% Auto-generated from data/segments.json + layout.json (+ transforms) — do not edit by hand.
% volume: 26Khu09
% mode: reading
% segments: 2268
% transform_rules: 0
% toc_marks: 268

% Volume layout defaults (layout.json ``layout``).
\csromanlayoutapply{1}{1.5}{21.6pt}{5pt}{6.3pt}{65pt}{2.5em}%

% Reading physical-page layout (layout.json ``page_layout_reading_mode``).
\csromanreadingpagelayoutvolume{1}{1.5}{21.6pt}{5pt}{6.3pt}{65pt}{2.5em}%
\csromanreadingpagelayoutenable

\csromanheader{2}{\textbf{ขุทฺทกนิกาย}}
\csromanmatikaanchor{mtk.0}
\csromantocmarknopage{chapter}{ปฏิสมฺภิทามคฺคปาฬิ}{mtk.0}
\bookmark[dest={mtk.0},level=0]{ปฏิสมฺภิทามคฺคปาฬิ}
\gambhira{\textbf{ปฏิสมฺภิทามคฺคปาฬิ}}
\namakkaram{นโม ตสฺส ภควโต อรหโต สมฺมาสมฺพุทฺธสฺส.}
\csromanmatikaanchor{mtk.1}
\csromantocmarknopage{chapter}{1. มหาวคฺค}{mtk.1}
\bookmark[dest={mtk.1},level=0]{1. มหาวคฺค}
\chapterhead{1. มหาวคฺค}
\csromanheader{2}{\textbf{มาติกา}}
\proseitem{1}{\csromanfolio{1}โสตาวธาเน ปญฺญา สุตมเย ญาณํ.}

\proseitem{2}{สุตฺวาน สํวเร ปญฺญา สีลมเย ญาณํ.}

\proseitem{3}{สํวริตฺวา สมาทหเน ปญฺญา สมาธิภาวนามเย ญาณํ.}

\proseitem{4}{ปจฺจยปริคฺคเห ปญฺญา ธมฺมฏฺฐิติญาณํ. 5.\csromanspacer{} อตีตานาคตปจฺจุปฺปนฺนานํ ธมฺมานํ สงฺขิปิตฺวา ววตฺถาเน ปญฺญา สมฺมสเน ญาณํ. 6.\csromanspacer{} ปจฺจุปฺปนฺนานํ ธมฺมานํ วิปริณามานุปสฺสเน ปญฺญา อุทยพฺพยานุปสฺสเน ญาณํ.}

\proseitem{7}{อารมฺมณํ ปฏิสงฺขา ภงฺคานุปสฺสเน ปญฺญา วิปสฺสเน ญาณํ.}

\proseitem{8}{ภยตุปฏฺฐาเน ปญฺญา อาทีนเว ญาณํ. 9.\csromanspacer{} มุญฺจิตุกมฺยตาปฏิสงฺขาสนฺติฏฺฐนา ปญฺญา สงฺขารุเปกฺขาสุ ญาณํ. 10.\csromanspacer{} พหิทฺธา วุฏฺฐานวิวฏฺฏเน ปญฺญา โคตฺรภุญาณํ. 11.\csromanspacer{} ทุภโต วุฏฺฐานวิวฏฺฏเน ปญฺญา มคฺเค ญาณํ. 12.\csromanspacer{} ปโยคปฺปฏิปฺปสฺสทฺธิปญฺญา ผเล ญาณํ.}

\proseitem{13}{\csromanfolio{2}ฉินฺนวฏุมานุปสฺสเน\footnote{ฉินฺนมนุปสฺสเน (สฺยา), ฉินฺนวฏฺฏมนุปสฺสเน (สี-ฏฺฐ)} ปญฺญา วิมุตฺติญาณํ. 14.\csromanspacer{} ตทา สมุทาคเต\footnote{สมุปาคเต (สฺยา)} ธมฺเม ปสฺสเน ปญฺญา ปจฺจเวกฺขเณ ญาณํ. 15.\csromanspacer{} อชฺฌตฺตววตฺถาเน ปญฺญา วตฺถุนานตฺเต ญาณํ. 16.\csromanspacer{} พหิทฺธาววตฺถาเน ปญฺญา โคจรนานตฺเต ญาณํ. 17.\csromanspacer{} จริยาววตฺถาเน ปญฺญา จริยานานตฺเต ญาณํ. 18.\csromanspacer{} จตุธมฺมววตฺถาเน ปญฺญา ภูมินานตฺเต ญาณํ. 19.\csromanspacer{} นวธมฺมววตฺถาเน ปญฺญา ธมฺมนานตฺเต ญาณํ. 20.\csromanspacer{} อภิญฺญาปญฺญา ญาตฏฺเฐ ญาณํ. 21.\csromanspacer{} ปริญฺญาปญฺญา ตีรณฏฺเฐ ญาณํ. 22.\csromanspacer{} ปหาเน ปญฺญา ปริจฺจาคฏฺเฐ ญาณํ. 23.\csromanspacer{} ภาวนาปญฺญา เอกรสฏฺเฐ ญาณํ. 24.\csromanspacer{} สจฺฉิกิริยาปญฺญา ผสฺสนฏฺเฐ\footnote{ผุสนฏฺเฐ (ก)} ญาณํ. 25.\csromanspacer{} อตฺถนานตฺเต ปญฺญา อตฺถปฏิสมฺภิเท ญาณํ. 26.\csromanspacer{} ธมฺมนานตฺเต ปญฺญา ธมฺมปฏิสมฺภิเท ญาณํ. 27.\csromanspacer{} นิรุตฺตินานตฺเต ปญฺญา นิรุตฺติปฏิสมฺภิเท ญาณํ. 28.\csromanspacer{} ปฏิภานนานตฺเต\footnote{ปฏิภาณนานตฺเต (สฺยา)} ปญฺญา ปฏิภานปฏิสมฺภิเท ญาณํ. 29.\csromanspacer{} วิหารนานตฺเต ปญฺญา วิหารฏฺเฐ ญาณํ. 30.\csromanspacer{} สมาปตฺตินานตฺเต ปญฺญา สมาปตฺตฏฺเฐ ญาณํ. 31.\csromanspacer{} วิหารสมาปตฺตินานตฺเต ปญฺญา วิหารสมาปตฺตฏฺเฐ ญาณํ. 32.\csromanspacer{} อวิกฺเขปปริสุทฺธตฺตา อาสวสมุจฺเฉเท ปญฺญา อานนฺตริกสมาธิมฺหิ ญาณํ. 33.\csromanspacer{} ทสฺสนาธิปเตยฺยํ สนฺโต จ วิหาราธิคโม ปณีตาธิมุตฺตตา ปญฺญา อรณวิหาเร ญาณํ.}

\csromanheader{2}{มาติกา}
\proseitem{34}{\csromanfolio{3}ทฺวีหิ พเลหิ สมนฺนาคตตฺตา ตโย จ สงฺขารานํ ปฏิปฺปสฺสทฺธิยา โสฬสหิ ญาณจริยาหิ นวหิ สมาธิจริยาหิ วสิภาวตา ปญฺญา นิโรธสมาปตฺติยา ญาณํ. 35.\csromanspacer{} สมฺปชานสฺส ปวตฺตปริยาทาเน ปญฺญา ปรินิพฺพาเน ญาณํ. 36.\csromanspacer{} สพฺพธมฺมานํ สมฺมา สมุจฺเฉเท นิโรเธ จ อนุปฏฺฐานตา ปญฺญา สมสีสฏฺเฐ ญาณํ. 37.\csromanspacer{} ปุถุนานตฺตเตชปริยาทาเน ปญฺญา สลฺเลขฏฺเฐ ญาณํ. 38.\csromanspacer{} อสลฺลีนตฺตปหิตตฺตปคฺคหฏฺเฐ ปญฺญา วีริยารมฺเภ ญาณํ. 39.\csromanspacer{} นานาธมฺมปฺปกาสนตา ปญฺญา อตฺถสนฺทสฺสเน ญาณํ. 40.\csromanspacer{} สพฺพธมฺมานํ เอกสงฺคหตานานตฺเตกตฺตปฏิเวเธ ปญฺญา ทสฺสนวิสุทฺธิญาณํ. 41.\csromanspacer{} วิทิตตฺตา ปญฺญา ขนฺติญาณํ. 42.\csromanspacer{} ผุฏฺฐตฺตา ปญฺญา ปริโยคาหเณ\footnote{ปริโยคาหเน (สฺยา)} ญาณํ. 43.\csromanspacer{} สโมทหเน ปญฺญา ปเทสวิหาเร ญาณํ. 44.\csromanspacer{} อธิปตตฺตา ปญฺญา สญฺญาวิวฏฺเฏ ญาณํ. 45.\csromanspacer{} นานตฺเต ปญฺญา เจโตวิวฏฺเฏ ญาณํ. 46.\csromanspacer{} อธิฏฺฐาเน ปญฺญา จิตฺตวิวฏฺเฏ ญาณํ. 47.\csromanspacer{} สุญฺญเต ปญฺญา ญาณวิวฏฺเฏ ญาณํ. 48.\csromanspacer{} โวสคฺเค\footnote{โวสฺสคฺเค (พหูสุ)} ปญฺญา วิโมกฺขวิวฏฺเฏ ญาณํ. 49.\csromanspacer{} ตถฏฺเฐ ปญฺญา สจฺจวิวฏฺเฏ ญาณํ. 50.\csromanspacer{} กายมฺปิ จิตฺตมฺปิ เอกววตฺถานตา สุขสญฺญญฺจ ลหุสญฺญญฺจ อธิฏฺฐานวเสน อิชฺฌนฏฺเฐ ปญฺญา อิทฺธิวิเธ ญาณํ. 51.\csromanspacer{} วิตกฺกวิปฺผารวเสน นานตฺเตกตฺตสทฺทนิมิตฺตานํ ปริโยคาหเณ ปญฺญา โสตธาตุวิสุทฺธิญฺญาณํ.}

\proseitem{52}{\csromanfolio{4}ติณฺณนฺนํ จิตฺตานํ วิปฺผารตฺตา อินฺทฺริยานํ ปสาทวเสน นานตฺเตกตฺตวิญฺญาณจริยา ปริโยคาหเณ ปญฺญา เจโตปริยญาณํ. 53.\csromanspacer{} ปจฺจยปฺปวตฺตานํ ธมฺมานํ นานตฺเตกตฺตกมฺมวิปฺผารวเสน ปริโยคาหเณ ปญฺญา ปุพฺเพนิวาสานุสฺสติญาณํ. 54.\csromanspacer{} โอภาสวเสน นานตฺเตกตฺตรูปนิมิตฺตานํ ทสฺสนฏฺเฐ ปญฺญา ทิพฺพจกฺขุญาณํ. 55.\csromanspacer{} จตุสฏฺฐิยา อากาเรหิ ติณฺณนฺนํ อินฺทฺริยานํ วสิภาวตา ปญฺญา อาสวานํ ขเย ญาณํ. 56.\csromanspacer{} ปริญฺญฏฺเฐ ปญฺญา ทุกฺเข ญาณํ. 57.\csromanspacer{} ปหานฏฺเฐ ปญฺญา สมุทเย ญาณํ. 58.\csromanspacer{} สจฺฉิกิริยฏฺเฐ ปญฺญา นิโรเธ ญาณํ. 59.\csromanspacer{} ภาวนฏฺเฐ ปญฺญา มคฺเค ญาณํ. 60.\csromanspacer{} ทุกฺเข ญาณํ. 61.\csromanspacer{} ทุกฺขสมุทเย ญาณํ. 62.\csromanspacer{} ทุกฺขนิโรเธ ญาณํ. 63.\csromanspacer{} ทุกฺขนิโรธคามินิยา ปฏิปทาย ญาณํ. 64.\csromanspacer{} อตฺถปฏิสมฺภิเท ญาณํ. 65.\csromanspacer{} ธมฺมปฏิสมฺภิเท ญาณํ. 66.\csromanspacer{} นิรุตฺติปฏิสมฺภิเท ญาณํ. 67.\csromanspacer{} ปฏิภานปฏิสมฺภิเท ญาณํ. 68.\csromanspacer{} อินฺทฺริยปโรปริยตฺตญาณํ. 69.\csromanspacer{} สตฺตานํ อาสยานุสเย ญาณํ. 70.\csromanspacer{} ยมกปาฏิหีเร ญาณํ. 71.\csromanspacer{} มหากรุณาสมาปตฺติยา ญาณํ. 72.\csromanspacer{} สพฺพญฺญุตญฺญาณํ.}

\proseitem{73}{\csromanfolio{5}อนาวรณญาณํ.}

\prose{อิมานิ เตสตฺตติ ญาณานิ, อิเมสํ เตสตฺตติยา ญาณานํ สตฺตสฏฺฐิ ญาณานิ สาวกสาธารณานิ, ฉ ญาณานิ อสาธารณานิ สาวเกหิ.\csromanspacer{} มาติกา นิฏฺฐิตา.}
\csromansectionrule

\csromanmatikaanchor{mtk.2}
\csromanmatikaanchor{mtk.3}
\csromantocmarknopage{section}{1. ญาณกถา}{mtk.2}
\bookmark[dest={mtk.2},level=1]{1. ญาณกถา}
\csromantocmarkref{subsection}{1. สุตมยญาณนิทฺเทส}{mtk.3}
\bookmark[dest={mtk.3},level=2]{1. สุตมยญาณนิทฺเทส}
\csromanheader{2}{1. สุตมยญาณนิทฺเทส}
\csromanheader{2}{1. กถํ โสตาวธาเน ปญฺญา สุตมเย ญาณํ\csromandash{}}
\prose{“อิเม ธมฺมา อภิญฺเญยฺยา”ติ โสตาวธานํ, ตํปชานนา ปญฺญา สุตมเย ญาณํ. (1)}

\prose{“อิเม ธมฺมา ปริญฺเญยฺยา”ติ โสตาวธานํ, ตํปชานนา ปญฺญา สุตมเย ญาณํ. (2)}

\prose{“อิเม ธมฺมา ปหาตพฺพา”ติ โสตาวธานํ, ตํปชานนา ปญฺญา สุตมเย ญาณํ. (3)}

\prose{“อิเม ธมฺมา ภาเวตพฺพา”ติ โสตาวธานํ, ตํปชานนา ปญฺญา สุตมเย ญาณํ. (4)}

\prose{“อิเม ธมฺมา สจฺฉิกาตพฺพา”ติ โสตาวธานํ, ตํปชานนา ปญฺญา สุตมเย ญาณํ. (5)}

\prose{“อิเม ธมฺมา หานภาคิยา”ติ โสตาวธานํ, ตํปชานนา ปญฺญา สุตมเย ญาณํ. (6)}

\prose{“อิเม ธมฺมา ฐิติภาคิยา”ติ โสตาวธานํ, ตํปชานนา ปญฺญา สุตมเย ญาณํ. (7)}

\prose{“อิเม ธมฺมา วิเสสภาคิยา”ติ โสตาวธานํ, ตํปชานนา ปญฺญา สุตมเย ญาณํ. (8)}

\prose{“อิเม ธมฺมา นิพฺเพธภาคิยา”ติ โสตาวธานํ, ตํปชานนา ปญฺญา สุตมเย ญาณํ. (9)}

\prose{\csromanfolio{6}“สพฺเพ สงฺขารา อนิจฺจา”ติ โสตาวธานํ, ตํปชานนา ปญฺญา สุตมเย ญาณํ. (10)}

\prose{“สพฺเพ สงฺขารา ทุกฺขา”ติ โสตาวธานํ, ตํปชานนา ปญฺญา สุตมเย ญาณํ. (11)}

\prose{“สพฺเพ ธมฺมา อนตฺตา”ติ โสตาวธานํ, ตํปชานนา ปญฺญา สุตมเย ญาณํ. (12)}

\prose{“อิทํ ทุกฺขํ อริยสจฺจนฺ”ติ โสตาวธานํ, ตํปชานนา ปญฺญา สุตมเย ญาณํ. (13)}

\prose{“อิทํ ทุกฺขสมุทยํ\footnote{ทุกฺขสมุทโย (สฺยา)} อริยสจฺจนฺ”ติ โสตาวธานํ, ตํปชานนา ปญฺญา สุตมเย ญาณํ. (14)}

\prose{“อิทํ ทุกฺขนิโรธํ\footnote{ทุกฺขนิโรโธ (สฺยา) อฏฺฐกถา โอโลเกตพฺพา.} อริยสจฺจนฺ”ติ โสตาวธานํ, ตํปชานนา ปญฺญา สุตมเย ญาณํ. (15)}

\prose{“อิทํ ทุกฺขนิโรธคามินี ปฏิปทา อริยสจฺจนฺ”ติ โสตาวธานํ, ตํปชานนา ปญฺญา สุตมเย ญาณํ. (16)}

\proseitem{2}{กถํ “อิเม ธมฺมา อภิญฺเญยฺยา”ติ โสตาวธานํ, ตํปชานนา ปญฺญา สุตมเย ญาณํ\csromandash{}}

\prose{เอโก ธมฺโม อภิญฺเญยฺโย, สพฺเพ สตฺตา อาหารฏฺฐิติกา.\csromanspacer{} เทฺว ธมฺมา อภิญฺเญยฺยา, เทฺว ธาตุโย.\csromanspacer{} ตโย ธมฺมา อภิญฺเญยฺยา, ติสฺโส ธาตุโย.\csromanspacer{} จตฺตาโร ธมฺมา อภิญฺเญยฺยา, จตฺตาริ อริยสจฺจานิ.\csromanspacer{} ปญฺจ ธมฺมา อภิญฺเญยฺยา, ปญฺจ วิมุตฺตายตนานิ.\csromanspacer{} ฉ ธมฺมา อภิญฺเญยฺยา, ฉ อนุตฺตริยานิ.\csromanspacer{} สตฺต ธมฺมา อภิญฺเญยฺยา, สตฺต นิทฺทสวตฺถูนิ.\csromanspacer{} อฏฺฐ ธมฺมา อภิญฺเญยฺยา, อฏฺฐ อภิภายตนานิ.\csromanspacer{} นว ธมฺมา อภิญฺเญยฺยา, นว อนุปุพฺพวิหารา.\csromanspacer{} ทส ธมฺมา อภิญฺเญยฺยา, ทส นิชฺชรวตฺถูนิ. (10)}

\proseitem{3}{สพฺพํ ภิกฺขเว อภิญฺเญยฺยํ.\csromanspacer{} กิญฺจ ภิกฺขเว สพฺพํ อภิญฺเญยฺยํ? จกฺขุ\footnote{จกฺขุํ (สฺยา, ก)} ภิกฺขเว อภิญฺเญยฺยํ, รูปา อภิญฺเญยฺยา, จกฺขุวิญฺญาณํ อภิญฺเญยฺยํ, จกฺขุสมฺผสฺโส อภิญฺเญยฺโย, ยมฺปิทํ\footnote{ยมิทํ (ก) สํ 2. 258 ปิฏฺเฐ ปสฺสิตพฺพา.} จกฺขุสมฺผสฺสปจฺจยา อุปฺปชฺชติ เวทยิตํ สุขํ วา ทุกฺขํ วา อทุกฺขมสุขํ วา, ตมฺปิ อภิญฺเญยฺยํ.\csromanspacer{} โสตํ อภิญฺเญยฺยํ, \csromanfolio{7}สทฺทา อภิญฺเญยฺยา.\csromanspacer{} ฆานํ อภิญฺเญยฺยํ, คนฺธา อภิญฺเญยฺยา.\csromanspacer{} ชิวฺหา อภิญฺเญยฺยา, รสา อภิญฺเญยฺยา.\csromanspacer{} กาโย อภิญฺเญยฺโย, โผฏฺฐพฺพา อภิญฺเญยฺยา.\csromanspacer{} มโน อภิญฺเญยฺโย, ธมฺมา อภิญฺเญยฺยา, มโนวิญฺญาณํ อภิญฺเญยฺยํ, มโนสมฺผสฺโส อภิญฺเญยฺโย, ยมฺปิทํ มโนสมฺผสฺสปจฺจยา อุปฺปชฺชติ เวทยิตํ สุขํ วา ทุกฺขํ วา อทุกฺขมสุขํ วา, ตมฺปิ อภิญฺเญยฺยํ. (30)}

\prose{รูปํ อภิญฺเญยฺยํ, เวทนา อภิญฺเญยฺยา, สญฺญา อภิญฺเญยฺยา, สงฺขารา อภิญฺเญยฺยา, วิญฺญาณํ อภิญฺเญยฺยํ. (5)}

\prose{จกฺขุ อภิญฺเญยฺยํ, โสตํ อภิญฺเญยฺยํ, ฆานํ อภิญฺเญยฺยํ, ชิวฺหา อภิญฺเญยฺยา, กาโย อภิญฺเญยฺโย, มโน อภิญฺเญยฺโย.\csromanspacer{}\csromanspacer{}รูปา อภิญฺเญยฺยา, สทฺทา อภิญฺเญยฺยา, คนฺธา อภิญฺเญยฺยา, รสา อภิญฺเญยฺยา, โผฏฺฐพฺพา อภิญฺเญยฺยา, ธมฺมา อภิญฺเญยฺยา.\csromanspacer{}\csromanspacer{}จกฺขุวิญฺญาณํ อภิญฺเญยฺยํ, โสตวิญฺญาณํ อภิญฺเญยฺยํ, ฆานวิญฺญาณํ อภิญฺเญยฺยํ, ชิวฺหาวิญฺญาณํ อภิญฺเญยฺยํ, กายวิญฺญาณํ อภิญฺเญยฺยํ, มโนวิญฺญาณํ อภิญฺเญยฺยํ.\csromanspacer{}\csromanspacer{}จกฺขุสมฺผสฺโส อภิญฺเญยฺโย, โสตสมฺผสฺโส อภิญฺเญยฺโย, ฆานสมฺผสฺโส อภิญฺเญยฺโย, ชิวฺหาสมฺผสฺโส อภิญฺเญยฺโย, กายสมฺผสฺโส อภิญฺเญยฺโย, มโนสมฺผสฺโส อภิญฺเญยฺโย.\csromanspacer{}\csromanspacer{}จกฺขุสมฺผสฺสชา เวทนา อภิญฺเญยฺยา, โสตสมฺผสฺสชา เวทนา อภิญฺเญยฺยา, ฆานสมฺผสฺสชา เวทนา อภิญฺเญยฺยา, ชิวฺหาสมฺผสฺสชา เวทนา อภิญฺเญยฺยา, กายสมฺผสฺสชา เวทนา อภิญฺเญยฺยา, มโนสมฺผสฺสชา เวทนา อภิญฺเญยฺยา.\csromanspacer{}\csromanspacer{}รูปสญฺญา อภิญฺเญยฺยา, สทฺทสญฺญา อภิญฺเญยฺยา, คนฺธสญฺญา อภิญฺเญยฺยา, รสสญฺญา อภิญฺเญยฺยา, โผฏฺฐพฺพสญฺญา อภิญฺเญยฺยา, ธมฺมสญฺญา อภิญฺเญยฺยา.\csromanspacer{}\csromanspacer{}รูปสญฺเจตนา อภิญฺเญยฺยา, สทฺทสญฺเจตนา อภิญฺเญยฺยา, คนฺธสญฺเจตนา อภิญฺเญยฺยา, รสสญฺเจตนา อภิญฺเญยฺยา, โผฏฺฐพฺพสญฺเจตนา อภิญฺเญยฺยา, ธมฺมสญฺเจตนา อภิญฺเญยฺยา.\csromanspacer{}\csromanspacer{}รูปตณฺหา อภิญฺเญยฺยา, สทฺทตณฺหา อภิญฺเญยฺยา, คนฺธตณฺหา อภิญฺเญยฺยา, รสตณฺหา อภิญฺเญยฺยา, โผฏฺฐพฺพตณฺหา อภิญฺเญยฺยา, ธมฺมตณฺหา อภิญฺเญยฺยา.\csromanspacer{}\csromanspacer{}รูปวิตกฺโก อภิญฺเญยฺโย, สทฺทวิตกฺโก อภิญฺเญยฺโย, คนฺธวิตกฺโก อภิญฺเญยฺโย, รสวิตกฺโก อภิญฺเญยฺโย, โผฏฺฐพฺพวิตกฺโก อภิญฺเญยฺโย, ธมฺมวิตกฺโก อภิญฺเญยฺโย.\csromanspacer{}\csromanspacer{}รูปวิจาโร อภิญฺเญยฺโย, สทฺทวิจาโร \csromanfolio{8}อภิญฺเญยฺโย, คนฺธวิจาโร อภิญฺเญยฺโย, รสวิจาโร อภิญฺเญยฺโย, โผฏฺฐพฺพวิจาโร อภิญฺเญยฺโย, ธมฺมวิจาโร อภิญฺเญยฺโย. (60)}

\proseitem{4}{ปถวีธาตุ อภิญฺเญยฺยา, อาโปธาตุ อภิญฺเญยฺยา, เตโชธาตุ อภิญฺเญยฺยา, วาโยธาตุ อภิญฺเญยฺยา, อากาสธาตุ อภิญฺเญยฺยา, วิญฺญาณธาตุ อภิญฺเญยฺยา. (6)}

\prose{ปถวีกสิณํ อภิญฺเญยฺยํ, อาโปกสิณํ อภิญฺเญยฺยํ, เตโชกสิณํ อภิญฺเญยฺยํ, วาโยกสิณํ อภิญฺเญยฺยํ, นีลกสิณํ อภิญฺเญยฺยํ, ปีตกสิณํ อภิญฺเญยฺยํ, โลหิตกสิณํ อภิญฺเญยฺยํ, โอทาตกสิณํ อภิญฺเญยฺยํ, อากาสกสิณํ อภิญฺเญยฺยํ, วิญฺญาณกสิณํ อภิญฺเญยฺยํ. (10)}

\prose{เกสา อภิญฺเญยฺยา, โลมา อภิญฺเญยฺยา, นขา อภิญฺเญยฺยา, ทนฺตา อภิญฺเญยฺยา, ตโจ อภิญฺเญยฺโย, มํสํ อภิญฺเญยฺยํ, นฺหารู\footnote{นหารู (สฺยา)} อภิญฺเญยฺยา, อฏฺฐี อภิญฺเญยฺยา, อฏฺฐิมิญฺชา อภิญฺเญยฺยา\footnote{อฏฺฐิมญฺชํ อภิญฺเญยฺยํ (สฺยา, ก)}, วกฺกํ อภิญฺเญยฺยํ, หทยํ อภิญฺเญยฺยํ, ยกนํ อภิญฺเญยฺยํ, กิโลมกํ อภิญฺเญยฺยํ, ปิหกํ อภิญฺเญยฺยํ, ปปฺผาสํ อภิญฺเญยฺยํ, อนฺตํ อภิญฺเญยฺยํ, อนฺตคุณํ อภิญฺเญยฺยํ, อุทริยํ อภิญฺเญยฺยํ, กรีสํ อภิญฺเญยฺยํ, ปิตฺตํ อภิญฺเญยฺยํ, เสมฺหํ อภิญฺเญยฺยํ, ปุพฺโพ อภิญฺเญยฺโย, โลหิตํ อภิญฺเญยฺยํ, เสโท อภิญฺเญยฺโย, เมโท อภิญฺเญยฺโย, อสฺสุ อภิญฺเญยฺยํ, วสา อภิญฺเญยฺยา, เขโฬ อภิญฺเญยฺโย, สิงฺฆาณิกา อภิญฺเญยฺยา, ลสิกา อภิญฺเญยฺยา, มุตฺตํ อภิญฺเญยฺยํ, มตฺถลุงฺคํ อภิญฺเญยฺยํ. (32)}

\prose{จกฺขายตนํ อภิญฺเญยฺยํ, รูปายตนํ อภิญฺเญยฺยํ.\csromanspacer{} โสตายตนํ อภิญฺเญยฺยํ, สทฺทายตนํ อภิญฺเญยฺยํ.\csromanspacer{} ฆานายตนํ อภิญฺเญยฺยํ, คนฺธายตนํ อภิญฺเญยฺยํ.\csromanspacer{} ชิวฺหายตนํ อภิญฺเญยฺยํ, รสายตนํ อภิญฺเญยฺยํ.กายายตนํ อภิญฺเญยฺยํ, โผฏฺฐพฺพายตนํ อภิญฺเญยฺยํ.\csromanspacer{} มนายตนํ อภิญฺเญยฺยํ, ธมฺมายตนํ อภิญฺเญยฺยํ. (12)}

\prose{จกฺขุธาตุ อภิญฺเญยฺยา, รูปธาตุ อภิญฺเญยฺยา, จกฺขุวิญฺญาณธาตุ อภิญฺเญยฺยา.\csromanspacer{} โสตธาตุ อภิญฺเญยฺยา, สทฺทธาตุ อภิญฺเญยฺยา, โสตวิญฺญาณธาตุ อภิญฺเญยฺยา.\csromanspacer{} ฆานธาตุ อภิญฺเญยฺยา, คนฺธธาตุ อภิญฺเญยฺยา, ฆานวิญฺญาณธาตุ อภิญฺเญยฺยา.\csromanspacer{} ชิวฺหาธาตุ อภิญฺเญยฺยา, \csromanfolio{9}รสธาตุ อภิญฺเญยฺยา, ชิวฺหาวิญฺญาณธาตุ อภิญฺเญยฺยา.\csromanspacer{} กายธาตุ อภิญฺเญยฺยา, โผฏฺฐพฺพธาตุ อภิญฺเญยฺยา, กายวิญฺญาณธาตุ อภิญฺเญยฺยา.\csromanspacer{} มโนธาตุ อภิญฺเญยฺยา, ธมฺมธาตุ อภิญฺเญยฺยา, มโนวิญฺญาณธาตุ อภิญฺเญยฺยา. (18)}

\prose{จกฺขุนฺทฺริยํ อภิญฺเญยฺยํ, โสตินฺทฺริยํ อภิญฺเญยฺยํ, ฆานินฺทฺริยํ อภิญฺเญยฺยํ, ชิวฺหินฺทฺริยํ อภิญฺเญยฺยํ, กายินฺทฺริยํ อภิญฺเญยฺยํ, มนินฺทฺริยํ อภิญฺเญยฺยํ, ชีวิตินฺทฺริยํ อภิญฺเญยฺยํ, อิตฺถินฺทฺริยํ อภิญฺเญยฺยํ, ปุริสินฺทฺริยํ อภิญฺเญยฺยํ, สุขินฺทฺริยํ อภิญฺเญยฺยํ, ทุกฺขินฺทฺริยํ อภิญฺเญยฺยํ, โสมนสฺสินฺทฺริยํ อภิญฺเญยฺยํ, โทมนสฺสินฺทฺริยํ อภิญฺเญยฺยํ, อุเปกฺขินฺทฺริยํ อภิญฺเญยฺยํ, สทฺธินฺทฺริยํ อภิญฺเญยฺยํ, วีริยินฺทฺริยํ\footnote{วิริยินฺทฺริยํ (สฺยา)} อภิญฺเญยฺยํ, สตินฺทฺริยํ อภิญฺเญยฺยํ, สมาธินฺทฺริยํ อภิญฺเญยฺยํ, ปญฺญินฺทฺริยํ อภิญฺเญยฺยํ, อนญฺญาตญฺญสฺสามีตินฺทฺริยํ อภิญฺเญยฺยํ, อญฺญินฺทฺริยํ อภิญฺเญยฺยํ, อญฺญาตาวินฺทฺริยํ อภิญฺเญยฺยํ. (22)}

\proseitem{5}{กามธาตุ อภิญฺเญยฺยา, รูปธาตุ อภิญฺเญยฺยา, อรูปธาตุ อภิญฺเญยฺยา.\csromanspacer{}\csromanspacer{}กามภโว อภิญฺเญยฺโย, รูปภโว อภิญฺเญยฺโย, อรูปภโว อภิญฺเญยฺโย.\csromanspacer{}\csromanspacer{}สญฺญาภโว อภิญฺเญยฺโย, อสญฺญาภโว อภิญฺเญยฺโย, เนวสญฺญานาสญฺญาภโว อภิญฺเญยฺโย.\csromanspacer{}\csromanspacer{}เอกโวการภโว อภิญฺเญยฺโย, จตุโวการภโว อภิญฺเญยฺโย, ปญฺจโวการภโว อภิญฺเญยฺโย. (12)}

\proseitem{6}{ปฐมํ ฌานํ\footnote{ปฐมชฺฌานํ (สฺยา) เอวมีทิเสสุ ฐาเนสุ.} อภิญฺเญยฺยํ, ทุติยํ ฌานํ อภิญฺเญยฺยํ, ตติยํ ฌานํ อภิญฺเญยฺยํ, จตุตฺถํ ฌานํ อภิญฺเญยฺยํ.\csromanspacer{}\csromanspacer{}เมตฺตาเจโตวิมุตฺติ อภิญฺเญยฺยา, กรุณาเจโตวิมุตฺติ อภิญฺเญยฺยา, มุทิตาเจโตวิมุตฺติ อภิญฺเญยฺยา, อุเปกฺขาเจโตวิมุตฺติ อภิญฺเญยฺยา.\csromanspacer{}\csromanspacer{}อากาสานญฺจายตนสมาปตฺติ อภิญฺเญยฺยา, วิญฺญาณญฺจายตนสมาปตฺติ อภิญฺเญยฺยา, อากิญฺจญฺญายตนสมาปตฺติ อภิญฺเญยฺยา, เนวสญฺญานาสญฺญายตนสมาปตฺติ อภิญฺเญยฺยา. (12)}

\prose{อวิชฺชา อภิญฺเญยฺยา, สงฺขารา อภิญฺเญยฺยา, วิญฺญาณํ อภิญฺเญยฺยํ, นามรูปํ อภิญฺเญยฺยํ, สฬายตนํ อภิญฺเญยฺยํ, ผสฺโส อภิญฺเญยฺโย, เวทนา อภิญฺเญยฺยา, ตณฺหา อภิญฺเญยฺยา, อุปาทานํ อภิญฺเญยฺยํ, ภโว อภิญฺเญยฺโย, ชาติ อภิญฺเญยฺยา, ชรามรณํ อภิญฺเญยฺยํ. (12)}

\proseitem{7}{\csromanfolio{10}ทุกฺขํ อภิญฺเญยฺยํ, ทุกฺขสมุทโย อภิญฺเญยฺโย, ทุกฺขนิโรโธ อภิญฺเญยฺโย, ทุกฺขนิโรธคามินี ปฏิปทา อภิญฺเญยฺยา.\csromanspacer{} รุปํ อภิญฺเญยฺยํ, รูปสมุทโย อภิญฺเญยฺโย, รูปนิโรโธ อภิญฺเญยฺโย, รูปนิโรธคามินี ปฏิปทา อภิญฺเญยฺยา.\csromanspacer{} เวทนา อภิญฺเญยฺยา ฯเปฯ.\csromanspacer{} สญฺญา อภิญฺเญยฺยา ฯเปฯ.\csromanspacer{} สงฺขารา อภิญฺเญยฺยา ฯเปฯ.\csromanspacer{} วิญฺญาณํ อภิญฺเญยฺยํ.\csromanspacer{} จกฺขุ อภิญฺเญยฺยํ ฯเปฯ.\csromanspacer{} ชรามรณํ อภิญฺเญยฺยํ, ชรามรณสมุทโย อภิญฺเญยฺโย, ชรามรณนิโรโธ อภิญฺเญยฺโย, ชรามรณนิโรธคามินี ปฏิปทา อภิญฺเญยฺยา. (808)}

\prose{ทุกฺขสฺส ปริญฺญฏฺโฐ อภิญฺเญยฺโย, ทุกฺขสมุทยสฺส ปหานฏฺโฐ อภิญฺเญยฺโย, ทุกฺขนิโรธสฺส สจฺฉิกิริยฏฺโฐ อภิญฺเญยฺโย, ทุกฺขนิโรธคามินิยา ปฏิปทาย ภาวนฏฺโฐ อภิญฺเญยฺโย.\csromanspacer{} รูปสฺส ปริญฺญฏฺโฐ อภิญฺเญยฺโย, รูปสมุทยสฺส ปหานฏฺโฐ อภิญฺเญยฺโย, รูปนิโรธสฺส สจฺฉิกิริยฏฺโฐ อภิญฺเญยฺโย, รูปนิโรธคามินิยา ปฏิปทาย ภาวนฏฺโฐ อภิญฺเญยฺโย.\csromanspacer{} เวทนาย ฯเปฯ.\csromanspacer{} สญฺญาย.\csromanspacer{} สงฺขารานํ.\csromanspacer{} วิญฺญาณสฺส.\csromanspacer{} จกฺขุสฺส ฯเปฯ.\csromanspacer{} ชรามรณสฺส ปริญฺญฏฺโฐ อภิญฺเญยฺโย, ชรามรณสมุทยสฺส ปหานฏฺโฐ อภิญฺเญยฺโย, ชรามรณนิโรธสฺส สจฺฉิกิริยฏฺโฐ อภิญฺเญยฺโย, ชรามรณนิโรธคามินิยา ปฏิปทาย ภาวนฏฺโฐ อภิญฺเญยฺโย. (808)}

\prose{ทุกฺขสฺส ปริญฺญาปฏิเวธฏฺโฐ อภิญฺเญยฺโย, ทุกฺขสมุทยสฺส ปหานปฏิเวธฏฺโฐ อภิญฺเญยฺโย, ทุกฺขนิโรธสฺส สจฺฉิกิริยาปฏิเวธฏฺโฐ อภิญฺเญยฺโย, ทุกฺขนิโรธคามินิยา ปฏิปทาย ภาวนาปฏิเวธฏฺโฐ อภิญฺเญยฺโย.\csromanspacer{} รูปสฺส ปริญฺญาปฏิเวธฏฺโฐ อภิญฺเญยฺโย, รูปสมุทยสฺส ปหานปฏิเวธฏฺโฐ อภิญฺเญยฺโย, รูปนิโรธสฺส สจฺฉิกิริยาปฏิเวธฏฺโฐ อภิญฺเญยฺโย, รูปนิโรธคามินิยา ปฏิปทาย ภาวนาปฏิเวธฏฺโฐ อภิญฺเญยฺโย.\csromanspacer{} เวทนาย ฯเปฯ.\csromanspacer{} สญฺญาย.\csromanspacer{} สงฺขารานํ.\csromanspacer{} วิญฺญาณสฺส.\csromanspacer{} จกฺขุสฺส ฯเปฯ.\csromanspacer{} ชรามรณสฺส ปริญฺญาปฏิเวธฏฺโฐ อภิญฺเญยฺโย, ชรามรณสมุทยสฺส ปหานปฏิเวธฏฺโฐ อภิญฺเญยฺโย, ชรามรณนิโรธสฺส สจฺฉิกิริยาปฏิเวธฏฺโฐ อภิญฺเญยฺโย, ชรามรณนิโรธคามินิยา ปฏิปทาย ภาวนาปฏิเวธฏฺโฐ อภิญฺเญยฺโย. (808)}

\proseitem{8}{ทุกฺขํ อภิญฺเญยฺยํ, ทุกฺขสมุทโย อภิญฺเญยฺโย, ทุกฺขนิโรโธ อภิญฺเญยฺโย, ทุกฺขสฺส สมุทยนิโรโธ อภิญฺเญยฺโย, ทุกฺขสฺส ฉนฺทราคนิโรโธ อภิญฺเญยฺโย, ทุกฺขสฺส อสฺสาโท อภิญฺเญยฺโย, ทุกฺขสฺส \csromanfolio{11}อาทีนโว อภิญฺเญยฺโย, ทุกฺขสฺส นิสฺสรณํ อภิญฺเญยฺยํ.\csromanspacer{}\csromanspacer{}รูปํ อภิญฺเญยฺยํ, รูปสมุทโย อภิญฺเญยฺโย, รูปนิโรโธ อภิญฺเญยฺโย, รูปสฺส สมุทยนิโรโธ อภิญฺเญยฺโย, รูปสฺส ฉนฺทราคนิโรโธ อภิญฺเญยฺโย, รูปสฺส อสฺสาโท อภิญฺเญยฺโย, รูปสฺส อาทีนโว อภิญฺเญยฺโย, รูปสฺส นิสฺสรณํ อภิญฺเญยฺยํ.\csromanspacer{} เวทนา อภิญฺเญยฺยา ฯเปฯ.\csromanspacer{} สญฺญา อภิญฺเญยฺยา.\csromanspacer{} สงฺขารา อภิญฺเญยฺยา.\csromanspacer{} วิญฺญาณํ อภิญฺเญยฺยํ.\csromanspacer{} จกฺขุ อภิญฺเญยฺยํ ฯเปฯ.\csromanspacer{} ชรามรณํ อภิญฺเญยฺยํ, ชรามรณสมุทโย อภิญฺเญยฺโย, ชรามรณนิโรโธ อภิญฺเญยฺโย, ชรามรณสฺส สมุทยนิโรโธ อภิญฺเญยฺโย, ชรามรณสฺส ฉนฺทราคนิโรโธ อภิญฺเญยฺโย, ชรามรณสฺส อสฺสาโท อภิญฺเญยฺโย, ชรามรณสฺส อาทีนโว อภิญฺเญยฺโย, ชรามรณสฺส นิสฺสรณํ อภิญฺเญยฺยํ. (1616)}

\prose{ทุกฺขํ อภิญฺเญยฺยํ, ทุกฺขสมุทโย อภิญฺเญยฺโย, ทุกฺขนิโรโธ อภิญฺเญยฺโย, ทุกฺขนิโรธคามินี ปฏิปทา อภิญฺเญยฺยา, ทุกฺขสฺส อสฺสาโท อภิญฺเญยฺโย, ทุกฺขสฺส อาทีนโว อภิญฺเญยฺโย, ทุกฺขสฺส นิสฺสรณํ อภิญฺเญยฺยํ.\csromanspacer{}\csromanspacer{}รูปํ อภิญฺเญยฺยํ, รูปสมุทโย อภิญฺเญยฺโย, รูปนิโรโธ อภิญฺเญยฺโย, รูปนิโรธคามินี ปฏิปทา อภิญฺเญยฺยา, รูปสฺส อสฺสาโท อภิญฺเญยฺโย, รูปสฺส อาทีนโว อภิญฺเญยฺโย, รูปสฺส นิสฺสรณํ อภิญฺเญยฺยํ.\csromanspacer{} เวทนา อภิญฺเญยฺยา ฯเปฯ.\csromanspacer{} สญฺญา อภิญฺเญยฺยา.\csromanspacer{} สงฺขารา อภิญฺเญยฺยา.\csromanspacer{} วิญฺญาณํ อภิญฺเญยฺยํ.\csromanspacer{} จกฺขุ อภิญฺเญยฺยํ ฯเปฯ.\csromanspacer{} ชรามรณํ อภิญฺเญยฺยํ.\csromanspacer{} ชรามรณสมุทโย อภิญฺเญยฺโย, ชรามรณนิโรโธ อภิญฺเญยฺโย, ชรามรณนิโรธคามินี ปฏิปทา อภิญฺเญยฺยา, ชรามรณสฺส อสฺสาโท อภิญฺเญยฺโย, ชรามรณสฺส อาทีนโว อภิญฺเญยฺโย, ชรามรณสฺส นิสฺสรณํ อภิญฺเญยฺยํ. (1414)}

\proseitem{9}{อนิจฺจานุปสฺสนา อภิญฺเญยฺยา, ทุกฺขานุปสฺสนา อภิญฺเญยฺยา, อนตฺตานุปสฺสนา อภิญฺเญยฺยา, นิพฺพิทานุปสฺสนา อภิญฺเญยฺยา, วิราคานุปสฺสนา อภิญฺเญยฺยา, นิโรธานุปสฺสนา อภิญฺเญยฺยา, ปฏินิสฺสคฺคานุปสฺสนา อภิญฺเญยฺยา.\csromanspacer{}\csromanspacer{}รูเป อนิจฺจานุปสฺสนา อภิญฺเญยฺยา, รูเป ทุกฺขานุปสฺสนา อภิญฺเญยฺยา, รูเป อนตฺตานุปสฺสนา อภิญฺเญยฺยา, รูเป นิพฺพิทานุปสฺสนา อภิญฺเญยฺยา, รูเป วิราคานุปสฺสนา อภิญฺเญยฺยา, รูเป นิโรธานุปสฺสนา อภิญฺเญยฺยา, รูเป ปฏินิสฺสคฺคานุปสฺสนา อภิญฺเญยฺยา.\csromanspacer{} เวทนาย ฯเปฯ.\csromanspacer{} สญฺญาย.\csromanspacer{} สงฺขาเรสุ.\csromanspacer{} วิญฺญาเณ.\csromanspacer{} จกฺขุสฺมิํ ฯเปฯ.\csromanspacer{} ชรามรเณ \csromanfolio{12}อนิจฺจานุปสฺสนา อภิญฺเญยฺยา, ชรามรเณ ทุกฺขานุปสฺสนา อภิญฺเญยฺยา, ชรามรเณ อนตฺตานุปสฺสนา อภิญฺเญยฺยา, ชรามรเณ นิพฺพิทานุปสฺสนา อภิญฺเญยฺยา, ชรามรเณ วิราคานุปสฺสนา อภิญฺเญยฺยา, ชรามรเณ นิโรธานุปสฺสนา อภิญฺเญยฺยา, ชรามรเณ ปฏินิสฺสคฺคานุปสฺสนา อภิญฺเญยฺยา. (1414)}

\proseitem{10}{อุปฺปาโท อภิญฺเญยฺโย, ปวตฺตํ อภิญฺเญยฺยํ, นิมิตฺตํ อภิญฺเญยฺยํ, อายูหนา\footnote{อายุหนา (สฺยา) เอวมุปริปิ.} อภิญฺเญยฺยา, ปฏิสนฺธิ อภิญฺเญยฺยา, คติ อภิญฺเญยฺยา, นิพฺพตฺติ อภิญฺเญยฺยา, อุปปตฺติ อภิญฺเญยฺยา, ชาติ อภิญฺเญยฺยา, ชรา อภิญฺเญยฺยา, พฺยาธิ อภิญฺเญยฺโย, มรณํ อภิญฺเญยฺยํ, โสโก อภิญฺเญยฺโย, ปริเทโว อภิญฺเญยฺโย, อุปายาโส อภิญฺเญยฺโย. (15)}

\prose{อนุปฺปาโท อภิญฺเญยฺโย, อปฺปวตฺตํ อภิญฺเญยฺยํ, อนิมิตฺตํ อภิญฺเญยฺยํ, อนายูหนา\footnote{อนายุหนา (สฺยา) เอวมุปริปิ.} อภิญฺเญยฺยา, อปฺปฏิสนฺธิ อภิญฺเญยฺยา, อคติ อภิญฺเญยฺยา, อนิพฺพตฺติ อภิญฺเญยฺยา, อนุปปตฺติ อภิญฺเญยฺยา, อชาติ อภิญฺเญยฺยา, อชรา อภิญฺเญยฺยา, อพฺยาธิ อภิญฺเญยฺโย, อมตํ อภิญฺเญยฺยํ, อโสโก อภิญฺเญยฺโย, อปริเทโว อภิญฺเญยฺโย, อนุปายาโส อภิญฺเญยฺโย. (15)}

\prose{อุปฺปาโท อภิญฺเญยฺโย, อนุปฺปาโท อภิญฺเญยฺโย.\csromanspacer{} ปวตฺตํ อภิญฺเญยฺยํ, อปฺปวตฺตํ อภิญฺเญยฺยํ.\csromanspacer{} นิมิตฺตํ อภิญฺเญยฺยํ, อนิมิตฺตํ อภิญฺเญยฺยํ.\csromanspacer{} อายูหนา อภิญฺเญยฺยา, อนายูหนา อภิญฺเญยฺยา.\csromanspacer{} ปฏิสนฺธิ อภิญฺเญยฺยา, อปฺปฏิสนฺธิ อภิญฺเญยฺยา.\csromanspacer{} คติ อภิญฺเญยฺยา, อคติ อภิญฺเญยฺยา.\csromanspacer{} นิพฺพตฺติ อภิญฺเญยฺยา, อนิพฺพตฺติ อภิญฺเญยฺยา.\csromanspacer{} อุปปตฺติ อภิญฺเญยฺยา, อนุปปตฺติ อภิญฺเญยฺยา.\csromanspacer{} ชาติ อภิญฺเญยฺยา, อชาติ อภิญฺเญยฺยา.\csromanspacer{} ชรา อภิญฺเญยฺยา, อชรา อภิญฺเญยฺยา.\csromanspacer{} พฺยาธิ อภิญฺเญยฺโย, อพฺยาธิ อภิญฺเญยฺโย.\csromanspacer{} มรณํ อภิญฺเญยฺยํ, อมตํ อภิญฺเญยฺยํ.\csromanspacer{} โสโก อภิญฺเญยฺโย, อโสโก อภิญฺเญยฺโย.\csromanspacer{} ปริเทโว อภิญฺเญยฺโย, อปริเทโว อภิญฺเญยฺโย.\csromanspacer{} อุปายาโส อภิญฺเญยฺโย, อนุปายาโส อภิญฺเญยฺโย. (30)}

\prose{อุปฺปาโท ทุกฺขนฺติ อภิญฺเญยฺยํ, ปวตฺตํ ทุกฺขนฺติ อภิญฺเญยฺยํ, นิมิตฺตํ ทุกฺขนฺติ อภิญฺเญยฺยํ, อายูหนา ทุกฺขนฺติ อภิญฺเญยฺยํ, ปฏิสนฺธิ ทุกฺขนฺติ อภิญฺเญยฺยํ \csromanfolio{13}คติ ทุกฺขนฺติ อภิญฺเญยฺยํ, นิพฺพตฺติ ทุกฺขนฺติ อภิญฺเญยฺยํ, อุปปตฺติ ทุกฺขนฺติ อภิญฺเญยฺยํ, ชาติ ทุกฺขนฺติ อภิญฺเญยฺยํ, ชรา ทุกฺขนฺติ อภิญฺเญยฺยํ, พฺยาธิ ทุกฺขนฺติ อภิญฺเญยฺยํ, มรณํ ทุกฺขนฺติ อภิญฺเญยฺยํ, โสโก ทุกฺขนฺติ อภิญฺเญยฺยํ, ปริเทโว ทุกฺขนฺติ อภิญฺเญยฺยํ, อุปายาโส ทุกฺขนฺติ อภิญฺเญยฺยํ. (15)}

\prose{อนุปฺปาโท สุขนฺติ อภิญฺเญยฺยํ, อปฺปวตฺตํ สุขนฺติ อภิญฺเญยฺยํ, อนิมิตฺตํ สุขนฺติ อภิญฺเญยฺยํ, อนายูหนา สุขนฺติ อภิญฺเญยฺยํ, อปฺปฏิสนฺธิ สุขนฺติ อภิญฺเญยฺยํ, อคติ สุขนฺติ อภิญฺเญยฺยํ, อนิพฺพตฺติ สุขนฺติ อภิญฺเญยฺยํ, อนุปปตฺติ สุขนฺติ อภิญฺเญยฺยํ, อชาติ สุขนฺติ อภิญฺเญยฺยํ, อชรา สุขนฺติ อภิญฺเญยฺยํ, อพฺยาธิ สุขนฺติ อภิญฺเญยฺยํ, อมตํ สุขนฺติ อภิญฺเญยฺยํ, อโสโก สุขนฺติ อภิญฺเญยฺยํ, อปริเทโว สุขนฺติ อภิญฺเญยฺยํ, อนุปายาโส สุขนฺติ อภิญฺเญยฺยํ. (15)}

\prose{อุปฺปาโท ทุกฺขํ, อนุปฺปาโท สุขนฺติ อภิญฺเญยฺยํ, ปวตฺตํ ทุกฺขํ, อปฺปวตฺตํ สุขนฺติ อภิญฺเญยฺยํ, นิมิตฺตํ ทุกฺขํ, อนิมิตฺตํ สุขนฺติ อภิญฺเญยฺยํ, อายูหนา ทุกฺขํ, อนายูหนา สุขนฺติ อภิญฺเญยฺยํ, ปฏิสนฺธิ ทุกฺขํ, อปฺปฏิสนฺธิ สุขนฺติ อภิญฺเญยฺยํ, คติ ทุกฺขํ, อคติ สุขนฺติ อภิญฺเญยฺยํ, นิพฺพตฺติ ทุกฺขํ, อนิพฺพตฺติ สุขนฺติ อภิญฺเญยฺยํ, อุปปตฺติ ทุกฺขํ, อนุปปตฺติ สุขนฺติ อภิญฺเญยฺยํ, ชาติ ทุกฺขํ, อชาติ สุขนฺติ อภิญฺเญยฺยํ, ชรา ทุกฺขํ, อชรา สุขนฺติ อภิญฺเญยฺยํ, พฺยาธิ ทุกฺขํ, อพฺยาธิ สุขนฺติ อภิญฺเญยฺยํ, มรณํ ทุกฺขํ, อมตํ สุขนฺติ อภิญฺเญยฺยํ, โสโก ทุกฺขํ, อโสโก สุขนฺติ อภิญฺเญยฺยํ, ปริเทโว ทุกฺขํ, อปริเทโว สุขนฺติ อภิญฺเญยฺยํ, อุปายาโส ทุกฺขํ, อนุปายาโส สุขนฺติ อภิญฺเญยฺยํ. (30)}

\prose{อุปฺปาโท ภยนฺติ อภิญฺเญยฺยํ, ปวตฺตํ ภยนฺติ อภิญฺเญยฺยํ, นิมิตฺตํ ภยนฺติ อภิญฺเญยฺยํ, อายูหนา ภยนฺติ อภิญฺเญยฺยํ, ปฏิสนฺธิ ภยนฺติ อภิญฺเญยฺยํ, คติ ภยนฺติ อภิญฺเญยฺยํ, นิพฺพตฺติ ภยนฺติ อภิญฺเญยฺยํ, อุปปตฺติ ภยนฺติ อภิญฺเญยฺยํ, ชาติ ภยนฺติ อภิญฺเญยฺยํ.\csromanspacer{} ชรา ภยนฺติ อภิญฺเญยฺยํ, พฺยาธิ ภยนฺติ อภิญฺเญยฺยํ, มรณํ ภยนฺติ อภิญฺเญยฺยํ, โสโก ภยนฺติ อภิญฺเญยฺยํ, ปริเทโว ภยนฺติ อภิญฺเญยฺยํ, อุปายาโส ภยนฺติ อภิญฺเญยฺยํ. (15)}

\prose{อนุปฺปาโท เขมนฺติ อภิญฺเญยฺยํ, อปฺปวตฺตํ เขมนฺติ อภิญฺเญยฺยํ, อนิมิตฺตํ เขมนฺติ อภิญฺเญยฺยํ, อนายูหนา เขมนฺติ อภิญฺเญยฺยํ, อปฺปฏิสนฺธิ \csromanfolio{14}เขมนฺติ อภิญฺเญยฺยํ, อคติ เขมนฺติ อภิญฺเญยฺยํ, อนิพฺพตฺติ เขมนฺติ อภิญฺเญยฺยํ, อนุปปตฺติ เขมนฺติ อภิญฺเญยฺยํ, อชาติ เขมนฺติ อภิญฺเญยฺยํ, อชรา เขมนฺติ อภิญฺเญยฺยํ, อพฺยาธิ เขมนฺติ อภิญฺเญยฺยํ, อมตํ เขมนฺติ อภิญฺเญยฺยํ, อโสโก เขมนฺติ อภิญฺเญยฺยํ, อปริเทโว เขมนฺติ อภิญฺเญยฺยํ, อนุปายาโส เขมนฺติ อภิญฺเญยฺยํ. (15)}

\prose{อุปฺปาโท ภยํ, อนุปฺปาโท เขมนฺติ อภิญฺเญยฺยํ, ปวตฺตํ ภยํ, อปฺปวตฺตํ เขมนฺติ อภิญฺเญยฺยํ, นิมิตฺตํ ภยํ, อนิมิตฺตํ เขมนฺติ อภิญฺเญยฺยํ, อายูหนา ภยํ, อนายูหนา เขมนฺติ อภิญฺเญยฺยํ, ปฏิสนฺธิ ภยํ, อปฺปฏิสนฺธิ เขมนฺติ อภิญฺเญยฺยํ, คติ ภยํ, อคติ เขมนฺติ อภิญฺเญยฺยํ, นิพฺพตฺติ ภยํ, อนิพฺพตฺติ เขมนฺติ อภิญฺเญยฺยํ, อุปปตฺติ ภยํ, อนุปปตฺติ เขมนฺติ อภิญฺเญยฺยํ, ชาติ ภยํ, อชาติ เขมนฺติ อภิญฺเญยฺยํ, ชรา ภยํ, อชรา เขมนฺติ อภิญฺเญยฺยํ, พฺยาธิ ภยํ, อพฺยาธิ เขมนฺติ อภิญฺเญยฺยํ, มรณํ ภยํ, อมตํ เขมนฺติ อภิญฺเญยฺยํ, โสโก ภยํ, อโสโก เขมนฺติ อภิญฺเญยฺยํ, ปริเทโว ภยํ อปริเทโว เขมนฺติ อภิญฺเญยฺยํ, อุปายาโส ภยํ, อนุปายาโส เขมนฺติ อภิญฺเญยฺยํ. (30)}

\prose{อุปฺปาโท สามิสนฺติ อภิญฺเญยฺยํ, ปวตฺตํ สามิสนฺติ อภิญฺเญยฺยํ, นิมิตฺตํ สามิสนฺติ อภิญฺเญยฺยํ, อายูหนา สามิสนฺติ อภิญฺเญยฺยํ, ปฏิสนฺธิ สามิสนฺติ อภิญฺเญยฺยํ, คติ สามิสนฺติ อภิญฺเญยฺยํ, นิพฺพตฺติ สามิสนฺติ อภิญฺเญยฺยํ, อุปปตฺติ สามิสนฺติ อภิญฺเญยฺยํ, ชาติ สามิสนฺติ อภิญฺเญยฺยํ, ชรา สามิสนฺติ อภิญฺเญยฺยํ, พฺยาธิ สามิสนฺติ อภิญฺเญยฺยํ, มรณํ สามิสนฺติ อภิญฺเญยฺยํ, โสโก สามิสนฺติ อภิญฺเญยฺยํ, ปริเทโว สามิสนฺติ อภิญฺเญยฺยํ, อุปายาโส สามิสนฺติ อภิญฺเญยฺยํ. (15)}

\prose{อนุปฺปาโท นิรามิสนฺติ อภิญฺเญยฺยํ, อปฺปวตฺตํ นิรามิสนฺติ อภิญฺเญยฺยํ, อนิมิตฺตํ นิรามิสนฺติ อภิญฺเญยฺยํ, อนายูหนา นิรามิสนฺติ อภิญฺเญยฺยํ, อปฺปฏิสนฺธิ นิรามิสนฺติ อภิญฺเญยฺยํ, อคติ นิรามิสนฺติ อภิญฺเญยฺยํ, อนิพฺพตฺติ นิรามิสนฺติ อภิญฺเญยฺยํ, อนุปปตฺติ นิรามิสนฺติ อภิญฺเญยฺยํ, อชาติ นิรามิมิสนฺติ อภิญฺเญยฺยํ, อชรา นิรามิสนฺติ อภิญฺเญยฺยํ, อพฺยาธิ นิรามิสนฺติ อภิญฺเญยฺยํ, อมตํ นิรามิสนฺติ อภิญฺเญยฺยํ, อโสโก นิรามิสนฺติ อภิญฺเญยฺยํ, อปริเทโว นิรามิสนฺติ อภิญฺเญยฺยํ, อนุปายาโส นิรามิสนฺติ อภิญฺเญยฺยํ. (15)}

\prose{\csromanfolio{15}อุปฺปาโท สามิสํ, อนุปฺปาโท นิรามิสนฺติ อภิญฺเญยฺยํ, ปวตฺตํ สามิสํ, อปฺปวตฺตํ นิรามิสนฺติ อภิญฺเญยฺยํ, นิมิตฺตํ สามิสํ, อนิมิตฺตํ นิรามิสนฺติ อภิญฺเญยฺยํ, อายูหนา สามิสํ, อนายูหนา นิรามิสนฺติ อภิญฺเญยฺยํ, ปฏิสนฺธิ สามิสํ, อปฺปฏิสนฺธิ นิรามิสนฺติ อภิญฺเญยฺยํ, คติ สามิสํ, อคติ นิรามิสนฺติ อภิญฺเญยฺยํ, นิพฺพตฺติ สามิสํ, อนิพฺพตฺติ นิรามิสนฺติ อภิญฺเญยฺยํ, อุปปตฺติ สามิสํ, อนุปปตฺติ นิรามิสนฺติ อภิญฺเญยฺยํ, ชาติ สามิสํ, อชาติ นิรามิสนฺติ อภิญฺเญยฺยํ, ชรา สามิสํ, อชรา นิรามิสนฺติ อภิญฺเญยฺยํ, พฺยาธิ สามิสํ, อพฺยาธิ นิรามิสนฺติ อภิญฺเญยฺยํ, มรณํ สามิสํ, อมตํ นิรามิสนฺติ อภิญฺเญยฺยํ, โสโก สามิสํ, อโสโก นิรามิสนฺติ อภิญฺเญยฺยํ, ปริเทโว สามิสํ, อปริเทโว นิรามิสนฺติ อภิญฺเญยฺยํ, อุปายาโส สามิสํ, อนุปายาโส นิรามิสนฺติ อภิญฺเญยฺยํ. (30)}

\prose{อุปฺปาโท สงฺขาราติ อภิญฺเญยฺยํ, ปวตฺตํ สงฺขาราติ อภิญฺเญยฺยํ, นิมิตฺตํ สงฺขาราติ อภิญฺเญยฺยํ, อายูหนา สงฺขาราติ อภิญฺเญยฺยํ, ปฏิสนฺธิ สงฺขาราติ อภิญฺเญยฺยํ, คติ สงฺขาราติ อภิญฺเญยฺยํ, นิพฺพตฺติ สงฺขาราติ อภิญฺเญยฺยํ, อุปปตฺติ สงฺขาราติ อภิญฺเญยฺยํ, ชาติ สงฺขาราติ อภิญฺเญยฺยํ, ชรา สงฺขาราติ อภิญฺเญยฺยํ, พฺยาธิ สงฺขาราติ อภิญฺเญยฺยํ, มรณํ สงฺขาราติ อภิญฺเญยฺยํ, โสโก สงฺขาราติ อภิญฺเญยฺยํ, ปริเทโว สงฺขาราติ อภิญฺเญยฺยํ, อุปายาโส สงฺขาราติ อภิญฺเญยฺยํ. (15)}

\prose{อนุปฺปาโท นิพฺพานนฺติ อภิญฺเญยฺยํ, อปฺปวตฺตํ นิพฺพานนฺติ อภิญฺเญยฺยํ, อนิมิตฺตํ นิพฺพานนฺติ อภิญฺเญยฺยํ, อนายูหนา นิพฺพานนฺติ อภิญฺเญยฺยํ, อปฺปฏิสนฺธิ นิพฺพานนฺติ อภิญฺเญยฺยํ, อคติ นิพฺพานนฺติ อภิญฺเญยฺยํ, อนิพฺพตฺติ นิพฺพานนฺติ อภิญฺเญยฺยํ, อนุปปตฺติ นิพฺพานนฺติ อภิญฺเญยฺยํ, อชาติ นิพฺพานนฺติ อภิญฺเญยฺยํ, อชรา นิพฺพานนฺติ อภิญฺเญยฺยํ, อพฺยาธิ นิพฺพานนฺติ อภิญฺเญยฺยํ, อมตํ นิพฺพานนฺติ อภิญฺเญยฺยํ, อโสโก นิพฺพานนฺติ อภิญฺเญยฺยํ.\csromanspacer{} อปริเทโว นิพฺพานนฺติ อภิญฺเญยฺยํ, อนุปายาโส นิพฺพานนฺติ อภิญฺเญยฺยํ. (15)}

\prose{อุปฺปาโท สงฺขารา, อนุปฺปาโท นิพฺพานนฺติ อภิญฺเญยฺยํ, ปวตฺตํ สงฺขารา, อปฺปวตฺตํ นิพฺพานนฺติ อภิญฺเญยฺยํ, นิมิตฺตํ สงฺขารา, อนิมิตฺตํ นิพฺพานนฺติ อภิญฺเญยฺยํ, อายูหนา สงฺขารา, อนายูหนา นิพฺพานนฺติ อภิญฺเญยฺยํ, ปฏิสนฺธิ สงฺขารา, อปฺปฏิสนฺธิ นิพฺพานนฺติ อภิญฺเญยฺยํ, คติ สงฺขารา, อคติ นิพฺพานนฺติ อภิญฺเญยฺยํ, นิพฺพตฺติ สงฺขารา, อนิพฺพตฺติ นิพฺพานนฺติ อภิญฺเญยฺยํ, อุปปตฺติ สงฺขารา, อนุปปตฺติ นิพฺพานนฺติ อภิญฺเญยฺยํ, ชาติ สงฺขารา, อชาติ นิพฺพานนฺติ อภิญฺเญยฺยํ, ชรา สงฺขารา, อชรา \csromanfolio{16}นิพฺพานนฺติ อภิญฺเญยฺยํ, พฺยาธิ สงฺขารา, อพฺยาธิ นิพฺพานนฺติ อภิญฺเญยฺยํ, มรณํ สงฺขารา, อมตํ นิพฺพานนฺติ อภิญฺเญยฺยํ, โสโก สงฺขารา, อโสโก นิพฺพานนฺติ อภิญฺเญยฺยํ, ปริเทโว สงฺขารา, อปริเทโว นิพฺพานนฺติ อภิญฺเญยฺยํ, อุปายาโส สงฺขารา, อนุปายาโส นิพฺพานนฺติ อภิญฺเญยฺยํ. (30) ปฐมภาณวาโร.}

\proseitem{11}{ปริคฺคหฏฺโฐ อภิญฺเญยฺโย, ปริวารฏฺโฐ อภิญฺเญยฺโย, ปริปูรฏฺโฐ\footnote{ปริปูรณฏฺโฐ (ก), ปริปูริฏฺโฐ (สี-ฏฺฐ)} อภิญฺเญยฺโย, เอกคฺคฏฺโฐ อภิญฺเญยฺโย, อวิกฺเขปฏฺโฐ อภิญฺเญยฺโย, ปคฺคหฏฺโฐ อภิญฺเญยฺโย, อวิสารฏฺโฐ อภิญฺเญยฺโย, อนาวิลฏฺโฐ อภิญฺเญยฺโย, อนิญฺชนฏฺโฐ อภิญฺเญยฺโย, เอกตฺตุปฏฺฐานวเสน จิตฺตสฺส ฐิตฏฺโฐ อภิญฺเญยฺโย, อารมฺมณฏฺโฐ อภิญฺเญยฺโย, โคจรฏฺโฐ อภิญฺเญยฺโย, ปหานฏฺโฐ อภิญฺเญยฺโย, ปริจฺจาคฏฺโฐ อภิญฺเญยฺโย, วุฏฺฐานฏฺโฐ อภิญฺเญยฺโย, วิวฏฺฏนฏฺโฐ\footnote{นิวตฺตนฏฺโฐ (ก)} อภิญฺเญยฺโย, สนฺตฏฺโฐ อภิญฺเญยฺโย, ปณีตฏฺโฐ อภิญฺเญยฺโย, วิมุตฺตฏฺโฐ อภิญฺเญยฺโย, อนาสวฏฺโฐ อภิญฺเญยฺโย, ตรณฏฺโฐ อภิญฺเญยฺโย, อนิมิตฺตฏฺโฐ อภิญฺเญยฺโย, อปฺปณิหิตฏฺโฐ อภิญฺเญยฺโย, สุญฺญตฏฺโฐ อภิญฺเญยฺโย, เอกรสฏฺโฐ อภิญฺเญยฺโย, อนติวตฺตนฏฺโฐ อภิญฺเญยฺโย, ยุคนทฺธฏฺโฐ\footnote{ยุคนนฺธนฏฺโฐ (ก)} อภิญฺเญยฺโย, นิยฺยานฏฺโฐ อภิญฺเญยฺโย, เหตุฏฺโฐ\footnote{เหตฏฺโฐ (สฺยา)} อภิญฺเญยฺโย, ทสฺสนฏฺโฐ อภิญฺเญยฺโย, อาธิปเตยฺยฏฺโฐ อภิญฺเญยฺโย. (31)}

\proseitem{12}{สมถสฺส อวิกฺเขปฏฺโฐ อภิญฺเญยฺโย, วิปสฺสนาย อนุปสฺสนฏฺโฐ อภิญฺเญยฺโย, สมถวิปสฺสนานํ เอกรสฏฺโฐ อภิญฺเญยฺโย, ยุคนทฺธสฺส อนติวตฺตนฏฺโฐ อภิญฺเญยฺโย. (4)}

\prose{สิกฺขาย สมาทานฏฺโฐ อภิญฺเญยฺโย, อารมฺมณสฺส โคจรฏฺโฐ อภิญฺเญยฺโย, ลีนสฺส จิตฺตสฺส ปคฺคหฏฺโฐ อภิญฺเญยฺโย, อุทฺธตสฺส จิตฺตสฺส นิคฺคหฏฺโฐ อภิญฺเญยฺโย, อุโภวิสุทฺธานํ อชฺฌุเปกฺขนฏฺโฐ อภิญฺเญยฺโย, วิเสสาธิคมฏฺโฐ อภิญฺเญยฺโย, อุตฺตริปฏิเวธฏฺโฐ อภิญฺเญยฺโย, สจฺจาภิสมยฏฺโฐ อภิญฺเญยฺโย, นิโรเธ ปติฏฺฐาปกฏฺโฐ อภิญฺเญยฺโย. (9)}

\prose{\csromanfolio{17}สทฺธินฺทฺริยสฺส อธิโมกฺขฏฺโฐ อภิญฺเญยฺโย, วีริยินฺทฺริยสฺส ปคฺคหฏฺโฐ อภิญฺเญยฺโย, สตินฺทฺริยสฺส อุปฏฺฐานฏฺโฐ อภิญฺเญยฺโย, สมาธินฺทฺริยสฺส อวิกฺเขปฏฺโฐ อภิญฺเญยฺโย, ปญฺญินฺทฺริยสฺส ทสฺสนฏฺโฐ อภิญฺเญยฺโย. (5)}

\prose{สทฺธาพลสฺส อสฺสทฺธิเย อกมฺปิยฏฺโฐ อภิญฺเญยฺโย, วีริยพลสฺส โกสชฺเช อกมฺปิยฏฺโฐ อภิญฺเญยฺโย, สติพลสฺส ปมาเท อกมฺปิยฏฺโฐ อภิญฺเญยฺโย, สมาธิพลสฺส อุทฺธจฺเจ อกมฺปิยฏฺโฐ อภิญฺเญยฺโย, ปญฺญาพลสฺส อวิชฺชาย อกมฺปิยฏฺโฐ อภิญฺเญยฺโย. (5)}

\prose{สติสมฺโพชฺฌงฺคสฺส อุปฏฺฐานฏฺโฐ อภิญฺเญยฺโย, ธมฺมวิจยสมฺโพชฺฌงฺคสฺส ปวิจยฏฺโฐ อภิญฺเญยฺโย, วีริยสมฺโพชฺฌงฺคสฺส ปคฺคหฏฺโฐ อภิญฺเญยฺโย, ปีติสมฺโพชฺฌงฺคสฺส ผรณฏฺโฐ อภิญฺเญยฺโย, ปสฺสทฺธิสมฺโพชฺฌงฺคสฺส อุปสมฏฺโฐ อภิญฺเญยฺโย, สมาธิสมฺโพชฺฌงฺคสฺส อวิกฺเขปฏฺโฐ อภิญฺเญยฺโย, อุเปกฺขาสมฺโพชฺฌงฺคสฺส ปฏิสงฺขานฏฺโฐ อภิญฺเญยฺโย. (7)}

\prose{สมฺมาทิฏฺฐิยา ทสฺสนฏฺโฐ อภิญฺเญยฺโย, สมฺมาสงฺกปฺปสฺส อภินิโรปนฏฺโฐ อภิญฺเญยฺโย, สมฺมาวาจาย ปริคฺคหฏฺโฐ อภิญฺเญยฺโย, สมฺมากมฺมนฺตสฺส สมุฏฺฐานฏฺโฐ อภิญฺเญยฺโย, สมฺมา-อาชีวสฺส โวทานฏฺโฐ อภิญฺเญยฺโย, สมฺมาวายามสฺส ปคฺคหฏฺโฐ อภิญฺเญยฺโย, สมฺมาสติยา อุปฏฺฐานฏฺโฐ อภิญฺเญยฺโย, สมฺมาสมาธิสฺส อวิกฺเขปฏฺโฐ อภิญฺเญยฺโย. (8)}

\proseitem{13}{อินฺทฺริยานํ อาธิปเตยฺยฏฺโฐ อภิญฺเญยฺโย, พลานํ อกมฺปิยฏฺโฐ อภิญฺเญยฺโย, โพชฺฌงฺคานํ นิยฺยานฏฺโฐ อภิญฺเญยฺโย, มคฺคสฺส เหตุฏฺโฐ อภิญฺเญยฺโย, สติปฏฺฐานานํ อุปฏฺฐานฏฺโฐ อภิญฺเญยฺโย, สมฺมปฺปธานานํ ปทหนฏฺโฐ อภิญฺเญยฺโย, อิทฺธิปาทานํ อิชฺฌนฏฺโฐ อภิญฺเญยฺโย, สจฺจานํ ตถฏฺโฐ อภิญฺเญยฺโย, ปโยคานํ\footnote{มคฺคานํ (สพฺพตฺถ) อฏฺฐกถา ปสฺสิตพฺพา.} ปฏิปฺปสฺสทฺธฏฺโฐ อภิญฺเญยฺโย, ผลานํ สจฺฉิกิริยฏฺโฐ อภิญฺเญยฺโย. (10)}

\prose{วิตกฺกสฺส อภินิโรปนฏฺโฐ อภิญฺเญยฺโย, วิจารสฺส อุปวิจารฏฺโฐ อภิญฺเญยฺโย, ปีติยา ผรณฏฺโฐ อภิญฺเญยฺโย, สุขสฺส อภิสนฺทนฏฺโฐ อภิญฺเญยฺโย, จิตฺตสฺส เอกคฺคฏฺโฐ อภิญฺเญยฺโย. (5)}

\prose{อาวชฺชนฏฺโฐ อภิญฺเญยฺโย, วิชานนฏฺโฐ อภิญฺเญยฺโย, ปชานนฏฺโฐ อภิญฺเญยฺโย, สญฺชานนฏฺโฐ อภิญฺเญยฺโย, เอโกทฏฺโฐ \csromanfolio{18}อภิญฺเญยฺโย, อภิญฺญาย ญาตฏฺโฐ อภิญฺเญยฺโย, ปริญฺญาย ตีรณฏฺโฐ อภิญฺเญยฺโย, ปหานสฺส ปริจฺจาคฏฺโฐ อภิญฺเญยฺโย, ภาวนาย เอกรสฏฺโฐ อภิญฺเญยฺโย, สจฺฉิกิริยาย ผสฺสนฏฺโฐ อภิญฺเญยฺโย, ขนฺธานํ ขนฺธฏฺโฐ อภิญฺเญยฺโย, ธาตูนํ ธาตุฏฺโฐ\footnote{ธาตฏฺโฐ (สฺยา)} อภิญฺเญยฺโย, อายตนานํ อายตนฏฺโฐ อภิญฺเญยฺโย, สงฺขตานํ สงฺขตฏฺโฐ อภิญฺเญยฺโย, อสงฺขตสฺส อสงฺขตฏฺโฐ อภิญฺเญยฺโย. (15)}

\proseitem{14}{จิตฺตฏฺโฐ อภิญฺเญยฺโย, จิตฺตานนฺตริยฏฺโฐ อภิญฺเญยฺโย, จิตฺตสฺส วุฏฺฐานฏฺโฐ อภิญฺเญยฺโย, จิตฺตสฺส วิวฏฺฏนฏฺโฐ อภิญฺเญยฺโย, จิตฺตสฺส เหตุฏฺโฐ อภิญฺเญยฺโย, จิตฺตสฺส ปจฺจยฏฺโฐ อภิญฺเญยฺโย, จิตฺตสฺส วตฺถุฏฺโฐ\footnote{วตฺถฏฺโฐ (สฺยา)} อภิญฺเญยฺโย, จิตฺตสฺส ภูมฏฺโฐ\footnote{ภุมฺมฏฺโฐ (สฺยา, สี-ฏฺฐ) (เอกตฺเต อุปนิพนฺธนฏฺโฐ อภิญฺเญยฺโย) (ก) อฏฺฐกถา ปสฺสิตพฺพา.} อภิญฺเญยฺโย, จิตฺตสฺส อารมฺมณฏฺโฐ อภิญฺเญยฺโย, จิตฺตสฺส โคจรฏฺโฐ อภิญฺเญยฺโย, จิตฺตสฺส จริยฏฺโฐ อภิญฺเญยฺโย, จิตฺตสฺส คตฏฺโฐ อภิญฺเญยฺโย, จิตฺตสฺส อภินีหารฏฺโฐ อภิญฺเญยฺโย, จิตฺตสฺส นิยฺยานฏฺโฐ อภิญฺเญยฺโย, จิตฺตสฺส นิสฺสรณฏฺโฐ อภิญฺเญยฺโย. (15)}

\proseitem{15}{เอกตฺเต อาวชฺชนฏฺโฐ อภิญฺเญยฺโย, เอกตฺเต วิชานนฏฺโฐ อภิญฺเญยฺโย, เอกตฺเต ปชานนฏฺโฐ อภิญฺเญยฺโย, เอกตฺเต สญฺชานนฏฺโฐ อภิญฺเญยฺโย, เอกตฺเต เอโกทฏฺโฐ อภิญฺเญยฺโย, ( ) เอกตฺเต ปกฺขนฺทนฏฺโฐ อภิญฺเญยฺโย, เอกตฺเต ปสีทนฏฺโฐ อภิญฺเญยฺโย, เอกตฺเต สนฺติฏฺฐนฏฺโฐ อภิญฺเญยฺโย, เอกตฺเต วิมุจฺจนฏฺโฐ อภิญฺเญยฺโย, เอกตฺเต “เอตํ สนฺตนฺ”ติ ปสฺสนฏฺโฐ อภิญฺเญยฺโย, เอกตฺเต ยานีกตฏฺโฐ อภิญฺเญยฺโย, เอกตฺเต วตฺถุกตฏฺโฐ อภิญฺเญยฺโย, เอกตฺเต อนุฏฺฐิตฏฺโฐ อภิญฺเญยฺโย, เอกตฺเต ปริจิตฏฺโฐ อภิญฺเญยฺโย, เอกตฺเต สุสมารทฺธฏฺโฐ อภิญฺเญยฺโย, เอกตฺเต ปริคฺคหฏฺโฐ อภิญฺเญยฺโย, เอกตฺเต ปริวารฏฺโฐ อภิญฺเญยฺโย, เอกตฺเต ปริปูรฏฺโฐ อภิญฺเญยฺโย, เอกตฺเต สโมธานฏฺโฐ อภิญฺเญยฺโย, เอกตฺเต อธิฏฺฐานฏฺโฐ อภิญฺเญยฺโย, เอกตฺเต อาเสวนฏฺโฐ อภิญฺเญยฺโย, เอกตฺเต ภาวนฏฺโฐ อภิญฺเญยฺโย, เอกตฺเต พหุลีกมฺมฏฺโฐ อภิญฺเญยฺโย, เอกตฺเต สุสมุคฺคตฏฺโฐ อภิญฺเญยฺโย, เอกตฺเต สุวิมุตฺตฏฺโฐ \csromanfolio{19}อภิญฺเญยฺโย, เอกตฺเต พุชฺฌนฏฺโฐ อภิญฺเญยฺโย, เอกตฺเต อนุพุชฺฌนฏฺโฐ อภิญฺเญยฺโย, เอกตฺเต ปฏิพุชฺฌนฏฺโฐ อภิญฺเญยฺโย, เอกตฺเต สมฺพุชฺฌนฏฺโฐ อภิญฺเญยฺโย, เอกตฺเต โพธนฏฺโฐ อภิญฺเญยฺโย, เอกตฺเต อนุโพธนฏฺโฐ อภิญฺเญยฺโย, เอกตฺเต ปฏิโพธนฏฺโฐ อภิญฺเญยฺโย, เอกตฺเต สมฺโพธนฏฺโฐ อภิญฺเญยฺโย, เอกตฺเต โพธิปกฺขิยฏฺโฐ อภิญฺเญยฺโย, เอกตฺเต อนุโพธิปกฺขิยฏฺโฐ อภิญฺเญยฺโย, เอกตฺเต ปฏิโพธิปกฺขิยฏฺโฐ อภิญฺเญยฺโย, เอกตฺเต สมฺโพธิปกฺขิยฏฺโฐ อภิญฺเญยฺโย, เอกตฺเต โชตนฏฺโฐ อภิญฺเญยฺโย, เอกตฺเต อุชฺโชตนฏฺโฐ อภิญฺเญยฺโย, เอกตฺเต อนุโชตนฏฺโฐ อภิญฺเญยฺโย, เอกตฺเต ปฏิโชตนฏฺโฐ อภิญฺเญยฺโย, เอกตฺเต สญฺโชตนฏฺโฐ อภิญฺเญยฺโย. (42)}

\proseitem{16}{ปตาปนฏฺโฐ\footnote{ปกาสนฏฺโฐ (ก)} อภิญฺเญยฺโย, วิโรจนฏฺโฐ อภิญฺเญยฺโย, กิเลสานํ สนฺตาปนฏฺโฐ อภิญฺเญยฺโย, อมลฏฺโฐ อภิญฺเญยฺโย, วิมลฏฺโฐ อภิญฺเญยฺโย, นิมฺมลฏฺโฐ อภิญฺเญยฺโย, สมฏฺโฐ อภิญฺเญยฺโย, สมยฏฺโฐ อภิญฺเญยฺโย, วิเวกฏฺโฐ อภิญฺเญยฺโย, วิเวกจริยฏฺโฐ อภิญฺเญยฺโย, วิราคฏฺโฐ อภิญฺเญยฺโย, วิราคจริยฏฺโฐ อภิญฺเญยฺโย, นิโรธฏฺโฐ อภิญฺเญยฺโย, นิโรธจริยฏฺโฐ อภิญฺเญยฺโย, โวสคฺคฏฺโฐ\footnote{โปสฺสคฺคฏฺโฐ (สฺยา, ก)} อภิญฺเญยฺโย, โวสคฺคจริยฏฺโฐ อภิญฺเญยฺโย วิมุตฺตฏฺโฐ อภิญฺเญยฺโย, วิมุตฺติจริยฏฺโฐ อภิญฺเญยฺโย. (18)}

\prose{ฉนฺทฏฺโฐ อภิญฺเญยฺโย, ฉนฺทสฺส มูลฏฺโฐ อภิญฺเญยฺโย, ฉนฺทสฺส ปาทฏฺโฐ อภิญฺเญยฺโย, ฉนฺทสฺส ปธานฏฺโฐ อภิญฺเญยฺโย, ฉนฺทสฺส อิชฺฌนฏฺโฐ อภิญฺเญยฺโย, ฉนฺทสฺส อธิโมกฺขฏฺโฐ อภิญฺเญยฺโย, ฉนฺทสฺส ปคฺคหฏฺโฐ อภิญฺเญยฺโย, ฉนฺทสฺส อุปฏฺฐานฏฺโฐ อภิญฺเญยฺโย, ฉนฺทสฺส อวิกฺเขปฏฺโฐ อภิญฺเญยฺโย, ฉนฺทสฺส ทสฺสนฏฺโฐ อภิญฺเญยฺโย. (10)}

\prose{วีริยฏฺโฐ อภิญฺเญยฺโย, วีริยสฺส มูลฏฺโฐ อภิญฺเญยฺโย, วีริยสฺส ปาทฏฺโฐ อภิญฺเญยฺโย, วีริยสฺส ปธานฏฺโฐ อภิญฺเญยฺโย, วีริยสฺส อิชฺฌนฏฺโฐ อภิญฺเญยฺโย, วีริยสฺส อธิโมกฺขฏฺโฐ อภิญฺเญยฺโย, วีริยสฺส ปคฺคหฏฺโฐ อภิญฺเญยฺโย, วีริยสฺส อุปฏฺฐานฏฺโฐ อภิญฺเญยฺโย, วีริยสฺส อวิกฺเขปฏฺโฐ อภิญฺเญยฺโย, วีริยสฺส ทสฺสนฏฺโฐ อภิญฺเญยฺโย. (10)}

\prose{\csromanfolio{20}จิตฺตฏฺโฐ อภิญฺเญยฺโย, จิตฺตสฺส มูลฏฺโฐ อภิญฺเญยฺโย, จิตฺตสฺส ปาทฏฺโฐ อภิญฺเญยฺโย, จิตฺตสฺส ปธานฏฺโฐ อภิญฺเญยฺโย, จิตฺตสฺส อิชฺฌนฏฺโฐ อภิญฺเญยฺโย, จิตฺตสฺส อธิโมกฺขฏฺโฐ อภิญฺเญยฺโย, จิตฺตสฺส ปคฺคหฏฺโฐ อภิญฺเญยฺโย, จิตฺตสฺส อุปฏฺฐานฏฺโฐ อภิญฺเญยฺโย, จิตฺตสฺส อวิกฺเขปฏฺโฐ อภิญฺเญยฺโย, จิตฺตสฺส ทสฺสนฏฺโฐ อภิญฺเญยฺโย. (10)}

\prose{วีมํสฏฺโฐ อภิญฺเญยฺโย, วีมํสาย มูลฏฺโฐ อภิญฺเญยฺโย, วีมํสาย ปาทฏฺโฐ อภิญฺเญยฺโย, วีมํสาย ปธานฏฺโฐ อภิญฺเญยฺโย, วีมํสาย อิชฺฌนฏฺโฐ อภิญฺเญยฺโย, วีมํสาย อธิโมกฺขฏฺโฐ อภิญฺเญยฺโย, วีมํสาย ปคฺคหฏฺโฐ อภิญฺเญยฺโย, วีมํสาย อุปฏฺฐานฏฺโฐ อภิญฺเญยฺโย, วีมํสาย อวิกฺเขปฏฺโฐ อภิญฺเญยฺโย, วีมํสาย ทสฺสนฏฺโฐ อภิญฺเญยฺโย. (10)}

\proseitem{17}{ทุกฺขฏฺโฐ อภิญฺเญยฺโย\csromandash{}ทุกฺขสฺส ปีฬนฏฺโฐ อภิญฺเญยฺโย, ทุกฺขสฺส สงฺขตฏฺโฐ อภิญฺเญยฺโย, ทุกฺขสฺส สนฺตาปฏฺโฐ อภิญฺเญยฺโย, ทุกฺขสฺส วิปริณามฏฺโฐ อภิญฺเญยฺโย.\csromanspacer{}\csromanspacer{}สมุทยฏฺโฐ อภิญฺเญยฺโย\csromandash{} สมุทยสฺส อายูหนฏฺโฐ\footnote{อายุหนฏฺโฐ (สฺยา)} อภิญฺเญยฺโย, สมุทยสฺส นิทานฏฺโฐ อภิญฺเญยฺโย, สมุทยสฺส สญฺโญคฏฺโฐ อภิญฺเญยฺโย, สมุทยสฺส ปลิโพธฏฺโฐ อภิญฺเญยฺโย.\csromanspacer{}\csromanspacer{}นิโรธฏฺโฐ อภิญฺเญยฺโย\csromandash{}นิโรธสฺส นิสฺสรณฏฺโฐ อภิญฺเญยฺโย, นิโรธสฺส วิเวกฏฺโฐ อภิญฺเญยฺโย, นิโรธสฺส อสงฺขตฏฺโฐ อภิญฺเญยฺโย, นิโรธสฺส อมตฏฺโฐ อภิญฺเญยฺโย.\csromanspacer{}\csromanspacer{}มคฺคฏฺโฐ อภิญฺเญยฺโย\csromandash{}มคฺคสฺส นิยฺยานฏฺโฐ อภิญฺเญยฺโย, มคฺคสฺส เหตุฏฺโฐ อภิญฺเญยฺโย, มคฺคสฺส ทสฺสนฏฺโฐ อภิญฺเญยฺโย, มคฺคสฺส อาธิปเตยฺยฏฺโฐ อภิญฺเญยฺโย. (16)}

\prose{ตถฏฺโฐ อภิญฺเญยฺโย, อนตฺตฏฺโฐ อภิญฺเญยฺโย, สจฺจฏฺโฐ อภิญฺเญยฺโย, ปฏิเวธฏฺโฐ อภิญฺเญยฺโย, อภิชานนฏฺโฐ อภิญฺเญยฺโย, ปริชานนฏฺโฐ อภิญฺเญยฺโย, ธมฺมฏฺโฐ อภิญฺเญยฺโย, ธาตุฏฺโฐ อภิญฺเญยฺโย, ญาตฏฺโฐ อภิญฺเญยฺโย, สจฺฉิกิริยฏฺโฐ อภิญฺเญยฺโย, ผสฺสนฏฺโฐ อภิญฺเญยฺโย, อภิสมยฏฺโฐ อภิญฺเญยฺโย. (12)}

\proseitem{18}{เนกฺขมฺมํ อภิญฺเญยฺยํ, อพฺยาปาโท อภิญฺเญยฺโย, อาโลกสญฺญา อภิญฺเญยฺยา, อวิกฺเขโป อภิญฺเญยฺโย, ธมฺมววตฺถานํ อภิญฺเญยฺยํ, ญาณํ อภิญฺเญยฺยํ, ปาโมชฺชํ\footnote{ปามุชฺชํ (สฺยา)} อภิญฺเญยฺยํ. (7)}

\prose{\csromanfolio{21}ปฐมํ ฌานํ อภิญฺเญยฺยํ, ทุติยํ ฌานํ อภิญฺเญยฺยํ, ตติยํ ฌานํ อภิญฺเญยฺยํ, จตุตฺถํ ฌานํ อภิญฺเญยฺยํ, อากาสานญฺจายตนสมาปตฺติ อภิญฺเญยฺยา, วิญฺญาณญฺจายตนสมาปตฺติ อภิญฺเญยฺยา, อากิญฺจญฺญายตนสมาปตฺติ อภิญฺเญยฺยา, เนวสญฺญานาสญฺญายตนสมาปตฺติ อภิญฺเญยฺยา. (8)}

\prose{อนิจฺจานุปสฺสนา อภิญฺเญยฺยา, ทุกฺขานุปสฺสนา อภิญฺเญยฺยา, อนตฺตานุปสฺสนา อภิญฺเญยฺยา, นิพฺพิทานุปสฺสนา อภิญฺเญยฺยา, วิราคานุปสฺสนา อภิญฺเญยฺยา, นิโรธานุปสฺสนา อภิญฺเญยฺยา, ปฏินิสฺสคฺคานุปสฺสนา อภิญฺเญยฺยา, ขยานุปสฺสนา อภิญฺเญยฺยา, วยานุปสฺสนา อภิญฺเญยฺยา, วิปริณามานุปสฺสนา อภิญฺเญยฺยา, อนิมิตฺตานุปสฺสนา อภิญฺเญยฺยา, อปฺปณิหิตานุปสฺสนา อภิญฺเญยฺยา, สุญฺญตานุปสฺสนา อภิญฺเญยฺยา, อธิปญฺญาธมฺมวิปสฺสนา อภิญฺเญยฺยา, ยถาภูตญาณทสฺสนํ อภิญฺเญยฺยํ, อาทีนวานุปสฺสนา อภิญฺเญยฺยา, ปฏิสงฺขานุปสฺสนา อภิญฺเญยฺยา, วิวฏฺฏนานุปสฺสนา อภิญฺเญยฺยา. (18)}

\proseitem{19}{โสตาปตฺติมคฺโค อภิญฺเญยฺโย, โสตาปตฺติผลสมาปตฺติ อภิญฺเญยฺยา, สกทาคามิมคฺโค อภิญฺเญยฺโย, สกทาคามิผลสมาปตฺติ อภิญฺเญยฺยา, อนาคามิมคฺโค อภิญฺเญยฺโย, อนาคามิผลสมาปตฺติ อภิญฺเญยฺยา อรหตฺตมคฺโค อภิญฺเญยฺโย, อรหตฺตผลสมาปตฺติ อภิญฺเญยฺยา. (8)}

\prose{อธิโมกฺขฏฺเฐน สทฺธินฺทฺริยํ อภิญฺเญยฺยํ, ปคฺคหฏฺเฐน วีริยินฺทฺริยํ อภิญฺเญยฺยํ, อุปฏฺฐานฏฺเฐน สตินฺทฺริยํ อภิญฺเญยฺยํ, อวิกฺเขปฏฺเฐน สมาธินฺทฺริยํ อภิญฺเญยฺยํ, ทสฺสนฏฺเฐน ปญฺญินฺทฺริยํ อภิญฺเญยฺยํ.}

\prose{อสฺสทฺธิเย อกมฺปิยฏฺเฐน สทฺธาพลํ อภิญฺเญยฺยํ, โกสชฺเช อกมฺปิยฏฺเฐน วีริยพลํ อภิญฺเญยฺยํ, ปมาเท อกมฺปิยฏฺเฐน สติพลํ อภิญฺเญยฺยํ, อุทฺธจฺเจ อกมฺปิยฏฺเฐน สมาธิพลํ อภิญฺเญยฺยํ, อวิชฺชาย อกมฺปิยฏฺเฐน ปญฺญาพลํ อภิญฺเญยฺยํ.}

\prose{อุปฏฺฐานฏฺเฐน สติสมฺโพชฺฌงฺโค อภิญฺเญยฺโย, ปวิจยฏฺเฐน ธมฺมวิจยสมฺโพชฺฌงฺโค อภิญฺเญยฺโย, ปคฺคหฏฺเฐน วีริยสมฺโพชฺฌงฺโค อภิญฺเญยฺโย, ผรณฏฺเฐน ปีติสมฺโพชฺฌงฺโค อภิญฺเญยฺโย, อุปสมฏฺเฐน ปสฺสทฺธิสมฺโพชฺฌงฺโค อภิญฺเญยฺโย, อวิกฺเขปฏฺเฐน สมาธิสมฺโพชฺฌงฺโค อภิญฺเญยฺโย, ปฏิสงฺขานฏฺเฐน อุเปกฺขาสมฺโพชฺฌงฺโค อภิญฺเญยฺโย.}

\prose{\csromanfolio{22}ทสฺสนฏฺเฐน สมฺมาทิฏฺฐิ อภิญฺเญยฺยา, อภินิโรปนฏฺเฐน สมฺมาสงฺกปฺโป อภิญฺเญยฺโย, ปริคฺคหฏฺเฐน สมฺมาวาจา อภิญฺเญยฺยา, สมุฏฺฐานฏฺเฐน สมฺมากมฺมนฺโต อภิญฺเญยฺโย, โวทานฏฺเฐน สมฺมา-อาชีโว อภิญฺเญยฺโย, ปคฺคหฏฺเฐน สมฺมาวายาโม อภิญฺเญยฺโย, อุปฏฺฐานฏฺเฐน สมฺมาสติ อภิญฺเญยฺยา, อวิกฺเขปฏฺเฐน สมฺมาสมาธิ อภิญฺเญยฺโย.}

\prose{อาธิปเตยฺยฏฺเฐน อินฺทฺริยา อภิญฺเญยฺยา, อกมฺปิยฏฺเฐน พลา อภิญฺเญยฺยา, นิยฺยานฏฺเฐน โพชฺฌงฺคา อภิญฺเญยฺยา, เหตุฏฺเฐน มคฺโค อภิญฺเญยฺโย, อุปฏฺฐานฏฺเฐน สติปฏฺฐานา อภิญฺเญยฺยา, ปทหนฏฺเฐน สมฺมปฺปธานา อภิญฺเญยฺยา, อิชฺฌนฏฺเฐน อิทฺธิปาทา อภิญฺเญยฺยา, ตถฏฺเฐน สจฺจา อภิญฺเญยฺยา. (33)}

\prose{อวิกฺเขปฏฺเฐน สมโถ อภิญฺเญยฺโย, อนุปสฺสนฏฺเฐน วิปสฺสนา อภิญฺเญยฺยา, เอกรสฏฺเฐน สมถวิปสฺสนา อภิญฺเญยฺยา, อนติวตฺตนฏฺเฐน ยุคนทฺธํ อภิญฺเญยฺยํ. (4)}

\prose{สํวรฏฺเฐน สีลวิสุทฺธิ อภิญฺเญยฺยา, อวิกฺเขปฏฺเฐน จิตฺตวิสุทฺธิ อภิญฺเญยฺยา, ทสฺสนฏฺเฐน ทิฏฺฐิวิสุทฺธิ อภิญฺเญยฺยา, มุตฺตฏฺเฐน วิโมกฺโข อภิญฺเญยฺโย, ปฏิเวธฏฺเฐน วิชฺชา อภิญฺเญยฺยา, ปริจฺจาคฏฺเฐน วิมุตฺติ อภิญฺเญยฺยา, สมุจฺเฉทฏฺเฐน ขเย ญาณํ อภิญฺเญยฺยํ, ปฏิปฺปสฺสทฺธฏฺเฐน อนุปฺปาเท ญาณํ อภิญฺเญยฺยํ. (8)}

\proseitem{20}{ฉนฺโท มูลฏฺเฐน อภิญฺเญยฺโย, มนสิกาโร สมุฏฺฐานฏฺเฐน อภิญฺเญยฺโย, ผสฺโส สโมธานฏฺเฐน อภิญฺเญยฺโย, เวทนา สโมสรณฏฺเฐน อภิญฺเญยฺยา, สมาธิ ปมุขฏฺเฐน อภิญฺเญยฺโย, สติ อาธิปเตยฺยฏฺเฐน อภิญฺเญยฺยา, ปญฺญา ตทุตฺตรฏฺเฐน อภิญฺเญยฺยา, วิมุตฺติ สารฏฺเฐน อภิญฺเญยฺยา, อมโตคธํ นิพฺพานํ ปริโยสานฏฺเฐน อภิญฺเญยฺยํ. (9)}

\prose{เย เย ธมฺมา อภิญฺญาตา โหนฺติ, เต เต ธมฺมา ญาตา โหนฺติ, ตํญาตฏฺเฐน ญาณํ, ปชานนฏฺเฐน ปญฺญา, เตน วุจฺจติ “อิเม ธมฺมา อภิญฺเญยฺยาติ โสตาวธานํ, ตํปชานนา ปญฺญา สุตมเย ญาณนฺ”ติ. (1)}

\vspace{\tipitakaparskip}
\csromancenter{ทุติยภาณวาโร.}
\proseitem{21}{\csromanfolio{23}กถํ “อิเม ธมฺมา ปริญฺเญยฺยา”ติ โสตาวธานํ, ตํปชานนา ปญฺญา สุตมเย ญาณํ\csromandash{}}

\prose{เอโก ธมฺโม ปริญฺเญยฺโย, ผสฺโส สาสโว อุปาทานิโย.\csromanspacer{} เทฺว ธมฺมา ปริญฺเญยฺยา, นามญฺจ รูปญฺจ.\csromanspacer{} ตโย ธมฺมา ปริญฺเญยฺยา, ติสฺโส เวทนา.\csromanspacer{} จตฺตาโร ธมฺมา ปริญฺเญยฺยา, จตฺตาโร อาหารา.\csromanspacer{} ปญฺจ ธมฺมา ปริญฺเญยฺยา, ปญฺจุปาทานกฺขนฺธา.\csromanspacer{} ฉ ธมฺมา ปริญฺเญยฺยา, ฉ อชฺฌตฺติกานิ อายตนานิ.\csromanspacer{} สตฺต ธมฺมา ปริญฺเญยฺยา, สตฺต วิญฺญาณฏฺฐิติโย.\csromanspacer{} อฏฺฐ ธมฺมา ปริญฺเญยฺยา, อฏฺฐ โลกธมฺมา.\csromanspacer{} นว ธมฺมา ปริญฺเญยฺยา, นว สตฺตาวาสา.\csromanspacer{} ทส ธมฺมา ปริญฺเญยฺยา, ทสายตนานิ. (10)}

\prose{สพฺพํ ภิกฺขเว ปริญฺเญยฺยํ.\csromanspacer{} กิญฺจ ภิกฺขเว สพฺพํ ปริญฺเญยฺยํ? จกฺขุ ภิกฺขเว ปริญฺเญยฺยํ, รูปา ปริญฺเญยฺยา, จกฺขุวิญฺญาณํ ปริญฺเญยฺยํ, จกฺขุสมฺผสฺโส ปริญฺเญยฺโย, ยมฺปิทํ จกฺขุสมฺผสฺสปจฺจยา อุปฺปชฺชติ เวทยิตํ สุขํ วา ทุกฺขํ วา อทุกฺขมสุขํ วา, ตมฺปิ ปริญฺเญยฺยํ.\csromanspacer{} โสตํ ปริญฺเญยฺยํ, สทฺทา ปริญฺเญยฺยา ฯเปฯ.\csromanspacer{} ฆานํ ปริญฺเญยฺยํ, คนฺธา ปริญฺเญยฺยา.\csromanspacer{} ชิวฺหา ปริญฺเญยฺยา, รสา ปริญฺเญยฺยา.\csromanspacer{} กาโย ปริญฺเญยฺโย, โผฏฺฐพฺพา ปริญฺเญยฺยา ฯเปฯ.\csromanspacer{} มโน ปริญฺเญยฺโย, ธมฺมา ปริญฺเญยฺยา.\csromanspacer{} มโนวิญฺญาณํ ปริญฺเญยฺยํ, มโนสมฺผสฺโส ปริญฺเญยฺโย.\csromanspacer{} ยมฺปิทํ มโนสมฺผสฺสปจฺจยา อุปฺปชฺชติ เวทยิตํ สุขํ วา ทุกฺขํ วา อทุกฺขมสุขํ วา, ตมฺปิ ปริญฺเญยฺยํ.\csromanspacer{}\csromanspacer{}รูปํ ปริญฺเญยฺยํ.\csromanspacer{} เวทนา ปริญฺเญยฺยา.\csromanspacer{} สญฺญา ปริญฺเญยฺยา.\csromanspacer{} สงฺขารา ปริญฺเญยฺยา.\csromanspacer{} วิญฺญาณํ ปริญฺเญยฺยํ.\csromanspacer{} จกฺขุ ปริญฺเญยฺยํ ฯเปฯ.\csromanspacer{} ชรามรณํ ปริญฺเญยฺยํ ฯเปฯ อมโตคธํ นิพฺพานํ ปริโยสานฏฺเฐน ปริญฺเญยฺยํ.\csromanspacer{} เยสํ เยสํ ธมฺมานํ ปฏิลาภตฺถาย วายมนฺตสฺส, เต เต ธมฺมา ปฏิลทฺธา โหนฺติ.\csromanspacer{} เอวํ เต ธมฺมา ปริญฺญาตา เจว โหนฺติ ตีริตา จ.}

\proseitem{22}{เนกฺขมฺมํ ปฏิลาภตฺถาย วายมนฺตสฺส เนกฺขมฺมํ ปฏิลทฺธํ โหติ, เอวํ โส ธมฺโม ปริญฺญาโต เจว โหติ ตีริโต จ.\csromanspacer{} อพฺยาปาทํ ปฏิลาภตฺถาย วายมนฺตสฺส อพฺยาปาโท ปฏิลทฺโธ โหติ, เอวํ โส ธมฺโม ปริญฺญาโต เจว โหติ ตีริโต จ.\csromanspacer{} อาโลกสญฺญํ ปฏิลาภตฺถาย วายมนฺตสฺส อาโลกสญฺญา ปฏิลทฺธา โหติ, เอวํ โส ธมฺโม ปริญฺญาโต เจว โหติ ตีริโต จ.\csromanspacer{} อวิกฺเขปํ ปฏิลาภตฺถาย วายมนฺตสฺส อวิกฺเขโป ปฏิลทฺโธ โหติ, เอวํ โส \csromanfolio{24}ธมฺโม ปริญฺญาโต เจว โหติ ตีริโต จ.\csromanspacer{} ธมฺมววตฺถานํ ปฏิลาภตฺถาย วายมนฺตสฺส ธมฺมววตฺถานํ ปฏิลทฺธํ โหติ, เอวํ โส ธมฺโม ปริญฺญาโต เจว โหติ ตีริโต จ.\csromanspacer{} ญาณํ ปฏิลาภตฺถาย วายมนฺตสฺส ญาณํ ปฏิลทฺธํ โหติ, เอวํ โส ธมฺโม ปริญฺญาโต เจว โหติ ตีริโต จ.\csromanspacer{} ปาโมชฺชํ ปฏิลาภตฺถาย วายมนฺตสฺส ปาโมชฺชํ ปฏิลทฺธํ โหติ, เอวํ โส ธมฺโม ปริญฺญาโต เจว โหติ ตีริโต จ. (7)}

\prose{ปฐมํ ฌานํ ปฏิลาภตฺถาย วายมนฺตสฺส ปฐมํ ฌานํ ปฏิลทฺธํ โหติ, เอวํ โส ธมฺโม ปริญฺญาโต เจว โหติ ตีริโต จ.\csromanspacer{} ทุติยํ ฌานํ ฯเปฯ.\csromanspacer{} ตติยํ ฌานํ.\csromanspacer{} จตุตฺถํ ฌานํ ปฏิลาภตฺถาย วายมนฺตสฺส จตุตฺถํ ฌานํ ปฏิลทฺธํ โหติ, เอวํ โส ธมฺโม ปริญฺญาโต เจว โหติ ตีริโต จ.\csromanspacer{} อากาสานญฺจายตนสมาปตฺติํ ปฏิลาภตฺถาย วายมนฺตสฺส อากาสานญฺจายตนสมาปตฺติ ปฏิลทฺธา โหติ, เอวํ โส ธมฺโม ปริญฺญาตฺโต เจว โหติ ตีริโต จ.\csromanspacer{} วิญฺญาณญฺจายตนสมาปตฺติํ ปฏิลาภตฺถาย วายมนฺตสฺส วิญฺญาณญฺจายตนสมาปตฺติ ปฏิลทฺธา โหติ, เอวํ โส ธมฺโม ปริญฺญาโต เจว โหติ ตีริโต จ.\csromanspacer{} อากิญฺจญฺญายตนสมาปตฺติํ ปฏิลาภตฺถาย วายมนฺตสฺส อากิญฺจญฺญายตนสมาปตฺติ ปฏิลทฺธา โหติ, เอวํ โส ธมฺโม ปริญฺญาโต เจว โหติ ตีริโต จ.\csromanspacer{} เนวสญฺญานาสญฺญายตนสมาปตฺติํ ปฏิลาภตฺถาย วายมนฺตสฺส เนวสญฺญานาสญฺญายตนสมาปตฺติ ปฏิลทฺธา โหติ, เอวํ โส ธมฺโม ปริญฺญาโต เจว โหติ ตีริโต จ. (8)}

\prose{อนิจฺจานุปสฺสนํ ปฏิลาภตฺถาย วายมนฺตสฺส อนิจฺจานุปสฺสนา ปฏิลทฺธา โหติ, เอวํ โส ธมฺโม ปริญฺญาโต เจว โหติ ตีริโต จ.\csromanspacer{} ทุกฺขานุปสฺสนํ ปฏิลาภตฺถาย วายมนฺตสฺส ทุกฺขานุปสฺสนา ปฏิลทฺธา โหติ, เอวํ โส ธมฺโม ปริญฺญาโต เจว โหติ ตีริโต จ.\csromanspacer{} อนตฺตานุปสฺสนํ ปฏิลาภตฺถาย วายมนฺตสฺส อนตฺตานุปสฺสนา ปฏิลทฺธา โหติ, เอวํ โส ธมฺโม ปริญฺญาโต เจว โหติ ตีริโต จ.\csromanspacer{}\csromanspacer{}นิพฺพิทานุปสฺสนํ ปฏิลาภตฺถาย วายมนฺตสฺส นิพฺพิทานุปสฺสนา ปฏิลทฺธา โหติ, เอวํ โส ธมฺโม ปริญฺญาโต เจว โหติ ตีริโต จ.\csromanspacer{} วิราคานุปสฺสนํ ปฏิลาภตฺถาย วายมนฺตสฺส วิราคานุปสฺสนา ปฏิลทฺธา โหติ, เอวํ โส ธมฺโม ปริญฺญาโต เจว โหติ ตีริโต จ. \csromanfolio{25}นิโรธานุปสฺสนํ ปฏิลาภตฺถาย วายมนฺตสฺส นิโรธานุปสฺสนา ปฏิลทฺธา โหติ, เอวํ โส ธมฺโม ปริญฺญาโต เจว โหติ ตีริโต จ.\csromanspacer{} ปฏินิสฺสคฺคานุปสฺสนํ ปฏิลาภตฺถาย วายมนฺตสฺส ปฏินิสฺสคฺคานุปสฺสนา ปฏิลทฺธา โหติ, เอวํ โส ธมฺโม ปริญฺญาโต เจว โหติ ตีริโต จ.\csromanspacer{}\csromanspacer{}ขยานุปสฺสนํ ปฏิลาภตฺถาย วายมนฺตสฺส ขยานุปสฺสนา ปฏิลทฺธา โหติ, เอวํ โส ธมฺโม ปริญฺญาโต เจว โหติ ตีริโต จ.\csromanspacer{} วยานุปสฺสนํ ปฏิลาภตฺถาย วายมนฺตสฺส วยานุปสฺสนา ปฏิลทฺธา โหติ, เอวํ โส ธมฺโม ปริญฺญาโต เจว โหติ ตีริโต จ.\csromanspacer{} วิปริณามานุปสฺสนํ ปฏิลาภตฺถาย วายมนฺตสฺส วิปริณามานุปสฺสนา ปฏิลทฺธา โหติ, เอวํ โส ธมฺโม ปริญฺญาโต เจว โหติ ตีริโต จ.\csromanspacer{}\csromanspacer{}อนิมิตฺตานุปสฺสนํ ปฏิลาภตฺถาย วายมนฺตสฺส อนิมิตฺตานุปสฺสนา ปฏิลทฺธา โหติ, เอวํ โส ธมฺโม ปริญฺญาโต เจว โหติ ตีริโต จ.\csromanspacer{} อปฺปณิหิตานุปสฺสนํ ปฏิลาภตฺถาย วายมนฺตสฺส อปฺปณิหิตานุปสฺสนา ปฏิลทฺธา โหติ, เอวํ โส ธมฺโม ปริญฺญาโต เจว โหติ ตีริโต จ.\csromanspacer{} สุญฺญตานุปสฺสนํ ปฏิลาภตฺถาย วายมนฺตสฺส สุญฺญตานุปสฺสนา ปฏิลทฺธา โหติ, เอวํ โส ธมฺโม ปริญฺญาโต เจว โหติ ตีริโต จ.}

\prose{อธิปญฺญาธมฺมวิปสฺสนํ ปฏิลาภตฺถาย วายมนฺตสฺส อธิปญฺญาธมฺมวิปสฺสนา ปฏิลทฺธา โหติ, เอวํ โส ธมฺโม ปริญฺญาโต เจว โหติ ตีริโต จ.\csromanspacer{} ยถาภูตญาณทสฺสนํ ปฏิลาภตฺถาย วายมนฺตสฺส ยถาภูตญาณทสฺสนํ ปฏิลทฺธํ โหติ, เอวํ โส ธมฺโม ปริญฺญาโต เจว โหติ ตีริโต จ.\csromanspacer{} อาทีนวานุปสฺสนํ ปฏิลาภตฺถาย วายมนฺตสฺส อาทีนวานุปสฺสนา ปฏิลทฺธา โหติ, เอวํ โส ธมฺโม ปริญฺญาโต เจว โหติ ตีริโต จ.\csromanspacer{} ปฏิสงฺขานุปสฺสนํ ปฏิลาภตฺถาย วายมนฺตสฺส ปฏิสงฺขานุปสฺสนา ปฏิลทฺธา โหติ, เอวํ โส ธมฺโม ปริญฺญาโต เจว โหติ ตีริโต จ.\csromanspacer{} วิวฏฺฏนานุปสฺสนํ ปฏิลาภตฺถาย วายมนฺตสฺส วิวฏฺฏนานุปสฺสนา ปฏิลทฺธา โหติ, เอวํ โส ธมฺโม ปริญฺญาโต เจว โหติ ตีริโต จ. (18)}

\prose{โสตาปตฺติมคฺคํ ปฏิลาภตฺถาย วายมนฺตสฺส โสตาปตฺติมคฺโค ปฏิลทฺโธ โหติ.\csromanspacer{} เอวํ โส ธมฺโม ปริญฺญาโต เจว โหติ ตีริโต จ.\csromanspacer{} สกทาคามิมคฺคํ ปฏิลาภตฺถาย วายมนฺตสฺส สกทาคามิมคฺโค ปฏิลทฺโธ โหติ, เอวํ โส ธมฺโม ปริญฺญาโต เจว โหติ ตีริโต จ. \csromanfolio{26}อนาคามิมคฺคํ ปฏิลาภตฺถาย วายมนฺตสฺส อนาคามิมคฺโค ปฏิลทฺโธ โหติ, เอวํ โส ธมฺโม ปริญฺญาโต เจว โหติ ตีริโต จ.\csromanspacer{} อรหตฺตมคฺคํ ปฏิลาภตฺถาย วายมนฺตสฺส อรหตฺตมคฺโค ปฏิลทฺโธ โหติ, เอวํ โส ธมฺโม ปริญฺญาโต เจว โหติ ตีริโต จ. (4)}

\prose{เยสํ เยสํ ธมฺมานํ ปฏิลาภตฺถาย วายมนฺตสฺส, เต เต ธมฺมา ปฏิลทฺธา โหนฺติ.\csromanspacer{} เอวํ เต ธมฺมา ปริญฺญาตา เจว โหนฺติ ตีริตา จ.\csromanspacer{} ตํญาตฏฺเฐน ญาณํ, ปชานนฏฺเฐน ปญฺญา, เตน วุจฺจติ “อิเม ธมฺมา ปริญฺเญยฺยาติ โสตาวธานํ, ตํปชานนา ปญฺญา สุตมเย ญาณนฺ”ติ. (2)}

\proseitem{23}{กถํ “อิเม ธมฺมา ปหาตพฺพา”ติ โสตาวธานํ, ตํปชานนา ปญฺญา สุตมเย ญาณํ\csromandash{}}

\prose{เอโก ธมฺโม ปหาตพฺโพ, อสฺมิมาโน.\csromanspacer{}\csromanspacer{}เทฺว ธมฺมา ปหาตพฺพา, อวิชฺชา จ ภวตณฺหา จ.\csromanspacer{}\csromanspacer{}ตโย ธมฺมา ปหาตพฺพา, ติสฺโส ตณฺหา.\csromanspacer{}\csromanspacer{}จตฺตาโร ธมฺมา ปหาตพฺพา, จตฺตาโร โอฆา.\csromanspacer{}\csromanspacer{}ปญฺจ ธมฺมา ปหาตพฺพา, ปญฺจ นีวรณานิ.\csromanspacer{}\csromanspacer{}ฉ ธมฺมา ปหาตพฺพา, ฉ ตณฺหากายา.\csromanspacer{}\csromanspacer{}สตฺต ธมฺมา ปหาตพฺพา, สตฺตานุสยา.\csromanspacer{}\csromanspacer{}อฏฺฐ ธมฺมา ปหาตพฺพา, อฏฺฐ มิจฺฉตฺตา.\csromanspacer{}\csromanspacer{}นว ธมฺมา ปหาตพฺพา, นว ตณฺหามูลกา.\csromanspacer{}\csromanspacer{}ทส ธมฺมา ปหาตพฺพา, ทส มิจฺฉตฺตา.}

\proseitem{24}{เทฺว ปหานานิ สมุจฺเฉทปฺปหานํ, ปฏิปฺปสฺสทฺธิปฺปหานํ.\csromanspacer{} สมุจฺเฉทปฺปหานญฺจ โลกุตฺตรํ ขยคามิมคฺคํ ภาวยโต, ปฏิปฺปสฺสทฺธิปฺปหานญฺจ ผลกฺขเณ.\csromanspacer{}\csromanspacer{}ตีณิ ปหานานิ, กามานเมตํ นิสฺสรณํ ยทิทํ เนกฺขมฺมํ, รูปานเมตํ นิสฺสรณํ ยทิทํ อารุปฺปํ.\csromanspacer{} ยํ โข ปน กิญฺจิ ภูตํ สงฺขตํ ปฏิจฺจสมุปฺปนฺนํ นิโรโธ, ตสฺส นิสฺสรณํ.\csromanspacer{} เนกฺขมฺมํ ปฏิลทฺธสฺส กามา ปหีนา เจว โหนฺติ ปริจฺจตฺตา จ, อารุปฺปํ ปฏิลทฺธสฺส รูปา ปหีนา เจว โหนฺติ ปริจฺจตฺตา จ, นิโรธํ ปฏิลทฺธสฺส สงฺขารา ปหีนา เจว โหนฺติ ปริจฺจตฺตา จ.\csromanspacer{}\csromanspacer{}จตฺตาริ ปหานานิ, ทุกฺขสจฺจํ ปริญฺญาปฏิเวธํ ปฏิวิชฺฌนฺโต ปชหาติ, สมุทยสจฺจํ ปหานปฏิเวธํ ปฏิวิชฺฌนฺโต ปชหาติ, นิโรธสจฺจํ สจฺฉิกิริยาปฏิเวธํ ปฏิวิชฺฌนฺโต ปชหาติ, มคฺคสจฺจํ ภาวนาปฏิเวธํ ปฏิวิชฺฌนฺโต ปชหาติ.\csromanspacer{}\csromanspacer{}ปญฺจ ปหานานิ วิกฺขมฺภนปฺปหานํ ตทงฺคปฺปหานํ สมุจฺเฉทปฺปหานํ ปฏิปฺปสฺสทฺธิปฺปหานํ นิสฺสรณปฺปหานํ.\csromanspacer{} วิกฺขมฺภนปฺปหานญฺจ นีวรณานํ ปฐมํ ฌานํ ภาวยโต, ตทงฺคปฺปหานญฺจ ทิฏฺฐิคตานํ \csromanfolio{27}นิพฺเพธภาคิยํ สมาธิํ ภาวยโต, สมุจฺเฉทปฺปหานญฺจ โลกุตฺตรํ ขยคามิมคฺคํ ภาวยโต, ปฏิปฺปสฺสทฺธิปฺปหานญฺจ ผลกฺขเณ, นิสฺสรณปฺปหานญฺจ นิโรโธ นิพฺพานํ.}

\prose{สพฺพํ ภิกฺขเว ปหาตพฺพํ.\csromanspacer{} กิญฺจ ภิกฺขเว สพฺพํ ปหาตพฺพํ? จกฺขุ ภิกฺขเว ปหาตพฺพํ, รูปา ปหาตพฺพา, จกฺขุวิญฺญาณํ ปหาตพฺพํ, จกฺขุสมฺผสฺโส ปหาตพฺโพ.\csromanspacer{} ยมฺปิทํ จกฺขุสมฺผสฺสปจฺจยา อุปฺปชฺชติ เวทยิตํ สุขํ วา ทุกฺขํ วา อทุกฺขมสุขํ วา, ตมฺปิ ปหาตพฺพํ.\csromanspacer{} โสตํ ปหาตพฺพํ, สทฺทา ปหาตพฺพา ฯเปฯ.\csromanspacer{} ฆานํ ปหาตพฺพํ, คนฺธา ปหาตพฺพา.\csromanspacer{} ชิวฺหา ปหาตพฺพา, รสา ปหาตพฺพา.\csromanspacer{} กาโย ปหาตพฺโพ, โผฏฺฐพฺพา ปหาตพฺพา.\csromanspacer{} มโน ปหาตพฺโพ, ธมฺมา ปหาตพฺพา.\csromanspacer{} มโนวิญฺญาณํ ปหาตพฺพํ, มโนสมฺผสฺโส ปหาตพฺโพ.\csromanspacer{} ยมฺปิทํ มโนสมฺผสฺสปจฺจยา อุปฺปชฺชติ เวทยิตํ สุขํ วา ทุกฺขํ วา อทุกฺขมสุขํ วา, ตมฺปิ ปหาตพฺพํ.\csromanspacer{} รูปํ ปสฺสนฺโต ปชหาติ, เวทนํ ปสฺสนฺโต ปชหาติ, สญฺญํ ปสฺสนฺโต ปชหาติ, สงฺขาเร ปสฺสนฺโต ปชหาติ, วิญฺญาณํ ปสฺสนฺโต ปชหาติ.\csromanspacer{} จกฺขุํ ฯเปฯ.\csromanspacer{} ชรามรณํ ฯเปฯ อมโตคธํ นิพฺพานํ ปริโยสานฏฺเฐน ปสฺสนฺโต ปชหาติ.\csromanspacer{} เย เย ธมฺมา ปหีนา โหนฺติ, เต เต ธมฺมา ปริจฺจตฺตา โหนฺติ, ตํญาตฏฺเฐน ญาณํ, ปชานนฏฺเฐน ปญฺญา, เตน วุจฺจติ “อิเม ธมฺมา ปหาตพฺพาติ โสตาวธานํ, ตํปชานนา ปญฺญา สุตมเย ญาณนฺ”ติ. (3) ตติยภาณวาโร.}

\proseitem{25}{กถํ “อิเม ธมฺมา ภาเวตพฺพา”ติ โสตาวธานํ, ตํปชานนา ปญฺญา สุตมเย ญาณํ\csromandash{}}

\prose{เอโก ธมฺโม ภาเวตพฺโพ, กายคตาสติ สาตสหคตา.\csromanspacer{}\csromanspacer{}เทฺว ธมฺมา ภาเวตพฺพา, สมโถ จ วิปสฺสนา จ.\csromanspacer{}\csromanspacer{}ตโย ธมฺมา ภาเวตพฺพา, ตโย สมาธี.\csromanspacer{}\csromanspacer{}จตฺตาโร ธมฺมา ภาเวตพฺพา, จตฺตาโร สติปฏฺฐานา.\csromanspacer{}\csromanspacer{}ปญฺจ ธมฺมา ภาเวตพฺพา, ปญฺจงฺคิโก สมาธิ\footnote{สมฺมาสมาธิ (สฺยา)}.\csromanspacer{}\csromanspacer{}ฉ ธมฺมา ภาเวตพฺพา, ฉ อนุสฺสติฏฺฐานานิ.\csromanspacer{}\csromanspacer{}สตฺต ธมฺมา ภาเวตพฺพา, สตฺต โพชฺฌงฺคา.\csromanspacer{}\csromanspacer{}อฏฺฐ ธมฺมา ภาเวตพฺพา, อริโย อฏฺฐงฺคิโก มคฺโค.\csromanspacer{}\csromanspacer{}นว ธมฺมา ภาเวตพฺพา, นว ปาริสุทฺธิปธานิยงฺคานิ.\csromanspacer{}\csromanspacer{}ทส ธมฺมา ภาเวตพฺพา, ทส กสิณายตนานิ.}

\proseitem{26}{\csromanfolio{28}เทฺว ภาวนา โลกิยา จ ภาวนา, โลกุตฺตรา จ ภาวนา.\csromanspacer{} ติสฺโส ภาวนา, รูปาวจรกุสลานํ ธมฺมานํ ภาวนา อรูปาวจรกุสลานํ ธมฺมานํ ภาวนา อปริยาปนฺนกุสลานํ ธมฺมานํ ภาวนา.\csromanspacer{} รูปาวจรกุสลานํ ธมฺมานํ ภาวนา อตฺถิ หีนา, อตฺถิ มชฺฌา อตฺถิ ปณีตา, อรูปาวจรกุสลานํ ธมฺมานํ ภาวนา อตฺถิ หีนา, อตฺถิ มชฺฌา, อตฺถิ ปณีตา, อปริยาปนฺนกุสลานํ ธมฺมานํ ภาวนา ปณีตา.}

\proseitem{27}{จตสฺโส ภาวนา, ทุกฺขสจฺจํ ปริญฺญาปฏิเวธํ ปฏิวิชฺฌนฺโต ภาเวติ, สมุทยสจฺจํ ปหานปฺปฏิเวธํ ปฏิวิชฺฌนฺโต ภาเวติ, นิโรธสจฺจํ สจฺฉิกิริยาปฏิเวธํ ปฏิวิชฺฌนฺโต ภาเวติ, มคฺคสจฺจํ ภาวนาปฏิเวธํ ปฏิวิชฺฌนฺโต ภาเวติ, อิมา จตสฺโส ภาวนา. (4)}

\prose{อปราปิ จตสฺโส ภาวนา, เอสนาภาวนา, ปฏิลาภภาวนา, เอกรสภาวนา, อาเสวนาภาวนา.\csromanspacer{} กตมา เอสนาภาวนา? สพฺเพสํ สมาธิํ สมาปชฺชนฺตานํ ตตฺถ ชาตา ธมฺมา เอกรสา โหนฺตีติ อยํ เอสนาภาวนา.\csromanspacer{}\csromanspacer{}กตมา ปฏิลาภภาวนา? สพฺเพสํ สมาธิํ สมาปนฺนานํ ตตฺถ ชาตา ธมฺมา อญฺญมญฺญํ นาติวตฺตนฺตีติ อยํ ปฏิลาภภาวนา.\csromanspacer{}\csromanspacer{}กตมา เอกรสภาวนา? อธิโมกฺขฏฺเฐน สทฺธินฺทฺริยํ ภาวยโต สทฺธินฺทฺริยสฺส วเสน จตฺตาริ อินฺทฺริยานิ เอกรสา โหนฺตีติ อินฺทฺริยานํ เอกรสฏฺเฐน ภาวนา, ปคฺคหฏฺเฐน วีริยินฺทฺริยํ ภาวยโต วีริยินฺทฺริยสฺส วเสน จตฺตาริ อินฺทฺริยานิ เอกรสา โหนฺตีติ อินฺทฺริยานํ เอกรสฏฺเฐน ภาวนา.\csromanspacer{} อุปฏฺฐานฏฺเฐน สตินฺทฺริยํ ภาวยโต สตินฺทฺริยสฺส วเสน จตฺตาริ อินฺทฺริยานิ เอกรสา โหนฺตีติ อินฺทฺริยานํ เอกรสฏฺเฐน ภาวนา.\csromanspacer{} อวิกฺเขปฏฺเฐน สมาธินฺทฺริยํ ภาวยโต สมาธินฺทฺริยสฺส วเสน จตฺตาริ อินฺทฺริยานิ เอกรสา โหนฺตีติ อินฺทฺริยานํ เอกรสฏฺเฐน ภาวนา, ทสฺสนฏฺเฐน ปญฺญินฺทฺริยํ ภาวยโต ปญฺญินฺทฺริยสฺส วเสน จตฺตาริ อินฺทฺริยานิ เอกรสา โหนฺตีติ อินฺทฺริยานํ เอกรสฏฺเฐน ภาวนา.}

\prose{อสฺสทฺธิเย อกมฺปิยฏฺเฐน สทฺธาพลํ ภาวยโต สทฺธาพลสฺส วเสน จตฺตาริ พลานิ เอกรสา โหนฺตีติ พลานํ เอกรสฏฺเฐน ภาวนา, โกสชฺเช อกมฺปิยฏฺเฐน วีริยพลํ ภาวยโต วีริยพลสฺส วเสน \csromanfolio{29}จตฺตาริ พลานิ เอกรสา โหนฺตีติ พลานํ เอกรสฏฺเฐน ภาวนา, ปมาเท อกมฺปิยฏฺเฐน สติพลํ ภาวยโต สติพลสฺส วเสน จตฺตาริ พลานิ เอกรสา โหนฺตีติ พลานํ เอกรสฏฺเฐน ภาวนา, อุทฺธจฺเจ อกมฺปิยฏฺเฐน สมาธิพลํ ภาวยโต สมาธิพลสฺส วเสน จตฺตาริ พลานิ เอกรสา โหนฺตีติ พลานํ เอกรสฏฺเฐน ภาวนา, อวิชฺชาย อกมฺปิยฏฺเฐน ปญฺญาพลํ ภาวยโต ปญฺญาพลสฺส วเสน จตฺตาริ พลานิ เอกรสา โหนฺตีติ พลานํ เอกรสฏฺเฐน ภาวนา.}

\prose{อุปฏฺฐานฏฺเฐน สติสมฺโพชฺฌงฺคํ ภาวยโต สติสมฺโพชฺฌงฺคสฺส วเสน ฉ โพชฺฌงฺคา เอกรสา โหนฺตีติ โพชฺฌงฺคานํ เอกรสฏฺเฐน ภาวนา, ปวิจยฏฺเฐน ธมฺมวิจยสมฺโพชฺฌงฺคํ ภาวยโต ธมฺมวิจยสมฺโพชฺฌงฺคสฺส วเสน ฉ โพชฺฌงฺคา เอกรสา โหนฺตีติ โพชฺฌงฺคานํ เอกรสฏฺเฐน ภาวนา, ปคฺคหฏฺเฐน วีริยสมฺโพชฺฌงฺคํ ภาวยโต วีริยสมฺโพชฺฌงฺคสฺส วเสน ฉ โพชฺฌงฺคา เอกรสา โหนฺตีติ โพชฺฌงฺคานํ เอกรสฏฺเฐน ภาวนา, ผรณฏฺเฐน ปีติสมฺโพชฺฌงฺคํ ภาวยโต ปีติสมฺโพชฺฌงฺคสฺส วเสน ฉ โพชฺฌงฺคา เอกรสา โหนฺตีติ โพชฺฌงฺคานํ เอกรสฏฺเฐน ภาวนา, อุปสมฏฺเฐน ปสฺสทฺธิสมฺโพชฺฌงฺคํ ภาวยโต ปสฺสทฺธิสมฺโพชฺฌงฺคสฺส วเสน ฉ โพชฺฌงฺคา เอกรสา โหนฺตีติ โพชฺฌงฺคานํ เอกรสฏฺเฐน ภาวนา, อวิกฺเขปฏฺเฐน สมาธิสมฺโพชฺฌงฺคํ ภาวยโต สมาธิสมฺโพชฺฌงฺคสฺส วเสน ฉ โพชฺฌงฺคา เอกรสา โหนฺตีติ โพชฺฌงฺคานํ เอกรสฏฺเฐน ภาวนา, ปฏิสงฺขานฏฺเฐน อุเปกฺขาสมฺโพชฺฌงฺคํ ภาวยโต อุเปกฺขาสมฺโพชฺฌงฺคสฺส วเสน ฉ โพชฺฌงฺคา เอกรสา โหนฺตีติ โพชฺฌงฺคานํ เอกรสฏฺเฐน ภาวนา.}

\prose{ทสฺสนฏฺเฐน สมฺมาทิฏฺฐิํ ภาวยโต สมฺมาทิฏฺฐิยา วเสน สตฺต มคฺคงฺคา เอกรสา โหนฺตีติ มคฺคงฺคานํ เอกรสฏฺเฐน ภาวนา, อภินิโรปนฏฺเฐน สมฺมาสงฺกปฺปํ ภาวยโต สมฺมาสงฺกปฺปสฺส วเสน สตฺต มคฺคงฺคา เอกรสา โหนฺตีติ มคฺคงฺคานํ เอกรสฏฺเฐน ภาวนา, ปริคฺคหฏฺเฐน สมฺมาวาจํ ภาวยโต สมฺมาวาจาย วเสน สตฺต มคฺคงฺคา เอกรสา โหนฺตีติ มคฺคงฺคานํ เอกรสฏฺเฐน ภาวนา, สมุฏฺฐานฏฺเฐน สมฺมากมฺมนฺตํ ภาวยโต สมฺมากมฺมนฺตสฺส วเสน สตฺต มคฺคงฺคา เอกรสา โหนฺตีติ มคฺคงฺคานํ เอกรสฏฺเฐน ภาวนา, โวทานฏฺเฐน สมฺมา-อาชีวํ ภาวยโต สมฺมา-อาชีวสฺส วเสน สตฺต มคฺคงฺคา เอกรสา โหนฺตีติ มคฺคงฺคานํ เอกรสฏฺเฐน \csromanfolio{30}ภาวนา, ปคฺคหฏฺเฐน สมฺมาวายามํ ภาวยโต สมฺมาวายามสฺส วเสน สตฺต มคฺคงฺคา เอกรสา โหนฺตีติ มคฺคงฺคานํ เอกรสฏฺเฐน ภาวนา, อุปฏฺฐานฏฺเฐน สมฺมาสติํ ภาวยโต สมฺมาสติยา วเสน สตฺต มคฺคงฺคา เอกรสา โหนฺตีติ มคฺคงฺคานํ เอกรสฏฺเฐน ภาวนา, อวิกฺเขปฏฺเฐน สมฺมาสมาธิํ ภาวยโต สมฺมาสมาธิสฺส วเสน สตฺต มคฺคงฺคา เอกรสา โหนฺตีติ มคฺคงฺคานํ เอกรสฏฺเฐน ภาวนา.\csromanspacer{} อยํ เอกรสาภาวนา.}

\prose{กตมา อาเสวนาภาวนา? อิธ ภิกฺขุ ปุพฺพณฺหสมยํ อาเสวติ, มชฺฌนฺหิกสมยมฺปิ\footnote{มชฺฌนฺติกสมยมฺปิ (สฺยา, ก)} อาเสวติ, สายนฺหสมยมฺปิ อาเสวติ, ปุเรภตฺตมฺปิ อาเสวติ, ปจฺฉาภตฺตมฺปิ อาเสวติ, ปุริเมปิ ยาเม อาเสวติ, มชฺฌิเมปิ ยาเม อาเสวติ, ปจฺฉิเมปิ ยาเม อาเสวติ, รตฺติมฺปิ อาเสวติ, ทิวาปิ อาเสวติ, รตฺตินฺทิวาปิ\footnote{รตฺติทิวาปิ (ก)} อาเสวติ, กาเฬปิ อาเสวติ, ชุเณฺหปิ อาเสวติ, วสฺเสปิ อาเสวติ, เหมนฺเตปิ อาเสวติ, คิเมฺหปิ อาเสวติ, ปุริเมปิ วโยขนฺเธ อาเสวติ, มชฺฌิเมปิ วโยขนฺเธ อาเสวติ, ปจฺฉิเมปิ วโยขนฺเธ อาเสวติ.\csromanspacer{} อยํ อาเสวนาภาวนา.\csromanspacer{} อิมา จตสฺโส ภาวนา.}

\proseitem{28}{อปราปิ จตสฺโส ภาวนา, ตตฺถ ชาตานํ ธมฺมานํ อนติวตฺตนฏฺเฐน ภาวนา, อินฺทฺริยานํ เอกรสฏฺเฐน ภาวนา, ตทุปควีริยวาหนฏฺเฐน ภาวนา, อาเสวนฏฺเฐน ภาวนา.\csromanspacer{} กถํ ตตฺถ ชาตานํ ธมฺมานํ อนติวตฺตนฏฺเฐน ภาวนา? กามจฺฉนฺทํ ปชหโต เนกฺขมฺมวเสน ชาตา ธมฺมา อญฺญมญฺญํ นาติวตฺตนฺตีติ ตตฺถ ชาตานํ ธมฺมานํ อนติวตฺตนฏฺเฐน ภาวนา, พฺยาปาทํ ปชหโต อพฺยาปาทวเสน ชาตา ธมฺมา อญฺญมญฺญํ นาติวตฺตนฺตีติ ตตฺถ ชาตานํ ธมฺมานํ อนติวตฺตนฏฺเฐน ภาวนา, ถินมิทฺธํ\footnote{ถีนมิทฺธํ (สฺยา, สี-ฏฺฐ)} ปชหโต อาโลกสญฺญาวเสน ชาตา ธมฺมา อญฺญมญฺญํ นาติวตฺตนฺตีติ ตตฺถ ชาตานํ ธมฺมานํ อนติวตฺตนฏฺเฐน ภาวนา, อุทฺธจฺจํ ปชหโต อวิกฺเขปวเสน ชาตา ธมฺมา อญฺญมญฺญํ นาติวตฺตนฺตีติ ตตฺถ ชาตานํ ธมฺมานํ อนติวตฺตนฏฺเฐน ภาวนา, วิจิกิจฺฉํ ปชหโต ธมฺมววตฺถานวเสน ชาตา ธมฺมา อญฺญมญฺญํ นาติวตฺตนฺตีติ ตตฺถ ชาตานํ ธมฺมานํ อนติวตฺตนฏฺเฐน ภาวนา, อวิชฺชํ ปชหโต ญาณวเสน ชาตา ธมฺมา อญฺญมญฺญํ นาติวตฺตนฺตีติ ตตฺถ ชาตานํ ธมฺมานํ อนติวตฺตนฏฺเฐน \csromanfolio{31}ภาวนา, อรติํ ปชหโต ปาโมชฺชวเสน ชาตา ธมฺมา อญฺญมญฺญํ นาติวตฺตนฺตีติ ตตฺถ ชาตานํ ธมฺมานํ อนติวตฺตนฏฺเฐน ภาวนา, นีวรเณ ปชหโต ปฐมชฺฌานวเสน ชาตา ธมฺมา อญฺญมญฺญํ นาติวตฺตนฺตีติ ตตฺถ ชาตานํ ธมฺมานํ อนติวตฺตนฏฺเฐน ภาวนา, วิตกฺกวิจาเร ปชหโต ทุติยชฺฌานวเสน ชาตา ธมฺมา อญฺญมญฺญํ นาติวตฺตนฺตีติ ตตฺถ ชาตานํ ธมฺมานํ อนติวตฺตนฏฺเฐน ภาวนา, ปีติํ ปชหโต ตติยชฺฌานวเสน ชาตา ธมฺมา อญฺญมญฺญํ นาติวตฺตนฺตีติ ตตฺถ ชาตานํ ธมฺมานํ อนติวตฺตนฏฺเฐน ภาวนา, สุขทุกฺเข ปชหโต จตุตฺถชฺฌานวเสน ชาตา ธมฺมา อญฺญมญฺญํ นาติวตฺตนฺตีติ ตตฺถ ชาตานํ ธมฺมานํ อนติวตฺตนฏฺเฐน ภาวนา.}

\prose{รูปสญฺญํ ปฏิฆสญฺญํ นานตฺตสญฺญํ ปชหโต อากาสานญฺจายตนสมาปตฺติวเสน ชาตา ธมฺมา อญฺญมญฺญํ นาติวตฺตนฺตีติ ตตฺถ ชาตานํ ธมฺมานํ อนติวตฺตนฏฺเฐน ภาวนา, อากาสานญฺจายตนสญฺญํ ปชหโต วิญฺญาณญฺจายตนสมาปตฺติวเสน ชาตา ธมฺมา อญฺญมญฺญํ นาติวตฺตนฺตีติ ตตฺถ ชาตานํ ธมฺมานํ อนติวตฺตนฏฺเฐน ภาวนา, วิญฺญาณญฺจายตนสญฺญํ ปชหโต อากิญฺจญฺญายตนสมาปตฺติวเสน ชาตา ธมฺมา อญฺญมญฺญํ นาติวตฺตนฺตีติ ตตฺถ ชาตานํ ธมฺมานํ อนติวตฺตนฏฺเฐน ภาวนา, อากิญฺจญฺญายตนสญฺญํ ปชหโต เนวสญฺญานาสญฺญายตนสมาปตฺติวเสน ชาตา ธมฺมา อญฺญมญฺญํ นาติวตฺตนฺตีติ ตตฺถ ชาตานํ ธมฺมานํ อนติวตฺตนฏฺเฐน ภาวนา.}

\prose{นิจฺจสญฺญํ ปชหโต อนิจฺจานุปสฺสนาวเสน ชาตา ธมฺมา อญฺญมญฺญํ นาติวตฺตนฺตีติ ตตฺถ ชาตานํ ธมฺมานํ อนติวตฺตนฏฺเฐน ภาวนา, สุขสญฺญํ ปชหโต ทุกฺขานุปสฺสนาวเสน ชาตา ธมฺมา อญฺญมญฺญํ นาติวตฺตนฺตีติ ตตฺถ ชาตานํ ธมฺมานํ อนติวตฺตนฏฺเฐน ภาวนา, อตฺตสญฺญํ ปชหโต อนตฺตานุปสฺสนาวเสน ชาตา ธมฺมา อญฺญมญฺญํ นาติวตฺตนฺตีติ ตตฺถ ชาตานํ ธมฺมานํ อนติวตฺตนฏฺเฐน ภาวนา, นนฺทิํ ปชหโต นิพฺพิทานุปสฺสนาวเสน ชาตา ธมฺมา อญฺญมญฺญํ นาติวตฺตนฺตีติ ตตฺถ ชาตานํ ธมฺมานํ อนติวตฺตนฏฺเฐน ภาวนา, ราคํ ปชหโต วิราคานุปสฺสนาวเสน ชาตา ธมฺมา อญฺญมญฺญํ นาติวตฺตนฺตีติ ตตฺถ ชาตานํ ธมฺมานํ อนติวตฺตนฏฺเฐน ภาวนา, สมุทยํ ปชหโต นิโรธานุปสฺสนาวเสน ชาตา ธมฺมา อญฺญมญฺญํ นาติวตฺตนฺตีติ ตตฺถ ชาตานํ ธมฺมานํ อนติวตฺตนฏฺเฐน \csromanfolio{32}ภาวนา, อาทานํ ปชหโต ปฏินิสฺสคฺคานุปสฺสนาวเสน ชาตา ธมฺมา อญฺญมญฺญํ นาติวตฺตนฺตีติ ตตฺถ ชาตานํ ธมฺมานํ อนติวตฺตนฏฺเฐน ภาวนา, ฆนสญฺญํ ปชหโต ขยานุปสฺสนาวเสน ชาตา ธมฺมา อญฺญมญฺญํ นาติวตฺตนฺตีติ ตตฺถ ชาตานํ ธมฺมานํ อนติวตฺตนฏฺเฐน ภาวนา, อายูหนํ ปชหโต วยานุปสฺสนาวเสน ชาตา ธมฺมา อญฺญมญฺญํ นาติวตฺตนฺตีติ ตตฺถ ชาตานํ ธมฺมานํ อนติวตฺตนฏฺเฐน ภาวนา.\csromanspacer{} ธุวสญฺญํ ปชหโต วิปริณามานุปสฺสนาวเสน ชาตา ธมฺมา อญฺญมญฺญํ นาติวตฺตนฺตีติ ตตฺถ ชาตานํ ธมฺมานํ อนติวตฺตนฏฺเฐน ภาวนา, นิมิตฺตํ ปชหโต อนิมิตฺตานุปสฺสนาวเสน ชาตา ธมฺมา อญฺญมญฺญํ นาติวตฺตนฺตีติ ตตฺถ ชาตานํ ธมฺมานํ อนติวตฺตนฏฺเฐน ภาวนา, ปณิธิํ ปชหโต อปฺปณิหิตานุปสฺสนาวเสน ชาตา ธมฺมา อญฺญมญฺญํ นาติวตฺตนฺตีติ ตตฺถ ชาตานํ ธมฺมานํ อนติวตฺตนฏฺเฐน ภาวนา, อภินิเวสํ ปชหโต สุญฺญตานุปสฺสนาวเสน ชาตา ธมฺมา อญฺญมญฺญํ นาติวตฺตนฺตีติ ตตฺถ ชาตานํ ธมฺมานํ อนติวตฺตนฏฺเฐน ภาวนา, สาราทานาภินิเวสํ ปชหโต อธิปญฺญาธมฺมวิปสฺสนาวเสน ชาตา ธมฺมา อญฺญมญฺญํ นาติวตฺตนฺตีติ ตตฺถ ชาตานํ ธมฺมานํ อนติวตฺตนฏฺเฐน ภาวนา, สมฺโมหาภินิเวสํ ปชหโต ยถาภูตญาณทสฺสนวเสน ชาตา ธมฺมา อญฺญมญฺญํ นาติวตฺตนฺตีติ ตตฺถ ชาตานํ ธมฺมานํ อนติวตฺตนฏฺเฐน ภาวนา, อาลยาภินิเวสํ ปชหโต อาทีนวานุปสฺสนาวเสน ชาตา ธมฺมา อญฺญมญฺญํ นาติวตฺตนฺตีติ ตตฺถ ชาตานํ ธมฺมานํ อนติวตฺตนฏฺเฐน ภาวนา, อปฺปฏิสงฺขํ ปชหโต ปฏิสงฺขานุปสฺสนาวเสน ชาตา ธมฺมา อญฺญมญฺญํ นาติวตฺตนฺตีติ ตตฺถ ชาตานํ ธมฺมานํ อนติวตฺตนฏฺเฐน ภาวนา, สญฺโญคาภินิเวสํ ปชหโต วิวฏฺฏนานุปสฺสนาวเสน ชาตา ธมฺมา อญฺญมญฺญํ นาติวตฺตนฺตีติ ตตฺถ ชาตานํ ธมฺมานํ อนติวตฺตนฏฺเฐน ภาวนา.}

\prose{ทิฏฺเฐกฏฺเฐ กิเลเส ปชหโต โสตาปตฺติมคฺควเสน ชาตา ธมฺมา อญฺญมญฺญํ นาติวตฺตนฺตีติ ตตฺถ ชาตานํ ธมฺมานํ อนติวตฺตนฏฺเฐน ภาวนา, โอฬาริเก กิเลเส ปชหโต สกทาคามิมคฺควเสน ชาตา ธมฺมา อญฺญมญฺญํ นาติวตฺตนฺตีติ ตตฺถ ชาตานํ ธมฺมานํ อนติวตฺตนฏฺเฐน ภาวนา, อณุสหคเต กิเลเส ปชหโต อนาคามิมคฺควเสน ชาตา ธมฺมา อญฺญมญฺญํ นาติวตฺตนฺตีติ ตตฺถ ชาตานํ ธมฺมานํ อนติวตฺตนฏฺเฐน ภาวนา, สพฺพกิเลเส ปชหโต อรหตฺตมคฺควเสน ชาตา ธมฺมา อญฺญมญฺญํ \csromanfolio{33}นาติวตฺตนฺตีติ ตตฺถ ชาตานํ ธมฺมานํ อนติวตฺตนฏฺเฐน ภาวนา, เอวํ ตตฺถ ชาตานํ ธมฺมานํ อนติวตฺตนฏฺเฐน ภาวนา.}

\prose{กถํ อินฺทฺริยานํ เอกรสฏฺเฐน ภาวนา? กามจฺฉนฺทํ ปชหโต เนกฺขมฺมวเสน ปญฺจินฺทฺริยานิ เอกรสา โหนฺตีติ อินฺทฺริยานํ เอกรสฏฺเฐน ภาวนา, พฺยาปาทํ ปชหโต อพฺยาปาทวเสน ปญฺจินฺทฺริยานิ เอกรสา โหนฺตีติ อินฺทฺริยานํ เอกรสฏฺเฐน ภาวนา ฯเปฯ สพฺพกิเลเส ปชหโต อรหตฺตมคฺควเสน ปญฺจินฺทฺริยานิ เอกรสา โหนฺตีติ อินฺทฺริยานํ เอกรสฏฺเฐน ภาวนา, เอวํ อินฺทฺริยานํ เอกรสฏฺเฐน ภาวนา.}

\prose{กถํ ตทุปควีริยวาหนฏฺเฐน ภาวนา? กามจฺฉนฺทํ ปชหโต เนกฺขมฺมวเสน วีริยํ วาเหตีติ ตทุปควีริยวาหนฏฺเฐน ภาวนา, พฺยาปาทํ ปชหโต อพฺยาปาทวเสน วีริยํ วาเหตีติ ตทุปควีริยวาหนฏฺเฐน ภาวนา ฯเปฯ สพฺพกิเลเส ปชหโต อรหตฺตมคฺควเสน วีริยํ วาเหตีติ ตทุปควีริยวาหนฏฺเฐน ภาวนา, เอวํ ตทุปควีริยวาหนฏฺเฐน ภาวนา.}

\prose{กถํ อาเสวนฏฺเฐน ภาวนา? กามจฺฉนฺทํ ปชหนฺโต เนกฺขมฺมํ อาเสวตีติ อาเสวนฏฺเฐน ภาวนา, พฺยาปาทํ ปชหนฺโต อพฺยาปาทํ อาเสวตีติ อาเสวนฏฺเฐน ภาวนา ฯเปฯ สพฺพกิเลเส ปชหนฺโต อรหตฺตมคฺคํ อาเสวตีติ อาเสวนฏฺเฐน ภาวนา, เอวํ อาเสวนฏฺเฐน ภาวนา.\csromanspacer{} อิมา จตสฺโส ภาวนา.}

\prose{รูปํ ปสฺสนฺโต ภาเวติ, เวทนํ ปสฺสนฺโต ภาเวติ, สญฺญํ ปสฺสนฺโต ภาเวติ, สงฺขาเร ปสฺสนฺโต ภาเวติ, วิญฺญาณํ ปสฺสนฺโต ภาเวติ, จกฺขุํ ฯเปฯ ชรามรณํ ฯเปฯ อมโตคธํ นิพฺพานํ ปริโยสานฏฺเฐน ปสฺสนฺโต ภาเวติ.\csromanspacer{} เย เย ธมฺมา ภาวิตา โหนฺติ, เต เต ธมฺมา เอกรสา โหนฺติ, ตํญาตฏฺเฐน ญาณํ, ปชานนฏฺเฐน ปญฺญา, เตน วุจฺจติ “อิเม ธมฺมา ภาเวตพฺพาติ โสตาวธานํ, ตํปชานนา ปญฺญา สุตมเย ญาณนฺ”ติ. (4)}

\vspace{\tipitakaparskip}
\csromancenter{จตุตฺถภาณวาโร.}
\proseitem{29}{\csromanfolio{34}กถํ “อิเม ธมฺมา สจฺฉิกาตพฺพา”ติ โสตาวธานํ, ตํปชานนา ปญฺญา สุตมเย ญาณํ\csromandash{}}

\prose{เอโก ธมฺโม สจฺฉิกาตพฺโพ, อกุปฺปา เจโตวิมุตฺติ.\csromanspacer{}\csromanspacer{}เทฺว ธมฺมา สจฺฉิกาตพฺพา, วิชฺชา จ วิมุตฺติ จ.\csromanspacer{}\csromanspacer{}ตโย ธมฺมา สจฺฉิกาตพฺพา, ติสฺโส วิชฺชา.\csromanspacer{}\csromanspacer{}จตฺตาโร ธมฺมา สจฺฉิกาตพฺพา, จตฺตาริ สามญฺญผลานิ.\csromanspacer{}\csromanspacer{}ปญฺจ ธมฺมา สจฺฉิกาตพฺพา, ปญฺจ ธมฺมกฺขนฺธา.\csromanspacer{}\csromanspacer{}ฉ ธมฺมา สจฺฉิกาตพฺพา, ฉ อภิญฺญา.\csromanspacer{}\csromanspacer{}สตฺต ธมฺมา สจฺฉิกาตพฺพา, สตฺต ขีณาสวพลานิ.\csromanspacer{}\csromanspacer{}อฏฺฐ ธมฺมา สจฺฉิกาตพฺพา, อฏฺฐ วิโมกฺขา.\csromanspacer{}\csromanspacer{}นว ธมฺมา สจฺฉิกาตพฺพา, นว อนุปุพฺพนิโรธา.\csromanspacer{}\csromanspacer{}ทส ธมฺมา สจฺฉิกาตพฺพา, ทส อเสกฺขา ธมฺมา.}

\prose{สพฺพํ ภิกฺขเว สจฺฉิกาตพฺพํ.\csromanspacer{} กิญฺจ ภิกฺขเว สพฺพํ สจฺฉิกาตพฺพํ? จกฺขุ ภิกฺขเว สจฺฉิกาตพฺพํ, รูปา สจฺฉิกาตพฺพา, จกฺขุวิญฺญาณํ สจฺฉิกาตพฺพํ, จกฺขุสมฺผสฺโส สจฺฉิกาตพฺโพ, ยมฺปิทํ จกฺขุสมฺผสฺสปจฺจยา อุปฺปชฺชติ เวทยิตํ สุขํ วา ทุกฺขํ วา อทุกฺขมสุขํ วา, ตมฺปิ สจฺฉิกาตพฺพํ.\csromanspacer{} โสตํ สจฺฉิกาตพฺพํ, สทฺทา สจฺฉิกาตพฺพา ฯเปฯ.\csromanspacer{} ฆานํ สจฺฉิกาตพฺพํ, คนฺธา สจฺฉิกาตพฺพา.\csromanspacer{} ชิวฺหา สจฺฉิกาตพฺพา, รสา สจฺฉิกาตพฺพา.\csromanspacer{} กาโย สจฺฉิกาตพฺโพ, โผฏฺฐพฺพา สจฺฉิกาตพฺพา.\csromanspacer{} มโน สจฺฉิกาตพฺโพ, ธมฺมา สจฺฉิกาตพฺพา, มโนวิญฺญาณํ สจฺฉิกาตพฺพํ, มโนสมฺผสฺโส สจฺฉิกาตพฺโพ, ยมฺปิทํ มโนสมฺผสฺสปจฺจยา อุปฺปชฺชติ เวทยิตํ สุขํ วา ทุกฺขํ วา อทุกฺขมสุขํ วา, ตมฺปิ สจฺฉิกาตพฺพํ.\csromanspacer{} รูปํ ปสฺสนฺโต สจฺฉิกโรติ, เวทนํ ปสฺสนฺโต สจฺฉิกโรติ, สญฺญํ ปสฺสนฺโต สจฺฉิกโรติ, สงฺขาเร ปสฺสนฺโต สจฺฉิกโรติ, วิญฺญาณํ ปสฺสนฺโต สจฺฉิกโรติ.\csromanspacer{} จกฺขุํ ฯเปฯ.\csromanspacer{} ชรามรณํ ฯเปฯ อมโตคธํ นิพฺพานํ ปริโยสานฏฺเฐน ปสฺสนฺโต สจฺฉิกโรติ.\csromanspacer{} เย เย ธมฺมา สจฺฉิกตา โหนฺติ, เต เต ธมฺมา ผสฺสิตา\footnote{ผุสิตา (สฺยา)} โหนฺติ.\csromanspacer{} ตํญาตฏฺเฐน ญาณํ, ปชานนฏฺเฐน ปญฺญา, เตน วุจฺจติ “อิเม ธมฺมา สจฺฉิกาตพฺพาติ โสตาวธานํ, ตํปชานนา ปญฺญา สุตมเย ญาณนฺ”ติ. (5)}

\proseitem{30}{กถํ “อิเม ธมฺมา หานภาคิยา”, “อิเม ธมฺมา ฐิติภาคิยา”, “อิเม ธมฺมา วิเสสภาคิยา”, “อิเม ธมฺมา นิพฺเพธภาคิยา”ติ โสตาวธานํ, ตํปชานนา ปญฺญา สุตมเย ญาณํ\csromandash{}}

\prose{\csromanfolio{35}ปฐมสฺส ฌานสฺส ลาภิํ กามสหคตา สญฺญามนสิการา สมุทาจรนฺติ, หานภาคิโย ธมฺโม, ตทนุธมฺมตา สติ สนฺติฏฺฐติ, ฐิติภาคิโย ธมฺโม, อวิตกฺกสหคตา สญฺญามนสิการา สมุทาจรนฺติ, วิเสสภาคิโย ธมฺโม, นิพฺพิทาสหคตา สญฺญามนสิการา สมุทาจรนฺติ วิราคูปสํหิตา, นิพฺเพธภาคิโย ธมฺโม.}

\prose{ทุติยสฺส ฌานสฺส ลาภิํ วิตกฺกสหคตา สญฺญามนสิการา สมุทาจรนฺติ, หานภาคิโย ธมฺโม, ตทนุธมฺมตา สติ สนฺติฏฺฐติ, ฐิติภาคิโย ธมฺโม, อุเปกฺขาสุขสหคตา สญฺญามนสิการา สมุทาจรนฺติ, วิเสสภาคิโย ธมฺโม, นิพฺพิทาสหคตา สญฺญามนสิการา สมุทาจรนฺติ วิราคูปสํหิตา, นิพฺเพธภาคิโย ธมฺโม.}

\prose{ตติยสฺส ฌานสฺส ลาภิํ ปีติสุขสหคตา สญฺญามนสิการา สมุทาจรนฺติ, หานภาคิโย ธมฺโม, ตทนุธมฺมตา สติ สนฺติฏฺฐติ, ฐิติภาคิโย ธมฺโม, อทุกฺขมสุขสหคตา สญฺญามนสิการา สมุทาจรนฺติ, วิเสสภาคิโย ธมฺโม, นิพฺพิทาสหคตา สญฺญามนสิการา สมุทาจรนฺติ วิราคูปสํหิตา, นิพฺเพธภาคิโย ธมฺโม.}

\prose{จตุตฺถสฺส ฌานสฺส ลาภิํ อุเปกฺขาสหคตา สญฺญามนสิการา สมุทาจรนฺติ, หานภาคิโย ธมฺโม, ตทนุธมฺมตา สติ สนฺติฏฺฐติ, ฐิติภาคิโย ธมฺโม, อากาสานญฺจายตนสหคตา สญฺญามนสิการา สมุทาจรนฺติ, วิเสสภาคิโย ธมฺโม, นิพฺพิทาสหคตา สญฺญามนสิการา สมุทาจรนฺติ วิราคูปสํหิตา, นิพฺเพธภาคิโย ธมฺโม.}

\prose{อากาสานญฺจายตนสฺส ลาภิํ รูปสหคตา สญฺญามนสิการา สมุทาจรนฺติ, หานภาคิโย ธมฺโม, ตทนุธมฺมตา สติ สนฺติฏฺฐติ, ฐิติภาคิโย ธมฺโม, วิญฺญาณญฺจายตนสหคตา สญฺญามนสิการา สมุทาจรนฺติ, วิเสสภาคิโย ธมฺโม, นิพฺพิทาสหคตา สญฺญามนสิการา สมุทาจรนฺติ วิราคูปสํหิตา, นิพฺเพธภาคิโย ธมฺโม.}

\prose{วิญฺญาณญฺจายตนสฺส ลาภิํ อากาสานญฺจายตนสหคตา สญฺญามนสิการา สมุทาจรนฺติ, หานภาคิโย ธมฺโม, ตทนุธมฺมตา สติ สนฺติฏฺฐติ, ฐิติภาคิโย ธมฺโม, อากิญฺจญฺญายตนสหคตา สญฺญามนสิการา สมุทาจรนฺติ, วิเสสภาคิโย ธมฺโม, นิพฺพิทาสหคตา สญฺญามนสิการา สมุทาจรนฺติ วิราคูปสํหิตา, นิพฺเพธภาคิโย ธมฺโม.}

\prose{\csromanfolio{36}อากิญฺจญฺญายตนสฺส ลาภิํ วิญฺญาณญฺจายตนสหคตา สญฺญามนสิการา สมุทาจรนฺติ, หานภาคิโย ธมฺโม, ตทนุธมฺมตา สติ สนฺติฏฺฐติ, ฐิติภาคิโย ธมฺโม.\csromanspacer{} เนวสญฺญานาสญฺญายตนสหคตา สญฺญามนสิการา สมุทาจรนฺติ, วิเสสภาคิโย ธมฺโม, นิพฺพิทาสหคตา สญฺญามนสิการา สมุทาจรนฺติ วิราคูปสํหิตา, นิพฺเพธภาคิโย ธมฺโม.\csromanspacer{} ตํญาตฏฺเฐน ญาณํ ปชานนฏฺเฐน ปญฺญา, เตน วุจฺจติ “อิเม ธมฺมา หานภาคิยา”, “อิเม ธมฺมา ฐิติภาคิยา”, “อิเม ธมฺมา วิเสสภาคิยา”, “อิเม ธมฺมา นิพฺเพธภาคิยา”ติ โสตาวธานํ, ตํปชานนา ปญฺญา สุตมเย ญาณํ. (4- 9)}

\proseitem{31}{กถํ “สพฺเพ สงฺขารา อนิจฺจา”, “สพฺเพ สงฺขารา ทุกฺขา”, “สพฺเพ ธมฺมา อนตฺตา”ติ โสตาวธานํ, ตํปชานนา ปญฺญา สุตมเย ญาณํ\csromandash{} รูปํ อนิจฺจํ ขยฏฺเฐน, ทุกฺขํ ภยฏฺเฐน, อนตฺตา อสารกฏฺเฐนาติ โสตาวธานํ, ตํปชานนา ปญฺญา สุตมเย ญาณํ.\csromanspacer{} เวทนา.\csromanspacer{} สญฺญา.\csromanspacer{} สงฺขารา.\csromanspacer{} วิญฺญาณํ.\csromanspacer{} จกฺขุ ฯเปฯ.\csromanspacer{} ชรามรณํ อนิจฺจํ ขยฏฺเฐน, ทุกฺขํ ภยฏฺเฐน, อนตฺตา อสารกฏฺเฐนาติ โสตาวธนํ, ตํปชานนา ปญฺญา สุตมเย ญาณํ.\csromanspacer{} ตํญาตฏฺเฐน ญาณํ, ปชานนฏฺเฐน ปญฺญา, เตน วุจฺจติ “สพฺเพ สงฺขารา อนิจฺจา”, “สพฺเพ สงฺขารา ทุกฺขา”, “สพฺเพ ธมฺมา อนตฺตา”ติ โสตาวธานํ, ตํปชานนา ปญฺญา สุตมเย ญาณํ. (3-12)}

\proseitem{32}{กถํ “อิทํ ทุกฺขํ อริยสจฺจํ”, “อิทํ ทุกฺขสมุทยํ อริยสจฺจํ”, “อิทํ ทุกฺขนิโรธํ อริยสจฺจํ”, “อิทํ ทุกฺขนิโรธคามินี ปฏิปทา อริยสจฺจนฺ”ติ โสตาวธานํ, ตํปชานนา ปญฺญา สุตมเย ญาณํ\csromandash{}}

\proseitem{33}{ตตฺถ กตมํ ทุกฺขํ อริยสจฺจํ? ชาติปิ ทุกฺขา, ชาราปิ ทุกฺขา, มรณมฺปิ ทุกฺขํ, โสกปริเทวทุกฺขโทมนสฺสุปายาสาปิ ทุกฺขา, อปฺปิเยหิ สมฺปโยโค ทุกฺโข, ปิเยหิ วิปฺปโยโค ทุกฺโข, ยมฺปิจฺฉํ น ลภติ, ตมฺปิ ทุกฺขํ, สํขิตฺเตน ปญฺจุปาทานกฺขนฺธา\footnote{ปญฺจุปาทานกฺขนฺธาปิ (ก) สติปฏฺฐานสุตฺเต สจฺจวิภงฺเค จ ปสฺสิตพฺพํ.} ทุกฺขา.}

\prose{ตตฺถ กตมา ชาติ? ยา เตสํ เตสํ สตฺตานํ ตมฺหิ ตมฺหิ สตฺตนิกาเย ชาติ สญฺชาติ โอกฺกนฺติ อภินิพฺพตฺติ ขนฺธานํ ปาตุภาโว อายตนานํ ปฏิลาโภ, อยํ วุจฺจติ ชาติ.}

\prose{\csromanfolio{37}ตตฺถ กตมา ชรา? ยา เตสํ เตสํ สตฺตานํ ตมฺหิ ตมฺหิ สตฺตนิกาเย ชรา ชีรณตา ขณฺฑิจฺจํ ปาลิจฺจํ วลิตฺตจตา อายุโน สํหานิ อินฺทฺริยานํ ปริปาโก, อยํ วุจฺจติ ชรา.}

\prose{ตตฺถ กตมํ มรณํ? ยา เตสํ เตสํ สตฺตานํ ตมฺหา ตมฺหา สตฺตนิกายา จุติ จวนตา เภโท อนฺตรธานํ มจฺจุ มรณํ กาลกิริยา\footnote{กาลํกิริยา (ก)} ขนฺธานํ เภโท กเฬวรสฺส นิกฺเขโป ชีวิตินฺทฺริยสฺสุปจฺเฉโท, อิทํ วุจฺจติ มรณํ.}

\prose{ตตฺถ กตโม โสโก? ญาติพฺยสเนน\footnote{โย ญาติพฺยสเนน (สฺยา) ปสฺส สจฺจวิภงฺเค 104 ปิฏฺเฐ.} วา ผุฏฺฐสฺส, โภคพฺยสเนน วา ผุฏฺฐสฺส, โรคพฺยสเนน วา ผุฏฺฐสฺส, สีลพฺยสเนน วา ผุฏฺฐสฺส, ทิฏฺฐิพฺยสเนน วา ผุฏฺฐสฺส, อญฺญตรญฺญตเรน พฺยสเนน สมนฺนาคตสฺส, อญฺญตรญฺญตเรน ทุกฺขธมฺเมน ผุฏฺฐสฺส โสโก โสจนา โสจิตตฺตํ อนฺโตโสโก อนฺโตปริโสโก เจตโส ปริชฺฌายนา โทมนสฺสํ โสกสลฺลํ, อยํ วุจฺจติ โสโก.}

\prose{ตตฺถ กตโม ปริเทโว? ญาติพฺยสเนน2 วา ผุฏฺฐสฺส, โภคพฺยสเนน วา ผุฏฺฐสฺส, โรคพฺยสเนน วา ผุฏฺฐสฺส, สีลพฺยสเนน วา ผุฏฺฐสฺส, ทิฏฺฐิพฺยสเนน วา ผุฏฺฐสฺส, อญฺญตรญฺญตเรน พฺยสเนน สมนฺนาคตสฺส, อญฺญตรญฺญตเรน ทุกฺขธมฺเมน ผุฏฺฐสฺส อาเทโว ปริเทโว อาเทวนา ปริเทวนา อาเทวิตตฺตํ ปริเทวิตตฺตํ วาจา ปลาโป วิปฺปลาโป ลาลปฺโป ลาลปฺปนา ลาลปฺปิตตฺตํ, อยํ วุจฺจติ ปริเทโว.}

\prose{ตตฺถ กตมํ ทุกฺขํ? ยํ กายิกํ อสาตํ กายิกํ ทุกฺขํ, กายสมฺผสฺสชํ อสาตํ ทุกฺขํ เวทยิตํ, กายสมฺผสฺสชา อสาตา ทุกฺขา เวทนา, อิทํ วุจฺจติ ทุกฺขํ.}

\prose{ตตฺถ กตมํ โทมนสฺสํ? ยํ เจตสิกํ อสาตํ เจตสิกํ ทุกฺขํ, เจโตสมฺผสฺสชํ อสาตํ ทุกฺขํ เวทยิตํ, เจโตสมฺผสฺสชา อสาตา ทุกฺขา เวทนา, อิทํ วุจฺจติ โทมนสฺสํ.}

\prose{ตตฺถ กตโม อุปายาโส? ญาติพฺยสเนน วา ผุฏฺฐสฺส, โภคพฺยสเนน วา ผุฏฺฐสฺส, โรคพฺยสเนน วา ผุฏฺฐสฺส, สีลพฺยสเนน วา ผุฏฺฐสฺส, ทิฏฺฐิพฺยสเนน วา ผุฏฺฐสฺส, อญฺญตรญฺญตเรน พฺยสเนน สมนฺนาคตสฺส, \csromanfolio{38}อญฺญตรญฺญตเรน ทุกฺขธมฺเมน ผุฏฺฐสฺส อายาโส อุปายาโส อายาสนา อุปายาสนา อายาสิตตฺตํ อุปายาสิตตฺตํ, อยํ วุจฺจติ อุปายาโส.}

\prose{ตตฺถ กตโม อปฺปิเยหิ สมฺปโยโค ทุกฺโข? อิธ ย’สฺส เต โหนฺติ อนิฏฺฐา อกนฺตา อมนาปา รูปา สทฺทา คนฺธา รสา โผฏฺฐพฺพา, เย วา ปนสฺส เต โหนฺติ อนตฺถกามา อหิตกามา อผาสุกามา อโยคกฺเขมกามา, ยา เตหิ สทฺธิํ สงฺคติ สมาคโม สโมธานํ มิสฺสีภาโว, อยํ วุจฺจติ อปฺปิเยหิ สมฺปโยโค ทุกฺโข.}

\prose{ตตฺถ กตโม ปิเยหิ วิปฺปโยโค ทุกฺโข? อิธ ย’สฺส เต โหนฺติ อิฏฺฐา กนฺตา มนาปา รูปา สทฺทา คนฺธา รสา โผฏฺฐพฺพา, เย วา ปนสฺส เต โหนฺติ อตฺถกามา หิตกามา ผาสุกามา โยคกฺเขมกามา มาตา วา ปิตา วา ภาตา วา ภคินี วา มิตฺตา วา อมจฺจา วา ญาตี วา สาโลหิตา วา, ยา เตหิ สทฺธิํ อสงฺคติ อสมาคโม อสโมธานํ อมิสฺสีภาโว, อยํ วุจฺจติ ปิเยหิ วิปฺปโยโค ทุกฺโข.}

\prose{ตตฺถ กตมํ ยมฺปิจฺฉํ น ลภติ, ตมฺปิ ทุกฺขํ? ชาติธมฺมานํ สตฺตานํ เอวํ อิจฺฉา อุปฺปชฺชติ “อโห วต มยํ น ชาติธมฺมา อสฺสาม, น จ วต โน ชาติ อาคจฺเฉยฺยา”ติ.\csromanspacer{} น โข ปเนตํ อิจฺฉาย ปตฺตพฺพํ, อิทมฺปิ ยมฺปิจฺฉํ น ลภติ, ตมฺปิ ทุกฺขํ.\csromanspacer{} ชราธมฺมานํ สตฺตานํ ฯเปฯ.\csromanspacer{} พฺยาธิธมฺมานํ สตฺตานํ.\csromanspacer{} มรณธมฺมานํ สตฺตานํ.\csromanspacer{} โสกปริเทวทุกฺขโทมนสฺสุปายาส- ธมฺมานํ สตฺตานํ เอวํ อิจฺฉา อุปฺปชฺชติ “อโห วต มยํ น โสกปริเทวทุกฺขโทมนสฺสุปายาสธมฺมา อสฺสาม, น จ วต โน โสกปริเทวทุกฺขโทมนสฺสุปายาสา อาคจฺเฉยฺยุนฺ”ติ, น โข ปเนตํ อิจฺฉาย ปตฺตพฺพํ, อิทมฺปิ ยมฺปิจฺฉํ น ลภติ, ตมฺปิ ทุกฺขํ.}

\prose{ตตฺถ กตเม สํขิตฺเตน ปญฺจุปาทานกฺขนฺธา ทุกฺขา? เสยฺยถิทํ\footnote{เสยฺยถีทํ (สฺยา, สี-ฏฺฐ)}, รูปุปาทานกฺขนฺโธ\footnote{รูปูปาทานกฺขนฺโธ (สฺยา) เอวมุปริปิ.} เวทนุปาทานกฺขนฺโธ สญฺญุปาทานกฺขนฺโธ สงฺขารุปาทานกฺขนฺโธ วิญฺญาณุปาทานกฺขนฺโธ, อิเม วุจฺจนฺติ สํขิตฺเตน ปญฺจุปาทานกฺขนฺธา ทุกฺขา.\csromanspacer{} อิทํ วุจฺจติ ทุกฺขํ อริยสจฺจํ.}

\proseitem{34}{\csromanfolio{39}ตตฺถ กตมํ ทุกฺขสมุทยํ อริยสจฺจํ? ยายํ ตณฺหา โปโนภวิกา\footnote{โปโนพฺภคิกา (สฺยา)} นนฺทิราคสหคตา ตตฺร ตตฺราภินนฺทินี.\csromanspacer{} เสยฺยถิทํ, กามตณฺหา ภวตณฺหา วิภวตณฺหา, สา โข ปเนสา ตณฺหา กตฺถ อุปฺปชฺชมานา อุปฺปชฺชติ, กตฺถ นิวิสมานา นิวิสติ? ยํ โลเก ปิยรูปํ สาตรูปํ, เอตฺเถสา ตณฺหา อุปฺปชฺชมานา อุปฺปชฺชติ, เอตฺถ นิวิสมานา นิวิสติ.\csromanspacer{} กิญฺจ โลเก ปิยรูปํ สาตรูปํ? จกฺขุ โลเก ปิยรูปํ สาตรูปํ, เอตฺเถสา ตณฺหา อุปฺปชฺชมานา อุปฺปชฺชติ, เอตฺถ นิวิสมานา นิวิสติ.\csromanspacer{} โสตํ โลเก.\csromanspacer{} ฆานํ โลเก.\csromanspacer{} ชิวฺหา โลเก.\csromanspacer{} กาโย โลเก.\csromanspacer{} มโน โลเก ปิยรูปํ สาตรูปํ, เอตฺเถสา ตณฺหา อุปฺปชฺชมานา อุปฺปชฺชติ, เอตฺถ นิวิสมานา นิวิสติ.\csromanspacer{}\csromanspacer{}รูปา โลเก ปิยรูปํ สาตรูปํ, เอตฺเถสา ตณฺหา อุปฺปชฺชมานา อุปฺปชฺชติ, เอตฺถ นิวิสมานา นิวิสติ.\csromanspacer{} สทฺทา โลเก ปิยรูปํ สาตรูปํ ฯเปฯ.\csromanspacer{} ธมฺมา โลเก.\csromanspacer{}\csromanspacer{}จกฺขุวิญฺญาณํ โลเก ฯเปฯ.\csromanspacer{} มโนวิญฺญาณํ โลเก.\csromanspacer{}\csromanspacer{}จกฺขุสมฺผสฺโส โลเก ฯเปฯ.\csromanspacer{} มโนสมฺผสฺโส โลเก.\csromanspacer{}\csromanspacer{}จกฺขุสมฺผสฺสชา เวทนา โลเก ฯเปฯ.\csromanspacer{} มโนสมฺผสฺสชา เวทนา โลเก.\csromanspacer{}\csromanspacer{}รูปสญฺญา โลเก ฯเปฯ.\csromanspacer{} ธมฺมสญฺญา โลเก.\csromanspacer{}\csromanspacer{}รูปสญฺเจตนา โลเก ฯเปฯ.\csromanspacer{} ธมฺมสญฺเจตนา โลเก.\csromanspacer{}\csromanspacer{}รูปตณฺหา โลเก ฯเปฯ.\csromanspacer{} ธมฺมตณฺหา โลเก.\csromanspacer{}\csromanspacer{}รูปวิตกฺโก โลเก ฯเปฯ.\csromanspacer{} ธมฺมวิตกฺโก โลเก.\csromanspacer{}\csromanspacer{}รูปวิจาโร โลเก ฯเปฯ.\csromanspacer{} ธมฺมวิจาโร โลเก ปิยรูปํ สาตรูปํ, เอตฺเถสา ตณฺหา อุปฺปชฺชมานา อุปฺปชฺชติ, เอตฺถ นิวิสมานา นิวิสติ.\csromanspacer{} อิทํ วุจฺจติ ทุกฺขสมุทยํ อริยสจฺจํ.}

\proseitem{35}{ตตฺถ กตมํ ทุกฺขนิโรธํ อริยสจฺจํ? โย ตสฺสาเยว ตณฺหาย อเสสวิราคนิโรโธ จาโค ปฏินิสฺสคฺโค มุตฺติ อนาลโย, สา โข ปเนสา ตณฺหา กตฺถ ปหียมานา ปหียติ, กตฺถ นิรุชฺฌมานา นิรุชฺฌติ? ยํ โลเก ปิยรูปํ สาตรูปํ, เอตฺเถสา ตณฺหา ปหียมานา ปหียติ, เอตฺถ นิรุชฺฌมานา นิรุชฺฌติ.\csromanspacer{} กิญฺจ โลเก ปิยรูปํ สาตรูปํ? จกฺขุ โลเก ปิยรูปํ สาตรูปํ, เอตฺเถสา ตณฺหา ปหียมานา ปหียติ, เอตฺถ นิรุชฺฌมานา นิรุชฺฌติ ฯเปฯ.\csromanspacer{} ธมฺมวิจาโร โลเก ปิยรูปํ สาตรูปํ, เอตฺเถสา ตณฺหา ปหียมานา ปหียติ, เอตฺถ นิรุชฺฌมานา นิรุชฺฌติ.\csromanspacer{} อิทํ วุจฺจติ ทุกฺขนิโรธํ อริยสจฺจํ.}

\proseitem{36}{\csromanfolio{40}ตตฺถ กตมํ ทุกฺขนิโรธคามินี ปฏิปทา อริยสจฺจ? อยเมว อริโย อฏฺฐงฺคิโก มคฺโค.\csromanspacer{} เสยฺยถิทํ, สมฺมาทิฏฺฐิ สมฺมาสงฺกปฺโป สมฺมาวาจา สมฺมากมฺมนฺโต สมฺมา-อาชีโว สมฺมาวายาโม สมฺมาสติ สมฺมาสมาธิ.}

\prose{ตตฺถ กตมา สมฺมาทิฏฺฐิ? ทุกฺเข ญาณํ, ทุกฺขสมุทเย ญาณํ, ทุกฺขนิโรเธ ญาณํ, ทุกฺขนิโรธคามินิยา ปฏิปทาย ญาณํ, อยํ วุจฺจติ สมฺมาทิฏฺฐิ.}

\prose{ตตฺถ กตโม สมฺมาสงฺกปฺโป? เนกฺขมฺมสงฺกปฺโป อพฺยาปาทสงฺกปฺโป อวิหิํสาสงฺกปฺโป, อยํ วุจฺจติ สมฺมาสงฺกปฺโป.}

\prose{ตตฺถ กตมา สมฺมาวาจา? มุสาวาทา เวรมณี\footnote{เวรมณิ (ก)}, ปิสุณาย วาจาย เวรมณี, ผรุสาย วาจาย เวรมณี, สมฺผปฺปลาปา เวรมณี, อยํ วุจฺจติ สมฺมาวาจา.}

\prose{ตตฺถ กตโม สมฺมากมฺมนฺโต? ปาณาติปาตา เวรมณี, อทินฺนาทานา เวรมณี, กาเมสุมิจฺฉาจารา เวรมณี, อยํ วุจฺจติ สมฺมากมฺมนฺโต.}

\prose{ตตฺถ กตโม สมฺมา-อาชีโว? อิธ อริยภาวโก มิจฺฉา-อาชีวํ ปหาย สมา-อาชีเวน ชีวิกํ\footnote{ชีวิตํ (ก)} กปฺเปติ, อยํ วุจฺจติ สมฺมา-อาชีโว.}

\prose{ตตฺถ กตโม สมฺมาวายาโม? อิธ ภิกฺขุ อนุปฺปนฺนานํ ปาปกานํ อกุสลานํ ธมฺมานํ อนุปฺปาทาย ฉนฺทํ ชเนติ วายมติ, วีริยํ อารภติ, จิตฺตํ ปคฺคณฺหาติ ปทหติ.\csromanspacer{} อุปฺปนฺนานํ ปาปกานํ อกุสลานํ ธมฺมานํ ปหานาย ฯเปฯ.\csromanspacer{} อนุปฺปนฺนานํ กุสลานํ ธมฺมานํ อุปฺปาทาย ฯเปฯ.\csromanspacer{} อุปฺปนฺนานํ กุสลานํ ธมฺมานํ ฐิติยา อสมฺโมสาย ภิยฺโยภาวาย เวปุลฺลาย ภาวนาย ปาริปูริยา ฉนฺทํ ชเนติ วายมติ, วีริยํ อารภติ, จิตฺตํ ปคฺคณฺหาติ ปทหติ, อยํ วุจฺจติ สมฺมาวายาโม.}

\prose{ตตฺถ กตมา สมฺมาสติ? อิธ ภิกฺขุ กาเย กายานุปสฺสี วิหรติ อาตาปี สมฺปชาโน สติมา วิเนยฺย โลเก อภิชฺฌาโทมนสฺสํ.\csromanspacer{} เวทนาสุ ฯเปฯ.\csromanspacer{} จิตฺเต ฯเปฯ.\csromanspacer{} ธมฺเมสุ ธมฺมานุปสฺสี วิหรติ อาตาปี สมฺปชาโน สติมา วิเนยฺย โลเก อภิชฺฌาโทมนสฺสํ, อยํ วุจฺจติ สมฺมาสติ.}

\prose{ตตฺถ กตโม สมฺมาสมาธิ? อิธ ภิกฺขุ วิวิจฺเจว กาเมหิ วิวิจฺจ อกุสเลหิ ธมฺเมหิ สวิตกฺกํ สวิจารํ วิเวกชํ ปีติสุขํ ปฐมํ ฌานํ \csromanfolio{41}อุปสมฺปชฺช วิหรติ.\csromanspacer{} วิตกฺกวิจารานํ วูปสมา อชฺฌตฺตํ สมฺปสาทนํ เจตโส เอโกทิภาวํ อวิตกฺกํ อวิจารํ สมาธิชํ ปีติสุขํ ทุติยํ ฌานํ อุปสมฺปชฺช วิหรติ.\csromanspacer{} ปีติยา จ วิราคา อุเปกฺขโก จ วิหรติ สโต จ สมฺปชาโน สุขญฺจ กาเยน ปฏิสํเวเทติ, ยํ ตํ อริยา อาจิกฺขนฺติ “อุเปกฺขโก สติมา สุขวิหารี”ติ, ตติยํ ฌานํ อุปสมฺปชฺช วิหรติ.\csromanspacer{} สุขสฺส จ ปหานา ทุกฺขสฺส จ ปหานา ปุพฺเพว โสมนสฺสโทมนสฺสานํ อตฺถงฺคมา อทุกฺขมสุขํ อุเปกฺขาสติปาริสุทฺธิํ จตุตฺถํ ฌานํ อุปสมฺปชฺช วิหรติ.\csromanspacer{} อยํ วุจฺจติ สมฺมาสมาธิ.\csromanspacer{} อิทํ วุจฺจติ ทุกฺขนิโรธคามินี ปฏิปทา อริยสจฺจํ.\csromanspacer{}\csromanspacer{}ตํญาตฏฺเฐน ญาณํ, ปชานนฏฺเฐน ปญฺญา, เตน วุจฺจติ “อิทํ ทุกฺขํ อริยสจฺจํ”, “อิทํ ทุกฺขสมุทยํ อริยสจฺจํ”, “อิทํ ทุกฺขนิโรธํ อริยสจฺจํ”, “อิทํ ทุกฺขนิโรธคามินี ปฏิปทา อริยสจฺจนฺ”ติ โสตาวธานํ, ตํปชานนา ปญฺญา สุตมเย ญาณํ.\csromanspacer{} เอวํ โสตาวธาเน ปญฺญา สุตมเย ญาณํ. (4-16)}

\csromanheader{2}{สุตมยญาณนิทฺเทโส ปฐโม.}
\csromansectionrule
\csromanmatikaanchor{mtk.4}
\csromantocmarkref{subsection}{2. สีลมยญาณนิทฺเทส}{mtk.4}
\bookmark[dest={mtk.4},level=2]{2. สีลมยญาณนิทฺเทส}
\csromanheader{2}{2. สีลมยญาณนิทฺเทส}
\proseitem{37}{กถํ สุตฺวาน สํวเร ปญฺญา สีลมเย ญาณํ\csromandash{}ปญฺจ สีลานิ ปริยนฺตปาริสุทฺธิสีลํ อปริยนฺตปาริสุทฺธิสีลํ ปริปุณฺณปาริสุทฺธิสีลํ อปรามฏฺฐปาริสุทฺธิสีลํ ปฏิปฺปสฺสทฺธิปาริสุทฺธิสีลนฺติ.}

\prose{ตตฺถ กตมํ ปริยนฺตปาริสุทฺธิสีลํ? อนุปสมฺปนฺนานํ ปริยนฺตสิกฺขาปทานํ, อิทํ ปริยนฺตปาริสุทฺธิสีลํ.\csromanspacer{}\csromanspacer{}กตมํ อปริยนฺตปาริสุทฺธิสีลํ? อุปสมฺปนฺนานํ อปริยนฺตสิกฺขาปทานํ, อิทํ อปริยนฺตปาริสุทฺธิสีลํ.\csromanspacer{}\csromanspacer{}กตมํ ปริปุณฺณปาริสุทฺธิสีลํ? ปุถุชฺชนกลฺยาณกานํ กุสลธมฺเม ยุตฺตานํ เสกฺขปริยนฺเต\footnote{เสขปริยนฺเต (ก)} ปริปูรการีนํ กาเย จ ชีวิเต จ อนเปกฺขานํ ปริจฺจตฺตชีวิตานํ, อิทํ ปริปุณฺณปาริสุทฺธิสีลํ.\csromanspacer{}\csromanspacer{}กตมํ อปรามฏฺฐปาริสุทฺธิสีลํ? สตฺตนฺนํ เสกฺขานํ, อิทํ อปรามฏฺฐปาริสุทฺธิสีลํ.\csromanspacer{}\csromanspacer{}กตมํ ปฏิปฺปสฺสทฺธิปาริสุทฺธิสีลํ? ตถาคตสาวกานํ ขีณาสวานํ ปจฺเจกพุทฺธานํ ตถาคตานํ อรหนฺตานํ สมฺมาสมฺพุทฺธานํ, อิทํ ปฏิปฺปสฺสทฺธิปาริสุทฺธิสีลํ.}

\proseitem{38}{\csromanfolio{42}อตฺถิ สีลํ ปริยนฺตํ, อตฺถิ สีลํ อปริยนฺตํ.\csromanspacer{} ตตฺถ กตมํ ตํ สีลํ ปริยนฺตํ? อตฺถิ สีลํ ลาภปริยนฺตํ, อตฺถิ สีลํ ยสปริยนฺตํ, อตฺถิ สีลํ ญาติปริยนฺตํ, อตฺถิ สีลํ องฺคปริยนฺตํ, อตฺถิ สีลํ ชีวิตปริยนฺตํ.}

\prose{กตมํ ตํ สีลํ ลาภปริยนฺตํ? อิเธกจฺโจ ลาภเหตุ ลาภปจฺจยา ลาภการณา ยถาสมาทินฺนํ\footnote{ยถาสมาทิณฺณํ (ก)} สิกฺขาปทํ วีติกฺกมติ, อิทํ ตํ สีลํ ลาภปริยนฺตํ.\csromanspacer{}\csromanspacer{}กตมํ ตํ สีลํ ยสปริยนฺตํ? อิเธกจฺโจ ยสเหตุ ยสปจฺจยา ยสการณา ยถาสมาทินฺนํ สิกฺขาปทํ วีติกฺกมติ, อิทํ ตํ สีลํ ยสปริยนฺตํ.\csromanspacer{}\csromanspacer{}กตมํ ตํ สีลํ ญาติปริยนฺตํ? อิเธกจฺโจ ญาติเหตุ ญาติปจฺจยา ญาติการณา ยถาสมาทินฺนํ สิกฺขาปทํ วีติกฺกมติ, อิทํ ตํ สีลํ ญาติปริยนฺตํ.\csromanspacer{}\csromanspacer{}กตมํ ตํ สีลํ องฺคปริยนฺตํ? อิเธกจฺโจ องฺคเหตุ องฺคปจฺจยา องฺคการณา ยถาสมาทินฺนํ สิกฺขาปทํ วีติกฺกมติ, อิทํ ตํ สีลํ องฺคปริยนฺตํ.\csromanspacer{}\csromanspacer{}กตมํ ตํ สีลํ ชีวิตปริยนฺตํ? อิเธกจฺโจ ชีวิตเหตุ ชีวิตปจฺจยา ชีวิตการณา ยถาสมาทินฺนํ สิกฺขาปทํ วีติกฺกมติ, อิทํ ตํ สีลํ ชีวิตปริยนฺตํ.\csromanspacer{} เอวรูปานิ สีลานิ ขณฺฑานิ ฉิทฺทานิ สพลานิ กมฺมาสานิ น ภุชิสฺสานิ น วิญฺญุปฺปสตฺถานิ\footnote{น วิญฺญูปสตฺถานิ (ก)} ปรามฏฺฐานิ อสมาธิสํวตฺตนิกานิ\footnote{น สมาธิสํวตฺตนิกานิ (ก)} น อวิปฺปฏิสารวตฺถุกานิ น ปาโมชฺชวตฺถุกานิ น ปีติวตฺถุกานิ น ปสฺสทฺธิวตฺถุกานิ น สุขวตฺถุกานิ น สมาธิวตฺถุกานิ น ยถาภูตญาณทสฺสนวตฺถุกานิ น เอกนฺตนิพฺพิทาย น วิราคาย น นิโรธาย น อุปสมาย น อภิญฺญาย น สมฺโพธาย น นิพฺพานาย สํวตฺตนฺติ, อิทํ ตํ สีลํ ปริยนฺตํ.}

\prose{กตมํ ตํ สีลํ อปริยนฺตํ? อตฺถิ สีลํ น ลาภปริยนฺตํ, อตฺถิ สีลํ น ยสปริยนฺตํ, อตฺถิ สีลํ น ญาติปริยนฺตํ, อตฺถิ สีลํ น องฺคปริยนฺตํ, อตฺถิ สีลํ น ชีวิตปริยนฺตํ.}

\prose{กตมํ ตํ สีลํ น ลาภปริยนฺตํ? อิเธกจฺโจ ลาภเหตุ ลาภปจฺจยา ลาภการณา ยถาสมาทินฺนํ สิกฺขาปทํ วีติกฺกมาย จิตฺตมฺปิ น อุปฺปาเทติ, กิํ โส วีติกฺกมิสฺสติ, อิทํ ตํ สีลํ น ลาภปริยนฺตํ.\csromanspacer{}\csromanspacer{}กตมํ ตํ สีลํ น ยสปริยนฺตํ? อิเธกจฺโจ ยสเหตุ ยสปจฺจยา ยสการณา ยถาสมาทินฺนํ สิกฺขาปทํ วีติกฺกมาย จิตฺตมฺปิ น \csromanfolio{43}อุปฺปาเทติ, กิํ โส วีติกฺกมิสฺสติ, อิทํ ตํ สีลํ น ยสปริยนฺตํ.\csromanspacer{}\csromanspacer{}กตมํ ตํ สีลํ น ญาติปริยนฺตํ? อิเธกจฺโจ ญาติเหตุ ญาติปจฺจยา ญาติการณา ยถาสมาทินฺนํ สิกฺขาปทํ วีติกฺกมาย จิตฺตมฺปิ น อุปฺปาเทติ, กิํ โส วีติกฺกมิสฺสติ, อิทํ ตํ สีลํ น ญาติปริยนฺตํ.\csromanspacer{}\csromanspacer{}กตมํ ตํ สีลํ น องฺคปริยนฺตํ? อิเธกจฺโจ องฺคเหตุ องฺคปจฺจยา องฺคการณา ยถาสมาทินฺนํ สิกฺขาปทํ วีติกฺกมาย จิตฺตมฺปิ น อุปฺปาเทติ, กิํ โส วีติกฺกมิสฺสติ, อิทํ ตํ สีลํ น องฺคปริยนฺตํ.\csromanspacer{}\csromanspacer{}กตมํ ตํ สีลํ น ชีวิตปริยนฺตํ? อิเธกจฺโจ ชีวิตเหตุ ชีวิตปจฺจยา ชีวิตการณา ยถาสมาทินฺนํ สิกฺขาปทํ วีติกฺกมาย จิตฺตมฺปิ น อุปฺปาเทติ, กิํ โส วีติกฺกมิสฺสติ, อิทํ ตํ สีลํ น ชีวิตปริยนฺตํ.\csromanspacer{} เอวรูปานิ สีลานิ อขณฺฑานิ อจฺฉิทฺทานิ อสพลานิ อกมฺมาสานิ ภุชิสฺสานิ วิญฺญุปฺปสตฺถานิ อปรามฏฺฐานิ สมาธิสํวตฺตนิกานิ อวิปฺปฏิสารวตฺถุกานิ ปาโมชฺชวตฺถุกานิ ปีติวตฺถุกานิ ปสฺสทฺธิวตฺถุกานิ สฺ อุขวตฺถุกานิ สมาธิวตฺถุกานิ ยถาภูตญาณทสฺสนวตฺถุกานิ เอกนฺตนิพฺพิทาย วิราคาย นิโรธาย อุปสมาย อภิญฺญาย สมฺโพธาย นิพฺพานาย สํวตฺตนฺติ, อิทํ ตํ สีลํ อปริยนฺตํ.}

\proseitem{39}{กิํ สีลํ, กติ สีลานิ, กิํสมุฏฺฐานํ สีลํ, กติธมฺมสโมธานํ สีลํ.}

\prose{\textbf{กิํ สีลนฺ}ติ เจตนา สีลํ, เจตสิกํ สีลํ, สํวโร สีลํ, อวีติกฺกโม สีลํ.\csromanspacer{}\csromanspacer{}\textbf{กติ สีลานี}ติ ตีณิ สีลานิ, กุสลสีลํ, อกุสลสีลํ, อพฺยากตสีลํ.\csromanspacer{}\csromanspacer{}กฺ\textbf{อิํสมุฏฺฐานํ สีลนฺ}ติ กุสลจิตฺตสมุฏฺฐานํ กุสลสีลํ, อกุสลจิตฺตสมุฏฺฐานํ อกุสลสีลํ, อพฺยากตจิตฺตสมุฏฺฐานํ อพฺยากตสีลํ.\csromanspacer{}\csromanspacer{}\textbf{กติธมฺมสโมธานํ สีลนฺ}ติ สํวรสโมธานํ สีลํ, อวีติกฺกมสโมธานํ สีลํ, ตถาภาเว ชาตเจตนาสโมธานํ สีลํ.}

\proseitem{40}{ปาณาติปาตํ สํวรฏฺเฐน สีลํ, อวีติกฺกมฏฺเฐน สีลํ.\csromanspacer{} อทินฺนาทานํ สํวรฏฺเฐน สีลํ, อวีติกฺกมฏฺเฐน สีลํ.\csromanspacer{} กาเมสุมิจฺฉาจารํ สํวรฏฺเฐน สีลํ, อวีติกฺกมฏฺเฐน สีลํ.\csromanspacer{} มุสาวาทํ สํวรฏฺเฐน สีลํ, อวีติกฺกมฏฺเฐน สีลํ.\csromanspacer{} ปิสุณํ วาจํ\footnote{ปิสุณาวาจํ (สฺยา, ก) ที 1. 4 ปิฏฺเฐ ปสฺสิตพฺพา.} สํวรฏฺเฐน สีลํ, อวีติกฺกมฏฺเฐน สีลํ.\csromanspacer{} ผรุสํ วาจํ\footnote{ผรุสวาจํ (สฺยา, ก)} สํวรฏฺเฐน สีลํ, อวีติกฺกมฏฺเฐน สีลํ.\csromanspacer{} สมฺผปฺปลาปํ สํวรฏฺเฐน สีลํ, อวีติกฺกมฏฺเฐน \csromanfolio{44}สีลํ.\csromanspacer{} อภิชฺฌํ สํวรฏฺเฐน สีลํ, อวีติกฺกมฏฺเฐน สีลํ.\csromanspacer{} พฺยาปาทํ สํวรฏฺเฐน สีลํ, อวีติกฺกมฏฺเฐน สีลํ.\csromanspacer{} มิจฺฉาทิฏฺฐิํ สํวรฏฺเฐน สีลํ, อวีติกฺกมฏฺเฐน สีลํ.}

\proseitem{41}{เนกฺขมฺเมน กามจฺฉนฺทํ สํวรฏฺเฐน สีลํ, อวีติกฺกมฏฺเฐน สีลํ.\csromanspacer{} อพฺยาปาเทน พฺยาปาทํ สํวรฏฺเฐน สีลํ, อวีติกฺกมฏฺเฐน สีลํ.\csromanspacer{} อาโลกสญฺญาย ถินมิทฺธํ.\csromanspacer{} อวิกฺเขปฏฺเฐน อุทฺธจฺจํ.\csromanspacer{} ธมฺมววตฺถาเนน วิจิกิจฺฉํ.\csromanspacer{} ญาเณน อวิชฺชํ.\csromanspacer{} ปาโมชฺเชน อรติํ.}

\prose{ปฐเมน ฌาเนน\footnote{ปฐมชฺฌาเนน (สฺยา) เอวมีทิเสสุ ฐาเนสุ.} นีวรเณ.\csromanspacer{} ทุติเยน ฌาเนน วิตกฺกวิจาเร.\csromanspacer{} ตติเยน ฌาเนน ปีติํ.\csromanspacer{} จตุตฺเถน ฌาเนน สุขทุกฺขํ.\csromanspacer{} อากาสานญฺจายตนสมาปตฺติยา รูปสญฺญํ ปฏิฆสญฺญํ นานตฺตสญฺญํ.\csromanspacer{} วิญฺญาณญฺจายตนสมาปตฺติยา อากาสานญฺจายตนสญฺญํ.\csromanspacer{} อากิญฺจญฺญายตนสมาปตฺติยา วิญฺญาณญฺจายตนสญฺญํ.\csromanspacer{} เนวสญฺญานาสญฺญายตนสมาปตฺติยา อากิญฺจญฺญายตนสญฺญํ.}

\prose{อนิจฺจานุปสฺสนาย นิจฺจสญฺญํ.\csromanspacer{} ทุกฺขานุปสฺสนาย สุขสญฺญํ.\csromanspacer{} อนตฺตานุปสฺสนาย อตฺตสญฺญํ.\csromanspacer{} นิพฺพิทานุปสฺสนาย นนฺทิํ.\csromanspacer{} วิราคานุปสฺสนาย ราคํ.\csromanspacer{} นิโรธานุปสฺสนาย สมุทยํ.\csromanspacer{} ปฏินิสฺสคฺคานุปสฺสนาย อาทานํ.\csromanspacer{} ขยานุปสฺสนาย ฆนสญฺญํ.\csromanspacer{} วยานุปสฺสนาย อายูหนํ.\csromanspacer{} วิปริณามานุปสฺสนาย ธุวสญฺญํ.\csromanspacer{} อนิมิตฺตานุปสฺสนาย นิมิตฺตํ.\csromanspacer{} อปฺปณิหิตานุปสฺสนาย ปณิธิํ.\csromanspacer{} สุญฺญตานุปสฺสนาย อภินิเวสํ.\csromanspacer{} อธิปญฺญาธมฺมวิปสฺสนาย สาราทานาภินิเวสํ.\csromanspacer{} ยถาภูตญาณทสฺสเนน สมฺโมหาภินิเวสํ.\csromanspacer{} อาทีนวานุปสฺสนาย อาลยาภินิเวสํ.\csromanspacer{} ปฏิสงฺขานุปสฺสานาย อปฺปฏิสงฺขํ.\csromanspacer{} วิวฏฺฏนานุปสฺสนาย สญฺโญคาภินิเวสํ.}

\prose{โสตาปตฺติมคฺเคน ทิฏฺเฐกฏฺเฐ กิเลเส.\csromanspacer{} สกทาคามิมคฺเคน โอฬาริเก กิเลเส.\csromanspacer{} อนาคามิมคฺเคน อณุสหคเต กิเลเส.\csromanspacer{} อรหตฺตมคฺเคน สพฺพกิเลเส สํวรฏฺเฐน สีลํ, อวีติกฺกมฏฺเฐน สีลํ.}

\prose{ปญฺจ สีลานิ ปาณาติปาตสฺส ปหานํ สีลํ, เวรมณี สีลํ, เจตนา สีลํ, สํวโร สีลํ, อวีติกฺกโม สีลํ.\csromanspacer{} เอวรูปานิ สีลานิ จิตฺตสฺส อวิปฺปฏิสาราย สํวตฺตนฺติ, ปาโมชฺชาย สํวตฺตนฺติ, ปีติยา สํวตฺตนฺติ, \csromanfolio{45}ปสฺสทฺธิยา สํวตฺตนฺติ, โสมนสฺสาย สํวตฺตนฺติ, อาเสวนาย สํวตฺตนฺติ, ภาวนาย สํวตฺตนฺติ, พหุลีกมฺมาย สํวตฺตนฺติ, อลงฺการาย สํวตฺตนฺติ, ปริกฺขาราย สํวตฺตนฺติ, ปริวาราย สํวตฺตนฺติ, ปาริปูริยา สํวตฺตนฺติ, เอกนฺตนิพฺพิทาย วิราคาย นิโรธาย อุปสมาย อภิญฺญาย สมฺโพธาย นิพฺพานาย สํวตฺตนฺติ.}

\prose{เอวรูปานํ สีลานํ สํวรปาริสุทฺธิ อธิสีลํ, สํวรปาริสุทฺธิยา ฐิตํ จิตฺตํ น วิกฺเขปํ คจฺฉติ, อวิกฺเขปปาริสุทฺธิ อธิจิตฺตํ, สํวรปาริสุทฺธิํ สมฺมา ปสฺสติ, อวิกฺเขปปาริสุทฺธิํ สมฺมา ปสฺสติ.\csromanspacer{} ทสฺสนปาริสุทฺธิ อธิปญฺญา.\csromanspacer{} โย ตตฺถ สํวรฏฺโฐ, อยํ อธิสีลสิกฺขา.\csromanspacer{} โย ตตฺถ อวิกฺเขปฏฺโฐ, อยํ อธิจิตฺตสิกฺขา.\csromanspacer{} โย ตตฺถ ทสฺสนฏฺโฐ, อยํ อธิปญฺญาสิกฺขา.}

\prose{อิมา ติสฺโส สิกฺขาโย อาวชฺชนฺโต สิกฺขติ, ชานนฺโต สิกฺขติ, ปสฺสนฺโต สิกฺขติ, ปจฺจเวกฺขนฺโต สิกฺขติ, จิตฺตํ อธิฏฺฐหนฺโต สิกฺขติ, สทฺธาย อธิมุจฺจนฺโต สิกฺขติ, วีริยํ ปคฺคณฺหนฺโต สิกฺขติ, สติํ อุปฏฺฐเปนฺโต สิกฺขติ, จิตฺตํ สมาทหนฺโต สิกฺขติ, ปญฺญาย ปชานนฺโต สิกฺขติ, อภิญฺเญยฺยํ อภิชานนฺโต สิกฺขติ, ปริญฺเญยฺยํ ปริชานนฺโต สิกฺขติ, ปหาตพฺพํ ปชหนฺโต สิกฺขติ, สจฺฉิกาตพฺพํ สจฺฉิกโรนฺโต สิกฺขติ, ภาเวตพฺพํ ภาเวนฺโต สิกฺขติ.}

\prose{ปญฺจ สีลานิ อทินฺนาทานสฺส.\csromanspacer{} กาเมสุมิจฺฉาจารสฺส.\csromanspacer{} มุสาวาทสฺส.\csromanspacer{} ปิสุณาย วาจาย.\csromanspacer{} ผรุสาย วาจาย.\csromanspacer{} สมฺผปฺปลาปสฺส.\csromanspacer{} อภิชฺฌาย.\csromanspacer{} พฺยาปาทสฺส.\csromanspacer{} มิจฺฉาทิฏฺฐิยา.}

\prose{เนกฺขมฺเมน กามจฺฉนฺทสฺส.\csromanspacer{} อพฺยาปาเทน พฺยาปาทสฺส.\csromanspacer{} อาโลกสญฺญาย ถินมิทฺธสฺส.\csromanspacer{} อวิกฺเขเปน อุทฺธจฺจสฺส.\csromanspacer{} ธมฺมววตฺถาเนน วิจิกิจฺฉาย.\csromanspacer{} ญาเณน อวิชฺชาย.\csromanspacer{} ปาโมชฺเชน อรติยา.}

\prose{ปฐเมน ฌาเนน นีวรณานํ.\csromanspacer{} ทุติเยน ฌาเนน วิตกฺกวิจารานํ.\csromanspacer{} ตติเยน ฌาเนน ปีติยา.\csromanspacer{} จตุตฺเถน ฌาเนน สุขทุกฺขานํ.\csromanspacer{} อากาสานญฺจายตนสมาปตฺติยา รูปสญฺญาย ปฏิฆสญฺญาย นานตฺตสญฺญาย.\csromanspacer{} วิญฺญาณญฺจายตนสมาปตฺติยา อากาสานญฺจายตนสญฺญาย.\csromanspacer{} อากิญฺจญฺญายตนสมาปตฺติยา วิญฺญณญฺจายตนสญฺญาย.\csromanspacer{} เนวสญฺญานาสญฺญายตนสมาปตฺติยา อากิญฺจญฺญายตนสญฺญาย.}

\prose{\csromanfolio{46}อนิจฺจานุปสฺสนาย นิจฺจสญฺญาย.\csromanspacer{} ทุกฺขานุปสฺสนาย สุขสญฺญาย, อนตฺตานุปสฺสนาย อตฺตสญฺญาย.\csromanspacer{} นิพฺพิทานุปสฺสนาย นนฺทิยา.\csromanspacer{} วิราคานุปสฺสนาย ราคสฺส.\csromanspacer{} นิโรธานุปสฺสนาย สมุทยสฺส.\csromanspacer{} ปฏินิสฺสคฺคานุปสฺสนาย อาทานสฺส.\csromanspacer{} ขยานุปสฺสนาย ฆนสญฺญาย.\csromanspacer{} วยานุปสฺสนาย อายูหนสฺส.\csromanspacer{} วิปริณามานุปสฺสนาย ธุวสญฺญาย.\csromanspacer{} อนิมิตฺตานุปสฺสนาย นิมิตฺตสฺส.\csromanspacer{} อปฺปณิหิตานุปสฺสนาย ปณิธิยา.\csromanspacer{} สุญฺญตานุปสฺสนาย อภินิเวสสฺส.\csromanspacer{} อธิปญฺญาธมฺมวิปสฺสนาย สาราคาภินิเวสสฺส.\csromanspacer{} ยถาภูตญาณทสฺสเนน สมฺโมหาภินิเวสสฺส.\csromanspacer{} อาทีนวานุปสฺสนาย อาลยาภินิเวสสฺส.\csromanspacer{} ปฏิสงฺขานุปสฺสนาย อปฺปฏิสงฺขาย.\csromanspacer{} วิวฏฺฏนานุปสฺสนาย สญฺโญคาภินิเวสสฺส.}

\prose{โสตาปตฺติมคฺเคน ทิฏฺเฐกฏฺฐานํ กิเลสานํ.\csromanspacer{} สกทาคามิมคฺเคน โอฬาริกานํ กิเลสานํ.\csromanspacer{} อนาคามิมคฺเคน อณุสหคตานํ กิเลสานํ.\csromanspacer{} อรหตฺตมคฺเคน สพฺพกิเลสานํ ปหานํ สีลํ, เวรมณี สีลํ, เจตนา สีลํ, สํวโร สีลํ, อวีติกฺกโม สีลํ.\csromanspacer{} เอวรูปานิ สีลานิ จิตฺตสฺส อวิปฺปฏิสาราย สํวตฺตนฺติ, ปาโมชฺชาย สํวตฺตนฺติ, ปีติยา สํวตฺตนฺติ, ปสฺสทฺธิยา สํวตฺตนฺติ, โสมนสฺสาย สํวตฺตนฺติ, อาเสวนาย สํวตฺตนฺติ, ภาวนาย สํวตฺตนฺติ, พหุลีกมฺมาย สํวตฺตนฺติ, อลงฺการาย สํวตฺตนฺติ, ปริกฺขาราย สํวตฺตนฺติ, ปริวาราย สํวตฺตนฺติ, ปาริปูริยา สํวตฺตนฺติ, เอกนฺตนิพฺพิทาย วิราคาย นิโรธาย อุปสมาย อภิญฺญาย สมฺโพธาย นิพฺพานาย สํวตฺตนฺติ.}

\proseitem{42}{เอวรูปานํ สีลานํ สํวรปาริสุทฺธิ อธิสีลํ.\csromanspacer{} สํวรปาริสุทฺธิยา ฐิตํ จิตฺตํ อวิกฺเขปํ คจฺฉติ, อวิกฺเขปปาริสุทฺธิ อธิจิตฺตํ.\csromanspacer{} สํวรปาริสุทฺธิํ สมฺมา ปสฺสติ, อวิกฺเขปปาริสุทฺธิํ สมฺมา ปสฺสติ.\csromanspacer{} ทสฺสนปาริสุทฺธิ อธิปญฺญา.\csromanspacer{} โย ตตฺถ สํวรฏฺโฐ, อยํ อธิสีลสิกฺขา.\csromanspacer{} โย ตตฺถ อวิกฺเขปฏฺโฐ, อยํ อธิจิตฺตสิกฺขา.\csromanspacer{} โย ตตฺถ ทสฺสนฏฺโฐ, อยํ อธิปญฺญาสิกฺขา.}

\prose{อิมา ติสฺโส สิกฺขาโย อาวชฺชนฺโต สิกฺขติ, ชานนฺโต สิกฺขติ, ปสฺสนฺโต สิกฺขติ, ปจฺจเวกฺขนฺโต สิกฺขติ, จิตฺตํ อธิฏฺฐหนฺโต สิกฺขติ, สทฺธาย อธิมุจฺจนฺโต สิกฺขติ, วีริยํ ปคฺคณฺหนฺโต สิกฺขติ, สติํ อุปฏฺฐเปนฺโต สิกฺขติ, จิตฺตํ สมาทหนฺโต สิกฺขติ, ปญฺญาย ปชานนฺโต สิกฺขติ, อภิญฺเญยฺยํ อภิชานนฺโต สิกฺขติ, ปริญฺเญยฺยํ ปริชานนฺโต สิกฺขติ, ปหาตพฺพํ ปชหนฺโต สิกฺขติ, สจฺฉิกาตพฺพํ สจฺฉิกโรนฺโต สิกฺขติ, \csromanfolio{47}ภาเวตพฺพํ ภาเวนฺโต สิกฺขติ, ตํญาตฏฺเฐน ญาณํ, ปชานนฏฺเฐน ปญฺญา, เตน วุจฺจติ สุตฺวาน สํวเร ปญฺญา สีลมเย ญาณํ.}

\csromanheader{2}{สีลมยญาณนิทฺเทโส ทุติโย.}
\csromansectionrule
\csromanmatikaanchor{mtk.5}
\csromantocmarkref{subsection}{3. สมาธิภาวนามยญาณนิทฺเทส}{mtk.5}
\bookmark[dest={mtk.5},level=2]{3. สมาธิภาวนามยญาณนิทฺเทส}
\csromanheader{2}{3. สมาธิภาวนามยญาณนิทฺเทส}
\proseitem{43}{กถํ สํวริตฺวา สมาทหเน ปญฺญา สมาธิภาวนามเย ญาณํ\csromandash{}เอโก สมาธิ จิตฺตสฺส เอกคฺคตา.\csromanspacer{}\csromanspacer{}เทฺว สมาธี, โลกิโย สมาธิ โลกุตฺตโร สมาธิ.\csromanspacer{}\csromanspacer{}ตโย สมาธี, สวิตกฺกสวิจาโร สมาธิ อวิตกฺกวิจารมตฺโต สมาธิ อวิตกฺก-อวิจาโร สมาธิ.\csromanspacer{}\csromanspacer{}จตฺตาโร สมาธี, หานภาคิโย สมาธิ ฐิติภาคิโย สมาธิ วิเสสภาคิโย สมาธิ นิพฺเพธภาคิโย สมาธิ.\csromanspacer{}\csromanspacer{}ปญฺจ สมาธี, ปีติผรณตา สุขผรณตา เจโตผรณตา อาโลกผรณตา ปจฺจเวกฺขณานิมิตฺตํ\footnote{ปจฺจเวกฺขนานิมิตฺตํ (ก)}.\csromanspacer{}\csromanspacer{}ฉ สมาธี, พุทฺธานุสฺสติวเสน จิตฺตสฺส เอกคฺคตา อวิกฺเขโป สมาธิ, ธมฺมานุสฺสติวเสน จิตฺตสฺส เอกคฺคตา อวิกฺเขโป สมาธิ, สํฆานุสฺสติวเสน จิตฺตสฺส เอกคฺคตา อวิกฺเขโป สมาธิ, สีลานุสฺสติวเสน จิตฺตสฺส เอกคฺคตา อวิกฺเขโป สมาธิ, จาคานุสฺสติวเสน จิตฺตสฺส เอกคฺคตา อวิกฺเขโป สมาธิ, เทวตานุสฺสติวเสน จิตฺตสฺส เอกคฺคตา อวิกฺเขโป สมาธิ.\csromanspacer{}\csromanspacer{}สตฺต สมาธี, สมาธิกุสลตา สมาธิสฺส สมาปตฺติกุสลตา สมาธิสฺส ฐิติกุสลตา สมาธิสฺส วุฏฺฐานกุสลตา สมาธิสฺส กลฺลตากุสลตา\footnote{กลฺลิตกุสลตา (สฺยา, ก)} สมาธิสฺส โคจรกุสลตา สมาธิสฺส อภินีหารกุสลตา.\csromanspacer{}\csromanspacer{}อฏฺฐ สมาธี, ปถวีกสิณวเสน จิตฺตสฺส เอกคฺคตา อวิกฺเขโป สมาธิ, อาโปกสิณวเสน ฯเปฯ.\csromanspacer{} เตโชกสิณวเสน.\csromanspacer{} วาโยกสิณวเสน.\csromanspacer{} นีลกสิณวเสน.\csromanspacer{} ปีตกสิณวเสน.\csromanspacer{} โลหิตกสิณวเสน.\csromanspacer{} โอทาตกสิณวเสน จิตฺตสฺส เอกคฺคตา อวิกฺเขโป สมาธิ.\csromanspacer{}\csromanspacer{}นว สมาธี, รูปาวจโร สมาธิ อตฺถิ หีโน อตฺถิ มชฺโฌ อตฺถิ ปณีโต, อรูปาวจโร สมาธิ อตฺถิ หีโน อตฺถิ มชฺโฌ อตฺถิ ปณีโต, สุญฺญโต สมาธิ อนิมิตฺโต สมาธิ อปฺปณิหิโต สมาธิ.\csromanspacer{}\csromanspacer{}ทส สมาธี, อุทฺธุมาตกสญฺญาวเสน \csromanfolio{48}จิตฺตสฺส เอกคฺคตา อวิกฺเขโป สมาธิ, วินีลกสญฺญาวเสน.\csromanspacer{} วิปุพฺพกสญฺญาวเสน.\csromanspacer{} วิจฺฉิทฺทกสญฺญาวเสน.\csromanspacer{} วิกฺขายิตกสญฺญาวเสน.\csromanspacer{} วิกฺขิตฺตกสญฺญาวเสน.\csromanspacer{} หตวิกฺขิตฺตกสญฺญาวเสน.\csromanspacer{} โลหิตกสญฺญาวเสน.\csromanspacer{} ปุฬวกสญฺญาวเสน\footnote{ปุฬุวก สญฺญาวเสน (ก)}.\csromanspacer{} อฏฺฐิกสญฺญาวเสน จิตฺตสฺส เอกคฺคตา อวิกฺเขโป สมาธิ.\csromanspacer{} อิเม ปญฺจปญฺญาส สมาธี.}

\proseitem{44}{อปิ จ ปญฺจวีสติ สมาธิสฺส สมาธิฏฺฐา\csromandash{}ปริคฺคหฏฺเฐน สมาธิ, ปริวารฏฺเฐน สมาธิ, ปริปูรฏฺเฐน สมาธิ, เอกคฺคฏฺเฐน สมาธิ, อวิกฺเขปฏฺเฐน สมาธิ, อนาวิลฏฺเฐน สมาธิ, อนิญฺชนฏฺเฐน สมาธิ, วิมุตฺตฏฺเฐน สมาธิ, เอกตฺตุปฏฺฐานวเสน จิตฺตสฺส ฐิตตฺตา สมาธิ.\csromanspacer{} สมํ เอสตีติ สมาธิ, วิสมํ เนสตีติ สมาธิ, สมํ เอสิตตฺตา สมาธิ, วิสมํ เนสิตตฺตา สมาธิ, สมํ อาทิยตีติ สมาธิ, วิสมํ นาทิยตีติ สมาธิ, สมํ อาทินฺนตฺตา สมาธิ, วิสมํ อนาทินฺนตฺตา สมาธิ, สมํ ปฏิปชฺชตีติ สมาธิ, วิสมํ นปฺปฏิปชฺชตีติ สมาธิ, สมํ ปฏิปนฺนตฺตา สมาธิ, วิสมํ นปฺปฏิปนฺนตฺตา สมาธิ, สมํ ฌายตีติ สมาธิ, วิสมํ ฌาเปตีติ สมาธิ, สมํ ฌาตตฺตา สมาธิ, วิสมํ ฌาปิตตฺตา สมาธิ.\csromanspacer{} สโม จ หิโต จ สุโข จาติ สมาธิ.\csromanspacer{} อิเม ปญฺจวีสติ สมาธิสฺส สมาธิฏฺฐา.\csromanspacer{} ตํญาตฏฺเฐน ญาณํ, ปชานนฏฺเฐน ปญฺญา, เตน วุจฺจติ สํวริตฺวา สมาทหเน ปญฺญา สมาธิภาวนามเย ญาณํ.}

\csromanheader{2}{สมาธิภาวนามยญาณนิทฺเทโส ตติโย.}
\csromansectionrule
\csromanmatikaanchor{mtk.6}
\csromantocmarkref{subsection}{4. ธมฺมฏฺฐิติญาณนิทฺเทส}{mtk.6}
\bookmark[dest={mtk.6},level=2]{4. ธมฺมฏฺฐิติญาณนิทฺเทส}
\csromanheader{2}{4. ธมฺมฏฺฐิติญาณนิทฺเทส}
\proseitem{45}{กถํ ปจฺจยปริคฺคเห ปญฺญา ธมฺมฏฺฐิติญาณํ\csromandash{}อวิชฺชา สงฺขารานํ อุปฺปาทฏฺฐิติ จ ปวตฺตฏฺฐิติ จ นิมิตฺตฏฺฐิติ จ อายูหนฏฺฐิติ จ สญฺโญคฏฺฐิติ จ ปลิโพธฏฺฐิติ จ สมุทยฏฺฐิติ จ เหตุฏฺฐิติ จ ปจฺจยฏฺฐิติ จ, อิเมหิ นวหากาเรหิ อวิชฺชา ปจฺจโย, สงฺขารา ปจฺจยสมุปฺปนฺนา, อุโภเปเต ธมฺมา ปจฺจยสมุปฺปนฺนาติ ปจฺจยปริคฺคเห ปญฺญา ธมฺมฏฺฐิติญาณํ.\csromanspacer{} อตีตมฺปิ อทฺธานํ, อนาคตมฺปิ อทฺธานํ อวิชฺชา สงฺขารานํ อุปฺปาทฏฺฐิติ จ ปวตฺตฏฺฐิติ จ นิมิตฺตฏฺฐิติ จ อายูหนฏฺฐิติ จ สญฺโญคฏฺฐิติ จ ปลิโพธฏฺฐิติ จ สมุทยฏฺฐิติ จ \csromanfolio{49}เหตุฏฺฐิติ จ ปจฺจยฏฺฐิติ จ, อิเมหิ นวหากาเรหิ อวิชฺชา ปจฺจโย, สงฺขารา ปจฺจยสมุปฺปนฺนา, อุโภเปเต ธมฺมา ปจฺจยสมุปฺปนฺนาติ ปจฺจยปริคฺคเห ปญฺญา ธมฺมฏฺฐิติญาณํ.}

\prose{สงฺขารา วิญฺญาณสฺส ฯเปฯ.\csromanspacer{} วิญฺญาณํ นามรูปสฺส.\csromanspacer{} นามรูปํ สฬายตนสฺส, สฬายตนํ ผสฺสสฺส.\csromanspacer{} ผสฺโส เวทนาย.\csromanspacer{} เวทนา ตณฺหาย.\csromanspacer{} ตณฺหา อุปาทานสฺส.\csromanspacer{} อุปาทานํ ภวสฺส.\csromanspacer{} ภโว ชาติยา.\csromanspacer{} ชาติ ชรามรณสฺส อุปฺปาทฏฺฐิติ จ ปวตฺตฏฺฐิติ จ นิมิตฺตฏฺฐิติ จ อายูหนฏฺฐิติ จ สญฺโญคฏฺฐิติ จ ปลิโพธฏฺฐิติ จ สมุทยฏฺฐิติ จ เหตุฏฺฐิติ จ ปจฺจยฏฺฐิติ จ, อิเมหิ นวหากาเรหิ ชาติ ปจฺจโย, ชรามรณํ ปจฺจยสมุปฺปนฺนํ, อุโภเปเต ธมฺมา ปจฺจยสมุปฺปนฺนาติ ปจฺจยปริคฺคเห ปญฺญา ธมฺมฏฺฐิติญาณํ.\csromanspacer{} อตีตมฺปิ อทฺธานํ.\csromanspacer{} อนาคตมฺปิ อทฺธานํ ชาติ ชรามรณสฺส อุปฺปาทฏฺฐิติ จ ปวตฺตฏฺฐิติ จ นิมิตฺตฏฺฐิติ จ อายูหนฏฺฐิติ จ สญฺโญคฏฺฐิติ จ ปลิโพธฏฺฐิติ จ สมุทยฏฺฐิติ จ เหตุฏฺฐิติ จ ปจฺจยฏฺฐิติ จ, อิเมหิ นวหากาเรหิ ชาติ ปจฺจโย, ชรามรณ ปจฺจยสมุปฺปนฺนํ, อุโภเปเต ธมฺมา ปจฺจยสมุปฺปนฺนาติ ปจฺจยปริคฺคเห ปญฺญา ธมฺมฏฺฐิติญาณํ.}

\proseitem{46}{อวิชฺชา เหตุ, สงฺขารา เหตุสมุปฺปนฺนา, อุโภเปเต ธมฺมา เหตุสมุปฺปนฺนาติ ปจฺจยปริคฺคเห ปญฺญา ธมฺมฐิติญาณํ.\csromanspacer{} อตีตมฺปิ อทฺธานํ.\csromanspacer{} อนาคตมฺปิ อทฺธานํ อวิชฺชา เหตุ, สงฺขารา เหตุสมุปฺปนฺนา, อุโภเปเต ธมฺมา เหตุสมุปฺปนฺนาติ ปจฺจยปริคฺคเห ปญฺญา ธมฺมฏฺฐิติญาณํ.}

\prose{สงฺขารา เหตุ, วิญฺญาณํ เหตุสมุปฺปนฺนํ ฯเปฯ.\csromanspacer{} วิญฺญาณํ เหตุ, นามรูปํ เหตุสมุปฺปนฺนํ.\csromanspacer{} นามรูปํ เหตุ, สฬายตนํ เหตุสมุปฺปนฺนํ.\csromanspacer{} สฬายตนํ เหตุ, ผสฺโส เหตุสมุปฺปนฺโน.\csromanspacer{} ผสฺโส เหตุ, เวทนา เหตุสมุปฺปนฺนา.\csromanspacer{} เวทนา เหตุ, ตณฺหา เหตุสมุปฺปนฺนา.\csromanspacer{} ตณฺหา เหตุ, อุปาทานํ เหตุสมุปฺปนฺนํ.\csromanspacer{} อุปาทานํ เหตุ, ภโว เหตุสมุปฺปนฺโน.\csromanspacer{} ภโว เหตุ, ชาติ เหตุสมุปฺปนฺนา.\csromanspacer{} ชาติ เหตุ, ชรามรณํ เหตุสมุปฺปนฺนํ, อุโภเปเต ธมฺมา เหตุสมุปฺปนฺนาติ ปจฺจยปริคฺคเห ปญฺญา ธมฺมฏฺฐิติญาณํ.\csromanspacer{} อตีตมฺปิ อทฺธานํ.\csromanspacer{} อนาคตมฺปิ อทฺธานํ ชาติ เหตุ, ชรามรณํ เหตุสมุปฺปนฺนํ, อุโภเปเต ธมฺมา เหตุสมุปฺปนฺนาติ ปจฺจยปริคฺคเห ปญฺญา ธมฺมฏฺฐิติญาณํ.}

\prose{อวิชฺชา ปฏิจฺจา, สงฺขารา ปฏิจฺจสมุปฺปนฺนา, อุโภเปเต ธมฺมา ปฏิจฺจสมุปฺปนฺนาติ ปจฺจยปริคฺคเห ปญฺญา ธมฺมฏฺฐิติญาณํ.\csromanspacer{} อตีตมฺปิ อทฺธานํ.\csromanspacer{} อนาคตมฺปิ อทฺธานํ \csromanfolio{50}อวิชฺชา ปฏิจฺจา, สงฺขารา ปฏิจฺจสมุปฺปนฺนา, อุโภเปเต ธมฺมา ปฏิจฺจสมุปฺปนฺนาติ ปจฺจยปริคฺคเห ปญฺญา ธมฺมฏฺฐิติญาณํ.}

\prose{สงฺขารา ปฏิจฺจา, วิญฺญาณํ ปฏิจฺจสมุปฺปนฺนํ ฯเปฯ.\csromanspacer{} วิญฺญาณํ ปฏิจฺจํ, นามรูปํ ปฏิจฺจสมุปฺปนฺนํ.\csromanspacer{} นามรูปํ ปฏิจฺจํ, สฬายตนํ ปฏิจฺจสมุปฺปนฺนํ.\csromanspacer{} สฬายตนํ ปฏิจฺจํ, ผสฺโส ปฏิจฺจสมุปฺปนฺโน.\csromanspacer{} ผสฺโส ปฏิจฺโจ, เวทนา ปฏิจฺจสมุปฺปนฺนา.\csromanspacer{} เวทนา ปฏิจฺจา, ตณฺหา ปฏิจฺจสมุปฺปนฺนา.\csromanspacer{} ตณฺหา ปฏิจฺจา, อุปาทานํ ปฏิจฺจสมุปฺปนฺนํ.\csromanspacer{} อุปาทานํ ปฏิจฺจํ, ภโว ปฏิจฺจสมุปฺปนฺโน.\csromanspacer{} ภโว ปฏิจฺโจ, ชาติ ปฏิจฺจสมุปฺปนฺนา.\csromanspacer{} ชาติ ปฏิจฺจา, ชรามรณํ ปฏิจฺจสมุปฺปนฺนํ, อุโภเปเต ธมฺมา ปฏิจฺจสมุปฺปนฺนาติ ปจฺจยปริคฺคเห ปญฺญา ธมฺมฏฺฐิติญาณํ.\csromanspacer{} อตีตมฺปิ อทฺธานํ.\csromanspacer{} อนาคตมฺปิ อทฺธานํ ชาติ ปฏิจฺจา, ชรามรณํ ปฏิจฺจสมุปฺปนฺนํ, อุโภเปเต ธมฺมา ปฏิจฺจสมุปฺปนฺนาติ ปจฺจยปริคฺคเห ปญฺญา ธมฺมฏฺฐิติญาณํ.}

\prose{อวิชฺชา ปจฺจโย, สงฺขารา ปจฺจยสมุปฺปนฺนา, อุโภเปเต ธมฺมา ปจฺจยสมุปฺปนฺนาติ ปจฺจยปริคฺคเห ปญฺญา ธมฺมฏฺฐิติญาณํ.\csromanspacer{} อตีตมฺปิ อทฺธานํ.\csromanspacer{} อนาคตมฺปิ อทฺธานํ อวิชฺชา ปจฺจโย, สงฺขารา ปจฺจยสมุปฺปนฺนา, อุโภเปเต ธมฺมา ปจฺจยสมุปฺปนฺนาติ ปจฺจยปริคฺคเห ปญฺญา ธมฺมฏฺฐิติญาณํ.}

\prose{สงฺขารา ปจฺจยา, วิญฺญาณํ ปจฺจยสมุปฺปนฺนํ ฯเปฯ.\csromanspacer{} วิญฺญาณํ ปจฺจโย, นามรูปํ ปจฺจยสมุปฺปนฺนํ.\csromanspacer{} นามรูปํ ปจฺจโย, สฬายตนํ ปจฺจยสมุปฺปนฺนํ.\csromanspacer{} สฬายตนํ ปจฺจโย, ผสฺโส ปจฺจยสมุปฺปนฺโน.\csromanspacer{} ผสฺโส ปจฺจโย, เวทนา ปจฺจยสมุปฺปนฺนา.\csromanspacer{} เวทนา ปจฺจโย, ตณฺหา ปจฺจยสมุปฺปนฺนา.\csromanspacer{} ตณฺหา ปจฺจโย, อุปาทานํ ปจฺจยสมุปฺปนฺนํ.\csromanspacer{} อุปาทานํ ปจฺจโย, ภโว ปจฺจยสมุปฺปนฺโน.\csromanspacer{} ภโว ปจฺจโย, ชาติ ปจฺจยสมุปฺปนฺนา.\csromanspacer{} ชาติ ปจฺจโย, ชรามรณํ ปจฺจยสมุปฺปนฺนํ, อุโภเปเต ธมฺมา ปจฺจยสมุปฺปนฺนาติ ปจฺจยปริคฺคเห ปญฺญา ธมฺมฏฺฐิติญาณํ.\csromanspacer{} อตีตมฺปิ อทฺธานํ.\csromanspacer{} อนาคตมฺปิ อทฺธานํ ชาติ ปจฺจโย, ชรามรณํ ปจฺจยสมุปฺปนฺนํ, อุโภเปเต ธมฺมา ปจฺจยสมุปฺปนฺนาติ ปจฺจยปริคฺคเห ปญฺญา ธมฺมฏฺฐิติญาณํ.}

\proseitem{47}{ปุริมกมฺมภวสฺมิํ โมโห อวิชฺชา, อายูหนา สงฺขารา, นิกนฺติ ตณฺหา, อุปคมนํ อุปาทานํ, เจตนา ภโว, อิเม ปญฺจ ธมฺมา ปุริมกมฺมภวสฺมิํ อิธ ปฏิสนฺธิยา ปจฺจยา.\csromanspacer{} อิธ ปฏิสนฺธิ วิญฺญาณํ, โอกฺกนฺติ นามรูปํ, ปสาโท อายตนํ, ผุฏฺโฐ ผสฺโส, เวทยิตํ เวทนา, อิเม ปญฺจ ธมฺมา อิธุปปตฺติภวสฺมิํ ปุเรกตสฺส กมฺมสฺส ปจฺจยา.\csromanspacer{} อิธ ปริปกฺกตฺตา อายตนานํ โมโห อวิชฺชา, อายูหนา สงฺขารา, นิกนฺติ ตณฺหา, อุปคมนํ อุปาทานํ, \csromanfolio{51}เจตนา ภโว, อิเม ปญฺจ ธมฺมา อิธ กมฺมภวสฺมิํ อายติํ ปฏิสนฺธิยา ปจฺจยา.\csromanspacer{} อายติํ ปฏิสนฺธิ วิญฺญาณํ, โอกฺกนฺติ นามรูปํ, ปสาโท อายตนํ, ผุฏฺโฐ ผสฺโส, เวทยิตํ เวทนา, อิเม ปญฺจ ธมฺมา อายติํ อุปปตฺติภวสฺมิํ อิธ กตสฺส กมฺมสฺส ปจฺจยา.\csromanspacer{} อิติเม จตุสงฺเขเป ตโย อทฺเธ ติสนฺธิํ วีสติยา อากาเรหิ ปฏิจฺจสมุปฺปาทํ ชานาติ ปสฺสติ อญฺญาติ ปฏิวิชฺฌติ, ตํญาตฏฺเฐน ญาณํ, ปชานนฏฺเฐน ปญฺญา, เตน วุจฺจติ ปจฺจยปริคฺคเห ปญฺญา ธมฺมฏฺฐิติญาณํ.}

\csromanheader{2}{ธมฺมฏฺฐิติญาณนิทฺเทโส จตุตฺโถ.}
\csromansectionrule
\csromanmatikaanchor{mtk.7}
\csromantocmarkref{subsection}{5. สมฺมสนญาณนิทฺเทส}{mtk.7}
\bookmark[dest={mtk.7},level=2]{5. สมฺมสนญาณนิทฺเทส}
\csromanheader{2}{5. สมฺมสนญาณนิทฺเทส}
\proseitem{48}{กถํ อตีตานาคตปจฺจุปฺปนฺนานํ ธมฺมานํ สงฺขิปิตฺวา ววตฺถาเน ปญฺญา สมฺมสเน ญาณํ\csromandash{}}

\prose{ยํ กิญฺจิ รูปํ อตีตานาคตปจฺจุปฺปนฺนํ อชฺฌตฺตํ วา พหิทฺธา วา โอฬาริกํ วา สุขุมํ วา หีนํ วา ปณีตํ วา ยํ ทูเร สนฺติเก วา, สพฺพํ รูปํ อนิจฺจโต ววตฺเถติ เอกํ สมฺมสนํ, ทุกฺขโต ววตฺเถติ เอกํ สมฺมสนํ, อนตฺตโต ววตฺเถติ เอกํ สมฺมสนํ.}

\prose{ยา กาจิ เวทนา ฯเปฯ.\csromanspacer{} ยา กาจิ สญฺญา.\csromanspacer{} เย เกจิ สงฺขารา.\csromanspacer{} ยํ กิญฺจิ วิญฺญาณํ อตีตานาคตปจฺจุปฺปนฺนํ อชฺฌตฺตํ วา พหิทฺธา วา โอฬาริกํ วา สุขุมํ วา หีนํ วา ปณีตํ วา ยํ ทูเร สนฺติเก วา, สพฺพํ วิญฺญาณํ อนิจฺจโต ววตฺเถติ เอกํ สมฺมสนํ, ทุกฺขโต ววตฺเถติ เอกํ สมฺมสนํ, อนตฺตโต ววตฺเถติ เอกํ สมฺมสนํ.}

\prose{จกฺขุํ ฯเปฯ.\csromanspacer{} ชรามรณํ อตีตานาคตปจฺจุปฺปนฺนํ อนิจฺจโต ววตฺเถติ เอกํ สมฺมสนํ, ทุกฺขโต ววตฺเถติ เอกํ สมฺมสนํ, อนตฺตโต ววตฺเถติ เอกํ สมฺมสนํ.}

\prose{“รูปํ อตีตานาคตปจฺจุปฺปนฺนํ อนิจฺจํ ขยฏฺเฐน, ทุกฺขํ ภยฏฺเฐน, อนตฺตา อสารกฏฺเฐนา”ติ สงฺขิปิตฺวา ววตฺถาเน ปญฺญา สมฺมสเน ญาณํ. “เวทนา ฯเปฯ.\csromanspacer{} สญฺญา ฯเปฯ.\csromanspacer{} สงฺขารา ฯเปฯ.\csromanspacer{} วิญฺญาณํ ฯเปฯ.\csromanspacer{} จกฺขุ ฯเปฯ.\csromanspacer{} ชรามรณํ อตีตานาคตปจฺจุปฺปนฺนํ อนิจฺจํ ขยฏฺเฐน ทุกฺขํ ภยฏฺเฐน อนตฺตา อสารกฏฺเฐนา”ติ สงฺขิปิตฺวา ววตฺถาเน ปญฺญา สมฺมสเน ญาณํ.}

\prose{\csromanfolio{52}“รูปํ อตีตานาคตปจฺจุปฺปนฺนํ อนิจฺจํ สงฺขตํ ปฏิจฺจสมุปฺปนฺนํ ขยธมฺมํ วยธมฺมํ วิราคธมฺมํ นิโรธธมฺมนฺ”ติ สงฺขิปิตฺวา ววตฺถาเน ปญฺญา สมฺมสเน ญาณํ. “เวทนา ฯเปฯ.\csromanspacer{} สญฺญา.\csromanspacer{} สงฺขารา.\csromanspacer{} วิญฺญาณํ.\csromanspacer{}\csromanspacer{}จกฺขุ ฯเปฯ.\csromanspacer{} ชรามรณํ อตีตานาคตปจฺจุปฺปนฺนํ อนิจฺจํ สงฺขตํ ปฏิจฺจสมุปฺปนฺนํ ขยธมฺมํ วยธมฺมํ วิราคธมฺมํ นิโรธธมฺมนฺ”ติ สงฺขิปิตฺวา ววตฺถาเน ปญฺญา สมฺมสเน ญาณํ.}

\prose{“ชาติปจฺจยา ชรามรณํ, อสติ ชาติยา นตฺถิ ชรามรณนฺ”ติ สงฺขิปิตฺวา ววตฺถาเน ปญฺญา สมฺมสเน ญาณํ.\csromanspacer{} อตีตมฺปิ อทฺธานํ.\csromanspacer{} อนาคตมฺปิ อทฺธานํ “ชาติปจฺจยา ชรามรณํ, อสติ ชาติยา นตฺถิ ชรามรณนฺ”ติ สงฺขิปิตฺวา ววตฺถาเน ปญฺญา สมฺมสเน ญณํ.\csromanspacer{}\csromanspacer{}ภวปฺปจฺจยา ชาติ, อสติ ฯเปฯ.\csromanspacer{} อุปาทานปจฺจยา ภโว, อสติ ฯเปฯ.\csromanspacer{} ตณฺหาปจฺจยา อุปาทานํ, อสติ ฯเปฯ.\csromanspacer{} เวทนาปจฺจยา ตณฺหา, อสติ ฯเปฯ.\csromanspacer{} ผสฺสปจฺจยา เวทนา, อสติ ฯเปฯ.\csromanspacer{} สฬายตนปจฺจยา ผสฺโส, อสติ ฯเปฯ.\csromanspacer{} นามรูปปจฺจยา สฬายตนํ, อสติ ฯเปฯ.\csromanspacer{} วิญฺญาณปจฺจยา นามรูปํ, อสติ ฯเปฯ.\csromanspacer{} สงฺขารปจฺจยา วิญฺญาณํ, อสติ ฯเปฯ. “อวิชฺชาปจฺจยา สงฺขารา, อสติ อวิชฺชาย นตฺถิ สงฺขารา”ติ สงฺขิปิตฺวา ววตฺถาเน ปญฺญา สมฺมสเน ญาณํ. “อตีตมฺปิ อทฺธานํ.\csromanspacer{} อนาคตมฺปิ อทฺธานํ อวิชฺชาปจฺจยา สงฺขารา, อสติ อวิชฺชาย นตฺถิ สงฺขารา”ติ สงฺขิปิตฺวา ววตฺถาเน ปญฺญา สมฺมสเน ญาณํ, ตํญาตฏฺเฐน ญาณํ, ปชานนฏฺเฐน ปญฺญา, เตน วุจฺจติ อตีตานาคตปจฺจุปฺปนฺนานํ ธมฺมานํ สงฺขิปิตฺวา ววตฺถาเน ปญฺญา สมฺมสเน ญาณํ.}

\csromanheader{2}{สมฺมสนญาณนิทฺเทโส ปญฺจโม.}
\csromansectionrule
\csromanmatikaanchor{mtk.8}
\csromantocmarkref{subsection}{6. อุทยพฺพยญาณนิทฺเทส}{mtk.8}
\bookmark[dest={mtk.8},level=2]{6. อุทยพฺพยญาณนิทฺเทส}
\csromanheader{2}{6. อุทยพฺพยญาณนิทฺเทส}
\proseitem{49}{กถํ ปจฺจุปฺปนฺนานํ ธมฺมานํ วิปริณามานุปสฺสเน ปญฺญา อุทยพฺพยานุปสฺสเน ญาณํ\csromandash{}ชาตํ รูปํ ปจฺจุปฺปนฺนํ, ตสฺส นิพฺพตฺติลกฺขณํ อุทโย, วิปริณามลกฺขณํ วโย, อนุปสฺสนา ญาณํ.\csromanspacer{} ชาตา เวทนา.\csromanspacer{} ชาตา สญฺญา.\csromanspacer{} ชาตา สงฺขารา.\csromanspacer{} ชาตํ วิญฺญาณํ.\csromanspacer{} ชาตํ จกฺขุ ฯเปฯ.\csromanspacer{} ชาโต ภโว ปจฺจุปฺปนฺโน, ตสฺส นิพฺพตฺติลกฺขณํ อุทโย, วิปริณามลกฺขณํ วโย, อนุปสฺสนา ญาณํ.}

\proseitem{50}{ปญฺจนฺนํ ขนฺธานํ อุทยํ ปสฺสนฺโต กติ ลกฺขณานิ ปสฺสติ, วยํ ปสฺสนฺโต กติ ลกฺขณานิ ปสฺสติ, อุทยพฺพยํ ปสฺสนฺโต กติ \csromanfolio{53}ลกฺขณานิ ปสฺสติ? ปญฺจนฺนํ ขนฺธานํ อุทยํ ปสฺสนฺโต ปญฺจวีสติ ลกฺขณานิ ปสฺสติ, วยํ ปสฺสนฺโต ปญฺจวีสติ ลกฺขณานิ ปสฺสติ, อุทยพฺพยํ ปสฺสนฺโต ปญฺญาส ลกฺขณานิ ปสฺสติ.}

\prose{รูปกฺขนฺธสฺส อุทยํ ปสฺสนฺโต กติ ลกฺขณานิ ปสฺสติ, วยํ ปสฺสนฺโต กติ ลกฺขณานิ ปสฺสติ, อุทยพฺพยํ ปสฺสนฺโต กติ ลกฺขณานิ ปสฺสติ? เวทนากฺขนฺธสฺส ฯเปฯ.\csromanspacer{} สญฺญากฺขนฺธสฺส ฯเปฯ.\csromanspacer{} สงฺขารกฺขนฺธสฺส ฯเปฯ.\csromanspacer{} วิญฺญาณกฺขนฺธสฺส อุทยํ ปสฺสนฺโต กติ ลกฺขณานิ ปสฺสติ, วยํ ปสฺสนฺโต กติ ลกฺขณานิ ปสฺสติ, อุทยพฺพยํ ปสฺสนฺโต กติ ลกฺขณานิ ปสฺสติ? รูปกฺขนฺธสฺส อุทยํ ปสฺสนฺโต ปญฺจ ลกฺขณานิ ปสฺสติ, วยํ ปสฺสนฺโต ปญฺจ ลกฺขณานิ ปสฺสติ, อุทยพฺพยํ ปสฺสนฺโต ทส ลกฺขณานิ ปสฺสติ.\csromanspacer{}\csromanspacer{}เวทนากฺขนฺธสฺส ฯเปฯ.\csromanspacer{} สญฺญากฺขนฺธสฺส.\csromanspacer{} สงฺขารกฺขนฺธสฺส.\csromanspacer{} วิญฺญาณกฺขนฺธสฺส อุทยํ ปสฺสนฺโต ปญฺจ ลกฺขณานิ ปสฺสติ, วยํ ปสฺสนฺโต ปญฺจ ลกฺขณานิ ปสฺสติ, อุทยพฺพยํ ปสฺสนฺโต ทส ลกฺขณานิ ปสฺสติ.}

\prose{รูปกฺขนฺธสฺส อุทยํ ปสฺสนฺโต กตมานิ ปญฺจ ลกฺขณานิ ปสฺสติ? “อวิชฺชาสมุทยา รูปสมุทโย”ติ ปจฺจยสมุทยฏฺเฐน รูปกฺขนฺธสฺส อุทยํ ปสฺสติ, “ตณฺหาสมุทยา รูปสมุทโย”ติ ปจฺจยสมุทยฏฺเฐน รูปกฺขนฺธสฺส อุทยํ ปสฺสติ, “กมฺมสมุทยา รูปสมุทโย”ติ ปจฺจยสมุทยฏฺเฐน รูปกฺขนฺธสฺส อุทยํ ปสฺสติ, “อาหารสมุทยา รูปสมุทโย”ติ ปจฺจยสมุทยฏฺเฐน รูปกฺขนฺธสฺส อุทยํ ปสฺสติ.\csromanspacer{} นิพฺพตฺติลกฺขณํ ปสฺสนฺโตปิ รูปกฺขนฺธสฺส อุทยํ ปสฺสติ.\csromanspacer{} รูปกฺขนฺธสฺส อุทยํ ปสฺสนฺโต อิมานิ ปญฺจ ลกฺขณานิ ปสฺสติ.}

\prose{วยํ ปสฺสนฺโต กตมานิ ปญฺจ ลกฺขณานิ ปสฺสติ? “อวิชฺชานิโรธา รูปนิโรโธ”ติ ปจฺจยนิโรธฏฺเฐน รูปกฺขนฺธสฺส วยํ ปสฺสติ, “ตณฺหา นิโรธา รูปนิโรโธ”ติ ปจฺจยนิโรธฏฺเฐน รูปกฺขนฺธสฺส วยํ ปสฺสติ, “กมฺมนิโรธา รูปนิโรโธ”ติ ปจฺจยนิโรธฏฺเฐน รูปกฺขนฺธสฺส วยํ ปสฺสติ, “อาหารนิโรธา รูปนิโรโธ”ติ ปจฺจยนิโรธฏฺเฐน รูปกฺขนฺธสฺส วยํ ปสฺสติ, วิปริณามลกฺขณํ ปสฺสนฺโตปิ รูปกฺขนฺธสฺส วยํ ปสฺสติ.\csromanspacer{} รูปกฺขนฺธสฺส วยํ ปสฺสนฺโต อิมานิ ปญฺจ ลกฺขณานิ ปสฺสติ.\csromanspacer{} อุทยพฺพยํ ปสฺสนฺโต อิมานิ ทส ลกฺขณานิ ปสฺสติ.}

\prose{\csromanfolio{54}เวทนากฺขนฺธสฺส อุทยํ ปสฺสนฺโต กตมานิ ปญฺจ ลกฺขณานิ ปสฺสติ? “อวิชฺชาสมุทยา เวทนาสมุทโย”ติ ปจฺจยสมุทยฏฺเฐน เวทนากฺขนฺธสฺส อุทยํ ปสฺสติ, “ตณฺหาสมุทยา เวทนาสมุทโย”ติ ปจฺจยสมุทยฏฺเฐน เวทนากฺขนฺธสฺส อุทยํ ปสฺสติ, “กมฺมสมุทยา เวทนาสมุทโย”ติ ปจฺจยสมุทยฏฺเฐน เวทนากฺขนฺธสฺส อุทยํ ปสฺสติ, “ผสฺสสมุทยา เวทนาสมุทโย”ติ ปจฺจยสมุทยฏฺเฐน เวทนากฺขนฺธสฺส อุทยํ ปสฺสติ.\csromanspacer{} นิพฺพตฺติลกฺขณํ ปสฺสนฺโตปิ เวทนากฺขนฺธสฺส อุทยํ ปสฺสติ.\csromanspacer{} เวทนากฺขนฺธสฺส อุทยํ ปสฺสนฺโต อิมานิ ปญฺจ ลกฺขณานิ ปสฺสติ.}

\prose{วยํ ปสฺสนฺโต กตมานิ ปญฺจ ลกฺขณานิ ปสฺสติ? “อวิชฺชานิโรธา เวทนานิโรโธ”ติ ปจฺจยนิโรธฏฺเฐน เวทนากฺขนฺธสฺส วยํ ปสฺสติ, “ตณฺหานิโรธา เวทนานิโรโธ”ติ ปจฺจยนิโรธฏฺเฐน เวทนากฺขนฺธสฺส วยํ ปสฺสติ, “กมฺมนิโรธา เวทนานิโรโธ”ติ ปจฺจยนิโรธฏฺเฐน เวทนากฺขนฺธสฺส วยํ ปสฺสติ, “ผสฺสนิโรธา เวทนานิโรโธ”ติ ปจฺจยนิโรธฏฺเฐน เวทนากฺขนฺธสฺส วยํ ปสฺสติ.\csromanspacer{} วิปริณามลกฺขณํ ปสฺสนฺโตปิ เวทนากฺขนฺธสฺส วยํ ปสฺสติ.\csromanspacer{} เวทนากฺขนฺธสฺส วยํ ปสฺสนฺโต อิมานิ ปญฺจ ลกฺขณานิ ปสฺสติ.\csromanspacer{}\csromanspacer{}อุทยพฺพยํ ปสฺสนฺโต อิมานิ ทส ลกฺขณานิ ปสฺสติ.}

\prose{สญฺญากฺขนฺธสฺส ฯเปฯ.\csromanspacer{} สงฺขารกฺขนฺธสฺส ฯเปฯ.\csromanspacer{} วิญฺญาณกฺขนฺธสฺส อุทยํ ปสฺสนฺโต กตมานิ ปญฺจ ลกฺขณานิ ปสฺสติ? “อวิชฺชาสมุทยา วิญฺญาณสมุทโย”ติ ปจฺจยสมุทยฏฺเฐน วิญฺญาณกฺขนฺธสฺส อุทยํ ปสฺสติ, “ตณฺหาสมุทยา วิญฺญาณสมุทโย”ติ ปจฺจยสมุทยฏฺเฐน วิญฺญาณกฺขนฺธสฺส อุทยํ ปสฺสติ, “กมฺมสมุทยา วิญฺญาณสมุทโย”ติ ปจฺจยสมุทยฏฺเฐน วิญฺญาณกฺขนฺธสฺส อุทยํ ปสฺสติ, “นามรูปสมุทยา วิญฺญาณสมุทโย”ติ ปจฺจยสมุทยฏฺเฐน วิญฺญาณกฺขนฺธสฺส อุทยํ ปสฺสติ.\csromanspacer{} นิพฺพตฺติลกฺขณํ ปสฺสนฺโตปิ วิญฺญาณกฺขนฺธสฺส อุทยํ ปสฺสติ.\csromanspacer{} วิญฺญาณกฺขนฺธสฺส อุทยํ ปสฺสนฺโต อิมานิ ปญฺจ ลกฺขณานิ ปสฺสติ.}

\prose{วยํ ปสฺสนฺโต กตมานิ ปญฺจ ลกฺขณานิ ปสฺสติ? “อวิชฺชานิโรธา วิญฺญาณนิโรโธ”ติ ปจฺจยนิโรธฏฺเฐน วิญฺญาณกฺขนฺธสฺส วยํ ปสฺสติ, “ตณฺหานิโรธา วิญฺญาณนิโรโธ”ติ ปจฺจยนิโรธฏฺเฐน วิญฺญาณกฺขนฺธสฺส วยํ ปสฺสติ, “กมฺมนิโรธา วิญฺญาณนิโรโธ”ติ ปจฺจยนิโรธฏฺเฐน วิญฺญาณกฺขนฺธสฺส วยํ ปสฺสติ, “นามรูปนิโรธา วิญฺญาณนิโรโธ”ติ ปจฺจยนิโรธฏฺเฐน วิญฺญาณกฺขนฺธสฺส วยํ ปสฺสติ, วิปริณามลกฺขณํ ปสฺสนฺโตปิ \csromanfolio{55}วิญฺญาณกฺขนฺธสฺส วยํ ปสฺสติ.\csromanspacer{} วิญฺญาณกฺขนฺธสฺส วยํ ปสฺสนฺโต อิมานิ ปญฺจ ลกฺขณานิ ปสฺสติ.\csromanspacer{}\csromanspacer{}อุทยพฺพยํ ปสฺสนฺโต อิมานิ ทส ลกฺขณานิ ปสฺสติ.}

\prose{ปญฺจนฺนํ ขนฺธานํ อุทยํ ปสฺสนฺโต อิมานิ ปญฺจวีสติ ลกฺขณานิ ปสฺสติ, วยํ ปสฺสนฺโต อิมานิ ปญฺจวีสติ ลกฺขณานิ ปสฺสติ, อุทยพฺพยํ ปสฺสนฺโต อิมานิ ปญฺญาส ลกฺขณานิ ปสฺสติ.\csromanspacer{} ตํญาตฏฺเฐน ญาณํ, ปชานนฏฺเฐน ปญฺญา, เตน วุจฺจติ ปจฺจุปฺปนฺนานํ ธมฺมานํ วิปริณามานุปสฺสเน ปญฺญา อุทยพฺพยานุปสฺสเน ญาณํ.\csromanspacer{}\csromanspacer{}รูปกฺขนฺโธ\footnote{รูปกฺขนฺธา (สฺยา)} อาหารสมุทโย, เวทนาสญฺญา สงฺขารา ตโย\footnote{สงฺขาราติ เสสา (สฺยา)} ขนฺธา ผสฺสสมุทยา.\csromanspacer{} วิญฺญาณกฺขนฺโธ นามรูปสมุทโย.}

\csromanheader{2}{อุทยพฺพยญาณนิทฺเทสฺโส ฉฏฺโฐ.}
\csromansectionrule
\csromanmatikaanchor{mtk.9}
\csromantocmarkref{subsection}{7. ภงฺคานุปสฺสนาญาณนิทฺเทส}{mtk.9}
\bookmark[dest={mtk.9},level=2]{7. ภงฺคานุปสฺสนาญาณนิทฺเทส}
\csromanheader{2}{7. ภงฺคานุปสฺสนาญาณนิทฺเทส}
\proseitem{51}{กถํ อารมฺมณํ ปฏิสงฺขา ภงฺคานุปสฺสเน ปญฺญา วิปสฺสเน ญาณํ\csromandash{}รูปารมฺมณตา จิตฺตํ อุปฺปชฺชิตฺวา ภิชฺชติ, ตํ อารมฺมณํ ปฏิสงฺขา ตสฺส จิตฺตสฺส ภงฺค อนุปสฺสติ.\csromanspacer{}\csromanspacer{}\textbf{อนุปสฺสตี}ติ กถํ อนุปสฺสติ? อนิจฺจโต อนุปสฺสติ, โน นิจฺจโต.\csromanspacer{} ทุกฺขโต อนุปสฺสติ, โน สุขโต.\csromanspacer{} อนตฺตโต อนุปสฺสติ, โน อตฺตโต.\csromanspacer{} นิพฺพินฺทติ, โน นนฺทติ.\csromanspacer{} วิรชฺชติ, โน รชฺชติ.\csromanspacer{} นิโรเธติ, โน สมุเทติ.\csromanspacer{} ปฏินิสฺสชฺชติ, โน อาทิยติ.}

\proseitem{52}{อนิจฺจโต อนุปสฺสนฺโต นิจฺจสญฺญํ ปชหาติ, ทุกฺขโต อนุปสฺสนฺโต สุขสญฺญํ ปชหาติ, อนตฺตโต อนุปสฺสนฺโต อตฺตสญฺญํ ปชหาติ, นิพฺพินฺทนฺโต นนฺทิํ ปชหาติ, วิรชฺชนฺโต ราคํ ปชหาติ, นิโรเธนฺโต สมุทยํ ปชหาติ, ปฏินิสฺสชฺชนฺโต อาทานํ ปชหาติ.}

\prose{เวทนารมฺมณตา ฯเปฯ.\csromanspacer{} สญฺญารมฺมณตา.\csromanspacer{} สงฺขารารมฺมณตา.\csromanspacer{} วิญฺญาณารมฺมณตา.\csromanspacer{}\csromanspacer{}จกฺขุ ฯเปฯ.\csromanspacer{} ชรามรณารมฺมณตา จิตฺตํ อุปฺปชฺชิตฺวา ภิชฺชติ, ตํ อารมฺมณํ ปฏิสงฺขา ตสฺส จิตฺตสฺส ภงฺคํ อนุปสฺสติ.\csromanspacer{}\csromanspacer{}\textbf{อนุปสฺสตี}ติ กถํ อนุปสฺสติ? อนิจฺจโต อนุปสฺสติ, โน นิจฺจโต.\csromanspacer{} ทุกฺขโต อนุปสฺสติ, \csromanfolio{56}โน สุขโต.\csromanspacer{} อนตฺตโต อนุปสฺสติ, โน อตฺตโต.\csromanspacer{} นิพฺพินฺทติ, โน นนฺทติ.\csromanspacer{} วิรชฺชติ, โน รชฺชติ.\csromanspacer{} นิโรเธติ, โน สมุเทติ.\csromanspacer{} ปฏินิสฺสชฺชติ, โน อาทิยติ.}

\prose{อนิจฺจโต อนุปสฺสนฺโต นิจฺจสญฺญํ ปชหาติ, ทุกฺขโต อนุปสฺสนฺโต สุขสญฺญํ ปชหาติ, อนตฺตโต อนุปสฺสนฺโต อตฺตสญฺญํ ปชหาติ, นิพฺพินฺทนฺโต นนฺทิํ ปชหาติ, วิรชฺชนฺโต ราคํ ปชหาติ, นิโรเธนฺโต สมุทยํ ปชหาติ, ปฏินิสฺสชฺชนฺโต อาทานํ ปชหาติ.}

\csromangathameasure{วตฺถุสงฺกมนา เจว, ปญฺญาย จ วิวฏฺฏนา. \\ อาวชฺชนาพลญฺเจว, ปฏิสงฺขา วิปสฺสนา. \\ อารมฺมณ-อนฺวเยน, อุโภ เอกววตฺถนา. \\ นิโรเธ อธิมุตฺตตา, วยลกฺขณวิปสฺสนา. \\ อารมฺมณญฺจ ปฏิสงฺขา, ภงฺคญฺจ อนุปสฺสติ. \\ สุญฺญโต จ อุปฏฺฐานํ, อธิปญฺญา วิปสฺสนา. \\ กุสโล ตีสุ อนุปสฺสนาสุ, จตสฺโส จ วิปสฺสนาสุ. \\ ตโย อุปฏฺฐาเน กุสลตา, นานาทิฏฺฐีสุ น กมฺปตีติ.}
\csromangathagroup{\csromangathastanza{วตฺถุสงฺกมนา เจว, ปญฺญาย จ วิวฏฺฏนา. \\ อาวชฺชนาพลญฺเจว, ปฏิสงฺขา วิปสฺสนา.}\csromangathastanza{อารมฺมณ-อนฺวเยน, อุโภ เอกววตฺถนา\footnote{เอกววตฺถานา (ก)}. \\ นิโรเธ อธิมุตฺตตา, วยลกฺขณวิปสฺสนา.}\csromangathastanza{อารมฺมณญฺจ ปฏิสงฺขา, ภงฺคญฺจ อนุปสฺสติ. \\ สุญฺญโต จ อุปฏฺฐานํ, อธิปญฺญา วิปสฺสนา.}\csromangathastanza{กุสโล ตีสุ อนุปสฺสนาสุ, จตสฺโส จ\footnote{จตูสุ จ (สฺยา)} วิปสฺสนาสุ. \\ ตโย อุปฏฺฐาเน กุสลตา, นานาทิฏฺฐีสุ น กมฺปตีติ.}}

\prose{ตํญาตฏฺเฐน ญาณํ, ปชานนฏฺเฐน ปญฺญา, เตน วุจฺจติ อารมฺมณํ ปฏิสงฺขา ภงฺคานุปสฺสเน ปญฺญา วิปสฺสเน ญาณํ.}

\csromanheader{2}{ภงฺคานุปสฺสนาญาณนิทฺเทโส สตฺตโม.}
\csromansectionrule
\csromanmatikaanchor{mtk.10}
\csromantocmarkref{subsection}{8. อาทีนวญาณนิทฺเทส}{mtk.10}
\bookmark[dest={mtk.10},level=2]{8. อาทีนวญาณนิทฺเทส}
\csromanheader{2}{8. อาทีนวญาณนิทฺเทส}
\proseitem{53}{กถํ ภยตุปฏฺฐาเน ปญฺญา อาทีนเว ญาณํ\csromandash{}อุปฺปาโท ภยนฺติ ภยตุปฏฺฐาเน ปญฺญา อาทีนเว ญาณํ, ปวตฺตํ ภยนฺติ ภยตุปฏฺฐาเน ปญฺญา อาทีนเว ญาณํ, นิมิตฺตํ ภยนฺติ ฯเปฯ, อายูหนา ภยนฺติ ฯเปฯ, ปฏิสนฺธิ ภยนฺติ.\csromanspacer{} คติ ภยนฺติ.\csromanspacer{} นิพฺพตฺติ ภยนฺติ.\csromanspacer{} อุปปตฺติ ภยนฺติ.\csromanspacer{} ชาติ ภยนฺติ.\csromanspacer{} ชรา ภยนฺติ.\csromanspacer{} พฺยาธิ ภยนฺติ.\csromanspacer{} มรณํ ภยนฺติ.\csromanspacer{} โสโก ภยนฺติ.\csromanspacer{} ปริเทโว ภยนฺติ.\csromanspacer{} อุปายาโส ภยนฺติ ภยตุปฏฺฐาเน ปญฺญา อาทีนเว ญาณํ.}

\prose{\csromanfolio{57}อนุปฺปาโท เขมนฺติ สนฺติปเท ญาณํ, อปฺปวตฺตํ เขมนฺติ สนฺติปเท ญาณํ ฯเปฯ อนุปายาโส เขมนฺติ สนฺติปเท ญาณํ.}

\prose{อุปฺปาโท ภยํ, อนุปฺปาโท เขมนฺติ สนฺติปเท ญาณํ, ปวตฺตํ ภยํ, อปฺปวตฺตํ เขมนฺติ สนฺติปเท ญาณํ ฯเปฯ อุปายาโส ภยํ, อนุปายาโส เขมนฺติ สนฺติปเท ญาณํ.}

\prose{อุปฺปาโท ทุกฺขนฺติ ภยตุปฏฺฐาเน ปญฺญา อาทีนเว ญาณํ, ปวตฺตํ ทุกฺขนฺติ ภยตุปฏฺฐาเน ปญฺญา อาทีนเว ญาณํ ฯเปฯ อุปายาโส ทุกฺขนฺติ ภยตุปฏฺฐาเน ปญฺญา อาทีนเว ญาณํ.}

\prose{อนุปฺปาโท สุขนฺติ สนฺติปเท ญาณํ, อปฺปวตฺตํ สุขนฺติ สนฺติปเท ญาณํ ฯเปฯ อนุปายาโส สุขนฺติ สนฺติปเท ญาณํ.}

\prose{อุปฺปาโท ทุกฺขํ, อนุปฺปาโท สุขนฺติ สนฺติปเท ญาณํ, ปวตฺตํ ทุกฺขํ, อปฺปวตฺตํ สุขนฺติ สนฺติปเท ญาณํ ฯเปฯ อุปายาโส ทุกฺขํ, อนุปายาโส สุขนฺติ สนฺติปเท ญาณํ.}

\prose{อุปฺปาโท สามิสนฺติ ภยตุปฏฺฐาเน ปญฺญา อาทีนเว ญาณํ, ปวตฺตํ สามิสนฺติ ภยตุปฏฺฐาเน ปญฺญา อาทีนเว ญาณํ ฯเปฯ อุปายาโส สามิสนฺติ ภยตุปฏฺฐาเน ปญฺญา อาทีนเว ญาณํ.}

\prose{อนุปฺปาโท นิรามิสนฺติ สนฺติปเท ญาณํ, อปฺปวตฺตํ นิรามิสนฺติ สนฺติปเท ญาณํ ฯเปฯ อนุปายาโส นิรามิสนฺติ สนฺติปเท ญาณํ.}

\prose{อุปฺปาโท สามิสํ, อนุปฺปาโท นิรามิสนฺติ สนฺติปเท ญาณํ, ปวตฺตํ สามิสํ, อปฺปวตฺตํ นิรามิสนฺติ สนฺติปเท ญาณํ ฯเปฯ อุปายาโส สามิสํ, อนุปายาโส นิรามิสนฺติ สนฺติปเท ญาณํ.}

\prose{อุปฺปาโท สงฺขาราติ ภยตุปฏฺฐาเน ปญฺญา อาทีนเว ญาณํ, ปวตฺตํ สงฺขาราติ ภยตุปฏฺฐาเน ปญฺญา อาทีนเว ญาณํ ฯเปฯ อุปายาโส สงฺขาราติ ภยตุปฏฺฐาเน ปญฺญา อาทีนเว ญาณํ.}

\prose{อนุปฺปาโท นิพฺพานนฺติ สนฺติปเท ญาณํ, อปฺปวตฺตํ นิพฺพานนฺติ สนฺติปเท ญาณํ ฯเปฯ อนุปายาโส นิพฺพานนฺติ สนฺติปเท ญาณํ.}

\prose{\csromanfolio{58}อุปฺปาโท สงฺขารา, อนุปฺปาโท นิพฺพานนฺติ สนฺติปเท ญาณํ, ปวตฺตํ สงฺขารา, อปฺปวตฺตํ นิพฺพานนฺติ สนฺติปเท ญาณํ ฯเปฯ อุปายาโส สงฺขารา, อนุปายาโส นิพฺพานนฺติ สนฺติปเท ญาณํ.}

\csromangathameasure{อุปฺปาทญฺจ ปวตฺตญฺจ, นิมิตฺตํ ทุกฺขนฺติ ปสฺสติ. \\ อายูหนํ ปฏิสนฺธิํ, ญาณํ อาทีนเว อิทํ. \\ อนุปฺปาทํ อปฺปวตฺตํ, อนิมิตฺตํ สุขนฺติ จ. \\ อนายูหนํ อปฺปฏิสนฺธิํ, ญาณํ สนฺติปเท อิทํ. \\ อิทํ อาทีนเว ญาณํ, ปญฺจฐาเนสุ ชายติ. \\ ปญฺจฐาเน สนฺติปเท, ทส ญาเณ ปชานาติ. \\ ทฺวินฺนํ ญาณานํ กุสลตา, นานาทิฏฺฐีสุ น กมฺปตีติ.}
\csromangathagroup{\csromangathastanza{อุปฺปาทญฺจ ปวตฺตญฺจ, นิมิตฺตํ ทุกฺขนฺติ ปสฺสติ. \\ อายูหนํ ปฏิสนฺธิํ, ญาณํ อาทีนเว อิทํ.}\csromangathastanza{อนุปฺปาทํ อปฺปวตฺตํ, อนิมิตฺตํ สุขนฺติ จ. \\ อนายูหนํ อปฺปฏิสนฺธิํ, ญาณํ สนฺติปเท อิทํ.}\csromangathastanza{อิทํ อาทีนเว ญาณํ, ปญฺจฐาเนสุ ชายติ. \\ ปญฺจฐาเน สนฺติปเท, ทส ญาเณ ปชานาติ. \\ ทฺวินฺนํ ญาณานํ กุสลตา, นานาทิฏฺฐีสุ น กมฺปตีติ.}}

\prose{ตํญาตฏฺเฐน ญาณํ, ปชานนฏฺเฐน ปญฺญา, เตน วุจฺจติ ภยตุปฏฺฐาเน ปญฺญา อาทีนเว ญาณํ.}

\csromanheader{2}{อาทีนวญาณนิทฺเทโส อฏฺฐโม.}
\csromansectionrule
\csromanmatikaanchor{mtk.11}
\csromantocmarkref{subsection}{9. สงฺขารุเปกฺขาญาณนิทฺเทส}{mtk.11}
\bookmark[dest={mtk.11},level=2]{9. สงฺขารุเปกฺขาญาณนิทฺเทส}
\csromanheader{2}{9. สงฺขารุเปกฺขาญาณนิทฺเทส}
\proseitem{54}{กถํ มุญฺจิตุกมฺยตาปฏิสงฺขาสนฺติฏฺฐนา\footnote{มุจฺจิตุกมฺยตาปฏิสงฺขาสนฺติฏฺฐนา (ก)} ปญฺญา สงฺขารุเปกฺขาสุ ญาณํ\csromandash{}อุปฺปาทํ มุญฺจิตุกมฺยตาปฏิสงฺขาสนฺติฏฺฐนา\footnote{อุปฺปาทมุญฺจิตุกมฺยตา... (สฺยา)} ปญฺญา สงฺขารุเปกฺขาสุ ญาณํ, ปวตฺตํ มุญฺจิตุกมฺยตาปฏิสงฺขาสนฺติฏฺฐนา ปญฺญา สงฺขารุเปกฺขาสุ ญาณํ, นิมิตฺตํ มุญฺจตุกมฺยตา ฯเปฯ อายูหนํ มุญฺจิตุกมฺยตา.\csromanspacer{} ปฏิสนฺธิํ มุญฺจิตุกมฺยตา.\csromanspacer{} คติํ มุญฺจิตุกมฺยตา.\csromanspacer{} นิพฺพตฺตํ มุญฺจิตุกมฺยตา.\csromanspacer{} อุปปตฺติํ มุญฺจิตุกมฺยตา.\csromanspacer{} ชาติํ มุญฺจิตุกมฺยตา.\csromanspacer{} ชรํ มุญฺจิตุกมฺยตา.\csromanspacer{} พฺยาธิํ มุญฺจิตุกมฺยตา.\csromanspacer{} มรณํ มุญฺจิตุกมฺยตา.\csromanspacer{} โสกํ มุญฺจิตุกมฺยตา.\csromanspacer{} ปริเทวํ มุญฺจิตุกมฺยตา ฯเปฯ อุปายาสํ มุญฺจิตุกมฺยตาปฏิสงฺขาสนฺติฏฺฐนา ปญฺญา สงฺขารุเปกฺขาสุ ญาณํ.}

\prose{อุปฺปาโท ทุกฺขนฺติ มุญฺจิตุกมฺยตาปฏิสงฺขาสนฺติฏฺฐนา ปญฺญา สงฺขารุเปกฺขาสุ ญาณํ, ปวตฺตํ ทุกฺขนฺติ มุญฺจิตุกมฺยตาปฏิสงฺขาสนฺติฏฺฐนา ปญฺญา สงฺขารุเปกฺขาสุ ญาณํ ฯเปฯ อุปายาโส ทุกฺขนฺติ มุญฺจิตุกมฺยตาปฏิสงฺขาสนฺติฏฺฐนา ปญฺญา สงฺขารุเปกฺขาสุ ญาณํ.}

\prose{\csromanfolio{59}อุปฺปาโท ภยนฺติ มุญฺจิตุกมฺยตาปฏิสงฺขาสนฺติฏฺฐนา ปญฺญา สงฺขารุเปกฺขาสุ ญาณํ, ปวตฺตํ ภยนฺติ มุญฺจิตุกมฺยตาปฏิสงฺขาสนฺติฏฺฐนา ปญฺญา สงฺขารุเปกฺขาสุ ญาณํ ฯเปฯ อุปายาโส ภยนฺติ มุญฺจิตุกมฺยตาปฏิสงฺขาสนฺติฏฺฐนา ปญฺญา สงฺขารุเปกฺขาสุ ญาณํ.}

\prose{อุปฺปาโท สามิสนฺติ มุญฺจิตุกมฺยตาปฏิสงฺขาสนฺติฏฺฐนา ปญฺญา สงฺขารุเปกฺขาสุ ญาณํ, ปวตฺตํ สามิสนฺติ มุญฺจิตุกมฺยตาปฏิสงฺขาสนฺติฏฺฐนา ปญฺญา สงฺขารุเปกฺขาสุ ญาณํ ฯเปฯ อุปายาโส สามิสนฺติ มุญฺจิตุกมฺยตาปฏิสงฺขาสนฺติฏฺฐนา ปญฺญา สงฺขารุเปกฺขาสุ ญาณํ.}

\prose{อุปฺปาโท สงฺขาราติ มุญฺจิตุกมฺยตาปฏิสงฺขาสนฺติฏฺฐนา ปญฺญา สงฺขารุเปกฺขาสุ ญาณํ, ปวตฺตํ สงฺขาราติ มุญฺจิตุกมฺยตาปฏิสงฺขาสนฺติฏฺฐนา ปญฺญา สงฺขารุเปกฺขาสุ ญาณํ ฯเปฯ อุปายาโส สงฺขาราติ มุญฺจิตุกมฺยตาปฏิสงฺขาสนฺติฏฺฐนา ปญฺญา สงฺขารุเปกฺขาสุ ญาณํ.}

\prose{อุปฺปาโท สงฺขารา, เต สงฺขาเร อชฺฌุเปกฺขตีติ สงฺขารุเปกฺขา.\csromanspacer{} เย จ สงฺขารา ยา จ อุเปกฺขา, อุโภเปเต สงฺขารา, เต สงฺขาเร อชฺฌุเปกฺขตีติ สงฺขารุเปกฺขา.\csromanspacer{} ปวตฺตํ สงฺขารา ฯเปฯ.\csromanspacer{} นิมิตฺตํ สงฺขารา.\csromanspacer{} อายูหนา สงฺขารา.\csromanspacer{} ปฏิสนฺธิ สงฺขารา.\csromanspacer{} คติ สงฺขารา.\csromanspacer{} นิพฺพตฺติ สงฺขารา.\csromanspacer{} อุปปตฺติ สงฺขารา.\csromanspacer{} ชาติ สงฺขารา.\csromanspacer{} ชรา สงฺขารา.\csromanspacer{} พฺยาธิ สงฺขารา.\csromanspacer{} มรณํ สงฺขารา.\csromanspacer{} โสโก สงฺขารา.\csromanspacer{} ปริเทโว สงฺขารา ฯเปฯ.\csromanspacer{} อุปายาโส สงฺขารา, เต สงฺขาเร อชฺฌุเปกฺขตีติ สงฺขารุเปกฺขา.\csromanspacer{} เย จ สงฺขารา ยา จ อุเปกฺขา, อุโภเปเต สงฺขารา, เต สงฺขาเร อชฺฌุเปกฺขตีติ สงฺขารุเปกฺขา.}

\proseitem{55}{กติหากาเรหิ สงฺขารุเปกฺขาย จิตฺตสฺส อภินีหาโร โหติ? อฏฺฐหากาเรหิ สงฺขารุเปกฺขาย จิตฺตสฺส อภินีหาโร โหติ.\csromanspacer{}\csromanspacer{}ปุถุชฺชนสฺส กติหากาเรหิ สงฺขารุเปกฺขาย จิตฺตสฺส อภินีหาโร โหติ? เสกฺขสฺส กติหากาเรหิ สงฺขารุเปกฺขาย จิตฺตสฺส อภินีหาโร โหติ? วีตราคสฺส กติหากาเรหิ สงฺขารุเปกฺขาย จิตฺตสฺส อภินีหาโร โหติ? ปุถุชฺชนสฺส ทฺวีหากาเรหิ สงฺขารุเปกฺขาย จิตฺตสฺส อภินีหาโร โหติ, เสกฺขสฺส ตีหากาเรหิ สงฺขารุเปกฺขาย จิตฺตสฺส อภินีหาโร โหติ, วีตราคสฺส ตีหากาเรหิ สงฺขารุเปกฺขาย จิตฺตสฺส อภินีหาโร โหติ.}

\prose{\csromanfolio{60}ปุถุชฺชนสฺส กตเมหิ ทฺวีหากาเรหิ สงฺขารุเปกฺขาย จิตฺตสฺส อภินีหาโร โหติ? ปุถุชฺชโน สงฺขารุเปกฺขํ อภินนฺทติ วา วิปสฺสติ วา.\csromanspacer{} ปุถุชฺชนสฺส อิเมหิ ทฺวีหากาเรหิ สงฺขารุเปกฺขาย จิตฺตสฺส อภินีหาโร โหติ.\csromanspacer{}\csromanspacer{}เสกฺขสฺส กตเมหิ ตีหากาเรหิ สงฺขารุเปกฺขาย จิตฺตสฺส อภินีหาโร โหติ? เสกฺโข สงฺขารุเปกฺขํ อภินนฺทติ วา วิปสฺสติ วา ปฏิสงฺขาย วา ผลสมาปตฺติํ สมาปชฺชติ เสกฺขสฺส อิเมหิ ตีหากาเรหิ สงฺขารุเปกฺขาย จิตฺตสฺส อภินีหาโร โหติ.\csromanspacer{}\csromanspacer{}วีตราคสฺส กตเมหิ ตีหากาเรหิ สงฺขารุเปกฺขาย จิตฺตสฺส อภินีหาโร โหติ? วีตราโค สงฺขารุเปกฺขํ วิปสฺสติ วา ปฏิสงฺขาย วา ผลสมาปตฺติํ สมาปชฺชติ, ตทชฺฌุเปกฺขิตฺวา สุญฺญตวิหาเรน วา อนิมิตฺตวิหาเรน วา อปฺปณิหิตหาเรน วา วิหรติ.\csromanspacer{} วีตราคสฺส อิเมหิ ตีหากาเรหิ สงฺขารุเปกฺขาย จิตฺตสฺส อภินีหาโร โหติ.}

\proseitem{56}{กถํ ปุถุชฺชนสฺส จ เสกฺขสฺส จ สงฺขารุเปกฺขาย จิตฺตสฺส อภินีหาโร เอกตฺตํ โหติ? ปุถุชฺชนสฺส สงฺขารุเปกฺขํ อภินนฺทโต จิตฺตํ กิลิสฺสติ, ภาวนาย ปริปนฺโถ โหติ, ปฏิเวธสฺส อนฺตราโย โหติ, อายติํ ปฏิสนฺธิยา ปจฺจโย โหติ.\csromanspacer{} เสกฺขสฺสปิ สงฺขารุเปกฺขํ อภินนฺทโต จิตฺตํ กิลิสฺสติ, ภาวนาย ปริปนฺโถ โหติ, อุตฺตริปฏิเวธสฺส อนฺตราโย โหติ, อายติํ ปฏิสนฺธิยา ปจฺจโย โหติ.\csromanspacer{} เอวํ ปุถุชฺชนสฺส จ เสกฺขสฺส จ สงฺขารุเปกฺขาย จิตฺตสฺส อภินีหาโร เอกตฺตํ โหติ อภินนฺทฏฺเฐน.}

\prose{กถํ ปุถุชฺชนสฺส จ เสกฺขสฺส จ วีตราคสฺส จ สงฺขารุเปกฺขาย จิตฺตสฺส อภินีหาโร เอกตฺตํ โหติ? ปุถุชฺชโน สงฺขารุเปกฺขํ อนิจฺจโตปิ ทุกฺขโตปิ อนตฺตโตปิ วิปสฺสติ, เสกฺโขปิ สงฺขารุเปกฺขํ อนิจฺจโตปิ ทุกฺขโตปิ อนตฺตโตปิ วิปสฺสติ, วีตราโคปิ สงฺขารุเปกฺขํ อนิจฺจโตปิ ทุกฺขโตปิ อนตฺตโตปิ วิปสฺสติ.\csromanspacer{} เอวํ ปุถุชฺชนสฺส จ เสกฺขสฺส จ วีตราคสฺส จ สงฺขารุเปกฺขาย จิตฺตสฺส อภินีหาโร เอกตฺตํ โหติ อนุปสฺสนฏฺเฐน.}

\prose{กถํ ปุถุชฺชนสฺส จ เสกฺขสฺส จ วีตราคสฺส จ สงฺขารุเปกฺขาย จิตฺตสฺส อภินีหาโร นานตฺตํ โหติ? ปุถุชฺชนสฺส สงฺขารุเปกฺขา กุสลา โหติ, เสกฺขสฺสปิ สงฺขารุเปกฺขา กุสลา โหติ, วีตราคสฺส \csromanfolio{61}สงฺขารุเปกฺขา อพฺยากตา โหติ.\csromanspacer{} เอวํ ปุถุชฺชนสฺส จ เสกฺขสฺส จ วีตราคสฺส จ สงฺขารุเปกฺขาย จิตฺตสฺส อภินีหาโร นานตฺตํ โหติ กุสลาพฺยากตฏฺเฐน.}

\prose{กถํ ปุถุชฺชนสฺส จ เสกฺขสฺส จ วีตราคสฺส จ สงฺขารุเปกฺขาย จิตฺตสฺส อภินีหาโร นานตฺตํ โหติ? ปุถุชฺชนสฺส สงฺขารุเปกฺขา กิญฺจิ กาเล สุวิทิตา โหติ, กิญฺจิ กาเล น สุวิทิตา โหติ.\csromanspacer{} เสกฺขสฺสปิ สงฺขารุเปกฺขา กิญฺจิ กาเล สุวิทิตา โหติ, กิญฺจิ กาเล น สุวิทิตา โหติ.\csromanspacer{} วีตราคสฺส สงฺขารุเปกฺขา อจฺจนฺตํ สุวิทิตา โหติ.\csromanspacer{} เอวํ ปุถุชฺชนสฺส จ เสกฺขสฺส จ วีตราคสฺส จ สงฺขารุเปกฺขาย จิตฺตสฺส อภินีหาโร นานตฺตํ โหติ วิทิตฏฺเฐน จ อวิทิตฏฺเฐน จ.}

\prose{กถํ ปุถุชฺชนสฺส จ เสกฺขสฺส จ วีตราคสฺส จ สงฺขารุเปกฺขาย จิตฺตสฺส อภินีหาโร นานตฺตํ โหติ? ปุถุชฺชโน สงฺขารุเปกฺขํ อติตฺตตฺตา วิปสฺสติ, เสกฺโขปิ สงฺขารุเปกฺขํ อติตฺตตฺตา วิปสฺสติ, วีตราโค สงฺขารุเปกฺขํ ติตฺตตฺตา วิปสฺสติ.\csromanspacer{} เอวํ ปุถุชฺชนสฺส จ เสกฺขสฺส จ วีตราคสฺส จ สงฺขารุเปกฺขาย จิตฺตสฺส อภินีหาโร นานตฺตํ โหติ ติตฺตฏฺเฐน จ อติตฺตฏฺเฐน จ.}

\prose{กถํ ปุถุชฺชนสฺส จ เสกฺขสฺส จ วีตราคสฺส จ สงฺขารุเปกฺขาย จิตฺตสฺส อภินีหาโร นานตฺตํ โหติ? ปุถุชฺชโน สงฺขารุเปกฺขํ ติณฺณํ สํโยชนานํ ปหานาย โสตาปตฺติมคฺคํ ปฏิลาภตฺถาย วิปสฺสติ, เสกฺโข สงฺขารุเปกฺขํ ติณฺณํ สํโยชนานํ ปหีนตฺตา อุตฺตริปฏิลาภตฺถาย วิปสฺสติ, วีตราโค สงฺขารุเปกฺขํ สพฺพกิเลสานํ ปหีนตฺตา ทิฏฺฐธมฺมสุขวิหารตฺถาย วิปสฺสติ.\csromanspacer{} เอวํ ปุถุชฺชนสฺส จ เสกฺขสฺส จ วีตราคสฺส จ สงฺขารุเปกฺขาย จิตฺตสฺส อภินีหาโร นานตฺตํ โหติ ปหีนฏฺเฐน จ อปฺปหีนฏฺเฐน จ.}

\prose{กถํ เสกฺขสฺส จ วีตราคสฺส จ สงฺขารุเปกฺขาย จิตฺตสฺส อภินีหาโร นานตฺตํ โหติ? เสกฺโข สงฺขารุเปกฺขํ อภินนฺทติ วา วิปสฺสติ วา ปฏิสงฺขาย วา ผลสมาปตฺติํ สมาปชฺชติ, วีตราโค สงฺขารุเปกฺขํ วิปสฺสติ วา ปฏิสงฺขาย วา ผลสมาปตฺติํ สมาปชฺชติ, ตทชฺฌุเปกฺขิตฺวา สุญฺญตวิหาเรน วา อนิมิตฺตวิหาเรน วา อปฺปณิหิตวิหาเรน วา วิหรติ.\csromanspacer{} เอวํ เสกฺขสฺส จ วีตราคสฺส จ สงฺขารุเปกฺขาย จิตฺตสฺส อภินีหาโร นานตฺตํ โหติ วิหารสมาปตฺตฏฺเฐน.}

\proseitem{57}{\csromanfolio{62}กติ สงฺขารุเปกฺขา สมถวเสน อุปฺปชฺชนฺติ, กติ สงฺขารุเปกฺขา วิปสฺสนาวเสน อุปฺปชฺชนฺติ? อฏฺฐ สงฺขารุเปกฺขา สมถวเสน อุปฺปชฺชนฺติ, ทส สงฺขารุเปกฺขา วิปสฺสนาวเสน อุปฺปชฺชนฺติ.}

\prose{กตมา อฏฺฐ สงฺขารุเปกฺขา สมถวเสน อุปฺปชฺชนฺติ? ปฐมํ ฌานํ ปฏิลาภตฺถาย นีวรเณ ปฏิสงฺขา สนฺติฏฺฐนา ปญฺญา สงฺขารุเปกฺขาสุ ญาณํ.\csromanspacer{} ทุติยํ ฌานํ ปฏิลาภตฺถาย วิตกฺกวิจาเร ปฏิสงฺขา สนฺติฏฺฐนา ปญฺญา สงฺขารุเปกฺขาสุ ญาณํ.\csromanspacer{} ตติยํ ฌานํ ปฏิลาภตฺถาย ปีติํ ปฏิสงฺขา สนฺติฏฺฐนา ปญฺญา สงฺขารุเปกฺขาสุ ญาณํ.\csromanspacer{} จตุตฺถํ ฌานํ ปฏิลาภตฺถาย สุขทุกฺเข ปฏิสงฺขา สนฺติฏฺฐนา ปญฺญา สงฺขารุเปกฺขาสุ ญาณํ.\csromanspacer{}\csromanspacer{}อากาสานญฺจายตนสมาปตฺติํ ปฏิลาภตฺถาย รูปสญฺญํ ปฏิฆสญฺญํ นานตฺตสญฺญํ ปฏิสงฺขา สนฺติฏฺฐนา ปญฺญา สงฺขารุเปกฺขาสุ ญาณํ.\csromanspacer{} วิญฺญาณญฺจายตนสมาปตฺติํ ปฏิลาภตฺถาย อากาสานญฺจายตนสญฺญํ ปฏิสงฺขา สนฺติฏฺฐนา ปญฺญา สงฺขารุเปกฺขาสุ ญาณํ.\csromanspacer{} อากิญฺจญฺญายตนสมาปตฺติํ ปฏิลาภตฺถาย วิญฺญาณญฺจายตนสญฺญํ ปฏิสงฺขา สนฺติฏฺฐนา ปญฺญา สงฺขารุเปกฺขาสุ ญาณํ.\csromanspacer{} เนวสญฺญานาสญฺญายตนสมาปตฺติํ ปฏิลาภตฺถาย อากิญฺจญฺญายตนสญฺญํ ปฏิสงฺขา สนฺติฏฺฐนา ปญฺญา สงฺขารุเปกฺขาสุ ญาณํ.\csromanspacer{} อิมา อฏฺฐ สงฺขารุเปกฺขา สมถวเสน อุปฺปชฺชนฺติ.}

\prose{กตมา ทส สงฺขารุเปกฺขา วิปสฺสนาวเสน อุปฺปชฺชนฺติ? โสตาปตฺติมคฺคํ ปฏิลาภตฺถาย อุปฺปาทํ ปวตฺตํ นิมิตฺตํ อายูหนํ ปฏิสนฺธิํ คติํ นิพฺพตฺติํ อุปปตฺติํ ชาติํ ชรํ พฺยาธิํ มรณํ โสกํ ปริเทวํ อุปายาสํ ปฏิสงฺขา สนฺติฏฺฐนา ปญฺญา สงฺขารุเปกฺขาสุ ญาณํ.\csromanspacer{} โสตาปตฺติผลสมาปตฺตตฺถาย อุปฺปาทํ ปวตฺตํ นิมิตฺตํ อายูหนํ ปฏิสนฺธิํ ฯเปฯ ปฏิสงฺขา สนฺติฏฺฐนา ปญฺญา สงฺขารุเปกฺขาสุ ญาณํ.\csromanspacer{} สกทาคามิมคฺคํ ปฏิลาภตฺถาย ฯเปฯ.\csromanspacer{} สกทาคามิผลสมาปตฺตตฺถาย.\csromanspacer{} อนาคามิมคฺคํ ปฏิลาภตฺถาย ฯเปฯ.\csromanspacer{} อนาคามิผลสมาปตฺตตฺถาย.\csromanspacer{} อรหตฺตมคฺคํ ปฏิลาภตฺถาย อุปฺปาทํ ปวตฺตํ นิมิตฺตํ อายูหนํ ปฏิสนฺธิํ คติํ นิพฺพตฺติํ อุปปตฺติํ ชาติํ ชรํ พฺยาธิํ มรณํ โสกํ ปริเทวํ อุปายาสํ ปฏิสงฺขา สนฺติฏฺฐนา ปญฺญา สงฺขารุเปกฺขาสุ ญาณํ.\csromanspacer{} อรหตฺตผลสมาปตฺตตฺถาย ฯเปฯ.\csromanspacer{} สุญฺญตวิหารสมาปตฺตตฺถาย ฯเปฯ.\csromanspacer{} อนิมิตฺตวิหารสมาปตฺตตฺถาย อุปฺปาทํ ปวตฺตํ นิมิตฺตํ อายูหนํ ปฏิสนฺธิํ ฯเปฯ.\csromanspacer{} ปฏิสงฺขา สนฺติฏฺฐนา ปญฺญา สงฺขารุเปกฺขาสุ ญาณํ.\csromanspacer{} อิมา ทส สงฺขารุเปกฺขา วิปสฺสนาวเสน อุปฺปชฺชนฺติ.}

\proseitem{58}{\csromanfolio{63}กติ สงฺขารุเปกฺขา กุสลา, กติ อกุสลา, กติ อพฺยากตา? ปนฺนรส สงฺขารุเปกฺขา กุสลา, ติสฺโส สงฺขารุเปกฺขา อพฺยากตา, นตฺถิ สงฺขารุเปกฺขา อกุสลา.}

\csromangathameasure{ปฏิสงฺขา สนฺติฏฺฐนา ปญฺญา, อฏฺฐ จิตฺตสฺส โคจรา. \\ ปุถุชฺชนสฺส เทฺว โหนฺติ, ตโย เสกฺขสฺส โคจรา. \\ ตโย จ วีตราคสฺส, เยหิ จิตฺตํ วิวฏฺฏติ. \\ อฏฺฐ สมาธิสฺส ปจฺจยา, ทส ญาณสฺส โคจรา. \\ อฏฺฐารส สงฺขารุเปกฺขา, ติณฺณํ วิโมกฺขาน ปจฺจยา. \\ อิเม อฏฺฐารสาการา, ปญฺญา ยสฺส ปริจฺจิตา. \\ กุสโล สงฺขารุเปกฺขาสุ, นานาทิฏฺฐีสุ น กมฺปตีติ.}
\csromangathagroup{\csromangathastanza{ปฏิสงฺขา สนฺติฏฺฐนา ปญฺญา, อฏฺฐ จิตฺตสฺส โคจรา. \\ ปุถุชฺชนสฺส เทฺว โหนฺติ, ตโย เสกฺขสฺส โคจรา.}\csromangathastanza{ตโย จ วีตราคสฺส, เยหิ จิตฺตํ วิวฏฺฏติ. \\ อฏฺฐ สมาธิสฺส ปจฺจยา, ทส ญาณสฺส โคจรา.}\csromangathastanza{อฏฺฐารส สงฺขารุเปกฺขา, ติณฺณํ วิโมกฺขาน ปจฺจยา. \\ อิเม อฏฺฐารสาการา, ปญฺญา ยสฺส ปริจฺจิตา. \\ กุสโล สงฺขารุเปกฺขาสุ, นานาทิฏฺฐีสุ น กมฺปตีติ.}}

\prose{ตํญาตฏฺเฐน ญาณํ, ปชานนฏฺเฐน ปญฺญา, เตน วุจฺจติ มุญฺจิตุกมฺยตาปฏิสงฺขาสนฺติฏฺฐนา ปญฺญา สงฺขารุเปกฺขาสุ ญาณํ.}

\csromanheader{2}{สงฺขารุเปกฺขาญาณนิทฺเทโส นวโม.}
\csromansectionrule
\csromanmatikaanchor{mtk.12}
\csromantocmarkref{subsection}{10. โคตฺรภุญาณนิทฺเทส}{mtk.12}
\bookmark[dest={mtk.12},level=2]{10. โคตฺรภุญาณนิทฺเทส}
\csromanheader{2}{10. โคตฺรภุญาณนิทฺเทส}
\proseitem{59}{กถํ พหิทฺธา วุฏฺฐานวิวฏฺฏเน ปญฺญา โคตฺรภุญาณํ\csromandash{} อุปฺปาทํ อภิภุยฺยตีติ โคตฺรภุ, ปวตฺตํ อภิภุยฺยตีติ โคตฺรภุ, นิมิตฺตํ อภิภุยฺยตีติ โคตฺรภุ, อายูหนํ อภิภุยฺยตีติ โคตฺรภุ, ปฏิสนฺธิํ อภิภุยฺยตีติ โคตฺรภุ, คติํ อภิภุยฺยตีติ โคตฺรภุ, นิพฺพตฺติํ อภิภุยฺยตีติ โคตฺรภุ, อุปปตฺติํ อภิภุยฺยตีติ โคตฺรภุ, ชาติํ อภิภุยฺยตีติ โคตฺรภุ, ชรํ อภิภุยฺยตีติ โคตฺรภุ, พฺยาธิํ อภิภุยฺยตีติ โคตฺรภุ, มรณํ อภิภุยฺยตีติ โคตฺรภุ, โสกํ อภิภุยฺยตีติ โคตฺรภุ, ปริเทวํ อภิภุยฺยตีติ โคตฺรภุ, อุปายาสํ อภิภุยฺยตีติ โคตฺรภุ, พหิทฺธา สงฺขารนิมิตฺตํ อภิภุยฺยตีติ โคตฺรภุ.\csromanspacer{}\csromanspacer{}อนุปฺปาทํ ปกฺขนฺทตีติ โคตฺรภุ, อปฺปวตฺตํ ปกฺขนฺทตีติ โคตฺรภุ ฯเปฯ นิโรธํ นิพฺพานํ ปกฺขนฺทตีติ โคตฺรภุ.}

\prose{อุปฺปาทํ อภิภุยฺยิตฺวา อนุปฺปาทํ ปกฺขนฺทตีติ โคตฺรภุ, ปวตฺตํ อภิภุยฺยิตฺวา อปฺปวตฺตํ ปกฺขนฺทตีติ โคตฺรภุ, นิมิตฺตํ อภิภุยฺยิตฺวา อนิมิตฺตํ ปกฺขนฺทตีติ \csromanfolio{64}โคตฺรภุ ฯเปฯ พหิทฺธา สงฺขารนิมิตฺตํ อภิภุยฺยิตฺวา นิโรธํ นิพฺพานํ ปกฺขนฺทตีติ โคตฺรภุ.}

\prose{อุปฺปาทา วุฏฺฐาตีติ โคตฺรภุ, ปวตฺตา วุฏฺฐาตีติ โคตฺรภุ, นิมิตฺตา วุฏฺฐาตีติ โคตฺรภุ, อายูหนา วุฏฺฐาตีติ โคตฺรภุ, ปฏิสนฺธิยา วุฏฺฐาตีติ โคตฺรภุ, คติยา วุฏฺฐาตีติ โคตฺรภุ, นิพฺพตฺติยา วุฏฺฐาตีติ โคตฺรภุ, อุปปตฺติยา วุฏฺฐาตีติ โคตฺรภุ, ชาติยา วุฏฺฐาตีติ โคตฺรภุ, ชราย วุฏฺฐตีติ โคตฺรภุ, พฺยาธิมฺหา วุฏฺฐาตีติ โคตฺรภุ, มรณา วุฏฺฐาตีติ โคตฺรภุ, โสกา วุฏฺฐาตีติ โคตฺรภุ, ปริเทวา วุฏฺฐาตีติ โคตฺรภุ, อุปายาสา วุฏฺฐาตีติ โคตฺรภุ, พหิทฺธา สงฺขารนิมิตฺตา วุฏฺฐาตีติ โคตฺรภุ.\csromanspacer{}\csromanspacer{}อนุปฺปาทํ ปกฺขนฺทตีติ โคตฺรภุ, อปฺปวตฺตํ ปกฺขนฺทตีติ โคตฺรภุ ฯเปฯ นิโรธํ นิพฺพานํ ปกฺขนฺทตีติ โคตฺรภุ.}

\prose{อุปฺปาทา วุฏฺฐหิตฺวา\footnote{วุฏฺฐิตฺวา (สฺยา, ก)} อนุปฺปาทํ ปกฺขนฺทตีติ โคตฺรภุ, ปวตฺตา วุฏฺฐหิตฺวา อปฺปวตฺตํ ปกฺขนฺทตีติ โคตฺรภุ, นิมิตฺตา วุฏฺฐหิตฺวา อนิมิตฺตํ ปกฺขนฺทตีติ โคตฺรภุ, อายูหนา วุฏฺฐหิตฺวา อนายูหนํ ปกฺขนฺทตีติ โคตฺรภุ, ปฏิสนฺธิยา วุฏฺฐหิตฺวา อปฺปฏิสนฺธิํ ปกฺขนฺทตีติ โคตฺรภุ, คติยา วุฏฺฐหิตฺวา อคติํ ปกฺขนฺทตีติ โคตฺรภุ, นิพฺพตฺติยา วุฏฺฐหิตฺวา อนิพฺพตฺติํ ปกฺขนฺทตีติ โคตฺรภุ, อุปปตฺติยา วุฏฺฐหิตฺวา อนุปปตฺติํ ปกฺขนฺทตีติ โคตฺรภุ, ชาติยา วุฏฺฐหิตฺวา อชาติํ ปกฺขนฺทตีติ โคตฺรภุ, ชราย วุฏฺฐหิตฺวา อชรํ ปกฺขนฺทตีติ โคตฺรภุ, พฺยาธิมฺหา วุฏฺฐหิตฺวา อพฺยาธิํ ปกฺขนฺทตีติ โคตฺรภุ, มรณา วุฏฺฐหิตฺวา อมตํ ปกฺขนฺทตีติ โคตฺรภุ, โสกา วุฏฺฐหิตฺวา อโสกํ ปกฺขนฺทตีติ โคตฺรภุ, ปริเทวา วุฏฺฐหิตฺวา อปริเทวํ ปกฺขนฺทตีติ โคตฺรภุ, อุปายาสา วุฏฺฐหิตฺวา อนุปายาสํ ปกฺขนฺทตีติ โคตฺรภุ, พหิทฺธา สงฺขารนิมิตฺตา วุฏฺฐหิตฺวา นิโรธํ นิพฺพานํ ปกฺขนฺทตีติ โคตฺรภุ.}

\prose{อุปฺปาทา วิวฏฺฏตีติ โคตฺรภุ, ปวตฺตา วิวฏฺฏตีติ โคตฺรภุ ฯเปฯ พหิทฺธา สงฺขารนิมิตฺตา วิวฏฺฏตีติ โคตฺรภุ.\csromanspacer{}\csromanspacer{}อนุปฺปาทํ ปกฺขนฺทตีติ โคตฺรภุ, อปฺปวตฺตํ ปกฺขนฺทตีติ โคตฺรภุ ฯเปฯ นิโรธํ นิพฺพานํ ปกฺขนฺทตีติ โคตฺรภุ.}

\prose{อุปฺปาทา วิวฏฺฏิตฺวา อนุปฺปาทํ ปกฺขนฺทตีติ โคตฺรภุ, ปวตฺตา วิวฏฺฏิตฺวา อปฺปวตฺตํ ปกฺขนฺทตีติ โคตฺรภุ ฯเปฯ พหิทฺธา สงฺขารนิมิตฺตา วิวฏฺฏิตฺวา นิโรธํ นิพฺพานํ ปกฺขนฺทตีติ โคตฺรภุ.}

\proseitem{60}{\csromanfolio{65}กติ โคตฺรภู ธมฺมา สมถวเสน อุปฺปชฺชนฺติ, กติ โคตฺรภู ธมฺมา วิปสฺสนาวเสน อุปฺปชฺชนฺติ? อฏฺฐ โคตฺรภู ธมฺมา สมถวเสน อุปฺปชฺชนฺติ, ทส โคตฺรภู ธมฺมา วิปสฺสนาวเสน อุปฺปชฺชนฺติ.}

\prose{กตเม อฏฺฐ โคตฺรภู ธมฺมา สมถวเสน อุปฺปชฺชนฺติ? ปฐมํ ฌานํ ปฏิลาภตฺถาย นีวรเณ อภิภุยฺยตีติ โคตฺรภุ, ทุติยํ ฌานํ ปฏิลาภตฺถาย วิตกฺกวิจาเร อภิภุยฺยตีติ โคตฺรภุ, ตติยํ ฌานํ ปฏิลาภตฺถาย ปีติํ อภิภุยฺยตีติ โคตฺรภุ, จตุตฺถํ ฌานํ ปฏิลาภตฺถาย สุขทุกฺเข อภิภุยฺยตีติ โคตฺรภุ, อากาสานญฺจายตนสมาปตฺติํ ปฏิลาภตฺถาย รูปสญฺญํ ปฏิฆสญฺญํ นานตฺตสญฺญํ อภิภุยฺยตีติ โคตฺรภุ, วิญฺญาณญฺจายตนสมาปตฺติํ ปฏิลาภตฺถาย อากาสานญฺจายตนสญฺญํ อภิภุยฺยตีติ โคตฺรภุ, อากิญฺจญฺญายตนสมาปตฺติํ ปฏิลาภตฺถาย วิญฺญาณญฺจายตนสญฺญํ อภิภุยฺยตีติ โคตฺรภุ, เนวสญฺญานาสญฺญายตนสมาปตฺติํ ปฏิลาภตฺถาย อากิญฺจญฺญายตนสญฺญํ อภิภุยฺยตีติ โคตฺรภุ.\csromanspacer{} อิเม อฏฺฐ โคตฺรภู ธมฺมา สมถวเสน อุปฺปชฺชนฺติ.}

\prose{กตเม ทส โคตฺรภู ธมฺมา วิปสฺสนาวเสน อุปฺปชฺชนฺติ? โสตาปตฺติมคฺคํ ปฏิลาภตฺถาย อุปฺปาทํ ปวตฺตํ นิมิตฺตํ อายูหนํ ปฏิสนฺธิํ คติํ นิพฺพตฺติํ อุปปตฺติํ ชาติํ ชรํ พฺยาธิํ มรณํ โสกํ ปริเทวํ อุปายาสํ พหิทฺธา สงฺขารนิมิตฺตํ อภิภุยฺยตีติ โคตฺรภุ, โสตาปตฺติผลสมาปตฺตตฺถาย อุปฺปาทํ ปวตฺตํ นิมิตฺตํ อายูหนํ ปฏิสนฺธิํ ฯเปฯ อภิภุยฺยตีติ โคตฺรภุ.\csromanspacer{} สกทาคามิมคฺคํ ปฏิลาภตฺถาย ฯเปฯ สกทาคามิผลสมาปตฺตตฺถาย.\csromanspacer{} อนาคามิมคฺคํ ปฏิลาภตฺถาย.\csromanspacer{} อนาคามิผลสมาปตฺตตฺถาย.\csromanspacer{} อรหตฺตมคฺคํ ปฏิลาภตฺถาย อุปฺปาทํ ปวตฺตํ นิมิตฺตํ อายูหนํ ปฏิสนฺธิํ คติํ นิพฺพตฺติํ อุปปตฺติํ ชาติํ ชรํ พฺยาธิํ มรณํ โสกํ ปริเทวํ อุปายาสํ พหิทฺธา สงฺขารนิมิตฺตํ อภิภุยฺยตีติ โคตฺรภุ, อรหตฺตผลสมาปตฺตตฺถาย.\csromanspacer{} สุญฺญตวิหารสมาปตฺตตฺถาย.\csromanspacer{} อนิมิตฺตวิหารสมาปตฺตตฺถาย อุปฺปาทํ ปวตฺตํ นิมิตฺตํ อายูหนํ ปฏิสนฺธิํ ฯเปฯ อภิภุยฺยตีติ โคตฺรภุ.\csromanspacer{} อิเม ทส โคตฺรภู ธมฺมา วิปสฺสนาวเสน อุปฺปชฺชนฺติ.}

\prose{กติ โคตฺรภู ธมฺมา กุสลา, กติ อกุสลา, กติ อพฺยากตา? ปนฺนรส โคตฺรภู ธมฺมา กุสลา, ตโย โคตฺรภู ธมฺมา อพฺยากตา, นตฺถิ โคตฺรภู ธมฺมา อกุสลาติ.}

\csromangathameasure{สามิสญฺจ นิรามิสํ, ปณิหิตญฺจ อปฺปณิหิตํ. \\ สญฺญุตฺตญฺจ วิสญฺญุตฺตํ, วุฏฺฐิตญฺจ อวุฏฺฐิตํ. \\ อฏฺฐ สมาธิสฺส ปจฺจยา, ทส ญาณสฺส โคจรา. \\ อฏฺฐารส โคตฺรภู ธมฺมา, ติณฺณํ วิโมกฺขาน ปจฺจยา. \\ อิเม อฏฺฐารสาการา, ปญฺญา ยสฺส ปริจฺจิตา. \\ กุสโล วิวฏฺเฏ วุฏฺฐาเน, นานาทิฏฺฐีสุ น กมฺปตีติ.}
\csromangathagroup{\csromangathastanza{\csromanfolio{66}สามิสญฺจ นิรามิสํ, ปณิหิตญฺจ อปฺปณิหิตํ. \\ สญฺญุตฺตญฺจ วิสญฺญุตฺตํ, วุฏฺฐิตญฺจ อวุฏฺฐิตํ.}\csromangathastanza{อฏฺฐ สมาธิสฺส ปจฺจยา, ทส ญาณสฺส โคจรา. \\ อฏฺฐารส โคตฺรภู ธมฺมา, ติณฺณํ วิโมกฺขาน ปจฺจยา.}\csromangathastanza{อิเม อฏฺฐารสาการา, ปญฺญา ยสฺส ปริจฺจิตา. \\ กุสโล วิวฏฺเฏ วุฏฺฐาเน, นานาทิฏฺฐีสุ น กมฺปตีติ.}}

\prose{ตํญาตฏฺเฐน ญาณํ, ปชานนฏฺเฐน ปญฺญา, เตน วุจฺจติ พหิทฺธา วุฏฺฐานวิวฏฺฏเน ปญฺญา โคตฺรภุญาณํ.}

\csromanheader{2}{โคตฺรภุญาณนิทฺเทโส ทสโม.}
\csromansectionrule
\csromanmatikaanchor{mtk.13}
\csromantocmarkref{subsection}{11. มคฺคญาณนิทฺเทส}{mtk.13}
\bookmark[dest={mtk.13},level=2]{11. มคฺคญาณนิทฺเทส}
\csromanheader{2}{11. มคฺคญาณนิทฺเทส}
\proseitem{61}{กถํ ทุภโต วุฏฺฐานวิวฏฺฏเน ปญฺญา มคฺเค ญาณํ\csromandash{} โสตาปตฺติมคฺคกฺขเณ ทสฺสนฏฺเฐน สมฺมาทิฏฺฐิ มิจฺฉาทิฏฺฐิยา วุฏฺฐาติ, ตทนุวตฺตกกิเลเสหิ จ ขนฺเธหิ จ วุฏฺฐาติ, พหิทฺธา จ สพฺพนิมิตฺเตหิ วุฏฺฐาติ.\csromanspacer{} เตน วุจฺจติ ทุภโต วุฏฺฐานวิวฏฺฏเน ปญฺญา มคฺเค ญาณํ.}

\prose{อภินิโรปนฏฺเฐน สมฺมาสงฺกปฺโป มิจฺฉาสงฺกปฺปา วุฏฺฐาติ, ตทนุวตฺตากกิเลเสหิ จ ขนฺเธหิ จ วุฏฺฐาติ, พหิทฺธา จ สพฺพนิมิตฺเตหิ วุฏฺฐาติ.\csromanspacer{} เตน วุจฺจติ ทุภโต วุฏฺฐานวิวฏฺฏเน ปญฺญา มคฺเค ญาณํ.}

\prose{ปริคฺคหฏฺเฐน สมฺมาวาจา มิจฺฉาวาจาย วุฏฺฐาติ, ตทนุวตฺตกกิเลเสหิ จ ขนฺเธหิ จ วุฏฺฐาติ, พหิทฺธา จ สพฺพนิมิตฺเตหิ วุฏฺฐาติ.\csromanspacer{} เตน วุจฺจติ ทุภโต วุฏฺฐานวิวฏฺฏเน ปญฺญา มคฺเค ญาณํ.}

\prose{สมุฏฺฐานฏฺเฐน สมฺมากมฺมนฺโต มิจฺฉากมฺมนฺตา วุฏฺฐาติ, ตทนุวตฺตกกิเลเสหิ จ ขนฺเธหิ จ วุฏฺฐาติ, พหิทฺธา จ สพฺพนิมิตฺเตหิ วุฏฺฐาติ.\csromanspacer{} เตน วุจฺจติ ทุภโต วุฏฺฐานวิวฏฺฏเน ปญฺญา มคฺเค ญาณํ.}

\prose{โวทานฏฺเฐน สมฺมา-อาชีโว มิจฺฉา-อาชีวา วุฏฺฐาติ, ตทนุวตฺตกกิเลเสหิ จ ขนฺเธหิ จ วุฏฺฐาติ, พหิทฺธา จ สพฺพนิมิตฺเตหิ วุฏฺฐาติ.\csromanspacer{} เตน วุจฺจติ ทุภโต วุฏฺฐานวิวฏฺฏเน ปญฺญา มคฺเค ญาณํ.}

\prose{\csromanfolio{67}ปคฺคหฏฺเฐน สมฺมาวายาโม มิจฺฉาวายามา วุฏฺฐาติ, ตทนุวตฺตกกิเลเสหิ จ ขนฺเธหิ จ วุฏฺฐาติ, พหิทฺธา จ สพฺพนิมิตฺเตหิ วุฏฺฐาติ.\csromanspacer{} เตน วุจฺจติ ทุภโต วุฏฺฐานวิวฏฺฏเน ปญฺญา มคฺเค ญาณํ.}

\prose{อุปฏฺฐานฏฺเฐน สมฺมาสติ มิจฺฉาสติยา วุฏฺฐาติ, ตทนุวตฺตกกิเลเสหิ จ ขนฺเธหิ จ วุฏฺฐาติ, พหิทฺธา จ สพฺพนิมิตฺเตหิ วุฏฺฐาติ.\csromanspacer{} เตน วุจฺจติ ทุภโต วุฏฺฐานวิวฏฺฏเน ปญฺญา มคฺเค ญาณํ.}

\prose{อวิกฺเขปฏฺเฐน สมฺมาสมาธิ มิจฺฉาสมาธิโต วุฏฺฐาติ, ตทนุวตฺตกกิเลเสหิ จ ขนฺเธหิ จ วุฏฺฐาติ, พหิทฺธา จ สพฺพนิมิตฺเตหิ วุฏฺฐาติ.\csromanspacer{} เตน วุจฺจติ ทุภโต วุฏฺฐานวิวฏฺฏเน ปญฺญา มคฺเค ญาณํ.}

\prose{สกทาคามิมคฺคกฺขเณ ทสฺสนฏฺเฐน สมฺมาทิฏฺฐิ ฯเปฯ.\csromanspacer{} อวิกฺเขปฏฺเฐน สมฺมาสมาธิ โอฬาริกา กามราคสญฺโญชนา ปฏิฆสญฺโญชนา โอฬาริกา กามราคานุสยา ปฏิฆานุสยา วุฏฺฐาติ, ตทนุวตฺตกกิเลเสหิ จ ขนฺเธหิ จ วุฏฺฐาติ, พหิทฺธา จ สพฺพนิมิตฺเตหิ วุฏฺฐาติ.\csromanspacer{} เตน วุจฺจติ ทุภโต วุฏฺฐานวิวฏฺฏเน ปญฺญา มคฺเค ญาณํ.}

\prose{อนาคามิมคฺคกฺขเณ ทสฺสนฏฺเฐน สมฺมาทิฏฺฐิ ฯเปฯ.\csromanspacer{} อวิกฺเขปฏฺเฐน สมฺมาสมาธิ อณุสหคตา กามราคสญฺโญชนา ปฏิฆสญฺโญชนา อณุสหคตา กามราคานุสยา ปฏิฆานุสยา วุฏฺฐาติ, ตทนุวตฺตกกิเลเสหิ จ ขนฺเธหิ จ วุฏฺฐาติ, พหิทฺธา จ สพฺพนิมิตฺเตหิ วุฏฺฐาติ.\csromanspacer{} เตน วุจฺจติ ทุภโต วุฏฺฐานวิวฏฺฏเน ปญฺญา มคฺเค ญาณํ.}

\prose{อรหตฺตมคฺคกฺขเณ ทสฺสนฏฺเฐน สมฺมาทิฏฺฐิ ฯเปฯ.\csromanspacer{} อวิกฺเขปฏฺเฐน สมฺมาสมาธิ รูปราคา อรูปราคา มานา อุทฺธจฺจา อวิชฺชาย มานานุสยา ภวราคานุสยา อวิชฺชานุสยา วุฏฺฐาติ, ตทนุวตฺตกกิเลเสหิ จ ขนฺเธหิ จ วุฏฺฐาติ, พหิทฺธา จ สพฺพนิมิตฺเตหิ วุฏฺฐาติ.\csromanspacer{} เตน วุจฺจติ ทุภโต วุฏฺฐานวิวฏฺฏเน ปญฺญา มคฺเค ญาณํ.}

\csromanhangingbegin
\hangingheaditem{62}{อชาตํ ฌาเปติ ชาเตน, ฌานํ เตน ปวุจฺจติ.}
\hangingbody{ฌานวิโมกฺเข กุสลตา, นานาทิฏฺฐีสุ น กมฺปติ.}
\hangingbody{สมาทหิตฺวา ยถา เจ วิปสฺสติ,}
\hangingbody{วิปสฺสมาโน ตถา เจ สมาทเห.}
\hangingbody{วิปสฺสนา จ สมโถ ตทา อหุ,}
\hangingbody{สมานภาคา ยุคนทฺธา วตฺตเร.}
\csromanhangingend

\csromangathameasure{ทุกฺขา สงฺขารา สุโข, นิโรโธ อิติ ทสฺสนํ. \\ ทุภโต วุฏฺฐิตา ปญฺญา, ผสฺเสติ อมตํ ปทํ. \\ วิโมกฺขจริยํ ชานาติ, นานตฺเตกตฺตโกวิโท. \\ ทฺวินฺนํ ญาณานํ กุสลตา, นานาทิฏฺฐีสุ น กมฺปตีติ.}
\csromangathagroup{\csromangathastanza{\csromanfolio{68}ทุกฺขา สงฺขารา สุโข, นิโรโธ อิติ ทสฺสนํ\footnote{นิโรโธติ ทสฺสนํ(สฺยา, ก)}. \\ ทุภโต วุฏฺฐิตา ปญฺญา, ผสฺเสติ อมตํ ปทํ.}\csromangathastanza{วิโมกฺขจริยํ ชานาติ, นานตฺเตกตฺตโกวิโท. \\ ทฺวินฺนํ ญาณานํ กุสลตา, นานาทิฏฺฐีสุ น กมฺปตีติ.}}

\prose{ตํญาตฏฺเฐน ญาณํ, ปชานนฏฺเฐน ปญฺญา.\csromanspacer{} เตน วุจฺจติ ทุภโต วุฏฺฐานวิวฏฺฏเน ปญฺญา มคฺเค ญาณํ.}

\csromanheader{2}{มคฺคญาณนิทฺเทโส เอกาทสโม.}
\csromansectionrule
\csromanmatikaanchor{mtk.14}
\csromantocmarkref{subsection}{12. ผลญาณนิทฺเทส}{mtk.14}
\bookmark[dest={mtk.14},level=2]{12. ผลญาณนิทฺเทส}
\csromanheader{2}{12. ผลญาณนิทฺเทส}
\proseitem{63}{กถํ ปโยคปฺปฏิปฺปสฺสทฺธิปญฺญา ผเล ญาณํ\csromandash{} โสตาปตฺติมคฺคกฺขเณ ทสฺสนฏฺเฐน สมฺมาทิฏฺฐิ มิจฺฉาทิฏฺฐิยา วุฏฺฐาติ, ตทนุวตฺตกกิเลเสหิ จ ขนฺเธหิ จ วุฏฺฐาติ, พหิทฺธา จ สพฺพนิมิตฺเตหิ วุฏฺฐาติ.\csromanspacer{} ตํปโยคปฺปฏิปฺปสฺสทฺธตฺตา อุปฺปชฺชติ สมฺมาทิฏฺฐิ, มคฺคสฺเสตํ ผลํ.}

\prose{อภินิโรปนฏฺเฐน สมฺมาสงฺกปฺโป มิจฺฉาสงฺกปฺปา วุฏฺฐาติ, ตทนุวตฺตกกิเลเสหิ จ ขนฺเธหิ จ วุฏฺฐาติ, พหิทฺธา จ สพฺพนิมิตฺเตหิ วุฏฺฐาติ.\csromanspacer{} ตํปโยคปฺปฏิปฺปสฺสทฺธตฺตา อุปฺปชฺชติ สมฺมาสงฺกปฺโป, มคฺคสฺเสตํ ผลํ.}

\prose{ปริคฺคหฏฺเฐน สมฺมาวาจา มิจฺฉาวาจาย วุฏฺฐาติ, ตทนุวตฺตกกิเลเสหิ จ ขนฺเธหิ จ วุฏฺฐาติ, พหิทฺธา จ สพฺพนิมิตฺเตหิ วุฏฺฐาติ.\csromanspacer{} ตํปโยคปฺปฏิปฺปสฺสทฺธตฺตา อุปฺปชฺชติ สมฺมาวาจา, มคฺคสฺเสตํ ผลํ.}

\prose{สมุฏฺฐานฏฺเฐน สมฺมากมฺมนฺโต มิจฺฉากมฺมนฺตา วุฏฺฐาติ, ตทนุวตฺตกกิเลเสหิ จ ขนฺเธหิ จ วุฏฺฐาติ, พหิทฺธา จ สพฺพนิมิตฺเตหิ วุฏฺฐาติ.\csromanspacer{} ตํปโยคปฺปฏิปฺปสฺสทฺธตฺตา อุปฺปชฺชติ สมฺมากมฺมนฺโต, มคฺคสฺเสตํ ผลํ.}

\prose{โวทานฏฺเฐน สมฺมา-อาชีโว มิจฺฉา-อาชีวา วุฏฺฐาติ, ตทนุวตฺตกกิเลเสหิ จ ขนฺเธหิ จ วุฏฺฐาติ, พหิทฺธา จ สพฺพนิมิตฺเตหิ วุฏฺฐาติ.\csromanspacer{} ตํปโยคปฺปฏิปฺปสฺสทฺธตฺตา อุปฺปชฺชติ สมฺมา-อาชีโว, มคฺคสฺเสตํ ผลํ.}

\prose{\csromanfolio{69}ปคฺคหฏฺเฐน สมฺมาวายาโม มิจฺฉาวายามา วุฏฺฐาติ, ตทนุวตฺตกกิเลเสหิ จ ขนฺเธหิ จ วุฏฺฐาติ, พหิทฺธา จ สพฺพนิมิตฺเตหิ วุฏฺฐาติ.\csromanspacer{} ตํปโยคปฺปฏิปฺปสฺสทฺธตฺตา อุปฺปชฺชติ สมฺมาวายาโม, มคฺคสฺเสตํ ผลํ.}

\prose{อุปฏฺฐานฏฺเฐน สมฺมาสติ มิจฺฉาสติยา วุฏฺฐาติ, ตทนุวตฺตกกิเลเสหิ จ ขนฺเธหิ จ วุฏฺฐาติ, พหิทฺธา จ สพฺพนิมิตฺเตหิ วุฏฺฐาติ.\csromanspacer{} ตํปโยคปฺปฏิปฺปสฺสทฺธตฺตา อุปฺปชฺชติ สมฺมาสติ, มคฺคสฺเสตํ ผลํ.}

\prose{อวิกฺเขปฏฺเฐน สมฺมาสมาธิ มิจฺฉาสมาธิโต วุฏฺฐาติ, ตทนุวตฺตกกิเลเสหิ จ ขนฺเธหิ จ วุฏฺฐาติ, พหิทฺธา จ สพฺพนิมิตฺเตหิ วุฏฺฐาติ.\csromanspacer{} ตํปโยคปฺปฏิปฺปสฺสทฺธตฺตา อุปฺปชฺชติ สมฺมาสมาธิ, มคฺคสฺเสตํ ผลํ.}

\prose{สกทาคามิมคฺคกฺขเณ ทสฺสนฏฺเฐน สมฺมาทิฏฺฐิ ฯเปฯ.\csromanspacer{} อวิกฺเขปฏฺเฐน สมฺมาสมาธิ โอฬาริกา กามราคสญฺโญชนา ปฏิฆสญฺโญชนา, โอฬาริกา กามราคานุสยา ปฏิฆานุสยา วุฏฺฐาติ, ตทนุวตฺตกกิเลเสหิ จ ขนฺเธหิ จ วุฏฺฐาติ, พหิทฺธา จ สพฺพนิมิตฺเตหิ วุฏฺฐาติ.\csromanspacer{} ตํปโยคปฺปฏิปฺปสฺสทฺธตฺตา อุปฺปชฺชติ สมฺมาสมาธิ, มคฺคสฺเสตํ ผลํ.}

\prose{อนาคามิมคฺคกฺขเณ ทสฺสนฏฺเฐน สมฺมาทิฏฺฐิ ฯเปฯ.\csromanspacer{} อวิกฺเขปฏฺเฐน สมฺมาสมาธิ อณุสหคตา กามราคสญฺโญชนา ปฏิฆสญฺโญชนา, อณุสหคตา กามราคานุสยา ปฏิฆานุสยา วุฏฺฐาติ, ตทนุวตฺตกกิเลเสหิ จ ขนฺเธหิ จ วุฏฺฐาติ, พหิทฺธา จ สพฺพนิมิตฺเตหิ วุฏฺฐาติ.\csromanspacer{} ตํปโยคปฺปฏิปฺปสฺสทฺธตฺตา อุปฺปชฺชติ สมฺมาสมาธิ, มคฺคสฺเสตํ ผลํ.}

\prose{อรหตฺตมคฺคกฺขเณ ทสฺสนฏฺเฐน สมฺมาทิฏฺฐิ ฯเปฯ.\csromanspacer{} อวิกฺเขปฏฺเฐน สมฺมาสมาธิ รูปราคา อรูปราคา มานา อุทฺธจฺจา อวิชฺชาย, มานานุสยา ภวราคานุสยา อวิชฺชานุสยา วุฏฺฐาติ, ตทนุวตฺตกกิเลเสหิ จ ขนฺเธหิ จ วุฏฺฐาติ, พหิทฺธา จ สพฺพนิมิตฺเตหิ วุฏฺฐาติ.\csromanspacer{} ตํปโยคปฺปฏิปฺปสฺสทฺธตฺตา อุปฺปชฺชติ สมฺมาสมาธิ, มคฺคสฺเสตํ ผลํ.\csromanspacer{} ตํญาตฏฺเฐน ญาณํ, ปชานนฏฺเฐน ปญฺญา.\csromanspacer{} เตน วุจฺจติ ปโยคปฺปฏิปฺปสฺสทฺธิปญฺญา ผเล ญาณํ.}

\vspace{\tipitakaparskip}
\csromancenter{ผลญาณนิทฺเทโส ทฺวาทสโม.}
\csromanmatikaanchor{mtk.15}
\csromantocmarkref{subsection}{13. วิมุตฺติญาณนิทฺเทส}{mtk.15}
\bookmark[dest={mtk.15},level=2]{13. วิมุตฺติญาณนิทฺเทส}
\csromanheader{2}{13. วิมุตฺติญาณนิทฺเทส}
\proseitem{64}{\csromanfolio{70}กถํ ฉินฺนวฏุมานุปสฺสเน ปญฺญา วิมุตฺติญาณํ\csromandash{} โสตาปตฺติมคฺเคน สกฺกายทิฏฺฐิ วิจิกิจฺฉา สีลพฺพตปรามาโส, ทิฏฺฐานุสโย วิจิกิจฺฉานุสโย อตฺตโน จิตฺตสฺส อุปกฺกิเลสา สมฺมา สมุจฺฉินฺนา โหนฺติ, อิเมหิ ปญฺจหิ อุปกฺกิเลเสหิ สปริยุฏฺฐาเนหิ จิตฺตํ วิมุตฺตํ โหติ สุวิมุตฺตํ.\csromanspacer{} ตํวิมุตฺติญาตฏฺเฐน ญาณํ, ปชานนฏฺเฐน ปญฺญา.\csromanspacer{} เตน วุจฺจติ ฉินฺนวฏุมานุปสฺสเน ปญฺญา วิมุตฺติญาณํ.}

\prose{สกทาคามิมคฺเคน โอฬาริกํ กามราคสญฺโญชนํ ปฏิฆสญฺโญชนํ, โอฬาริโก กามราคานุสโย ปฏิฆานุสโย อตฺตโน จิตฺตสฺส อุปกฺกิเลสา สมฺมา สมุจฺฉินฺนา โหนฺติ, อิเมหิ จตูหิ อุปกฺกิเลเสหิ สปริยุฏฺฐาเนหิ จิตฺตํ วิมุตฺตํ โหติ สุวิมุตฺตํ.\csromanspacer{} ตํวิมุตฺติญาตฏฺเฐน ญาณํ, ปชานนฏฺเฐน ปญฺญา.\csromanspacer{} เตน วุจฺจติ ฉินฺนวฏุมานุปสฺสเน ปญฺญา วิมุตฺติญาณํ.}

\prose{อนาคามิมคฺเคน อณุสหคตํ กามราคสญฺโญชนํ ปฏิฆสญฺโญชนํ, อณุสหคโต กามราคานุสโย ปฏิฆานุสโย อตฺตโน จิตฺตสฺส อุปกฺกิเลสา สมฺมา สมุจฺฉินฺนา โหนฺติ, อิเมหิ จตูหิ อุปกฺกิเลเสหิ สปริยุฏฺฐาเนหิ จิตฺตํ วิมุตฺตํ โหติ สุวิมุตฺตํ.\csromanspacer{} ตํวิมุตฺติญาตฏฺเฐน ญาณํ, ปชานนฏฺเฐน ปญฺญา.\csromanspacer{} เตน วุจฺจติ ฉินฺนวฏุมานุปสฺสเน ปญฺญา วิมุตฺติญาณํ.}

\prose{อรหตฺตมคฺเคน รูปราโค อรูปราโค มาโน อุทฺธจฺจํ อวิชฺชา, มานานุสโย ภวราคานุสโย อวิชฺชานุสโย อตฺตโน จิตฺตสฺส อุปกฺกิเลสา สมฺมา สมุจฺฉินฺนา โหนฺติ, อิเมหิ อฏฺฐหิ อุปกฺกิเลเสหิ สปริยุฏฺฐาเนหิ จิตฺตํ วิมุตฺตํ โหติ สุวิมุตฺตํ.\csromanspacer{} ตํวิมุตฺติญาตฏฺเฐน ญาณํ, ปชานนฏฺเฐน ปญฺญา.\csromanspacer{} เตน วุจฺจติ ฉินฺนวฏุมานุปสฺสเน ปญฺญา วิมุตฺติญาณํ.\csromanspacer{} ตํญาตฏฺเฐน ญาณํ, ปชานนฏฺเฐน ปญฺญา.\csromanspacer{} เตน วุจฺจติ ฉินฺนวฏุมานุปสฺสเน ปญฺญา วิมุตฺติญาณํ.}

\csromanheader{2}{วิมุตฺติญาณนิทฺเทโส เตรสโม.}
\csromansectionrule
\csromanmatikaanchor{mtk.16}
\csromantocmarkref{subsection}{14. ปจฺจเวกฺขณญาณนิทฺเทส}{mtk.16}
\bookmark[dest={mtk.16},level=2]{14. ปจฺจเวกฺขณญาณนิทฺเทส}
\csromanheader{2}{14. ปจฺจเวกฺขณญาณนิทฺเทส}
\proseitem{65}{กถํ ตทา สมุทาคเต ธมฺเม ปสฺสเน ปญฺญา ปจฺจเวกฺขเณ ญาณํ\csromandash{}โสตาปตฺติมคฺคกฺขเณ ทสฺสนฏฺเฐน สมฺมาทิฏฺฐิ ตทา สมุทาคตา. \csromanfolio{71}อภินิโรปนฏฺเฐน สมฺมาสงฺกปฺโป ตทา สมุทาคโต.\csromanspacer{} ปริคฺคหฏฺเฐน สมฺมาวาจา ตทา สมุทาคตา.\csromanspacer{} สมุฏฺฐานฏฺเฐน สมฺมากมฺมนฺโต ตทา สมุทาคโต.\csromanspacer{} โวทานฏฺเฐน สมฺมา-อาชีโว ตทา สมุทาคโต.\csromanspacer{} ปคฺคหฏฺเฐน สมฺมาวายาโม ตทา สมุทาคโต.\csromanspacer{} อุปฏฺฐานฏฺเฐน สมฺมาสติ ตทา สมุทาคตา.\csromanspacer{} อวิกฺเขปฏฺเฐน สมฺมาสมาธิ ตทา สมุทาคโต.}

\prose{อุปฏฺฐานฏฺเฐน สติสมฺโพชฺฌงฺโค ตทา สมุทาคโต.\csromanspacer{} ปวิจยฏฺเฐน ธมฺมวิจยสมฺโพชฺฌงฺโค ตทา สมุทาคโต.\csromanspacer{} ปคฺคหฏฺเฐน วีริยสมฺโพชฺฌงฺโค ตทา สมุทาคโต.\csromanspacer{} ผรณฏฺเฐน ปีติสมฺโพชฺฌงฺโค ตทา สมุทาคโต.\csromanspacer{} อุปสมฏฺเฐน ปสฺสทฺธิสมฺโพชฺฌงฺโค ตทา สมุทาคโต.\csromanspacer{} อวิกฺเขปฏฺเฐน สมาธิสมฺโพชฺฌงฺโค ตทา สมุทาคโต.\csromanspacer{} ปฏิสงฺขานฏฺเฐน อุเปกฺขาสมฺโพชฺฌงฺโค ตทา สมุทาคโต.}

\prose{อสฺสทฺธิเย อกมฺปิยฏฺเฐน สทฺธาพลํ ตทา สมุทาคตํ, โกสชฺเช อกมฺปิยฏฺเฐน วีริยพลํ ตทา สมุทาคตํ.\csromanspacer{} ปมาเท อกมฺปิยฏฺเฐน สติพลํ ตทา สมุทาคตํ.\csromanspacer{} อุทฺธจฺเจ อกมฺปิยฏฺเฐน สมาธิพลํ ตทา สมุทาคตํ.\csromanspacer{} อวิชฺชาย อกมฺปิยฏฺเฐน ปญฺญาพลํ ตทา สมุทาคตํ.}

\prose{อธิโมกฺขฏฺเฐน สทฺธินฺทฺริยํ ตทา สมุทาคตํ.\csromanspacer{} ปคฺคหฏฺเฐน วีริยินฺทฺริยํ ตทา สมุทาคตํ.\csromanspacer{} อุปฏฺฐานฏฺเฐน สตินฺทฺริยํ ตทา สมุทาคตํ.\csromanspacer{} อวิกฺเขปฏฺเฐน สมาธินฺทฺริยํ ตทา สมุทาคตํ.\csromanspacer{} ทสฺสนฏฺเฐน ปญฺญินฺทฺริยํ ตทา สมุทาคตํ.}

\prose{อาธิปเตยฺยฏฺเฐน อินฺทฺริยา ตทา สมุทาคตา.\csromanspacer{} อกมฺปิยฏฺเฐน พลา ตทา สมุทาคตา.\csromanspacer{} นิยฺยานฏฺเฐน สมฺโพชฺฌงฺคา ตทา สมุทาคตา.\csromanspacer{} เหตุฏฺเฐน มคฺโค ตทา สมุทาคโต.\csromanspacer{} อุปฏฺฐานฏฺเฐน สติปฏฺฐานา ตทา สมุทาคตา.\csromanspacer{} ปทหนฏฺเฐน สมฺมปฺปธานา ตทา สมุทาคตา.\csromanspacer{} อิชฺฌนฏฺเฐน อิทฺธิปาทา ตทา สมุทาคตา.\csromanspacer{} ตถฏฺเฐน สจฺจา ตทา สมุทาคตา.\csromanspacer{} อวิกฺเขปฏฺเฐน สมโถ ตทา สมุทาคโต.\csromanspacer{} อนุปสฺสนฏฺเฐน วิปสฺสนา ตทา สมุทาคตา.\csromanspacer{} เอกรสฏฺเฐน สมถวิปสฺสนา ตทา สมุทาคตา.\csromanspacer{} อนติวตฺตนฏฺเฐน ยุคนทฺธํ ตทา สมุทาคตํ.\csromanspacer{} สํวรฏฺเฐน สีลวิสุทฺธิ ตทา สมุทาคตา.\csromanspacer{} อวิกฺเขปฏฺเฐน จิตฺตวิสุทฺธิ ตทา สมุทาคตา.\csromanspacer{} ทสฺสนฏฺเฐน ทิฏฺฐิวิสุทฺธิ ตทา สมุทาคตา.\csromanspacer{} วิมุตฺตฏฺเฐน วิโมกฺโข ตทา สมุทาคโต.\csromanspacer{} ปฏิเวธฏฺเฐน วิชฺชา ตทา สมุทาคตา.\csromanspacer{} ปริจฺจาคฏฺเฐน วิมุตฺติ ตทา สมุทาคตา.\csromanspacer{} สมุจฺเฉทฏฺเฐน ขเย ญาณํ ตทา สมุทาคตํ.}

\prose{\csromanfolio{72}ฉนฺโท มูลฏฺเฐน ตทา สมุทาคโต.\csromanspacer{} มนสิกาโร สมุฏฺฐานฏฺเฐน ตทา สมุทาคโต.\csromanspacer{} ผสฺโส สโมธานฏฺเฐน ตทา สมุทาคโต.\csromanspacer{} เวทนา สโมสรณฏฺเฐน ตทา สมุทาคตา.\csromanspacer{} สมาธิ ปมุขฏฺเฐน ตทา สมุทาคโต.\csromanspacer{} สติ อาธิปเตยฺยฏฺเฐน ตทา สมุทาคตา.\csromanspacer{} ปญฺญา ตทุตฺตรฏฺเฐน ตทา สมุทาคตา.\csromanspacer{} วิมุตฺติ สารฏฺเฐน ตทา สมุทาคตา.\csromanspacer{} อมโตคธํ นิพฺพานํ ปริโยสานฏฺเฐน ตทา สมุทาคตํ.\csromanspacer{} วุฏฺฐหิตฺวา ปจฺจเวกฺขติ, อิเม ธมฺมา ตทา สมุทาคตา.}

\prose{โสตาปตฺติผลกฺขเณ ทสฺสนฏฺเฐน สมฺมาทิฏฺฐิ ตทา สมุทาคตา.\csromanspacer{} อภินิโรปนฏฺเฐน สมฺมาสงฺกปฺโป ตทา สมุทาคโต ฯเปฯ.\csromanspacer{} ปฏิปฺปสฺสทฺธฏฺเฐน อนุปฺปาเท ญาณํ ตทา สมุทาคตํ.\csromanspacer{}\csromanspacer{}ฉนฺโท มูลฏฺเฐน ตทา สมุทาคโต.\csromanspacer{} มนสิกาโร สมุฏฺฐานฏฺเฐน ตทา สมุทาคโต.\csromanspacer{} ผสฺโส สโมธานฏฺเฐน ตทา สมุทาคโต.\csromanspacer{} เวทนา สโมสรณฏฺเฐน ตทา สมุทาคตา.\csromanspacer{} สมาธิ ปมุขฏฺเฐน ตทา สมุทาคโต.\csromanspacer{} สติ อาธิปเตยฺยฏฺเฐน ตทา สมุทาคตา.\csromanspacer{} ปญฺญา ตทุตฺตรฏฺเฐน ตทา สมุทาคตา.\csromanspacer{} วิมุตฺติ สารฏฺเฐน ตทา สมุทาคตา.\csromanspacer{} อมโตคธํ นิพฺพานํ ปริโยสานฏฺเฐน ตทา สมุทาคตํ.\csromanspacer{} วุฏฺฐหิตฺวา ปจฺจเวกฺขติ, อิเม ธมฺมา ตทา สมุทาคตา.}

\prose{สกทาคามิมคฺคกฺขเณ ฯเปฯ.\csromanspacer{} สกทาคามิผลกฺขเณ ฯเปฯ.\csromanspacer{} อนาคามิมคฺคกฺขเณ ฯเปฯ.\csromanspacer{} อนาคามิผลกฺขเณ ฯเปฯ.\csromanspacer{} อรหตฺตมคฺคกฺขเณ ทสฺสนฏฺเฐน สมฺมาทิฏฺฐิ ตทา สมุทาคตา ฯเปฯ.\csromanspacer{} สมุจฺเฉทฏฺเฐน ขเย ญาณํ ตทา สมุทาคตํ.\csromanspacer{}\csromanspacer{}ฉนฺโท มูลฏฺเฐน ตทา สมุทาคโต ฯเปฯ อมโตคธํ นิพฺพานํ ปริโยสานฏฺเฐน ตทา สมุทาคตํ.\csromanspacer{} วุฏฺฐหิตฺวา ปจฺจเวกฺขติ, อิเม ธมฺมา ตทา สมุทาคตา.}

\prose{อรหตฺตผลกฺขเณ ทสฺสนฏฺเฐน สมฺมาทิฏฺฐิ ตทา สมุทาคตา ฯเปฯ.\csromanspacer{} ปฏิปฺปสฺสทฺธฏฺเฐน อนุปฺปาเท ญาณํ ตทา สมุทาคตํ.\csromanspacer{} ฉนฺโท มูลฏฺเฐน ตทา สมุทาคโต ฯเปฯ อมโตคธํ นิพฺพานํ ปริโยสานฏฺเฐน ตทา สมุทาคตํ.\csromanspacer{} วุฏฺฐหิตฺวา ปจฺจเวกฺขติ, อิเม ธมฺมา ตทา สมุทาคตา, ตํญาตฏฺเฐน ญาณํ, ปชานนฏฺเฐน ปญฺญา.\csromanspacer{} เตน วุจฺจติ ตทา สมุทาคเต ธมฺเม ปสฺสเน ปญฺญา ปจฺจเวกฺขเณ ญาณํ.}

\vspace{\tipitakaparskip}
\csromancenter{ปจฺจเวกฺขณญาณนิทฺเทโส จุทฺทสโม.}
\csromanmatikaanchor{mtk.17}
\csromantocmarkref{subsection}{15. วตฺถุนานตฺตญาณนิทฺเทส}{mtk.17}
\bookmark[dest={mtk.17},level=2]{15. วตฺถุนานตฺตญาณนิทฺเทส}
\csromanheader{2}{15. วตฺถุนานตฺตญาณนิทฺเทส}
\proseitem{66}{\csromanfolio{73}กถํ อชฺฌตฺตววตฺถาเน ปญฺญา วตฺถุนานตฺเต ญาณํ, กถํ อชฺฌตฺตธมฺเม ววตฺเถติ\csromandash{}จกฺขุํ อชฺฌตฺตํ ววตฺเถติ, โสตํ อชฺฌตฺตํ ววตฺเถติ, ฆานํ อชฺฌตฺตํ ววตฺเถติ, ชิวฺหํ อชฺฌตฺตํ ววตฺเถติ, กายํ อชฺฌตฺตํ ววตฺเถติ, มนํ อชฺฌตฺตํ ววตฺเถติ.}

\prose{กถํ จกฺขุํ อชฺฌตฺตํ ววตฺเถติ? จกฺขุ อวิชฺชาสมฺภูตนฺติ ววตฺเถติ, จกฺขุ ตณฺหาสมฺภูตนฺติ ววตฺเถติ, จกฺขุ กมฺมสมฺภูตนฺติ ววตฺเถติ, จกฺขุ อาหารสมฺภูตนฺติ ววตฺเถติ, จกฺขุ จตุนฺนํ มหาภูตานํ อุปาทายาติ ววตฺเถติ, จกฺขุ อุปฺปนฺนนฺติ ววตฺเถติ, จกฺขุ สมุทาคตนฺติ ววตฺเถติ.\csromanspacer{} จกฺขุ อหุตฺวา สมฺภูตํ, หุตฺวา น ภวิสฺสตีติ ววตฺเถติ, จกฺขุํ อนฺตวนฺตโต ววตฺเถติ, จกฺขุ อทฺธุคํ อสสฺสตํ วิปริณามธมฺมนฺติ ววตฺเถติ, จกฺขุ อนิจฺจํ สงฺขตํ ปฏิจฺจสมุปฺปนฺนํ ขยธมฺมํ วยธมฺมํ วิราคธมฺมํ นิโรธธมฺมนฺติ ววตฺเถติ, จกฺขุํ อนิจฺจโต ววตฺเถติ, โน นิจฺจโต, ทุกฺขโต ววตฺเถติ, โน สุขโต, อนตฺตโต ววตฺเถติ, โน อตฺตโต, นิพฺพินฺทติ, โน นนฺทติ, วิรชฺชติ, โน รชฺชติ, นิโรเธติ, โน สมุเทติ, ปฏินิสฺสชฺชติ, โน อาทิยติ.\csromanspacer{}\csromanspacer{}อนิจฺจโต ววตฺเถนฺโต นิจฺจสญฺญํ ปชหาติ, ทุกฺขโต ววตฺเถนฺโต สุขสญฺญํ ปชหาติ, อนตฺตโต ววตฺเถนฺโต อตฺตสญฺญํ ปชหาติ, นิพฺพินฺทนฺโต นนฺทิํ ปชหาติ, วิรชฺชนฺโต ราคํ ปชหาติ, นิโรเธนฺโต สมุทยํ ปชหาติ, ปฏินิสฺสชฺชนฺโต อาทานํ ปชหาติ.\csromanspacer{} เอวํ จกฺขุํ อชฺฌตฺตํ ววตฺเถติ.}

\prose{กถํ โสตํ อชฺฌตฺตํ ววตฺเถติ? โสตํ อวิชฺชาสมฺภูตนฺติ ววตฺเถติ ฯเปฯ เอวํ โสตํ อชฺฌตฺตํ ววตฺเถติ.\csromanspacer{}\csromanspacer{}กถํ ฆานํ อชฺฌตฺตํ ววตฺเถติ? ฆานํ อวิชฺชาสมฺภูตนฺติ ววตฺเถติ ฯเปฯ เอวํ ฆานํ อชฺฌตฺตํ ววตฺเถติ.\csromanspacer{}\csromanspacer{}กถํ ชิวฺหํ อชฺฌตฺตํ ววตฺเถติ? ชิวฺหา อวิชฺชาสมฺภูตาติ ววตฺเถติ, ชิวฺหา ตณฺหาสมฺภูตาติ ววตฺเถติ, ชิวฺหา กมฺมสมฺภูตาติ ววตฺเถติ, ชิวฺหา อาหารสมฺภูตาติ ววตฺเถติ, ชิวฺหา จตุนฺนํ มหาภูตานํ อุปาทายาติ ววตฺเถติ, ชิวฺหา อุปฺปนฺนาติ ววตฺเถติ, ชิวฺหา สมุทาคตาติ ววตฺเถติ, ชิวฺหา อหุตฺวา สมฺภูตา, หุตฺวา น ภวิสฺสตีติ ววตฺเถติ, ชิวฺหํ อนฺตวนฺตโต ววตฺเถติ, ชิวฺหา อทฺธุวา อสสฺสตา วิปริณามธมฺมาติ ววตฺเถติ, ชิวฺหา อนิจฺจา สงฺขตา ปฏิจฺจสมุปฺปนฺนา ขยธมฺมา วยธมฺมา วิราคธมฺมา นิโรธธมฺมาติ ววตฺเถติ, ชิวฺหํ อนิจฺจโต ววตฺเถติ, \csromanfolio{74}โน นิจฺจโต ฯเปฯ ปฏินิสฺสชฺชติ, โน อาทิยติ.\csromanspacer{}\csromanspacer{}อนิจฺจโต ววตฺเถนฺโต นิจฺจสญฺญํ ปชหาติ ฯเปฯ ปฏินิสฺสชฺชนฺโต อาทานํ ปชหาติ.\csromanspacer{} เอวํ ชิวฺหํ อชฺฌตฺตํ ววตฺเถติ.}

\prose{กถํ กายํ อชฺฌตฺตํ ววตฺเถติ? กาโย อวิชฺชาสมฺภูโตติ ววตฺเถติ, กาโย ตณฺหาสมฺภูโตติ ววตฺเถติ, กาโย กมฺมสมฺภูโตติ ววตฺเถติ, กาโย อาหารสมฺภูโตติ ววตฺเถติ, กาโย จตุนฺนํ มหาภูตานํ อุปาทายาติ ววตฺเถติ, กาโย อุปฺปนฺโนติ ววตฺเถติ, กาโย สมุทาคโตติ ววตฺเถติ, กาโย อหุตฺวา สมฺภูโต, หุตฺวา น ภวิสฺสตีติ ววตฺเถติ, กายํ อนฺตวนฺตโต ววตฺเถติ, กาโย อทฺธุโว อสสฺสโต วิปริณามธมฺโมติ ววตฺเถติ, กาโย อนิจฺโจ สงฺขโต ปฏิจฺจสมุปฺปนฺโน ขยธมฺโม วยธมฺโม วิราคธมฺโม นิโรธธมฺโมติ ววตฺเถติ, กายํ อนิจฺจโต ววตฺเถติ, โน นิจฺจโต, ทุกฺขโต ววตฺเถติ, โน สุขโต ฯเปฯ ปฏินิสฺสชฺชติ, โน อาทิยติ.\csromanspacer{}\csromanspacer{}อนิจฺจโต ววตฺเถนฺโต นิจฺจสญฺญํ ปชหาติ, ทุกฺขโต ววตฺเถนฺโต สุขสญฺญํ ปชหาติ ฯเปฯ ปฏินิสฺสชฺชนฺโต อาทานํ ปชหาติ.\csromanspacer{} เอวํ กายํ อชฺฌตฺตํ ววตฺเถติ.}

\prose{กถํ มนํ อชฺฌตฺตํ ววตฺเถติ? มโน อวิชฺชาสมฺภูโตติ ววตฺเถติ, มโน ตณฺหาสมฺภูโตติ ววตฺเถติ, มโน กมฺมสมฺภูโตติ ววตฺเถติ, มโน อาหารสมฺภูโตติ ววตฺเถติ, มโน อุปฺปนฺโนติ ววตฺเถติ, มโน สมุทาคโตติ ววตฺเถติ, มโน อหุตฺวา สมฺภูโต, หุตฺวา น ภวิสฺสตีติ ววตฺเถติ, มนํ อนฺตวนฺตโต ววตฺเถติ, มโน อทฺธุโว อสสฺสโต วิปริณามธมฺโมติ ววตฺเถติ, มโน อนิจฺโจ สงฺขโต ปฏิจฺจสมุปฺปนฺโน ขยธมฺโม วยธมฺโม วิราคธมฺโม นิโรธธมฺโมติ ววตฺเถติ.\csromanspacer{} มนํ อนิจฺจโต ววตฺเถติ, โน นิจฺจโต, ทุกฺขโต ววตฺเถติ, โน สุขโต, อนตฺตโต ววตฺเถติ, โน อตฺตโต นิพฺพินฺทติ, โน นนฺทติ, วิรชฺชติ, โน รชฺชติ, นิโรเธติ, โน สมุเทติ, ปฏินิสฺสชฺชติ, โน อาทิยติ.\csromanspacer{}\csromanspacer{}อนิจฺจโต ววตฺเถนฺโต นิจฺจสญฺญํ ปชหาติ.\csromanspacer{} ทุกฺขโต ววตฺเถนฺโต สุขสญฺญํ ปชหาติ, อนตฺตโต ววตฺเถนฺโต อตฺตสญฺญํ ปชหาติ, นิพฺพินฺทนฺโต นนฺทิํ ปชหาติ, วิรชฺชนฺโต ราคํ ปชหาติ, นิโรเธนฺโต สมุทยํ ปชหาติ, ปฏินิสฺสชฺชนฺโต อาทานํ ปชหาติ.\csromanspacer{} เอวํ มนํ อชฺฌตฺตํ ววตฺเถติ.\csromanspacer{} เอวํ อชฺฌตฺตธมฺเม ววตฺเถติ.\csromanspacer{} ตํญาตฏฺเฐน \csromanfolio{75}ญาณํ ปชานนฏฺเฐน ปญฺญา.\csromanspacer{} เตน วุจฺจติ อชฺฌตฺตววตฺถาเน ปญฺญา วตฺถุนานตฺเต ญาณํ.}

\csromanheader{2}{วตฺถุนานตฺตญาณนิทฺเทโส ปนฺนรสโม.}
\csromansectionrule
\csromanmatikaanchor{mtk.18}
\csromantocmarkref{subsection}{16. โคจรนานตฺตญาณนิทฺเทส}{mtk.18}
\bookmark[dest={mtk.18},level=2]{16. โคจรนานตฺตญาณนิทฺเทส}
\csromanheader{2}{16. โคจรนานตฺตญาณนิทฺเทส}
\proseitem{67}{กถํ พหิทฺธา ววตฺถาเน ปญฺญา โคจรนานตฺเต ญาณํ\csromandash{}กถํ พหิทฺธา ธมฺเม ววตฺเถติ, รูเป พหิทฺธา ววตฺเถติ, สทฺเท พหิทฺธา ววตฺเถติ, คนฺเธ พหิทฺธา ววตฺเถติ, รเส พหิทฺธา ววตฺเถติ, โผฏฺฐพฺเพ พหิทฺธา ววตฺเถติ, ธมฺเม พหิทฺธา ววตฺเถติ?}

\prose{กถํ รูเป พหิทฺธา ววตฺเถติ? รูปา อวิชฺชาสมฺภูตาติ ววตฺเถติ, รูปา ตณฺหาสมฺภูตาติ ววตฺเถติ, รูปา กมฺมสมฺภูตาติ ววตฺเถติ, รูปา อาหารสมฺภูตาติ ววตฺเถติ.\csromanspacer{} รูปา จตุนฺนํ มหาภูตานํ อุปาทายาติ ววตฺเถติ, รูปา อุปฺปนฺนาติ ววตฺเถติ, รูปา สมุทาคตาติ ววตฺเถติ, รูปา อหุตฺวา สมฺภูตา, หุตฺวา น ภวิสฺสนฺตีติ ววตฺเถติ, รูเป อนฺตวนฺตโต ววตฺเถติ, รูปา อทฺธุวา อสสฺสตา วิปริณามธมฺมาติ ววตฺเถติ.\csromanspacer{} รูปา อนิจฺจา สงฺขตา ปฏิจฺจสมุปฺปนฺนา ขยธมฺมา วยธมฺมา วิราคธมฺมา นิโรธธมฺมาติ ววตฺเถติ, รูเป อนิจฺจโต ววตฺเถติ, โน นิจฺจโต, ทุกฺขโต ววตฺเถติ โน สุขโต, อนตฺตโต ววตฺเถติ, โน อตฺตโต, นิพฺพินฺทติ, โน นนฺทติ, วิรชฺชติ โน รชฺชติ, นิโรเธติ, โน สมุเทติ, ปฏินิสฺสชฺชติ, โน อาทิยติ.\csromanspacer{}\csromanspacer{}อนิจฺจโต ววตฺเถนฺโต นิจฺจสญฺญํ ปชหาติ, ทุกฺขโต ววตฺเถนฺโต สุขสญฺญํ ปชหาติ, อนตฺตโต ววตฺเถนฺโต อตฺตสญฺญํ ปชหาติ, นิพฺพินฺทนฺโต นนฺทิํ ปชหาติ, วิรชฺชนฺโต ราคํ ปชหาติ, นิโรเธนฺโต สมุทยํ ปชหาติ, ปฏินิสฺสชฺชนฺโต อาทานํ ปชหาติ.\csromanspacer{} เอวํ รูเป พหิทฺธา ววตฺเถติ.}

\prose{กถํ สทฺเท พหิทฺธา ววตฺเถติ? สทฺทา จตุนฺนํ มหาภูตานํ อุปาทายาติ ววตฺเถติ, สทฺทา อุปฺปนฺนาติ ววตฺเถติ, สทฺทา สมุทาคตาติ ววตฺเถติ.\csromanspacer{} สทฺธา อหุตฺวา สมฺภูตา, หุตฺวา น ภวิสฺสนฺตีติ ววตฺเถติ, สทฺเท อนฺตวนฺตโต ววตฺเถติ, สทฺทา \csromanfolio{76}อทฺธุวา อสสฺสตา วิปริณามธมฺมาติ ววตฺเถติ, สทฺทา อนิจฺจา สงฺขตา ปฏิจฺจสมุปฺปนฺนา ขยธมฺมา วยธมฺมา วิราคธมฺมา นิโรธธมฺมาติ ววตฺเถติ, สทฺเท อนิจฺจโต ววตฺเถติ, โน นิจฺจโต ฯเปฯ.\csromanspacer{} เอวํ สทฺเท พหิทฺธา ววตฺเถติ.}

\prose{กถํ คนฺเธ พหิทฺธา ววตฺเถติ? คนฺธา อวิชฺชาสมฺภูตาติ ววตฺเถติ, คนฺธา ตณฺหาสมฺภูตาติ ววตฺเถติ ฯเปฯ.\csromanspacer{} เอวํ คนฺเธ พหิทฺธา ววตฺเถติ.\csromanspacer{}\csromanspacer{}กถํ รเส พหิทฺธา ววตฺเถติ? รสา อวิชฺชาสมฺภูตาติ ววตฺเถติ, รสา ตณฺหาสมฺภูตาติ ววตฺเถติ ฯเปฯ.\csromanspacer{} เอวํ รเส พหิทฺธา ววตฺเถติ.\csromanspacer{}\csromanspacer{}กถํ โผฏฺฐพฺเพ พหิทฺธา ววตฺเถติ? โผฏฺฐพฺพา อวิชฺชาสมฺภูตาติ ววตฺเถติ, โผฏฺฐพฺพา ตณฺหาสมฺภูตาติ ววตฺเถติ, โผฏฺฐพฺพา กมฺมสมฺภูตาติ ววตฺเถติ, โผฏฺฐพฺพา อาหารสมฺภูตาติ ววตฺเถติ, โผฏฺฐพฺพา อุปฺปนฺนาติ ววตฺเถติ.\csromanspacer{} โผฏฺฐพฺพา สมุทาคตาติ ววตฺเถติ ฯเปฯ.\csromanspacer{} เอวํ โผฏฺฐพฺเพ พหิทฺธา ววตฺเถติ.}

\prose{กถํ ธมฺเม พหิทฺธา ววตฺเถติ? ธมฺมา อวิชฺชาสมฺภูตาติ ววตฺเถติ, ธมฺมา ตณฺหาสมฺภูตาติ ววตฺเถติ, ธมฺมา กมฺมสมฺภูตาติ ววตฺเถติ, ธมฺมา อาหารสมฺภูตาติ ววตฺเถติ, ธมฺมา อุปฺปนฺนาติ ววตฺเถติ, ธมฺมา สมุทาคตาติ ววตฺเถติ, ธมฺมา อหุตฺวา สมฺภูตา, หุตฺวา น ภวิสฺสนฺตีติ ววตฺเถติ, ธมฺเม อนฺตวนฺตโต ววตฺเถติ, ธมฺมา อทฺธุวา อสสฺสตา วิปริณามธมฺมาติ ววตฺเถติ, ธมฺมา อนิจฺจา สงฺขตา ปฏิจฺจสมุปฺปนฺนา ขยธมฺมา วยธมฺมา วิราคธมฺมา นิโรธธมฺมาติ ววตฺเถติ, ธมฺเม อนิจฺจโต ววตฺเถติ, โน นิจฺจโต, ทุกฺขโต ววตฺเถติ, โน สุขโต, อนตฺตโต ววตฺเถติ, โน อตฺตโต, นิพฺพินฺทติ, โน นนฺทติ, วิรชฺชติ, โน รชฺชติ, นิโรเธติ, โน สมุเทติ, ปฏินิสฺสชฺชติ, โน อาทิยติ.\csromanspacer{}\csromanspacer{}อนิจฺจโต ววตฺเถนฺโต นิจฺจสญฺญํ ปชหาติ, ทุกฺขโต ววตฺเถนฺโต สุขสญฺญํ ปชหาติ, อนตฺตโต ววตฺเถนฺโต อตฺตสญฺญํ ปชหาติ, นิพฺพินฺทนฺโต นนฺทิํ ปชหาติ, วิรชฺชนฺโต ราคํ ปชหาติ, นิโรเธนฺโต สมุทยํ ปชหาติ, ปฏินิสฺสชฺชนฺโต อาทานํ ปชหาติ.\csromanspacer{} เอวํ ธมฺเม พหิทฺธา ววตฺเถติ.\csromanspacer{} ตํญาตฏฺเฐน ญาณํ, ปชานนฏฺเฐน ปญฺญา.\csromanspacer{} เตน วุจฺจติ พหิทฺธา ววตฺถาเน ปญฺญา โคจรนานตฺเต ญาณํ.}

\vspace{\tipitakaparskip}
\csromancenter{โคจรนานตฺตญาณนิทฺเทโส โสฬสโม.}
\csromanmatikaanchor{mtk.19}
\csromantocmarkref{subsection}{17. จริยานานตฺตญาณนิทฺเทส}{mtk.19}
\bookmark[dest={mtk.19},level=2]{17. จริยานานตฺตญาณนิทฺเทส}
\csromanheader{2}{17. จริยานานตฺตญาณนิทฺเทส}
\proseitem{68}{\csromanfolio{77}กถํ \textbf{จริยา}ววตฺถาเน ปญฺญา \textbf{จริยา}นานตฺเต ญาณํ\csromandash{}\textbf{จริยา}ติ ติสฺโส \textbf{จริยา}โย วิญฺญาณ\textbf{จริยา}, อญฺญาณ\textbf{จริยา}, ญาณ\textbf{จริยา}.}

\prose{กตมา วิญฺญาณ\textbf{จริยา}? ทสฺสนตฺถาย อาวชฺชนกิริยาพฺยากตา วิญฺญาณ\textbf{จริยา} รูเปสุ, ทสฺสนฏฺโฐ จกฺขุวิญฺญาณํ วิญฺญาณ\textbf{จริยา} รูเปสุ, ทิฏฺฐตฺตา อภินิโรปนา วิปากมโนธาตุ\footnote{อภินิโรปนวิปากมโนธาตุ (สฺยา)} วิญฺญาณ\textbf{จริยา} รูเปสุ, อภินิโรปิตตฺตา วิปากมโนวิญฺญาณธาตุ วิญฺญาณ\textbf{จริยา} รูเปสุ.\csromanspacer{}\csromanspacer{}สวนตฺถาย อาวชฺชนกิริยาพฺยากตา วิญฺญาณ\textbf{จริยา} สทฺเทสุ, สวนฏฺโฐ โสตวิญฺญาณํ วิญฺญาณ\textbf{จริยา} สทฺเทสุ.\csromanspacer{} สุตตฺตา อภินิโรปนา วิปากมโนธาตุ วิญฺญาณ\textbf{จริยา} สทฺเทสุ, อภินิโรปิตตฺตา วิปากมโนวิญฺญาณธาตุ วิญฺญาณ\textbf{จริยา} สทฺเทสุ.\csromanspacer{}\csromanspacer{}ฆายนตฺถาย อาวชฺชนกิริยาพฺยากตา วิญฺญาณ\textbf{จริยา} คนฺเธสุ, ฆายนฏฺโฐ\footnote{ฆายนตฺโถ (สฺยา, ก)} ฆานวิญฺญาณํ วิญฺญาณ\textbf{จริยา} คนฺเธสุ, ฆายิตตฺตา อภินิโรปนา วิปากมโนธาตุ วิญฺญาณ\textbf{จริยา} คนฺเธสุ, อภินิโรปิตตฺตา วิปากมโนวิญฺญาณธาตุ วิญฺญาณ\textbf{จริยา} คนฺเธสุ.\csromanspacer{}\csromanspacer{}สายนตฺถาย อาวชฺชนกิริยาพฺยากตา วิญฺญาณ\textbf{จริยา} รเสสุ, สายนฏฺโฐ ชิวฺหาวิญฺญาณํ วิญฺญาณ\textbf{จริยา} รเสสุ, สายิตตฺตา อภินิโรปนา วิปากมโนธาตุ วิญฺญาณ\textbf{จริยา} รเสสุ, อภินิโรปิตตฺตา วิปากมโนวิญฺญาณธาตุ วิญฺญาณ\textbf{จริยา} รเสสุ.\csromanspacer{}\csromanspacer{}ผุสนตฺถาย อาวชฺชนกิริยาพฺยากตา วิญฺญาณ\textbf{จริยา} โผฏฺฐพฺเพสุ, ผุสนฏฺโฐ กายวิญฺญาณํ วิญฺญาณ\textbf{จริยา} โผฏฺฐพฺเพสุ, ผุฏฺฐตฺตา อภินิโรปนา วิปากมโนธาตุ วิญฺญาณ\textbf{จริยา} โผฏฺฐพฺเพสุ, อภินิโรปิตตฺตา วิปากมโนวิญฺญาณธาตุ วิญฺญาณ\textbf{จริยา} โผฏฺฐพฺเพสุ.\csromanspacer{}\csromanspacer{}วิชานนตฺถาย อาวชฺชนกิริยาพฺยากตา วิญฺญาณ\textbf{จริยา} ธมฺเมสุ, วิชานนฏฺโฐ มโนวิญฺญาณํ วิญฺญาณ\textbf{จริยา} ธมฺเมสุ.\csromanspacer{} วิญฺญาตตฺตา อภินิโรปนา วิปากมโนธาตุ วิญฺญาณ\textbf{จริยา} ธมฺเมสุ, อภินิโรปิตตฺตา วิปากมโนวิญฺญาณธาตุ วิญฺญาณ\textbf{จริยา} ธมฺเมสุ.}

\proseitem{69}{\textbf{วิญฺญาณจริยา}ติ เกนฏฺเฐน วิญฺญาณ\textbf{จริยา}? นีราคา จรตีติ วิญฺญาณ\textbf{จริยา}, นิทฺโทสา จรตีติ วิญฺญาณ\textbf{จริยา}, นิมฺโมหา จรตีติ วิญฺญาณ\textbf{จริยา}, นิมฺมานา จรตีติ วิญฺญาณ\textbf{จริยา}, นิทฺทิฏฺฐิ จรตีติ \csromanfolio{78}วิญฺญาณจริยา, นิ-อุทฺธจฺจา จรตีติ วิญฺญาณจริยา.\csromanspacer{} นิพฺพิจิกิจฺฉา จรตีติ วิญฺญาณจริยา, นานุสยา จรตีติ วิญฺญาณจริยา.\csromanspacer{} ราควิปฺปยุตฺตา จรตีติ วิญฺญาณจริยา, โทสวิปฺปยุตฺตา จรตีติ วิญฺญาณจริยา, โมหวิปฺปยุตฺตา จรตีติ วิญฺญาณจริยา, มานวิปฺปยุตฺตา จรตีติ วิญฺญาณจริยา, ทิฏฺฐิวิปฺปยุตฺตา จรตีติ วิญฺญาณจริยา, อุทฺธจฺจวิปฺปยุตฺตา จรตีติ วิญฺญาณจริยา, วิจิกิจฺฉาวิปฺปยุตฺตา จรตีติ วิญฺญาณจริยา, อนุสยวิปฺปยุตฺตา จรตีติ วิญฺญาณจริยา, กุสเลหิ กมฺเมหิ สมฺปยุตฺตา จรตีติ วิญฺญาณจริยา, อกุสเลหิ กมฺเมหิ วิปฺปยุตฺตา จรตีติ วิญฺญาณจริยา, สาวชฺเชหิ กมฺเมหิ วิปฺปยุตฺตา จรตีติ วิญฺญาณจริยา, อนวชฺเชหิ กมฺเมหิ สมฺปยุตฺตา จรตีติ วิญฺญาณจริยา, กเณฺหหิ กมฺเมหิ วิปฺปยุตฺตา จรตีติ วิญฺญาณจริยา, สุกฺเกหิ กมฺเมหิ สมฺปยุตฺตา จรตีติ วิญฺญาณจริยา, สุขุทฺรเยหิ กมฺเมหิ สมฺปยุตฺตา จรตีติ วิญฺญาณจริยา, ทุกฺขุทฺรเยหิ กมฺเมหิ วิปฺปยุตฺตา จรตีติ วิญฺญาณจริยา, สุขวิปาเกหิ กมฺเมหิ สมฺปยุตฺตา จรตีติ วิญฺญาณจริยา, ทุกฺขวิปาเกหิ กมฺเมหิ วิปฺปยุตฺตา จรตีติ วิญฺญาณจริยา, วิญฺญาเต จรตีติ วิญฺญาณจริยา, วิญฺญาณสฺส เอวรูปา จริยา โหตีติ วิญฺญาณจริยา, ปกติปริสุทฺธมิทํ จิตฺตํ นิกฺกิเลสฏฺเฐนาติ วิญฺญาณจริยา.\csromanspacer{} อยํ วิญฺญาณจริยา.}

\prose{กตมา อญฺญาณจริยา? มนาปิเยสุ รูเปสุ ราคสฺส ชวนตฺถาย อาวชฺชนกิริยาพฺยากตา วิญฺญาณจริยา, ราคสฺส ชวนา อญฺญาณจริยา.\csromanspacer{} อมนาปิเยสุ รูเปสุ โทสสฺส ชวนตฺถาย อาวชฺชนกิริยาพฺยากตา วิญฺญาณจริยา, โทสสฺส ชวนา อญฺญาณจริยา.\csromanspacer{} ตทุภเยน อสมเปกฺขนสฺมิํ วตฺถุสฺมิํ โมหสฺส ชวนตฺถาย อาวชฺชนกิริยาพฺยากตา วิญฺญาณจริยา, โมหสฺส ชวนา อญฺญาณจริยา, วินิพนฺธสฺส มานสฺส ชวนตฺถาย อาวชฺชนกิริยาพฺยากตา วิญฺญาณจริยา, มานสฺส ชวนา อญฺญาณจริยา.\csromanspacer{} ปรามฏฺฐาย ทิฏฺฐิยา ชวนตฺถาย อาวชฺชนกิริยาพฺยากตา วิญฺญาณจริยา, ทิฏฺฐิยา ชวนา อญฺญาณจริยา.\csromanspacer{} วิกฺเขปคตสฺส อุทฺธจฺจสฺส ชวนตฺถาย อาวชฺชนกิริยาพฺยากตา วิญฺญาณจริยา, อุทฺธจฺจสฺส ชวนา อญฺญาณจริยา.\csromanspacer{} อนิฏฺฐงฺคตาย วิจิกิจฺฉาย ชวนตฺถาย อาวชฺชนกิริยาพฺยากตา วิญฺญาณจริยา, วิจิกิจฺฉาย ชวนา อญฺญาณจริยา.\csromanspacer{} ถามคตสฺส อนุสยสฺส ชวนตฺถาย อาวชฺชนกิริยาพฺยากตา วิญฺญาณจริยา, อนุสยสฺส ชวนา อญฺญาณจริยา.}

\prose{\csromanfolio{79}มนาปิเยสุ สทฺเทสุ ฯเปฯ.\csromanspacer{} มนาปิเยสุ คนฺเธสุ ฯเปฯ.\csromanspacer{} มนาปิเยสุ รเสสุ ฯเปฯ.\csromanspacer{} มนาปิเยสุ โผฏฺฐพฺเพสุ ฯเปฯ.\csromanspacer{} มนาปิเยสุ ธมฺเมสุ ราคสฺส ชวนตฺถาย อาวชฺชนกิริยาพฺยากตา วิญฺญาณจริยา, ราคสฺส ชวนา อญฺญาณจริยา.\csromanspacer{} อมนาปิเยสุ ธมฺเมสุ โทสสฺส ชวนตฺถาย อาวชฺชนกิริยาพฺยากตา วิญฺญาณจริยา, โทสสฺส ชวนา อญฺญาณจริยา.\csromanspacer{} ตทุภเยน อสมเปกฺขนสฺมิํ วตฺถุสฺมิํ โมหสฺส ชวนตฺถาย อาวชฺชนกิริยาพฺยากตา วิญฺญาณจริยา, โมหสฺส ชวนา อญฺญาณจริยา.\csromanspacer{} วินิพนฺธสฺส มานสฺส ชวนตฺถาย อาวชฺชนกิริยาพฺยากตา วิญฺญาณจริยา, มานสฺส ชวนา อญฺญาณจริยา.\csromanspacer{} ปรามฏฺฐาย ทิฏฺฐิยา ชวนตฺถาย อาวชฺชนกิริยาพฺยากตา วิญฺญาณจริยา, ทิฏฺฐิยา ชวนา อญฺญาณจริยา.\csromanspacer{} วิกฺเขปคตสฺส อุทฺธจฺจสฺส ชวนตฺถาย อาวชฺชนกิริยาพฺยากตา วิญฺญาณจริยา, อุทฺธจฺจสฺส ชวนา อญฺญาณจริยา.\csromanspacer{} อนิฏฺฐงฺคตาย วิจิกิจฺฉาย ชวนตฺถาย อาวชฺชนกิริยาพฺยากตา วิญฺญาณจริยา, วิจิกิจฺฉาย ชวนา อญฺญาณจริยา.\csromanspacer{} ถามคตสฺส อนุสยสฺส ชวนตฺถาย อาวชฺชนกิริยาพฺยากตา วิญฺญาณจริยา, อนุสยสฺส ชวนา อญฺญาณจริยา.}

\proseitem{70}{\textbf{อญฺญาณจริยา}ติ เกนฏฺเฐน อญฺญาณจริยา? สราคา จรตีติ อญฺญาณจริยา, สโทสา จรตีติ อญฺญาณจริยา, สโมหา จรตีติ อญฺญาณจริยา, สมานา จรตีติ อญฺญาณจริยา, สทิฏฺฐิ จรตีติ อญฺญาณจริยา, ส-อุทฺธจฺจา จรตีติ อญฺญาณจริยา, สวิจิกิจฺฉา จรตีติ อญฺญาณจริยา, สานุสยา จรตีติ อญฺญาณจริยา.\csromanspacer{} ราคสมฺปยุตฺตา จรตีติ อญฺญาณจริยา, โทสสมฺปยุตฺตา จรตีติ อญฺญาณจริยา, โมหสมฺปยุตฺตา จรตีติ อญฺญาณจริยา, มานสมฺปยุตฺตา จรตีติ อญฺญาณจริยา, ทิฏฺฐิสมฺปยุตฺตา จรตีติ อญฺญาณจริยา, อุทฺธจฺจสมฺปยุตฺตา จรตีติ อญฺญาณจริยา.\csromanspacer{} วิจิกิจฺฉาสมฺปยุตฺตา จรตีติ อญฺญาณจริยา, อนุสยสมฺปยุตฺตา จรตีติ อญฺญาณจริยา.\csromanspacer{} กุสเลหิ กมฺเมหิ วิปฺปยุตฺตา จรตีติ อญฺญาณจริยา, อกุสเลหิ กมฺเมหิ สมฺปยุตฺตา จรตีติ อญฺญาณจริยา, สาวชฺเชหิ กมฺเมหิ สมฺปยุตฺตา จรตีติ อญฺญาณจริยา, อนวชฺเชหิ กมฺเมหิ วิปฺปยุตฺตา จรตีติ อญฺญาณจริยา, กเณฺหหิ กมฺเมหิ สมฺปยุตฺตา จรตีติ อญฺญาณจริยา, สุกฺเกหิ กมฺเมหิ วิปฺปยุตฺตา จรตีติ อญฺญาณจริยา, สุขุทฺรเยหิ กมฺเมหิ วิปฺปยุตฺตา จรตีติ อญฺญาณจริยา, ทุกฺขุทฺรเยหิ \csromanfolio{80}กมฺเมหิ สมฺปยุตฺตา จรตีติ อญฺญาณจริยา, สุขวิปาเกหิ กมฺเมหิ วิปฺปยุตฺตา จรตีติ อญฺญาณจริยา, ทุกฺขวิปาเกหิ กมฺเมหิ สมฺปยุตฺตา จรตีติ อญฺญาณจริยา, อญฺญาเต จรตีติ อญฺญาณจริยา, อญฺญาณสฺส เอวรูปา จริยา โหตีติ อญฺญาณจริยา.\csromanspacer{} อยํ อญฺญาณจริยา.}

\proseitem{71}{กตมา ญาณจริยา? อนิจฺจานุปสฺสนตฺถาย อาวชฺชนกิริยาพฺยากตา วิญฺญาณจริยา, อนิจฺจานุปสฺสนา ญาณจริยา.\csromanspacer{} ทุกฺขานุปสฺสนตฺถาย อาวชฺชนกิริยาพฺยากตา วิญฺญาณจริยา, ทุกฺขานุปสฺสนา ญาณจริยา.\csromanspacer{} อนตฺตานุปสฺสนตฺถาย อาวชฺชนกิริยาพฺยากตา วิญฺญาณจริยา, อนตฺตานุปสฺสนา ญาณจริยา.\csromanspacer{} นิพฺพิทานุปสฺสนตฺถาย ฯเปฯ.\csromanspacer{} วิราคานุปสฺสนตฺถาย.\csromanspacer{} นิโรธานุปสฺสนตฺถาย.\csromanspacer{} ปฏินิสฺสคฺคานุปสฺสนตฺถาย.\csromanspacer{} ขยานุปสฺสนตฺถาย.\csromanspacer{} วยานุปสฺสนตฺถาย.\csromanspacer{} วิปริณามานุปสฺสนตฺถาย.\csromanspacer{} อนิมิตฺตานุปสฺสนตฺถาย.\csromanspacer{} อปฺปณิหิตานุปสฺสนตฺถาย.\csromanspacer{} สุญฺญตานุปสฺสนตฺถาย.\csromanspacer{} อธิปญฺญาธมฺมานุปสฺสนตฺถาย.\csromanspacer{} ยถาภูตญาณทสฺสนตฺถาย.\csromanspacer{} อาทีนวานุปสฺสนตฺถาย.\csromanspacer{} ปฏิสงฺขานุปสฺสนตฺถาย.\csromanspacer{} อาวชฺชนกิริยาพฺยากตา วิญฺญาณจริยา, ปฏิสงฺขานุปสฺสนา ญาณจริยา.\csromanspacer{} วิวฏฺฏนานุปสฺสนา ญาณจริยา.\csromanspacer{} โสตาปตฺติมคฺโค ญาณจริยา.\csromanspacer{} โสตาปตฺติผลสมาปตฺติ ญาณจริยา.\csromanspacer{} สกทาคามิมคฺโค ญาณจริยา.\csromanspacer{} สกทาคามิผลสมาปตฺติ ญาณจริยา.\csromanspacer{} อนาคามิมคฺโค ญาณจริยา.\csromanspacer{} อนาคามิผลสมาปตฺติ ญาณจริยา.\csromanspacer{} อรหตฺตมคฺโค ญาณจริยา.\csromanspacer{} อรหตฺตผลสมาปตฺติ ญาณจริยา.}

\prose{\textbf{ญาณจริยา}ติ เกนฏฺเฐน ญาณจริยา? นีราคา จรตีติ ญาณจริยา นิทฺโทสา จรตีติ ญาณจริยา ฯเปฯ นานุสยา จรตีติ ญาณจริยา, ราควิปฺปยุตฺตา จรตีติ ญาณจริยา, โทสวิปฺปยุตฺตา จรตีติ ญาณจริยา, โมหวิปฺปยุตฺตา จรตีติ ญาณจริยา, มานวิปฺปยุตฺตา ฯเปฯ ทิฏฺฐิวิปฺปยุตฺตา.\csromanspacer{} อุทฺธจฺจวิปฺปยุตฺตา.\csromanspacer{} วิจิกิจฺฉาวิปฺปยุตฺตา.\csromanspacer{} อนุสยวิปฺปยุตฺตา.\csromanspacer{} กุสเลหิ กมฺเมหิ สมฺปยุตฺตา.\csromanspacer{} อกุสเลหิ กมฺเมหิ วิปฺปยุตฺตา.\csromanspacer{} สาวชฺเชหิ กมฺเมหิ วิปฺปยุตฺตา.\csromanspacer{} อนวชฺเชหิ กมฺเมหิ สมฺปยุตฺตา.\csromanspacer{} กเณฺหหิ กมฺเมหิ วิปฺปยุตฺตา.\csromanspacer{} สุกฺเกหิ กมฺเมหิ สมฺปยุตฺตา.\csromanspacer{} สุขุทฺรเยหิ กมฺเมหิ สมฺปยุตฺตา.\csromanspacer{} ทุกฺขุทฺรเยหิ กมฺเมหิ วิปฺปยุตฺตา.\csromanspacer{} สุขวิปาเกหิ กมฺเมหิ สมฺปยุตฺตา จรตีติ ญาณจริยา, ทุกฺขวิปาเกหิ กมฺเมหิ วิปฺปยุตฺตา \csromanfolio{81}จรตีติ ญาณจริยา, ญาเต จรตีติ ญาณจริยา, ญาณสฺส เอวรูปา จริยา โหตีติ ญาณจริยา, อยํ ญาณจริยา, อญฺญา วิญฺญาณจริยา, อญฺญา อญฺญาณจริยา, อญฺญา ญาณจริยาติ.\csromanspacer{} ตํญาตฏฺเฐน ญาณํ, ปชานนฏฺเฐน ปญฺญา.\csromanspacer{} เตน วุจฺจติ จริยาววตฺถาเน ปญฺญา จริยานานตฺเต ญาณํ.}

\csromanheader{2}{จริยานานตฺตญาณนิทฺเทโส สตฺตรสโม.}
\csromansectionrule
\csromanmatikaanchor{mtk.20}
\csromantocmarkref{subsection}{18. ภูมินานตฺตญาณนิทฺเทส}{mtk.20}
\bookmark[dest={mtk.20},level=2]{18. ภูมินานตฺตญาณนิทฺเทส}
\csromanheader{2}{18. ภูมินานตฺตญาณนิทฺเทส}
\proseitem{72}{กถํ จตุธมฺมววตฺถาเน ปญฺญา ภูมินานตฺเต ญาณํ\csromandash{} จตสฺโส ภูมิโย กามาวจรา ภูมิ, รูปาวจรา ภูมิ, อรูปาวจรา ภูมิ, อปริยาปนฺนา ภูมิ.\csromanspacer{} กตมา กามาวจรา ภูมิ? เหฏฺฐโต อวีจินิรยํ ปริยนฺตํ กริตฺวา อุปริโต ปรนิมฺมิตวสวตฺตี เทเว\footnote{ปรนิมฺมิตวสวตฺติเทเว (ก)} อนฺโต กริตฺวา ยํ เอตสฺมิํ อนฺตเร เอตฺถาวจรา เอตฺถ ปริยาปนฺนา ขนฺธธาตุ-อายตนา รูปํ เวทนา สญฺญา สงฺขารา วิญฺญาณํ.\csromanspacer{} อยํ กามาวจรา ภูมิ.}

\prose{กตมา รูปาวจรา ภูมิ? เหฏฺฐโต พฺรหฺมโลกํ ปริยนฺตํ กริตฺวา อุปริโต อกนิฏฺเฐ เทเว อนฺโต กริตฺวา ยํ เอตสฺมิํ อนฺตเร เอตฺถาวจรา เอตฺถ ปริยาปนฺนา สมาปนฺนสฺส วา อุปปนฺนสฺส วา ทิฏฺฐธมฺมสุขวิหาริสฺส วา จิตฺตเจตสิกา ธมฺมา.\csromanspacer{} อยํ รูปาวจรา ภูมิ.}

\prose{กตมา อรูปาวจรา ภูมิ? เหฏฺฐโต อากาสานญฺจายตนูปเค เทเว ปริยนฺตํ กริตฺวา อุปริโต เนวสญฺญานาสญฺญายตนูปเค เทเว อนฺโต กริตฺวา ยํ เอตสฺมิํ อนฺตเร เอตฺถาวจรา เอตฺถ ปริยาปนฺนา สมาปนฺนสฺส วา อุปปนฺนสฺส วา ทิฏฺฐธมฺมสุขวิหาริสฺส วา จิตฺตเจตสิกา ธมฺมา.\csromanspacer{} อยํ อรูปาวจรา ภูมิ.}

\prose{กตมา อปริยาปนฺนา ภูมิ? อปริยาปนฺนา มคฺคา จ มคฺคผลานิ จ อสงฺขตา จ ธาตุ.\csromanspacer{} อยํ อปริยาปนฺนา ภูมิ.\csromanspacer{} อิมา จตสฺโส ภูมิโย.}

\prose{อปราปิ จตสฺโส ภูมิโย จตฺตาโร สติปฏฺฐานา, จตฺตาโร สมฺมปฺปธานา, จตฺตาโร อิทฺธิปาทา, จตฺตาริ ฌานานิ, จตสฺโส อปฺปมญฺญาโย, จตสฺโส อรูปสมาปตฺติโย, จตสฺโส ปฏิสมฺภิทา, \csromanfolio{82}จตสฺโส ปฏิปทา, จตฺตาริ อารมฺมณานิ, จตฺตาโร อริยวํสา, จตฺตาริ สงฺคหวตฺถูนิ, จตฺตาริ จกฺกานิ, จตฺตาริ ธมฺมปทานิ.\csromanspacer{} อิมา จตสฺโส ภูมิโย.\csromanspacer{} ตํญาตฏฺเฐน ญาณํ, ปชานนฏฺเฐน ปญฺญา.\csromanspacer{} เตน วุจฺจติ จตุธมฺมววตฺถาเน ปญฺญา ภูมินานตฺเต ญาณํ.}

\csromanheader{2}{ภูมินานตฺตญาณนิทฺเทโส อฏฺฐารสโม.}
\csromansectionrule
\csromanmatikaanchor{mtk.21}
\csromantocmarkref{subsection}{19. ธมฺมนานตฺตญาณนิทฺเทส}{mtk.21}
\bookmark[dest={mtk.21},level=2]{19. ธมฺมนานตฺตญาณนิทฺเทส}
\csromanheader{2}{19. ธมฺมนานตฺตญาณนิทฺเทส}
\proseitem{73}{กถํ นวธมฺมววตฺถาเน ปญฺญา ธมฺมนานตฺเต ญาณํ\csromandash{} กถํ ธมฺเม ววตฺเถติ? กามาวจเร ธมฺเม กุสลโต ววตฺเถติ, อกุสลโต ววตฺเถติ, อพฺยากตโต ววตฺเถติ.\csromanspacer{} รูปาวจเร ธมฺเม กุสลโต ววตฺเถติ, อพฺยากตโต ววตฺเถติ.\csromanspacer{} อรูปาวจเร ธมฺเม กุสลโต ววตฺเถติ, อพฺยากตโต ววตฺเถติ.\csromanspacer{} อปริยาปนฺเน ธมฺเม กุสลโต ววตฺเถติ, อพฺยากตโต ววตฺเถติ.}

\prose{กถํ กามาวจเร ธมฺเม กุสลโต ววตฺเถติ, อกุสลโต ววตฺเถติ, อพฺยากตโต ววตฺเถติ? ทส กุสลกมฺมปเถ กุสลโต ววตฺเถติ.\csromanspacer{} ทส อกุสลกมฺมปเถ อกุสลโต ววตฺเถติ, รูปญฺจ วิปากญฺจ กิริยญฺจ อพฺยากตโต ววตฺเถติ.\csromanspacer{} เอวํ กามาวจเร ธมฺเม กุสลโต ววตฺเถติ อกุสลโต ววตฺเถติ อพฺยากตโต ววตฺเถติ.}

\prose{กถํ รูปาวจเร ธมฺเม กุสลโต ววตฺเถติ, อพฺยากตโต ววตฺเถติ? อิธฏฺฐสฺส จตฺตาริ ฌานานิ กุสลโต ววตฺเถติ, ตตฺรูปปนฺนสฺส จตฺตาริ ฌานานิ อพฺยากตโต ววตฺเถติ.\csromanspacer{} เอวํ รูปาวจเร ธมฺเม กุสลโต ววตฺเถติ, อพฺยากตโต ววเถติ.}

\prose{กถํ อรูปาวจเร ธมฺเม กุสลโต ววตฺเถติ, อพฺยากตโต ววตฺเถติ? อิธฏฺฐสฺส จตสฺโส อรูปาวจรสมาปตฺติโย กุสลโต ววตฺเถติ, ตตฺรูปปนฺนสฺส จตสฺโส อรูปาวจรสมาปตฺติโย อพฺยากตโต ววตฺเถติ.\csromanspacer{} เอวํ อรูปาวจเร ธมฺเม กุสลโต ววตฺเถติ อพฺยากตโต ววตฺเถติ.}

\prose{\csromanfolio{83}กถํ อปริยาปนฺเน ธมฺเม กุสลโต ววตฺเถติ, อพฺยากตโต ววตฺเถติ? จตฺตาโร อริยมคฺเค กุสลโต ววตฺเถติ, จตฺตาริ จ สามญฺญผลานิ นิพฺพานญฺจ อพฺยากตโต ววตฺเถติ.\csromanspacer{} เอวํ อปริยาปนฺเน ธมฺเม กุสลโต ววตฺเตติ อพฺยากตโต ววตฺเถติ.\csromanspacer{} เอวํ ธมฺเม ววตฺเถติ.}

\prose{นว ปาโมชฺชมูลกา ธมฺมา, อนิจฺจโต มนสิกโรโต ปาโมชฺชํ ชายติ, ปมุทิตสฺส ปีติ ชายติ, ปีติมนสฺส กาโย ปสฺสมฺภติ, ปสฺสทฺธกาโย ทุขํ เวเทติ, สุขิโน จิตฺตํ สมาธิยติ, สมาหิเต จิตฺเต ยถาภูตํ ปชานาติ ปสฺสติ, ยถาภูตํ ชานํ ปสฺสํ นิพฺพินฺทติ.\csromanspacer{} นิพฺพินฺทํ วิรชฺชติ, วิราคา วิมุจฺจติ.\csromanspacer{} ทุกฺขโต มนสิกโรโต ปาโมชฺชํ ชายติ ฯเปฯ.\csromanspacer{} อนตฺตโต มนสิกโรโต ปาโมชฺชํ ชายติ ฯเปฯ.}

\prose{รูปํ อนิจฺจโต มนสิกโรโต ปาโมชฺชํ ชายติ ฯเปฯ.\csromanspacer{} รูปํ ทุกฺขโต มนสิกโรโต ฯเปฯ.\csromanspacer{} เวทนํ.\csromanspacer{} สญฺญํ.\csromanspacer{} สงฺขาเร.\csromanspacer{} วิญฺญาณํ.\csromanspacer{} จกฺขุํ ฯเปฯ.\csromanspacer{} ชรามรณํ อนิจฺจโต มนสิกโรโต ปาโมชฺชํ ชายติ ฯเปฯ.\csromanspacer{} ชรามรณํ ทุกฺขโต มนสิกโรโต ปาโมชฺชํ ชายติ ฯเปฯ.\csromanspacer{} ชรามรณํ อนตฺตโต มนสิกโรโต ปาโมชฺชํ ชายติ, ปมุทิตสฺส ปีติ ชายติ, ปีติมนสฺส กาโย ปสฺสมฺภติ, ปสฺสทฺธกาโย สุขํ เวเทติ, สุขิโน จิตฺตํ สมาธิยติ, สมาหิเต จิตฺเต ยถาภูตํ ปชานาติ ปสฺสติ, ยถาภูตํ ชานํ ปสฺสํ นิพฺพินฺทติ, นิพฺพินฺทํ วิรชฺชติ, วิราคา วิมุจฺจติ, อิเม นว ปาโมชฺชมูลกา ธมฺมา.}

\proseitem{74}{นว โยนิโสมนสิการมูลกา ธมฺมา, อนิจฺจโต โยนิโส มนสิกโรโต ปาโมชฺชํ ชายติ, ปมุทิตสฺส ปีติ ชายติ, ปีติมนสฺส กาโย ปสฺสมฺภติ, ปสฺสทฺธกาโย สุขํ เวเทติ, สุขิโน จิตฺตํ สมาธิยติ, สมาหิเตน จิตฺเตน “อิทํ ทุกฺขนฺ”ติ ยถาภูตํ ปชานาติ, “อยํ ทุกฺขสมุทโย”ติ ยถาภูตํ ปชานาติ, “อยํ ทุกฺขนิโรโธ”ติ ยยาภูตํ ปชานาติ, “อยํ ทุกฺขนิโรธคามินี ปฏิปทา”ติ ยถาภูตํ ปชานาติ.\csromanspacer{}\csromanspacer{}ทุกฺขโต โยนิโส มนสิกโรโต ปาโมชฺชํ ชายติ, ปมุทิตสฺส ปีติ ชายติ, ปีติมนสฺส กาโย ปสฺสมฺภติ, ปสฺสทฺธกาโย สุขํ เวเทติ, สุขิโน จิตฺตํ สมาธิยติ, สมาหิเตน จิตฺเตน “อิทํ ทุกฺขนฺ”ติ ยถาภูตํ ปชานาติ, “อยํ ทุกฺขสมุทโย”ติ ยถาภูตํ ปชานาติ, “อยํ ทุกฺขนิโรโธ”ติ ยถาภูตํ ปชานาติ, \csromanfolio{84}“อยํ ทุกฺขนิโรธคามินี ปฏิปทา”ติ ยถาภูตํ ปชานาติ.\csromanspacer{} อนตฺตโต โยนิโส มนสิกโรโต ปาโมชฺชํ ชายติ ฯเปฯ.}

\prose{รูปํ อนิจฺจโต โยนิโส มนสิกโรโต ปาโมชฺชํ ชายติ ฯเปฯ.\csromanspacer{} รูปํ ทุกฺขโต โยนิโส มนสิกโรโต ปาโมชฺชํ ชายติ ฯเปฯ.\csromanspacer{} รูปํ อนตฺตโต โยนิโส มนสิกโรโต ปาโมชฺชํ ชายติ ฯเปฯ.\csromanspacer{} เวทนํ.\csromanspacer{} สญฺญํ.\csromanspacer{} สงฺขาเร.\csromanspacer{} วิญฺญาณํ.\csromanspacer{} จกฺขุํ ฯเปฯ.\csromanspacer{} ชรามรณํ อนิจฺจโต โยนิโส มนสิกโรโต ปาโมชฺชํ ชายติ ฯเปฯ.\csromanspacer{} ชรามรณํ ทุกฺขโต โยนิโส มนสิกโรโต ปาโมชฺชํ ชายติ ฯเปฯ.\csromanspacer{} ชรามรณํ อนตฺตโต โยนิโส มนสิกโรโต ปาโมชฺชํ ชายติ, ปมุทิตสฺส ปีติ ชายติ, ปีติมนสฺส กาโย ปสฺสมฺภติ, ปสฺสทฺธกาโย สุขํ เวเทติ, สุขิโน จิตฺตํ สมาธิยติ, สมาหิเตน จิตฺเตน “อิทํ ทุกฺขนฺ”ติ ยถาภูตํ ปชานาติ, “อยํ ทุกฺขสมุทโย”ติ ยถาภูตํ ปชานาติ, “อยํ ทุกฺขนิโรโธ”ติ ยถาภูตํ ปชานาติ, “อยํ ทุกฺขนิโรธคามินี ปฏิปทา”ติ ยถาภูตํ ปชานาติ.\csromanspacer{} อิเม นว โยนิโสมนสิการมูลกา ธมฺมา.}

\prose{นว นานตฺตา, ธาตุนานตฺตํ ปฏิจฺจ อุปฺปชฺชติ ผสฺสนานตฺตํ, ผสฺสนานตฺตํ ปฏิจฺจ อุปฺปชฺชติ เวทนานานตฺตํ, เวทนานานตฺตํ ปฏิจฺจ อุปฺปชฺชติ สญฺญานานตฺตํ, สญฺญานานตฺตํ ปฏิจฺจ อุปฺปชฺชติ สงฺกปฺปนานตฺตํ, สงฺกปฺปนานตฺตํ ปฏิจฺจ อุปฺปชฺชติ ฉนฺทนานตฺตํ, ฉนฺทนานตฺตํ ปฏิจฺจ อุปฺปชฺชติ ปริฬาหนานตฺตํ, ปริฬาหนานตฺตํ ปฏิจฺจ อุปฺปชฺชติ ปริเยสนานานตฺตํ, ปริเยสนานานตฺตํ ปฏิจฺจ อุปฺปชฺชติ ลาภนานตฺตํ.\csromanspacer{} อิเม นว นานตฺตา.\csromanspacer{} ตํญาตฏฺเฐน ญาณํ, ปชานนฏฺเฐน ปญฺญา.\csromanspacer{} เตน วุจฺจติ นวธมฺมววตฺถาเน ปญฺญา ธมฺมนานตฺเต ญาณํ.}

\csromanheader{2}{ธมฺมนานตฺตญาณนิทฺเทโส เอกูนวีสติโม.}
\csromansectionrule
\csromanmatikaanchor{mtk.22}
\csromantocmarkref{subsection}{20-24. ญาณปญฺจกนิทฺเทส}{mtk.22}
\bookmark[dest={mtk.22},level=2]{20-24. ญาณปญฺจกนิทฺเทส}
\csromanheader{2}{20-24. ญาณปญฺจกนิทฺเทส}
\proseitem{75}{กถํ อภิญฺญาปญฺญา ญาตฏฺเฐ ญาณํ, ปริญฺญาปญฺญา ตีรณฏฺเฐ ญาณํ, ปหาเน ปญฺญา ปริจฺจาคฏฺเฐ ญาณํ, ภาวนาปญฺญา เอกรสฏฺเฐ ญาณํ, สจฺฉิกิริยาปญฺญา ผสฺสนฏฺเฐ ญาณํ\csromandash{}เย เย ธมฺมา อภิญฺญาตา โหนฺติ, เต เต ธมฺมา ญาตา โหนฺติ.\csromanspacer{} เย เย ธมฺมา ปริญฺญาตา โหนฺติ, \csromanfolio{85}เต เต ธมฺมา ตีริตา โหนฺติ.\csromanspacer{} เย เย ธมฺมา ปหีนา โหนฺติ, เต เต ธมฺมา ปริจฺจตฺตา โหนฺติ.\csromanspacer{} เย เย ธมฺมา ภาวิตา โหนฺติ, เต เต ธมฺมา เอกรสา โหนฺติ.\csromanspacer{} เย เย ธมฺมา สจฺฉิกตา โหนฺติ, เต เต ธมฺมา ผสฺสิตา โหนฺติ.\csromanspacer{} ตํญาตฏฺเฐน ญาณํ, ปชานนฏฺเฐน ปญฺญา.\csromanspacer{} เตน วุจฺจติ อภิญฺญาปญฺญา ญาตฏฺเฐ ญาณํ, ปริญฺญาปญฺญา ตีรณฏฺเฐ ญาณํ.\csromanspacer{} ปหาเน ปญฺญา ปริจฺจาคฏฺเฐ ญาณํ.\csromanspacer{} ภาวนาปญฺญา เอกรสฏฺเฐ ญาณํ.\csromanspacer{} สจฺฉิกิริยาปญฺญา ผสฺสนฏฺเฐ ญาณํ.}

\csromanheader{2}{ญาณปญฺจกนิทฺเทโส จตุวีสติโม.}
\csromansectionrule
\csromanmatikaanchor{mtk.23}
\csromantocmarkref{subsection}{25-28. ปฏิสมฺภิทาญาณนิทฺเทส}{mtk.23}
\bookmark[dest={mtk.23},level=2]{25-28. ปฏิสมฺภิทาญาณนิทฺเทส}
\csromanheader{2}{25-28. ปฏิสมฺภิทาญาณนิทฺเทส}
\proseitem{76}{กถํ อตฺถนานตฺเต ปญฺญา อตฺถปฏิสมฺภิเท ญาณํ, ธมฺมนานตฺเต ปญฺญา ธมฺมปฏิสมฺภิเท ญาณํ, นิรุตฺตินานตฺเต ปญฺญา นิรุตฺติปฏิสมฺภิเท ญาณํ, ปฏิภานนานตฺเต ปญฺญา ปฏิภานปฏิสมฺภิเท ญาณํ\csromandash{}สทฺธินฺทฺริยํ ธมฺโม, วีริยินฺทฺริยํ ธมฺโม, สตินฺทฺริยํ ธมฺโม, สมาธินฺทฺริยํ ธมฺโม, ปญฺญินฺทฺริยํ ธมฺโม.\csromanspacer{} อญฺโญ สทฺธินฺทฺริยํ ธมฺโม, อญฺโญ วีริยินฺทฺริยํ ธมฺโม, อญฺโญ สตินฺทฺริยํ ธมฺโม, อญฺโญ สมาธินฺทฺริยํ ธมฺโม, อญฺโญ ปญฺญินฺทฺริยํ ธมฺโม.\csromanspacer{} เยน ญาเณน อิเม นานา ธมฺมา ญาตา, เตเนว ญาเณน อิเม นานา ธมฺมา ปฏิวิทิตาติ.\csromanspacer{} เตน วุจฺจติ ธมฺมนานตฺเต ปญฺญา ธมฺมปฏิสมฺภิเท ญาณํ. (1)}

\prose{อธิโมกฺขฏฺโฐ อตฺโถ, ปคฺคหฏฺโฐ อตฺโถ, อุปฏฺฐานฏฺโฐ อตฺโถ, อวิกฺเขปฏฺโฐ อตฺโถ, ทสฺสนฏฺโฐ อตฺโถ.\csromanspacer{} อญฺโญ อธิโมกฺขฏฺโฐ อตฺโถ, อญฺโญ ปคฺคหฏฺโฐ อตฺโถ, อญฺโญ อุปฏฺฐานฏฺโฐ อตฺโถ, อญฺโญ อวิกฺเขปฏฺโฐ อตฺโถ, อญฺโญ ทสฺสนฏฺโฐ อตฺโถ.\csromanspacer{} เยน ญาเณน อิเม นานา อตฺถา ญาตา, เตเนว ญาเณน อิเม นานา อตฺถา ปฏิวิทิตาติ.\csromanspacer{} เตน วุจฺจติ อตฺถนานตฺเถ ปญฺญา อตฺถปฏิสมฺภิเท ญาณํ. (2)}

\prose{ปญฺจ ธมฺเม สนฺทสฺเสตุํ พฺยญฺชนนิรุตฺตาภิลาปา, ปญฺจ อตฺเถ สนฺทสฺเสตุํ พฺยญฺชนนิรุตฺตาภิลาปา.\csromanspacer{} อญฺญา ธมฺมนิรุตฺติโย, อญฺญา อตฺถนิรุตฺติโย.\csromanspacer{} เยน ญาเณน อิมา นานา นิรุตฺติโย ญาตา, เตเนว ญาเณน อิมา นานา นิรุตฺติโย ปฏิวิทิตาติ.\csromanspacer{} เตน วุจฺจติ นิรุตฺตินานตฺเต ปญฺญา นิรุตฺติปฏิสมฺภิเท ญาณํ. (3)}

\prose{\csromanfolio{86}ปญฺจสุ ธมฺเมสุ ญาณานิ, ปญฺจสุ อตฺเถสุ ญาณานิ, ทสสุ นิรุตฺตีสุ ญาณานิ.\csromanspacer{} อญฺญานิ ธมฺเมสุ ญาณานิ, อญฺญานิ อตฺเถสุ ญาณานิ, อญฺญานิ นิรุตฺตีสุ ญาณานิ.\csromanspacer{} เยน ญาเณน อิเม นานา ญาณา ญาตา, เตเนว ญาเณน อิเม นานา ญาณา ปฏิวิทิตาติ.\csromanspacer{} เตน วุจฺจติ ปฏิภานนานตฺเต ปญฺญา ปฏิภานปฏิสมฺภิเท ญาณํ. (4)}

\proseitem{77}{สทฺธาพลํ ธมฺโม, วีริยพลํ ธมฺโม, สติพลํ ธมฺโม, สมาธิพลํ ธมฺโม, ปญฺญาพลํ ธมฺโม.\csromanspacer{} อญฺโญ สทฺธาพลํ ธมฺโม, อญฺโญ วีริยพลํ ธมฺโม, อญฺโญ สติพลํ ธมฺโม, อญฺโญ สมาธิพลํ ธมฺโม, อญฺโญ ปญฺญาพลํ ธมฺโม.\csromanspacer{} เยน ญาเณน อิเม นานา ธมฺมา ญาตา, เตเนว ญาเณน อิเม นานา ธมฺมา ปฏิวิทิตาติ.\csromanspacer{} เตน วุจฺจติ ธมฺมนานตฺเต ปญฺญา ธมฺมปฏิสมฺภิเท ญาณํ. (1)}

\prose{อสฺสทฺธิเย อกมฺปิยฏฺโฐ อตฺโถ, โกสชฺเช อกมฺปิยฏฺโฐ อตฺโถ, ปมาเท อกมฺปิยฏฺโฐ อตฺโถ, อุทฺธจฺเจ อกมฺปิยฏฺโฐ อตฺโถ, อวิชฺชาย อกมฺปิยฏฺโฐ อตฺโถ.\csromanspacer{} อญฺโญ อสฺสทฺธิเย อกมฺปิยฏฺโฐ อตฺโถ, อญฺโญ โกสชฺเช อกมฺปิยฏฺโฐ อตฺโถ, อญฺโญ ปมาเท อกมฺปิยฏฺโฐ อตฺโถ, อญฺโญ อุทฺธจฺเจ อกมฺปิยฏฺโฐ อตฺโถ, อญฺโญ อวิชฺชาย อกมฺปิยฏฺโฐ อตฺโถ.\csromanspacer{} เยน ญาเณน อิเม นานา อตฺถา ญาตา, เตเนว ญาเณน อิเม นานา อตฺถา ปฏิวิทิตาติ, เตน วุจฺจติ อตฺถนานตฺเต ปญฺญา อตฺถปฏิสมฺภิเท ญาณํ. (2)}

\prose{ปญฺจ ธมฺเม สนฺทสฺเสตุํ พฺยญฺชนนิรุตฺตาภิลาปา, ปญฺจ อตฺเถ สนฺทสฺเสตุํ พฺยญฺชนนิรุตฺตาภิลาปา.\csromanspacer{} อญฺญา ธมฺมนิรุตฺติโย, อญฺญา อตฺถนิรุตฺติโย.\csromanspacer{} เยน ญาเณน อิมา นานา นิรุตฺติโย ญาตา, เตเนว ญาเณน อิมา นานา นิรุตฺติโย ปฏิวิทิตาติ, เตน วุจฺจติ นิรุตฺตินานตฺเต ปญฺญา นิรุตฺติปฏิสมฺภิเท ญาณํ. (3)}

\prose{ปญฺจสุ\footnote{ปญฺจ (?)} ธมฺเมสุ ญาณานิ, ปญฺจสุ1 อตฺเถสุ ญาณานิ, ทสสุ\footnote{ทส (?)} นิรุตฺตีสุ ญาณานิ.\csromanspacer{} อญฺญานิ ธมฺเมสุ ญาณานิ, อญฺญานิ อตฺเถสุ ญาณานิ, อญฺญานิ นิรุตฺตีสุ ญาณานิ.\csromanspacer{} เยน ญาเณน อิเม นานา ญาณา ญาตา, เตเนว ญาเณน อิเม นานา ญาณา ปฏิวิทิตาติ, เตน วุจฺจติ ปฏิภานนานตฺเต ปญฺญา ปฏิภานปฏิสมฺภิเท ญาณํ. (4)}

\prose{\csromanfolio{87}สติสมฺโพชฺฌงฺโค ธมฺโม, ธมฺมวิจยสมฺโพชฺฌงฺโค ธมฺโม, วีริยสมฺโพชฺฌงฺโค ธมฺโม, ปีติสมฺโพชฺฌงฺโค ธมฺโม, ปสฺสทฺธิสมฺโพชฺฌงฺโค ธมฺโม, สมาธิสมฺโพชฺฌงฺโค ธมฺโม, อุเปกฺขาสมฺโพชฺฌงฺโค ธมฺโม.\csromanspacer{} อญฺโญ สติสมฺโพชฺฌงฺโค ธมฺโม, อญฺโญ ธมฺมวิจยสมฺโพชฺฌงฺโค ธมฺโม, อญฺโญ วีริยสมฺโพชฺฌงฺโค ธมฺโม, อญฺโญ ปีติสมฺโพชฺฌงฺโค ธมฺโม, อญฺโญ ปสฺสทฺธิสมฺโพชฺฌงฺโค ธมฺโม, อญฺโญ สมาธิสมฺโพชฺฌงฺโค ธมฺโม, อญฺโญ อุเปกฺขาสมฺโพชฺฌงฺโค ธมฺโม, เยน ญาเณน อิเม นานา ธมฺมา ญาตา, เตเนว ญาเณน อิเม นานา ธมฺมา ปฏิวิทิตาติ, เตน วุจฺจติ ธมฺมนานตฺเต ปญฺญา ธมฺมปฏิสมฺภิเท ญาณํ. (1)}

\prose{อุปฏฺฐานฏฺโฐ อตฺโถ, ปวิจยฏฺโฐ อตฺโถ, ปคฺคหฏฺโฐ อตฺโถ, ผรณฏฺโฐ อตฺโถ, อุปสมฏฺโฐ อตฺโถ, อวิกฺเขปฏฺโฐ อตฺโถ, ปฏิสงฺขานฏฺโฐ อตฺโถ.\csromanspacer{} อญฺโญ อุปฏฺฐานฏฺโฐ อตฺโถ, อญฺโญ ปวิจยฏฺโฐ อตฺโถ, อญฺโญ ปคฺคหฏฺโฐ อตฺโต, อญฺโญ ผรณฏฺโฐ อตฺโถ, อญฺโญ อุปสมฏฺโฐ อตฺโถ, อญฺโญ อวิกฺเขปฏฺโฐ อตฺโถ, อญฺโญ ปฏิสงฺขานฏฺโฐ อตฺโถ.\csromanspacer{} เยน ญาเณน อิเม นานา อตฺถา ญาตา.\csromanspacer{} เตเนว ญาเณน อิเม นานา อตฺถา ปฏิวิทิตาติ.\csromanspacer{} เตน วุจฺจติ อตฺถนานตฺเต ปญฺญา อตฺถปฏิสมฺภิเท ญาณํ. (2)}

\prose{สตฺต ธมฺเม สนฺทสฺเสตุํ พฺยญฺชนนิรุตฺตาภิลาปา, สตฺต อตฺเถ สนฺทสฺเสตุํ พฺยญฺชนนิรุตฺตาภิลาปา.\csromanspacer{} อญฺญา ธมฺมนิรุตฺติโย, อญฺญา อตฺถนิรุตฺติโย.\csromanspacer{} เยน ญาเณน อิมา นานา นิรุตฺติโย ญาตา, เตเนว ญาเณน อิมา นานา นิรุตฺติโย ปฏิวิทิตาติ, เตน วุจฺจติ นิรุตฺตินานตฺเต ปญฺญา นิรุตฺติปฏิสมฺภิเท ญาณํ. (3)}

\prose{สตฺตสุ ธมฺเมสุ ญาณานิ, สตฺตสุ อตฺเถสุ ญาณานิ, จุทฺทสสุ นิรุตฺตีสุ ญาณานิ.\csromanspacer{} อญฺญานิ ธมฺเมสุ ญาณานิ, อญฺญานิ อตฺเถสุ ญาณานิ, อญฺญานิ นิรุตฺตีสุ ญาณานิ, เยน ญาเณน อิเม นานา ญาณา ญาตา, เตเนว ญาเณน อิเม นานา ญาณา ปฏิวิทิตาติ, เตน วุจฺจติ ปฏิภานนานตฺเต ปญฺญา ปฏิภานปฏิสมฺภิเท ญาณํ. (4)}

\prose{สมฺมาทิฏฺฐิ ธมฺโม, สมฺมาสงฺกปฺโป ธมฺโม, สมฺมาวาจา ธมฺโม, สมฺมากมฺมนฺโต ธมฺโม, สมฺมา-อาชีโว ธมฺโม, สมฺมาวายาโม ธมฺโม, สมฺมาสติ ธมฺโม, สมฺมาสมาธิ ธมฺโม.\csromanspacer{} อญฺโญ สมฺมาทิฏฺฐิ ธมฺโม, อญฺโญ \csromanfolio{88}สมฺมาสงฺกปฺโป ธมฺโม, อญฺโญ สมฺมาวาจา ธมฺโม, อญฺโญ สมฺมากมฺมนฺโต ธมฺโม, อญฺโญ สมฺมา-อาชีโว ธมฺโม, อญฺโญ สมฺมาวายาโม ธมฺโม, อญฺโญ สมฺมาสติ ธมฺโม, อญฺโญ สมฺมาสมาธิ ธมฺโม.\csromanspacer{} เยน ญาเณน อิเม นานา ธมฺมา ญาตา, เตเนว ญาเณน อิเม นานา ธมฺมา ปฏิวิทิตาติ, เตน วุจฺจติ ธมฺมนานตฺเต ปญฺญา ธมฺมปฏิสมฺภิเท ญาณํ.}

\prose{ทสฺสนฏฺโฐ อตฺโถ, อภินิโรปนฏฺโฐ อตฺโถ, ปริคฺคหฏฺโฐ อตฺโถ, สมุฏฺฐานฏฺโฐ อตฺโถ, โวทานฏฺโฐ อตฺโถ, ปคฺคหฏฺโฐ อตฺโถ, อุปฏฺฐานฏฺโฐ อตฺโถ, อวิกฺเขปฏฺโฐ อตฺโถ.\csromanspacer{} อญฺโญ ทสฺสนฏฺโฐ อตฺโถ, อญฺโญ อภินิโรปนฏฺโฐ อตฺโถ, อญฺโญ ปริคฺคหฏฺโฐ อตฺโถ, อญฺโญ สมุฏฺฐานฏฺโฐ อตฺโถ, อญฺโญ โวทานฏฺโฐ อตฺโถ, อญฺโญ ปคฺคหฏฺโฐ อตฺโถ, อญฺโญ อุปฏฺฐานฏฺโฐ อตฺโถ, อญฺโญ อวิกฺเขปฏฺโฐ อตฺโถ.\csromanspacer{} เยน ญาเณน อิเม นานา อตฺถา ญาตา, เตเนว ญาเณน อิเม นานา อตฺถา ปฏิวิทิตาติ, เตน วุจฺจติ อตฺถนานตฺเต ปญฺญา อตฺถปฏิสมฺภิเท ญาณํ.}

\prose{อฏฺฐ ธมฺเม สนฺทสฺเสตุํ พฺยญฺชนนิรุตฺตาภิลาปา, อฏฺฐ อตฺเถ สนฺทสฺเสตุํ พฺยญฺชนนิรุตฺตาภิลาปา.\csromanspacer{} อญฺญา ธมฺมนิรุตฺติโย, อญฺญา อตฺถนิรุตฺติโย เยน ญาเณน อิมา นานา นิรุตฺติโย ญาตา, เตเนว ญาเณน อิมา นานา นิรุตฺติโย ปฏิวิทิตาติ, เตน วุจฺจติ นิรุตฺตินานตฺเต ปญฺญา นิรุตฺติปฏิสมฺภิเท ญาณํ.}

\prose{อฏฺฐสุ ธมฺเมสุ ญาณานิ, อฏฺฐสุ อตฺเถสุ ญาณานิ, โสฬสสุ นิรุตฺตีสุ ญาณานิ.\csromanspacer{} อญฺญานิ ธมฺเมสุ ญาณานิ, อญฺญานิ อตฺเถสุ ญาณานิ อญฺญานิ นิรุตฺตีสุ ญาณานิ.\csromanspacer{} เยน ญาเณน อิเม นานา ญาณา ญาตา, เตเนว ญาเณน อิเม นานา ญาณา ปฏิวิทิตาติ, เตน วุจฺจติ ปฏิภานนานตฺเต ปญฺญา ปฏิภานปฏิสมฺภิเท ญาณํ.\csromanspacer{} ตํญาตฏฺเฐน ญาณํ, ปชานนฏฺเฐน ปญฺญา, เตน วุจฺจติ อตฺถนานตฺเต ปญฺญา อตฺถปฏิสมฺภิเท ญาณํ, ธมฺมนานตฺเต ปญฺญา ธมฺมปฏิสมฺภิเท ญาณํ, นิรุตฺตินานตฺเต ปญฺญา นิรุตฺติปฏิสมฺภิเท ญาณํ, ปฏิภานนานตฺเต ปญฺญา ปฏิภานปฏิสมฺภิเท ญาณํ.}

\vspace{\tipitakaparskip}
\csromancenter{ปฏิสมฺภิทาญาณนิทฺเทโส อฏฺฐวีสติโม.}
\csromanmatikaanchor{mtk.24}
\csromantocmarkref{subsection}{29-31. ญาณตฺตยนิทฺเทส}{mtk.24}
\bookmark[dest={mtk.24},level=2]{29-31. ญาณตฺตยนิทฺเทส}
\csromanheader{2}{29-31. ญาณตฺตยนิทฺเทส}
\proseitem{78}{\csromanfolio{89}กถํ วิหารนานตฺเต ปญฺญา วิหารฏฺเฐ ญาณํ, สมาปตฺตินานตฺเต ปญฺญา สมาปตฺตฏฺเฐ ญาณํ, วิหารสมาปตฺตินานตฺเต ปญฺญา วิหารสมาปตฺตฏฺเฐ ญาณํ\csromandash{}นิมิตฺตํ ภยโต สมฺปสฺสมาโน อนิมิตฺเต อธิมุตฺตตฺตา ผุสฺส ผุสฺส วยํ ปสฺสติ, อนิมิตฺโต วิหาโร.\csromanspacer{} ปณิธิํ ภยโต สมฺปสฺสมาโน อปฺปณิหิเต อธิมุตฺตตฺตา ผุสฺส ผุสฺส วยํ ปสฺสติ, อปฺปณิหิโต วิหาโร.\csromanspacer{} อภินิเวสํ ภยโต สมฺปสฺสมาโน สุญฺญเต อธิมุตฺตตฺตา ผุสฺส ผุสฺส วยํ ปสฺสติ, สุญฺญโต วิหาโร.}

\prose{นิมิตฺตํ ภยโต สมฺปสฺสมาโน อนิมิตฺเต อธิมุตฺตตฺตา ปวตฺตํ อชฺฌุเปกฺขิตฺวา นิโรธํ นิพฺพานํ อนิมิตฺตํ อาวชฺชิตฺวา สมาปชฺชติ, อนิมิตฺตา สมาปตฺติ.\csromanspacer{} ปณิธิํ ภยโต สมฺปสฺสมาโน อปฺปณิหิเต อธิมุตฺตตฺตา ปวตฺตํ อชฺฌุเปกฺขิตฺวา นิโรธํ นิพฺพานํ อปฺปณิหิตํ อาวชฺชิตฺวา สมาปชฺชติ, อปฺปณิหิตา สมาปตฺติ.\csromanspacer{} อภินิเวสํ ภยโต สมฺปสฺสมาโน สุญฺญเต อธิมุตฺตตฺตา ปวตฺตํ อชฺฌุเปกฺขิตฺวา นิโรธํ นิพฺพานํ สุญฺญตํ อาวชฺชิตฺวา สมาปชฺชติ, สุญฺญตา สมาปตฺติ.}

\prose{นิมิตฺตํ ภยโต สมฺปสฺสมาโน อนิมิตฺเต อธิมุตฺตตฺตา ผุสฺส ผุสฺส วยํ ปสฺสติ, ปวตฺตํ อชฺฌุเปกฺขิตฺวา นิโรธํ นิพฺพานํ อนิมิตฺตํ อาวชฺชิตฺวา สมาปชฺชติ, อนิมิตฺตวิหารสมาปตฺติ.\csromanspacer{} ปณิธิํ ภยโต สมฺปสฺสมาโน อปฺปณิหิเต อธิมุตฺตตฺตา ผุสฺส ผุสฺส วยํ ปสฺสติ, ปวตฺตํ อชฺฌุเปกฺขิตฺวา นิโรธํ นิพฺพานํ อปฺปณิหิตํ อาวชฺชิตฺวา สมาปชฺชติ, อปฺปณิหิตวิหารสมาปตฺติ.\csromanspacer{} อภินิเวสํ ภยโต สมฺปสฺสมาโน สุญฺญเต อธิมุตฺตตฺตา ผุสฺส ผุสฺส วยํ ปสฺสติ, ปวตฺตํ อชฺฌุเปกฺขิตฺวา นิโรธํ นิพฺพานํ สุญฺญตํ อาวชฺชิตฺวา สมาปชฺชติ, สุญฺญตวิหารสมาปตฺติ.}

\proseitem{79}{รูปนิมิตฺตํ ภยโต สมฺปสฺสมาโน อนิมิตฺเต อธิมุตฺตตฺตา ผุสฺส ผุสฺส วยํ ปสฺสติ, อนิมิตฺโต วิหาโร.\csromanspacer{} รูปปณิธํ ภยโต สมฺปสฺสมาโน อปฺปณิหิเต อธิมุตฺตตฺตา ผุสฺส ผุสฺส วยํ ปสฺสติ, อปฺปณิหิโต วิหาโร.\csromanspacer{} รูปาภินิเวสํ ภยโต สมฺปสฺสมาโน สุญฺญเต อธิมุตฺตตฺตา ผุสฺส ผุสฺส วยํ ปสฺสติ, สุญฺญโต วิหาโร.}

\prose{รูปนิมิตฺตํ ภยโต สมฺปสฺสมาโน อนิมิตฺเต อธิมุตฺตตฺตา ปวตฺตํ อชฺฌุเปกฺขิตฺวา นิโรธํ นิพฺพานํ อนิมิตฺตํ อาวชฺชิตฺวา สมาปชฺชติ, อนิมิตฺตา \csromanfolio{90}สมาปตฺติ.\csromanspacer{} รูปปณิธิํ ภยโต สมฺปสฺสมาโน อปฺปณิหิเต อธิมุตฺตตฺตา ปวตฺตํ อชฺฌุเปกฺขิตฺวา นิโรธํ นิพฺพานํ อปฺปณิหิตํ อาวชฺชิตฺวา สมาปชฺชติ, อปฺปณิหิตา สมาปตฺติ.\csromanspacer{} รูปาภินิเวสํ ภยโต สมฺปสฺสมาโน สุญฺญเต อธิมุตฺตตฺตา ปวตฺตํ อชฺฌุเปกฺขิตฺวา นิโรธํ นิพฺพานํ สุญฺญตํ อาวชฺชิตฺวา สมาปชฺชติ, สุญฺญตา สมาปตฺติ.}

\prose{รูปนิมิตฺตํ ภยโต สมฺปสฺสมาโน อนิมิตฺเต อธิมุตฺตตฺตา ผุสฺส ผุสฺส วยํ ปสฺสติ, ปวตฺตํ อชฺฌุเปกฺขิตฺวา นิโรธํ นิพฺพานํ อนิมิตฺตํ อาวชฺชิตฺวา สมาปชฺชติ, อนิมิตฺตวิหารสมาปตฺติ.\csromanspacer{} รูปปณิธิํ ภยโต สมฺปสฺสมาโน อปฺปณิหิเต อธิมุตฺตตฺตา ผุสฺส ผุสฺส วยํ ปสฺสติ, ปวตฺตํ อชฺฌุเปกฺขิตฺวา นิโรธํ นิพฺพานํ อปฺปณิหิตํ อาวชฺชิตฺวา สมาปชฺชติ, อปฺปณิหิตวิหารสมาปตฺติ.\csromanspacer{} รูปาภินิเวสํ ภยโต สมฺปสฺสมาโน สุญฺญเต อธิมุตฺตตฺตา ผุสฺส ผุสฺส วยํ ปสฺสติ, ปวตฺตํ อชฺฌุเปกฺขิตฺวา นิโรธํ นิพฺพานํ สุญฺญตํ อาวชฺชิตฺวา สมาปชฺชติ, สุญฺญตวิหารสมาปตฺติ.}

\prose{เวทนานิมิตฺตํ ฯเปฯ.\csromanspacer{} สญฺญานิมิตฺตํ.\csromanspacer{} สงฺขารนิมิตฺตํ.\csromanspacer{} วิญฺญาณนิมิตฺตํ.\csromanspacer{} จกฺขุนิมิตฺตํ ฯเปฯ.\csromanspacer{} ชรามรณนิมิตฺตํ ภยโต สมฺปสฺสมาโน อนิมิตฺเต อธิมุตฺตตฺตา ผุสฺส ผุสฺส วยํ ปสฺสติ, อนิมิตฺโต วิหาโร.\csromanspacer{} ชรามรณปณิธิํ ภยโต สมฺปสฺสมาโน อปฺปณิหิเต อธิมุตฺตตฺตา ผุสฺส ผุสฺส วยํ ปสฺสติ, อปฺปณิหิโต วิหาโร.\csromanspacer{} ชรามรณาภินิเวสํ ภยโต สมฺปสฺสมาโน สุญฺญเต อธิมุตฺตตฺตา ผุสฺส ผุสฺส วยํ ปสฺสติ, สุญฺญโต วิหาโร.}

\prose{ชรามรณนิมิตฺตํ ภยโต สมฺปสฺสมาโน อนิมิตฺเต อธิมุตฺตตฺตา ปวตฺตํ อชฺฌุเปกฺขิตฺวา นิโรธํ นิพฺพานํ อนิมิตฺตํ อาวชฺชิตฺวา สมาปชฺชติ, อนิมิตฺตา สมาปตฺติ.\csromanspacer{} ชรามรณปณิธิํ ภยโต สมฺปสฺสมาโน อปฺปณิหิเต อธิมุตฺตตฺตา ปวตฺตํ อชฺฌุเปกฺขิตฺวา นิโรธํ นิพฺพานํ อปฺปณิหิตํ อาวชฺชิตฺวา สมาปชฺชติ, อปฺปณิหิตา สมาปตฺติ.\csromanspacer{} ชรามรณาภินิเวสํ ภยโต สมฺปสฺสมาโน สุญฺญเต อธิมุตฺตตฺตา ปวตฺตํ อชฺฌุเปกฺขิตฺวา นิโรธํ นิพฺพานํ สุญฺญตํ อาวชฺชิตฺวา สมาปชฺชติ, สุญฺญตา สมาปตฺติ.}

\prose{ชรามรณนิมิตฺตํ ภยโต สมฺปสฺสมาโน อนิมิตฺเต อธิมุตฺตตฺตา ผุสฺส ผุสฺส วยํ ปสฺสติ, ปวตฺตํ อชฺฌุเปกฺขิตฺวา นิโรธํ นิพฺพานํ อนิมิตฺตํ อาวชฺชิตฺวา สมาปชฺชติ, อนิมิตฺตวิหารสมาปตฺติ.\csromanspacer{} ชรามรณปณิธิํ \csromanfolio{91}ภยโต สมฺปสฺสมาโน อปฺปณิหิเต อธิมุตฺตตฺตา ผุสฺส ผุสฺส วยํ ปสฺสติ, ปวตฺตํ อชฺฌุเปกฺขิตฺวา นิโรธํ นิพฺพานํ อปฺปณิหิตํ อาวชฺชิตฺวา สมาปชฺชติ, อปฺปณิหิตวิหารสมาปตฺติ.\csromanspacer{} ชรามรณาภินิเวสํ ภยโต สมฺปสฺสมาโน สุญฺญเต อธิมุตฺตตฺตา ผุสฺส ผุสฺส วยํ ปสฺสติ, ปวตฺตํ อชฺฌุเปกฺขิตฺวา นิโรธํ นิพฺพานํ สุญฺญตํ อาวชฺชิตฺวา สมาปชฺชติ, สุญฺญตวิหารสมาปตฺติ.\csromanspacer{} อญฺโญ อนิมิตฺโต วิหาโร, อญฺโญ อปฺปณิหิโต วิหาโร, อญฺโญ สุญฺญโต วิหาโร.\csromanspacer{} อญฺญา อนิมิตฺตา สมาปตฺติ, อญฺญา อปฺปณิหิตา สมาปตฺติ, อญฺญา สุญฺญตา สมาปตฺติ.\csromanspacer{} อญฺญา อนิมิตฺตวิหารสมาปตฺติ, อญฺญา อปฺปณิหิตวิหารสมาปตฺติ, อญฺญา สุญฺญตวิหารสมาปตฺติ.\csromanspacer{} ตํญาตฏฺเฐน ญาณํ, ปชานนฏฺเฐน ปญฺญา.\csromanspacer{} เตน วุจฺจติ วิหารนานตฺเต ปญฺญา วิหารฏฺเฐ ญาณํ.\csromanspacer{} สมาปตฺตินานตฺเต ปญฺญา สมาปตฺตฏฺเฐ ญาณํ.\csromanspacer{} วิหารสมาปตฺตินานตฺเต ปญฺญา วิหารสมาปตฺตฏฺเฐ ญาณํ.}

\csromanheader{2}{ญาณตฺตยนิทฺเทโส เอกติํสติโม.}
\csromansectionrule
\csromanmatikaanchor{mtk.25}
\csromantocmarkref{subsection}{32. อานนฺตริก สมาธิญาณนิทฺเทส}{mtk.25}
\bookmark[dest={mtk.25},level=2]{32. อานนฺตริก สมาธิญาณนิทฺเทส}
\csromanheader{2}{32. อานนฺตริกสมาธิญาณนิทฺเทส}
\proseitem{80}{กถํ อวิกฺเขปปริสุทฺธตฺตา อาสวสมุจฺเฉเท ปญฺญา อานนฺตริกสมาธิมฺหิ ญาณํ\csromandash{}เนกฺขมฺมวเสน จิตฺตสฺส เอกคฺคตา อวิกฺเขโป สมาธิ, ตสฺส สมาธิสฺส วเสน อุปฺปชฺชติ ญาณํ, เตน ญาเณน อาสวา ขียนฺติ, อิติ ปฐมํ สมโถ, ปจฺฉา ญาณํ, เตน ญาเณน อาสวานํ ขโย โหติ.\csromanspacer{} เตน วุจฺจติ อวิกฺเขปปริสุทฺธตฺตา อาสวสมุจฺเฉเท ปญฺญา อานนฺตริกสมาธิมฺหิ ญาณํ.}

\prose{\textbf{อาสวา}ติ กตเม เต อาสวา? กามาสโว ภวาสโว ทิฏฺฐาสโว อวิชฺชาสโว.\csromanspacer{} กตฺเถเต อาสวา ขียนฺติ? โสตาปตฺติมคฺเคน อนวเสโส ทิฏฺฐาสโว ขียติ, อปายคมนีโย กามาสโว ขียติ, อปายคมนีโย ภวาสโว ขียติ, อปายคมนีโย อวิชฺชาสโว ขียติ, เอตฺเถเต อาสวา ขียนฺติ.\csromanspacer{} สกทาคามิมคฺเคน โอฬาริโก กามาสโว ขียติ, ตเทกฏฺโฐ ภวาสโว ขียติ, ตเทกฏฺโฐ อวิชฺชาสโว ขียติ, เอตฺเถเต อาสวา ขียนฺติ.\csromanspacer{} อนาคามิมคฺเคน อนวเสโส กามาสโว ขียติ, ตเทกฏฺโฐ ภวาสโว ขียติ, \csromanfolio{92}ตเทกฏฺโฐ อวิชฺชาสโว ขียติ, เอตฺเถเต อาสวา ขียนฺติ.\csromanspacer{} อรหตฺตมคฺเคน อนวเสโส ภวาสโว ขียติ, อนวเสโส อวิชฺชาสโว ขียติ, เอตฺเถเต อาสวา ขียนฺติ.}

\prose{อพฺยาปาทวเสน.\csromanspacer{} อาโลกสญฺญาวเสน.\csromanspacer{} อวิกฺเขปวเสน.\csromanspacer{} ธมฺมววตฺถานวเสน.\csromanspacer{} ญาณวเสน.\csromanspacer{} ปาโมชฺชวเสน.\csromanspacer{}\csromanspacer{}ปฐมชฺฌานวเสน.\csromanspacer{} ทุติยชฺฌานวเสน.\csromanspacer{} ตติยชฺฌานวเสน.\csromanspacer{} จตุตฺถชฺฌานวเสน.\csromanspacer{}\csromanspacer{}อากาสานญฺจายตนสมาปตฺติวเสน.\csromanspacer{} วิญฺญาณญฺจายตนสมาปตฺติวเสน.\csromanspacer{} อากิญฺจญฺญายตนสมาปตฺติวเสน.\csromanspacer{} เนวสญฺญานาสญฺญายตนสมาปตฺติวเสน.\csromanspacer{}\csromanspacer{}ปถวีกสิณวเสน.\csromanspacer{} อาโปกสิณวเสน.\csromanspacer{} เตโชกสิณวเสน.\csromanspacer{} วาโยกสิณวเสน.\csromanspacer{} นีลกสิณวเสน.\csromanspacer{} ปีตกสิณวเสน.\csromanspacer{} โลหิตกสิณวเสน.\csromanspacer{} โอทาตกสิณวเสน.\csromanspacer{} อากาสกสิณวเสน.\csromanspacer{} วิญฺญาณกสิณวเสน.\csromanspacer{}\csromanspacer{}พุทฺธานุสฺสติวเสน.\csromanspacer{} ธมฺมานุสฺสติวเสน.\csromanspacer{} สํฆานุสฺสติวเสน.\csromanspacer{} สีลานุสฺสติวเสน.\csromanspacer{} จาคานุสฺสติวเสน.\csromanspacer{} เทวตานุสฺสติวเสน.\csromanspacer{} อานาปานสฺสติวเสน.\csromanspacer{} มรณสฺสติวเสน.\csromanspacer{} กายคตาสติวเสน.\csromanspacer{} อุปสมานุสฺสติวเสน.\csromanspacer{}\csromanspacer{}อุทฺธุมาตกสญฺญาวเสน.\csromanspacer{} วินีลกสญฺญาวเสน.\csromanspacer{} วิปุพฺพกสญฺญาวเสน.\csromanspacer{} วิจฺฉิทฺทกสญฺญาวเสน.\csromanspacer{} วิกฺขายิตกสญฺญาวเสน.\csromanspacer{} วิกฺขิตฺตกสญฺญาวเสน.\csromanspacer{} หตวิกฺขิตฺตกสญฺญาวเสน.\csromanspacer{} โลหิตกสญฺญาวเสน.\csromanspacer{} ปุฬวกสญฺญาวเสน.\csromanspacer{} อฏฺฐิกสญฺญาวเสน.}

\proseitem{81}{ทีฆํ อสฺสาสวเสน, ทีฆํ ปสฺสาสวเสน.\csromanspacer{} รสฺสํ อสฺสาสวเสน, รสฺสํ ปสฺสาสวเสน.\csromanspacer{} สพฺพกายปฏิสํเวที อสฺสาสวเสน, สพฺพกายปฏิสํเวที ปสฺสาสวเสน.\csromanspacer{} ปสฺสมฺภยํ กายสงฺขารํ อสฺสาสวเสน, ปสฺสมฺภยํ กายสงฺขารํ ปสฺสาสวเสน.\csromanspacer{}\csromanspacer{}ปีติปฏิสํเวที อสฺสาสวเสน, ปีติปฏิสํเวที ปสฺสาสวเสน.\csromanspacer{} สุขปฏิสํเวที อสฺสาสวเสน, สุขปฏิสํเวที ปสฺสาสวเสน.\csromanspacer{} จิตฺตสงฺขารปฏิสํเวที อสฺสาสวเสน, จิตฺตสงฺขารปฏิสํเวที ปสฺสาสวเสน.\csromanspacer{} ปสฺสมฺภยํ จิตฺตสงฺขารํ อสฺสาสวเสน, ปสฺสมฺภยํ จิตฺตาสงฺขารํ ปสฺสาสวเสน.\csromanspacer{}\csromanspacer{}จิตฺตปฏิสํเวที อสฺสาสวเสน, จิตฺตปฏิสํเวที ปสฺสาสวเสน.\csromanspacer{} อภิปฺปโมทยํ จิตฺตํ อสฺสาสวเสน, อภิปฺปโมทยํ จิตฺตํ ปสฺสาสวเสน.\csromanspacer{} สมาทหํ จิตฺตํ ฯเปฯ.\csromanspacer{} วิโมจยํ จิตฺตํ.\csromanspacer{}\csromanspacer{}อนิจฺจานุปสฺสี.\csromanspacer{} วิราคานุปสฺสี.\csromanspacer{} นิโรธานุปสฺสี.\csromanspacer{} ปฏินิสฺสคฺคานุปสฺสี อสฺสาสวเสน, ปฏินิสฺสคฺคานุปสฺสี \csromanfolio{93}ปสฺสาสวเสน จิตฺตสฺส เอกคฺคตา อวิกฺเขโป สมาธิ.\csromanspacer{} ตสฺส สมาธิสฺส วเสน อุปฺปชฺชติ ญาณํ.\csromanspacer{} เตน ญาเณน อาสวา ขียนฺติ, อิติ ปฐมํ สมโถ, ปจฺฉา ญาณํ.\csromanspacer{} เตน ญาเณน อาสวานํ ขโย โหติ, เตน วุจฺจติ อวิกฺเขปปริสุทฺธตฺตา อาสวสมุจฺเฉเท ปญฺญา อานนฺตริกสมาธิมฺหิ ญาณํ.}

\prose{\textbf{อาสวา}ติ กตเม เต อาสวา? กามาสโว ภวาสโว ทิฏฺฐาสโว อวิชฺชาสโว.\csromanspacer{} กตฺเถเต อาสวา ขียนฺติ? โสตาปตฺติมคฺเคน อนวเสโส ทิฏฺฐาสโว ขียติ, อปายคมนีโย กามาสโว ขียติ, อปายคมนีโย ภวาสโว ขียติ, อปายคมนีโย อวิชฺชาสโว ขียติ, เอตฺเถเต อาสวา ขียนฺติ.\csromanspacer{} สกทาคามิมคฺเคน โอฬาริโก กามาสโว ขียติ, ตเทกฏฺโฐ ภวาสโว ขียติ, ตเทกฏฺโฐ อวิชฺชาสโว ขียติ, เอตฺเถเต อาสวา ขียนฺติ.\csromanspacer{} อนาคามิมคฺเคน อนวเสโส กามาสโว ขียติ, ตเทกฏฺโฐ ภวาสโว ขียติ, ตเทกฏฺโฐ อวิชฺชาสโว ขียติ, เอตฺเถเต อาสวา ขียนฺติ.\csromanspacer{} อรหตฺตมคฺเคน อนวเสโส ภวาสโว ขียติ, อนวเสโส อวิชฺชาสโว ขียติ, เอตฺเถเต อาสวา ขียนฺติ.\csromanspacer{} ตํญาตฏฺเฐน ญาณํ, ปชานนฏฺเฐน ปญฺญา, เตน วุจฺจติ อวิกฺเขปปริสุทฺธตฺตา อาสวสมุจฺเฉเท ปญฺญา อานนฺตริกสมาธิมฺหิ ญาณํ.}

\csromanheader{2}{อานนฺตริกสมาธิญาณนิทฺเทโส ทฺวตฺติํสติโม.}
\csromansectionrule
\csromanmatikaanchor{mtk.26}
\csromantocmarkref{subsection}{33. อรณวิหารญาณนิทฺเทส}{mtk.26}
\bookmark[dest={mtk.26},level=2]{33. อรณวิหารญาณนิทฺเทส}
\csromanheader{2}{33. อรณวิหารญาณนิทฺเทส}
\proseitem{82}{กถํ ทสฺสนาธิปเตยฺยํ สนฺโต จ วิหาราธิคโม ปณีตาธิมุตฺตตา ปญฺญา อรณวิหาเร ญาณํ\csromandash{}\textbf{ทสฺสนาธิปเตยฺยนฺ}ติ อนิจฺจานุปสฺสนา ทสฺสนาธิปเตยฺยํ, ทุกฺขานุปสฺสนา ทสฺสนาธิปเตยฺยํ, อนตฺตานุปสฺสนา ทสฺสนาธิปเตยฺยํ, รูเป อนิจฺจานุปสฺสนา ทสฺสนาธิปเตยฺยํ, รูเป ทุกฺขานุปสฺสนา ทสฺสนาธิปเตยฺยํ, รูเป อนตฺตานุปสฺสนา ทสฺสนาธิปเตยฺยํ, เวทนาย ฯเปฯ สญฺญาย.\csromanspacer{} สงฺขาเรสุ.\csromanspacer{} วิญฺญาเณ.\csromanspacer{} จกฺขุสฺมิํ ฯเปฯ.\csromanspacer{} ชรามรเณ อนิจฺจานุปสฺสนา ทสฺสนาธิปเตยฺยํ, ชรามรเณ ทุกฺขานุปสฺสนา ทสฺสนาธิปเตยฺยํ, ชรามรเณ อนตฺตานุปสฺสนา ทสฺสนาธิปเตยฺยํ.}

\prose{\csromanfolio{94}\textbf{สนฺโต จ วิหาราธิคโม}ติ สุญฺญโต วิหาโร สนฺโต วิหาราธิคโม, อนิมิตฺโต วิหาโร สนฺโต วิหาราธิคโม, อปฺปณิหิโต วิหาโร สนฺโต วิหาราธิคโม.}

\prose{\textbf{ปณีตาธิมุตฺตตา}ติ สุญฺญเต อธิมุตฺตตา ปณีตาธิมุตฺตตา, อนิมิตฺเต อธิมุตฺตตา ปณีตาธิมุตฺตตา, อปฺปณิหิเต อธิมุตฺตตา ปณีตาธิมุตฺตตา.}

\prose{\textbf{อรณวิหาโร}ติ ปฐมํ ฌานํ อรณวิหาโร, ทุติยํ ฌานํ อรณวิหาโร, ตติยํ ฌานํ อรณวิหาโร, จตุตฺถํ ฌานํ อรณวิหาโร อากาสานญฺจายตนสมาปตฺติ อรณวิหาโร ฯเปฯ เนวสญฺญานาสญฺญายตน- สมาปตฺติ อรณวิหาโร.}

\prose{\textbf{อรณวิหาโร}ติ เกนฏฺเฐน อรณวิหาโร, ปฐเมน ฌาเนน นีวรเณ หรตีติ อรณวิหาโร, ทุติเยน ฌาเนน วิตกฺกวิจาเร หรตีติ อรณวิหาโร, ตติเยน ฌาเนน ปีติํ หรตีติ อรณวิหาโร, จตุตฺเถน ฌาเนน สุขทุกฺเข หรตีติ อรณวิหาโร.\csromanspacer{}\csromanspacer{}อากาสานญฺจายตนสมาปตฺติยา รูปสญฺญํ ปฏิฆสญฺญํ นานตฺตสญฺญํ หรตีติ อรณวิหาโร, วิญฺญาณญฺจายตนสมาปตฺติยา อากาสานญฺจายตนสญฺญํ หรตีติ อรณวิหาโร, อากิญฺจญฺญายตนสมาปตฺติยา วิญฺญาณญฺจายตนสญฺญํ หรตีติ อรณวิหาโร, เนวสญฺญานาสญฺญายตนสมาปตฺติยา อากิญฺจญฺญายตนสญฺญํ หรตีติ อรณวิหาโร, อยํ อรณวิหาโร.\csromanspacer{} ตํญาตฏฺเฐน ญาณํ, ปชานนฏฺเฐน ปญฺญา เตน วุจฺจติ ทสฺสนาธิปเตยฺยํ สนฺโต จ วิหาราธิคโม ปณีตาธิมุตฺตตา ปญฺญา อรณวิหาเร ญาณํ.}

\csromanheader{2}{อรณวิหารญาณนิทฺเทโส เตตฺติํสติโม.}
\csromansectionrule
\csromanmatikaanchor{mtk.27}
\csromantocmarkref{subsection}{34. นิโรธสมาปตฺติญาณนิทฺเทส}{mtk.27}
\bookmark[dest={mtk.27},level=2]{34. นิโรธสมาปตฺติญาณนิทฺเทส}
\csromanheader{2}{34. นิโรธสมาปตฺติญาณนิทฺเทส}
\proseitem{83}{กถํ ทฺวีหิ พเลหิ สมนฺนาคตตฺตา ตโย จ สงฺขารานํ ปฏิปฺปสฺสทฺธิยา โสฬสหิ ญาณจริยาหิ นวหิ สมาธิจริยาหิ วสิภาวตา ปญฺญา นิโรธสมาปตฺติยา ญาณํ\csromandash{}}

\prose{\csromanfolio{95}\textbf{ทฺวีหิ พเลหี}ติ เทฺว พลานิ สมถพลํ วิปสฺสนาพลํ.\csromanspacer{} กตมํ สมถพลํ? เนกฺขมฺมวเสน จิตฺตสฺเสกคฺคตา อวิกฺเขโป สมถพลํ, อพฺยาปาทวเสน จิตฺตสฺเสกคฺคตา อวิกฺเขโป สมถพลํ, อาโลกสญฺญาวเสน จิตฺตสฺเสกคฺคตา อวิกฺเขโป สมถพลํ, อวิกฺเขปวเสน จิตฺตสฺเสกคฺคตา อวิกฺเขโป สมถพลํ ฯเปฯ ปฏินิสฺสคฺคานุปสฺสี อสฺสาสวเสน จิตฺตสฺเสกคฺคตา อวิกฺเขโป สมถพลํ, ปฏินิสฺสคฺคานุปสฺสี ปสฺสาสวเสน จิตฺตสฺเสกคฺคตา อวิกฺเขโป สมถพลํ.}

\prose{\textbf{สมถพลนฺ}ติ เกนฏฺเฐน สมถพลํ? ปฐเมน ฌาเนน นีวรเณ น กมฺปตีติ สมถพลํ, ทุติเยน ฌาเนน วิตกฺกวิจาเร น กมฺปตีติ สมถพลํ, ตติเยน ฌาเนน ปีติยา น กมฺปตีติ สมถพลํ, จตุตฺเถน ฌาเนน สุขทุกฺเข น กมฺปตีติ สมถพลํ.\csromanspacer{}\csromanspacer{}อากาสานญฺจายตนสมาปตฺติยา รูปสญฺญาย ปฏิฆสญฺญาย นานตฺตสญฺญาย น กมฺปตีติ สมถพลํ, วิญฺญาณญฺจายตนสมาปตฺติยา อากาสานญฺจายตนสญฺญาย น กมฺปตีติ สมถพลํ, อากิญฺจญฺญายตนสมาปตฺติยา วิญฺญาณญฺจายตนสญฺญาย น กมฺปตีติ สมถพลํ, เนวสญฺญานาสญฺญายตนสมาปตฺติยา อากิญฺจญฺญายตนสญฺญาย น กมฺปตีติ สมถพลํ, อุทฺธจฺเจ จ อุทฺธจฺจสหคตกิเลเส จ ขนฺเธ จ น กมฺปติ น จลติ น เวธตีติ สมถพลํ.\csromanspacer{} อิทํ สมถพลํ.}

\prose{กตมํ วิปสฺสนาพลํ? อนิจฺจานุปสฺสนา วิปสฺสนาพลํ, ทุกฺขานุปสฺสนา วิปสฺสนาพลํ, อนตฺตานุปสฺสนา วิปสฺสนาพลํ, นิพฺพิทานุปสฺสนา วิปสฺสนาพลํ, วิราคานุปสฺสนา วิปสฺสนาพลํ, นิโรธานุปสฺสนา วิปสฺสนาพลํ, ปฏินิสฺสคฺคานุปสฺสนา วิปสฺสนาพลํ.\csromanspacer{} รูเป อนิจฺจานุปสฺสนา วิปสฺสนาพลํ ฯเปฯ รูเป ปฏินิสฺสคฺคานุปสฺสนา วิปสฺสนาพลํ.\csromanspacer{}\csromanspacer{}เวทนาย ฯเปฯ.\csromanspacer{} สญฺญาย.\csromanspacer{} สงฺขาเรสุ.\csromanspacer{} วิญฺญาเณ.\csromanspacer{} จกฺขุสฺมิํ ฯเปฯ.\csromanspacer{} ชรามรเณ อนิจฺจานุปสฺสนา วิปสฺสนาพลํ ฯเปฯ ชรามรเณ ปฏินิสฺสคฺคานุปสฺสนา วิปสฺสนาพลํ.}

\prose{\textbf{วิปสฺสนาพลนฺ}ติ เกนฏฺเฐน วิปสฺสนาพลํ? อนิจฺจานุปสฺสนาย นิจฺจสญฺญาย น กมฺปตีติ วิปสฺสนาพลํ, ทุกฺขานุปสฺสนาย สุขสญฺญาย น กมฺปตีติ วิปสฺสนาพลํ, อนตฺตานุปสฺสนาย อตฺตสญฺญาย น กมฺปตีติ วิปสฺสนาพลํ, นิพฺพิทานุปสฺสนาย นนฺทิยา น กมฺปตีติ วิปสฺสนาพลํ, วิราคานุปสฺสนาย ราเค น กมฺปตีติ วิปสฺสนาพลํ, นิโรธานุปสฺสนาย สมุทเย น กมฺปตีติ วิปสฺสนาพลํ, ปฏินิสฺสคฺคานุปสฺสนาย อาทาเน น \csromanfolio{96}กมฺปตีติ วิปสฺสนาพลํ, อวิชฺชาย จ อวิชฺชาสหคตกิเลเส จ ขนฺเธ จ น กมฺปติ น จลติ น เวธตีติ วิปสฺสนาพลํ.\csromanspacer{} อิทํ วิปสฺสนาพลํ.}

\prose{\textbf{ตโย จ สงฺขารานํ ปฏิปฺปสฺสทฺธิยา}ติ กตเมสํ ติณฺณนฺนํ สงฺขารานํ ปฏิปฺปสฺสทฺธิยา? ทุติยํ ฌานํ สมาปนฺนสฺส วิตกฺกวิจารา วจีสงฺขารา ปฏิปฺปสฺสทฺธา โหนฺติ, จตุตฺถํ ฌานํ สมาปนฺนสฺส อสฺสาสปสฺสาสา กายสงฺขารา ปฏิปฺปสฺสทฺธา โหนฺติ, สญฺญาเวทยิตนิโรธํ สมาปนฺนสฺส สญฺญา จ เวทนา จ จิตฺตสงฺขารา ปฏิปฺปสฺสทฺธา โหนฺติ.\csromanspacer{} อิเมสํ ติณฺณนฺนํ สงฺขารานํ ปฏิปฺปสฺสทฺธิยา.}

\proseitem{84}{\textbf{โสฬสหิ ญาณจริยาหี}ติ กตมาหิ โสฬสหิ ญาณจริยาหิ? อนิจฺจานุปสฺสนา ญาณจริยา, ทุกฺขานุปสฺสนา ญาณจริยา, อนตฺตานุปสฺสนา ญาณจริยา, นิพฺพิทานุปสฺสนา ญาณจริยา, วิราคานุปสฺสนา ญาณจริยา, นิโรธานุปสฺสนา ญาณจริยา, ปฏินิสฺสคฺคานุปสฺสนา ญาณจริยา, วิวฏฺฏนานุปสฺสนา ญาณจริยา.\csromanspacer{} โสตาปตฺติมคฺโค ญาณจริยา, โสตาปตฺติผลสมาปตฺติ ญาณจริยา, สกทาคามิมคฺโค ญาณจริยา, สกทาคามิผลสมาปตฺติ ญาณจริยา, อนาคามิมคฺโค ญาณจริยา, อนาคามิผลสมาปตฺติ ญาณจริยา, อรหตฺตมคฺโค ญาณจริยา, อรหตฺตผลสมาปตฺติ ญาณจริยา.\csromanspacer{} อิมาหิ โสฬสหิ ญาณจริยาหิ.}

\proseitem{85}{\textbf{นวหิ สมาธิจริยาหี}ติ กตมาหิ นวหิ สมาธิจริยาหิ? ปฐมํ ฌานํ สมาธิจริยา, ทุติยํ ฌานํ สมาธิจริยา, ตติยํ ฌานํ สมาธิจริยา, จตุตฺถํ ฌานํ สมาธิจริยา.\csromanspacer{}\csromanspacer{}อากาสานญฺจายตนสมาปตฺติ ฯเปฯ.\csromanspacer{} วิญฺญาณญฺจายตนสมาปตฺติ.\csromanspacer{} อากิญฺจญฺญายตนสมาปตฺติ.\csromanspacer{} เนวสญฺญานาสญฺญายตนสมาปตฺติ สมาธิจริยา.\csromanspacer{}\csromanspacer{}ปฐมํ ฌานํ ปฏิลาภตฺถาย วิตกฺโก จ วิจาโร จ ปีติ จ สุขญฺจ จิตฺเตกคฺคตา จ ฯเปฯ.\csromanspacer{} เนวสญฺญานาสญฺญายตนสมาปตฺติํ ปฏิลาภตฺถาย วิตกฺโก จ วิจาโร จ ปีติ จ สุขญฺจ จิตฺเตกคฺคตา จ.\csromanspacer{} อิมาหิ นวหิ สมาธิจริยาหิ.}

\prose{\textbf{วสี}ติ ปญฺจ วสิโย อาวชฺชนวสี สมาปชฺชนวสี อธิฏฺฐานวสี วุฏฺฐานวสี ปจฺจเวกฺขณาวสี.\csromanspacer{} ปฐมํ ฌานํ ยตฺถิจฺฉกํ ยทิจฺฉกํ ยาวติจฺฉกํ อาวชฺชติ, อาวชฺชนาย ทนฺธายิตตฺตํ นตฺถีติ อาวชฺชนวสี.\csromanspacer{} ปฐมํ ฌานํ ยตฺถิจฺฉกํ ยทิจฺฉกํ ยาวติจฺฉกํ สมาปชฺชติ, สมาปชฺชนาย ทนฺธายิตตฺตํ นตฺถีติ สมาปชฺชนวสี.\csromanspacer{} ปฐมํ ฌานํ ยตฺถิจฺฉกํ ยทิจฺฉกํ ยาวติจฺฉกํ \csromanfolio{97}อธิฏฺฐาติ, อธิฏฺฐาเน ทนฺธายิตตฺตํ นตฺถีติ อธิฏฺฐานวสี\textbf{.}\csromanspacer{} ปฐมํ ฌานํ ยตฺถิจฺฉกํ ยทิจฺฉกํ ยาวติจฺฉกํ วุฏฺฐาติ, วุฏฺฐาเน ทนฺธายิตตฺตํ นตฺถีติ วุฏฺฐานวสี\textbf{.}\csromanspacer{} ปฐมํ ฌานํ ยตฺถิจฺฉกํ ยทิจฺฉกํ ยาวติจฺฉกํ ปจฺจเวกฺขติ, ปจฺจเวกฺขณาย ทนฺธายิตตฺตํ นตฺถีติ ปจฺจเวกฺขณาวสี\textbf{.}}

\prose{ทุติยํ ฌานํ ฯเปฯ.\csromanspacer{} เนวสญฺญานาสญฺญายตนสมาปตฺติํ ยตฺถิจฺฉกํ ยทิจฺฉกํ ยาวติจฺฉกํ อาวชฺชติ, อาวชฺชนาย ทนฺธายิตตฺตํ นตฺถีติ อาวชฺชนวสี.\csromanspacer{} เนวสญฺญานาสญฺญายตนสมาปตฺติํ ยตฺถิจฺฉกํ ยทิจฺฉกํ ยาวติจฺฉกํ สมาปชฺชติ ฯเปฯ อธิฏฺฐาติ, วุฏฺฐาติ, ปจฺจเวกฺขติ, ปจฺจเวกฺขณาย ทนฺธายิตตฺตํ นตฺถีติ ปจฺจเวกฺขณาวสี.\csromanspacer{} อิมา ปญฺจ วสิโย.\csromanspacer{} ตํญาตฏฺเฐน ญาณํ ปชานนฏฺเฐน ปญฺญา.\csromanspacer{} เตน วุจฺจติ ทฺวีหิ พเลหิ สมนฺนาคตตฺตา ตโย จ สงฺขารานํ ปฏิปสฺสทฺธิยา โสฬสหิ ญาณจริยาหิ นวหิ สมาธิจริยาหิ วสีภาวตา ปญฺญา นิโรธสมาปตฺติยา ญาณํ.}

\csromanheader{2}{นิโรธสมาปตฺติญาณนิทฺเทโส จตุตฺติํสติโม.}
\csromansectionrule
\csromanmatikaanchor{mtk.28}
\csromantocmarkref{subsection}{35. ปรินิพฺพานญาณนิทฺเทส}{mtk.28}
\bookmark[dest={mtk.28},level=2]{35. ปรินิพฺพานญาณนิทฺเทส}
\csromanheader{2}{35. ปรินิพฺพานญาณนิทฺเทส}
\proseitem{86}{กถํ สมฺปชานสฺส ปวตฺตปริยาทาเน ปญฺญา ปรินิพฺพาเน ญาณํ\csromandash{}อิธ สมฺปชาโน เนกฺขมฺเมน กามจฺฉนฺทสฺส ปวตฺตํ ปริยาทิยติ.\csromanspacer{} อพฺยาปาเทน พฺยาปาทสฺส ปวตฺตํ ปริยาทิยติ.\csromanspacer{} อาโลกสญฺญาย ถินมิทฺธสฺส ปวตฺตํ ปริยาทิยติ.\csromanspacer{} อวิกฺเขเปน อุทฺธจฺจสฺส ปวตฺตํ ปริยาทิยติ.\csromanspacer{} ธมฺมววตฺถาเนน วิจิกิจฺฉาย ฯเปฯ.\csromanspacer{} ญาเณน อวิชฺชาย.\csromanspacer{} ปาโมชฺเชน อรติยา.\csromanspacer{} ปฐเมน ฌาเนน นีวรณานํ ปวตฺตํ ปริยาทิยติ ฯเปฯ.\csromanspacer{} อรหตฺตมคฺเคน สพฺพกิเลสานํ ปวตฺตํ ปริยาทิยติ.}

\prose{อถวา ปน สมฺปชานสฺส อนุปาทิเสสาย นิพฺพานธาตุยา ปรินิพฺพายนฺตสฺส อิทํ เจว จกฺขุปวตฺตํ ปริยาทิยติ, อญฺญญฺจ จกฺขุปวตฺตํ น อุปฺปชฺชติ.\csromanspacer{} อิทํ เจว โสตปวตฺตํ ฯเปฯ.\csromanspacer{} ฆานปวตฺตํ.\csromanspacer{} ชิวฺหาปวตฺตํ.\csromanspacer{} กายปวตฺตํ.\csromanspacer{} มโนปวตฺตํ ปริยาทิยติ, อญฺญญฺจ มโนปวตฺตํ น อุปฺปชฺชติ.\csromanspacer{} อิทํ สมฺปชานสฺส \csromanfolio{98}ปวตฺตปริยาทาเน ปญฺญา ปรินิพฺพาเน ญาณํ\textbf{.}\csromanspacer{} ตํญาตฏฺเฐน ญาณํ, ปชานนฏฺเฐน ปญฺญา\textbf{.}\csromanspacer{} เตน วุจฺจติ สมฺปชานสฺส ปวตฺตปริยาทาเน ปญฺญา ปรินิพฺพาเน ญาณํ\textbf{.}}

\csromanheader{2}{ปรินิพฺพานญาณนิทฺเทโส ปญฺจติํสติโม\textbf{.}}
\csromansectionrule
\csromanmatikaanchor{mtk.29}
\csromantocmarkref{subsection}{36. สมสีสฏฺฐญาณนิทฺเทส}{mtk.29}
\bookmark[dest={mtk.29},level=2]{36. สมสีสฏฺฐญาณนิทฺเทส}
\csromanheader{2}{36. สมสีสฏฺฐญาณนิทฺเทส}
\proseitem{87}{กถํ สพฺพธมฺมานํ สมฺมา สมุจฺเฉเท นิโรเธ จ อนุปฏฺฐานตา ปญฺญา สมสีสฏฺเฐ ญาณํ\csromandash{}\textbf{สพฺพธมฺมานนฺ}ติ ปญฺจกฺขนฺธา ทฺวาทสายตนานิ อฏฺฐารส ธาตุโย กุสลา ธมฺมา อกุสลา ธมฺมา อพฺยากตา ธมฺมา กามาวจรา ธมฺมา รูปาวจรา ธมฺมา อรูปาวจรา ธมฺมา อปริยาปนฺนา ธมฺมา\textbf{.}\csromanspacer{} \textbf{สมฺมา สมุจฺเฉเท}ติ เนกฺขมฺเมน กามจฺฉนฺทํ สมฺมา สมุจฺฉินฺทติ, อพฺยาปาเทน พฺยาปาทํ สมฺมา สมุจฺฉินฺทติ\textbf{.}\csromanspacer{} อาโลกสญฺญาย ถินมิทฺธํ สมฺมา สมุจฺฉินฺทติ\textbf{.}\csromanspacer{} อวิกฺเขเปน อุทฺธจฺจํ สมฺมา สมุจฺฉินฺทติ\textbf{.}\csromanspacer{} ธมฺมววตฺถาเนน วิจิกิจฺฉํ สมฺมา สมุจฺฉินฺทติ\textbf{.}\csromanspacer{} ญาเณน อวิชฺชํ สมฺมา สมุจฺฉินฺทติ\textbf{.}\csromanspacer{} ปาโมชฺเชน อรติํ สมฺมา สมุจฺฉินฺทติ\textbf{.}\csromanspacer{} ปฐเมน ฌาเนน นีวรเณ สมฺมา สมุจฺฉินฺทติ ฯเปฯ\textbf{.}\csromanspacer{} อรหตฺตมคฺเคน สพฺพกิเลเส สมฺมา สมุจฺฉินฺทติ\textbf{.}}

\prose{\textbf{นิโรเธ}ติ เนกฺขมฺเมน กามจฺฉนฺทํ นิโรเธติ, อพฺยาปาเทน พฺยาปาทํ นิโรเธติ\textbf{.}\csromanspacer{} อาโลกสญฺญาย ถินมิทฺธํ นิโรเธติ\textbf{.}\csromanspacer{} อวิกฺเขเปน อุทฺธจฺจํ นิโรเธติ, ธมฺมววตฺถาเนน วิจิกิจฺฉํ นิโรเธติ\textbf{.}\csromanspacer{} ญาเณน อวิชฺชํ นิโรเธติ\textbf{.}\csromanspacer{} ปาโมชฺเชน อรติํ นิโรเธติ\textbf{.}\csromanspacer{} ปฐเมน ฌาเนน นีวรเณ นิโรเธติ ฯเปฯ อรหตฺตมคฺเคน สพฺพกิเลเส นิโรเธติ\textbf{.}}

\prose{\textbf{อนุปฏฺฐานตา}ติ เนกฺขมฺมํ ปฏิลทฺธสฺส กามจฺฉนฺโท น อุปฏฺฐาติ\textbf{.}\csromanspacer{} อพฺยาปาทํ ปฏิลทฺธสฺส พฺยาปาโท น อุปฏฺฐาติ\textbf{.}\csromanspacer{} อาโลกสญฺญํ ปฏิลทฺธสฺส ถินมิทฺธํ น อุปฏฺฐาติ\textbf{.}\csromanspacer{} อวิกฺเขปํ ปฏิลทฺธสฺส อุทฺธจฺจํ น อุปฏฺฐาติ\textbf{.}\csromanspacer{} ธมฺมววตฺถานํ ปฏิลทฺธสฺส วิจิกิจฺฉา น อุปฏฺฐาติ\textbf{.}\csromanspacer{} ญาณํ ปฏิลทฺธสฺส อวิชฺชา น อุปฏฺฐาติ\textbf{.}\csromanspacer{} ปาโมชฺชํ ปฏิลทฺธสฺส อรติ น อุปฏฺฐาติ\textbf{.}\csromanspacer{} ปฐมํ ฌานํ ปฏิลทฺธสฺส นีวรณา น อุปฏฺฐหนฺติ ฯเปฯ\textbf{.}\csromanspacer{} อรหตฺตมคฺคํ ปฏิลทฺธสฺส สพฺพกิเลสา น อุปฏฺฐหนฺติ\textbf{.}}

\prose{\textbf{สมนฺ}ติ กามจฺฉนฺทสฺส ปหีนตฺตา เนกฺขมฺมํ สมํ\textbf{.}\csromanspacer{} พฺยาปาทสฺส ปหีนตฺตา อพฺยาปาโท สมํ\textbf{.}\csromanspacer{} ถินมิทฺธสฺส ปหีนตฺตา อาโลกสญฺญา สมํ\textbf{.} \csromanfolio{99}อุทฺธจฺจสฺส ปหีนตฺตา อวิกฺเขโป สมํ.\csromanspacer{} วิจิกิจฺฉาย ปหีนตฺตา ธมฺมววตฺถานํ สมํ.\csromanspacer{} อวิชฺชาย ปหีนตฺตา ญาณํ สมํ.\csromanspacer{} อรติยา ปหีนตฺตา ปาโมชฺชํ สมํ.\csromanspacer{} นีวรณานํ ปหีนตฺตา ปฐมํ ฌานํ สมํ ฯเปฯ.\csromanspacer{} สพฺพกิเลสานํ ปหีนตฺตา อรหตฺตมคฺโค สมํ.}

\prose{\textbf{สีสนฺ}ติ เตรส สีสานิ.\csromanspacer{} ปลิโพธสีสญฺจ ตณฺหา, วินิพนฺธนสีสญฺจ มาโน, ปรามาสสีสญฺจ ทิฏฺฐิ, วิกฺเขปสีสญฺจ อุทฺธจฺจํ, สํกิเลสสีสญฺจ อวิชฺชา, อธิโมกฺขสีสญฺจ สทฺธา, ปคฺคหสีสญฺจ วีริยํ, อุปฏฺฐานสีสญฺจ สติ, อวิกฺเขปสีสญฺจ สมาธิ, ทสฺสนสีสญฺจ ปญฺญา, ปวตฺตสีสญฺจ ชีวิตินฺทฺริยํ, โคจรสีสญฺจ วิโมกฺโข, สงฺขารสีสญฺจ นิโรโธ.\csromanspacer{} ตํญาตฏฺเฐน ญาณํ ปชานนฏฺเฐน ปญฺญา.\csromanspacer{} เตน วุจฺจติ สพฺพธมฺมานํ สมฺมา สมุจฺเฉเท นิโรเธ จ อนุปฏฺฐานตา ปญฺญา สมสีสฏฺเฐ ญาณํ.}

\csromanheader{2}{สมสีสฏฺฐญาณนิทฺเทโส ฉตฺติํสติโม.}
\csromansectionrule
\csromanmatikaanchor{mtk.30}
\csromantocmarkref{subsection}{37. สลฺเลขฏฺฐญาณนิทฺเทส}{mtk.30}
\bookmark[dest={mtk.30},level=2]{37. สลฺเลขฏฺฐญาณนิทฺเทส}
\csromanheader{2}{37. สลฺเลขฏฺฐญาณนิทฺเทส}
\proseitem{88}{กถํ ปุถุนานตฺตเตชปริยาทาเน ปญฺญา สลฺเลขฏฺเฐ\footnote{ปุถุนานตฺเตกตฺตเตชปริยาทาเน (ก)} ญาณํ\csromandash{} \textbf{ปุถู}ติ ราโค ปุถุ, โทโส ปุถุ, โมโห ปุถุ, โกโธ ฯเปฯ.\csromanspacer{} อุปนาโห.\csromanspacer{} มกฺโข.\csromanspacer{} ปฬาโส.\csromanspacer{} อิสฺสา.\csromanspacer{} มจฺฉริยํ.\csromanspacer{} มายา.\csromanspacer{} สาเฐยฺยํ.\csromanspacer{} ถมฺโภ.\csromanspacer{} สารมฺโภ.\csromanspacer{} มาโน.\csromanspacer{} อติมาโน.\csromanspacer{} มโท.\csromanspacer{} ปมาโท.\csromanspacer{} สพฺเพ กิเลสา.\csromanspacer{} สพฺเพ ทุจฺจริตา.\csromanspacer{} สพฺเพ อภิสงฺขารา.\csromanspacer{} สพฺเพ ภวคามิกมฺมา.}

\prose{\textbf{นานตฺเตกตฺตนฺ}ติ กามจฺฉนฺโท นานตฺตํ, เนกฺขมฺมํ เอกตฺตํ.\csromanspacer{} พฺยาปาโท นานตฺตํ อพฺยาปาโท เอกตฺตํ.\csromanspacer{} ถินมิทฺธํ นานตฺตํ, อาโลกสญฺญา เอกตฺตํ.\csromanspacer{} อุทฺธจฺจํ นานตฺตํ, อวิกฺเขโป เอกตฺตํ.\csromanspacer{} วิจิกิจฺฉา นานตฺตํ, ธมฺมววตฺถานํ เอกตฺตํ.\csromanspacer{} อวิชฺชา นานตฺตํ, ญาณํ เอกตฺตํ.\csromanspacer{} อรติ นานตฺตํ, ปาโมชฺชํ เอกตฺตํ.\csromanspacer{} นีวรณา นานตฺตํ, ปฐมํ ฌานํ เอกตฺตํ ฯเปฯ สพฺเพ กิเลสา นานตฺตํ, อรหตฺตมคฺโค เอกตฺตํ.}

\prose{\textbf{เตโช}ติ ปญฺจ เตชา จรณเตโช คุณเตโช ปญฺญาเตโช ปุญฺญเตโช ธมฺมเตโช.\csromanspacer{} จรณเตเชน เตชิตตฺตา ทุสฺสีลฺยเตชํ ปริยาทิยติ, คุณเตเชน เตชิตตฺตา อคุณเตชํ ปริยาทิยติ, ปญฺญาเตเชน เตชิตตฺตา ทุปฺปญฺญเตชํ ปริยาทิยติ, ปุญฺญเตเชน \csromanfolio{100}เตชิตตฺตา อปุญฺญเตชํ ปริยาทิยติ, ธมฺมเตเชน เตชิตตฺตา อธมฺมเตชํ ปริยาทิยติ.}

\prose{\textbf{สลฺเลโข}ติ กามจฺฉนฺโท อสลฺเลโข, เนกฺขมฺมํ สลฺเลโข.\csromanspacer{} พฺยาปาโท อสลฺเลโข, อพฺยาปาโท สลฺเลโข.\csromanspacer{} ถินมิทฺธํ อสลฺเลโข, อาโลกสญฺญา สลฺเลโข.\csromanspacer{} อุทฺธจฺจํ อสลฺเลโข, อวิกฺเขโป สลฺเลโข.\csromanspacer{} วิจิกิจฺฉา อสลฺเลโข, ธมฺมววตฺถานํ สลฺเลโข.\csromanspacer{} อวิชฺชา อสลฺเลโข, ญาณํ สลฺเลโข.\csromanspacer{} อรติ อสลฺเลโข, ปาโมชฺชํ สลฺเลโข.\csromanspacer{} นีวรณา อสลฺเลโข, ปฐมํ ฌานํ สลฺเลโข ฯเปฯ.\csromanspacer{} สพฺพกิเลสา อสลฺเลโข, อรหตฺตมคฺโค สลฺเลโข.\csromanspacer{} ตํญาตฏฺเฐน ญาณํ, ปชานนฏฺเฐน ปญฺญา.\csromanspacer{} เตน วุจฺจติ ปุถุนานตฺตเตชปริยาทาเน ปญฺญา สลฺเลขฏฺเฐ ญาณํ.}

\csromanheader{2}{สลฺเลขฏฺฐญาณนิทฺเทโส สตฺตติํสติโม.}
\csromansectionrule
\csromanmatikaanchor{mtk.31}
\csromantocmarkref{subsection}{38. วีริยารมฺภญาณนิทฺเทส}{mtk.31}
\bookmark[dest={mtk.31},level=2]{38. วีริยารมฺภญาณนิทฺเทส}
\csromanheader{2}{38. วีริยารมฺภญาณนิทฺเทส}
\proseitem{89}{กถํ อสลฺลีนตฺตปหิตตฺตปคฺคหฏฺเฐ ปญฺญา วีริยารมฺเภ ญาณํ\csromandash{}อนุปฺปนฺนานํ ปาปกานํ อกุสลานํ ธมฺมานํ อนุปฺปาทาย อสลฺลีนตฺตปหิตตฺตปคฺคหฏฺเฐ ปญฺญา วีริยารมฺเภ ญาณํ.\csromanspacer{} อุปฺปนฺนานํ ปาปกานํ อกุสลานํ ธมฺมานํ ปหานาย อสลฺลีนตฺตปหิตตฺตปคฺคหฏฺเฐ ปญฺญา วีริยารมฺเภ ญาณํ.\csromanspacer{} อนุปฺปนฺนานํ กุสลานํ ธมฺมานํ อุปฺปาทาย อสลฺลีนตฺตปหิตตฺตปคฺคหฏฺเฐ ปญฺญา วีริยารมฺเภ ญาณํ.\csromanspacer{} อุปฺปนฺนานํ กุสลานํ ธมฺมานํ ฐิติยา อสมฺโมสาย ภิโยภาวาย เวปุลฺลาย ภาวนาย ปาริปูริยา อสลฺลีนตฺตปหิตตฺตปคฺคหฏฺเฐ ปญฺญา วีริยารมฺเภ ญาณํ.}

\prose{อนุปฺปนฺนสฺส กามจฺฉนฺทสฺส อนุปฺปาทาย อสลฺลีนตฺตปหิตตฺตปคฺคหฏฺเฐ ปญฺญา วีริยารมฺเภ ญาณํ.\csromanspacer{} อุปฺปนฺนสฺส กามจฺฉนฺทสฺส ปหานาย อสลฺลีนตฺตปหิตตฺตปคฺคหฏฺเฐ ปญฺญา วีริยารมฺเภ ญาณํ.\csromanspacer{} อนุปฺปนฺนสฺส เนกฺขมฺมสฺส อุปฺปาทาย อสลฺลีนตฺตปหิตตฺตปคฺคหฏฺเฐ ปญฺญา วีริยารมฺเภ ญาณํ.\csromanspacer{} อุปฺปนฺนสฺส เนกฺขมฺมสฺส ฐิติยา อสมฺโมสาย ภิยฺโยภาวาย เวปุลฺลาย ภาวนาย ปาริปูริยา อสลฺลีนตฺตปหิตตฺตปคฺคหฏฺเฐ ปญฺญา วีริยารมฺเภ ญาณํ ฯเปฯ.}

\prose{อนุปฺปนฺนานํ สพฺพกิเลสานํ อนุปฺปาทาย อสลฺลีนตฺตปหิตตฺตปคฺคหฏฺเฐ ปญฺญา วีริยารมฺเภ ญาณํ.\csromanspacer{} อุปฺปนฺนานํ สพฺพกิเลสานํ ปหานาย \csromanfolio{101}อสลฺลีนตฺตปหิตตฺตปคฺคหฏฺเฐ ปญฺญา วีริยารมฺเภ ญาณํ ฯเปฯ.\csromanspacer{} อนุปฺปนฺนสฺส อรหตฺตมคฺคสฺส อุปฺปาทาย อสลฺลีนตฺตปหิตตฺตปคฺคหฏฺเฐ ปญฺญา วีริยารมฺเภ ญาณํ.\csromanspacer{} อุปฺปนฺนสฺส อรหตฺตมคฺคสฺส ฐิติยา อสมฺโมสาย ภิยฺโยภาวาย เวปุลฺลาย ภาวนาย ปาริปูริยา อสลฺลีนตฺตปหิตตฺตปคฺคหฏฺเฐ ปญฺญา วีริยารมฺเภ ญาณํ.\csromanspacer{} ตํญาตฏฺเฐน ญาณํ, ปชานนฏฺเฐน ปญฺญา.\csromanspacer{} เตน วุจฺจติ อสลฺลีนตฺตปหิตตฺตปคฺคหฏฺเฐ ปญฺญา วีริยารมฺเภ ญาณํ.}

\csromanheader{2}{วีริยารมฺภญาณนิทฺเทโส อฏฺฐติํสติโม.}
\csromansectionrule
\csromanmatikaanchor{mtk.32}
\csromantocmarkref{subsection}{39. อตฺถสนฺทสฺสนญาณนิทฺเทส}{mtk.32}
\bookmark[dest={mtk.32},level=2]{39. อตฺถสนฺทสฺสนญาณนิทฺเทส}
\csromanheader{2}{39. อตฺถสนฺทสฺสนญาณนิทฺเทส}
\proseitem{90}{กถํ นานาธมฺมปฺปกาสนตา ปญฺญา อตฺถสนฺทสฺสเน ญาณํ\csromandash{} \textbf{นานาธมฺมา}ติ ปญฺจกฺขนฺธา ทฺวาทสายตนานิ อฏฺฐารส ธาตุโย กุสลา ธมฺมา อกุสลา ธมฺมา อพฺยากตา ธมฺมา กามาวจรา ธมฺมา รูปาวจรา ธมฺมา อรูปาวจรา ธมฺมา อปริยาปนฺนา ธมฺมา.}

\prose{\textbf{ปกาสนตา}ติ รูปํ อนิจฺจโต ปกาเสติ, รูปํ ทุกฺขโต ปกาเสติ, รูปํ อนตฺตโต ปกาเสติ.\csromanspacer{} เวทนํ ฯเปฯ.\csromanspacer{} สญฺญํ.\csromanspacer{} สงฺขาร.\csromanspacer{} วิญฺญาณํ.\csromanspacer{} จกฺขุํ ฯเปฯ.\csromanspacer{} ชรามรณํ อนิจฺจโต ปกาเสติ, ชรามรณํ ทุกฺขโต ปกาเสติ, ชรามรณํ อนตฺตโต ปกาเสติ.}

\prose{\textbf{อตฺถสนฺทสฺสเน}ติ กามจฺฉนฺทํ ปชหนฺโต เนกฺขมฺมตฺถํ สนฺทสฺเสติ, พฺยาปาทํ ปชหนฺโต อพฺยาปาทตฺถํ สนฺทสฺเสติ, ถินมิทฺธํ ปชหนฺโต อาโลกสญฺญตฺถํ สนฺทสฺเสติ, อุทฺธจฺจํ ปชหนฺโต อวิกฺเขปตฺถํ สนฺทสฺเสติ, วิจิกิจฺฉํ ปชหนฺโต ธมฺมววตฺถานตฺถํ สนฺทสฺเสติ, อวิชฺชํ ปชหนฺโต ญาณตฺถํ สนฺทสฺเสติ, อรติํ ปชหนฺโต ปาโมชฺชตฺถํ สนฺทสฺเสติ, นีวรเณ ปชหนฺโต ปฐมฌานตฺถํ สนฺทสฺเสติ ฯเปฯ สพฺพกิเลเส ปชหนฺโต อรหตฺตมคฺคตฺถํ สนฺทสฺเสติ.\csromanspacer{} ตํญาตฏฺเฐน ญาณํ, ปชานนฏฺเฐน ปญฺญา.\csromanspacer{} เตน วุจฺจติ นานาธมฺมปกาสนตา ปญฺญา อตฺถสนฺทสฺสเน ญาณํ.}

\vspace{\tipitakaparskip}
\csromancenter{อตฺถสนฺทสฺสนญาณนิทฺเทโส นวติํสติโม.}
\csromanmatikaanchor{mtk.33}
\csromantocmarkref{subsection}{40. ทสฺสนวิสุทฺธิญาณนิทฺเทส}{mtk.33}
\bookmark[dest={mtk.33},level=2]{40. ทสฺสนวิสุทฺธิญาณนิทฺเทส}
\csromanheader{2}{40. ทสฺสนวิสุทฺธิญาณนิทฺเทส}
\proseitem{91}{\csromanfolio{102}กถํ สพฺพธมฺมานํ เอกสงฺคหตานานตฺเตกตฺตปฏิเวเธ ปญฺญา ทสฺสนวิสุทฺธิญาณํ\csromandash{}\textbf{สพฺพธมฺมานนฺ}ติ ปญฺจกฺขนฺธา ฯเปฯ อปริยาปนฺนา ธมฺมา.}

\prose{\textbf{เอกสงฺคหตา}ติ ทฺวาทสหิ อากาเรหิ สพฺเพ ธมฺมา เอกสงฺคหิตา, ตถฏฺเฐน อนตฺตฏฺเฐน สจฺจฏฺเฐน ปฏิเวธฏฺเฐน อภิชานนฏฺเฐน ปริชานนฏฺเฐน ธมฺมฏฺเฐน ธาตุฏฺเฐน ญาตฏฺเฐน สจฺฉิกิริยฏฺเฐน ผุสนฏฺเฐน อภิสมยฏฺเฐน.\csromanspacer{} อิเมหิ ทฺวาทสหิ อากาเรหิ สพฺเพ ธมฺมา เอกสงฺคหิตา.}

\prose{\textbf{นานตฺเตกตฺตนฺ}ติ กามจฺฉนฺโท นานตฺตํ, เนกฺขมฺมํ เอกตฺตํ ฯเปฯ สพฺพกิเลสา นานตฺตํ, อรหตฺตมคฺโค เอกตฺตํ.}

\prose{\textbf{ปฏิเวเธ}ติ ทุกฺขสจฺจํ ปริญฺญาปฏิเวธํ ปฏิวิชฺฌติ, สมุทยสจฺจํ ปหานปฏิเวธํ ปฏิวิชฺฌติ, นิโรธสจฺจํ สจฺฉิกิริยาปฏิเวธํ ปฏิวิชฺฌติ, มคฺคสจฺจํ ภาวนาปฏิเวธํ ปฏิวิชฺฌติ.}

\prose{\textbf{ทสฺสนวิสุทฺธี}ติ โสตาปตฺติมคฺคกฺขเณ ทสฺสนํ วิสุชฺฌติ, โสตาปตฺติผลกฺขเณ ทสฺสนํ วิสุทฺธํ.\csromanspacer{} สกทาคามิมคฺคกฺขเณ ทสฺสนํ วิสุชฺฌติ, สกทาคามิผลกฺขเณ ทสฺสนํ วิสุทฺธํ, อนาคามิมคฺคกฺขเณ ทสฺสนํ วิสุชฺฌติ, อนาคามิผลกฺขเณ ทสฺสนํ วิสุทฺธํ.\csromanspacer{} อรหตฺตมคฺคกฺขเณ ทสฺสนํ วิสุชฺฌติ, อรหตฺตผลกฺขเณ ทสฺสนํ วิสุทฺธํ.\csromanspacer{} ตํญาตฏฺเฐน ญาณํ ปชานนฏฺเฐน ปญฺญา.\csromanspacer{} เตน วุจฺจติ สพฺพธมฺมานํ เอกสงฺคหตานานตฺเตกตฺตปฏิเวเธ ปญฺญา ทสฺสนวิสุทฺธิญาณํ.}

\csromanheader{2}{ทสฺสนวิสุทฺธิญาณนนิทฺเทโส จตฺตาลีสโม.}
\csromansectionrule
\csromanmatikaanchor{mtk.34}
\csromantocmarkref{subsection}{41. ขนฺติญาณนิทฺเทส}{mtk.34}
\bookmark[dest={mtk.34},level=2]{41. ขนฺติญาณนิทฺเทส}
\csromanheader{2}{41. ขนฺติญาณนิทฺเทส}
\proseitem{92}{กถํ วิทิตตฺตา ปญฺญา ขนฺติญาณํ\csromandash{}รูปํ อนิจฺจโต วิทิตํ, รูปํ ทุกฺขโต วิทิตํ, รูปํ อนตฺตโต วิทิตํ.\csromanspacer{} ยํ ยํ วิทิตํ, ตํ ตํ ขมตีติ วิทิตตฺตา ปญฺญา ขนฺติญาณํ.\csromanspacer{} เวทนา ฯเปฯ.\csromanspacer{} สญฺญา.\csromanspacer{} สงฺขารา.\csromanspacer{} วิญฺญาณํ.\csromanspacer{} จกฺขุ ฯเปฯ.\csromanspacer{} ชรามรณํ อนิจฺจโต วิทิตํ, ชรามรณํ ทุกฺขโต วิทิตํ, อรามรณํ อนตฺตโต วิทิตํ.\csromanspacer{} ยํ ยํ วิทิตํ, ตํ ตํ ขมตีติ วิทิตตฺตา ปญฺญา ขนฺติญาณํ. \csromanfolio{103}ตํญาตฏฺเฐน ญาณํ, ปชานนฏฺเฐน ปญฺญา.\csromanspacer{} เตน วุจฺจติ วิทิตตฺตา ปญฺญา ขนฺติญาณํ.\csromanspacer{} ขนฺติญาณนิทฺเทโส เอกจตฺตาลีสโม.}

\csromanmatikaanchor{mtk.35}
\csromantocmarkref{subsection}{42. ปริโยคาหณญาณนิทฺเทส}{mtk.35}
\bookmark[dest={mtk.35},level=2]{42. ปริโยคาหณญาณนิทฺเทส}
\csromanheader{2}{42. ปริโยคาหณญาณนิทฺเทส}
\proseitem{93}{กถํ ผุฏฺฐตฺตา ปญฺญา ปริโยคาหเณ ญาณํ\csromandash{}รูปํ อนิจฺจโต ผุสติ, รูปํ ทุกฺขโต ผุสติ, รูปํ อนตฺตโต ผุสติ.\csromanspacer{} ยํ ยํ ผุสติ, ตํ ตํ ปริโยคหตีติ ผุฏฺฐตฺตา ปญฺญา ปริโยคาหเณ ญาณํ.\csromanspacer{} เวทนํ ฯเปฯ.\csromanspacer{} สญฺญํ.\csromanspacer{} สงฺขาเร.\csromanspacer{} วิญฺญาณํ.\csromanspacer{} จกฺขุํ ฯเปฯ.\csromanspacer{} ชรามรณํ อนิจฺจโต ผุสติ, ทุกฺขโต ผุสติ, อนตฺตโต ผุสติ.\csromanspacer{} ยํ ยํ ผุสติ, ตํ ตํ ปริโยคหตีติ ผุฏฺฐตฺตา ปญฺญา ปริโยคาหเณ ญาณํ.\csromanspacer{} ตํญาตฏฺเฐน ญาณํ.\csromanspacer{} ปชานนฏฺเฐน ปญฺญา.\csromanspacer{} เตน วุจฺจติ ผุฏฺฐตฺตา ปญฺญา ปริโยคาหเณ ญาณํ.}

\csromanheader{2}{ปริโยคาหณญาณนิทฺเทโส เทฺวจตฺตาลีสโม.}
\csromansectionrule
\csromanmatikaanchor{mtk.36}
\csromantocmarkref{subsection}{43. ปเทสวิหารญาณนิทฺเทส}{mtk.36}
\bookmark[dest={mtk.36},level=2]{43. ปเทสวิหารญาณนิทฺเทส}
\csromanheader{2}{43. ปเทสวิหารญาณนิทฺเทส}
\proseitem{94}{กถํ สโมทหเน ปญฺญา ปเทสวิหาเร ญาณํ\csromandash{}มิจฺฉาทิฏฺฐิ ปจฺจยาปิ เวทยิตํ, มิจฺฉาทิฏฺฐิวูปสมปจฺจยาปิ เวทยิตํ.\csromanspacer{} สมฺมาทิฏฺฐิปจฺจยาปิ เวทยิตํ, สมฺมาทิฏฺฐิวูปสมปจฺจยาปิ เวทยิตํ.\csromanspacer{} มิจฺฉาสงฺกปฺปปจฺจยาปิ เวทยิตํ, มิจฺฉาสงฺกปฺปวูปสมปจฺจยาปิ เวทยิตํ.\csromanspacer{} สมฺมาสงฺกปฺปปจฺจยาปิ เวทยิตํ, สมฺมาสงฺกปฺปวูปสมปจฺจยาปิ เวทยิตํ ฯเปฯ.\csromanspacer{} มิจฺฉาวิมุตฺติปจฺจยาปิ เวทยิตํ, มิจฺฉาวิมุตฺติวูปสมปจฺจยาปิ เวทยิตํ.\csromanspacer{} สมฺมาวิมุตฺติปจฺจยาปิ เวทยิตํ, สมฺมาวิมุตฺติวูปสมปจฺจยาปิ เวทยิตํ.\csromanspacer{} ฉนฺทปจฺจยาปิ เวทยิตํ, ฉนฺทวูปสมปจฺจยาปิ เวทยิตํ.\csromanspacer{} วิตกฺกปจฺจยาปิ เวทยิตํ.\csromanspacer{} วิตกฺกวูปสมปจฺจยาปิ เวทยิตํ.\csromanspacer{} สญฺญาปจฺจยาปิ เวทยิตํ, สญฺญาวูปสมปจฺจยาปิ เวทยิตํ.}

\prose{ฉนฺโท จ อวูปสนฺโต โหติ, วิตกฺโก จ อวูปสนฺโต โหติ, สญฺญา จ อวูปสนฺตา โหติ, ตปฺปจฺจยาปิ เวทยิตํ.\csromanspacer{} ฉนฺโท จ \csromanfolio{104}วูปสนฺโต โหติ, วิตกฺโก จ อวูปสนฺโต โหติ, สญฺญา จ อวูปสนฺตา โหติ, ตปฺปจฺจยาปิ เวทยิตํ.\csromanspacer{} ฉนฺโท จ วูปสนฺโต โหติ, วิตกฺโก จ วูปสนฺโต โหติ, สญฺญา จ อวูปสนฺตา โหติ, ตปฺปจฺจยาปิ เวทยิตํ.\csromanspacer{} ฉนฺโท จ วูปสนฺโต โหติ, วิตกฺโก จ วูปสนฺโต โหติ, สญฺญา จ วูปสนฺตา โหติ.\csromanspacer{} ตปฺปจฺจยาปิ เวทยิตํ.\csromanspacer{} อปฺปตฺตสฺส ปตฺติยา อตฺถิ อายวํ, ตสฺมิมฺปิ ฐาเน อนุปฺปตฺเต ตปฺปจฺจยาปิ เวทยิตํ.\csromanspacer{} ตํญาตฏฺเฐน ญาณํ, ปชานนฏฺเฐน ปญฺญา.\csromanspacer{} เตน วุจฺจติ สโมทหเน ปญฺญา ปเทสวิหาเร ญาณํ.}

\csromanheader{2}{ปเทสวิหารญาณนิทฺเทโส เตจตฺตาลีสโม.}
\csromansectionrule
\csromanmatikaanchor{mtk.37}
\csromantocmarkref{subsection}{44-49. ฉวิวฏฺฏญาณนิทฺเทส}{mtk.37}
\bookmark[dest={mtk.37},level=2]{44-49. ฉวิวฏฺฏญาณนิทฺเทส}
\csromanheader{2}{44-49. ฉวิวฏฺฏญาณนิทฺเทส}
\proseitem{95}{กถํ อธิปตตฺตา ปญฺญา สญฺญาวิวฏฺเฏ ญาณํ\csromandash{} เนกฺขมฺมาธิปตตฺตา ปญฺญา กามจฺฉนฺทโต สญฺญาย วิวฏฺฏตีติ อธิปตตฺตา ปญฺญา สญฺญาวิวฏฺเฏ ญาณํ, อพฺยาปาทาธิปตตฺตา ปญฺญา พฺยาปาทโต สญฺญาย วิวฏฺฏตีติ อธิปตตฺตา ปญฺญา สญฺญาวิวฏฺเฏ ญาณํ, อาโลกสญฺญาธิปตตฺตา ปญฺญา ถินมิทฺธโต สญฺญาย วิวฏฺฏตีติ อธิปตตฺตา ปญฺญา สญฺญาวิวฏฺเฏ ญาณํ, อวิกฺเขปาธิปตตฺตา ปญฺญา อุทฺธจฺจโต สญฺญาย วิวฏฺฏตีติ อธิปตตฺตา ปญฺญา สญฺญาวิวฏฺเฏ ญาณํ, ธมฺมววตฺถานาธิปตตฺตา ปญฺญา วิจิกิจฺฉาย สญฺญาย วิวฏฺฏตีติ อธิปตตฺตา ปญฺญา สญฺญาวิวฏฺเฏ ญาณํ, ญาณาธิปตตฺตา ปญฺญา อวิชฺชาย สญฺญาย วิวฏฺฏตีติ อธิปตตฺตา ปญฺญา สญฺญาวิวฏฺเฏ ญาณํ.\csromanspacer{} ปาโมชฺชาธิปตตฺตา ปญฺญา อรติยา สญฺญาย วิวฏฺฏตีติ อธิปตตฺตา ปญฺญา สญฺญาวิวฏฺเฏ ญาณํ, ปฐมชฺฌานาธิปตตฺตา ปญฺญา นีวรเณหิ สญฺญาย วิวฏฺฏตีติ อธิปตตฺตา ปญฺญา สญฺญาวิวฏฺเฏ ญาณํ ฯเปฯ อรหตฺตมคฺคาธิปตตฺตา ปญฺญา สพฺพกิเลเสหิ สญฺญาย วิวฏฺฏตีติ อธิปตตฺตา ปญฺญา สญฺญาวิวฏฺเฏ ญาณํ.\csromanspacer{} ตํญาตฏฺเฐน ญาณํ, ปชานนฏฺเฐน ปญฺญา.\csromanspacer{} เตน วุจฺจติ อธิปตตฺตา ปญฺญา สญฺญาวิวฏฺเฏ ญาณํ.}

\proseitem{96}{กถํ นานตฺเต ปญฺญา เจโตวิวฏฺเฏ ญาณํ\csromandash{}กามจฺฉนฺโท นานตฺตํ, เนกฺขมฺมํ เอกตฺตํ.\csromanspacer{} เนกฺขมฺเมกตฺตํ เจตยโต กามจฺฉนฺทโต จิตฺตํ วิวฏฺฏตีติ \csromanfolio{105}นานตฺเต ปญฺญา เจโตวิวฏฺเฏ ญาณํ.\csromanspacer{} พฺยาปาโท นานตฺตํ, อพฺยาปาโท เอกตฺตํ.\csromanspacer{} อพฺยาปาเทกตฺตํ เจตยโต พฺยาปาทโต จิตฺตํ วิวฏฺฏตีติ นานตฺเต ปญฺญา เจโตวิวฏฺเฏ ญาณํ.\csromanspacer{} ถินมิทฺธํ นานตฺตํ, อาโลกสญฺญา เอกตฺตํ.\csromanspacer{} อาโลกสญฺเญกตฺตํ เจตยโต ถินมิทฺธโต จิตฺตํ วิวฏฺฏตีติ นานตฺเต ปญฺญา เจโตวิวฏฺเฏ ญาณํ ฯเปฯ.\csromanspacer{} สพฺพกิเลสา นานตฺตํ, อรหตฺตมคฺโค เอกตฺตํ.\csromanspacer{} อรหตฺตมคฺเคกตฺตํ เจตยโต สพฺพกิเลเสหิ จิตฺตํ วิวฏฺฏตีติ นานตฺเต ปญฺญา เจโตวิวฏฺเฏ ญาณํ, ตํญาตฏฺเฐน ญาณํ, ปชานนฏฺเฐน ปญฺญา.\csromanspacer{} เตน วุจฺจติ นานตฺเต ปญฺญา เจโตวิวฏฺเฏ ญาณํ.}

\proseitem{97}{กถํ อธิฏฺฐาเน ปญฺญา จิตฺตวิวฏฺเฏ ญาณํ\csromandash{}กามจฺฉนฺทํ ปชหนฺโต เนกฺขมฺมวเสน จิตฺตํ อธิฏฺฐาตีติ อธิฏฺฐาเน ปญฺญา จิตฺตวิวฏฺเฏ ญาณํ.\csromanspacer{} พฺยาปาทํ ปชหนฺโต อพฺยาปาทวเสน จิตฺตํ อธิฏฺฐาตีติ อธิฏฺฐาเน ปญฺญา จิตฺตวิวฏฺเฏ ญาณํ.\csromanspacer{} ถินมิทฺธํ ปชหนฺโต อาโลกสญฺญาวเสน จิตฺตํ อธิฏฺฐาตีติ อธิฏฺฐาเน ปญฺญา จิตฺตวิวฏฺเฏ ญาณํ ฯเปฯ.\csromanspacer{} สพฺพกิเลเส ปชหนฺโต อรหตฺตมคฺควเสน จิตฺตํ อธิฏฺฐาตีติ อธิฏฺฐาเน ปญฺญา จิตฺตวิวฏฺเฏ ญาณํ, ตํ ญาตฏฺเฐน ญาณํ, ปชานนฏฺเฐน ปญฺญา.\csromanspacer{} เตน วุจฺจติ อธิฏฺฐาเน ปญฺญา จิตฺตวิวฏฺเฏ ญาณํ.}

\proseitem{98}{กถํ สุญฺญเต ปญฺญา ญาณวิวฏฺเฏ ญาณํ\csromandash{}จกฺขุ สุญฺญํ อตฺเตน วา อตฺตนิเยน วา นิจฺเจน วา ธุเวน วา สสฺสเตน วา อวิปริณามธมฺเมน วาติ ยถาภูตํ ชานโต\footnote{ปชานโต (สฺยา)} ปสฺสโต จกฺขาภินิเวสโต\footnote{กามาภินิเวสโต (สฺยา) อฏฺฐกถา โอโลเกตพฺพา.} ญาณํ วิวฏฺฏตีติ สุญฺญเต ปญฺญา ญาณวิวฏฺเฏ ญาณํ.\csromanspacer{} โสตํ สุญฺญํ ฯเปฯ.\csromanspacer{} ฆานํ สุญฺญํ.\csromanspacer{} ชิวฺหา สุญฺญา.\csromanspacer{} กาโย สุญฺโญ.\csromanspacer{} มโน สุญฺโญ อตฺเตน วา อตฺตนิเยน วา นิจฺเจน วา ธุเวน วา สสฺสเตน วา อวิปริณามธมฺเมน วาติ ยถาภูตํ ชานโต ปสฺสโต มนาภินิเวสโต ญาณํ วิวฏฺฏตีติ สุญฺญเต ปญฺญา ญาณวิวฏฺเฏ ญาณํ, ตํญาตฏฺเฐน ญาณํ, ปชานนฏฺเฐน ปญฺญา.\csromanspacer{} เตน วุจฺจติ สุญฺญเต ปญฺญา ญาณวิวฏฺเฏ ญาณํ.}

\proseitem{99}{กถํ โวสคฺเค ปญฺญา วิโมกฺขวิวฏฺเฏ ญาณํ\csromandash{}เนกฺขมฺเมน กามจฺฉนฺทํ โวสชฺชตีติ โวสคฺเค ปญฺญา วิโมกฺขวิวฏฺเฏ ญาณํ, อพฺยาปาเทน พฺยาปาทํ โวสชฺชตีติ โวสคฺเค ปญฺญา วิโมกฺขวิวฏฺเฏ ญาณํ, อาโลกสญฺญาย ถินมิทฺธํ โวสชฺชตีติ โวสคฺเค ปญฺญา วิโมกฺขวิวฏฺเฏ \csromanfolio{106}ญาณํ, อวิกฺเขเปน อุทฺธจฺจํ โวสชฺชตีติ โวสคฺเค ปญฺญา วิโมกฺขวิวฏฺเฏ ญาณํ, ธมฺมววตฺถาเนน วิจิกิจฺฉํ โวสชฺชตีติ โวสคฺเค ปญฺญา วิโมกฺขวิวฏฺเฏ ญาณํ ฯเปฯ อรหตฺตมคฺเคน สพฺพกิเลเส โวสชฺชตีติ โวสคฺเค ปญฺญา วิโมกฺขวิวฏฺเฏ ญาณํ, ตํญาตฏฺเฐน ญาณํ, ปชานนฏฺเฐน ปญฺญา.\csromanspacer{} เตน วุจฺจติ โวสคฺเค ปญฺญา วิโมกฺขวิวฏฺเฏ ญาณํ.}

\proseitem{100}{กถํ ตถฏฺเฐ ปญฺญา สจฺจวิวฏฺเฏ ญาณํ\csromandash{}ทุกฺขสฺส ปีฬนฏฺฐํ สงฺขตฏฺฐํ สนฺตาปฏฺฐํ วิปริณามฏฺฐํ ปริชานนฺโต วิวฏฺฏตีติ ตถฏฺเฐ ปญฺญา สจฺจวิวฏฺเฏ ญาณํ, สมุทยสฺส อายูหนฏฺฐํ นิทานฏฺฐํ สญฺโญคฏฺฐํ ปลิโพธฏฺฐํ ปชหนฺโต วิวฏฺฏตีติ ตถฏฺเฐ ปญฺญา สจฺจวิวฏฺเฏ ญาณํ, นิโรธสฺส นิสฺสรณฏฺฐํ วิเวกฏฺฐํ อสงฺขตฏฺฐํ อมตฏฺฐํ สจฺฉิกโรนฺโต วิวฏฺฏตีติ ตถฏฺเฐ ปญฺญา สจฺจวิวฏฺเฏ ญาณํ, มคฺคสฺส นิยฺยานฏฺฐํ เหตุฏฺฐํ ทสฺสนฏฺฐํ อาธิปเตยฺยฏฺฐํ ภาเวนฺโต วิวฏฺฏตีติ ตถฏฺเฐ ปญฺญา สจฺจวิวฏฺเฏ ญาณํ.}

\prose{สญฺญาวิวฏฺโฏ เจโตวิวฏฺโฏ จิตฺตวิวฏฺโฏ ญาณวิวฏฺโฏ วิโมกฺขวิวฏฺโฏ สจฺจวิวฏฺโฏ.\csromanspacer{} สญฺชานนฺโต วิวฏฺฏตีติ สญฺญาวิวฏฺโฏ, เจตยนฺโต วิวฏฺฏตีติ เจโตวิวฏฺโฏ, วิชานนฺโต วิวฏฺฏตีติ จิตฺตวิวฏฺโฏ, ญาณํ กโรนฺโต วิวฏฺฏตีติ ญาณวิวฏฺโฏ, โวสชฺชนฺโต วิวฏฺฏตีติ วิโมกฺขวิวฏฺโฏ, ตถฏฺเฐ วิวฏฺฏตีติ สจฺจวิวฏฺโฏ.}

\prose{ยตฺถ สญฺญาวิวฏฺโฏ, ตตฺถ เจโตวิวฏฺโฏ.\csromanspacer{} ยตฺถ เจโตวิวฏฺโฏ, ตตฺถ สญฺญาวิวฏฺโฏ.\csromanspacer{} ยตฺถ สญฺญาวิวฏฺโฏ เจโตวิวฏฺโฏ, ตตฺถ จิตฺตวิวฏฺโฏ.\csromanspacer{} ยตฺถ จิตฺตวิวฏฺโฏ, ตตฺถ สญฺญาวิวฏฺโฏ เจโตวิวฏฺโฏ.\csromanspacer{} ยตฺถ สญฺญาวิวฏฺโฏ เจโตวิวฏฺโฏ จิตฺตวิวฏฺโฏ, ตตฺถ ญาณวิวฏฺโฏ.\csromanspacer{} ยตฺถ ญาณวิวฏฺโฏ, ตตฺถ สญฺญาวิวฏฺโฏ เจโตวิวฏฺโฏ จิตฺตวิวฏฺโฏ.\csromanspacer{} ยตฺถ สญฺญาวิวฏฺโฏ เจโตวิวฏฺโฏ จิตฺตวิวฏฺโฏ ญาณวิวฏฺโฏ, ตตฺถ วิโมกฺขวิวฏฺโฏ.\csromanspacer{} ยตฺถ วิโมกฺขวิวฏฺโฏ, ตตฺถ สญฺญาวิวฏฺโฏ เจโตวิวฏฺโฏ จิตฺตวิวฏฺโฏ ญาณวิวฏฺโฏ.\csromanspacer{} ยตฺถ สญฺญาวิวฏฺโฏ เจโตวิวฏฺโฏ จิตฺตวิวฏฺโฏ ญาณวิวฏฺโฏ วิโมกฺขวิวฏฺโฏ, ตตฺถ สจฺจวิวฏฺโฏ.\csromanspacer{} ยตฺถ สจฺจวิวฏฺโฏ, ตตฺถ สญฺญาวิวฏฺโฏ เจโตวิวฏฺโฏ จิตฺตวิวฏฺโฏ ญาณวิวฏฺโฏ วิโมกฺขวิวฏฺโฏ.\csromanspacer{} ตํญาตฏฺเฐน ญาณํ, ปชานนฏฺเฐน ปญฺญา.\csromanspacer{} เตน วุจฺจติ ตถฏฺเฐ ปญฺญา สจฺจวิวฏฺเฏ ญาณํ.}

\vspace{\tipitakaparskip}
\csromancenter{ฉวิวฏฺฏญาณนิทฺเทโส นวจตฺตาลีสโม.}
\csromanmatikaanchor{mtk.38}
\csromantocmarkref{subsection}{50. อิทฺธิวิธญาณนิทฺเทส}{mtk.38}
\bookmark[dest={mtk.38},level=2]{50. อิทฺธิวิธญาณนิทฺเทส}
\csromanheader{2}{50. อิทฺธิวิธญาณนิทฺเทส}
\proseitem{101}{\csromanfolio{107}กถํ กายมฺปิ จิตฺตมฺปิ เอกววตฺถานตา สุขสญฺญญฺจ ลหุสญฺญญฺจ อธิฏฺฐานวเสน อิชฺฌนฏฺเฐ ปญฺญา อิทฺธิวิเธ ญาณํ\csromandash{}อิธ ภิกฺขุ ฉนฺทสมาธิปธานสงฺขารสมนฺนาคตํ อิทฺธิปาทํ ภาเวติ, วีริยสมาธิปธานสงฺขารสมนฺนาคตํ อิทฺธิปาทํ ภาเวติ, จิตฺตสมาธิปธานสงฺขารสมนฺนาคตํ อิทฺธิปาทํ ภาเวติ, วีมํสาสมาธิปธานสงฺขารสมนฺนาคตํ อิทฺธิปาทํ ภาเวติ.\csromanspacer{} โส อิเมสุ จตูสุ อิทฺธิปาเทสุ จิตฺตํ ปริภาเวติ ปริทเมติ, มุทุํ กโรติ กมฺมนิยํ.\csromanspacer{} โส อิเมสุ จตูสุ อิทฺธิปาเทสุ จิตฺตํ ปริภาเวตฺวา ปริทเมตฺวา มุทุํ กริตฺวา กมฺมนิยํ กายมฺปิ จิตฺเต สโมทหติ, จิตฺตมฺปิ กาเย สโมทหติ, กายวเสน จิตฺตํ ปริณาเมติ, จิตฺตวเสน กายํ ปริณาเมติ, กายวเสน จิตฺตํ อธิฏฺฐาติ, จิตฺตวเสน กายํ อธิฏฺฐาติ, กายวเสน จิตฺตํ ปริณาเมตฺวา จิตฺตวเสน กายํ ปริณาเมตฺวา กายวเสน จิตฺตํ อธิฏฺฐหิตฺวา จิตฺตวเสน กายํ อธิฏฺฐหิตฺวา สุขสญฺญญฺจ ลหุสญฺญญฺจ กาเย โอกฺกมิตฺวา วิหรติ, โส ตถาภาวิเตน จิตฺเตน ปริสุทฺเธน ปริโยทาเตน อิทฺธิวิธญาณาย จิตฺตํ อภินีหรติ อภินินฺนาเมติ, โส อเนกวิหิตํ อิทฺธิวิธํ ปจฺจนุโภติ.}

\proseitem{102}{เอโกปิ หุตฺวา พหุธา โหติ, พหุธาปิ หุตฺวา เอโก โหติ, อาวิภาวํ ติโรภาวํ ติโรกุฏฺฏํ\footnote{ติโรกุฑฺฑํ (สฺยา)} ติโรปาการํ ติโรปพฺพตํ อสชฺชมาโน คจฺฉติ เสยฺยถาปิ อากาเส.\csromanspacer{} ปถวิยาปิ อุมฺมุชฺชนิมุชฺชํ กโรติ เสยฺยถาปิ อุทเก.\csromanspacer{} อุทเกปิ อภิชฺชมาเน คจฺฉติ เสยฺยถาปิ ปถวิยํ.\csromanspacer{} อากาเสปิ ปลฺลงฺเกน กมติ\footnote{จงฺกมติ (สฺยา) ที 1. 74 ปิฏฺเฐ ปสฺสิตพฺพา.} เสยฺยถาปิ ปกฺขี สกุโณ.\csromanspacer{} อิเมปิ จนฺทิมสูริเย เอวํมหิทฺธิเก เอวํมหานุภาเว ปาณินา ปรามสติ\footnote{ปริมสติ (สฺยา)} ปริมชฺชติ, ยาว พฺรหฺมโลกาปิ กาเยน วสํ วตฺเตติ.\csromanspacer{} ตํญาตฏฺเฐน ญาณํ, ปชานนฏฺเฐน ปญฺญา.\csromanspacer{} เตน วุจฺจติ กามมฺปิ จิตฺตมฺปิ เอกววตฺถานตา สุขสญฺญญฺจ ลหุสญฺญญฺจ อธิฏฺฐานวเสน อิชฺฌนฏฺเฐ ปญฺญา อิทฺธิวิเธ ญาณํ.}

\vspace{\tipitakaparskip}
\csromancenter{อิทฺธิวิธญาณนิทฺเทโส ปญฺญาสโม.}
\csromanmatikaanchor{mtk.39}
\csromantocmarkref{subsection}{51. โสตธาตุวิสุทฺธิญาณนิทฺเทส}{mtk.39}
\bookmark[dest={mtk.39},level=2]{51. โสตธาตุวิสุทฺธิญาณนิทฺเทส}
\csromanheader{2}{51. โสตธาตุวิสุทฺธิญาณนิทฺเทส}
\proseitem{103}{\csromanfolio{108}กถํ วิตกฺกวิปฺผารวเสน นานตฺเตกตฺตสทฺทนิมิตฺตานํ ปริโยคาหเณ ปญฺญา โสตธาตุวิสุทฺธิญาณํ\csromandash{}อิธ ภิกฺขุ ฉนฺทสมาธิ ฯเปฯ วีริยสมาธิ.\csromanspacer{} จิตฺตสมาธิ.\csromanspacer{} วีมํสาสมาธิปธานสงฺขาร- สมนฺนาคตํ อิทฺธิปาทํ ภาเวติ.\csromanspacer{} โส อิเมสุ จตูสุ อิทฺธิปาเทสุ จิตฺตํ ปริภาเวติ ปริทเมติ, มุทุํ กโรติ กมฺมนิยํ.\csromanspacer{} โส อิเมสุ จตูสุ อิทฺธิปาเทสุ จิตฺตํ ปริภาเวตฺวา ปริทเมตฺวา มุทุํ กริตฺวา กมฺมนิยํ ทูเรปิ สทฺทานํ สทฺทนิมิตฺตํ มนสิ กโรติ, สนฺติเกปิ สทฺทานํ สทฺทนิมิตฺตํ มนสิ กโรติ, โอฬาริกานมฺปิ สทฺทานํ สทฺทนิมิตฺตํ มนสิ กโรติ, สุขุมานมฺปิ สทฺทานํ สทฺทนิมิตฺตํ มนสิ กโรติ, สณฺหสณฺหานมฺปิ สทฺทานํ สทฺทนิมิตฺตํ มนสิ กโรติ, ปุรตฺถิมายปิ ทิสาย สทฺทานํ สทฺทนิมิตฺตํ มนสิ กโรติ, ปจฺฉิมายปิ ทิสาย สทฺทานํ สทฺทนิมิตฺตํ มนสิ กโรติ, อุตฺตรายปิ ทิสาย สทฺทานํ สทฺทนิมิตฺตํ มนสิ กโรติ, ทกฺขิณายปิ ทิสาย สทฺทานํ สทฺทนิมิตฺตํ มนสิ กโรติ, ปุรตฺถิมายปิ อนุทิสาย สทฺทานํ สทฺทนิมิตฺตํ มนสิ กโรติ, ปจฺฉิมายปิ อนุทิสาย สทฺทานํ สทฺทนิมิตฺตํ มนสิ กโรติ, อุตฺตรายปิ อนุทิสาย สทฺทานํ สทฺทนิมิตฺตํ มนสิ กโรติ, ทกฺขิณายปิ อนุทิสาย สทฺทานํ สทฺทนิมิตฺตํ มนสิ กโรติ, เหฏฺฐิมยปิ ทิสาย สทฺทานํ สทฺทนิมิตฺตํ มนสิ กโรติ, อุปริมายปิ ทิสาย สทฺทานํ สทฺทนิมิตฺตํ มนสิ กโรติ.\csromanspacer{} โส ตถาภาวิเตน จิตฺเตน ปริสุทฺเธน ปริโยทาเตน โสตธาตุวิสุทฺธิญาณาย จิตฺตํ อภินีหรติ อภินินฺนาเมติ, โส ทิพฺพาย โสตธาตุยา วิสุทฺธาย อติกฺกนฺตมานุสิกาย อุโภ สทฺเท สุณาติ ทิพฺเพ จ มานุเส จ เย ทูเร สนฺติเก จ.\csromanspacer{} ตํญาตฏฺเฐน ญาณํ, ปชานนฏฺเฐน ปญฺญา.\csromanspacer{} เตน วุจฺจติ วิตกฺกวิปฺผารวเสน นานตฺเตกตฺตสทฺทนิมิตฺตานํ ปริโยคาหเณ ปญฺญา โสตธาตุวิสุทฺธิญาณํ.}

\csromanheader{2}{โสตธาตุวิสุทฺธิญาณนิทฺเทโส เอกปญฺญาสโม.}
\csromansectionrule
\csromanmatikaanchor{mtk.40}
\csromantocmarkref{subsection}{52. เจโตปริยญาณนิทฺเทส}{mtk.40}
\bookmark[dest={mtk.40},level=2]{52. เจโตปริยญาณนิทฺเทส}
\csromanheader{2}{52. เจโตปริยญาณนิทฺเทส}
\proseitem{104}{กถํ ติณฺณํ จิตฺตานํ วิปฺผารตฺตา อินฺทฺริยานํ ปสาทวเสน นานตฺเตกตฺตวิญฺญาณจริยาปริโยคาหเณ ปญฺญา เจโตปริยญาณํ\csromandash{}อิธ \csromanfolio{109}ภิกฺขุ ฉนฺทสมาธิปธานสงฺขารสมนฺนาคตํ อิทฺธิปาทํ ภาเวติ, วีริยสมาธฺ ฯเปฯ จิตฺตสมาธิ ฯเปฯ วีมํสาสมาธิปธานสงฺขาร- สมนฺนาคตํ อิทฺธิปาทํ ภาเวติ.\csromanspacer{} โส อิเมสุ จตูสุ อิทฺธิปาเทสุ จิตฺตํ ปริภาเวติ ปริทเมติ, มุทุํ กโรติ กมฺมนิยํ.\csromanspacer{} โส อิเมสุ จตูสุ อิทฺธิปาเทสุ จิตฺตํ ปริภาเวตฺวา ปริทเมตฺวา มุทุํ กริตฺวา กมฺมนิยํ เอวํ ปชานาติ “อิทํ รูปํ โสมนสฺสินฺทฺริยสมุฏฺฐิตํ, อิทํ รูปํ โทมนสฺสินฺทฺริยสมุฏฺฐิตํ, อิทํ รูปํ อุเปกฺขินฺทฺริยสมุฏฺฐิตนฺ”ติ.\csromanspacer{} โส ตถาภาวิเตน จิตฺเตน ปริสุทฺเธน ปริโยทาเตน เจโตปริยญาณาย จิตฺตํ อภินีหรติ อภินินฺนาเมติ.\csromanspacer{} โส ปรสตฺตานํ ปรปุคฺคลานํ เจตสา เจโต ปริจฺจ ปชานาติ สราคํ วา จิตฺตํ “สราคํ จิตฺตนฺ”ติ ปชานาติ, วีตราคํ วา จิตฺตํ “วีตราคํ จิตฺตนฺ”ติ ปชานาติ, สโทสํ วา จิตฺตํ ฯเปฯ วีตโทสํ วา จิตฺตํ.\csromanspacer{} สโมหํ วา จิตฺตํ.\csromanspacer{} วีตโมหํ วา จิตฺตํ.\csromanspacer{} สํขิตฺตํ วา จิตฺตํ.\csromanspacer{} วิกฺขิตฺตํ วา จิตฺตํ.\csromanspacer{} มหคฺคตํ วา จิตฺตํ.\csromanspacer{} อมหคฺคตํ วา จิตฺตํ.\csromanspacer{} ส-อุตฺตรํ วา จิตฺตํ.\csromanspacer{} อนุตฺตรํ วา จิตฺตํ.\csromanspacer{} สมาหิตํ วา จิตฺตํ.\csromanspacer{} อสมาหิตํ วา จิตฺตํ.\csromanspacer{} วิมุตฺตํ วา จิตฺตํ.\csromanspacer{} อวิมุตฺตํ วา จิตฺตํ “อวิมุตฺตํ จิตฺตนฺ”ติ ปชานาติ.\csromanspacer{} ตํญาตฏฺเฐน ญาณํ, ปชานนฏฺเฐน ปญฺญา.\csromanspacer{} เตน วุจฺจติ ติณฺณํ จิตฺตานํ วิปฺผารตฺตา อินฺทฺริยานํ ปสาทวเสน นานตฺเตกตฺตวิญฺญาณจริยาปริโยคาหเณ ปญฺญา เจโตปริยญาณํ.}

\csromanheader{2}{เจโตปริยญาณนิทฺเทโส เทฺวปญฺญาสโม.}
\csromansectionrule
\csromanmatikaanchor{mtk.41}
\csromantocmarkref{subsection}{53. ปุพฺเพนิวาสานุสฺสติญาณนิทฺเทส}{mtk.41}
\bookmark[dest={mtk.41},level=2]{53. ปุพฺเพนิวาสานุสฺสติญาณนิทฺเทส}
\csromanheader{2}{53. ปุพฺเพนิวาสานุสฺสติญาณนิทฺเทส}
\proseitem{105}{กถํ ปจฺจยปวตฺตานํ ธมฺมานํ นานตฺเตกตฺตกมฺมวิปฺผารวเสน ปริโยคาหเณ ปญฺญา ปุพฺเพนิวาสานุสฺสติญาณํ\csromandash{}อิธ ภิกฺขุ ฉนฺทสมาธิ ฯเปฯ มุทุํ กริตฺวา กมฺมนิยํ เอวํ ปชานาติ “อิมสฺมิํ สติ อิทํ โหติ, อิมสฺสุปฺปาทา อิทํ อุปฺปชฺชติ ยทิทํ อวิชฺชาปจฺจยา สงฺขารา, สงฺขารปจฺจยา วิญฺญาณํ, วิญฺญาณปจฺจยา นามรูปํ, นามรูปปจฺจยา สฬายตนํ, สฬายตนปจฺจยา ผสฺโส, ผสฺสปจฺจยา เวทนา, เวทนาปจฺจยา ตณฺหา, ตณฺหาปจฺจยา อุปาทานํ, อุปาทานปจฺจยา ภโว, ภวปจฺจยา ชาติ, ชาติปจฺจยา ชรามรณํ โสกปริเทวทุกฺขโทมนสฺสุปายาสา สมฺภวนฺติ, เอวเมตสฺส เกวลสฺส ทุกฺขกฺขนฺธสฺส สมุทโย โหติ.}

\prose{\csromanfolio{110}โส ตถาภาวิเตน จิตฺเตน ปริสุทฺเธน ปริโยทาเตน ปุพฺเพนิวาสานุสฺสติญาณาย จิตฺตํ อภินีหรติ อภินินฺนาเมติ.\csromanspacer{} โส อเนกวิหิตํ ปุพฺเพนิวาสํ อนุสฺสรติ.\csromanspacer{} เสยฺยถิทํ, เอกมฺปิ ชาติํ เทฺวปิ ชาติโย ติสฺโสปิ ชาติโย จตสฺโสปิ ชาติโย ปญฺจปิ ชาติโย ทสปิ ชาติโย วีสมฺปิ ชาติโย ติํสมฺปิ ชาติโย จตฺตาลีสมฺปิ ชาติโย ปญฺญาสมฺปิ ชาติโย, ชาติสตมฺปิ ชาติสหสฺสมฺปิ ชาติสตสหสฺสมฺปิ, อเนเกปิ สํวฏฺฏกปฺเป อเนเกปิ วิวฏฺฏกปฺเป อเนเกปิ สํวฏฺฏวิวฏฺฏกปฺเป, อมุตฺราสิํ เอวํนาโม เอวํโคตฺโต เอวํวณฺโณ เอวมาหาโร เอวํสุขทุกฺขปฺปฏิสํเวที เอวมายุปริยนฺโต.\csromanspacer{} โส ตโต จุโต อมุตฺร อุทปาทิํ.\csromanspacer{} ตตฺราปาสิํ เอวํนาโม เอวํโคตฺโต เอวํวณฺโณ เอวมาหาโร เอวํสุขทุกฺขปฺปฏิสํเวที เอวมายุปริยนฺโต.\csromanspacer{} โส ตโต จุโต อิธูปปนฺโน”ติ, อิติ สาการํ ส-อุทฺเทสํ อเนกวิหิตํ ปุพฺเพนิวาสํ อนุสฺสรติ.\csromanspacer{} ตํญาตฏฺเฐน ญาณํ, ปชานนฏฺเฐน ปญฺญา.\csromanspacer{} เตน วุจฺจติ ปจฺจยปวตฺตานํ ธมฺมานํ นานตฺเตกตฺตกมฺมวิปฺผารวเสน ปริโยคาหเณ ปญฺญา ปุพฺเพนิวาสานุสฺสติญาณํ.}

\csromanheader{2}{ปุพฺเพนิวาสานุสฺสติญาณนิทฺเทโส เตปญฺญาสโม.}
\csromansectionrule
\csromanmatikaanchor{mtk.42}
\csromantocmarkref{subsection}{54. ทิพฺพจกฺขุญาณนิทฺเทส}{mtk.42}
\bookmark[dest={mtk.42},level=2]{54. ทิพฺพจกฺขุญาณนิทฺเทส}
\csromanheader{2}{54. ทิพฺพจกฺขุญาณนิทฺเทส}
\proseitem{106}{กถํ โอภาสวเสน นานตฺเตกตฺตรูปนิมิตฺตานํ ทสฺสนฏฺเฐ ปญฺญา ทิพฺพจกฺขุญาณํ\csromandash{}อิธ ภิกฺขุ ฉนฺทสมาธิปธานสงฺขารสมนฺนาคตํ อิทฺธิปาทํ ภาเวติ.\csromanspacer{} วีริยสมาธิ ฯเปฯ.\csromanspacer{} จิตฺตสมาธิ ฯเปฯ.\csromanspacer{} วีมํสาสมาธิปธานสงฺขารสมนฺนาคตํ อิทฺธิปาทํ ภาเวติ.\csromanspacer{} โส อิเมสุ จตูสุ อิทฺธิปาเทสุ จิตฺตํ ปริภาเวติ ปริทเมติ, มุทุํ กโรติ กมฺมนิยํ.\csromanspacer{} โส อิเมสุ จตูสุ อิทฺธิปาเทสุ จิตฺตํ ปริภาเวตฺวา ปริทเมตฺวา มุทุํ กริตฺวา กมฺมนิยํ อาโลกสญฺญํ มนสิ กโรติ, ทิวาสญฺญํ อธิฏฺฐาติ, ยถา ทิวา ตถา รตฺติํ, ยถา รตฺติํ ตถา ทิวา.\csromanspacer{} อิติ วิวเฏน เจตสา อปริโยนทฺเธน สปฺปภาสํ จิตฺตํ ภาเวติ.\csromanspacer{} โส ตถาภาวิเตน จิตฺเตน ปริสุทฺเธน ปริโยทาเตน สตฺตานํ จุตูปปาตญาณาย จิตฺตํ อภินีหรติ อภินินฺนาเมติ.\csromanspacer{} โส ทิพฺเพน จกฺขุนา วิสุทฺเธน อติกฺกนฺตมานุสเกน สตฺเต ปสฺสติ จวมาเน อุปปชฺชมาเน หีเน ปณีเต สุวณฺเณ \csromanfolio{111}ทุพฺพณฺเณ สุคเต ทุคฺคเต, ยถากมฺมูปเค สตฺเต ปชานาติ “อิเม วต โภนฺโต สตฺตา กายทุจฺจริเตน สมนฺนาคตา, วจีทุจฺจริเตน สมนฺนาคตา, มโนทุจฺจริเตน สมนฺนาคตา, อริยานํ อุปวาทกา, มิจฺฉาทิฏฺฐิกา มิจฺฉาทิฏฺฐิกมฺมสมาทานา, เต กายสฺส เภทา ปรํ มรณา อปายํ ทุคฺคติํ วินิปาตํ นิรยํ อุปปนฺนา.\csromanspacer{} อิเม วา ปน โภนฺโต สตฺตา กายสุจริเตน สมนฺนาคตา, วจีสุจริเตน สมนฺนาคตา, มโนสุจริเตน สมนฺนาคตา, อริยานํ อนุปวาทกา, สมฺมาทิฏฺฐิกา สมฺมาทิฏฺฐิกมฺมสมาทานา, เต กายสฺส เภทา ปรํ มรณา สุคติํ สคฺคํ โลกํ อุปปนฺนา”ติ, อิติ ทิพฺเพน จกฺขุนา วิสุทฺเธน อติกฺกนฺตมานุสเกน สตฺเต ปสฺสติ จวมาเน อุปปชฺชมาเน หีเน ปณีเต สุวณฺเณ ทุพฺพณฺเณ สุคเต ทุคฺคเต, ยถากมฺมูปเค สตฺเต ปชานาติ.\csromanspacer{} ตํญาตฏฺเฐน ญาณํ, ปชานนฏฺเฐน ปญฺญา.\csromanspacer{} เตน วุจฺจติ โอภาสวเสน นานตฺเตกตฺตรูปนิมิตฺตานํ ทสฺสงฏฺเฐ ปญฺญา ทิพฺพจกฺขุญาณํ.}

\csromanheader{2}{ทิพฺพจกฺขุญาณนิทฺเทโส จตุปญฺญาสโม.}
\csromansectionrule
\csromanmatikaanchor{mtk.43}
\csromantocmarkref{subsection}{55. อาสวกฺขยญาณนิทฺเทส}{mtk.43}
\bookmark[dest={mtk.43},level=2]{55. อาสวกฺขยญาณนิทฺเทส}
\csromanheader{2}{55. อาสวกฺขยญาณนิทฺเทส}
\proseitem{107}{กถํ จตุสฏฺฐิยา อากาเรหิ ติณฺณนฺนํ อินฺทฺริยานํ วสิภาวตา ปญฺญา อาสวานํ ขเย ญาณํ\csromandash{}กตเมสํ ติณฺณนฺนํ อินฺทฺริยานํ? อนญฺญาตญฺญสฺสามีตินฺทฺริยสฺส อญฺญินฺทฺริยสฺส อญฺญาตาวินฺทฺริยสฺส.}

\prose{อนญฺญาตญฺญสฺสามีตินฺทฺริยํ กติ ฐานานิ คจฺฉติ, อญฺญินฺทฺริยํ กติ ฐานานิ คจฺฉติ, อญฺญาตาวินฺทฺริยํ กติ ฐานานิ คจฺฉติ? อนญฺญาตญฺญสฺสามีตินฺทฺริยํ เอกํ ฐานํ คจฺฉติ, โสตาปตฺติมคฺคํ.\csromanspacer{} อญฺญินฺทฺริยํ ฉ ฐานานิ คจฺฉติ, โสตาปตฺติผลํ สกทาคามิมคฺคํ สกทาคามิผลํ อนาคามิมคฺคํ อนาคามิผลํ อรหตฺตมคฺคํ.\csromanspacer{} อญฺญาตาวินฺทฺริยํ เอกํ ฐานํ คจฺฉติ, อรหตฺตผลํ.}

\prose{โสตาปตฺติมคฺคกฺขเณ อนญฺญาตญฺญสฺสามีตินฺทฺริยสฺส สทฺธินฺทฺริยํ อธิโมกฺขปริวารํ โหติ, วีริยินฺทฺริยํ ปคฺคหปริวารํ โหติ, สตินฺทฺริยํ อุปฏฺฐานปริวารํ โหติ, สมาธินฺทฺริยํ อวิกฺเขปปริวารํ โหติ, ปญฺญินฺทฺริยํ ทสฺสนปริวารํ โหติ, มนินฺทฺริยํ วิชานนปริวารํ โหติ, โสมนสฺสินฺทฺริยํ \csromanfolio{112}อภิสนฺทนปริวารํ โหติ, ชีวิตินฺทฺริยํ ปวตฺตสนฺตตาธิปเตยฺยปริวารํ โหติ.\csromanspacer{} โสตาปตฺติมคฺคกฺขเณ ชาตา ธมฺมา ฐเปตฺวา จิตฺตสมุฏฺฐานํ รูปํ สพฺเพว กุสลา โหนฺติ, สพฺเพว อนาสวา โหนฺติ, สพฺเพว นิยฺยานิกา โหนฺติ, สพฺเพว อปจยคามิโน โหนฺติ, สพฺเพว โลกุตฺตรา โหนฺติ, สพฺเพว นิพฺพานารมฺมณา โหนฺติ.\csromanspacer{} โสตาปตฺติมคฺคกฺขเณ อนญฺญาตญฺญสฺสามีตินฺทฺริยสฺส อิมานิ อฏฺฐินฺทฺริยานิ สหชาตปริวารา โหนฺติ, อญฺญมญฺญปริวารา โหนฺติ, นิสฺสยปริวารา โหนฺติ, สมฺปยุตฺตปริวารา โหนฺติ, สหคตา โหนฺติ, สหชาตา โหนฺติ, สํสฏฺฐา โหนฺติ, สมฺปยุตฺตา โหนฺติ.\csromanspacer{} เตว ตสฺส อาการา เจว โหนฺติ ปริวารา จ.}

\prose{โสตาปตฺติผลกฺขเณ อญฺญินฺทฺริยสฺส สทฺธินฺทฺริยํ อธิโมกฺขปริวารํ โหติ, วีริยินฺทฺริยํ ปคฺคหปริวารํ โหติ, สตินฺทฺริยํ อุปฏฺฐานปริวารํ โหติ, สมาธินฺทฺริยํ อวิกฺเขปปริวารํ โหติ, ปญฺญินฺทฺริยํ ทสฺสนปริวารํ โหติ, มนินฺทฺริยํ วิชานนปริวารํ โหติ, โสมนสฺสินฺทฺริยํ อภิสนฺทนปริวารํ โหติ, ชีวิตินฺทฺริยํ ปวตฺตสนฺตตาธิปเตยฺยปริวารํ โหติ.\csromanspacer{} โสตาปตฺติผลกฺขเณ ชาตา ธมฺมา สพฺเพว อพฺยากตา โหนฺติ, ฐเปตฺวา จิตฺตสมุฏฺฐานํ รูปํ สพฺเพว อนาสวา โหนฺติ, สพฺเพว โลกุตฺตรา โหนฺติ, สพฺเพว นิพฺพานารมฺมณา โหนฺติ.\csromanspacer{} โสตาปตฺติผลกฺขเณ อญฺญินฺทฺริยสฺส อิมานิ อฏฺฐินฺทฺริยานิ สหชาตปริวารา โหนฺติ, อญฺญมญฺญปริวารา โหนฺติ, นิสฺสยปริวารา โหนฺติ, สมฺปยุตฺตปริวารา โหนฺติ, สหคตา โหนฺติ, สหชาตา โหนฺติ, สํสฏฺฐา โหนฺติ, สมฺปยุตฺตา โหนฺติ, เตว ตสฺส อาการา เจว โหนฺติ ปริวารา จ.}

\prose{สกทาคามิมคฺคกฺขเณ ฯเปฯ.\csromanspacer{} สกทาคามิผลกฺขเณ ฯเปฯ.\csromanspacer{} อนาคามิมคฺคกฺขเณ ฯเปฯ.\csromanspacer{} อนาคามิผลกฺขเณ ฯเปฯ.\csromanspacer{} อรหตฺตมคฺคกฺขเณ อญฺญินฺทฺริยสฺส สทฺธินฺทฺริยํ อธิโมกฺขปริวารํ โหติ ฯเปฯ.ชีวิตินฺทฺริยํ ปวตฺตสนฺตตาธิปเตยฺยปริวารํ โหติ.\csromanspacer{} อรหตฺตมคฺคกฺขเณ ชาตา ธมฺมา ฐเปตฺวา จิตฺตสมุฏฺฐานํ รูปํ สพฺเพว กุสลา โหนฺติ, สพฺเพว อนาสวา โหนฺติ, สพฺเพว นิยฺยานิกา โหนฺติ, สพฺเพว อปจยคามิโน โหนฺติ, สพฺเพว โลกุตฺตรา โหนฺติ, สพฺเพว นิพฺพานารมฺมณา โหนฺติ.\csromanspacer{} อรหตฺตมคฺคกฺขเณ อญฺญินฺทฺริยสฺส อิมานิ อฏฺฐินฺทฺริยานิ สหชาตปริวารา โหนฺติ, อญฺญมญฺญปริวารา โหนฺติ, นิสฺสยปริวารา โหนฺติ, สมฺปยุตฺตปริวารา โหนฺติ, \csromanfolio{113}สหคตา โหนฺติ, สหชาตา โหนฺติ, สํสฏฺฐา โหนฺติ, สมฺปยุตฺตา โหนฺติ, เตว ตสฺส อาการา เจว โหนฺติ ปริวารา จ.}

\prose{อรหตฺตผลกฺขเณ อญฺญาตาวินฺทฺริยสฺส สทฺธินฺทฺริยํ อธิโมกฺขปริวารํ โหติ, วีริยินฺทฺริยํ ปคฺคหปริวารํ โหติ, สตินฺทฺริยํ อุปฏฺฐานปริวารํ โหติ, สมาธินฺทฺริยํ อวิกฺเขปปริวารํ โหติ, ปญฺญินฺทฺริยํ ทสฺสนปริวารํ โหติ, มนินฺทฺริยํ วิชานนปริวารํ โหติ, โสมนสฺสินฺทฺริยํ อภิสนฺทนปริวารํ โหติ, ชีวิตินฺทฺริยํ ปวตฺตสนฺตตาธิปเตยฺยปริวารํ โหติ.\csromanspacer{} อรหตฺตผลกฺขเณ ชาตา ธมฺมา สพฺเพว อพฺยากตา โหนฺติ, ฐเปตฺวา จิตฺตสมุฏฺฐานํ รูปํ สพฺเพว อนาสวา โหนฺติ, สพฺเพว โลกุตฺตรา โหนฺติ, สพฺเพว นิพฺพานารมฺมณา โหนฺติ.\csromanspacer{} อรหตฺตผลกฺขเณ อญฺญาตาวินฺทฺริยสฺส อิมานิ อฏฺฐินฺทฺริยานิ สหชาตปริวารา โหนฺติ, อญฺญมญฺญปริวารา โหนฺติ, นิสฺสยปริวารา โหนฺติ, สมฺปยุตฺตปริวารา โหนฺติ, สหคตา โหนฺติ, สหชาตา โหนฺติ, สํสฏฺฐา โหนฺติ, สมฺปยุตฺตา โหนฺติ, เตว ตสฺส อาการา เจว โหนฺติ ปริวารา จ.\csromanspacer{} อิติ อิมานิ อฏฺฐฏฺฐกานิ จตุสฏฺฐิ โหนฺติ.}

\prose{\textbf{อาสวา}ติ กตเม เต อาสวา? กามาสโว ภวาสโว ทิฏฺฐาสโว อวิชฺชาสโว.\csromanspacer{} กตฺเถเต อาสวา ขียนฺติ? โสตาปตฺติมคฺเคน อนวเสโส ทิฏฺฐาสโว ขียติ, อปายคมนีโย กามาสโว ขียติ, อปายคมนีโย ภวาสโว ขียติ, อปายคมนีโย อวิชฺชาสโว ขียติ, เอตฺเถเต อาสวา ขียนฺติ.\csromanspacer{} สกทาคามิมคฺเคน โอฬาริโก กามาสโว ขียติ, ตเทกฏฺโฐ ภวาสโว ขียติ, ตเทกฏฺโฐ อวิชฺชาสโว ขียติ, เอตฺเถเต อาสวา ขียนฺติ.\csromanspacer{} อนาคามิมคฺเคน อนวเสโส กามาสโว ขียติ, ตเทกฏฺโฐ ภวาสโว ขียติ, ตเทกฏฺโฐ อวิชฺชาสโค ขียติ, เอตฺเถเต อาสวา ขียนฺติ.\csromanspacer{} อรหตฺตมคฺเคน อนวเสโส ภวาสโว ขียติ, อนวเสโส อวิชฺชาสโว ขียติ, เอตฺเถเต อาสวา ขียนฺติ.\csromanspacer{} ตํญาตฏฺเฐน ญาณํ, ปชานนฏฺเฐน ปญฺญา, เตน วุจฺจติ จตุสฏฺฐิยา อากาเรหิ ติณฺณนฺนํ อินฺทฺริยานํ วสิภาวตา ปญฺญา อาสวานํ ขเย ญาณํ.}

\vspace{\tipitakaparskip}
\csromancenter{อาสวกฺขยญาณนิทฺเทโส ปญฺจปญฺญาสโม.}
\csromanmatikaanchor{mtk.44}
\csromantocmarkref{subsection}{56-63. สจฺจญาณจตุกฺกทฺวยนิทฺเทส}{mtk.44}
\bookmark[dest={mtk.44},level=2]{56-63. สจฺจญาณจตุกฺกทฺวยนิทฺเทส}
\csromanheader{2}{56-63. สจฺจญาณจตุกฺกทฺวยนิทฺเทส}
\proseitem{108}{\csromanfolio{114}กถํ ปริญฺญฏฺเฐ ปญฺญา ทุกฺเข ญาณํ, ปหานฏฺเฐ ปญฺญา สมุทเย ญาณํ, สจฺฉิกิริยฏฺเฐ ปญฺญา นิโรเธ ญาณํ, ภาวนฏฺเฐ ปญฺญา มคฺเค ญาณํ\csromandash{}ทุกฺขสฺส ปีฬนฏฺโฐ สงฺขตฏฺโฐ สนฺตาปฏฺโฐ วิปริณามฏฺโฐ ปริญฺญาตฏฺโฐ.\csromanspacer{} สมุทยสฺส อายูหนฏฺโฐ นิทานฏฺโฐ สญฺโญคฏฺโฐ ปลิโพธฏฺโฐ ปหานฏฺโฐ.\csromanspacer{} นิโรธสฺส นิสฺสรณฏฺโฐ วิเวกฏฺโฐ อสงฺขตฏฺโฐ อมตฏฺโฐ สจฺฉิกิริยฏฺโฐ.\csromanspacer{} มคฺคสฺส นิยฺยานฏฺโฐ เหตุฏฺโฐ ทสฺสนฏฺโฐ อาธิปเตยฺยฏฺโฐ ภาวนฏฺโฐ.\csromanspacer{} ตํญาตฏฺเฐน ญาณํ, ปชานนฏฺเฐน ปญฺญา, เตน วุจฺจติ ปริญฺญฏฺเฐ ปญฺญา ทุกฺเข ญาณํ.\csromanspacer{} ปหานฏฺเฐ ปญฺญา สมุทเย ญาณํ.\csromanspacer{} สจฺฉิกิริยฏฺเฐ ปญฺญา นิโรเธ ญาณํ.\csromanspacer{} ภาวนฏฺเฐ ปญฺญา มคฺเค ญาณํ.}

\proseitem{109}{กถํ ทุกฺเข ญาณํ.\csromanspacer{} ทุกฺขสมุทเย ญาณํ.\csromanspacer{} ทุกฺขนิโรเธ ญาณํ.\csromanspacer{} ทุกฺขนิโรธคามินิยา ปฏิปทาย ญาณํ\csromandash{}มคฺคสมงฺคิสฺส ญาณํ ทุกฺเขเปตํ ญาณํ, ทุกฺขสมุทเยเปตํ ญาณํ, ทุกฺขนิโรเธเปตํ ญาณํ, ทุกฺขนิโรธคามินิยา ปฏิปทายเปตํ ญาณํ.}

\prose{ตตฺถ กตมํ ทุกฺเข ญาณํ? ทุกฺขํ อารพฺภ ยา อุปฺปชฺชติ ปญฺญา ปชานนา วิจโย ปวิจโย ธมฺมวิจโย สลฺลกฺขณา อุปลกฺขณา ปจฺจุปลกฺขณา ปณฺฑิจฺจํ โกสลฺลํ เนปุญฺญํ เวภพฺยา จินฺตา อุปปริกฺขา ภูริ เมธา ปริณายิกา วิปสฺสนา สมฺปชญฺญํ ปโตโท ปญฺญา ปญฺญินฺทฺริยํ ปญฺญาพลํ ปญฺญาสตฺถํ ปญฺญาปาสาโท ปญฺญา-อาโลโก ปญฺญา-โอภาโส ปญฺญาปชฺโชโต ปญฺญารตนํ อโมโห ธมฺมวิจโย สมฺมาทิฏฺฐิ, อิทํ วุจฺจติ ทุกฺเข ญาณํ.\csromanspacer{} ทุกฺขสมุทยํ อารพฺภ ฯเปฯ ทุกฺขนิโรธํ อารพฺภ ฯเปฯ ทุกฺขนิโรธคามินิํ ปฏิปทํ อารพฺภ ยา อุปฺปชฺชติ ปญฺญา ปชานนา ฯเปฯ อโมโห ธมฺมวิจโย สมฺมาทิฏฺฐิ, อิทํ วุจฺจติ ทุกฺขนิโรธคามินิยา ปฏิปทาย ญาณํ, ตํญาตฏฺเฐน ญาณํ, ปชานนฏฺเฐน ปญฺญา, เตน วุจฺจติ ทุกฺเข ญาณํ.\csromanspacer{} ทุกฺขสมุทเย ญาณํ.\csromanspacer{} ทุกฺขนิโรเธ ญาณํ.\csromanspacer{} ทุกฺขนิโรธคามินิยา ปฏิปทาย ญาณํ.}

\vspace{\tipitakaparskip}
\csromancenter{สจฺจญาณจตุกฺกทฺวยนิทฺเทโส เตสฏฺฐิโม.}
\csromanmatikaanchor{mtk.45}
\csromantocmarkref{subsection}{64-67. สุทฺธิกปฏิสมฺภิทาญาณนิทฺเทส}{mtk.45}
\bookmark[dest={mtk.45},level=2]{64-67. สุทฺธิกปฏิสมฺภิทาญาณนิทฺเทส}
\csromanheader{2}{64-67. สุทฺธิกปฏิสมฺภิทาญาณนิทฺเทส}
\proseitem{110}{\csromanfolio{115}กถํ อตฺถปฏิสมฺภิเท ญาณํ, ธมฺมปฏิสมฺภิเท ญาณํ, นิรุตฺติปฏิสมฺภิเท ญาณํ, ปฏิภานปฏิสมฺภิเท ญาณํ\csromandash{}อตฺเถสุ ญาณํ อตฺถปฏิสมฺภิทา, ธมฺเมสุ ญาณํ ธมฺมปฏิสมฺภิทา, นิรุตฺตีสุ ญาณํ นิรุตฺติปฏิสมฺภิทา, ปฏิภาเนสุ ญาณํ ปฏิภานปฏิสมฺภิทา.\csromanspacer{} อตฺถนานตฺเต ปญฺญา อตฺถปฏิสมฺภิเท ญาณํ, ธมฺมนานตฺเต ปญฺญา ธมฺมปฏิสมฺภิเท ญาณํ, นิรุตฺตินานตฺเต ปญฺญา นิรุตฺติปฏิสมฺภิเท ญาณํ, ปฏิภานนานตฺเต ปญฺญา ปฏิภานปฏิสมฺภิเท ญาณํ.\csromanspacer{}\csromanspacer{}อตฺถววตฺถาเน ปญฺญา อตฺถปฏิสมฺภิเท ญาณํ, ธมฺมววตฺถาเน ปญฺญา ธมฺมปฏิสมฺภิเท ญาณํ, นิรุตฺติววตฺถาเน ปญฺญา นิรุตฺติปฏิสมฺภิเท ญาณํ, ปฏิภานววตฺถาเน ปญฺญา ปฏิภานปฏิสมฺภิเท ญาณํ.}

\prose{อตฺถสลฺลกฺขเณ ปญฺญา อตฺถปฏิสมฺภิเท ญาณํ, ธมฺมสลฺลกฺขเณ ปญฺญา ธมฺมปฏิสมฺภิเท ญาณํ, นิรุตฺติสลฺลกฺขเณ ปญฺญา นิรุตฺติปฏิสมฺภิเท ญาณํ, ปฏิภานสลฺลกฺขเณ ปญฺญา ปฏิภานปฏิสมฺภิเท ญาณํ.\csromanspacer{}\csromanspacer{}อตฺถูปลกฺขเณ ปญฺญา อตฺถปฏิสมฺภิเท ญาณํ, ธมฺมูปลกฺขเณ ปญฺญา ธมฺมปฏิสมฺภิเท ญาณํ, นิรุตฺตูปลกฺขเณ ปญฺญา นิรุตฺติปฏิสมฺภิเท ญาณํ.\csromanspacer{} ปฏิภานูปลกฺขเณ ปญฺญา ปฏิภานปฏิสมฺภิเท ญาณํ.}

\prose{อตฺถปฺปเภเท ปญฺญา อตฺถปฏิสมฺภิเท ญาณํ, ธมฺมปฺปเภเท ปญฺญา ธมฺมปฏิสมฺภิเท ญาณํ, นิรุตฺติปฺปเภเท ปญฺญา นิรุตฺติปฏิสมฺภิเท ญาณํ, ปฏิภานปฺปเภเท ปญฺญา ปฏิภานปฏิสมฺภิเท ญาณํ.\csromanspacer{}\csromanspacer{}อตฺถปฺปภาวเน ปญฺญา อตฺถปฏิสมฺภิเท ญาณํ, ธมฺมปฺปภาวเน ปญฺญา ธมฺมปฏิสมฺภิเท ญาณํ, นิรุตฺติปฺปภาวเน ปญฺญา นิรุตฺติปฏิสมฺภิเท ญาณํ, ปฏิภานปฺปภาวเน ปญฺญา ปฏิภานปฏิสมฺภิเท ญาณํ.}

\prose{อตฺถโชตเน ปญฺญา อตฺถปฏิสมฺภิเท ญาณํ, ธมฺมโชตเน ปญฺญา ธมฺมปฏิสมฺภิเท ญาณํ, นิรุตฺติโชตเน ปญฺญา นิรุตฺติปฏิสมฺภิเท ญาณํ, ปฏิภานโชตเน ปญฺญา ปฏิภานปฏิสมฺภิเท ญาณํ.\csromanspacer{}\csromanspacer{}อตฺถวิโรจเน ปญฺญา อตฺถปฏิสมฺภิเท ญาณํ, ธมฺมวิโรจเน ปญฺญา ธมฺมปฏิสมฺภิเท ญาณํ, นิรุตฺติวิโรจเน ปญฺญา นิรุตฺติปฏิสมฺภิเท ญาณํ, ปฏิภานวิโรจเน ปญฺญา ปฏิภานปฏิสมฺภิเท ญาณํ.\csromanspacer{}\csromanspacer{}อตฺถปฺปกาสเน ปญฺญา อตฺถปฏิสมฺภิเท ญาณํ, ธมฺมปฺปกาสเน ปญฺญา ธมฺมปฏิสมฺภิเท ญาณํ, นิรุตฺถิปฺปกาสเน ปญฺญา นิรุตฺติปฏิสมฺภิเท ญาณํ, ปฏิภานปฺปกาสเน ปญฺญา ปฏิภานปฏิสมฺภิเท ญาณํ. \csromanfolio{116}ตํญาตฏฺเฐน ญาณํ, ปชานนฏฺเฐน ปญฺญา, เตน วุจฺจติ อตฺถปฏิสมฺภิเท ญาณํ, ธมฺมปฏิสมฺภิเท ญาณํ, นิรุตฺติปฏิสมฺภิเท ญาณํ, ปฏิภานปฏิสมฺภิเท ญาณํ.}

\csromanheader{2}{สุทฺธิกปฏิสมฺภิทาญาณนิทฺเทโส สตฺตสฏฺฐิโม.}
\csromansectionrule
\csromanmatikaanchor{mtk.46}
\csromantocmarkref{subsection}{68. อินฺทฺริยปโรปริยตฺตญาณนิทฺเทส}{mtk.46}
\bookmark[dest={mtk.46},level=2]{68. อินฺทฺริยปโรปริยตฺตญาณนิทฺเทส}
\csromanheader{2}{68. อินฺทฺริยปโรปริยตฺตญาณนิทฺเทส}
\proseitem{111}{กตมํ \textbf{ตถาคตสฺส อินฺทฺริยปโรปริยตฺตญาณํ}? อิธ ตถาคโต สตฺเต ปสฺสติ อปฺปรชกฺเข มหารชกฺเข ติกฺขินฺทฺริเย มุทินฺทฺริเย สฺวากาเร ทฺวากาเร สุวิญฺญาปเย ทุวิญฺญาปเย อปฺเปกจฺเจ ปรโลกวชฺชภยทสฺสาวิโน อปฺเปกจฺเจ นปรโลกวชฺชภยทสฺสาวิโน.}

\prose{\textbf{อปฺปรชกฺเข มหารชกฺเข}ติ สทฺโธ ปุคฺคโล อปฺปรชกฺโข, อสฺสทฺโธ ปุคฺคโล มหารชกฺโข.\csromanspacer{} อารทฺธวีริโย ปุคฺคโล อปฺปรชกฺโข, กุสีโต ปุคฺคโล มหารชกฺโข.\csromanspacer{} อุปฏฺฐิตสฺสติ ปุคฺคโล อปฺปรชกฺโข, มุฏฺฐสฺสติ ปุคฺคโล มหารชกฺโข.\csromanspacer{} สมาหิโต ปุคฺคโล อปฺปรชกฺโข, อสมาหิโต ปุคฺคโล มหารชกฺโข.\csromanspacer{} ปญฺญวา ปุคฺคโล อปฺปรชกฺโข, ทุปฺปญฺโญ ปุคฺคโล มหารชกฺโข.}

\prose{\textbf{ติกฺขินฺทฺริเย มุทินฺทฺริเย}ติ สทฺโธ ปุคฺคโล ติกฺขินฺทฺริโย, อสฺสทฺโธ ปุคฺคโล มุทินฺทฺริโย.\csromanspacer{} อารทฺธวีริโย ปุคฺคโล ติกฺขินฺทฺริโย, กุสีโต ปุคฺคโล มุทินฺทฺริโย.\csromanspacer{} อุปฏฺฐิตสฺสติ ปุคฺคโล ติกฺขินฺทฺริโย, มุฏฺฐสฺสติ ปุคฺคโล มุทินฺทฺริโย.\csromanspacer{} สมาหิโต ปุคฺคโล ติกฺขินฺทฺริโย, อสมาหิโต ปุคฺคโล มุทินฺทฺริโย.\csromanspacer{} ปญฺญวา ปุคฺคโล ติกฺขินฺทฺริโย, ทุปฺปญฺโญ ปุคฺคโล มุทินฺทฺริโย.}

\prose{\textbf{สฺวากาเร ทฺวากาเร}ติ สทฺโธ ปุคฺคโล สฺวากาโร, อสฺสทฺโธ ปุคฺคโล ทฺวากาโร.\csromanspacer{} อารทฺธวีริโย ปุคฺคโล สฺวากาโร, กุสีโต ปุคฺคโล ทฺวากาโร.\csromanspacer{} อุปฏฺฐิตสฺสติ ปุคฺคโล สฺวากาโร, มุฏฺฐสฺสติ ปุคฺคโล ทฺวากาโร.\csromanspacer{} สมาหิโต ปุคฺคโล สฺวากาโร, อสมาหิโต ปุคฺคโล ทฺวากาโร.\csromanspacer{} ปญฺญวา ปุคฺคโล สฺวากาโร, ทุปฺปญฺโญ ปุคฺคโล ทฺวากาโร.}

\prose{\textbf{สุวิญฺญาปเย ทุวิญฺญาปเย}ติ สทฺโธ ปุคฺคโล สุวิญฺญาปโย, อสฺสทฺโธ ปุคฺคโล ทุวิญฺญาปโย.\csromanspacer{} อารทฺธวีริโย ปุคฺคโล สุวิญฺญาปโย, กุสีโต ปุคฺคโล ทุวิญฺญาปโย.\csromanspacer{} อุปฏฺฐิตสฺสติ ปุคฺคโล สุวิญฺญาปโย, มุฏฺฐสฺสติ ปุคฺคโล ทุวิญฺญาปโย.\csromanspacer{} สมาหิโต ปุคฺคโล \csromanfolio{117}สุวิญฺญาปโย, อสมาหิโต ปุคฺคโล ทุวิญฺญาปโย.\csromanspacer{} ปญฺญวา ปุคฺคโล สุวิญฺญาปโย, ทุปฺปญฺโญ ปุคฺคโล ทุวิญฺญาปโย.}

\prose{\textbf{อปฺเปกจฺเจ ปรโลกวชฺชภยทสฺสาวิโน อปฺเปกจฺเจ นปรโลกวชฺชภยทสฺสาวิโน}ติ สทฺโธ ปุคฺคโล ปรโลกวชฺชภยทสฺสาวี, อสฺสทฺโธ ปุคฺคโล นปรโลกวชฺชภยทสฺสาวี.\csromanspacer{} อารทฺธวีริโย ปุคฺคโล ปรโลกวชฺชภยทสฺสาวี, กุสีโต ปุคฺคโล นปรโลกวชฺชภยทสฺสาวี.\csromanspacer{} อุปฏฺฐิตสฺสติ ปุคฺคโล ปรโลกวชฺชภยทสฺสาวี, มุฏฺฐสฺสติ ปุคฺคโล นปรโลกวชฺชภยทสฺสาวี.\csromanspacer{} สมาหิโต ปุคฺคโล ปรโลกวชฺชภยทสฺสาวี, อสมาหิโต ปุคฺคโล นปรโลกวชฺชภยทสฺสาวี.\csromanspacer{} ปญฺญวา ปุคฺคโล ปรโลกวชฺชภยทสฺสาวี, ทุปฺปญฺโญ ปุคฺคโล นปรโลกวชฺชภยทสฺสาวี.}

\proseitem{112}{\textbf{โลโก}ติ ขนฺธโลโก ธาตุโลโก อายตนโลโก วิปตฺติภวโลโก วิปตฺติสมฺภวโลโก สมฺปตฺติภวโลโก สมฺปตฺติสมฺภวโลโก.}

\prose{เอโก โลโก, สพฺเพ สตฺตา อาหารฏฺฐิติกา.\csromanspacer{} เทฺว โลกา, นามญฺจ รูปญฺจ.\csromanspacer{} ตโย โลกา, ติสฺโส เวทนา.\csromanspacer{} จตฺตาโร โลกา, จตฺตาโร อาหารา.\csromanspacer{} ปญฺจ โลกา, ปญฺจุปาทานกฺขนฺธา.\csromanspacer{} ฉ โลกา, ฉ อชฺฌตฺติกานิ อายตนานิ.\csromanspacer{} สตฺต โลกา, สตฺต วิญฺญาณฏฺฐิติโย.\csromanspacer{} อฏฺฐ โลกา, อฏฺฐ โลกธมฺมา.\csromanspacer{} นว โลกา, นว สตฺตาวาสา.\csromanspacer{} ทส โลกา, ทสายตนานิ.\csromanspacer{} ทฺวาทส โลกา, ทฺวาทสายตนานิ.\csromanspacer{} อฏฺฐารส โลกา, อฏฺฐารส ธาตุโย.}

\prose{\textbf{วชฺชนฺ}ติ สพฺเพ กิเลสา วชฺชา, สพฺเพ ทุจฺจริตา วชฺชา, สพฺเพ อภิสงฺขารา วชฺชา, สพฺเพ ภวคามิกมฺมา วชฺชา, อิติ อิมสฺมิญฺจ โลเก อิมสฺมิญฺจ วชฺเช ติพฺพา ภยสญฺญา ปจฺจุปฏฺฐิตา โหติ เสยฺยถาปิ อุกฺขิตฺตาสิเก วธเก.\csromanspacer{} อิเมหิ ปญฺญาสาย อากาเรหิ อิมานิ ปญฺจินฺทฺริยานิ ชานาติ ปสฺสติ อญฺญาติ ปฏิวิชฺฌติ, อิทํ ตถาคตสฺส อินฺทฺริยปโรปริยตฺตญาณํ.}

\csromanheader{2}{อินฺทฺริยปโรปริยตฺตญาณนิทฺเทโส อฏฺฐสฏฺฐิโม.}
\csromansectionrule
\csromanmatikaanchor{mtk.47}
\csromantocmarkref{subsection}{69. อาสยานุสยญาณนิทฺเทส}{mtk.47}
\bookmark[dest={mtk.47},level=2]{69. อาสยานุสยญาณนิทฺเทส}
\csromanheader{2}{69. อาสยานุสยญาณนิทฺเทส}
\proseitem{113}{กตมํ \textbf{ตถาคตสฺส สตฺตานํ อาสยานุสเย ญาณํ}? อิธ ตถาคโต สตฺตานํ อาสยํ ชานาติ, อนุสยํ ชานาติ, จริตํ \csromanfolio{118}ชานาติ, อธิมุตฺติํ ชานาติ, ภพฺพาภพฺเพ สตฺเต ปชานาติ.\csromanspacer{} กตโม\footnote{กตโม จ (สฺยา, ก)} สตฺตานํ อาสโย? “สสฺสโต โลโก”ติ วา “อสสฺสโต โลโก”ติ วา “อนฺตวา โลโก”ติ วา “อนนฺตวา โลโก”ติ วา “ตํ ชีวํ ตํ สรีรนฺ”ติ วา “อญฺญํ ชีวํ อญฺญํ สรีรนฺ”ติ วา “โหติ ตถาคโต ปรํ มรณา”ติ วา “น โหติ ตถาคโต ปรํ มรณา”ติ วา “โหติ จ น จ โหติ ตถาคโต ปรํ มรณา”ติ วา “เนว โหติ น น โหติ ตถาคโต ปรํ มรณา”ติ วา.\csromanspacer{} อิติ ภวทิฏฺฐิสนฺนิสฺสิตา วา สตฺตา โหนฺติ วิภวทิฏฺฐิสนฺนิสฺสิตา วา.}

\prose{เอเต วา ปน อุโภ อนฺเต อนุปคมฺม อิทปฺปจฺจยตาปฏิจฺจสมุปฺปนฺเนสุ ธมฺเมสุ อนุโลมิกา ขนฺติ ปฏิลทฺธา โหติ, ยถาภูตํ วา ญาณํ.\csromanspacer{} กามํ เสวนฺตญฺเญว ชานาติ “อยํ ปุคฺคโล กามครุโก กามาสโย กามาธิมุตฺโต”ติ.\csromanspacer{} กามํ เสวนฺตญฺเญว ชานาติ “อยํ ปุคฺคโล เนกฺขมฺมครุโก เนกฺขมฺมาสโย เนกฺขมฺมาธิมุตฺโต”ติ.\csromanspacer{} เนกฺขมฺมํ เสวนฺตญฺเญว ชานาติ “อยํ ปุคฺคโล (เนกฺขมฺมครุโก เนกฺขมฺมาสโย เนกฺขมฺมาธิมุตฺโต”ติ.\csromanspacer{} เนกฺขมฺมํ เสวนฺตญฺเญว ชานาติ “อยํ ปุคฺคโล กามครุโก กามาสโย กามาธิมุตฺโต”ติ.\csromanspacer{} พฺยาปาทํ เสวนฺตญฺเญว ชานาติ)\footnote{( ) เอตฺถนฺตเร ปาฐา นตฺถิ สฺยามโปตฺถเก, เอวมุปริปิ.} “อยํ ปุคฺคโล พฺยาปาทครุโก พฺยาปาทาสโย พฺยาปาทาธิมุตฺโต”ติ.\csromanspacer{} พฺยาปาทํ เสวนฺตญฺเญว ชานาติ “อยํ ปุคฺคโล อพฺยาปาทครุโก อพฺยาปาทาสโย อพฺยาปาทาธิมุตฺโต”ติ.\csromanspacer{} อพฺยาปาทํ เสวนฺตญฺเญว ชานาติ “อยํ ปุคฺคโล (อพฺยาปาทครุโก อพฺยาปาทาสโย อพฺยาปาทาธิมุตฺโต”ติ.\csromanspacer{} อพฺยาปาทํ เสวนฺตญฺเญว ชานาติ “อยํ ปุคฺคโล พฺยาปาทครุโก พฺยาปาทาสโย พฺยาปาทาธิมุตฺโต”ติ.\csromanspacer{} ถินมิทฺธํ เสวนฺตญฺเญว ชานาติ) “อยํ ปุคฺคโล ถินมิทฺธครุโก ถินมิทฺธาสโย ถินมิทฺธาธิมุตฺโต”ติ.\csromanspacer{} ถินมิทฺธํ เสวนฺตญฺเญว ชานาติ “อยํ ปุคฺคโล อาโลกสญฺญาครุโก อาโลกสญฺญาสโย อาโลกสญฺญาธิมุตฺโต”ติ.\csromanspacer{} อาโลกสญฺญํ เสวนฺตญฺเญว ชานาติ “อยํ ปุคฺคโล (อาโลกสญฺญาครุโก อาโลกสญฺญาสโย อาโลกสญฺญาธิมุตฺโต”ติ.\csromanspacer{} อาโลกสญฺญํ เสวนฺตญฺเญว ชานาติ “อยํ ปุคฺคโล ถินมิทฺธครุโก ถินมิทฺธาสโย ถินมิทฺธาธิมุตฺโต”ติ,) อยํ สตฺตานํ อาสโย.}

\proseitem{114}{\csromanfolio{119}กตโม จ สตฺตานํ อนุสโย? สตฺตานุสยา กามราคานุสโย ปฏิฆานุสโย มานานุสโย ทิฏฺฐานุสโย วิจิกิจฺฉานุสโย ภวราคานุสโย อวิชฺชานุสโย.\csromanspacer{} ยํ โลเก ปิยรูปํ สาตรูปํ, เอตฺถ สตฺตานํ กามราคานุสโย อนุเสติ.\csromanspacer{} ยํ โลเก อปฺปิยรูปํ อสาตรูปํ, เอตฺถ สตฺตานํ ปฏิฆานุสโย อนุเสติ.\csromanspacer{} อิติ อิเมสุ ทฺวีสุ ธมฺเมสุ อวิชฺชา อนุปติตา, ตเทกฏฺโฐ มาโน จ ทิฏฺฐิ จ วิจิกิจฺฉา จ ทฏฺฐพฺพา.\csromanspacer{} อยํ สตฺตานํ อนุสโย.}

\prose{กตมญฺจ สตฺตานํ จริตํ? ปุญฺญาภิสงฺขาโร อปุญฺญาภิสงฺขาโร อาเนญฺชาภิสงฺขาโร ปริตฺตภูมโก วา มหาภูมโก วา.\csromanspacer{} อิทํ สตฺตานํ จริตํ.}

\proseitem{115}{กตมา จ สตฺตานํ อธิมุตฺติ? สนฺติ สตฺตา หีนาธิมุตฺติกา, สนฺติ สตฺตา ปณีตาธิมุตฺติกา.\csromanspacer{}\csromanspacer{}หีนาธิมุตฺติกา สตฺตา หีนาธิมุตฺติเก สตฺเต เสวนฺติ ภชนฺติ ปยิรุปาสนฺติ, ปณีตาธิมุตฺติกา สตฺตา ปณีตาธิมุตฺติเก สตฺเต เสวนฺติ ภชนฺติ ปยิรุปาสนฺติ.\csromanspacer{} อตีตมฺปิ อทฺธานํ หีนาธิมุตฺติกา สตฺตา หีนาธิมุตฺติเก สตฺเต เสวิํสุ ภชิํสุ ปยิรุปาสิํสุ ปณีตาธิมุตฺติกา สตฺตา ปณีตาธิมุตฺติเก สตฺเต เสวิํสุ ภชิํสุ ปยิรุปาสิํสุ.\csromanspacer{} อนาคตมฺปิ อทฺธานํ หีนาธิมุตฺติกา สตฺตา หีนาธิมุตฺติเก สตฺเต เสวิสฺสนฺติ ภชิสฺสนฺติ ปยิรุปาสสฺสนฺติ, ปณีตาธิมุตฺติกา สตฺตา ปณีภาธิมุตฺติเก สตฺเถ เสวิสฺสนฺติ ภชิสฺสนฺติ ปยิรุปาสิสฺสนฺติ.\csromanspacer{} อยํ สตฺตานํ อธิมุตฺติ.}

\prose{กตเม สตฺตา อภพฺพา? เย เต สตฺตา กมฺมาวรเณน สมนฺนาคตา กิเลสาวรเณน สมนฺนาคตา วิปากาวรเณน สมนฺนาคตา อสฺสทฺธา อจฺฉนฺทิกา ทุปฺปญฺญา อภพฺพา นิยามํ โอกฺกมิตุํ กุสเลสุ ธมฺเมสุ สมฺมตฺตํ, อิเม เต สตฺตา อภพฺพา.}

\prose{กตเม สตฺตา ภพฺพา? เย เต สตฺตา น กมฺมาวรเณน สมนฺนาคตา น กิเลสาวรเณน สมนฺนาคตา น วิปากาวรเณน สมนฺนาคตา สทฺธา ฉนฺทิกา ปญฺญวนฺโต ภพฺพา นิยามํ โอกฺกมิตุํ กุสเลสุ ธมฺเมสุ สมฺมตฺตํ.\csromanspacer{} อิเม เต สตฺตา ภพฺพา, อิทํ ตถาคตสฺส สตฺตานํ อาสยานุสเย ญาณํ.}

\vspace{\tipitakaparskip}
\csromancenter{อาสยานุสยญาณนิทฺเทโส นวสฏฺฐิโม.}
\csromanmatikaanchor{mtk.48}
\csromantocmarkref{subsection}{70. ยมกปาฏิหีรญาณนิทฺเทส}{mtk.48}
\bookmark[dest={mtk.48},level=2]{70. ยมกปาฏิหีรญาณนิทฺเทส}
\csromanheader{2}{70. ยมกปาฏิหีรญาณนิทฺเทส}
\proseitem{116}{\csromanfolio{120}กตมํ \textbf{ตถาคตสฺส ยมกปาฏิหีเร ญาณํ?} อิธ ตถาคโต ยมกปาฏิหีรํ กโรติ อสาธารณํ สาวเกหิ, อุปริมกายโต อคฺคิกฺขนฺโธ ปวตฺตติ, เหฏฺฐิมกายโต อุทกธารา ปวตฺตติ.\csromanspacer{} เหฏฺฐิมกายโต อคฺคิกฺขนฺโธ ปวตฺตติ, อุปริมกายโต อุทกธารา ปวตฺตติ.\csromanspacer{}\csromanspacer{}ปุรตฺถิมกายโต อคฺคิกฺขนฺโธ ปวตฺตติ, ปจฺฉิมกายโต อุทกธารา ปวตฺตติ.\csromanspacer{} ปจฺฉิมกายโต อคฺคิกฺขนฺโธ ปวตฺตติ, ปุรตฺติมกายโต อุทกธารา ปวตฺตติ.\csromanspacer{}\csromanspacer{}ทกฺขิณ-อกฺขิโต อคฺคิกฺขนฺโธ ปวตฺตติ, วาม-อกฺขิโต อุทกธารา ปวตฺตติ.\csromanspacer{} วาม-อกฺขิโต อคฺคิกฺขนฺโธ ปวตฺตติ, ทกฺขิณ-อกฺขิโต อุทกธารา ปวตฺตติ.\csromanspacer{}\csromanspacer{}ทกฺขิณกณฺณโสตโต อคฺคิกฺขนฺโธ ปวตฺตติ, วามกณฺณโสตโต อุทกธารา ปวตฺตติ, วามกณฺณโสตโต อคฺคิกฺขนฺโธ ปวตฺตติ, ทกฺขิณกณฺณโสตโต อุทกธารา ปวตฺตติ.\csromanspacer{}\csromanspacer{}ทกฺขิณนาสิกาโสตโต อคฺคิกฺขนฺโธ ปวตฺตติ, วามนาสิกาโสตโต อุทกธารา ปวตฺตติ.\csromanspacer{} วามนาสิกาโสตโต อคฺคิกฺขนฺโธ ปวตฺตติ, ทกฺขิณนาสิกาโสตโต อุทกธารา ปวตฺตติ.\csromanspacer{}\csromanspacer{}ทกฺขิณ-อํสกูฏโต อคฺคิกฺขนฺโธ ปวตฺตติ, วาม-อํสกูฏโต อุทกธารา ปวตฺตติ.\csromanspacer{} วาม-อํสกูฏโต อคฺคิกฺขนฺโธ ปวตฺตติ, ทกฺขิณ-อํสกูฏโต อุทกธารา ปวตฺตติ.\csromanspacer{}\csromanspacer{}ทกฺขิณหตฺถโต อคฺคิกฺขนฺโธ ปวตฺตติ, วามหตฺถโต อุทกธารา ปวตฺตติ.\csromanspacer{} วามหตฺถโต อคฺคิกฺขนฺโธ ปวตฺตติ, ทกฺขิณหตฺถโต อุทกธารา ปวตฺตติ.\csromanspacer{}\csromanspacer{}ทกฺขิณปสฺสโต อคฺคิกฺขนฺโธ ปวตฺตติ, วามปสฺสโต อุทกธารา ปวตฺตติ.\csromanspacer{} วามปสฺสโต อคฺคิกฺขนฺโธ ปวตฺตติ, ทกฺขิณปสฺสโต อุทกธารา ปวตฺตติ.\csromanspacer{}\csromanspacer{}ทกฺขิณปาทโต อคฺคิกฺขนฺโธ ปวตฺตติ, วามปาทโต อุทกธารา ปวตฺตติ.\csromanspacer{} วามปาทโต อคฺคิกฺขนฺโธ ปวตฺตติ, ทกฺขิณปาทโต อุทกธารา ปวตฺตติ.\csromanspacer{}\csromanspacer{}องฺคุลงฺคุเลหิ อคฺคิกฺขนฺโธ ปวตฺตติ, องฺคุลนฺตริกาหิ อุทกธารา ปวตฺตติ, องฺคุลนฺตริกาหิ อคฺคิกฺขนฺโธ ปวตฺตติ, องฺคุลงฺคุเลหิ อุทกธารา ปวตฺตติ.\csromanspacer{}\csromanspacer{}เอเกกโลมโต อคฺคิกฺขนฺโธ ปวตฺตติ, เอเกกโลมโต อุทกธารา ปวตฺตติ.\csromanspacer{} โลมกูปโต โลมกูปโต อคฺคิกฺขนฺโธ ปวตฺตติ, โลมกูปโต โลมกูปโต อุทกธารา ปวตฺตติ. \csromanfolio{121}ฉนฺนํ วณฺณานํ นีลานํ ปีตกานํ โลหิตกานํ โอทาตานํ มญฺชิฏฺฐานํ\footnote{มญฺเชฏฺฐานํ (สฺยา, ก)} ปภสฺสรานํ.}

\prose{ภควา จงฺกมติ, นิมฺมิโต ติฏฺฐติ วา นิสีทติ วา เสยฺยํ วา กปฺเปติ.\csromanspacer{} ภควา ติฏฺฐติ, นิมฺมิโต จงฺกมติ วา นิสีทติ วา เสยฺยํ วา กปฺเปติ.\csromanspacer{} ภควา นิสีทติ, นิมฺมิโต จงฺกมติ วา ติฏฺฐติ วา เสยฺยํ วา กปฺเปติ.\csromanspacer{} ภควา เสยฺยํ กปฺเปติ, นิมฺมิโต จงฺกมติ วา ติฏฺฐติ วา นิสีทติ วา.\csromanspacer{} นิมฺมิโต จงฺกมติ, ภควา ติฏฺฐติ วา นิสีทติ วา เสยฺยํ วา กปฺเปติ.\csromanspacer{} นิมฺมิโต ติฏฺฐติ, ภควา จงฺกมติ วา นิสีทติ วา เสยฺยํ วา กปฺเปติ.\csromanspacer{} นิมฺมิโต นิสีทติ, ภควา จงฺกมติ วา ติฏฺฐติ วา เสยฺยํ วา กปฺเปติ.\csromanspacer{} นิมฺมิโต เสยฺยํ กปฺเปติ, ภควา จงฺกมติ วา ติฏฺฐติ วา นิสีทติ วา.\csromanspacer{} อิทํ ตถาคตสฺส ยมกปาฏิหีเร ญาณํ.}

\csromanheader{2}{ยมกปาฏิหีรญาณนิทฺเทโส สตฺตติโม.}
\csromansectionrule
\csromanmatikaanchor{mtk.49}
\csromantocmarkref{subsection}{71. มหากรุณาญาณนิทฺเทส}{mtk.49}
\bookmark[dest={mtk.49},level=2]{71. มหากรุณาญาณนิทฺเทส}
\csromanheader{2}{71. มหากรุณาญาณนิทฺเทส}
\proseitem{117}{กตมํ \textbf{ตถาคตสฺส มหาครุณาสมาปตฺติยา ญาณํ?} พหุเกหิ อากาเรหิ ปสฺสนฺตานํ พุทฺธานํ ภควนฺตานํ สตฺเตตุ มหากรุณา โอกฺกมติ, “อาทิตฺโต โลกสนฺนิวาโส”ติ ปสฺสนฺตานํ พุทฺธานํ ภควนฺตานํ สตฺเตสุ มหากรุณา โอกฺกมติ. “อุยฺยุตฺโต โลกสนฺนิวาโส”ติ ปสฺสนฺตานํ พุทฺธานํ ภควนฺตานํ สตฺเตสุ มหากรุณา โอกฺกมติ, “ปยาโต โลกสนฺนิวาโส”ติ ปสฺสนฺตานํ พุทฺธานํ ภควนฺตานํ สตฺเตสุ มหากรุณา โอกฺกมติ. “กุมฺมคฺคปฺปฏิปนฺโน\footnote{กุมฺมคฺคํ ปฏิปนฺโน (สฺยา)} โลกสนฺนิวาโส”ติ ปสฺสนฺตานํ พุทฺธานํ ภควนฺตานํ สตฺเตสุ มหากรุณา โอกฺกมติ, “อุปนียติ โลโก อทฺธุโว”ติ ปสฺสนฺตานํ พุทฺธานํ ภควนฺตานํ สตฺเตสุ มหากรุณา โอกฺกมติ. “อตาโณ\footnote{อตฺตาโณ (สฺยา)} โลโก อนภิสฺสโร”ติ ปสฺสนฺตานํ พุทฺธานํ ภควนฺตานํ สตฺเตสุ มหากรุณา โอกฺกมติ, “อสฺสโก โลโก สพฺพํ ปหาย คมนียนฺ”ติ ปสฺสนฺตานํ พุทฺธานํ ภควนฺตานํ สตฺเตสุ มหากรุณา โอกฺกมติ, “อูโน โลโก อติตฺโต ตณฺหาทาโส”ติ ปสฺสนฺตานํ พุทฺธานํ ภควนฺตานํ สตฺเตสุ มหากรุณา โอกฺกมติ, “อตายโน โลกสนฺนิวาโส”ติ \csromanfolio{122}ปสฺสนฺตานํ ฯเปฯ “อเลโณ โลกสนฺนิวาโส”ติ ปสฺสนฺตานํ ฯเปฯ “อสรโณ โลกสนฺนิวาโส”ติ ปสฺสนฺตานํ ฯเปฯ “อสรณีภูโต โลกสนฺนิวาโส”ติ ปสฺสนฺตานํ ฯเปฯ.}

\prose{“อุทฺธโต โลโก อวูปสนฺโต”ติ ปสฺสนฺตานํ ฯเปฯ “สสลฺโลโลกสนฺนิวาโส วิทฺโธ ปุถุสลฺเลหิ, ตสฺส นตฺถญฺโญ โกจิ สลฺลานํ อุทฺธตา อญฺญตฺร มยา”ติ ปสฺสนฺตานํ ฯเปฯ “อวิชฺชนฺธการาวรโณ โลกสนฺนิวาโส อณฺฑภูโต กิเลสปญฺชรปกฺขิตฺโต, ตสฺส นตฺถญฺโญ โกจิ อาโลกํ ทสฺเสตา อญฺญตฺร มยา”ติ ปสฺสนฺตานํ ฯเปฯ “อวิชฺชาคโต โลกสนฺนิวาโส อณฺฑภูโต ปริโยนทฺโธ ตนฺตากุลกชาโต\footnote{ตนฺตากุลชาโต (สฺยา)} กุลาคณฺฐิกชาโต\footnote{คุฬาคุณฺฐิกชาโต (สฺยา), คุลาคุณฺฑิกชาโต (ก, สี-ฏฺฐ) ที 2. 47 ปิฏฺเฐ ปสฺสิตพฺพา.} มุญฺชปพฺพชภูโต อปายํ ทุคฺคติํ วินิปาตํ สํสารํ นาติวตฺตตี”ติ ปสฺสนฺตานํ ฯเปฯ “อวิชฺชาวิสโทสสํลิตฺโต โลกสนฺนิวาโส กิเลสกลลีภูโต”ติ ปสฺสนฺตานํ ฯเปฯ “ราคโทสโมห- ชฏาชฏิโต โลกสนฺนิวาโส, ตสฺส นตฺถญฺโญ โกจิ ชฏํ วิชเฏตา อญฺญตฺร มยา”ติ ปสฺสนฺตานํ ฯเปฯ.}

\prose{“ตณฺหาสํฆาฏปฏิมุกฺโก โลกสนฺนิวาโส”ติ ปสฺสนฺตานํ ฯเปฯ “ตณฺหาชาเลน โอตฺถโฏ\footnote{โอตฺถโต (ก)} โลกสนฺนิวาโส”ติ ปสฺสนฺตานํ ฯเปฯ “ตณฺหาโสเตน วุยฺหติ โลกสนฺนิวาโส”ติ ปสฺสนฺตานํ ฯเปฯ “ตณฺหาสญฺโญชเนน สญฺญุตฺโต โลกสนฺนิวาโส”ติ ปสฺสนฺตานํ ฯเปฯ “ตณฺหานุสเยน อนุสโฏ โลกสนฺนิวาโส”ติ ปสฺสนฺตานํ ฯเปฯ “ตณฺหาสนฺตาเปน สนฺตปฺปติ โลกสนฺนิวาโส”ติ ปสฺสนฺตานํ ฯเปฯ “ตณฺหาปริฬาเหน ปริฑยฺหติ โลกสนฺนิวาโส”ติ ปสฺสนฺตานํ ฯเปฯ.}

\prose{“ทิฏฺฐิสํฆาฏปฏิมุกฺโก โลกสนฺนิวาโส”ติ ปสฺสนฺตานํ ฯเปฯ “ทิฏฺฐิชาเลน โอตฺถโฏ โลกสนฺนิวาโส”ติ ปสฺสนฺตานํ ฯเปฯ “ทิฏฺฐิโสเตน วุยฺหติ โลกสนฺนิวาโส”ติ ปสฺสนฺตานํ ฯเปฯ “ทิฏฺฐิสญฺโญชเนน สญฺญุตฺโต โลกสนฺนิวาโส”ติ ปสฺสนฺตานํ ฯเปฯ “ทิฏฺฐานุสเยน อนุสโฏ โลกสนฺนิวาโส”ติ ปสฺสนฺตานํ ฯเปฯ “ทิฏฺฐิสนฺตาเปน สนฺตปฺปติ โลกสนฺนิวาโส”ติ ปสฺสนฺตานํ ฯเปฯ “ทิฏฺฐิปริฬาเหน ปริฑยฺหติ โลกสนฺนิวาโส”ติ ปสฺสนฺตานํวฺ ฯเปฯ. \csromanfolio{123}“ชาติยา อนุคโต โลกสนฺนิวาโส”ติ ปสฺสนฺตานํ ฯเปฯ “ชราย อนุสโฏ โลกสนฺนิวาโส”ติ ปสฺสนฺตานํ ฯเปฯ “พฺยาธินา อภิภูโต โลกสนฺนิวาโส”ติ ปสฺสนฺตานํ ฯเปฯ “มรเณน อพฺภาหโต โลกสนฺนิวาโส”ติ ปสฺสนฺตานํ ฯเปฯ “ทุกฺเข ปติฏฺฐิโต โลกสนฺนิวาโส”ติ ปสฺสนฺตานํ ฯเปฯ.}

\prose{“ตณฺหาย อุฑฺฑิโต โลกสนฺนิวาโส”ติ ปสฺสนฺตานํ ฯเปฯ “ชราปาการปริกฺขิตฺโต โลกสนฺนิวาโส”ติ ปสฺสนฺตานํ ฯเปฯ “มจฺจุปาเสน ปริกฺขิตฺโต โลกสนฺนิวาโส”ติ ปสฺสนฺตานํ ฯเปฯ “มหาพนฺธนพทฺโธ โลกสนฺนิวาโส ราคพนฺธเนน โทสพนฺธเนน โมหพนฺธเนน มานพนฺธเนน ทิฏฺฐิพนฺธเนน กิเลสพนฺธเนน ทุจฺจริตพนฺธเนน, ตสฺส นตฺถญฺโญ โกจิ พนฺธนํ โมเจตา อญฺญตฺร มยา”ติ ปสฺสนฺตานํ ฯเปฯ “มหาสมฺพาธปฺปฏิปนฺโน โลกสนฺนิวาโส, ตสฺส นตฺถญฺโญ โกจิ โอกาสํ ทสฺเสตา อญฺญตฺร มยา”ติ ปสฺสนฺตานํ. “มหาปลิโพเธน ปลิพุทฺโธ โลกสนฺนิวาโส, ตสฺส นตฺถญฺโญ โกจิ ปลิโพธํ เฉตา อญฺญตฺร มยา”ติ ปสฺสนฺตานํ ฯเปฯ “มหาปปาเต ปติโต โลกสนฺนิวาโส, ตสฺส นตฺถญฺโญ โกจิ ปปาตา อุทฺธตา อญฺญตฺร มยา”ติ ปสฺสนฺตานํ ฯเปฯ “มหากนฺตารปฺปฏิปนฺโน โลกสนฺนิวาโส, ตสฺส นตฺถญฺโญ โกจิ กนฺตารํ ตาเรตา อญฺญตฺร มยา”ติ ปสฺสนฺตานํ ฯเปฯ “มหาสํสารปฺปฏิปนฺโน โลกสนฺนิวาโส, ตสฺส นตฺถญฺโญ โกจิ สํสารา โมเจตา อญฺญตฺร มยา”ติ ปสฺสนฺตานํ ฯเปฯ “มหาวิทุคฺเค สมฺปริวตฺตติ โลกสนฺนิวาโส, ตสฺส นตฺถญฺโญ โกจิ วิทุคฺคา อุทฺธตา อญฺญตฺร มยา”ติ ปสฺสนฺตานํ ฯเปฯ “มหาปลิเป\footnote{มหาปลฺเลเป (สฺยา)} ปลิปนฺโน โลกสนฺนิวาโส, ตสฺส นตฺถญฺโญ โกจิ ปลิปา อุทฺธตา อญฺญตฺร มยา”ติ ปสฺสนฺตานํ ฯเปฯ.}

\prose{“อพฺภาหโต โลกสนฺนิวาโส”ติ ปสฺสนฺตานํ ฯเปฯ “อาทิตฺโต โลกสนฺนิวาโส ราคคฺคินา โทสคฺคินา โมหคฺคินา ชาติยา ชราย มรเณน โสเกหิ ปริเทเวหิ ทุกฺเขหิ โทมนสฺเสหิ อุปายาเสหิ, ตสฺส นตฺถญฺโญ โกจิ นิพฺพาเปตา อญฺญตฺร มยา”ติ ปสฺสนฺตานํ ฯเปฯ “อุนฺนีตโก โลกสนฺนิวาโส หญฺญติ นิจฺจมตาโณ ปตฺตทณฺโฑ ตกฺกโร”ติ ปสฺสนฺตานํ ฯเปฯ “วชฺชพนฺธนพทฺโธ โลกสนฺนิวาโส อาฆาตนปจฺจุปฏฺฐิโต, ตสฺส นตฺถญฺโญ โกจิ พนฺธนํ โมเจตา \csromanfolio{124}อญฺญตฺร มยา”ติ ปสฺสนฺตานํ ฯเปฯ “อนาโถ โลกสนฺนิวาโส ปรมการุญฺญปฺปตฺโต, ตสฺส นตฺถญฺโญ โกจิ ตาเยตา อญฺญตฺร มยา”ติ ปสฺสนฺตานํ ฯเปฯ “ทุกฺขาภิตุนฺโน\footnote{ทุกฺขาภิตุณฺโณ (ก)} โลกสนฺนิวาโส จิรรตฺตํ ปีฬิโต”ติ ปสฺสนฺตานํ ฯเปฯ “คธิโต โลกสนฺนิวาโส นิจฺจํ ปิปาสิโต”ติ ปสฺสนฺตานํ ฯเปฯ.}

\prose{“อนฺโธ โลกสนฺนิวาโส อจกฺขุโก”ติ ปสฺสนฺตานํ ฯเปฯ “หตเนตฺโต โลกสนฺนิวาโส อปริณายโก”ติ ปสฺสนฺตานํ ฯเปฯ “วิปถปกฺขนฺโท\footnote{วิปถํ ปกฺขนฺโต (สฺยา)} โลกสนฺนิวาโส อญฺชสาปรทฺโธ, ตสฺส นตฺถญฺโญ โกจิ อริยปถํ อาเนตา อญฺญตฺร มยา”ติ ปสฺสนฺตานํ ฯเปฯ “มโหฆปกฺขนฺโท โลกสนฺนิวาโส, ตสฺส นตฺถญฺโญ โกจิ โอฆา อุทฺธตา อญฺญตฺร มยา”ติ ปสฺสนฺตานํ ฯเปฯ.}

\proseitem{118}{“ทฺวีหิ ทิฏฺฐิคเตหิ ปริยุฏฺฐิโต โลกสนฺนิวาโส”ติ ปสฺสนฺตานํ ฯเปฯ “ตีหิ ทุจฺจริเตหิ วิปฺปฏิปนฺโน โลกสนฺนิวาโส”ติ ปสฺสนฺตานํ ฯเปฯ “จตูหิ โยเคหิ ยุตฺโต โลกสนฺนิวาโส จตุโยคโยชิโต”ติ ปสฺสนฺตานํ ฯเปฯ “จตูหิ คนฺเถหิ\footnote{คณฺเฐหิ (สฺยา)} คนฺถิโต โลกสนฺนิวาโส”ติ ปสฺสนฺตานํ ฯเปฯ “จตูหิ อุปาทาเนหิ อุปาทียติ โลกสนฺนิวาโส”ติ ปสฺสนฺตานํ ฯเปฯ “ปญฺจคติสมารุโฬฺห โลกสนฺนิวาโส”ติ ปสฺสนฺตานํ ฯเปฯ “ปญฺจหิ กามคุเณหิ รชฺชติ โลกสนฺนิวาโส”ติ ปสฺสนฺตานํ ฯเปฯ “ปญฺจหิ นีวรเณหิ โอตฺถโฏ โลกสนฺนิวาโส”ติ ปสฺสนฺตานํ ฯเปฯ “ฉหิ วิวาทมูเลหิ วิวทติ โลกสนฺนิวาโส”ติ ปสฺสนฺตานํ ฯเปฯ “ฉหิ ตณฺหากาเยหิ รชฺชติ โลกสนฺนิวาโส”ติ ปสฺสนฺตานํ ฯเปฯ “ฉหิ ทิฏฺฐิคเตหิ ปริยุฏฺฐิโต โลกสนฺนิวาโส”ติ ปสฺสนฺตานํ ฯเปฯ “สตฺตหิ อนุสเยหิ อนุสโฏ โลกสนฺนิวาโส”ติ ปสฺสนฺตานํ ฯเปฯ “สตฺตหิ สญฺโญชเนหิ สญฺญุตฺโต โลกสนฺนิวาโส”ติ ปสฺสนฺตานํ ฯเปฯ “สตฺตหิ มาเนหิ อุนฺนโต\footnote{อุณฺณโต (สฺยา, ก)} โลกสนฺนิวาโส”ติ ปสฺสนฺตานํ ฯเปฯ “อฏฺฐหิ โลกธมฺเมหิ สมฺปริวตฺตติ โลกสนฺนิวาโส”ติ ปสฺสนฺตานํ ฯเปฯ “อฏฺฐหิ มิจฺฉตฺเตหิ นิยฺยาโต โลกสนฺนิวาโส”ติ ปสฺสนฺตานํ ฯเปฯ “อฏฺฐหิ ปุริสโทเสหิ ทุสฺสติ โลกสนฺนิวาโส”ติ \csromanfolio{125}ปสฺสนฺตานํ ฯเปฯ “นวหิ อาฆาตวตฺถูหิ อาฆาติโต โลกสนฺนิวาโส”ติ ปสฺสนฺตานํ ฯเปฯ “นววิธมาเนหิ อุนฺนโต โลกสนฺนิวาโส”ติ ปสฺสนฺตานํ ฯเปฯ “นวหิ ตณฺหามูลเกหิ ธมฺเมหิ รชฺชติ โลกสนฺนีวาโส”ติ ปสฺสนฺตานํ ฯเปฯ “ทสหิ กิเลสวตฺถูหิ กิลิสฺสติ โลกสนฺนิวาโส”ติ ปสฺสนฺตานํ ฯเปฯ “ทสหิ อาฆาตวตฺถูหิ อาฆาติโต โลกสนฺนิวาโส”ติ ปสฺสนฺตานํ ฯเปฯ “ทสหิ อกุสลกมฺมปเถหิ สมนฺนาคโต โลกสนฺนิวาโส”ติ ปสฺสนฺตานํ ฯเปฯ “ทสหิ สญฺโญชเนหิ สญฺญุตฺโต โลกสนฺนิวาโส”ติ ปสฺสนฺตานํ ฯเปฯ “ทสหิ มิจฺฉตฺเตหิ นิยฺยาโต โลกสนฺนิวาโส”ติ ปสฺสนฺตานํ ฯเปฯ “ทสวตฺถุกาย มิจฺฉาทิฏฺฐิยา สมนฺนาคโต โลกสนฺนิวาโส”ติ ปสฺสนฺตานํ ฯเปฯ “ทสวตฺถุกาย อนฺตคฺคาหิกาย ทิฏฺฐิยา สมนฺนาคโต โลกสนฺนิวาโส”ติ ปสฺสนฺตานํ ฯเปฯ “อฏฺฐสตตณฺหาปปญฺจสเตหิ ปปญฺจิโต โลกสนฺนิวาโส”ติ ปสฺสนฺตานํ พุทฺธานํ ภควนฺตานํ สตฺเตสุ มหากรุณา โอกฺกมติ.\csromanspacer{} ทฺวาสฏฺฐิยา ทิฏฺฐิคเตหิ ปริยุฏฺฐิโต โลกสนฺนิวาโสติ ปสฺสนฺตานํ พุทฺธานํ ภควนฺตานํ สตฺเตสุ มหากรุณา โอกฺกมติ.}

\prose{“อหญฺจมฺหิ ติณฺโณ, โลโก จ อติณฺโณ.\csromanspacer{} อหญฺจมฺหิ มุตฺโต, โลโก จ อมุตฺโต.\csromanspacer{} อหญฺจมฺหิ ทนฺโต, โลโก จ อทนฺโต.\csromanspacer{} อหญฺจมฺหิ สนฺโต, โลโก จ อสนฺโต.\csromanspacer{} อหญฺจมฺหิ อสฺสตฺโถ, โลโก จ อนสฺสตฺโถ.\csromanspacer{} อหญฺจมฺหิ ปรินิพฺพุโต, โลโก จ อปรินิพฺพุโต.\csromanspacer{} ปโหมิ ขฺวาหํ ติณฺโณ ตาเรตุํ, มุตฺโต โมเจตุํ, ทนฺโต ทเมตุํ, สนฺโต สเมตุํ, อสฺสตฺโถ อสฺสาเสตุํ, ปรินิพฺพุโต ปเร จ ปรินิพฺพาเปตุนฺ”ติ ปสฺสนฺตานํ พุทฺธานํ ภควนฺตานํ สตฺเตสุ มหากรุณา โอกฺกมติ.\csromanspacer{} อิทํ ตถาคตสฺส มหากรุณาสมาปตฺติยา ญาณํ.}

\csromanheader{2}{มหากรุณาญาณนิทฺเทโส เอกสตฺตติโม.}
\csromansectionrule
\csromanmatikaanchor{mtk.50}
\csromantocmarkref{subsection}{72-73. สพฺพญฺญุตญฺญาณนิทฺเทส}{mtk.50}
\bookmark[dest={mtk.50},level=2]{72-73. สพฺพญฺญุตญฺญาณนิทฺเทส}
\csromanheader{2}{72-73. สพฺพญฺญุตญฺญาณนิทฺเทส}
\proseitem{119}{กตมํ \textbf{ตถาคตสฺส สพฺพญฺญุตญฺญาณํ?} สพฺพํ สงฺขตมสงฺขตํ อนวเสสํ ชานาตีติ สพฺพญฺญุตญฺญาณํ, ตตฺถ อาวรณํ นตฺถีติ อนาวรณญาณํ\footnote{อนาวรณํ ญาณํ (สฺยา) เอวมุปริปิ.}.}

\proseitem{120}{\csromanfolio{126}อตีตํ สพฺพํ ชานาตีติ สพฺพญฺญุตญฺญาณํ, ตตฺถ อาวรณํ นตฺถีติ อนาวรณญาณํ.\csromanspacer{} อนาคตํ สพฺพํ ชานาตีติ สพฺพญฺญุตญฺญาณํ, ตตฺถ อาวรณํ นตฺถีติ อนาวรณญาณํ.\csromanspacer{} ปจฺจุปฺปนฺนํ สพฺพํ ชานาตีติ สพฺพญฺญุตญฺญาณํ, ตตฺถ อาวรณํ นตฺถีติ อนาวรณญาณํ.}

\prose{จกฺขุ เจว รูปา จ, เอวํ ตํ สพฺพํ ชานาตีติ สพฺพญฺญุตญฺญาณํ, ตตฺถ อาวรณํ นตฺถีติ อนาวรณญาณํ.\csromanspacer{} โสตญฺเจว สทฺทา จ ฯเปฯ.\csromanspacer{} ฆานญฺเจว คนฺธา จ, ชิวฺหา เจว รสา จ.\csromanspacer{} กาโย เจว โผฏฺฐพฺพา จ.\csromanspacer{} มโน เจว ธมฺมา จ, เอวํ ตํ สพฺพํ ชานาตีติ สพฺพญฺญุตญฺญาณํ, ตตฺถ อาวรณํ นตฺถีติ อนาวรณญาณํ.}

\prose{ยาวตา อนิจฺจฏฺฐํ ทุกฺขฏฺฐํ อนตฺตฏฺฐํ, ตํ สพฺพํ ชานาตีติ สพฺพญฺญุตญฺญาณํ, ตตฺถ อาวรณํ นตฺถีติ อนาวรณญาณํ.\csromanspacer{} ยาวตา รูปสฺส อนิจฺจฏฺฐํ ทุกฺขฏฺฐํ อนตฺตฏฺฐํ, ตํ สพฺพํ ชานาตีติ สพฺพญฺญุตญฺญาณํ, ตตฺถ อาวรณํ นตฺถีติ อนาวรณญาณํ.\csromanspacer{} ยาวตา เวทนาย ฯเปฯ.\csromanspacer{} ยาวตา สญฺญาย ฯเปฯ.\csromanspacer{} ยาวตา สงฺขารานํ ฯเปฯ.\csromanspacer{} ยาวตา วิญฺญาณสฺส ฯเปฯ.\csromanspacer{} ยาวตา จกฺขุสฺส ฯเปฯ.\csromanspacer{} ชรามรณสฺส อนิจฺจฏฺฐํ ทุกฺขฏฺฐํ อนตฺตฏฺฐํ, ตํ สพฺพํ ชานาตีติ สพฺพญฺญุตญฺญาณํ, ตตฺถ อาวรณํ นตฺถีติ อนาวรณญาณํ.}

\prose{ยาวตา อภิญฺญาย อภิญฺญฏฺฐํ, ตํ สพฺพํ ชานาตีติ สพฺพญฺญุตญฺญาณํ, ตตฺถ อาวรณํ นตฺถีติ อนาวรณญาณํ.\csromanspacer{} ยาวตา ปริญฺญาย ปริญฺญฏฺฐํ ฯเปฯ.\csromanspacer{} ยาวตา ปหานสฺส\footnote{ปหานาย (สฺยา)} ปหานฏฺฐํ ฯเปฯ.\csromanspacer{} ยาวตา ภาวนาย ภาวนฏฺฐํ ฯเปฯ.\csromanspacer{} ยาวตา สจฺฉิกิริยาย สจฺฉิกิริยฏฺฐํ, ตํ สพฺพํ ชานาตีติ สพฺพญฺญุตญฺญาณํ, ตตฺถ อาวรณํ นตฺถีติ อนาวรณญาณํ.}

\prose{ยาวตา ขนฺธานํ ขนฺธฏฺฐํ, ตํ สพฺพํ ชานาตีติ สพฺพญฺญุตญฺญาณํ, ตตฺถ อาวรณํ นตฺถีติ อนาวรณญาณํ.\csromanspacer{} ยาวตา ธาตูนํ ธาตุฏฺฐํ ฯเปฯ.\csromanspacer{} ยาวตา อายตนานํ อายตนฏฺฐํ ฯเปฯ.\csromanspacer{} ยาวตา สงฺขตานํ สงฺขตฏฺฐํ ฯเปฯ.\csromanspacer{} ยาวตา อสงฺขตสฺส อสงฺขตฏฺฐํ, ตํ สพฺพํ ชานาตีติ สพฺพญฺญุตญฺญาณํ, ตตฺถ อาวรณํ นตฺถีติ อนาวรณญาณํ.}

\prose{ยาวตา กุสเล ธมฺเม สพฺพํ\footnote{สพฺเพ (อฏฺฐกถา)} ชานาตีติ สพฺพญฺญุตญฺญาณํ, ตตฺถ อาวรณํ นตฺถีติ อนาวรณญาณํ.\csromanspacer{} ยาวตา อกุสเล ธมฺเม ฯเปฯ.\csromanspacer{} ยาวตา อพฺยากเต ธมฺเม.\csromanspacer{} ยาวตา กามาวจเร ธมฺเม.\csromanspacer{} ยาวตา \csromanfolio{127}รูปาวจเร ธมฺเม.\csromanspacer{} ยาวตา อรูปาวจเร ธมฺเม.\csromanspacer{} ยาวตา อปริยาปนฺเน ธมฺเม สพฺพํ\footnote{สพฺเพ (อฏฺฐกถา)} ชานาตีติ สพฺพญฺญุตญฺญาณํ, ตตฺถ อาวรณํ นตฺถีติ อนาวรณญาณํ.}

\prose{ยาวตา ทุกฺขสฺส ทุกฺขฏฺฐํ, ตํ สพฺพํ ชานาตีติ สพฺพญฺญุตญฺญาณํ, ตตฺถ อาวรณํ นตฺถีติ อนาวรณญาณํ.\csromanspacer{} ยาวตา สมุทยสฺส สมุทยฏฺฐํ ฯเปฯ.\csromanspacer{} ยาวตา นิโรธสฺส นิโรธฏฺฐํ ฯเปฯ.\csromanspacer{} ยาวตา มคฺคสฺส มคฺคฏฺฐํ, ตํ สพฺพํ ชานาตีติ สพฺพญฺญุตญฺญาณํ.\csromanspacer{} ตตฺถ อาวรณํ นตฺถีติ อนาวรณญาณํ.}

\prose{ยาวตา อตฺถปฏิสมฺภิทาย อตฺถปฏิสมฺภิทฏฺฐํ, ตํ สพฺพํ ชานาตีติ สพฺพญฺญุตญฺญาณํ, ตตฺถ อาวรณํ นตฺถีติ อนาวรณญาณํ.\csromanspacer{} ยาวตา ธมฺมปฏิสมฺภิทาย ธมฺมปฏิสมฺภิทฏฺฐํ ฯเปฯ.\csromanspacer{} ยาวตา นิรุตฺติปฏิสมฺภิทาย นิรุตฺติปฏิสมฺภิทฏฺฐํ ฯเปฯ.\csromanspacer{} ยาวตา ปฏิภานปฏิสมฺภิทาย ปฏิภานปฏิสมฺภิทฏฺฐํ, ตํ สพฺพํ ชานาตีติ สพฺพญฺญุตญฺญาณํ, ตตฺถ อาวรณํ นตฺถีติ อนาวรณญาณํ.}

\prose{ยาวตา อินฺทฺริยปโรปริยตฺตญาณํ, ตํ สพฺพํ ชานาตีติ สพฺพญฺญุตญฺญาณํ, ตตฺถ อาวรณํ นตฺถีติ อนาวรณญาณํ.\csromanspacer{} ยาวตา สตฺตานํ อาสยานุสเย ญาณํ ฯเปฯ.\csromanspacer{} ยาวตา ยมกปาฏิหีเร ญาณํ ฯเปฯ.\csromanspacer{} ยาวตา มหากรุณาสมาปตฺติยา ญาณํ, ตํ สพฺพํ ชานาตีติ สพฺพญฺญุตญฺญาณํ, ตตฺถ อาวรณํ นตฺถีติ อนาวรณญาณํ.}

\prose{ยาวตา สเทวกสฺส โลกสฺส สมารกสฺส สพฺรหฺมกสฺส สสฺสมณพฺรหฺมณิยา ปชาย สเทวมนุสฺสาย ทิฏฺฐํ สุตํ มุตํ วิญฺญาตํ ปตฺตํ ปริเยสิตํ อนุวิจริตํ มนสา, ตํ สพฺพํ ชานาตีติ สพฺพญฺญุตญฺญาณํ, ตตฺถ อาวรณํ นตฺถีติ อนาวรณญาณํ.}

\csromanhangingbegin
\hangingheaditem{121}{น ตสฺส อทฺทิฏฺฐมิธตฺถิ กิญฺจิ,}
\hangingbody{อโถ อวิญฺญาตมชานิตพฺพํ.}
\hangingbody{สพฺพํ อภิญฺญาสิ ยทตฺถิ เนตฺยํ,}
\hangingbody{ตถาคโต เตน สมนฺตจกฺขูติ2.}
\csromanhangingend

\prose{\textbf{สมนฺตจกฺขู}ติ เกนฏฺเฐน สมนฺตจกฺขุ? จุทฺทส พุทฺธญาณานิ ทุกฺเข ญาณํ พุทฺธญาณํ, ทุกฺขสมุทเย ญาณํ พุทฺธญาณํ, ทุกฺขนิโรเธ ญาณํ พุทฺธญาณํ, \csromanfolio{128}ทุกฺขนิโรธคามินิยา ปฏิปทาย ญาณํ พุทฺธญาณํ, อตฺถปฏิสมฺภิเท ญาณํ พุทฺธญาณํ, ธมฺมปฏิสมฺภิเท ญาณํ พุทฺธญาณํ, นิรุตฺติปฏิสมฺภิเท ญาณํ พุทฺธญาณํ, ปฏิภานปฏิสมฺภิเท ญาณํ พุทฺธญาณํ, อินฺทฺริยปโรปริยตฺตญาณํ พุทฺธญาณํ, สตฺตานํ อาสยานุสเย ญาณํ พุทฺธญาณํ, ยมกปาฏิหีเร ญาณํ พุทฺธญาณํ, มหากรุณาสมาปตฺติยา ญาณํ พุทฺธญาณํ, สพฺพญฺญุตญฺญาณํ พุทฺธญาณํ, อนาวรณญาณํ พุทฺธญาณํ.\csromanspacer{} อิมานิ จุทฺทส พุทฺธญาณานิ.\csromanspacer{} อิเมสํ จุทฺทสนฺนํ พุทฺธญาณานํ อฏฺฐ ญาณานิ สาวกสาธารณานิ, ฉ ญาณานิ อสาธารณานิ สาวเกหิ.}

\prose{ยาวตา ทุกฺขสฺส ทุกฺขฏฺโฐ, สพฺโพ ญาโต, อญฺญาโต ทุกฺขฏฺโฐ นตฺถีติ สพฺพญฺญุตญฺญาณํ, ตตฺถ อาวรณํ นตฺถีติ อนาวรณญาณํ.\csromanspacer{} ยาวตา ทุกฺขสฺส ทุกฺขฏฺโฐ, สพฺโพ ญาโต สพฺโพ ทิฏฺโฐ สพฺโพ วิทิโต สพฺโพ สจฺฉิกโต สพฺโพ ผสฺสิโต ปญฺญาย, อผสฺสิโต ปญฺญาย ทุกฺขฏฺโฐ นตฺถีติ สพฺพญฺญุตญฺญาณํ, ตตฺถ อาวรณํ นตฺถีติ อนาวรณญาณํ.\csromanspacer{} ยาวตา สมุทยสฺส สมุทยฏฺโฐ.\csromanspacer{} ยาวตา นิโรธสฺส นิโรธฏฺโฐ.\csromanspacer{} ยาวตา มคฺคสฺส มคฺคฏฺโฐ ฯเปฯ.\csromanspacer{} ยาวตา อตฺถปฏิสมฺภิทาย อตฺถปฏิสมฺภิทฏฺโฐ.\csromanspacer{} ยาวตา ธมฺมปฏิสมฺภิทาย ธมฺมปฏิสมฺภิทฏฺโฐ.\csromanspacer{} ยาวตา นิรุตฺติปฏิสมฺภิทาย นิรุตฺติปฏิสมฺภิทฏฺโฐ.\csromanspacer{} ยาวตา ปฏิภานปฏิสมฺภิทาย ปฏิภานปฏิสมฺภิทฏฺโฐ, สพฺโพ ญาโต สพฺโพ ทิฏฺโฐ สพฺโพ วิทิโต สพฺโพ สจฺฉิกโต สพฺโพ ผสฺสิโต ปญฺญาย, อผสฺสิโต ปญฺญาย ปฏิภานปฏิสมฺภิทฏฺโฐ นตฺถีติ สพฺพญฺญุตญฺญาณํ, ตตฺถ อาวรณํ นตฺถีติ อนาวรณญาณํ.}

\prose{ยาวตา อินฺทฺริยปโรปริยตฺตญาณํ.\csromanspacer{} ยาวตา สตฺตานํ อาสยานุสเย ญาณํ.\csromanspacer{} ยาวตา ยมกปาฏิหีเร ญาณํ.\csromanspacer{} ยาวตา มหากรุณาสมาปตฺติยา ญาณํ, สพฺพํ ญาตํ สพฺพํ ทิฏฺฐํ สพฺพํ วิทิตํ สพฺพํ สจฺฉิกตํ สพฺพํ ผสฺสิตํ ปญฺญาย, อผสฺสิตํ ปญฺญาย มหากรุณาสมาปตฺติยา ญาณํ นตฺถีติ สพฺพญฺญุตญฺญาณํ, ตตฺถ อาวรณํ นตฺถีติ อนาวรณญาณํ.}

\prose{ยาวตา สเทวกสฺส โลกสฺส สมารกสฺส สพฺรหฺมกสฺส สสฺสมณพฺราหฺมณิยา ปชาย สเทวมนุสฺสาย ทิฏฺฐํ สุตํ มุตํ วิญฺญาตํ ปตฺตํ ปริเยสิตํ อนุวิจริตํ มนสา, สพฺพํ ญาตํ สพฺพํ ทิฏฺฐํ สพฺพํ วิทิตํ สพฺพํ สจฺฉิกตํ สพฺพํ ผสฺสิตํ ปญฺญาย, อผสฺสิตํ ปญฺญาย นตฺถีติ สพฺพญฺญุตญฺญาณํ, ตตฺถ อาวรณํ นตฺถีติ อนาวรณญาณํ.}

\csromanmatikaanchor{mtk.51}
\csromantocmarknopage{section}{2. ทิฏฺฐิกถา}{mtk.51}
\bookmark[dest={mtk.51},level=1]{2. ทิฏฺฐิกถา}
\csromanheader{1}{2. ทิฏฺฐิกถา}
\csromangathameasure{น ตสฺส อทฺทิฏฺฐมิธตฺถิ กิญฺจิ, \\ อโถ อวิญฺญาตมชานิตพฺพํ. \\ สพฺพํ อภิญฺญาสิ ยทตฺถิ เนยฺยํ, \\ ตถาคโต เตน สมนฺตจกฺขูติ.}
\csromangathagroup{\csromangathastanza{\csromanfolio{129}น ตสฺส อทฺทิฏฺฐมิธตฺถิ กิญฺจิ, \\ อโถ อวิญฺญาตมชานิตพฺพํ. \\ สพฺพํ อภิญฺญาสิ ยทตฺถิ เนยฺยํ, \\ ตถาคโต เตน สมนฺตจกฺขูติ.}}

\nitthitamruled{ญาณกถา นิฏฺฐิตา.}
\csromanheader{2}{2. ทิฏฺฐิกถา}
\proseitem{122}{กา ทิฏฺฐิ, กติ ทิฏฺฐิฏฺฐานานิ, กติ ทิฏฺฐิปริยุฏฺฐานานิ, กติ ทิฏฺฐิโย กติ ทิฏฺฐาภินิเวสา, กตโม ทิฏฺฐิฏฺฐานสมุคฺฆาโตติ.}

\prose{\textbf{กา ทิฏฺฐี}ติ \textbf{อภินิเวสปรามาโส ทิฏฺฐิ}.\csromanspacer{}\csromanspacer{}\textbf{กติ ทิฏฺฐิฏฺฐานานี}ติ อฏฺฐ ทิฏฺฐิฏฺฐานานิ.\csromanspacer{}\csromanspacer{}กฺ\textbf{อติ ทิฏฺฐิปริยุฏฺฐานานี}ติ อฏฺฐารส ทิฏฺฐิปริยุฏฺฐานานิ.\csromanspacer{}\csromanspacer{}\textbf{กติ ทิฏฺฐิโย}ติ โสฬส ทิฏฺฐิโย.\csromanspacer{}\csromanspacer{}\textbf{กติ ทิฏฺฐาภินิเวสา}ติ ตีณิ สตํ ทิฏฺฐาภินิเวสา.\csromanspacer{}\csromanspacer{}\textbf{กตโม ทิฏฺฐิฏฺฐานสมุคฺฆาโต}ติ โสตาปตฺติมคฺโค ทิฏฺฐิฏฺฐานสมุคฺฆาโต.}

\proseitem{123}{กถํ \textbf{อภินิเวสปรามาโส ทิฏฺฐิ}\footnote{กติ อภินิเวสปรามาโส ทิฏฺฐีติ (สฺยา)}? รูปํ “เอตํ มม, เอโสหมสฺมิ, เอโส เม อตฺตา”ติ \textbf{อภินิเวสปรามาโส ทิฏฺฐิ}.\csromanspacer{} เวทนํ เอตํ มม ฯเปฯ.\csromanspacer{} สญฺญํ เอตํ มม ฯเปฯ.\csromanspacer{} สงฺขาเร เอตํ มม ฯเปฯ.\csromanspacer{} วิญฺญาณํ “เอตํ มม, เอโสหมสฺมิ, เอโส เม อตฺตา”ติ \textbf{อภินิเวสปรามาโส ทิฏฺฐิ}.\csromanspacer{} จกฺขุํ เอตํ มม.\csromanspacer{} โสตํ เอตํ มม.\csromanspacer{} ฆานํ เอตํ มม.\csromanspacer{} ชิวฺหํ เอตํ มม.\csromanspacer{} กายํ เอตํ มม.\csromanspacer{} มนํ “เอตํ มม, เอโสหมสฺมิ, เอโส เม อตฺตา”ติ \textbf{อภินิเวสปรามาโส ทิฏฺฐิ}.\csromanspacer{}\csromanspacer{}รูเป\footnote{รูปํ (สฺยา) ตถา ปญฺจสุ อารมฺมเณสุ เอกวจเนน.} เอตํ มม.\csromanspacer{} สทฺเท เอตํ มม.\csromanspacer{} คนฺเธ เอตํ มม.\csromanspacer{} รเส เอตํ มม.\csromanspacer{} โผฏฺฐพฺเพ เอตํ มม.\csromanspacer{} ธมฺเม “เอตํ มม, เอโสหมสฺมิ, เอโส เม อตฺตา”ติ \textbf{อภินิเวสปรามาโส ทิฏฺฐิ}.\csromanspacer{}\csromanspacer{}จกฺขุวิญฺญาณํ เอตํ มม.\csromanspacer{} โสตวิญฺญาณํ เอตํ มม.\csromanspacer{} ฆานวิญฺญาณํ เอตํ มม.\csromanspacer{} ชิวฺหาวิญฺญาณํ เอตํ มม.\csromanspacer{} กายวิญฺญาณํ เอตํ มม.\csromanspacer{} มโนวิญฺญาณํ “เอตํ มม, เอโสหมสฺมิ, เอโส เม อตฺตา”ติ \textbf{อภินิเวสปรามาโส ทิฏฺฐิ}.\csromanspacer{}\csromanspacer{}จกฺขุสมฺผสฺสํ เอตํ มม.\csromanspacer{} โสตสมฺผสฺสํ เอตํ มม.\csromanspacer{} ฆานสมฺผสฺสํ เอตํ มม.\csromanspacer{} ชิวฺหาสมฺผสฺสํ เอตํ มม. \csromanfolio{130}กายสมฺผสฺสํ เอตํ มม.\csromanspacer{} มโนสมฺผสฺสํ “เอตํ มม, เอโสหมสฺมิ, เอโส เม อตฺตา”ติ อภินิเวสปรามาโส ทิฏฺฐิ.\csromanspacer{}\csromanspacer{}จกฺขุสมฺผสฺสชํ เวทนํ.\csromanspacer{} โสตสมฺผสฺสชํ เวทนํ.\csromanspacer{} ฆานสมฺผสฺสชํ เวทนํ.\csromanspacer{} ชิวฺหาสมฺผสฺสชํ เวทนํ.\csromanspacer{} กายสมฺผสฺสชํ เวทนํ.\csromanspacer{} มโนสมฺผสฺสชํ เวทนํ “เอตํ มม, เอโสหมสฺมิ, เอโส เม อตฺตา”ติ อภินิเวสปรามาโส ทิฏฺฐิ.}

\prose{รูปสญฺญํ เอตํ มม.\csromanspacer{} สทฺทสญฺญํ เอตํ มม.\csromanspacer{} คนฺธสญฺญํ เอตํ มม.\csromanspacer{} รสสญฺญํ เอตํ มม.\csromanspacer{} โผฏฺฐพฺพสญฺญํ เอตํ มม.\csromanspacer{} ธมฺมสญฺญํ “เอตํ มม, เอโสหมสฺมิ, เอโส เม อตฺตา”ติ อภินิเวสปรามาโส ทิฏฺฐิ.\csromanspacer{}\csromanspacer{}รูปสญฺเจตนํ เอตํ มม.\csromanspacer{} สทฺทสญฺเจตนํ เอตํ มม.\csromanspacer{} คนฺธสญฺเจตนํ เอตํ มม.\csromanspacer{} รสสญฺเจตนํ เอตํ มม.\csromanspacer{} โผฏฺฐพฺพสญฺเจตนํ เอตํ มม.\csromanspacer{} ธมฺมสญฺเจตนํ “เอตํ มม, เอโสหมสฺมิ, เอโส เม อตฺตา”ติ อภินิเวสปรามาโส ทิฏฺฐิ.\csromanspacer{}\csromanspacer{}รูปตณฺหํ เอตํ มม.\csromanspacer{} สทฺทตณฺหํ เอตํ มม.\csromanspacer{} คนฺธตณฺหํ เอตํ มม.\csromanspacer{} รสตณฺหํ เอตํ มม.\csromanspacer{} โผฏฺฐพฺพตณฺหํ เอตํ มม.\csromanspacer{} ธมฺมตณฺหํ “เอตํ มม, เอโสหมสฺมิ, เอโส เม อตฺตา”ติ อภินิเวสปรามาโส ทิฏฺฐิ.\csromanspacer{}\csromanspacer{}รูปวิตกฺกํ เอตํ มม.\csromanspacer{} สทฺทวิตกฺกํ เอตํ มม.\csromanspacer{} คนฺธวิตกฺกํ เอตํ มม.\csromanspacer{} รสวิตกฺกํ เอตํ มม.\csromanspacer{} โผฏฺฐพฺพวิตกฺกํ เอตํ มม.\csromanspacer{} ธมฺมวิตกฺกํ “เอตํ มม, เอโสหมสฺมิ, เอโส เม อตฺตา”ติ อภินิเวสปรามาโส ทิฏฺฐิ.\csromanspacer{}\csromanspacer{}รูปวิจารํ เอตํ มม.\csromanspacer{} สทฺทวิจารํ เอตํ มม.\csromanspacer{} คนฺธวิจารํ เอตํ มม.\csromanspacer{} รสวิจารํ เอตํ มม.\csromanspacer{} โผฏฺฐพฺพวิจารํ เอตํ มม.\csromanspacer{} ธมฺมวิจารํ “เอตํ มม, เอโสหมสฺมิ, เอโส เม อตฺตา”ติ อภินิเวสปรามาโส ทิฏฺฐิ.}

\prose{ปถวีธาตุํ เอตํ มม.\csromanspacer{} อาโปธาตุํ เอตํ มม.\csromanspacer{} เตโชธาตุํ เอตํ มม.\csromanspacer{} วาโยธาตุํ เอตํ มม.\csromanspacer{} อากาสธาตุํ เอตํ มม.\csromanspacer{} วิญฺญาณธาตุํ “เอตํ มม, เอโสหมสฺมิ, เอโส เม อตฺตา”ติ อภินิเวสปรามาโส ทิฏฺฐิ.\csromanspacer{}\csromanspacer{}ปถวีกสิณํ เอตํ มม.\csromanspacer{} อาโปกสิณํ.\csromanspacer{} เตโชกสิณํ.\csromanspacer{} วาโยกสิณํ.\csromanspacer{} นีลกสิณํ.\csromanspacer{} ปีตกสิณํ.\csromanspacer{} โลหิตกสิณํ.\csromanspacer{} โอทาตกสิณํ.\csromanspacer{} อากาสกสิณํ.\csromanspacer{} วิญฺญาณกสิณํ “เอตํ มม, เอโสหมสฺมิ, เอโส เม อตฺตา”ติ อภินิเวสปรามาโส ทิฏฺฐิ.}

\prose{เกสํ เอตํ มม.\csromanspacer{} โลมํ เอตํ มม.\csromanspacer{} นขํ เอตํ มม.\csromanspacer{} ทนฺตํ เอตํ มม.\csromanspacer{} ตจํ เอตํ มม.\csromanspacer{} มํสํ เอตํ มม.\csromanspacer{} นฺหารุํ เอตํ มม.\csromanspacer{} อฏฺฐิํ เอตํ มม.\csromanspacer{} อฏฺฐิมิญฺชํ เอตํ มม.\csromanspacer{} วกฺกํ เอตํ มม.\csromanspacer{} หทยํ เอตํ มม.\csromanspacer{} ยกนํ เอตํ มม.\csromanspacer{} กิโลมกํ เอตํ มม.\csromanspacer{} ปิหกํ เอตํ มม.\csromanspacer{} ปปฺผาสํ เอตํ มม.\csromanspacer{} อนฺตํ เอตํ มม.\csromanspacer{} อนฺตคุณํ เอตํ มม.\csromanspacer{} อุทริยํ เอตํ มม.\csromanspacer{} กรีสํ เอตํ มม.\csromanspacer{} ปิตฺตํ เอตํ มม.\csromanspacer{} เสมฺหํ เอตํ มม.\csromanspacer{} ปุพฺพํ เอตํ}

\vspace{\tipitakaparskip}
\csromancenter{ทิฏฺฐิกถา มม.\csromanspacer{} โลหิตํ เอตํ มม.\csromanspacer{} เสทํ เอตํ มม.\csromanspacer{} เมทํ เอตํ มม.\csromanspacer{} อสฺสุํ เอตํ มม.\csromanspacer{} วสํ เอตํ มม.\csromanspacer{} เขฬํ เอตํ มม.\csromanspacer{} สิงฺฆาณิกํ เอตํ มม.\csromanspacer{} ลสิกํ เอตํ มม.\csromanspacer{} มุตฺตํ เอตํ มม.\csromanspacer{} มตฺถลุงฺคํ “เอตํ มม, เอโสหมสฺมิ, เอโส เม อตฺตา”ติ อภินิเวสปรามาโส ทิฏฺฐิ.}
\prose{\csromanfolio{131}จกฺขายตนํ เอตํ มม.\csromanspacer{} รูปายตนํ เอตํ มม.\csromanspacer{} โสตายตนํ เอตํ มม.\csromanspacer{} สทฺทายตนํ เอตํ มม.\csromanspacer{} ฆานายตนํ เอตํ มม.\csromanspacer{} คนฺธายตนํ เอตํ มม.\csromanspacer{} ชิวฺหายตนํ เอตํ มม.\csromanspacer{} รสายตนํ เอตํ มม.\csromanspacer{} กายายตนํ เอตํ มม.\csromanspacer{} โผฏฺฐพฺพายตนํ เอตํ มม.\csromanspacer{} มนายตนํ เอตํ มม.\csromanspacer{} ธมฺมายตนํ เอตํ มม.}

\prose{จกฺขุธาตุํ เอตํ มม.\csromanspacer{} รูปธาตุํ เอตํ มม.\csromanspacer{} จกฺขุวิญฺญาณธาตุํ เอตํ มม.\csromanspacer{} โสตธาตุํ เอตํ มม.\csromanspacer{} สทฺทธาตุํ เอตํ มม.\csromanspacer{} โสตวิญฺญาณธาตุํ เอตํ มม.\csromanspacer{} ฆานธาตุํ เอตํ มม.\csromanspacer{} คนฺธธาตุํ เอตํ มม.\csromanspacer{} ฆานวิญฺญาณธาตุํ เอตํ มม.\csromanspacer{} ชิวฺหาธาตุํ เอตํ มม.\csromanspacer{} รสธาตุํ เอตํ มม.\csromanspacer{} ชิวฺหาวิญฺญาณธาตุํ เอตํ มม.\csromanspacer{} กายธาตุํ เอตํ มม.\csromanspacer{} โผฏฺฐพฺพธาตุํ เอตํ มม.\csromanspacer{} กายวิญฺญาณธาตุํ เอตํ มม.\csromanspacer{} มโนธาตุํ เอตํ มม.\csromanspacer{} ธมฺมธาตุํ เอตํ มม.\csromanspacer{} มโนวิญฺญาณธาตุํ “เอตํ มม, เอโสหมสฺมิ, เอโส เม อตฺตา”ติ อภินิเวสปรามาโส ทิฏฺฐิ.}

\prose{จกฺขุนฺทฺริยํ เอตํ มม.\csromanspacer{} โสตินฺทฺริยํ เอตํ มม.\csromanspacer{} ฆานินฺทฺริยํ เอตํ มม.\csromanspacer{} ชิวฺหินฺทฺริยํ เอตํ มม.\csromanspacer{} กายินฺทฺริยํ เอตํ มม, มนินฺทฺริยํ เอตํ มม.\csromanspacer{} ชีวิตินฺทฺริยํ เอตํ มม.\csromanspacer{} อิตฺถินฺทฺริยํ เอตํ มม.\csromanspacer{} ปุริสินฺทฺริยํ เอตํ มม.\csromanspacer{} สุขินฺทฺริยํ เอตํ มม.\csromanspacer{} ทุกฺขินฺทฺริยํ เอตํ มม.\csromanspacer{} โสมนสฺสินฺทฺริยํ เอตํ มม.\csromanspacer{} โทมนสฺสินฺทฺริยํ เอตํ มม.\csromanspacer{} อุเปกฺขินฺทฺริยํ เอตํ มม.\csromanspacer{} สทฺธินฺทฺริยํ เอตํ มม.\csromanspacer{} วีริยินฺทฺริยํ เอตํ มม.\csromanspacer{} สตินฺทฺริยํ เอตํ มม.\csromanspacer{} สมาธินฺทฺริยํ เอตํ มม.\csromanspacer{} ปญฺญินฺทฺริยํ “เอตํ มม, เอโสหมสฺมิ, เอโส เม อตฺตา”ติ อภินิเวสปรามาโส ทิฏฺฐิ.}

\prose{กามธาตุํ เอตํ มม.\csromanspacer{} รูปธาตุํ เอตํ มม.\csromanspacer{} อรูปธาตุํ เอตํ มม.\csromanspacer{} กามภวํ เอตํ มม.\csromanspacer{} รูปภวํ เอตํ มม.\csromanspacer{} อรูปภวํ เอตํ มม.\csromanspacer{} สญฺญาภวํ เอตํ มม.\csromanspacer{} อสญฺญาภวํ เอตํ มม.\csromanspacer{} เนวสญฺญานาสญฺญาภวํ เอตํ มม.\csromanspacer{} เอกโวการภวํ เอตํ มม.\csromanspacer{} จตุโวการภวํ เอตํ มม.\csromanspacer{} ปญฺจโวการภวํ เอตํ มม.\csromanspacer{}\csromanspacer{}ปฐมชฺฌานํ เอตํ มม.\csromanspacer{} ทุติยชฺฌานํ เอตํ มม.\csromanspacer{} ตติยชฺฌานํ เอตํ มม.\csromanspacer{} จตุตฺถชฺฌานํ เอตํ มม.\csromanspacer{} เมตฺตํ เจโตวิมุตฺติํ เอตํ มม.\csromanspacer{} กรุณํ เจโตวิมุตฺติํ เอตํ มม.\csromanspacer{} มุทิตํ เจโตวิมุตฺติํ เอตํ มม.\csromanspacer{} อุเปกฺขํ เจโตวิมุตฺติํ เอตํ มม.\csromanspacer{}\csromanspacer{}อากาสานญฺจายตนสมาปตฺติํ เอตํ มม. \csromanfolio{132}วิญฺญาณญฺจายตนสมาปตฺติํ เอตํ มม.\csromanspacer{} อากิญฺจญฺญายตนสมาปตฺติํ เอตํ มม.\csromanspacer{} เนวสญฺญานาสญฺญายตนสมาปตฺติํ “เอตํ มม, เอโสหมสฺมิ, เอโส เม อตฺตา”ติ อภินิเวสปรามาโส ทิฏฺฐิ.}

\prose{อวิชฺชํ เอตํ มม.\csromanspacer{} สงฺขาเร เอตํ มม.\csromanspacer{} วิญฺญาณํ เอตํ มม.\csromanspacer{} นามรูปํ เอตํ มม.\csromanspacer{} สฬายตนํ เอตํ มม.\csromanspacer{} ผสฺสํ เอตํ มม.\csromanspacer{} เวทนํ เอตํ มม.\csromanspacer{} ตณฺหํ เอตํ มม.\csromanspacer{} อุปาทานํ เอตํ มม.\csromanspacer{} ภวํ เอตํ มม.\csromanspacer{} ชาติํ เอตํ มม.\csromanspacer{} ชรามรณํ “เอตํ มม, เอโสหมสฺมิ, เอโส เม อตฺตา”ติ อภินิเวสปรามาโส ทิฏฺฐิ.\csromanspacer{} เอวํ อภินิเวสปรามาโส ทิฏฺฐิ.}

\proseitem{124}{กตมานิ \textbf{อฏฺฐ ทิฏฺฐิฏฺฐานานิ?} ขนฺธาปิ ทิฏฺฐิฏฺฐานํ, อวิชฺชาปิ ทิฏฺฐิฏฺฐานํ, ผสฺโสปิ ทิฏฺฐิฏฺฐานํ, สญฺญาปิ ทิฏฺฐิฏฺฐานํ, วิตกฺโกปิ\footnote{วิตกฺกาปิ (สฺยา)} ทิฏฺฐิฏฺฐานํ, อโยนิโส มนสิกาโรปิ ทิฏฺฐิฏฺฐานํ, ปาปมิตฺโตปิ ทิฏฺฐิฏฺฐานํ, ปรโตโฆโสปิ ทิฏฺฐิฏฺฐานํ.}

\prose{ขนฺธา เหตุ ขนฺธา ปจฺจโย ทิฏฺฐิฏฺฐานํ อุปาทาย สมุฏฺฐานฏฺเฐน, เอวํ ขนฺธาปิ ทิฏฺฐิฏฺฐานํ.\csromanspacer{} อวิชฺชา เหตุ อวิชฺชา ปจฺจโย ทิฏฺฐิฏฺฐานํ อุปาทาย สมุฏฺฐานฏฺเฐน, เอวํ อวิชฺชาปิ ทิฏฺฐิฏฺฐานํ.\csromanspacer{} ผสฺโส เหตุ ผสฺโส ปจฺจโย ทิฏฺฐิฏฺฐานํ อุปาทาย สมุฏฺฐานฏฺเฐน, เอวํ ผสฺโสปิ ทิฏฺฐิฏฺฐานํ.\csromanspacer{} สญฺญา เหตุ สญฺญา ปจฺจโย ทิฏฺฐิฏฺฐานํ อุปาทาย สมุฏฺฐานฏฺเฐน, เอวํ สญฺญาปิ ทิฏฺฐิฏฺฐานํ.\csromanspacer{} วิตกฺโก เหตุ วิตกฺโก ปจฺจโย ทิฏฺฐิฏฺฐานํ อุปาทาย สมุฏฺฐานฏฺเฐน, เอวํ วิตกฺโกปิ ทิฏฺฐิฏฺฐานํ.\csromanspacer{} อโยนิโส มนสิกาโร เหตุ อโยนิโส มนสิกาโร ปจฺจโย ทิฏฺฐิฏฺฐานํ อุปาทาย สมุฏฺฐานฏฺเฐน, เอวํ อโยนิโส มนสิกาโรปิ ทิฏฺฐิฏฺฐานํ.\csromanspacer{} ปาปมิตฺโต เหตุ ปาปมิตฺโต ปจฺจโย ทิฏฺฐิฏฺฐานํ อุปาทาย สมุฏฺฐานฏฺเฐน, เอวํ ปาปมิตฺโตปิ ทิฏฺฐิฏฺฐานํ.\csromanspacer{} ปรโตโฆโส เหตุ ปรโตโฆโส ปจฺจโย ทิฏฺฐิฏฺฐานํ อุปาทาย สมุฏฺฐานฏฺเฐน, เอวํ ปรโตโฆโสปิ ทิฏฺฐิฏฺฐานํ.\csromanspacer{} อิมานิ อฏฺฐ ทิฏฺฐิฏฺฐานานิ.}

\proseitem{125}{กตมานิ \textbf{อฏฺฐารส ทิฏฺฐิปริยุฏฺฐานานิ}? ยา ทิฏฺฐิ ทิฏฺฐิคตํ ทิฏฺฐิคหนํ ทิฏฺฐิกนฺตารํ ทิฏฺฐิวิสูกํ ทิฏฺฐิวิปฺผนฺทิตํ ทิฏฺฐิสญฺโญชนํ ทิฏฺฐิสลฺลํ ทิฏฺฐิสมฺพาโธ ทิฏฺฐิปลิโพโธ ทิฏฺฐิพนฺธนํ ทิฏฺฐิปปาโต ทิฏฺฐานุสโย ทิฏฺฐิสนฺตาโป}

\vspace{\tipitakaparskip}
\csromancenter{ทิฏฺฐิกถา ทิฏฺฐิปริฬาโห ทิฏฺฐิคนฺโถ ทิฏฺฐุปาทานํ ทิฏฺฐาภินิเวโส ทิฏฺฐิปรามาโส อิมานิ อฏฺฐารส ทิฏฺฐิปริยุฏฺฐานานิ.}
\proseitem{126}{\csromanfolio{133}กตมา \textbf{โสฬส ทิฏฺฐิโย}? อสฺสาททิฏฺฐิ อตฺตานุทิฏฺฐิ มิจฺฉาทิฏฺฐิ สกฺกายทิฏฺฐิ สกฺกายวตฺถุกา สสฺสตทิฏฺฐิ สกฺกายวตฺถุกา อุจฺเฉททิฏฺฐิ อนฺตคฺคาหิกา ทิฏฺฐิ ปุพฺพนฺตานุทิฏฺฐิ อปรนฺตานุทิฏฺฐิ สญฺโญชนิกา ทิฏฺฐิ “อหนฺ”ติ มานวินิพนฺธาทฺùฐิ “มมนฺ”ติ มานวินิพนฺธา ทิฏฺฐิ อตฺตวาทปฏิสํยุตฺตา ทิฏฺฐิ โลกวาทปฏิสํยุตฺตา ทิฏฺฐิ ภวทิฏฺฐิ วิภวทิฏฺฐิ, อิมา \textbf{โสฬส ทิฏฺฐิโย}.}

\proseitem{127}{กตเม \textbf{ตีณิ สตํ ทิฏฺฐาภินิเวสา}? อสฺสาททิฏฺฐิยา กติหากาเรหิ อภินิเวโส โหติ, อตฺตานุทิฏฺฐิยา กติหากาเรหิ อภินิเวโส โหติ, มิจฺฉาทิฏฺฐิยา กติหากาเรหิ อภินิเวโส โหติ, สกฺกายทิฏฺฐิยา กติหากาเรหิ อภินิเวโส โหติ, สกฺกายวตฺถุกาย สสฺสตทิฏฺฐิยา กติหากาเรหิ อภินิเวโส โหติ, สกฺกายวตฺถุกาย อุจฺเฉททิฏฺฐิยา กติหากาเรหิ อภินิเวโส โหติ, อนฺตคฺคาหิกาย ทิฏฺฐิยา กติหากาเรหิ อภินิเวโส โหติ, ปุพฺพนฺตานุทิฏฺฐิยา กติหากาเรหิ อภินิเวโส โหติ, อปรนฺตานุทิฏฺฐิยา กติหากาเรหิ อภินิเวโส โหติ, สญฺโญชนิกาย ทิฏฺฐิยา กติหากาเรหิ อภินิเวโส โหติ, “อหนฺ”ติ มานวินิพนฺธาย ทิฏฺฐิยา กติหากาเรหิ อภินิเวโส โหติ, “มมนฺ”ติ มานวินิพนฺธาย ทิฏฺฐิยา กติหากาเรหิ อภินิเวโส โหติ, อตฺตวาทปฏิสํยุตฺตาย ทิฏฺฐิยา กติหากาเรหิ อภินิเวโส โหติ, โลกวาทปฏิสํยุตฺตาย ทิฏฺฐิยา กติหากาเรหิ อภินิเวโส โหติ, ภวทิฏฺฐิยา กติหากาเรหิ อภินิเวโส โหติ, วิภวทิฏฺฐิยา กติหากาเรหิ อภินิเวโส โหติ?}

\prose{อสฺสาททิฏฺฐิยา ปญฺจติํสาย อากาเรหิ อภินิเวโส โหติ, อตฺตานุทิฏฺฐิยา วีสติยา อากาเรหิ อภินิเวโส โหติ, มิจฺฉาทิฏฺฐิยา ทสหากาเรหิ อภินิเวโส โหติ, สกฺกายทิฏฺฐิยา วีสติยา อากาเรหิ อภินิเวโส โหติ, สกฺกายวตฺถุกาย สสฺสตทิฏฺฐิยา ปนฺนรสหิ อากาเรหิ อภินิเวโส โหติ, สกฺกายวตฺถุกาย อุจฺเฉททิฏฺฐิยา ปญฺจหากาเรหิ อภินิเวโส โหติ, อนฺตคฺคาหิกาย ทิฏฺฐิยา ปญฺญาสาย อากาเรหิ อภินิเวโส \csromanfolio{134}โหติ, ปุพฺพนฺตานุทิฏฺฐิยา อฏฺฐารสหิ อากาเรหิ อภินิเวโส โหติ, อปรนฺตานุทิฏฺฐิยา จตุจตฺตาลีสาย อากาเรหิ อภินิเวโส โหติ, สญฺโญชนิกาย ทิฏฺฐิยา อฏฺฐารสหิ อากาเรหิ อภินิเวโส โหติ, “อหนฺ”ติ มานวินิพนฺธาย ทิฏฺฐิยา อฏฺฐารสหิ อากาเรหิ อภินิเวโส โหติ, “มมนฺ”ติ มานวินิพนฺธาย ทิฏฺฐิยา อฏฺฐารสหิ อากาเรหิ อภินิเวโส โหติ, อตฺตวาทปฏิสํยุตฺตาย ทิฏฺฐิยา วีสติยา อากาเรหิ อภินิเวโส โหติ, โลกวาทปฏิสํยุตฺตาย ทิฏฺฐิยา อฏฺฐหิ อากาเรหิ อภินิเวโส โหติ, ภวทิฏฺฐิยา เอเกน อากาเรน\footnote{เอกูนวีสติยา อากาเรหิ (สฺยา)} อภินิเวโส โหติ, วิภวทิฏฺฐิยา เอเกน อากาเรน1 อภินิเวโส โหติ.}

\csromanmatikaanchor{mtk.52}
\csromantocmarkref{subsection}{1. อสฺสาททิฏฺฐินิทฺเทส}{mtk.52}
\bookmark[dest={mtk.52},level=2]{1. อสฺสาททิฏฺฐินิทฺเทส}
\csromanheader{2}{1. อสฺสาททิฏฺฐินิทฺเทส}
\proseitem{128}{อสฺสาททิฏฺฐิยา กตเมหิ ปญฺจติํสาย อากาเรหิ อภินิเวโส โหติ? ยํ รูปํ ปฏิจฺจ อุปฺปชฺชติ สุขํ โสมนสฺสํ, “อยํ รูปสฺส อสฺสาโท”ติ อภินิเวสปรามาโส ทิฏฺฐิ.\csromanspacer{} ทิฏฺฐิ น อสฺสาโท, อสฺสาโท น ทิฏฺฐิ.\csromanspacer{} อญฺญา ทิฏฺฐิ, อญฺโญ อสฺสาโท.\csromanspacer{} ยา จ ทิฏฺฐิ โย จ อสฺสาโท, อยํ วุจฺจติ อสฺสาททิฏฺฐิ.}

\prose{อสฺสาททิฏฺฐิ มิจฺฉาทิฏฺฐิ ทิฏฺฐิวิปตฺติ, ตาย ทิฏฺฐิวิปตฺติยา สมนฺนาคโต ปุคฺคโล ทิฏฺฐิวิปนฺโน.\csromanspacer{} ทิฏฺฐิวิปนฺโน ปุคฺคโล น เสวิตพฺโพ น ภชิตพฺโพ น ปยิรุปาสิตพฺโพ, ตํ กิสฺส เหตุ, ทิฏฺฐิ หิสฺส ปาปิกา.\csromanspacer{} โย ทิฏฺฐิยา ราโค\footnote{ยา ทิฏฺฐิ โย ราโค (สฺยา)}, โส น ทิฏฺฐิ, ทิฏฺฐิ น ราโค.\csromanspacer{} อญฺญา ทิฏฺฐิ, อญฺโญ ราโค.\csromanspacer{} ยา จ ทิฏฺฐิ โย จ ราโค, อยํ วุจฺจติ ทิฏฺฐิราโค.\csromanspacer{} ตาย จ ทิฏฺฐิยา เตน จ ราเคน สมนฺนาคโต ปุคฺคโล ทิฏฺฐิราครตฺโต.\csromanspacer{} ทิฏฺฐิราครตฺเต ปุคฺคเล ทินฺนํ ทานํ น มหปฺผลํ โหติ น มหานิสํสํ.\csromanspacer{} ตํ กิสฺส เหตุ, ทิฏฺฐิ หิสฺส ปาปิกา, อสฺสาททิฏฺฐิ มิจฺฉาทิฏฺฐิ.}

\prose{มิจฺฉาทิฏฺฐิกสฺส ปุริสปุคฺคลสฺส เทฺวว คติโย นิรโย วา ติรจฺฉานโยนิ วา.\csromanspacer{} มิจฺฉาทิฏฺฐิกสฺส ปุริสปุคฺคลสฺส ยญฺเจว กายกมฺมํ ยถาทิฏฺฐิ สมตฺตํ สมาทินฺนํ, ยญฺจ วจีกมฺมํ ฯเปฯ ยญฺจ มโนกมฺมํ ยถาทิฏฺฐิ สมตฺตํ สมาทินฺนํ, ยา จ เจตนา ยา จ ปตฺถนา โย จ ปณิธิ เย จ สงฺขารา,}

\vspace{\tipitakaparskip}
\csromancenter{ทิฏฺฐิกถา สพฺเพ เต ธมฺมา อนิฏฺฐาย อกนฺตาย อมนาปาย อหิตาย ทุกฺขาย สํวตฺตนฺติ.\csromanspacer{} ตํ กิสฺส เหตุ, ทิฏฺฐิ หิสฺส ปาปิกา.\csromanspacer{} เสยฺยถาปิ นิมฺพพีชํ วา โกสาตกีพีชํ วา ติตฺตกาลาพุพีชํ วา อลฺลาย ปถวิยา นิกฺขิตฺตํ ยญฺเจว ปถวิรสํ อุปาทิยติ, ยญฺจ อาโปรสํ อุปาทิยติ, สพฺพํ ตํ ติตฺตกตฺตาย กฏุกตฺตาย อสาตตฺตาย\footnote{อสารตาย (สฺยา)} สํวตฺตติ.\csromanspacer{} ตํ กิสฺส เหตุ, พีชํ หิสฺส ปาปิกํ.\csromanspacer{} เอวเมวํ มิจฺฉาทิฏฺฐิกสฺส ปุริสปุคฺคลสฺส ยญฺเจว กายกมฺมํ ยถาทิฏฺฐิ สมตฺตํ สมาทินฺนํ, ยญฺจ วจีกมฺมํ ฯเปฯ ยญฺจ มโนกมฺมํ ยถาทิฏฺฐิ สมตฺตํ สมาทินฺนํ, ยา จ เจตนา ยา จ ปตฺถนา โย จ ปณิธิ เย จ สงฺขารา, สพฺเพ เต ธมฺมา อนิฏฺฐาย อกนฺตาย อมนาปาย อหิตาย ทุกฺขาย สํวตฺตนฺติ.\csromanspacer{} ตํ กิสฺส เหตุ, ทิฏฺฐิ หิสฺส ปาปิกา.\csromanspacer{} อสฺสาททิฏฺฐิ มิจฺฉาทิฏฺฐิ.}
\prose{\csromanfolio{135}มิจฺฉาทิฏฺฐิ ทิฏฺฐิคตํ ทิฏฺฐิคหนํ ทิฏฺฐิกนฺตารํ ทิฏฺฐิวิสูกํ ทิฏฺฐิวิปฺผนฺทิตํ ทิฏฺฐิสญฺโญชนํ ทิฏฺฐิสลฺลํ ทิฏฺฐิสมฺพาโธ ทิฏฺฐิปลิโพโธ ทิฏฺฐิพนฺธนํ ทิฏฺฐิปปาโต ทิฏฺฐานุสโย ทิฏฺฐิสนฺตาโป ทิฏฺฐิปริฬาโห ทิฏฺฐิคนฺโถ ทิฏฺฐุปาทานํ ทิฏฺฐาภินิเวโส ทิฏฺฐิปรามาโส, อิเมหิ อฏฺฐารสหิ อากาเรหิ ปริยุฏฺฐิตจิตฺตสฺส สญฺโญโค.}

\proseitem{129}{อตฺถิ สญฺโญชนานิ เจว ทิฏฺฐิโย จ, อตฺถิ สญฺโญชนานิ น จ ทิฏฺฐิโย.\csromanspacer{} กตมานิ สญฺโญชนานิ เจว ทิฏฺฐิโย จ? สกฺกายทิฏฺฐิ สีลพฺพตปรามาโส, อิมานิ สญฺโญชนานิ เจว ทิฏฺฐิโย จ.\csromanspacer{}\csromanspacer{}กตมานิ สญฺโญชนานิ น จ ทิฏฺฐิโย? กามราคสญฺโญชนํ ปฏิฆสญฺโญชนํ มานสญฺโญชนํ วิจิกิจฺฉาสญฺโญชนํ ภวราคสญฺโญชนํ อิสฺสาสญฺโญชนํ มจฺฉริยสญฺโญชนํ อนุนยสญฺโญชนํ อวิชฺชาสญฺโญชนํ, อิมานิ สญฺโญชนานิ น จ ทิฏฺฐิโย.}

\prose{ยํ เวทนํ ปฏิจฺจ ฯเปฯ.\csromanspacer{} ยํ สญฺญํ ปฏิจฺจ ฯเปฯ.\csromanspacer{} ยํ สงฺขาเร\footnote{เย สงฺขาเร (ก) สํ 2. 24 ปิฏฺเฐ ปสฺสิตพฺพา.} ปฏิจฺจ ฯเปฯ.\csromanspacer{} ยํ วิญฺญาณํ ปฏิจฺจ.\csromanspacer{}\csromanspacer{}ยํ จกฺขุํ ปฏิจฺจ.\csromanspacer{} ยํ โสตํ ปฏิจฺจ.\csromanspacer{} ยํ ฆานํ ปฏิจฺจ.\csromanspacer{} ยํ ชิวฺหํ ปฏิจฺจ.\csromanspacer{} ยํ กายํ ปฏิจฺจ.\csromanspacer{} ยํ มนํ ปฏิจฺจ.\csromanspacer{} ยํ รูเป ปฏิจฺจ.\csromanspacer{} ยํ สทฺเท ปฏิจฺจ.\csromanspacer{} ยํ คนฺเธ ปฏิจฺจ.\csromanspacer{} ยํ รเส ปฏิจฺจ.\csromanspacer{} ยํ โผฏฺฐพฺเพ ปฏิจฺจ.\csromanspacer{} ยํ ธมฺเม ปฏิจฺจ.\csromanspacer{}\csromanspacer{}ยํ จกฺขุวิญฺญาณํ ปฏิจฺจ.\csromanspacer{} ยํ โสตวิญฺญาณํ ปฏิจฺจ.\csromanspacer{} ยํ ฆานวิญฺญาณํ ปฏิจฺจ.\csromanspacer{} ยํ ชิวฺหาวิญฺญาณํ ปฏิจฺจ.\csromanspacer{} ยํ กายวิญฺญาณํ ปฏิจฺจ.\csromanspacer{} ยํ มโนวิญฺญาณํ \csromanfolio{136}ปฏิจฺจ.\csromanspacer{}\csromanspacer{}ยํ จกฺขุสมฺผสฺสํ ปฏิจฺจ.\csromanspacer{} ยํ โสตสมฺผสฺสํ ปฏิจฺจ.\csromanspacer{} ยํ ฆานสมฺผสฺสํ ปฏิจฺจ.\csromanspacer{} ยํ ชิวฺหาสมฺผสฺสํ ปฏิจฺจ.\csromanspacer{} ยํ กายสมฺผสฺสํ ปฏิจฺจ.\csromanspacer{} ยํ มโนสมฺผสฺสํ ปฏิจฺจ.\csromanspacer{}\csromanspacer{}ยํ จกฺขุสมฺผสฺสชํ เวทนํ ปฏิจฺจ.\csromanspacer{} ยํ โสตสมฺผสฺสชํ เวทนํ ปฏิจฺจ.\csromanspacer{} ยํ ฆานสมฺผสฺสชํ เวทนํ ปฏิจฺจ.\csromanspacer{} ยํ ชิวฺหาสมฺผสฺสชํ เวทนํ ปฏิจฺจ.\csromanspacer{} ยํ กายสมฺผสฺสชํ เวทนํ ปฏิจฺจ.\csromanspacer{} ยํ มโนสมฺผสฺสชํ เวทนํ ปฏิจฺจ อุปฺปชฺชติ สุขํ โสมนสฺสํ “อยํ มโนสมฺผสฺสชาย เวทนาย อสฺสาโท”ติ อภินิเวสปรามาโส ทิฏฺฐิ.\csromanspacer{} ทิฏฺฐิ น อสฺสาโท, อสฺสาโท น ทิฏฺฐิ.\csromanspacer{} อญฺญา ทิฏฺฐิ, อญฺโญ อสฺสาโท.\csromanspacer{} ยา จ ทิฏฺฐิ โย จ อสฺสาโท, อยํ วุจฺจติ อสฺสาททิฏฺฐิ.}

\prose{อสฺสาททิฏฺฐิ มิจฺฉาทิฏฺฐิ ทิฏฺฐิวิปตฺติ, ตาย ทิฏฺฐิวิปตฺติยา สมานฺนาคโต ปุคฺคโล ทิฏฺฐิวิปนฺโน.\csromanspacer{} ทิฏฺฐิวิปนฺโน ปุคฺคโล น เสวิตพฺโพ น ภชิตพฺโพ น ปยิรุปาสิตพฺโพ.\csromanspacer{} ตํ กิสฺส เหตุ, ทิฏฺฐิ หิสฺส ปาปิกา.\csromanspacer{} โย ทิฏฺฐิยา ราโค, โส น ทิฏฺฐิ, ทิฏฺฐิ น ราโค.\csromanspacer{} อญฺญา ทิฏฺฐิ, อญฺโญ ราโค.\csromanspacer{} ยา จ ทิฏฺฐิ โย จ ราโค, อยํ วุจฺจติ ทิฏฺฐิราโค.\csromanspacer{} ตาย จ ทิฏฺฐิยา เตน จ ราเคน สมนฺนาคโต ปุคฺคโล ทิฏฺฐิราครตฺโต.\csromanspacer{} ทิฏฺฐิราครตฺเต ปุคฺคเล ทินฺนํ ทานํ น มหปฺผลํ โหติ น มหานิสํสํ.\csromanspacer{} ตํ กิสฺส เหสุ, ทิฏฺฐิ หิสฺส ปาปิกา.\csromanspacer{} อสฺสาททิฏฺฐิ มิจฺฉาทิฏฺฐิ.}

\prose{มิจฺฉาทิฏฺฐิกสฺส ปุริสปุคฺคลสฺส เทฺวว คติโย นิรโย วา ติรจฺฉานโยนิ วา.\csromanspacer{} มิจฺฉาทิฏฺฐิกสฺส ปุริสปุคฺคลสฺส ยญฺเจว กายกมฺมํ ยถาทิฏฺฐิ สมตฺตํ สมาทินฺนํ, ยญฺจ วจีกมฺมํ ฯเปฯ ยญฺจ มโนกมฺมํ ยถาทิฏฺฐิ สมตฺตํ สมาทินฺนํ, ยา จ เจตนา ยา จ ปตฺถนา โย จ ปณิธิ เย จ สงฺขารา, สพฺเพ เต ธมฺมา อนิฏฺฐาย อกนฺตาย อมนาปาย อหิตาย ทุกฺขาย สํวตฺตนฺติ.\csromanspacer{} ตํ กิสฺส เหตุ, ทิฏฺฐิ หิสฺส ปาปิกา.\csromanspacer{} เสยฺยถาปิ นิมฺพพีชํ วา โกสาตกีพีชํ วา ติตฺตกาลาพุพีชํ วา อลฺลาย ปถวิยา นิกฺขิตฺตํ ยญฺเจว ปถวิรสํ อุปาทิยติ, ยญฺจ อาโปรสํ อุปาทิยติ, สพฺพํ ตํ ติตฺตกตฺตาย กฏุกตฺตาย อสาตตฺตาย สํวตฺตติ.\csromanspacer{} ตํ กิสฺส เหตุ, พีชํ หิสฺส ปาปิกํ.\csromanspacer{} เอวเมวํ มิจฺฉาทิฏฺฐิกสฺส ปุริสปุคฺคลสฺส ยญฺเจว กายกมฺมํ ยถาทิฏฺฐิ สมตฺตํ สมาทินฺนํ, ยญฺจ วจีกมฺมํ ฯเปฯ ยญฺจ มโนกมฺมํ ยถาทิฏฺฐิ สมตฺตํ สมาทินฺนํ, ยา จ เจตนา ยา จ ปตฺถนา โย จ ปณิธิ เย จ สงฺขารา, สพฺเพ เต ธมฺมา อนิฏฺฐาย อกนฺตาย}

\vspace{\tipitakaparskip}
\csromancenter{ทิฏฺฐิกถา อมนาปาย อหิตาย ทุกฺขาย สํวตฺตนฺติ.\csromanspacer{} ตํ กิสฺส เหตุ, ทิฏฺฐิ หิสฺส ปาปิกา.\csromanspacer{} อสฺสาททิฏฺฐิ มิจฺฉาทิฏฺฐิ.}
\prose{\csromanfolio{137}มิจฺฉาทิฏฺฐิ ทิฏฺฐิคตํ ทิฏฺฐิคหนํ ฯเปฯ ทิฏฺฐาภินิเวโส ทิฏฺฐิปรามาโส, อิเมหิ อฏฺฐารสหิ อากาเรหิ ปริยุฏฺฐิตจิตฺตสฺส สญฺโญโค.}

\prose{อตฺถิ สญฺโญชนานิ เจว ทิฏฺฐิโย จ, อตฺถิ สญฺโญชนานิ น จ ทิฏฺฐิโย.\csromanspacer{} กตมานิ สญฺโญชนานิ เจว ทิฏฺฐิโย จ? สกฺกายทิฏฺฐิ สีลพฺพตปรามาโส, อิมานิ สญฺโญชนานิ เจว ทิฏฺฐิโย จ.\csromanspacer{}\csromanspacer{}กตมานิ สญฺโญชนานิ น จ ทิฏฺฐิโย? กามราคสญฺโญชนํ ปฏิฆสญฺโญชนํ มานสญฺโญชนํ วิจิกิจฺฉาสญฺโญชนํ ภวราคสญฺโญชนํ อิสฺสาสญฺโญชนํ มจฺฉริยสญฺโญชนํ อนุนยสญฺโญชนํ อวิชฺชาสญฺโญชนํ, อิมานิ สญฺโญชนานิ น จ ทิฏฺฐิโย.\csromanspacer{} อสฺสาททิฏฺฐิยา อิเมหิ ปญฺจติํสาย อากาเรหิ อภินิเวโส โหติ.}

\csromanheader{2}{อสฺสาททิฏฺฐินิทฺเทโส ปฐโม.}
\csromansectionrule
\csromanmatikaanchor{mtk.53}
\csromantocmarkref{subsection}{2. อตฺตานุทิฏฺฐินิทฺเทส}{mtk.53}
\bookmark[dest={mtk.53},level=2]{2. อตฺตานุทิฏฺฐินิทฺเทส}
\csromanheader{2}{2. อตฺตานุทิฏฺฐินิทฺเทส}
\proseitem{130}{อตฺตานุทิฏฺฐิยา กตเมหิ วีสติยา อากาเรหิ อภินิเวโส โหติ? อิธ อสฺสุตวา ปุถุชฺชโน อริยานํ อทสฺสาวี อริยธมฺมสฺส อโกวิโท อริยธมฺเม อวินีโต สปฺปุริสานํ อทสฺสาวี สปฺปุริสธมฺมสฺส อโกวิโท สปฺปุริสธมฺเม อวินีโต รูปํ อตฺตโต สมนุปสฺสติ, รูปวนฺตํ วา อตฺตานํ, อตฺตนิ วา รูปํ, รูปสฺมิํ วา อตฺตานํ.\csromanspacer{}\csromanspacer{}เวทนํ ฯเปฯ.\csromanspacer{} สญฺญํ ฯเปฯ.\csromanspacer{} สงฺขาเร ฯเปฯ.\csromanspacer{} วิญฺญาณํ อตฺตโต สมนุปสฺสติ, วิญฺญาณวนฺตํ วา อตฺตานํ, อตฺตนิ วา วิญฺญาณํ, วิญฺญาณสฺมิํ วา อตฺตานํ.}

\proseitem{131}{กถํ \textbf{รูปํ อตฺตโต สมนุปสฺสติ}? อิเธกจฺโจ ปถวีกสิณํ อตฺตโต สมนุปสฺสติ “ยํ ปถวีกสิณํ, โส อหํ.\csromanspacer{} โย อหํ, ตํ ปถวีกสิณนฺ”ติ ปถวีกสิณญฺจ อตฺตญฺจ อทฺวยํ สมนุปสฺสติ.\csromanspacer{} เสยฺยถาปิ เตลปฺปทีปสฺส ฌายโต ยา อจฺจิ, โส วณฺโณ.\csromanspacer{} โย วณฺโณ, สา อจฺจีติ อจฺจิญฺจ วณฺณญฺจ อทฺวยํ สมนุปสฺสติ.\csromanspacer{} เอวเมวํ อิเธกจฺโจ ปถวีกสิณํ อตฺตโต สมนุปสฺสติ “ยํ ปถวีกสิณํ, โส อหํ.\csromanspacer{} โย \csromanfolio{138}อหํ, ตํ ปถวีกสิณนฺ”ติ ปถวีกสิณญฺจ อตฺตญฺจ อทฺวยํ สมนุปสฺสติ.\csromanspacer{} อภินิเวสปรามาโส ทิฏฺฐิ, ทิฏฺฐิ น วตฺถุ, วตฺถุ น ทิฏฺฐิ.\csromanspacer{} อญฺญา ทิฏฺฐิ, อญฺญํ วตฺถุ.\csromanspacer{} ยา จ ทิฏฺฐิ ยญฺจ วตฺถุ, อยํ ปฐมา รูปวตฺถุกา อตฺตานุทิฏฺฐิ.\csromanspacer{} อตฺตานุทิฏฺฐิ มิจฺฉาทิฏฺฐิ ทิฏฺฐิวิปตฺติ ฯเปฯ อตฺตานุทิฏฺฐิ มิจฺฉาทิฏฺฐิ.\csromanspacer{} มิจฺฉาทิฏฺฐิกสฺส ปุริสปุคฺคลสฺส เทฺวว คติโย ฯเปฯ อิมานิ สญฺโญชนานิ น จ ทิฏฺฐิโย.}

\prose{อิเธกจฺโจ อาโปกสิณํ.\csromanspacer{} เตโชกสิณํ.\csromanspacer{} วาโยกสิณํ.\csromanspacer{} นีลกสิณํ.\csromanspacer{} ปีตกสิณํ.\csromanspacer{} โลหิตกสิณํ.\csromanspacer{} โอทาตกสิณํ อตฺตโต สมนุปสฺสติ “ยํ โอทาตกสิณํ, โส อหํ.\csromanspacer{} โย อหํ, ตํ โอทาตกสิณนฺ”ติ โอทาตกสิณญฺจ, อตฺตญฺจ อทฺวยํ สมนุปสฺสติ.\csromanspacer{} เสยฺยถาปิ เตลปฺปทีปสฺส ฌายโต ยา อจฺจิ, โส วณฺโณ.\csromanspacer{} โย วณฺโณ, สา อจฺจีติ อจฺจิญฺจ วณฺณญฺจ อทฺวยํ สมนุปสฺสติ.\csromanspacer{} เอวเมวํ อิเธกจฺโจ ฯเปฯ.\csromanspacer{} โอทาตกสิณญฺจ อตฺตญฺจ อทฺวยํ สมนุปสฺสติ.\csromanspacer{} อภินิเวสปรามาโส ทิฏฺฐิ, ทิฏฺฐิ น วตฺถุ, วตฺถุ น ทิฏฺฐิ.\csromanspacer{} อญฺญา ทิฏฺฐิ, อญฺญํ วตฺถุ.\csromanspacer{} ยา จ ทิฏฺฐิ ยญฺจ วตฺถุ, อยํ ปฐมา รูปวตฺถุกา อตฺตานุทิฏฺฐิ.\csromanspacer{} อตฺตานุทิฏฺฐิ มิจฺฉาทิฏฺฐิ ทิฏฺฐิวิปตฺติ ฯเปฯ อิมานิ สญฺโญชนานิ น จ ทิฏฺฐิโย.\csromanspacer{} เอวํ รูปํ อตฺตโต สมนุปสฺสติ. (1)}

\prose{กถํ \textbf{รูปวนฺตํ อตฺตานํ สมนุปสฺสติ}? อิเธกจฺโจ เวทนํ สญฺญํ สงฺขาเร วิญฺญาณํ อตฺตโต สมนุปสฺสติ, ตสฺส เอวํ โหติ “อยํ โข เม อตฺตา, โส โข ปน เม อยํ อตฺตา อิมินา รูเปน รูปวา”ติ \textbf{รูปวนฺตํ อตฺตานํ สมนุปสฺสติ}.\csromanspacer{} เสยฺยถาปิ รุกฺโข ฉายาสมฺปนฺโน อสฺส.\csromanspacer{} ตเมนํ ปุริโส เอวํ วเทยฺย “อยํ รุกฺโข, อยํ ฉายา.\csromanspacer{} อญฺโญ รุกฺโข, อญฺญา ฉายา.\csromanspacer{} โส โข ปนายํ รุกฺโข อิมาย ฉายาย ฉายาวา”ติ ฉายาวนฺตํ รุกฺขํ สมนุปสฺสติ.\csromanspacer{} เอวเมวํ อิเธกจฺโจ เวทนํ.\csromanspacer{} สญฺญํ.\csromanspacer{} สงฺขาเร.\csromanspacer{} วิญฺญาณํ อตฺตโต สมนุปสฺสติ, ตสฺส เอวํ โหติ “อยํ โข เม อตฺตา, โส โข ปน อยํ อตฺตา อิมินา รูเปน รูปวา”ติ \textbf{รูปวนฺตํ อตฺตานํ สมนุปสฺสติ}.\csromanspacer{} อภินิเวสปรามาโส ทิฏฺฐิ, ทิฏฺฐิ น วตฺถุ, วตฺถุ น ทิฏฺฐิ.\csromanspacer{} อญฺญา ทิฏฺฐิ, อญฺญํ วตฺถุ.\csromanspacer{} ยา จ ทิฏฺฐิ ยญฺจ วตฺถุ, อยํ ทุติยา รูปวตฺถุกา อตฺตานุทิฏฺฐิ.\csromanspacer{} อตฺตานุทิฏฺฐิ มิจฺฉาทิฏฺฐิ ทิฏฺฐิวิปตฺติ ฯเปฯ อิมานิ สญฺโญชนานิ น จ ทิฏฺฐิโย.\csromanspacer{} เอวํ \textbf{รูปวนฺตํ อตฺตานํ สมนุปสฺสติ}. (2)}

\csromanheader{2}{2. ทิฏฺฐิกถา}
\prose{\csromanfolio{139}กถํ \textbf{อตฺตนิ รูปํ สมนุปสฺสติ}? อิเธกจฺโจ เวทนํ สญฺญํ สงฺขาเร วิญฺญาณํ อตฺตโต สมนุปสฺสติ, ตสฺส เอวํ โหติ “อยํ โข เม อตฺตา, อิมสฺมิํ จ ปน อตฺตนิ อิทํ รูปนฺ”ติ \textbf{อตฺตนิ รูปํ สมนุปสฺสติ}.\csromanspacer{} เสยฺยถาปิ ปุปฺผํ คนฺธสมฺปนฺนํ อสฺส.\csromanspacer{} ตเมนํ ปุริโส เอวํ วเทยฺย “อิทํ ปุปฺผํ, อยํ คนฺโธ.\csromanspacer{} อญฺญํ ปุปฺผํ, อญฺโญ คนฺโธ.\csromanspacer{} โส โข ปนายํ คนฺโธ อิมสฺมิํ ปุปฺเผ”ติ ปุปฺผสฺมิํ คนฺธํ สมนุปสฺสติ.\csromanspacer{} เอวเมวํ อิเธกจฺโจ เวทนํ สญฺญํ สงฺขาเร วิญฺญาณํ อตฺตโต สมนุปสฺสติ.\csromanspacer{} ตสฺส เอวํ โหติ “อยํ โข เม อตฺตา, อิมสฺมิํ จ ปน อตฺตนิ อิทํ รูปนฺ”ติ \textbf{อตฺตนิ รูปํ สมนุปสฺสติ}.\csromanspacer{} อภินิเวสปรามาโส ทิฏฺฐิ.\csromanspacer{} ทิฏฺฐิ น วตฺถุ, วตฺถุ น ทิฏฺฐิ, อญฺญา ทิฏฺฐิ, อญฺญํ วตฺถุ.\csromanspacer{} ยา จ ทิฏฺฐิ ยญฺจ วตฺถุ, อยํ ตติยา รูปวตฺถุกา อตฺตานุทิฏฺฐิ.\csromanspacer{} อตฺตานุทิฏฺฐิ มิจฺฉาทิฏฺฐิ ทิฏฺฐิวิปตฺติ ฯเปฯ อิมานิ สญฺโญชนานิ น จ ทิฏฺฐิโย.\csromanspacer{} เอวํ \textbf{อตฺตนิ รูปํ สมนุปสฺสติ}. (3)}

\prose{กถํ \textbf{รูปสฺมิํ อตฺตานํ สมนุปสฺสติ}? อิเธกจฺโจ เวทนํ สญฺญํ สงฺขาเร วิญฺญาณํ อตฺตโต สมนุปสฺสติ, ตสฺส เอวํ โหติ “อยํ โข เม อตฺตา, โส โข ปน เม อยํ อตฺตา อิมสฺมิํ รูเป”ติ \textbf{รูปสฺมิํ อตฺตานํ สมนุปสฺสติ}.\csromanspacer{} เสยฺยถาปิ มณิ กรณฺฑเก ปกฺขิตฺโต อสฺส.\csromanspacer{} ตเมนํ ปุริโส เอวํ วเทยฺย “อยํ มณิ, อยํ กรณฺฑโก.\csromanspacer{} อญฺโญ มณิ, อญฺโญ กรณฺฑโก.\csromanspacer{} โส โข ปนายํ มณิ อิมสฺมิํ กรณฺฑเก”ติ กรณฺฑกสฺมิํ มณิํ สมนุปสฺสติ.\csromanspacer{} เอวเมวํ อิเธกจฺโจ เวทนํ สญฺญํ สงฺขาเร วิญฺญาณํ อตฺตโต สมนุปสฺสติ.\csromanspacer{} ตสฺส เอวํ โหติ “อยํ โข เม อตฺตา, โส โข ปน เม อยํ อตฺตา อิมสฺมิํ รูเป”ติ \textbf{รูปสฺมิํ อตฺตานํ สมนุปสฺสติ}, อภินิเวสปรามาโส ทิฏฺฐิ, ทิฏฺฐิ น วตฺถุ, วตฺถุ น ทิฏฺฐิ.\csromanspacer{} อญฺญา ทิฏฺฐิ, อญฺญํ วตฺถุ.\csromanspacer{} ยา จ ทิฏฺฐิ ยญฺจ วตฺถุ.\csromanspacer{} อยํ จตุตฺถา รูปวตฺถุกา อตฺตานุทิฏฺฐิ.\csromanspacer{} อตฺตานุทิฏฺฐิ มิจฺฉาทิฏฺฐิ ทิฏฺฐิวิปตฺติ ฯเปฯ อิมานิ สญฺโญชนานิ น จ ทิฏฺฐิโย.\csromanspacer{} เอวํ \textbf{รูปสฺมิํ อตฺตานํ สมนุปสฺสติ}. (4)}

\proseitem{132}{กถํ \textbf{เวทนํ อตฺตโต สมนุปสฺสติ?} อิเธกจฺโจ จกฺขุสมฺผสฺสชํ เวทนํ.\csromanspacer{} โสตสมฺผสฺสชํ เวทนํ.\csromanspacer{} ฆานสมฺผสฺสชํ เวทนํ.\csromanspacer{} ชิวฺหาสมฺผสฺสชํ เวทนํ.\csromanspacer{} กายสมฺผสฺสชํ เวทนํ.\csromanspacer{} มโนสมฺผสฺสชํ เวทนํ อตฺตโต สมนุปสฺสติ “ยา มโนสมฺผสฺสชา เวทนา, โส อหํ.\csromanspacer{} โย อหํ, สา มโนสมฺผสฺสชา เวทนา”ติ มโนสมฺผสฺสชํ เวทนญฺจ \csromanfolio{140}อตฺตญฺจ อทฺวยํ สมนุปสฺสติ.\csromanspacer{} เสยฺยถาปิ เตลปฺปทีปสฺส ฌายโต ยา อจฺจิ, โส วณฺโณ.\csromanspacer{} โย วณฺโณ, สา อจฺจีติ อจฺจิญฺจ วณฺณญฺจ อทฺวยํ สมนุปสฺสติ.\csromanspacer{} เอวเมวํ อิเธกจฺโจ มโนสมฺผสฺสชํ เวทนํ อตฺตโต สมนุปสฺสติ “ยา มโนสมฺผสฺสชา เวทนา, โส อหํ.\csromanspacer{} โย อหํ, สา มโนสมฺผสฺสชา เวทนา”ติ มโนสมฺผสฺสชํ เวทนญฺจ อตฺตญฺจ อทฺวยํ สมนุปสฺสติ.\csromanspacer{} อภินิเวสปรามาโส ทิฏฺฐิ, ทิฏฺฐิ น วตฺถุ, วตฺถุ น ทิฏฺฐิ.\csromanspacer{} อญฺญา ทิฏฺฐิ, อญฺญํ วตฺถุ.\csromanspacer{} ยา จ ทิฏฺฐิ ยญฺจ วตฺถุ.\csromanspacer{} อยํ ปฐมา เวทนาวตฺถุกา อตฺตานุทิฏฺฐิ.\csromanspacer{} อตฺตานุทิฏฺฐิ มิจฺฉาทิฏฺฐิ ทิฏฺฐิวิปตฺติ ฯเปฯ อิมานิ สญฺโญชนานิ น จ ทิฏฺฐิโย.\csromanspacer{} เอวํ เวทนํ อตฺตโต สมนุปสฺสติ. (1, 5)}

\prose{กถํ \textbf{เวทนาวนฺตํ อตฺตานํ สมนุปสฺสติ}? อิเธกจฺโจ สญฺญํ สงฺขาเร วิญฺญาณํ รูปํ อตฺตโต สมนุปสฺสติ.\csromanspacer{} ตสฺส เอวํ โหติ “อยํ โข เม อตฺตา, โส โข ปน เม อยํ อตฺตา, อิมาย เวทนาย เวทนาวา”ติ \textbf{เวทนาวนฺตํ อตฺตานํ สมนุปสฺสติ}.\csromanspacer{} เสยฺยถาปิ รุกฺโข ฉายาสมฺปนฺโน อสฺส, ตเมนํ ปุริโส เอวํ วเทยฺย “อยํ รุกฺโข, อยํ ฉายา.\csromanspacer{} อญฺโญ รุกฺโข, อญฺญา ฉายา.\csromanspacer{} โส โข ปนายํ รุกฺโข อิมาย ฉายาย ฉายาวา”ติ ฉายาวนฺตํ รุกฺขํ สมนุปสฺสติ.\csromanspacer{} เอวเมวํ อิเธกจฺโจ สญฺญํ สงฺขาเร วิญฺญาณํ รูปํ อตฺตโต สมนุปสฺสติ, ตสฺส เอวํ โหติ “อยํ โข เม อตฺตา, โส โข ปน เม อยํ อตฺตา อิมาย เวทนาย เวทนาวา”ติ \textbf{เวทนาวนฺตํ อตฺตานํ สมนุปสฺสติ}.\csromanspacer{} อภินิเวสปรามาโส ทิฏฺฐิ, ทิฏฺฐิ น วตฺถุ, วตฺถุ น ทิฏฺฐิ.\csromanspacer{} อญฺญา ทิฏฺฐิ, อญฺญํ วตฺถุ, ยา จ ทิฏฺฐิ ยญฺจ วตฺถุ.\csromanspacer{} อยํ ทุติยา เวทนาวตฺถุกา อตฺตานุทิฏฺฐิ.\csromanspacer{} อตฺตานุทิฏฺฐิ มิจฺฉาทิฏฺฐิ ทิฏฺฐิวิปตฺติ ฯเปฯ อิมานิ สญฺโญชนานิ น จ ทิฏฺฐิโย.\csromanspacer{} เอวํ \textbf{เวทนาวนฺตํ อตฺตานํ สมนุปสฺสติ}. (2, 6)}

\prose{กถํ \textbf{อตฺตนิ เวทนํ สมนุปสฺสติ}? อิเธกจฺโจ สญฺญํ สงฺขาเร วิญฺญาณํ รูปํ อตฺตโต สมนุปสฺสติ, ตสฺส เอวํ โหติ “อยํ โข เม อตฺตา, อิมสฺมิญฺจ ปน อตฺตนิ อยํ เวทนา”ติ \textbf{อตฺตนิ เวทนํ สมนุปสฺสติ}.\csromanspacer{} เสยฺยถาปิ ปุปฺผํ คนฺธสมฺปนฺนํ อสฺส, ตเมนํ ปุริโส เอวํ วเทยฺย “อิทํ ปุปฺผํ, อยํ คนฺโธ.\csromanspacer{} อญฺญํ ปุปฺผํ, อญฺโญ คนฺโธ.\csromanspacer{} โส โข ปนายํ คนฺโธ อิมสฺมิํ ปุปฺเผ”ติ ปุปฺผสฺมิํ คนฺธํ สมนุปสฺสติ.\csromanspacer{} เอวเมวํ อิเธกจฺโจ สญฺญํ สงฺขาเร วิญฺญาณํ รูปํ อตฺตโต สมนุปสฺสติ, ตสฺส เอวํ โหติ “อยํ โข เม อตฺตา, อิมสฺมิํ จ ปน อตฺตนิ อยํ เวทนา”ติ อตฺตนิ เวทนํ}

\vspace{\tipitakaparskip}
\csromancenter{ทิฏฺฐิกถา สมนุปสฺสติ, อภินิเวสปรามาโส ทิฏฺฐิ, ทิฏฺฐิ น วตฺถุ, วตฺถุ น ทิฏฺฐิ.\csromanspacer{} อญฺญา ทิฏฺฐิ, อญฺญํ วตฺถุ.\csromanspacer{} ยา จ ทิฏฺฐิ ยญฺจ วตฺถุ.\csromanspacer{} อยํ ตติยา เวทนาวตฺถุกา อตฺตานุทิฏฺฐิ.\csromanspacer{} อตฺตานุทิฏฺฐิ มิจฺฉาทิฏฺฐิ ทิฏฺฐิวิปตฺติ ฯเปฯ อิมานิ สญฺโญชนานิ น จ ทิฏฺฐิโย.\csromanspacer{} เอวํ อตฺตนิ เวทนํ สมนุปสฺสติ. (3, 7)}
\prose{\csromanfolio{141}กถํ \textbf{เวทนาย อตฺตานํ สมนุปสฺสติ}? อิเธกจฺโจ สญฺญํ สงฺขาเร วิญฺญาณํ รูปํ อตฺตโต สมนุปสฺสติ, ตสฺส เอวํ โหติ “อยํ โข เม อตฺตา, โส โข ปน เม อยํ อตฺตา อิมาย เวทนายา”ติ \textbf{เวทนาย อตฺตานํ สมนุปสฺสติ}.\csromanspacer{} เสยฺยถาปิ มณิ กรณฺฑเก ปกฺขิตฺโต อสฺส.\csromanspacer{} ตเมนํ ปุริโส เอวํ วเทยฺย “อยํ มณิ, อยํ กรณฺฑโก.\csromanspacer{} อญฺโญ มณิ, อญฺโญ กรณฺฑโก.\csromanspacer{} โส โข ปนายํ มณิ อิมสฺมิํ กรณฺฑเก”ติ กรณฺฑกสฺมิํ มณิํ สมนุปสฺสติ.\csromanspacer{} เอวเมวํ อิเธกจฺโจ สญฺญํ สงฺขาเร วิญฺญาณํ รูปํ อตฺตโต สมนุปสฺสติ, ตสฺส เอวํ โหติ “อยํ โข เม อตฺตา, โส โข ปน เม อยํ อตฺตา อิมาย เวทนายา”ติ \textbf{เวทนาย อตฺตานํ สมนุปสฺสติ}.\csromanspacer{} อภินิเวสปรามาโส ทิฏฺฐิ, ทิฏฺฐิ น วตฺถุ, วตฺถุ น ทิฏฺฐิ.\csromanspacer{} อญฺญา ทิฏฺฐิ, อญฺญํ วตฺถุ.\csromanspacer{} ยา จ ทิฏฺฐิ ยญฺจ วตฺถุ.\csromanspacer{} อยํ จตุตฺถา เวทนาวตฺถุกา อตฺตานุทิฏฺฐิ.\csromanspacer{} อตฺตานุทิฏฺฐิ มิจฺฉาทิฏฺฐิ ทิฏฺฐิวิปตฺติ ฯเปฯ อิมานิ สญฺโญชนานิ น จ ทิฏฺฐิโย.\csromanspacer{} เอวํ \textbf{เวทนาย อตฺตานํ สมนุปสฺสติ}. (4, 8)}

\proseitem{133}{กถํ \textbf{สญฺญํ อตฺตโต สมนุปสฺสติ}? อิเธกจฺโจ จกฺขุสมฺผสฺสชํ สญฺญํ.\csromanspacer{} โสตสมฺผสฺสชํ สญฺญํ.\csromanspacer{} ฆานสมฺผสฺสชํ สญฺญํ.\csromanspacer{} ชิวฺหาสมฺผสฺสชํ สญฺญํ.\csromanspacer{} กายสมฺผสฺสชํ สญฺญํ.\csromanspacer{} มโนสมฺผสฺสชํ \textbf{สญฺญํ อตฺตโต สมนุปสฺสติ} “ยา มโนสมฺผสฺสชา สญฺญา, โส อหํ.\csromanspacer{} โย อหํ, สา มโนสมฺผสฺสชา สญฺญา”ติ มโนสมฺผสฺสชํ สญฺญญฺจ อตฺตญฺจ อทฺวยํ สมนุปสฺสติ.\csromanspacer{} เสยฺยถาปิ เตลปฺปทีปสฺส ฌายโต ยา อจฺจิ, โส วณฺโณ.\csromanspacer{} โย วณฺโณ, สา อจฺจีติ อจฺจิญฺจ วณฺณญฺจ อทฺวยํ สมนุปสฺสติ.\csromanspacer{} เอวเมวํ อิเธกจฺโจ มโนสมฺผสฺสชํ \textbf{สญฺญํ อตฺตโต สมนุปสฺสติ} “ยา มโนสมฺผสฺสชา สญฺญา, โส อหํ.\csromanspacer{} โย อหํ, สา มโนสมฺผสฺสชา สญฺญา”ติ มโนสมฺผสฺสชํ สญฺญญฺจ อตฺตญฺจ อทฺวยํ สมนุปสฺสติ, อภินิเวสปรามาโส ทิฏฺฐิ, ทิฏฺฐิ น วตฺถุ, วตฺถุ น ทิฏฺฐิ.\csromanspacer{} อญฺญา ทิฏฺฐิ, อญฺญํ วตฺถุ.\csromanspacer{} ยา จ ทิฏฺฐิ ยญฺจ วตฺถุ.\csromanspacer{} อยํ ปฐมา สญฺญาวตฺถุกา อตฺตานุทิฏฺฐิ. \csromanfolio{142}อตฺตานุทิฏฺฐิ มิจฺฉาทิฏฺฐิ.\csromanspacer{} ทิฏฺฐิวิปตฺติ ฯเปฯ อิมานิ สญฺโญชนานิ น จ ทิฏฺฐิโย.\csromanspacer{} เอวํ สญฺญํ อตฺตโต สมนุปสฺสติ. (1, 9)}

\prose{กถํ \textbf{สญฺญาวนฺตํ อตฺตานํ สมนุปสฺสติ}? อิเธกจฺโจ สงฺขาเร วิญฺญาณํ รูปํ เวทนํ อตฺตโต สมนุปสฺสติ, ตสฺส เอวํ โหติ “อยํ โข เม อตฺตา, โส โข ปน เม อยํ อตฺตา อิมาย สญฺญาย สญฺญาวา”ติ \textbf{สญฺญาวนฺตํ อตฺตานํ สมนุปสฺสติ}.\csromanspacer{} เสยฺยถาปิ รุกฺโข ฉายาสมฺปนฺโน อสฺส, ตเมนํ ปุริโส เอวํ วเทยฺย “อยํ รุกฺโข, อยํ ฉายา.\csromanspacer{} อญฺโญ รุกฺโข, อญฺญา ฉายา.\csromanspacer{} โส โข ปนายํ รุกฺโข อิมาย ฉายาย ฉายาวา”ติ ฉายาวนฺตํ รุกฺขํ สมนุปสฺสติ.\csromanspacer{} เอวเมวํ อิเธกจฺโจ สงฺขาเร วิญฺญาณํ รูปํ เวทนํ อตฺตโต สมนุปสฺสติ, ตสฺส เอวํ โหติ “อยํ โข เม อตฺตา, โส โข ปน เม อยํ อตฺตา อิมาย สญฺญาย สญฺญาวา”ติ \textbf{สญฺญาวนฺตํ อตฺตานํ สมนุปสฺสติ}, อภินิเวสปรามาโส ทิฏฺฐิ, ทิฏฺฐิ น วตฺถุ, วตฺถุ น ทิฏฺฐิ.\csromanspacer{} อญฺญา ทิฏฺฐิ, อญฺญํ วตฺถุ.\csromanspacer{} ยา จ ทิฏฺฐิ ยญฺจ วตฺถุ.\csromanspacer{} อยํ ทุติยา สญฺญาวตฺถุกา อตฺตานุทิฏฺฐิ.\csromanspacer{} อตฺตานุทิฏฺฐิ มิจฺฉาทิฏฺฐิ ฯเปฯ อิมานิ สญฺโญชนานิ น จ ทิฏฺฐิโย.\csromanspacer{} เอวํ \textbf{สญฺญาวนฺตํ อตฺตานํ สมนุปสฺสติ}. (2, 10)}

\prose{กถํ \textbf{อตฺตนิ สญฺญํ สมนุปสฺสติ}? อิเธกจฺโจ สงฺขาเร วิญฺญาณํ รูปํ เวทนํ อตฺตโต สมนุปสฺสติ, ตสฺส เอวํ โหติ “อยํ โข เม อตฺตา, อิมสฺมิํ จ ปน อตฺตนิ อยํ สญฺญา”ติ \textbf{อตฺตนิ สญฺญํ สมนุปสฺสติ}.\csromanspacer{} เสยฺยถาปิ ปุปฺผํ คนฺธสมฺปนฺนํ อสฺส, ตเมนํ ปุริโส เอวํ วเทยฺย “อิทํ ปุปฺผํ อยํ คนฺโธ.\csromanspacer{} อญฺญํ ปุปฺผํ, อญฺโญ คนฺโธ.\csromanspacer{} โส โข ปนายํ คนฺโธ อิมสฺมิํ ปุปฺเผ”ติ ปุปฺผสฺมิํ คนฺธํ สมนุปสฺสติ.\csromanspacer{} เอวเมวํ อิเธกจฺโจ สงฺขาเร วิญฺญาณํ รูปํ เวทนํ อตฺตโต สมนุปสฺสติ, ตสฺส เอวํ โหติ “อยํ โข เม อตฺตา, อิมสฺมิํ จ ปน อตฺตนิ อยํ สญฺญา”ติ \textbf{อตฺตนิ สญฺญํ สมนุปสฺสติ}, อภินิเวสปรามาโส ทิฏฺฐิ, ทิฏฺฐิ น วตฺถุ, วตฺถุ น ทิฏฺฐิ.\csromanspacer{} อญฺญา ทิฏฺฐิ, อญฺญํ วตฺถุ.\csromanspacer{} ยา จ ทิฏฺฐิ ยญฺจ วตฺถุ.\csromanspacer{} อยํ ตติยา สญฺญาวตฺถุกา อตฺตานุทิฏฺฐิ.\csromanspacer{} อตฺตานุทิฏฺฐิ มิจฺฉาทิฏฺฐิ ฯเปฯ อิมานิ สญฺโญชนานิ น จ ทิฏฺฐิโย.\csromanspacer{} เอวํ \textbf{อตฺตนิ สญฺญํ สมนุปสฺสติ}. (3, 11)}

\prose{กถํ \textbf{สญฺญาย อตฺตานํ สมนุปสฺสติ?} อิเธกจฺโจ สงฺขาเร วิญฺญาณํ รูปํ เวทนํ อตฺตโต สมนุปสฺสติ, ตสฺส เอวํ โหติ “อยํ โข เม อตฺตา, โส โข ปน เม อยํ อตฺตา อิมาย สญฺญายา”ติ}

\vspace{\tipitakaparskip}
\csromancenter{ทิฏฺฐิกถา สญฺญาย อตฺตานํ สมนุปสฺสติ.\csromanspacer{} เสยฺยถาปิ มณิ กรณฺฑเก ปกฺขิตฺโต อสฺส, ตเมนํ ปุริโส เอวํ วเทยฺย “อยํ มณิ, อยํ กรณฺฑโก.\csromanspacer{} อญฺโญ มณิ, อญฺโญ กรณฺฑโก.\csromanspacer{} โส โข ปนายํ มณิ อิมสฺมิํ กรณฺฑเก”ติ กรณฺฑกสฺมิํ, มณิํ สมนุปสฺสติ.\csromanspacer{} เอวเมวํ อิเธกจฺโจ สงฺขาเร วิญฺญาณํ รูปํ เวทนํ อตฺตโต สมนุปสฺสติ, ตสฺส เอวํ โหติ “อยํ โข เม อตฺตา, โส โข ปน เม อยํ อตฺตา อิมาย สญฺญายา”ติ สญฺญาย อตฺตานํ สมนุปสฺสติ, อภินิเวสปรามาโส ทิฏฺฐิ, ทิฏฺฐิ น วตฺถุ, วตฺถุ น ทิฏฺฐิ.\csromanspacer{} อญฺญา ทิฏฺฐิ, อญฺญํ วตฺถุ.\csromanspacer{} ยา จ ทิฏฺฐิ ยญฺจ วตฺถุ.\csromanspacer{} อยํ จตุตฺถา สญฺญาวตฺถุกา อตฺตานุทิฏฺฐิ.\csromanspacer{} อตฺตานุทิฏฺฐิ มิจฺฉาทิฏฺฐิ ฯเปฯ อิมานิ สญฺโญชนานิ น จ ทิฏฺฐิโย.\csromanspacer{} เอวํ สญฺญาย อตฺตานํ สมนุปสฺสติ. (4, 12)}
\proseitem{134}{\csromanfolio{143}กถํ \textbf{สงฺขาเร อตฺตโต สมนุปสฺสติ}? อิเธกจฺโจ จกฺขุสมฺผสฺสชํ เจตนํ.\csromanspacer{} โสตสมฺผสฺสชํ เจตนํ.\csromanspacer{} ฆานสมฺผสฺสชํ เจตนํ.\csromanspacer{} ชิวฺหาสมฺผสฺสชํ เจตนํ.\csromanspacer{} กายสมฺผสฺสชํ เจตนํ.\csromanspacer{} มโนสมฺผสฺสชํ เจตนํ อตฺตโต สมนุปสฺสติ “ยา มโนสมฺผสฺสชา เจตนา, โส อหํ.\csromanspacer{} โย อหํ, สา มโนสมฺผสฺสชา เจตนา”ติ มโนสมฺผสฺสชํ เจตนญฺจ อตฺตญฺจ อทฺวยํ สมนุปสฺสติ.\csromanspacer{} เสยฺยถาปิ เตลปฺปทีปสฺส ฌายโต ยา อจฺจิ, โส วณฺโณ.\csromanspacer{} โย วณฺโณ, สา อจฺจีติ อจฺจิญฺจ วณฺณญฺจ อทฺวยํ สมนุปสฺสติ.\csromanspacer{} เอวเมวํ อิเธกจฺโจ มโนสมฺผสฺสชํ เจตนํ อตฺตโต สมนุปสฺสติ “ยา มโนสมฺผสฺสชา เจตนา, โส อหํ.\csromanspacer{} โย อหํ, สา มโนสมฺผสฺสชา เจตนา”ติ มโนสมฺผสฺสชํ เจตนญฺจ อตฺตญฺจ อทฺวยํ สมนุปสฺสติ, อภินิเวสปรามาโส ทิฏฺฐิ, ทิฏฺฐิ น วตฺถุ, วตฺถุ น ทิฏฺฐิ.\csromanspacer{} อญฺญา ทิฏฺฐิ, อญฺญํ วตฺถุ.\csromanspacer{} ยา จ ทิฏฺฐิ ยญฺจ วตฺถุ.\csromanspacer{} อยํ ปฐมา สงฺขารวตฺถุกา อตฺตานุทิฏฺฐิ.\csromanspacer{} อตฺตานุทิฏฺฐิ มิจฺฉาทิฏฺฐิ ฯเปฯ อิมานิ สญฺโญชนานิ น จ ทิฏฺฐิโย.\csromanspacer{} เอวํ \textbf{สงฺขาเร อตฺตโต สมนุปสฺสติ}. (1, 13)}

\prose{กถํ \textbf{สงฺขารวนฺตํ อตฺตานํ สมนุปสฺสติ}? อิเธกจฺโจ วิญฺญาณํ รูปํ เวทนํ สญฺญํ อตฺตโต สมนุปสฺสติ, ตสฺส เอวํ โหติ “อยํ โข เม อตฺตา, โส โข ปน เม อยํ อตฺตา อิเมหิ สงฺขาเรหิ สงฺขารวา”ติ \textbf{สงฺขารวนฺตํ อตฺตานํ สมนุปสฺสติ}.\csromanspacer{} เสยฺยถาปิ รุกฺโข ฉายาสมฺปนฺโน อสฺส, ตเมนํ ปุริโส เอวํ วเทยฺย “อยํ รุกฺโข, อยํ ฉายา.\csromanspacer{} อญฺโญ รุกฺโข, อญฺญา ฉายา.\csromanspacer{} โส โข ปนายํ \csromanfolio{144}รุกฺโข อิมาย ฉายาย ฉายาวา”ติ ฉายาวนฺตํ รุกฺขํ สมนุปสฺสติ.\csromanspacer{} เอวเมวํ อิเธกจฺโจ วิญฺญาณํ รูปํ เวทนํ สญฺญํ อตฺตโต สมนุปสฺสติ, ตสฺส เอวํ โหติ “อยํ โข เม อตฺตา, โส โข ปน เม อยํ อตฺตา อิเมหิ สงฺขาเรหิ สงฺขารวา”ติ สงฺขารวนฺตํ อตฺตานํ สมนุปสฺสติ, อภินิเวสปรามาโส ทิฏฺฐิ, ทิฏฺฐิ น วตฺถุ, วตฺถุ น ทิฏฺฐิ.\csromanspacer{} อญฺญา ทิฏฺฐิ, อญฺญํ วตฺถุ.\csromanspacer{} ยา จ ทิฏฺฐิ ยญฺจ วตฺถุ.\csromanspacer{} อยํ ทุติยา สงฺขารวตฺถุกา อตฺตานุทิฏฺฐิ.\csromanspacer{} อตฺตานุทิฏฺฐิ มิจฺฉาทิฏฺฐิ ฯเปฯ อิมานิ สญฺโญชนานิ น จ ทิฏฺฐิโย.\csromanspacer{} เอวํ สงฺขารวนฺตํ อตฺตานํ สมนุปสฺสติ. (1, 14)}

\prose{กถํ \textbf{อตฺตนิ สงฺขาเร สมนุปสฺสติ}? อิเธกจฺโจ วิญฺญาณํ รูปํ เวทนํ สญฺญํ อตฺตโต สมนุปสฺสติ, ตสฺส เอวํ โหติ “อยํ โข เม อตฺตา, อิมสฺมิํ จ ปน อตฺตนิ อิเม สงฺขารา”ติ \textbf{อตฺตนิ สงฺขาเร สมนุปสฺสติ}.\csromanspacer{} เสยฺยถาปิ, ปุปฺผํ คนฺธสมฺปนฺนํ อสฺส, ตเมนํ ปุริโส เอวํ วเทยฺย “อิทํ ปุปฺผํ, อยํ คนฺโธ.\csromanspacer{} อญฺญํ ปุปฺผํ, อญฺโญ คนฺโธ.\csromanspacer{} โส โข ปนายํ คนฺโธ อิมสฺมิํ ปุปฺเผ”ติ ปุปฺผสฺมิํ คนฺธํ สมนุปสฺสติ.\csromanspacer{} เอวเมวํ อิเธกจฺโจ วิญฺญาณํ รูปํ เวทนํ สญฺญํ อตฺตโต สมนุปสฺสติ, ตสฺส เอวํ โหติ “อยํ โข เม อตฺตา, อิมสฺมิํ จ ปน อตฺตนิ อิเม สงฺขารา”ติ \textbf{อตฺตนิ สงฺขาเร สมนุปสฺสติ}, อภินิเวสปรามาโส ทิฏฺฐิ, ทิฏฺฐิ น วตฺถุ.\csromanspacer{} วตฺถุ น ทิฏฺฐิ.\csromanspacer{} อญฺญา ทิฏฺฐิ, อญฺญํ วตฺถุ.\csromanspacer{} ยา จ ทิฏฺฐิ ยญฺจ วตฺถุ.\csromanspacer{} อยํ ตติยา สงฺขารวตฺถุกา อตฺตานุทิฏฺฐิ.\csromanspacer{} อตฺตานุทิฏฺฐิ มิจฺฉาทิฏฺฐิ ฯเปฯ อิมานิ สญฺโญชนานิ น จ ทิฏฺฐิโย.\csromanspacer{} เอวํ \textbf{อตฺตนิ สงฺขาเร สมนุปสฺสติ}. (3, 15)}

\prose{กถํ \textbf{สขาเรสุ อตฺตานํ สมนุปสฺสติ}? อิเธกจฺโจ วิญฺญาณํ รูปํ เวทนํ สญฺญํ อตฺตโต สมนุปสฺสติ, ตสฺส เอวํ โหติ “อยํ โข เม อตฺตา, โส โข ปน เม อยํ อตฺตา อิเมสุ สงฺขาเรสู”ติ สงฺขาเรสุ อตฺตานํ สมนุปสฺสติ.\csromanspacer{} เสยฺยถาปิ มณิ กรณฺฑเก ปกฺขิตฺโต อสฺส, ตเมนํ ปุริโส เอวํ วเทยฺย “อยํ มณิ, อยํ กรณฺฑโก.\csromanspacer{} อญฺโญ มณิ, อญฺโญ กรณฺฑโก.\csromanspacer{} โส โข ปนายํ มณิ อิมสฺมิํ กรณฺฑเก”ติ กรณฺฑกสฺมิํ มณิํ สมนุปสฺสติ.\csromanspacer{} เอวเมวํ อิเธกจฺโจ วิญฺญาณํ รูปํ เวทนํ สญฺญํ อตฺตโต สมนุปสฺสติ, ตสฺส เอวํ โหติ “อยํ โข เม อตฺตา, โส โข ปน เม อยํ อตฺตา อิเมสุ สงฺขาเรสูติ สงฺขาเรสุ อตฺตานํ สมนุปสฺสติ, อภินิเวสปรามาโส ทิฏฺฐิ, ทิฏฺฐิ น วตฺถุ, วตฺถุ น ทิฏฺฐิ.\csromanspacer{} อญฺญา ทิฏฺฐิ, อญฺญํ วตฺถุ.\csromanspacer{} ยา จ ทิฏฺฐิ ยญฺจ}

\vspace{\tipitakaparskip}
\csromancenter{ทิฏฺฐิกถา วตฺถุ อยํ จตุตฺถา สงฺขารวตฺถุกา อตฺตานุทิฏฺฐิ.\csromanspacer{} อตฺตานุทิฏฺฐิ มิจฺฉาทิฏฺฐิ ฯเปฯ อิมานิ สญฺโญชนานิ น จ ทิฏฺฐิโย.\csromanspacer{} เอวํ สงฺขาเรสุ อตฺตานํ สมนุปสฺสติ. (4, 16)}
\proseitem{135}{\csromanfolio{145}กถํ \textbf{วิญฺญาณํ อตฺตโต สมนุปสฺสติ}? อิเธกจฺโจ จกฺขุวิญฺญาณํ.\csromanspacer{} โสตวิญฺญาณํ.\csromanspacer{} ฆานวิญฺญาณํ.\csromanspacer{} ชิวฺหาวิญฺญาณํ.\csromanspacer{} กายวิญฺญาณํ.\csromanspacer{} มโน\textbf{วิญฺญาณํ อตฺตโต สมนุปสฺสติ} “ยํ มโนวิญฺญาณํ, โส อหํ.\csromanspacer{} โย อหํ, ตํ มโนวิญฺญาณนฺ”ติ มโนวิญฺญาณญฺจ, อตฺตญฺจ อทฺวยํ สมนุปสฺสติ.\csromanspacer{} เสยฺยถาปิ เตลปฺปทีปสฺส ฌายโต ยา อจฺจิ, โส วณฺโณ.\csromanspacer{} โย วณฺโณ, สา อจฺจีติ อจฺจิญฺจ วณฺณญฺจ อทฺวยํ สมนุปสฺสติ.\csromanspacer{} เอวเมวํ อิเธกจฺโจ มโน\textbf{วิญฺญาณํ อตฺตโต สมนุปสฺสติ} “ยํ มโนวิญฺญาณํ, โส อหํ.\csromanspacer{} โย อหํ, ตํ มโนวิญฺญาณนฺ”ติ มโนวิญฺญาณญฺจ อตฺตญฺจ อทฺวยํ สมนุปสฺสติ, อภินิเวสปรามาโส ทิฏฺฐิ, ทิฏฺฐิ น วตฺถุ, วตฺถุ น ทิฏฺฐิ, อญฺญา ทิฏฺฐิ, อญฺญํ วตฺถุ, ยา จ ทิฏฺฐิ, ยญฺจ วตฺถุ.\csromanspacer{} อยํ ปฐมา วิญฺญาณวตฺถุกา อตฺตานุทิฏฺฐิ.\csromanspacer{} อตฺตานุทิฏฺฐิ มิจฺฉาทิฏฺฐิ ฯเปฯ อิมานิ สญฺโญชนานิ น จ ทิฏฺฐิโย.\csromanspacer{} เอวํ \textbf{วิญฺญาณํ อตฺตโต สมนุปสฺสติ}. (1, 17)}

\prose{กถํ \textbf{วิญฺญาณวนฺตํ อตฺตานํ สมนุปสฺสติ}? อิเธกจฺโจ รูปํ เวทนํ สญฺญํ สงฺขาเร อตฺตโต สมนุปสฺสติ, ตสฺส เอวํ โหติ “อยํ โข เม อตฺตา, โส โข ปน เม อยํ อตฺตา อิมินา วิญฺญาเณน วิญฺญาณวา”ติ \textbf{วิญฺญาณวนฺตํ อตฺตานํ สมนุปสฺสติ}.\csromanspacer{} เสยฺยถาปิ รุกฺโข ฉายาสมฺปนฺโน อสฺส, ตเมนํ ปุริโส เอวํ วเทยฺย “อยํ รุกฺโข, อยํ ฉายา.\csromanspacer{} อญฺโญ รุกฺโข, อญฺญา ฉายา.\csromanspacer{} โส โข ปนายํ รุกฺโข อิมาย ฉายาย ฉายาวา”ติ ฉายาวนฺตํ รุกฺขํ สมนุปสฺสติ.\csromanspacer{} เอวเมวํ อิเธกจฺโจ รูปํ เวทนํ สญฺญํ สงฺขาเร อตฺตโต สมนุปสฺสติ, ตสฺส เอวํ โหติ “อยํ โข เม อตฺตา, โส โข ปน เม อยํ อตฺตา อิมินา วิญฺญาเณน วิญฺญาณวา”ติ \textbf{วิญฺญาณวนฺตํ อตฺตานํ สมนุปสฺสติ}, อภินิเวสปรามาโส ทิฏฺฐิ, ทิฏฺฐิ น วตฺถุ, วตฺถุ น ทิฏฺฐิ, อญฺญา ทิฏฺฐิ, อญฺญํ วตฺถุ, ยา จ ทิฏฺฐิ ยญฺจ วตฺถุ.\csromanspacer{} อยํ ทุติยา วิญฺญาณวตฺถุกา อตฺตานุทิฏฺฐิ, อตฺตานุทิฏฺฐิ มิจฺฉาทิฏฺฐิ ฯเปฯ อิมานิ สญฺโญชนานิ น จ ทิฏฺฐิโย.\csromanspacer{} เอวํ \textbf{วิญฺญาณวนฺตํ อตฺตานํ สมนุปสฺสติ}. (2, 18)}

\prose{กถํ \textbf{อตฺตนิ วิญฺญาณํ สมนุปสฺสติ}? อิเธกจฺโจ รูปํ เวทนํ สญฺญํ สงฺขาเร อตฺตโต สมนุปสฺสติ, ตสฺส เอว โหติ “อยํ โข เม \csromanfolio{146}อตฺตา, อิมสฺมิญฺจ ปน อตฺตนิ อิทํ วิญฺญาณนฺ”ติ อตฺตนิ วิญฺญาณํ สมนุปสฺสติ.\csromanspacer{} เสยฺยถาปิ ปุปฺผํ คนฺธสมฺปนฺนํ อสฺส, ตเมนํ ปุริโส เอวํ วเทยฺย “อิทํ ปุปฺผํ, อยํ คนฺโธ.\csromanspacer{} อญฺญํ ปุปฺผํ, อญฺโญ คนฺโธ.\csromanspacer{} โส โข ปนายํ คนฺโธ อิมสฺมิํ ปุปฺเผ”ติ ปุปฺผสฺมิํ คนฺธํ สมนุปสฺสติ.\csromanspacer{} เอวเมวํ อิเธกจฺโจ รูปํ เวทนํ สญฺญํ สงฺขาเร อตฺตโต สมนุปสฺสติ, ตสฺส เอวํ โหติ “อยํ โข เม อตฺตา, อิมสฺมิญฺจ ปน อตฺตนิ อิทํ วิญฺญาณนฺ”ติ อตฺตนิ วิญฺญาณํ สมนุปสฺสติ.\csromanspacer{} อภินิเวสปรามาโส ทิฏฺฐิ, ทิฏฺฐิ น วตฺถุ, วตฺถุ น ทิฏฺฐิ, อญฺญา ทิฏฺฐิ, อญฺญํ วตฺถุ, ยา จ ทิฏฺฐิ ยญฺจ วตฺถุ.\csromanspacer{} อยํ ตติยา วิญฺญาณวตฺถุกา อตฺตานุทิฏฺฐิ.\csromanspacer{} อตฺตานุทิฏฺฐิ มิจฺฉาทิฏฺฐิ ฯเปฯ อิมานิ สญฺโญชนานิ น จ ทิฏฺฐิโย.\csromanspacer{} เอวํ อตฺตนิ วิญฺญาณํ สมนุปสฺสติ. (3, 19)}

\prose{กถํ \textbf{วิญฺญาณสฺมิํ อตฺตานํ สมนุปสฺสติ?} อิเธกจฺโจ รูปํ เวทนํ สญฺญํ สงฺขาเร อตฺตโต สมนุปสฺสติ, ตสฺส เอวํ โหติ “อยํ โข เม อตฺตา, โส โข ปน เม อยํ อตฺตา อิมสฺมิํ วิญฺญาเณ”ติ วิญฺญาณสฺมิํ อตฺตานํ สมนุปสฺสติ.\csromanspacer{} เสยฺยถาปิ มณิ กรณฺฑเก ปกฺขิตฺโต อสฺส, ตเมนํ ปุริโส เอวํ วเทยฺย “อยํ มณิ, อยํ กรณฺฑโก.\csromanspacer{} อญฺโญ มณิ, อญฺโญ กรณฺฑโก.\csromanspacer{} โส โข ปนายํ มณิ อิมสฺมิํ กรณฺฑเก”ติ กรณฺฑกสฺมิํ มณิํ สมนุปสฺสติ.\csromanspacer{} เอวเมวํ อิเธกจฺโจ รูปํ เวทนํ สญฺญํ สงฺขาเร อตฺตโต สมนุปสฺสติ, ตสฺส เอวํ โหติ “อยํ โข เม อตฺตา, โส โข ปน เม อยํ อตฺตา อิมสฺมิํ วิญฺญาเณ”ติ วิญฺญาณสฺมิํ อตฺตานํ สมนุปสฺสติ, อภินิเวสปรามาโส ทิฏฺฐิ, ทิฏฺฐิ น วตฺถุ, วตฺถุ น ทิฏฺฐิ, อญฺญา ทิฏฺฐิ, อญฺญํ วตฺถุ, ยา จ ทิฏฺฐิ ยญฺจ วตฺถุ.\csromanspacer{} อยํ จตุตฺถา วิญฺญาณวตฺถุกา อตฺตานุทิฏฺฐิ.\csromanspacer{} อตฺตานุทิฏฺฐิ มิจฺฉาทิฏฺฐิ ฯเปฯ อิมานิ สญฺโญชนานิ น จ ทิฏฺฐิโย.\csromanspacer{} เอวํ วิญฺญาณสฺมิํ อตฺตานํ สมนุปสฺสติ.\csromanspacer{} อตฺตานุทิฏฺฐิยา อิเมหิ วีสติยา อากาเรหิ อภินิเวโส โหติ. (4, 20)}

\csromanheader{2}{อตฺตานุทิฏฺฐินิทฺเทโส ทุติโย.}
\csromansectionrule
\csromanmatikaanchor{mtk.54}
\csromantocmarkref{subsection}{3. มิจฺฉาทิฏฺฐินิทฺเทส}{mtk.54}
\bookmark[dest={mtk.54},level=2]{3. มิจฺฉาทิฏฺฐินิทฺเทส}
\csromanheader{2}{3. มิจฺฉาทิฏฺฐินิทฺเทส}
\proseitem{136}{มิจฺฉาทิฏฺฐิยา กตเมหิ ทสหากาเรหิ อภินิเวโส โหติ? “นตฺถิ ทินฺนนฺ”ติ วตฺถุ\footnote{วตฺถุํ (สฺยา) เอวมุปริปิ.} เอวํวาโท มิจฺฉาภินิเวสปรามาโส\footnote{มิจฺฉาทิฏฺฐาภินิเวสปรามาโส (สฺยา)} ทิฏฺฐิ,}

\vspace{\tipitakaparskip}
\csromancenter{ทิฏฺฐิกถา ทิฏฺฐิ น วตฺถุ, วตฺถุ น ทิฏฺฐิ, อญฺญา ทิฏฺฐิ, อญฺญํ วตฺถุ, ยา จ ทิฏฺฐิ ยญฺจ วตฺถุ.\csromanspacer{} อยํ ปฐมา มิจฺฉาวตฺถุกา มิจฺฉาทิฏฺฐิ.\csromanspacer{} มิจฺฉาทิฏฺฐิ ทิฏฺฐิวิปตฺติ ฯเปฯ อิมานิ สญฺโญชนานิ น จ ทิฏฺฐิโย. “นตฺถิ ยิฏฺฐนฺ”ติ วตฺถุ ฯเปฯ “นตฺถิ หุตนฺ”ติ วตฺถุ. “นตฺถิ สุกตทุกฺกฏานํ กมฺมานํ ผลํ วิปาโก”ติ วตฺถุ. “นตฺถิ อยํ โลโก”ติ วตฺถุ. “นตฺถิ ปโร โลโก”ติ วตฺถุ. “นตฺถิ มาตา”ติ วตฺถุ. “นตฺถิ ปิตา”ติ วตฺถุ. “นตฺถิ สตฺตา โอปปาติกา”ติ วตฺถุ. “นตฺถิ โลเก สมณพฺราหฺมณา สมฺมคฺคตา\footnote{สมคฺคตา (ก)} สมฺมาปฏิปนฺนา เย อิมญฺจ โลกํ ปรญฺจ โลกํ สยํ อภิญฺญา สจฺฉิกตฺวา ปเวเทนฺตี”ติ วตฺถุ เอวํวาโท มิจฺฉาภินิเวสปรามาโส ทิฏฺฐิ, ทิฏฺฐิ น วตฺถุ, วตฺถุ น ทิฏฺฐิ, อญฺญา ทิฏฺฐิ, อญฺญํ วตฺถุ, ยา จ ทิฏฺฐิ, ยญฺจ วตฺถุ, อยํ ทสมา มิจฺฉาวตฺถุกา มิจฺฉาทิฏฺฐิ มิจฺฉาทิฏฺฐิ ทิฏฺฐิวิปตฺติ ฯเปฯ.\csromanspacer{} มิจฺฉาทิฏฺฐิกสฺส ปุริสปุคฺคลสฺส เทฺวว คติโย ฯเปฯ อิมานิ สญฺโญชนานิ น จ ทิฏฺฐิโย.\csromanspacer{} มิจฺฉาทิฏฺฐิยา อิเมหิ ทสหากาเรหิ อภินิเวโส โหติ.}
\csromanmatikaanchor{mtk.55}
\csromantocmarkref{subsection}{4. สกฺกายทิฏฺฐินิทฺเทส}{mtk.55}
\bookmark[dest={mtk.55},level=2]{4. สกฺกายทิฏฺฐินิทฺเทส}
\csromanheader{2}{4. สกฺกายทิฏฺฐินิทฺเทส}
\proseitem{137}{\csromanfolio{147}สกฺกายทิฏฺฐิยา กตเมหิ วีสติยา อากาเรหิ อภินิเวโส โหติ? อิธ อสฺสุตวา ปุถุชฺชโน อริยานํ อทสฺสาวี อริยธมฺมสฺส อโกวิโท อริยธมฺเม อวินีโต สปฺปุริสานํ อทสฺสาวี สปฺปุริสธมฺมสฺส อโกวิโท สปฺปุริสธมฺเม อวินีโต \textbf{รูปํ อตฺตโต สมนุปสฺสติ}, รูปวนฺตํ วา อตฺตานํ, อตฺตนิ วา รูปํ, รูปสฺมิํ วา อตฺตานํ.\csromanspacer{} เวทนํ.\csromanspacer{} สญฺญํ.\csromanspacer{} สงฺขาเร.\csromanspacer{} วิญฺญาณํ อตฺตโต สมนุปสฺสติ, วิญฺญาณวนฺตํ วา อตฺตานํ, อตฺตนิ วา วิญฺญาณํ วิญฺญาณสฺมิํ วา อตฺตานํ.}

\prose{กถํ \textbf{รูปํ อตฺตโต สมนุปสฺสติ}? อิเธกจฺโจ ปถวีกสิณํ ฯเปฯ.\csromanspacer{} โอทาตกสิณํ อตฺตโต สมนุปสฺสติ “ยํ โอทาตกสิณํ, โส อหํ.\csromanspacer{} โย อหํ, ตํ โอทาตกสิณนฺ”ติ โอทาตกสิณญฺจ อตฺตญฺจ อทฺวยํ สมนุปสฺสติ.\csromanspacer{} เสยฺยถาปิ เตลปฺปทีปสฺส ฌายโต ฯเปฯ.\csromanspacer{} เอวเมวํ อิเธกจฺโจ โอทาตกสิณํ อตฺตโต สมนุปสฺสติ, อภินิเวสปรามาโส ทิฏฺฐิ ฯเปฯ.\csromanspacer{} อยํ ปฐมา รูปวตฺถุกา สกฺกายทิฏฺฐิ, สกฺกายทิฏฺฐิ มิจฺฉาทิฏฺฐิ ฯเปฯ อิมานิ สญฺโญชนานิ น จ ทิฏฺฐิโย.\csromanspacer{} เอวํ \textbf{รูปํ อตฺตโต สมนุปสฺสติ} ฯเปฯ.\csromanspacer{} สกฺกายทิฏฺฐิยา อิเมหิ วีสติยา อากาเรหิ อภินิเวโส โหติ.}

\csromanmatikaanchor{mtk.56}
\csromantocmarkref{subsection}{5. สสฺสตทิฏฺฐินิทฺเทส}{mtk.56}
\bookmark[dest={mtk.56},level=2]{5. สสฺสตทิฏฺฐินิทฺเทส}
\csromanheader{2}{5. สสฺสตทิฏฺฐินิทฺเทส}
\proseitem{138}{\csromanfolio{148}สกฺกายวตฺถุกาย สสฺสตทิฏฺฐิยา กตเมหิ ปนฺนรสหิ อากาเรหิ อภินิเวโส โหติ? อิธ อสฺสุตวา ปุถุชฺชโน อริยานํ อทสฺสาวี อริยธมฺมสฺส อโกวิโท อริยธมฺเม อวินีโต สปฺปุริสานํ อทสฺสาวี สปฺปุริสธมฺมสฺส อโกวิโท สปฺปุริสธมฺเม อวินีโต รูปวนฺตํ วา อตฺตานํ, อตฺตนิ วา รูปํ, รูปสฺมิํ วา อตฺตานํ.\csromanspacer{} เวทนาวนฺตํ วา อตฺตานํ.\csromanspacer{} สญฺญาวนฺตํ วา อตฺตานํ.\csromanspacer{} สงฺขารวนฺตํ วา อตฺตานํ.\csromanspacer{} วิญฺญาณวนฺตํ วา อตฺตานํ.\csromanspacer{} อตฺตนิ วา วิญฺญาณํ, วิญฺญาณสฺมิํ วา อตฺตานํ.}

\prose{กถํ \textbf{รูปวนฺตํ อตฺตานํ สมนุปสฺสติ}? อิเธกจฺโจ เวทนํ สญฺญํ สงฺขาเร วิญฺญาณํ อตฺตโต สมนุปสฺสติ, ตสฺส เอวํ โหติ “อยํ โข เม อตฺตา, โส โข ปน เม อยํ อตฺตา อิมินา รูเปน รูปวา”ติ \textbf{รูปวนฺตํ อตฺตานํ สมนุปสฺสติ}.\csromanspacer{} เสยฺยถาปิ รุกฺโข ฉายาสมฺปนฺโน อสฺส, ตเมนํ ปุริโส เอวํ วเทยฺย “อยํ รุกฺโข, อยํ ฉายา.\csromanspacer{} อญฺโญ รุกฺโข, อญฺญา ฉายา.\csromanspacer{} โส โข ปนายํ รุกฺโข อิมาย ฉายาย ฉายาวา”ติ ฉายาวนฺตํ รุกฺขํ สมนุปสฺสติ.\csromanspacer{} เอวเมวํ อิเธกจฺโจ เวทนํ ฯเปฯ.\csromanspacer{} อยํ ปฐมา สกฺกายวตฺถุกา สสฺสตทิฏฺฐิ, สสฺสตทิฏฺฐิ มิจฺฉาทิฏฺฐิ ฯเปฯ อิมานิ สญฺโญชนานิ น จ ทิฏฺฐิโย.\csromanspacer{} เอวํ \textbf{รูปวนฺตํ อตฺตานํ สมนุปสฺสติ} ฯเปฯ.\csromanspacer{} สกฺกายวตฺถุกาย สสฺสตทิฏฺฐิยา อิเมหิ ปนฺนรสหิ อากาเรหิ อภินิเวโส โหติ.}

\csromanmatikaanchor{mtk.57}
\csromantocmarkref{subsection}{6. อุจฺเฉททิฏฺฐินิทฺเทส}{mtk.57}
\bookmark[dest={mtk.57},level=2]{6. อุจฺเฉททิฏฺฐินิทฺเทส}
\csromanheader{2}{6. อุจฺเฉททิฏฺฐินิทฺเทส}
\proseitem{139}{สกฺกายวตฺถุกาย อุจฺเฉททิฏฺฐิยา กตเมหิ ปญฺจหิ อากาเรหิ อภินิเวโส โหติ? อิธ อสฺสุตวา ปุถุชฺชโน อริยานํ อทสฺสวี อริยธมฺมสฺส อโกวิโท อริยธมฺเม อวินีโต สปฺปุริสานํ อทสฺสาวี สปฺปุริสธมฺมสฺส อโกวิโท สปฺปุริสธมฺเม อวินีโต รูปํ อตฺตโต สมนุปสฺสติ.\csromanspacer{} เวทนํ อตฺตโต สมนุปสฺสติ.\csromanspacer{} สญฺญํ อตฺตโต สมนุปสฺสติ.\csromanspacer{} สงฺขาเร อตฺตโต สมนุปสฺสติ.\csromanspacer{} วิญฺญาณํ อตฺตโต สมนุปสฺสติ.}

\prose{กถํ \textbf{รูปํ อตฺตโต สมนุปสฺสติ?} อิเธกจฺโจ ปถวีกสิณํ ฯเปฯ.\csromanspacer{} โอทาตกสิณํ อตฺตโต สมนุปสฺสติ “ยํ โอทาตกสิณํ, โส}

\vspace{\tipitakaparskip}
\csromancenter{ทิฏฺฐิกถา อหํ.\csromanspacer{} โย อหํ, ตํ โอทาตกสิณนฺ”ติ โอทาตกสิณญฺจ อตฺตญฺจ อทฺวยํ สมนุปสฺสติ.\csromanspacer{} เสยฺยถาปิ เตลปฺปทีปสฺส ฌายโต ฯเปฯ.\csromanspacer{} อยํ ปฐมา สกฺกายวตฺถุกา อุจฺเฉททิฏฺฐิ, อุจฺเฉททิฏฺฐิ มิจฺฉาทิฏฺฐิ ฯเปฯ อิมานิ สญฺโญชนานิ น จ ทิฏฺฐิโย.\csromanspacer{} เอวํ รูปํ อตฺตโต สมนุปสฺสติ ฯเปฯ.\csromanspacer{} สกฺกายวตฺถุกาย อุจฺเฉททิฏฺฐิยา อิเมหิ ปญฺจหิ อากาเรหิ อภินิเวโส โหติ.}
\csromanmatikaanchor{mtk.58}
\csromantocmarkref{subsection}{7. อนฺตคฺคาหิกาทิฏฺฐินิทฺเทส}{mtk.58}
\bookmark[dest={mtk.58},level=2]{7. อนฺตคฺคาหิกาทิฏฺฐินิทฺเทส}
\csromanheader{2}{7. อนฺตคฺคาหิกาทิฏฺฐินิทฺเทส}
\proseitem{140}{\csromanfolio{149}อนฺตคฺคาหิกาย ทิฏฺฐิยา กตเมหิ ปญฺญาสาย อากาเรหิ อภินิเวโส โหติ? “สสฺสโต โลโก”ติ อนฺตคฺคาหิกาย ทิฏฺฐิยา กติหากาเรหิ อภินิเวโส โหติ, “อสสฺสโต โลโก”ติ อนฺตคฺคาหิกาย ทิฏฺฐิยา กติหากาเรหิ อภินิเวโส โหติ, “อนฺตวา โลโก”ติ อนฺตคฺคาหิกาย ทิฏฺฐิยา. “อนนฺตวา โลโก”ติ อนฺตคฺคาหิกาย ทิฏฺฐิยา. “ตํ ชีวํ ตํ สรีรนฺ”ติ อนฺตคฺคาหิกาย ทิฏฺฐิยา. “อญฺญํ ชีวํ อญฺญํ สรีรนฺ”ติ อนฺตคฺคาหิกาย ทิฏฺฐิยา. “โหติ ตถาคโต ปรํ มรณา”ติ ฯเปฯ “น โหติ ตถาคโต ปรํ มรณา”ติ. “โหติ จ น จ โหติ ตถาคโต ปรํ มรณา”ติ. “เนว โหติ น น โหติ ตถาคโต ปรํ มรณา”ติ อนฺตคฺคาหิกาย ทิฏฺฐิยา กติหากาเรหิ อภินิเวโส โหติ?}

\prose{“สสฺสโต โลโก”ติ อนฺตคฺคาหิกาย ทิฏฺฐิยา ปญฺจหากาเรหิ อภินิเวโส โหติ ฯเปฯ “เนว โหติ น น โหติ ตถาคโต ปรํ มรณา”ติ อนฺตคฺคาหิกาย ทิฏฺฐิยา ปญฺจหากาเรหิ อภินิเวโส โหติ.}

\prose{(ก) “สสฺสโต โลโก”ติ อนฺตคฺคาหิกาย ทิฏฺฐิยา กตเมหิ ปญฺจหากาเรหิ อภินิเวโส โหติ? “รูปํ โลโก เจว สสฺสตญฺจา”ติ อภินิเวสปรามาโส ทิฏฺฐิ, ตาย ทิฏฺฐิยา โส อนฺโต คหิโตติ อนฺตคฺคาหิกา ทิฏฺฐิ, ทิฏฺฐิ น วตฺถุ, วตฺถุ น ทิฏฺฐิ, อญฺญา ทิฏฺฐิ, อญฺญํ วตฺถุ, ยา จ ทิฏฺฐิ ยญฺจ วตฺถุ.\csromanspacer{} อยํ ปฐมา “สสฺสโต โลโก”ติ อนฺตคฺคาหิกา ทิฏฺฐิ, อนฺตคฺคาหิกา ทิฏฺฐิ มิจฺฉาทิฏฺฐิ ฯเปฯ อิมานิ สญฺโญชนานิ น จ ทิฏฺฐิโย.}

\prose{“เวทนา โลโก เจว สสฺสตา จา”ติ ฯเปฯ “สญฺญา โลโก เจว อสฺสตา จา”ติ ฯเปฯ “สงฺขารา โลโก เจว สสฺสตา จา”ติ ฯเปฯ \csromanfolio{150}“วิญฺญาณํ โลโก เจว สสฺสตญฺจา”ติ อภินิเวสปรามาโส ทิฏฺฐิ, ตาย ทิฏฺฐิยา โส อนฺโต คหิโตติ อนฺตคฺคาหิกา ทิฏฺฐิ, ทิฏฺฐิ น วตฺถุ, วตฺถุ น ทิฏฺฐิ, อญฺญา ทิฏฺฐิ, อญฺญํ วตฺถุ, ยา จ ทิฏฺฐิ, ยญฺจ วตฺถุ.\csromanspacer{} อยํ ปญฺจมี “สสฺสโต โลโก”ติ อนฺตคฺคาหิกา ทิฏฺฐิ.\csromanspacer{} อนฺตคฺคาหิกา ทิฏฺฐิ มิจฺฉาทิฏฺฐิ ฯเปฯ อิมานิ สญฺโญชนานิ น จ ทิฏฺฐิโย. “สสฺสโต โลโก”ติ อนฺตคฺคาหิกาย ทิฏฺฐิยา อิเมหิ ปญฺจหากาเรหิ อภินิเวโส โหติ. (5)}

\prose{(ข) “อสสฺสโต โลโก”ติ อนฺตคฺคาหิกาย ทิฏฺฐิยา กตเมหิ ปญฺจหากาเรหิ อภินิเวโส โหติ? “รูปํ โลโก เจว อสสฺสตญฺจา”ติ อภินิเวสปรามาโส ทิฏฺฐิ, ตาย ทิฏฺฐิยา โส อนฺโต คหิโตติ อนฺตคฺคาหิกา ทิฏฺฐิ ฯเปฯ.\csromanspacer{} อยํ ปฐมา “อสสฺสโต โลโก”ติ อนฺตคฺคาหิกา ทิฏฺฐิ, อนฺตคฺคาหิกา ทิฏฺฐิ มิจฺฉาทิฏฺฐิ ทิฏฺฐิวิปตฺติ ฯเปฯ อิมานิ สญฺโญชนานิ น จ ทิฏฺฐิโย.}

\prose{“เวทนา โลโก เจว อสสฺสตา จา”ติ ฯเปฯ “สญฺญา โลโก เจว อสสฺสตา จา”ติ ฯเปฯ “สงฺขารา โลโก เจว อสสฺสตา จา”ติ ฯเปฯ “วิญฺญาณํ โลโก เจว อสสฺสตญฺจา”ติ อภินิเวสปรามาโส ทิฏฺฐิ.\csromanspacer{} ตาย ทิฏฺฐิยา โส อนฺโต คหิโตติ อนฺตคฺคาหิกา ทิฏฺฐิ, อนฺตคฺคาหิกา ทิฏฺฐิ มิจฺฉาทิฏฺฐิ ฯเปฯ อิมานิ สญฺโญชนานิ น จ ทิฏฺฐิโย. “อสสฺสโต โลโก”ติ อนฺตคฺคาหิกาย ทิฏฺฐิยา อิเมหิ ปญฺจหากาเรหิ อภินิเวโส โหติ. (5, 10)}

\prose{(ค) “อนฺตวา โลโก”ติ อนฺตคฺคาหิกาย ทิฏฺฐิยา กตเมหิ ปญฺจหากาเรหิ อภินิเวโส โหติ? อิเธกจฺโจ ปริตฺตํ โอกาสํ นีลกโต ผรติ, ตสฺส เอวํ โหติ “อนฺตวา อยํ โลโก ปริวฏุโม”ติ อนฺตสญฺญี โหติ. “ยํ ผรติ, ตํ วตฺถุ เจว โลโก จ.\csromanspacer{} เยน ผรติ, โส อตฺตา เจว โลโก จา”ติ อภินิเวสปรามาโส ทิฏฺฐิ, ตาย ทิฏฺฐิยา โส อนฺโต คหิโตติ อนฺตคฺคาหิกา ทิฏฺฐิ, ทิฏฺฐิ น วตฺถุ, วตฺถุ น ทิฏฺฐิ, อญฺญา ทิฏฺฐิ, อญฺญํ วตฺถุ, ยา จ ทิฏฺฐิ ยญฺจ วตฺถุ.\csromanspacer{} อยํ ปฐมา “อนฺตวา โลโก”ติ อนฺตคฺคาหิกา ทิฏฺฐิ, อนฺตคฺคาหิกา ทิฏฺฐิ มิจฺฉาทิฏฺฐิ ฯเปฯ อิมานิ สญฺโญชนานิ น จ ทิฏฺฐิโย.}

\csromanheader{2}{2. ทิฏฺฐิกถา}
\prose{\csromanfolio{151}อิเธกจฺโจ ปริตฺตํ โอกาสํ ปีตกโต ผรติ.\csromanspacer{} โลหิตกโต ผรติ.\csromanspacer{} โอทาตกโต ผรติ.\csromanspacer{} โอภาสกโต ผรติ, ตสฺส เอวํ โหติ “อนฺตวา อยํ โลโก ปริวฏุโม”ติ อนฺตสญฺญี โหติ. “ยํ ผรติ, ตํ วตฺถุ เจว โลโก จ.\csromanspacer{} เยน ผรติ, โส อตฺตา เจว โลโก จา”ติ อภินิเวสปรามาโส ทิฏฺฐิ, ตาย ทิฏฺฐิยา โส อนฺโต คหิโตติ อนฺตคฺคาหิกา ทิฏฺฐิ ฯเปฯ “อนฺตวา โลโก”ติ อนฺตคฺคาหิกาย ทิฏฺฐิยา อิเมหิ ปญฺจหากาเรหิ อภินิเวโส โหติ. (5, 15)}

\prose{(ฆ) “อนนฺตวา โลโก”ติ อนฺตคฺคาหิกาย ทิฏฺฐิยา กตเมหิ ปญฺจหากาเรหิ อภินิเวโส โหติ\csromandash{}อิเธกจฺโจ วิปุลํ โอกาสํ นีลกโต ผรติ, ตสฺส เอวํ โหติ “อนนฺตวา อยํ โลโก อปริยนฺโต”ติ อนนฺตสญฺญี โหติ. “ยํ ผรติ, ตํ วตฺถุ เจว โลโก จ.\csromanspacer{} เยน ผรติ, โส อตฺตา เจว โลโก จา”ติ อภินิเวสปรามาโส ทิฏฺฐิ, ตาย ทิฏฺฐิยา โส อนฺโต คหิโตติ อนฺตคฺคาหิกา ทิฏฺฐิ, ทิฏฺฐิ น วตฺถุ, วตฺถุ น ทิฏฺฐิ, อญฺญา ทิฏฺฐิ, อญฺญํ วตฺถุ, ยา จ ทิฏฺฐิ ยญฺจ วตฺถุ.\csromanspacer{} อยํ ปฐมา “อนนฺตวา โลโก”ติ อนฺตคฺคาหิกา ทิฏฺฐิ, อนฺตคฺคาหิกา ทิฏฺฐิ มิจฺฉาทิฏฺฐิ ฯเปฯ อิมานิ สญฺโญชนานิ น จ ทิฏฺฐิโย.}

\prose{อิเธกจฺโจ วิปุลํ โอกาสํ ปีตกโต ผรติ.\csromanspacer{} โลหิตกโต ผรติ.\csromanspacer{} โอทาตกโต ผรติ.\csromanspacer{} โอภาสกโต ผรติ, ตสฺส เอวํ โหติ “อนนฺตวา อยํ โลโก อปริยนฺโต”ติ อนนฺตสญฺญี โหติ. “ยํ ผรติ, ตํ วตฺถุ เจว โลโก จ.\csromanspacer{} เยน ผรติ, โส อตฺตา เจว โลโก จา”ติ อภินิเวสปรามาโส ทิฏฺฐิ, ตาย ทิฏฺฐิยา โส อนฺโต คหิโตติ อนฺตคฺคาหิกา ทิฏฺฐิ ฯเปฯ “อนนฺตวา โลโก”ติ อนฺตคฺคาหิกาย ทิฏฺฐิยา อิเมหิ ปญฺจหากาเรหิ อภินิเวโส โหติ. (5, 20)}

\prose{(\^{}อ) “ตํ ชีวํ ตํ สรีรนฺ”ติ อนฺตคฺคาหิกาย ทิฏฺฐิยา กตเมหิ ปญฺจหากาเรหิ อภินิเวโส โหติ? “รูปํ ชีวญฺเจว สรีรญฺจ, ยํ ชีวํ, ตํ สรีรํ.\csromanspacer{} ยํ สรีรํ, ตํ ชีวนฺ”ติ อภินิเวสปรามาโส ทิฏฺฐิ, ตาย ทิฏฺฐิยา โส อนฺโต คหิโตติ อนฺตคฺคาหิกา ทิฏฺฐิ, ทิฏฺฐิ น วตฺถุ, วตฺถุ น ทิฏฺฐิ, อญฺญา ทิฏฺฐิ, อญฺญํ วตฺถุ, ยา จ ทิฏฺฐิ ยญฺจ วตฺถุ.\csromanspacer{} อยํ ปฐมา “ตํ ชีวํ ตํ สรีรนฺ”ติ อนฺตคฺคาหิกา ทิฏฺฐิ, อนฺตคฺคาหิกา ทิฏฺฐิ มิจฺฉาทิฏฺฐิ ฯเปฯ อิมานิ สญฺโญชนานิ น จ ทิฏฺฐิโย.}

\prose{\csromanfolio{152}“เวทนา ชีวา เจว สรีรญฺจ.\csromanspacer{} สญฺญา ชีวา เจว สรีรญฺจ.\csromanspacer{} สงฺขารา ชีวา เจว สรีรญฺจ.\csromanspacer{} วิญฺญาณํ ชีวํ เจว สรีรญฺจ.\csromanspacer{} ยํ ชีวํ ตํ สรีรํ, ยํ สรีรํ ตํ ชีวนฺ”ติ อภินิเวสปรามาโส ทิฏฺฐิ, ตาย ทิฏฺฐิยา โส อนฺโต คหิโตติ อนฺตคฺคาหิกา ทิฏฺฐิ ฯเปฯ “ตํ ชีวํ ตํ สรีรนฺ”ติ อนฺตคฺคาหิกาย ทิฏฺฐิยา อิเมหิ ปญฺจหากาเรหิ อภินิเวโส โหติ. (5, 25)}

\prose{(จ) “อญฺญํ ชีวํ อญฺญํ สรีรนฺ”ติ อนฺตคฺคาหิกาย ทิฏฺฐิยา กตเมหิ ปญฺจหากาเรหิ อภินิเวโส โหติ? “รูปํ สรีรํ น ชีวํ, ชีวํ น สรีรํ.\csromanspacer{} อญฺญํ ชีวํ อญฺญํ สรีรนฺ”ติ อภินิเวสปรามาโส ทิฏฺฐิ, ตาย ทิฏฺฐิยา โส อนฺโต คหิโตติ อนฺตคฺคาหิกา ทิฏฺฐิ, ทิฏฺฐิ น วตฺถุ, วตฺถุ น ทิฏฺฐิ, อญฺญา ทิฏฺฐิ, อญฺญํ วตฺถุ, ยา จ ทิฏฺฐิ ยญฺจ วตฺถุ.\csromanspacer{} อยํ ปฐมา “อญฺญํ ชีวํ อญฺญํ สรีรนฺ”ติ อนฺตคฺคาหิกา ทิฏฺฐิ, อนฺตคฺคาหิกา ทิฏฺฐิ มิจฺฉาทิฏฺฐิ ฯเปฯ อิมานิ สญฺโญชนานิ น จ ทิฏฺฐิโย.}

\prose{“เวทนา สรีรํ น ชีวํ.\csromanspacer{} สญฺญา สรีรํ น ชีวํ.\csromanspacer{} สงฺขารา สรีรํ น ชีวํ.\csromanspacer{} วิญฺญาณํ สรีรํ น ชีวํ.\csromanspacer{} ชีวํ น สรีรํ.\csromanspacer{} อญฺญํ ชีวํ, อญฺญํ สรีรนฺ”ติ อภินิเวสปรามาโส ทิฏฺฐิ, ตาย ทิฏฺฐิยา โส อนฺโต คหิโตติ อนฺตคฺคาหิกา ทิฏฺฐิ ฯเปฯ “อญฺญํ ชีวํ อญฺญํ สรีรนฺ”ติ อนฺตคฺคาหิกาย ทิฏฺฐิยา อิเมหิ ปญฺจหากาเรหิ อภินิเวโส โหติ. (5, 30)}

\prose{(ฉ) “โหติ ตถาคโต ปรํ มรณา”ติ อนฺตคฺคาหิกาย ทิฏฺฐิยา กตเมหิ ปญฺจหากาเรหิ อภินิเวโส โหติ? “รูปํ อิเธว มรณธมฺมํ, ตถาคโต กายสฺส เภทา โหติปิ ติฏฺฐติปิ อุปฺปชฺชติปิ นิพฺพตฺตติปี”ติ อภินิเวสปรามาโส ทิฏฺฐิ, ตาย ทิฏฺฐิยา โส อนฺโต คหิโตติ อนฺตคฺคาหิกา ทิฏฺฐิ, ทิฏฺฐิ น วตฺถุ, วตฺถุ น ทิฏฺฐิ, อญฺญา ทิฏฺฐิ, อญฺญํ วตฺถุ, ยา จ ทิฏฺฐิ ยญฺจ วตฺถุ.\csromanspacer{} อยํ ปฐมา “โหติ ตถาคโต ปรํ มรณา”ติ อนฺตคฺคาหิกา ทิฏฺฐิ, อนฺตคฺคาหิกา ทิฏฺฐิ มิจฺฉาทิฏฺฐิ ฯเปฯ อิมานิ สญฺโญชนานิ น จ ทิฏฺฐิโย.}

\prose{“เวทนา อิเธว มรณธมฺมา.\csromanspacer{} สญฺญา อิเธว มรณธมฺมา.\csromanspacer{} สงฺขารา อิเธว มรณธมฺมา.\csromanspacer{} วิญฺญาณํ อิเธว มรณธมฺมํ, ตถาคโต กายสฺส เภทา โหติปิ ติฏฺฐติปิ อุปฺปชฺชติปิ นิพฺพตฺตติปี”ติ อภินิเวสปรามาโส ทิฏฺฐิ, ตาย ทิฏฺฐิยา โส อนฺโต คหิโตติ อนฺตคฺคาหิกา ทิฏฺฐิ ฯเปฯ “โหติ}

\vspace{\tipitakaparskip}
\csromancenter{ทิฏฺฐิกถา ตถาคโต ปรํ มรณา”ติ อนฺตคฺคาหิกาย ทิฏฺฐิยา อิเมหิ ปญฺจหากาเรหิ อภินิเวโส โหติ. (5, 35)}
\prose{\csromanfolio{153}(ช) “น โหติ ตถาคโต ปรํ มรณา”ติ อนฺตคฺคาหิกาย ทิฏฺฐิยา กตเมหิ ปญฺจหากาเรหิ อภินิเวโส โหติ? “รูปํ อิเธว มรณธมฺมํ, ตถาคโตปิ กายสฺส เภทา อุจฺฉิชฺชติ วินสฺสติ น โหติ ตถาคโต ปรํ มรณา”ติ อภินิเวสปรามาโส ทิฏฺฐิ, ตาย ทิฏฺฐิยา โส อนฺโต คหิโตติ อนฺตคฺคาหิกา ทิฏฺฐิ, ทิฏฺฐิ น วตฺถุ, วตฺถุ น ทิฏฺฐิ, อญฺญา ทิฏฺฐิ, อญฺญํ วตฺถุ.\csromanspacer{} ยา จ ทิฏฺฐิ ยญฺจ วตฺถุ.\csromanspacer{} อยํ ปฐมา “น โหติ ตถาคโต ปรํ มรณา”ติ อนฺตคฺคาหิกา ทิฏฺฐิ, อนฺตคฺคาหิกา ทิฏฺฐิ มิจฺฉาทิฏฺฐิ ฯเปฯ อิมานิ สญฺโญชนานิ น จ ทิฏฺฐิโย.}

\prose{“เวทนา อิเธว มรณธมฺมา.\csromanspacer{} สญฺญา อิเธว มรณธมฺมา.\csromanspacer{} สงฺขารา อิเธว มรณธมฺมา.\csromanspacer{} วิญฺญาณํ อิเธว มรณธมฺมํ, ตถาคโตปิ กายสฺส เภทา อุจฺฉิชฺชติ วินสฺสติ น โหติ ตถาคโต ปรํ มรณา”ติ อภินิเวสปรามาโส ทิฏฺฐิ.\csromanspacer{} ตาย ทิฏฺฐิยา โส อนฺโต คหิโตติ อนฺตคฺคาหิกา ทิฏฺฐิ ฯเปฯ “น โหติ ตถาคโต ปรํ มรณา”ติ อนฺตคฺคาหิกาย ทิฏฺฐิยา อิเมหิ ปญฺจหากาเรหิ อภินิเวโส โหติ. (5, 40)}

\prose{(ฌ) “โหติ จ น จ โหติ ตถาคโต ปรํ มรณา”ติ อนฺตคฺคาหิกาย ทิฏฺฐิยา กตเมหิ ปญฺจหากาเรหิ อภินิเวโส โหติ? “รูปํ อิเธว มรณธมฺมํ, ตถาคโต กายสฺส เภทา โหติ จ น จ โหตี”ติ อภินิเวสปรามาโส ทิฏฺฐิ, ตาย ทิฏฺฐิยา โส อนฺโต คหิโตติ อนฺตคฺคาหิกา ทิฏฺฐิ, ทิฏฺฐิ น วตฺถุ, วตฺถุ น ทิฏฺฐิ, อญฺญา ทิฏฺฐิ, อญฺญํ วตฺถุ, ยา จ ทิฏฺฐิ ยญฺจ วตฺถุ.\csromanspacer{} อยํ ปฐมา “โหติ จ น จ โหติ ตถาคโต ปรํ มรณา”ติ อนฺตคฺคาหิกา ทิฏฺฐิ.\csromanspacer{} อนฺตคฺคาหิกา ทิฏฺฐิ มิจฺฉาทิฏฺฐิ ฯเปฯ อิมานิ สญฺโญชนานิ น จ ทิฏฺฐิโย.}

\prose{“เวทนา อิเธว มรณธมฺมา.\csromanspacer{} สญฺญา อิเธว มรณธมฺมา.\csromanspacer{} สงฺขารา อิเธว มรณธมฺมา.\csromanspacer{} วิญฺญาณํ อิเธว มรณธมฺมํ, ตถาคโต กายสฺส เภทา โหติ จ น จ โหตี”ติ อภินิเวสปรามาโส ทิฏฺฐิ, ตาย ทิฏฺฐิยา โส อนฺโต คหิโตติ อนฺตคฺคาหิกา ทิฏฺฐิ ฯเปฯ “โหติ จ น จ โหติ ตถาคโต ปรํ มรณา”ติ อนฺตคฺคาหิกาย ทิฏฺฐิยา อิเมหิ ปญฺจหากาเรหิ อภินิเวโส โหติ. (5, 45)}

\prose{\csromanfolio{154}(ญ) “เนว โหติ น น โหติ ตถาคโต ปรํ มรณา”ติ อนฺตคฺคาหิกาย ทิฏฺฐิยา กตเมหิ ปญฺจหากาเรหิ อภินิเวโส โหติ? “รูปํ อิเธว มรณธมฺมํ, ตถาคโต กายสฺส เภทา ปรํ มรณา เนว โหติ น น โหตี”ติ อภินิเวสปรามาโส ทิฏฺฐิ, ตาย ทิฏฺฐิยา โส อนฺโต คหิโตติ อนฺตคฺคาหิกา ทิฏฺฐิ, ทิฏฺฐิ น วตฺถุ, วตฺถุ น ทิฏฺฐิ, อญฺญา ทิฏฺฐิ, อญฺญํ วตฺถุ, ยา จ ทิฏฺฐิ, ยญฺจ วตฺถุ.\csromanspacer{} อยํ ปฐมา “เนว โหติ น น โหติ ตถาคโต ปรํ มรณา”ติ อนฺตคฺคาหิกา ทิฏฺฐิ, อนฺตคฺคาหิกา ทิฏฺฐิ มิจฺฉาทิฏฺฐิ ฯเปฯ อิมานิ สญฺโญชนานิ น จ ทิฏฺฐิโย.}

\prose{“เวทนา อิเธว มรณธมฺมา.\csromanspacer{} สญฺญา อิเธว มรณธมฺมา.\csromanspacer{} สงฺขารา อิเธว มรณธมฺมา.\csromanspacer{} วิญฺญาณํ อิเธว มรณธมฺมํ, ตถาคโต กายสฺส เภทา ปรํ มรณา เนว โหติ น น โหตี”ติ อภินิเวสปรามาโส ทิฏฺฐิ, ตาย ทิฏฺฐิยา โส อนฺโต คหิโตติ อนฺตคฺคาหิกา ทิฏฺฐิ, ทิฏฺฐิ น วตฺถุ, วตฺถุ น ทิฏฺฐิ, อญฺญา ทิฏฺฐิ, อญฺญํ วตฺถุ, ยา จ ทิฏฺฐิ ยญฺจ วตฺถุ.\csromanspacer{} อยํ ปญฺจมี “เนว โหติ น น โหติ ตถาคโต ปรํ มรณา”ติ อนฺตคฺคาหิกา ทิฏฺฐิ, อนฺตคฺคาหิกา ทิฏฺฐิ มิจฺฉาทิฏฺฐิ ฯเปฯ อิมานิ สญฺโญชนานิ น จ ทิฏฺฐิโย. “เนว โหติ น น โหติ ตถาคโต ปรํ มรณา”ติ อนฺตคฺคาหิกาย ทิฏฺฐิยา อิเมหิ ปญฺจหากาเรหิ อภินิเวโส โหติ.\csromanspacer{} อนฺตคฺคาหิกาย ทิฏฺฐิยา อิเมหิ ปญฺญาสาย อากาเรหิ อภินิเวโส โหติ. (50)}

\csromanheader{2}{อนฺตคฺคาหิกาทิฏฺฐินิทฺเทโส สตฺตโม.}
\csromansectionrule
\csromanmatikaanchor{mtk.59}
\csromantocmarkref{subsection}{8. ปุพฺพนฺตานุทิฏฺฐินิทฺเทส}{mtk.59}
\bookmark[dest={mtk.59},level=2]{8. ปุพฺพนฺตานุทิฏฺฐินิทฺเทส}
\csromanheader{2}{8. ปุพฺพนฺตานุทิฏฺฐินิทฺเทส}
\proseitem{141}{ปุพฺพนฺตานุทิฏฺฐิยา กตเมหิ อฏฺฐารสหิ อากาเรหิ อภินิเวโส โหติ? จตฺตาโร สสฺสตวาทา, จตฺตาโร เอกจฺจสสฺสติกา, จตฺตาโร อนฺตานนฺติกา, จตฺตาโร อมราวิกฺเขปิกา, เทฺว อธิจฺจสมุปฺปนฺนิกา.\csromanspacer{} ปุพฺพนฺตานุทิฏฺฐิยา อิเมหิ อฏฺฐารสหิ อากาเรหิ อภินิเวโส โหติ.}

\csromanmatikaanchor{mtk.60}
\csromantocmarkref{subsection}{9. อปรนฺตานุทิฏฺฐินิทฺเทส}{mtk.60}
\bookmark[dest={mtk.60},level=2]{9. อปรนฺตานุทิฏฺฐินิทฺเทส}
\csromanheader{2}{9. อปรนฺตานุทิฏฺฐินิทฺเทส}
\proseitem{142}{อปรนฺตานุทิฏฺฐิยา กตเมหิ จตุจตฺตาลีสาย อากาเรหิ อภินิเวโส โหติ? โสฬส สญฺญีวาทา, อฏฺฐ อสญฺญีวาทา, อฏฺฐ}

\vspace{\tipitakaparskip}
\csromancenter{ทิฏฺฐิกถา เนวสญฺญีนาสญฺญีวาทา, สตฺต อุจฺเฉทวาทา, ปญฺจ ทิฏฺฐธมฺมนิพฺพานวาทา.\csromanspacer{} อปรนฺตานุทิฏฺฐิยา อิเมหิ จตุจตฺตาลีสาย อากาเรหิ อภินิเวโส โหติ.}
\csromanmatikaanchor{mtk.61}
\csromantocmarkref{subsection}{10-12. สญฺโญชนิกาทิทิฏฺฐินิทฺเทส}{mtk.61}
\bookmark[dest={mtk.61},level=2]{10-12. สญฺโญชนิกาทิทิฏฺฐินิทฺเทส}
\csromanheader{2}{10-12. สญฺโญชนิกาทิทิฏฺฐินิทฺเทส}
\proseitem{143}{\csromanfolio{155}สญฺโญชนิกาย ทิฏฺฐิยา กตเมหิ อฏฺฐารสหิ อากาเรหิ อภินิเวโส โหติ? ยา ทิฏฺฐิ ทิฏฺฐิคตํ ทิฏฺฐิคหนํ ฯเปฯ ทิฏฺฐาภินิเวโส ทิฏฺฐิปรามาโส.\csromanspacer{} สญฺโญชนิกาย ทิฏฺฐิยา อิเมหิ อฏฺฐารสหิ อากาเรหิ อภินิเวโส โหติ.}

\proseitem{144}{“อหนฺ”ติ มานวินิพนฺธาย ทิฏฺฐิยา กตเมหิ อฏฺฐารสหิ อากาเรหิ อภินิเวโส โหติ? “จกฺขุ อหนฺ”ติ อภินิเวสปรามาโส “อหนฺ”ติ มานวินิพนฺธา ทิฏฺฐิ, ทิฏฺฐิ น วตฺถุ, วตฺถุ น ทิฏฺฐิ, อญฺญา ทิฏฺฐิ, อญฺญํ วตฺถุ, ยา จ ทิฏฺฐิ, ยญฺจ วตฺถุ, อยํ ปฐมา “อหนฺ”ติ มานวินิพนฺธา ทิฏฺฐิ.\csromanspacer{} มานวินิพนฺธา ทิฏฺฐิ มิจฺฉาทิฏฺฐิ ฯเปฯ อิมานิ สญฺโญชนานิ น จ ทิฏฺฐิโย.}

\prose{“โสตํ อหนฺ”ติ ฯเปฯ “ฆานํ อหนฺ”ติ ฯเปฯ “ชิวฺหา อหนฺ”ติ ฯเปฯ “กาโย อหนฺ”ติ ฯเปฯ “มโน อหนฺ”ติ ฯเปฯ “รูปา อหนฺ”ติ ฯเปฯ “ธมฺมา อหนฺ”ติ. “จกฺขุวิญฺญาณํ อหนฺ”ติ ฯเปฯ “มโนวิญฺญาณํ อหนฺ”ติ อภินิเวสปรามาโส “อหนฺ”ติ มานวินิพนฺธา ทิฏฺฐิ, ทิฏฺฐิ น วตฺถุ, วตฺถุ น ทิฏฺฐิ, อญฺญา ทิฏฺฐิ, อญฺญํ วตฺถุ, ยา จ ทิฏฺฐิ ยญฺจ วตฺถุ, อยํ อฏฺฐารสมี “อหนฺ”ติ มานวินิพนฺธา ทิฏฺฐิ, มานวินิพนฺธา ทิฏฺฐิ มิจฺฉาทิฏฺฐิ ฯเปฯ อิมานิ สญฺโญชนานิ น จ ทิฏฺฐิโย. “อหนฺ”ติ มานวินิพนฺธาย ทิฏฺฐิยา อิเมหิ อฏฺฐารสหิ อากาเรหิ อภินิเวโส โหติ.}

\proseitem{145}{“มมนฺ”ติ มานวินิพนฺธาย ทิฏฺฐิยา กตเมหิ อฏฺฐารสหิ อากาเรหิ อภินิเวโส โหติ? “จกฺขุ มมนฺ”ติ อภินิเวสปรามาโส “มมนฺ”ติ มานวินิพนฺธา ทิฏฺฐิ, ทิฏฺฐิ น วตฺถุ, วตฺถุ น ทิฏฺฐิ, อญฺญา ทิฏฺฐิ, อญฺญํ วตฺถุ, ยา จ ทิฏฺฐิ ยญฺจ วตฺถุ, อยํ ปฐมา “มมนฺ”ติ มานวินิพนฺธา ทิฏฺฐิ.\csromanspacer{} มานวินิพนฺธา ทิฏฺฐิ มิจฺฉาทิฏฺฐิ ฯเปฯ อิมานิ สญฺโญชนานิ น จ ทิฏฺฐิโย.}

\prose{“โสตํ มมนฺ”ติ ฯเปฯ “ฆานํ มมนฺ”ติ ฯเปฯ “ชิวฺหา มมนฺ”ติ ฯเปฯ “กาโย มมนฺ”ติ ฯเปฯ “มโน มมนฺ”ติ ฯเปฯ “รูปา มมนฺ”ติ ฯเปฯ “ธมฺมา มมนฺ”ติ ฯเปฯ “จกฺขุวิญฺญาณํ มมนฺ”ติ ฯเปฯ “มโนวิญฺญาณํ มมนฺ”ติ อภินิเวสปรามาโส \csromanfolio{156}“มมนฺ”ติ มานวินิพนฺธา ทิฏฺฐิ, ทิฏฺฐิ น วตฺถุ, วตฺถุ น ทิฏฺฐิ, อญฺญา ทิฏฺฐิ, อญฺญํ วตฺถุ, ยา จ ทิฏฺฐิ, ยญฺจ วตฺถุ, อยํ อฏฺฐารสมี “มมนฺ”ติ มานวินิพนฺธา ทิฏฺฐิ, มานวินิพนฺธา ทิฏฺฐิ มิจฺฉาทิฏฺฐิ ฯเปฯ อิมานิ สญฺโญชนานิ น จ ทิฏฺฐิโย. “มมนฺ”ติ มานวินิพนฺธาย ทิฏฺฐิยา อิเมหิ อฏฺฐารสหิ อากาเรหิ อภินิเวโส โหติ.}

\csromanmatikaanchor{mtk.62}
\csromantocmarkref{subsection}{13. อตฺตวาทปฏิสํยุตฺตทิฏฺฐินิทฺเทส}{mtk.62}
\bookmark[dest={mtk.62},level=2]{13. อตฺตวาทปฏิสํยุตฺตทิฏฺฐินิทฺเทส}
\csromanheader{2}{13. อตฺตวาทปฏิสํยุตฺตทิฏฺฐินิทฺเทส}
\proseitem{146}{อตฺตวาทปฏิสํยุตฺตาย ทิฏฺฐิยา กตเมหิ วีสติยา อากาเรหิ อภินิเวโส โหติ? อิธ อสฺสุตวา ปุถุชฺชโน อริยานํ อทสฺสาวี อริยธมฺมสฺส อโกวิโท อริยธมฺเม อวินีโต สปฺปุริสานํ อทสฺสาวี สปฺปุริสธมฺมสฺส อโกวิโท สปฺปุริสธมฺเม อวินีโต รูปํ อตฺตโต สมนุปสฺสติ, รูปวนฺตํ วา อตฺตานํ, อตฺตนิ วา รูปํ, รูปสฺมิํ วา อตฺตานํ ฯเปฯ เวทนํ.\csromanspacer{} สญฺญํ.\csromanspacer{} สงฺขาเร.\csromanspacer{} วิญฺญาณํ อตฺตโต สมนุปสฺสติ, วิญฺญาณวนฺตํ วา อตฺตานํ, อตฺตนิ วา วิญฺญาณํ, วิญฺญาณสฺมิํ วา อตฺตานํ ฯเปฯ.}

\prose{กถํ \textbf{รูปํ อตฺตโต สมนุปสฺสติ?} อิเธกจฺโจ ปถวีกสิณํ ฯเปฯ.\csromanspacer{} โอทาตกสิณํ อตฺตโต สมนุปสฺสติ. “ยํ โอทาตกสิณํ, โส อหํ.\csromanspacer{} โย อหํ, ตํ โอทาตกสิณนฺ”ติ โอทาตกสิณญฺจ อตฺตญฺจ อทฺวยํ สมนุปสฺสติ.\csromanspacer{} เสยฺยถาปิ เตลปฺปทีปสฺส ฌายโต ยา อจฺจิ, โส วณฺโณ.\csromanspacer{} โย วณฺโณ, สา อจฺจีติ อจฺจิญฺจ วณฺณญฺจ อทฺวยํ สมนุปสฺสติ.\csromanspacer{} เอวเมวํ อิเธกจฺโจ โอทาตกสิณํ อตฺตโต สมนุปสฺสติ ฯเปฯ อยํ ปฐมา รูปวตฺถุกา อตฺตวาทปฏิสํยุตฺตา ทิฏฺฐิ.\csromanspacer{} อตฺตวาทปฏิสํยุตฺตา ทิฏฺฐิ มิจฺฉาทิฏฺฐิ ฯเปฯ อิมานิ สญฺโญชนานิ น จ ทิฏฺฐิโย, เอวํ รูปํ อตฺตโต สมนุปสฺสติ ฯเปฯ.\csromanspacer{} อตฺตวาทปฏิสํยุตฺตาย ทิฏฺฐิยา อิเมหิ วีสติยา อากาเรหิ อภินิเวโส โหติ.}

\csromanmatikaanchor{mtk.63}
\csromantocmarkref{subsection}{14. โลกวาทปฏิสํยุตฺตทิฏฺฐินิทฺเทส}{mtk.63}
\bookmark[dest={mtk.63},level=2]{14. โลกวาทปฏิสํยุตฺตทิฏฺฐินิทฺเทส}
\csromanheader{2}{14. โลกวาทปฏิสํยุตฺตทิฏฺฐินิทฺเทส}
\proseitem{147}{โลกวาทปฏิสํยุตฺตาย ทิฏฺฐิยา กตเมหิ อฏฺฐหิ อากาเรหิ อภินิเวโส โหติ? “สสฺสโต อตฺตา จ โลโก จา”ติ อภินิเวสปรามาโส โลกวาทปฏิสํยุตฺตา ทิฏฺฐิ, ทิฏฺฐิ น วตฺถุ, วตฺถุ น ทิฏฺฐิ, อญฺญา ทิฏฺฐิ, อญฺญํ วตฺถุ, ยา จ ทิฏฺฐิ ยญฺจ วตฺถุ.\csromanspacer{} อยํ ปฐมา โลกวาทปฏิสํยุตฺตา ทิฏฺฐิ.\csromanspacer{} โลกวาทปฏิสํยุตฺตา ทิฏฺฐิ มิจฺฉาทิฏฺฐิ ฯเปฯ อิมานิ สญฺโญชนานิ น จ ทิฏฺฐิโย.}

\csromanheader{2}{2. ทิฏฺฐิกถา}
\prose{\csromanfolio{157}“อสสฺสโต อตฺตา จ โลโก จา”ติ ฯเปฯ “สสฺสโต จ อสสฺสโต จ อตฺตา จ โลโก จา”ติ ฯเปฯ “เนว สสฺสโต นาสสฺสโต อตฺตา จ โลโก จา”ติ. “อนฺตวา อตฺตา จ โลโก จา”ติ. “อนนฺตวา อตฺตา จ โลโก จา”ติ. “อนฺตวา จ อนนฺตวา จ อตฺตา จ โลโก จา”ติ. “เนว อนฺตวา น อนนฺตวา อตฺตา จ โลโก จา”ติ อภินิเวสปรามาโส โลกวาทปฏิสํยุตฺตา ทิฏฺฐิ.\csromanspacer{} ทิฏฺฐิ น วตฺถุ, วตฺถุ น ทิฏฺฐิ, อญฺญา ทิฏฺฐิ, อญฺญํ วตฺถุ, ยา จ ทิฏฺฐิ ยญฺจ วตฺถุ.\csromanspacer{} อยํ อฏฺฐมี โลกวาทปฏิสํยุตฺตา ทิฏฺฐิ.\csromanspacer{} โลกวาทปฏิสํยุตฺตา ทิฏฺฐิ มิจฺฉาทิฏฺฐิ ฯเปฯ อิมานิ สญฺโญชนานิ น จ ทิฏฺฐิโย.\csromanspacer{} โลกวาทปฏิสํยุตฺตาย ทิฏฺฐิยา อิเมหิ อฏฺฐหิ อากาเรหิ อภินิเวโส โหติ.}

\csromanmatikaanchor{mtk.64}
\csromantocmarkref{subsection}{15-16. ภว-วิภวทิฏฺฐินิทฺเทส}{mtk.64}
\bookmark[dest={mtk.64},level=2]{15-16. ภว-วิภวทิฏฺฐินิทฺเทส}
\csromanheader{2}{15-16. ภว-วิภวทิฏฺฐินิทฺเทส}
\proseitem{148}{โอลียนาภินิเวโส ภวทิฏฺฐิ, อติธาวนาภินิเวโส วิภวทิฏฺฐิ, อสฺสาททิฏฺฐิยา ปญฺจติํสาย อากาเรหิ อภินิเวโส กติ ภวทิฏฺฐิโย, กติ วิภวทิฏฺฐิโย.\csromanspacer{} อตฺตานุทิฏฺฐิยา วีสติยา อากาเรหิ อภินิเวโส กติ ภวทิฏฺฐิโย, กติ วิภวทิฏฺฐิโย ฯเปฯ.\csromanspacer{} โลกวาทปฏิสํยุตฺตาย ทิฏฺฐิยา อฏฺฐหิ อากาเรหิ อภินิเวโส กติ ภวทิฏฺฐิโย, กติ วิภวทิฏฺฐิโย?}

\prose{อสฺสาททิฏฺฐิยา ปญฺจาติํสาย อากาเรหิ อภินิเวโส สิยา ภวทิฏฺฐิโย, สิยา วิภวทิฏฺฐิโย.\csromanspacer{} อตฺตานุทิฏฺฐิยา วีสติยา อากาเรหิ อภินิเวโส ปนฺนรสภวทิฏฺฐิโย, ปญฺจ วิภวทิฏฺฐิโย.\csromanspacer{} มิจฺฉาทิฏฺฐิยา ทสหิ อากาเรหิ อภินิเวโส, สพฺพาว ตา วิภวทิฏฺฐิโย.\csromanspacer{} สกฺกายทิฏฺฐิยา วีสติยา อากาเรหิ อภินิเวโส ปนฺนรส ภวทิฏฺฐิโย, ปญฺจ วิภวทิฏฺฐิโย.\csromanspacer{} สกฺกายวตฺถุกาย สสฺสตทิฏฺฐิยา ปนฺนรสหิ อากาเรหิ อภินิเวโส, สพฺพาว ตา ภวทิฏฺฐิโย.\csromanspacer{} สกฺกายวตฺถุกาย อุจฺเฉททิฏฺฐิยา ปญฺจหากาเรหิ อภินิเวโส, สพฺพาว ตา วิภวทิฏฺฐิโย.}

\prose{“สสฺสโต โลโก”ติ อนฺตคฺคาหิกาย ทิฏฺฐิยา ปญฺจหากาเรหิ อภินิเวโส, สพฺพาว ตา ภวทิฏฺฐิโย, “อสสฺสโต โลโก”ติ อนฺตคฺคาหิกาย ทิฏฺฐิยา ปญฺจหากาเรหิ อภินิเวโส, สพฺพาว ตา \csromanfolio{158}วิภวทิฏฺฐิโย. “อนฺตวา โลโก”ติ อนฺตคฺคาหิกาย ทิฏฺฐิยา ปญฺจหากาเรหิ อภินิเวโส สิยา ภวทิฏฺฐิโย, สิยา วิภวทิฏฺฐิโย. “อนนฺตวา โลโก”ติ อนฺตคฺคาหิกาย ทิฏฺฐิยา ปญฺจหากาเรหิ อภินิเวโส สิยา ภวทิฏฺฐิโย, สิยา วิภวทิฏฺฐิโย. “ตํ ชีวํ ตํ สรีรนฺ”ติ อนฺตคฺคาหิกาย ทิฏฺฐิยา ปญฺจหากาเรหิ อภินิเวโส, สพฺพาว ตา วิภวทิฏฺฐิโย. “อญฺญํ ชีวํ อญฺญํ สรีรนฺ”ติ อนฺตคฺคาหิกาย ทิฏฺฐิยา ปญฺจหากาเรหิ อภินิเวโส, สพฺพาว ตา ภวทิฏฺฐิโย. “โหติ ตถาคโต ปรํ มรณา”ติ อนฺตคฺคาหิกาย ทิฏฺฐิยา ปญฺจหากาเรหิ อภินิเวโส, สพฺพาว ตา ภวทิฏฺฐิโย. “น โหติ ตถาคโต ปรํ มรณา”ติ อนฺตคฺคาหิกาย ทิฏฺฐิยา ปญฺจหากาเรหิ อภินิเวโส, สพฺพาว ตา วิภวทิฏฺฐิโย. “โหติ จ น จ โหติ ตถาคโต ปรํ มรณา”ติ อนฺตคฺคาหิกาย ทิฏฺฐิยา ปญฺจหากาเรหิ อภินิเวโส สิยา ภวทิฏฺฐิโย, สิยา วิภวทิฏฺฐิโย. “เนว โหติ น น โหติ ตถาคโต ปรํ มรณา”ติ อนฺตคฺคาหิกาย ทิฏฺฐิยา ปญฺจหากาเรหิ อภินิเวโส สิยา ภวทิฏฺฐิโย, สิยา วิภวทิฏฺฐิโย.}

\prose{ปุพฺพนฺตานุทิฏฺฐิยา อฏฺฐารสหิ อากาเรหิ อภินิเวโส สิยา ภวทิฏฺฐิโย, สิยา วิภวทิฏฺฐิโย.\csromanspacer{} อปรนฺตานุทิฏฺฐิยา จตุจตฺตารีสาย อากาเรหิ อภินิเวโส สิยา ภวทิฏฺฐิโย, สิยา วิภวทิฏฺฐิโย.\csromanspacer{} สญฺโญชนิกาย ทิฏฺฐิยา อฏฺฐารสหิ อากาเรหิ อภินิเวโส สิยา ภวทิฏฺฐิโย, สิยา วิภวทิฏฺฐิโย. “อหนฺ”ติ มานวินิพนฺธาย ทิฏฺฐิยา อฏฺฐารสหิ อากาเรหิ อภินิเวโส สพฺพาว ตา วิภวทิฏฺฐิโย. “มมนฺ”ติ มานวินิพนฺธาย ทิฏฺฐิยา อฏฺฐารสหิ อากาเรหิ อภินิเวโส สพฺพาว ตา ภวทิฏฺฐิโย.\csromanspacer{} อตฺตวาทปฏิสํยุตฺตาย ทิฏฺฐิยา วีสติยา อากาเรหิ อภินิเวโส ปนฺนรส ภวทิฏฺฐิโย, ปญฺจ วิภวทิฏฺฐิโย.\csromanspacer{} โลกวาทปฏิสํยุตฺตาย ทิฏฺฐิยา อฏฺฐหิ อากาเรหิ อภินิเวโส สิยา ภวทิฏฺฐิโย, สิยา วิภวทิฏฺฐิโย.}

\prose{สพฺพาว ตา ทิฏฺฐิโย อสฺสาททิฏฺฐิโย.\csromanspacer{} สพฺพาว ตา ทิฏฺฐิโย อตฺตานุทิฏฺฐิโย.\csromanspacer{} สพฺพาว ตา ทิฏฺฐิโย มิจฺฉาทิฏฺฐิโย.\csromanspacer{} สพฺพาว ตา ทิฏฺฐิโย สกฺกายทิฏฺฐิโย สพฺพาว ตา ทิฏฺฐิโย อนฺตคฺคาหิกา ทิฏฺฐิโย.\csromanspacer{} สพฺพาว ตา ทิฏฺฐิโย สญฺโญชนิกา ทิฏฺฐิโย.\csromanspacer{} สพฺพาว ตา ทิฏฺฐิโย อตฺตวาทปฏิสํยุตฺตา ทิฏฺฐิโย.}

\csromanheader{2}{2. ทิฏฺฐิกถา}
\csromangathameasure{ภวญฺจ ทิฏฺฐิํ วิภวญฺจ ทิฏฺฐิํ, \\ เอตํ ทฺวยํ ตกฺกิกา นิสฺสิตาเส. \\ เตสํ นิโรธมฺหิ น ห’ตฺถิ ญาณํ, \\ ยตฺถายํ โลโก วิปรีตสญฺญีติ.}
\csromangathagroup{\csromangathastanza{\csromanfolio{159}ภวญฺจ ทิฏฺฐิํ วิภวญฺจ ทิฏฺฐิํ, \\ เอตํ ทฺวยํ ตกฺกิกา นิสฺสิตาเส. \\ เตสํ นิโรธมฺหิ น ห’ตฺถิ ญาณํ, \\ ยตฺถายํ โลโก วิปรีตสญฺญีติ.}}

\proseitem{149}{ทฺวีหิ ภิกฺขเว ทิฏฺฐิคเตหิ ปริยุฏฺฐิตา เทวมนุสฺสา โอลียนฺติ\footnote{โอลิยนฺติ (สฺยา, ก) ขุ 1. 224 ปิฏฺเฐ ปสฺสิตพฺพา.} เอเก, อติธาวนฺติ เอเก, จกฺขุมนฺโต จ ปสฺสนฺติ.\csromanspacer{} กถญฺจ ภิกฺขเว โอลียนฺติ เอเก? ภวารามา ภิกฺขเว เทวมนุสฺสา ภวรตา ภวสมฺมุทิตา, เตสํ ภวนิโรธาย ธมฺเม เทสิยมาเน จิตฺตํ น ปกฺขนฺทติ น ปสีทติ น สนฺติฏฺฐติ นาธิมุจฺจติ.\csromanspacer{} เอวํ โข ภิกฺขเว โอลียนฺติ เอเก.}

\prose{กถญฺจ ภิกฺขเว อติธาวนฺติ เอเก? ภเวเนว โข ปเนเก อฏฺฏียมานา\footnote{อฏฺฏิยมานา (สฺยา, ก) ขุ 1. 224 ปิฏฺเฐ ปสฺสิตพฺพา.} หรายมานา ชิคุจฺฉมานา วิภวํ อภินนฺทนฺติ, ยโต กิร โภ อยํ อตฺตา กายสฺส เภทา ปรํ มรณา อุจฺฉิชฺชติ วินสฺสติ น โหติ ปรํ มรณา, เอตํ สนฺตํ เอตํ ปณีตํ เอตํ ยาถาวนฺติ.\csromanspacer{} เอวํ โข ภิกฺขเว อติธาวนฺติ เอเก.}

\prose{กถญฺจ ภิกฺขเว จกฺขุมนฺโต จ ปสฺสนฺติ? อิธ ภิกฺขเว ภิกฺขุ ภูตํ ภูตโต ปสฺสติ, ภูตํ ภูตโต ทิสฺวา ภูตสฺส นิพฺพิทาย วิราคาย นิโรธาย ปฏิปนฺโน โหติ.\csromanspacer{} เอวํ โข ภิกฺขเว จกฺขุมนฺโต จ ปสฺสนฺติ.}

\csromangathameasure{“โย ภูตํ ภูตโต ทิสฺวา, ภูตสฺส จ อติกฺกมํ. \\ ยถาภูเต’ธิมุจฺจติ, ภวตณฺหา ปริกฺขยา. \\ ส เว ภูตปริญฺญาโต, วีตตโณฺห ภวาภเว. \\ ภูตสฺส วิภวา ภิกฺขุ, นาคจฺฉติ ปุนพฺภวนฺ”ติ.}
\csromangathagroup{\csromangathastanza{“โย\footnote{อยํ คาถา ขุ 1. 225 ปิฏฺเฐ ทิสฺสติ.} ภูตํ ภูตโต ทิสฺวา, ภูตสฺส จ อติกฺกมํ. \\ ยถาภูเต’ธิมุจฺจติ, ภวตณฺหา ปริกฺขยา.}\csromangathastanza{ส เว ภูตปริญฺญาโต, วีตตโณฺห ภวาภเว. \\ ภูตสฺส วิภวา ภิกฺขุ, นาคจฺฉติ ปุนพฺภวนฺ”ติ.}}

\proseitem{150}{ตโย ปุคฺคลา วิปนฺนทิฏฺฐี, ตโย ปุคฺคลา สมฺปนฺนทิฏฺฐี.\csromanspacer{} กตเม ตโย ปุคฺคลา วิปนฺนทิฏฺฐี? ติตฺถิโย จ ติตฺถิยสาวโก จ โย จ มิจฺฉาทิฏฺฐิโก.\csromanspacer{} อิเม ตโย ปุคฺคลา วิปนฺนทิฏฺฐี.}

\prose{กตเม ตโย ปุคฺคลา สมฺปนฺนทิฏฺฐี? ตถาคโต จ ตถาคตสาวโก จ โย จ สมฺมาทิฏฺฐิโก.\csromanspacer{} อิเม ตโย ปุคฺคลา สมฺปนฺนทิฏฺฐี.}

\csromangathameasure{โกธโน อุปนาหี จ, ปาปมกฺขี จ โย นโร. \\ วิปนฺนทิฏฺฐิ มายาวี, ตํ ชญฺญา วสโล อิติ. \\ อกฺโกธโน อนุปนาหี, วิสุทฺโธ สุทฺธตํ คโต. \\ สมฺปนฺนทิฏฺฐิ เมธาวี, ตํ ชญฺญา อริโย อิตีติ.}
\csromangathagroup{\csromangathastanza{\csromanfolio{160}โกธโน อุปนาหี จ, ปาปมกฺขี จ โย นโร. \\ วิปนฺนทิฏฺฐิ มายาวี, ตํ ชญฺญา วสโล อิติ\footnote{อิเม (สฺยา) ขุ 1. 298 ปิฏฺเฐ ปสฺสิตพฺพา.}.}\csromangathastanza{อกฺโกธโน อนุปนาหี, วิสุทฺโธ\footnote{อมกฺขี (สฺยา) อฏฺฐกถา โอโลเกตพฺพา. 3. สกฺกายทิฏฺฐิปฺปมุเขน (สฺยา)} สุทฺธตํ คโต. \\ สมฺปนฺนทิฏฺฐิ เมธาวี, ตํ ชญฺญา อริโย อิตีติ.}}

\prose{ติสฺโส วิปนฺนทิฏฺฐิโย, ติสฺโส สมฺปนฺนทิฏฺฐิโย.\csromanspacer{} กตมา ติสฺโส วิปนฺนทิฏฺฐิโย? \textbf{“เอตํ มมา”}ติ วิปนฺนทิฏฺฐิ, “เอโสหมสฺมี”ติ วิปนฺนทิฏฺฐิ, “เอโส เม อตฺตา”ติ วิปนฺนทิฏฺฐิ.\csromanspacer{} อิมา ติสฺโส วิปนฺนทิฏฺฐิโย.}

\prose{กตมา ติสฺโส สมฺปนฺนทิฏฺฐิโย? “เนตํ มมา”ติ สมฺปนฺนทิฏฺฐิ, “เนโสหมสฺมี”ติ สมฺปนฺนทิฏฺฐิ, “น เมโส อตฺตา”ติ สมฺปนฺนทิฏฺฐิ.\csromanspacer{} อิมา ติสฺโส สมฺปนฺนทิฏฺฐิโย.}

\prose{\textbf{“เอตํ มมา”}ติ กา ทิฏฺฐิ, กติ ทิฏฺฐิโย, กตมนฺตานุคฺคหิตา ตา ทิฏฺฐิโย. \textbf{“เอโสหมสฺมี”}ติ กา ทิฏฺฐิ, กติ ทิฏฺฐิโย, กตมนฺตานุคฺคหิตา ตา ทิฏฺฐิโย. \textbf{“เอโส เม อตฺตา”}ติ กา ทิฏฺฐิ, กติ ทิฏฺฐิโย, กตมนฺตานุคฺคหิตา ตา ทิฏฺฐิโย?}

\prose{\textbf{“เอตํ มมา”}ติ ปุพฺพนฺตานุทิฏฺฐิ, อฏฺฐารส ทิฏฺฐิโย, ปุพฺพนฺตานุคฺคหิตา ตา ทิฏฺฐิโย. \textbf{“เอโสหมสฺมี”}ติ อปรนฺตานุทิฏฺฐิ, จตุจตฺตารีสํ ทิฏฺฐิโย, อปรนฺตานุคฺคหิตา ตา ทิฏฺฐิโย. \textbf{“เอโส เม อตฺตา”}ติ วีสติวตฺถุกา อตฺตานุทิฏฺฐิ วีสติวตฺถุกา สกฺกายทิฏฺฐิ, สกฺกายทิฏฺฐิปฺปมุขานิ3 ทฺวาสฏฺฐิ ทิฏฺฐิคตานิ, ปุพฺพนฺตาปรนฺตานุคฺคหิตา ตา ทิฏฺฐิโย.}

\proseitem{151}{เย เกจิ ภิกฺขเว มยิ นิฏฺฐํ คตา, สพฺเพ เต ทิฏฺฐิสมฺปนฺนา.\csromanspacer{} เตสํ ทิฏฺฐิสมฺปนฺนานํ ปญฺจนฺนํ อิธ นิฏฺฐา, ปญฺจนฺนํ อิธ วิหาย นิฏฺฐา.\csromanspacer{} กตเมสํ ปญฺจนฺนํ อิธ นิฏฺฐา? สตฺตกฺขตฺตุปรมสฺส, โกลํโกลสฺส, เอกพีชิสฺส สกทาคามิสฺส โย จ ทิฏฺเฐว ธมฺเม อรหา.\csromanspacer{} อิเมสํ ปญฺจนฺนํ อิธ นิฏฺฐา.}

\prose{กตเมสํ ปญฺจนฺนํ อิธ วิหาย นิฏฺฐา? อนฺตราปรินิพฺพายิสฺส, อุปหจฺจปรินิพฺพายิสฺส, อสงฺขารปรินิพฺพายิสฺส, สสงฺขารปรินิพฺพายิสฺส, อุทฺธํโสตสฺส อกนิฏฺฐคามิโน.\csromanspacer{} อิเมสํ ปญฺจนฺนํ อิธ วิหาย นิฏฺฐา.}

\csromanmatikaanchor{mtk.65}
\csromantocmarknopage{section}{3. อานาปานสฺสติกถา}{mtk.65}
\bookmark[dest={mtk.65},level=1]{3. อานาปานสฺสติกถา}
\csromanheader{1}{3. อานาปานสฺสติกถา}
\prose{\csromanfolio{161}เย เกจิ ภิกฺขเว มยิ นิฏฺฐํ คตา, สพฺเพ เต ทิฏฺฐิสมฺปนฺนา, เตสํ ทิฏฺฐิสมฺปนฺนานํ อิเมสํ ปญฺจนฺนํ อิธ นิฏฺฐา, อิเมสํ ปญฺจนฺนํ อิธ วิหาย นิฏฺฐา.}

\prose{เย เกจิ ภิกฺขเว มยิ อเวจฺจปฺปสนฺนา, สพฺเพ เต โสตาปนฺนา, เตสํ โสตาปนฺนานํ ปญฺจนฺนํ อิธ นิฏฺฐา, ปญฺจนฺนํ อิธ วิหาย นิฏฺฐา, กตเมสํ ปญฺจนฺนํ อิธ นิฏฺฐา? สตฺตกฺขตฺตุปรมสฺส โกลํโกลสฺส เอกพีชิสฺส สกทาคามิสฺส โย จ ทิฏฺเฐว ธมฺเม อรหา, อิเมสํ ปญฺจนฺนํ อิธ นิฏฺฐา.}

\prose{กตเมสํ ปญฺจนฺนํ อิธ วิหาย นิฏฺฐา? อนฺตราปรินิพฺพายิสฺส อุปหจฺจปรินิพฺพายิสฺส อสงฺขารปรินิพฺพายิสฺส สสงฺขารปรินิพฺพายิสฺส อุทฺธํโสตสฺส อกนิฏฺฐคามิโน, อิเมสํ ปญฺจนฺนํ อิธ วิหาย นิฏฺฐา.}

\prose{เย เกจิ ภิกฺขเว มยิ อเวจฺจปฺปสนฺนา, สพฺเพ เต โสตาปนฺนา, เตสํ โสตาปนฺนานํ อิเมสํ ปญฺจนฺนํ อิธ นิฏฺฐา, อิเมสํ ปญฺจนฺนํ อิธ วิหาย นิฏฺฐาติ.}

\vspace{\tipitakaparskip}
\csromancenter{ภววิภวทิฏฺฐินิทฺเทโส โสฬสโม.}
\nitthitamruled{ทิฏฺฐิกถา นิฏฺฐิตา.}
\csromanheader{2}{3. อานาปานสฺสติกถา}
\csromanmatikaanchor{mtk.66}
\csromantocmarkref{subsection}{1. คณนวาร}{mtk.66}
\bookmark[dest={mtk.66},level=2]{1. คณนวาร}
\csromanheader{2}{1. คณนวาร}
\proseitem{152}{โสฬสวตฺถุกํ อานาปานสฺสติสมาธิํ\footnote{อานาปานสติสมาธิํ (สี-ฏฺฐ)} ภาวยโต สมธิกานิ เทฺว ญาณสตานิ อุปฺปชฺชนฺติ, อฏฺฐ ปริปนฺเถ\footnote{ปริพนฺเธ (ก)} ญาณานิ, อฏฺฐ จ อุปกาเร ญาณานิ, อฏฺฐารส อุปกฺกิเลเส ญาณานิ, เตรส โวทาเน ญาณานิ, พาตฺติํส สโตการิสฺส\footnote{สโตการีสุ (สฺยา)} ญาณานิ, จตุวีสติ สมาธิวเสน ญาณานิ, เทฺวสตฺตติ วิปสฺสนาวเสน ญาณานิ, อฏฺฐ นิพฺพิทาญาณานิ, อฏฺฐ นิพฺพิทานุโลมญาณานิ, อฏฺฐ นิพฺพิทาปฏิปฺปสฺสทฺธิญาณานิ, เอกวีสติ วิมุตฺติสุเข ญาณานิ.}

\prose{กตมานิ อฏฺฐ ปริปนฺเถ ญาณานิ, อฏฺฐ จ อุปกาเร ญาณานิ? กามจฺฉนฺโท สมาธิสฺส ปริปนฺโถ, เนกฺขมฺมํ สมาธิสฺส อุปการํ.\csromanspacer{} พฺยาปาโท \csromanfolio{162}สมาธิสฺส ปริปนฺโถ, อพฺยาปาโท สมาธิสฺส อุปการํ.\csromanspacer{} ถินมิทฺธํ สมาธิสฺส ปริปนฺโถ, อาโลกสญฺญา สมาธิสฺส อุปการํ.\csromanspacer{} อุทฺธจฺจํ สมาธิสฺส ปริปนฺโถ, อวิกฺเขโป สมาธิสฺส อุปการํ.\csromanspacer{} วิจิกิจฺฉา สมาธิสฺส ปริปนฺโถ, ธมฺมววตฺถานํ สมาธิสฺส อุปการํ.\csromanspacer{} อวิชฺชา สมาธิสฺส ปริปนฺโถ, ญาณํ สมาธิสฺส อุปการํ.\csromanspacer{} อรติ สมาธิสฺส ปริปนฺโถ, ปาโมชฺชํ สมาธิสฺส อุปการํ.\csromanspacer{} สพฺเพปิ อกุสลา ธมฺมา สมาธิสฺส ปริปนฺถา, สพฺเพปิ กุสลา ธมฺมา สมาธิสฺส อุปการา.\csromanspacer{} อิมานิ อฏฺฐ ปริปนฺเถ ญาณานิ, อฏฺฐ จ อุปกาเร ญาณานิ.}

\csromanheader{2}{คณนวาโร ปฐโม.}
\csromansectionrule
\csromanmatikaanchor{mtk.67}
\csromantocmarkref{subsection}{2. โสฬสญาณนิทฺเทส}{mtk.67}
\bookmark[dest={mtk.67},level=2]{2. โสฬสญาณนิทฺเทส}
\csromanheader{2}{2. โสฬสญาณนิทฺเทส}
\proseitem{153}{อิเมหิ โสฬสหิ อากาเรหิ อุทุจิตํ จิตฺตํ สมุทุจิตํ จิตฺตํ เอกตฺเต สนฺติฏฺฐติ, นีวรเณหิ วิสุชฺฌติ.\csromanspacer{} กตเม เต เอกตฺตา? เนกฺขมฺมํ เอกตฺตํ, อพฺยาปาโท เอกตฺตํ, อาโลกสญฺญา เอกตฺตํ, อวิกฺเขโป เอกตฺตํ, ธมฺมววตฺถานํ เอกตฺตํ, ญาณํ เอกตฺตํ, ปาโมชฺชํ เอกตฺตํ, สพฺเพปิ กุสลา ธมฺมา เอกตฺตา.}

\prose{\textbf{นีวรณา}ติ กตเม เต นีวรณา? กามจฺฉนฺโท นีวรณํ, พฺยาปาโท นีวรณํ, ถินมิทฺธํ นีวรณํ, อุทฺธจฺจกุกฺกุจฺจํ นีวรณํ, วิจิกิจฺฉา นีวรณํ, อวิชฺชา นีวรณํ, อรติ นีวรณํ, สพฺเพปิ อกุสลา ธมฺมา นีวรณา.}

\prose{\textbf{นีวรณา}ติ เกนฏฺเฐน นีวรณา? นิยฺยานาวรณฏฺเฐน นีวรณา.\csromanspacer{} กตเม เต นิยฺยานา? เนกฺขมฺมํ อริยานํ นิยฺยานํ, เตน จ เนกฺขมฺเมน อริยา นิยฺยนฺติ, กามจฺฉนฺโท นิยฺยานาวรณํ, เตน จ กามจฺฉนฺเทน นิวุตตฺตา เนกฺขมฺมํ อริยานํ นิยฺยานํ นปฺปชานาตีติ กามจฺฉนฺโท นิยฺยานาวรณํ.\csromanspacer{}\csromanspacer{}อพฺยาปาโท อริยานํ นิยฺยานํ, เตน จ อพฺยาปาเทน อริยา นิยฺยนฺติ, พฺยาปาโท นิยฺยานาวรณํ, เตน จ พฺยาปาเทน นิวุตตฺตา อพฺยาปาทํ อริยานํ นิยฺยานํ นปฺปชานาตีติ พฺยาปาโท นิยฺยานาวรณํ.\csromanspacer{}\csromanspacer{}อาโลกสญฺญา อริยานํ นิยฺยานํ, ตาย จ อาโลกสญฺญาย อริยา นิยฺยนฺติ, ถินมิทฺธํ นิยฺยานาวรณํ, เตน จ ถินมิทฺเธน นิวุตตฺตา อาโลกสญฺญํ อริยานํ นิยฺยานํ นปฺปชานาตีติ ถินมิทฺธํ นิยฺยานาวรณํ.\csromanspacer{}\csromanspacer{}อวิกฺเขโป อริยานํ นิยฺยานํ, เตน จ อวิกฺเขเปน อริยา นิยฺยนฺติ, อุทฺธจฺจํ}

\vspace{\tipitakaparskip}
\csromancenter{อานาปานสฺสติกถา นิยฺยานาวรณํ, เตน จ อุทฺธจฺเจน นิวุตตฺตา อวิกฺเขปํ อริยานํ นิยฺยานํ นปฺปชานาตีติ อุทฺธจฺจํ นิยฺยานาวรณํ.\csromanspacer{}\csromanspacer{}ธมฺมววตฺถานํ อริยานํ นิยฺยานํ, เตน จ ธมฺมววตฺถาเนน อริยา นิยฺยนฺติ, วิจิกิจฺฉา นิยฺยานาวรณํ, ตาย จ วิจิกิจฺฉาย นิวุตตฺตา ธมฺมววตฺถานํ อริยานํ นิยฺยานํ นปฺปชานาตีติ วิจิกิจฺฉา นิยฺยานาวรณํ.\csromanspacer{}\csromanspacer{}ญาณํ อริยานํ นิยฺยานํ, เตน จ ญาเณน อริยา นิยฺยนฺติ, อวิชฺชา นิยฺยานาวรณํ, ตาย จ อวิชฺชาย นิวุตตฺตา ญาณํ อริยานํ นิยฺยานํ นปฺปชานาตีติ อวิชฺชา นิยฺยานาวรณํ.\csromanspacer{}\csromanspacer{}ปาโมชฺชํ อริยานํ นิยฺยานํ, เตน จ ปาโมชฺเชน อริยา นิยฺยนฺติ, อรติ นิยฺยานาวรณํ, ตาย จ อรติยา นิวุตตฺตา ปาโมชฺชํ อริยานํ นิยฺยานํ นปฺปชานาตีติ อรติ นิยฺยานาวรณํ.\csromanspacer{}\csromanspacer{}สพฺเพปิ กุสลา ธมฺมา อริยานํ นิยฺยานํ, เตหิ จ กุสเลหิ ธมฺเมหิ อริยา นิยฺยนฺติ.\csromanspacer{} สพฺเพปิ อกุสลา ธมฺมา นิยฺยานาวรณา, เตหิ จ อกุสเลหิ ธมฺเมหิ นิวุตตฺตา กุสเล ธมฺเม อริยานํ นิยฺยานํ นปฺปชานาตีติ สพฺเพปิ อกุสลา ธมฺมา นิยฺยานาวรณา.}
\csromanheader{2}{โสฬสญาณนิทฺเทโส ทุติโย.}
\csromansectionrule
\csromanmatikaanchor{mtk.68}
\csromantocmarkref{subsection}{3. อุปกฺกิเลสญาณนิทฺเทส}{mtk.68}
\bookmark[dest={mtk.68},level=2]{3. อุปกฺกิเลสญาณนิทฺเทส}
\csromanheader{2}{3. อุปกฺกิเลสญาณนิทฺเทส}
\csromanheader{2}{\textbf{ปฐมจฺฉกฺก}}
\proseitem{154}{\csromanfolio{163}อิเมหิ จ ปน นีวรเณหิ วิสุทฺธจิตฺตสฺส โสฬสวตฺถุกํ อานาปานสฺสติสมาธิํ ภาวยโต ขณิกสโมธานา กตเม อฏฺฐารส อุปกฺกิเลสา อุปฺปชฺชนฺติ? อสฺสาสาทิมชฺฌปริโยสานํ สติยา อนุคจฺฉโต อชฺฌตฺตวิกฺเขปคตํ จิตฺตํ สมาธิสฺส ปริปนฺโถ, ปสฺสาสาทิมชฺฌปริโยสานํ สติยา อนุคจฺฉโต พหิทฺธาวิกฺเขปคตํ จิตฺตํ สมาธิสฺส ปริปนฺโถ, อสฺสาสปฏิกงฺขนา นิกนฺติตณฺหาจริยา สมาธิสฺส ปริปนฺโถ, ปสฺสาสปฏิกงฺขนา นิกนฺติตณฺหาจริยา สมาธิสฺส ปริปนฺโถ, อสฺสาเสนาภิตุนฺนสฺส ปสฺสาสปฏิลาเภ มุจฺฉนา สมาธิสฺส ปริปนฺโถ, ปสฺสาเสนาภิตุนฺนสฺส อสฺสาสปฏิลาเภ มุจฺฉนา สมาธิสฺส ปริปนฺโถ.}

\csromangathameasure{อนุคจฺฉนา จ อสฺสาสํ, ปสฺสาสํ, อนุคจฺฉนา.}
\csromangathagroup{\csromangathastanza{\csromanfolio{164}อนุคจฺฉนา จ อสฺสาสํ, ปสฺสาสํ, อนุคจฺฉนา.}}

\prose{สติ อชฺฌตฺตวิกฺเขปา-กงฺขนา พหิทฺธาวิกฺเขปปตฺถนา\footnote{วิกฺเขปปนฺถนา (สฺยา)}.}

\csromangathameasure{อสฺสาเสนาภิตุนฺนสฺส, ปสฺสาสปฏิลาเภ มุจฺฉนา. \\ ปสฺสาเสนาภิตุนฺนสฺส, อสฺสาสปฏิลาเภ มุจฺฉนา. \\ ฉ เอเต อุปกฺกิเลสา, อานาปานสฺสติสมาธิสฺส. \\ เยหิ วิกฺขิปฺปมานสฺส, โน จ จิตฺตํ วิมุจฺจติ. \\ วิโมกฺขํ อปฺปชานนฺตา, เต โหนฺติ ปรปตฺติยาติ.}
\csromangathagroup{\csromangathastanza{อสฺสาเสนาภิตุนฺนสฺส, ปสฺสาสปฏิลาเภ มุจฺฉนา. \\ ปสฺสาเสนาภิตุนฺนสฺส, อสฺสาสปฏิลาเภ มุจฺฉนา.}\csromangathastanza{ฉ เอเต อุปกฺกิเลสา, อานาปานสฺสติสมาธิสฺส. \\ เยหิ วิกฺขิปฺปมานสฺส\footnote{วิกมฺปมานสฺส (สฺยา)}, โน จ จิตฺตํ วิมุจฺจติ. \\ วิโมกฺขํ อปฺปชานนฺตา, เต โหนฺติ ปรปตฺติยาติ.}}

\csromanheader{2}{\textbf{ทุติยจฺฉกฺก}}
\proseitem{155}{นิมิตฺตํ อาวชฺชโต อสฺสาเส จิตฺตํ วิกมฺปติ, สมาธิสฺส ปริปนฺโถ.\csromanspacer{} อสฺสาสํ อาวชฺชโต นิมิตฺเต จิตฺตํ วิกมฺปติ, สมาธิสฺส ปริปนฺโถ.\csromanspacer{} นิมิตฺตํ อาวชฺชโต ปสฺสาเส จิตฺตํ วิกมฺปติ, สมาธิสฺส ปริปนฺโถ.\csromanspacer{} ปสฺสาสํ อาวชฺชโต นิมิตฺเต จิตฺตํ วิกมฺปติ, สมาธิสฺส ปริปนฺโถ.\csromanspacer{} อสฺสาสํ อาวชฺชโต ปสฺสาเส จิตฺตํ วิกมฺปติ, สมาธิสฺส ปริปนฺโถ.\csromanspacer{} ปสฺสาสํ อาวชฺชโต อสฺสาเส จิตฺตํ วิกมฺปติ, สมาธิสฺส ปริปนฺโถ.}

\csromangathameasure{นิมิตฺตํ อาวชฺชมานสฺส, อสฺสาเส วิกฺขิปฺปเต มโน. \\ อสฺสาสํ อาวชฺชมานสฺส, นิมิตฺเต จิตฺตํ วิกมฺปติ. \\ นิมิตฺตํ อาวชฺชมานสฺส, ปสฺสาเส วิกฺขิปฺปเต มโน. \\ ปสฺสาสํ อาวชฺชมานสฺส, นิมิตฺเต จิตฺตํ วิกมฺปติ. \\ อสฺสาสํ อาวชฺชมานสฺส, ปสฺสาเส วิกฺขิปฺปเต มโน. \\ ปสฺสาสํ อาวชฺชมานสฺส, อสฺสาเส จิตฺตํ วิกมฺปติ. \\ ฉ เอเต อุปกฺกิเลสา, อานาปานสฺสติสมาธิสฺส. \\ เยหิ วิกฺขิปฺปมานสฺส, โน จ จิตฺตํ วิมุจฺจติ. \\ วิโมกฺขํ อปฺปชานนฺตา, เต โหนฺติ ปรปตฺติยาติ.}
\csromangathagroup{\csromangathastanza{นิมิตฺตํ อาวชฺชมานสฺส, อสฺสาเส วิกฺขิปฺปเต มโน. \\ อสฺสาสํ อาวชฺชมานสฺส, นิมิตฺเต จิตฺตํ วิกมฺปติ.}\csromangathastanza{นิมิตฺตํ อาวชฺชมานสฺส, ปสฺสาเส วิกฺขิปฺปเต มโน. \\ ปสฺสาสํ อาวชฺชมานสฺส, นิมิตฺเต จิตฺตํ วิกมฺปติ.}\csromangathastanza{อสฺสาสํ อาวชฺชมานสฺส, ปสฺสาเส วิกฺขิปฺปเต มโน. \\ ปสฺสาสํ อาวชฺชมานสฺส, อสฺสาเส จิตฺตํ วิกมฺปติ.}\csromangathastanza{ฉ เอเต อุปกฺกิเลสา, อานาปานสฺสติสมาธิสฺส. \\ เยหิ วิกฺขิปฺปมานสฺส, โน จ จิตฺตํ วิมุจฺจติ. \\ วิโมกฺขํ อปฺปชานนฺตา, เต โหนฺติ ปรปตฺติยาติ.}}

\csromanheader{2}{\textbf{ตติยจฺฉกฺก}}
\proseitem{156}{อตีตานุธาวนํ จิตฺตํ วิกฺเขปานุปติตํ สมาธิสฺส ปริปนฺโถ.\csromanspacer{} อนาคตปฏิกงฺขนํ จิตฺตํ วิกมฺปิตํ สมาธิสฺส ปริปนฺโถ.\csromanspacer{} ลีนํ จิตฺตํ}

\vspace{\tipitakaparskip}
\csromancenter{อานาปานสฺสติกถา โกสชฺชานุปติตํ สมาธิสฺส ปริปนฺโถ.\csromanspacer{} อติปคฺคหิตํ จิตฺตํ อุทฺธจฺจานุปติตํ สมาธิสฺส ปริปนฺโถ.\csromanspacer{} อภินตํ จิตฺตํ ราคานุปติตํ สมาธิสฺส ปริปนฺโถ.\csromanspacer{} อปนตํ จิตฺตํ พฺยาปาทานุปติตํ สมาธิสฺส ปริปนฺโถ.}
\vspace{0.5\baselineskip}
\csromangathameasure{อตีตานุธาวนํ จิตฺตํ, อนาคตปฏิกงฺขนํ ลีนํ. \\ อติปคฺคหิตํ อภินตํ, อปนตํ จิตฺตํ น สมาธิยติ. \\ ฉ เอเต อุปกฺกิเลสา, อานาปานสฺสติสมาธิสฺส. \\ เยหิ อุปกฺกิลิฏฺฐสงฺกปฺโป, อธิจิตฺตํ นปฺปชานาตีติ.}
\csromangathagroup{\csromangathastanza{\csromanfolio{165}อตีตานุธาวนํ จิตฺตํ, อนาคตปฏิกงฺขนํ ลีนํ. \\ อติปคฺคหิตํ อภินตํ, อปนตํ จิตฺตํ น สมาธิยติ.}\csromangathastanza{ฉ เอเต อุปกฺกิเลสา, อานาปานสฺสติสมาธิสฺส. \\ เยหิ อุปกฺกิลิฏฺฐสงฺกปฺโป, อธิจิตฺตํ นปฺปชานาตีติ.}}

\proseitem{157}{อสฺสาสาทิมชฺฌปริโยสานํ สติยา อนุคจฺฉโต อชฺฌตฺตํ วิกฺเขปคเตน จิตฺเตน กาโยปิ จิตฺตมฺปิ สารทฺธา จ โหนฺติ อิญฺชิตา จ ผนฺทิตา จ, ปสฺสาสาทิมชฺฌปริโยสานํ สติยา อนุคจฺฉโต พหิทฺธาวิกฺเขปคเตน จิตฺเตน กาโยปิ จิตฺตมฺปิ สารทฺธา จ โหนฺติ อิญฺชิตา จ ผนฺทิตา จ, อสฺสาสปฏิกงฺขนาย นิกนฺติยา ตณฺหาจริยาย กาโยปิ จิตฺตมฺปิ สารทฺธา จ โหนฺติ อิญฺชิตา จ ผนฺทิตา จ, ปสฺสาสปฏิกงฺขนาย นิกนฺติยา ตณฺหาจริยาย กาโยปิ จิตฺตมฺปิ สารทฺธา จ โหนฺติ อิญฺชิตา จ ผนฺทิตา จ, อสฺสาเสนาภิตุนฺนสฺส ปสฺสาสปฏิลาเภ มุจฺฉิตตฺตา กาโยปิ จิตฺตมฺปิ สารทฺธา จ โหนฺติ อิญฺชิตา จ ผนฺทิตา จ, ปสฺสาเสนาภิตุนฺนสฺส อสฺสาสปฏิลาเภ มุจฺฉิตตฺตา กาโยปิ จิตฺตมฺปิ สารทฺธา จ โหนฺติ อิญฺชิตา จ ผนฺทิตา จ.}

\prose{นิมิตฺตํ อาวชฺชโต อสฺสาเส จิตฺตํ วิกมฺปิตตฺตา กาโยปิ จิตฺตมฺปิ สารทฺธา จ โหนฺติ อิญฺชิตา จ ผนฺทิตา จ, อสฺสาสํ อาวชฺชโต นิมิตฺเต จิตฺตํ วิกมฺปิตตฺตา กาโยปิ จิตฺตมฺปิ สารทฺธา จ โหนฺติ อิญฺชิตา จ ผนฺทิตา จ.\csromanspacer{} นิมิตฺตํ อาวชฺชโต ปสฺสาเส จิตฺตํ วิกมฺปิตตฺตา กาโยปิ จิตฺตมฺปิ สารทฺธา จ โหนฺติ อิญฺชิตา จ ผนฺทิตา จ, ปสฺสาสํ อาวชฺชโต นิมิตฺเต จิตฺตํ วิกมฺปิตตฺตา กาโยปิ จิตฺตมฺปิ สารทฺธา จ โหนฺติ อิญฺชิตา จ ผนฺทิตา จ.\csromanspacer{} อสฺสาสํ อาวชฺชโต ปสฺสาเส จิตฺตํ วิกมฺปิตตฺตา กาโยปิ จิตฺตมฺปิ สารทฺธา จ โหนฺติ อิญฺชิตา จ ผนฺทิตา จ, ปสฺสาสํ อาวชฺชโต อสฺสาเส จิตฺตํ วิกมฺปิตตฺตา กาโยปิ จิตฺตมฺปิ สารทฺธา จ โหนฺติ อิญฺชิตา จ ผนฺทิตา จ.}

\prose{อตีตานุธาวเนน จิตฺเตน วิกฺเขปานุปติเตน กาโยปิ จิตฺตมฺปิ สารทฺธา จ โหนฺติ อิญฺชิตา จ ผนฺทิตา จ, อนาคตปฏิกงฺขเนน จิตฺเตน วิกมฺปิเตน กาโยปิ จิตฺตมฺปิ สารทฺธา จ โหนฺติ อิญฺชิตา จ ผนฺทิตา จ.\csromanspacer{} ลีเนน \csromanfolio{166}จิตฺเตน โกสชฺชานุปติเตน กาโยปิ จิตฺตมฺปิ สารทฺธา จ โหนฺติ อิญฺชิตา จ ผนฺทิตา จ, อติปคฺคหิเตน จิตฺเตน อุทฺธจฺจานุปติเตน กาโยปิ จิตฺตมฺปิ สารทฺธา จ โหนฺติ อิญฺชิตา จ ผนฺทิตา จ.\csromanspacer{} อภินเตน จิตฺเตน ราคานุปติเตน กาโยปิ จิตฺตมฺปิ สารทฺธา จ โหนฺติ อิญฺชิตา จ ผนฺทิตา จ, อปนเตน จิตฺเตน พฺยาปาทานุปติเตน กาโยปิ จิตฺตมฺปิ สารทฺธา จ โหนฺติ อิญฺชิตา จ ผนฺทิตา จ.}

\csromangathameasure{อานาปานสฺสติ ยสฺส, ปริปุณฺณา อภาวิตา. \\ กาโยปิ อิญฺชิโต โหติ, จิตฺตมฺปิ โหติ อิญฺชิตํ. \\ กาโยปิ ผนฺทิโต โหติ, จิตฺตมฺปิ โหติ ผนฺทิตํ. \\ อานาปานสฺสติ ยสฺส, ปริปุณฺณา สุภาวิตา. \\ กาโยปิ อนิญฺชิโต โหติ, จิตฺตมฺปิ โหติ อนิญฺชิตํ. \\ กาโยปิ อผนฺทิโต โหติ, จิตฺตมฺปิ โหติ อผนฺทิตนฺติ.}
\csromangathagroup{\csromangathastanza{อานาปานสฺสติ ยสฺส, ปริปุณฺณา อภาวิตา. \\ กาโยปิ อิญฺชิโต โหติ, จิตฺตมฺปิ โหติ อิญฺชิตํ.}\csromangathastanza{กาโยปิ ผนฺทิโต โหติ, จิตฺตมฺปิ โหติ ผนฺทิตํ. \\ อานาปานสฺสติ ยสฺส, ปริปุณฺณา สุภาวิตา.}\csromangathastanza{กาโยปิ อนิญฺชิโต โหติ, จิตฺตมฺปิ โหติ อนิญฺชิตํ. \\ กาโยปิ อผนฺทิโต โหติ, จิตฺตมฺปิ โหติ อผนฺทิตนฺติ.}}

\prose{อิเมหิ จ ปน นีวรเณหิ วิสุทฺธจิตฺตสฺส โสฬสวตฺถุกํ อานาปานสฺสติสมาธิํ ภาวยโต ขณิกสโมธานา อิเม อฏฺฐารส อุปกฺกิเลสา อุปฺปชฺชนฺติ.}

\csromanheader{2}{อุปกฺกิเลสญาณนิทฺเทโส ตติโย.}
\csromansectionrule
\csromanmatikaanchor{mtk.69}
\csromantocmarkref{subsection}{4. โวทานญาณนิทฺเทส}{mtk.69}
\bookmark[dest={mtk.69},level=2]{4. โวทานญาณนิทฺเทส}
\csromanheader{2}{4. โวทานญาณนิทฺเทส}
\proseitem{158}{กตมานิ เตรส โวทาเน ญาณานิ? อตีตานุธาวนํ จิตฺตํ วิกฺเขปานุปติตํ, ตํ วิวชฺชยิตฺวา เอกฏฺฐาเน สมาทหติ, เอวมฺปิ จิตฺตํ น วิกฺเขปํ คจฺฉติ.\csromanspacer{} อนาคตปฏิกงฺขนํ จิตฺตํ วิกมฺปิตํ, ตํ วิวชฺชยิตฺวา ตตฺเถว อธิโมเจติ, เอวมฺปิ จิตฺตํ น วิกฺเขปํ คจฺฉติ.\csromanspacer{} ลีนํ จิตฺตํ โกสชฺชานุปติตํ, ตํ ปคฺคณฺหิตฺวา โกสชฺชํ ปชหติ, เอวมฺปิ จิตฺตํ น วิกฺเขปํ คจฺฉติ.\csromanspacer{} อติปคฺคหิตํ จิตฺตํ อุทฺธจฺจานุปติตํ, ตํ วินิคฺคณฺหิตฺวา อุทฺธจฺจํ ปชหติ, เอวมฺปิ จิตฺตํ น วิกฺเขปํ คจฺฉติ.\csromanspacer{} อภินตํ จิตฺตํ ราคานุปติตํ, ตํ สมฺปชาโน หุตฺวา ราคํ ปชหติ, เอวมฺปิ จิตฺตํ น วิกฺเขปํ คจฺฉติ.\csromanspacer{} อปนตํ จิตฺตํ พฺยาปาทานุปติตํ, ตํ สมฺปชาโน หุตฺวา พฺยาปาทํ ปชหติ, เอวมฺปิ จิตฺตํ น วิกฺเขปํ คจฺฉติ.\csromanspacer{} อิเมหิ ฉหิ ฐาเนหิ ปริสุทฺธํ จิตฺตํ ปริโยทาตํ เอกตฺตคตํ โหติ.}

\csromanheader{2}{3. อานาปานสฺสติกถา}
\prose{\csromanfolio{167}กตเม เต เอกตฺตา? ทานโวสคฺคุปฏฺฐาเนกตฺตํ, สมถนิมิตฺตุปฏฺฐาเนกตฺตํ, วยลกฺขณุปฏฺฐาเนกตฺตํ, นิโรธุปฏฺฐาเนกตฺตํ.\csromanspacer{} ทานโวสคฺคุปฏฺฐาเนกตฺตํ จาคาธิมุตฺตานํ.\csromanspacer{} สมถนิมิตฺตุปฏฺฐาเนกตฺตญฺจ อธิจิตฺตมนุยุตฺตานํ.\csromanspacer{} วยลกฺขณุปฏฺฐาเนกตฺตญฺจ วิปสฺสกานํ.\csromanspacer{} นิโรธุปฏฺฐาเนกตฺตญฺจ อริยปุคฺคลานํ.\csromanspacer{} อิเมหิ จตูหิ ฐาเนหิ เอกตฺตคตํ จิตฺตํ ปฏิปทาวิสุทฺธิปกฺขนฺทญฺเจว โหติ, อุเปกฺขานุพฺรูหิตญฺจ ญาเณน จ สมฺปหํสิตํ.}

\prose{ปฐมสฺส ฌานสฺส โก อาทิ, กิํ มชฺเฌ, กิํ ปริโยสานํ? ปฐมสฺส ฌานสฺส ปฏิปทาวิสุทฺธิ อาทิ, อุเปกฺขานุพฺรูหนา มชฺเฌ, สมฺปหํสนา ปริโยสานํ.\csromanspacer{}\csromanspacer{}ปฐมสฺส ฌานสฺส ปฏิปทาวิสุทฺธิ อาทิ.\csromanspacer{} อาทิสฺส กติ ลกฺขณานิ? อาทิสฺส ตีณิ ลกฺขณานิ, โย ตสฺส ปริปนฺโถ, ตโต จิตฺตํ วิสุชฺฌติ, วิสุทฺธตฺตา จิตฺตํ มชฺฌิมํ สมถนิมิตฺตํ ปฏิปชฺชติ, ปฏิปนฺนตฺตา ตตฺถ จิตฺตํ ปกฺขนฺทติ, ยญฺจ ปริปนฺถโต จิตฺตํ วิสุชฺฌติ, ยญฺจ วิสุทฺธตฺตา จิตฺตํ มชฺฌิมํ สมถนิมิตฺตํ ปฏิปชฺชติ, ยญฺจ ปฏิปนฺนตฺตา ตตฺถ จิตฺตํ ปกฺขนฺทติ, ปฐมสฺส ฌานสฺส ปฏิปทาวิสุทฺธิ อาทิ, อาทิสฺส อิมานิ ตีณิ ลกฺขณานิ, เตน วุจฺจติ ปฐมํ ฌานํ อาทิกลฺยาณญฺเจว โหติ ลกฺขณสมฺปนฺนญฺจ.}

\prose{ปฐมสฺส ฌานสฺส อุเปกฺขานุพฺรูหนา มชฺเฌ.\csromanspacer{} มชฺฌสฺส กติ ลกฺขณานิ? มชฺฌสฺส ตีณิ ลกฺขณานิ, วิสุทฺธํ จิตฺตํ อชฺฌุเปกฺขติ, สมถปฏิปนฺนํ อชฺฌุเปกฺขติ, เอกตฺตุปฏฺฐานํ อชฺฌุเปกฺขติ, ยญฺจ วิสุทฺธํ จิตฺตํ อชฺฌุเปกฺขติ, ยญฺจ สมถปฏิปนฺนํ อชฺฌุเปกฺขติ, ยญฺจ เอกตฺตุปฏฺฐานํ อชฺฌุเปกฺขติ, ปฐมสฺส ฌานสฺส อุเปกฺขานุพฺรูหนา มชฺเฌ, มชฺฌสฺส อิมานิ ตีณิ ลกฺขณานิ, เตน วุจฺจติ ปฐมํ ฌานํ มชฺเฌกลฺยาณญฺเจว โหติ ลกฺขณสมฺปนฺนญฺจ.}

\prose{ปฐมสฺส ฌานสฺส สมฺปหํสนา ปริโยสานํ.\csromanspacer{} ปริโยสานสฺส กติ ลกฺขณานิ? ปริโยสานสฺส จตฺตาริ ลกฺขณานิ, ตตฺถ ชาตานํ ธมฺมานํ อนติวตฺตนฏฺเฐน สมฺปหํสนา, อินฺทฺริยานํ เอกรสฏฺเฐน สมฺปหํสนา, ตทุปควีริยวาหนฏฺเฐน สมฺปหํสนา, อาเสวนฏฺเฐน สมฺปหํสนา, ปฐมสฺส ฌานสฺส สมฺปหํสนา ปริโยสานํ, ปริโยสานสฺส อิมานิ จตฺตาริ ลกฺขณานิ, เตน วุจฺจติ ปฐมํ ฌานํ ปริโยสานกลฺยาณญฺเจว โหติ ลกฺขณสมฺปนฺนญฺจ, เอวํ ติวตฺตคตํ จิตฺตํ ติวิธกลฺยาณกํ ทสลกฺขณสมฺปนฺนํ วิตกฺกสมฺปนฺนญฺเจว โหติ วิจารสมฺปนฺนญฺจ ปีติสมฺปนฺนญฺจ สุขสมฺปนฺนญฺจ จิตฺตสฺส อธิฏฺฐานสมฺปนฺนญฺจ สทฺธาสมฺปนฺนญฺจ วีริยสมฺปนฺนญฺจ สติสมฺปนฺนญฺจ สมาธิสมฺปนฺนญฺจ ปญฺญาสมฺปนฺนญฺจ.}

\prose{\csromanfolio{168}ทุติยสฺส ฌานสฺส โก อาทิ, กิํ มชฺเฌ, กิํ ปริโยสานํ? ทุติยสฺส ฌานสฺส ปฏิปทาวิสุทฺธิ อาทิ, อุเปกฺขานุพฺรูหนา มชฺเฌ, สมฺปหํสนา ปริโยสานํ ฯเปฯ เอวํ ติวตฺตคตํ จิตฺตํ ติวิธกลฺยาณกํ ทสลกฺขณสมฺปนฺนํ ปีติสมฺปนฺนญฺเจว โหติ สุขสมฺปนฺนญฺจ จิตฺตสฺส อธิฏฺฐานสมฺปนฺนญฺจ สทฺธาสมฺปนฺนญฺจ วีริยสมฺปนฺนญฺจ สติสมฺปนฺนญฺจ สมาธิสมฺปนฺนญฺจ ปญฺญาสมฺปนฺนญฺจ.}

\prose{ตติยสฺส ฌานสฺส โก อาทิ, กิํ มชฺเฌ, กิํ ปริโยสานํ ฯเปฯ เอวํ ติวตฺตคตํ จิตฺตํ ติวิธกลฺยาณกํ ทสลกฺขณสมฺปนฺนํ สุขสมฺปนฺนญฺเจว โหติ จิตฺตสฺส อธิฏฺฐานสมฺปนฺนญฺจ สทฺธาสมฺปนฺนญฺจ วีริยสมฺปนฺนญฺจ สติสมฺปนฺนญฺจ สมาธิสมฺปนฺนญฺจ ปญฺญาสมฺปนฺนญฺจ.}

\prose{จตุตฺถสฺส ฌานสฺส โก อาทิ, กิํ มชฺเฌ, กิํ ปริโยสานํ ฯเปฯ เอวํ ติวตฺตคตํ จิตฺตํ ติวิธกลฺยาณกํ ทสลกฺขณสมฺปนฺนญฺจ อุเปกฺขาสมฺปนญฺเจว โหติ จิตฺตสฺส อธิฏฺฐานสมฺปนฺนญฺจ สทฺธาสมฺปนฺนญฺจ วีริยสมฺปนฺนญฺจ สติสมฺปนฺนญฺจ สมาธิสมฺปนฺนญฺจ ปญฺญาสมฺปนฺนญฺจ.}

\prose{อากาสานญฺจายตนสมาปตฺติยา ฯเปฯ.\csromanspacer{} วิญฺญาณญฺจายตนสมาปตฺติยา.\csromanspacer{} อากิญฺจญฺญายตนสมาปตฺติยา.\csromanspacer{} เนวสญฺญานาสญฺญายตนสมาปตฺติยา โก อาทิ, กิํ มชฺเฌ, กิํปริโยสานํ ฯเปฯ เอวํ ติวตฺตคตํ จิตฺตํ ติวิธกลฺยาณกํ ทสลกฺขณสมฺปนฺนํ อุเปกฺขาสมฺปนฺนญฺเจว โหติ จิตฺตสฺส อธิฏฺฐานสมฺปนฺนญฺจ ฯเปฯ ปญฺญาสมฺปนฺนญฺจ.}

\prose{อนิจฺจานุปสฺสนาย โก อาทิ, กิํ มชฺเฌ, กิํ ปริโยสานํ ฯเปฯ เอวํ ติวตฺตคตํ จิตฺตํ ติวิธกลฺยาณกํ ทสลกฺขณสมฺปนฺนํ วิตกฺกสมฺปนฺนญฺเจว โหติ วิจารสมฺปนฺนญฺจ ปีติสมฺปนฺนญฺจ สุขสมฺปนฺนญฺจ จิตฺตสฺส อธิฏฺฐานสมฺปนฺนญฺจ สทฺธาสมฺปนฺนญฺจ วีริยสมฺปนฺนญฺจ สติสมฺปนฺนญฺจ สมาธิสมฺปนฺนญฺจ ปญฺญาสมฺปนฺนญฺจ.\csromanspacer{}\csromanspacer{}ทุกฺขานุปสฺสนาย ฯเปฯ.\csromanspacer{} อนตฺตานุปสฺสนาย.\csromanspacer{} นิพฺพิทานุปสฺสนาย.\csromanspacer{} วิราคานุปสฺสนาย.\csromanspacer{} นิโรธานุปสฺสนาย.\csromanspacer{} ปฏินิสฺสคฺคานุปสฺสนาย.\csromanspacer{} ขยานุปสฺสนาย.\csromanspacer{} วยานุปสฺสนาย.\csromanspacer{} วิปริณามานุปสฺสนาย.\csromanspacer{} อนิมิตฺตานุปสฺสนาย.\csromanspacer{} อปฺปณิหิตานุปสฺสนาย.\csromanspacer{} สุญฺญตานุปสฺสนาย.\csromanspacer{} อธิปญฺญาธมฺมวิปสฺสนาย.\csromanspacer{} ยถาภูตญาณทสฺสนาย.\csromanspacer{} อาทีนวานุปสฺสนาย.\csromanspacer{} ปฏิสงฺขานุปสฺสนาย.\csromanspacer{} วิวฏฺฏนานุปสฺสนาย.}

\prose{โสตาปตฺติมคฺคสฺส.\csromanspacer{} สกทาคามิมคฺคสฺส.\csromanspacer{} อนาคามิมคฺคสฺส.\csromanspacer{} อรหตฺตมคฺคสฺส โก อาทิ, กิํ มชฺเฌ, กิํ ปริโยสานํ? อรหตฺตมคฺคสฺส}

\vspace{\tipitakaparskip}
\csromancenter{อานาปานสฺสติกถา ปฏิปทาวิสุทฺธิ อาทิ, อุเปกฺขานุพฺรูหนา มชฺเฌ, สมฺปหํสนา ปริโยสานํ.\csromanspacer{} อรหตฺตมคฺคสฺส ปฏิปทาวิสุทฺธิ อาทิ.\csromanspacer{} อาทิสฺส กติ ลกฺขณานิ? อาทิสฺส ตีณิ ลกฺขณานิ, โย ตสฺส ปริปนฺโถ, ตโต จิตฺตํ วิสุชฺฌติ, วิสุทฺธตฺตา จิตฺตํ มชฺฌิมํ สมถนิมิตฺตํ ปฏิปชฺชติ.\csromanspacer{} ปฏิปนฺนตฺตา ตตฺถ จิตฺตํ ปกฺขนฺทติ, ยญฺจ ปริปนฺถโต จิตฺตํ วิสุชฺฌติ, ยญฺจ วิสุทฺธตฺตา จิตฺตํ มชฺฌิมํ สมถนิมิตฺตํ ปฏิปชฺชติ, ยญฺจ ปฏิปนฺนตฺตา ตตฺถ จิตฺตํ ปกฺขนฺทติ, อรหตฺตมคฺคสฺส ปฏิปทาวิสุทฺธิ อาทิ, อาทิสฺส อิมานิ ตีณิ ลกฺขณานิ.\csromanspacer{} เตน วุจฺจติ อรหตฺตมคฺโค อาทิกลฺยาโณ เจว โหติ ลกฺขณสมฺปนฺโน จ.}
\prose{\csromanfolio{169}อรหตฺตมคฺคสฺส อุเปกฺขานุพฺรูหนา มชฺเฌ.\csromanspacer{} มชฺฌสฺส กติ ลกฺขณานิ? มชฺฌสฺส ตีณิ ลกฺขณานิ, วิสุทฺธํ จิตฺตํ อชฺฌูเปกฺขติ.\csromanspacer{} สมถปฏิปนฺนํ อชฺฌุเปกฺขติ, เอกตฺตุปฏฺฐานํ อชฺฌุเปกฺขติ, ยญฺจ วิสุทฺธํ จิตฺตํ อชฺฌุเปกฺขติ, ยญฺจ สมถปฏิปนฺนํ อชฺฌุเปกฺขติ, ยญฺจ เอกตฺตุปฏฺฐานํ อชฺฌุเปกฺขติ.\csromanspacer{} เตน วุจฺจติ อรหตฺตมคฺโค มชฺเฌกลฺยาโณ เจว โหติ ลกฺขณสมฺปนฺโน จ.}

\prose{อรหตฺตมคฺคสฺส สมฺปหํสนา ปริโยสานํ.\csromanspacer{} ปริโยสานสฺส กติ ลกฺขณานิ? ปริโยสานสฺส จตฺตาริ ลกฺขณานิ, ตตฺถ ชาตานํ ธมฺมานํ อนติวตฺตนฏฺเฐน สมฺปหํสนา, อินฺทฺริยานํ เอกรสฏฺเฐน สมฺปหํสนา, ตทุปควีริยวาหนฏฺเฐน สมฺปหํสนา, อาเสวนฏฺเฐน สมฺปหํสนา.\csromanspacer{} อรหตฺตมคฺคสฺส สมฺปหํสนา ปริโยสานํ, ปริโยสานสฺส อิมานิ จตฺตาริ ลกฺขณานิ, เตน วุจฺจติ อรหตฺตมคฺโค ปริโยสานกลฺยาโณ เจว โหติ ลกฺขณสมฺปนฺโน จ.\csromanspacer{} เอวํ ติวตฺตคตํ จิตฺตํ ติวิธกลฺยาณกํ ทสลกฺขณสมฺปนฺนํ วิตกฺกสมฺปนฺนญฺเจว โหติ วิจารสมฺปนฺนญฺจ ปีติสมฺปนฺนญฺจ สุขสมฺปนฺนญฺจ จิตฺตสฺส อธิฏฺฐานสมฺปนฺนญฺจ สทฺธาสมฺปนฺนญฺจ วีริยสมฺปนฺนญฺจ สติสมฺปนฺนญฺจ สมาธิสมฺปนฺนญฺจ ปญฺญาสมฺปนฺนญฺจ.}

\csromanhangingbegin
\hangingheaditem{159}{นิมิตฺตํ อสฺสาสปสฺสาสา, อนารมฺมณเมกจิตฺตสฺส.}
\hangingbody{อชานโต จ ตโย ธมฺเม, ภาวนา นุปลพฺภติ.}
\hangingbody{นิมิตฺตํ อสฺสาสปสฺสาสา, อนารมฺมณเมกจิตฺตสฺส.}
\hangingbody{ชานโต จ ตโย ธมฺเม, ภาวนา อุปลพฺภตีติ.}
\csromanhangingend

\prose{กถํ อิเม ตโย ธมฺมา เอกจิตฺตสฺส อารมฺมณา น โหนฺติ, น จิเม ตโย ธมฺมา อวิทิตา โหนฺติ, น จ จิตฺตํ วิกฺเขปํ คจฺฉติ, ปธานญฺจ ปญฺญายติ, ปโยคญฺจ สาเธติ, วิเสสมธิคจฺฉติ? เสยฺยถาปิ \csromanfolio{170}รุกฺโข สเม ภูมิภาเค นิกฺขิตฺโต, ตเมนํ ปุริโส กกเจน ฉินฺเทยฺย, รุกฺเข ผุฏฺฐกกจทนฺตานํ วเสน ปุริสสฺส สติ อุปฏฺฐิตา โหติ, น อาคเต วา คเต วา กกจทนฺเต มนสิ กโรติ, น อาคตา วา คตา วา กกจทนฺตา อวิทิตา โหนฺติ, ปธานญฺจ ปญฺญายติ, ปโยคญฺจ สาเธติ.\csromanspacer{} ยถา รุกฺโข สเม ภูมิภาเค นิกฺขิตฺโต, เอวํ อุปนิพนฺธนา นิมิตฺตํ.\csromanspacer{} ยถา กกจทนฺตา, เอวํ อสฺสาสปสฺสาสา.\csromanspacer{} ยถา รุกฺเข ผุฏฺฐกกจทนฺตานํ วเสน ปุริสสฺส สติ อุปฏฺฐิตา โหติ, น อาคเต วา คเต วา กกจทนฺเต มนสิ กโรติ, น อาคตา วา คตา วา กกจทนฺตา อวิทิตา โหนฺติ, ปธานญฺจ ปญฺญายติ, ปโยคญฺจ สาเธติ.\csromanspacer{} เอวเมวํ ภิกฺขุ นาสิกคฺเค วา มุขนิมิตฺเต วา สติํ อุปฏฺฐเปตฺวา นิสินฺโน โหติ.\csromanspacer{} น อาคเต วา คเต วา อสฺสาสปสฺสาเส มนสิ กโรติ, น อาคตา วา คตา วา อสฺสาสปสฺสาสา อวิทิตา โหนฺติ, ปธานญฺจ ปญฺญายติ, ปโยคญฺจ สาเธติ, วิเสสมธิคจฺฉติ ปธานญฺจ.}

\prose{กตมํ ปธานํ? อารทฺธวีริยสฺส กาโยปิ จิตฺตมฺปิ กมฺมนิยํ โหติ.\csromanspacer{} อิทํ ปธานํ.\csromanspacer{}\csromanspacer{}กตโม ปโยโค? อารทฺธวีริยสฺส อุปกฺกิเลสา ปหียนฺติ, วิตกฺกา วูปสมฺมนฺติ.\csromanspacer{} อยํ ปโยโค.\csromanspacer{}\csromanspacer{}กตโม วิเสโส? อารทฺธวีริยสฺส สญฺโญชนา ปหียนฺติ, อนุสยา พฺยนฺตีโหนฺติ\footnote{อนุสยา พฺยาสนฺติ (สฺยา)}.\csromanspacer{} อยํ วิเสโส.\csromanspacer{} เอวํ อิเม ตโย ธมฺมา เอกจิตฺตสฺส อารมฺมณา น โหนฺติ, น จิเม ตโย ธมฺมา อวิทิตา โหนฺติ, น จ จิตฺตํ วิกฺเขปํ คจฺฉติ, ปธานญฺจ ปญฺญายติ, ปโยคญฺจ สาเธติ.\csromanspacer{} วิเสสมธิคจฺฉติ.}

\csromanhangingbegin
\hangingheaditem{160}{อานาปานสฺสติ ยสฺส, ปริปุณฺณา สุภาวิตา.}
\hangingbody{อนุปุพฺพํ ปริจิตา, ยถา พุทฺเธน เทสิตา.}
\hangingbody{โส อิมํ โลกํ ปภาเสติ, อพฺภา มุตฺโตว จนฺทิมาติ.}
\csromanhangingend

\prose{\textbf{อานนฺ}ติ อสฺสาโส, โน ปสฺสาโส.\csromanspacer{} \textbf{อาปานนฺ}ติ\footnote{อปานนฺติ (ก) 69 ปิฏฺเฐ อฏฺฐกถา โอโลเกตพฺพา.} ปสฺสาโส, โน อสฺสาโส.\csromanspacer{} อสฺสาสวเสน อุปฏฺฐานํ สติ, ปสฺสาสวเสน อุปฏฺฐานํ สติ.}

\prose{โย อสฺสสติ, ตสฺสุปฏฺฐาติ.\csromanspacer{} โย ปสฺสสติ ตสฺสุปฏฺฐาติ.\csromanspacer{} \textbf{ปริปุณฺณา}ติ ปริคฺคหฏฺเฐน ปริปุณฺณา, ปริวารฏฺเฐน ปริปุณฺณา, ปริปูรฏฺเฐน ปริปุณฺณา.\csromanspacer{} \textbf{สุภาวิตา}ติ จตสฺโส ภาวนา.\csromanspacer{} ตตฺถ ชาตานํ ธมฺมานํ}

\vspace{\tipitakaparskip}
\csromancenter{อานาปานสฺสติกถา อนติวตฺตนฏฺเฐน ภาวนา, อินฺทฺริยานํ เอกรสฏฺเฐน ภาวนา, ตทุปควีริยวาหนฏฺเฐน ภาวนา.\csromanspacer{} อาเสวนฏฺเฐน ภาวนา.\csromanspacer{} ตสฺสิเม จตฺตาโร ภาวนฏฺฐา ยานีกตา โหนฺติ วตฺถุกตา อนุฏฺฐิตา ปฺ\textbf{อริจิตา} สุสมารทฺธา.}
\prose{\csromanfolio{171}\textbf{ยานีกตา}ติ ยตฺถ ยตฺถ อากงฺขติ, ตตฺถ ตตฺถ วสิปฺปตฺโต โหติ พลปฺปตฺโต เวสารชฺชปฺปตฺโต.\csromanspacer{} ตสฺส เม เต ธมฺมา อาวชฺชนปฏิพทฺธา โหนฺติ อากงฺขปฏิพทฺธา มนสิการปฏิพทฺธา จิตฺตุปฺปาทปฏิพทฺธา.\csromanspacer{} เตน วุจฺจติ “ยานีกตา”ติ.\csromanspacer{}\csromanspacer{}\textbf{วตฺถุกตา}ติ ยสฺมิํ ยสฺมิํ วตฺถุสฺมิํ จิตฺตํ สฺวาธิฏฺฐิตํ โหติ, ตสฺมิํ ตสฺมิํ วตฺถุสฺมิํ สติ สูปฏฺฐิตา\footnote{สุปฏฺฐิตา (ก)} โหติ.\csromanspacer{} ยสฺมิํ ยสฺมิํ วา ปน วตฺถุสฺมิํ สติ สูปฏฺฐิตา โหติ, ตสฺมิํ ตสฺมิํ วตฺถุสฺมิํ จิตฺตํ สฺวาธิฏฺฐิตํ โหติ.\csromanspacer{} เตน วุจฺจติ “วตฺถุกตา”ติ.\csromanspacer{}\csromanspacer{}\textbf{อนุฏฺฐิตา}ติ วตฺถุสฺมิํ เยน เยน จิตฺตํ อภินีหรติ, เตน เตน สติ อนุปริวตฺตติ.\csromanspacer{} เยน เยน วา ปน สติ อนุปริวตฺตติ, เตน เตน จิตฺตํ อภินีหรติ.\csromanspacer{} เตน วุจฺจติ “อนุฏฺฐิตา”ติ.\csromanspacer{}\csromanspacer{}ปฺ\textbf{อริจิตา}ติ ปริคฺคหฏฺเฐน ปฺ\textbf{อริจิตา}, ปริวารฏฺเฐน ปฺ\textbf{อริจิตา}, ปริปูรฏฺเฐน ปฺ\textbf{อริจิตา}, สติยา ปริคฺคณฺหนฺโต ชินาติ ปาปเก อกุสเล ธมฺเม.\csromanspacer{} เตน วุจฺจติ “ปฺ\textbf{อริจิตา}”ติ.\csromanspacer{}\csromanspacer{}\textbf{สุสมารทฺธา}ติ จตฺตาโร สุสมารทฺธา.\csromanspacer{} ตตฺถ ชาตานํ ธมฺมานํ อนติวตฺตนฏฺเฐน สุสมารทฺธา, อินฺทฺริยานํ เอกรสฏฺเฐน สุสมารทฺธา, ตทุปควีริยวาหนฏฺเฐน สุสมารทฺธา, ตปฺปจฺจนีกานํ กิเลสานํ สุสมูหตตฺตา\footnote{สุสมุหตตฺตา (ก)} สุสมารทฺธา.}

\proseitem{161}{\textbf{สุสมนฺ}ติ อตฺถิ สมํ อตฺถิ สุสมํ.\csromanspacer{} กตมํ สมํ? เย ตตฺถ ชาตา อนวชฺชา กุสลา โพธิปกฺขิยา, อิทํ สมํ.\csromanspacer{} กตมํ สุสมํ? ยํ เตสํ เตสํ ธมฺมานํ อารมฺมณํ นิโรโธ นิพฺพานํ, อิทํ สุสมํ.\csromanspacer{} อิติ อิทญฺจ สมํ, อิทญฺจ สุสมํ ญาตํ โหติ ทิฏฺฐิ วิทิตํ สจฺฉิกตํ ผสฺสิตํ ปญฺญาย, อารทฺธํ โหติ วีริยํ อสลฺลีนํ, อุปฏฺฐิตา สติ อสมฺมุฏฺฐา\footnote{อปมุฏฺฐา (สฺยา)}, ปสฺสทฺโธ กาโย อสารทฺโธ, สมาหิตํ จิตฺตํ เอกคฺคํ.\csromanspacer{} เตน วุจฺจติ “สุสมารทฺธา”ติ.}

\prose{\textbf{อนุปุพฺพํ ปริจิตา}ติ ทีฆํ อสฺสาสวเสน ปุริมา ปุริมา ปฺ\textbf{อริจิตา}, ปจฺฉิมา ปจฺฉิมา อนุปฺ\textbf{อริจิตา}.\csromanspacer{} ทีฆํ ปสฺสาสวเสน ปุริมา ปุริมา ปฺ\textbf{อริจิตา}, ปจฺฉิมา ปจฺฉิมา อนุปฺ\textbf{อริจิตา}.\csromanspacer{} รสฺสํ อสฺสาสวเสน ปุริมา ปุริมา \csromanfolio{172}ปริจิตา, ปจฺฉิมา ปจฺฉิมา อนุปริจิตา.\csromanspacer{} รสฺสํ ปสฺสาสวเสน ปุริมา ปุริมา ปริจิตา, ปจฺฉิมา ปจฺฉิมา อนุปริจิตา ฯเปฯ.\csromanspacer{} ปฏินิสฺสคฺคานุปสฺสี อสฺสาสวเสน ปุริมา ปุริมา ปริจิตา, ปจฺฉิมา ปจฺฉิมา อนุปริจิตา.\csromanspacer{} ปฏินิสฺสคฺคานุปสฺสี ปสฺสาสวเสน ปุริมา ปุริมา ปริจิตา, ปจฺฉิมา ปจฺฉิมา อนุปริจิตา.\csromanspacer{} สพฺพาปิ โสฬสวตฺถุกา อานาปานสฺสติโย อญฺญมญฺญํ ปริจิตา เจว โหนฺติ อนุปริจิตา จ.\csromanspacer{} เตน วุจฺจติ “อนุปุพฺพปริจิตา”ติ.}

\prose{\textbf{ยถา}ติ ทส ยถตฺถา.\csromanspacer{} อตฺตทมถตฺโถ ยถตฺโถ, อตฺตสมถตฺโถ ยถตฺโถ, อตฺตปรินิพฺพาปนตฺโถ ยถตฺโถ, อภิญฺญตฺโถ ยถตฺโถ, ปริญฺญตฺโถ ยถตฺโถ, ปหานตฺโถ ยถตฺโถ, ภาวนตฺโถ ยถตฺโถ, สจฺฉิกิริยตฺโถ ยถตฺโถ, สจฺจาภิสมยตฺโถ ยถตฺโถ, นิโรเธ ปติฏฺฐาปกตฺโถ ยถตฺโถ.}

\prose{\textbf{พุทฺโธ}ติ โย โส ภควา สยมฺภู อนาจริยโก ปุพฺเพ อนนุสฺสุเตสุ ธมฺเมสุ สามํ สจฺจานิ อภิสมฺพุชฺฌิ, ตตฺถ จ สพฺพญฺญุตํ ปาปุณิ, พเลสุ จ วสีภาวํ.}

\proseitem{162}{\textbf{พุทฺโธ}ติ เกนฏฺเฐน \textbf{พุทฺโธ}? พุชฺฌิตา สจฺจานีติ \textbf{พุทฺโธ}.\csromanspacer{} โพเธตา ปชายาติ \textbf{พุทฺโธ}.\csromanspacer{} สพฺพญฺญุตาย \textbf{พุทฺโธ}.\csromanspacer{} สพฺพทสฺสาวิตาย \textbf{พุทฺโธ}.\csromanspacer{} อนญฺญเนยฺยตาย \textbf{พุทฺโธ}.\csromanspacer{} วิสวิตาย\footnote{วิกติตาย (สฺยา)} \textbf{พุทฺโธ}.\csromanspacer{} ขีณาสวสงฺขาเตน \textbf{พุทฺโธ}.\csromanspacer{} นิรุปเลปสงฺขาเตน\footnote{นิรุปกฺกิเลสสงฺขาเตน (สฺยา)} \textbf{พุทฺโธ}.\csromanspacer{} เอกนฺตวีตราโคติ \textbf{พุทฺโธ}.\csromanspacer{} เอกนฺตวีตโทโสติ \textbf{พุทฺโธ}.\csromanspacer{} เอกนฺตวีตโมโหติ \textbf{พุทฺโธ}.\csromanspacer{} เอกนฺตนิกฺกิเลโสติ \textbf{พุทฺโธ}.\csromanspacer{} เอกายนมคฺคํ คโตติ \textbf{พุทฺโธ}.\csromanspacer{} เอโก อนุตฺตรํ สมฺมาสมฺโพธิํ อภิสมฺพุทฺโธติ \textbf{พุทฺโธ}.\csromanspacer{} อพุทฺธิวิหตตฺตา พุทฺธิปฏิลาภา \textbf{พุทฺโธ}. “\textbf{พุทฺโธ}”ติ เนตํ นามํ มาตรา กตํ, น ปิตรา กตํ, น ภาตรา กตํ, น ภคินิยา กตํ, น มิตฺตามจฺเจหิ กตํ, น ญาติสาโลหิเตหิ กตํ, น สมณพฺราหฺมเณหิ กตํ, น เทวตาหิ กตํ, วิโมกฺขนฺติกเมตํ พุทฺธานํ ภควนฺตานํ โพธิยา มุเล สห สพฺพญฺญุตญาณสฺส ปฏิลาภา สจฺฉิกา ปญฺญตฺติ, ยทิทํ \textbf{พุทฺโธ}ติ.\csromanspacer{}\csromanspacer{}\textbf{เทสิตา}ติ อตฺตทมถตฺโถ ยถตฺโถ ยถา พุทฺเธน เทสิโต.\csromanspacer{} อตฺตสมถตฺโถ ยถตฺโถ ยถา พุทฺเทน เทสิโต.\csromanspacer{} อตฺตปรินิพฺพาปนตฺโถ ยถตฺโถ ยถา พุทฺเธน เทสิโต ฯเปฯ นิโรเธ ปติฏฺฐาปกตฺโถ ยถตฺโถ ยถา พุทฺเธน เทสิโต.}

\csromanheader{2}{3. อานาปานสฺสติกถา}
\prose{\csromanfolio{173}\textbf{โส}ติ คหฏฺโฐ วา โหติ ปพฺพชิโต วา.\csromanspacer{}\csromanspacer{}\textbf{โลโก}ติ ขนฺธโลโก ธาตุโลโก อายตนโลโก วิปตฺติภวโลโก วิปตฺติสมฺภวโลโก สมฺปตฺติภวโลโก สมฺปตฺติสมฺภวโลโก.\csromanspacer{} เอโก โลโก, สพฺเพ สตฺตา อาหารฏฺฐิติกา ฯเปฯ.\csromanspacer{} อฏฺฐารส โลกา อฏฺฐารส ธาตุโย.\csromanspacer{}\csromanspacer{}\textbf{ปภาเสตี}ติ อตฺตทมตฺถตฺถํ ยถตฺถํ อภิสมฺพุทฺธตฺตา โส อิมํ โลกํ โอภาเสติ ภาเสติ ปภาเสติ.\csromanspacer{} อตฺตสมถตฺถํ ยถตฺถํ อภิสมฺพุทฺธตฺตา โส อิมํ โลกํ โอภาเสติ ภาเสติ ปภาเสติ.\csromanspacer{} อตฺตปรินิพฺพาปนตฺถํ ยถตฺถํ อภิสมฺพุทฺธตฺตา โส อิมํ โลกํ โอภาเสติ ภาเสติ ปภาเสติ ฯเปฯ นิโรเธ ปติฏฺฐาปกตฺถํ ยถตฺถํ อภิสมฺพุทฺธตฺตา โส อิมํ โลกํ โอภาเสติ ภาเสติ ปภาเสติ.}

\prose{\textbf{อพฺภา มุตฺโตว จนฺทิมา}ติ ยถา อพฺภา, เอวํ กิเลสา.\csromanspacer{} ยถา จนฺโท เอวํ อริยญาณํ.\csromanspacer{} ยถา จนฺทิมา เทวปุตฺโต, เอวํ ภิกฺขุ.\csromanspacer{} ยถา จนฺโท อพฺภา มุตฺโต มหิกา มุตฺโต ธูมรชา มุตฺโต ราหุคหณา\footnote{ราหุปาณา (สฺยา)} วิปฺปมุตฺโต ภาสเต จ ตปเต จ วิโรจเต\footnote{วิโรจติ (สฺยา)} จ.\csromanspacer{} เอวเมวํ ภิกฺขุ สพฺพกิเลเสหิ วิปฺปมุตฺโต ภาสเต จ ตปเต จ วิโรจเต จ, เตน วุจฺจติ “อพฺภา มุตฺโตว จนฺทิมา”ติ.\csromanspacer{} อิมานิ เตรส โวทาเน ญาณานิ.}

\vspace{\tipitakaparskip}
\csromancenter{โวทานญาณนิทฺเทโส จตุตฺโถ.}
\csromanheader{2}{ภาณวาโร.}
\csromansectionrule
\csromanmatikaanchor{mtk.70}
\csromantocmarkref{subsection}{5. สโตการิญาณนิทฺเทส}{mtk.70}
\bookmark[dest={mtk.70},level=2]{5. สโตการิญาณนิทฺเทส}
\csromanheader{2}{5. สโตการิญาณนิทฺเทส}
\proseitem{163}{กตมานิ พาตฺติํส สโตการิสฺส ญาณานิ? อิธ ภิกฺขุ อรญฺญคโต วา รุกฺขมูลคโต วา สุญฺญาคารคโต วา นิสีทติ ปลฺลงฺกํ อาภุชิตฺวา อุชุํ กายํ ปณิธาย ปริมุขํ สติํ อุปฏฺฐเปตฺวา, โส สโตว อสฺสสติ, สโต ปสฺสสติ\csromandash{}ทีฆํ วา อสฺสสนฺโต “ทีฆํ อสฺสสามี”ติ ปชานาติ, ทีฆํ วา ปสฺสสนฺโต “ทีฆํ ปสฺสสามี”ติ ปชานาติ.\csromanspacer{} รสฺสํ วา อสฺสสนฺโต “รสฺสํ อสฺสสามี”ติ ปชานาติ, รสฺสํ วา ปสฺสสนฺโต “รสฺสํ ปสฺสสามี”ติ ปชานาติ. “สพฺพกายปฏิสํเวที อสฺสสิสฺสามี”ติ สิกฺขติ, “สพฺพกายปฏิสํเวที ปสฺสสิสฺสามี”ติ \csromanfolio{174}สิกฺขติ. “ปสฺสมฺภยํ กายสงฺขารํ อสฺสสิสฺสามี”ติ สิกฺขติ, “ปสฺสมฺภยํ กายสงฺขารํ ปสฺสสิสฺสามี”ติ สิกฺขติ. “ปีติปฏิสํเวที ฯเปฯ.\csromanspacer{} สุขปฏิสํเวที.\csromanspacer{} จิตฺตสงฺขารปฏิสํเวที.\csromanspacer{} ปสฺสมฺภยํ จิตฺตสงฺขารํ.\csromanspacer{} จิตฺตปฏิสํเวที.\csromanspacer{} อภิปฺปโมทยํ จิตฺตํ.\csromanspacer{} สมาทหํ จิตฺตํ.\csromanspacer{} วิโมจยํ จิตฺตํ.\csromanspacer{} อนิจฺจานุปสฺสี.\csromanspacer{} วิราคานุปสฺสี.\csromanspacer{} นิโรธานุปสฺสี.\csromanspacer{} ปฏินิสฺสคฺคานุปสฺสี อสฺสสิสฺสามี”ติ สิกฺขติ, “ปฏินิสฺสคฺคานุปสฺสี ปสฺสสิสฺสามี”ติ สิกฺขติ.}

\proseitem{164}{\textbf{อิธา}ติ อิมิสฺสา ทิฏฺฐิยา อิมิสฺสา ขนฺติยา อิมิสฺสา รุจิยา อิมสฺมิํ อาทาเย อิมสฺมิํ ธมฺเม อิมสฺมิํ วินเย อิมสฺมิํ ธมฺมวินเย อิมสฺมิํ ปาวจเน อิมสฺมิํ พฺรหฺมจริเย อิมสฺมิํ สตฺถุสาสเน, เตน วุจฺจติ “อิธา”ติ.\csromanspacer{} พฺ\textbf{หิกฺขู}ติ ปุถุชฺชนกลฺยาณโก วา โหติ ภิกฺขุ เสกฺโข วา อรหา วา อกุปฺปธมฺโม.\csromanspacer{} \textbf{อรญฺญนฺ}ติ นิกฺขมิตฺวา พหิ อินฺทขีลา, สพฺพเมตํ อรญฺญํ.\csromanspacer{} \textbf{รุกฺขมูลนฺ}ติ ยตฺถ ภิกฺขุโน อาสนํ ปญฺญตฺตํ โหติ มญฺโจ วา ปีฐํ วา ภิสิ วา ตฏฺฏิกา วา จมฺมขณฺโฑ วา ติณสนฺถโร\footnote{ติณสณฺฐโร (สฺยา)} วา ปณฺณสนฺถโร วา ปลาลสนฺถโร\footnote{ปลาสสณฺฐโร (สฺยา)} วา, ตตฺถ ภิกฺขุ จงฺกมติ วา ติฏฺฐติ วา นิสีทติ วา เสยฺยํ วา กปฺเปติ.\csromanspacer{} \textbf{สุญฺญนฺ}ติ เกนจิ อนากิณฺณํ โหติ คหฏฺเฐหิ วา ปพฺพชิเตหิ วา.\csromanspacer{} \textbf{อคารนฺ}ติ\footnote{อาคารนฺติ (สฺยา)} วิหาโร อฑฺฒโยโค ปาสาโท หมฺมิยํ คุหา.\csromanspacer{} \textbf{นิสีทติ ปลฺลงฺกํ อาภุชิตฺวา}ติ นิสินฺโน โหติ ปลฺลงฺกํ อาภุชิตฺวา.\csromanspacer{} \textbf{อุชุํ กายํ ปณิธายา}ติ อุชุโก โหติ กาโย ฐิโต\footnote{ปนิธิโต (สฺยา)} สุปณิหิโต.\csromanspacer{} \textbf{ปริมุขํ สติํ อุปฏฺฐเปตฺวา}ติ \textbf{ปรี}ติ ปริคฺคหฏฺโฐ.\csromanspacer{} มฺ\textbf{อุขนฺ}ติ นิยฺยานฏฺโฐ.\csromanspacer{} \textbf{สตี}ติ อุปฏฺฐานฏฺโฐ, เตน วุจฺจติ “ปริมุขํ สติํ อุปฏฺฐเปตฺวา”ติ.}

\proseitem{165}{\textbf{สโตว อสฺสสติ, สโต ปสฺสสตี}ติ พาตฺติํสาย อากาเรหิ สโต การี โหติ.\csromanspacer{} ทีฆํ อสฺสาสวเสน จิตฺตสฺส เอกคฺคตํ อวิกฺเขปํ ปชานโต สติ อุปฏฺฐิตา โหติ, ตาย สติยา เตน ญาเณน สโต การี โหติ.\csromanspacer{} ทีฆํ ปสฺสาสวเสน จิตฺตสฺส เอกคฺคตํ อวิกฺเขปํ ปชานโต สติ อุปฏฺฐิตา โหติ, ตาย สติยา เตน ญาเณน สโต การี โหติ.\csromanspacer{} รสฺสํ อสฺสาสวเสน จิตฺตสฺส เอกคฺคตํ อวิกฺเขปํ ปชานโต สติ อุปฏฺฐิตา โหติ, ตาย สติยา เตน ญาเณน สโต การี โหติ.\csromanspacer{} รสฺสํ ปสฺสาสวเสน จิตฺตสฺส}

\vspace{\tipitakaparskip}
\csromancenter{อานาปานสฺสติกถา เอกคฺคตํ อวิกฺเขปํ ปชานโต สติ อุปฏฺฐิตา โหติ, ตาย สติยา เตน ญาเณน สโต การี โหติ ฯเปฯ.\csromanspacer{} ปฏินิสฺสคฺคานุปสฺสี อสฺสาสวเสน.\csromanspacer{} ปฏินิสฺสคฺคานุปสฺสี ปสฺสาสวเสน จิตฺตสฺส เอกคฺคตํ อวิกฺเขปํ ปชานโต สติ อุปฏฺฐิตา โหติ, ตาย สติยา เตน ญาเณน สโต การี โหติ.}
\csromanheader{2}{\textbf{ปฐมจตุกฺกนิทฺเทส}}
\proseitem{166}{\csromanfolio{175}กถํ ทีฆํ อสฺสสนฺโต “ทีฆํ อสฺสสามี”ติ ปชานาติ, ทีฆํ ปสฺสสนฺโต “ทีฆํ ปสฺสสามี”ติ ปชานาติ? ทีฆํ อสฺสาสํ อทฺธานสงฺขาเต อสฺสสติ, ทีฆํ ปสฺสาสํ อทฺธานสงฺขาเต ปสฺสสติ, ทีฆํ อสฺสาสปสฺสาสํ อทฺธานสงฺขาเต อสฺสสติปิ อสฺสสติปิ, ทีฆํ อสฺสาสปสฺสาสํ อทฺธานสงฺขาเต อสฺสสโตปิ ปสฺสสโตปิ ฉนฺโท อุปฺปชฺชติ.\csromanspacer{} ฉนฺทวเสน ตโต สุขุมตรํ ทีฆํ อสฺสาสํ อทฺธานสงฺขาเต อสฺสสติ.\csromanspacer{} ฉนฺทวเสน ตโต สุขุมตรํ ทีฆํ ปสฺสาสํ อทฺธานสงฺขาเต ปสฺสสติ.\csromanspacer{} ฉนฺทวเสน ตโต สุขุมตรํ ทีฆํ อสฺสาสปสฺสาสํ อทฺธานสงฺขาเต อสฺสสติปิ ปสฺสสติปิ.\csromanspacer{} ฉนฺทวเสน ตโต สุขุมตรํ ทีฆํ อสฺสาสปสฺสาสํ อทฺธานสงฺขาเต อสฺสสโตปิ ปสฺสสโตปิ ปาโมชฺชํ อุปฺปชฺชติ.\csromanspacer{} ปาโมชฺชวเสน ตโต สุขุมตรํ ทีฆํ อสฺสาสํ อทฺธานสงฺขาเต อสฺสสติ.\csromanspacer{} ปาโมชฺชวเสน ตโต สุขุมตรํ ทีฆํ ปสฺสาสํ อทฺธานสงฺขาเต ปสฺสสติ.\csromanspacer{} ปาโมชฺชวเสน ตโต สุขุมตรํ ทีฆํ อสฺสาสปสฺสาสํ อทฺธานสงฺขาเต อสฺสสติปิ ปสฺสสติปิ.\csromanspacer{} ปาโมชฺชวเสน ตโต สุขุมตรํ ทีฆํ อสฺสาสปสฺสาสํ อทฺธานสงฺขาเต อสฺสสโตปิ ปสฺสสโตปิ ทีฆํ อสฺสาสปสฺสาสาปิ จิตฺตํ วิวตฺตติ, อุเปกฺขา สณฺฐาติ.\csromanspacer{} อิเมหิ นวหากาเรหิ ทีฆํ อสฺสาสปสฺสาสา กาโย, อุปฏฺฐานํ สติ, อนุปสฺสนา ญาณํ.\csromanspacer{} กาโย อุปฏฺฐานํ, โน สติ.\csromanspacer{} สติ อุปฏฺฐานญฺเจว สติ จ.\csromanspacer{} ตาย สติยา เตน ญาเณน ตํ กายํ อนุปสฺสติ, เตน วุจฺจติ “กาเย กายานุปสฺสนาสติปฏฺฐานภาวนา”ติ.}

\proseitem{167}{อ\textbf{นุปสฺสตี}ติ กถํ ตํ กายํ อนุปสฺสติ? อนิจฺจโต อนุปสฺสติ, โน นิจฺจโต.\csromanspacer{} ทุกฺขโต อนุปสฺสติ, โน สุขโต.\csromanspacer{} อนตฺตโต อนุปสฺสติ, โน อตฺตโต.\csromanspacer{} นิพฺพินฺทติ, โน นนฺทติ.\csromanspacer{} วิรชฺชติ, โน รชฺชติ.\csromanspacer{} นิโรเธติ, โน สมุเทติ.\csromanspacer{} ปฏินิสฺสชฺชติ, โน อาทิยติ. \csromanfolio{176}อนิจฺจโต อนุปสฺสนฺโต นิจฺจสญฺญํ ปชหติ.\csromanspacer{} ทุกฺขโต อนุปสฺสนฺโต สุขสญฺญํ ปชหติ.\csromanspacer{} อนตฺตโต อนุปสฺสนฺโต อตฺตสญฺญํ ปชหติ.\csromanspacer{} นิพฺพินฺทนฺโต นนฺทิํ ปชหติ.\csromanspacer{} วิรชฺชนฺโต ราคํ ปชหติ.\csromanspacer{} นิโรเธนฺโต สมุทยํ ปชหติ.\csromanspacer{} ปฏินิสฺสชฺชนฺโต อาทานํ ปชหติ.\csromanspacer{} เอวํ ตํ กายํ อนุปสฺสติ.}

\prose{\textbf{ภาวนา}ติ จตสฺโส ภาวนา.\csromanspacer{} ตตฺถ ชาตานํ ธมฺมานํ อนติวตฺตนฏฺเฐน ภาวนา, อินฺทฺริยานํ เอกรสฏฺเฐน ภาวนา, ตทุปควีริยวาหนฏฺเฐน ภาวนา, อาเสวนฏฺเฐน ภาวนา, ทีฆํ อสฺสาสปสฺสาสวเสน จิตฺตสฺส เอกคฺคตํ อวิกฺเขปํ ปชานโต วิทิตา เวทนา อุปฺปชฺชนฺติ, วิทิตา อุปฏฺฐหนฺติ, วิทิตา อพฺภตฺถํ คจฺฉนฺติ, วิทิตา สญฺญา อุปฺปชฺชนฺติ, วิทิตา อุปฏฺฐหนฺติ, วิทิตา อพฺภตฺถํ คจฺฉนฺติ, วิทิตา วิตกฺกา อุปฺปชฺชนฺติ, วิทิตา อุปฏฺฐนฺติ, วิทิตา อพฺภตฺถํ คจฺฉนฺติ.}

\prose{กถํ วิทิตา เวทนา อุปฺปชฺชนฺติ, วิทิตา อุปฏฺฐหนฺติ, วิทิตา อพฺภตฺถํ คจฺฉนฺติ? กถํ เวทนาย อุปฺปาโท วิทิโต โหติ? “อวิชฺชาสมุทยา เวทนาสมุทโย”ติ ปจฺจยสมุทยฏฺเฐน เวทนาย อุปฺปาโท วิทิโต โหติ. “ตณฺหาสมุทยา เวทนาสมุทโย”ติ. “กมฺมสมุทยา เวทนาสมุทโย”ติ. “ผสฺสสมุทยา เวทนาสมุทโย”ติ ปจฺจยสมุทยฏฺเฐน เวทนาย อุปฺปาโท วิทิโต โหติ.\csromanspacer{} นิพฺพตฺติลกฺขณํ ปสฺสโตปิ เวทนาย อุปฺปาโท วิทิโต โหติ.\csromanspacer{} เอวํ เวทนาย อุปฺปาโท วิทิโต โหติ.}

\prose{กถํ เวทนาย อุปฏฺฐานํ วิทิตํ โหติ? อนิจฺจโต มนสิกโรโต ขยตุปฏฺฐานํ วิทิตํ โหติ, ทุกฺขโต มนสิกโรโต ภยตุปฏฺฐานํ วิทิตํ โหติ, อนตฺตโต มนสิกโรโต สุญฺญตุปฏฺฐานํ วิทิตํ โหติ.\csromanspacer{} เอวํ เวทนาย อุปฏฺฐานํ วิทิตํ โหติ.}

\prose{กถํ เวทนาย อตฺถงฺคโม วิทิโต โหติ? “อวิชฺชานิโรธา เวทนานิโรโธ”ติ ปจฺจยนิโรธฏฺเฐน เวทนาย อตฺถงฺคโม วิทิโต โหติ, “ตณฺหานิโรธา เวทนานิโรโธ”ติ ฯเปฯ “กมฺมนิโรธา เวทนานิโรโธ”ติ ฯเปฯ “ผสฺสนิโรธา เวทนานิโรโธ”ติ ปจฺจยนิโรธฏฺเฐน เวทนาย อตฺถงฺคโม วิทิโต โหติ, วิปริณามลกฺขณํ ปสฺสโตปิ เวทนาย อตฺถงฺคโม วิทิโต โหติ.\csromanspacer{} เอวํ เวทนาย อตฺถงฺคโม วิทิโต โหติ.\csromanspacer{} เอวํ วิทิตา เวทนา อุปฺปชฺชนฺติ, วิทิตา อุปฏฺฐหนฺติ, วิทิตา อพฺภตฺถํ คจฺฉนฺติ.}

\csromanheader{2}{3. อานาปานสฺสติกถา}
\prose{\csromanfolio{177}กถํ วิทิตา สญฺญา อุปฺปชฺชนฺติ, วิทิตา อุปฏฺฐหนฺติ, วิทิตา อพฺภตฺถํ คจฺฉนฺติ, กถํ สญฺญาย อุปฺปาโท วิทิโต โหติ? “อวิชฺชาสมุทยา สญฺญาสมุทโย”ติ ปจฺจยสมุทยฏฺเฐน สญฺญาย อุปฺปาโท วิทิโต โหติ, “ตณฺหาสมุทยา สญฺญาสมุทโย”ติ ฯเปฯ “กมฺมสมุทยา สญฺญาสมุทโย”ติ ฯเปฯ “ผสฺสสมุทยา สญฺญาสมุทโย”ติ ปจฺจยสมุทยฏฺเฐน สญฺญาย อุปฺปาโท วิทิโต โหติ, นิพฺพตฺติลกฺขณํ ปสฺสโตปิ สญฺญาย อุปฺปโท วิทิโต โหติ.\csromanspacer{} เอวํ สญฺญาย อุปฺปาโท วิทิโต โหติ.}

\prose{กถํ สญฺญาย อุปฏฺฐานํ วิทิตํ โหติ? อนิจฺจโต มนสิกโรโต ขยตุปฏฺฐานํ วิทิตํ โหติ, ทุกฺขโต มนสิกโรโต ภยตุปฏฺฐานํ วิทิตํ โหติ, อนตฺตโต มนสิกโรโต สุญฺญาตุปฏฺฐานํ วิทิตํ โหติ.\csromanspacer{} เอวํ สญฺญาย อุปฏฺฐานํ วิทิตํ โหติ.}

\prose{กถํ สญฺญาย อตฺถงฺคโม วิทิโต โหติ? “อวิชฺชานิโรธา สญฺญานิโรโธ”ติ ปจฺจยนิโรธฏฺเฐน สญฺญาย อตฺถงฺคโม วิทิโต โหติ, “ตณฺหานิโรธา สญฺญานิโรโธ”ติ ฯเปฯ “กมฺมนิโรธา สญฺญานิโรโธ”ติ ฯเปฯ “ผสฺสนิโรธา สญฺญานิโรโธ”ติ ปจฺจยนิโรธฏฺเฐน สญฺญาย อตฺถงฺคโม วิทิโต โหติ, วิปริณามลกฺขณํ ปสฺสโตปิ สญฺญาย อตฺถงฺคโม วิทิโต โหติ.\csromanspacer{} เอวํ สญฺญาย อตฺถงฺคโม วิทิโต โหติ.\csromanspacer{} เอวํ วิทิตา สญฺญา อุปฺปชฺชนฺติ, วิทิตา อุปฏฺฐหนฺติ, วิทิตา อพฺภตฺถํ คจฺฉนฺติ.}

\prose{กถํ วิทิตา วิตกฺกา อุปฺปชฺชนฺติ, วิทิตา อุปฏฺฐหนฺติ, วิทิตา อพฺภตฺถํ คจฺฉนฺติ, กถํ วิตกฺกานํ อุปฺปาโท วิทิโต โหติ? “อวิชฺชาสมุทยา วิตกฺกสมุทโย”ติ ปจฺจยสมุทยฏฺเฐน วิตกฺกานํ อุปฺปาโท วิทิโต โหติ, “ตณฺหาสมุทยา วิตกฺกสมุทโย”ติ ฯเปฯ “กมฺมสมุทยา วิตกฺกสมุทโย”ติ ฯเปฯ “สญฺญาสมุทยา วิตกฺกสมุทโย”ติ ปจฺจยสมุทยฏฺเฐน วิตกฺกานํ อุปฺปาโท วิทิโต โหติ, นิพฺพตฺติลกฺขณํ ปสฺสโตปิ วิตกฺกานํ อุปฺปาโท วิทิโต โหติ.\csromanspacer{} เอวํ วิตกฺกานํ อุปฺปาโท วิทิโต โหติ.}

\prose{กถํ วิตกฺกานํ อุปฏฺฐานํ วิทิตํ โหติ? อนิจฺจโต มนสิกโรโต ขยตุปฏฺฐานํ วิทิตํ โหติ, ทุกฺขโต มนสิกโรโต ภยตุปฏฺฐานํ วิทิตํ โหติ, อนตฺตโต มนสิกโรโต สุญฺญตุปฏฺฐานํ วิทิตํ โหติ.\csromanspacer{} เอวํ วิตกฺกานํ อุปฏฺฐานํ วิทิตํ โหติ.}

\prose{\csromanfolio{178}กถํ วิตกฺกานํ อตฺถงฺคโม วิทิโต โหติ? “อวิชฺชานิโรธา วิตกฺกนิโรโธ”ติ ปจฺจยนิโรธฏฺเฐน วิตกฺกานํ อตฺถงฺคโม วิทิโต โหติ, “ตณฺหานิโรธา วิตกฺกนิโรโธ”ติ ฯเปฯ “กมฺมนิโรธา วิตกฺกนิโรโธ”ติ ฯเปฯ “สญฺญานิโรธา วิตกฺกนิโรโธ”ติ ปจฺจยนิโรธฏฺเฐน วิตกฺกานํ อตฺถงฺคโม วิทิโต โหติ, วิปริณามลกฺขณํ ปสฺสโตปิ วิตกฺกานํ อตฺถงฺคโม วิทิโต โหติ.\csromanspacer{} เอวํ วิตกฺกานํ อตฺถงฺคโม วิทิโต โหติ.\csromanspacer{} เอวํ วิทิตา วิตกฺกา อุปฺปชฺชนฺติ, วิทิตา อุปฏฺฐหนฺติ, วิทิตา อพฺภตฺถํ คจฺฉนฺติ.}

\proseitem{168}{ทีฆํ อสฺสาสปสฺสาสวเสน จิตฺตสฺส เอกคฺคตํ อวิกฺเขปํ ปชานนฺโต อินฺทฺริยานิ สโมธาเนติ, โคจรญฺจ ปชานาติ, สมตฺถญฺจ ปฏิวิชฺฌติ ฯเปฯ มคฺคํ สโมธาเนติ, ธมฺเม สโมธาเนติ, โคจรญฺจ ปชานาติ, สมตฺตญฺจ ปฏิวิชฺฌติ.}

\prose{\textbf{อินฺทฺริยานิ สโมธาเนตี}ติ กถํ อินฺทฺริยานิ สโมธาเนติ? อธิโมกฺขฏฺเฐน สทฺธินฺทฺริยํ สโมธาเนติ, ปคฺคหฏฺเฐน วีริยินฺทฺริยํ สโมธาเนติ, อุปฏฺฐานฏฺเฐน สตินฺทฺริยํ สโมธาเนติ, อวิกฺเขปฏฺเฐน สมาธินฺทฺริยํ สโมธาเนติ, ทสฺสนฏฺเฐน ปญฺญินฺทฺริยํ สโมธาเนติ.\csromanspacer{} อยํ ปุคฺคโล อิมานิ อินฺทฺริยานิ อิมสฺมิํ อารมฺมเณ สโมธาเนติ.\csromanspacer{} เตน วุจฺจติ “อินฺทฺริยานิ สโมธาเนตี”ติ.}

\prose{\textbf{โคจรญฺจ ปชานาตี}ติ ยํ ตสฺส อารมฺมณํ, ตํ ตสฺส โคจรํ.\csromanspacer{} ยํ ตสฺส โคจรํ, ตํ ตสฺส อารมฺมณํ.\csromanspacer{} \textbf{ปชานาตี}ติ ปุคฺคโล ปชานนา ปญฺญา.}

\prose{\textbf{สมนฺ}ติ อารมฺมณสฺส อุปฏฺฐานํ สมํ, จิตฺตสฺส อวิกฺเขโป สมํ, จิตฺตสฺส อธิฏฺฐานํ สมํ, จิตฺตสฺส โวทานํ สมํ.\csromanspacer{} \textbf{อตฺโถ}ติ อนวชฺชฏฺโฐ นิเกฺลสฏฺโฐ โวทานฏฺโฐ ปรมฏฺโฐ.\csromanspacer{} \textbf{ปฏิวิชฺฌตี}ติ อารมฺมณสฺส อุปฏฺฐานฏฺฐํ ปฏิวิชฺฌติ, จิตฺตสฺส อวิกฺเขปฏฺฐํ ปฏิวิชฺฌติ, จิตฺตสฺส อธิฏฺฐานฏฺฐํ ปฏิวิชฺฌติ, จิตฺตสฺส โวทานฏฺฐํ ปฏิวิชฺฌติ, เตน วุจฺจติ “สมตฺถญฺจ ปฏิวิชฺฌตี”ติ.}

\prose{\textbf{พลานิ สโมธาเนตี}ติ กถํ พลานิ สโมธาเนติ? อสฺสทฺธิเย อกมฺปิยฏฺเฐน สทฺธาพลํ สโมธาเนติ, โกสชฺเช อกมฺปิยฏฺเฐน วิริยพลํ สโมธาเนติ, ปมาเท อกมฺปิยฏฺเฐน สติพลํ สโมธาเนติ, อุทฺธจฺเจ อกมฺปิยฏฺเฐน สมาธิพลํ สโมธาเนติ, อวิชฺชาย อกมฺปิยฏฺเฐน ปญฺญาพลํ สโมธาเนติ.\csromanspacer{} อยํ ปุคฺคโล อิมานิ พลานิ อิมสฺมิํ อารมฺมเณ}

\vspace{\tipitakaparskip}
\csromancenter{อานาปานสฺสติกถา สโมธาเนติ, เตน วุจฺจติ “พลานิ สโมธาเนตี”ติ.\csromanspacer{}\csromanspacer{}โคจรญฺจ ปชานาตีติ ฯเปฯ.\csromanspacer{} เตน วุจฺจติ “สมตฺถญฺจ ปฏิวิชฺฌตี”ติ.}
\prose{\csromanfolio{179}\textbf{โพชฺฌงฺเค สโมธาเนตี}ติ กถํ โพชฺฌงฺค สโมธาเนติ? อุปฏฺฐานฏฺเฐน สติสมฺโพชฺฌงฺค สโมธาเนติ, ปวิจยฏฺเฐน ธมฺมวิจยสมฺโพชฺฌงฺคํ สโมธาเนติ, ปคฺคหฏฺเฐน วีริยสมฺโพชฺฌงฺคํ สโมธาเนติ, ผรณฏฺเฐน ปีติสมฺโพชฺฌงฺคํ สโมธาเนติ, อุปสมฏฺเฐน ปสฺสทฺธิสมฺโพชฺฌงฺคํ สโมธาเนติ, อวิกฺเขปฏฺเฐน สมาธิสมฺโพชฺฌงฺค สโมธาเนติ, ปฏิสงฺขานฏฺเฐน อุเปกฺขาสมฺโพชฺฌงฺคํ สโมธาเนติ.\csromanspacer{} อยํ ปุคฺคโล อิเม โพชฺฌงฺเค อิมสฺมิํ อารมฺมเณ สโมธาเนติ.\csromanspacer{} เตน วุจฺจติ “โพชฺฌงฺเค สโมธาเนตี”ติ.\csromanspacer{}\csromanspacer{}โคจรญฺจ ปชานาตีติ ฯเปฯ.\csromanspacer{} เตน วุจฺจติ “สมตฺถญฺจ ปฏิวิชฺฌตี”ติ.}

\prose{\textbf{มคฺคํ สโมธาเนตี}ติ กถํ มคฺคํ สโมธาเนติ? ทสฺสนฏฺเฐน สมฺมาทิฏฺฐิํ สโมธาเนติ, อภินิโรปนฏฺเฐน สมฺมาสงฺกปฺปํ สโมธาเนติ, ปริคฺคหฏฺเฐน สมฺมาวาจํ สโมธาเนติ, สมุฏฺฐานฏฺเฐน สมฺมากมฺมนฺตํ สโมธาเนติ, โวทานฏฺเฐน สมฺมา-อาชีวํ สโมธาเนติ, ปคฺคหฏฺเฐน สมฺมาวายามํ สโมธาเนติ, อุปฏฺฐานฏฺเฐน สมฺมาสติํ สโมธาเนติ, อวิกฺเขปฏฺเฐน สมฺมาสมาธิํ สโมธาเนติ.\csromanspacer{} อยํ ปุคฺคโล อิมํ มคฺคํ อิมสฺมิํ อารมฺมเณ สโมธาเนติ, เตน วุจฺจติ “มคฺคํ สโมธาเนตี”ติ.\csromanspacer{}\csromanspacer{}โคจรญฺจ ปชานาตีติ ฯเปฯ.\csromanspacer{} เตน วุจฺจติ “สมตฺถญฺจ ปฏิวิชฺฌตี”ติ.}

\prose{\textbf{ธมฺเม สโมธาเนตี}ติ กถํ ธมฺเม สโมธาเนติ? อาธิปเตยฺยฏฺเฐน อินฺทฺริยานิ สโมธาเนติ, อกมฺปิยฏฺเฐน พลานิ สโมธาเนติ, นิยฺยานฏฺเฐน โพชฺฌงฺเค สโมธาเนติ, เหตุฏฺเฐน มคฺคํ สโมธาเนติ, อุปฏฺฐานฏฺเฐน สติปฏฺฐานํ สโมธาเนติ, ปทหนฏฺเฐน สมฺมปฺปธานํ สโมธาเนติ, อิชฺฌนฏฺเฐน อิทฺธิปาทํ สโมธาเนติ, ตถฏฺเฐน สจฺจํ สโมธาเนติ, อวิกฺเขปฏฺเฐน สมถํ สโมธาเนติ, อนุปสฺสนฏฺเฐน วิปสฺสนํ สโมธาเนติ, เอกรสฏฺเฐน สมถวิปสฺสนํ สโมธาเนติ, อนติวตฺตนฏฺเฐน ยุคนทฺธํ สโมธาเนติ, สํวรฏฺเฐน สีลวิสุทฺธิํ สโมธาเนติ, อวิกฺเขปฏฺเฐน จิตฺตวิสุทฺธิํ สโมธาเนติ, ทสฺสนฏฺเฐน ทิฏฺฐิวิสุทฺธิํ สโมธาเนติ, วิมุตฺตฏฺเฐน วิโมกฺขํ สโมธาเนติ, ปฏิเวธฏฺเฐน วิชฺชํ สโมธาเนติ, ปริจฺจาคฏฺเฐน วิมุตฺติํ สโมธาเนติ, สมุจฺเฉทฏฺเฐน ขเย ญาณํ สโมธาเนติ, ปฏิปฺปสฺสทฺธฏฺเฐน อนุปฺปาเท ญาณํ สโมธาเนติ, \csromanfolio{180}ฉนฺทํ มูลฏฺเฐน สโมธาเนติ, มนสิการํ สมุฏฺฐานฏฺเฐน สโมธาเนติ, ผสฺสํ สโมธานฏฺเฐน สโมธาเนติ, เวทนํ สโมสรณฏฺเฐน สโมธาเนติ, สมาธิํ ปมุขฏฺเฐน สโมธาเนติ, สติํ อาธิปเตยฺยฏฺเฐน สโมธาเนติ, ปญฺญํ ตทุตฺตรฏฺเฐน สโมธาเนติ, วิมุตฺติํ สารฏฺเฐน สโมธาเนติ, อมโตคธํ นิพฺพานํ ปริโยสานฏฺเฐน สโมธาเนติ.\csromanspacer{} อยํ ปุคฺคโล อิเม ธมฺเม อิมสฺมิํ อารมฺมเณ สโมธาเนติ, เตน วุจฺจติ “ธมฺเม สโมธาเนตี”ติ.}

\prose{\textbf{โคจรญฺจ ปชานาตี}ติ ยํ ตสฺส อารมฺมณํ, ตํ ตสฺส โคจรํ.\csromanspacer{} ยํ ตสฺส โคจรํ, ตํ ตสฺส อารมฺมณํ.\csromanspacer{} \textbf{ปชานาตี}ติ ปุคฺคโล ปชานนา ปญฺญา.\csromanspacer{} \textbf{สมนฺ}ติ อารมฺมณสฺส อุปฏฺฐานํ สมํ, จิตฺตสฺส อวิกฺเขโป สมํ, จิตฺตสฺส อธิฏฺฐานํ สมํ, จิตฺตสฺส โวทานํ สมํ.\csromanspacer{} \textbf{อตฺโถ}ติ อนวชฺชฏฺโฐ นิเกฺลสฏฺโฐ โวทานฏฺโฐ ปรมฏฺโฐ.\csromanspacer{} \textbf{ปฏิวิชฺฌตี}ติ อารมฺมณสฺส อุปฏฺฐานฏฺฐํ ปฏิวิชฺฌติ, จิตฺตสฺส อวิกฺเขปฏฺฐํ ปฏิวิชฺฌติ, จิตฺตสฺส อธิฏฺฐานฏฺฐํ ปฏิวิชฺฌติ, จิตฺตสฺส โวทานฏฺฐํ ปฏิวิชฺฌติ, เตน วุจฺจติ “สมตฺถญฺจ ปฏิวิชฺฌตี”ติ. (1)}

\proseitem{169}{กถํ รสฺสํ อสฺสสนฺโต “รสฺสํ อสฺสสามี”ติ ปชานาติ, รสฺสํ ปสฺสสนฺโต “รสฺสํ ปสฺสสามี”ติ ปชานาติ? รสฺสํ อสฺสาสํ อิตฺตรสงฺขาเต อสฺสสติ, รสฺสํ ปสฺสาสํ อิตฺตรสงฺขาเต ปสฺสสติ, รสฺสํ อสฺสาสปสฺสาสํ อิตฺตรสงฺขาเต อสฺสสติปิ ปสฺสสติปิ.\csromanspacer{} รสฺสํ อสฺสาสปสฺสาสํ อิตฺตรสงฺขาเต อสฺสสโตปิ ปสฺสสโตปิ ฉนฺโท อุปฺปชฺชติ, ฉนฺทวเสน ตโต สุขุมตรํ รสฺสํ อสฺสาสํ อิตฺตรสงฺขาเต อสฺสสติ, ฉนฺทวเสน ตโต สุขุมตรํ รสฺสํ ปสฺสาสํ อิตฺตรสงฺขาเต ปสฺสสติ, ฉนฺทวเสน ตโต สุขุมตรํ รสฺสํ อสฺสาสปสฺสาสํ อิตฺตรสงฺขาเต อสฺสสติปิ ปสฺสสติปิ.\csromanspacer{} ฉนฺทวเสน ตโต สุขุมตรํ รสฺสํ อสฺสาสปสฺสาสํ อิตฺตรสงฺขาเต อสฺสสโตปิ ปสฺสสโตปิ ปาโมชฺชํ อุปฺปชฺชติ, ปาโมชฺชวเสน ตโต สุขุมตรํ รสฺสํ อสฺสาสํ อิตฺตรสงฺขาเต อสฺสสติ, ปาโมชฺชวเสน ตโต สุขุมตรํ รสฺสํ ปสฺสาสํ อิตฺตรสงฺขาเต ปสฺสสติ, ปาโมชฺชวเสน ตโต สุขุมตรํ รสฺสํ อสฺสาสปสฺสาสํ อิตฺตรสงฺขาเต อสฺสสติปิ ปสฺสสติปิ.\csromanspacer{} ปาโมชฺชวเสน ตโต สุขุมตรํ รสฺสํ อสฺสาสปสฺสาสํ อิตฺตรสงฺขาเต อสฺสสโตปิ ปสฺสสโตปิ, รสฺสา อสฺสาสปสฺสาสา จิตฺตํ วิวตฺตติ, อุเปกฺขา สณฺฐาติ.\csromanspacer{} อิเมหิ นวหากาเรหิ รสฺสา}

\vspace{\tipitakaparskip}
\csromancenter{อานาปานสฺสติกถา อสฺสาสปสฺสาสา กาโย, อุปฏฺฐานํ สติ, อนุปสฺสนา ญาณํ.\csromanspacer{} \textbf{กาโย} อุปฏฺฐานํ โน สติ, สติ อุปฏฺฐานญฺเจว สติ จ.\csromanspacer{} ตาย สติยา เตน ญาเณน ตํ กายํ อนุปสฺสติ, เตน วุจฺจติ “กาเย กายานุปสฺสนาสติปฏฺฐานพฺ\textbf{หาวนา}”ติ.}
\prose{\csromanfolio{181}\textbf{อนุปสฺสตี}ติ กถํ ตํ กายํ อนุปสฺสติ ฯเปฯ.\csromanspacer{} เอวํ ตํ กายํ อนุปสฺสติ.\csromanspacer{} พฺ\textbf{หาวนา}ติ จตสฺโส พฺ\textbf{หาวนา} ฯเปฯ อาเสวนฏฺเฐน พฺ\textbf{หาวนา}.\csromanspacer{} รสฺสํ อสฺสาสปสฺสาสวเสน จิตฺตสฺส เอกคฺคตํ อวิกฺเขปํ ปชานโต วิทิตา เวทนา อุปฺปชฺชนฺติ ฯเปฯ.\csromanspacer{} รสฺสํ อสฺสาสปสฺสาสวเสน จิตฺตสฺส เอกคฺคตํ อวิกฺเขปํ ปชานนฺโต อินฺทฺริยานิ สโมธาเนติ ฯเปฯ เตน วุจฺจติ “สมตฺถญฺจ ปฏิวิชฺฌตี”ติ. (2)}

\proseitem{170}{กถํ “สพฺพกายปฏิสํเวที อสฺสสิสฺสามี”ติ สิกฺขติ, “สพฺพกายปฏิสํเวที ปสฺสสิสฺสามี”ติ สิกฺขติ? \textbf{กาโย}ติ เทฺว กายา นามกาโย จ รูปกาโย จ.\csromanspacer{} กตโม นามกาโย? เวทนา สญฺญา เจตนา ผสฺโส มนสิกาโร นามญฺจ นามกาโย จ, เย จ วุจฺจนฺติ จิตฺตสงฺขารา, อยํ นามกาโย.\csromanspacer{}\csromanspacer{}กตโม รูปกาโย? จตฺตาโร จ มหาภูตา จตุนฺนญฺจ มหาภูตานํ อุปาทายรูปํ อสฺสาโส จ ปสฺสาโส จ นิมิตฺตญฺจ อุปนิพนฺธนา, เย จ วุจฺจนฺติ กายสงฺขารา, อยํ รูปกาโย.}

\prose{กถํ เต กายา ปฏิวิทิตา โหนฺติ? ทีฆํ อสฺสาสวเสน จิตฺตสฺส เอกคฺคตํ อวิกฺเขปํ ปชานโต สติ อุปฏฺฐิตา โหติ, ตาย สติยา เตน ญาเณน เต กายา ปฏิวิทิตา โหนฺติ.\csromanspacer{} ทีฆํ ปสฺสาสวเสน จิตฺตสฺส เอกคฺคตํ อวิกฺเขปํ ปชานโต สติ อุปฏฺฐิตา โหติ, ตาย สติยา เตน ญาเณน เต กายา ปฏิวิทิตา โหนฺติ.\csromanspacer{} รสฺสํ อสฺสาสวเสน จิตฺตสฺส เอกคฺคตํ อวิกฺเขปํ ปชานโต สติ อุปฏฺฐิตา โหติ, ตาย สติยา เตน ญาเณน เต กายา ปฏิวิทิตา โหนฺติ.\csromanspacer{} รสฺสํ ปสฺสาสวเสน จิตฺตสฺส เอกคฺคตํ อวิกฺเขปํ ปชานโต สติ อุปฏฺฐิตา โหติ, ตาย สติยา เตน ญาเณน เต กายา ปฏิวิทิตา โหนฺติ.}

\prose{อาวชฺชโต เต กายา ปฏิวิทิตา โหนฺติ, ปชานโต เต กายา ปฏิวิทิตา โหนฺติ, ปสฺสโต เต กายา ปฏิวิทิตา โหนฺติ, \csromanfolio{182}ปจฺจเวกฺขโต เต กายา ปฏิวิทิตา โหนฺติ, จิตฺตํ อธิฏฺฐหโต เต กายา ปฏิวิทิตา โหนฺติ, สทฺธาย อธิมุจฺจโต เต กายา ปฏิวิทิตา โหนฺติ, วีริยํ ปคฺคณฺหโต เต กายา ปฏิวิทิตา โหนฺติ, สติํ อุปฏฺฐาปยโต เต กายา ปฏิวิทิตา โหนฺติ, จิตฺตํ สมาทหโต เต กายา ปฏิวิทิตา โหนฺติ, ปญฺญาย ปชานโต เต กายา ปฏิวิทิตา โหนฺติ, อภิญฺเญยฺยํ อภิชานโต เต กายา ปฏิวิทิตา โหนฺติ, ปริญฺเญยฺยํ ปริชานโต เต กายา ปฏิวิทิตา โหนฺติ, ปหาตพฺพํ ปชหโต เต กายา ปฏิวิทิตา โหนฺติ, ภาเวตพฺพํ ภาวยโต เต กายา ปฏิวิทิตา โหนฺติ, สจฺฉิกาตพฺพํ สจฺฉิกโรโต เต กายา ปฏิวิทิตา โหนฺติ.\csromanspacer{} เอวํ เต กายา ปฏิวิทิตา โหนฺติ, สพฺพกายปฏิสํเวที อสฺสาสปสฺสาสา กาโย, อุปฏฺฐานํ สติ, อนุปสฺสนา ญาณํ.\csromanspacer{} กาโย อุปฏฺฐานํ, โน สติ, สติ อุปฏฺฐานญฺเจว สติ จ.\csromanspacer{} ตาย สติยา เตน ญาเณน ตํ กายํ อนุปสฺสติ, เตน วุจฺจติ “กาเย กายานุปสฺสนาสติปฏฺฐานพฺ\textbf{หาวนา}”ติ.}

\prose{\textbf{อนุปสฺสตี}ติ กถํ ตํ กายํ อนุปสฺสติ ฯเปฯ.\csromanspacer{} เอวํ ตํ กายํ อนุปสฺสติ.\csromanspacer{} พฺ\textbf{หาวนา}ติ จตสฺโส พฺ\textbf{หาวนา} ฯเปฯ.\csromanspacer{} อาเสวนฏฺเฐน พฺ\textbf{หาวนา}.}

\prose{สพฺพกายปฏิสํเวที อสฺสาสปสฺสาสานํ\footnote{อสฺสาสปสฺสาสา (สฺยา)} สํวรฏฺเฐน สีลวิสุทฺธิ, อวิกฺเขปฏฺเฐน จิตฺตวิสุทฺธิ, ทสฺสนฏฺเฐน ทิฏฺฐิวิสุทฺธิ.\csromanspacer{} โย ตตฺถ สํวรฏฺโฐ, อยํ อธิสีลสิกฺขา.\csromanspacer{} โย ตตฺถ อวิกฺเขปฏฺโฐ, อยํ อธิจิตฺตสิกฺขา.\csromanspacer{} โย ตตฺถ ทสฺสนฏฺโฐ, อยํ อธิปญฺญาสิกฺขา.\csromanspacer{} อิมา ติสฺโส สิกฺขาโย อาวชฺชนฺโต สิกฺขติ, ชานนฺโต สิกฺขติ, ปสฺสนฺโต สิกฺขติ, ปจฺจเวกฺขนฺโต สิกฺขติ, จิตฺตํ อธิฏฺฐหนฺโต สิกฺขติ, สทฺธาย อธิมุจฺจนฺโต สิกฺขติ, วีริยํ ปคฺคณฺหนฺโต สิกฺขติ, สติํ อุปฏฺฐเปนฺโต สิกฺขติ, จิตฺตํ สมาทหนฺโต สิกฺขติ, ปญฺญาย ปชานนฺโต สิกฺขติ, อภิญฺเญยฺยํ อภิชานนฺโต สิกฺขติ, ปริญฺเญยฺยํ ปริชานนฺโต สิกฺขติ, ปหาตพฺพํ ปชหนฺโต สิกฺขติ, ภาเวตพฺพํ ภาเวนฺโต สิกฺขติ, สิจฺฉิกาตพฺพํ สจฺฉิกโรนฺโต สิกฺขติ.}

\csromanheader{2}{3. อานาปานสฺสติกถา}
\prose{\csromanfolio{183}สพฺพกายปฏิสํเวที อสฺสาสปสฺสาสวเสน จิตฺตสฺส เอกคฺคตํ อวิกฺเขปํ ปชานโต วิทิตา เวทนา อุปฺปชฺชนฺติ ฯเปฯ.\csromanspacer{} สพฺพกายปฏิสํเวที อสฺสาสปสฺสาสวเสน จิตฺตสฺส เอกคฺคตํ อวิกฺเขปํ ปชานนฺโต อินฺทฺริยานิ สโมธาเนติ ฯเปฯ เตน วุจฺจติ “สมตฺถญฺจ ปฏิวิชฺฌตี”ติ. (3)}

\proseitem{171}{กถํ “ปสฺสมฺภยํ กายสงฺขารํ อสฺสสิสฺสามี”ติ สิกฺขติ, “ปสฺสมฺภยํ กายสงฺขารํ ปสฺสสิสฺสามี”ติ สิกฺขติ? กตโม กายสงฺขาโร? ทีฆํ อสฺสาสา กายิกา เอเต ธมฺมา กายปฏิพทฺธา กายสงฺขารา, เต กายสงฺขาเร ปสฺสมฺเภนฺโต นิโรเธนฺโต วูปสเมนฺโต สิกฺขติ.\csromanspacer{} ทีฆํ ปสฺสาสา กายิกา เอเต ธมฺมา กายปฏิพทฺธา กายสงฺขารา, เต กายสงฺขาเร ปสฺสมฺเภนฺโต นิโรเธนฺโต วูปสเมนฺโต สิกฺขติ.\csromanspacer{} รสฺสํ อสฺสาสา.\csromanspacer{} รสฺสํ ปสฺสาสา.\csromanspacer{} สพฺพกายปฏิสํเวที อสฺสาสา.\csromanspacer{} สพฺพกายปฏิสํเวที ปสฺสาสา กายิกา เอเต ธมฺมา กายปฏิพทฺธา กายสงฺขารา, เต กายสงฺขาเร ปสฺสมฺเภนฺโต นิโรเธนฺโต วูปสเมนฺโต สิกฺขติ.}

\prose{ยถารูเปหิ กายสงฺขาเรหิ ยา กายสฺส อานมนา วินมนา สนฺนมนา ปณมนา อิญฺชนา ผนฺทนา จลนา ปกมฺปนา, “ปสฺสมฺภยํ กายสงฺขารํ อสฺสสิสฺสามี”ติ สิกฺขติ, “ปสฺสมฺภยํ กายสงฺขารํ ปสฺสสิสฺสามี”ติ สิกฺขติ.\csromanspacer{} ยถารูเปหิ กายสงฺขาเรหิ ยา กายสฺส น อานมนา น วินมนา น สนฺนมนา น ปณมนา อนิญฺชนา อผนฺทนา อจลนา อกมฺปนา, “สนฺตํ สุขุมํ ปสฺสมฺภยํ กายสงฺขารํ อสฺสสิสฺสามี”ติ สิกฺขติ, “ปสฺสมฺภยํ กายสงฺขารํ ปสฺสสิสฺสามี”ติ สิกฺขติ.}

\prose{อิติ กิร “ปสฺสมฺภยํ กายสงฺขารํ อสฺสสิสฺสามี”ติ สิกฺขติ, “ปสฺสมฺภยํ กายสงฺขารํ ปสฺสสิสฺสามี”ติ สิกฺขติ.\csromanspacer{} เอวํ สนฺเต วาตูปลทฺธิยา จ ปภาวนา น โหติ, อสฺสาสปสฺสาสานญฺจ ปภาวนา น โหติ, อานาปานสฺสติยา จ ปภาวนา น โหติ, อานาปานสฺสติสมาธิสฺส จ ปภาวนา น โหติ, น จ นํ ตํ สมาปตฺติํ ปณฺฑิตา สมาปชฺชนฺติปิ วุฏฺฐหนฺติปิ.}

\prose{อิติ กิร “ปสฺสมฺภยํ กายสงฺขารํ อสฺสสิสฺสามี”ติ สิกฺขติ, “ปสฺสมฺภยํ กายสงฺขารํ ปสฺสสิสฺสามี”ติ สิกฺขติ.\csromanspacer{} เอวํ สนฺเต วาตูปลทฺธิยา จ ปภาวนา โหติ, อสฺสาสปสฺสาสานญฺจ ปภาวนา โหติ, อานาปานสฺสติยา จ ปภาวนา โหติ, อานาปานสฺสติสมาธิสฺส \csromanfolio{184}จ ปพฺ\textbf{หาวนา} โหติ, ตญฺจ นํ สมาปตฺติํ ปณฺฑิตา สมาปชฺชนฺติปิ วุฏฺฐหนฺติปิ.\csromanspacer{} ยถา กถํ วิย เสยฺยถาปิ กํเส อาโกฏิเต ปฐมํ โอฬาริกา สทฺทา ปวตฺตนฺติ, โอฬาริกานํ สทฺทานํ นิมิตฺตํ สุคฺคหิตตฺตา สุมนสิกตตฺตา สูปธาริตตฺตา นิรุทฺเธปิ โอฬาริเก สทฺเท อถ ปจฺฉา สุขุมกา สทฺทา ปวตฺตนฺติ.\csromanspacer{} สุขุมกานํ สทฺทานํ นิมิตฺตํ สุคฺคหิตตฺตา สุมนสิกตตฺตา สูปธาริตตฺตา นิรุทฺเธปิ สุขุมเก สทฺเท อถ ปจฺฉา สุขุมสทฺทนิมิตฺตารมฺมณตาปิ จิตฺตํ ปวตฺตติ.\csromanspacer{} เอวเมวํ ปฐมํ โอฬาริกา อสฺสาสปสฺสาสา ปวตฺตนฺติ, โอฬาริกานํ อสฺสาสปสฺสาสานํ นิมิตฺตํ สุคฺคหิตตฺตา สุมนสิกตตฺตา สูปธาริตตฺตา นิรุทฺเธปิ โอฬาริเก อสฺสาสปสฺสาเส อถ ปจฺฉา สุขุมกา อสฺสาสปสฺสาสา ปวตฺตนฺติ, สุขุมกานํ อสฺสาสปสฺสาสานํ นิมิตฺตํ สุคฺคหิตตฺตา สุมนสิกตตฺตา สูปธาริตตฺตา นิรุทฺเธปิ สุขุมเก อสฺสาสปสฺสาเส อถ ปจฺฉา สุขุมก-อสฺสาสปสฺสาสานํ นิมิตฺตารมฺมณตาปิ จิตฺตํ น วิกฺเขปํ คจฺฉติ.}

\prose{เอวํ สนฺเต วาตูปลทฺธิยา จ ปพฺ\textbf{หาวนา} โหติ, อสฺสาสปสฺสาสานญฺจ ปพฺ\textbf{หาวนา} โหติ, อานาปานสฺสติยา จ ปพฺ\textbf{หาวนา} โหติ, อานาปานสฺสติสมาธิสฺส จ ปพฺ\textbf{หาวนา} โหติ, ตญฺจ นํ สมาปตฺติํ ปณฺฑิตา สมาปชฺชนฺติปิ วุฏฺฐหนฺติปิ.\csromanspacer{} ปสฺสมฺภยํ กายสงฺขารํ อสฺสาสปสฺสาสา กาโย, อุปฏฺฐานํ สติ, อนุปสฺสนา ญาณํ.\csromanspacer{} กาโย อุปฏฺฐานํ, โน สติ, สติ อุปฏฺฐานญฺเจว สติ จ.\csromanspacer{} ตาย สติยา เตน ญาเณน ตํ กายํ อนุปสฺสติ, เตน วุจฺจติ “กาเย กายานุปสฺสนาสติปฏฺฐานพฺ\textbf{หาวนา}”ติ.}

\prose{\textbf{อนุปสฺสตี}ติ กถํ ตํ กายํ อนุปสฺสติ ฯเปฯ.\csromanspacer{} เอวํ ตํ กายํ อนุปสฺสติ.\csromanspacer{} พฺ\textbf{หาวนา}ติ จตสฺโส พฺ\textbf{หาวนา} ฯเปฯ อาเสวนฏฺเฐน พฺ\textbf{หาวนา}.\csromanspacer{} ปสฺสมฺภยํ กายสงฺขารํ อสฺสาสปสฺสาสานํ สํวรฏฺเฐน สีลวิสุทฺธิ, อวิกฺเขปฏฺเฐน จิตฺตวิสุทฺธิ, ทสฺสนฏฺเฐน ทิฏฺฐิวิสุทฺธิ.\csromanspacer{} โย ตตฺถ สํวรฏฺโฐ, อยํ อธิสีลสิกฺขา.\csromanspacer{} โย ตตฺถ อวิกฺเขปฏฺโฐ, อยํ อธิจิตฺตสิกฺขา.\csromanspacer{} โย ตตฺถ ทสฺสนฏฺโฐ, อยํ อธิปญฺญาสิกฺขา.\csromanspacer{} อิมา ติสฺโส สิกฺขาโย อาวชฺชนฺโต สิกฺขติ ฯเปฯ สจฺฉิกาตพฺพํ สจฺฉิกโรนฺโต สิกฺขติ, ปสฺสมฺภยํ กายสงฺขารํ อสฺสาสปสฺสาสวเสน จิตฺตสฺส เอกคฺคตํ อวิกฺเขปํ ปชานโต วิทิตา เวทนา อุปฺปชฺชนฺติ ฯเปฯ.\csromanspacer{} ปสฺสมฺภยํ กายสงฺขารํ อสฺสาสปสฺสาสวเสน จิตฺตสฺส เอกคฺคตํ อวิกฺเขปํ}

\vspace{\tipitakaparskip}
\csromancenter{อานาปานสฺสติกถา ปชานนฺโต อินฺทฺริยานิ สโมธาเนติ ฯเปฯ เตน วุจฺจติ “สมตฺถญฺจ ปฏิวิชฺฌตี”ติ. (4)}
\prose{\csromanfolio{185}อฏฺฐ อนุปสฺสนาญาณานิ, อฏฺฐ จ อุปฏฺฐานานุสฺสติโย, จตฺตาริ สุตฺตนฺติกวตฺถูนิ กาเย กายานุปสฺสนาย.}

\csromanheader{2}{ภาณวาโร.}
\csromansectionrule
\csromanheader{2}{\textbf{ทุติยจตุกฺกนิทฺเทส}}
\proseitem{172}{กถํ “ปีติปฏิสํเวที อสฺสสิสฺสามี”ติ สิกฺขติ, “ปีติปฏิสํเวที ปสฺสสิสฺสามี”ติ สิกฺขติ? กตมา ปีติ? ทีฆํ อสฺสาสวเสน จิตฺตสฺส เอกคฺคตํ อวิกฺเขปํ ปชานโต อุปฺปชฺชติ ปีติ ปาโมชฺชํ, ยา ปีติ ปาโมชฺชํ อาโมทนา ปโมทนา หาโส ปหาโส วิตฺติ โอทคฺยํ อตฺตมนตา จิตฺตสฺส.\csromanspacer{} ทีฆํ ปสฺสาสวเสน จิตฺตสฺส เอกคฺคตํ อวิกฺเขปํ ปชานโต อุปฺปชฺชติ ปีติ ปาโมชฺชํ ฯเปฯ.\csromanspacer{} รสฺสํ อสฺสาสวเสน.\csromanspacer{} รสฺสํ ปสฺสาสวเสน.\csromanspacer{} สพฺพกายปฏิสํเวที อสฺสาสวเสน.\csromanspacer{} สพฺพกายปฏิสํเวที ปสฺสาสวเสน.\csromanspacer{} ปสฺสมฺภยํ กายสงฺขารํ อสฺสาสวเสน.\csromanspacer{} ปสฺสมฺภยํ กายสงฺขารํ ปสฺสาสวเสน จิตฺตสฺส เอกคฺคตํ อวิกฺเขปํ ปชานโต อุปฺปชฺชติ ปีติ ปาโมชฺชํ.\csromanspacer{} ยา ปีติ ปาโมชฺชํ อาโมทนา ปโมทนา หาโส ปหาโส วิตฺติ โอทคฺยํ อตฺตมนตา จิตฺตสฺส, อยํ ปีติ.}

\prose{กถํ สา ปีติ ปฏิวิทิตา โหติ? ทีฆํ อสฺสาสปสฺสาสวเสน จิตฺตสฺส เอกคฺคตํ อวิกฺเขปํ ปชานโต สติ อุปฏฺฐิตา โหติ, ตาย สติยา เตน ญาเณน สา ปีติ ปฏิวิทิตา โหติ.\csromanspacer{} ทีฆํ ปสฺสาสวเสน จิตฺตสฺส เอกคฺคตํ อวิกฺเขปํ ปชานโต สติ อุปฏฺฐิตา โหติ, ตาย สติยา เตน ญาเณน สา ปีติ ปฏิวิทิตา โหติ.\csromanspacer{} รสฺสํ อสฺสาสวเสน ฯเปฯ.\csromanspacer{} รสฺสํ ปสฺสาสวเสน.\csromanspacer{} สพฺพกายปฏิสํเวที อสฺสาสวเสน.\csromanspacer{} สพฺพกายปฏิสํเวที ปสฺสาสวเสน.\csromanspacer{} ปสฺสมฺภยํ กายสงฺขารํ อสฺสาสวเสน.\csromanspacer{} ปสฺสมฺภยํ กายสงฺขารํ ปสฺสาสวเสน จิตฺตสฺส เอกคฺคตํ อวิกฺเขปํ ปชานโต สติ อุปฏฺฐิตา โหติ, ตาย สติยา เตน ญาเณน สา ปีติ ปฏิวิทิตา โหติ.\csromanspacer{} อาวชฺชโต สา ปีติ ปฏิวิทิตา โหติ, ชานโต.\csromanspacer{} ปสฺสโต.\csromanspacer{} ปจฺจเวกฺขโต.\csromanspacer{} จิตฺตํ อธิฏฺฐหโต. \csromanfolio{186}สทฺธาย อธิมุจฺจโต.\csromanspacer{} วีริยํ ปคฺคณฺหโต.\csromanspacer{} สติํ อุปฏฺฐาปยโต.\csromanspacer{} จิตฺตํ สมาทหโต.\csromanspacer{} ปญฺญาย ปชานโต.\csromanspacer{} อภิญฺเญยฺยํ อภิชานโต.\csromanspacer{} ปริญฺเญยฺยํ ปริชานโต.\csromanspacer{} ปหาตพฺพํ ปชหโต.\csromanspacer{} ภาเวตพฺพํ ภาวยโต.\csromanspacer{} สจฺฉิกาตพฺพํ สจฺฉิกโรโต สา ปีติ ปฏิวิทิตา โหติ.\csromanspacer{} เอวํ สา ปีติ ปฏิวิทิตา โหติ.}

\prose{ปีติปฏิสํเวที อสฺสาสปสฺสาสวเสน เวทนา, อุปฏฺฐานํ สติ, อนุปสฺสนา ญาณํ.\csromanspacer{} เวทนา อุปฏฺฐานํ, โน สติ, สติ อุปฏฺฐานญฺเจว สติ จ.ตาย สติยา เตน ญาเณน ตํ เวทนํ อนุปสฺสติ, เตน วุจฺจติ “เวทนาสุ เวทนานุปสฺสนาสติปฏฺฐานภาวนา”ติ.}

\prose{\textbf{อนุปสฺสตี}ติ กถํ ตํ เวทนํ อนุปสฺสติ, อนิจฺจโต อนุปสฺสติ ฯเปฯ.\csromanspacer{} เอวํ ตํ เวทนํ อนุปสฺสติ.\csromanspacer{} \textbf{ภาวนา}ติ จตสฺโส ภาวนา ฯเปฯ อาเสวนฏฺเฐน ภาวนา.\csromanspacer{} ปีติปฏิสํเวที อสฺสาสปสฺสาสานํ สํวรฏฺเฐน สีลวิสุทฺธิ ฯเปฯ ปีติปฏิสํเวที อสฺสาสปสฺสาสวเสน จิตฺตสฺส เอกคฺคตํ อวิกฺเขปํ ปชานโต ฯเปฯ ปชานนฺโต อินฺทฺริยานิ สโมธาเนติ, เตน วุจฺจติ “สมตฺถญฺจ ปฏิวิชฺฌตี”ติ. (1)}

\proseitem{173}{กถํ “สุขปฏิสํเวที อสฺสสิสฺสามี”ติ สิกฺขติ, “สุขปฏิสํเวที ปสฺสสิสฺสามี”ติ สิกฺขติ? \textbf{สุขนฺ}ติ เทฺว สุขานิ กายิกญฺจ สุขํ, เจตสิกญฺจ สุขํ.\csromanspacer{} กตมํ กายิกํ สุขํ? ยํ กายิกํ สาตํ กายิกํ สุขํ กายสมฺผสฺสชํ สาตํ สุขํ เวทยิตํ กายสมฺผสฺสชา สาตา สุขา เวทนา, อิทํ กายิกํ สุขํ.\csromanspacer{} กตมํ เจตสิกํ สุขํ? ยํ เจตสิกํ สาตํ เจตสิกํ สุขํ เจโตสมฺผสฺสชํ สาตํ สุขํ เวทยิตํ เจโตสมฺผสฺสชา สาตา สุขา เวทนา, อิทํ เจตสิกํ สุขํ.}

\prose{กถํ เต สุขา ปฏิวิทิตา โหนฺติ? ทีฆํ อสฺสาสวเสน จิตฺตสฺส เอกคฺคตํ อวิกฺเขปํ ปชานโต สติ อุปฏฺฐิตา โหติ, ตาย สติยา เตน ญาเณน เต สุขา ปฏิวิทิตา โหนฺติ.\csromanspacer{} ทีฆํ ปสฺสาสวเสน จิตฺตสฺส เอกคฺคตํ อวิกฺเขปํ ปชานโต สติ อุปฏฺฐิตา โหติ, ตาย สติยา เตน ญาเณน เต สุขา ปฏิวิทิตา โหนฺติ ฯเปฯ สจฺฉิกาตพฺพํ สจฺฉิกโรโต เต สุขา ปฏิวิทิตา โหนฺติ.\csromanspacer{} เอวํ เต สุขา ปฏิวิทิตา โหนฺติ.\csromanspacer{} สุขปฏิสํเวที อสฺสาสปสฺสาสวเสน เวทนา, อุปฏฺฐานํ สติ, อนุปสฺสนา ญาณํ.\csromanspacer{} เวทนา อุปฏฺฐานํ, โน สติ, สติ อุปฏฺฐานญฺเจว}

\vspace{\tipitakaparskip}
\csromancenter{อานาปานสฺสติกถา สติ จ, ตาย สติยา เตน ญาเณน ตํ เวทนํ อนุปสฺสติ, เตน วุจฺจติ “เวทนาสุ เวทนานุปสฺสนาสติปฏฺฐานภาวนา”ติ.}
\prose{\csromanfolio{187}\textbf{อนุปสฺสตี}ติ กถํ ตํ เวทนํ อนุปสฺสติ, อนิจฺจโต อนุปสฺสติ ฯเปฯ.\csromanspacer{} เอวํ ตํ เวทนํ อนุปสฺสติ.\csromanspacer{} \textbf{ภาวนา}ติ จตสฺโส ภาวนา ฯเปฯ อาเสวนฏฺเฐน ภาวนา.\csromanspacer{} สุขปฏิสํเวที อสฺสาสปสฺสาสานํ สํวรฏฺเฐน สีลวิสุทฺธิ ฯเปฯ.\csromanspacer{} สุขปฏิสํเวที อสฺสาสปสฺสาสวเสน จิตฺตสฺส เอกคฺคตํ อวิกฺเขปํ ปชานโต ฯเปฯ.\csromanspacer{} ปชานนฺโต อินฺทฺริยานิ สโมธาเนติ, เตน วุจฺจติ “สมตฺถญฺจ ปฏิวิชฺฌตี”ติ. (2)}

\proseitem{174}{กถํ “จิตฺตสงฺขารปฏิสํเวที อสฺสสิสฺสามี”ติ สิกฺขติ, “จิตฺตสงฺขารปฏิสํเวที ปสฺสสิสฺสามี”ติ สิกฺขติ? กตโม จิตฺตสงฺขาโร? ทีฆํ อสฺสาสวเสน สญฺญา จ เวทนา จ เจตสิกา, เอเต ธมฺมา จิตฺตปฏิพทฺธา จิตฺตสงฺขารา.\csromanspacer{} ทีฆํ ปสฺสาสวเสน สญฺญา จ เวทนา จ เจตสิกา, เอเต ธมฺมา จิตฺตปฏิพทฺธา จิตฺตสงฺขารา ฯเปฯ.\csromanspacer{} สุขปฏิสํเวที อสฺสาสวเสน.\csromanspacer{} สุขปฏิสํเวที ปสฺสาสวเสน สญฺญา จ เวทนา จ เจตสิกา, เอเต ธมฺมา จิตฺตปฏิพทฺธา จิตฺตสงฺขารา, อยํ จิตฺตสงฺขาโร.}

\prose{กถํ เต จิตฺตสงฺขารา ปฏิวิทิตา โหนฺติ? ทีฆํ อสฺสาสวเสน จิตฺตสฺส เอกคฺคตํ อวิกฺเขปํ ปชานโต สติ อุปฏฺฐิตา โหติ, ตาย สติยา เตน ญาเณน เต จิตฺตสงฺขารา ปฏิวิทิตา โหนฺติ.\csromanspacer{} ทีฆํ ปสฺสาสวเสน จิตฺตสฺส เอกคฺคตํ อวิกฺเขปํ ปชานโต สติ อุปฏฺฐิตา โหติ, ตาย สติยา เตน ญาเณน เต จิตฺตสงฺขารา ปฏิวิทิตา โหนฺติ ฯเปฯ.\csromanspacer{} สจฺฉิกาตพฺพํ สจฺฉิกโรโต เต จิตฺตสงฺขารา ปฏิวิทิตา โหนฺติ.\csromanspacer{} เอวํ เต จิตฺตสงฺขารา ปฏิวิทิตา โหนฺติ.\csromanspacer{} จิตฺตสงฺขารปฏิสํเวที อสฺสาสปสฺสาสวเสน เวทนา, อุปฏฺฐานํ สติ, อนุปสฺสนา ญาณํ.\csromanspacer{} เวทนา อุปฏฺฐานํ, โน สติ, สติ อุปฏฺฐานญฺเจว สติ จ, ตาย สติยา เตน ญาเณน ตํ เวทนํ อนุปสฺสติ, เตน วุจฺจติ “เวทนาสุ เวทนานุปสฺสนาสติปฏฺฐานภาวนา”ติ.}

\prose{\textbf{อนุปสฺสตี}ติ กถํ ตํ เวทนํ อนุปสฺสติ, อนิจฺจโต อนุปสฺสติ ฯเปฯ.\csromanspacer{} เอวํ ตํ เวทนํ อนุปสฺสติ.\csromanspacer{} \textbf{ภาวนา}ติ จตสฺโส ภาวนา ฯเปฯ อาเสวนฏฺเฐน ภาวนา.\csromanspacer{} จิตฺตสงฺขารปฏิสํเวที อสฺสาสปสฺสาสานํ สํวรฏฺเฐน สีลวิสุทฺธิ ฯเปฯ.\csromanspacer{} จิตฺตสงฺขารปฏิสํเวที อสฺสาสปสฺสาสวเสน \csromanfolio{188}จิตฺตสฺส เอกคฺคตํ อวิกฺเขปํ ปชานโต ฯเปฯ ปชานนฺโต อินฺทฺริยานิ สโมธาเนติ, เตน วุจฺจติ “สมตฺถญฺจ ปฏิวิชฺฌตี”ติ. (3)}

\proseitem{175}{กถํ “ปสฺสมฺภยํ จิตฺตสงฺขารํ อสฺสสิสฺสามี”ติ สิกฺขติ, “ปสฺสมฺภยํ จิตฺตสงฺขารํ ปสฺสสิสฺสามี”ติ สิกฺขติ? กตโม จิตฺตสงฺขาโร? ทีฆํ อสฺสาสวเสน สญฺญา จ เวทนา จ เจตสิกา, เอเต ธมฺมา จิตฺตปฏิพทฺธา จิตฺตสงฺขารา, เต จิตฺตสงฺขาเร ปสฺสมฺเภนฺโต นิโรเธนฺโต วูปสเมนฺโต สิกฺขติ.\csromanspacer{} ทีฆํ ปสฺสาสวเสน สญฺญา จ เวทนา จ เจตสิกา, เอเต ธมฺมา จิตฺตปฏิพทฺธา จิตฺตสงฺขารา, เต จิตฺตสงฺขาเร ปสฺสมฺเภนฺโต นิโรเธนฺโต วูปสเมนฺโต สิกฺขติ.\csromanspacer{} จิตฺตสงฺขารปฏิสํเวที อสฺสาสวเสน.\csromanspacer{} จิตฺตสงฺขารปฏิสํเวที ปสฺสาสวเสน สญฺญา จ เวทนา จ เจตสิกา, เอเต ธมฺมา จิตฺตปฏิพทฺธา จิตฺตสงฺขารา, เต จิตฺตสงฺขาเร ปสฺสมฺเภนฺโต นิโรเธนฺโต วูปสเมนฺโต สิกฺขติ.\csromanspacer{} ปสฺสมฺภยํ จิตฺตสงฺขารํ อสฺสาสปสฺสาสวเสน เวทนา, อุปฏฺฐานํ สติ, อนุปสฺสนา ญาณํ.\csromanspacer{} เวทนา อุปฏฺฐานํ, โน สติ, สติ อุปฏฺฐานญฺเจว สติ จ.\csromanspacer{} ตาย สติยา เตน ญาเณน ตํ เวทนํ อนุปสฺสติ, เตน วุจฺจติ “เวทนาสุ เวทนานุปสฺสนาสติปฏฺฐาน\textbf{ภาวนา}”ติ.}

\prose{\textbf{อนุปสฺสตี}ติ กถํ ตํ เวทนํ อนุปสฺสติ ฯเปฯ.\csromanspacer{} เอวํ ตํ เวทนํ อนุปสฺสติ, \textbf{ภาวนา}ติ จตสฺโส \textbf{ภาวนา} ฯเปฯ อาเสวนฏฺเฐน \textbf{ภาวนา}.\csromanspacer{} ปสฺสมฺภยํ จิตฺตสงฺขารํ อสฺสาสปสฺสาสานํ สํวรฏฺเฐน สีลวิสุทฺธิ ฯเปฯ.\csromanspacer{} ปสฺสมฺภยํ จิตฺตสงฺขารํ อสฺสาสปสฺสาสวเสน จิตฺตสฺส เอกคฺคตํ อวิกฺเขปํ ปชานโต ฯเปฯ.\csromanspacer{} ปชานนฺโต อินฺทฺริยานิ สโมธาเนติ, เตน วุจฺจติ “สมตฺถญฺจ ปฏิวิชฺฌตี”ติ. (4)}

\prose{อฏฺฐ อนุปสฺสนาญาณานิ, อฏฺฐ จ อุปฏฺฐานานุสฺสติโย, จตฺตาริ สุตฺตนฺติกวตฺถูนิ เวทนาสุ เวทนานุปสฺสนาย.}

\csromanheader{2}{\textbf{ตติยจตุกฺกนิทฺเทส}}
\proseitem{176}{กถํ “จิตฺตปฏิสํเวที อสฺสสิสฺสามี”ติ สิกฺขติ, “จิตฺตปฏิสํเวที ปสฺสสิสฺสามี”ติ สิกฺขติ? กตมํ ตํ จิตฺตํ? ทีฆํ อสฺสาสวเสน วิญฺญาณํ จิตฺตํ.\csromanspacer{} ยํ จิตฺตํ มโน มานสํ หทยํ ปณฺฑรํ มโน มนายตนํ มนินฺทฺริยํ วิญฺญาณํ วิญฺญาณกฺขนฺโธ ตชฺชา มโนวิญฺญาณธาตุ.\csromanspacer{} ทีฆํ ปสฺสาสวเสน ฯเปฯ.\csromanspacer{} ปสฺสมฺภยํ จิตฺตสงฺขารํ อสฺสาสวเสน.\csromanspacer{} ปสฺสมฺภยํ จิตฺตสงฺขารํ}

\vspace{\tipitakaparskip}
\csromancenter{อานาปานสฺสติกถา ปสฺสาสวเสน วิญฺญาณํ จิตฺตํ.\csromanspacer{} ยํ จิตฺตํ มโน มานสํ หทยํ ปณฺฑรํ มโน มนายตนํ มนินฺทฺริยํ วิญฺญาณํ วิญฺญาณกฺขนฺโธ ตชฺชา มโนวิญฺญาณธาตุ, อิทํ จิตฺตํ.}
\prose{\csromanfolio{189}กถํ ตํ จิตฺตํ ปฏิวิทิตํ โหติ? ทีฆํ อสฺสาสวเสน จิตฺตสฺส เอกคฺคตํ อวิกฺเขปํ ปชานโต สติ อุปฏฺฐิตา โหติ, ตาย สติยา เตน ญาเณน ตํ จิตฺตํ ปฏิวิทิตํ โหติ.\csromanspacer{} ทีฆํ ปสฺสาสวเสน จิตฺตสฺส เอกคฺคตํ อวิกฺเขปํ ปชานโต สติ อุปฏฺฐิตา โหติ, ตาย สติยา เตน ญาเณน ตํ จิตฺตํ ปฏิวิทิตํ โหติ ฯเปฯ สจฺฉิกาตพฺพํ สจฺฉิกโรโต ตํ จิตฺตํ ปฏิวิทิตํ โหติ.\csromanspacer{} เอวํ ตํ จิตฺตํ ปฏิวิทิตํ โหติ.\csromanspacer{} จิตฺตปฏิสํเวที อสฺสาสปสฺสาสวเสน วิญฺญาณํ จิตฺตํ, อุปฏฺฐานํ สติ, อนุปสฺสนา ญาณํ.\csromanspacer{} จิตฺตํ อุปฏฺฐานํ, โน สติ, สติ อุปฏฺฐานญฺเจว สติ จ.\csromanspacer{} ตาย สติยา เตน ญาเณน ตํ จิตฺตํ อนุปสฺสติ, เตน วุจฺจติ “จิตฺเต จิตฺตานุปสฺสนาสติปฏฺฐานพฺ\textbf{หาวนา}”ติ.}

\prose{\textbf{อนุปสฺสตี}ติ กถํ ตํ จิตฺตํ อนุปสฺสติ ฯเปฯ.\csromanspacer{} เอวํ ตํ จิตฺตํ อนุปสฺสติ.\csromanspacer{} พฺ\textbf{หาวนา}ติ จตสฺโส พฺ\textbf{หาวนา} ฯเปฯ อาเสวนฏฺเฐน พฺ\textbf{หาวนา}.\csromanspacer{} จิตฺตปฏิสํเวที อสฺสาสปสฺสาสานํ สํวรฏฺเฐน สีลวิสุทฺธิ ฯเปฯ.\csromanspacer{} จิตฺตปฏิสํเวที อสฺสาสปสฺสาสวเสน จิตฺตสฺส เอกคฺคตํ อวิกฺเขปํ ปชานโต ฯเปฯ ปชานนฺโต อินฺทฺริยานิ สโมธาเนติ, เตน วุจฺจติ “สมตฺถญฺจ ปฏิวิชฺฌตี”ติ. (1)}

\proseitem{177}{กถํ “อภิปฺปโมทยํ จิตฺตํ อสฺสสิสฺสามี”ติ สิกฺขติ, “อภิปฺปโมทยํ จิตฺตํ ปสฺสสิสฺสามี”ติ สิกฺขติ? กตโม จิตฺตสฺส อภิปฺปโมโท? ทีฆํ อสฺสาสวเสน จิตฺตสฺส เอกคฺคตํ อวิกฺเขปํ ปชานโต อุปฺปชฺชติ จิตฺตสฺส อภิปฺปโมโท, ยา จิตฺตสฺส อาโมทนา ปโมทนา หาโส ปหาโส วิตฺติ โอทคฺยํ อตฺตมนตา จิตฺตสฺส.\csromanspacer{} ทีฆํ อสฺสาสปสฺสาสวเสน จิตฺตสฺส เอกคฺคตํ อวิกฺเขปํ ปชานโต อุปฺปชฺชติ จิตฺตสฺส อภิปฺปโมโท, ยา จิตฺตสฺส อาโมทนา ปโมทนา หาโส ปหาโส วิตฺติ โอทคฺยํ อตฺตมนตา จิตฺตสฺส ฯเปฯ.\csromanspacer{} จิตฺตปฏิสํเวที อสฺสาสวเสน.\csromanspacer{} จิตฺตปฏิสํเวที ปสฺสาสวเสน จิตฺตสฺส เอกคฺคตํ อวิกฺเขปํ ปชานโต อุปฺปชฺชติ จิตฺตสฺส อภิปฺปโมโท, ยา จิตฺตสฺส อาโมทนา ปโมทนา หาโส ปหาโส วิตฺติ โอทคฺยํ อตฺตมนตา จิตฺตสฺส, อยํ จิตฺตสฺส อภิปฺปโมโท.\csromanspacer{} อภิปฺปโมทยํ จิตฺตํ อสฺสาสปสฺสาสวเสน วิญฺญาณํ จิตฺตํ, อุปฏฺฐานํ \csromanfolio{190}สติ, อนุปสฺสนา ญาณํ.\csromanspacer{} จิตฺตํ อุปฏฺฐานํ, โน สติ, สติ อุปฏฺฐานญฺเจว สติ จ.\csromanspacer{} ตาย สติยา เตน ญาเณน ตํ จิตฺตํ อนุปสฺสติ, เตน วุจฺจติ “จิตฺเต จิตฺตานุปสฺสนาสติปฏฺฐานพฺ\textbf{หาวนา}”ติ.}

\prose{\textbf{อนุปสฺสตี}ติ กถํ ตํ จิตฺตํ อนุปสฺสติ ฯเปฯ.\csromanspacer{} เอวํ ตํ จิตฺตํ อนุปสฺสตีติ.\csromanspacer{} พฺ\textbf{หาวนา}ติ จตสฺโส พฺ\textbf{หาวนา} ฯเปฯ อาเสวนฏฺเฐน พฺ\textbf{หาวนา}.\csromanspacer{} อภิปฺปโมทยํ จิตฺตํ อสฺสาสปสฺสาสานํ สํวรฏฺเฐน สีลวิสุทฺธิ ฯเปฯ.\csromanspacer{} อภิปฺปโมทยํ จิตฺตํ อสฺสาสปสฺสาสวเสน จิตฺตสฺส เอกคฺคตํ อวิกฺเขปํ ปชานโต ฯเปฯ ปชานนฺโต อินฺทฺริยานิ สโมธาเนติ, เตน วุจฺจติ “สมตฺถญฺจ ปฏิวิชฺฌตี”ติ. (2)}

\proseitem{178}{กถํ “สมาทหํ จิตฺตํ อสฺสสิสฺสามี”ติ สิกฺขติ, “สมาทหํ จิตฺตํ ปสฺสสิสฺสามี”ติ สิกฺขติ? กตโม สมาธิ\footnote{กตมญฺจ สมาธินฺทฺริยํ (สฺยา)}? ทีฆํ อสฺสาสวเสน จิตฺตสฺส เอกคฺคตา อวิกฺเขโป สมาธิ, ยา จิตฺตสฺส ฐิติ สณฺฐิติ อวฏฺฐิติ\footnote{อธิฏฺฐิติ (สฺยา)} อวิสาหาโร อวิกฺเขโป อวิสาหฏมานสตา\footnote{อวิสาหตมานสตา (สฺยา, ก)} สมโถ สมาธินฺทฺริยํ สมาธิพลํ สมฺมาสมาธิ.\csromanspacer{} ทีฆํ ปสฺสาสวเสน จิตฺตสฺส เอกคฺคตา อวิกฺเขโป สมาธิ ฯเปฯ สมาทหํ จิตฺตํ อสฺสาสวเสน ฯเปฯ สมาทหํ จิตฺตํ ปสฺสาสวเสน จิตฺตสฺส เอกคฺคตา อวิกฺเขโป สมาธิ, ยา จิตฺตสฺส ฐิติ สณฺฐิติ อวฏฺฐิติ อวิสาหาโร อวิกฺเขโป อวิสาหฏมานสตา สมโถ สมาธินฺทฺริยํ สมาธิพลํ สมฺมาสมาธิ, อยํ สมาธิ.\csromanspacer{} สมาทหํ จิตฺตํ อสฺสาสปสฺสาสวเสน วิญฺญาณํ จิตฺตํ, อุปฏฺฐานํ สติ, อนุปสฺสนา ญาณํ.\csromanspacer{} จิตฺตํ อุปฏฺฐานํ, โน สติ.\csromanspacer{} สติ อุปฏฺฐานญฺเจว สติ จ.\csromanspacer{} ตาย สติยา เตน ญาเณน ตํ จิตฺตํ อนุปสฺสติ, เตน วุจฺจติ “จิตฺเต จิตฺตานุปสฺสนาสติปฏฺฐานพฺ\textbf{หาวนา}”ติ.}

\prose{\textbf{อนุปสฺสตี}ติ กถํ ตํ จิตฺตํ อนุปสฺสติ ฯเปฯ.\csromanspacer{} เอวํ ตํ จิตฺตํ อนุปสฺสติ.\csromanspacer{} พฺ\textbf{หาวนา}ติ จตสฺโส พฺ\textbf{หาวนา} ฯเปฯ อาเสวนฏฺเฐน พฺ\textbf{หาวนา}.\csromanspacer{} สมาทหํ จิตฺตํ อสฺสาสปสฺสาสานํ สํวรฏฺเฐน สีลวิสุทฺธิ ฯเปฯ สมาทหํ จิตฺตํ อสฺสาสปสฺสาสวเสน จิตฺตสฺส เอกคฺคตํ อวิกฺเขปํ ปชานโต ฯเปฯ ปชานนฺโต อินฺทฺริยานิ สโมธาเนติ, เตน วุจฺจติ “สมตฺถญฺจ ปฏิวิชฺฌตี”ติ. (3)}

\csromanheader{2}{3. อานาปานสฺสติกถา}
\proseitem{179}{\csromanfolio{191}กถํ “วิโมจยํ จิตฺตํ อสฺสสิสฺสามี”ติ สิกฺขติ, “วิโมจยํ จิตฺตํ ปสฺสสิสฺสามี”ติ สิกฺขติ? “ราคโต วิโมจยํ จิตฺตํ อสฺสสิสฺสามี”ติ สิกฺขติ, “ราคโต วิโมจยํ จิตฺตํ ปสฺสสิสฺสามี”ติ สิกฺขติ. “โทสโต วิโมจยํ จิตฺตํ อสฺสสิสฺสามี”ติ สิกฺขติ, “โทสโต วิโมจยํ จิตฺตํ ปสฺสสิสฺสามี”ติ สิกฺขติ. “โมหโต วิโมจยํ จิตฺตํ อสฺสสิสฺสามี”ติ สิกฺขติ ฯเปฯ “มานโต วิโมจยํ จิตฺตํ.\csromanspacer{} ทิฏฺฐิยา วิโมจยํ จิตฺตํ.\csromanspacer{} วิจิกิจฺฉาย วิโมจยํ จิตฺตํ.\csromanspacer{} ถินโต วิโมจยํ จิตฺตํ.\csromanspacer{} อุทฺธจฺจโต วิโมจยํ จิตฺตํ.\csromanspacer{} อหิริกโต วิโมจยํ จิตฺตํ.\csromanspacer{} อโนตฺตปฺปโต วิโมจยํ จิตฺตํ อสฺสสิสฺสามี”ติ สิกฺขติ, “อโนตฺตปฺปโต วิโมจยํ จิตฺตํ ปสฺสสิสฺสามี”ติ สิกฺขติ.\csromanspacer{} วิโมจยํ จิตฺตํ อสฺสาสปสฺสาสวเสน วิญฺญาณํ จิตฺตํ, อุปฏฺฐานํ สติ ฯเปฯ.}

\prose{\textbf{อนุปสฺสตี}ติ กถํ ตํ จิตฺตํ อนุปสฺสติ ฯเปฯ.\csromanspacer{} เอวํ ตํ จิตฺตํ อนุปสฺสติ.\csromanspacer{} พฺ\textbf{หาวนา}ติ จตสฺโส พฺ\textbf{หาวนา} ฯเปฯ อาเสวนฏฺเฐน พฺ\textbf{หาวนา}.\csromanspacer{} วิโมจยํ จิตฺตํ อสฺสาสปสฺสาสานํ สํวรฏฺเฐน สีลวิสุทฺธิ ฯเปฯ.\csromanspacer{} วิโมจยํ จิตฺตํ อสฺสาสปสฺสาสวเสน จิตฺตสฺส เอกคฺคตํ อวิกฺเขปํ ปชานโต ฯเปฯ ปชานนฺโต อินฺทฺริยานิ สโมธาเนติ, เตน วุจฺจติ “สมตฺถญฺจ ปฏิวิชฺฌตี”ติ. (4)}

\prose{อฏฺฐ อนุปสฺสนาญาณานิ, อฏฺฐ จ อุปฏฺฐานานุสฺสติโย, จตฺตาริ สุตฺตนฺติกวตฺถูนิ จิตฺเต จิตฺตานุปสฺสนาย.}

\csromanheader{2}{\textbf{จตุตฺถจตุกฺกนิทฺเทส}}
\proseitem{180}{กถํ “อนิจฺจานุปสฺสี อสฺสสิสฺสามี”ติ สิกฺขติ, “อนิจฺจานุปสฺสี ปสฺสสิสฺสามี”ติ สิกฺขติ? \textbf{อนิจฺจนฺ}ติ กิํ อนิจฺจํ? ปญฺจกฺขนฺธา อนิจฺจา.\csromanspacer{} เกนฏฺเฐน อนิจฺจา? อุปฺปาทวยฏฺเฐน อนิจฺจา.\csromanspacer{} ปญฺจนฺนํ ขนฺธานํ อุทยํ ปสฺสนฺโต กติ ลกฺขณานิ ปสฺสติ, วยํ ปสฺสนฺโต กติ ลกฺขณานิ ปสฺสติ, อุทยพฺพยํ ปสฺสนฺโต กติ ลกฺขณานิ ปสฺสติ? ปญฺจนฺนํ ขนฺธานํ อุทยํ ปสฺสนฺโต ปญฺจวีสติ ลกฺขณานิ ปสฺสติ, วยํ ปสฺสนฺโต ปญฺจวีสติ ลกฺขณานิ ปสฺสติ, ปญฺจนฺนํ ขนฺธานํ อุทยพฺพยํ ปสฺสนฺโต อิมานิ ปญฺญาส ลกฺขณานิ ปสฺสติ.}

\prose{“รูเป อนิจฺจานุปสฺสี อสฺสสิสฺสามี”ติ สิกฺขติ, “รูเป อนิจฺจานุปสฺสี ปสฺสสิสฺสามี”ติ สิกฺขติ, “เวทนาย ฯเปฯ.\csromanspacer{} สญฺญาย.\csromanspacer{} สงฺขาเรสุ.\csromanspacer{} วิญฺญาเณ.\csromanspacer{} จกฺขุสฺมิํ ฯเปฯ.\csromanspacer{} ชรามรเณ อนิจฺจานุปสฺสี อสฺสสิสฺสามี”ติ สิกฺขติ, \csromanfolio{192}“ชรามรเณ อนิจฺจานุปสฺสี ปสฺสสิสฺสามี”ติ สิกฺขติ.\csromanspacer{} อนิจฺจานุปสฺสี อสฺสาสปสฺสาสวเสน ธมฺมา, อุปฏฺฐานํ สติ, อนุปสฺสนา ญาณํ.\csromanspacer{} ธมฺมา อุปฏฺฐานํ, โน สติ, สติ อุปฏฺฐานญฺเจว สติ จ.\csromanspacer{} ตาย สติยา เตน ญาเณน เต ธมฺเม อนุปสฺสติ, เตน วุจฺจติ “ธมฺเมสุ ธมฺมานุปสฺสนาสติปฏฺฐานพฺ\textbf{หาวนา}”ติ.}

\prose{\textbf{อนุปสฺสตี}ติ กถํ เต ธมฺเม อนุปสฺสติ ฯเปฯ.\csromanspacer{} เอวํ เต ธมฺเม อนุปสฺสติ.\csromanspacer{} พฺ\textbf{หาวนา}ติ จตสฺโส พฺ\textbf{หาวนา} ฯเปฯ อาเสวนฏฺเฐน พฺ\textbf{หาวนา} อนิจฺจานุปสฺสี อสฺสาสปสฺสาสานํ สํวรฏฺเฐน สีลวิสุทฺธิ ฯเปฯ.\csromanspacer{} อนิจฺจานุปสฺสี อสฺสาสปสฺสาสวเสน จิตฺตสฺส เอกคฺคตํ อวิกฺเขปํ ปชานโต ฯเปฯ.\csromanspacer{} ปชานนฺโต อินฺทฺริยานิ สโมธาเนติ, เตน วุจฺจติ “สมตฺถญฺจ ปฏิวิชฺฌตี”ติ. (1)}

\prose{กถํ “วิราคานุปสฺสี อสฺสสิสฺสามี”ติ สิกฺขติ, “วิราคานุปสฺสี ปสฺสสิสฺสามี”ติ สิกฺขติ? รูเป อาทีนวํ ทิสฺวา รูปวิราเค ฉนฺทชาโต โหติ สทฺธาธิมุตฺโต, จิตฺตญฺจสฺส สฺวาธิฏฺฐิตํ. “รูเป วิราคานุปสฺสี อสฺสสิสฺสามี”ติ สิกฺขติ, “รูเป วิราคานุปสฺสี ปสฺสสิสฺสามี”ติ สิกฺขติ, เวทนาย ฯเปฯ สญฺญาย.\csromanspacer{} สงฺขาเรสุ.\csromanspacer{} วิญฺญาเณ.\csromanspacer{} จกฺขุสฺมิํ ฯเปฯ.\csromanspacer{} ชรามรเณ อาทีนวํ ทิสฺวา ชรามรณวิราเค ฉนฺทชาโต โหติ สทฺธาธิมุตฺโต, จิตฺตญฺจสฺส สฺวาธิฏฺฐิตํ. “ชรามรเณ วิราคานุปสฺสี อสฺสสิสฺสามี”ติ สิกฺขติ, “ชรามรเณ วิราคานุปสฺสี ปสฺสสิสฺสามี”ติ สิกฺขติ.\csromanspacer{} วิราคานุปสฺสี อสฺสาสปสฺสาสวเสน ธมฺมา, อุปฏฺฐานํ สติ, อนุปสฺสนา ญาณํ.\csromanspacer{} ธมฺมา อุปฏฺฐานํ, โน สติ, สติ อุปฏฺฐานญฺเจว สติ จ.\csromanspacer{} ตาย สติยา เตน ญาเณน เต ธมฺเม อนุปสฺสติ, เตน วุจฺจติ “ธมฺเมสุ ธมฺมานุปสฺสนาสติปฏฺฐานพฺ\textbf{หาวนา}”ติ.}

\prose{\textbf{อนุปสฺสตี}ติ กถํ เต ธมฺเม อนุปสฺสติ ฯเปฯ.\csromanspacer{} เอวํ เต ธมฺเม อนุปสฺสติ.\csromanspacer{} พฺ\textbf{หาวนา}ติ จตสฺโส พฺ\textbf{หาวนา} ฯเปฯ อาเสวนฏฺเฐน พฺ\textbf{หาวนา}, วิราคานุปสฺสี อสฺสาสปสฺสาสานํ สํวรฏฺเฐน สีลวิสุทฺธิ ฯเปฯ.\csromanspacer{} วิราคานุปสฺสี อสฺสาสปสฺสาสวเสน จิตฺตสฺส เอกคฺคตํ อวิกฺเขปํ ปชานโต ฯเปฯ ปชานนฺโต อินฺทฺริยานิ สโมธาเนติ, เตน วุจฺจติ “สมตฺถญฺจ ปฏิวิชฺฌตี”ติ. (2)}

\csromanheader{2}{3. อานาปานสฺสติกถา}
\prose{\csromanfolio{193}กถํ “นิโรธานุปสฺสี อสฺสสิสฺสามี”ติ สิกฺขติ, “นิโรธานุปสฺสี ปสฺสสิสฺสามี”ติ สิกฺขติ? รูเป อาทีนวํ ทิสฺวา รูปนิโรเธ ฉนฺทชาโต โหติ สทฺธาธิมุตฺโต, จิตฺตญฺจสฺส สฺวาธิฏฺฐิตํ. “รูเป นิโรธานุปสฺสี อสฺสสิสฺสามี”ติ สิกฺขติ, “รูเป นิโรธานุปสฺสี ปสฺสสิสฺสามี”ติ สิกฺขติ, เวทนาย ฯเปฯ สญฺญาย.\csromanspacer{} สงฺขาเรสุ.\csromanspacer{} วิญฺญาเณ.\csromanspacer{} จกฺขุสฺมิํ ฯเปฯ.\csromanspacer{} ชรามรเณ อาทีนวํ ทิสฺวา ชรามรณนิโรเธ ฉนฺทชาโต โหติ สทฺธาธิมุตฺโต, จิตฺตญฺจสฺส สฺวาธิฏฺฐิตํ. “ชรามรเณ นิโรธานุปสฺสี อสฺสสิสฺสามี”ติ สิกฺขติ, “ชรามรเณ นิโรธานุปสฺสี ปสฺสสิสฺสามี”ติ สิกฺขติ.}

\proseitem{181}{กติหากาเรหิ อวิชฺชาย อาทีนโว โหติ, กติหากาเรหิ อวิชฺชา นิรุชฺฌติ? ปญฺจหากาเรหิ อวิชฺชาย อาทีนโว โหติ, อฏฺฐหากาเรหิ อวิชฺชา นิรุชฺฌติ.}

\prose{กตเมหิ ปญฺจหากาเรหิ อวิชฺชาย อาทีนโว โหติ? อนิจฺจฏฺเฐน อวิชฺชาย อาทีนโว โหติ, ทุกฺขฏฺเฐน อวิชฺชาย อาทีนโว โหติ, อนตฺตฏฺเฐน อวิชฺชาย อาทีนโว โหติ, สนฺตาปฏฺเฐน อวิชฺชาย อาทีนโว โหติ, วิปริณามฏฺเฐน อวิชฺชาย อาทีนโว โหติ, อิเมหิ ปญฺจหากาเรหิ อวิชฺชาย อาทีนโว โหติ.}

\prose{กตเมหิ อฏฺฐหากาเรหิ อวิชฺชา นิรุชฺฌติ? นิทานนิโรเธน อวิชฺชา นิรุชฺฌติ, สมุทยนิโรเธน อวิชฺชา นิรุชฺฌติ, ชาตินิโรเธน อวิชฺชา นิรุชฺฌติ, ปภวนิโรเธน\footnote{อาหารนิโรเธน (สฺยา)} อวิชฺชา นิรุชฺฌติ, เหตุนิโรเธน อวิชฺชา นิรุชฺฌติ, ปจฺจยนิโรเธน อวิชฺชา นิรุชฺฌติ, ญาณุปฺปาเทน อวิชฺชา นิรุชฺฌติ, นิโรธุปฏฺฐาเนน อวิชฺชา นิรุชฺฌติ, อิเมหิ อฏฺฐหากาเรหิ อวิชฺชา นิรุชฺฌติ.\csromanspacer{} อิเมหิ ปญฺจหากาเรหิ อวิชฺชาย อาทีนวํ ทิสฺวา อิเมหิ อฏฺฐหากาเรหิ อวิชฺชานิโรเธ ฉนฺทชาโต โหติ สทฺธาธิมุตฺโต, จิตฺตญฺจสฺส สฺวาธิฏฺฐิตํ. “อวิชฺชาย นิโรธานุปสฺสี อสฺสสิสฺสามี”ติ สิกฺขติ, “อวิชฺชาย นิโรธานุปสฺสี ปสฺสสิสฺสามี”ติ สิกฺขติ.}

\prose{กติหากาเรหิ สงฺขาเรสุ อาทีนโว โหติ, กติหากาเรหิ สงฺขารา นิรุชฺฌนฺติ ฯเปฯ.\csromanspacer{} กติหากาเรหิ วิญฺญาเณ อาทีนโว โหติ, กติหากาเรหิ วิญฺญาณํ นิรุชฺฌติ.\csromanspacer{} กติหากาเรหิ นามรูเป \csromanfolio{194}อาทีนโว โหติ, กติหากาเรหิ นามรูปํ นิรุชฺฌติ.\csromanspacer{} กติหากาเรหิ สฬายตเน อาทีนโว โหติ, กติหากาเรหิ สฬายตนํ นิรุชฺฌติ.\csromanspacer{} กติหากาเรหิ ผสฺเส อาทีนโว โหติ, กติหากาเรหิ ผสฺโส นิรุชฺฌติ.\csromanspacer{} กติหากาเรหิ เวทนาย อาทีนโว โหติ, กติหากาเรหิ เวทนา นิรุชฺฌติ.\csromanspacer{} กติหากาเรหิ ตณฺหาย อาทีนโว โหติ, กติหากาเรหิ ตณฺหา นิรุชฺฌติ.\csromanspacer{} กติหากาเรหิ อุปาทาเน อาทีนโว โหติ, กติหากาเรหิ อุปาทานํ นิรุชฺฌติ.\csromanspacer{} กติหากาเรหิ ภเว อาทีนโว โหติ, กติหากาเรหิ ภโว นิรุชฺฌติ.\csromanspacer{} กติหากาเรหิ ชาติยา อาทีนโว โหติ, กติหากาเรหิ ชาติ นิรุชฺฌติ.\csromanspacer{} กติหากาเรหิ ชรามรเณ อาทีนโว โหติ, กติหากาเรหิ ชรามรณํ นิรุชฺฌติ? ปญฺจหากาเรหิ ชรามรเณ อาทีนโว โหติ, อฏฺฐหากาเรหิ ชรามรณํ นิรุชฺฌติ.}

\prose{กตเมหิ ปญฺจหากาเรหิ ชรามรเณ อาทินโว โหติ? อนิจฺจฏฺเฐน ชรามรเณ อาทีนโว โหติ, ทุกฺขฏฺเฐน ฯเปฯ อนตฺตฏฺเฐน ฯเปฯ สนฺตาปฏฺเฐน ฯเปฯ วิปริณามฏฺเฐน ชรามรเณ อาทีนโว โหติ.\csromanspacer{} อิเมหิ ปญฺจหากาเรหิ ชรามรเณ อาทีนโว โหติ.}

\prose{กตเมหิ อฏฺฐหากาเรหิ ชรามรณํ นิรุชฺฌติ? นิทานนิโรเธน ชรามรณํ นิรุชฺฌติ, สมุทยนิโรเธน ฯเปฯ ชาตินิโรเธน ฯเปฯ ปภวนิโรเธน\footnote{ภวนิโรเธน (สฺยา)}.\csromanspacer{} เหตุนิโรเธน.\csromanspacer{} ปจฺจยนิโรเธน.\csromanspacer{} ญาณุปฺปาเทน ฯเปฯ.\csromanspacer{} นิโรธุปฏฺฐาเนน ชรามรณํ นิรุชฺฌติ.\csromanspacer{} อิเมหิ อฏฺฐหากาเรหิ ชรามรณํ นิรุชฺฌติ.\csromanspacer{} อิเมหิ ปญฺจหากาเรหิ ชรามรเณ อาทีนวํ ทิสฺวา อิเมหิ อฏฺฐหากาเรหิ ชรามรณนิโรเธ ฉนฺทชาโต โหติ สทฺธาธิมุตฺโต, จิตฺตญฺจสฺส สฺวาธิฏฺฐิตํ. “ชรามรเณ นิโรธานุปสฺสี อสฺสสิสฺสามี”ติ สิกฺขติ, “ชรามรเณ นิโรธานุปสฺสี ปสฺสสิสฺสามี”ติ สิกฺขติ.\csromanspacer{} นิโรธานุปสฺสี อสฺสาสปสฺสาสวเสน ธมฺมา, อุปฏฺฐานํ สติ, อนุปสฺสนา ญาณํ.\csromanspacer{} ธมฺมา อุปฏฺฐานํ, โน สติ, สติ อุปฏฺฐานญฺเจว สติ จ.\csromanspacer{} ตาย สติยา เตน ญาเณน เต ธมฺเม อนุปสฺสติ, เตน วุจฺจติ “ธมฺเมสุ ธมฺมานุปสฺสนาสติปฏฺฐานภาวนา”ติ.}

\csromanheader{2}{3. อานาปานสฺสติกถา}
\prose{\csromanfolio{195}\textbf{อนุปสฺสตี}ติ กถํ เต ธมฺเม อนุปสฺสติ ฯเปฯ.\csromanspacer{} เอวํ เต ธมฺเม อนุปสฺสติ.\csromanspacer{} \textbf{ภาวนา}ติ จตสฺโส พฺ\textbf{หาวนา} ฯเปฯ อาเสวนฏฺเฐน พฺ\textbf{หาวนา}.\csromanspacer{} นิโรธานุปสฺสี อสฺสาสปสฺสาสานํ สํวรฏฺเฐน สีลวิสุทฺธิ ฯเปฯ นิโรธานุปสฺสี อสฺสาสปสฺสาสวเสน จิตฺตสฺส เอกคฺคตํ อวิกฺเขปํ ปชานโต ฯเปฯ ปชานนฺโต อินฺทฺริยานิ สโมธาเนติ, เตน วุจฺจติ “สมตฺถญฺจ ปฏิวิชฺฌตี”ติ. (3)}

\proseitem{182}{กถํ “ปฏินิสฺสคฺคานุปสฺสี อสฺสสิสฺสามี”ติ สิกฺขติ, “ปฏินิสฺสคฺคานุปสฺสี ปสฺสสิสฺสามี”ติ สิกฺขติ? \textbf{ปฏินิสฺสคฺคา}ติ เทฺว ปฏินิสฺสคฺคา ปริจฺจาคปฏินิสฺสคฺโค จ ปกฺขนฺทนปฏินิสฺสคฺโค จ.\csromanspacer{} รูปํ ปริจฺจชตีติ ปริจฺจาคปฏินิสฺสคฺโค, รูปนิโรเธ นิพฺพาเน จิตฺตํ ปกฺขนฺทตีติ ปกฺขนฺทนปฏินิสฺสคฺโค. “รูเป ปฏินิสฺสคฺคานุปสฺสี อสฺสสิสฺสามี”ติ สิกฺขติ, “รูเป ปฏินิสฺสคฺคานุปสฺสี ปสฺสสิสฺสามี”ติ สิกฺขติ.\csromanspacer{} เวทนํ ฯเปฯ.\csromanspacer{} สญฺญํ.\csromanspacer{} สงฺขาเร.\csromanspacer{} วิญฺญาณํ.\csromanspacer{} จกฺขุํ ฯเปฯ.\csromanspacer{} ชรามรณํ ปริจฺจชตีติ ปริจฺจาคปฏินิสฺสคฺโค, ชรามรณนิโรเธ นิพฺพาเน จิตฺตํ ปกฺขนฺทตีติ ปกฺขนฺทนปฏินิสฺสคฺโค, “ชรามรเณ ปฏินิสฺสคฺคานุปสฺสี อสฺสสิสฺสามี”ติ สิกฺขติ, “ชรามรเณ ปฏินิสฺสคฺคานุปสฺสี ปสฺสสิสฺสามี”ติ สิกฺขติ, ปฏินิสฺสคฺคานุปสฺสี อสฺสาสปสฺสาสวเสน ธมฺมา, อุปฏฺฐานํ สติ, อนุปสฺสนา ญาณํ.\csromanspacer{} ธมฺมา อุปฏฺฐานํ, โน สติ, สติ อุปฏฺฐานญฺเจว สติ จ.\csromanspacer{} ตาย สติยา เตน ญาเณน เต ธมฺเม อนุปสฺสติ, เตน วุจฺจติ “ธมฺเมสุ ธมฺมานุปสฺสนาสติปฏฺฐานพฺ\textbf{หาวนา}”ติ.}

\prose{\textbf{อนุปสฺสตี}ติ กถํ เต ธมฺเม อนุปสฺสติ อนิจฺจโต อนุปสฺสติ, โน นิจฺจโต ฯเปฯ ปฏินิสฺสชฺชติ, โน อาทิยติ.\csromanspacer{} อนิจฺจโต อนุปสฺสนฺโต นิจฺจสญฺญํ ปชหติ ฯเปฯ ปฏินิสฺสชฺชนฺโต อาทานํ ปชหติ, เอวํ เต ธมฺเม อนุปสฺสติ.\csromanspacer{} \textbf{ภาวนา}ติ จตสฺโส พฺ\textbf{หาวนา}.\csromanspacer{} ตตฺถ ชาตานํ ธมฺมานํ อนติวตฺตนฏฺเฐน พฺ\textbf{หาวนา} ฯเปฯ อาเสวนฏฺเฐน พฺ\textbf{หาวนา}, ปฏินิสฺสคฺคานุปสฺสี อสฺสาสปสฺสาสานํ สํวรฏฺเฐน สีลวิสุทฺธิ, อวิกฺเขปฏฺเฐน จิตฺตวิสุทฺธิ, ทสฺสนฏฺเฐน ทิฏฺฐิวิสุทฺธิ, โย ตตฺถ สํวรฏฺโฐ, อยํ อธิสีลสิกฺขา.\csromanspacer{} โย ตตฺถ อวิกฺเขปฏฺโฐ, อยํ อธิจิตฺตสิกฺขา.\csromanspacer{} โย ตตฺถ ทสฺสนฏฺโฐ, อยํ อธิปญฺญาสิกฺขา.\csromanspacer{} อิมา ติสฺโส สิกฺขาโย อาวชฺชนฺโต สิกฺขติ, ชานนฺโต สิกฺขติ ฯเปฯ สจฺฉิกาตพฺพํ สจฺฉิกโรนฺโต สิกฺขติ.}

\prose{\textbf{ปฏินิสฺสคฺคา}นุปสฺสี อสฺสาสปสฺสาสวเสน จิตฺตสฺส เอกคฺคตํ อวิกฺเขปํ ปชานโต วิทิตา เวทนา อุปชฺชนฺติ, วิทิตา อุปฏฺฐหนฺติ, วิทิตา อพฺภตฺถํ \csromanfolio{196}คจฺฉนฺติ ฯเปฯ.\csromanspacer{} ปฏินิสฺสคฺคานุปสฺสี อสฺสาสปสฺสาสวเสน จิตฺตสฺส เอกคฺคตํ อวิกฺเขปํ ปชานนฺโต อินฺทฺริยานิ สโมธาเนติ, โคจรญฺจ ปชานาติ, สมตฺถญฺจ ปฏิวิชฺฌติ, พลานิ สโมธาเนติ, โพชฺฌงฺเค สโมธาเนติ, มคฺคํ สโมธาเนติ, ธมฺเม สโมธาเนติ, โคจรญฺจ ปชานาติ, สมตฺถญฺจ ปฏิวิชฺฌติ.}

\prose{\textbf{อินฺทฺริยานิ สโมธาเนตี}ติ กถํ อินฺทฺริยานิ สโมธาเนติ, อธิโมกฺขฏฺเฐน สทฺธินฺทฺริยํ สโมธาเนติ ฯเปฯ เตน วุจฺจติ “สมตฺถญฺจ ปฏิวิชฺฌตี”ติ. (4)}

\prose{อฏฺฐ อนุปสฺสเน ญาณานิ, อฏฺฐ จ อุปฏฺฐานานุสฺสติโย, จตฺตาริ สุตฺตนฺติกวตฺถูนิ ธมฺเมสุ ธมฺมานุปสฺสนาย.\csromanspacer{} อิมานิ พาตฺติํส สโตการิสฺส ญาณานิ.}

\csromanheader{2}{สโตการิญาณนิทฺเทโส ปญฺจโม.}
\csromansectionrule
\csromanmatikaanchor{mtk.71}
\csromantocmarkref{subsection}{6. ญาณราสิฉกฺกนิทฺเทส}{mtk.71}
\bookmark[dest={mtk.71},level=2]{6. ญาณราสิฉกฺกนิทฺเทส}
\csromanheader{2}{6. ญาณราสิฉกฺกนิทฺเทส}
\proseitem{183}{กตมานิ จตุวีสติ สมาธิวเสน ญาณานิ? ทีฆํ อสฺสาสวเสน จิตฺตสฺส เอกคฺคตา อวิกฺเขโป สมาธิ, ทีฆํ ปสฺสาสวเสน จิตฺตสฺส เอกคฺคตา อวิกฺเขโป สมาธิ ฯเปฯ.\csromanspacer{} วิโมจยํ จิตฺตํ อสฺสาสวเสน จิตฺตสฺส เอกคฺคตา อวิกฺเขโป สมาธิ.\csromanspacer{} วิโมจยํ จิตฺตํ ปสฺสาสวเสน จิตฺตสฺส เอกคฺคตา อวิกฺเขโป สมาธิ.\csromanspacer{} อิมานิ จตุวีสติ สมาธิวเสน ญาณานิ. (1)}

\prose{กตมานิ เทฺวสตฺตติ วิปสฺสนาวเสน ญาณานิ? ทีฆํ อสฺสาสํ อนิจฺจโต อนุปสฺสนฏฺเฐน วิปสฺสนา, ทุกฺขโต อนุปสฺสนฏฺเฐน วิปสฺสนา, อนตฺตโต อนุปสฺสนฏฺเฐน วิปสฺสนา.\csromanspacer{} ทีฆํ ปสฺสาสํ อนิจฺจโต อนุปสฺสนฏฺเฐน วิปสฺสนา, ทุกฺขโต อนุปสฺสนฏฺเฐน วิปสฺสนา, อนตฺตโต อนุปสฺสนฏฺเฐน วิปสฺสนา ฯเปฯ วิโมจยํ จิตฺตํ อสฺสาสํ.\csromanspacer{} วิโมจยํ จิตฺตํ ปสฺสาสํ อนิจฺจโต อนุปสฺสนฏฺเฐน วิปสฺสนา, ทุกฺขโต อนุปสฺสนฏฺเฐน วิปสฺสนา, อนตฺตโต อนุปสฺสนฏฺเฐน วิปสฺสนา.\csromanspacer{} อิมานิ เทฺวสตฺตติ วิปสฺสนาวเสน ญาณานิ. (2)}

\csromanheader{2}{3. อานาปานสฺสติกถา}
\prose{\csromanfolio{197}กตมานิ อฏฺฐ นิพฺพิทาญาณานิ? อนิจฺจานุปสฺสี อสฺสสํ ยถาภูตํ ชานาติ ปสฺสตีติ นิพฺพิทาญาณํ, อนิจฺจานุปสฺสี ปสฺสสํ ยถาภูตํ ชานาติ ปสฺสตีติ นิพฺพิทาญาณํ ฯเปฯ ปฏินิสฺสคฺคานุปสฺสี อสฺสสํ ยถาภูตํ ชานาติ ปสฺสตีติ นิพฺพิทาญาณํ, ปฏินิสฺสคฺคานุปสฺสี ปสฺสสํ ยถาภูตํ ชานาติ ปสฺสตีติ นิพฺพิทาญาณํ.\csromanspacer{} อิมานิ อฏฺฐ นิพฺพิทาญาณานิ. (3)}

\prose{กตมานิ อฏฺฐ นิพฺพิทานุโลเม ญาณานิ? อนิจฺจานุปสฺสี อสฺสสํ ภยตุปฏฺฐาเน ปญฺญา นิพฺพิทานุโลเม ญาณํ.\csromanspacer{} อนิจฺจานุปสฺสี ปสฺสสํ ภยตุปฏฺฐาเน ปญฺญา นิพฺพิทานุโลเม ญาณํ ฯเปฯ.\csromanspacer{} ปฏินิสฺสคฺคานุปสฺสี อสฺสสํ ภยตุปฏฺฐาเน ปญฺญา นิพฺพิทานุโลเม ญาณํ.\csromanspacer{} ปฏินิสฺสคฺคานุปสฺสี ปสฺสสํ ภยตุปฏฺฐาเน ปญฺญา นิพฺพิทานุโลเม ญาณํ.\csromanspacer{} อิมานิ อฏฺฐ นิพฺพิทานุโลเม ญาณานิ. (4)}

\prose{กตมานิ อฏฺฐ นิพฺพิทาปฏิปฺปสฺสทฺธิญาณานิ? อนิจฺจานุปสฺสี อสฺสสํ ปฏิสงฺขา สนฺติฏฺฐนา ปญฺญา นิพฺพิทาปฏิปฺปสฺสทฺธิญาณํ.\csromanspacer{} อนิจฺจานุปสฺสี ปสฺสสํ ปฏิสงฺขา สนฺติฏฺฐนา ปญฺญา นิพฺพิทาปฏิปฺปสฺสทฺธิญาณํ ฯเปฯ.\csromanspacer{} ปฏินิสฺสคฺคานุปสฺสี อสฺสสํ ปฏิสงฺขา สนฺติฏฺฐนา ปญฺญา นิพฺพิทาปฏิปฺปสฺสทฺธิญาณํ.\csromanspacer{} ปฏินิสฺสคฺคานุปสฺสี ปสฺสสํ ปฏิสงฺขา สนฺติฏฺฐนา ปญฺญา นิพฺพิทาปฏิปฺปสฺสทฺธิญาณํ.\csromanspacer{} อิมานิ อฏฺฐ นิพฺพิทาปฏิปฺปสฺสทฺธิญาณานิ. (5)}

\prose{กตมานิ เอกวีสติ วิมุตฺติสุเข ญาณานิ? โสตาปตฺติมคฺเคน สกฺกายทิฏฺฐิยา ปหีนตฺตา สมุจฺฉินฺนตฺตา อุปฺปชฺชติ วิมุตฺติสุเข ญาณํ.\csromanspacer{} วิจิกิจฺฉาย ปหีนตฺตา สมุจฺฉินฺนตฺตา อุปฺปชฺชติ วิมุตฺติสุเข ญาณํ.\csromanspacer{} สีลพฺพตปรามาสสฺส ฯเปฯ.\csromanspacer{} ทิฏฺฐานุสยสฺส วิจิกิจฺฉานุสยสฺส ปหีนตฺตา สมุจฺฉินฺนตฺตา อุปฺปชฺชติ วิมุตฺติสุเข ญาณํ.\csromanspacer{} สกทาคามิมคฺเคน โอฬาริกสฺส กามราคสญฺโญชนสฺส ฯเปฯ.\csromanspacer{} ปฏิฆสญฺโญชนสฺส โอฬาริกสฺส กามราคานุสยสฺส ปฏิฆานุสยสฺส ปหีนตฺตา สมุจฺฉินฺนตฺตา อุปฺปชฺชติ วิมุตฺติสุเข ญาณํ.\csromanspacer{} อนาคามิมคฺเคน อณุสหคตสฺส กามราคสญฺโญชนสฺส ฯเปฯ.\csromanspacer{} ปฏิฆสญฺโญชนสฺส อณุสหคตสฺส กามราคานุสยสฺส ปฏิฆานุสยสฺส ปหีนตฺตา สมุจฺฉินฺนตฺตา อุปฺปชฺชติ วิมุตฺติสุเข ญาณํ.\csromanspacer{} อรหตฺตมคฺเคน รูปราคสฺส ฯเปฯ.\csromanspacer{} อรูปราคสฺส มานสฺส อุทฺธจฺจสฺส อวิชฺชาย มานานุสยสฺส ภวราคานุสยสฺส อวิชฺชานุสยสฺส ปหีนตฺตา สมุจฺฉินฺนตฺตา อุปฺปชฺชติ}

\prose{\csromanfolio{198}วิมุตฺติสุเข ญาณํ.\csromanspacer{} อิมานิ เอกวีสติ วิมุตฺติสุเข ญาณานิ.\csromanspacer{} โสฬสวตฺถุกํ อานาปานสฺสติสมาธิํ ภาวยโต สมธิกานิ อิมานิ เทฺว ญาณสตานิ อุปฺปชฺชนฺติ. (6)}

\nitthitamruled{อานาปานสฺสติกถา\footnote{อานาปานกถา (สฺยา, ก)} นิฏฺฐิตา.}
\csromanmatikaanchor{mtk.72}
\csromantocmarknopage{section}{4. อินฺทฺริยกถา}{mtk.72}
\bookmark[dest={mtk.72},level=1]{4. อินฺทฺริยกถา}
\csromanheader{1}{4. อินฺทฺริยกถา}
\csromanmatikaanchor{mtk.73}
\csromantocmarkref{subsection}{1. ปฐมสุตฺตนฺตนิทฺเทส}{mtk.73}
\bookmark[dest={mtk.73},level=2]{1. ปฐมสุตฺตนฺตนิทฺเทส}
\csromanheader{2}{1. ปฐมสุตฺตนฺตนิทฺเทส}
\proseitem{184}{เอวํ เม สุตํ\csromandash{}เอกํ สมยํ ภควา สาวตฺถิยํ วิหรติ เชตวเน อนาถปิณฺฑิกสฺส อาราเม.\csromanspacer{} ตตฺร โข ภควา ภิกฺขู อามนฺเตสิ “ภิกฺขโว”ติ. “ภทนฺเต”ติ\footnote{“ภทฺทนฺเต”ติ (ก)} เต ภิกฺขู ภควโต ปจฺจสฺโสสุํ.\csromanspacer{} ภควา เอตทโวจ\csromandash{}}

\prose{ปญฺจิมานิ ภิกฺขเว อินฺทฺริยานิ.\csromanspacer{} กตมานิ ปญฺจ? สทฺธินฺทฺริยํ วีริยินฺทฺริยํ สตินฺทฺริยํ สมาธินฺทฺริยํ ปญฺญินฺทฺริยํ.\csromanspacer{} อิมานิ โข ภิกฺขเว ปญฺจินฺทฺริยานิ\footnote{สํ 3. 170 ปิฏฺเฐปิ.}.}

\proseitem{185}{อิมานิ ปญฺจินฺทฺริยานิ กติหากาเรหิ วิสุชฺฌนฺติ? อิมานิ ปญฺจินฺทฺริยานิ ปนฺนรสหิ อากาเรหิ วิสุชฺฌนฺติ.\csromanspacer{} อสฺสทฺเธ ปุคฺคเล ปริวชฺชยโต สทฺเธ ปุคฺคเล เสวโต ภชโต ปยิรุปาสโต ปสาทนีเย สุตฺตนฺเต ปจฺจเวกฺขโต อิเมหิ ตีหากาเรหิ สทฺธินฺทฺริยํ วิสุชฺฌติ.\csromanspacer{}\csromanspacer{}กุสีเต ปุคฺคเล ปริวชฺชยโต อารทฺธวีริเย ปุคฺคเล เสวโต ภชโต ปยิรุปาสโต สมฺมปฺปธาเน ปจฺจเวกฺขโต อิเมหิ ตีหากาเรหิ วีริยินฺทฺริยํ วิสุชฺฌติ.\csromanspacer{}\csromanspacer{}มุฏฺฐสฺสตี ปุคฺคเล ปริวชฺชยโต อุปฏฺฐิตสฺสตี ปุคฺคเล เสวโต ภชโต ปยิรุปาสโต สติปฏฺฐาเน ปจฺจเวกฺขโต อิเมหิ ตีหากาเรหิ สตินฺทฺริยํ วิสุชฺฌติ.\csromanspacer{}\csromanspacer{}อสมาหิเต ปุคฺคเล ปริวชฺชยโต สมาหิเต ปุคฺคเล เสวโต ภชโต ปยิรุปาสโต ฌานวิโมกฺเข ปจฺจเวกฺขโต อิเมหิ ตีหากาเรหิ สมาธินฺทฺริยํ วิสุชฺฌติ.\csromanspacer{}\csromanspacer{}ทุปฺปญฺเญ ปุคฺคเล ปริวชฺชยโต ปญฺญวนฺเต ปุคฺคเล เสวโต ภชโต ปยิรุปาสโต คมฺภีรญาณจริยํ ปจฺจเวกฺขโต อิเมหิ ตีหากาเรหิ ปญฺญินฺทฺริยํ วิสุชฺฌติ.\csromanspacer{} อิติ อิเม ปญฺจ ปุคฺคเล ปริวชฺชยโต ปญฺจ}

\csromanheader{2}{4. อินฺทฺริยกถา}
\prose{\csromanfolio{199}ปุคฺคเล เสวโต ภชโต ปยิรุปาสโต ปญฺจ สุตฺตนฺตกฺขนฺเธ ปจฺจเวกฺขโต อิเมหิ ปนฺนรสหิ อากาเรหิ อิมานิ ปญฺจินฺทฺริยานิ วิสุชฺฌนฺติ.}

\prose{กติหากาเรหิ ปญฺจินฺทฺริยานิ ภาวิยนฺติ, กติหากาเรหิ ปญฺจนฺนํ อินฺทฺริยานํ ภาวนา โหติ? ทสหากาเรหิ ปญฺจินฺทฺริยานิ ภาวิยนฺติ, ทสหากาเรหิ ปญฺจนฺนํ อินฺทฺริยานํ ภาวนา โหติ.\csromanspacer{} อสฺสทฺธิยํ ปชหนฺโต สทฺธินฺทฺริยํ ภาเวติ, สทฺธินฺทฺริยํ ภาเวนฺโต อสฺสทฺธิยํ ปชหติ, โกสชฺชํ ปชหนฺโต วีริยินฺทฺริยํ ภาเวติ, วีริยินฺทฺริยํ ภาเวนฺโต โกสชฺชํ ปชหติ, ปมาทํ ปชหนฺโต สตินฺทฺริยํ ภาเวติ, สตินฺทฺริยํ ภาเวนฺโต ปมาทํ ปชหติ, อุทฺธจฺจํ ปชหนฺโต สมาธินฺทฺริยํ ภาเวติ, สมาธินฺทฺริยํ ภาเวนฺโต อุทฺธจฺจํ ปชหติ, อวิชฺชํ ปชหนฺโต ปญฺญินฺทฺริยํ ภาเวติ, ปญฺญินฺทฺริยํ ภาเวนฺโต อวิชฺชํ ปชหติ.\csromanspacer{} อิเมหิ ทสหากาเรหิ ปญฺจินฺทฺริยานิ ภาวิยนฺติ, อิเมหิ ทสหากาเรหิ ปญฺจนฺนํ อินฺทฺริยานํ ภาวนา โหติ.}

\prose{กติหากาเรหิ ปญฺจินฺทฺริยานิ ภาวิตานิ โหนฺติ สุภาวิตานิ.\csromanspacer{} ทสหากาเรหิ ปญฺจินฺทฺริยานิ ภาวิตานิ โหนฺติ สุภาวิตานิ.\csromanspacer{} อสฺสทฺธิยสฺส ปหีนตฺตา สุปฺปหีนตฺตา สทฺธินฺทฺริยํ ภาวิตํ โหติ สุภาวิตํ, สทฺธินฺทฺริยสฺส ภาวิตตฺตา สุภาวิตตฺตา อสฺสทฺธิยํ ปหีนํ โหติ สุปฺปหีนํ, โกสชฺชสฺส ปหีนตฺตา สุปฺปหีนตฺตา วีริยินฺทฺริยํ ภาวิตํ โหติ สุภาวิตํ, วีริยินฺทฺริยสฺส ภาวิตตฺตา สุภาวิตตฺตา โกสชฺชํ ปหีนํ โหติ สุปฺปหีนํ, ปมาทสฺส ปหีนตฺตา สุปฺปหีนตฺตา สตินฺทฺริยํ ภาวิตํ โหติ สุภาวิตํ, สตินฺทฺริยสฺส ภาวิตตฺตา สุภาวิตตฺตา ปมาโท ปหีโน โหติ สุปฺปหีโน, อุทฺธจฺจสฺส ปหีนตฺตา สุปฺปหีนตฺตา สมาธินฺทฺริยํ ภาวิตํ โหติ สุภาวิตํ, สมาธินฺทฺริยสฺส ภาวิตตฺตา สุภาวิตตฺตา อุทฺธจฺจํ ปหีนํ โหติ สุปฺปหีนํ, อวิชฺชาย ปหีนตฺตา ปญฺญินฺทฺริยํ ภาวิตํ โหติ สุภาวิตํ, ปญฺญินฺทฺริยสฺส ภาวิตตฺตา สุภาวิตตฺตา อวิชฺชา ปหีนา โหติ สุปฺปหีนา.\csromanspacer{} อิเมหิ ทสหากาเรหิ ปญฺจินฺทฺริยานิ ภาวิตานิ โหนฺติ สุภาวิตานิ.}

\proseitem{186}{กติหากาเรหิ ปญฺจินฺทฺริยานิ ภาวิยนฺติ, กติหากาเรหิ ปญฺจินฺทฺริยานิ ภาวิตานิ เจว โหนฺติ สุภาวิตานิ จ ปฏิปฺปสฺสทฺธานิ จ สุปฺปฏิปฺปสฺสทฺธานิ จ? จตูหากาเรหิ ปญฺจินฺทฺริยานิ ภาวิยนฺติ, จตูหากาเรหิ ปญฺจินฺทฺริยานิ ภาวิตานิ เจว โหนฺติ สุภาวิตานิ จ ปฏิปฺปสฺสทฺธานิ จ สุปฺปฏิปฺปสฺสทฺธานิ จ.\csromanspacer{} โสตาปตฺติมคฺคกฺขเณ ปญฺจินฺทฺริยานิ ภาวิยนฺติ, \csromanfolio{200}โสตาปตฺติผลกฺขเณ ปญฺจินฺทฺริยานิ ภาวิตานิ เจว โหนฺติ สุภาวิตานิ จ ปฏิปฺปสฺสทฺธานิ จ สุปฺปฏิปฺปสฺสทฺธานิ จ.\csromanspacer{} สกทาคามิมคฺคกฺขเณ ปญฺจินฺทฺริยานิ ภาวิยนฺติ, สกทาคามิผลกฺขเณ ปญฺจินฺทฺริยานิ ภาวิตานิ เจว โหนฺติ สุภาวิตานิ จ ปฏิปฺปสฺสทฺธานิ จ สุปฺปฏิปฺปสฺสทฺธานิ จ.\csromanspacer{} อนาคามิมคฺคกฺขเณ ปญฺจินฺทฺริยานิ ภาวิยนฺติ, อนาคามิผลกฺขเณ ปญฺจินฺทฺริยานิ ภาวิตานิ เจว โหนฺติ สุภาวิตานิ จ ปฏิปฺปสฺสทฺธานิ จ สุปฺปฏิปฺปสฺสทฺธานิ จ.\csromanspacer{} อรหตฺตมคฺคกฺขเณ ปญฺจินฺทฺริยานิ ภาวิยนฺติ, อรหตฺตผลกฺขเณ ปญฺจินฺทฺริยานิ ภาวิตานิ เจว โหนฺติ สุภาวิตานิ จ ปฏิปฺปสฺสทฺธานิ จ สุปฺปฏิปฺปสฺสทฺธานิ จ.\csromanspacer{} อิติ จตสฺโส มคฺควิสุทฺธิโย จตสฺโส ผลวิสุทฺธิโย จตสฺโส สมุจฺเฉทวิสุทฺธิโย จตสฺโส ปฏิปฺปสฺสทฺธิวิสุทฺธิโย.\csromanspacer{} อิเมหิ จตูหากาเรหิ ปญฺจินฺทฺริยานิ ภาวิยนฺติ, อิเมหิ จตูหากาเรหิ ปญฺจินฺทฺริยานิ ภาวิตานิ เจว โหนฺติ สุภาวิตานิ จ ปฏิปฺปสฺสทฺธานิ จ สุปฺปฏิปฺปสฺสทฺธานิ จ.}

\prose{กตินํ ปุคฺคลานํ อินฺทฺริยภาวนา, กติ ปุคฺคลา ภาวิตินฺทฺริยา? อฏฺฐนฺนํ ปุคฺคลานํ อินฺทฺริยภาวนา, ตโย ปุคฺคลา ภาวิตินฺทฺริยา.\csromanspacer{}\csromanspacer{}กตเมสํ อฏฺฐนฺนํ ปุคฺคลานํ อินฺทฺริยภาวนา? สตฺตนฺนญฺจ เสกฺขานํ ปุถุชฺชนกลฺยาณกสฺส จ, อิเมสํ อฏฺฐนฺนํ ปุคฺคลานํ อินฺทฺริยภาวนา.\csromanspacer{}\csromanspacer{}กตเม ตโย ปุคฺคลา ภาวิตินฺทฺริยา, สวเนน พุทฺโธ ตถาคตสฺส สาวโก ขีณาสโว ภาวิตินฺทฺริโย, สยํ ภูตฏฺเฐน ปจฺเจกสมฺพุทฺโธ ภาวิตินฺทฺริโย, อปฺปเมยฺยฏฺเฐน ตถาคโต อรหํ สมฺมาสมฺพุทฺโธ ภาวิตินฺทฺริโย.\csromanspacer{} อิเม ตโย ปุคฺคลา ภาวิตินฺทฺริยา.\csromanspacer{} อิติ อิเมสํ อฏฺฐนฺนํ ปุคฺคลานํ อินฺทฺริยภาวนา.\csromanspacer{} อิเม ตโย ปุคฺคลา ภาวิตินฺทฺริยา.}

\csromanheader{2}{สุตฺตนฺตนิทฺเทโส ปฐโม.}
\csromansectionrule
\csromanmatikaanchor{mtk.74}
\csromantocmarkref{subsection}{2. ทุติยสุตฺตนฺตนิทฺเทส}{mtk.74}
\bookmark[dest={mtk.74},level=2]{2. ทุติยสุตฺตนฺตนิทฺเทส}
\csromanheader{2}{2. ทุติยสุตฺตนฺตนิทฺเทส}
\proseitem{187}{สาวตฺถินิทานํ.\csromanspacer{}\csromanspacer{}ปญฺจิมานิ ภิกฺขเว อินฺทฺริยานิ.\csromanspacer{} กตมานิ ปญฺจ? สทฺธินฺทฺริยํ วีริยินฺทฺริยํ สตินฺทฺริยํ สมาธินฺทฺริยํ ปญฺญินฺทฺริยํ.\csromanspacer{} เย หิ เกจิ ภิกฺขเว สมณา วา พฺราหฺมณา วา อิเมสํ ปญฺจนฺนํ อินฺทฺริยานํ สมุทยญฺจ อตฺถงฺคมญฺจ อสฺสาทญฺจ อาทีนวญฺจ นิสฺสรณญฺจ ยถาภูตํ นปฺปชานนฺติ, น เม เต ภิกฺขเว สมณา วา พฺราหฺมณา วา สมเณสุ วา\footnote{เจว (ก) สํ 3. 171 ปิฏฺเฐ ปสฺสิตพฺพา.} สมณสมฺมตา,}

\csromanheader{2}{4. อินฺทฺริยกถา}
\prose{\csromanfolio{201}พฺราหฺมเณสุ วา พฺราหฺมณสมฺมตา, น จ ปเนเต อายสฺมนฺตา\footnote{อายสฺมนฺโต สํ 1. 433 ปิฏฺเฐ; สํ 3. 171 ปิฏฺเฐปิ.} สามญฺญตฺถํ วา พฺรหฺมญฺญตฺถํ วา ทิฏฺเฐว ธมฺเม สยํ อภิญฺญา สจฺฉิกตฺวา อุปสมฺปชฺช วิหรนฺติ.\csromanspacer{} เย จ โข เกจิ\footnote{เย เกจิ โข (สฺยา), เย หิ เกจิ โข (ก) สํ 3. 171 ปิฏฺเฐ ปสฺสิตพฺพา.} ภิกฺขเว สมณา วา พฺราหฺมณา วา อิเมสํ ปญฺจนฺนํ อินฺทฺริยานํ สมุทยญฺจ อตฺถงฺคมญฺจ อสฺสาทญฺจ อาทีนวญฺจ นิสฺสรณญฺจ ยถาภูตํ ปชานนฺติ, เต โข เม ภิกฺขเว สมณา วา พฺราหฺมณา วา สมเณสุ วา สมณสมฺมตา, พฺราหฺมเณสุ วา พฺราหฺมณสมฺมตา, เต จ ปนายสฺมนฺตา สามญฺญตฺถญฺจ พฺราหฺมญฺญตฺถญฺจ ทิฏฺเฐว ธมฺเม สยํ อภิญฺญา สจฺฉิกตฺวา อุปสมฺปชฺช วิหรนฺตีติ.}

\proseitem{188}{กติหากาเรหิ ปญฺจนฺนํ อินฺทฺริยานํ สมุทโย โหติ, กติหากาเรหิ ปญฺจนฺนํ อินฺทฺริยานํ สมุทยํ ปชานาติ.\csromanspacer{} กติหากาเรหิ ปญฺจนฺนํ อินฺทฺริยานํ อตฺถงฺคโม โหติ, กติหากาเรหิ ปญฺจนฺนํ อินฺทฺริยานํ อตฺถงฺคมํ ปชานาติ.\csromanspacer{} กติหากาเรหิ ปญฺจนฺนํ อินฺทฺริยานํ อสฺสาโท โหติ, กติหากาเรหิ ปญฺจนฺนํ อินฺทฺริยานํ อสฺสาทํ ปชานาติ.\csromanspacer{} กติหากาเรหิ ปญฺจนฺนํ อินฺทฺริยานํ อาทีนโว โหติ, กติหากาเรหิ ปญฺจนฺนํ อินฺทฺริยานํ อาทีนวํ ปชานาติ.\csromanspacer{} กติหากาเรหิ ปญฺจนฺนํ อินฺทฺริยานํ นิสฺสรณํ โหติ, กติหากาเรหิ ปญฺจนฺนํ อินฺทฺริยานํ นิสฺสรณํ ปชานาติ?}

\prose{จตฺตารีสาย อากาเรหิ ปญฺจนฺนํ อินฺทฺริยานํ สมุทโย โหติ, จตฺตารีสาย อากาเรหิ ปญฺจนฺนํ อินฺทฺริยานํ สมุทยํ ปชานาติ.\csromanspacer{} จตฺตารีสาย อากาเรหิ ปญฺจนฺนํ อินฺทฺริยานํ อตฺถงฺคโม โหติ, จตฺตารีสาย อากาเรหิ ปญฺจนฺนํ อินฺทฺริยานํ อตฺถงฺคมํ ปชานาติ.\csromanspacer{} ปญฺจวีสติยา อากาเรหิ ปญฺจนฺนํ อินฺทฺริยานํ อสฺสาโท โหติ, ปญฺจวีสติยา อากาเรหิ ปญฺจนฺนํ อินฺทฺริยานํ อสฺสาทํ ปชานาติ.\csromanspacer{} ปญฺจวีสติยา อากาเรหิ ปญฺจนฺนํ อินฺทฺริยานํ อาทีนโว โหติ, ปญฺจวีสติยา อากาเรหิ ปญฺจนฺนํ อินฺทฺริยานํ อาทีนวํ ปชานาติ.\csromanspacer{} อสีติสตํ อากาเรหิ ปญฺจนฺนํ อินฺทฺริยานํ นิสฺสรณํ โหติ, อสีติสตํ อากาเรหิ\footnote{อสีติสตากาเรหิ (สฺยา)} ปญฺจนฺนํ อินฺทฺริยานํ นิสฺสรณํ ปชานาติ.}

\prose{\csromanfolio{202}กตเมหิ จตฺตารีสาย อากาเรหิ ปญฺจนฺนํ อินฺทฺริยานํ สมุทโย โหติ, กตเมหิ จตฺตารีสาย อากาเรหิ ปญฺจนฺนํ อินฺทฺริยานํ สมุทยํ ปชานาติ? อธิโมกฺขตฺถาย อาวชฺชนายา สมุทโย สทฺธินฺทฺริยสฺส สมุทโย โหติ, อธิโมกฺขวเสน ฉนฺทสฺส สมุทโย สทฺธินฺทฺริยสฺส สมุทโย โหติ, อธิโมกฺขวเสน มนสิการสฺส สมุทโย สทฺธินฺทฺริยสฺส สมุทโย โหติ.\csromanspacer{} สทฺธินฺทฺริยสฺส วเสน เอกตฺตุปฏฺฐานํ สทฺธินฺทฺริยสฺส สมุทโย โหติ.\csromanspacer{}\csromanspacer{}ปคฺคหตฺถาย อาวชฺชนาย สมุทโย วีริยินฺทฺริยสฺส สมุทโย โหติ, ปคฺคหวเสน ฉนฺทสฺส สมุทโย วีริยินฺทฺริยสฺส สมุทโย โหติ, ปคฺคหวเสน มนสิการสฺส สมุทโย วีริยินฺทฺริยสฺส สมุทโย โหติ.\csromanspacer{} วีริยินฺทฺริยสฺส วเสน เอกตฺตุปฏฺฐานํ วีริยินฺทฺริยสฺส สมุทโย โหติ.\csromanspacer{}\csromanspacer{}อุปฏฺฐานตฺถาย อาวชฺชนาย สมุทโย สตินฺทฺริยสฺส สมุทโย โหติ, อุปฏฺฐานวเสน ฉนฺทสฺส สมุทโย สตินฺทฺริยสฺส สมุทโย โหติ, อุปฏฺฐานวเสน มนสิการสฺส สมุทโย สตินฺทฺริยสฺส สมุทโย โหติ.\csromanspacer{} สตินฺทฺริยสฺส วเสน เอกตฺตุปฏฺฐานํ สตินฺทฺริยสฺส สมุทโย โหติ.\csromanspacer{}\csromanspacer{}อวิกฺเขปตฺถาย อาวชฺชนาย สมุทโย สมาธินฺทฺริยสฺส สมุทโย โหติ, อวิกฺเขปวเสน ฉนฺทสฺส สมุทโย สมาธินฺทฺริยสฺส สมุทโย โหติ, อวิกฺเขปวเสน มนสิการสฺส สมุทโย สมาธินฺทฺริยสฺส สมุทโย โหติ.\csromanspacer{} สมาธินฺทฺริยสฺส วเสน เอกตฺตุปฏฺฐานํ สมาธินฺทฺริยสฺส สมุทโย โหติ.\csromanspacer{}\csromanspacer{}ทสฺสนตฺถาย อาวชฺชนาย สมุทโย ปญฺญินฺทฺริยสฺส สมุทโย โหติ, ทสฺสนวเสน ฉนฺทสฺส สมุทโย ปญฺญินฺทฺริยสฺส สมุทโย โหติ, ทสฺสนวเสน มนสิการสฺส สมุทโย ปญฺญินฺทฺริยสฺส สมุทโย โหติ.\csromanspacer{} ปญฺญินฺทฺริยสฺส วเสน เอกตฺตุปฏฺฐานํ ปญฺญินฺทฺริยสฺส สมุทโย โหติ.}

\prose{อธิโมกฺขตฺถาย อาวชฺชนาย สมุทโย สทฺธินฺทฺริยสฺส สมุทโย โหติ, ปคฺคหตฺถาย อาวชฺชนาย สมุทโย วีริยินฺทฺริยสฺส สมุทโย โหติ, อุปฏฺฐานตฺถาย อาวชฺชนาย สมุทโย สตินฺทฺริยสฺส สมุทโย โหติ.\csromanspacer{} อวิกฺเขปตฺถาย อาวชฺชนาย สมุทโย สมาธินฺทฺริยสฺส สมุทโย โหติ, ทสฺสนตฺถาย อาวชฺชนาย สมุทโย ปญฺญินฺทฺริยสฺส สมุทโย โหติ.\csromanspacer{}\csromanspacer{}อธิโมกฺขวเสน ฉนฺทสฺส สมุทโย สทฺธินฺทฺริยสฺส สมุทโย โหติ, ปคฺคหวเสน ฉนฺทสฺส สมุทโย วีริยินฺทฺริยสฺส สมุทโย โหติ, อุปฏฺฐานวเสน ฉนฺทสฺส สมุทโย สตินฺทฺริยสฺส}

\csromanheader{2}{4. อินฺทฺริยกถา}
\prose{\csromanfolio{203}สมุทโย โหติ, อวิกฺเขปวเสน ฉนฺทสฺส สมุทโย สมาธินฺทฺริยสฺส สมุทโย โหติ, ทสฺสนวเสน ฉนฺทสฺส สมุทโย ปญฺญินฺทฺริยสฺส สมุทโย โหติ.\csromanspacer{}\csromanspacer{}อธิโมกฺขวเสน มนสิการสฺส สมุทโย สทฺธินฺทฺริยสฺส สมุทโย โหติ, ปคฺคหวเสน มนสิการสฺส สมุทโย วีริยินฺทฺริยสฺส สมุทโย โหติ, อุปฏฺฐานวเสน มนสิการสฺส สมุทโย สตินฺทฺริยสฺส สมุทโย โหติ, อวิกฺเขปวเสน มนสิการสฺส สมุทโย สมาธินฺทฺริยสฺส สมุทโย โหติ, ทสฺสนวเสน มนสิการสฺส สมุทโย ปญฺญินฺทฺริยสฺส สมุทโย โหติ, สทฺธินฺทฺริยสฺส วเสน เอกตฺตุปฏฺฐานํ สทฺธินฺทฺริยสฺส สมุทโย โหติ, วีริยินฺทฺริยสฺส วเสน เอกตฺตุปฏฺฐานํ วีริยินฺทฺริยสฺส สมุทโย โหติ, สตินฺทฺริยสฺส วเสน เอกตฺตุปฏฺฐานํ สตินฺทฺริยสฺส สมุทโย โหติ, สมาธินฺทฺริยสฺส วเสน เอกตฺตุปฏฺฐานํ สมาธินฺทฺริยสฺส สมุทโย โหติ, ปญฺญินฺทฺริยสฺส วเสน เอกตฺตุปฏฺฐานํ ปญฺญินฺทฺริยสฺส สมุทโย โหติ.}

\prose{อิเมหิ จตฺตารีสาย อากาเรหิ ปญฺจนฺนํ อินฺทฺริยานํ สมุทโย โหติ, อิเมหิ จตฺตารีสาย อากาเรหิ ปญฺจนฺนํ อินฺทฺริยานํ สมุทยํ ปชานาติ.}

\prose{กตเมหิ จตฺตารีสาย อากาเรหิ ปญฺจนฺนํ อินฺทฺริยานํ อตฺถงฺคโม โหติ, กตเมหิ จตฺตารีสาย อากาเรหิ ปญฺจนฺนํ อินฺทฺริยานํ อตฺถงฺคมํ ปชานาติ? อธิโมกฺขตฺถาย อาวชฺชนาย อตฺถงฺคโม สทฺธินฺทฺริยสฺส อตฺถงฺคโม โหติ, อธิโมกฺขวเสน ฉนฺทสฺส อตฺถงฺคโม สทฺธินฺทฺริยสฺส อตฺถงฺคโม โหติ, อธิโมกฺขวเสน มนสิการสฺส อตฺถงฺคโม สทฺธินฺทฺริยสฺส อตฺถงฺคโม โหติ.\csromanspacer{} สทฺธินฺทฺริยสฺส วเสน เอกตฺต-อนุปฏฺฐานํ\footnote{เอกตฺตํ อนุปฏฺฐานํ (สฺยา)} สทฺธินฺทฺริยสฺส วเสน อตฺถงฺคโม โหติ.\csromanspacer{}\csromanspacer{}ปคฺคหตฺถาย อาวชฺชนาย อตฺถงฺคโม วีริยินฺทฺริยสฺส อตฺถงฺคโม โหติ, ปคฺคหวเสน ฉนฺทสฺส อตฺถงฺคโม วีริยินฺทฺริยสฺส อตฺถงฺคโม โหติ.\csromanspacer{} ปคฺคหวเสน มนสิการสฺส อตฺถงฺคโม วีริยินฺทฺริยสฺส อตฺถงฺคโม โหติ, วีริยินฺทฺริยสฺส วเสน เอกตฺต-อนุปฏฺฐานํ วีริยินฺทฺริยสฺส อตฺถงฺคโม โหติ.\csromanspacer{}\csromanspacer{}อุปฏฺฐานตฺถาย อาวชฺชนาย อตฺถงฺคโม สตินฺทฺริยสฺส อตฺถงฺคโม โหติ, อุปฏฺฐานวเสน ฉนฺทสฺส อตฺถงฺคโม สตินฺทฺริยสฺส อตฺถงฺคโม โหติ, อุปฏฺฐานวเสน มนสิการสฺส อตฺถงฺคโม สตินฺทฺริยสฺส อตฺถงฺคโม โหติ.\csromanspacer{} สตินฺทฺริยสฺส}

\prose{\csromanfolio{204}วเสน เอกตฺต-อนุปฏฺฐานํ สตินฺทฺริยสฺส อตฺถงฺคโม โหติ.\csromanspacer{}\csromanspacer{}อวิกฺเขปตฺถาย อาวชฺชนาย อตฺถงฺคโม สมาธินฺทฺริยสฺส อตฺถงฺคโม โหติ, อวิกฺเขปวเสน ฉนฺทสฺส อตฺถงฺคโม สมาธินฺทฺริยสฺส อตฺถงฺคโม โหติ, อวิกฺเขปวเสน มนสิการสฺส อตฺถงฺคโม สมาธินฺทฺริยสฺส อตฺถงฺคโม โหติ.\csromanspacer{} สมาธินฺทฺริยสฺส วเสน เอกตฺต-อนุปฏฺฐานํ สมาธินฺทฺริยสฺส อตฺถงฺคโม โหติ.\csromanspacer{}\csromanspacer{}ทสฺสนตฺถาย อาวชฺชนาย อตฺถงฺคโม ปญฺญินฺทฺริยสฺส อตฺถงฺคโม โหติ, ทสฺสนวเสน ฉนฺทสฺส อตฺถงฺคโม ปญฺญินฺทฺริยสฺส อตฺถงฺคโม โหติ, ทสฺสนวเสน มนสิการสฺส อตฺถงฺคโม ปญฺญินฺทฺริยสฺส อตฺถงฺคโม โหติ.\csromanspacer{} ปญฺญินฺทฺริยสฺส วเสน เอกตฺต-อนุปฏฺฐานํ ปญฺญินฺทฺริยสฺส อตฺถงฺคโม โหติ.}

\prose{อธิโมกฺขตฺถาย อาวชฺชนาย อตฺถงฺคโม สทฺธินฺทฺริยสฺส อตฺถงฺคโม โหติ, ปคฺคหตฺถาย อาวชฺชนาย อตฺถงฺคโม วีริยินฺทฺริยสฺส อตฺถงฺคโม โหติ, อุปฏฺฐานตฺถาย อาวชฺชนาย อตฺถงฺคโม สตินฺทฺริยสฺส อตฺถงฺคโม โหติ, อวิกฺเขปตฺถาย อาวชฺชนาย อตฺถงฺคโม สมาธินฺทฺริยสฺส อตฺถงฺคโม โหติ, ทสฺสนตฺถาย อาวชฺชนาย อตฺถงฺคโม ปญฺญินฺทฺริยสฺส อตฺถงฺคโม โหติ.\csromanspacer{}\csromanspacer{}อธิโมกฺขวเสน ฉนฺทสฺส อตฺถงฺคโม สทฺธินฺทฺริยสฺส อตฺถงฺคโม โหติ, ปคฺคหวเสน ฉนฺทสฺส อตฺถงฺคโม วีริยินฺทฺริยสฺส อตฺถงฺคโม โหติ, อุปฏฺฐานวเสน ฉนฺทสฺส อตฺถงฺคโม สตินฺทฺริยสฺส อตฺถงฺคโม โหติ, อวิกฺเขปวเสน ฉนฺทสฺส อตฺถงฺคโม สมาธินฺทฺริยสฺส อตฺถงฺคโม โหติ, ทสฺสนวเสน ฉนฺทสฺส อตฺถงฺคโม ปญฺญินฺทฺริยสฺส อตฺถงฺคโม โหติ.\csromanspacer{}\csromanspacer{}อธิโมกฺขวเสน มนสิการสฺส อตฺถงฺคโม สทฺธินฺทฺริยสฺส อตฺถงฺคโม โหติ, ปคฺคหวเสน มนสิการสฺส อตฺถงฺคโม วีริยินฺทฺริยสฺส อตฺถงฺคโม โหติ, อุปฏฺฐานวเสน มนสิการสฺส อตฺถงฺคโม สตินฺทฺริยสฺส อตฺถงฺคโม โหติ, อวิกฺเขปวเสน มนสิการสฺส อตฺถงฺคโม สมาธินฺทฺริยสฺส อตฺถงฺคโม โหติ, ทสฺสนวเสน มนสิการสฺส อตฺถงฺคโม ปญฺญินฺทฺริยสฺส อตฺถงฺคโม โหติ.\csromanspacer{}\csromanspacer{}สทฺธินฺทฺริยสฺส วเสน เอกตฺต-อนุปฏฺฐานํ สทฺธินฺทฺริยสฺส อตฺถงฺคโม โหติ.\csromanspacer{} วีริยินฺทฺริยสฺส วเสน เอกตฺต-อนุปฏฺฐานํ วีริยินฺทฺริยสฺส อตฺถงฺคโม โหติ.\csromanspacer{} สตินฺทฺริยสฺส วเสน เอกตฺต-อนุปฏฺฐานํ สตินฺทฺริยสฺส อตฺถงฺคโม โหติ.\csromanspacer{} สมาธินฺทฺริยสฺส วเสน เอกตฺต-อนุปฏฺฐานํ สมาธินฺทฺริยสฺส อตฺถงฺคโม โหติ.\csromanspacer{} ปญฺญินฺทฺริยสฺส วเสน เอกตฺต-อนุปฏฺฐานํ ปญฺญินฺทฺริยสฺส อตฺถงฺคโม โหติ.}

\csromanheader{2}{4. อินฺทฺริยกถา}
\prose{\csromanfolio{205}อิเมหิ จตฺตารีสาย อากาเรหิ ปญฺจนฺนํ อินฺทฺริยานํ อตฺถงฺคโม โหติ.\csromanspacer{} อิเมหิ จตฺตารีสาย อากาเรหิ ปญฺจนฺนํ อินฺทฺริยานํ อตฺถงฺคมํ ปชานาติ.}

\csromanheader{2}{\textbf{ก. อสฺสาทนิทฺเทส}}
\proseitem{189}{กตเมหิ ปญฺจวีสติยา อากาเรหิ ปญฺจนฺนํ อินฺทฺริยานํ อสฺสาโท โหติ, กตเมหิ ปญฺจวีสติยา อากาเรหิ ปญฺจนฺนํ อินฺทฺริยานํ อสฺสาทํ ปชานาติ? อสฺสทฺธิยสฺส อนุปฏฺฐานํ สทฺธินฺทฺริยสฺส อสฺสาโท โหติ.\csromanspacer{} อสฺสทฺธิยปริฬาหสฺส อนุปฏฺฐานํ สทฺธินฺทฺริยสฺส อสฺสาโท โหติ.\csromanspacer{} อธิโมกฺขจริยาย เวสารชฺชํ สทฺธินฺทฺริยสฺส อสฺสาโท โหติ.\csromanspacer{} สนฺโต จ วิหาราธิคโม สทฺธินฺทฺริยสฺส อสฺสาโท โหติ.\csromanspacer{} ยํ สทฺธินฺทฺริยํ ปฏิจฺจ อุปฺปชฺชติ สุขํ โสมนสฺสํ, อยํ สทฺธินฺทฺริยสฺส อสฺสาโท.}

\prose{โกสชฺชสฺส อนุปฏฺฐานํ วีริยินฺทฺริยสฺส อสฺสาโท โหติ.\csromanspacer{} โกสชฺชปริฬาหสฺส อนุปฏฺฐานํ วีริยินฺทฺริยสฺส อสฺสาโท โหติ.\csromanspacer{} ปคฺคหจริยาย เวสารชฺชํ วีริยินฺทฺริยสฺส อสฺสาโท โหติ.\csromanspacer{} สนฺโต จ วิหาราธิคโม วีริยินฺทฺริยสฺส อสฺสาโท โหติ.\csromanspacer{} ยํ วีริยินฺทฺริยํ ปฏิจฺจ อุปฺปชฺชติ สุขํ โสมนสฺสํ, อยํ วีริยินฺทฺริยสฺส อสฺสาโท.}

\prose{ปมาทสฺส อนุปฏฺฐานํ สตินฺทฺริยสฺส อสฺสาโท โหติ.\csromanspacer{} ปมาทปริฬาหสฺส อนุปฏฺฐานํ สตินฺทฺริยสฺส อสฺสาโท โหติ.\csromanspacer{} อุปฏฺฐานจริยาย เวสารชฺชํ สตินฺทฺริยสฺส อสฺสาโท โหติ.\csromanspacer{} สนฺโต จ วิหาราธิคโม สตินฺทฺริยสฺส อสฺสาโท โหติ.\csromanspacer{} ยํ สตินฺทฺริยํ ปฏิจฺจ อุปฺปชฺชติ สุขํ โสมนสฺสํ, อยํ สตินฺทฺริยสฺส อสฺสาโท.}

\prose{อุทฺธจฺจสฺส อนุปฏฺฐานํ สมาธินฺทฺริยสฺส อสฺสาโท โหติ.\csromanspacer{} อุทฺธจฺจปริฬาหสฺส อนุปฏฺฐานํ สมาธินฺทฺริยสฺส อสฺสาโท โหติ.\csromanspacer{} อวิกฺเขปจริยาย เวสารชฺชํ สมาธินฺทฺริยสฺส อสฺสาโท โหติ.\csromanspacer{} สนฺโต จ วิหาราธิคโม สมาธินฺทฺริยสฺส อสฺสาโท โหติ.\csromanspacer{} ยํ สมาธินฺทฺริยํ ปฏิจฺจ อุปฺปชฺชติ สุขํ โสมนสฺส, อยํ สมาธินฺทฺริยสฺส อสฺสาโท.}

\prose{อวิชฺชาย อนุปฏฺฐานํ ปญฺญินฺทฺริยสฺส อสฺสาโท โหติ.\csromanspacer{} อวิชฺชาปริฬาหสฺส อนุปฏฺฐานํ ปญฺญินฺทฺริยสฺส อสฺสาโท โหติ.\csromanspacer{} ทสฺสนจริยาย เวสารชฺชํ ปญฺญินฺทฺริยสฺส อสฺสาโท โหติ.\csromanspacer{} สนฺโต จ วิหาราธิคโม \csromanfolio{206}ปญฺญินฺทฺริยสฺส อสฺสาโท โหติ.\csromanspacer{} ยํ ปญฺญินฺทฺริยํ ปฏิจฺจ อุปฺปชฺชติ สุขํ โสมนสฺสํ, อยํ ปญฺญินฺทฺริยสฺส อสฺสาโท.}

\prose{อิเมหิ ปญฺจวีสติยา อากาเรหิ ปญฺจนฺนํ อินฺทฺริยานํ อสฺสาโท โหติ, อิเมหิ ปญฺจวีสติยา อากาเรหิ ปญฺจนฺนํ อินฺทฺริยานํ อสฺสาทํ ปชานาติ.}

\csromanheader{2}{\textbf{ข. อาทีนวนิทฺเทส}}
\proseitem{190}{กตเมหิ ปญฺจวีสติยา อากาเรหิ ปญฺจนฺนํ อินฺทฺริยานํ อาทีนโว โหติ, กตเมหิ ปญฺจวีสติยา อากาเรหิ ปญฺจนฺนํ อินฺทฺริยานํ อาทีนวํ ปชานาติ? อสฺสทฺธิยสฺส อุปฏฺฐานํ สทฺธินฺทฺริยสฺส อาทีนโว โหติ.\csromanspacer{} อสฺสทฺธิยปริฬาหสฺส อุปฏฺฐานํ สทฺธินฺทฺริยสฺส อาทีนโว โหติ.\csromanspacer{} อนิจฺจฏฺเฐน สทฺธินฺทฺริยสฺส อาทีนโว โหติ.\csromanspacer{} ทุกฺขฏฺเฐน สทฺธินฺทฺริยสฺส อาทีนโว โหติ.\csromanspacer{} อนตฺตฏฺเฐน สทฺธินฺทฺริยสฺส อาทีนโว โหติ.}

\prose{โกสชฺชสฺส อุปฏฺฐานํ วีริยินฺทฺริยสฺส อาทีนโว โหติ.\csromanspacer{} โกสชฺชปริฬาหสฺส อุปฏฺฐานํ วีริยินฺทฺริยสฺส อาทีนโว โหติ.\csromanspacer{} อนิจฺจฏฺเฐน วีริยินฺทฺริยสฺส อาทีนโว โหติ.\csromanspacer{} ทุกฺขฏฺเฐน ฯเปฯ.\csromanspacer{} อนตฺตฏฺเฐน วีริยินฺทฺริสฺส อาทีนโว โหติ.}

\prose{ปมาทสฺส อุปฏฺฐานํ สตินฺทฺริยสฺส อาทีนโว โหติ.\csromanspacer{} ปมาทปริฬาหสฺส อุปฏฺฐานํ สตินฺทฺริยสฺส อาทีนโว โหติ.\csromanspacer{} อนิจฺจฏฺเฐน สตินฺทฺริยสฺส อาทีนโว โหติ.\csromanspacer{} ทุกฺขฏฺเฐน ฯเปฯ.\csromanspacer{} อนตฺตฏฺเฐน สตินฺทฺริยสฺส อาทีนโว โหติ.}

\prose{อุทฺธจฺจสฺส อุปฏฺฐานํ สมาธินฺทฺริยสฺส อาทีนโว โหติ.\csromanspacer{} อุทฺธจฺจปริฬาหสฺส อุปฏฺฐานํ สมาธินฺทฺริยสฺส อาทีนโว โหติ.\csromanspacer{} อนิจฺจฏฺเฐน สมาธินฺทฺริยสฺส อาทีนโว โหติ.\csromanspacer{} ทุกฺขฏฺเฐน ฯเปฯ.\csromanspacer{} อนตฺตฏฺเฐน สมาธินฺทฺริยสฺส อาทีนโว โหติ.}

\prose{อวิชฺชาย อุปฏฺฐานํ ปญฺญินฺทฺริยสฺส อาทีนโว โหติ.\csromanspacer{} อวิชฺชาปริฬาหสฺส อุปฏฺฐานํ ปญฺญินฺทฺริยสฺส อาทีนโว โหติ.\csromanspacer{} อนิจฺจฏฺเฐน ปญฺญินฺทฺริยสฺส อาทีนโว โหติ.\csromanspacer{} ทุกฺขฏฺเฐน ปญฺญินฺทฺริยสฺส อาทีนโว โหติ.\csromanspacer{} อนตฺตฏฺเฐน ปญฺญินฺทฺริยสฺส อาทีนโว โหติ.}

\csromanheader{2}{4. อินฺทฺริยกถา}
\prose{\csromanfolio{207}อิเมหิ ปญฺจวีสติยา อากาเรหิ ปญฺจนฺนํ อินฺทฺริยานํ อาทีนโว โหติ, อิเมหิ ปญฺจวีสติยา อากาเรหิ ปญฺจนฺนํ อินฺทฺริยานํ อาทีนวํ ปชานาติ.}

\csromanheader{2}{8. นิสฺสรณนิทฺเทส}
\proseitem{191}{กตเมหิ อสีติสตํ อากาเรหิ ปญฺจนฺนํ อินฺทฺริยานํ นิสฺสรณํ โหติ, กตเมหิ อสีติสตํ อากาเรหิ ปญฺจนฺนํ อินฺทฺริยานํ นิสฺสรณํ ปชานาติ? อธิโมกฺขฏฺเฐน สทฺธินฺทฺริยํ อสฺสทฺธิยา นิสฺสฏํ โหติ.\csromanspacer{} อสฺสทฺธิยปริฬาหา นิสฺสฏํ โหติ, ตทนุวตฺตกกิเลเสหิ จ ขนฺเธหิ จ นิสฺสฏํ โหติ, พหิทฺธา จ สพฺพนิมิตฺเตหิ นิสฺสฏํ โหติ, ตโต ปณีตตรสทฺธินฺทฺริยสฺส ปฏิลาภา ปุริมตรสทฺธินฺทฺริยา นิสฺสฏํ โหติ.\csromanspacer{} ปคฺคหฏฺเฐน วีริยินฺทฺริยํ โกสชฺชา นิสฺสฏํ โหติ, โกสชฺชปริฬาหา นิสฺสฏํ โหติ, ตทนุวตฺตกกิเลเสหิ จ ขนฺเธหิ จ นิสฺสฏํ โหติ, พหิทฺธา จ สพฺพนิมิตฺเตหิ นิสฺสฏํ โหติ, ตโต ปณีตตรวีริยินฺทฺริยสฺส ปฏิลาภา ปุริมตรวีริยินฺทฺริยา นิสฺสฏํ โหติ.\csromanspacer{} อุปฏฺฐานฏฺเฐน สตินฺทฺริยํ ปมาทา นิสฺสฏํ โหติ, ปมาทปริฬาหา นิสฺสฏํ โหติ, ตทนุวตฺตกกิเลเสหิ จ ขนฺเธหิ จ นิสฺสฏํ โหติ, พหิทฺธา จ สพฺพนิมิตฺเตหิ นิสฺสฏํ โหติ, ตโต ปณีตตรสตินฺทฺริยสฺส ปฏิลาภา ปุริมตรสตินฺทฺริยา นิสฺสฏํ โหติ.\csromanspacer{} อวิกฺเขปฏฺเฐน สมาธินฺทฺริยํ อุทฺธจฺจา นิสฺสฏํ โหติ, อุทฺธจฺจปริฬาหา นิสฺสฏํ โหติ, ตทนุวตฺตกกิเลเสหิ จ ขนฺเธหิ จ นิสฺสฏํ โหติ, พหิทฺธา จ สพฺพนิมิตฺเตหิ นิสฺสฏํ โหติ, ตโต ปณีตตรสมาธินฺทฺริยสฺส ปฏิลาภา ปุริมตรสมาธินฺทฺริยา นิสฺสฏํ โหติ.\csromanspacer{} ทสฺสนฏฺเฐน ปญฺญินฺทฺริยํ อวิชฺชาย นิสฺสฏํ โหติ, อวิชฺชาปริฬาหา นิสฺสฏํ โหติ, ตทนุวตฺตกกิเลเสหิ จ ขนฺเธหิ จ นิสฺสฏํ โหติ, พหิทฺธา จ สพฺพนิมิตฺเตหิ นิสฺสฏํ โหติ, ตโต ปณีตตรปญฺญินฺทฺริยสฺส ปฏิลาภา ปุริมตรปญฺญินฺทฺริยา นิสฺสฏํ โหติ.}

\proseitem{192}{ปุพฺพภาเค ปญฺจหิ อินฺทฺริเยหิ ปฐมชฺฌานวเสน ปญฺจินฺทฺริยานิ นิสฺสฏานิ โหนฺติ, ปฐเม ฌาเน ปญฺจหินฺทฺริเยหิ ทุติยชฺฌานวเสน ปญฺจินฺทฺริยานิ นิสฺสฏานิ โหนฺติ.\csromanspacer{} ทุติเย ฌาเน ปญฺจหินฺทฺริเยหิ ตติยชฺฌานวเสน ปญฺจินฺทฺริยานิ นิสฺสฏานิ โหนฺติ, ตติเย ฌาเน ปญฺจหินฺทฺริเยหิ จตุตฺถชฺฌานวเสน ปญฺจินฺทฺริยานิ นิสฺสาฏานิ โหนฺติ, จตุตฺเถ ฌาเน ปญฺจหินฺทฺริเยหิ อากาสานญฺจายตนสมาปตฺติวเสน ปญฺจินฺทฺริยานิ นิสฺสฏานิ \csromanfolio{208}โหนฺติ, อากาสานญฺจายตนสมาปตฺติยา ปญฺจหินฺทฺริเยหิ วิญฺญาณญฺจายตนสมาปตฺติวเสน ปญฺจินฺทฺริยานิ นิสฺสฏานิ โหนฺติ, วิญฺญาณญฺจายตนสมาปตฺติยา ปญฺจหินฺทฺริเยหิ อากิญฺจญฺญายตนสมาปตฺติวเสน ปญฺจินฺทฺริยานิ นิสฺสฏานิ โหนฺติ, อากิญฺจญฺญายตนสมาปตฺติยา ปญฺจหินฺทฺริเยหิ เนวสญฺญานาสญฺญายตนสมาปตฺติวเสน ปญฺจินฺทฺริยานิ นิสฺสฏานิ โหนฺติ, เนวสญฺญานาสญฺญายตนสมาปตฺติยา ปญฺจหินฺทฺริเยหิ อนิจฺจานุปสฺสนาวเสน ปญฺจินฺทฺริยานิ นิสฺสฏานิ โหนฺติ. (8)}

\prose{อนิจฺจานุปสฺสนาย ปญฺจหินฺทฺริเยหิ ทุกฺขานุปสฺสนาวเสน ปญฺจินฺทฺริยานิ นิสฺสฏานิ โหนฺติ, ทุกฺขานุปสฺสนาย ปญฺจหินฺทฺริเยหิ อนตฺตานุปสฺสนาวเสน ปญฺจินฺทฺริยานิ นิสฺสฏานิ โหนฺติ, อนตฺตานุปสฺสนาย ปญฺจหินฺทฺริเยหิ นิพฺพิทานุปสฺสนาวเสน ปญฺจินฺทฺริยานิ นิสฺสฏานิ โหนฺติ, นิพฺพิทานุปสฺสนาย ปญฺจหินฺทฺริเยหิ วิราคานุปสฺสนาวเสน ปญฺจินฺทฺริยานิ นิสฺสฏานิ โหนฺติ, วิราคานุปสฺสนาย ปญฺจหินฺทฺริเยหิ นิโรธานุปสฺสนาวเสน ปญฺจินฺทฺริยานิ นิสฺสฏานิ โหนฺติ, นิโรธานุปสฺสนาย ปญฺจหินฺทฺริเยหิ ปฏินิสฺสคฺคานุปสฺสนาวเสน ปญฺจินฺทฺริยานิ นิสฺสฏานิ โหนฺติ, ปฏินิสฺสคฺคานุปสฺสนาย ปญฺจหินฺทฺริเยหิ ขยานุปสฺสนาวเสน ปญฺจินฺทฺริยานิ นิสฺสฏานิ โหนฺติ, ขยานุปสฺสนาย ปญฺจหินฺทฺริเยหิ วยานุปสฺสนาวเสน ปญฺจินฺทฺริยานิ นิสฺสฏานิ โหนฺติ, วยานุปสฺสนาย ปญฺจหินฺทฺริเยหิ วิปริณามานุปสฺสนาวเสน ปญฺจินฺทฺริยานิ นิสฺสฏานิ โหนฺติ, วิปริณามานุปสฺสนาย ปญฺจหินฺทฺริเยหิ อนิมิตฺตานุปสฺสนาวเสน ปญฺจินฺทฺริยานิ นิสฺสฏานิ โหนฺติ, อนิมิตฺตานุปสฺสนาย ปญฺจหินฺทฺริเยหิ อปฺปณิหิตานุปสฺสนาวเสน ปญฺจินฺทฺริยานิ นิสฺสฏานิ โหนฺติ, อปฺปณิหิตานุปสฺสนาย ปญฺจหินฺทฺริเยหิ สุญฺญตานุปสฺสนาวเสน ปญฺจินฺทฺริยานิ นิสฺสฏานิ โหนฺติ, สุญฺญตานุปสฺสนาย ปญฺจหินฺทฺริเยหิ อธิปญฺญาธมฺมวิปสฺสนาวเสน ปญฺจินฺทฺริยานิ นิสฺสฏานิ โหนฺติ.\csromanspacer{} อธิปญฺญาธมฺมวิปสฺสนาย ปญฺจหินฺทฺริเยหิ ยถาภูตญาณทสฺสนวเสน ปญฺจินฺทฺริยานิ นิสฺสฏานิ โหนฺติ, ยถาภูตญาณทสฺสเน ปญฺจหินฺทฺริเยหิ อาทีนวานุปสฺสนาวเสน ปญฺจินฺทฺริยานิ นิสฺสฏานิ โหนฺติ, อาทีนวานุปสฺสนาย ปญฺจหินฺทฺริเยหิ ปฏิสงฺขานุปสฺสนาวเสน ปญฺจินฺทฺริยานิ นิสฺสฏานิ โหนฺติ, ปฏิสงฺขานุปสฺสนาย ปญฺจหินฺทฺริเยหิ วิวฏฺฏนานุปสฺสนาวเสน ปญฺจินฺทฺริยานิ นิสฺสฏานิ โหนฺติ.\csromanspacer{} วิวฏฺฏนานุปสฺสนาย ปญฺจหินฺทฺริเยหิ โสตาปตฺติมคฺควเสน ปญฺจินฺทฺริยานิ นิสฺสฏานิ โหนฺติ. (18, 26)}

\csromanheader{2}{4. อินฺทฺริยกถา}
\prose{\csromanfolio{209}โสตาปตฺติมคฺเค ปญฺจหินฺทฺริเยหิ โสตาปตฺติผลสมาปตฺติวเสน ปญฺจินฺทฺริยานิ นิสฺสฏานิ โหนฺติ, โสตาปตฺติผลสมาปตฺติยา ปญฺจหินฺทฺริเยหิ สกทาคามิมคฺควเสน ปญฺจินฺทฺริยานิ นิสฺสฏานิ โหนฺติ, สกทาคามิมคฺเค ปญฺจหินฺทฺริเยหิ สกทาคามิผลสมาปตฺติวเสน ปญฺจินฺทฺริยานิ นิสฺสฏานิ โหนฺติ, สกทามิผลสมาปตฺติยา ปญฺจหินฺทฺริเยหิ อนาคามิมคฺควเสน ปญฺจินฺทฺริยานิ นิสฺสฏานิ โหนฺติ, อนาคามิมคฺเค ปญฺจหินฺทฺริเยหิ อนาคามิผลสมาปตฺติวเสน ปญฺจินฺทฺริยานิ นิสฺสฏานิ โหนฺติ, อนาคามิผลสมาปตฺติยา ปญฺจหินฺทฺริเยหิ อรหตฺตมคฺควเสน ปญฺจินฺทฺริยานิ นิสฺสฏานิ โหนฺติ, อรหตฺตมคฺเค ปญฺจหินฺทฺริเยหิ อรหตฺตผลสมาปตฺติวเสน ปญฺจินฺทฺริยานิ นิสฺสฏานิ โหนฺติ. (8, 34)}

\prose{เนกฺขมฺเม ปญฺจินฺทฺริยานิ กามจฺฉนฺทโต นิสฺสฏานิ โหนฺติ, อพฺยาปาเท ปญฺจินฺทฺริยานิ พฺยาปาทโต นิสฺสฏานิ โหนฺติ, อาโลกสญฺญาย ปญฺจินฺทฺริยานิ ถินมิทฺธโต นิสฺสฏานิ โหนฺติ, อวิกฺเขเป ปญฺจินฺทฺริยานิ อุทฺธจฺจโต นิสฺสฏานิ โหนฺติ, ธมฺมววตฺถาเน ปญฺจินฺทฺริยานิ วิจิกิจฺฉาย นิสฺสฏานิ โหนฺติ, ญาเณ ปญฺจินฺทฺริยานิ อวิชฺชาย นิสฺสฏานิ โหนฺติ, ปาโมชฺเช ปญฺจินฺทฺริยานิ อรติยา นิสฺสฏานิ โหนฺติ.}

\proseitem{193}{ปฐเม ฌาเน ปญฺจินฺทฺริยานิ นีวรเณหิ นิสฺสฏานิ โหนฺติ, ทุติเย ฌาเน ปญฺจินฺทฺริยานิ วิตกฺกวิจาเรหิ นิสฺสฏานิ โหนฺติ, ตติเย ฌาเน ปญฺจินฺทฺริยานิ ปีติยา นิสฺสฏานิ โหนฺติ, จตุตฺเถ ฌาเน ปญฺจินฺทฺริยานิ สุขทุกฺเขหิ นิสฺสฏานิ โหนฺติ, อากาสานญฺจายตนสมาปตฺติยา ปญฺจินฺทฺริยานิ รูปสญฺญาย ปฏิฆสญฺญาย นานตฺตสญฺญาย นิสฺสฏานิ โหนฺติ, วิญฺญาณญฺจายตนสมาปตฺติยา ปญฺจินฺทฺริยานิ อากาสานญฺจายตนสญฺญาย นิสฺสฏานิ โหนฺติ, อากิญฺจญฺญายตนสมาปตฺติยา ปญฺจินฺทฺริยานิ วิญฺญาณญฺจายตนสญฺญาย นิสฺสฏานิ โหนฺติ, เนวสญฺญานาสญฺญายตนสมาปตฺติยา ปญฺจินฺทฺริยานิ อากิญฺจญฺญายตนสญฺญาย นิสฺสฏานิ โหนฺติ.}

\prose{อนิจฺจานุปสฺสนาย ปญฺจินฺทฺริยานิ นิจฺจสญฺญาย นิสฺสฏานิ โหนฺติ, ทุกฺขานุปสฺสนาย ปญฺจินฺทฺริยานิ สุขสญฺญาย นิสฺสฏานิ โหนฺติ, อนตฺตานุปสฺสนาย ปญฺจินฺทฺริยานิ อตฺตสญฺญาย นิสฺสฏานิ โหนฺติ, นิพฺพิทานุปสฺสนาย ปญฺจินฺทฺริยานิ นนฺทิยา นิสฺสฏานิ โหนฺติ, วิราคานุปสฺสนาย ปญฺจินฺทฺริยานิ ราคโต นิสฺสฏานิ โหนฺติ, นิโรธานุปสฺสนาย ปญฺจินฺทฺริยานิ สมุทยโต \csromanfolio{210}นิสฺสฏานิ โหนฺติ, ปฏินิสฺสคฺคานุปสฺสนาย ปญฺจินฺทฺริยานิ อาทานโต นิสฺสฏานิ โหนฺติ, ขยานุปสฺสนาย ปญฺจินฺทฺริยานิ ฆนสญฺญาย นิสฺสฏานิ โหนฺติ, วยานุปสฺสนาย ปญฺจินฺทฺริยานิ อายูหนโต นิสฺสฏานิ โหนฺติ, วิปริณามานุปสฺสนาย ปญฺจินฺทฺริยานิ ธุวสญฺญาย นิสฺสฏานิ โหนฺติ, อนิมิตฺตานุปสฺสนาย ปญฺจินฺทฺริยานิ นิมิตฺตโต นิสฺสฏานิ โหนฺติ, อปฺปณิหิตานุปสฺสนาย ปญฺจินฺทฺริยานิ ปณิธิยา นิสฺสฏานิ โหนฺติ, สุญฺญตานุปสฺสนาย ปญฺจินฺทฺริยานิ อภินิเวสโต นิสฺสฏานิ โหนฺติ.\csromanspacer{} อธิปญฺญาธมฺมวิปสฺสนาย ปญฺจินฺทฺริยานิ สาราทานาภินิเวสโต นิสฺสฏานิ โหนฺติ, ยถาภูตญาณทสฺสเน ปญฺจินฺทฺริยานิ สมฺโมหาภินิเวสโต นิสฺสฏานิ โหนฺติ, อาทีนวานุปสฺสนาย ปญฺจินฺทฺริยานิ อาลยาภินิเวสโต นิสฺสฏานิ โหนฺติ, ปฏิสงฺขานุปสฺสนาย ปญฺจินฺทฺริยานิ อปฺปฏิสงฺขาย นิสฺสฏานิ โหนฺติ, วิวฏฺฏนานุปสฺสนาย ปญฺจินฺทฺริยานิ สญฺโญคาภินิเวสโต นิสฺสฏานิ โหนฺติ.}

\prose{โสตาปตฺติมคฺเค ปญฺจินฺทฺริยานิ ทิฏฺเฐกฏฺเฐหิ กิเลเสหิ นิสฺสฏานิ โหนฺติ, สกทาคามิมคฺเค ปญฺจินฺทฺริยานิ โอฬาริเกหิ กิเลเสหิ นิสฺสฏานิ โหนฺติ, อนาคามิมคฺเค ปญฺจินฺทฺริยานิ อณุสหคเตหิ กิเลเสหิ นิสฺสฏานิ โหนฺติ, อรหตฺตมคฺเค ปญฺจินฺทฺริยานิ สพฺพกิเลเสหิ นิสฺสฏานิ โหนฺติ, สพฺเพสญฺเญว ขีณาสวานํ ตตฺถ ตตฺถ ปญฺจินฺทฺริยานิ นิสฺสฏานิ เจว โหนฺติ สุนิสฺสฏานิ จ ปฏิปฺปสฺสทฺธานิ จ สุปฺปฏิปฺปสฺสทฺธานิ จ.}

\prose{อิเมหิ อสีติสตํ อากาเรหิ ปญฺจนฺนํ อินฺทฺริยานํ นิสฺสรณํ โหติ, อิเมหิ อสีติสตํ อากาเรหิ ปญฺจนฺนํ อินฺทฺริยานํ นิสฺสรณํ ปชานาติ.}

\vspace{\tipitakaparskip}
\csromancenter{สุตฺตนฺตนิทฺเทโส ทุติโย.}
\csromanheader{2}{ปฐมภาณวาโร.}
\csromansectionrule
\csromanmatikaanchor{mtk.75}
\csromantocmarkref{subsection}{3. ตติยสุตฺตนฺตนิทฺเทส}{mtk.75}
\bookmark[dest={mtk.75},level=2]{3. ตติยสุตฺตนฺตนิทฺเทส}
\csromanheader{2}{3. ตติยสุตฺตนฺตนิทฺเทส}
\proseitem{194}{สาวตฺถินิทานํ\footnote{สํ 3. 172 ปิฏฺเฐ ปสฺสิตพฺพา.}.\csromanspacer{}\csromanspacer{}ปญฺจิมานิ ภิกฺขเว อินฺทฺริยานิ.\csromanspacer{} กตมานิ ปญฺจ? สทฺธินฺทฺริยํ วีริยินฺทฺริยํ สตินฺทฺริยํ สมาธินฺทฺริยํ ปญฺญินฺทฺริยํ.\csromanspacer{} กตฺถ จ ภิกฺขเว สทฺธินฺทฺริยํ ทฏฺฐพฺพํ? จตูสุ โสตาปตฺติยงฺคสุ เอตฺถ สทฺธินฺทฺริยํ ทฏฺฐพฺพํ.\csromanspacer{} กตฺถ จ ภิกฺขเว วีริยินฺทฺริยํ ทฏฺฐพฺพํ? จตูสุ สมฺมปฺปธาเนสุ เอตฺถ วีริยินฺทฺริยํ ทฏฺฐพฺพํ.\csromanspacer{} กตฺถ}

\csromanheader{2}{4. อินฺทฺริยกถา}
\prose{\csromanfolio{211}จ ภิกฺขเว สตินฺทฺริยํ ทฏฺฐพฺพํ? จตูสุ สติปฏฺฐาเนสุ เอตฺถ สตินฺทฺริยํ ทฏฺฐพฺพํ, กตฺถ จ ภิกฺขเว สมาธินฺทฺริยํ ทฏฺฐพฺพํ? จตูสุ ฌาเนสุ เอตฺถ สมาธินฺทฺริยํ ทฏฺฐพฺพํ.\csromanspacer{} กตฺถ จ ภิกฺขเว ปญฺญินฺทฺริยํ ทฏฺฐพฺพํ? จตูสุ อริยสจฺเจสุ เอตฺถ ปญฺญินฺทฺริยํ ทฏฺฐพฺพํ.}

\prose{จตูสุ โสตาปตฺติยงฺเคสุ สทฺธินฺทฺริยสฺส วเสน กติหากาเรหิ ปญฺจินฺทฺริยานิ ทฏฺฐพฺพานิ, จตูสุ สมฺมปฺปธาเนสุ วีริยินฺทฺริยสฺส วเสน กติหากาเรหิ ปญฺจินฺทฺริยานิ ทฏฺฐพฺพานิ, จตูสุ สติปฏฺฐาเนสุ สตินฺทฺริยสฺส วเสน กติหากาเรหิ ปญฺจินฺทฺริยานิ ทฏฺฐพฺพานิ, จตูสุ ฌาเนสุ สมาธินฺทฺริยสฺส วเสน กติหากาเรหิ ปญฺจินฺทฺริยานิ ทฏฺฐพฺพานิ, จตูสุ อริยสจฺเจสุ ปญฺญินฺทฺริยสฺส วเสน กติหากาเรหิ ปญฺจินฺทฺริยานิ ทฏฺฐพฺพานิ?}

\prose{จตูสุ โสตาปตฺติยงฺเคสุ สทฺธินฺทฺริยสฺส วเสน วีสติยา อากาเรหิ ปญฺจินฺทฺริยานิ ทฏฺฐพฺพานิ, จตูสุ สมฺมปฺปธาเนสุ วีริยินฺทฺริยสฺส วเสน วีสติยา อากาเรหิ ปญฺจินฺทฺริยานิ ทฏฺฐพฺพานิ, จตูสุ สติปฏฺฐาเนสุ สตินฺทฺริยสฺส วเสน วีสติยา อากาเรหิ ปญฺจินฺทฺริยานิ ทฏฺฐพฺพานิ, จตูสุ ฌาเนสุ สมาธินฺทฺริยสฺส วเสน วีสติยา อากาเรหิ ปญฺจินฺทฺริยานิ ทฏฺฐพฺพานิ, จตูสุ อริยสจฺเจสุ ปญฺญินฺทฺริยสฺส วเสน วีสติยา อากาเรหิ ปญฺจินฺทฺริยานิ ทฏฺฐพฺพานิ.}

\csromanheader{2}{\textbf{ก. ปเภทคณนนิทฺเทส}}
\proseitem{195}{จตูสุ โสตาปตฺติยงฺเคสุ สทฺธินฺทฺริยสฺส วเสน กตเมหิ วีสติยา อากาเรหิ ปญฺจินฺทฺริยานิ ทฏฺฐพฺพานิ? สปฺปุริสสํเสเว โสตาปตฺติยงฺเค อธิโมกฺขาธิปเตยฺยฏฺเฐน สทฺธินฺทฺริยํ ทฏฺฐพฺพํ, สทฺธินฺทฺริยสฺส วเสน ปคฺคหฏฺเฐน วีริยินฺทฺริยํ ทฏฺฐพฺพํ, อุปฏฺฐานฏฺเฐน สตินฺทฺริยํ ทฏฺฐพฺพํ, อวิกฺเขปฏฺเฐน สมาธินฺทฺริยํ ทฏฺฐพฺพํ, ทสฺสนฏฺเฐน ปญฺญินฺทฺริยํ ทฏฺฐพฺพํ, สทฺธมฺมสวเน โสตาปตฺติยงฺเค โยนิโสมนสิกาเร โสตาปตฺติยงฺเค ธมฺมานุธมฺมปฏิปตฺติยา โสตาปตฺติยงฺเค อธิโมกฺขาธิปเตยฺยฏฺเฐน สทฺธินฺทฺริยํ ทฏฺฐพฺพํ, สทฺธินฺทฺริยสฺส วเสน ปคฺคหฏฺเฐน วีริยินฺทฺริยํ ทฏฺฐพฺพํ, อุปฏฺฐานฏฺเฐน สตินฺทฺริยํ ทฏฺฐพฺพํ, อวิกฺเขปฏฺเฐน สมาธินฺทฺริยํ ทฏฺฐพฺพํ, ทสฺสนฏฺเฐน ปญฺญินฺทฺริยํ ทฏฺฐพฺพํ.\csromanspacer{} จตูสุ โสตาปตฺติยงฺเคสุ สทฺธินฺทฺริยสฺส วเสน อิเมหิ วีสติยา อากาเรหิ ปญฺจินฺทฺริยานิ ทฏฺฐพฺพานิ.}

\prose{\csromanfolio{212}จตูสุ สมฺมปฺปธาเนสุ วีริยินฺทฺริยสฺส วเสน กตเมหิ วีสติยา อากาเรหิ ปญฺจินฺทฺริยานิ ทฏฺฐพฺพานิ? อนุปฺปนฺนานํ ปาปกานํ อกุสลานํ ธมฺมานํ อนุปฺปาทาย สมฺมปฺปธาเน ปคฺคหาธิปเตยฺยฏฺเฐน วีริยินฺทฺริยํ ทฏฺฐพฺพํ, วีริยินฺทฺริยสฺส วเสน อุปฏฺฐานฏฺเฐน สตินฺทฺริยํ ทฏฺฐพฺพํ, อวิกฺเขปฏฺเฐน สมาธินฺทฺริยํ ทฏฺฐพฺพํ, ทสฺสนฏฺเฐน ปญฺญินฺทฺริยํ ทฏฺฐพฺพํ, อธิโมกฺขฏฺเฐน สทฺธินฺทฺริยํ ทฏฺฐพฺพํ.\csromanspacer{} อุปฺปนฺนานํ ปาปกานํ อกุสลานํ ธมฺมานํ ปหานาย สมฺมปฺปธาเน ฯเปฯ.\csromanspacer{} อนุปฺปนฺนานํ กุสลานํ ธมฺมานํ อุปฺปาทาย สมฺมปฺปธาเน ฯเปฯ.\csromanspacer{} อุปฺปนฺนานํ กุสลานํ ธมฺมานํ ฐิติยา อสมฺโมสาย ภิยฺโยภาวาย เวปุลฺลาย ภาวนาย ปาริปูริยา สมฺมปฺปธาเน ปคฺคหาธิปเตยฺยฏฺเฐน วีริยินฺทฺริยํ ทฏฺฐพฺพํ, วีริยินฺทฺริยสฺส วเสน อุปฏฺฐานฏฺเฐน สตินฺทฺริยํ ทฏฺฐพฺพํ, อวิกฺเขปฏฺเฐน สมาธินฺทฺริยํ ทฏฺฐพฺพํ, ทสฺสนฏฺเฐน ปญฺญินฺทฺริยํ ทฏฺฐพฺพํ, อธิโมกฺขฏฺเฐน สทฺธินฺทฺริยํ ทฏฺฐพฺพํ.\csromanspacer{} จตูสุ สมฺมปฺปธาเนสุ วีริยินฺทฺริยสฺส วเสน อิเมหิ วีสติยา อากาเรหิ ปญฺจินฺทฺริยานิ ทฏฺฐพฺพานิ.}

\prose{จตูสุ สติปฏฺฐาเนสุ สตินฺทฺริยสฺส วเสน กตเมหิ วีสติยา อากาเรหิ ปญฺจินฺทฺริยานิ ทฏฺฐพฺพานิ? กาเย กายานุปสฺสนาสติปฏฺฐาเน อุปฏฺฐานาธิปเตยฺยฏฺเฐน สตินฺทฺริยํ ทฏฺฐพฺพํ, สตินฺทฺริยสฺส วเสน อวิกฺเขปฏฺเฐน สมาธินฺทฺริยํ ทฏฺฐพฺพํ, ทสฺสนฏฺเฐน ปญฺญินฺทฺริยํ ทฏฺฐพฺพํ, อธิโมกฺขฏฺเฐน สทฺธินฺทฺริยํ ทฏฺฐพฺพํ, ปคฺคหฏฺเฐน วีริยินฺทฺริยํ ทฏฺฐพฺพํ.\csromanspacer{} เวทนาสุ เวทนานุปสฺสนาสติปฏฺฐาเน ฯเปฯ.\csromanspacer{} จิตฺเต จิตฺตานุปสฺสนาสติปฏฺฐาเน.\csromanspacer{} ธมฺเมสุ ธมฺมานุปสฺสนาสติปฏฺฐาเน อุปฏฺฐานาธิปเตยฺยฏฺเฐน สตินฺทฺริยํ ทฏฺฐพฺพํ, สตินฺทฺริยสฺส วเสน อวิกฺเขปฏฺเฐน สมาธินฺทฺริยํ ทฏฺฐพฺพํ, ทสฺสนฏฺเฐน ปญฺญินฺทฺริยํ ทฏฺฐพฺพํ, อธิโมกฺขฏฺเฐน สทฺธินฺทฺริยํ ทฏฺฐพฺพํ, ปคฺคหฏฺเฐน วีริยินฺทฺริยํ ทฏฺฐพฺพํ.\csromanspacer{} จตูสุ สติปฏฺฐาเนสุ สตินฺทฺริยสฺส วเสน อิเมหิ วีสติยา อากาเรหิ ปญฺจินฺทฺริยานิ ทฏฺฐพฺพานิ.}

\prose{จตูสุ ฌาเนสุ สมาธินฺทฺริยสฺส วเสน กตเมหิ วีสติยา อากาเรหิ ปญฺจินฺทฺริยานิ ทฏฺฐพฺพานิ? ปฐเม ฌาเน อวิกฺเขปาธิปเตยฺยฏฺเฐน สมาธินฺทฺริยํ ทฏฺฐพฺพํ, สมาธินฺทฺริยสฺส วเสน ทสฺสนฏฺเฐน ปญฺญินฺทฺริยํ ทฏฺฐพฺพํ, อธิโมกฺขฏฺเฐน สทฺธินฺทฺริยํ ทฏฺฐพฺพํ, ปคฺคหฏฺเฐน วีริยินฺทฺริยํ ทฏฺฐพฺพํ, อุปฏฺฐานฏฺเฐน สตินฺทฺริยํ ทฏฺฐพฺพํ.\csromanspacer{} ทุติเย ฌาเน ฯเปฯ.\csromanspacer{} ตติเย ฌาเน.\csromanspacer{} จตุตฺเถ ฌาเน อวิกฺเขปาธิปเตยฺยฏฺเฐน สมาธินฺทฺริยํ ทฏฺฐพฺพํ, สมาธินฺทฺริยสฺส วเสน ทสฺสนฏฺเฐน ปญฺญินฺทฺริยํ ทฏฺฐพฺพํ, อธิโมกฺขฏฺเฐน สทฺธินฺทฺริยํ ทฏฺฐพฺพํ, ปคฺคหฏฺเฐน}

\csromanheader{2}{4. อินฺทฺริยกถา}
\prose{\csromanfolio{213}วีริยินฺทฺริยํ ทฏฺฐพฺพํ, อุปฏฺฐานฏฺเฐน สตินฺทฺริยํ ทฏฺฐพฺพํ.\csromanspacer{} จตูสุ ฌาเนสุ สมาธินฺทฺริยสฺส วเสน อิเมหิ วีสติยา อากาเรหิ ปญฺจินฺทฺริยานิ ทฏฺฐพฺพานิ.}

\prose{จตูสุ อริยสจฺเจสุ ปญฺญินฺทฺริยสฺส วเสน กตเมหิ วีสติยา อากาเรหิ ปญฺจินฺทฺริยานิ ทฏฺฐพฺพานิ? ทุกฺเข อริยสจฺเจ ทสฺสนาธิปเตยฺยฏฺเฐน ปญฺญินฺทฺริยํ ทฏฺฐพฺพํ, ปญฺญินฺทฺริยสฺส วเสน อธิโมกฺขฏฺเฐน สทฺธินฺทฺริยํ ทฏฺฐพฺพํ, ปคฺคหฏฺเฐน วีริยินฺทฺริยํ ทฏฺฐพฺพํ.\csromanspacer{} อุปฏฺฐานฏฺเฐน สตินฺทฺริยํ ทฏฺฐพฺพํ.\csromanspacer{} อวิกฺเขปฏฺเฐน สมาธินฺทฺริยํ ทฏฺฐพฺพํ.\csromanspacer{} ทุกฺขสมุทเย อริยสจฺเจ ฯเปฯ.\csromanspacer{} ทุกฺขนิโรเธ อริยสจฺเจ.\csromanspacer{} ทุกฺขนิโรธคามินิยา ปฏิปทาย อริยสจฺเจ ทสฺสนาธิปเตยฺยฏฺเฐน ปญฺญินฺทฺริยํ ทฏฺฐพฺพํ, ปญฺญินฺทฺริยสฺส วเสน อธิโมกฺขฏฺเฐน สทฺธินฺทฺริยํ ทฏฺฐพฺพํ, ปคฺคหฏฺเฐน วีริยินฺทฺริยํ ทฏฺฐพฺพํ, อุปฏฺฐานฏฺเฐน สตินฺทฺริยํ ทฏฺฐพฺพํ, อวิกฺเขปฏฺเฐน สมาธินฺทฺริยํ ทฏฺฐพฺพํ.\csromanspacer{} จตูสุ อริยสจฺเจสุ ปญฺญินฺทฺริยสฺส วเสน อิเมหิ วีสติยา อากาเรหิ ปญฺจินฺทฺริยานิ ทฏฺฐพฺพานิ.}

\csromanheader{2}{\textbf{ข. จริยวาร}}
\proseitem{196}{จตูสุ โสตาปตฺติยงฺเคสุ สทฺธินฺทฺริยสฺส วเสน กติหากาเรหิ ปญฺจนฺนํ อินฺทฺริยานํ จริยา ทฏฺฐพฺพา? จตูสุ สมฺมปฺปธาเนสุ ฯเปฯ.\csromanspacer{} จตูสุ สติปฏฺฐาเนสุ.\csromanspacer{} จตูสุ ฌาเนสุ.\csromanspacer{} จตูสุ อริยสจฺเจสุ ปญฺญินฺทฺริยสฺส วเสน กติหากาเรหิ ปญฺจนฺนํ อินฺทฺริยานํ จริยา ทฏฺฐพฺพา.\csromanspacer{}\csromanspacer{}จตูสุ โสตาปตฺติยงฺเคสุ สทฺธินฺทฺริยสฺส วเสน วีสติยา อากาเรหิ ปญฺจนฺนํ อินฺทฺริยานํ จริยา ทฏฺฐพฺพา.\csromanspacer{} จตูสุ สมฺมปฺปธาเนสุ ฯเปฯ.\csromanspacer{} จตูสุ สติปฏฺฐาเนสุ จตูสุ ฌาเนสุ จตูสุ อริยสจฺเจสุ ปญฺญินฺทฺริยสฺส วเสน วีสติยา อากาเรหิ ปญฺจนฺนํ อินฺทฺริยานํ จริยา ทฏฺฐพฺพา.}

\prose{จตูสุ โสตาปตฺติยงฺเคสุ สทฺธินฺทฺริยสฺส วเสน กตเมหิ วีสติยา อากาเรหิ ปญฺจนฺนํ อินฺทฺริยานํ จริยา ทฏฺฐพฺพา? สปฺปุริสสํเสเว โสตาปตฺติยงฺเค อธิโมกฺขาธิปเตยฺยฏฺเฐน สทฺธินฺทฺริยสฺส จริยา ทฏฺฐพฺพา, สทฺธินฺทฺริยสฺส วเสน ปคฺคหฏฺเฐน วีริยินฺทฺริยสฺส จริยา ทฏฺฐพฺพา, อุปฏฺฐานฏฺเฐน สตินฺทฺริยสฺส จริยา ทฏฺฐพฺพา, อวิกฺเขปฏฺเฐน สมาธินฺทฺริยสฺส จริยา ทฏฺฐพฺพา.\csromanspacer{} ทสฺสนฏฺเฐน ปญฺญินฺทฺริยสฺส จริยา ทฏฺฐพฺพา.\csromanspacer{} สทฺธมฺมสวเน โสตาปตฺติยงฺเค ฯเปฯ.\csromanspacer{} โยนิโส มนสิกาเร โสตาปตฺติยงฺเค ธมฺมานุธมฺมปฏิปตฺติยา โสตาปตฺติยงฺเค อธิโมกฺขาธิปเตยฺยฏฺเฐน สทฺธินฺทฺริยสฺส จริยา ทฏฺฐพฺพา, สทฺธินฺทฺริยสฺส วเสน ปคฺคหฏฺเฐน วีริยินฺทฺริยสฺส \csromanfolio{214}จริยา ทฏฺฐพฺพา, อุปฏฺฐานฏฺเฐน สตินฺทฺริยสฺส จริยา ทฏฺฐพฺพา, อวิกฺเขปฏฺเฐน สมาธินฺทฺริยสฺส จริยา ทฏฺฐพฺพา, ทสฺสนฏฺเฐน ปญฺญินฺทฺริยสฺส จริยา ทฏฺฐพฺพา.\csromanspacer{} จตูสุ โสตาปตฺติยงฺเคสุ สทฺธินฺทฺริยสฺส วเสน อิเมหิ วีสติยา อากาเรหิ ปญฺจนฺนํ อินฺทฺริยานํ จริยา ทฏฺฐพฺพา.}

\prose{จตูสุ สมฺมปฺปธาเนสุ วีริยินฺทฺริยสฺส วเสน กตเมหิ วีสติยา อากาเรหิ ปญฺจนฺนํ อินฺทฺริยานํ จริยา ทฏฺฐพฺพา? อนุปฺปนฺนานํ ปาปกานํ อกุสลานํ ธมฺมานํ อนุปฺปาทาย สมฺมปฺปธาเน ปคฺคหาธิปเตยฺยฏฺเฐน วีริยินฺทฺริยสฺส จริยา ทฏฺฐพฺพา, วีริยินฺทฺริยสฺส วเสน อุปฏฺฐานฏฺเฐน สตินฺทฺริยสฺส จริยา ทฏฺฐพฺพา, อวิกฺเขปฏฺเฐน สมาธินฺทฺริยสฺส จริยา ทฏฺฐพฺพา, ทสฺสนฏฺเฐน ปญฺญินฺทฺริยสฺส จริยา ทฏฺฐพฺพา, อธิโมกฺขฏฺเฐน สทฺธินฺทฺริยสฺส จริยา ทฏฺฐพฺพา.\csromanspacer{} อุปฺปนฺนานํ ปาปกานํ อกุสลานํ ธมฺมานํ ปหานาย สมฺมปฺปธาเน ฯเปฯ.\csromanspacer{} อนุปฺปนฺนานํ กุสลานํ ธมฺมานํ อุปฺปาทาย สมฺมปฺปธาเน ฯเปฯ.\csromanspacer{} อุปฺปนฺนานํ กุสลานํ ธมฺมานํ ฐิติยา อสมฺโมสาย ภิยฺโยภาวาย เวปุลฺลาย ภาวนาย ปาริปูริยา สมฺมปฺปธาเน ปคฺคหาธิปเตยฺยฏฺเฐน วีริยินฺทฺริยสฺส จริยา ทฏฺฐพฺพา ฯเปฯ.\csromanspacer{} จตูสุ สมฺมปฺปธาเนสุ วีริยินฺทฺริยสฺสวเสน อิเมหิ วีสติยา อากาเรหิ ปญฺจนฺนํ อินฺทฺริยานํ จริยา ทฏฺฐพฺพา.}

\prose{จตูสุ สติปฏฺฐาเนสุ สตินฺทฺริยสฺส วเสน กตเมหิ วีสติยา อากาเรหิ ปญฺจนฺนํ อินฺทฺริยานํ จริยา ทฏฺฐพฺพา? กาเย กายานุปสฺสนาสติปฏฺฐาเน อุปฏฺฐานาธิปเตยฺยฏฺเฐน สตินฺทฺริยสฺส จริยา ทฏฺฐพฺพา, สตินฺทฺริยสฺส วเสน อวิกฺเขปฏฺเฐน สมาธินฺทฺริยสฺส จริยา ทฏฺฐพฺพา, ทสฺสนฏฺเฐน ปญฺญินฺทฺริยสฺส จริยา ทฏฺฐพฺพา, อธิโมกฺขฏฺเฐน สทฺธินฺทฺริยสฺส จริยา ทฏฺฐพฺพา, ปคฺคหฏฺเฐน วีริยินฺทฺริยสฺส จริยา ทฏฺฐพฺพา.\csromanspacer{} เวทนาสุ เวทนานุปสฺสนาสติปฏฺฐาเน ฯเปฯ.\csromanspacer{} จิตฺเต จิตฺตานุปสฺสนาสติปฏฺฐาเน ฯเปฯ.\csromanspacer{} ธมฺเมสุ ธมฺมานุปสฺสนา สติปฏฺฐาเน อุปฏฺฐานาธิปเตยฺยฏฺเฐน สตินฺทฺริยสฺส จริยา ทฏฺฐพฺพา ฯเปฯ.\csromanspacer{} จตูสุ สติปฏฺฐาเนสุ สตินฺทฺริยสฺส วเสน อิเมหิ วีสติยา อากาเรหิ ปญฺจนฺนํ อินฺทฺริยานํ จริยา ทฏฺฐพฺพา.}

\prose{จตูสุ ฌาเนสุ สมาธินฺทฺริยสฺส วเสน กตเมหิ วีสติยา อากาเรหิ ปญฺจนฺนํ อินฺทฺริยานํ จริยา ทฏฺฐพฺพา? ปฐเม ฌาเน อวิกฺเขปาธิปเตยฺยฏฺเฐน สมาธินฺทฺริยสฺส จริยา ทฏฺฐพฺพา, สมาธินฺทฺริยสฺส วเสน ทสฺสนฏฺเฐน ปญฺญินฺทฺริยสฺส จริยา ทฏฺฐพฺพา, อธิโมกฺขฏฺเฐน สทฺธินฺทฺริยสฺส จริยา ทฏฺฐพฺพา, ปคฺคหฏฺเฐน วีริยินฺทฺริยสฺส จริยา ทฏฺฐพฺพา, อุปฏฺฐานฏฺเฐน}

\csromanheader{2}{4. อินฺทฺริยกถา}
\prose{\csromanfolio{215}สตินฺทฺริยสฺส จริยา ทฏฺฐพฺพา.\csromanspacer{} ทุติเย ฌาเน ฯเปฯ.\csromanspacer{} ตติเย ฌาเน ฯเปฯ.\csromanspacer{} จตุตฺเถ ฌาเน อวิกฺเขปาธิปเตยฺยฏฺเฐน สมาธินฺทฺริยสฺส จริยา ทฏฺฐพฺพา ฯเปฯ.\csromanspacer{} จตูสุ ฌาเนสุ สมาธินฺทฺริยสฺส วเสน อิเมหิ วีสติยา อากาเรหิ ปญฺจนฺนํ อินฺทฺริยานํ จริยา ทฏฺฐพฺพา.}

\prose{จตูสุ อริยสจฺเจสุ ปญฺญินฺทฺริยสฺส วเสน กตเมหิ วีสติยา อากาเรหิ ปญฺจนฺนํ อินฺทฺริยานํ จริยา ทฏฺฐพฺพา? ทุกฺเข อริยสจฺเจ ทสฺสนาธิปเตยฺยฏฺเฐน ปญฺญินฺทฺริยสฺส จริยา ทฏฺฐพฺพา, ปญฺญินฺทฺริยสฺส วเสน อธิโมกฺขฏฺเฐน สทฺธินฺทฺริยสฺส จริยา ทฏฺฐพฺพา, ปคฺคหฏฺเฐน วีริยินฺทฺริยสฺส จริยา ทฏฺฐพฺพา, อุปฏฺฐานฏฺเฐน สตินฺทฺริยสฺส จริยา ทฏฺฐพฺพา, อวิกฺเขปฏฺเฐน สมาธินฺทฺริยสฺส จริยา ทฏฺฐพฺพา.\csromanspacer{} ทุกฺขสมุทเย อริยสจฺเจ ฯเปฯ.\csromanspacer{} ทุกฺขนิโรเธ อริยสจฺเจ.\csromanspacer{} ทุกฺขนิโรธคามินิยา ปฏิปทาย อริยสจฺเจ ทสฺสนาธิปเตยฺยฏฺเฐน ปญฺญินฺทฺริยสฺส จริยา ทฏฺฐพฺพา, ปญฺญินฺทฺริยสฺส วเสน อธิโมกฺขฏฺเฐน สทฺธินฺทฺริยสฺส จริยา ทฏฺฐพฺพา, ปคฺคหฏฺเฐน วีริยินฺทฺริยสฺส จริยา ทฏฺฐพฺพา, อุปฏฺฐานฏฺเฐน สตินฺทฺริยสฺส จริยา ทฏฺฐพฺพา, อวิกฺเขปฏฺเฐน สมาธินฺทฺริยสฺส จริยา ทฏฺฐพฺพา.\csromanspacer{} จตูสุ อริยสจฺเจสุ ปญฺญินฺทฺริยสฺส วเสน อิเมหิ วีสติยา อากาเรหิ ปญฺจนฺนํ อินฺทฺริยานํ จริยา ทฏฺฐพฺพา.}

\csromanheader{2}{\textbf{จารวิหารนิทฺเทส}}
\proseitem{197}{จาโร จ วิหาโร จ อนุพุทฺโธ โหติ ปฏิวิทฺโธ.\csromanspacer{} ยถาจรนฺตํ ยถาวิหรนฺตํ วิญฺญู สพฺรหฺมจารี คมฺภีเรสุ ฐาเนสุ โอกปฺเปยฺยุํ อทฺธา อยมายสฺมา ปตฺโต วา ปาปุณิสฺสติ วา.}

\prose{\textbf{จริยา}ติ อฏฺฐ จริยาโย อิริยาปถจริยา, อายตนจริยา, สติจริยา, สมาธิจริยา, ญาณจริยา, มคฺคจริยา, ปตฺติจริยา, โลกตฺถจริยาติ.\csromanspacer{} \textbf{อิริยาปถจริยา}ติ จตูสุ อิริยาปเถสุ.\csromanspacer{} \textbf{อายตนจริยา}ติ ฉสุ อชฺฌตฺติกพาหิเรสุ อายตเนสุ.\csromanspacer{} \textbf{สติจริยา}ติ จตูสุ สติปฏฺฐาเนสุ.\csromanspacer{} \textbf{สมาธิจริยา}ติ จตูสุ ฌาเนสุ.\csromanspacer{} \textbf{ญาณจริยา}ติ จตูสุ อริยสจฺเจสุ.\csromanspacer{} \textbf{มคฺคจริยา}ติ จตูสุ อริยมคฺเคสุ.\csromanspacer{} \textbf{ปตฺติจริยา}ติ จตูสุ สามญฺญผเลสุ.\csromanspacer{} \textbf{โลกตฺถจริยา}ติ ตถาคเตสุ อรหนฺเตสุ สมฺมาสมฺพุทฺเธสุ, ปเทเส ปจฺเจกพุทฺเธสุ, ปเทเส สาวเกสุ.\csromanspacer{} \textbf{อิริยาปถจริยา} จ ปณิธิสมฺปนฺนานํ.\csromanspacer{} \textbf{อายตนจริยา} จ \csromanfolio{216}อินฺทฺริเยสุ คุตฺตทฺวารานํ.\csromanspacer{} สติจริยา จ อปฺปมาทวิหารีนํ.\csromanspacer{} สมาธิจริยา จ อธิจิตฺตมนุยุตฺตานํ.\csromanspacer{} ญาณจริยา จ พุทฺธิสมฺปนฺนานํ.\csromanspacer{} มคฺคจริยา จ สมฺมาปฏิปนฺนานํ.\csromanspacer{} ปตฺติจริยา จ อธิคตผลานํ.\csromanspacer{} โลกตฺถจริยา จ ตถาคตานํ อรหนฺตานํ สมฺมาสมฺพุทฺธานํ, ปเทเส ปจฺเจกพุทฺธานํ, ปเทเส สาวกานํ.\csromanspacer{} อิมา อฏฺฐ จริยาโย.}

\prose{อปราปิ อฏฺฐ จริยาโย อธิมุจฺจนฺโต สทฺธาย จรติ, ปคฺคณฺหนฺโต วีริเยน จรติ, อุปฏฺฐาเปนฺโต สติยา จรติ, อวิกฺเขปํ กโรนฺโต สมาธินา จรติ, ปชานนฺโต ปญฺญาย จรติ, วิชานนฺโต วิญฺญาณจริยาย จรติ, “เอวํ ปฏิปนฺนสฺส กุสลา ธมฺมา อายาเปนฺตี”ติ อายตนจริยาย จรติ, “เอวํ ปฏิปนฺโน วิเสสมธิคจฺฉตี”ติ วิเสสจริยาย จรติ.\csromanspacer{} อิมา อฏฺฐ จริยาโย.}

\prose{อปราปิ อฏฺฐ จริยาโย ทสฺสนจริยา จ สมฺมาทิฏฺฐิยา, อภินิโรปนจริยา จ สมฺมาสงฺกปฺปสฺส, ปริคฺคหจริยา จ สมฺมาวาจาย, สมุฏฺฐานจริยา จ สมฺมากมฺมนฺตสฺส, โวทานจริยา จ สมฺมา-อาชีวสฺส, ปคฺคหจริยา จ สมฺมาวายามสฺส, อุปฏฺฐานจริยา จ สมฺมาสติยา, อวิกฺเขปจริยา จ สมฺมาสมาธิสฺส.\csromanspacer{} อิมา อฏฺฐ จริยาโย.}

\prose{\textbf{วิหาโร}ติ อธิมุจฺจนฺโต สทฺธาย วิหรติ, ปคฺคณฺหนฺโต วีริเยน วิหรติ, อุปฏฺฐาเปนฺโต สติยา วิหรติ, อวิกฺเขปํ กโรนฺโต สมาธินา วิหรติ, ปชานนฺโต ปญฺญาย วิหรติ.}

\prose{\textbf{อนุพุทฺโธ}ติ สทฺธินฺทฺริยสฺส อธิโมกฺขฏฺโฐ อนุพุทฺโธ โหติ, วีริยินฺทฺริยสฺส ปคฺคหฏฺโฐ อนุพุทฺโธ โหติ, สตินฺทฺริยสฺส อุปฏฺฐานฏฺโฐ อนุพุทฺโธ โหติ, สมาธินฺทฺริยสฺส อวิกฺเขปฏฺโฐ อนุพุทฺโธ โหติ, ปญฺญินฺทฺริยสฺส ทสฺสนฏฺโฐ อนุพุทฺโธ โหติ.}

\prose{\textbf{ปฏิวิทฺโธ}ติ สทฺธินฺทฺริยสฺส อธิโมกฺขฏฺโฐ ปฏิวิทฺโธ โหติ, วีริยินฺทฺริยสฺส ปคฺคหฏฺโฐ ปฏิวิทฺโธ โหติ, สตินฺทฺริยสฺส อุปฏฺฐานฏฺโฐ ปฏิวิทฺโธ โหติ, สมาธินฺทฺริยสฺส อวิกฺเขปฏฺโฐ ปฏิวิทฺโธ โหติ, ปญฺญินฺทฺริยสฺส ทสฺสนฏฺโฐ ปฏิวิทฺโธ โหติ.\csromanspacer{}\csromanspacer{}\textbf{ยถาจรนฺตนฺ}ติ เอวํ สทฺธาย จรนฺตํ เอวํ วีริเยน จรนฺตํ เอวํ สติยา จรนฺตํ เอวํ สมาธินา จรนฺตํ เอวํ ปญฺญาย จรนฺตํ.\csromanspacer{}\csromanspacer{}\textbf{ยถาวิหรนฺต}นฺติ เอวํ สทฺธาย วิหรนฺตํ เอวํ วีริเยน วิหรนฺตํ เอวํ สติยา วิหรนฺตํ เอวํ สมาธินา วิหรนฺตํ เอวํ ปญฺญาย วิหรนฺตํ.\csromanspacer{} \textbf{วิญฺญู}ติ}

\csromanheader{2}{4. อินฺทฺริยกถา}
\prose{\csromanfolio{217}วิญฺญู วิภาวี เมธาวี ปณฺฑิตา พุทฺธิสมฺปนฺนา.\csromanspacer{} \textbf{สพฺรหฺมจารี}ติ เอกกมฺมํ เอกุทฺเทโส สมสิกฺขตา.\csromanspacer{} \textbf{คมฺภีเรสุ ฐาเนสู}ติ คมฺภีรานิ ฐานานิ วุจฺจนฺติ ฌานา จ วิโมกฺขา จ สมาธี จ สมาปตฺติโย จ มคฺคา จ ผลานิ จ อภิญฺญาโย จ ปฏิสมฺภิทา จ.\csromanspacer{} \textbf{โอกปฺเปยฺยุนฺ}ติ สทฺทเหยฺยุํ อธิมุจฺเจยฺยุํ.\csromanspacer{} \textbf{อทฺธา}ติ เอกํสวจนเมตํ นิสฺสํสยวจนเมตํ นิกฺกงฺขวจนเมตํ อเทฺวชฺฌวจนเมตํ อเทฺวฬฺหกวจนเมตํ นิโยควจนเมตํ อปณฺณกวจนเมตํ อวตฺถาปนวจนเมตํ “อทฺธา”ติ.\csromanspacer{} \textbf{อายสฺมา}ติ ปิยวจนเมตํ ครุวจนเมตํ สคารวสปฺปติสฺสาธิวจนเมตํ “อายสฺมา”ติ.\csromanspacer{} \textbf{ปตฺโต วา}ติ อธิคโต วา\textbf{.}\csromanspacer{}\textbf{ ปาปุณิสฺสติ วา}ติ อธิคมิสฺสติ วา.}

\csromanheader{2}{สุตฺตนฺตนิทฺเทโส ตติโย.}
\csromansectionrule
\csromanmatikaanchor{mtk.76}
\csromantocmarkref{subsection}{4. จตุตฺถสุตฺตนฺตนิทฺเทส}{mtk.76}
\bookmark[dest={mtk.76},level=2]{4. จตุตฺถสุตฺตนฺตนิทฺเทส}
\csromanheader{2}{4. จตุตฺถสุตฺตนฺตนิทฺเทส}
\proseitem{198}{ปุริมนิทานํ\footnote{ปริปุณฺณนิทานํ (สฺยา, ก)}.\csromanspacer{}\csromanspacer{}ปญฺจิมานิ ภิกฺขเว อินฺทฺริยานิ.\csromanspacer{} กตมานิ ปญฺจ? สทฺธินฺทฺริยํ, วีริยินฺทฺริยํ, สตินฺทฺริยํ, สมาธินฺทฺริยํ, ปญฺญินฺทฺริยํ.\csromanspacer{} อิมานิ โข ภิกฺขเว ปญฺจินฺทฺริยานิ.\csromanspacer{} อิมานิ ปญฺจินฺทฺริยานิ กติหากาเรหิ เกนฏฺเฐน ทฏฺฐพฺพานิ? อิมานิ ปญฺจินฺทฺริยานิ ฉหากาเรหิ เตนฏฺเฐน ทฏฺฐพฺพานิ อาธิปเตยฺยฏฺเฐน, อาทิวิโสธนฏฺเฐน, อธิมตฺตฏฺเฐน, อธิฏฺฐานฏฺเฐน, ปริยาทานฏฺเฐน, ปติฏฺฐาปกฏฺเฐน.}

\csromanheader{2}{\textbf{ก. อาธิปเตยฺยฏฺฐนิทฺเทส}}
\proseitem{199}{กถํ \textbf{อาธิปเตยฺยฏฺเฐน อินฺทฺริยานิ ทฏฺฐพฺพานิ?} อสฺสทฺธิยํ ปชหโต อธิโมกฺขาธิปเตยฺยฏฺเฐน สทฺธินฺทฺริยํ ทฏฺฐพฺพํ, สทฺธินฺทฺริยสฺส วเสน ปคฺคหฏฺเฐน วีริยินฺทฺริยํ ทฏฺฐพฺพํ, อุปฏฺฐานฏฺเฐน สตินฺทฺริยํ ทฏฺฐพฺพํ, อวิกฺเขปฏฺเฐน สมาธินฺทฺริยํ ทฏฺฐพฺพํ, ทสฺสนฏฺเฐน ปญฺญินฺทฺริยํ ทฏฺฐพฺพํ.\csromanspacer{} โกสชฺชํ ปชหโต ปคฺคหาธิปเตยฺยฏฺเฐน วีริยินฺทฺริยํ ทฏฺฐพฺพํ, วีริยินฺทฺริยสฺส วเสน อุปฏฺฐานฏฺเฐน สตินฺทฺริยํ ทฏฺฐพฺพํ, อวิกฺเขปฏฺเฐน สมาธินฺทฺริยํ ทฏฺฐพฺพํ, ทสฺสนฏฺเฐน ปญฺญินฺทฺริยํ ทฏฺฐพฺพํ, อธิโมกฺขฏฺเฐน สทฺธินฺทฺริยํ ทฏฺฐพฺพํ.\csromanspacer{} ปมาทํ ปชหโต อุปฏฺฐานาธิปเตยฺยฏฺเฐน สตินฺทฺริยํ ทฏฺฐพฺพํ, สตินฺทฺริยสฺส วเสน อวิกฺเขปฏฺเฐน สมาธินฺทฺริยํ ทฏฺฐพฺพํ, ทสฺสนฏฺเฐน ปญฺญินฺทฺริยํ ทฏฺฐพฺพํ, อธิโมกฺขฏฺเฐน สทฺธินฺทฺริยํ ทฏฺฐพฺพํ.\csromanspacer{} ปคฺคหฏฺเฐน \csromanfolio{218}วีริยินฺทฺริยํ ทฏฺฐพฺพํ.\csromanspacer{} อุทฺธจฺจํ ปชหโต อวิกฺเขปาธิปเตยฺยฏฺเฐน สมาธินฺทฺริยํ ทฏฺฐพฺพํ, สมาธินฺทฺริยสฺส วเสน ทสฺสนฏฺเฐน ปญฺญินฺทฺริยํ ทฏฺฐพฺพํ อธิโมกฺขฏฺเฐน สทฺธินฺทฺริยํ ทฏฺฐพฺพํ, ปคฺคหฏฺเฐน วีริยินฺทฺริยํ ทฏฺฐพฺพํ.\csromanspacer{} อุปฏฺฐานฏฺเฐน สตินฺทฺริยํ ทฏฺฐพฺพํ.\csromanspacer{} อวิชฺชํ ปชหโต ทสฺสนาธิปเตยฺยฏฺเฐน ปญฺญินฺทฺริยํ ทฏฺฐพฺพํ.\csromanspacer{} ปญฺญินฺทฺริยสฺส วเสน อธิโมกฺขฏฺเฐน สทฺธินฺทฺริยํ ทฏฺฐพฺพํ.\csromanspacer{} ปคฺคหฏฺเฐน วีริยินฺทฺริยํ ทฏฺฐพฺพํ.\csromanspacer{} อุปฏฺฐานฏฺเฐน สตินฺทฺริยํ ทฏฺฐพฺพํ.\csromanspacer{} อวิกฺเขปฏฺเฐน สมาธินฺทฺริยํ ทฏฺฐพฺพํ.\csromanspacer{} ทสฺสนฏฺเฐน ปญฺญินฺทฺริยํ ทฏฺฐพฺพํ.}

\prose{กามจฺฉนฺทํ ปชหโต เนกฺขมฺมวเสน อธิโมกฺขาธิปเตยฺยฏฺเฐน สทฺธินฺทฺริยํ ทฏฺฐพฺพํ, สทฺธินฺทฺริยสฺส วเสน ปคฺคหฏฺเฐน วีริยินฺทฺริยํ ทฏฺฐพฺพํ.\csromanspacer{} อุปฏฺฐานฏฺเฐน สตินฺทฺริยํ ทฏฺฐพฺพํ, อวิกฺเขปฏฺเฐน สมาธินฺทฺริยํ ทฏฺฐพฺพํ, ทสฺสนฏฺเฐน ปญฺญินฺทฺริยํ ทฏฺฐพฺพํ.\csromanspacer{} กามจฺฉนฺทํ ปชหโต เนกฺขมฺมวเสน ปคฺคหาธิปเตยฺยฏฺเฐน วีริยินฺทฺริยํ ทฏฺฐพฺพํ, วีริยินฺทฺริยสฺส วเสน อุปฏฺฐานฏฺเฐน สตินฺทฺริยํ ทฏฺฐพฺพํ, อวิกฺเขปฏฺเฐน สมาธินฺทฺริยํ ทฏฺฐพฺพํ, ทสฺสนฏฺเฐน ปญฺญินฺทฺริยํ ทฏฺฐพฺพํ, อธิโมกฺขฏฺเฐน สทฺธินฺทฺริยํ ทฏฺฐพฺพํ.\csromanspacer{} กามจฺฉนฺทํ ปชหโต เนกฺขมฺมวเสน อุปฏฺฐานาธิปเตยฺยฏฺเฐน สตินฺทฺริยํ ทฏฺฐพฺพํ, สตินฺทฺริยสฺส วเสน อวิกฺเขปฏฺเฐน สมาธินฺทฺริยํ ทฏฺฐพฺพํ, ทสฺสนฏฺเฐน ปญฺญินฺทฺริยํ ทฏฺฐพฺพํ, อธิโมกฺขฏฺเฐน สทฺธินฺทฺริยํ ทฏฺฐพฺพํ, ปคฺคหฏฺเฐน วีริยินฺทฺริยํ ทฏฺฐพฺพํ.\csromanspacer{} กามจฺฉนฺทํ ปชหโต เนกฺขมฺมวเสน อวิกฺเขปาธิปเตยฺยฏฺเฐน สมาธินฺทฺริยํ ทฏฺฐพฺพํ, สมาธินฺทฺริยสฺส วเสน ทสฺสนฏฺเฐน ปญฺญินฺทฺริยํ ทฏฺฐพฺพํ, อธิโมกฺขฏฺเฐน สทฺธินฺทฺริยํ ทฏฺฐพฺพํ, ปคฺคหฏฺเฐน วีริยินฺทฺริยํ ทฏฺฐพฺพํ, อุปฏฺฐานฏฺเฐน สตินฺทฺริยํ ทฏฺฐพฺพํ.\csromanspacer{} กามจฺฉนฺทํ ปชหโต เนกฺขมฺมวเสน ทสฺสนาธิปเตยฺยฏฺเฐน ปญฺญินฺทฺริยํ ทฏฺฐพฺพํ, ปญฺญินฺทฺริยสฺส วเสน อธิโมกฺขฏฺเฐน สทฺธินฺทฺริยํ ทฏฺฐพฺพํ, ปคฺคหฏฺเฐน วีริยินฺทฺริยํ ทฏฺฐพฺพํ, อุปฏฺฐานฏฺเฐน สตินฺทฺริยํ ทฏฺฐพฺพํ, อวิกฺเขปฏฺเฐน สมาธินฺทฺริยํ ทฏฺฐพฺพํ.}

\prose{พฺยาปาทํ ปชหโต อพฺยาปาทวเสน ฯเปฯ.\csromanspacer{} ถนมิทฺธํ ปชหโต อาโลกสญฺญาวเสน ฯเปฯ.\csromanspacer{} สพฺพกิเลเส ปชหโต อรหตฺตมคฺควเสน อธิโมกฺขาธิปเตยฺยฏฺเฐน สทฺธินฺทฺริยํ ทฏฺฐพฺพํ, สทฺธินฺทฺริยสฺส วเสน ปคฺคหฏฺเฐน วีริยินฺทฺริยํ ทฏฺฐพฺพํ, อุปฏฺฐานฏฺเฐน สตินฺทฺริยํ ทฏฺฐพฺพํ, อวิกฺเขปฏฺเฐน สมาธินฺทฺริยํ ทฏฺฐพฺพํ, ทสฺสนฏฺเฐน ปญฺญินฺทฺริยํ ทฏฺฐพฺพํ ฯเปฯ สพฺพกิเลเส ปชหโต อรหตฺตมคฺควเสน ทสฺสนาธิปเตยฺยฏฺเฐน ปญฺญินฺทฺริยํ ทฏฺฐพฺพํ, ปญฺญินฺทฺริยสฺส วเสน อธิโมกฺขฏฺเฐน สทฺธินฺทฺริยํ ทฏฺฐพฺพํ, ปคฺคหฏฺเฐน วีริยินฺทฺริยํ ทฏฺฐพฺพํ,}

\csromanheader{2}{4. อินฺทฺริยกถา}
\prose{\csromanfolio{219}อุปฏฺฐานฏฺเฐน สตินฺทฺริยํ ทฏฺฐพฺพํ, อวิกฺเขปฏฺเฐน สมาธินฺทฺริยํ ทฏฺฐพฺพํ.\csromanspacer{} เอวํ อาธิปเตยฺยฏฺเฐน อินฺทฺริยานิ ทฏฺฐพฺพานิ. (1)}

\csromanheader{2}{\textbf{ข. อาทิวิโสธนฏฺฐนิทฺเทส}}
\proseitem{200}{กถํ \textbf{อาทิวิโสธนฏฺเฐน อินฺทฺริยานิ ทฏฺฐพฺพานิ?} อธิโมกฺขฏฺเฐน สทฺธินฺทฺริยํ อสฺสทฺธิยสํวรฏฺเฐน\footnote{อสฺสทฺธิยํ สํวรฏฺเฐน (สฺยา, ก) เอวมีทิเสสุ ปเทสุปิ.} สีลวิสุทฺธิ สทฺธินฺทฺริยสฺส อาทิวิโสธนา.\csromanspacer{} ปคฺคหฏฺเฐน วีริยินฺทฺริยํ โกสชฺชสํวรฏฺเฐน สีลวิสุทฺธิ วีริยินฺทฺริยสฺส อาทิวิโสธนา.\csromanspacer{} อุปฏฺฐานฏฺเฐน สตินฺทฺริยํ ปมาทสํวรฏฺเฐน สีลวิสุทฺธิ สตินฺทฺริยสฺส อาทิวิโสธนา.\csromanspacer{} อวิกฺเขปฏฺเฐน สมาธินฺทฺริยํ อุทฺธจฺจสํวรฏฺเฐน สีลวิสุทฺธิ สมาธินฺทฺริยสฺส อาทิวิโสธนา.\csromanspacer{} ทสฺสนฏฺเฐน ปญฺญินฺทฺริยํ อวิชฺชาสํวรฏฺเฐน สีลวิสุทฺธิ ปญฺญินฺทฺริยสฺส อาทิวิโสธนา.\csromanspacer{} เนกฺขมฺเม ปญฺจินฺทฺริยานิ กามจฺฉนฺทสํวรฏฺเฐน สีลวิสุทฺธิ ปญฺจนฺนํ อินฺทฺริยานํ อาทิวิโสธนา.\csromanspacer{} อพฺยาปาเท ปญฺจินฺทฺริยานิ พฺยาปาทสํวรฏฺเฐน สีลวิสุทฺธิ ปญฺจนฺนํ อินฺทฺริยานํ อาทิวิโสธนา ฯเปฯ.\csromanspacer{} อรหตฺตมคฺเค ปญฺจินฺทฺริยานิ สพฺพกิเลสสํวรฏฺเฐน สีลวิสุทฺธิ ปญฺจนฺนํ อินฺทฺริยานํ อาทิวิโสธนา.\csromanspacer{} เอวํ อาทิวิโสธนฏฺเฐน อินฺทฺริยานิ ทฏฺฐพฺพานิ. (2)}

\csromanheader{2}{\textbf{ค. อธิมตฺตฏฺฐนิทฺเทส}}
\proseitem{201}{กถํ \textbf{อธิมตฺตฏฺเฐน อินฺทฺริยานิ ทฏฺฐพฺพานิ?} สทฺธินฺทฺริยสฺส ภาวนาย ฉนฺโท อุปฺปชฺชติ, ฉนฺทวเสน สทฺธาวเสน สทฺธินฺทฺริยํ อธิมตฺตํ โหติ.\csromanspacer{} ฉนฺทวเสน ปาโมชฺชํ อุปฺปชฺชติ, ปาโมชฺชวเสน สทฺธาวเสน สทฺธินฺทฺริยํ อธิมตฺตํ โหติ.\csromanspacer{} ปาโมชฺชวเสน ปีติ อุปฺปชฺชติ, ปีติวเสน สทฺธาวเสน สทฺธินฺทฺริยํ อธิมตฺตํ โหติ.\csromanspacer{} ปีติวเสน ปสฺสทฺธิ อุปฺปชฺชติ, ปสฺสทฺธิวเสน สทฺธาวเสน สทฺธินฺทฺริยํ อธิมตฺตํ โหติ.\csromanspacer{} ปสฺสทฺธิวเสน สุขํ อุปฺปชฺชติ, สุขวเสน สทฺธาวเสน สทฺธินฺทฺริยํ อธิมตฺตํ โหติ.\csromanspacer{} สุขวเสน โอภาโส อุปฺปชฺชติ, โอภาสวเสน สทฺธาวเสน สทฺธินฺทฺริยํ อธิมตฺตํ โหติ.\csromanspacer{} โอภาสวเสน สํเวโค อุปฺปชฺชติ, สํเวควเสน สทฺธาวเสน สทฺธินฺทฺริยํ อธิมตฺตํ โหติ.\csromanspacer{} สํเวเชตฺวา จิตฺตํ สมาทหติ, สมาธิวเสน สทฺธาวเสน สทฺธินฺทฺริยํ อธิมตฺตํ โหติ.\csromanspacer{} ตถาสมาหิตํ จิตฺตํ สาธุกํ ปคฺคณฺหาติ, ปคฺคหวเสน สทฺธาวเสน สทฺธินฺทฺริยํ อธิมตฺตํ โหติ.\csromanspacer{} ตถาปคฺคหิตํ จิตฺตํ สาธุกํ อชฺฌุเปกฺขติ, อุเปกฺขาวเสน สทฺธาวเสน สทฺธินฺทฺริยํ อธิมตฺตํ โหติ.\csromanspacer{} อุเปกฺขาวเสน \csromanfolio{220}นานตฺตกิเลเสหิ จิตฺตํ วิมุจฺจติ, วิโมกฺขวเสน สทฺธาวเสน สทฺธินฺทฺริยํ อธิมตฺตํ โหติ.\csromanspacer{} วิมุตฺตตฺตา เต ธมฺมา เอกรสา โหนฺติ, เอกรสฏฺเฐน ภาวนาวเสน สทฺธาวเสน สทฺธินฺทฺริยํ อธิมตฺตํ โหติ.\csromanspacer{} ภาวิตตฺตา ตโต ปณีตตเร วิวฏฺฏนฺติ, วิวฏฺฏนาวเสน สทฺธาวเสน สทฺธินฺทฺริยํ อธิมตฺตํ โหติ.\csromanspacer{} วิวฏฺฏิตตฺตา ตโต โวสชฺชติ\footnote{โวสฺสชฺชติ (สฺยา, ก)}, โวสคฺควเสน สทฺธาวเสน สทฺธินฺทฺริยํ อธิมตฺตํ โหติ.\csromanspacer{} โวสชฺชิตตฺตา ตโต นิรุชฺฌนฺติ, นิโรธวเสน สทฺธาวเสน สทฺธินฺทฺริยํ อธิมตฺตํ โหติ.\csromanspacer{} นิโรธวเสน เทฺว โวสคฺคา ปริจฺจาคโวสคฺโค จ ปกฺขนฺทนโวสคฺโค จ.\csromanspacer{} กิเลเส จ ขนฺเธ จ ปริจฺจชตีติ ปริจฺจาคโวสคฺโค.\csromanspacer{} นิโรธนิพฺพานธาตุยา จิตฺตํ ปกฺขนฺทตีติ ปกฺขนฺทนโวสคฺโค.\csromanspacer{} นิโรธวเสน อิเม เทฺว โวสคฺคา.}

\prose{อสฺสทฺธิยสฺส ปหานาย ฉนฺโท อุปฺปชฺชติ ฯเปฯ.\csromanspacer{} อสฺสทฺธิยปริฬาหสฺส ปหานาย ฉนฺโท อุปฺปชฺชติ.\csromanspacer{} ทิฏฺเฐกฏฺฐานํ กิเลสานํ ปหานาย ฉนฺโท อุปฺปชฺชติ.\csromanspacer{} โอฬาริกานํ กิเลสานํ ปหานาย ฉนฺโท อุปฺปชฺชติ.\csromanspacer{} อณุสหคตานํ กิเลสานํ ปหานาย ฉนฺโท อุปฺปชฺชติ.\csromanspacer{} สพฺพกิเลสานํ ปหานาย ฉนฺโท อุปฺปชฺชติ, ฉนฺทวเสน สทฺธาวเสน สทฺธินฺทฺริยํ อธิมตฺตํ โหติ ฯเปฯ.\csromanspacer{} วีริยินฺทฺริยสฺส ภาวนาย ฉนฺโท อุปฺปชฺชติ ฯเปฯ.\csromanspacer{} โกสชฺชสฺส ปหานาย ฉนฺโท อุปฺปชฺชติ.\csromanspacer{} โกสชฺชปริฬาหสฺส ปหานาย ฉนฺโท อุปฺปชฺชติ.\csromanspacer{} ทิฏฺเฐกฏฺฐานํ กิเลสานํ ปหานาย ฉนฺโท อุปฺปชฺชติ ฯเปฯ.\csromanspacer{} สพฺพกิเลสานํ ปหานาย ฉนฺโท อุปฺปชฺชติ.\csromanspacer{} สตินฺทฺริยสฺส ภาวนาย ฉนฺโท อุปฺปชฺชติ ฯเปฯ.\csromanspacer{} ปมาทสฺส ปหานาย ฉนฺโท อุปฺปชฺชติ.\csromanspacer{} ปมาทปริฬาหสฺส ปหานาย ฉนฺโท อุปฺปชฺชติ ฯเปฯ.\csromanspacer{} สพฺพกิเลสานํ ปหานาย ฉนฺโท อุปฺปชฺชติ.\csromanspacer{} สมาธินฺทฺริยสฺส ภาวนาย ฉนฺโท อุปฺปชฺชติ ฯเปฯ.\csromanspacer{} อุทฺธจฺจสฺส ปหานาย ฉนฺโท อุปฺปชฺชติ.\csromanspacer{} อุทฺธจฺจปริฬาหสฺส ปหานาย ฉนฺโท อุปฺปชฺชติ ฯเปฯ.\csromanspacer{} สพฺพกิเลสานํ ปหานาย ฉนฺโท อุปฺปชฺชติ.\csromanspacer{} ปญฺญินฺทฺริยสฺส ภาวนาย ฉนฺโท อุปฺปชฺชติ ฯเปฯ.\csromanspacer{} อวิชฺชาย ปหานาย ฉนฺโท อุปฺปชฺชติ.\csromanspacer{} อวิชฺชาปริฬาหสฺส ปหานาย ฉนฺโท อุปฺปชฺชติ.\csromanspacer{} ทิฏฺเฐกฏฺฐานํ กิเลสานํ ปหานาย ฉนฺโท อุปฺปชฺชติ.\csromanspacer{} โอฬาริกานํ กิเลสานํ ปหานาย ฉนฺโท อุปฺปชฺชติ.\csromanspacer{} อณุสหคตานํ กิเลสานํ ปหานาย ฉนฺโท อุปฺปชฺชติ.\csromanspacer{} สพฺพกิเลสานํ ปหานาย ฉนฺโท อุปฺปชฺชติ.\csromanspacer{} ฉนฺทวเสน ปญฺญาวเสน ปญฺญินฺทฺริยํ อธิมตฺตํ โหติ.\csromanspacer{} ฉนฺทวเสน ปาโมชฺชํ อุปฺปชฺชติ, ปาโมชฺชวเสน}

\csromanheader{2}{4. อินฺทฺริยกถา}
\prose{\csromanfolio{221}ปญฺญาวเสน ปญฺญินฺทฺริยํ อธิมตฺตํ โหติ.\csromanspacer{} ปาโมชฺชวเสน ปีติ อุปฺปชฺชติ, ปีติวเสน ปญฺญาวเสน ปญฺญินฺทฺริยํ อธิมตฺตํ โหติ.\csromanspacer{} ปีติวเสน ปสฺสทฺธิ อุปฺปชฺชติ, ปสฺสทฺธิวเสน ปญฺญาวเสน ปญฺญินฺทฺริยํ อธิมตฺตํ โหติ.\csromanspacer{} ปสฺสทฺธิวเสน สุขํ อุปฺปชฺชติ, สุขวเสน ปญฺญาวเสน ปญฺญินฺทฺริยํ อธิมตฺตํ โหติ.\csromanspacer{} สุขวเสน โอภาโส อุปฺปชฺชติ, โอภาสวเสน ปญฺญาวเสน ปญฺญินฺทฺริยํ อธิมตฺตํ โหติ.\csromanspacer{} โอภาสวเสน สํเวโค อุปฺปชฺชติ, สํเวควเสน ปญฺญาวเสน ปญฺญินฺทฺริยํ อธิมตฺตํ โหติ.\csromanspacer{} สํเวเชตฺวา จิตฺตํ สมาทหติ, สมาธิวเสน ปญฺญาวเสน ปญฺญินฺทฺริยํ อธิมตฺตํ โหติ.\csromanspacer{} ตถาสมาหิตํ จิตฺตํ สาธุกํ ปคฺคณฺหาติ, ปคฺคหวเสน ปญฺญาวเสน ปญฺญินฺทฺริยํ อธิมตฺตํ โหติ.\csromanspacer{} ตถาปคฺคหิตํ จิตฺตํ สาธุกํ อชฺฌุเปกฺขติ, อุเปกฺขาวเสน ปญฺญาวเสน ปญฺญินฺทฺริยํ อธิมตฺตํ โหติ.\csromanspacer{} อุเปกฺขาวเสน นานตฺตกิเลเสหิ จิตฺตํ วิมุจฺจติ, วิโมกฺขวเสน ปญฺญาวเสน ปญฺญินฺทฺริยํ อธิมตฺตํ โหติ.\csromanspacer{} วิมุตฺตตฺตา เต ธมฺมา เอกรสา โหนฺติ, ภาวนาวเสน\footnote{เอกรสฏฺเฐน ภาวนาวเสน (สฺยา, ก) อฏฺฐกถา โอโลเกตพฺโพ.} ปญฺญาวเสน ปญฺญินฺทฺริยํ อธิมตฺตํ โหติ.\csromanspacer{} ภาวิตตฺตา ตโต ปณีตตเร วิวฏฺฏนฺติ, วิวฏฺฏนาวเสน ปญฺญาวเสน ปญฺญินฺทฺริยํ อธิมตฺตํ โหติ.\csromanspacer{} วิวฏฺฏิตตฺตา ตโต โวสชฺชติ, โวสคฺควเสน ปญฺญาวเสน ปญฺญินฺทฺริยํ อธิมตฺตํ โหติ.\csromanspacer{} โวสชฺชิตตฺตา ตโต นิรุชฺฌนฺติ, นิโรธวเสน ปญฺญาวเสน ปญฺญินฺทฺริยํ อธิมตฺตํ โหติ.\csromanspacer{} นิโรธวเสน เทฺว โวสคฺคา ปริจฺจาคโวสคฺโค จ ปกฺขนฺทนโวสคฺโค จ.\csromanspacer{} กิเลเส จ ขนฺเธ จ ปริจฺจชตีติ ปริจฺจาคโวสคฺโค.\csromanspacer{} นิโรธนิพฺพานธาตุยา จิตฺตํ ปกฺขนฺทตีติ ปกฺขนฺทนโวสคฺโค.\csromanspacer{} นิโรธวเสน อิเม เทฺว โวสคฺคา.\csromanspacer{} เอวํ อธิมตฺตฏฺเฐน อินฺทฺริยานิ ทฏฺฐพฺพานิ.}

\csromanheader{2}{ทุติยภาณวาโร.}
\csromansectionrule
\csromanheader{2}{\textbf{ฆ. อธิฏฺฐานฏฺฐนิทฺเทส}}
\proseitem{202}{กถํ \textbf{อธิฏฺฐานฏฺเฐน อินฺทฺริยานิ ทฏฺฐพฺพานิ?} สทฺธินฺทฺริยสฺส ภาวนาย ฉนฺโท อุปฺปชฺชติ, ฉนฺทวเสน สทฺธาวเสน สทฺธินฺทฺริยํ อธิฏฺฐาติ.\csromanspacer{} ฉนฺทวเสน ปาโมชฺชํ อุปฺปชฺชติ, ปาโมชฺชวเสน สทฺธาวเสน สทฺธินฺทฺริยํ อธิฏฺฐาติ ฯเปฯ.\csromanspacer{} เอวํ อธิฏฺฐานฏฺเฐน อินฺทฺริยานิ ทฏฺฐพฺพานิ.}

\vspace{\tipitakaparskip}
\csromancenter{\textbf{\^{}อ.}\csromanspacer{}\textbf{ ปริยาทานฏฺฐนิทฺเทส}}
\prose{\csromanfolio{222}กถํ \textbf{ปริยาทานฏฺเฐน อินฺทฺริยานิ ทฏฺฐพฺพานิ?} อธิโมกฺขฏฺเฐน สทฺธินฺทฺริยํ อสฺสทฺธิยํ ปริยาทิยติ, อสฺสทฺธิยปริฬาหํ ปริยาทิยติ.\csromanspacer{} ปคฺคหฏฺเฐน วีริยินฺทฺริยํ โกสชฺชํ ปริยาทิยติ, โกสชฺชปริฬาหํ ปริยาทิยติ.\csromanspacer{} อุปฏฺฐานฏฺเฐน สตินฺทฺริยํ ปมาทํ ปริยาทิยติ, ปมาทปริฬาหํ ปริยาทิยติ.\csromanspacer{} อวิกฺเขปฏฺเฐน สมาธินฺทฺริยํ อุทฺธจฺจํ ปริยาทิยติ, อุทฺธจฺจปริฬาหํ ปริยาทิยติ.\csromanspacer{} ทสฺสนฏฺเฐน ปญฺญินฺทฺริยํ อวิชฺชํ ปริยาทิยติ, อวิชฺชาปริฬาหํ ปริยาทิยติ.\csromanspacer{} เนกฺขมฺเม ปญฺจินฺทฺริยานิ กามจฺฉนฺทํ ปริยาทิยนฺติ.\csromanspacer{} อพฺยาปาเท ปญฺจินฺทฺริยานิ พฺยาปาทํ ปริยาทิยนฺติ.\csromanspacer{} อาโลกสญฺญาย ปญฺจินฺทฺริยานิ ถินมิทฺธํ ปริยาทิยนฺติ.\csromanspacer{} อวิกฺเขเป ปญฺจินฺทฺริยานิ อุทฺธจฺจํ ปริยาทิยนฺติ ฯเปฯ.\csromanspacer{} อรหตฺตมคฺเค ปญฺจินฺทฺริยานิ สพฺพกิเลเส ปริยาทิยนฺติ.\csromanspacer{} เอวํ ปริยาทานฏฺเฐน อินฺทฺริยานิ ทฏฺฐพฺพานิ.}

\csromanheader{2}{\textbf{จ. ปติฏฺฐาปกฏฺฐนิทฺเทส}}
\proseitem{203}{กถํ \textbf{ปติฏฺฐาปกฏฺเฐน อินฺทฺริยานิ ทฏฺฐพฺพานิ?} สทฺโธ สทฺธินฺทฺริยํ อธิโมกฺเข ปติฏฺฐาเปติ, สทฺธสฺส สทฺธินฺทฺริยํ อธิโมกฺเข ปติฏฺฐาเปติ.\csromanspacer{} วีริยวา วีริยินฺทฺริยํ ปคฺคเห ปติฏฺฐาเปติ, วีริยวโต วีริยินฺทฺริยํ ปคฺคเห ปติฏฺฐาเปติ.\csromanspacer{} สติมา สตินฺทฺริยํ อุปฏฺฐาเน ปติฏฺฐาเปติ, สติมโต สตินฺทฺริยํ อุปฏฺฐาเน ปติฏฺฐาเปติ.\csromanspacer{} สมาหิโต สมาธินฺทฺริยํ อวิกฺเขเป ปติฏฺฐาเปติ, สมาหิตสฺส สมาธินฺทฺริยํ อวิกฺเขเป ปติฏฺฐาเปติ.\csromanspacer{} ปญฺญวา ปญฺญินฺทฺริยํ ทสฺสเน ปติฏฺฐาเปติ, ปญฺญวโต ปญฺญินฺทฺริยํ ทสฺสเน ปติฏฺฐาเปติ.\csromanspacer{} โยคาวจโร ปญฺจินฺทฺริยานิ เนกฺขมฺเม ปติฏฺฐาเปติ, โยคาวจรสฺส ปญฺจินฺทฺริยานิ เนกฺขมฺเม ปติฏฺฐาเปนฺติ.\csromanspacer{} โยคาวจโร ปญฺจินฺทฺริยานิ อพฺยาปาเท ปติฏฺฐาเปติ, โยคาวจรสฺส ปญฺจินฺทฺริยานิ อพฺยาปาเท ปติฏฺฐาเปนฺติ.\csromanspacer{} โยคาวจโร ปญฺจินฺทฺริยานิ อาโลกสญฺญาย ปติฏฺฐาเปติ, โยคาวจรสฺส ปญฺจินฺทฺริยานิ อาโลกสญฺญาย ปติฏฺฐาเปนฺติ.\csromanspacer{} โยคาวจโร ปญฺจินฺทฺริยานิ อวิกฺเขเป ปติฏฺฐาเปติ, โยคาวจรสฺส ปญฺจินฺทฺริยานิ อวิกฺเขเป ปติฏฺฐาเปนฺติ ฯเปฯ.\csromanspacer{} โยคาวจโร ปญฺจินฺทฺริยานิ อรหตฺตมคฺเค ปติฏฺฐาเปติ, โยคาวจรสฺส ปญฺจินฺทฺริยานิ อรหตฺตมคฺเค ปติฏฺฐาเปนฺติ.\csromanspacer{} เอวํ ปติฏฺฐาปกฏฺเฐน อินฺทฺริยานิ ทฏฺฐพฺพานิ.\csromanspacer{} สุตฺตนฺตนิทฺเทโส จตุตฺโถ.}

\csromanheader{2}{4. อินฺทฺริยกถา}
\csromanmatikaanchor{mtk.77}
\csromantocmarkref{subsection}{5. อินฺทฺริยสโมธาน}{mtk.77}
\bookmark[dest={mtk.77},level=2]{5. อินฺทฺริยสโมธาน}
\csromanheader{2}{5. อินฺทฺริยสโมธาน}
\proseitem{204}{\csromanfolio{223}ปุถุชฺชโน สมาธิํ ภาเวนฺโต กติหากาเรหิ อุปฏฺฐานกุสโล โหติ, เสกฺโข สมาธิํ ภาเวนฺโต กติหากาเรหิ อุปฏฺฐานกุสโล โหติ, วีตราโค สมาธิํ ภาเวนฺโต กติหากาเรหิ อุปฏฺฐานกุสโล โหติ? ปุถุชฺชโน สมาธิํ ภาเวนฺโต สตฺตหิ อากาเรหิ อุปฏฺฐานกุสโล โหติ, เสกฺโข สมาธิํ ภาเวนฺโต อฏฺฐหิ อากาเรหิ อุปฏฺฐานกุสโล โหติ, วีตราโค สมาธิํ ภาเวนฺโต ทสหิ อากาเรหิ อุปฏฺฐานกุสโล โหติ.}

\prose{ปุถุชฺชโน สมาธิํ ภาเวนฺโต กตเมหิ สตฺตหิ อากาเรหิ อุปฏฺฐานกุสโล โหติ\csromandash{}อาวชฺชิตตฺตา อารมฺมณูปฏฺฐานกุสโล โหติ, สมถนิมิตฺตูปฏฺฐานกุสโล โหติ, ปคฺคหนิมิตฺตูปฏฺฐานกุสโล โหติ, อวิกฺเขปูปฏฺฐานกุสโล โหติ, โอภาสูปฏฺฐานกุสโล โหติ, สมฺปหํสนูปฏฺฐานกุสโล โหติ, อุเปกฺขูปฏฺฐานกุสโล โหติ.\csromanspacer{} ปุถุชฺชโน สมาธิํ ภาเวนฺโต อิเมหิ สตฺตหิ อากาเรหิ อุปฏฺฐานกุสโล โหติ.}

\prose{เสกฺโข สมาธิํ ภาเวนฺโต กตเมหิ อฏฺฐหิ อากาเรหิ อุปฏฺฐานกุสโล โหติ\csromandash{}อาวชฺชิตตฺตา อารมฺมณูปฏฺฐานกุสโล โหติ, สมถนิมิตฺตูปฏฺฐานกุสโล โหติ, ปคฺคหนิมิตฺตูปฏฺฐานกุสโล โหติ, อวิกฺเขปูปฏฺฐานกุสโล โหติ, โอภาสูปฏฺฐานกุสโล โหติ, สมฺปหํสนูปฏฺฐานกุสโล โหติ, อุเปกฺขูปฏฺฐานกุสโล โหติ, เอกตฺตูปฏฺฐานกุสโล โหติ.\csromanspacer{} เสกฺโข สมาธิํ ภาเวนฺโต อิเมหิ อฏฺฐหิ อากาเรหิ อุปฏฺฐานกุสโล โหติ.}

\prose{วีตราโค สมาธิํ ภาเวนฺโต กตเมหิ ทสหิ อากาเรหิ อุปฏฺฐานกุสโล โหติ\csromandash{}อาวชฺชิตตฺตา อารมฺมณูปฏฺฐานกุสโล โหติ ฯเปฯ เอกตฺตูปฏฺฐานกุสโล โหติ, ญาณูปฏฺฐานกุสโล โหติ, วิมุตฺตูปฏฺฐานกุสโล โหติ.\csromanspacer{} วีตราโค สมาธิํ ภาเวนฺโต อิเมหิ ทสหิ อากาเรหิ อุปฏฺฐานกุสโล โหติ.}

\proseitem{205}{ปุถุชฺชโน วิปสฺสนํ ภาเวนฺโต กติหากาเรหิ อุปฏฺฐานกุสโล โหติ, กติหากาเรหิ อนุปฏฺฐานกุสโล โหติ.\csromanspacer{} เสกฺโข วิปสฺสนํ ภาเวนฺโต กติหากาเรหิ อุปฏฺฐานกุสโล \csromanfolio{224}โหติ, กติหาเรหิ อนุปฏฺฐานกุสโล โหติ.\csromanspacer{} วีตราโค วิปสฺสนํ ภาเวนฺโต กติหากาเรหิ อุปฏฺฐานกุสโล โหติ, กติหากาเรหิ อนุปฏฺฐานกุสโล โหติ?}

\prose{ปุถุชฺชโน วิปสฺสนํ ภาเวนฺโต นวหิ อากาเรหิ อุปฏฺฐานกุสโล โหติ, นวหิ อากาเรหิ อนุปฏฺฐานกุสโล โหติ.\csromanspacer{} เสกฺโข วิปสฺสนํ ภาเวนฺโต ทสหิ อากาเรหิ อุปฏฺฐานกุสโล โหติ, ทสหิ อากาเรหิ อนุปฏฺฐานกุสโล โหติ.\csromanspacer{} วีตราโค วิปสฺสนํ ภาเวนฺโต ทฺวาทสหิ อากาเรหิ อุปฏฺฐานกุสโล โหติ, ทฺวาทสหิ อากาเรหิ อนุปฏฺฐานกุสโล โหติ.}

\prose{ปุถุชฺชโน วิปสฺสนํ ภาเวนฺโต กตเมหิ นวหิ อากาเรหิ อุปฏฺฐานกุสโล โหติ, กตเมหิ นวหิ อากาเรหิ อนุปฏฺฐานกุสโล โหติ? อนิจฺจโต อุปฏฺฐานกุสโล โหติ, นิจฺจโต อนุปฏฺฐานกุสโล โหติ.\csromanspacer{} ทุกฺขโต อุปฏฺฐานกุสโล โหติ, สุขโต อนุปฏฺฐานกุสโล โหติ.\csromanspacer{} อนตฺตโต อุปฏฺฐานกุสโล โหติ, อตฺตโต อนุปฏฺฐานกุสโล โหติ.\csromanspacer{} ขยโต อุปฏฺฐานกุสโล โหติ, ฆนโต อนุปฏฺฐานกุสโล โหติ.\csromanspacer{} วยโต อุปฏฺฐานกุสโล โหติ, อายูหนานุปฏฺฐานกุสโล โหติ.\csromanspacer{} วิปริณามูปฏฺฐานกุสโล โหติ, ธุวโต อนุปฏฺฐานกุสโล โหติ.\csromanspacer{} อนิมิตฺตูปฏฺฐานกุสโล โหติ, นิมิตฺตานุปฏฺฐานกุสโล โหติ.\csromanspacer{} อปฺปณิหิตูปฏฺฐานกุสโล โหติ, ปณิธิ-อนุปฏฺฐานกุสโล\footnote{ปณิธานุปฏฺฐานกุสโล (สฺยา)} โหติ, สุญฺญตูปฏฺฐานกุสโล โหติ, อภินิเวสานุปฏฺฐานกุสโล โหติ.\csromanspacer{} ปุถุชฺชโน วิปสฺสนํ ภาเวนฺโต อิเมหิ นวหิ อากาเรหิ อุปฏฺฐานกุสโล โหติ, อิเมหิ นวหิ อากาเรหิ อนุปฏฺฐานกุสโล โหติ.}

\proseitem{206}{เสกฺโข วิปสฺสนํ ภาเวนฺโต กตเมหิ ทสหิ อากาเรหิ อุปฏฺฐานกุสโล โหติ, กตเมหิ ทสหิ อากาเรหิ อนุปฏฺฐานกุสโล โหติ? อนิจฺจโต อุปฏฺฐานกุสโล โหติ, นิจฺจโต อนุปฏฺฐานกุสโล โหติ ฯเปฯ.\csromanspacer{} สุญฺญตูปฏฺฐานกุสโล โหติ, อภินิเวสานุปฏฺฐานกุสโล โหติ.\csromanspacer{} ญาณูปฏฺฐานกุสโล โหติ อญาณ-อนุปฏฺฐานกุสโล โหติ, เสกฺโข วิปสฺสนํ ภาเวนฺโต}

\csromanheader{2}{4. อินฺทฺริยกถา}
\prose{\csromanfolio{225}อิเมหิ ทสหิ อากาเรหิ อุปฏฺฐานกุสโล โหติ, อิเมหิ ทสหิ อากาเรหิ อนุปฏฺฐานกุสโล โหติ.}

\prose{วีตราโค วิปสฺสนํ ภาเวนฺโต กตเมหิ ทฺวาทสหิ อากาเรหิ อุปฏฺฐานกุสโล โหติ, กตเมหิ ทฺวาทสหิ อากาเรหิ อนุปฏฺฐานกุสโล โหติ? อนิจฺจโต อุปฏฺฐานกุสโล โหติ, นิจฺจโต อนุปฏฺฐานกุสโล โหติ ฯเปฯ.\csromanspacer{} ญาณูปฏฺฐานกุสโล โหติ, อญาณานุปฏฺฐานกุสโล โหติ.\csromanspacer{} วิสญฺโญคูปฏฺฐานกุสโล โหติ, สญฺโญคานุปฏฺฐานกุสโล โหติ.\csromanspacer{} นิโรธูปฏฺฐานกุสโล โหติ, สงฺขารานุปฏฺฐานกุสโล โหติ.\csromanspacer{} วีตราโค วิปสฺสนํ ภาเวนฺโต อิเมหิ ทฺวาทสหิ อากาเรหิ อุปฏฺฐานกุสโล โหติ, อิเมหิ ทฺวาทสหิ อากาเรหิ อนุปฏฺฐานกุสโล โหติ.}

\prose{อาวชฺชิตตฺตา อารมฺมณูปฏฺฐานกุสลวเสน อินฺทฺริยานิ สโมธาเนติ โคจรญฺจ ปชานาติ.\csromanspacer{} สมตฺถญฺจ ปฏิวิชฺฌติ ฯเปฯ ธมฺเม สโมธาเนติ.\csromanspacer{} โคจรญฺจ ปชานาติ สมตฺถญฺจ ปฏิวิชฺฌติ.\csromanspacer{} \textbf{อินฺทฺริยานิ สโมธาเนตี}ติ กถํ อินฺทฺริยานิ สโมธาเนติ.\csromanspacer{} อธิโมกฺขฏฺเฐน สทฺธินฺทฺริยํ สโมธาเนติ ฯเปฯ.\csromanspacer{} สมถนิมิตฺตูปฏฺฐานกุสลวเสน.\csromanspacer{} ปคฺคหนิมิตฺตูปฏฺฐานกุสลวเสน.\csromanspacer{} อวิกฺเขปูปฏฺฐานกุสลวเสน.\csromanspacer{} โอภาสูปฏฺฐานกุสลวเสน.\csromanspacer{} สมฺปหํสนูปฏฺฐานกุสลวเสน.\csromanspacer{} อุเปกฺขูปฏฺฐานกุสลวเสน.\csromanspacer{} เอกตฺตูปฏฺฐานกุสลวเสน.\csromanspacer{} ญาณูปฏฺฐานกุสลวเสน.\csromanspacer{} วิมุตฺตูปฏฺฐานกุสลวเสน.}

\prose{อนิจฺจโต อุปฏฺฐานกุสลวเสน.\csromanspacer{} นิจฺจโต อนุปฏฺฐานกุสลวเสน.\csromanspacer{} ทุกฺขโต อุปฏฺฐานกุสลวเสน.\csromanspacer{} สุขโต อนุปฏฺฐานกุสลวเสน.\csromanspacer{} อนตฺตโต อุปฏฺฐานกุสลวเสน.\csromanspacer{} อตฺตโต อนุปฏฺฐานกุสลวเสน, ขยโต อุปฏฺฐานกุสลวเสน.\csromanspacer{} ฆนโต อนุปฏฺฐานกุสลวเสน.\csromanspacer{} วยโต อุปฏฺฐานกุสลวเสน.\csromanspacer{} อายูหนานุปฏฺฐานกุสลวเสน.\csromanspacer{} วิปริณามูปฏฺฐานกุสลวเสน.\csromanspacer{} ธุวโต อนุปฏฺฐานกุสลวเสน.\csromanspacer{} อนิมิตฺตูปฏฺฐานกุสลวเสน.\csromanspacer{} นิมิตฺตานุปฏฺฐานกุสลวเสน.\csromanspacer{} อปฺปณิหิตูปฏฺฐานกุสลวเสน.\csromanspacer{} ปณิธิ-อนุปฏฺฐานกุสลวเสน.\csromanspacer{} สุญฺญตูปฏฺฐานกุสลวเสน.\csromanspacer{} อภินิเวสานุปฏฺฐานกุสลวเสน.\csromanspacer{} ญาณูปฏฺฐานกุสลวเสน.\csromanspacer{} อญาณานุปฏฺฐานกุสลวเสน.\csromanspacer{} วิสญฺโญ รูปฏฺฐานกุสลวเสน.}

\prose{\csromanfolio{226}สญฺโญคานุปฏฺฐานกุสลวเสน.\csromanspacer{} นิโรธูปฏฺฐานกุสลวเสน.\csromanspacer{} สงฺขารานุปฏฺฐานกุสลวเสน.\csromanspacer{} อินฺทฺริยานิ สโมธาเนติ.\csromanspacer{} โคจรญฺจ ปชานาติ.\csromanspacer{} สมตฺถญฺจ ปฏิวิชฺฌติ.}

\proseitem{207}{จตุสฏฺฐิยา อากาเรหิ ติณฺณนฺนํ อินฺทฺริยานํ วสิภาวตา ปญฺญา อาสวานํ ขเย ญาณํ\csromandash{}กตเมสํ ติณฺณนฺนํ อินฺทฺริยานํ? อนญฺญาตญฺญสฺสามีตินฺทฺริยสฺส อญฺญินฺทฺริยสฺส อญฺญาตาวินฺทฺริยสฺส.}

\prose{อนญฺญาตญฺญสฺสามีตินฺทฺริยํ กติ ฐานานิ คจฺฉติ, อญฺญินฺทฺริยํ กติ ฐานานิ คจฺฉติ, อญฺญาตาวินฺทฺริยํ กติ ฐานานิ คจฺฉติ? อนญฺญาตญฺญสฺสามีตินฺทฺริยํ เอกํ ฐานํ คจฺฉติ, โสตาปตฺติมคฺคํ.\csromanspacer{} อญฺญินฺทฺริยํ ฉ ฐานานิ คจฺฉติ, โสตาปตฺติผลํ สกทาคามิมคฺคํ สกทามิผลํ อนาคามิมคฺคํ อนาคามิผลํ อรหตฺตมคฺคํ.\csromanspacer{} อญฺญาตาวินฺทฺริยํ เอกํ ฐานํ คจฺฉติ, อรหตฺตผลํ.}

\prose{โสตาปตฺติมคฺคกฺขเณ อนญฺญาตญฺญสฺสามีตินฺทฺริยสฺส สทฺธินฺทฺริยํ อธิโมกฺขปริวารํ โหติ.\csromanspacer{} วีริยินฺทฺริยํ ปคฺคหปริวารํ โหติ.\csromanspacer{} สตินฺทฺริยํ อุปฏฺฐานปริวารํ โหติ.\csromanspacer{} สมาธินฺทฺริยํ อวิกฺเขปปริวารํ โหติ.\csromanspacer{} ปญฺญินฺทฺริยํ ทสฺสนปริวารํ โหติ.\csromanspacer{} มนินฺทฺริยํ วิชานนปริวารํ โหติ.\csromanspacer{} โสมนสฺสินฺทฺริยํ อภิสนฺทนปริวารํ โหติ.\csromanspacer{} ชีวิตินฺทฺริยํ ปวตฺตสนฺตตาธิปเตยฺยปริวารํ โหติ.\csromanspacer{} โสตาปตฺติมคฺคกฺขเณ ชาตา ธมฺมา ฐเปตฺวา จิตฺตสมุฏฺฐานํ รูปํ สพฺเพว กุสลา โหนฺติ.\csromanspacer{} สพฺเพว อนาสวา โหนฺติ.\csromanspacer{} สพฺเพว นิยฺยานิกา โหนฺติ สพฺเพว อปจยคามิโน โหนฺติ.\csromanspacer{} สพฺเพว โลกุตฺตรา โหนฺติ.\csromanspacer{} สพฺเพว นิพฺพานารมฺมณา โหนฺติ.\csromanspacer{} โสตาปตฺติมคฺคกฺขเณ อนญฺญาตญฺญสฺสามีตินฺทฺริยสฺส อิมานิ อฏฺฐินฺทฺริยานิ สหชาตปริวารา โหนฺติ, อญฺญมญฺญปริวารา โหนฺติ, นิสฺสยปริวารา โหนฺติ, สมฺปยุตฺตปริวารา โหนฺติ, สหคตา โหนฺติ, สหชาตา โหนฺติ, สํสฏฺฐา โหนฺติ, สมฺปยุตฺตา โหนฺติ, เตว ตสฺส อาการา เจว โหนฺติ ปริวารา จ.}

\prose{โสตาปตฺติผลกฺขเณ ฯเปฯ.\csromanspacer{} อรหตฺตผลกฺขเณ อญฺญาตาวินฺทฺริยสฺส สทฺธินฺทฺริยํ อธิโมกฺขปริวารํ โหติ ฯเปฯ ชีวิตินฺทฺริยํ ปวตฺตสนฺตตาธิปเตยฺยปริวารํ โหติ, อรหตฺตผลกฺขเณ ชาตา ธมฺมา สพฺเพว อพฺยากตา โหนฺติ, ฐเปตฺวา จิตฺตสมุฏฺฐานํ รูปํ สพฺเพว อนาสวา โหนฺติ สพฺเพว โลกุตฺตรา โหนฺติ สพฺเพว นิพฺพานารมฺมณา โหนฺติ.\csromanspacer{} อรหตฺตผลกฺขเณ อญฺญาตาวินฺทฺริยสฺส อิมานิ อฏฺฐินฺทฺริยานิ สหชาตปริวารา}

\csromanheader{2}{4. อินฺทฺริยกถา}
\prose{\csromanfolio{227}โหนฺติ, เตว ตสฺส อาการา เจว โหนฺติ ปริวารา จ.\csromanspacer{} อิติ อิมานิ อฏฺฐ อฏฺฐกานิ จตุสฏฺฐิ โหนฺติ.}

\prose{\textbf{อาสวา}ติ กตเม เต อาสวา? กามาสโว ภวาสโว ทิฏฺฐาสโว อวิชฺชาสโว.\csromanspacer{} กตฺเถเต อาสวา ขียนฺติ\footnote{ขิยฺยนฺติ (ก)}? โสตาปตฺติมคฺเคน อนวเสโส ทิฏฺฐาสโว ขียติ, อปายคมนีโย กามาสโว ขียติ, อปายคมนีโย ภวาสโว ขียติ, อปายคมนีโย อวิชฺชาสโว ขียติ, เอตฺเถเต อาสวา ขียนฺติ.\csromanspacer{} สกทาคามิมคฺเคน โอฬาริโก กามาสโว ขียติ, ตเทกฏฺโฐ ภวาสโว ขียติ, ตเทกฏฺโฐ อวิชฺชาสโว ขียติ, เอตฺเถเต อาสวา ขียนฺติ.\csromanspacer{} อนาคามิมคฺเคน อนวเสโส กามาสโว ขียติ, ตเทกฏฺโฐ ภวาสโว ขียติ, ตเทกฏฺโฐ อวิชฺชาสโว ขียติ, เอตฺเถเต อาสวา ขียนฺติ.\csromanspacer{} อรหตฺตมคฺเคน อนวเสโส ภวาสโว ขียติ, อนวเสโส อวิชฺชาสโว ขียติ, เอตฺเถเต อาสวา ขียนฺติ.}

\csromanhangingbegin
\hangingheaditem{208}{น ตสฺส อทฺทิฏฺฐมิธตฺถิกิญฺจิ,}
\hangingbody{อโถ อวิญฺญาตมชานิตพฺพํ.}
\hangingbody{สพฺพํ อภิญฺญาสิ ยทตฺถิ เนยฺยํ,}
\hangingbody{ตถาคโต เตน สมนฺตจกฺขูติ.}
\csromanhangingend

\prose{\textbf{สมนฺตจกฺขู}ติ เกนฏฺเฐน สมนฺตจกฺขุ? จุทฺทส พุทฺธญาณานิ ทุกฺเข ญาณํ พุทฺธญาณํ, ทุกฺขสมุทเย ญาณํ พุทฺธญาณํ ฯเปฯ สพฺพญฺญุตญฺญาณํ พุทฺธญาณํ, อนาวรณญาณํ พุทฺธญาณํ, อิมานิ จุทฺทส พุทฺธญาณานิ.\csromanspacer{} อิเมสํ จุทฺทสนฺนํ พุทฺธญาณานํ อฏฺฐ ญาณานิ สาวกสาธารณานิ ฉ ญาณานิ อสาธารณานิ สาวเกหิ.}

\prose{ยาวตา ทุกฺขสฺส ทุกฺขฏฺโฐ ญาโต, อญฺญาโต ทุกฺขฏฺโฐ นตฺถีติ สมนฺตจกฺขุ.\csromanspacer{} ยํ สมนฺตจกฺขุ ตํ ปญฺญินฺทฺริยํ, ปญฺญินฺทฺริยสฺส วเสน อธิโมกฺขฏฺเฐน สทฺธินฺทฺริยํ, วคฺคหฏฺเฐน วีริยินฺทฺริยํ, อุปฏฺฐานฏฺเฐน สตินฺทฺริยํ, อวิกฺเขปฏฺเฐน สมาธินฺทฺริยํ.\csromanspacer{} ยาวตา ทุกฺขสฺส ทุกฺขฏฺโฐ ทิฏฺโฐ วิทิโต สจฺฉิกโต ผสฺสิโต ปญฺญาย, อผสฺสิโต ปญฺญาย ทุกฺขฏฺโฐ นตฺถีติ สมนฺตจกฺขุ.\csromanspacer{} ยํ สมนฺตจกฺขุ, ตํ ปญฺญินฺทฺริยํ, ปญฺญินฺทฺริยสฺส วเสน อธิโมกฺขฏฺเฐน สทฺธินฺทฺริยํ, ปคฺคหฏฺเฐน วีริยินฺทฺริยํ, อุปฏฺฐานฏฺเฐน สตินฺทฺริยํ, อวิกฺเขปฏฺเฐน สมาธินฺทฺริยํ. \csromanfolio{228}ยาวตา สมุทยสฺส สมุทยฏฺโฐ ฯเปฯ.\csromanspacer{} ยาวตา นิโรธสฺส นิโรธฏฺโฐ.\csromanspacer{} ยาวตา มคฺคสฺส มคฺคฏฺโฐ.\csromanspacer{} ยาวตา อตฺถปฏิสมฺภิทาย อตฺถปฏิสมฺภิทฏฺโฐ.\csromanspacer{} ยาวตา ธมฺมปฏิสมฺภิทาย ธมฺมปฏิสมฺภิทฏฺโฐ.\csromanspacer{} ยาวตา นิรุตฺติปฏิสมฺภิทาย นิรุตฺติปฏิสมฺภิทฏฺโฐ.\csromanspacer{} ยาวตา ปฏิภานปฏิสมฺภิทาย ปฏิภานปฏิสมฺภิทฏฺโฐ.\csromanspacer{} ยาวตา อินฺทฺริยปโรปริยตฺตญาณํ.\csromanspacer{} ยาวตา สตฺตานํ อาสยานุสเย ญาณํ.\csromanspacer{} ยาวตา ยมกปาฏิหีเร ญาณํ.\csromanspacer{} ยาวตา มหากรุณาสมาปตฺติยา ญาณํ.\csromanspacer{} ยาวตา สเทวกสฺส โลกสฺส สมารกสฺส สพฺรหฺมกสฺส สสฺสมณพฺราหฺมณิยา ปชาย สเทวมนุสฺสาย ทิฏฺฐํ สุตํ มุตํ วิญฺญาตํ ปตฺตํ ปริเยสิตํ อนุวิจริตํ มนสา, ตํ ญาตํ ทิฏฺฐํ วิทิตํ สจฺฉิกตํ ผสฺสิตํ ปญฺญาย, อผสฺสิตํ ปญฺญาย นตฺถีติ สมนฺตจกฺขุ.\csromanspacer{} ยํ สมนฺตจกฺขุ, ตํ ปญฺญินฺทฺริยํ, ปญฺญินฺทฺริยสฺส วเสน อธิโมกฺขฏฺเฐน สทฺธินฺทฺริยํ, ปคฺคหฏฺเฐน วีริยินฺทฺริยํ, อุปฏฺฐานฏฺเฐน สตินฺทฺริยํ, อวิกฺเขปฏฺเฐน สมาธินฺทฺริยํ.}

\prose{สทฺทหนฺโต ปคฺคณฺหาติ, ปคฺคณฺหนฺโต สทฺทหติ, สทฺทหนฺโต อุปฏฺฐาเปติ, อุปฏฺฐาเปนฺโต สทฺทหติ, สทฺทหนฺโต สมาทหติ, สมาทหนฺโต สทฺทหติ, สทฺทหนฺโต ปชานาติ, ปชานนฺโต สทฺทหติ.\csromanspacer{} ปคฺคณฺหนฺโต อุปฏฺฐาเปติ, อุปฏฺฐาเปนฺโต ปคฺคณฺหาติ, ปคฺคณฺหนฺโต สมาทหติ, สมาทหนฺโต ปคฺคณฺหาติ, ปคฺคณฺหนฺโต ปชานาติ, ปชานนฺโต ปคฺคณฺหาติ, ปคฺคณฺหนฺโต สทฺทหติ, สทฺทหนฺโต ปคฺคณฺหาติ.\csromanspacer{} อุปฏฺฐาเปนฺโต สมาทหติ, สมาทหนฺโต อุปฏฺฐาเปติ, อุปฏฺฐาเปนฺโต ปชานาติ, ปชานนฺโต อุปฏฺฐาเปติ, อุปฏฺฐาเปนฺโต สทฺทหติ, สทฺทหนฺโต อุปฏฺฐาเปติ, อุปฏฺฐาเปนฺโต ปคฺคณฺหาติ, ปคฺคณฺหนฺโต อุปฏฺฐาเปติ.\csromanspacer{} สมาทหนฺโต ปชานาติ, ปชานนฺโต สมาทหติ, สมาทหนฺโต สทฺทหติ, สทฺทหนฺโต สมาทหติ, สมาทหนฺโต ปคฺคณฺหาติ, ปคฺคณฺหนฺโต สมาทหติ, สมาทหนฺโต อุปฏฺฐาเปติ, อุปฏฺฐาเปนฺโต สมาทหติ.\csromanspacer{} ปชานนฺโต สทฺทหติ, สทฺทหนฺโต ปชานาติ, ปชานนฺโต ปคฺคณฺหาติ, ปคฺคณฺหนฺโต ปชานาติ, ปชานนฺโต อุปฏฺฐาเปติ, อุปฏฺฐาเปนฺโต ปชานาติ, ปชานนฺโต สมาทหติ, สมาทหนฺโต ปชานาติ.}

\prose{สทฺทหิตตฺตา ปคฺคหิตํ, ปคฺคหิตตฺตา สทฺทหิตํ, สทฺทหิตตฺตา อุปฏฺฐาปิตํ, อุปฏฺฐาปิตตฺตา สทฺทหิตํ, สทฺทหิตตฺตา สมาทหิตํ, สมาทหิตตฺตา สทฺทหิตํ สทฺทหิตตฺตา ปชานิตํ, ปชานิตตฺตา สทฺทหิตํ.\csromanspacer{} ปคฺคหิตตฺตา}

\csromanheader{2}{4. อินฺทฺริยกถา}
\prose{\csromanfolio{229}อุปฏฺฐาปิตํ, อุปฏฺฐาปิตตฺตา ปคฺคหิตํ, ปคฺคติตตฺตา สมาทหิตํ, สมาทหิตตฺตา ปคฺคหิตํ, ปคฺคหิตตฺตา ปชานิตํ, ปชานิตตฺตา ปคฺคหิตํ, ปคฺคหิตตฺตา สทฺทหิตํ, สทฺทหิตตฺตา ปคฺคหิตํ.\csromanspacer{} อุปฏฺฐาปิตตฺตา สมาทหิตํ, สมาทหิตตฺตา อุปฏฺฐาปิตํ, อุปฏฺฐาปิตตฺตา ปชานิตํ, ปชานิตตฺตา อุปฏฺฐาปิตํ, อุปฏฺฐาปิตตฺตา สทฺทหิตํ, สทฺทหิตตฺตา อุปฏฺฐาปิตํ, อุปฏฺฐาปิตตฺตา ปคฺคหิตํ, ปคฺคหิตตฺตา อุปฏฺฐาปิตํ.\csromanspacer{} สมาทหิตตฺตา ปชานิตํ, ปชานิตตฺตา สมาทหิตํ, สมาทหิตตฺตา สทฺทหิตํ, สทฺทหิตตฺตา สมาทหิตํ, สมาทหิตตฺตา ปคฺคหิตํ, ปคฺคหิตตฺตา สมาทหิตํ, สมาทหิตตฺตา อุปฏฺฐาปิตํ, อุปฏฺฐาปิตตฺตา สมาทหิตํ.\csromanspacer{} ปชานิตตฺตา สทฺทหิตํ, สทฺทหิตตฺตา ปชานิตํ, ปชานิตตฺตา ปคฺคหิตํ, ปคฺคหิตตฺตา ปชานิตํ, ปชานิตตฺตา อุปฏฺฐาปิตํ, อุปฏฺฐาปิตตฺตา ปชานิตํ, ปชานิตตฺตา สมาทหิตํ, สมาทหิตตฺตา ปชานิตํ.}

\prose{ยํ พุทฺธจกฺขุ, ตํ พุทฺธญาณํ.\csromanspacer{} ยํ พุทฺธญาณํ, ตํ พุทฺธจกฺขุ.\csromanspacer{} เยน จกฺขุนา ตถาคโต สตฺเต ปสฺสติ อปฺปรชกฺเข มหารชกฺเข ติกฺขินฺทฺริเย มุทินฺทฺริเย สฺวากาเร ทฺวากาเร สุวิญฺญาปเย ทุวิญฺญาปเย อปฺเปกจฺเจ ปรโลกวชฺชภยทสฺสาวิโน\footnote{ปรโลกวชฺชภยทสฺสาวิเน (สฺยา, ก)} อปฺเปกจฺเจ นปรโลกวชฺชภยทสฺสาวิโน.}

\prose{\textbf{อปฺปรชกฺเข มหารชกฺเข}ติ สทฺโธ ปุคฺคโล อปฺปรชกฺโข, อสฺสทฺโธ ปุคฺคโล มหารชกฺโข.\csromanspacer{} อารทฺธวีริโย ปุคฺคโล อปฺปรชกฺโข, กุสีโต ปุคฺคโล มหารชกฺโข.\csromanspacer{} อุปฏฺฐิตสฺสติ ปุคฺคโล อปฺปรชกฺโข, มุฏฺฐสฺสติ ปุคฺคโล มหารชกฺโข.\csromanspacer{} สมาหิโต ปุคฺคโล อปฺปรชกฺโข, อสมาหิโต ปุคฺคโล มหารชกฺโข.\csromanspacer{} ปญฺญวา ปุคฺคโล อปฺปรชกฺโข, ทุปฺปญฺโญ ปุคฺคโล มหารชกฺโข.}

\prose{\textbf{ติกฺขินฺทฺริเย มุทินฺทฺริเย}ติ สทฺโธ ปุคฺคโล ติกฺขินฺทฺริโย, อสฺสทฺโธ ปุคฺคโล มุทินฺทฺริโย ฯเปฯ.\csromanspacer{} ปญฺญวา ปุคฺคโล ติกฺขินฺทฺริโย, ทุปฺปญฺโญ ปุคฺคโล มุทินฺทฺริโย.\csromanspacer{} \textbf{สฺวากาเร ทฺวากาเร}ติ สทฺโธ ปุคฺคโล สฺวากาโร, อสฺสทฺโธ ปุคฺคโล ทฺวากาโร ฯเปฯ.\csromanspacer{} ปญฺญวา ปุคฺคโล สฺวากาโร, ทุปฺปญฺโญ ปุคฺคโล ทฺวากาโร.\csromanspacer{} \textbf{สุวิญฺญาปเย ทุวิญฺญาปเย}ติ สทฺโธ ปุคฺคโล สุวิญฺญาปโย, อสฺสทฺโธ ปุคฺคโล ทุวิญฺญาปโย ฯเปฯ.\csromanspacer{} ปญฺญวา ปุคฺคโล สุวิญฺญาปโย, ทุปฺปญฺโญ ปุคฺคโล ทุวิญฺญาปโย.}

\prose{\csromanfolio{230}\textbf{อปฺเปกจฺเจ ปรโลกวชฺชภยทสฺสาวิโน อปฺเปกจฺเจ นปรโลกวชฺชภยทสฺสาวิโน}ติ สทฺโธ ปุคฺคโล ปรโลกวชฺชภยทสฺสาวี, อสฺสทฺโธ ปุคฺคโล นปรโลกวชฺชภยทสฺสาวี.\csromanspacer{} อารทฺธวีริโย ปุคฺคโล ปรโลกวชฺชภยทสฺสาวี, กุสีโต ปุคฺคโล นปรโลกวชฺชภยทสฺสาวี ฯเปฯ.\csromanspacer{} ปญฺญวา ปุคฺคโล ปรโลกวชฺชภยทสฺสาวี, ทุปฺปญฺโญ ปุคฺคโล นปรโลกวชฺชภยทสฺสาวี.}

\prose{\textbf{โลโก}ติ ขนฺธโลโก ธาตุโลโก อายตนโลโก วิปตฺติภวโลโก วิปตฺติสมฺภวโลโก สมฺปตฺติภวโลโก สมฺปตฺติสมฺภวโลโก.}

\prose{เอโก โลโก, สพฺเพ สตฺตา อาหารฏฺฐิติกา.\csromanspacer{} เทฺว โลกา, นามญฺจ รูปญฺจ.\csromanspacer{} ตโย โลกา, ติสฺโส เวทนา.\csromanspacer{} จตฺตาโร โลกา, จตฺตาโร อาหารา.\csromanspacer{} ปญฺจ โลกา, ปญฺจุปาทานกฺขนฺธา.\csromanspacer{} ฉ โลกา, ฉ อชฺฌตฺติกานิ อายตนานิ.\csromanspacer{} สตฺต โลกา, สตฺต วิญฺญาณฏฺฐิติโย.\csromanspacer{} อฏฺฐ โลกา, อฏฺฐ โลกธมฺมา.\csromanspacer{} นว โลกา, นว สตฺตาวาสา.\csromanspacer{} ทส โลกา, ทสายตนานิ.\csromanspacer{} ทฺวาทส โลกา, ทฺวาทสายตนานิ.\csromanspacer{} อฏฺฐารส โลกา, อฏฺฐารส ธาตุโย.}

\prose{\textbf{วชฺชนฺ}ติ สพฺเพ กิเลสา วชฺชา, สพฺเพ ทุจฺจริตา วชฺชา, สพฺเพ อภิสงฺขารา วชฺชา, สพฺเพ ภวคามิกมฺมา วชฺชา.\csromanspacer{} อิติ อิมสฺมิญฺจ โลเก อิมสฺมิญฺจ วชฺเช ติพฺพา ภยสญฺญา ปจฺจุปฏฺฐิตา โหนฺติ เสยฺยถาปิ อุกฺขิตฺตาสิเก วธเก.\csromanspacer{} อิเมหิ ปญฺญาสาย อากาเรหิ อิมานิ ปญฺจินฺทฺริยานิ ชานาติ ปสฺสติ อญฺญาติ ปฏิวิชฺฌตีติ.}

\vspace{\tipitakaparskip}
\csromancenter{ตติยภาณวาโร.}
\nitthitamruled{อินฺทฺริยกถา นิฏฺฐิตา.}
\csromanmatikaanchor{mtk.78}
\csromantocmarkref{section}{5. วิโมกฺขกถา}{mtk.78}
\bookmark[dest={mtk.78},level=1]{5. วิโมกฺขกถา}
\csromanheader{1}{5. วิโมกฺขกถา}
\csromanheader{2}{\textbf{อุทฺเทส}}
\proseitem{209}{ปุริมนิทานํ, ตโยเม ภิกฺขเว วิโมกฺขา.\csromanspacer{} กตเม ตโย? สุญฺญโต วิโมกฺโข อนิมิตฺโต วิโมกฺโข อปฺปณิหิโต วิโมกฺโข, อิเม โข ภิกฺขเว ตโย วิโมกฺขา.}

\csromanheader{2}{5. วิโมกฺขกถา}
\prose{\csromanfolio{231}อปิ จ อฏฺฐสฏฺฐิ วิโมกฺขา\csromandash{}สุญฺญโต วิโมกฺโข, อนิมิตฺโต วิโมกฺโข, อปฺปณิหิโต วิโมกฺโข.\csromanspacer{} อชฺฌตฺตวุฏฺฐาโน วิโมกฺโข, พหิทฺธาวุฏฺฐาโน วิโมกฺโข, ทุภโต วุฏฺฐาโน วิโมกฺโข.\csromanspacer{} อชฺฌตฺตวุฏฺฐานา จตฺตาโร วิโมกฺขา, พหิทฺธาวุฏฺฐานา จตฺตาโร วิโมกฺขา, ทุภโต วุฏฺฐานา จตฺตาโร วิโมกฺขา.\csromanspacer{} อชฺฌตฺตวุฏฺฐานานํ อนุโลมา จตฺตาโร วิโมกฺขา, พหิทฺธาวุฏฺฐานานํ อนุโลมา จตฺตาโร วิโมกฺขา, ทุภโต วุฏฺฐานานํ อนุโลมา จตฺตาโร วิโมกฺขา.\csromanspacer{} อชฺฌตฺตวุฏฺฐานปฏิปฺปสฺสทฺธี\footnote{อชฺฌตฺตวุฏฺฐานา ปฏิปฺปสฺสทฺธี (สฺยา)} จตฺตาโร วิโมกฺขา, พหิทฺธาวุฏฺฐานปฏิปฺปสฺสทฺธี จตฺตาโร วิโมกฺขา, ทุภโต วุฏฺฐานปฏิปฺปสฺสทฺธี จตฺตาโร วิโมกฺขา.\csromanspacer{} รูปี รูปานิ ปสฺสตีติ วิโมกฺโข, อชฺฌตฺตํ อรูปสญฺญี พหิทฺธา รูปานิ ปสฺสตีติ วิโมกฺโข, “สุภนฺ”เตว อธิมุตฺโต โหตีติ วิโมกฺโข.\csromanspacer{} อากาสานญฺจายตนสมาปตฺติ วิโมกฺโข, วิญฺญาณญฺจายตนสมาปตฺติ วิโมกฺโข, อากิญฺจญฺญายตนสมาปตฺติ วิโมกฺโข.\csromanspacer{} เนวสญฺญานาสญฺญายตนสมาปตฺติ วิโมกฺโข, สญฺญาเวทยิตนิโรธสมาปตฺติ วิโมกฺโข.\csromanspacer{} สมยวิโมกฺโข อสมยวิโมกฺโข.\csromanspacer{} สามยิโก วิโมกฺโข, อสามยิโก วิโมกฺโข.\csromanspacer{} กุปฺโป วิโมกฺโข, อกุปฺโป วิโมกฺโข.\csromanspacer{} โลกิโก วิโมกฺโข, โลกุตฺตโร วิโมกฺโข.\csromanspacer{} สาสโว วิโมกฺโข, อนาสโว วิโมกฺโข.\csromanspacer{} สามิโส วิโมกฺโข, นิรามิโส วิโมกฺโข, นิรามิสา นิรามิสตโร วิโมกฺโข, ปณิหิโต วิโมกฺโข, อปฺปณิหิโต วิโมกฺโข, ปณิหิตปฺปฏิปฺปสฺสทฺธิ วิโมกฺโข.\csromanspacer{} สญฺญุตฺโต วิโมกฺโข, วิสญฺญุตฺโต วิโมกฺโข.\csromanspacer{} เอกตฺตวิโมกฺโข, นานตฺตวิโมกฺโข, สญฺญาวิโมกฺโข.\csromanspacer{} ญาณวิโมกฺโข, สีติสิยาวิโมกฺโข\footnote{สีตีสิยาวิโมกฺโข (สี-ฏฺฐ)}, ฌานวิโมกฺโข, อนุปาทา จิตฺตสฺส วิโมกฺโข. (75)}

\csromanheader{2}{\textbf{นิทฺเทส}}
\proseitem{210}{กตโม \textbf{สุญฺญโต วิโมกฺโข?} อิธ ภิกฺขุ อรญฺญคโต วา รุกฺขมูลคโต วา สุญฺญาคารคโต วา อิติ ปฏิสญฺจิกฺขติ “สุญฺญมิทํ อตฺเตน วา อตฺตนิเยน วา”ติ.\csromanspacer{} โส ตตฺถ อภินิเวสํ น กโรตีติ สุญฺญโต วิโมกฺโข.\csromanspacer{} อยํ สุญฺญโต วิโมกฺโข.}

\prose{\csromanfolio{232}กตโม \textbf{อนิมิตฺโต วิโมกฺโข?} อิธ ภิกฺขุ อรญฺญคโต วา รุกฺขมูลคโต วา สุญฺญาคารคโต วา อิติ ปฏิสญฺจิกฺขติ “สุญฺญมิทํ อตฺเตน วา อตฺตนิเยน วา”ติ, โส ตตฺถ นิมิตฺตํ น กโรตีติ อนิมิตฺโต วิโมกฺโข.\csromanspacer{} อยํ อนิมิตฺโต วิโมกฺโข.}

\prose{กตโม \textbf{อปฺปณิหิโต วิโมกฺโข?} อิธ ภิกฺขุ อรญฺญคโต วา รุกฺขมูลคโต วา สุญฺญาคารคโต วา อิติ ปฏิสญฺจิกฺขติ “สุญฺญมิทํ อตฺเตน วา อตฺตนิเยน วา”ติ, โส ตตฺถ ปณิธิํ น กโรตีติ อปฺปณิหิโต วิโมกฺโข.\csromanspacer{} อยํ อปฺปณิหิโต วิโมกฺโข.}

\prose{กตโม \textbf{อชฺฌตฺตวุฏฺฐาโน วิโมกฺโข?} จตฺตาริ ฌานานิ.\csromanspacer{} อยํ อชฺฌตฺตวุฏฺฐาโน วิโมกฺโข.\csromanspacer{} กตโม พหิทฺธาวุฏฺฐาโน วิโมกฺโข? จตสฺโส อรูปสมาปตฺติโย.\csromanspacer{} อยํ พหิทฺธาวุฏฺฐาโน วิโมกฺโข.\csromanspacer{} กตโม ทุภโต วุฏฺฐาโน วิโมกฺโข? จตฺตาโร อริยมคฺคา.\csromanspacer{} อยํ ทุภโต วุฏฺฐาโน วิโมกฺโข.}

\prose{กตเม \textbf{อชฺฌตฺตวุฏฺฐานา จตฺตาโร วิโมกฺขา?} ปฐมํ ฌานํ นีวรเณหิ วุฏฺฐาติ.\csromanspacer{} ทุติยํ ฌานํ วิตกฺกวิจาเรหิ วุฏฺฐาติ.\csromanspacer{} ตติยํ ฌานํ ปีติยา วุฏฺฐาติ.\csromanspacer{} จตุตฺถํ ฌานํ สุขทุกฺเขหิ วุฏฺฐาติ.\csromanspacer{} อิเม อชฺฌตฺตวุฏฺฐานา จตฺตาโร วิโมกฺขา. (4)}

\prose{กตเม \textbf{พหิทฺธาวุฏฺฐานา จตฺตาโร วิโมกฺขา?} อากาสานญฺจายตนสมาปตฺติ รูปสญฺญาย ปฏิฆสญฺญาย นานตฺตสญฺญาย วุฏฺฐาติ.\csromanspacer{} วิญฺญาณญฺจายตนสมาปตฺติ อากาสานญฺจายตนสญฺญาย วุฏฺฐาติ.\csromanspacer{} อากิญฺจญฺญายตนสมาปตฺติ วิญฺญาณญฺจายตนสญฺญาย วุฏฺฐาติ.\csromanspacer{} เนวสญฺญานาสญฺญายตนสมาปตฺติ อากิญฺจญฺญายตนสญฺญาย วุฏฺฐาติ.\csromanspacer{} อิเม พหิทฺธาวุฏฺฐานา จตฺตาโร วิโมกฺขา. (4, 8)}

\prose{กตเม \textbf{ทุภโต วุฏฺฐานา จตฺตาโร วิโมกฺขา?} โสตาปตฺติมคฺโค สกฺกายทิฏฺฐิวิจิกิจฺฉาสีลพฺพตปรามาสา ทิฏฺฐานุสยา วิจิกิจฺฉานุสยา วุฏฺฐาภิ, ตทนุวตฺตกกิเลเสหิ จ ขนฺเธหิ จ วุฏฺฐาติ, พหิทฺธา จ สพฺพนิมิตฺเตหิ วุฏฺฐาติ.\csromanspacer{} สกทาคามิมคฺโค โอฬาริกา กามราคสญฺโญชนา ปฏิฆสญฺโญชนา โอฬาริกา กามราคานุสยา ปฏิฆานุสยา วุฏฺฐาติ, ตทนุวตฺตกกิเลเสหิ จ ขนฺเธหิ จ วุฏฺฐาติ, พหิทฺธา จ สพฺพนิมิตฺเตหิ}

\csromanheader{2}{5. วิโมกฺขกถา}
\prose{\csromanfolio{233}วุฏฺฐาติ.\csromanspacer{} อนาคามิมคฺโค อณุสหคตา กามราคสญฺโญชนา ปฏิฆสญฺโญชนา อณุสหคตา กามราคานุสยา ปฏิฆานุสยา วุฏฺฐาติ, ตทนุวตฺตกกิเลเสหิ จ ขนฺเธหิ จ วุฏฺฐาติ, พหิทฺธา จ สพฺพนิมิตฺเตหิ วุฏฺฐาติ.\csromanspacer{} อรหตฺตมคฺโค รูปราคา อรูปราคา มานา อุทฺธจฺจา อวิชฺชาย มานานุสยา ภวราคานุสยา อวิชฺชานุสยา วุฏฺฐาติ, ตทนุวตฺตกกิเลเสหิ จ ขนฺเธหิ จ วุฏฺฐาติ, พหิทฺธา จ สพฺพนิมิตฺเตหิ วุฏฺฐาติ, อิเม ทุภโต วุฏฺฐานา จตฺตาโร วิโมกฺขา. (4, 12)}

\proseitem{211}{กตเม \textbf{อชฺฌตฺตวุฏฺฐานานํ อนุโลมา จตฺตาโร วิโมกฺขา?} ปฐมํ ฌานํ ปฏิลาภตฺถาย วิตกฺโก จ วิจาโร จ ปีติ จ สุขญฺจ จิตฺเตกคฺคตา จ.\csromanspacer{} ทุติยํ ฌานํ ปฏิลาภตฺถาย ฯเปฯ.\csromanspacer{} ตติยํ ฌานํ ปฏิลาภตฺถาย ฯเปฯ.\csromanspacer{} จตุตฺถํ ฌานํ ปฏิลาภตฺถาย วิตกฺโก จ วิจาโร จ ปีติ จ สุขญฺจ จิตฺเตกคฺคตา จ.\csromanspacer{} อิเม อชฺฌตฺตวุฏฺฐานานํ อนุโลมา จตฺตาโร วิโมกฺขา. (4, 16)}

\prose{กตเม \textbf{พหิทฺธาวุฏฺฐานานํ อนุโลมา จตฺตาโร วิโมกฺขา?} อากาสานญฺจายตนสมาปตฺติํ ปฏิลาภตฺถาย วิตกฺโก จ วิจาโร จ ปีติ จ สุขญฺจ จิตฺเตกคฺคตา จ.\csromanspacer{} วิญฺญาณญฺจายตนสมาปตฺติํ ฯเปฯ.\csromanspacer{} อากิญฺจญฺญายตนสมาปตฺติํ ฯเปฯ.\csromanspacer{} เนวสญฺญานาสญฺญายตนสมาปตฺติํ ปฏิลาภตฺถาย วิตกฺโก จ วิจาโร จ ปีติ จ สุขญฺจ จิตฺเตกคฺคตา จ.\csromanspacer{} อิเม พหิทฺธาวุฏฺฐานานํ อนุโลมา จตฺตาโร วิโมกฺขา. (4, 20)}

\prose{กตเม \textbf{ทุภโต วุฏฺฐานานํ อนุโลมา จตฺตาโร วิโมกฺขา?} โสตาปตฺติมคฺคํ ปฏิลาภตฺถาย อนิจฺจานุปสฺสนา, ทุกฺขานุปสฺสนา, อนตฺตานุปสฺสนา.\csromanspacer{} สกทาคามิมคฺคํ ปฏิลาภตฺถาย ฯเปฯ.\csromanspacer{} อนาคามิมคฺคํ ปฏิลาภตฺถาย ฯเปฯ.\csromanspacer{} อรหตฺตมคฺคํ ปฏิลาภตฺถาย อนิจฺจานุปสฺสนา, ทุกฺขานุปสฺสนา, อนตฺตานุปสฺสนา, อิเม ทุภโต วุฏฺฐานานํ อนุโลมา จตฺตาโร วิโมกฺขา. (4, 24)}

\prose{กตเม \textbf{อชฺฌตฺตวุฏฺฐานปฏิปฺปสฺสทฺธิ จตฺตาโร วิโมกฺขา?} ปฐมสฺส ฌานสฺส ปฏิลาโภ วา วิปาโก วา.\csromanspacer{} ทุติยสฺส ฌานสฺส ฯเปฯ.\csromanspacer{} ตติยสฺส ฌานสฺส ฯเปฯ.\csromanspacer{} จตุตฺถสฺส ฌานสฺส ปฏิลาโภ วา วิปาโก วา.\csromanspacer{} อิเม อชฺฌตฺตวุฏฺฐานปฏิปฺปสฺสทฺธี จตฺตาโร วิโมกฺขา. (4, 28)}

\prose{\csromanfolio{234}กตเม \textbf{พหิทฺธาวุฏฺฐานปฏิปฺปสฺสทฺธี จตฺตาโร วิโมกฺขา?} อากาสานญฺจายตนสมาปตฺติยา ปฏิลาโภ วา วิปาโก วา.\csromanspacer{} วิญฺญาณญฺจายตนสมาปตฺติยา ฯเปฯ.\csromanspacer{} อากิญฺจญฺญายตนสมาปตฺติยา ฯเปฯ.\csromanspacer{} เนวสญฺญานาสญฺญายตนสมาปตฺติยา ปฏิลาโภ วา วิปาโก วา.\csromanspacer{} อิเม พหิทฺธาวุฏฺฐานปฏิปฺปสฺสทฺธี จตฺตาโร วิโมกฺขา. (4,32)}

\prose{กตเม \textbf{ทุภโต วุฏฺฐานปฏิปฺปสฺสทฺธิ จตฺตาโร วิโมกฺขา?} โสตาปตฺติมคฺคสฺส โสตาปตฺติผลํ, สกทาคามิมคฺคสฺส สกทาคามิผลํ, อนาคามิมคฺคสฺส อนาคามิผลํ, อรหตฺตมคฺคสฺส อรหตฺตผลํ.\csromanspacer{} อิเม ทุภโต วุฏฺฐานปฏิปฺปสฺสทฺธี จตฺตาโร วิโมกฺขา. (4, 36)}

\proseitem{212}{กถํ \textbf{รูปี รูปานิ ปสฺสตีติ วิโมกฺโข?} อิเธกจฺโจ อชฺฌตฺตํ ปจฺจตฺตํ นีลนิมิตฺตํ มนสิ กโรติ, นีลสญฺญํ ปฏิลภติ, โส ตํ นิมิตฺตํ สุคฺคหิตํ กโรติ, สูปธาริตํ\footnote{สุปธาริตํ (ก)} อุปธาเรติ, สฺวาวตฺถิตํ อวตฺถาเปติ.\csromanspacer{} โส ตํ นิมิตฺตํ สุคฺคหิตํ กตฺวา สูปธาริตํ อุปธาเรตฺวา สฺวาวตฺถิตํ อวตฺถาเปตฺวา พหิทฺธา นีลนิมิตฺเต จิตฺตํ อุปสํหรติ, นีลสญฺญํ ปฏิลภติ.\csromanspacer{} โส ตํ นิมิตฺตํ สุคฺคหิตํ กโรติ, สูปธาริตํ อุปธาเรติ, สฺวาวตฺถิตํ อวตฺถาเปติ.\csromanspacer{} โส ตํ นิมิตฺตํ สุคฺคหิตํ กตฺวา อุปธาเรตฺวา สฺวาวตฺถิตํ อวตฺถาเปตฺวา อาเสวติ ภาเวติ พหุลีกโรติ.\csromanspacer{} ตสฺส เอวํ โหติ “อชฺฌตฺตญฺจ พหิทฺธา จ อุภยมิทํ รูปนฺ”ติ รูปสญฺญี โหติ.\csromanspacer{} อิเธกจฺโจ อชฺฌตฺตํ ปจฺจตฺตํ ปีตนิมิตฺตํ ฯเปฯ โลหิตนิมิตฺตํ ฯเปฯ โอทาตนิมิตฺตํ มนสิ กโรติ, โอทาตสญฺญํ ปฏิลภติ.\csromanspacer{} โส ตํ นิมิตฺตํ สุคฺคหิตํ กโรติ, สูปธาริตํ อุปธาเรติ, สฺวาวตฺถิตํ อวตฺถาเปติ.\csromanspacer{} โส ตํ นิมิตฺตํ สุคฺคหิตํ กตฺวา สูปธาริตํ อุปธาเรตฺวา สฺวาวตฺถิตํ อวตฺถาเปตฺวา พหิทฺธา โอทาตนิมิตฺเต จิตฺตํ อุปสํหรติ, โอทาตสญฺญํ ปฏิลภติ.\csromanspacer{} โส ตํ นิมิตฺตํ สุคฺคหิตํ กโรติ, สูปธาริตํ อุปธาเรติ, สฺวาวตฺถิตํ อวตฺถาเปติ.\csromanspacer{} โส ตํ นิมิตฺตํ สุคฺคหิตํ กตฺวา สูปธาริตํ อุปธาเรตฺวา สฺวาวตฺถิตํ อวตฺถาเปตฺวา อาเสวติ ภาเวติ พหุลีกโรติ.\csromanspacer{} ตสฺส เอวํ โหติ “อชฺฌตฺตญฺจ พหิทฺธา จ อุภยมิทํ รูปนฺ”ติ รูปสญฺญี โหติ.\csromanspacer{} เอวํ รูปี รูปานิ ปสฺสตีติ วิโมกฺโข. (37)}

\csromanheader{2}{5. วิโมกฺขกถา}
\prose{\csromanfolio{235}กถํ \textbf{อชฺฌตฺตํ อรูปสญฺญี พหิทฺธา รูปานิ ปสฺสตีติ วิโมกฺโข?} อิเธกจฺโจ อชฺฌตฺตํ ปจฺจตฺตํ นีลนิมิตฺตํ น มนสิ กโรติ, นีลสญฺญํ น ปฏิลภติ, พหิทฺธา นีลนิมิตฺเต จิตฺตํ อุปสํหรติ, นีลสญฺญํ ปฏิลภติ.\csromanspacer{} โส ตํ นิมิตฺตํ สุคฺคหิตํ กโรติ, สูปธาริตํ อุปธาเรติ, สฺวาวตฺถิตํ อวตฺถาเปติ.\csromanspacer{} โส ตํ นิมิตฺตํ สุคฺคหิตํ กตฺวา สูปธาริตํ อุปธาเรตฺวา สฺวาวตฺถิตํ อวตฺถาเปตฺวา อาเสวติ ภาเวติ พหุลีกโรติ.\csromanspacer{} ตสฺส เอวํ โหติ “อชฺฌตฺตํ อรูปํ พหิทฺธา รูปมิทนฺ”ติ รูปสญฺญี โหติ.\csromanspacer{} อิเธกจฺโจ อชฺฌตฺตํ ปจฺจตฺตํ ปีตนิมิตฺตํ ฯเปฯ โลหิตนิมิตฺตํ ฯเปฯ โอทาตนิมิตฺตํ น มนสิ กโรติ, โอทาตสญฺญํ น ปฏิลภติ.\csromanspacer{} พหิทฺธา โอทาตนิมิตฺเต จิตฺตํ อุปสํหรติ, โอทาตสญฺญํ ปฏิลภติ.\csromanspacer{} โส ตํ นิมิตฺตํ สุคฺคหิตํ กโรติ, สูปธาริตํ อุปธาเรติ, สฺวาวตฺถิตํ อวตฺถาเปติ.\csromanspacer{} โส ตํ นิมิตฺตํ สุคฺคหิตํ กตฺวา สูปธาริตํ อุปธาเรตฺวา สฺวาวตฺถิตํ อวตฺถาเปตฺวา อาเสวติ ภาเวติ พหุลีกโรติ.\csromanspacer{} ตสฺส เอวํ โหติ “อชฺฌตฺตํ อรูปํ พหิทฺธา รูปมิทนฺ”ติ รูปสญฺญี โหติ.\csromanspacer{} เอวํ อชฺฌตฺตํ อรูปสญฺญี พหิทฺธา รูปานิ ปสฺสตีติ วิโมกฺโข. (38)}

\prose{กถํ \textbf{“สุภนฺ”เตว อธิมุตฺโต โหตีติ วิโมกฺโข?} อิธ ภิกฺขุ เมตฺตาสหคเตน เจตสา เอกํ ทิสํ ผริตฺวา วิหรติ, ตถา ทุติยํ.\csromanspacer{} ตถา ตติยํ.\csromanspacer{} ตถา จตุตฺถํ.\csromanspacer{} อิติ อุทฺธมโธ ติริยํ สพฺพธิ สพฺพตฺตตาย สพฺพาวนฺตํ โลกํ เมตฺตาสหคเตน เจตสา วิปุเลน มหคฺคเตน อปฺปมาเณน อเวเรน อพฺยาปชฺเชน\footnote{อพฺยาปชฺเฌน (สฺยา)} ผริตฺวา วิหรติ, เมตฺตาย ภาวิตตฺตา สตฺตา อปฺปฏิกูลา โหนฺติ.\csromanspacer{} กรุณาสหคเตน เจเตสา ฯเปฯ กรุณาย ภาวิตตฺตา สตฺตา อปฺปฏิกูลา โหนฺติ.\csromanspacer{} มุทิตาสหคเตน เจตสา ฯเปฯ มุทิตาย ภาวิตตฺตา สตฺตา อปฺปฏิกูลา โหนฺติ.\csromanspacer{} อุเปกฺขาสหคเตน เจตสา เอกํ ทิสํ ผริตฺวา วิหรติ ฯเปฯ อุเปกฺขาย ภาวิตตฺตา สตฺตา อปฺปฏิกูลา โหนฺติ.\csromanspacer{} เอวํ “สุภนฺ”เตว อธิมุตฺโต โหตีติ วิโมกฺโข. (39)}

\proseitem{213}{กตโม \textbf{อากาสานญฺจายตนสมาปตฺติ วิโมกฺโข?} อิธ ภิกฺขุ สพฺพโส รูปสญฺญานํ สมติกฺกมา ปฏิฆสญฺญานํ อตฺถงฺคมา \csromanfolio{236}นานตฺตสญฺญานํ อมนสิการา “อนนฺโต อากาโส”ติ อากาสานญฺจายตนํ อุปสมฺปชฺช วิหรติ.\csromanspacer{} อยํ อากาสานญฺจายตนสมาปตฺติ วิโมกฺโข. (40)}

\prose{กตโม \textbf{วิญฺญาณญฺจายตนสมาปตฺติ วิโมกฺโข?} อิธ ภิกฺขุ สพฺพโส อากาสานญฺจายตนํ สมติกฺกมฺม “อนนฺตํ วิญฺญาณนฺ”ติ วิญฺญาณญฺจายตนํ อุปสมฺปชฺช วิหรติ.\csromanspacer{} อยํ วิญฺญาณญฺจายตนสมาปตฺติ วิโมกฺโข. (41)}

\prose{กตโม \textbf{อากิญฺจญฺญายตนสมาปตฺติ วิโมกฺโข?} อิธ ภิกฺขุ สพฺพโส วิญฺญาณญฺจายตนํ สมติกฺกมฺม “นตฺถิ กิญฺจี”ติ อากิญฺจญฺญายตนํ อุปสมฺปชฺช วิหรติ.\csromanspacer{} อยํ อากิญฺจญฺญายตนสมาปตฺติ วิโมกฺโข. (42)}

\prose{กตโม \textbf{เนวสญฺญานาสญฺญายตนสมาปตฺติ วิโมกฺโข?} อิธ ภิกฺขุ สพฺพโส อากิญฺจญฺญายตนํ สมติกฺกมฺม เนวสญฺญานาสญฺญายตนํ อุปสมฺปชฺช วิหรติ.\csromanspacer{} อยํ เนวสญฺญานาสญฺญายตนสมาปตฺติ วิโมกฺโข. (43)}

\prose{กตโม \textbf{สญฺญาเวทยิตนิโรธสมาปตฺติ วิโมกฺโข?} อิธ ภิกฺขุ สพฺพโส เนวสญฺญานาสญฺญายตนํ สมติกฺกมฺม สญฺญาเวทยิตนิโรธํ อุปสมฺปชฺช วิหรติ.\csromanspacer{} อยํ สญฺญาเวทยิตนิโรธสมาปตฺติ วิโมกฺโข. (44)}

\prose{กตโม \textbf{สมยวิโมกฺโข?} จตฺตาริ จ ฌานานิ, จตสฺโส จ อรูปสมาปตฺติโย.\csromanspacer{} อยํ สมยวิโมกฺโข. (45)}

\prose{กตโม \textbf{อสมยวิโมกฺโข?} จตฺตาโร จ อริยมคฺคา, จตฺตาริ จ สามญฺญผลานิ, นิพฺพานญฺจ.\csromanspacer{} อยํ อสมยวิโมกฺโข. (46)}

\prose{กตโม \textbf{สามยิโก วิโมกฺโข?} จตฺตาริ จ ฌานานิ, จตสฺโส จ อรูปสมาปตฺติโย.\csromanspacer{} อยํ สามยิโก วิโมกฺโข. (47)}

\prose{กตโม \textbf{อสามยิโก วิโมกฺโข?} จตฺตาโร จ อริยมคฺคา, จตฺตาริ จ สามญฺญผลานิ, นิพฺพานญฺจ.\csromanspacer{} อยํ อสามยิโก วิโมกฺโข. (48)}

\prose{กตโม \textbf{กุปฺโป วิโมกฺโข?} จตฺตาริ จ ฌานานิ, จตสฺโส จ อรูปสมาปตฺติโย.\csromanspacer{} อยํ กุปฺโป วิโมกฺโข. (49)}

\csromanheader{2}{5. วิโมกฺขกถา}
\prose{\csromanfolio{237}กตโม \textbf{อกุปฺโป วิโมกฺโข?} จตฺตาโร จ อริยมคฺคา, จตฺตาริ จ สามญฺญผลานิ, นิพฺพานญฺจ.\csromanspacer{} อยํ อกุปฺโป วิโมกฺโข. (50)}

\prose{กตโม \textbf{โลกิโก วิโมกฺโข?} จตฺตาริ จ ฌานานิ, จตสฺโส จ อรูปสมาปตฺติโย.\csromanspacer{} อยํ โลกิโก วิโมกฺโข. (51)}

\prose{กตโม \textbf{โลกุตฺตโร วิโมกฺโข?} จตฺตาโร จ อริยมคฺคา, จตฺตาริ จ สามญฺญผลานิ, นิพฺพานญฺจ.\csromanspacer{} อยํ โลกุตฺตโร วิโมกฺโข. (52)}

\prose{กตโม \textbf{สาสโว วิโมกฺโข?} จตฺตาริ จ ฌานานิ, จตสฺโส จ อรูปสมาปตฺติโย.\csromanspacer{} อยํ สาสโว วิโมกฺโข. (53)}

\prose{กตโม \textbf{อนาสโว วิโมกฺโข?} จตฺตาโร จ อริยมคฺคา, จตฺตาริ จ สามญฺญผลานิ, นิพฺพานญฺจ.\csromanspacer{} อยํ อนาสโว วิโมกฺโข. (54)}

\prose{กตโม \textbf{สามิโส วิโมกฺโข}? รูปปฺปฏิสญฺญุตฺโต วิโมกฺโข.\csromanspacer{} อยํ \textbf{สามิโส วิโมกฺโข}. (55)}

\prose{กตโม \textbf{นิรามิโส วิโมกฺโข?} อรูปปฺปฏิสญฺญุตฺโต วิโมกฺโข.\csromanspacer{} อยํ นิรามิโส วิโมกฺโข. (56)}

\prose{กตโม \textbf{นิรามิสา นิรามิสตโร วิโมกฺโข?} จตฺตาโร จ อริยมคฺคา, จตฺตาริ จ สามญฺญผลานิ, นิพฺพานญฺจ.\csromanspacer{} อยํ นิรามิสา นิรามิสตโร วิโมกฺโข. (57)}

\prose{กตโม \textbf{ปณิหิโต วิโมกฺโข?} จตฺตาริ จ ฌานานิ, จตสฺโส จ อรูปสมาปตฺติโย.\csromanspacer{} อยํ ปณิหิโต วิโมกฺโข.\csromanspacer{} กตโม อปฺ\textbf{ปณิหิโต วิโมกฺโข?} จตฺตาโร จ อริยมคฺคา, จตฺตาริ จ สามญฺญผลานิ, นิพฺพานญฺจ.\csromanspacer{} อยํ อปฺปณิหิโต วิโมกฺโข. (58)}

\prose{กตโม \textbf{ปณิหิตปฺปฏิปฺปสฺสทฺธิ วิโมกฺโข?} ปฐมสฺส ฌานสฺส ปฏิลาโภ วา วิปาโก วา ฯเปฯ.\csromanspacer{} เนวสญฺญานาสญฺญายตนสมาปตฺติยา ปฏิลาโภ วา วิปาโก วา.\csromanspacer{} อยํ ปณิหิตปฺปฏิปฺปสฺสทฺธิ วิโมกฺโข. (59)}

\prose{กตโม \textbf{สญฺญุตฺโต วิโมกฺโข?} จตฺตาริ จ ฌานานิ, จตสฺโส จ อรูปสมาปตฺติโย.\csromanspacer{} อยํ สญฺญุตฺโต วิโมกฺโข. (60)}

\prose{กตโม \textbf{วิสญฺญุตฺโต วิโมกฺโข?} จตฺตาโร จ อริยมคฺคา, จตฺตาริ จ สามญฺญผลานิ, นิพฺพานญฺจ.\csromanspacer{} อยํ วิสญฺญุตฺโต วิโมกฺโข. (61)}

\prose{\csromanfolio{238}กตโม \textbf{เอกตฺตวิโมกฺโข?} จตฺตาโร จ อริยมคฺคา, จตฺตาริ จ สามญฺญผลานิ, นิพฺพานญฺจ.\csromanspacer{} อยํ เอกตฺตวิโมกฺโข. (62)}

\prose{กตโม \textbf{นานตฺตวิโมกฺโข?} จตฺตาริ จ ฌานานิ, จตสฺโส จ อรูปสมาปตฺติโย.\csromanspacer{} อยํ นานตฺตวิโมกฺโข. (63)}

\proseitem{214}{กตโม \textbf{สญฺญาวิโมกฺโข?} \textbf{สิยา} เอโก สญฺญาวิโมกฺโข ทส สญฺญาวิโมกฺขา โหนฺติ.\csromanspacer{} ทส สญฺญาวิโมกฺขา เอโก สญฺญาวิโมกฺโข โหติ วตฺถุวเสน ปริยาเยน.}

\prose{\textbf{สิยา}ติ กถญฺจ สิยา? อนิจฺจานุปสฺสนญาณํ นิจฺจโต สญฺญาย มุจฺจตีติ สญฺญาวิโมกฺโข, ทุกฺขานุปสฺสนญาณํ สุขโต สญฺญาย มุจฺจตีติ สญฺญาวิโมกฺโข, อนตฺตานุปสฺสนญาณํ อตฺตโต สญฺญาย มุจฺจตีติ สญฺญาวิโมกฺโข, นิพฺพิทานุปสฺสนญาณํ นนฺทิยา สญฺญาย มุจฺจตีติ สญฺญาวิโมกฺโข, วิราคานุปสฺสนญาณํ ราคโต สญฺญาย มุจฺจตีติ สญฺญาวิโมกฺโข, นิโรธานุปสฺสนญาณํ สมุทยโต สญฺญาย มุจฺจตีติ สญฺญาวิโมกฺโข, ปฏินิสฺสคฺคานุปสฺสนญาณํ อาทานโต สญฺญาย มุจฺจตีติ สญฺญาวิโมกฺโข, อนิมิตฺตานุปสฺสนญาณํ นิมิตฺตโต สญฺญาย มุจฺจตีติ สญฺญาวิโมกฺโข, อปฺปณิหิตานุปสฺสนญาณํ ปณิธิยา สญฺญาย มุจฺจตีติ สญฺญาวิโมกฺโข, สุญฺญตานุปสฺสนญาณํ อภินิเวสโต สญฺญาย มุจฺจตีติ สญฺญาวิโมกฺโข.\csromanspacer{} เอวํ สิยา เอโก สญฺญาวิโมกฺโข ทส สญฺญาวิโมกฺขา โหนฺติ, ทส สญฺญาวิโมกฺขา เอโก สญฺญาวิโมกฺโข โหติ วตฺถุวเสน ปริยาเยน.}

\prose{รูเป อนิจฺจานุปสฺสนญาณํ นิจฺจโต สญฺญาย มุจฺจตีติ สญฺญาวิโมกฺโข ฯเปฯ รูเป สุญฺญตานุปสฺสนญาณํ อภินิเวสโต สญฺญาย มุจฺจตีติ สญฺญาวิโมกฺโข.\csromanspacer{} เอวํ สิยา เอโก สญฺญาวิโมกฺโข ทส สญฺญาวิโมกฺขา โหนฺติ.\csromanspacer{} ทส สญฺญาวิโมกฺขา เอโก สญฺญาวิโมกฺโข โหติ วตฺถุวเสน ปริยาเยน.}

\prose{เวทนาย ฯเปฯ.\csromanspacer{} สญฺญาย.\csromanspacer{} สงฺขาเรสุ.\csromanspacer{} วิญฺญาเณ.\csromanspacer{} จกฺขุสฺมิํ ฯเปฯ.\csromanspacer{} ชรามรเณ อนิจฺจานุปสฺสนญาณํ นิจฺจโต สญฺญาย มุจฺจตีติ สญฺญาวิโมกฺโข ฯเปฯ ชรามรเณ สุญฺญตานุปสฺสนญาณํ อภินิเวสโต สญฺญาย มุจฺจตีติ}

\csromanheader{2}{5. วิโมกฺขกถา}
\prose{\csromanfolio{239}สญฺญาวิโมกฺโข.\csromanspacer{} เอวํ สิยา เอโก สญฺญาวิโมกฺโข ทส สญฺญาวิโมกฺขา โหนฺติ.\csromanspacer{} ทส สญฺญาวิโมกฺขา เอโก สญฺญาวิโมกฺโข โหติ วตฺถุวเสน ปริยาเยน.\csromanspacer{} อยํ สญฺญาวิโมกฺโข. (64)}

\proseitem{215}{กตโม \textbf{ญาณวิโมกฺโข?} \textbf{สิยา} เอโก ญาณวิโมกฺโข ทส ญาณวิโมกฺขา โหนฺติ.\csromanspacer{} ทส ญาณวิโมกฺขา เอโก ญาณวิโมกฺโข โหติ วตฺถุวเสน ปริยาเยน.}

\prose{\textbf{สิยา}ติ กถญฺจ สิยา? อนิจฺจานุปสฺสนา ยถาภูตํ ญาณํ นิจฺจโต สมฺโมหา อญฺญาณา มุจฺจตีติ ญาณวิโมกฺโข, ทุกฺขานุปสฺสนา ยถาภูตํ ญาณํ สุขโต สมฺโมหา อญฺญาณา มุจฺจตีติ ญาณวิโมกฺโข, อนตฺตานุปสฺสนา ยถาภูตํ ญาณํ อตฺตโต สมฺโมหา อญฺญาณา มุจฺจตีติ ญาณวิโมกฺโข, นิพฺพิทานุปสฺสนา ยถาภูตํ ญาณํ นนฺทิยา สมฺโมหา อญฺญาณา มุจฺจตีติ ญาณวิโมกฺโข, วิราคานุปสฺสนา ยถาภูตํ ญาณํ ราคโต สมฺโมหา อญฺญาณา มุจฺจตีติ ญาณวิโมกฺโข, นิโรธานุปสฺสนา ยถาภูตํ ญาณํ สมุทยโต สมฺโมหา อญฺญาณา มุจฺจตีติ ญาณวิโมกฺโข, ปฏินิสฺสคฺคานุปสฺสนา ยถาภูตํ ญาณํ อาทานโต สมฺโมหา อญฺญาณา มุจฺจตีติ ญาณวิโมกฺโข, อนิมิตฺตานุปสฺสนา ยถาภูตํ ญาณํ นิมิตฺตโต สมฺโมหา อญฺญาณา มุจฺจตีติ ญาณวิโมกฺโข, อปฺปณิหิตานุปสฺสนา ยถาภูตํ ญาณํ ปณิธิยา สมฺโมหา อญฺญาณา มุจฺจตีติ ญาณวิโมกฺโข, สุญฺญตานุปสฺสนา ยถาภูตํ ญาณํ อภินิเวสโต สมฺโมหา อญฺญาณา มุจฺจตีติ ญาณวิโมกฺโข.\csromanspacer{} เอวํ สิยา เอโก ญาณวิโมกฺโข ทส ญาณวิโมกฺขา โหนฺติ, ทส ญาณวิโมกฺขา เอโก ญาณวิโมกฺโข โหติ วตฺถุวเสน ปริยาเยน.}

\prose{รูเป อนิจฺจานุปสฺสนา ยถาภูตํ ญาณํ นิจฺจโต สมฺโมหา อญฺญาณา มุจฺจตีติ ญาณวิโมกฺโข ฯเปฯ รูเป สุญฺญตานุปสฺสนา ยถาภูตํ ญาณํ อภินิเวสโต สมฺโมหา อญฺญาณา มุจฺจตีติ ญาณวิโมกฺโข.\csromanspacer{} เอวํ สิยา เอโก ญาณวิโมกฺโข ทส ญาณวิโมกฺขา โหนฺติ, ทส ญาณวิโมกฺขา เอโก ญาณวิโมกฺโข โหติ วตฺถุวเสน ปริยาเยน.}

\prose{\csromanfolio{240}เวทนาย ฯเปฯ.\csromanspacer{} สญฺญาย.\csromanspacer{} สงฺขาเรสุ.\csromanspacer{} วิญฺญาเณ.\csromanspacer{} จกฺขุสฺมิํ ฯเปฯ.\csromanspacer{} ชรามรเณ อนิจฺจานุปสฺสนา ยถาภูตํ ญาณํ นิจฺจโต สมฺโมหา อญฺญาณา มุจฺจตีติ ญาณวิโมกฺโข ฯเปฯ ชรามรเณ สุญฺญตานุปสฺสนา ยถาภูตํ ญาณํ อภินิเวสโต สมฺโมหา อญฺญาณา มุจฺจตีติ ญาณวิโมกฺโข.\csromanspacer{} เอวํ สิยา เอโก ญาณวิโมกฺโข ทส ญาณวิโมกฺขา โหนฺติ, ทส ญาณวิโมกฺขา เอโก ญาณวิโมกฺโข โหติ วตฺถุวเสน ปริยาเยน.\csromanspacer{} อยํ ญาณวิโมกฺโข. (65)}

\proseitem{216}{กตโม \textbf{สีติสิยาวิโมกฺโข?} \textbf{สิยา} เอโก สีติสิยาวิโมกฺโข ทส สีติสิยาวิโมกฺขา โหนฺติ, ทส สีติสิยาวิโมกฺขา เอโก สีติสิยาวิโมกฺโข โหติ วตฺถุวเสน ปริยาเยน.}

\prose{\textbf{สิยา}ติ กถญฺจ สิยา? อนิจฺจานุปสฺสนา อนุตฺตรํ สีติภาวญาณํ นิจฺจโต สนฺตาปปริฬาหทรถา มุจฺจตีติ สีติสิยาวิโมกฺโข.\csromanspacer{} ทุกฺขานุปสฺสนา อนุตฺตรํ สีติภาวญาณํ สุขโต สนฺตาปปริฬาหทรถา มุจฺจตีติ สีติสิยาวิโมกฺโข.\csromanspacer{} อนตฺตานุปสฺสนา อนุตฺตรํ สีติภาวญาณํ อตฺตโต สนฺตาปปริฬาหทรถา มุจฺจตีติ สีติสิยาวิโมกฺโข.\csromanspacer{} นิพฺพิทานุปสฺสนา อนุตฺตรํ สีติภาวญาณํ นนฺทิยา สนฺตาปปริฬาหทรถา มุจฺจตีติ สีติสิยาวิโมกฺโข.\csromanspacer{} วิราคานุปสฺสนา อนุตฺตรํ สีติภาวญาณํ ราคโต สนฺตาปปริฬาหทรถา มุจฺจตีติ สีติสิยาวิโมกฺโข.\csromanspacer{} นิโรธานุปสฺสนา อนุตฺตรํ สีติภาวญาณํ สมุทยโต สนฺตาปปริฬาหทรถา มุจฺจตีติ สีติสิยาวิโมกฺโข.\csromanspacer{} ปฏินิสฺสคฺคานุปสฺสนา อนุตฺตรํ สีติภาวญาณํ อาทานโต สนฺตาปปริฬาหทรถา มุจฺจตีติ สีกิสิยาวิโมกฺโข.\csromanspacer{} อนิมิตฺตานุปสฺสนา อนุตฺตรํ สีติภาวญาณํ นิมิตฺตโต สนฺตาปปริฬาหทรถา มุจฺจตีติ สีติสิยาวิโมกฺโข.\csromanspacer{} อปฺปณิหิตานุปสฺสนา อนุตฺตรํ สีติภาวญาณํ ปณิธิยา สนฺตาปปริฬาหทรถา มุจฺจตีติ สีติสิยาวิโมกฺโข.\csromanspacer{} สุญฺญตานุปสฺสนา อนุตฺตรํ สีติภาวญาณํ อภินิเวสโต สนฺตาปปริฬาหทรถา มุจฺจตีติ สีติสิยาวิโมกฺโข.\csromanspacer{} เอวํ สิยา เอโก สีติสิยาวิโมกฺโข ทส สีติสิยาวิโมกฺขา โหนฺติ, ทส สีติสิยาวิโมกฺขา เอโก สีติสิยาวิโมกฺโข โหติ วตฺถุวเสน ปริยาเยน.}

\csromanheader{2}{5. วิโมกฺขกถา}
\prose{\csromanfolio{241}รูเป อนิจฺจานุปสฺสนา อนุตฺตรํ สีติภาวญาณํ นิจฺจโต สนฺตาปปริฬาหทรถา มุจฺจตีติ สีติสิยาวิโมกฺโข ฯเปฯ รูเป สุญฺญาตานุปสฺสนา อนุตฺตรํ สีติภาวญาณํ อภินิเวสโต สนฺตาปปริฬาหทรถา มุจฺจตีติ สีติสิยาวิโมกฺโข.\csromanspacer{} เอวํ สิยา เอโก สีติสิยาวิโมกฺโข ทส สีติสิยาวิโมกฺขา โหนฺติ, ทส สีติสิยาวิโมกฺขา เอโก สีติสิยาวิโมกฺโข โหติ วตฺถุวเสน ปริยาเยน.}

\prose{เวทนาย ฯเปฯ.\csromanspacer{} สญฺญาย.\csromanspacer{} สงฺขาเรสุ.\csromanspacer{} วิญฺญาเณ.\csromanspacer{} จกฺขุสฺมิํ ฯเปฯ.\csromanspacer{} ชรามรเณ อนิจฺจานุปสฺสนา อนุตฺตรํ สีติภาวญาณํ นิจฺจโต สนฺตาปปริฬาหทรถา มุจฺจตีติ สีติสิยาวิโมกฺโข ฯเปฯ ชรามรเณ สุญฺญตานุปสฺสนา อนุตฺตรํ สีติภาวญาณํ อภินิเวสโต สนฺตาปปริฬาหทรถา มุจฺจตีติ สีติสิยาวิโมกฺโข.\csromanspacer{} เอวํ สิยา เอโก สีติสิยาวิโมกฺโข ทส สีติสิยาวิโมกฺขา โหนฺติ, ทส สีติสิยาวิโมกฺขา เอโก สีติสิยาวิโมกฺโข โหติ วตฺถุวเสน ปริยาเยน, อยํ สีติสียาวิโมกฺโข. (66)}

\proseitem{217}{กตโม \textbf{ฌานวิโมกฺโข?} เนกฺขมฺมํฌายตีติ\footnote{ชายตีติ (สฺยา)} ฌานํ, กามจฺฉนฺทํ ฌาเปตีติ ฌานํ, ฌายนฺโต\footnote{ชายนฺโต (สฺยา)} มุจฺจตีติ ฌานวิโมกฺโข, ฌาเปนฺโต มุจฺจตีติ ฌานวิโมกฺโข.\csromanspacer{} ฌายนฺตีติ ธมฺมา, ฌาเปตีติ กิเลเส, ฌาเต จ ฌาเป จ ชานาตีติ ฌานวิโมกฺโข.\csromanspacer{} อพฺยาปาโท ฌายตีติ ฌานํ พฺยาปาทํ ฌาเปตีติ ฌานํ.\csromanspacer{} ฌายนฺโต มุจฺจตีติ ฌานวิโมกฺโข, ฌาเปนฺโต มุจฺจตีติ ฌานวิโมกฺโข.\csromanspacer{} ฌายนฺตีติ ธมฺมา, ฌาเปตีติ กิเลเส, ฌาเต จ ฌาเป จ ชานาตีติ ฌานวิโมกฺโข.\csromanspacer{} อาโลกสญฺญา ฌายตีติ ฌานํ.\csromanspacer{} ถินมิทฺธํ ฌาเปตีติ ฌานํ ฯเปฯ.\csromanspacer{} อวิกฺเขโป ฌายตีติ ฌานํ, อุทฺธจฺจํ ฌาเปตีติ ฌานํ.\csromanspacer{} ธมฺมววตฺถานํ ฌายตีติ ฌานํ, วิจิกิจฺฉํ ฌาเปตีติ ฌานํ.\csromanspacer{} ญาณํ ฌายตีติ ฌานํ, อวิชฺชํ ฌาเปตีติ ฌานํ.\csromanspacer{} ปาโมชฺชํ ฌายตีติ ฌานํ, อรติํ ฌาเปตีติ ฌานํ.\csromanspacer{} ปฐมํ ฌานํ ฌายตีติ ฌานํ, นีวรเณ ฌาเปตีติ ฌานํ ฯเปฯ.\csromanspacer{} อรหตฺตมคฺโค ฌายตีติ ฌานํ สพฺพกิเลเส ฌาเปตีติ ฌานํ.\csromanspacer{} ฌายนฺโต มุจฺจตีติ ฌานวิโมกฺโข, ฌาเปนฺโต มุจฺจตีติ ฌานวิโมกฺโข.\csromanspacer{} ฌายนฺตีติ ธมฺมา, ฌาเปตีติ กิเลเส, ฌาเต จ ฌาเป จ ชานาตีติ ฌานวิโมกฺโข.\csromanspacer{} อยํ ฌานวิโมกฺโข. (67)}

\proseitem{218}{\csromanfolio{242}กตโม \textbf{อนุปาทา จิตฺตสฺส วิโมกฺโข?} \textbf{สิยา} เอโก อนุปาทา จิตฺตสฺส วิโมกฺโข ทส อนุปาทา จิตฺตสฺส วิโมกฺขา โหนฺติ, ทส อนุปาทา จิตฺตสฺส วิโมกฺขา เอโก อนุปาทา จิตฺตสฺส วิโมกฺโข โหติ วตฺถุวเสน ปริยาเยน.}

\prose{\textbf{สิยา}ติ กถญฺจ สิยา? อนิจฺจานุปสฺสนญาณํ นิจฺจโต อุปาทานา มุจฺจตีติ อนุปาทา จิตฺตสฺส วิโมกฺโข.\csromanspacer{} ทุกฺขานุปสฺสนญาณํ สุขโต อุปาทานา มุจฺจตีติ อนุปาทา จิตฺตสฺส วิโมกฺโข.\csromanspacer{} อนตฺตานุปสฺสนญาณํ อตฺตโต อุปาทานา มุจฺจตีติ อนุปาทา จิตฺตสฺส วิโมกฺโข.\csromanspacer{} นิพฺพิทานุปสฺสนญาณํ นนฺทิยา อุปาทานา มุจฺจตีติ อนุปาทา จิตฺตสฺส วิโมกฺโข.\csromanspacer{} วิราคานุปสฺสนญาณํ ราคโต อุปาทานา มุจฺจตีติ อนุปาทา จิตฺตสฺส วิโมกฺโข.\csromanspacer{} นิโรธานุปสฺสนญาณํ สมุทยโต อุปาทานา มุจฺจตีติ อนุปาทา จิตฺตสฺส วิโมกฺโข.\csromanspacer{} ปฏินิสฺสคฺคานุปสฺสนญาณํ อาทานโต อุปาทานา มุจฺจตีติ อนุปาทา จิตฺตสฺส วิโมกฺโข.\csromanspacer{} อนิมิตฺตานุปสฺสนญาณํ นิมิตฺตโต อุปาทานา มุจฺจตีติ อนุปาทา จิตฺตสฺส วิโมกฺโข.\csromanspacer{} อปฺปณิหิตานุปสฺสนญาณํ ปณิธิยา อุปาทานา มุจฺจตีติ อนุปาทา จิตฺตสฺส วิโมกฺโข.\csromanspacer{} สุญฺญตานุปสฺสนญาณํ อภินิเวสโต อุปาทานา มุจฺจตีติ อนุปาทา จิตฺตสฺส วิโมกฺโข.\csromanspacer{} เอวํ สิยา เอโก อนุปาทา จิตฺตสฺส วิโมกฺโข ทส อนุปาทา จิตฺตสฺส วิโมกฺขา โหนฺติ, ทส อนุปาทา จิตฺตสฺส วิโมกฺขา เอโก อนุปาทา จิตฺตสฺส วิโมกฺโข โหติ วตฺถุวเสน ปริยาเยน.}

\prose{รูเป อนิจฺจานุปสฺสนญาณํ นิจฺจโต อุปาทานา มุจฺจตีติ อนุปาทา จิตฺตสฺส วิโมกฺโข ฯเปฯ รูเป สุญฺญตานุปสฺสนญาณํ อภินิเวสโต อุปาทานา มุจฺจตีติ อนุปาทา จิตฺตสฺส วิโมกฺโข.\csromanspacer{} เอวํ สิยา เอโก อนุปาทา จิตฺตสฺส วิโมกฺโข ทส อนุปาทา จิตฺตสฺส วิโมกฺขา โหนฺติ, ทส อนุปาทา จิตฺตสฺส วิโมกฺขา เอโก อนุปาทา จิตฺตสฺส วิโมกฺโข โหติ วตฺถุวเสน ปริยาเยน.}

\prose{เวทนาย ฯเปฯ.\csromanspacer{} สญฺญาย.\csromanspacer{} สงฺขาเรสุ.\csromanspacer{} วิญฺญาเณ.\csromanspacer{} จกฺขุสฺมิํ ฯเปฯ.\csromanspacer{} ชรามรเณ อนิจฺจานุปสฺสนญาณํ นิจฺจโต อุปาทานา มุจฺจตีติ อนุปาทา จิตฺตสฺส วิโมกฺโข ฯเปฯ ชรามรเณ สุญฺญตานุปสฺสนญาณํ อภินิเวสโต อุปาทานา มุจฺจตีติ อนุปาทา จิตฺตสฺส วิโมกฺโข.\csromanspacer{} เอวํ สิยา เอโก อนุปาทา จิตฺตสฺส วิโมกฺโข ทส อนุปาทา จิตฺตสฺส วิโมกฺขา}

\csromanheader{2}{5. วิโมกฺขกถา}
\prose{\csromanfolio{243}โหนฺติ, ทส อนุปาทา จิตฺตสฺส วิโมกฺขา เอโก อนุปาทา จิตฺตสฺส วิโมกฺโข โหติ วตฺถุวเสน ปริยาเยน.}

\prose{อนิจฺจานุปสฺสนญาณํ กติหุปาทาเนหิ มุจฺจติ, ทุกฺขานุปสฺสนญาณํ กติหุปาทาเนหิ มุจฺจติ, อนตฺตานุปสฺสนญาณํ กติหุปาทาเนหิ มุจฺจติ, นิพฺพิทานุปสฺสนญาณํ ฯเปฯ วิราคานุปสฺสนญาณํ.\csromanspacer{} นิโรธานุปสฺสนญาณํ.\csromanspacer{} ปฏินิสฺสคฺคานุปสฺสนญาณํ.\csromanspacer{} อนิมิตฺตานุปสฺสนญาณํ.\csromanspacer{} อปฺปณิหิตานุปสฺสนญาณํ.\csromanspacer{} สุญฺญตานุปสฺสนญาณํ กติหุปาทาเนหิ มุจฺจตีติ?}

\prose{อนิจฺจานุปสฺสนญาณํ ตีหุปาทาเนหิ มุจฺจติ, ทุกฺขานุปสฺสนญาณํ เอกุปาทานา มุจฺจติ, อนตฺตานุปสฺสนญาณํ ตีหุปาทาเนหิ มุจฺจติ, นิพฺพิทานุปสฺสนญาณํ เอกุปาทานา มุจฺจติ, วิราคานุปสฺสนญาณํ เอกุปาทานา มุจฺจติ, นิโรธานุปสฺสนญาณํ จตูหุปาทาเนหิ มุจฺจติ, ปฏินิสฺสคฺคานุปสฺสนญาณํ จตูหุปาทาเนหิ มุจฺจติ, อนิมิตฺตานุปสฺสนญาณํ ตีหุปาทาเนหิ มุจฺจติ, อปฺปณิหิตานุปสฺสนญาณํ เอกุปาทานา มุจฺจติ, สุญฺญตานุปสฺสนญาณํ ตีหุปาทาเนหิ มุจฺจติ.}

\prose{อนิจฺจานุปสฺสนญาณํ กตเมหิ ตีหุปาทาเนหิ มุจฺจติ, ทิฏฺฐุปาทานา สีลพฺพตุปาทานา อตฺตวาทุปาทานา, อนิจฺจานุปสฺสนญาณํ อิเมหิ ตีหุปาทาเนหิ มุจฺจติ.\csromanspacer{}\csromanspacer{}ทุกฺขานุปสฺสนญาณํ กตมา เอกุปาทานา\footnote{เอกูปาทานา (สี-ฏฺฐ)} มุจฺจติ, กามุปาทานา.\csromanspacer{} ทุกฺขานุปสฺสนญาณํ อิทํ เอกุปาทานา\footnote{อิมา เอกุปาทานา (สฺยา) 159 ปิฏฺเฐ อฏฺฐกถา โอโลเกตพฺพา.} มุจฺจติ.\csromanspacer{}\csromanspacer{}อนตฺตานุปสฺสนญาณํ กตเมหิ ตีหุปาทาเนหิ มุจฺจติ, ทิฏฺฐุปาทานา สีลพฺพตุปาทานา อตฺตวาทุปาทานา.\csromanspacer{} อนตฺตานุปสฺสนญาณํ อิเมหิ ตีหุปาทาเนหิ มุจฺจติ.\csromanspacer{}\csromanspacer{}นิพฺพิทานุปสฺสนญาณํ กตมา เอกุปาทานา มุจฺจติ, กามุปาทานา.\csromanspacer{} นิพฺพิทานุปสฺสนญาณํ อิทํ เอกุปาทานา มุจฺจติ.\csromanspacer{}\csromanspacer{}วิราคานุปสฺสนญาณํ กตมา เอกุปาทานา มุจฺจติ, กามุปาทานา.\csromanspacer{} วิราคานุปสฺสนญาณํ อิทํ เอกุปาทานา มุจฺจติ.\csromanspacer{}\csromanspacer{}นิโรธานุปสฺสนญาณํ กตเมหิ จตูหุปาทาเนหิ มุจฺจติ, กามุปาทานา ทิฏฺฐุปาทานา สีลพฺพตุปาทานา อตฺตวาทุปาทานา.\csromanspacer{} นิโรธานุปสฺสนญาณํ อิเมหิ จตูหุปาทาเนหิ มุจฺจติ.\csromanspacer{}\csromanspacer{}ปฏินิสฺสคฺคานุปสฺสนญาณํ กตเมหิ จตูหุปาทาเนหิ มุจฺจติ, กามุปาทานา ทิฏฺฐุปาทานา สีลพฺพตุปาทานา อตฺตวาทุปาทานา.\csromanspacer{} ปฏินิสฺสคฺคานุปสฺสนญาณํ อิเมหิ จตูหุปาทาเนหิ มุจฺจติ.\csromanspacer{}\csromanspacer{}อนิมิตฺตานุปสฺสนญาณํ กตเมหิ ตีหุปาทาเนหิ มุจฺจติ, ทิฏฺฐุปาทานา สีลพฺพตุปาทานา}

\prose{\csromanfolio{244}อตฺตวาทุปาทานา.\csromanspacer{} อนิมิตฺตานุปสฺสนญาณํ อิเมหิ ตีหุปาทาเนหิ มุจฺจติ.\csromanspacer{}\csromanspacer{}อปฺปณิหิตานุปสฺสนญาณํ กตมา เอกุปาทานา มุจฺจติ, กามุปาทานา.\csromanspacer{} อปฺปณิหิตานุปสฺสนญาณํ อิทํ เอกุปาทานา มุจฺจติ.\csromanspacer{}\csromanspacer{}สุญฺญตานุปสฺสนญาณํ กตเมหิ ตีหุปาทาเนหิ มุจฺจติ, ทิฏฺฐุปาทานา สีลพฺพตุปาทานา อตฺตวาทุปาทานา.\csromanspacer{} สุญฺญตานุปสฺสนญาณํ อิเมหิ ตีหุปาทาเนหิ มุจฺจติ.}

\prose{ยญฺจ อนิจฺจานุปสฺสนญาณํ ยญฺจ อนตฺตานุปสฺสนญาณํ ยญฺจ อนิมิตฺตานุปสฺสนญาณํ ยญฺจ สุญฺญตานุปสฺสนญาณํ อิมานิ จตฺตาริ ญาณานิ ตีหุปาทาเนหิ มุจฺจนฺติ, ทิฏฺฐุปาทานา สีลพฺพตุปาทานา อตฺตวาทุปาทานา.\csromanspacer{}\csromanspacer{}ยญฺจ ทุกฺขานุปสฺสนญาณํ ยญฺจ นิพฺพิทานุปสฺสนญาณํ ยญฺจ วิราคานุปสฺสนญาณํ ยญฺจ อปฺปณิหิตานุปสฺสนญาณํ อิมานิ จตฺตาริ ญาณานิ เอกุปาทานา มุจฺจนฺติ กามุปาทานา.\csromanspacer{}\csromanspacer{}ยญฺจ นิโรธานุปสฺสนญาณํ ยญฺจ ปฏินิสฺสคฺคานุปสฺสนญาณํ อิมานิ เทฺว ญาณานิ จตูหุปาทาเนหิ มุจฺจนฺติ กามุปาทานา ทิฏฺฐุปาทานา สีลพฺพตุปาทานา อตฺตวาทุปาทานา.\csromanspacer{} อยํ อนุปาทา จิตฺตสฺส วิโมกฺโข. (68) วิโมกฺขกถาย ปฐมภาณวาโร.}

\proseitem{219}{ตีณิ โข ปนิมานิ วิโมกฺขมุขานิ โลกนิยฺยานาย สํวตฺตนฺติ, สพฺพสงฺขาเร ปริจฺเฉทปริวฏุมโต\footnote{ปริจฺเฉทปริวฏฺฏุมโต (สฺยา)} สมนุปสฺสนตาย อนิมิตฺตาย จ ธาตุยา จิตฺตสมฺปกฺขนฺทนตาย, สพฺพสงฺขาเรสุ มโนสมุตฺเตชนตาย อปฺปณิหิตาย จ ธาตุยา จิตฺตสมฺปกฺขนฺทนตาย, สพฺพธมฺเม ปรโต สมนุปสฺสนตาย สุญฺญตาย จ ธาตุยา จิตฺตสมฺปกฺขนฺทนตาย, อิมานิ ตีณิ วิโมกฺขมุขานิ โลกนิยฺยานาย สํวตฺตนฺติ.}

\prose{อนิจฺจโต มนสิกโรโต กถํ สงฺขารา อุปฏฺฐนฺติ, ทุกฺขโต มนสิกโรโต กถํ สงฺขารา อุปฏฺฐนฺติ, อนตฺตโต มนสิกโรโต กถํ สงฺขารา อุปฏฺฐนฺติ? อนิจฺจโต มนสิกโรโต ขยโต สงฺขารา อุปฏฺฐนฺติ, ทุกฺขโต มนสิกโรโต ภยโต สงฺขารา อุปฏฺฐนฺติ, อนตฺตโต มนสิกโรโต สุญฺญโต สงฺขารา อุปฏฺฐนฺติ.}

\csromanheader{2}{5. วิโมกฺขกถา}
\prose{\csromanfolio{245}อนิจฺจโต มนสิกโรโต กิํ พหุลํ จิตฺตํ โหติ, ทุกฺขโต มนสิกโรโต กิํ พหุลํ จิตฺตํ โหติ, อนตฺตโต มนสิกโรโต กิํ พหุลํ จิตฺตํ โหติ? อนิจฺจโต มนสิกโรโต อธิโมกฺขพหุลํ จิตฺตํ โหติ, ทุกฺขโต มนสิกโรโต ปสฺสทฺธิพหุลํ จิตฺตํ โหติ, อนตฺตโต มนสิกโรโต เวทพหุลํ จิตฺตํ โหติ.}

\prose{อนิจฺจโต มนสิกโรนฺโต อธิโมกฺขพหุโล กตมินฺทฺริยํ ปฏิลภติ, ทุกฺขโต มนสิกโรนฺโต ปสฺสทฺธิพหุโล กตมินฺทฺริยํ ปฏิลภติ, อนตฺตโต มนสิกโรนฺโต เวทพหุโล กตมินฺทฺริยํ ปฏิลภติ? อนิจฺจโต มนสิกโรนฺโต อธิโมกฺขพหุโล สทฺธินฺทฺริยํ ปฏิลภติ, ทุกฺขโต มนสิกโรนฺโต ปสฺสทฺธิพหุโล สมาธินฺทฺริยํ ปฏิลภติ, อนตฺตโต มนสิกโรนฺโต เวทพหุโล ปญฺญินฺทฺริยํ ปฏิลภติ.}

\prose{อนิจฺจโต มนสิกโรโต อธิโมกฺขพหุลสฺส กตมินฺทฺริยํ อาธิปเตยฺยํ โหติ, ภาวนาย กตินฺทฺริยานิ ตทนฺวยา\footnote{ตทนฺวยานิ (สฺยา)} โหนฺติ, สหชาตปจฺจยา โหนฺติ, อญฺญมญฺญปจฺจยา โหนฺติ, นิสฺสยปจฺจยา โหนฺติ, สมฺปยุตฺตปจฺจยา โหนฺติ, เอกรสา โหนฺติ, เกนฏฺเฐน ภาวนา, โก ภาเวติ.\csromanspacer{} ทุกฺขโต มนสิกโรโต ปสฺสทฺธิพหุลสฺส กตมินฺทฺริยํ อาธิปเตยฺยํ โหติ, ภาวนาย กตินฺทฺริยานิ ตทนฺวยา โหนฺติ, สหชาตปจฺจยา โหนฺติ, อญฺญมญฺญปจฺจยา โหนฺติ, นิสฺสยปจฺจยา โหนฺติ, สมฺปยุตฺตปจฺจยา โหนฺติ, เอกรสา โหนฺติ, เกนฏฺเฐน ภาวนา, โก ภาเวติ.\csromanspacer{} อนตฺตโต มนสิกโรโต เวทพหุลสฺส กตมินฺทฺริยํ อาธิปเตยฺยํ โหติ, ภาวนาย กตินฺทฺริยานิ ตทนฺวยา โหนฺติ, สหชาตปจฺจยา โหนฺติ, อญฺญมญฺญปจฺจยา โหนฺติ, นิสฺสยปจฺจยา โหนฺติ, สมฺปยุตฺตปจฺจยา โหนฺติ, เอกรสา โหนฺติ, เกนฏฺเฐน ภาวนา, โก ภาเวติ? อนิจฺจโต มนสิกโรโห อธิโมกฺขพหุลสฺส สทฺธินฺทฺริยํ อาธิปเตยฺยํ โหติ, ภาวนาย จตฺตารินฺทฺริยานิ ตทนฺวยา โหนฺติ, สหชาตปจฺจยา โหนฺติ, อญฺญมญฺญปจฺจยา โหนฺติ, นิสฺสยปจฺจยา โหนฺติ, สมฺปยุตฺตปจฺจยา โหนฺติ, เอกรสา โหนฺติ, เอกรสฏฺเฐน ภาวนา, โย สมฺมาปฏิปนฺโน, โส ภาเวติ, นตฺถิ มิจฺฉาปฏิปนฺนสฺส อินฺทฺริยภาวนา.\csromanspacer{}\csromanspacer{}ทุกฺขโต มนสิกโรโต ปสฺสทฺธิพหุลสฺส สมาธินฺทฺริยํ อาธิปเตยฺยํ โหติ, ภาวนาย จตฺตารินฺทฺริยานิ ตทนฺวยา}

\prose{\csromanfolio{246}โหนฺติ, สหชาตปจฺจยา โหนฺติ, อญฺญมญฺญปจฺจยา โหนฺติ, นิสฺสยปจฺจยา โหนฺติ, สมฺปยุตฺตปจฺจยา โหนฺติ, เอกรสา โหนฺติ, เอกรสฏฺเฐน ภาวนา, โย สมฺมาปฏิปนฺโน, โส ภาเวติ, นตฺถิ มิจฺฉาปฏิปนฺนสฺส อินฺทฺริยภาวนา.\csromanspacer{}\csromanspacer{}อนตฺตโต มนสิกโรโต เวทพหุลสฺส ปญฺญินฺทฺริยํ อาธิปเตยฺยํ โหติ, ภาวนาย จตฺตารินฺทฺริยานิ ตทนฺวยา โหนฺติ, สหชาตปจฺจยา โหนฺติ, อญฺญมญฺญปจฺจยา โหนฺติ, นิสฺสยปจฺจยา โหนฺติ, สมฺปยุตฺตปจฺจยา โหนฺติ, เอกรสา โหนฺติ, เอกรสฏฺเฐน ภาวนา, โย สมฺมาปฏิปนฺโน, โส ภาเวติ, นตฺถิ มิจฺฉาปฏิปนฺนสฺส อินฺทฺริยภาวนา.}

\proseitem{220}{อนิจฺจโต มนสิกโรโต อธิโมกฺขพหุลสฺส กตมินฺทฺริยํ อาธิปเตยฺยํ โหติ, ภาวนาย กตินฺทฺริยานิ ตทนฺวยา โหนฺติ, สหชาตปจฺจยา โหนฺติ, อญฺญมญฺญปจฺจยา โหนฺติ, นิสฺสยปจฺจยา โหนฺติ, สมฺปยุตฺตปจฺจยา โหนฺติ.\csromanspacer{} ปฏิเวธกาเล กตมินฺทฺริยํ อาธิปเตยฺยํ โหติ, ปฏิเวธาย กตินฺทฺริยานิ ตทนฺวยา โหนฺติ.\csromanspacer{} สหชาตปจฺจยา โหนฺติ, อญฺญมญฺญปจฺจยา โหนฺติ, นิสฺสยปจฺจยา โหนฺติ, สมฺปยุตฺตปจฺจยา โหนฺติ, เอกรสา โหนฺติ, เกนฏฺเฐน ภาวนา, เกนฏฺเฐน ปฏิเวโธ.\csromanspacer{}\csromanspacer{}ทุกฺขโต มนสิกโรโต ปสฺสทฺธิพหุลสฺส กตมินฺทฺริยํ อาธิปเตยฺยํ โหติ, ภาวนาย กตินฺทฺริยานิ ตทนฺวยา โหนฺติ, สหชาตปจฺจยา โหนฺติ, อญฺญมญฺญปจฺจยา โหนฺติ, นิสฺสยปจฺจยา โหนฺติ, สมฺปยุตฺตปจฺจยา โหนฺติ.\csromanspacer{} ปฏิเวธกาเล กตมินฺทฺริยํ อาธิปเตยฺยํ โหติ, ปฏิเวธาย กตินฺทฺริยานิ ตทนฺวยา โหนฺติ, สหชาตปจฺจยา โหนฺติ, อญฺญมญฺญปจฺจยา โหนฺติ, นิสฺสยปจฺจยา โหนฺติ, สมฺปยุตฺตปจฺจยา โหนฺติ, เอกรสา โหนฺติ, เกนฏฺเฐน ภาวนา, เกนฏฺเฐน ปฏิเวโธ.\csromanspacer{}\csromanspacer{}อนตฺตโต มนสิกโรโต เวทพหุลสฺส กตมินฺทฺริยํ อาธิปเตยฺยํ โหติ, ภาวนาย กตินฺทฺริยานิ ตทนฺวยา โหนฺติ, สหชาตปจฺจยา โหนฺติ, อญฺญมญฺญปจฺจยา โหนฺติ, นิสฺสยปจฺจยา โหนฺติ, สมฺปยุตฺตปจฺจยา โหนฺติ, ปฏิเวธกาเล กตมินฺทฺริยํ อาธิปเตยฺยํ โหติ, ปฏิเวธาย กตินฺทฺริยานิ ตทนฺวยา โหนฺติ.\csromanspacer{} สหชาตปจฺจยา โหนฺติ, อญฺญมญฺญปจฺจยา โหนฺติ, นิสฺสยปจฺจยา โหนฺติ, สมฺปยุตฺตปจฺจยา โหนฺติ, เอกรสา โหนฺติ, เกนฏฺเฐน ภาวนา, เกนฏฺเฐน ปฏิเวโธ?}

\csromanheader{2}{5. วิโมกฺขกถา}
\prose{\csromanfolio{247}อนิจฺจโต มนสิกโรโต อธิโมกฺขพหุลสฺส สทฺธินฺทฺริยํ อาธิปเตยฺยํ โหติ, ภาวนาย จตฺตารินฺทฺริยานิ ตทนฺวยา โหนฺติ, สหชาตปจฺจยา โหนฺติ, อญฺญมญฺญปจฺจยา โหนฺติ, นิสฺสยปจฺจยา โหนฺติ, สมฺปยุตฺตปจฺจยา โหนฺติ, ปฏิเวธกาเล ปญฺญินฺทฺริยํ อาธิปเตยฺยํ โหติ, ปฏิเวธาย จตฺตารินฺทฺริยานิ ตทนฺวยา โหนฺติ, สหชาตปจฺจยา โหนฺติ, อญฺญมญฺญปจฺจยา โหนฺติ, นิสฺสยปจฺจยา โหนฺติ, สมฺปยุตฺตปจฺจยา โหนฺติ, เอกรสา โหนฺติ, เอกรสฏฺเฐน ภาวนา, ทสฺสนฏฺเฐน ปฏิเวโธ.\csromanspacer{} เอวํ ปฏิวิชฺฌนฺโตปิ ภาเวติ, ภาเวนฺโตปิ ปฏิวิชฺฌติ.\csromanspacer{}\csromanspacer{}ทุกฺขโต มนสิกโรโต ปสฺสทฺธิพหุลสฺส สมาธินฺทฺริยํ อาธิปเตยฺยํ โหติ, ภาวนาย จตฺตารินฺทฺริยานิ ตทนฺวยา โหนฺติ, สหชาตปจฺจยา โหนฺติ, อญฺญมญฺญปจฺจยา โหนฺติ นิสฺสยปจฺจยา โหนฺติ, สมฺปยุตฺตปจฺจยา โหนฺติ, ปฏิเวธกาเล ปญฺญินฺทฺริยํ อาธิปเตยฺยํ โหติ, ปฏิเวธาย จตฺตารินฺทฺริยานิ ตทนฺวยา โหนฺติ, สหชาตปจฺจยา โหนฺติ, อญฺญมญฺญปจฺจยา โหนฺติ, นิสฺสยปจฺจยา โหนฺติ, สมฺปยุตฺตปจฺจยา โหนฺติ, เอกรสา โหนฺติ, เอกรสฏฺเฐน ภาวนา, ทสฺสนฏฺเฐน ปฏิเวโธ.\csromanspacer{} เอวํ ปฏิวิชฺฌนฺโตปิ ภาเวติ, ภาเวนฺโตปิ ปฏิวิชฺฌติ.\csromanspacer{}\csromanspacer{}อนตฺตโต มนสิกโรโต เวทพหุลสฺส ปญฺญินฺทฺริยํ อาธิปเตยฺยํ โหติ, ภาวนาย จตฺตารินฺทฺริยานิ ตทนฺวยา โหนฺติ, สหชาตปจฺจยา โหนฺติ, อญฺญมญฺญปจฺจยา โหนฺติ, นิสฺสยปจฺจยา โหนฺติ, สมฺปยุตฺตปจฺจยา โหนฺติ, ปฏิเวธกาเลปิ ปญฺญินฺทฺริยํ อาธิปเตยฺยํ โหติ, ปฏิเวธาย จตฺตารินฺทฺริยานิ ตทนฺวยา โหนฺติ, สหชาตปจฺจยา โหนฺติ, อญฺญมญฺญปจฺจยา โหนฺติ, นิสฺสยปจฺจยา โหนฺติ, สมฺปยุตฺตปจฺจยา โหนฺติ, เอกรสา โหนฺติ, เอกรสฏฺเฐน ภาวนา, ทสฺสนฏฺเฐน ปฏิเวโธ.\csromanspacer{} เอวํ ปฏิวิชฺฌนฺโตปิ ภาเวติ, ภาเวนฺโตปิ ปฏิวิชฺฌติ.}

\proseitem{221}{อนิจฺจโต มนสิกโรโต กตมินฺทฺริยํ อธิมตฺตํ โหติ, กตมินฺทฺริยสฺส อธิมตฺตตฺตา สทฺธาวิมุตฺโต\footnote{สทฺธาธิมุตฺโต (สฺยา)} โหติ.\csromanspacer{}\csromanspacer{}ทุกฺขโต มนสิกโรโต กตมินฺทฺริยํ อธิมตฺตํ โหติ, กตมินฺทฺริยสฺส อธิมตฺตตฺตา กายสกฺขี\footnote{กายสกฺขิ (สฺยา, ก) อภิ 3. 132 ปิฏฺเฐ ปสฺสิตพฺพา.} โหติ.\csromanspacer{}\csromanspacer{}อนตฺตโต มนสิกโรโต กตมินฺทฺริยํ อธิมตฺตํ โหติ.\csromanspacer{} กตมินฺทฺริยสฺส อธิมตฺตตฺตา ทิฏฺฐิปฺปตฺโต โหติ?}

\prose{\csromanfolio{248}อนิจฺจโต มนสิกโรโต สทฺธินฺทฺริยํ อธิมตฺตํ โหติ, สทฺธินฺทฺริยสฺส อธิมตฺตตฺตา สทฺธาวิมุตฺโต โหติ.\csromanspacer{}\csromanspacer{}ทุกฺขโต มนสิกโรโต สมาธินฺทฺริยํ อธิมตฺตํ โหติ, สมาธินฺทฺริยสฺส อธิมตฺตตฺตา กายสกฺขี โหติ.\csromanspacer{}\csromanspacer{}อนตฺตโต มนสิกโรโต ปญฺญินฺทฺริยํ อธิมตฺตํ โหติ, ปญฺญินฺทฺริยสฺส อธิมตฺตตฺตา ทิฏฺฐิปฺปตฺโต โหติ.}

\prose{สทฺทหนฺโต วิมุตฺโตติ สทฺธาวิมุตฺโต, ผุฏฺฐตฺตา สจฺฉิกโตติ\footnote{สจฺฉิกโรตีติ (สฺยา, อิ)} กายสกฺขี, ทิฏฺฐตฺตา ปตฺโตติ ทิฏฺฐิปฺปตฺโต.\csromanspacer{} สทฺทหนฺโต วิมุจฺจตีติ สทฺธาวิมุตฺโต.\csromanspacer{} ฌานผสฺสํ ปฐมํ ผุสติ, ปจฺฉา นิโรธํ นิพฺพานํ สจฺฉิกโรตีติ กายสกฺขี. “ทุกฺขา สงฺขารา, สุโข นิโรโธ”ติ ญาตํ โหติ ทิฏฺฐํ วิทิตํ สจฺฉิกตํ ผสฺสิตํ ปญฺญายาติ ทิฏฺฐิปฺปตฺโต.\csromanspacer{} โย จายํ ปุคฺคโล สทฺธาวิมุตฺโต, โย จ กายสกฺขี, โย จ ทิฏฺฐิปฺปตฺโต, สิยา อิเม ตโย ปุคฺคลา สทฺธาวิมุตฺตาปิ กายสกฺขีปิ ทิฏฺฐิปฺปตฺตาปิ วตฺถุวเสน ปริยาเยน.}

\prose{\textbf{สิยา}ติ กถญฺจ สิยา? อนิจฺจโต มนสิกโรโต สทฺธินฺทฺริยํ อธิมตฺตํ โหติ, สทฺธินฺทฺริยสฺส อธิมตฺตตฺตา สทฺธาวิมุตฺโต โหติ.\csromanspacer{} ทุกฺขโต มนสิกโรโต สทฺธินฺทฺริยํ อธิมตฺตํ โหติ, สทฺธินฺทฺริยสฺส อธิมตฺตตฺตา สทฺธาวิมุตฺโต โหติ.\csromanspacer{} อนตฺตโต มนสิกโรโต สทฺธินฺทฺริยํ อธิมตฺตํ โหติ, สทฺธินฺทฺริยสฺส อธิมตฺตตฺตา สทฺธาวิมุตฺโต โหติ.\csromanspacer{} เอวํ อิเม ตโย ปุคฺคลา สทฺธินฺทฺริยสฺส วเสน สทฺธาวิมุตฺตา.}

\prose{ทุกฺขโต มนสิกโรโต สมาธินฺทฺริยํ อธิมตฺตํ โหติ, สมาธินฺทฺริยสฺส อธิมตฺตตฺตา กายสกฺขี โหติ.\csromanspacer{} อนตฺตโต มนสิกโรโต สมาธินฺทฺริยํ อธิมตฺตํ โหติ, สมาธินฺทฺริยสฺส อธิมตฺตตฺตา กายสกฺขี โหติ.\csromanspacer{} อนิจฺจโต มนสิกโรโต สมาธินฺทฺริยํ อธิมตฺตํ โหติ, สมาธินฺทฺริยสฺส อธิมตฺตตฺตา กายสกฺขี โหติ.\csromanspacer{} เอวํ อิเม ตโย ปุคฺคลา สมาธินฺทฺริยสฺส วเสน กายสกฺขี.}

\prose{อนตฺตโต มนสิกโรโต ปญฺญินฺทฺริยํ อธิมตฺตํ โหติ, ปญฺญินฺทฺริยสฺส อธิมตฺตตฺตา ทิฏฺฐิปฺปตฺโต โหติ.\csromanspacer{} อนิจฺจโต มนสิกโรโต ปญฺญินฺทฺริยํ อธิมตฺตํ โหติ.\csromanspacer{} ปญฺญินฺทฺริยสฺส อธิมตฺตตฺตา ทิฏฺฐิปฺปตฺโต โหติ.\csromanspacer{} ทุกฺขโต มนสิกโรโต ปญฺญินฺทฺริยํ อธิมตฺตํ โหติ, ปญฺญินฺทฺริยสฺส อธิมตฺตตฺตา}

\csromanheader{2}{5. วิโมกฺขกถา}
\prose{\csromanfolio{249}ทิฏฺฐิปฺปตฺโต โหติ.\csromanspacer{} เอวํ อิเม ตโย ปุคฺคลา ปญฺญินฺทฺริยสฺส วเสน ทิฏฺฐิปฺปตฺตา.}

\prose{โย จายํ ปุคฺคโล สทฺธาวิมุตฺโต, โย จ กายสกฺขี โย จ ทิฏฺฐิปฺปตฺโต, เอวํ สิยา อิเม ตโย ปุคฺคลา สทฺธาวิมุตฺตาปิ กายสกฺขีปิ ทิฏฺฐิปฺปตฺตาปิ วตฺถุวเสน ปริยาเยน.\csromanspacer{}\csromanspacer{}โย จายํ ปุคฺคโล สทฺธาวิมุตฺโต, โย จ กายสกฺขี, โย จ ทิฏฺฐิปฺปตฺโต, สิยา อิเม ตโย ปุคฺคลา ฯเปฯ.\csromanspacer{} อญฺโญเยว สทฺธาวิมุตฺโต, อญฺโญ กายสกฺขี, อญฺโญ ทิฏฺฐิปฺปตฺโต.}

\prose{\textbf{สิยา}ติ กถญฺจ สิยา? อนิจฺจโต มนสิกโรโต สทฺธินฺทฺริยํ อธิมตฺตํ โหติ, สทฺธินฺทฺริยสฺส อธิมตฺตตฺตา สทฺธาวิมุตฺโต โหติ.\csromanspacer{} ทุกฺขโต มนสิ กโรโต สมาธินฺทฺริยํ อธิมตฺตํ โหติ.\csromanspacer{} สมาธินฺทฺริยสฺส อธิมตฺตตฺตา กายสกฺขี โหติ.\csromanspacer{} อนตฺตโต มนสิกโรโต ปญฺญินฺทฺริยํ อธิมตฺตํ โหติ, ปญฺญินฺทฺริยสฺส อธิมตฺตตฺตา ทิฏฺฐิปฺปตฺโต โหติ.\csromanspacer{} โย จายํ ปุคฺคโล สทฺธาวิมุตฺโต, โย จ กายสกฺขี, โย จ ทิฏฺฐิปฺปตฺโต, เอวํ สิยา อิเม ตโย ปุคฺคลา สทฺธาวิมุตฺตาปิ กายสกฺขีปิ ทิฏฺฐิปฺปตฺตาปิ วตฺถุวเสน ปริยาเยน.\csromanspacer{} อญฺโญเยว สทฺธาวิมุตฺโต, อญฺโญ กายสกฺขี อญฺโญ ทิฏฺฐิปฺปตฺโต.}

\prose{อนิจฺจโต มนสิกโรโต สทฺธินฺทฺริยํ อธิมตฺตํ โหติ, สทฺธินฺทฺริยสฺส อธิมตฺตตฺตา โสตาปตฺติมคฺคํ ปฏิลภติ, เตน วุจฺจติ สทฺธานุสารี.\csromanspacer{} จตฺตารินฺทฺริยานิ ตทนฺวยา โหนฺติ, สหชาตปจฺจยา โหนฺติ, อญฺญมญฺญปจฺจยา โหนฺติ, นิสฺสยปจฺจยา โหนฺติ, สมฺปยุตฺตปจฺจยา โหนฺติ.\csromanspacer{} สทฺธินฺทฺริยสฺส วเสน จตุนฺนํ อินฺทฺริยานํ ภาวนา โหติ.\csromanspacer{} เย หิ เกจิ สทฺธินฺทฺริยสฺส วเสน โสตาปตฺติมคฺคํ ปฏิลภนฺติ, สพฺเพ เต สทฺธานุสาริโน.}

\prose{อนิจฺจโต มนสิกโรโต สทฺธินฺทฺริยํ อธิมตฺตํ โหติ, สทฺธินฺทฺริยสฺส อธิมตฺตตฺตา โสตาปตฺติผลํ สจฺฉิกตํ โหติ, เตน วุจฺจติ สทฺธาวิมุตฺโต.\csromanspacer{} จตฺตารินฺทฺริยานิ ตทนฺวยา โหนฺติ, สหชาตปจฺจยา โหนฺติ, อญฺญมญฺญปจฺจยา โหนฺติ, นิสฺสยปจฺจยา โหนฺติ, สมฺปยุตฺตปจฺจยา โหนฺติ, สทฺธินฺทฺริยสฺส วเสน จตฺตารินฺทฺริยานิ ภาวิตานิ โหนฺติ}

\prose{\csromanfolio{250}สุภาวิตานิ.\csromanspacer{} เย หิ เกจิ สทฺธินฺทฺริยสฺส วเสน โสตาปตฺติผลํ สจฺฉิกตา, สพฺเพ เต สทฺธาวิมุตฺตา.}

\prose{อนิจฺจโต มนสิกโรโต สทฺธินฺทฺริยํ อธิมตฺตํ โหติ.\csromanspacer{} สทฺธินฺทฺริยสฺส อธิมตฺตตฺตา สกทาคามิมคฺคํ ปฏิลภติ ฯเปฯ สกทาคามิผลํ สจฺฉิกตํ โหติ.\csromanspacer{} อนาคามิมคฺคํ ปฏิลภติ, อนาคามิผลํ สจฺฉิกตํ โหติ.\csromanspacer{} อรหตฺตมคฺคํ ปฏิลภติ, อรหตฺตํ\footnote{อรหตฺตผลํ (สฺยา, ก) อฏฺฐกถา โอโลเกตพฺพา.} สจฺฉิกตํ โหติ, เตน วุจฺจติ สทฺธาวิมุตฺโต.\csromanspacer{} จตฺตารินฺทฺริยานิ ตทนฺวยา โหนฺติ ฯเปฯ สมฺปยุตฺตปจฺจยา โหนฺติ, สทฺธินฺทฺริยสฺส วเสน จตฺตารินฺทฺริยานิ ภาวิตานิ โหนฺติ สุภาวิตานิ.\csromanspacer{} เย หิ เกจิ สทฺธินฺทฺริยสฺส วเสน อรหตฺตํ สจฺฉิกตา, สพฺเพ เต สทฺธาวิมุตฺตา.}

\prose{ทุกฺขโต มนสิกโรโต สมาธินฺทฺริยํ อธิมตฺตํ โหติ, สมาธินฺทฺริยสฺส อธิมตฺตตฺตา โสตาปตฺติมคฺคํ ปฏิลภติ, เตน วุจฺจติ กายสกฺขี.\csromanspacer{} จตฺตารินฺทฺริยานิ ตทนฺวยา โหนฺติ, สหชาตปจฺจยา โหนฺติ, อญฺญมญฺญปจฺจยา โหนฺติ, นิสฺสยปจฺจยา โหนฺติ.\csromanspacer{} สมฺปยุตฺตปจฺจยา โหนฺติ, สมาธินฺทฺริยสฺส วเสน จตุนฺนํ อินฺทฺริยานํ ภาวนา โหติ, เย หิ เกจิ สมาธินฺทฺริยสฺส วเสน โสตาปตฺติมคฺคํ ปฏิลภนฺติ, สพฺเพ เต กายสกฺขี.}

\prose{ทุกฺขโต มนสิกโรโต สมาธินฺทฺริยํ อธิมตฺตํ โหติ, สมาธินฺทฺริยสฺส อธิมตฺตตฺตา โสตาปตฺติผลํ สจฺฉิกตํ โหติ.\csromanspacer{} สกทาคามิมคฺคํ ปฏิลภติ, สกทาคามิผลํ สจฺฉิกตํ โหติ.\csromanspacer{} อนาคามิมคฺคํ ปฏิลภติ, อนาคามิผลํ สจฺฉิกตํ โหติ.\csromanspacer{} อรหตฺตมคฺคํ ปฏิลภติ, อรหตฺตํ สจฺฉิกตํ โหติ, เตน วุจฺจติ กายสกฺขี.\csromanspacer{} จตฺตารินฺทฺริยานิ ตทนฺวยา โหนฺติ, สหชาตปจฺจยา โหนฺติ, อญฺญมญฺญปจฺจยา โหนฺติ, นิสฺสยปจฺจยา โหนฺติ, สมฺปยุตฺตปจฺจยา โหนฺติ.\csromanspacer{} สมาธินฺทฺริยสฺส วเสน จตฺตารินฺทฺริยานิ ภาวิตานิ โหนฺติ สุภาวิตานิ.\csromanspacer{} เย หิ เกจิ สมาธินฺทฺริยสฺส วเสน อรหตฺตํ สจฺฉิกตํ, สพฺเพ เต กายสกฺขี.}

\prose{อนตฺตโต มนสิกโรโต ปญฺญินฺทฺริยํ อธิมตฺตํ โหติ, ปญฺญินฺทฺริยสฺส อธิมตฺตตฺตา โสตาปตฺติมคฺคํ ปฏิลภติ, เตน วุจฺจติ ธมฺมานุสารี.\csromanspacer{} จตฺตารินฺทฺริยานิ ตทนฺวยา โหนฺติ ฯเปฯ สมฺปยุตฺตปจฺจยา โหนฺติ.\csromanspacer{} ปญฺญินฺทฺริยสฺส}

\csromanheader{2}{5. วิโมกฺขกถา}
\prose{\csromanfolio{251}วเสน จตุนฺนํ อินฺทฺริยานํ ภาวนา โหติ.\csromanspacer{} เย หิ เกจิ ปญฺญินฺทฺริยสฺส วเสน โสตาปตฺติมคฺคํ ปฏิลภนฺติ, สพฺเพ เต ธมฺมานุสาริโน.}

\prose{อนตฺตโต มนสิกโรโต ปญฺญินฺทฺริยํ อธิมตฺตํ โหติ, ปญฺญินฺทฺริยสฺส อธิมตฺตตฺตา โสตาปตฺติผลํ สจฺฉิกตํ โหติ, เตน วุจฺจติ ทิฏฺฐิปฺปตฺโต.\csromanspacer{} จตฺตารินฺทฺริยานิ ตทนฺวยา โหนฺติ ฯเปฯ สมฺปยุตฺตปจฺจยา โหนฺติ, ปญฺญินฺทฺริยสฺส วเสน จตฺตารินฺทฺริยานิ ภาวิตานิ โหนฺติ สุภาวิตานิ.\csromanspacer{} เย หิ เกจิ ปญฺญินฺทฺริยสฺส วเสน โสตาปตฺติผลํ สจฺฉิกตา, สพฺเพ เต ทิฏฺฐิปฺปตฺตา.}

\prose{อนตฺตโต มนสิกโรโต ปญฺญินฺทฺริยํ อธิมตฺตํ โหติ, ปญฺญินฺทฺริยสฺส อธิมตฺตตฺตา สกทาคามิมคฺคํ ปฏิลภติ, สกทาคามิผลํ สจฺฉิกตํ โหติ.\csromanspacer{} อนาคามิมคฺคํ ปฏิลภติ, อนาคามิผลํ สจฺฉิกตํ โหติ.\csromanspacer{} อรหตฺตมคฺคํ ปฏิลภติ, อรหตฺตํ สจฺฉิกตํ โหติ, เตน วุจฺจติ ทิฏฺฐิปฺปตฺโต.\csromanspacer{} จตฺตารินฺทฺริยานิ ตทนฺวยา โหนฺติ, สหชาตปจฺจยา โหนฺติ, อญฺญมญฺญปจฺจยา โหนฺติ, นิสฺสยปจฺจยา โหนฺติ, สมฺปยุตฺตปจฺจยา โหนฺติ, ปญฺญินฺทฺริยสฺส วเสน จตฺตารินฺทฺริยานิ ภาวิตานิ โหนฺติ สุภาวิตานิ.\csromanspacer{} เย หิ เกจิ ปญฺญินฺทฺริยสฺส วเสน อรหตฺตํ สจฺฉิกตา, สพฺเพ เต ทิฏฺฐิปฺปตฺตา.}

\proseitem{222}{เย หิ เกจิ เนกฺขมฺมํ ภาวิตา วา ภาเวนฺติ วา ภาวิสฺสนฺติ วา, อธิคตา วา อธิคจฺฉนฺติ วา อธิคมิสฺสนฺติ วา, ปตฺตา วา ปาปุณนฺติ วา ปาปุณิสฺสนฺติ วา, ปฏิลทฺธา วา ปฏิลภนฺติ วา ปฏิลภิสฺสนฺติ วา, ปฏิวิทฺธา วา ปฏิวิชฺฌนฺติ วา ปฏิวิชฺฌิสฺสนฺติ วา, สจฺฉิกตา วา สจฺฉิกโรนฺติ วา สจฺฉิกริสฺสนฺติ วา, ผสฺสิตา วา ผสฺสนฺติ วา ผสฺสิสฺสนฺติ วา, วสิปฺปตฺตา วา ปาปุณนฺติ วา ปาปุณิสฺสนฺติ วา, ปารมิปฺปตฺตา วา ปาปุณนฺติ วา ปาปุณิสฺสนฺติ วา, เวสารชฺชปฺปตฺตา วา ปาปุณนฺติ วา ปาปุณิสฺสนฺติ วา, สพฺเพ เต สทฺธินฺทฺริยสฺส วเสน สทฺธาวิมุตฺตา, สมาธินฺทฺริยสฺส วเสน กายสกฺขี, ปญฺญินฺทฺริยสฺส วเสน ทิฏฺฐิปฺปตฺตา.}

\prose{เย หิ เกจิ อพฺยาปาทํ ฯเปฯ อาโลกสญฺญํ.\csromanspacer{} อวิกฺเขปํ.\csromanspacer{} ธมฺมววตฺถานํ.\csromanspacer{} ญาณํ.\csromanspacer{} ปาโมชฺชํ.\csromanspacer{} ปฐมํ ฌานํ.\csromanspacer{} ทุติยํ ฌานํ.\csromanspacer{} ตติยํ ฌานํ.\csromanspacer{} จตุตฺถํ ฌานํ.\csromanspacer{} อากาสานญฺจายตนสมาปตฺติํ.\csromanspacer{} วิญฺญาณญฺจายตนสมาปตฺติํ.\csromanspacer{} อากิญฺจญฺญายตนาสมาปตฺติํ.\csromanspacer{} เนวสญฺญานาสญฺญายตนสมาปตฺติํ.\csromanspacer{} อนิจฺจานุปสฺสนํ.\csromanspacer{} ทุกฺขานุปสฺสนํ.\csromanspacer{} อนตฺตานุปสฺสนํ.\csromanspacer{} นิพฺพิทานุปสฺสนํ. \csromanfolio{252}วิราคานุปสฺสนํ.\csromanspacer{} นิโรธานุปสฺสนํ.\csromanspacer{} ปฏินิสฺสคฺคานุปสฺสนํ.\csromanspacer{} ขยานุปสฺสนํ.\csromanspacer{} วยานุปสฺสนํ.\csromanspacer{} วิปริณามานุปสฺสนํ.\csromanspacer{} อนิมิตฺตานุปสฺสนํ.\csromanspacer{} อปฺปณิหิตานุปสฺสนํ.\csromanspacer{} สุญฺญตานุปสฺสนํ.\csromanspacer{} อธิปญฺญาธมฺมวิปสฺสนํ.\csromanspacer{} ยถาภูตญาณทสฺสนํ.\csromanspacer{} อาทีนวานุปสฺสนํ.\csromanspacer{} ปฏิสงฺขานุปสฺสนํ.\csromanspacer{} วิวฏฺฏนานุปสฺสนํ.\csromanspacer{} โสตาปตฺติมคฺคํ.\csromanspacer{} สกทาคามิมคฺคํ.\csromanspacer{} อนาคามิมคฺคํ.\csromanspacer{} อรหตฺตมคฺคํ.}

\prose{เย หิ เกจิ จตฺตาโร สติปฏฺฐาเน.\csromanspacer{} จตฺตาโร สมฺมปฺปธาเน.\csromanspacer{} จตฺตาโร อิทฺธิปาเท.\csromanspacer{} ปญฺจินฺทฺริยานิ.\csromanspacer{} ปญฺจ พลานิ.\csromanspacer{} สตฺต โพชฺฌงฺเค.\csromanspacer{} อริยํ อฏฺฐงฺคิกํ มคฺคํ.\csromanspacer{} เย หิ เกจิ อฏฺฐ วิโมกฺเข ภาวิตา วา ภาเวนฺติ วา ภาวิสฺสนฺติ วา, อธิคตา วา อธิคจฺฉนฺติ วา อธิคมิสฺสนฺติ วา, ปตฺตา วา ปาปุณนฺติ วา ปาปุณิสฺสนฺติ วา, ปฏิลทฺธา วา ปฏิลภนฺติ วา ปฏิลภิสฺสนฺติ วา, ปฏิวิทฺธา วา ปฏิวิชฺฌนฺติ วา ปฏิวิชฺฌิสฺสนฺติ วา, สจฺฉิกตา วา สจฺฉิกโรนฺติ วา สจฺฉิกริสฺสนฺติ วา, ผสฺสิตา วา ผสฺสนฺติ วา ผสฺสิสฺสนฺติ วา, วสิปฺปตฺตา วา ปาปุณนฺติ วา ปาปุณิสฺสนฺติ วา, ปารมิปฺปตฺตา วา ปาปุณนฺติ วา ปาปุณิสฺสนฺติ วา, เวสารชฺชปฺปตฺตา วา ปาปุณนฺติ วา ปาปุณิสฺสนฺติ วา, สพฺเพ เต สทฺธินฺทฺริยสฺส วเสน สทฺธาวิมุตฺตา, สมาธินฺทฺริยสฺส วเสน กายสกฺขี, ปญฺญินฺทฺริยสฺส วเสน ทิฏฺฐิปฺปตฺตา.}

\prose{เย หิ เกจิ จตสฺโส ปฏิสมฺภิทา ปตฺตา วา ปาปุณนฺติ วา ปาปุณิสฺสนฺติ วา ฯเปฯ สพฺเพ เต สทฺธินฺทฺริยสฺส วเสน สทฺธาวิมุตฺตา, สมาธินฺทฺริยสฺส วเสน กายสกฺขี, ปญฺญินฺทฺริยสฺส วเสน ทิฏฺฐิปฺปตฺตา.}

\prose{เย หิ เกจิ ติสฺโส วิชฺชา ปฏิวิทฺธา วา ปฏิวิชฺฌนฺติ วา ปฏิวิชฺฌิสฺสนฺติ วา ฯเปฯ สพฺเพ เต สทฺธินฺทฺริยสฺส วเสน สทฺธาวิมุตฺตา, สมาธินฺทฺริยสฺส วเสน กายสกฺขี, ปญฺญินฺทฺริยสฺส วเสน ทิฏฺฐิปฺปตฺตา.}

\prose{เย หิ เกจิ ติสฺโส สิกฺขา สิกฺขิตา วา สิกฺขนฺติ วา สิกฺขิสฺสนฺติ วา, สจฺฉิกตา วา สจฺฉิกโรนฺติ วา สจฺฉิกริสฺสนฺติ วา, ผสฺสิตา วา ผสฺสนฺติ วา ผสฺสิสฺสนฺติ วา, วสิปฺปตฺตา วา ปาปุณนฺติ วา ปาปุณิสฺสนฺติ วา, ปารมิปฺปตฺตา วา ปาปุณนฺติ วา ปาปุณิสฺสนฺติ วา, เวสารชฺชปฺปตฺตา วา ปาปุณนฺติ วา ปาปุณิสฺสนฺติ วา, สพฺเพ เต สทฺธินฺทฺริยสฺส วเสน สทฺธาวิมุตฺตา, สมาธินฺทฺริยสฺส วเสน กายสกฺขี, ปญฺญินฺทฺริยสฺส วเสน ทิฏฺฐิปฺปตฺตา.}

\csromanheader{2}{5. วิโมกฺขกถา}
\prose{\csromanfolio{253}เย หิ เกจิ ทุกฺขํ ปริชานนฺติ.\csromanspacer{} สมุทยํ ปชหนฺติ.\csromanspacer{} นิโรธํ สจฺฉิกโรนฺติ.\csromanspacer{} มคฺคํ ภาเวนฺติ.\csromanspacer{} สพฺเพ เต สทฺธินฺทฺริยสฺส วเสน สทฺธาวิมุตฺตา, สมาธินฺทฺริยสฺส วเสน กายสกฺขี, ปญฺญินฺทฺริยสฺส วเสน ทิฏฺฐิปฺปตฺตา.}

\prose{กติหากาเรหิ สจฺจปฺปฏิเวโธ โหติ, กติหากาเรหิ สจฺจานิ ปฏิวิชฺฌติ? จตูหากาเรหิ สจฺจปฺปฏิเวโธ โหติ, จตูหากาเรหิ สจฺจานิ ปฏิวิชฺฌติ.\csromanspacer{} ทุกฺขสจฺจํ ปริญฺญาปฏิเวธํ ปฏิวิชฺฌติ, สมุทยสจฺจํ ปหานปฺปฏิเวธํ ปฏิวิชฺฌติ, นิโรธสจฺจํ สจฺฉิกิริยาปฏิเวธํ ปฏิวิชฺฌติ, มคฺคสจฺจํ ภาวนาปฏิเวธํ ปฏิวิชฺฌติ.\csromanspacer{} อิเมหิ จตูหากาเรหิ สจฺจปฺปฏิเวโธ โหติ.\csromanspacer{}\csromanspacer{}อิเมหิ จตูหากาเรหิ สจฺจานิ ปฏิวิชฺฌนฺโต สทฺธินฺทฺริยสฺส วเสน สทฺธาวิมุตฺโต, สมาธินฺทฺริยสฺส วเสน กายสกฺขี, ปญฺญินฺทฺริยสฺส วเสน ทิฏฺฐิปฺปตฺโต.}

\prose{กติหากาเรหิ สจฺจปฺปฏิเวโธ โหติ, กติหากาเรหิ สจฺจานิ ปฏิวิชฺฌติ? นวหากาเรหิ สจฺจปฺปฏิเวโธ โหติ, นวหากาเรหิ สจฺจานิ ปฏิวิชฺฌติ, ทุกฺขสจฺจํ ปริญฺญาปฏิเวธํ ปฏิวิชฺฌติ, สมุทยสจฺจํ ปหานปฺปฏิเวธํ ปฏิวิชฺฌติ.\csromanspacer{} นิโรธสจฺจํ สจฺฉิกิริยาปฏิเวธํ ปฏิวิชฺฌติ, มคฺคสจฺจํ ภาวนาปฏิเวธํ ปฏิวิชฺฌติ.\csromanspacer{} อภิญฺญาปฏิเวโธ จ สพฺพธมฺมานํ, ปริญฺญาปฏิเวโธ จ สพฺพสงฺขารานํ, ปหานปฺปฏิเวโธ จ สพฺพากุสลานํ, ภาวนาปฏิเวโธ จ จตุนฺนํ มคฺคานํ, สจฺฉิกิริยาปฏิเวโธ จ นิโรธสฺส.\csromanspacer{} อิเมหิ นวหากาเรหิ สจฺจปฺปฏิเวโธ โหติ.\csromanspacer{}\csromanspacer{}อิเมหิ นวหากาเรหิ สจฺจานิ ปฏิวิชฺฌนฺโต สทฺธินฺทฺริยสฺส วเสน สทฺธาวิมุตฺโต, สมาธินฺทฺริยสฺส วเสน กายสกฺขี, ปญฺญินฺทฺริยสฺส วเสน ทิฏฺฐิปฺปตฺโต.\csromanspacer{} ทุติยภาณวาโร.}

\proseitem{223}{อนิจฺจโต มนสิกโรโต กถํ สงฺขารา อุปฏฺฐนฺติ, ทุกฺขโต มนสิกโรโต กถํ สงฺขารา อุปฏฺฐนฺติ, อนตฺตโต มนสิกโรโต กถํ สงฺขารา อุปฏฺฐนฺติ? อนิจฺจโต มนสิกโรโต ขยโต สงฺขารา อุปฏฺฐนฺติ, ทุกฺขโต มนสิกโรโต ภยโต สงฺขารา อุปฏฺฐนฺติ, อนตฺตโต มนสิกโรโต สุญฺญโต สงฺขารา อุปฏฺฐนฺติ.}

\prose{อนิจฺจโต มนสิกโรโต กิํ พหุลํ จิตฺตํ โหติ, ทุกฺขโต มนสิกโรโต กิํ พหุลํ จิตฺตํ โหติ, อนตฺตโต มนสิกโรโต \csromanfolio{254}กิํ พหุลํ จิตฺตํ โหติ? อนิจฺจโต มนสิกโรโต อธิโมกฺขพหุลํ จิตฺตํ โหติ, ทุกฺขโต มนสิกโรโต ปสฺสทฺธิพหุลํ จิตฺตํ โหติ, อนตฺตโต มนสิกโรโต เวทพหุลํ จิตฺตํ โหติ.}

\prose{อนิจฺจโต มนสิกโรนฺโต อธิโมกฺขพหุโล กตมํ วิโมกฺขํ ปฏิลภติ.\csromanspacer{} ทุกฺขโต มนสิกโรนฺโต ปสฺสทฺธิพหุโล กตมํ วิโมกฺขํ ปฏิลภติ.\csromanspacer{} อนตฺตโต มนสิกโรนฺโต เวทพหุโล กตมํ วิโมกฺขํ ปฏิลภติ.\csromanspacer{}\csromanspacer{}อนิจฺจโต มนสิกโรนฺโต อธิโมกฺขพหุโล อนิมิตฺตํ วิโมกฺขํ ปฏิลภติ.\csromanspacer{} ทุกฺขโต มนสิกโรนฺโต ปสฺสทฺธิพหุโล อปฺปณิหิตํ วิโมกฺขํ ปฏิลภติ.\csromanspacer{} อนตฺตโต มนสิกโรนฺโต เวทพหุโล สุญฺญตํ วิโมกฺขํ ปฏิลภติ.}

\proseitem{224}{อนิจฺจโต มนสิกโรโต อธิโมกฺขพหุลสฺส กตโม วิโมกฺโข อาธิปเตยฺโย โหติ, ภาวนาย กติ วิโมกฺขา ตทนฺวยา โหนฺติ, สหชาตปจฺจยา โหนฺติ, อญฺญมญฺญปจฺจยา โหนฺติ, นิสฺสยปจฺจยา โหนฺติ, สมฺปยุตฺตปจฺจยา โหนฺติ, เอกรสา โหนฺติ.\csromanspacer{} เกนฏฺเฐน ภาวนา, โก ภาเวติ.\csromanspacer{} ทุกฺขโต มนสิกโรโต ปสฺสทฺธิพหุลสฺส กตโม วิโมกฺโข อาธิปเตยฺโย โหติ, ภาวนาย กติ วิโมกฺขา ตทนฺวยา โหนฺติ, สหชาตปจฺจยา โหนฺติ, อญฺญมญฺญปจฺจยา โหนฺติ, นิสฺสยปจฺจยา โหนฺติ, สมฺปยุตฺตปจฺจยา โหนฺติ, เอกรสา โหนฺติ.\csromanspacer{} เกนฏฺเฐน ภาวนา, โก ภาเวติ.\csromanspacer{} อนตฺตโต มนสิกโรโต เวทพหุลสฺส กตโม วิโมกฺโข อาธิปเตยฺโย โหติ, ภาวนาย กติ วิโมกฺขา ตทนฺวยา โหนฺติ, สหชาตปจฺจยา โหนฺติ, อญฺญมญฺญปจฺจยา โหนฺติ, นิสฺสยปจฺจยา โหนฺติ, สมฺปยุตฺตปจฺจยา โหนฺติ, เอกรสา โหนฺติ.\csromanspacer{} เกนฏฺเฐน ภาวนา, โก ภาเวติ?}

\prose{อนิจฺจโต มนสิกโรโต อธิโมกฺขพหุลสฺส อนิมิตฺโต วิโมกฺโข อาธิปเตยฺโย โหติ, ภาวนาย เทฺว วิโมกฺขา ตทนฺวยา โหนฺติ, สหชาตปจฺจยา โหนฺติ, อญฺญมญฺญปจฺจยา โหนฺติ, นิสฺสยปจฺจยา โหนฺติ, สมฺปยุตฺตปจฺจยา โหนฺติ, เอกรสา โหนฺติ.\csromanspacer{} เอกรสฏฺเฐน ภาวนา.\csromanspacer{} โย สมฺมาปฏิปนฺโน, โส ภาเวติ.\csromanspacer{} นตฺถิ มิจฺฉาปฏิปนฺนสฺส วิโมกฺขภาวนา.\csromanspacer{} ทุกฺขโต มนสิกโรโต ปสฺสทฺธิพหุลสฺส อปฺปณิหิโต วิโมกฺโข อาธิปเตยฺโย โหติ, ภาวนาย เทฺว วิโมกฺขา}

\csromanheader{2}{5. วิโมกฺขกถา}
\prose{\csromanfolio{255}ตทนฺวยา โหนฺติ, สหชาตปจฺจยา โหนฺติ, อญฺญมญฺญปจฺจยา โหนฺติ, นิสฺสยปจฺจยา โหนฺติ, สมฺปยุตฺตปจฺจยา โหนฺติ, เอกรสา โหนฺติ.\csromanspacer{} เอกรสฏฺเฐน ภาวนา.\csromanspacer{} โย สมฺมาปฏิปนฺโน, โส ภาเวติ.\csromanspacer{} นตฺถิ มิจฺฉาปฏิปนฺนสฺส วิโมกฺขภาวนา.\csromanspacer{} อนตฺตโต มนสิกโรโต เวทพหุลสฺส สุญฺญโต วิโมกฺโข อาธิปเตยฺโย โหติ, ภาวนาย เทฺว วิโมกฺขา ตทนฺวยา โหนฺติ, สหชาตปจฺจยา โหนฺติ, อญฺญมญฺญปจฺจยา โหนฺติ, นิสฺสยปจฺจยา โหนฺติ, สมฺปยุตฺตปจฺจยา โหนฺติ, เอกรสา โหนฺติ, เอกรสฏฺเฐน ภาวนา.\csromanspacer{} โย สมฺมาปฏิปนฺโน, โส ภาเวติ.\csromanspacer{} นตฺถิ มิจฺฉาปฏิปนฺนสฺส วิโมกฺขภาวนา.}

\prose{อนิจฺจโต มนสิกโรโต อธิโมกฺขพหุลสฺส กตโม วิโมกฺโข อาธิปเตยฺโย โหติ, ภาวนาย กติ วิโมกฺขา ตทนฺวยา โหนฺติ, สหชาตปจฺจยา โหนฺติ, อญฺญมญฺญปจฺจยา โหนฺติ, นิสฺสยปจฺจยา โหนฺติ, สมฺปยุตฺตปจฺจยา โหนฺติ, เอกรสา โหนฺติ.\csromanspacer{} ปฏิเวธกาเล กตโม วิโมกฺโข อาธิปเตยฺโย โหติ, ปฏิเวธาย กติ วิโมกฺขา ตทนฺวยา โหนฺติ, สหชาตปจฺจยา โหนฺติ, อญฺญมญฺญปจฺจยา โหนฺติ, นิสฺสยปจฺจยา โหนฺติ, สมฺปยุตฺตปจฺจยา โหนฺติ, เอกรสา โหนฺติ.\csromanspacer{} เกนฏฺเฐน ภาวนา, เกนฏฺเฐน ปฏิเวโธ.\csromanspacer{}\csromanspacer{}ทุกฺขโต มนสิกโรโต ปสฺสทฺธิพหุลสฺส กตโม วิโมกฺโข อาธิปเตยฺโย โหติ, ภาวนาย กติ วิโมกฺขา ตทนฺวยา โหนฺติ, สหชาตปจฺจยา โหนฺติ, อญฺญมญฺญปจฺจยา โหนฺติ, นิสฺสยปจฺจยา โหนฺติ, สมฺปยุตฺตปจฺจยา โหนฺติ, เอกรสา โหนฺติ.\csromanspacer{} ปฏิเวธกาเล กตโม วิโมกฺโข อาธิปเตยฺโย โหติ, ปฏิเวธาย กติ วิโมกฺขา ตทนฺวยา โหนฺติ, สหชาตปจฺจยา โหนฺติ, อญฺญมญฺญปจฺจยา โหนฺติ, นิสฺสยปจฺจยา โหนฺติ, สมฺปยุตฺตปจฺจยา โหนฺติ, เอกรสา โหนฺติ.\csromanspacer{} เกนฏฺเฐน ภาวนา, เกนฏฺเฐน ปฏิเวโธ.\csromanspacer{}\csromanspacer{}อนตฺตโต มนสิกโรโต เวทพหุลสฺส กตโม วิโมกฺโข อาธิปเตยฺโย โหติ, ภาวนาย กติ วิโมกฺขา ตทนฺวยา โหนฺติ, สหชาตปจฺจยา โหนฺติ, อญฺญมญฺญปจฺจยา โหนฺติ, นิสฺสยปจฺจยา โหนฺติ, สมฺปยุตฺตปจฺจยา โหนฺติ, เอกรสา โหนฺติ.\csromanspacer{} ปฏิเวธกาเล กตโม วิโมกฺโข อาธิปเตยฺโย โหติ, ปฏิเวธาย กติวิโมกฺขา ตทนฺวยา โหนฺติ, สหชาตปจฺจยา โหนฺติ, อญฺญมญฺญปจฺจยา}

\prose{\csromanfolio{256}โหนฺติ, นิสฺสยปจฺจยา โหนฺติ, สมฺปยุตฺตปจฺจยา โหนฺติ, เอกรสา โหนฺติ.\csromanspacer{} เกนฏฺเฐน ภาวนา, เกนฏฺเฐน ปฏิเวโธ?}

\proseitem{225}{อนิจฺจโต มนสิกโรโต อธิโมกฺขพหุลสฺส อนิมิตฺโต วิโมกฺโข อาธิปเตยฺโย โหติ, ภาวนาย เทฺว วิโมกฺขา ตทนฺวยา โหนฺติ, สหชาตปจฺจยา โหนฺติ, อญฺญมญฺญปจฺจยา โหนฺติ, นิสฺสยปจฺจยา โหนฺติ, สมฺปยุตฺตปจฺจยา โหนฺติ, เอกรสา โหนฺติ.\csromanspacer{} ปฏิเวธกาเลปิ อนิมิตฺโต วิโมกฺโข อาธิปเตยฺโย โหติ, ปฏิเวธาย เทฺว วิโมกฺขา ตทนฺวยา โหนฺติ, สหชาตปจฺจยา โหนฺติ, อญฺญมญฺญปจฺจยา โหนฺติ, นิสฺสยปจฺจยา โหนฺติ, สมฺปยุตฺตปจฺจยา โหนฺติ, เอกรสา โหนฺติ.\csromanspacer{} เอกรสฏฺเฐน ภาวนา, ทสฺสนฏฺเฐน ปฏิเวโธ.\csromanspacer{} เอวํ ปฏิวิชฺฌนฺโตปิ ภาเวติ, ภาเวนฺโตปิ ปฏิวิชฺฌติ.\csromanspacer{}\csromanspacer{}ทุกฺขโต มนสิกโรโต ปสฺสทฺธิพหุลสฺส อปฺปณิหิโต วิโมกฺโข อาธิปเตยฺโย โหติ, ภาวนาย เทฺว วิโมกฺขา ตทนฺวยา โหนฺติ, สหชาตปจฺจยา โหนฺติ, อญฺญมญฺญปจฺจยา โหนฺติ, นิสฺสยปจฺจยา โหนฺติ, สมฺปยุตฺตปจฺจยา โหนฺติ, เอกรสา โหนฺติ.\csromanspacer{} ปฏิเวธกาเลปิ อปฺปณิหิโต วิโมกฺโข อาธิปเตยฺโย โหติ, ปฏิเวธาย เทฺว วิโมกฺขา ตทนฺวยา โหนฺติ, สหชาตปจฺจยา โหนฺติ, อญฺญมญฺญปจฺจยา โหนฺติ, นิสฺสยปจฺจยา โหนฺติ, สมฺปยุตฺตปจฺจยา โหนฺติ, เอกรสา โหนฺติ.\csromanspacer{} เอกรสฏฺเฐน ภาวนา, ทสฺสนฏฺเฐน ปฏิเวโธ.\csromanspacer{} เอวํ ปฏิวิชฺฌนฺโตปิ ภาเวติ, ภาเวนฺโตปิ ปฏิวิชฺฌติ.\csromanspacer{}\csromanspacer{}อนตฺตโต มนสิกโรโต เวทพหุลสฺส สุญฺญโต วิโมกฺโข อาธิปเตยฺโย โหติ, ภาวนาย เทฺว วิโมกฺขา ตทนฺวยา โหนฺติ, สหชาตปจฺจยา โหนฺติ, อญฺญมญฺญปจฺจยา โหนฺติ, นิสฺสยปจฺจยา โหนฺติ, สมฺปยุตฺตปจฺจยา โหนฺติ.\csromanspacer{} เอกรสา โหนฺติ.\csromanspacer{} ปฏิเวธกาเลปิ สุญฺญโต วิโมกฺโข อาธิปเตยฺโย โหติ, ปฏิเวธาย เทฺว วิโมกฺขา ตทนฺวยา โหนฺติ, สหชาตปจฺจยา โหนฺติ, อญฺญมญฺญปจฺจยา โหนฺติ, นิสฺสยปจฺจยา โหนฺติ, สมฺปยุตฺตปจฺจยา โหนฺติ, เอกรสา โหนฺติ, เอกรสฏฺเฐน ภาวนา.\csromanspacer{} ทสฺสนฏฺเฐน ปฏิเวธา.\csromanspacer{} เอวํ ปฏิวิชฺฌนฺโตปิ ภาเวติ, ภาเวนฺโตปิ ปฏิวิชฺฌติ.}

\proseitem{226}{อนิจฺจโต มนสิกโรโต กตโม วิโมกฺโข อธิมตฺโต โหติ, กตมวิโมกฺขสฺส อธิมตฺตตฺตา สทฺธาวิมุตฺโต โหติ.}

\csromanheader{2}{5. วิโมกฺขกถา}
\prose{\csromanfolio{257}ทุกฺขโต มนสิกโรโต กตโม วิโมกฺโข อธิมตฺโต โหติ, กตมวิโมกฺขสฺส อธิมตฺตตฺตา กายสกฺขี โหติ.\csromanspacer{} อนตฺตโต มนสิกโรโต กตโม วิโมกฺโข อธิมตฺโต โหติ, กตมวิโมกฺขสฺส อธิมตฺตตฺตา ทิฏฺฐิปฺปตฺโต โหติ?}

\prose{อนิจฺจโต มนสิกโรโต อนิมิตฺโต วิโมกฺโข อธิมตฺโต โหติ, อนิมิตฺตวิโมกฺขสฺส อธิมตฺตตฺตา สทฺธาวิมุตฺโต โหติ.\csromanspacer{} ทุกฺขโต มนสิกโรโต อปฺปณิหิโต วิโมกฺโข อธิมตฺโต โหติ, อปฺปณิหิตวิโมกฺขสฺส อธิมตฺตตฺตา กายสกฺขี โหติ.\csromanspacer{} อนตฺตโต มนสิกโรโต สุญฺญโต วิโมกฺโข อธิมตฺโต โหติ, สุญฺญตวิโมกฺขสฺส อธิมตฺตตฺตา ทิฏฺฐิปฺปตฺโต โหติ.}

\prose{สทฺทหนฺโต วิมุตฺโตติ สทฺธาวิมุตฺโต, ผุฏฺฐตฺตา สจฺฉิกโตติ กายสกฺขี, ทิฏฺฐตฺตา ปตฺโตติ ทิฏฺฐิปฺปตฺโต.\csromanspacer{} สทฺทหนฺโต วิมุจฺจตีติ สทฺธาวิมุตฺโต.\csromanspacer{} ฌานผสฺสํ ปฐมํ ผุสติ, ปจฺฉา นิโรธํ นิพฺพานํ สจฺฉิกโรตีติ กายสกฺขี. “ทุกฺขา สงฺขารา, สุโข นิโรโธ”ติ วิญฺญาตํ โหติ ทิฏฺฐํ วิทิตํ สจฺฉิกตํ ผสฺสิตํ ปญฺญายาติ ทิฏฺฐิปฺปตฺโต ฯเปฯ.}

\prose{เย หิ เกจิ เนกฺขมฺมํ ภาวิตา วา ภาเวนฺติ วา ภาวิสฺสนฺติ วา ฯเปฯ สฺพฺเพ เต อนิมิตฺตวิโมกฺขสฺส วเสน สทฺธาวิมุตฺตา, อปฺปณิหิตวิโมกฺขสฺส วเสน กายสกฺขี, สุญฺญตวิโมกฺขสฺส วเสน ทิฏฺฐิปฺปตฺตา.}

\prose{เย หิ เกจิ อพฺยาปาทํ ฯเปฯ.\csromanspacer{} อาโลกสญฺญํ.\csromanspacer{} อวิกฺเขปํ ฯเปฯ.\csromanspacer{} เย หิ เกจิ ทุกฺขํ ปริชานนฺติ, สมุทยํ ปชหนฺติ, นิโรธํ สจฺฉิกโรนฺติ, มคฺคํ ภาเวนฺติ.\csromanspacer{} สพฺเพ เต อนิมิตฺตวิโมกฺขสฺส วเสน สทฺธาวิมุตฺตา, อปฺปณิหิตวิโมกฺขสฺส วเสน กายสกฺขี, สุญฺญตวิโมกฺขสฺส วเสน ทิฏฺฐิปฺปตฺตา.}

\prose{กติหากาเรหิ สจฺจปฺปฏิเวโธ โหติ, กติหากาเรหิ สจฺจานิ ปฏิวิชฺฌติ? จตูหากาเรหิ สจฺจปฺปฏิเวโธ โหติ, จตูหากาเรหิ สจฺจานิ ปฏิวิชฺฌติ.\csromanspacer{} ทุกฺขสจฺจํ ปริญฺญาปฏิเวธํ ปฏิวิชฺฌติ, สมุทยสจฺจํ ปหานปฺปฏิเวธํ ปฏิวิชฺฌติ.\csromanspacer{} นิโรธสจฺจํ สจฺฉิกิริยาปฏิเวธํ ปฏิวิชฺฌติ.\csromanspacer{} มคฺคสจฺจํ ภาวนาปฏิเวธํ ปฏิวิชฺฌติ.\csromanspacer{} อิเมหิ จตูหากาเรหิ สจฺจปฺปฏิเวโธ โหติ.\csromanspacer{}\csromanspacer{}อิเมหิ จตูหากาเรหิ สจฺจานิ ปฏิวิชฺฌนฺโต อนิมิตฺตวิโมกฺขสฺส วเสน สทฺธาวิมุตฺโต, อปฺปณิหิตวิโมกฺขสฺส วเสน กายสกฺขี, สุญฺญตวิโมกฺขสฺส วเสน ทิฏฺฐิปฺปตฺโต.}

\prose{\csromanfolio{258}กติหากาเรหิ สจฺจ`ปฏิเวโธ โหติ, กติหากาเรหิ สจฺจานิ ปฏิวิชฺฌติ? นวหากาเรหิ สจฺจปฺปฏิเวโธ โหติ, นวหากาเรหิ สจฺจานิ ปฏิวิชฺฌติ.\csromanspacer{} ทุกฺขสจฺจํ ปริญฺญาปฏิเวธํ ปฏิวิชฺฌติ ฯเปฯ.\csromanspacer{} สจฺฉิกิริยาปฏิเวโธ จ นิโรธสฺส.\csromanspacer{} อิเมหิ นวหากาเรหิ สจฺจปฺปฏิเวโธ โหติ.\csromanspacer{}\csromanspacer{}อิเมหิ นวหากาเรหิ สจฺจานิ ปฏิวิชฺฌนฺโต อนิมิตฺตวิโมกฺขสฺส วเสน สทฺธาวิมุตฺโต, อปฺปณิหิตวิโมกฺขสฺส วเสน กายสกฺขี, สุญฺญตวิโมกฺขสฺส วเสน ทิฏฺฐิปฺปตฺโต.}

\proseitem{227}{อนิจฺจโต มนสิกโรนฺโต กตเม ธมฺเม ยถาภูตํ ชานาติ ปสฺสติ, กถํ สมฺมาทสฺสนํ โหติ, กถํ ตทนฺวเยน สพฺเพ สงฺขารา อนิจฺจโต สุทิฏฺฐา โหนฺติ, กตฺถ กงฺขา ปหียติ.\csromanspacer{} ทุกฺขโต มนสิกโรนฺโต กตเม ธมฺเม ยถาภูตํ ชานาติ ปสฺสติ, กถํ สมฺมาทสฺสนํ โหติ, กถํ ตทนฺวเยน สพฺเพ สงฺขารา ทุกฺขโต สุทิฏฺฐา โหนฺติ, กตฺถ กงฺขา ปหียติ.\csromanspacer{} อนตฺตโต มนสิกโรนฺโต กตเม ธมฺเม ยถาภูตํ ชานาติ ปสฺสติ, กถํ สมฺมาทสฺสนํ โหติ, กถํ ตทนฺวเยน สพฺเพ ธมฺมา อนตฺตโต สุทิฏฺฐา โหนฺติ, กตฺถ กงฺขา ปหียติ?}

\prose{อนิจฺจโต มนสิกโรนฺโต นิมิตฺตํ ยถาภูตํ ชานาติ ปสฺสติ, เตน วุจฺจติ สมฺมาทสฺสนํ.\csromanspacer{} เอวํ ตทนฺวเยน สพฺเพ สงฺขารา อนิจฺจโต สุทิฏฺฐา โหนฺติ, เอตฺถ กงฺขา ปหียติ.\csromanspacer{} ทุกฺขโต มนสิกโรนฺโต ปวตฺตํ ยถาภูตํ ชานาติ ปสฺสติ, เตน วุจฺจติ สมฺมาทสฺสนํ.\csromanspacer{} เอวํ ตทนฺวเยน สพฺเพ สงฺขารา ทุกฺขโต สุทิฏฺฐา โหนฺติ, เอตฺถ กงฺขา ปหียติ.\csromanspacer{} อนตฺตโต มนสิกโรนฺโต นิมิตฺตญฺจ ปวตฺตญฺจ ยถาภูตํ ชานาติ ปสฺสติ, เตน วุจฺจติ สมฺมาทสฺสนํ.\csromanspacer{} เอวํ ตทนฺวเยน สพฺเพ ธมฺมา อนตฺตโต สุทิฏฺฐา โหนฺติ, เอตฺถ กงฺขา ปหียติ.}

\prose{ยญฺจ ยถาภูตํ ญาณํ ยญฺจ สมฺมาทสฺสนํ ยา จ กงฺขาวิตรณา, อิเม ธมฺมา นานตฺถา เจว นานาพฺยญฺชนา จ, อุทาหุ เอกตฺถา, พฺยญฺชนเมว นานนฺติ? ยญฺจ ยถาภูตํ ญาณํ ยญฺจ สมฺมาทสฺสนํ ยา จ กงฺขาวิตรณา, อิเม ธมฺมา เอกตฺถา, พฺยญฺชนเมว นานํ.}

\prose{อนิจฺจโต มนสิกโรโต กิํ ภยโต อุปฏฺฐาติ, ทุกฺขโต มนสิกโรโต กิํ ภยโต อุปฏฺฐาติ, อนตฺตโต มนสิกโรโต กิํ ภยโต อุปฏฺฐาติ? อนิจฺจโต มนสิกโรโต นิมิตฺตํ ภยโต}

\csromanheader{2}{5. วิโมกฺขกถา}
\prose{\csromanfolio{259}อุปฏฺฐาติ, ทุกฺขโต มนสิกโรโต ปวตฺตํ ภยโต อุปฏฺฐาติ, อนตฺตโต มนสิกโรโต นิมิตฺตญฺจ ปวตฺตญฺจ ภยโต อุปฏฺฐาติ.}

\prose{ยา จ ภยตุปฏฺฐาเน ปญฺญา ยญฺจ อาทีนเว ญาณํ ยา จ นิพฺพิทา, อิเม ธมฺมา นานตฺถา เจว นานาพฺยญฺชนา จ.\csromanspacer{} อุทาหุ เอกตฺถา, พฺยญฺชนเมว นานนฺติ.\csromanspacer{} ยา จ ภยตุปฏฺฐาเน ปญฺญา ยญฺจ อาทีนเว ญาณํ ยา จ นิพฺพิทา, อิเม ธมฺมา เอกตฺถา, พฺยญฺชนเมว นานํ.}

\prose{ยา จ อนตฺตานุปสฺสนา ยา จ สุญฺญตานุปสฺสนา, อิเม ธมฺมา นานตฺถา เจว นานาพฺยญฺชนา จ.\csromanspacer{} อุทาหุ เอกตฺถา, พฺยญฺชนเมว นานนฺติ? ยา จ อนตฺตานุปสฺสนา ยา จ สุญฺญตานุปสฺสนา, อิเม ธมฺมา เอกตฺถา, พฺยญฺชนเมว นานํ.}

\prose{อนิจฺจโต มนสิกโรโต กิํ ปฏิสงฺขา ญาณํ อุปฺปชฺชติ, ทุกฺขโต มนสิกโรโต กิํ ปฏิสงฺขา ญาณํ อุปฺปชฺชติ, อนตฺตโต มนสิกโรโต กิํ ปฏิสงฺขา ญาณํ อุปฺปชฺชติ? อนิจฺจโต มนสิกโรโต นิมิตฺตํ ปฏิสงฺขา ญาณํ อุปฺปชฺชติ, ทุกฺขโต มนสิกโรโต ปวตฺตํ ปฏิสงฺขา ญาณํ อุปฺปชฺชติ, อนตฺตโต มนสิกโรโต นิมิตฺตญฺจ ปวตฺตญฺจ ปฏิสงฺขา ญาณํ อุปฺปชฺชติ.}

\prose{ยา จ มุจฺจิตุกมฺยตา ยา จ ปฏิสงฺขานุปสฺสนา ยา จ สงฺขารุเปกฺขา, อิเม ธมฺมา นานตฺถา เจว นานาพฺยญฺชนา จ.\csromanspacer{} อุทาหุ เอกตฺถา, พฺยญฺชนเมว นานนฺติ? ยา จ มุจฺจิตุกมฺยตา ยา จ ปฏิสงฺขานุปสฺสนา ยา จ สงฺขารุเปกฺขา อิเม ธมฺมา เอกตฺถา, พฺยญฺชนเมว นานํ.}

\prose{อนิจฺจโต มนสิกโรโต กุโต จิตฺตํ วุฏฺฐาติ, กตฺถ จิตฺตํ ปกฺขนฺทติ.\csromanspacer{} ทุกฺขโต มนสิกโรโต กุโต จิตฺตํ วุฏฺฐาติ, กตฺถ จิตฺตํ ปกฺขนฺทติ.\csromanspacer{} อนตฺตโต มนสิกโรโต กุโต จิตฺตํ วุฏฺฐาติ, กตฺถ จิตฺตํ ปกฺขนฺทติ? อนิจฺจโต มนสิกโรโต นิมิตฺตา จิตฺตํ วุฏฺฐาติ, อนิมิตฺเต จิตฺตํ ปกฺขนฺทติ.\csromanspacer{} ทุกฺขโต มนสิกโรโต ปวตฺตา จิตฺตํ วุฏฺฐาติ, อปฺปวตฺเต จิตฺตํ ปกฺขนฺทติ.\csromanspacer{} อนตฺตโต มนสิกโรโต นิมิตฺตา จ ปวตฺตา จ จิตฺตํ วุฏฺฐาติ, อนิมิตฺเต อปฺปวตฺเต นิโรเธ นิพฺพานธาตุยา\footnote{นิโรธนิพฺพานธาตุยา (สี)} จิตฺตํ ปกฺขนฺทติ.}

\prose{ยา จ พหิทฺธาวุฏฺฐานวิวฏฺฏเน ปญฺญา เย จ โคตฺรภู ธมฺมา, อิเม ธมฺมา นานตฺถา เจว นานาพฺยญฺชนา จ.\csromanspacer{} อุทาหุ เอกตฺถา, พฺยญฺชนเมว นานนฺติ? ยา จ}

\prose{\csromanfolio{260}พหิทฺธาวุฏฺฐานวิวฏฺฏเน ปญฺญา เย จ โคตฺรภู ธมฺมา, อิเม ธมฺมา เอกตฺถา, พฺยญฺชนเมว นานํ.}

\prose{อนิจฺจโต มนสิกโรนฺโต กตเมน วิโมกฺเขน วิมุจฺจติ, ทุกฺขโต มนสิกโรนฺโต กตเมน วิโมกฺเขน วิมุจฺจติ, อนตฺตโต มนสิกโรนฺโต กตเมน วิโมกฺเขน วิมุจฺจติ? อนิจฺจโต มนสิกโรนฺโต อนิมิตฺตวิโมกฺเขน วิมุจฺจติ, ทุกฺขโต มนสิกโรนฺโต อปฺปณิหิตวิโมกฺเขน วิมุจฺจติ, อนตฺตโต มนสิกโรนฺโต สุญฺญตวิโมกฺเขน วิมุจฺจติ.}

\prose{ยา จ ทุภโต วุฏฺฐานวิวฏฺฏเน ปญฺญา ยญฺจ มคฺเค ญาณํ, อิเม ธมฺมา นานตฺถา เจว นานาพฺยญฺชนา จ.\csromanspacer{} อุทาหุ เอกตฺถา, พฺยญฺชนเมว นานนฺติ? ยา จ ทุภโต วุฏฺฐานวิวฏฺฏเน ปญฺญา ยญฺจ มคฺเค ญาณํ, อิเม ธมฺมา เอกตฺถา, พฺยญฺชนเมว นานํ.}

\proseitem{228}{กติหากาเรหิ ตโย วิโมกฺขา นานากฺขเณ\footnote{นานาขเณ (สฺยา, ก)} โหนฺติ กติหากาเรหิ ตโย วิโมกฺขา เอกกฺขเณ โหนฺติ? จตูหากาเรหิ ตโย วิโมกฺขา นานากฺขเณ โหนฺติ, สตฺตหากาเรหิ ตโย วิโมกฺขา เอกกฺขเณ โหนฺติ.}

\prose{กตเมหิ จตูหากาเรหิ ตโย วิโมกฺขา นานากฺขเณ โหนฺติ? อาธิปเตยฺยฏฺเฐน อธิฏฺฐานฏฺเฐน อภินีหารฏฺเฐน นิยฺยานฏฺเฐน.\csromanspacer{} กถํ อาธิปเตยฺยฏฺเฐน ตโย วิโมกฺขา นานากฺขเณ โหนฺติ? อนิจฺจโต มนสิกโรโต อนิมิตฺโต วิโมกฺโข อาธิปเตยฺโย โหติ, ทุกฺขโต มนสิกโรโต อปฺปณิหิโต วิโมกฺโข อาธิปเตยฺโย โหติ, อนตฺตโต มนสิกโรโต สุญฺญโต วิโมกฺโข อาธิปเตยฺโย โหติ.\csromanspacer{} เอวํ อาธิปเตยฺยฏฺเฐน ตโย วิโมกฺขา นานากฺขเณ โหนฺติ.}

\prose{กถํ อธิฏฺฐานฏฺเฐน ตโย วิโมกฺขา นานากฺขเณ โหนฺติ? อนิจฺจโต มนสิกโรนฺโต อนิมิตฺตวิโมกฺขสฺส วเสน จิตฺตํ อธิฏฺฐาติ, ทุกฺขโต มนสิกโรนฺโต อปฺปณิหิตวิโมกฺขสฺส วเสน จิตฺตํ อธิฏฺฐาติ, อนตฺตโต มนสิกโรนฺโต สุญฺญตวิโมกฺขสฺส วเสน จิตฺตํ อธิฏฺฐาติ.\csromanspacer{} เอวํ อธิฏฺฐานฏฺเฐน ตโย วิโมกฺขา นานากฺขเณ โหนฺติ.}

\csromanheader{2}{5. วิโมกฺขกถา}
\prose{\csromanfolio{261}กถํ อภินีหารฏฺเฐน ตโย วิโมกฺขา นานากฺขเณ โหนฺติ? อนิจฺจโต มนสิกโรนฺโต อนิมิตฺตวิโมกฺขสฺส วเสน จิตฺตํ อภินีหรติ, ทุกฺขโต มนสิกโรนฺโต อปฺปณิหิตวิโมกฺขสฺส วเสน จิตฺตํ อภินีหรติ, อนตฺตโต มนสิกโรนฺโต สุญฺญตวิโมกฺขสฺส วเสน จิตฺตํ อภินีหรติ.\csromanspacer{} เอวํ อภินีหารฏฺเฐน ตโย วิโมกฺขา นานากฺขเณ โหนฺติ.}

\prose{กถํ นิยฺยานฏฺเฐน ตโย วิโมกฺขา นานากฺขเณ โหนฺติ? อนิจฺจโต มนสิกโรนฺโต อนิมิตฺตวิโมกฺขสฺส วเสน นิโรธํ นิพฺพานํ นิยฺยาติ, ทุกฺขโต มนสิกโรนฺโต อปฺปณิหิตวิโมกฺขสฺส วเสน นิโรธํ นิพฺพานํ นิยฺยาติ, อนตฺตโต มนสิกโรนฺโต สุญฺญตวิโมกฺขสฺส วเสน นิโรธํ นิพฺพานํ นิยฺยาติ.\csromanspacer{} เอวํ นิยฺยานฏฺเฐน ตโย วิโมกฺขา นานากฺขเณ โหนฺติ.\csromanspacer{}\csromanspacer{}อิเมหิ จตูหากาเรหิ ตโย วิโมกฺขา นานากฺขเณ โหนฺติ.}

\prose{กตเมหิ สตฺตหากาเรหิ ตโย วิโมกฺขา เอกกฺขเณ โหนฺติ? สโมธานฏฺเฐน อธิคมนฏฺเฐน ปฏิลาภฏฺเฐน ปฏิเวธฏฺเฐน สจฺฉิกิริยฏฺเฐน ผสฺสนฏฺเฐน อภิสมยฏฺเฐน.\csromanspacer{} กถํ สโมธานฏฺเฐน อธิคมนฏฺเฐน ปฏิลาภฏฺเฐน ปฏิเวธฏฺเฐน สจฺฉิกิริยฏฺเฐน ผสฺสนฏฺเฐน อภิสมยฏฺเฐน ตโย วิโมกฺขา เอกกฺขเณ โหนฺติ? อนิจฺจโต มนสิกโรนฺโต นิมิตฺตา มุจฺจตีติ อนิมิตฺโต วิโมกฺโข.\csromanspacer{} ยโต มุจฺจติ, ตตฺถ น ปณิทหตีติ อปฺปณิหิโต วิโมกฺโข.\csromanspacer{} ยตฺถ น ปณิทหติ, เตน สุญฺโญติ สุญฺญโต วิโมกฺโข.\csromanspacer{} เยน สุญฺโญ, เตน นิมิตฺเตน อนิมิตฺโตติ อนิมิตฺโต วิโมกฺโข.\csromanspacer{} เอวํ สโมธานฏฺเฐน อธิคมนฏฺเฐน ปฏิลาภฏฺเฐน ปฏิเวธฏฺเฐน สจฺฉิกิริยฏฺเฐน ผสฺสนฏฺเฐน อภิสมยฏฺเฐน ตโย วิโมกฺขา เอกกฺขเณ โหนฺติ.}

\prose{ทุกฺขโต มนสิกโรนฺโต ปณิธิยา มุจฺจตีติ อปฺปณิหิโต วิโมกฺโข.\csromanspacer{} ยตฺถ น ปณิทหติ, เตน สุญฺโญติ สุญฺญโต วิโมกฺโข.\csromanspacer{} เยน สุญฺโญ, เตน นิมิตฺเตน อนิมิตฺโตติ อนิมิตฺโต วิโมกฺโข.\csromanspacer{} เยน นิมิตฺเตน อนิมิตฺโต, ตตฺถ น ปณิทหตีติ อปฺปณิหิโต วิโมกฺโข.\csromanspacer{} เอวํ สโมธานฏฺเฐน อธิคมนฏฺเฐน ปฏิลาภฏฺเฐน ปฏิเวธฏฺเฐน สจฺฉิกิริยฏฺเฐน ผสฺสนฏฺเฐน อภิสมยฏฺเฐน ตโย วิโมกฺขา เอกกฺขเณ โหนฺติ.}

\prose{\csromanfolio{262}อนตฺตโต มนสิกโรนฺโต อภินิเวสา มุจฺจตีติ สุญฺญโต \textbf{วิโมกฺโข}.\csromanspacer{} เยน สุญฺโญ, เตน นิมิตฺเตน อนิมิตฺโตติ อนิมิตฺโต \textbf{วิโมกฺโข}.\csromanspacer{} เยน นิมิตฺเตน อนิมิตฺโต, ตตฺถ น ปณิทหตีติ อปฺปณิหิโต \textbf{วิโมกฺโข}.\csromanspacer{} ยตฺถ น ปณิทหติ, เตน สุญฺโญติ สุญฺญโต \textbf{วิโมกฺโข}.\csromanspacer{} เอวํ สโมธานฏฺเฐน อธิคมนฏฺเฐน ปฏิลาภฏฺเฐน ปฏิเวธฏฺเฐน สจฺฉิกิริยฏฺเฐน ผสฺสนฏฺเฐน อภิสมยฏฺเฐน ตโย วิโมกฺขา เอกกฺขเณ โหนฺติ.\csromanspacer{} อิเมหิ สตฺตหากาเรหิ ตโย วิโมกฺขา เอกกฺขเณ โหนฺติ.}

\proseitem{229}{อตฺถิ \textbf{วิโมกฺโข}, อตฺถิ มุขํ, อตฺถิ วิโมกฺขมุขํ, อตฺถิ วิโมกฺขปจฺจนีกํ, อตฺถิ วิโมกฺขานุโลมํ, อตฺถิ วิโมกฺขวิวฏฺโฏ, อตฺถิ วิโมกฺขภาวนา, อตฺถิ วิโมกฺขปฺปฏิปฺปสฺสทฺธิ.}

\prose{กตโม \textbf{วิโมกฺโข}? สุญฺญโต \textbf{วิโมกฺโข} อนิมิตฺโต \textbf{วิโมกฺโข} อปฺปณิหิโต \textbf{วิโมกฺโข}.\csromanspacer{}\csromanspacer{}กตโม สุญฺญโต \textbf{วิโมกฺโข}? อนิจฺจานุปสฺสนญาณํ นิจฺจโต อภินิเวสา มุจฺจตีติ สุญฺญโต \textbf{วิโมกฺโข}.\csromanspacer{} ทุกฺขานุปสฺสนญาณํ สุขโต อภินิเวสา มุจฺจตีติ สุญฺญโต \textbf{วิโมกฺโข}.\csromanspacer{} อนตฺตานุปสฺสนญาณํ อตฺตโต อภินิเวสา มุจฺจตีติ สุญฺญโต \textbf{วิโมกฺโข}.\csromanspacer{} นิพฺพิทานุปสฺสนญาณํ นนฺทิยา อภินิเวสา มุจฺจตีติ สุญฺญโต \textbf{วิโมกฺโข}.\csromanspacer{} วิราคานุปสฺสนญาณํ ราคโต อภินิเวสา มุจฺจตีติ สุญฺญโต \textbf{วิโมกฺโข}.\csromanspacer{} นิโรธานุปสฺสนญาณํ สมุทยโต อภินิเวสา มุจฺจตีติ สุญฺญโต \textbf{วิโมกฺโข}.\csromanspacer{} ปฏินิสฺสคฺคานุปสฺสนญาณํ อาทานโต อภินิเวสา มุจฺจตีติ สุญฺญโต \textbf{วิโมกฺโข}.\csromanspacer{} อนิมิตฺตานุปสฺสนญาณํ นิมิตฺตโต อภินิเวสา มุจฺจตีติ สุญฺญโต \textbf{วิโมกฺโข}.\csromanspacer{} อปฺปณิหิตานุปสฺสนญาณํ ปณิธิยา อภินิเวสา มุจฺจตีติ สุญฺญโต \textbf{วิโมกฺโข}.\csromanspacer{} สุญฺญตานุปสฺสนญาณํ สพฺพาภินิเวเสหิ มุจฺจตีติ สุญฺญโต \textbf{วิโมกฺโข}.}

\prose{รูเป อนิจฺจานุปสฺสนญาณํ นิจฺจโต อภินิเวสา มุจฺจตีติ สุญฺญโต \textbf{วิโมกฺโข} ฯเปฯ รุเป สุญฺญตานุปสฺสนญาณํ สพฺพาภินิเวเสหิ มุจฺจตีติ สุญฺญโต \textbf{วิโมกฺโข}.\csromanspacer{}\csromanspacer{}เวทนาย ฯเปฯ.\csromanspacer{} สญฺญาย.\csromanspacer{} สงฺขาเรสุ.\csromanspacer{} วิญฺญาเณ.\csromanspacer{} จกฺขุสฺมิํ ฯเปฯ.\csromanspacer{} ชรามรเณ อนิจฺจานุปสฺสนญาณํ นิจฺจโต อภินิเวสา มุจฺจตีติ สุญฺญโต \textbf{วิโมกฺโข} ฯเปฯ ชรามรเณ สุญฺญตานุปสฺสนญาณํ}

\csromanheader{2}{5. วิโมกฺขกถา}
\prose{\csromanfolio{263}สพฺพาภินิเวเสหิ มุจฺจตีติ สุญฺญโต วิโมกฺโข.\csromanspacer{} อยํ สุญฺญโต วิโมกฺโข.}

\prose{กตโม อนิมิตฺโต วิโมกฺโข? อนิจฺจานุปสฺสนญาณํ นิจฺจโต นิมิตฺตา มุจฺจตีติ อนิมิตฺโต วิโมกฺโข.\csromanspacer{} ทุกฺขานุปสฺสนญาณํ สุขโต นิมิตฺตา มุจฺจตีติ อนิมิตฺโต วิโมกฺโข.\csromanspacer{} อนตฺตานุปสฺสนญาณํ อตฺตโต นิมิตฺตา มุจฺจตีติ อนิมิตฺโต วิโมกฺโข.\csromanspacer{} นิพฺพิทานุปสฺสนญาณํ นนฺทิยา นิมิตฺตา มุจฺจตีติ อนิมิตฺโต วิโมกฺโข.\csromanspacer{} วิราคานุปสฺสนญาณํ ราคโต นิมิตฺตา มุจฺจตีติ อนิมิตฺโต วิโมกฺโข.\csromanspacer{} นิโรธานุปสฺสนญาณํ สมุทยโต นิมิตฺตา มุจฺจตีติ อนิมิตฺโต วิโมกฺโข.\csromanspacer{} ปฏินิสฺสคฺคานุปสฺสนญาณํ อาทานโต นิมิตฺตา มุจฺจตีติ อนิมิตฺโต วิโมกฺโข.\csromanspacer{} อนิมิตฺตานุปสฺสนญาณํ สพฺพนิมิตฺเตหิ มุจฺจตีติ อนิมิตฺโต วิโมกฺโข.\csromanspacer{} อปฺปณิหิตานุปสฺสนญาณํ ปณิธิยา นิมิตฺตา มุจฺจตีติ อนิมิตฺโต วิโมกฺโข.\csromanspacer{} สุญฺญตานุปสฺสนญาณํ อภินิเวสโต นิมิตฺตา มุจฺจตีติ อนิมิตฺโต วิโมกฺโข.}

\prose{รูเป อนิจฺจานุปสฺสนญาณํ นิจฺจโต นิมิตฺตา มุจฺจตีติ อนิมิตฺโต วิโมกฺโข ฯเปฯ รูเป อนิมิตฺตานุปสฺสนญาณํ สพฺพนิมิตฺเตหิ มุจฺจตีติ อนิมิตฺโต วิโมกฺโข.\csromanspacer{} รูเป อปฺปณิหิตานุปสฺสนญาณํ ปณิธิยา นิมิตฺตา มุจฺจตีติ อนิมิตฺโต วิโมกฺโข.\csromanspacer{} รูเป สุญฺญตานุปสฺสนญาณํ อภินิเวสโต นิมิตฺตา มุจฺจตีติ อนิมิตฺโต วิโมกฺโข.\csromanspacer{}\csromanspacer{}เวทนาย ฯเปฯ.\csromanspacer{} สญฺญาย.\csromanspacer{} สงฺขาเรสุ.\csromanspacer{} วิญฺญาเณ.\csromanspacer{} จกฺขุสฺมิํ ฯเปฯ.\csromanspacer{} ชรามรเณ อนิจฺจานุปสฺสนญาณํ นิจฺจโต นิมิตฺตา มุจฺจตีติ อนิมิตฺโต วิโมกฺโข ฯเปฯ ชรามรเณ อนิจฺจานุปสฺสนญาณํ สพฺพนิมิตฺเตหิ มุจฺจตีติ อนิมิตฺโต วิโมกฺโข.\csromanspacer{} ชรามรเณ อปฺปณิหิตานุปสฺสนญาณํ ปณิธิยา นิมิตฺตา มุจฺจตีติ อนิมิตฺโต วิโมกฺโข.\csromanspacer{} ชรามรเณ สุญฺญตานุปสฺสนญาณํ อภินิเวสโต นิมิตฺตา มุจฺจตีติ อนิมิตฺโต วิโมกฺโข.\csromanspacer{} อยํ อนิมิตฺโต วิโมกฺโข.}

\prose{กตโม อปฺปณิหิโต วิโมกฺโข? อนิจฺจานุปสฺสนญาณํ นิจฺจโต ปณิธิยา มุจฺจตีติ อปฺปณิหิโต วิโมกฺโข.\csromanspacer{} ทุกฺขานุปสฺสนญาณํ สุขโต ปณิธิ มุจฺจตีติ อปฺปณิหิโต วิโมกฺโข.\csromanspacer{} อนตฺตานุปสฺสนญาณํ อตฺตโต ปณิธิ มุจฺจตีติ อปฺปณิหิโต วิโมกฺโข.\csromanspacer{} นิพฺพิทานุปสฺสนญาณํ นนฺทิยา ปณิธิ มุจฺจตีติ อปฺปณิหิโต วิโมกฺโข.}

\prose{\csromanfolio{264}วิราคานุปสฺสนญาณํ ราคโต ปณิธิ มุจฺจตีติ อปฺปณิหิโต วิโมกฺโข.\csromanspacer{} นิโรธานุปสฺสนญาณํ สมุทยโต ปณิธิ มุจฺจตีติ อปฺปณิหิโต วิโมกฺโข.\csromanspacer{} ปฏินิสฺสคฺคานุปสฺสนญาณํ อาทานโต ปณิธิ มุจฺจตีติ อปฺปณิหิโต วิโมกฺโข.\csromanspacer{} อนิมิตฺตานุปสฺสนญาณํ นิมิตฺตโต ปณิธิ มุจฺจตีติ อปฺปณิหิโต วิโมกฺโข.\csromanspacer{} อปฺปณิหิตานุปสฺสนญาณํ สพฺพปณิธีหิ มุจฺจตีติ อปฺปณิหิโต วิโมกฺโข.\csromanspacer{} สุญฺญตานุปสฺสนญาณํ อภินิเวสโต ปณิธิ มุจฺจตีติ อปฺปณิหิโต วิโมกฺโข.}

\prose{รูเป อนิจฺจานุปสฺสนญาณํ นิจฺจโต ปณิธิ มุจฺจตีติ อปฺปณิหิโต วิโมกฺโข ฯเปฯ รูเป อปฺปณิหิตานุปสฺสนญาณํ สพฺพปณิธีหิ มุจฺจตีติ อปฺปณิหิโต วิโมกฺโข.\csromanspacer{} รูเป สุญฺญตานุปสฺสนญาณํ อภินิเวสโต ปณิธิ มุจฺจตีติ อปฺปณิหิโต วิโมกฺโข.\csromanspacer{}\csromanspacer{}เวทนาย ฯเปฯ.\csromanspacer{} สญฺญาย.\csromanspacer{} สงฺขาเรสุ.\csromanspacer{} วิญฺญาเณ.\csromanspacer{} จกฺขุสฺมิํ ฯเปฯ.\csromanspacer{} ชรามรเณ อนิจฺจานุปสฺสนญาณํ นิจฺจโต ปณิธิ มุจฺจตีติ อปฺปณิหิโต วิโมกฺโข ฯเปฯ ชรามรเณ อปฺปณิหิตานุปสฺสนญาณํ สพฺพปณิธีหิ มุจฺจตีติ อปฺปณิหิโต วิโมกฺโข.\csromanspacer{} ชรามรเณ สุญฺญตานุปสฺสนญาณํ อภินิเวสโต ปณิธิ มุจฺจตีติ อปฺปณิหิโต วิโมกฺโข.\csromanspacer{} อยํ อปฺปณิหิโต วิโมกฺโข.\csromanspacer{} อยํ วิโมกฺโข. (1)}

\proseitem{230}{กตมํ \textbf{มุขํ?} เย ตตฺถ ชาตา อนวชฺชกุสลา โพธิปกฺขิยา ธมฺมา, อิทํ มุขํ. (2)}

\prose{กตมํ \textbf{วิโมกฺขมุขํ?} ยํ เตสํ ธมฺมานํ อารมฺมณํ นิโรโธ นิพฺพานํ, อิทํ วิโมกฺขมุขํ.\csromanspacer{}\csromanspacer{}วิโมกฺขญฺจ มุขญฺจ วิโมกฺขมุขํ, อิทํ วิโมกฺขมุขํ. (3)}

\prose{กตมํ \textbf{วิโมกฺขปจฺจนีกํ?} ตีณิ อกุสลมูลานิ วิโมกฺขปจฺจนีกานิ, ตีณิ ทุจฺจริตานิ วิโมกฺขปจฺจนีกานิ, สพฺเพปิ อกุสลา ธมฺมา วิโมกฺขปจฺจนีกา, อิทํ วิโมกฺขปจฺจนีกํ. (4)}

\prose{กตมํ \textbf{วิโมกฺขานุโลมํ?} ตีณิ กุสลมูลานิ วิโมกฺขานุโลมานิ, ตีณิ สุจริตานิ วิโมกฺขานุโลมานิ, สพฺเพปิ กุสลา ธมฺมา วิโมกฺขานุโลมา, อิทํ วิโมกฺขานุโลมํ. (5)}

\csromanheader{2}{5. วิโมกฺขกถา}
\prose{\csromanfolio{265}กตโม \textbf{วิโมกฺขวิวฏฺโฏ?} สญฺญาวิวฏฺโฏ เจโตวิวฏฺโฏ จิตฺตวิวฏฺโฏ ญาณวิวฏฺโฏ วิโมกฺขวิวฏฺโฏ สจฺจวิวฏฺโฏ.\csromanspacer{} สญฺชานนฺโต วิวฏฺฏตีติ สญฺญาวิวฏฺโฏ, เจตยนฺโต วิวฏฺฏตีติ เจโตวิวฏฺโฏ, วิชานนฺโต วิวฏฺฏตีติ จิตฺตวิวฏฺโฏ, ญาณํ กโรนฺโต วิวฏฺฏตีติ ญาณวิวฏฺโฏ, โวสชฺชนฺโต วิวฏฺฏตีติ วิโมกฺขวิวฏฺโฏ, ตถฏฺเฐ วิวฏฺฏตีติ สจฺจวิวฏฺโฏ.}

\prose{ยตฺถ สญฺญาวิวฏฺโฏ, ตตฺถ เจโตวิวฏฺโฏ.\csromanspacer{} ยตฺถ เจโตวิวฏฺโฏ, ตตฺถ สญฺญาวิวฏฺโฏ.\csromanspacer{} ยตฺถ สญฺญาวิวฏฺโฏ เจโตวิวฏฺโฏ, ตตฺถ จิตฺตวิวฏฺโฏ.\csromanspacer{} ยตฺถ จิตฺตวิวฏฺโฏ, ตตฺถ สญฺญาวิวฏฺโฏ เจโตวิวฏฺโฏ.\csromanspacer{} ยตฺถ สญฺญาวิวฏฺโฏ เจโตวิวฏฺโฏ จิตฺตวิวฏฺโฏ, ตตฺถ ญาณวิวฏฺโฏ.\csromanspacer{} ยตฺถ ญาณวิวฏฺโฏ, ตตฺถ สญฺญาวิวฏฺโฏ เจโตวิวฏฺโฏ จิตฺตวิวฏฺโฏ.\csromanspacer{} ยตฺถ สญฺญาวิวฏฺโฏ เจโตวิวฏฺโฏ จิตฺตวิวฏฺโฏ ญาณวิวฏฺโฏ, ตตฺถ วิโมกฺขวิวฏฺโฏ.\csromanspacer{} ยตฺถ วิโมกฺขวิวฏฺโฏ, ตตฺถ สญฺญาวิวฏฺโฏ เจโตวิวฏฺโฏ จิตฺตวิวฏฺโฏ ญาณวิวฏฺโฏ.\csromanspacer{} ยตฺถ สญฺญาวิวฏฺโฏ เจโตวิวฏฺโฏ จิตฺตวิวฏฺโฏ ญาณวิวฏฺโฏ วิโมกฺขวิวฏฺโฏ, ตตฺถ สจฺจวิวฏฺโฏ.\csromanspacer{} ยตฺถ สจฺจวิวฏฺโฏ, ตตฺถ สญฺญาวิวฏฺโฏ เจโตวิวฏฺโฏ จิตฺตวิวฏฺโฏ ญาณวิวฏฺโฏ วิโมกฺขวิวฏฺโฏ.\csromanspacer{} อยํ วิโมกฺขวิวฏฺโฏ. (6)}

\prose{กตมา \textbf{วิโมกฺขภาวนา?} ปฐมสฺส ฌานสฺส อาเสวนา ภาวนา พหุลีกมฺมํ, ทุติยสฺส ฌานสฺส อาเสวนา ภาวนา พหุลีกมฺมํ, ตติยสฺส ฌานสฺส อาเสวนา ภาวนา พหุลีกมฺมํ, จตุตฺถสฺส ฌานสฺส อาเสวนา ภาวนา พหุลีกมฺมํ, อากาสานญฺจายตนสมาปตฺติยา อาเสวนา ภาวนา พหุลีกมฺมํ, วิญฺญาณญฺจายตนสมาปตฺติยา ฯเปฯ อากิญฺจญฺญายตนสมาปตฺติยา.\csromanspacer{} เนวสญฺญานาสญฺญายตนสมาปตฺติยา อาเสวนา ภาวนา พหุลีกมฺมํ, โสตาปตฺติมคฺคสฺส อาเสวนา ภาวนา พหุลีกมฺมํ, สกทาคามิมคฺคสฺส อาเสวนา ภาวนา พหุลีกมฺมํ, อนาคามิมคฺคสฺส อาเสวนา ภาวนา พหุลีกมฺมํ, อรหตฺตมคฺคสฺส อาเสวนา ภาวนา พหุลีกมฺมํ.\csromanspacer{} อยํ วิโมกฺขภาวนา. (7)}

\prose{กตมา \textbf{วิโมกฺขปฺปฏิปฺปสฺสทฺธิ?} ปฐมสฺส ฌานสฺส ปฏิลาโภ วา วิปาโก วา, ทุติยสฺส ฌานสฺส ปฏิลาโภ วา วิปาโก วา, ตติยสฺส ฌานสฺส ฯเปฯ จตุตฺถสฺส ฌานสฺส.\csromanspacer{} อากาสานญฺจายตนสมาปตฺติยา.\csromanspacer{} วิญฺญาณญฺจายตนสมาปตฺติยา.\csromanspacer{} อากิญฺจญฺญายตนสมาปตฺติยา.\csromanspacer{} เนวสญฺญานาสญฺญายตนสมาปตฺติยา ปฏิลาโภ วา}

\prose{\csromanfolio{266}วิปาโก วา, โสตาปตฺติมคฺคสฺส โสตาปตฺติผลํ, สกทาคามิมคฺคสฺส สกทาคามิผลํ, อนาคามิมคฺคสฺส อนาคามิผลํ, อรหตฺตมคฺคสฺส อรหตฺตผลํ.\csromanspacer{} อยํ วิโมกฺขปฺปฏิปฺปสฺสทฺธีติ. (8)}

\nitthitam{วิโมกฺขกถา นิฏฺฐิตา.}
\csromanheader{2}{ตติยภาณวาโร.}
\csromansectionrule
\csromanmatikaanchor{mtk.79}
\csromantocmarkref{section}{6. คติกถา}{mtk.79}
\bookmark[dest={mtk.79},level=1]{6. คติกถา}
\csromanheader{1}{6. คติกถา}
\proseitem{231}{คติสมฺปตฺติยา ญาณสมฺปยุตฺเต กตินํ เหตูนํ ปจฺจยา อุปปตฺติ โหติ.\csromanspacer{} ขตฺติยมหาสาลานํ พฺราหฺมณมหาสาลานํ คหปติมหาสาลานํ กามาวจรานํ เทวานํ ญาณสมฺปยุตฺเต กตินํ เหตูนํ ปจฺจยา อุปปตฺติ โหติ.\csromanspacer{} รูปาวจรานํ เทวานํ กตินํ เหตูนํ ปจฺจยา อุปปตฺติ โหติ.\csromanspacer{} อรูปาวจรานํ เทวานํ กตินํ เหตูนํ ปจฺจยา อุปปตฺติ โหติ?}

\prose{คติสมฺปตฺติยา ญาณสมฺปยุตฺเต อฏฺฐนฺนํ เหตูนํ ปจฺจยา อุปปตฺติ โหติ.\csromanspacer{} ขตฺติยมหาสาลานํ พฺราหฺมณมหาสาลานํ คหปติมหาสาลานํ กามาวจรานํ เทวานํ ญาณสมฺปยุตฺเต อฏฺฐนฺนํ เหตูนํ ปจฺจยา อุปปตฺติ โหติ.\csromanspacer{} รูปาวจรานํ เทวานํ อฏฺฐนฺนํ เหตูนํ ปจฺจยา อุปปตฺติ โหติ, อรูปาวจรานํ เทวานํ อฏฺฐนฺนํ เหตูนํ ปจฺจยา อุปปตฺติ โหติ.}

\proseitem{232}{คติสมฺปตฺติยา ญาณสมฺปยุตฺเต กตเมสํ อฏฺฐนฺนํ เหตูนํ ปจฺจยา อุปปตฺติ โหติ? กุสลกมฺมสฺส ชวนกฺขเณ ตโย เหตู กุสลา ตสฺมิํ ขเณ ชาตเจตนาย สหชาตปจฺจยา โหนฺติ, เตน วุจฺจติ กุสลมูลปจฺจยาปิ สงฺขารา.\csromanspacer{} นิกนฺติกฺขเณ เทฺว เหตู อกุสลา ตสฺมิํ ขเณ ชาตเจตนาย สหชาตปจฺจยา โหนฺติ, เตน วุจฺจติ อกุสลมูลปจฺจยาปิ สงฺขารา.\csromanspacer{} ปฏิสนฺธิกฺขเณ ตโย เหตู อพฺยากตา ตสฺมิํ ขเณ ชาตเจตนาย สหชาตปจฺจยา โหนฺติ, เตน วุจฺจติ นามรูปปจฺจยาปิ วิญฺญาณํ, วิญฺญาณปจฺจยาปิ นามรูปํ.}

\prose{ปฏิสนฺธิกฺขเณ ปญฺจกฺขนฺธา สหชาตปจฺจยา โหนฺติ, อญฺญมญฺญปจฺจยา โหนฺติ, นิสฺสยปจฺจยา โหนฺติ, วิปฺปยุตฺตปจฺจยา โหนฺติ.\csromanspacer{} ปฏิสนฺธิกฺขเณ}

\csromanheader{2}{6. คติกถา}
\prose{\csromanfolio{267}จตฺตาโร มหาภูตา สหชาตปจฺจยา โหนฺติ, อญฺญมญฺญปจฺจยา โหนฺติ, นิสฺสยปจฺจยา โหนฺติ.\csromanspacer{} ปฏิสนฺธิกฺขเณ ตโย ชีวิตสงฺขารา สหชาตปจฺจยา โหนฺติ, อญฺญมญฺญปจฺจยา โหนฺติ, นิสฺสยปจฺจยา โหนฺติ, วิปฺปยุตฺตปจฺจยา โหนฺติ.\csromanspacer{} ปฏิสนฺธิกฺขเณ นามญฺจ รูปญฺจ สหชาตปจฺจยา โหนฺติ, อญฺญมญฺญปจฺจยา โหนฺติ, นิสฺสยปจฺจยา โหนฺติ, วิปฺปยุตฺตปจฺจยา โหนฺติ.\csromanspacer{} ปฏิสนฺธิกฺขเณ อิเม จุทฺทส ธมฺมา สหชาตปจฺจยา โหนฺติ, อญฺญมญฺญปจฺจยา โหนฺติ, นิสฺสยปจฺจยา โหนฺติ, วิปฺปยุตฺตปจฺจยา โหนฺติ.\csromanspacer{} ปฏิสนฺธิกฺขเณ จตฺตาโร ขนฺธา อรูปิโน สหชาตปจฺจยา โหนฺติ, อญฺญมญฺญปจฺจยา โหนฺติ, นิสฺสยปจฺจยา โหนฺติ, สมฺปยุตฺตปจฺจยา โหนฺติ.\csromanspacer{} ปฏิสนฺธิกฺขเณ ปญฺจินฺทฺริยานิ สหชาตปจฺจยา โหนฺติ, อญฺญมญฺญปจฺจยา โหนฺติ, นิสฺสยปฺปจฺจยา โหนฺติ, สมฺปยุตฺตปจฺจยา โหนฺติ.\csromanspacer{} ปฏิสนฺธิกฺขเณ ตโย เหตู สหชาตปจฺจยา โหนฺติ, อญฺญมญฺญปจฺจยา โหนฺติ, นิสฺสยปจฺจยา โหนฺติ, สมฺปยุตฺตปจฺจยา โหนฺติ.\csromanspacer{} ปฏิสนฺธิกฺขเณ นามญฺจ วิญฺญาณญฺจ สหชาตปจฺจยา โหนฺติ, อญฺญมญฺญปจฺจยา โหนฺติ, นิสฺสยปจฺจยา โหนฺติ, สมฺปยุตฺตปจฺจยา โหนฺติ.\csromanspacer{} ปฏิสนฺธิกฺขเณ อิเม จุทฺทส ธมฺมา สหชาตปจฺจยา โหนฺติ, อญฺญมญฺญปจฺจยา โหนฺติ, นิสฺสยปจฺจยา โหนฺติ, สมฺปยุตฺตปจฺจยา โหนฺติ.\csromanspacer{} ปฏิสนฺธิกฺขเณ อิเม อฏฺฐวีสติ ธมฺมา สหชาตปจฺจยา โหนฺติ, อญฺญมญฺญปจฺจยา โหนฺติ, นิสฺสยปจฺจยา โหนฺติ, วิปฺปยุตฺตปจฺจยา โหนฺติ.\csromanspacer{} คติสมฺปตฺติยา ญาณสมฺปยุตฺเต อิเมสํ อฏฺฐนฺนํ เหตูนฺนํ ปจฺจยา อุปปตฺติ โหติ. (1)}

\prose{ขตฺติยมหาสาลานํ พฺราหฺมณมหาสาลานํ คหปติมหาสาลานํ กามาวจรานํ เทวานํ ญาณสมฺปยุตฺเต กตเมสํ อฏฺฐนฺนํ เหตูนํ ปจฺจยา อุปปตฺติ โหติ? กุสลกมฺมสฺส ชวนกฺขเณ ตโย เหตู กุสลา ตสฺมิํ ขเณ ชาตเจตนาย สหชาตปจฺจยา โหนฺติ, เตน วุจฺจติ กุสลมูลปจฺจยาปิ สงฺขารา.\csromanspacer{} นิกนฺติกฺขเณ เทฺว เหตู อกุสลา ตสฺมิํ ขเณ ชาตเจตนาย สหชาตปจฺจยา โหนฺติ, เตน วุจฺจติ อกุสลมูลปจฺจยาปิ สงฺขารา.\csromanspacer{} ปฏิสนฺธิกฺขเณ ตโย เหตู อพฺยากตา ตสฺมิํ ขเณ ชาตเจตนาย สหชาตปจฺจยา โหนฺติ, เตน วุจฺจติ นามรูปปจฺจยาปิ วิญฺญาณํ, วิญฺญาณปจฺจยาปิ นามรูปํ.}

\prose{\csromanfolio{268}ปฏิสนฺธิกฺขเณ ปญฺจกฺขนฺธา สหชาตปจฺจยา โหนฺติ, อญฺญมญฺญปจฺจยา โหนฺติ, นิสฺสยปจฺจยา โหนฺติ, วิปฺปยุตฺตปจฺจยา โหนฺติ, ปฏิสนฺธิกฺขเณ จตฺตาโร มหาภูตา สหชาตปจฺจยา โหนฺติ.\csromanspacer{} อญฺญมญฺญปจฺจยา โหนฺติ, นิสฺสยปจฺจยา โหนฺติ.\csromanspacer{} ปฏิสนฺธิกฺขเณ ตโย ชีวิตสงฺขารา สหชาตปจฺจยา โหนฺติ, อญฺญมญฺญปจฺจยา โหนฺติ, นิสฺสยปจฺจยา โหนฺติ, วิปฺปยุตฺตปจฺจยา โหนฺติ.\csromanspacer{} ปฏิสนฺธิกฺขเณ นามญฺจ รูปญฺจ สหชาตปจฺจยา โหนฺติ, อญฺญมญฺญปจฺจยา โหนฺติ, นิสฺสยปจฺจยา โหนฺติ, วิปฺปยุตฺตปจฺจยา โหนฺติ.\csromanspacer{} ปฏิสนฺธิกฺขเณ อิเม จุทฺทส ธมฺมา สหชาตปจฺจยา โหนฺติ, อญฺญมญฺญปจฺจยา โหนฺติ, นิสฺสยปจฺจยา โหนฺติ, วิปฺปยุตฺตปจฺจยา โหนฺติ.\csromanspacer{} ปฏิสนฺธิกฺขเณ จตฺตาโร ขนฺธา อรูปิโน สหชาตปจฺจยา โหนฺติ, อญฺญมญฺญปจฺจยา โหนฺติ, นิสฺสยปจฺจยา โหนฺติ, สมฺปยุตฺตปจฺจยา โหนฺติ.\csromanspacer{} ปฏิสนฺธิกฺขเณ ปญฺจินฺทฺริยานิ สหชาตปจฺจยา โหนฺติ, อญฺญมญฺญปจฺจยา โหนฺติ, นิสฺสยปจฺจยา โหนฺติ, สมฺปยุตฺตปจฺจยา โหนฺติ.\csromanspacer{} ปฏิสนฺธิกฺขเณ ตโย เหตู สหชาตปจฺจยา โหนฺติ, อญฺญมญฺญปจฺจยา โหนฺติ, นิสฺสยปจฺจยา โหนฺติ, สมฺปยุตฺตปจฺจยา โหนฺติ.\csromanspacer{} ปฏิสนฺธิกฺขเณ นามญฺจ วิญฺญาณญฺจ สหชาตปจฺจยา โหนฺติ, อญฺญมญฺญปจฺจยา โหนฺติ, นิสฺสยปจฺจยา โหนฺติ, สมฺปยุตฺตปจฺจยา โหนฺติ.\csromanspacer{} ปฏิสนฺธิกฺขเณ อิเม จุทฺทส ธมฺมา สหชาตปจฺจยา โหนฺติ, อญฺญมญฺญปจฺจยา โหนฺติ, นิสฺสยปจฺจยา โหนฺติ, สมฺปยุตฺตปจฺจยา โหนฺติ.\csromanspacer{} ปฏิสนฺธิกฺขเณ อิเม อฏฺฐวีสติ ธมฺมา สหชาตปจฺจยา โหนฺติ, อญฺญมญฺญปจฺจยา โหนฺติ, นิสฺสยปจฺจยา โหนฺติ, วิปฺปยุตฺตปจฺจยา โหนฺติ, ขตฺติยมหาสาลานํ พฺรหฺมณมหาสาลานํ คหปติมหาสาลานํ กามาวจรานํ เทวานํ ญาณสมฺปยุตฺเต อิเมสํ อฏฺฐนฺนํ เหตูนํ ปจฺจยา อุปปตฺติ โหติ. (2)}

\prose{รูปาวจรานํ เทวานํ กตเมสํ อฏฺฐนฺนํ เหตูนํ ปจฺจยา อุปปตฺติ โหติ? กุสลกมฺมสฺส ชวนกฺขเณ ตโย เหตู กุสลา ฯเปฯ รูปาวจรานํ เทวานํ อิเมสํ อฏฺฐนฺนํ เหตูนํ ปจฺจยา อุปปตฺติ โหติ. (3)}

\prose{อรูปาวจรานํ เทวานํ กตเมสํ อฏฺฐนฺนํ เหตูนํ ปจฺจยา อุปปตฺติ โหติ? กุสลกมฺมสฺส ชวนกฺขเณ ตโย เหตู กุสลา ตสฺมิํ ขเณ ชาตเจตนาย สหชาตปจฺจยา โหนฺติ, เตน วุจฺจติ กุสลมูลปจฺจยาปิ สงฺขารา.\csromanspacer{} นิกนฺติกฺขเณ เทฺว เหตู อกุสลา ตสฺมิํ ขเณ}

\csromanheader{2}{6. คติกถา}
\prose{\csromanfolio{269}ชาตเจตนาย สหชาตปจฺจยา โหนฺติ, เตน วุจฺจติ อกุสลมูลปจฺจยาปิ สงฺขารา.\csromanspacer{} ปฏิสนฺธิกฺขเณ ตโย เหตู อพฺยากตา ตสฺมิํ ขเณ ชาตเจตนาย สหชาตปจฺจยา โหนฺติ, เตน วุจฺจติ นามปจฺจยาปิ วิญฺญาณํ วิญฺญาณปจฺจยาปิ นามํ.}

\prose{ปฏิสนฺธิกฺขเณ จตฺตาโร ขนฺธา อรูปิโน สหชาตปจฺจยา โหนฺติ, อญฺญมญฺญปจฺจยา โหนฺติ, นิสฺสยปจฺจยา โหนฺติ, สมฺปยุตฺตปจฺจยา โหนฺติ.\csromanspacer{} ปฏิสนฺธิกฺขเณ ปญฺจินฺทฺริยานิ สหชาตปจฺจยา โหนฺติ, อญฺญมญฺญปจฺจยา โหนฺติ, นิสฺสยปจฺจยา โหนฺติ, สมฺปยุตฺตปจฺจยา โหนฺติ.\csromanspacer{} ปฏิสนฺธิกฺขเณ ตโย เหตู สหชาตปจฺจยา โหนฺติ, อญฺญมญฺญปจฺจยา โหนฺติ, นิสฺสยปจฺจยา โหนฺติ, สมฺปยุตฺตปจฺจยา โหนฺติ.\csromanspacer{} ปฏิสนฺธิกฺขเณ นามญฺจ วิญฺญาณญฺจ สหชาตปจฺจยา โหนฺติ, อญฺญมญฺญปจฺจยา โหนฺติ, นิสฺสยปจฺจยา โหนฺติ, สมฺปยุตฺตปจฺจยา โหนฺติ.\csromanspacer{} ปฏิสนฺธิกฺขเณ อิเม จุทฺทส ธมฺมา สหชาตปจฺจยา โหนฺติ, อญฺญมญฺญปจฺจยา โหนฺติ, นิสฺสยปจฺจยา โหนฺติ, สมฺปยุตฺตปจฺจยา โหนฺติ, อรูปาวจรานํ เทวานํ อิเมสํ อฏฺฐนฺนํ เหตูนํ ปจฺจยา อุปปตฺติ โหติ. (4)}

\proseitem{233}{คติสมฺปตฺติยา ญาณวิปฺปยุตฺเต กตินํ เหตูนํ ปจฺจยา อุปปตฺติ โหติ.\csromanspacer{} ขตฺติยมหาสาลานํ พฺราหฺมณมหาสาลานํ คหปติมหาสาลานํ กามาวจรานํ เทวานํ ญาณวิปฺปยุตฺเต กตินํ เหตูนํ ปจฺจยา อุปปตฺติ โหติ?}

\prose{คติสมฺปตฺติยา ญาณวิปฺปยุตฺเต ฉนฺนํ เหตูนํ ปจฺจยา อุปปตฺติ โหติ.\csromanspacer{} ขตฺติยมหาสาลานํ พฺราหฺมณมหาสาลานํ คหปติมหาสาลานํ กามาวจรานํ เทวานํ ญาณวิปฺปยุตฺเต ฉนฺนํ เหตูนํ ปจฺจยา อุปปตฺติ โหติ.}

\prose{คติสมฺปตฺติยา ญาณวิปฺปยุตฺเต กตเมสํ ฉนฺนํ เหตูนํ ปจฺจยา อุปปตฺติ โหติ? กุสลกมฺมสฺส ชวนกฺขเณ เทฺว เหตู กุสลา ตสฺมิํ ขเณ ชาตเจตนาย สหชาตปจฺจยา โหนฺติ, เตน วุจฺจติ กุสลมูลปจฺจยาปิ สงฺขารา.\csromanspacer{} นิกนฺติกฺขเณ เทฺว เหตู อกุสลา ตสฺมิํ ขเณ ชาตเจตนาย สหชาตปจฺจยา โหนฺติ, เตน วุจฺจติ อกุสลมูลปจฺจยาปิ สงฺขารา.\csromanspacer{} ปฏิสนฺธิกฺขเณ เทฺว เหตู อพฺยากตา ตสฺมิํ ขเณ \csromanfolio{270}ชาตเจตนาย สหชาตปจฺจยา โหนฺติ, เตน วุจฺจติ นามรูปปจฺจยาปิ วิญฺญาณํ, วิญฺญาณปจฺจยาปิ นามรูปํ.}

\prose{ปฏิสนฺธิกฺขเณ ปญฺจกฺขนฺธา สหชาตปจฺจยา โหนฺติ, อญฺญมญฺญปจฺจยา โหนฺติ, นิสฺสยปจฺจยา โหนฺติ, วิปฺปยุตฺตปจฺจยา โหนฺติ.\csromanspacer{} ปฏิสนฺธิกฺขเณ จตฺตาโร มหาภูตา สหชาตปจฺจยา โหนฺติ, อญฺญมญฺญปจฺจยา โหนฺติ, นิสฺสยปจฺจยา โหนฺติ.\csromanspacer{} ปฏิสนฺธิกฺขเณ ตโย ชีวิตสงฺขารา สหชาตปจฺจยา โหนฺติ, อญฺญมญฺญปจฺจยา โหนฺติ, นิสฺสยปจฺจยา โหนฺติ, วิปฺปยุตฺตปจฺจยา โหนฺติ.\csromanspacer{} ปฏิสนฺธิกฺขเณ นามญฺจ รูปญฺจ สหชาตปจฺจยา โหนฺติ, อญฺญมญฺญปจฺจยา โหนฺติ, นิสฺสยปจฺจยา โหนฺติ, วิปฺปยุตฺตปจฺจยา โหนฺติ.\csromanspacer{} ปฏิสนฺธิกฺขเณ อิเม จุทฺทส ธมฺมา สหชาตปจฺจยา โหนฺติ, อญฺญมญฺญปจฺจยา โหนฺติ, นิสฺสยปจฺจยา โหนฺติ, วิปฺปยุตฺตปจฺจยา โหนฺติ.\csromanspacer{} ปฏิสนฺธิกฺขเณ จตฺตาโร ขนฺธา อรูปิโน สหชาตปจฺจยา โหนฺติ, อญฺญมญฺญปจฺจยา โหนฺติ, นิสฺสยปจฺจยา โหนฺติ, สมฺปยุตฺตปจฺจยา โหนฺติ.\csromanspacer{} ปฏิสนฺธิกฺขเณ จตฺตาริ อินฺทฺริยานิ สหชาตปจฺจยา โหนฺติ, อญฺญมญฺญปจฺจยา โหนฺติ, นิสฺสยปจฺจยา โหนฺติ, สมฺปยุตฺตปจฺจยา โหนฺติ.\csromanspacer{} ปฏิสนฺธิกฺขเณ เทฺว เหตู สหชาตปจฺจยา โหนฺติ, อญฺญมญฺญปจฺจยา โหนฺติ, นิสฺสยปจฺจยา โหนฺติ, สมฺปยุตฺตปจฺจยา โหนฺติ.\csromanspacer{} ปฏิสนฺธิกฺขเณ นามญฺจ วิญฺญาณญฺจ สหชาตปจฺจยา โหนฺติ, อญฺญมญฺญปจฺจยา โหนฺติ, นิสฺสยปจฺจยา โหนฺติ, สมฺปยุตฺตปจฺจยา โหนฺติ.\csromanspacer{} ปฏิสนฺธิกฺขเณ อิเม ทฺวาทส ธมฺมา สหชาตปจฺจยา โหนฺติ, อญฺญมญฺญปจฺจยา โหนฺติ, นิสฺสยปจฺจยา โหนฺติ, สมฺปยุตฺตปจฺจยา โหนฺติ.\csromanspacer{} ปฏิสนฺธิกฺขเณ อิเม ฉพฺพีสติ ธมฺมา สหชาตปจฺจยา โหนฺติ, อญฺญมญฺญปจฺจยา โหนฺติ, นิสฺสยปจฺจยา โหนฺติ, วิปฺปยุตฺตปจฺจยา โหนฺติ, คติสมฺปตฺติยา ญาณวิปฺปยุตฺเต อิเมสํ ฉนฺนํ เหตูนํ ปจฺจยา อุปปตฺติ โหติ. (1)}

\prose{ขตฺติยมหาสาลานํ พฺราหฺมณมหาสาลานํ คหปติมหาสาลานํ กามาวจรานํ เทวานํ ญาณวิปฺปยุตฺเต กตเมสํ ฉนฺนํ เหตูนํ ปจฺจยา อุปปตฺติ โหติ? กุสลกมฺมสฺส ชวนกฺขเณ เทฺว เหตู กุสลา ตสฺมิํ ขเณ ชาตเจตนาย สหชาตปจฺจยา โหนฺติ, เตน วุจฺจติ กุสลมูลปจฺจยาปิ สงฺขารา ฯเปฯ ขตฺติยมหาสาลานํ พฺราหฺมณมหาสาลานํ}

\csromanmatikaanchor{mtk.80}
\csromantocmarkref{section}{7. กมฺมกถา}{mtk.80}
\bookmark[dest={mtk.80},level=1]{7. กมฺมกถา}
\csromanheader{1}{7. กมฺมกถา}
\prose{\csromanfolio{271}คหปติมหาสาลานํ กามาวจรานํ เทวานํ ญาณวิปฺปยุตฺเต อิเมสํ ฉนฺนํ เหตูนํ ปจฺจยา อุปปตฺติ โหตีติ. (2)}

\nitthitamruled{คติกถา นิฏฺฐิตา.}
\csromanheader{2}{7. กมฺมกถา}
\proseitem{234}{อโหสิ กมฺมํ อโหสิ กมฺมวิปาโก.\csromanspacer{} อโหสิ กมฺมํ นาโหสิ กมฺมวิปาโก.\csromanspacer{} อโหสิ กมฺมํ อตฺถิ กมฺมวิปาโก.\csromanspacer{} อโหสิ กมฺมํ นตฺถิ กมฺมวิปาโก.\csromanspacer{} อโหสิ กมฺมํ ภวิสฺสติ กมฺมวิปาโก.\csromanspacer{} อโหสิ กมฺมํ น ภวิสฺสติ กมฺมวิปาโก. (อตีตกมฺมํ) (6)}

\prose{อตฺถิ กมฺมํ อตฺถิ กมฺมวิปาโก.\csromanspacer{} อตฺถิ กมฺมํ นตฺถิ กมฺมวิปาโก.\csromanspacer{} อตฺถิ กมฺมํ ภวิสฺสติ กมฺมวิปาโก.\csromanspacer{} อตฺถิ กมฺมํ น ภวิสฺสติ กมฺมวิปาโก. (ปจฺจุปฺปนฺนกมฺมํ) (4, 10)}

\prose{ภวิสฺสติ กมฺมํ ภวิสฺสติ กมฺมวิปาโก.\csromanspacer{} ภวิสฺสติ กมฺมํ น ภวิสฺสติ กมฺมวิปาโก. (อนาคตกมฺมํ) (2, 12)}

\proseitem{235}{อโหติ กุสลํ กมฺมํ อโหสิ กุสลสฺส กมฺมสฺส วิปาโก, อโหสิ กุสลํ กมฺมํ นาโหสิ กุสลสฺส กมฺมสฺส วิปาโก, อโหสิ กุสลํ กมฺมํ อตฺถิ กุสลสฺส กมฺมสฺส วิปาโก, อโหสิ กุสลํ กมฺมํ นตฺถิ กุสลสฺส กมฺมสฺส วิปาโก, อโหสิ กุสลํ กมฺมํ ภวิสฺสติ กุสลสฺส กมฺมสฺส วิปาโก, อโหสิ กุสลํ กมฺมํ น ภวิสฺสติ กุสลสฺส กมฺมสฺส วิปาโก.}

\prose{อตฺถิ กุสลํ กมฺมํ อตฺถิ กุสลสฺส กมฺมสฺส วิปาโก, อตฺถิ กุสลํ กมฺมํ นตฺถิ กุสลสฺส กมฺมสฺส วิปาโก, อตฺถิ กุสลํ กมฺมํ ภวิสฺสติ กุสลสฺส กมฺมสฺส วิปาโก, อตฺถิ กุสลํ กมฺมํ น ภวิสฺสติ กุสลสฺส กมฺมสฺส วิปาโก.}

\prose{ภวิสฺสติ กุสลํ กมฺมํ ภวิสฺสติ กุสลสฺส กมฺมสฺส วิปาโก, ภวิสฺสติ กุสลํ กมฺมํ น ภวิสฺสติ กุสลสฺส กมฺมสฺส วิปาโก.}

\prose{อโหสิ อกุสลํ กมฺมํ อโหสิ อกุสลสฺส กมฺมสฺส วิปาโก, อโหสิ อกุสลํ กมฺมํ นาโหสิ อกุสลสฺส กมฺมสฺส วิปาโก, \csromanfolio{272}อโหสิ อกุสลํ กมฺมํ อตฺถิ อกุสลสฺส กมฺมสฺส วิปาโก, อโหสิ อกุสลํ กมฺมํ นตฺถิ อกุสลสฺส กมฺมสฺส วิปาโก.\csromanspacer{} อโหสิ อกุสลํ กมฺมํ ภวิสฺสติ อกุสลสฺส กมฺมสฺส วิปาโก, อโหสิ อกุสลํ กมฺมํ น ภวิสฺสติ อกุสลสฺส กมฺมสฺส วิปาโก.}

\prose{อตฺถิ อกุสลํ กมฺมํ อตฺถิ อกุสลสฺส กมฺมสฺส วิปาโก, อตฺถิ อกุสลํ กมฺมํ นตฺถิ อกุสลสฺส กมฺมสฺส วิปาโก, อตฺถิ อกุสลํ กมฺมํ ภวิสฺสติ อกุสลสฺส กมฺมสฺส วิปาโก, อตฺถิ อกุสลํ กมฺมํ น ภวิสฺสติ อกุสลสฺส กมฺมสฺส วิปาโก.}

\prose{ภวิสฺสติ อกุสลํ กมฺมํ ภวิสฺสติ อกุสลสฺส กมฺมสฺส วิปาโก, ภวิสฺสติ อกุสลํ กมฺมํ น ภวิสฺสติ อกุสลสฺส กมฺมสฺส วิปาโก.}

\prose{อโหสิ สาวชฺชํ กมฺมํ ฯเปฯ.\csromanspacer{} อโหสิ อนวชฺชํ กมฺมํ ฯเปฯ.\csromanspacer{} อโหสิ กณฺหํ กมฺมํ ฯเปฯ.\csromanspacer{} อโหสิ สุกฺกํ กมฺมํ ฯเปฯ.\csromanspacer{} อโหสิ สุขุทฺรยํ กมฺมํ ฯเปฯ.\csromanspacer{} อโหสิ ทุกฺขุทฺรยํ กมฺมํ ฯเปฯ.\csromanspacer{} อโหสิ สุขวิปากํ กมฺมํ ฯเปฯ.\csromanspacer{} อโหสิ ทุกฺขวิปากํ กมฺมํ อโหสิ ทุกฺขวิปากสฺส กมฺมสฺส วิปาโก, อโหสิ ทุกฺขวิปากํ กมฺมํ นาโหสิ ทุกฺขวิปากสฺส กมฺมสฺส วิปาโก, อโหสิ ทุกฺขวิปากํ กมฺมํ อตฺถิ ทุกฺขวิปากสฺส กมฺมสฺส วิปาโก, อโหสิ ทุกฺขวิปากํ กมฺมํ นตฺถิ ทุกฺขวิปากสฺส กมฺมสฺส วิปาโก, อโหสิ ทุกฺขวิปากํ กมฺมํ ภวิสฺสติ ทุกฺขวิปากสฺส กมฺมสฺส วิปาโก, อโหสิ ทุกฺขวิปากํ กมฺมํ น ภวิสฺสติ ทุกฺขวิปากสฺส กมฺมสฺส วิปาโก.}

\prose{อตฺถิ ทุกฺขวิปากํ กมฺมํ อตฺถิ ทุกฺขวิปากสฺส กมฺมสฺส วิปาโก, อตฺถิ ทุกฺขวิปากํ กมฺมํ นตฺถิ ทุกฺขวิปากสฺส กมฺมสฺส วิปาโก, อตฺถิ ทุกฺข วิปากํ กมฺมํ ภวิสฺสติ ทุกฺขวิปากสฺส กมฺมสฺส วิปาโก, อตฺถิ ทุกฺขวิปากํ กมฺมํ น ภวิสฺสติ ทุกฺขวิปากสฺส กมฺมสฺส วิปาโก.}

\prose{ภวิสฺสติ ทุกฺขวิปากํ กมฺมํ ภวิสฺสติ ทุกฺขวิปากสฺส กมฺมสฺส วิปาโก, ภวิสฺสติ ทุกฺขวิปากํ กมฺมํ น ภวิสฺสติ ทุกฺขวิปากสฺส กมฺมสฺส วิปาโกติ.}

\nitthitam{กมฺมกถา นิฏฺฐิตา.}
\csromanmatikaanchor{mtk.81}
\csromantocmarkref{section}{8. วิปลฺลาสกถา}{mtk.81}
\bookmark[dest={mtk.81},level=1]{8. วิปลฺลาสกถา}
\csromanheader{1}{8. วิปลฺลาสกถา}
\csromanheader{2}{8. วิปลฺลาสกถา}
\proseitem{236}{\csromanfolio{273}ปุริมนิทานํ.\csromanspacer{}\csromanspacer{}จตฺตาโรเม ภิกฺขเว\footnote{อํ 1. 361 ปิฏฺเฐ ปสฺสิตพฺพา.} สญฺญาวิปลฺลาสา จิตฺตวิปลฺลาสา ทิฏฺฐิวิปลฺลาสา.\csromanspacer{} กตเม จตฺตาโร? อนิจฺเจ ภิกฺขเว “นิจฺจนฺ”ติ สญฺญาวิปลฺลาโส จิตฺตวิปลฺลาโส ทิฏฺฐิวิปลฺลาโส, ทุกฺเข ภิกฺขเว “สุขนฺ”ติ สญฺญาวิปลฺลาโส จิตฺตวิปลฺลาโส ทิฏฺฐิวิปลฺลาโส, อนตฺตนิ ภิกฺขเว “อตฺตา”ติ สญฺญาวิปลฺลาโส จิตฺตวิปลฺลาโส ทิฏฺฐิวิปลฺลาโส, อสุเภ ภิกฺขเว “สุภนฺ”ติ สญฺญาวิปลฺลาโส จิตฺตวิปฺปลฺลาโส ทิฏฺฐิวิปลฺลาโส, อิเม โข ภิกฺขเว จตฺตาโร สญฺญาวิปลฺลาสา จิตฺตวิปลฺลาสา ทิฏฺฐิวิปลฺลาสา.}

\prose{จตฺตาโรเม ภิกฺขเว นสญฺญาวิปลฺลาสา นจิตฺตวิปลฺลาสา นทิฏฺฐิปลฺลาสา.\csromanspacer{} กตเม จตฺตาโร? อนิจฺเจ ภิกฺขเว “อนิจฺจนฺ”ติ นสญฺญาวิปลฺลาโส นจิตฺตวิปลฺลาโส นทิฏฺฐิวิปลฺลาโส, ทุกฺเข ภิกฺขเว “ทุกฺขนฺ”ติ นสญฺญาวิปลฺลาโส นจิตฺตวิปลฺลาโส นทิฏฺฐิวิปลฺลาโส, อนตฺตนิ ภิกฺขเว “อนตฺตา”ติ นสญฺญาวิปลฺลาโส นจิตฺตวิปลฺลาโส นทิฏฺฐิวิปลฺลาโส, อสุเภ ภิกฺขเว “อสุภนฺ”ติ นสญฺญาวิปลฺลาโส นจิตฺตวิปลฺลาโส นทิฏฺฐิวิปลฺลาโส, อิเม โข ภิกฺขเว จตฺตโร นสญฺญาวิปลฺลาสา นจิตฺตวิปลฺลาสา นทิฏฺฐิวิปลฺลาสาติ.}

\csromangathameasure{อนิจฺเจ นิจฺจสญฺญิโน, ทุกฺเข จ สุขสญฺญิโน. \\ อนตฺตนิ จ “อตฺตา”ติ, อสุเภ สุภสญฺญิโน. \\ มิจฺฉาทิฏฺฐิหตา สตฺตา, ขิตฺตจิตฺตา วิสญฺญิโน. \\ เต โยคยุตฺตา มารสฺส, อโยคกฺเขมิโน ชนา. \\ สตฺตา คจฺฉนฺติ สํสารํ, ชาติมรณคามิโน. \\ ยทา จ พุทฺธา โลกสฺมิํ, อุปฺปชฺชนฺติ ปภงฺกรา. \\ เต อิมํ ธมฺมํ ปกาเสนฺติ, ทุกฺขูปสมคามินํ. \\ เตสํ สุตฺวาน สปฺปญฺญา, สจิตฺตํ ปจฺจลทฺธุ3 เต. \\ อนิจฺจํ อนิจฺจโต ทกฺขุํ, ทุกฺขมทฺทกฺขุ ทุกฺขโต. \\ อนตฺตนิ “อนตฺตา”ติ, อสุภํ อสุภตทฺทสุํ. \\ สมฺมาทิฏฺฐิสมาทานา, สพฺพํ ทุกฺขํ อุปจฺจคุนฺติ.}
\csromangathagroup{\csromangathastanza{อนิจฺเจ นิจฺจสญฺญิโน, ทุกฺเข จ สุขสญฺญิโน. \\ อนตฺตนิ จ “อตฺตา”ติ\footnote{อตฺตสญฺญิโน (สฺยา) อํ 1. 362 ปิฏฺเฐ ปสฺสิตพฺพา. 3. ปจฺจลทฺธุํ (สฺยา), ปจฺจลทฺธา (อํ 1. 362 ปิฏฺเฐ.)}, อสุเภ สุภสญฺญิโน.}\csromangathastanza{มิจฺฉาทิฏฺฐิหตา สตฺตา, ขิตฺตจิตฺตา วิสญฺญิโน. \\ เต โยคยุตฺตา มารสฺส, อโยคกฺเขมิโน ชนา.}\csromangathastanza{สตฺตา คจฺฉนฺติ สํสารํ, ชาติมรณคามิโน. \\ ยทา จ พุทฺธา โลกสฺมิํ, อุปฺปชฺชนฺติ ปภงฺกรา.}\csromangathastanza{เต อิมํ ธมฺมํ ปกาเสนฺติ, ทุกฺขูปสมคามินํ. \\ เตสํ สุตฺวาน สปฺปญฺญา, สจิตฺตํ ปจฺจลทฺธุ3 เต.}\csromangathastanza{อนิจฺจํ อนิจฺจโต ทกฺขุํ, ทุกฺขมทฺทกฺขุ ทุกฺขโต. \\ อนตฺตนิ “อนตฺตา”ติ, อสุภํ อสุภตทฺทสุํ. \\ สมฺมาทิฏฺฐิสมาทานา, สพฺพํ ทุกฺขํ อุปจฺจคุนฺติ.}}

\prose{\csromanfolio{274}อิเม จตฺตาโร วิปลฺลาสา ทิฏฺฐิสมฺปนฺนสฺส ปุคฺคลสฺส ปหีนา อปฺปหีนาติ? เกจิ ปหีนา เกจิ อปฺปหีนา.\csromanspacer{} อนิจฺเจ “นิจฺจนฺ”ติ สญฺญาวิปลฺลาโส จิตฺตวิปลฺลาโส ทิฏฺฐิวิปลฺลาโส ปหีโน.\csromanspacer{} ทุกฺเข “สุขนฺ”ติ สญฺญา อุปฺปชฺชติ, จิตฺตํ อุปฺปชฺชติ, ทิฏฺฐิวิปลฺลาโส ปหีโน.\csromanspacer{} อนตฺตนิ “อตฺตา”ติ สญฺญาวิปลฺลาโส จิตฺตวิปลฺลาโส ทิฏฺฐิวิปลฺลาโส ปหีโน, อสุเภ “สุภนฺ”ติ สญฺญา อุปฺปชฺชติ, จิตฺตํ อุปฺปชฺชติ, ทิฏฺฐิวิปลฺลาโส ปหีโน.\csromanspacer{} ทฺวีสุ วตฺถูสุ ฉ วิปลฺลาสา ปหีนา, ทฺวีสุ วตฺถูสุ เทฺว วิปลฺลาสา ปหีนา, จตฺตาโร วิปลฺลาสา อปฺปหีนา, จตูสุ วตฺถูสุ อฏฺฐ วิปลฺลาสา ปหีนา, จตฺตาโร วิปลฺลาสา อปฺปหีนาติ.}

\nitthitamruled{วิปลฺลาสกถา นิฏฺฐิตา.}
\csromanmatikaanchor{mtk.82}
\csromantocmarkref{section}{9. มคฺคกถา}{mtk.82}
\bookmark[dest={mtk.82},level=1]{9. มคฺคกถา}
\csromanheader{1}{9. มคฺคกถา}
\proseitem{237}{มฺ\textbf{อคฺโค}ติ เกนฏฺเฐน มฺ\textbf{อคฺโค}? โสตาปตฺติมคฺคกฺขเณ ทสฺสนฏฺเฐน สมฺมาทิฏฺฐิ มิจฺฉาทิฏฺฐิยา ปหานาย มฺ\textbf{อคฺโค} เจว เหตุ จ, สหชาตานํ ธมฺมานํ อุปตฺถมฺภนาย มฺ\textbf{อคฺโค} เจว เหตุ จ, กิเลสานํ ปริยาทานาย มฺ\textbf{อคฺโค} เจว เหตุ จ, ปฏิเวธาทิวิโสธนาย มฺ\textbf{อคฺโค} เจว เหตุ จ, จิตฺตสฺส อธิฏฺฐานาย มฺ\textbf{อคฺโค} เจว เหตุ จ, จิตฺตสฺส โวทานาย มฺ\textbf{อคฺโค} เจว เหตุ จ, วิเสสาธิคมาย มฺ\textbf{อคฺโค} เจว เหตุ จ, อุตฺตริปฏิเวธาย\footnote{อุตฺตริํ ปฏิเวธาย (สฺยา)} มฺ\textbf{อคฺโค} เจว เหตุ จ, สจฺจาภิสมยาย มฺ\textbf{อคฺโค} เจว เหตุ จ, นิโรเธ ปติฏฺฐาปนาย มฺ\textbf{อคฺโค} เจว เหตุ จ.}

\prose{อภินิโรปนฏฺเฐน สมฺมาสงฺกปฺโป มิจฺฉาสงฺกปฺปสฺส ปหานาย มฺ\textbf{อคฺโค} เจว เหตุ จ, สหชาตานํ ธมฺมานํ อุปตฺถมฺภนาย มฺ\textbf{อคฺโค} เจว เหตุ จ, กิเลสานํ ปริยาทานาย มฺ\textbf{อคฺโค} เจว เหตุ จ, ปฏิเวธาทิวิโสธนาย มฺ\textbf{อคฺโค} เจว เหตุ จ, จิตฺตสฺส อธิฏฺฐานาย มฺ\textbf{อคฺโค} เจว เหตุ จ, จิตฺตสฺส โวทานาย มฺ\textbf{อคฺโค} เจว เหตุ จ, วิเสสาธิคมาย มฺ\textbf{อคฺโค} เจว เหตุ จ,}

\csromanheader{2}{9. มคฺคกถา}
\prose{\csromanfolio{275}อุตฺตริปฏิเวธาย มคฺโค เจว เหตุ จ, สจฺจาภิสมยาย มคฺโค เจว เหตุ จ, นิโรเธ ปติฏฺฐาปนาย มคฺโค เจว เหตุ จ.}

\prose{ปริคฺคหฏฺเฐน สมฺมาวาจา มิจฺฉาวาจาย ปหานาย มคฺโค เจว เหตุ จ, สหชาตานํ ธมฺมานํ อุปตฺถมฺภนาย มคฺโค เจว เหตุ จ, กิเลสานํ ปริยาทานาย มคฺโค เจว เหตุ จ, ปฏิเวธาทิวิโสธนาย มคฺโค เจว เหตุ จ, จิตฺตสฺส อธิฏฺฐานาย มคฺโค เจว เหตุ จ, จิตฺตสฺส โวทานาย มคฺโค เจว เหตุ จ, วิเสสาธิคมาย มคฺโค เจว เหตุ จ, อุตฺตริปฏิเวธาย มคฺโค เจว เหตุ จ, สจฺจาภิสมยาย มคฺโค เจว เหตุ จ, นิโรเธ ปติฏฺฐาปนาย มคฺโค เจว เหตุ จ.}

\prose{สมุฏฺฐานฏฺเฐน สมฺมากมฺมนฺโต มิจฺฉากมฺมนฺตสฺส ปหานาย มคฺโค เจว เหตุ จ, สหชาตานํ ธมฺมานํ อุปตฺถมฺภนาย มคฺโค เจว เหตุ จ, กิเลสานํ ปริยาทานาย มคฺโค เจว เหตุ จ, ปฏิเวธาทิวิโสธนาย มคฺโค เจว เหตุ จ, จิตฺตสฺส อธิฏฺฐานาย มคฺโค เจว เหตุ จ, จิตฺตสฺส โวทานาย มคฺโค เจว เหตุ จ.\csromanspacer{} วิเสสาธิคมาย มคฺโค เจว เหตุ จ, อุตฺตริปฏิเวธาย มคฺโค เจว เหตุ จ, สจฺจาภิสมยาย มคฺโค เจว เหตุ จ, นิโรเธ ปติฏฺฐาปนาย มคฺโค เจว เหตุ จ.}

\prose{โวทานฏฺเฐน สมฺมา-อาชีโว มิจฺฉา-อาชีวสฺส ปหานาย มคฺโค เจว เหตุ จ ฯเปฯ.\csromanspacer{} ปคฺคหฏฺเฐน สมฺมาวายาโม มิจฺฉาวายามสฺส ปหานาย มคฺโค เจว เหตุ จ ฯเปฯ.\csromanspacer{} อุปฏฺฐานฏฺเฐน สมฺมาสติ มิจฺฉาสติยา ปหานาย มคฺโค เจว เหตุ จ ฯเปฯ.\csromanspacer{} อวิกฺเขปฏฺเฐน สมฺมาสมาธิ มิจฺฉาสมาธิสฺส ปหานาย มคฺโค เจว เหตุ จ, สหชาตานํ ธมฺมานํ อุปตฺถมฺภนาย มคฺโค เจว เหตุ จ, กิเลสานํ ปริยาทานาย มคฺโค เจว เหตุ จ, ปฏิเวธาทิวิโสธนาย มคฺโค เจว เหตุ จ, จิตฺตสฺส อธิฏฺฐานาย มคฺโค เจว เหตุ จ, จิตฺตสฺส โวทานาย มคฺโค เจว เหตุ จ, วิเสสาธิคมาย มคฺโค เจว เหตุ จ, อุตฺตริปฏิเวธาย มคฺโค เจว เหตุ จ สจฺจาภิสมยาย มคฺโค เจว เหตุ จ, นิโรเธ ปติฏฺฐาปนาย มคฺโค เจว เหตุ จ.}

\prose{สกทาคามิมคฺคกฺขเณ ทสฺสนฏฺเฐน สมฺมาทิฏฺฐิ ฯเปฯ.\csromanspacer{} อวิกฺเขปฏฺเฐน สมฺมาสมาธิ โอฬาริกสฺส กามราคสญฺโญชนสฺส ปฏิฆสญฺโญชนสฺส โอฬาริกสฺส กามราคานุสยสฺส ปฏิฆานุสยสฺส ปหานาย มคฺโค}

\prose{\csromanfolio{276}เจว เหตุ จ, สหชาตานํ ธมฺมานํ อุปตฺถมฺภนาย มคฺโค เจว เหตุ จ, กิเลสานํ ปริยาทานาย มคฺโค เจว เหตุ จ, ปฏิเวธาทิวิโสธนาย มคฺโค เจว เหตุ จ, จิตฺตสฺส อธิฏฺฐานาย มคฺโค เจว เหตุ จ, จิตฺตสฺส โวทานาย มคฺโค เจว เหตุ จ, วิเสสาธิคมาย มคฺโค เจว เหตุ จ, อุตฺตริปฏิเวธาย มคฺโค เจว เหตุ จ, สจฺจาภิสมยาย มคฺโค เจว เหตุ จ, นิโรเธ ปติฏฺฐาปนาย มคฺโค เจว เหตุ จ.}

\prose{อนาคามิมคฺคกฺขเณ ทสฺสนฏฺเฐน สมฺมาทิฏฺฐิ ฯเปฯ.\csromanspacer{} อวิกฺเขปฏฺเฐน สมฺมาสมาธิ อณุสหคตสฺส กามราคสญฺโญชนสฺส ปฏิฆสญฺโญชนสฺส, อณุสหคตสฺส กามราคานุสยสฺส ปฏิฆานุสยสฺส ปหานาย มคฺโค เจว เหตุ จ, สหชาตานํ ธมฺมานํ อุปตฺถมฺภนาย มคฺโค เจว เหตุ จ, กิเลสานํ ปริยาทานาย มคฺโค เจว เหตุ จ, ปฏิเวธาทิวิโสธนาย มคฺโค เจว เหตุ จ, จิตฺตสฺส อธิฏฺฐานาย มคฺโค เจว เหตุ จ, จิตฺตสฺส โวทานาย มคฺโค เจว เหตุ จ, วิเสสาธิคมาย มคฺโค เจว เหตุ จ, อุตฺตริปฏิเวธาย มคฺโค เจว เหตุ จ, สจฺจาภิสมยาย มคฺโค เจว เหตุ จ, นิโรเธ ปติฏฺฐาปนาย มคฺโค เจว เหตุ จ.}

\prose{อรหตฺตมคฺคกฺขเณ ทสฺสนฏฺเฐน สมฺมาทิฏฺฐิ ฯเปฯ.\csromanspacer{} อวิกฺเขปฏฺเฐน สมฺมาสมาธิ รูปราคสฺส อรูปราคสฺส มานสฺส อุทฺธจฺจสฺส อวิชฺชาย, มานานุสยสฺส ภวราคานุสยสฺส อวิชฺชานุสยสฺส ปหานาย มคฺโค เจว เหตุ จ, สหชาตานํ ธมฺมานํ อุปตฺถมฺภนาย มคฺโค เจว เหตุ จ, กิเลสานํ ปริยาทานาย มคฺโค เจว เหตุ จ, ปฏิเวธาทิวิโสธนาย มคฺโค เจว เหตุ จ, จิตฺตสฺส อธิฏฺฐานาย มคฺโค เจว เหตุ จ, จิตฺตสฺส โวทานาย มคฺโค เจว เหตุ จ, วิเสสาธิคมาย มคฺโค เจว เหตุ จ, อุตฺตริปฏิเวธาย มคฺโค เจว เหตุ จ, สจฺจาภิสมยาย มคฺโค เจว เหตุ จ, นิโรเธ ปติฏฺฐาปนาย มคฺโค เจว เหตุ จ.}

\prose{ทสฺสนมคฺโค สมฺมาทิฏฺฐิ, อภินิโรปนมคฺโค สมฺมาสงฺกปฺโป, ปริคฺคหมคฺโค สมฺมาวาจา, สมุฏฺฐานมคฺโค สมฺมากมฺมนฺโต, โวทานมคฺโค สมฺมา-อาชีโว, ปคฺคหมคฺโค สมฺมาวายาโม, อุปฏฺฐานมคฺโค สมฺมาสติ, อวิกฺเขปมคฺโค สมฺมาสมาธิ.\csromanspacer{} อุปฏฺฐานมคฺโค สติสมฺโพชฺฌงฺโค, ปวิจยมคฺโค ธมฺมวิจยสมฺโพชฺฌงฺโค, ปคฺคหมคฺโค วีริยสมฺโพชฺฌงฺโค, ผรณมคฺโค ปีติสมฺโพชฺฌงฺโค,}

\csromanmatikaanchor{mtk.83}
\csromantocmarkref{section}{10. มณฺฑเปยฺยกถา}{mtk.83}
\bookmark[dest={mtk.83},level=1]{10. มณฺฑเปยฺยกถา}
\csromanheader{1}{10. มณฺฑเปยฺยกถา}
\prose{\csromanfolio{277}อุปสมมคฺโค ปสฺสทฺธิสมฺโพชฺฌงฺโค, อวิกฺเขปมคฺโค สมาธิสมฺโพชฺฌงฺโค, ปฏิสงฺขานมคฺโค อุเปกฺขาสมฺโพชฺฌงฺโค.}

\prose{อสฺสทฺธิเย อกมฺปิยมคฺโค สทฺธาพลํ, โกสชฺเช อกมฺปิยมคฺโค วีริยพลํ, ปมาเท อกมฺปิยมคฺโค สติพลํ, อุทฺธจฺเจ อกมฺปิยมคฺโค สมาธิพลํ, อวิชฺชาย อกมฺปิยมคฺโค ปญฺญาพลํ, อธิโมกฺขมคฺโค สทฺธินฺทฺริยํ, ปคฺคหมคฺโค วีริยินฺทฺริยํ, อุปฏฺฐานมคฺโค สตินฺทฺริยํ, อวิกฺเขปมคฺโค สมาธินฺทฺริยํ, ทสฺสนมคฺโค ปญฺญินฺทฺริยํ.}

\prose{อาธิปเตยฺยฏฺเฐน อินฺทฺริยา มคฺโค, อกมฺปิยฏฺเฐน พลา มคฺโค, นิยฺยานฏฺเฐน โพชฺฌงฺคา มคฺโค, เหตุฏฺเฐน มคฺคงฺคา มคฺโค, อุปฏฺฐานฏฺเฐน สติปฏฺฐานา มคฺโค, ปทหนฏฺเฐน สมฺมปฺปธานา มคฺโค, อิชฺฌนฏฺเฐน อิทฺธิปาทา มคฺโค, ตถฏฺเฐน สจฺจานิ มคฺโค, อวิกฺเขปฏฺเฐน สมโถ มคฺโค, อนุปสฺสนฏฺเฐน วิปสฺสนา มคฺโค, เอกรสฏฺเฐน สมถวิปสฺสนา มคฺโค.\csromanspacer{} อนติวตฺตนฏฺเฐน ยุคนทฺธา มคฺโค.\csromanspacer{} สํวรฏฺเฐน สีลวิสุทฺธิ มคฺโค, อวิกฺเขปฏฺเฐน จิตฺตวิสุทฺธิ มคฺโค, ทสฺสนฏฺเฐน ทิฏฺฐิวิสุทฺธิ มคฺโค, มุตฺตฏฺเฐน วิโมกฺโข มคฺโค, ปฏิเวธฏฺเฐน วิชฺชา มคฺโค, ปริจฺจาคฏฺเฐน วิมุตฺติ มคฺโค, สมุจฺเฉทฏฺเฐน ขเย ญาณํ มคฺโค, ฉนฺโท มูลฏฺเฐน มคฺโค, มนสิกาโร สมุฏฺฐานฏฺเฐน มคฺโค, ผสฺโส สโมธานฏฺเฐน มคฺโค, เวทนา สโมสรณฏฺเฐน มคฺโค, สมาธิ ปมุขฏฺเฐน มคฺโค, สติ อาธิปเตยฺยฏฺเฐน มคฺโค, ปญฺญา ตทุตฺตรฏฺเฐน มคฺโค, วิมุตฺติ สารฏฺเฐน มคฺโค, อมโตคธํ นิพฺพานํ ปริโยสานฏฺเฐน มคฺโคติ.}

\nitthitamruled{มคฺคกถา นิฏฺฐิตา.}
\csromanheader{2}{10. มณฺฑเปยฺยกถา}
\proseitem{238}{มณฺฑเปยฺยมิทํ ภิกฺขเว พฺรหฺมจริยํ, สตฺถา สมฺมุขีภูโต.\csromanspacer{} ติธตฺตมณฺโฑ\footnote{ติวิโธ มณฺโฑ (พหูสุ) อฏฺฐกถา โอโลเกตพฺพา.} สตฺถริ สมฺมุขีภูเต เทสนามณฺโฑ ปฏิคฺคหมณฺโฑ พฺรหฺมจริยมณฺโฑ.}

\prose{กตโม \textbf{เทสนามณฺโฑ?} จตุนฺนํ อริยสจฺจานํ อาจิกฺขนา เทสนา ปญฺญาปนา ปฏฺฐปนา วิวรณา วิภชนา อุตฺตานีกมฺมํ\footnote{อุตฺตานิกมฺมํ (ก)}.\csromanspacer{} จตุนฺนํ สติปฏฺฐานานํ ฯเปฯ.\csromanspacer{} จตุนฺนํ สมฺมปฺปธานานํ.\csromanspacer{} จตุนฺนํ อิทฺธิปาทานํ.\csromanspacer{} ปญฺจนฺนํ อินฺทฺริยานํ.}

\prose{\csromanfolio{278}ปญฺจนฺนํ พลานํ.\csromanspacer{} สตฺตนฺนํ โพชฺฌงฺคานํ.\csromanspacer{} อริยสฺส อฏฺฐงฺคิกสฺส มคฺคสฺส อาจิกฺขนา เทสนา ปญฺญาปนา ปฏฺฐปนา วิวรณา วิภชนา อุตฺตานีกมฺมํ.\csromanspacer{} อยํ เทสนามณฺโฑ. (1)}

\prose{กตโม \textbf{ปฏิคฺคหมณฺโฑ?} ภิกฺขู ภิกฺขุนิโย อุปาสกา อุปาสิกาโย เทวา มนุสฺสา, เย วา ปนญฺเญปิ เกจิ วิญฺญาตาโร.\csromanspacer{} อยํ ปฏิคฺคหมณฺโฑ. (2)}

\prose{กตโม \textbf{พฺรหฺมจริยมณฺโฑ?} อยเมว อริโย อฏฺฐงฺคิโก มคฺโค.\csromanspacer{} เสยฺยถิทํ\csromandash{}สมฺมาทิฏฺฐิ สมฺมาสงฺกปฺโป สมฺมาวาจา สมฺมากมฺมนฺโต สมฺมา-อาชีโว สมฺมาวายาโม สมฺมาสติ สมฺมาสมาธิ.\csromanspacer{} อยํ พฺรหฺมจริยมณฺโฑ. (3)}

\proseitem{239}{อธิโมกฺขมณฺโฑ สทฺธินฺทฺริยํ, อสฺสทฺธิยํ กสโฏ, อสฺสทฺธิยํ กสฏํ ฉฑฺเฑตฺวา สทฺธินฺทฺริยสฺส อธิโมกฺขมณฺฑํ ปิวตีติ มณฺฑเปยฺยํ.\csromanspacer{} ปคฺคหมณฺโฑ วีริยินฺทฺริยํ, โกสชฺชํ กสโฏ, โกสชฺชํ กสฏํ ฉฑฺเฑตฺวา วีริยินฺทฺริยสฺส ปคฺคหมณฺฑํ ปิวตีติ มณฺฑเปยฺยํ.\csromanspacer{} อุปฏฺฐานมณฺโฑ สตินฺทฺริยํ, ปมาโท กสโฏ, ปมาทํ กสฏํ ฉฑฺเฑตฺวา สตินฺทฺริยสฺส อุปฏฺฐานมณฺฑํ ปิวตีติ มณฺฑเปยฺยํ.\csromanspacer{} อวิกฺเขปมณฺโฑ สมาธินฺทฺริยํ, อุทฺธจฺจํ กสโฏ, อุทฺธจฺจํ กสฏํ ฉฑฺเฑตฺวา สมาธินฺทฺริยสฺส อวิกฺเขปมณฺฑํ ปิวตีติ มณฺฑเปยฺยํ.\csromanspacer{} ทสฺสนมณฺโฑ ปญฺญินฺทฺริยํ, อวิชฺชา กสโฏ, อวิชฺชํ กสฏํ ฉฑฺเฑตฺวา ปญฺญินฺทฺริยสฺส ทสฺสนมณฺฑํ ปิวตีติ มณฺฑเปยฺยํ.}

\prose{อสฺสทฺธิเย อกมฺปิยมณฺโฑ สทฺธาพลํ, อสฺสทฺธิยํ กสโฏ, อสฺสทฺธิยํ กสฏํ ฉฑฺเฑตฺวา สทฺธาพลสฺส อสฺสทฺธิเย อกมฺปิยมณฺฑํ ปิวตีติ มณฺฑเปยฺยํ.\csromanspacer{} โกสชฺเช อกมฺปิยมณฺโฑ วิริยพลํ, โกสชฺชํ กสโฏ, โกสชฺชํ กสฏํ ฉฑฺเฑตฺวา วีริยพลสฺส โกสชฺเช อกมฺปิยมณฺฑํ ปิวตีติ มณฺฑเปยฺยํ.\csromanspacer{} ปมาเท อกมฺปิยมณฺโฑ สติพลํ, ปมาโท กสโฏ, ปมาทํ กสฏํ ฉฑฺเฑตฺวา สติพลสฺส ปมาเท อกมฺปิยมณฺฑํ ปิวตีติ มณฺฑเปยฺยํ.\csromanspacer{} อุทฺธจฺเจ อกมฺปิยมณฺโฑ สมาธิพลํ, อุทฺธจฺจํ กสโฏ, อุทฺธจฺจํ กสฏํ ฉฑฺเฑตฺวา สมาธิพลสฺส อุทฺธจฺเจ อกมฺปิยมณฺฑํ ปิวตีติ มณฺฑเปยฺยํ.\csromanspacer{} อวิชฺชาย อกมฺปิยมณฺโฑ ปญฺญาพลํ, อวิชฺชา กสโฏ, อวิชฺชํ กสฏํ ฉฑฺเฑตฺวา ปญฺญาพลสฺส อวิชฺชาย อกมฺปิยมณฺฑํ ปิวตีติ มณฺฑเปยฺยํ.}

\csromanheader{2}{10. มณฺฑเปยฺยกถา}
\prose{\csromanfolio{279}อุปฏฺฐานมณฺโฑ สติสมฺโพชฺฌงฺคา, ปมาโท กสโฏ, ปมาทํ กสฏํ ฉฑฺเฑตฺวา สติสมฺโพชฺฌงฺคสฺส อุปฏฺฐานมณฺฑํ ปิวตีติ มณฺฑเปยฺยํ.\csromanspacer{} ปวิจยมณฺโฑ ธมฺมวิจยสมฺโพชฺฌงฺโค, อวิชฺชา กสโฏ, อวิชฺชํ กสฏํ ฉฑฺเฑตฺวา ธมฺมวิจยสมฺโพชฺฌงฺคสฺส ปวิจยมณฺฑํ ปิวตีติ มณฺฑเปยฺยํ.\csromanspacer{} ปคฺคหมณฺโฑ วีริยสมฺโพชฺฌงฺโค, โกสชฺชํ กสโฏ, โกสชฺชํ กสฏํ ฉฑฺเฑตฺวา วีริยสมฺโพชฺฌงฺคสฺส ปคฺคหมณฺฑํ ปิวตีติ มณฺฑเปยฺยํ.\csromanspacer{} ผรณมณฺโฑ ปีติสมฺโพชฺฌงฺโค, ปริฬาโห กสโฏ, ปริฬาหํ กสฏํ ฉฑฺเฑตฺวา ปีติสมฺโพชฺฌงฺคสฺส ผรณมณฺฑํ ปิวตีติ มณฺฑเปยฺยํ.\csromanspacer{} อุปสมมณฺโฑ ปสฺสทฺธิสมฺโพชฺฌงฺโค, ทุฏฺฐุลฺลํ กสโฏ, ทุฏฺฐุลฺลํ กสฏํ ฉฑฺเฑตฺวา ปสฺสทฺธิสมฺโพชฺฌงฺคสฺส อุปสมมณฺฑํ ปิวตีติ มณฺฑเปยฺยํ.\csromanspacer{} อวิกฺเขปมณฺโฑ สมาธิสมฺโพชฺฌงฺโค, อุทฺธจฺจํ กสโฏ, อุทฺธจฺจํ กสฏํ ฉฑฺเฑตฺวา สมาธิสมฺโพชฺฌงฺคสฺส อวิกฺเขปมณฺฑํ ปิวตีติ มณฺฑเปยฺยํ.\csromanspacer{} ปฏิสงฺขานมณฺโฑ อุเปกฺขาสมฺโพชฺฌงฺโค, อปฺปฏิสงฺขา กสโฏ, อปฺปฏิสงฺขํ กสฏํ ฉฑฺเฑตฺวา อุเปกฺขาสมฺโพชฺฌงฺคสฺส ปฏิสงฺขานมณฺฑํ ปิวตีติ มณฺฑเปยฺยํ.}

\prose{ทสฺสนมณฺโฑ สมฺมาทิฏฺฐิ, มิจฺฉาทิฏฺฐิ กสโฏ, มิจฺฉาทิฏฺฐิํ กสฏํ ฉฑฺเฑตฺวา สมฺมาทิฏฺฐิยา ทสฺสนมณฺฑํ ปิวตีติ มณฺฑเปยฺยํ.\csromanspacer{} อภินิโรปนมณฺโฑ สมฺมาสงฺกปฺโป, มิจฺฉาสงฺกปฺโป กสโฏ, มิจฺฉาสงฺกปฺปํ กสฏํ ฉฑฺเฑตฺวา สมฺมาสงฺกปฺปสฺส อภินิโรปนมณฺฑํ ปิวตีติ มณฺฑเปยฺยํ.\csromanspacer{} ปริคฺคหมณฺโฑ สมฺมาวาจา, มิจฺฉาวาจา กสโฏ, มิจฺฉาวาจํ กสฏํ ฉฑฺเฑตฺวา สมฺมาวาจาย ปริคฺคหมณฺฑํ ปิวตีติ มณฺฑเปยฺยํ.\csromanspacer{} สมุฏฺฐานมณฺโฑ สมฺมากมฺมนฺโต, มิจฺฉากมฺมนฺโต กสโฏ, มิจฺฉากมฺมนฺตํ กสฏํ ฉฏฺเฏตฺวา สมฺมากมฺมนฺตสฺส สมุฏฺฐานมณฺฑํ ปิวตีติ มณฺฑเปยฺยํ.\csromanspacer{} โวทานมณฺโฑ สมฺมา-อาชีโว, มิจฺฉา-อาชีโว กสโฏ, มิจฺฉา-อาชีวํ กสฏํ ฉฑฺเฑตฺว ฉฑฺเฑตฺวา สมฺมา-อาชีวสฺส โวทานมณฺฑํ ปิวตีติ มณฺฑเปยฺยํ.\csromanspacer{} ปคฺคหมณฺโฑ สมฺมาวายาโม, มิจฺฉาวายาโม กสโฏ.\csromanspacer{} มิจฺฉาวายามํ กสฏํ ฉฑฺเฑตฺวา สมฺมาวายามสฺส ปคฺคหมณฺฑํ ปิวตีติ มณฺฑเปยฺยํ.\csromanspacer{} อุปฏฺฐานมณฺโฑ สมฺมาสติ, มิจฺฉาสติ กสโฏ, มิจฺฉาสติํ กสฏํ ฉฑฺเฑตฺวา สมฺมาสติยา อุปฏฺฐานมณฺฑํ ปิวตีติ มณฺฑเปยฺยํ.\csromanspacer{} อวิกฺเขปมณฺโฑ สมฺมาสมาธิ, มิจฺฉาสมาธิ กสโฏ, มิจฺฉาสมาธิํ กสฏํ ฉฑฺเฑตฺวา สมฺมาสมาธิสฺส อวิกฺเขปมณฺฑํ ปิวตีติ มณฺฑเปยฺยํ.}

\proseitem{240}{อตฺถิ มณฺโฑ, อตฺถิ เปยฺยํ, อตฺถิ กสโฏ.\csromanspacer{} อธิโมกฺขมณฺโฑ สทฺธินฺทฺริยํ, อสฺสทฺธิยํ กสโฏ, โย ตตฺถ อตฺถรโส ธมฺมรโส \csromanfolio{280}วิมุตฺติรโส, อิทํ เปยฺยํ.\csromanspacer{} ปคฺคหมณฺโฑ วีริยินฺทฺริยํ, โกสชฺชํ กสโฏ, โย ตตฺถ อตฺถรโส ธมฺมรโส วิมุตฺติรโส, อิทํ เปยฺยํ.\csromanspacer{} อุปฏฺฐานมณฺโฑ สตินฺทฺริยํ, ปมาโท กสโฏ, โย ตตฺถ อตฺถรโส ธมฺมรโส วิมุตฺติรโส, อิทํ เปยฺยํ.\csromanspacer{} อวิกฺเขปมณฺโฑ สมาธินฺทฺริยํ, อุทฺธจฺจํ กสโฏ, โย ตตฺถ อตฺถรโส ธมฺมรโส วิมุตฺติรโส, อิทํ เปยฺยํ.\csromanspacer{} ทสฺสนมณฺโฑ ปญฺญินฺทฺริยํ, อวิชฺชา กสโฏ, โย ตตฺถ อตฺถรโส ธมฺมรโส วิมุตฺติรโส, อิทํ เปยฺยํ.}

\prose{อสฺสทฺธิเย อกมฺปิยมณฺโฑ สทฺธาพลํ, อสฺสทฺธิยํ กสโฏ, โย ตตฺถ อตฺถรโส ธมฺมรโส วิมุตฺติรโส, อิทํ เปยฺยํ.\csromanspacer{} โกสชฺเช อกมฺปิยมณฺโฑ วีริยพลํ, โกสชฺชํ กสโฏ, โย ตตฺถ อตฺถรโส ธมฺมรโส วิมุตฺติรโส, อิทํ เปยฺยํ.\csromanspacer{} ปมาเท อกมฺปิยมณฺโฑ สติพลํ, ปมาโท กสโฏ, โย ตตฺถ อตฺถรโส ธมฺมรโส วิมุตฺติรโส, อิทํ เปยฺยํ.\csromanspacer{} อุทฺธจฺเจ อกมฺปิยมณฺโฑ สมาธิพลํ, อุทฺธจฺจํ กสโฏ, โย ตตฺถ อตฺถรโส ธมฺมรโส วิมุตฺติรโส, อิทํ เปยฺยํ.\csromanspacer{} อวิชฺชาย อกมฺปิยมณฺโฑ ปญฺญาพลํ, อวิชฺชา กสโฏ, โย ตตฺถ อตฺถรโส ธมฺมรโส วิมุตฺติรโส, อิทํ เปยฺยํ.}

\prose{อุปฏฺฐานมณฺโฑ สติสมฺโพชฺฌงฺโค, ปมาโท กสโฏ, โย ตตฺถ อตฺถรโส ธมฺมรโส วิมุตฺติรโส, อิทํ เปยฺยํ.\csromanspacer{} ปวิจยมณฺโฑ ธมฺมวิจยสมฺโพชฺฌงฺโค, อวิชฺชา กสโฏ, โย ตตฺถ อตฺถรโส ธมฺมรโส วิมุตฺติรโส, อิทํ เปยฺยํ.\csromanspacer{} ปคฺคหมณฺโฑ วีริยสมฺโพชฺฌงฺโค, โกสชฺชํ กสโฏ, โย ตตฺถ อตฺถรโส ธมฺมรโส วิมุตฺติรโส, อิทํ เปยฺยํ.\csromanspacer{} ผรณมณฺโฑ ปีติสมฺโพชฺฌงฺโค, ปริฬาโห กสโฏ, โย ตตฺถ อตฺถรโส ธมฺมรโส วิมุตฺติรโส, อิทํ เปยฺยํ.\csromanspacer{} อุปสมมณฺโฑ ปสฺสทฺธิสมฺโพชฺฌงฺโค, ทุฏฺฐุลฺลํ กสโฏ, โย ตตฺถ อตฺถรโส ธมฺมรโส วิมุตฺติรโส, อิทํ เปยฺยํ.\csromanspacer{} อวิกฺเขปมณฺโฑ สมาธิสมฺโพชฺฌงฺโค, อุทฺธจฺจํ กสโฏ, โย ตตฺถ อตฺถรโส ธมฺมรโส วิมุตฺติรโส, อิทํ เปยฺยํ.\csromanspacer{} ปฏิสงฺขานมณฺโฑ อุเปกฺขาสมฺโพชฺฌงฺโค, อปฏิสงฺขา กสโฏ, โย ตตฺถ อตฺถรโส ธมฺมรโส วิมุตฺติรโส, อิทํ เปยฺยํ.}

\prose{ทสฺสนมณฺโฑ สมฺมาทิฏฺฐิ, มิจฺฉาทิฏฺฐิ กสโฏ, โย ตตฺถ อตฺถรโส ธมฺมรโส วิมุตฺติรโส, อิทํ เปยฺยํ.\csromanspacer{} อภินิโรปนมณฺโฑ สมฺมาสงฺกปฺโป, มิจฺฉาสงฺกปฺโป กสโฏ, โย ตตฺถ อตฺถรโส ธมฺมรโส วิมุตฺติรโส,}

\csromanheader{2}{10. มณฺฑเปยฺยกถา}
\prose{\csromanfolio{281}อิทํ เปยฺยํ.\csromanspacer{} ปริคฺคหมณฺโฑ สมฺมาวาจา, มิจฺฉาวาจา กสโฏ, โย ตตฺถ อตฺถรโส ธมฺมรโส วิมุตฺติรโส, อิทํ เปยฺยํ.\csromanspacer{} สมุฏฺฐานมณฺโฑ สมฺมากมฺมนฺโต, มิจฺฉากมฺมนฺโต กสโฏ, โย ตตฺถ อตฺถรโส ธมฺมรโส วิมุตฺติรโส, อิทํ เปยฺยํ.\csromanspacer{} โวทานมณฺโฑ สมฺมา-อาชีโว, มิจฺฉา-อาชีโว กสโฏ, โย ตตฺถ อตฺถรโส ธมฺมรโส วิมุตฺติรโส, อิทํ เปยฺยํ.\csromanspacer{} ปคฺคหมณฺโฑ สมฺมาวายาโม, มิจฺฉาวายาโม กสโฏ, โย ตตฺถ อตฺถรโส ธมฺมรโส วิมุตฺติรโส, อิทํ เปยฺยํ.\csromanspacer{} อุปฏฺฐานมณฺโฑ สมฺมาสติ, มิจฺฉาสติ กสโฏ, โย ตตฺถ อตฺถรโส ธมฺมรโส วิมุตฺติรโส, อิทํ เปยฺยํ.\csromanspacer{} อวิกฺเขปมณฺโฑ สมฺมาสมาธิ, มิจฺฉาสมาธิ กสโฏ, โย ตตฺถ อตฺถรโส ธมฺมรโส วิมุตฺติรโส, อิทํ เปยฺยํ.}

\prose{ทสฺสนมณฺโฑ สมฺมาทิฏฺฐิ.\csromanspacer{} อภินิโรปนมณฺโฑ สมฺมาสงฺกปฺโป.\csromanspacer{} ปริคฺคหมณฺโฑ สมฺมาวาจา.\csromanspacer{} สมุฏฺฐานมณฺโฑ สมฺมากมฺมนฺโต.\csromanspacer{} โวทานมณฺโฑ สมฺมา-อาชีโว.\csromanspacer{} ปคฺคหมณฺโฑ สมฺมาวายาโม.\csromanspacer{} อุปฏฺฐานมณฺโฑ สมฺมาสติ.\csromanspacer{} อวิกฺเขปมณฺโฑ สมฺมาสมาธิ.}

\prose{อุปฏฺฐานมณฺโฑ สติสมฺโพชฺฌงฺโค, ปวิจยมณฺโฑ ธมฺมวิจยสมฺโพชฺฌงฺโค, ปคฺคหมณฺโฑ วีริยสมฺโพชฺฌงฺโค, ผรณมณฺโฑ ปีติสมฺโพชฺฌงฺโค, อุปสมมณฺโฑ ปสฺสทฺธิสมฺโพชฺฌงฺโค, อวิกฺเขปมณฺโฑ สมาธิสมฺโพชฺฌงฺโค, ปฏิสงฺขานมณฺโฑ อุเปกฺขาสมฺโพชฺฌงฺโค.}

\prose{อสฺสทฺธิเย อกมฺปิยมณฺโฑ สทฺธาพลํ, โกสชฺเช อกมฺปิยมณฺโฑ วีริยพลํ, ปมาเท อกมฺปิยมณฺโฑ สติพลํ, อุทฺธจฺเจ อกมฺปิยมณฺโฑ สมาธิพลํ, อวิชฺชาย อกมฺปิยมณฺโฑ ปญฺญาพลํ.}

\prose{อธิโมกฺขมณฺโฑ สทฺธินฺทฺริยํ, ปคฺคหมณฺโฑ วีริยินฺทฺริยํ, อุปฏฺฐานมณฺโฑ สตินฺทฺริยํ, อวิกฺเขปมณฺโฑ สมาธินฺทฺริยํ, ทสฺสนมณฺโฑ ปญฺญินฺทฺริยํ.}

\prose{อาธิปเตยฺยฏฺเฐน อินฺทฺริยา มณฺโฑ, อกมฺปิยฏฺเฐน พลา มณฺโฑ, นิยฺยานฏฺเฐน โพชฺฌงฺคา มณฺโฑ, เหตุฏฺเฐน มคฺโค มณฺโฑ, อุปฏฺฐานฏฺเฐน สติปฏฺฐานา มณฺโฑ, ปทหนฏฺเฐน สมฺมปฺปธานา มณฺโฑ, อิชฺฌนฏฺเฐน อิทฺธิปาทา มณฺโฑ, ตถฏฺเฐน สจฺจา มณฺโฑ, อวิกฺเขปฏฺเฐน สมโถ มณฺโฑ, อนุปสฺสนฏฺเฐน วิปสฺสนา มณฺโฑ, เอกรสฏฺเฐน สมถวิปสฺสนา มณฺโฑ, อนติวตฺตนฏฺเฐน ยุคนทฺธา มณฺโฑ, สํวรฏฺเฐน สีลวิสุทฺธิ มณฺโฑ, อวิกฺเขปฏฺเฐน จิตฺตวิสุทฺธิ มณฺโฑ, ทสฺสนฏฺเฐน ทิฏฺฐิวิสุทฺธิ มณฺโฑ, มุตฺตฏฺเฐน}

\prose{\csromanfolio{282}วิโมกฺโข มณฺโฑ, ปฏิเวธฏฺเฐน วิชฺชา มณฺโฑ, ปริจฺจาคฏฺเฐน วิมุตฺติ มณฺโฑ, สมุจฺเฉทฏฺเฐน ขเย ญาณํ มณฺโฑ, ปฏิปฺปสฺสทฺธฏฺเฐน อนุปฺปาเท ญาณํ มณฺโฑ.\csromanspacer{} ฉนฺโท มูลฏฺเฐน มณฺโฑ, มนสิกาโร สมุฏฺฐานฏฺเฐน มณฺโฑ, ผสฺโส สโมธานฏฺเฐน มณฺโฑ, เวทนา สโมสรณฏฺเฐน มณฺโฑ, สมาธิ ปมุขฏฺเฐน มณฺโฑ, สติ อาธิปเตยฺยฏฺเฐน มณฺโฑ, ปญฺญา ตทุตฺตรฏฺเฐน มณฺโฑ, วิมุตฺติ สารฏฺเฐน มณฺโฑ, อมโตคธํ นิพฺพานํ ปริโยสานฏฺเฐน มณฺโฑติ.}

\nitthitam{มณฺฑเปยฺยกถา นิฏฺฐิตา.}
\vspace{\tipitakaparskip}
\csromancenter{จตุตฺถภาณวาโร.}
\csromancenter{\textbf{มหาวคฺโค ปฐโม.}}
\titlehead{\textbf{ตสฺสุทฺทานํ}}
\csromangathameasure{ญาณทิฏฺฐี จ อสฺสาสา, อินฺทฺริยํ วิโมกฺขปญฺจมา. \\ คติกมฺมวิปลฺลาสา, มคฺโค มณฺเฑน เต ทสาติ. \\ เอส นิกายธเรหิ ฐปิโต อสโม ปฐโม ปวโร วรวคฺโคติ.}
\csromangathagroup{\csromangathastanza{ญาณทิฏฺฐี จ อสฺสาสา, อินฺทฺริยํ วิโมกฺขปญฺจมา. \\ คติกมฺมวิปลฺลาสา, มคฺโค มณฺเฑน เต ทสาติ. \\ เอส นิกายธเรหิ ฐปิโต อสโม\footnote{อสฺสโม (สฺยา)} ปฐโม ปวโร วรวคฺโคติ\footnote{วรมคฺโคติ (สฺยา)}.}}

\csromanmatikaanchor{mtk.84}
\csromantocmarknopage{chapter}{2. ยุคนทฺธวคฺค}{mtk.84}
\bookmark[dest={mtk.84},level=0]{2. ยุคนทฺธวคฺค}
\chapterheadpageread{2. ยุคนทฺธวคฺค}
\csromanmatikaanchor{mtk.85}
\csromantocmarknopage{section}{1. ยุคนทฺธกถา}{mtk.85}
\bookmark[dest={mtk.85},level=1]{1. ยุคนทฺธกถา}
\csromanheader{1}{1. ยุคนทฺธกถา}
\proseitem{1}{\csromanfolio{283}เอวํ เม สุตํ, เอกํ สมยํ อายสฺมา อานนฺโท โกสมฺพิยํ วิหรติ โฆสิตาราเม.\csromanspacer{} ตตฺร โข อายสฺมา อานนฺโท ภิกฺขู อามนฺเตสิ “อาวุโส ภิกฺขโว”ติ. “อาวุโส”ติ โข เต ภิกฺขู อายสฺมโต อานนฺทสฺส ปจฺจสฺโสสุํ.\csromanspacer{} อายสฺมา อานนฺโท เอตทโวจ\csromandash{}}

\prose{โย หิ โกจิ อาวุโส ภิกฺขุ วา ภิกฺขุนี วา มม สนฺติเก อรหตฺตปฺปตฺตํ\footnote{อรหตฺตํ (สฺยา), อรหตฺตปตฺติํ (อํ 1. 475 ปิฏฺเฐ.)} พฺยากโรติ สพฺพโส จตูหิ มคฺเคหิ, เอเตสํ วา อญฺญตเรน.\csromanspacer{} กตเมหิ จตูหิ?}

\prose{อิธาวุโส ภิกฺขุ สมถปุพฺพงฺคมํ วิปสฺสนํ ภาเวติ, ตสฺส สมถปุพฺพงฺคมํ วิปสฺสนํ ภาวยโต มคฺโค สญฺชายติ, โส ตํ มคฺคํ อาเสวติ ภาเวติ พหุลีกโรติ\footnote{พหุลิํ กโรติ (ก) อํ 1. 475 ปิฏฺเฐ ปสฺสิตพฺพา.}, ตสฺส ตํ มคฺคํ อาเสวโต ภาวยโต พหุลี กโรโต สญฺโญชนานิ ปหียนฺติ, อนุสยา พฺยนฺตีโหนฺติ.}

\prose{ปุน จปรํ อาวุโส ภิกฺขุ วิปสฺสนาปุพฺพงฺคมํ สมถํ ภาเวติ, ตสฺส วิปสฺสนาปุพฺพงฺคมํ สมถํ ภาวยโต มคฺโค สญฺชายติ, โส ตํ มคฺคํ อาเสวติ ภาเวติ พหุลีกโรติ, ตสฺส ตํ มคฺคํ อาเสวโต ภาวยโต พหุลีกโรโต สญฺโญชนานิ ปหียนฺติ, อนุสยา พฺยนฺตีโหนฺติ.}

\prose{ปุน จปรํ อาวุโส ภิกฺขุ สมถวิปสฺสนํ ยุคนทฺธํ\footnote{ยุคนนฺธํ (ก, สี-ฏฺฐ)} ภาเวติ, ตสฺส สมถวิปสฺสนํ ยุคนทฺธํ ภาวยโต มคฺโค สญฺชายติ, โส ตํ มคฺคํ อาเสวติ ภาเวติ พหุลีกโรติ, ตสฺส ตํ มคฺคํ อาเสวโต ภาวยโต พหุลีกโรโต สญฺโญชนานิ ปหียนฺติ, อนุสยา พฺยนฺตีโหนฺติ.}

\prose{ปุน จปรํ อาวุโส ภิกฺขุโน ธมฺมุทฺธจฺจวิคฺคหิตํ มานสํ โหติ, โหติ โส อาวุโส สมโย ยํ ตํ จิตฺตํ อชฺฌตฺตเมว\footnote{อชฺฌตฺตญฺเญว (สฺยา, ก)} สนฺติฏฺฐติ สนฺนิสีทติ \csromanfolio{284}เอโกทิ โหติ สมาธิยติ, ตสฺส มคฺโค สญฺชายติ, โส ตํ มคฺคํ อาเสวติ ภาเวติ พหุลีกโรติ, ตสฺส ตํ มคฺคํ อาเสวโต ภาวยโต พหุลีกโรโต สญฺโญชนานิ ปหียนฺติ, อนุสยา พฺยนฺตีโหนฺติ.}

\prose{โย หิ โกจิ อาวุโส ภิกฺขุ วา ภิกฺขุนี วา มม สนฺติเก อรหตฺตปฺปตฺตํ พฺยากโรติ สพฺพโส อิเมหิ จตูหิ มคฺเคหิ, เอเตสํ วา อญฺญตเรนาติ.}

\csromanmatikaanchor{mtk.86}
\csromantocmarkref{subsection}{1. สุตฺตนฺตนิทฺเทส}{mtk.86}
\bookmark[dest={mtk.86},level=2]{1. สุตฺตนฺตนิทฺเทส}
\csromanheader{2}{1. สุตฺตนฺตนิทฺเทส}
\proseitem{2}{กถํ \textbf{สมถปุพฺพงฺคมํ วิปสฺสนํ ภาเวติ?} เนกฺขมฺมวเสน จิตฺตสฺส เอกคฺคตา อวิกฺเขโป สมาธิ, ตตฺถ ชาเต ธมฺเม อนิจฺจโต อนุปสฺสนฏฺเฐน วิปสฺสนา, ทุกฺขโต อนุปสฺสนฏฺเฐน วิปสฺสนา, อนตฺตโต อนุปสฺสฏฺเฐน วิปสฺสนา.\csromanspacer{} อิติ ปฐมํ สมโถ, ปจฺฉา วิปสฺสนา.\csromanspacer{} เตน วุจฺจติ “สมถปุพฺพงฺคมํ วิปสฺสนํ ภาเวตี”ติ.\csromanspacer{}\csromanspacer{}\textbf{ภาเวตี}ติ จตสฺโส ภาวนา, ตตฺถ ชาตานํ ธมฺมานํ อนติวตฺตนฏฺเฐน ภาวนา, อินฺทฺริยานํ เอกรสฏฺเฐน ภาวนา, ตทุปควีริยวาหนฏฺเฐน ภาวนา, อาเสวนฏฺเฐน ภาวนา.}

\prose{\textbf{มคฺโค สญฺชายตี}ติ กถํ มคฺโค สญฺชายติ? ทสฺสนฏฺเฐน สมฺมาทิฏฺฐิ มคฺโค สญฺชายติ, อภินิโรปนฏฺเฐน สมฺมาสงฺกปฺโป มคฺโค สญฺชายติ, ปริคฺคหฏฺเฐน สมฺมาวาจา มคฺโค สญฺชายติ, สมุฏฺฐานฏฺเฐน สมฺมากมฺมนฺโต มคฺโค สญฺชายติ, โวทานฏฺเฐน สมฺมา-อาชีโว มคฺโค สญฺชายติ, ปคฺคหฏฺเฐน สมฺมาวายาโม มคฺโค สญฺชายติ, อุปฏฺฐานฏฺเฐน สมฺมาสติ มคฺโค สญฺชายติ, อวิกฺเขปฏฺเฐน สมฺมาสมาธิ มคฺโค สญฺชายติ.\csromanspacer{} เอวํ มคฺโค สญฺชายติ.}

\prose{\textbf{โส ตํ มคฺคํ อาเสวติ ภาเวติ พหุลีกโรติ อาเสวตี}ติ กถํ อาเสวติ? อาวชฺชนฺโต อาเสวติ, ชานนฺโต อาเสวติ, ปสฺสนฺโต อาเสวติ, ปจฺจเวกฺขนฺโต อาเสวติ, จิตฺตํ อธิฏฺฐหนฺโต อาเสวติ, สทฺธาย อธิมุจฺจนฺโต อาเสวติ, วีริยํ ปคฺคณฺหนฺโต อาเสวติ, สติํ อุปฏฺฐาเปนฺโต อาเสวติ, จิตฺตํ สมาทหนฺโต อาเสวติ, ปญฺญาย ปชานนฺโต อาเสวติ, อภิญฺเญยฺยํ อภิชานนฺโต อาเสวติ, ปริญฺเญยฺยํ ปริชานนฺโต อาเสวติ, ปหาตพฺพํ ปชหนฺโต}

\vspace{\tipitakaparskip}
\csromancenter{ยุคนทฺธกถา อาเสวติ, ภาเวตพฺพํ ภาเวนฺโต อาเสวติ, สจฺฉิกาตพฺพํ สจฺฉิกโรนฺโต อาเสวติ, เอวํ อาเสวติ.}
\prose{\csromanfolio{285}\textbf{ภาเวตี}ติ กถํ ภาเวติ? อาวชฺชนฺโต ภาเวติ, ชนนฺโต ภาเวติ, ปสฺสนฺโต ภาเวติ, ปจฺจเวกฺขนฺโต ภาเวติ, จิตฺตํ อธิฏฺฐหนฺโต ภาเวติ, สทฺธาย อธิมุจฺจนฺโต ภาเวติ วีริยํ ปคฺคณฺหนฺโต ภาเวติ, สติํ อุปฏฺฐาเปนฺโต ภาเวติ, จิตฺตํ สมาทหนฺโต ภาเวติ, ปญฺญาย ปชานนฺโต ภาเวติ.\csromanspacer{} อภิญฺเญยฺยํ อภิชานนฺโต ภาเวติ, ปริญฺเญยฺยํ ปริชานนฺโต ภาเวติ, ปหาตพฺพํ ปชหนฺโต ภาเวติ, ภาเวตพฺพํ ภาเวนฺโต ภาเวติ, สจฺฉิกาตพฺพํ สจฺฉิกโรนฺโต ภาเวติ, เอวํ ภาเวติ.}

\prose{\textbf{พหุลีกโรตี}ติ กถํ พหุลีกโรติ? อาวชฺชนฺโต พหุลีกโรติ, ชานนฺโต พหุลีกโรติ, ปสฺสนฺโต พหุลีกโรติ, ปจฺจเวกฺขนฺโต พหุลีกโรติ, จิตฺตํ อธิฏฺฐหนฺโต พหุลีกโรติ, สทฺธาย อธิมุจฺจนฺโต พหุลีกโรติ, วีริยํ ปคฺคณฺหนฺโต พหุลีกโรติ, สติํ อุปฏฺฐาเปนฺโต พหุลีกโรติ, จิตฺตํ สมาทหนฺโต พหุลีกโรติ, ปญฺญาย ปชานนฺโต พหุลีกโรติ, อภิญฺเญยฺยํ อภิชานนฺโต พหุลีกโรติ, ปริญฺเญยฺยํ ปริชานนฺโต พหุลีกโรติ, ปหาตพฺพํ ปชหนฺโต พหุลีกโรติ, ภาเวตพฺพํ ภาเวนฺโต พหุลีกโรติ, สจฺฉิกาตพฺพํ สจฺฉิกโรนฺโต พหุลีกโรติ, เอวํ พหุลีกโรติ.}

\prose{\textbf{ตสฺส ตํ มคฺคํ อาเสวโต ภาวยโต พหุลีกโรโต สญฺโญชนานิ ปหียนฺติ, อนุสยา พฺยนฺตีโหนฺตีติ} กถํ สญฺโญชนานิ ปหียนฺติ, อนุสยา พฺยนฺตีโหนฺติ? โสตาปตฺติมคฺเคน สกฺกายทิฏฺฐิ วิจิกิจฺฉา สีลพฺพตปรามาโส, อิมานิ ตีณิ สญฺโญชนานิ ปหียนฺติ, ทิฏฺฐานุสโย วิจิกิจฺฉานุสโย อิเม เทฺว อนุสยา พฺยนฺตีโหนฺติ.\csromanspacer{} สกทาคามิมคฺเคน โอฬาริกํ กามราคสญฺโญชนํ ปฏิฆสญฺโญชนํ อิมานิ เทฺว สญฺโญชนานิ ปหียนฺติ, โอฬาริโก กามราคานุสโย ปฏิฆานุสโย อิเม เทฺว อนุสยา พฺยนฺตีโหนฺติ.\csromanspacer{} อนาคามิมคฺเคน อณุสหคตํ กามราคสญฺโญชนํ ปฏิฆสญฺโญชนํ อิมานิ เทฺว สญฺโญชนานิ ปหียนฺติ, อณุสหคโต กามราคานุสโย ปฏิฆานุสโย อิเม เทฺว อนุสยา พฺยนฺตีโหนฺติ.\csromanspacer{} อรหตฺตมคฺเคน รูปราโค อรูปราโค มาโน อุทฺธจฺจํ อวิชฺชา อิมานิ ปญฺจ สญฺโญชนานิ ปหียนฺติ, มานานุสโย \csromanfolio{286}ภวราคานุสโย อวิชฺชานุสโย อิเม ตโย อนุสยา พฺยนฺตีโหนฺติ, เอวํ สญฺโญชนานิ ปหียนฺติ, อนุสยา พฺยนฺตีโหนฺติ.}

\proseitem{3}{อพฺยาปาทวเสน จิตฺตสฺส เอกคฺคตา อวิกฺเขโป สมาธิ ฯเปฯ.\csromanspacer{} อาโลกสญฺญาวเสน จิตฺตสฺส เอกคฺคตา อวิกฺเขโป สมาธิ ฯเปฯ.\csromanspacer{} ปฏินิสฺสคฺคานุปสฺสี อสฺสาสวเสน ปฏินิสฺสคฺคานุปสฺสี ปสฺสาสวเสน จิตฺตสฺส เอกคฺคตา อวิกฺเขโป สมาธิ.\csromanspacer{} ตตฺถ ชาเต ธมฺเม อนิจฺจโต อนุปสฺสนฏฺเฐน วิปสฺสนา, ทุกฺขโต อนุปสฺสนฏฺเฐน วิปสฺสนา, อนตฺตโต อนุปสฺสนฏฺเฐน วิปสฺสนา.\csromanspacer{} อิติ ปฐมํ สมโถ, ปจฺฉา วิปสฺสนา.\csromanspacer{} เตน วุจฺจติ “สมถปุพฺพงฺคมํ วิปสฺสนํ ภาเวตี”ติ.\csromanspacer{}\csromanspacer{}\textbf{ภาเวตี}ติ จตสฺโส ภาวนา.\csromanspacer{} ตตฺถ ชาตานํ ธมฺมานํ อนติวตฺตนฏฺเฐน ภาวนา, อินฺทฺริยานํ เอกรสฏฺเฐน ภาวนา, ตทุปควีริยวาหนฏฺเฐน ภาวนา, อาเสวนฏฺเฐน ภาวนา.}

\prose{\textbf{มคฺโค สญฺชายตี}ติ กถํ มคฺโค สญฺชายติ? ทสฺสนฏฺเฐน สมฺมาทิฏฺฐิ มคฺโค สญฺชายติ, อภินิโรปนฏฺเฐน สมฺมาสงฺกปฺโป มคฺโค สญฺชายติ ฯเปฯ อวิกฺเขปฏฺเฐน สมฺมาสมาธิ มคฺโค สญฺชายติ, เอวํ มคฺโค สญฺชายติ.}

\prose{\textbf{โส ตํ มคฺคํ อาเสวติ ภาเวติ พหุลีกโรติ อาเสวตี}ติ กถํ อาเสวติ? อาวชฺชนฺโต อาเสวติ ฯเปฯ สจฺฉิกาตพฺพํ สจฺฉิกโรนฺโต อาเสวติ, เอวํ อาเสวติ.\csromanspacer{}\csromanspacer{}\textbf{ภาเวตี}ติ กถํ ภาเวติ? อาวชฺชนฺโต ภาเวติ, ชานนฺโต ภาเวติ ฯเปฯ สจฺฉิกาตพฺพํ สจฺฉิกโรนฺโต ภาเวติ, เอวํ ภาเวติ.\csromanspacer{}\csromanspacer{}\textbf{พหุลีกโรตี}ติ กถํ พหุลีกโรติ? อาวชฺชนฺโต พหุลีกโรติ, ชานนฺโต พหุลีกโรติ ฯเปฯ สจฺฉิกาตพฺพํ สจฺฉิกโรนฺโต พหุลีกโรติ, เอวํ พหุลีกโรติ.}

\prose{\textbf{ตสฺส ตํ มคฺคํ อาเสวโต ภาวยโต พหุลีกโรโต สญฺโญชนานิ ปหียนฺติ, อนุสยา พฺยนฺตีโหนฺตี}ติ กถํ สญฺโญชนา ปหียนฺติ, อนุสยา พฺยนฺตีโหนฺติ? โสตาปตฺติมคฺเคน สกฺกายทิฏฺฐิ วิจิกิจฺฉา สีลพฺพตปรามาโส, อิมานิ ตีณิ สญฺโญชนานิ ปหียนฺติ, ทิฏฺฐานุสโย วิจิกิจฺฉานุสโย อิเม เทฺว อนุสยา พฺยนฺตีโหนฺติ.\csromanspacer{} สกทาคามิมคฺเคน โอฬาริกํ กามราคสญฺโญชนํ ปฏิฆสญฺโญชนํ อิมานิ เทฺว สญฺโญชนานิ ปหียนฺติ, โอฬาริโก กามราคานุสโย}

\vspace{\tipitakaparskip}
\csromancenter{ยุคนทฺธกถา ปฏิฆานุสโย อิเม เทฺว อนุสยา พฺยนฺตีโหนฺติ.\csromanspacer{} อนาคามิมคฺเคน อณุสหคตํ กามราคสญฺโญชนํ ปฏิฆสญฺโญชนํ อิมานิ เทฺว สญฺโญชนานิ ปหียนฺติ, อณุสหคโต กามราคานุสโย ปฏิฆานุสโย, อิเม เทฺว อนุสยา พฺยนฺตีโหนฺติ.\csromanspacer{} อรหตฺตมคฺเคน รูปราโค อรูปราโค มาโน อุทฺธจฺจํ อวิชฺชา อิมานิ ปญฺจ สญฺโญชนานิ ปหียนฺติ, มานานุสโย ภวราคานุสโย อวิชฺชานุสโย อิเม ตโย อนุสยา พฺยนฺตีโหนฺติ.\csromanspacer{} เอวํ สญฺโญชนานิ ปหียนฺติ, อนุสยา พฺยนฺตีโหนฺติ.\csromanspacer{} เอวํ สมถปุพฺพงฺคมํ วิปสฺสนํ ภาเวติ.}
\proseitem{4}{\csromanfolio{287}กถํ \textbf{วิปสฺสนาปุพฺพงฺคมํ สมถํ ภาเวติ?} อนิจฺจโต อนุปสฺสนฏฺเฐน วิปสฺสนา, ทุกฺขโต อนุปสฺสนฏฺเฐน วิปสฺสนา, อนตฺตโต อนุปสฺสนฏฺเฐน วิปสฺสนา.\csromanspacer{} ตตฺถ ชาตานํ ธมฺมานญฺจ โวสคฺคารมฺมณตา\footnote{โวสฺสคฺคารมฺมณตา (สฺยา, ก)} จิตฺตสฺส เอกคฺคตา อวิกฺเขโป สมาธิ.\csromanspacer{} อิติ ปฐมํ วิปสฺสนา, ปจฺฉา สมโถ.\csromanspacer{} เตน วุจฺจติ “วิปสฺสนาปุพฺพงฺคมํ สมถํ ภาเวตี”ติ.\csromanspacer{}\csromanspacer{}\textbf{ภาเวตี}ติ จตสฺโส ภาวนา, อาเสวนฏฺเฐน ภาวนา ฯเปฯ.\csromanspacer{} \textbf{มคฺโค สญฺชายตี}ติ กถํ มคฺโค สญฺชายติ ฯเปฯ เอวํ มคฺโค สญฺชายติ.\csromanspacer{} เอวํ สญฺโญชนานิ ปหียนฺติ, อนุสยา พฺยนฺตีโหนฺติ.}

\prose{รูปํ อนิจฺจโต อนุปสฺสนฏฺเฐน วิปสฺสนา, รูปํ ทุกฺขโต อนุปสฺสนฏฺเฐน วิปสฺสนา, รูปํ อนตฺตโต อนุปสฺสนฏฺเฐน วิปสฺสนา.\csromanspacer{} ตตฺถ ชาตานํ ธมฺมานญฺจ โวสคฺคารมฺมณตา จิตฺตสฺส เอกคฺคตา อวิกฺเขโป สมาธิ.\csromanspacer{} อิติ ปฐมํ วิปสฺสนา, ปจฺฉา สมโถ.\csromanspacer{} เตน วุจฺจติ “วิปสฺสนาปุพฺพงฺคมํ สมถํ ภาเวตี”ติ.\csromanspacer{} \textbf{ภาเวตี}ติ จตสฺโส ภาวนา, อาเสวนฏฺเฐน ภาวนา ฯเปฯ.\csromanspacer{} \textbf{มคฺโค สญฺชายตี}ติ กถํ มคฺโค สญฺชายติ ฯเปฯ เอวํ มคฺโค สญฺชายติ.\csromanspacer{} เอวํ สญฺโญชนานิ ปหียนฺติ, อนุสยา พฺยนฺตีโหนฺติ.}

\prose{เวทนํ ฯเปฯ.\csromanspacer{} สญฺญํ.\csromanspacer{} สงฺขาเร.\csromanspacer{} วิญฺญาณํ.\csromanspacer{} จกฺขุํ ฯเปฯ.\csromanspacer{} ชรามรณํ อนิจฺจโต อนุปสฺสนฏฺเฐน วิปสฺสนา, ชรามรณํ ทุกฺขโต ฯเปฯ.\csromanspacer{} อนตฺตโต อนุปสฺสนฏฺเฐน วิปสฺสนา.\csromanspacer{} ตตฺถ ชาตานํ ธมฺมานญฺจ โวสคฺคารมฺมณตา จิตฺตสฺส เอกคฺคตา อวิกฺเขโป สมาธิ.\csromanspacer{} อิติ ปฐมํ วิปสฺสนา, ปจฺฉา สมโถ.\csromanspacer{} เตน วุจฺจติ “วิปสฺสนาปุพฺพงฺคมํ สมถํ ภาเวตี”ติ.\csromanspacer{}\csromanspacer{}\textbf{ภาเวตี}ติ จตสฺโส ภาวนา, อาเสวนฏฺเฐน ภาวนา ฯเปฯ.\csromanspacer{} \textbf{มคฺโค สญฺชายตี}ติ \csromanfolio{288}กถํ มคฺโค สญฺชายติ ฯเปฯ เอวํ มคฺโค สญฺชายติ.\csromanspacer{} เอวํ สญฺโญชนานิ ปหียนฺติ, อนุสยา พฺยนฺตีโหนฺติ.\csromanspacer{} เอวํ วิปสฺสนาปุพฺพงฺคมํ สมถํ ภาเวติ.}

\proseitem{5}{กถํ \textbf{สมถวิปสฺสนํ ยุคนทฺธํ ภาเวติ?} โสฬสหิ อากาเรหิ สมถวิปสฺสนํ ยุคนทฺธํ ภาเวติ, อารมฺมณฏฺเฐน โคจรฏฺเฐน ปหานฏฺเฐน ปริจฺจาคฏฺเฐน วุฏฺฐานฏฺเฐน วิวฏฺฏนฏฺเฐน สนฺตฏฺเฐน ปณีตฏฺเฐน วิมุตฺตฏฺเฐน อนาสวฏฺเฐน ตรณฏฺเฐน อนิมิตฺตฏฺเฐน อปฺปณิหิตฏฺเฐน สุญฺญตฏฺเฐน เอกรสฏฺเฐน อนติวตฺตนฏฺเฐน ยุคนทฺธฏฺเฐน.}

\prose{กถํ \textbf{อารมฺมณฏฺเฐน สมถวิปสฺสนํ ยุคนทฺธํ ภาเวติ?} อุทฺธจฺจํ ปชหโต จิตฺตสฺส เอกคฺคตา อวิกฺเขโป สมาธิ นิโรธารมฺมโณ, อวิชฺชํ ปชหโต อนุปสฺสนฏฺเฐน วิปสฺสนา นิโรธารมฺมณา.\csromanspacer{} อิติ อารมฺมณฏฺเฐน สมถวิปสฺสนา เอกรสา โหนฺติ, ยุคนทฺธา โหนฺติ, อญฺญมญฺญํ นาติวตฺตนฺตีติ.\csromanspacer{} เตน วุจฺจติ “อารมฺมณฏฺเฐน สมถวิปสฺสนํ ยุคนทฺธํ ภาเวตี”ติ.\csromanspacer{} \textbf{ภาเวตี}ติ จตสฺโส ภาวนา.\csromanspacer{} อาเสวนฏฺเฐน ภาวนา ฯเปฯ.\csromanspacer{} \textbf{มคฺโค สญฺชายตี}ติ กถํ มคฺโค สญฺชายติ ฯเปฯ เอวํ มคฺโค สญฺชายติ.\csromanspacer{} เอวํ สญฺโญชนานิ ปหียนฺติ, อนุสยา พฺยนฺตีโหนฺติ.\csromanspacer{} เอวํ อารมฺมณฏฺเฐน สมถวิปสฺสนํ ยุคนทฺธํ ภาเวติ. (1)}

\prose{กถํ \textbf{โคจรฏฺเฐน สมถวิปสฺสนํ ยุคนทฺธํ ภาเวติ?} อุทฺธจฺจํ ปชหโต จิตฺตสฺส เอกคฺคตา อวิกฺเขโป สมาธิ นิโรธโคจโร, อวิชฺชํ ปชหโต อนุปสฺสนฏฺเฐน วิปสฺสนา นิโรธโคจรา.\csromanspacer{} อิติ โคจรฏฺเฐน สมถวิปสฺสนา เอกรสา โหนฺติ, ยุคนทฺธา โหนฺติ, อญฺญมญฺญํ นาติวตฺตนฺตีติ.\csromanspacer{} เตน วุจฺจติ “โคจรฏฺเฐน สมถวิปสฺสนํ ยุคนทฺธํ ภาเวตี”ติ. (2)}

\prose{กถํ \textbf{ปหานฏฺเฐน สมถวิปสฺสนํ ยุคนทฺธํ ภาเวติ?} อุทฺธจฺจสหคตกิเลเส จ ขนฺเธ จ ปชหโต จิตฺตสฺส เอกคฺคตา อวิกฺเขโป สมาธิ นิโรธโคจโร, อวิชฺชาสหคตกิเลเส จ ขนฺเธ จ ปชหโต อนุปสฺสนฏฺเฐน วิปสฺสนา นิโรธโคจรา.\csromanspacer{} อิติ ปหานฏฺเฐน สมถวิปสฺสนา เอกรสา โหนฺติ, ยุคนทฺธา โหนฺติ, อญฺญมญฺญํ นาติวตฺตนฺตีติ.\csromanspacer{} เตน วุจฺจติ “ปหานฏฺเฐน สมถวิปสฺสนํ ยุคนทฺธํ ภาเวตี”ติ. (3)}

\prose{กถํ \textbf{ปริจฺจาคฏฺเฐน สมถวิปสฺสนํ ยุคนทฺธํ ภาเวติ?} อุทฺธจฺจสหคตกิเลเส จ ขนฺเธ จ ปริจฺจชโต จิตฺตสฺส เอกคฺคตา อวิกฺเขโป สมาธิ นิโรธโคจโร, อวิชฺชาสหคตกิเลเส จ ขนฺเธ จ ปริจฺจชโต}

\vspace{\tipitakaparskip}
\csromancenter{ยุคนทฺธกถา อนุปสฺสนฏฺเฐน วิปสฺสนา นิโรธโคจรา.\csromanspacer{} อิติ ปริจฺจาคฏฺเฐน สมถวิปสฺสนา เอกรสา โหนฺติ, ยุคนทฺธา โหนฺติ, อญฺญมญฺญํ นาติวตฺตนฺตีติ.\csromanspacer{} เตน วุจฺจติ “ปริจฺจาคฏฺเฐน สมถวิปสฺสนํ ยุคนทฺธํ ภาเวตี”ติ. (4)}
\prose{\csromanfolio{289}กถํ \textbf{วุฏฺฐานฏฺเฐน สมถวิปสฺสนํ ยุคนทฺธํ ภาเวติ?} อุทฺธจฺจสหคตกิเลเสหิ จ ขนฺเธหิ จ วุฏฺฐหโต จิตฺตสฺส เอกคฺคตา อวิกฺเขโป สมาธิ นิโรธโคจโร, อวิชฺชาสหคตกิเลเสหิ จ ขนฺเธหิ จ วุฏฺฐหโต อนุปสฺสนฏฺเฐน วิปสฺสนา นิโรธโคจรา.\csromanspacer{} อิติ วุฏฺฐานฏฺเฐน สมถวิปสฺสนา เอกรสา โหนฺติ, ยุคนทฺธา โหนฺติ, อญฺญมญฺญํ นาติวตฺตนฺตีติ.\csromanspacer{} เตน วุจฺจติ “วุฏฺฐานฏฺเฐน สมถวิปสฺสนํ ยุคนทฺธํ ภาเวตี”ติ. (5)}

\prose{กถํ \textbf{วิวฏฺฏนฏฺเฐน สมถวิปสฺสนํ ยุคนทฺธํ ภาเวติ?} อุทฺธจฺจสหคตกิเลเสหิ จ ขนฺเธหิ จ วิวฏฺฏโต จิตฺตสฺส เอกคฺคตา อวิกฺเขโป สมาธิ นิโรธโคจโร, อวิชฺชาสหคตกิเลเสหิ จ ขนฺเธหิ จ วิวฏฺฏโต อนุปสฺสนฏฺเฐน วิปสฺสนา นิโรธโคจรา.\csromanspacer{} อิติ วิวฏฺฏนฏฺเฐน สมถวิปสฺสนา เอกรสา โหนฺติ, ยุคนทฺธา โหนฺติ, อญฺญมญฺญํ นาติวตฺตนฺตีติ.\csromanspacer{} เตน วุจฺจติ “วิวฏฺฏนฏฺเฐน สมถวิปสฺสนํ ยุคนทฺธํ ภาเวตี”ติ. (6)}

\prose{กถํ \textbf{สนฺตฏฺเฐน สมถวิปสฺสนํ ยุคนทฺธํ ภาเวติ?} อุทฺธจฺจํ ปชหโต จิตฺตสฺส เอกคฺคตา อวิกฺเขโป สมาธิ สนฺโต โหติ นิโรธโคจโร, อวิชฺชํ ปชหโต อนุปสฺสนฏฺเฐน วิปสฺสนา สนฺตา โหติ นิโรธโคจรา.\csromanspacer{} อิติ สนฺตฏฺเฐน สมถวิปสฺสนา เอกรสา โหนฺติ, ยุคนทฺธา โหนฺติ, อญฺญมญฺญํ นาติวตฺตนฺตีติ.\csromanspacer{} เตน วุจฺจติ “สนฺตฏฺเฐน สมถวิปสฺสนํ ยุคนทฺธํ ภาเวตี”ติ. (7)}

\prose{กถํ \textbf{ปณีตฏฺเฐน สมถวิปสฺสนํ ยุคนทฺธํ ภาเวติ?} อุทฺธจฺจํ ปชหโต จิตฺตสฺส เอกคฺคตา อวิกฺเขโป สมาธิ ปณีโต โหติ นิโรธโคจโร, อวิชฺชํ ปชหโต อนุปสฺสนฏฺเฐน วิปสฺสนา ปณีตา โหติ นิโรธโคจรา.\csromanspacer{} อิติ ปณีตฏฺเฐน สมถวิปสฺสนา เอกรสา โหนฺติ, ยุคนทฺธา โหนฺติ, อญฺญมญฺญํ นาติวตฺตนฺตีติ.\csromanspacer{} เตน วุจฺจติ “ปณีตฏฺเฐน สมถวิปสฺสนํ ยุคนทฺธํ ภาเวตี”ติ. (8)}

\prose{กถํ \textbf{วิมุตฺตฏฺเฐน สมถวิปสฺสนํ ยุคนทฺธํ ภาเวติ?} อุทฺธจฺจํ ปชหโต จิตฺตสฺส เอกคฺคตา อวิกฺเขโป สมาธิ กามาสวา วิมุตฺโต โหติ นิโรธโคจโร, อวิชฺชํ ปชหโต อนุปสฺสนฏฺเฐน วิปสฺสนา อวิชฺชาสวา \csromanfolio{290}วิมุตฺตา โหติ นิโรธโคจรา.\csromanspacer{} อิติ ราควิราคา เจโตวิมุตฺติ, อวิชฺชาวิราคา ปญฺญาวิมุตฺติ วิมุตฺตฏฺเฐน สมถวิปสฺสนา เอกรสา โหนฺติ, ยุคนทฺธา โหนฺติ, อญฺญมญฺญํ นาติวตฺตนฺตีติ.\csromanspacer{} เตน วุจฺจติ “วิมุตฺตฏฺเฐน สมถวิปสฺสนํ ยุคนทฺธํ พฺ\textbf{หาเวตี}”ติ. (9)}

\prose{กถํ \textbf{อนาสวฏฺเฐน สมถวิปสฺสนํ ยุคนทฺธํ ภาเวติ?} อุทฺธจฺจํ ปชหโต จิตฺตสฺส เอกคฺคตา อวิกฺเขโป สมาธิ กามาสเวน อนาสโว โหติ นิโรธโคจโร, อวิชฺชํ ปชหโต อนุปสฺสนฏฺเฐน วิปสฺสนา อวิชฺชาสเวน อนาสวา โหติ นิโรธโคจรา.\csromanspacer{} อิติ อนาสวฏฺเฐน สมถวิปสฺสนา เอกรสา โหนฺติ, ยุคนทฺธา โหนฺติ, อญฺญมญฺญํ นาติวตฺตนฺตีติ.\csromanspacer{} เตน วุจฺจติ “อนาสวฏฺเฐน สมถวิปสฺสนํ ยุคนทฺธํ ภาเวตี”ติ. (10)}

\prose{กถํ \textbf{ตรณฏฺเฐน สมถวิปสฺสนํ ยุคนทฺธํ ภาเวติ?} อุทฺธจฺจสหคตกิเลเส จ ขนฺเธ จ ตรโต จิตฺตสฺส เอกคฺคตา อวิกฺเขโป สมาธิ นิโรธโคจโร, อวิชฺชาสหคตกิเลเส จ ขนฺเธ จ ตรโต อนุปสฺสนฏฺเฐน วิปสฺสนา นิโรธโคจรา.\csromanspacer{} อิติ ตรณฏฺเฐน สมถวิปสฺสนา เอกรสา โหนฺติ, ยุคนทฺธา โหนฺติ, อญฺญมญฺญํ นาติวตฺตนฺตีติ.\csromanspacer{} เตน วุจฺจติ “ตรณฏฺเฐน สมถวิปสฺสนํ ยุคนทฺธํ ภาเวตี”ติ. (11)}

\prose{กถํ \textbf{อนิมิตฺตฏฺเฐน สมถวิปสฺสนํ ยุคนทฺธํ ภาเวติ?} อุทฺธจฺจํ ปชหโต จิตฺตสฺส เอกคฺคตา อวิกฺเขโป สมาธิ สพฺพนิมิตฺเตหิ อนิมิตฺโต โหติ นิโรธโคจโร, อวิชฺชํ ปชหโต อนุปสฺสนฏฺเฐน วิปสฺสนา สพฺพนิมิตฺเตหิ อนิมิตฺตา โหติ นิโรธโคจรา.\csromanspacer{} อิติ อนิมิตฺตฏฺเฐน สมถวิปสฺสนา เอกรสา โหนฺติ, ยุคนทฺธา โหนฺติ, อญฺญมญฺญํ นาติวตฺตนฺตีติ.\csromanspacer{} เตน วุจฺจติ “อนิมิตฺตฏฺเฐน สมถวิปสฺสนํ ยุคนทฺธํ ภาเวตี”ติ. (12)}

\prose{กถํ \textbf{อปฺปณิหิตฏฺเฐน สมถวิปสฺสนํ ยุคนทฺธํ ภาเวติ?} อุทฺธจฺจํ ปชหโต จิตฺตสฺส เอกคฺคตา อวิกฺเขโป สมาธิ สพฺพปณิธีหิ อปฺปณิหิโต โหติ นิโรธโคจโร, อวิชฺชํ ปชหโต อนุปสฺสนฏฺเฐน วิปสฺสนา สพฺพปณิธีหิ อปฺปณิหิตา โหติ นิโรธโคจรา.\csromanspacer{} อิติ อปฺปณิหิตฏฺเฐน สมถวิปสฺสนา เอกรสา โหนฺติ, ยุคนทฺธา โหนฺติ, อญฺญมญฺญํ นาติวตฺตนฺตีติ.\csromanspacer{} เตน วุจฺจติ “อปฺปณิหิตฏฺเฐน สมถวิปสฺสนํ ยุคนทฺธํ ภาเวตี”ติ. (13)}

\csromanheader{2}{1. ยุคนทฺธกถา}
\prose{\csromanfolio{291}กถํ \textbf{สุญฺญตฏฺเฐน สมถวิปสฺสนํ ยุคนทฺธํ ภาเวติ?} อุทฺธจฺจํ ปชหโต จิตฺตสฺส เอกคฺคตา อวิกฺเขโป สมาธิ สพฺพาภินิเวเสหิ สุญฺโญ โหติ นิโรธโคจโร, อวิชฺชํ ปชหโต อนุปสฺสนฏฺเฐน วิปสฺสนา สพฺพาภินิเวเสหิ สุญฺญา โหติ นิโรธโคจรา.\csromanspacer{} อิติ สุญฺญตฏฺเฐน สมถวิปสฺสนา เอกรสา โหนฺติ, ยุคนทฺธา โหนฺติ, อญฺญมญฺญํ นาติวตฺตนฺตีติ.\csromanspacer{} เตน วุจฺจติ “สุญฺญตฏฺเฐน สมถวิปสฺสนํ ยุคนทฺธํ พฺ\textbf{หาเวตี}”ติ.\csromanspacer{} พฺ\textbf{หาเวตี}ติ จตสฺโส ภาวนา, อาเสวนฏฺเฐน ภาวนา ฯเปฯ.\csromanspacer{} มฺ\textbf{อคฺโค สญฺชายตี}ติ กถํ มคฺโค สญฺชายติ ฯเปฯ เอวํ มคฺโค สญฺชายติ.\csromanspacer{} เอวํ สญฺโญชนานิ ปหียนฺติ, อนุสยา พฺยนฺตีโหนฺติ.\csromanspacer{} เอวํ สุญฺญตฏฺเฐน สมถวิปสฺสนํ ยุคนทฺธํ ภาเวติ.\csromanspacer{} อิเมหิ โสฬสหิ อากาเรหิ สมถวิปสฺสนํ ยุคนทฺธํ ภาเวติ, เอวํ สมถวิปสฺสนํ ยุคนทฺธํ ภาเวติ. (14) สุตฺตนฺตนิทฺเทโส.}
\csromansectionrule

\csromanmatikaanchor{mtk.87}
\csromantocmarkref{subsection}{2. ธมฺมุทฺธจฺจวารนิทฺเทส}{mtk.87}
\bookmark[dest={mtk.87},level=2]{2. ธมฺมุทฺธจฺจวารนิทฺเทส}
\csromanheader{2}{2. ธมฺมุทฺธจฺจวารนิทฺเทส}
\proseitem{6}{กถํ ธมฺมฺ\textbf{อุทฺธจฺจวิคฺคหิตํ มานสํ โหติ?} อนิจฺจโต มนสิกโรโต โอภาโส อุปฺปชฺชติ, “โอภาโส ธมฺโม”ติ โอภาสํ อาวชฺชติ.\csromanspacer{} ตโต วิกฺเขโป อุทฺธจฺจํ.\csromanspacer{} เตน อุทฺธจฺเจน วิคฺคหิตมานโส อนิจฺจโต อุปฏฺฐานํ ยถาภูตํ นปฺปชานาติ, ทุกฺขโต อุปฏฺฐานํ ยถาภูตํ นปฺปชานาติ, อนตฺตโต อุปฏฺฐานํ ยถาภูตํ นปฺปชานาติ, เตน วุจฺจติ ธมฺมุทฺธจฺจวิคฺคหิตมานโส.\csromanspacer{} โหติ โส สมโย, ยํ ตํ จิตฺตํ อชฺฌตฺตเมว สนฺติฏฺฐติ สนฺนิสีทติ เอโกทิ โหติ สมาธิยติ.\csromanspacer{} \textbf{ตสฺส มคฺโค สญฺชายตี}ติ กถํ มคฺโค สญฺชายติ ฯเปฯ เอวํ มคฺโค สญฺชายติ, เอวํ สญฺโญชนานิ ปหียนฺติ, อนุสยา พฺยนฺตีโหนฺติ.}

\prose{อนิจฺจโต มนสิกโรโต ญาณํ อุปฺปชฺชติ, ปีติ อุปฺปชฺชติ, ปสฺสทฺธิ อุปฺปชฺชติ, สุขํ อุปฺปชฺชติ, อธิโมกฺโข อุปฺปชฺชติ, ปคฺคโห อุปฺปชฺชติ, อุปฏฺฐานํ อุปฺปชฺชติ, อุเปกฺขา อุปฺปชฺชติ, นิกนฺติ อุปฺปชฺชติ, “นิกนฺติ ธมฺโม”ติ นิกนฺติํ อาวชฺชติ.\csromanspacer{} ตโต วิกฺเขโป อุทฺธจฺจํ.\csromanspacer{} เตน อุทฺธจฺเจน วิคฺคหิตมานโส อนิจฺจโต อุปฏฺฐานํ ยถาภูตํ นปฺปชานาติ, ทุกฺขโต อุปฏฺฐานํ ยถาภูตํ นปฺปชานาติ, อนตฺตโต อุปฏฺฐานํ ยถาภูตํ นปฺปชานาติ, เตน วุจฺจติ \csromanfolio{292}ธมฺมุทฺธจฺจวิคฺคหิตมานโส.\csromanspacer{} โหติ โส สมโย, ยํ ตํ จิตฺตํ อชฺฌตฺตเมว สนฺติฏฺฐติ สนฺนิสีทติ เอโกทิ โหติ สมาธิยติ.\csromanspacer{} \textbf{ตสฺส มคฺโค สญฺชายตี}ติ กถํ มคฺโค สญฺชายติ ฯเปฯ เอวํ มคฺโค สญฺชายติ, เอวํ สญฺโญชนานิ ปหียนฺติ, อนุสยา พฺยนฺตีโหนฺติ.}

\prose{ทุกฺขโต มนสิกโรโต ฯเปฯ.\csromanspacer{} อนตฺตโต มนสิกโรโต โอภาโส อุปฺปชฺชติ ฯเปฯ ญาณํ อุปฺปชฺชติ, ปีติ อุปฺปชฺชติ, ปสฺสทฺธิ อุปฺปชฺชติ, สุขํ อุปฺปชฺชติ, อธิโมกฺโข อุปฺปชฺชติ, ปคฺคโห อุปฺปชฺชติ, อุปฏฺฐานํ อุปฺปชฺชติ, อุเปกฺขา อุปฺปชฺชติ, นิกนฺติ อุปฺปชฺชติ, นิกนฺติ ธมฺโมติ นิกนฺติํ อาวชฺชติ.\csromanspacer{} ตโต วิกฺเขโป อุทฺธจฺจํ.\csromanspacer{} เตน อุทฺธจฺเจน วิคฺคหิตมานโส อนตฺตโต อุปฏฺฐานํ, อนิจฺจโต อุปฏฺฐานํ, ทุกฺขโต อุปฏฺฐานํ ยถาภูตํ นปฺปชานาติ, เตน วุจฺจติ ธมฺมุทฺธจฺจวิคฺคหิตมานโส ฯเปฯ เอวํ สญฺโญชนานิ ปหียนฺติ, อนุสยา พฺยนฺตีโหนฺติ.}

\prose{รูปํ อนิจฺจโต มนสิกโรโต ฯเปฯ รูปํ ทุกฺขโต มนสิกโรโต.\csromanspacer{} รูปํ อนตฺตโต มนสิกโรโต.\csromanspacer{} เวทนํ ฯเปฯ.\csromanspacer{} สญฺญํ.\csromanspacer{} สงฺขาเร, วิญฺญาณํ.\csromanspacer{} จกฺขุํ ฯเปฯ.\csromanspacer{} ชรามรณํ อนิจฺจโต มนสิกโรโต ฯเปฯ ชรามรณํ ทุกฺขโต มนสิกโรโต.\csromanspacer{} ชรามรณํ อนตฺตโต มนสิกโรโต โอภาโส อุปฺปชฺชติ ฯเปฯ ญาณํ อุปฺปชฺชติ, ปีติ อุปฺปชฺชติ, ปสฺสทฺธิ อุปฺปชฺชติ, สุขํ อุปฺปชฺชติ, อธิโมกฺโข อุปฺปชฺชติ, ปคฺคโห อุปฺปชฺชติ, อุปฏฺฐานํ อุปฺปชฺชติ, อุเปกฺโข อุปฺปชฺชติ, นิกนฺติ อุปฺปชฺชติ, “นิกนฺติ ธมฺโม”ติ นิกนฺติํ อาวชฺชติ.\csromanspacer{} ตโต วิกฺเขโป อุทฺธจฺจํ.\csromanspacer{} เตน อุทฺธจฺเจน วิคฺคหิตมานโส.\csromanspacer{} ชรามรณํ อนตฺตโต อุปฏฺฐานํ ยถาภูตํ นปฺปชานาติ.\csromanspacer{} ชรามรณํ อนิจฺจโต อุปฏฺฐานํ ยถาภูตํ นปฺปชานาติ, ชรามรณํ ทุกฺขโต อุปฏฺฐานํ ยถาภูตํ นปฺปชานาติ.\csromanspacer{} เตน วุจฺจติ ธมฺมุทฺธจฺจวิคฺคหิตมานโส.\csromanspacer{} โหติ โส สมโย, ยํ ตํ จิตฺตํ อชฺฌตฺตเมว สนฺติฏฺฐติ สนฺนิสีทติ เอโกทิ โหติ สมาธิยติ.\csromanspacer{} \textbf{ตสฺส มคฺโค สญฺชายตี}ติ กถํ มคฺโค สญฺชายติ ฯเปฯ เอวํ มคฺโค สญฺชายติ.\csromanspacer{} เอวํ สญฺโญชนานิ ปหียนฺติ, อนุสยา พฺยนฺตีโหนฺติ.\csromanspacer{} เอวํ ธมฺมุทฺธจฺจวิคฺคหิตํ มานสํ โหติ.}

\proseitem{7}{โอภาเส เจว ญาเณ จ, ปีติยา จ วิกมฺปติ.}

\csromangathameasure{ปสฺสทฺธิยา สุเข เจว, เยหิ จิตฺตํ ปเวธติ.}
\csromangathagroup{\csromangathastanza{ปสฺสทฺธิยา สุเข เจว, เยหิ จิตฺตํ ปเวธติ.}}

\csromanmatikaanchor{mtk.88}
\csromantocmarknopage{section}{2. สจฺจกถา}{mtk.88}
\bookmark[dest={mtk.88},level=1]{2. สจฺจกถา}
\csromanheader{1}{2. สจฺจกถา}
\csromangathameasure{อธิโมกฺเข จ ปคฺคาเห, อุปฏฺฐาเน จ กมฺปติ. \\ อุเปกฺขาวชฺชนาย เจว, อุเปกฺขาย จ นิกนฺติยา. \\ อิมานิ ทส ฐานานิ, ปญฺญา ยสฺส ปริจฺจิตา. \\ ธมฺมุทฺธจฺจกุสโล โหติ, น จ สมฺโมห คจฺฉติ. \\ วิกฺขิปฺปติ เจว กิลิสฺสติ จ, จวติ จิตฺตภาวนา. \\ วิกฺขิปติ น กิลิสฺสติ, ภาวนา ปริหายติ. \\ วิกฺขิปติ น กิลิสฺสติ, ภาวนา น ปริหายติ. \\ น จ วิกฺขิปเต จิตฺตํ น กิลิสฺสติ, น จวติ จิตฺตภาวนา.}
\csromangathagroup{\csromangathastanza{\csromanfolio{293}อธิโมกฺเข จ ปคฺคาเห, อุปฏฺฐาเน จ กมฺปติ. \\ อุเปกฺขาวชฺชนาย เจว, อุเปกฺขาย จ นิกนฺติยา.}\csromangathastanza{อิมานิ ทส ฐานานิ, ปญฺญา ยสฺส ปริจฺจิตา. \\ ธมฺมุทฺธจฺจกุสโล โหติ, น จ สมฺโมห คจฺฉติ.}\csromangathastanza{วิกฺขิปฺปติ เจว กิลิสฺสติ จ, จวติ จิตฺตภาวนา. \\ วิกฺขิปติ น กิลิสฺสติ, ภาวนา ปริหายติ.}\csromangathastanza{วิกฺขิปติ น กิลิสฺสติ, ภาวนา น ปริหายติ. \\ น จ วิกฺขิปเต จิตฺตํ น กิลิสฺสติ, น จวติ จิตฺตภาวนา.}}

\prose{อิเมหิ จตูหิ ฐาเนหิ จิตฺตสฺส สงฺเขปวิกฺเขปวิคฺคหิตํ\footnote{สงฺเขปํ วิกฺเขปํ วิคฺคหิตํ (สฺยา)} ทส ฐาเน สมฺปชานาตีติ.}

\nitthitamruled{ยุคนทฺธกถา นิฏฺฐิตา.}
\csromanheader{2}{2. สจฺจกถา}
\proseitem{8}{ปุริมนิทานํ.\csromanspacer{}\csromanspacer{}จตฺตาริมานิ ภิกฺขเว\footnote{สํ 3. 377 ปิฏฺเฐปิ.} ตถานิ อวิตถานิ อนญฺญถานิ.\csromanspacer{} กตมานิ จตฺตาริ\csromandash{}“อิทํ ทุกฺขนฺ”ติ ภิกฺขเว ตถเมตํ อวิตถเมตํ อนญฺญถเมตํ, “อยํ ทุกฺขสมุทโย”ติ ตถเมตํ อวิตถเมตํ อนญฺญถเมตํ, “อยํ ทุกฺขนิโรโธ”ติ ตถเมตํ อวิตถเมตํ อนญฺญถเมตํ, “อยํ ทุกฺขนิโรธคามินี ปฏิปทา”ติ ตถเมตํ อวิตถเมตํ อนญฺญถเมตํ.\csromanspacer{} อิมานิ โข ภิกฺขเว จตฺตาริ ตถานิ อวิตถานิ อนญฺญถานิ.}

\csromanmatikaanchor{mtk.89}
\csromantocmarkref{subsection}{1. ปฐมสุตฺตนฺตนิทฺเทส}{mtk.89}
\bookmark[dest={mtk.89},level=2]{1. ปฐมสุตฺตนฺตนิทฺเทส}
\csromanheader{2}{1. ปฐมสุตฺตนฺตนิทฺเทส}
\prose{กถํ ทุกฺขํ ตถฏฺเฐน สจฺจํ? จตฺตาโร ทุกฺขสฺส ทุกฺขฏฺฐา ตถา อวิตถา อนญฺญถา, ทุกฺขสฺส ปีฬนฏฺโฐ สงฺขตฏฺโฐ สนฺตาปฏฺโฐ วิปริณามฏฺโฐ, อิเม จตฺตาโร ทุกฺขสฺส ทุกฺขฏฺฐา ตถา อวิตถา อนญฺญถา, เอวํ ทุกฺขํ ตถฏฺเฐน สจฺจํ.}

\prose{กถํ สมุทโย ตถฏฺเฐน สจฺจํ? จตฺตาโร สมุทยสฺส สมุทยฏฺฐา ตถา อวิตถา อนญฺญถา, สมุทยสฺส อายูหนฏฺโฐ นิทานฏฺโฐ}

\prose{\csromanfolio{294}สํโยคฏฺโฐ ปลิโพธฏฺโฐ, อิเม จตฺตาโร สมุทยสฺส สมุทยฏฺฐา ตถา อวิตถา อนญฺญถา, เอวํ สมุทโย ตถฏฺเฐน สจฺจํ.}

\prose{กถํ นิโรโธ ตถฏฺเฐน สจฺจํ? จตฺตาโร นิโรธสฺส นิโรธฏฺฐา ตถา อวิตถา อนญฺญถา, นิโรธสฺส นิสฺสรณฏฺโฐ วิเวกฏฺโฐ อสงฺขตฏฺโฐ อมตฏฺโฐ, อิเม จตฺตาโร นิโรธสฺส นิโรธฏฺฐา ตถา อวิตถา อนญฺญถา, เอวํ นิโรโธ ตถฏฺเฐน สจฺจํ.}

\prose{กถํ มคฺโค ตถฏฺเฐน สจฺจํ? จตฺตาโร มคฺคสฺส มคฺคฏฺฐา ตถา อวิตถา อนญฺญถา, มคฺคสฺส นิยฺยานฏฺโฐ เหตุฏฺโฐ ทสฺสนฏฺโฐ อาธิปเตยฺยฏฺโฐ, อิเม จตฺตาโร มคฺคสฺส มคฺคฏฺฐา ตถา อวิตถา อนญฺญถา, เอวํ มคฺโค ตถฏฺเฐน สจฺจํ.}

\proseitem{9}{กติหากาเรหิ จตฺตาริ สจฺจานิ เอกปฺปฏิเวธานิ? จตูหากาเรหิ จตฺตาริ สจฺจานิ เอกปฺปฏิเวธานิ ตถฏฺเฐน อนตฺตฏฺเฐน สจฺจฏฺเฐน ปฏิเวธฏฺเฐน, อิเมหิ จตูหากาเรหิ จตฺตาริ สจฺจานิ เอกสงฺคหิตานิ, ยํ เอกสงฺคหิตํ ตํ เอกตฺตํ, เอกตฺตํ เอเกน ญาเณน ปฏิวิชฺฌตีติ จตฺตาริ สจฺจานิ เอกปฺปฏิเวธานิ.}

\prose{กถํ ตถฏฺเฐน จตฺตาริ สจฺจานิ เอกปฺปฏิเวธานิ? จตูหากาเรหิ ตถฏฺเฐน จตฺตาริ สจฺจานิ เอกปฺปฏิเวธานิ ทุกฺขสฺส ทุกฺขฏฺโฐ ตถฏฺโฐ, สมุทยสฺส สมุทยฏฺโฐ ตถฏฺโฐ, นิโรธสฺส นิโรธฏฺโฐ ตถฏฺโฐ, มคฺคสฺส มคฺคฏฺโฐ ตถฏฺโฐ, อิเมหิ จตูหากาเรหิ ตถฏฺเฐน จตฺตาริ สจฺจานิ เอกสงฺคหิตานิ, ยํ เอกสงฺคหิตํ ตํ เอกตฺตํ, เอกตฺตํ เอเกน ญาเณน ปฏิวิชฺฌตีติ จตฺตาริ สจฺจานิ เอกปฺปฏิเวธานิ.}

\prose{กถํ อนตฺตฏฺเฐน จตฺตาริ สจฺจานิ เอกปฺปฏิเวธานิ? จตูหากาเรหิ อนตฺตฏฺเฐน จตฺตาริ สจฺจานิ เอกปฺปฏิเวธานิ, ทุกฺขสฺส ทุกฺขฏฺโฐ อนตฺตฏฺโฐ, สมุทยสฺส สมุทยฏฺโฐ อนตฺตฏฺโฐ, นิโรธสฺส นิโรธฏฺโฐ อนตฺตฏฺโฐ, มคฺคสฺส มคฺคฏฺโฐ อนตฺตฏฺโฐ, อิเมหิ จตูหากาเรหิ อนตฺตฏฺเฐน จตฺตาริ สจฺจานิ เอกสงฺคหิตานิ, ยํ เอกสงฺคหิตํ ตํ เอกตฺตํ, เอกตฺตํ เอเกน ญาเณน ปฏิวิชฺฌตีติ จตฺตาริ สจฺจานิ เอกปฺปฏิเวธานิ.}

\prose{กถํ สจฺจฏฺเฐน จตฺตาริ สจฺจานิ เอกปฺปฏิเวธานิ? จตูหากาเรหิ สจฺจฏฺเฐน จตฺตาริ สจฺจานิ เอกปฺปฏิเวธานิ.\csromanspacer{} ทุกฺขสฺส ทุกฺขฏฺโฐ สจฺจฏฺโฐ.}

\vspace{\tipitakaparskip}
\csromancenter{สจฺจกถา สมุทยสฺส สมุทยฏฺโฐ สจฺจฏฺโฐ.\csromanspacer{} นิโรธสฺส นิโรธฏฺโฐ สจฺจฏฺโฐ.\csromanspacer{} มคฺคสฺส มคฺคฏฺโฐ สจฺจฏฺโฐ.\csromanspacer{} อิเมหิ จตูหากาเรหิ สจฺจฏฺเฐน จตฺตาริ สจฺจานิ เอกสงฺคหิตานิ.\csromanspacer{} ยํ เอกสงฺคหิตํ ตํ เอกตฺตํ, เอกตฺตํ เอเกน ญาเณน ปฏิวิชฺฌตีติ จตฺตาริ สจฺจานิ เอกปฺปฏิเวธานิ.}
\prose{\csromanfolio{295}กถํ ปฏิเวธฏฺเฐน จตฺตาริ สจฺจานิ เอกปฺปฏิเวธานิ? จตูหากาเรหิ ปฏิเวธฏฺเฐน จตฺตาริ สจฺจานิ เอกปฺปฏิเวธานิ.\csromanspacer{} ทุกฺขสฺส ทุกฺขฏฺโฐ ปฏิเวธฏฺโฐ สมุทยสฺส สมุทยฏฺโฐ ปฏิเวธฏฺโฐ.\csromanspacer{} นิโรธสฺส นิโรธฏฺโฐ ปฏิเวธฏฺโฐ.\csromanspacer{} มคฺคสฺส มคฺคฏฺโฐ ปฏิเวธฏฺโฐ.\csromanspacer{} อิเมหิ จตูหากาเรหิ ปฏิเวธฏฺเฐน จตฺตาริ สจฺจานิ เอกสงฺคหิตานิ.\csromanspacer{} ยํ เอกสงฺคหิตํ ตํ เอกตฺตํ, เอกตฺตํ เอเกน ญาเณน ปฏิวิชฺฌตีติ จตฺตาริ สจฺจานิ เอกปฺปฏิเวธานิ.}

\proseitem{10}{กติหากาเรหิ จตฺตาริ สจฺจานิ เอกปฺปฏิเวธานิ? ยํ อนิจฺจํ, ตํ ทุกฺขํ.\csromanspacer{} ยํ ทุกฺขํ, ตํ อนิจฺจํ.\csromanspacer{} ยํ อนิจฺจญฺจ ทุกฺขญฺจ, ตํ อนตฺตา.\csromanspacer{} ยํ อนิจฺจญฺจ ทุกฺขญฺจ อนตฺตา จ, ตํ ตถํ.\csromanspacer{} ยํ อนิจฺจญฺจ ทุกฺขญฺจ อนตฺตา จ ตถญฺจ, ตํ สจฺจํ.\csromanspacer{} ยํ อนิจฺจญฺจ ทุกฺขญฺจ อนตฺตา จ ตถญฺจ สจฺจญฺจ, ตํ เอกสงฺคหิตํ.\csromanspacer{} ยํ เอกสงฺคหิตํ ตํ เอกตฺตํ, เอกตฺตํ เอเกน ญาเณน ปฏิวิชฺฌตีติ จตฺตาริ สจฺจานิ เอกปฺปฏิเวธานิ.}

\prose{กติหากาเรหิ จตฺตาริ สจฺจานิ เอกปฺปฏิเวธานิ? นวหากาเรหิ จตฺตาริ สจฺจานิ เอกปฺปฏิเวธานิ, ตถฏฺเฐน อนตฺตฏฺเฐน สจฺจฏฺเฐน ปฏิเวธฏฺเฐน อภิญฺญฏฺเฐน ปริญฺญฏฺเฐน ปหานฏฺเฐน ภาวนฏฺเฐน สจฺฉิกิริยฏฺเฐน.\csromanspacer{} อิเมหิ นวหากาเรหิ จตฺตาริ สจฺจานิ เอกสงฺคหิตานิ.\csromanspacer{} ยํ เอกสงฺคหิตํ ตํ เอกตฺตํ, เอกตฺตํ เอเกน ญาเณน ปฏิวิชฺฌตีติ จตฺตาริ สจฺจานิ เอกปฺปฏิเวธานิ.}

\prose{กถํ ตถฏฺเฐน จตฺตาริ สจฺจานิ เอกปฺปฏิเวธา? นวหากาเรหิ ตถฏฺเฐน จตฺตาริ สจฺจานิ เอกปฺปฏิเวธานิ.\csromanspacer{} ทุกฺขสฺส ทุกฺขฏฺโฐ ตถฏฺโฐ.\csromanspacer{} สมุทยสฺส สมุทยฏฺโฐ ตถฏฺโฐ.\csromanspacer{} นิโรธสฺส นิโรธฏฺโฐ ตถฏฺโฐ.\csromanspacer{} มคฺคสฺส มคฺคฏฺโฐ ตถฏฺโฐ.\csromanspacer{} อภิญฺญาย อภิญฺญฏฺโฐ ตถฏฺโฐ.\csromanspacer{} ปริญฺญาย ปริญฺญฏฺโฐ ตถฏฺโฐ.\csromanspacer{} ปหานสฺส ปหานฏฺโฐ ตถฏฺโฐ.\csromanspacer{} ภาวนาย ภาวนฏฺโฐ ตถฏฺโฐ.\csromanspacer{} สจฺฉิกิริยาย สจฺฉิกิริยฏฺโฐ ตถฏฺโฐ.\csromanspacer{} อิเมหิ นวหากาเรหิ ตถฏฺเฐน จตฺตาริ สจฺจานิ เอกสงฺคหิตานิ.\csromanspacer{} ยํ เอกสงฺคหิตํ ตํ เอกตฺตํ, เอกตฺตํ เอเกน ญาเณน ปฏิวิชฺฌตีติ จตฺตาริ สจฺจานิ เอกปฺปฏิเวธานิ.}

\prose{\csromanfolio{296}กถํ อนตฺตฏฺเฐน สจฺจฏฺเฐน ปฏิเวธฏฺเฐน จตฺตาริ สจฺจานิ เอกปฺปฏิเวธานิ? นวหากาเรหิ ปฏิเวธฏฺเฐน จตฺตาริ สจฺจานิ เอกปฺปฏิเวธานิ.\csromanspacer{} ทุกฺขสฺส ทุกฺขฏฺโฐ ปฏิเวธฏฺโฐ.\csromanspacer{} สมุทยสฺส สมุทยฏฺโฐ ปฏิเวธฏฺโฐ.\csromanspacer{} นิโรธสฺส นิโรธฏฺโฐ ปฏิเวธฏฺโฐ.\csromanspacer{} มคฺคสฺส มคฺคฏฺโฐ ปฏิเวธฏฺโฐ.\csromanspacer{} อภิญฺญาย อภิญฺญฏฺโฐ ปฏิเวธฏฺโฐ.\csromanspacer{} ปริญฺญาย ปริญฺญฏฺโฐ ปฏิเวธฏฺโฐ.\csromanspacer{} ปหานสฺส ปหานฏฺโฐ ปฏิเวธฏฺโฐ.\csromanspacer{} ภาวนาย ภาวนฏฺโฐ ปฏิเวธฏฺโฐ.\csromanspacer{} สจฺฉิกิริยาย สจฺฉิกิริยฏฺโฐ ปฏิเวธฏฺโฐ.\csromanspacer{} อิเมหิ นวหากาเรหิ ปฏิเวธฏฺเฐน จตฺตาริ สจฺจานิ เอกสงฺคหิตานิ.\csromanspacer{} ยํ เอกสงฺคหิตํ ตํ เอกตฺตํ, เอกตฺตํ เอเกน ญาเณน ปฏิวิชฺฌตีติ จตฺตาริ สจฺจานิ เอกปฺปฏิเวธานิ.}

\proseitem{11}{กติหากาเรหิ จตฺตาริ สจฺจานิ เอกปฺปฏิเวธานิ? ทฺวาทสหากาเรหิ จตฺตาริ สจฺจานิ เอกปฺปฏิเวธานิ, ตถฏฺเฐน อนตฺตฏฺเฐน สจฺจฏฺเฐน ปฏิเวธฏฺเฐน อภิชานนฏฺเฐน ปริชานนฏฺเฐน ธมฺมฏฺเฐน ตถฏฺเฐน ญาตฏฺเฐน สจฺฉิกิริยฏฺเฐน ผสฺสนฏฺเฐน อภิสมยฏฺเฐน.\csromanspacer{} อิเมหิ ทฺวาทสหิ อากาเรหิ จตฺตาริ สจฺจานิ เอกสงฺคหิตานิ.\csromanspacer{} ยํ เอกสงฺคหิตํ ตํ เอกตฺตํ, เอกตฺตํ เอเกน ญาเณน ปฏิวิชฺฌตีติ จตฺตาริ สจฺจานิ เอกปฺปฏิเวธานิ.}

\prose{กถํ ตถฏฺเฐน จตฺตาริ สจฺจานิ เอกปฺปฏิเวธานิ? โสฬสหากาเรหิ ตถฏฺเฐน จตฺตาริ สจฺจานิ เอกปฺปฏิเวธานิ.\csromanspacer{} ทุกฺขสฺส ปีฬนฏฺโฐ สงฺขตฏฺโฐ สนฺตาปฏฺโฐ วิปริณามฏฺโฐ ตถฏฺโฐ.\csromanspacer{} สมุทยสฺส อายูหนฏฺโฐ นิทานฏฺโฐ สํโยคฏฺโฐ ปลิโพธฏฺโฐ ตถฏฺโฐ.\csromanspacer{} นิโรธสฺส นิสฺสรณฏฺโฐ วิเวกฏฺโฐ อสงฺขตฏฺโฐ อมตฏฺโฐ ตถฏฺโฐ.\csromanspacer{} มคฺคสฺส นิยฺยานฏฺโฐ เหตุฏฺโฐ ทสฺสนฏฺโฐ อาธิปเตยฺยฏฺโฐ ตถฏฺโฐ.\csromanspacer{} อิเมหิ โสฬสหิ อากาเรหิ ตถฏฺเฐน จตฺตาริ สจฺจานิ เอกสงฺคหิตานิ.\csromanspacer{} ยํ เอกสงฺคหิตํ ตํ เอกตฺตํ, เอกตฺตํ เอเกน ญาเณน ปฏิวิชฺฌตีติ จตฺตาริ สจฺจานิ เอกปฺปฏิเวธานิ.}

\prose{กถํ อนตฺตฏฺเฐน ฯเปฯ.\csromanspacer{} สจฺจฏฺเฐน ปฏิเวธฏฺเฐน อภิชานนฏฺเฐน ปริชานนฏฺเฐน ธมฺมฏฺเฐน ตถฏฺเฐน ญาตฏฺเฐน สจฺฉิกิริยฏฺเฐน ผสฺสนฏฺเฐน อภิสมยฏฺเฐน จตฺตาริ สจฺจานิ เอกปฺปฏิเวธานิ.\csromanspacer{} โสฬสหิ อากาเรหิ อภิสมยฏฺเฐน จตฺตาริ สจฺจานิ เอกปฺปฏิเวธานิ.\csromanspacer{} ทุกฺขสฺส ปีฬนฏฺโฐ สงฺขตฏฺโฐ สนฺตาปฏฺโฐ วิปริณามฏฺโฐ อภิสมยฏฺโฐ.\csromanspacer{} สมุทยสฺส อายูหนฏฺโฐ นิทานฏฺโฐ สํโยคฏฺโฐ ปลิโพธฏฺโฐ อภิสมยฏฺโฐ.\csromanspacer{} นิโรธสฺส}

\vspace{\tipitakaparskip}
\csromancenter{สจฺจกถา นิสฺสรณฏฺโฐ วิเวกฏฺโฐ อสงฺขตฏฺโฐ อมตฏฺโฐ อภิสมยฏฺโฐ.\csromanspacer{} มคฺคสฺส นิยฺยานฏฺโฐ เหตุฏฺโฐ ทสฺสนฏฺโฐ อาธิปเตยฺยฏฺโฐ อภิสมยฏฺโฐ.\csromanspacer{} อิเมหิ โสฬสหิ อากาเรหิ อภิสมยฏฺเฐน จตฺตาริ สจฺจานิ เอกสงฺคหิตานิ.\csromanspacer{} ยํ เอกสงฺคหิตํ ตํ เอกตฺตํ, เอกตฺตํ เอเกน ญาเณน ปฏิวิชฺฌตีติ จตฺตาริ สจฺจานิ เอกปฺปฏิเวธานิ.}
\proseitem{12}{\csromanfolio{297}สจฺจานํ กติ ลกฺขณานิ? สจฺจานํ เทฺว ลกฺขณานิ สงฺขตลกฺขณญฺจ อสงฺขตลกฺขณญฺจ, สจฺจานํ อิมานิ เทฺว ลกฺขณานิ.}

\prose{สจฺจานํ กติ ลกฺขณานิ? สจฺจานํ ฉ ลกฺขณานิ.\csromanspacer{} สงฺขตานํ สจฺจานํ อุปฺปาโท ปญฺญายติ, วโย ปญฺญายติ, ฐิตานํ อญฺญถตฺตํ ปญฺญายติ.\csromanspacer{} อสงฺขตสฺส สจฺจสฺส น อุปฺปาโท ปญฺญายติ, น วโย ปญฺญายติ, น ฐิตสฺส อญฺญถตฺตํ ปญฺญายติ, สจฺจานํ อิมานิ ฉ ลกฺขณานิ.}

\prose{สจฺจานํ กติ ลกฺขณานิ? สจฺจานํ ทฺวาทส ลกฺขณานิ, ทุกฺขสจฺจสฺส อุปฺปาโท ปญฺญายติ, วโย ปญฺญายติ, ฐิตสฺส อญฺญถตฺตํ ปญฺญายติ, สมุทยสจฺจสฺส อุปฺปาโท ปญฺญายติ, วโย ปญฺญายติ, ฐิตสฺส อญฺญถตฺตํ ปญฺญายติ.\csromanspacer{} มคฺคสจฺจสฺส อุปฺปาโท ปญฺญายติ, วโย ปญฺญายติ, ฐิตสฺส อญฺญถตฺตํ ปญฺญายติ.\csromanspacer{} นิโรธสจฺจสฺส น อุปฺปาโท ปญฺญายติ, น วโย ปญฺญายติ, น ฐิตสฺส อญฺญถตฺตํ ปญฺญายติ, สจฺจานํ อิมานิ ทฺวาทส ลกฺขณานิ.}

\prose{จตุนฺนํ สจฺจานํ กติ กุสลา, กติ อกุสลา, กติ อพฺยากตา? สมุทยสจฺจํ อกุสลํ, มคฺคสจฺจํ กุสลํ, นิโรธสจฺจํ อพฺยากตํ.\csromanspacer{} ทุกฺขสจฺจํ สิยา กุสลํ, สิยา อกุสลํ, สิยา อพฺยากตํ.}

\prose{\textbf{สิยา} ตีณิ สจฺจานิ เอกสจฺเจน สงฺคหิตานิ.\csromanspacer{} เอกสจฺจํ ตีหิ สจฺเจหิ สงฺคหิตํ วตฺถุวเสน ปริยาเยน.\csromanspacer{} \textbf{สิยา}ติ กถญฺจ สิยา? ยํ ทุกฺขสจฺจํ อกุสลํ, สมุทยสจฺจํ อกุสลํ, เอวํ อกุสลฏฺเฐน เทฺว สจฺจานิ เอกสจฺเจน สงฺคหิตานิ.\csromanspacer{} เอกสจฺจํ ทฺวีหิ สจฺเจหิ สงฺคหิตํ, ยํ ทุกฺขสจฺจํ กุสลํ, มคฺคสจฺจํ กุสลํ เอวํ กุสลฏฺเฐน เทฺว สจฺจานิ เอกสจฺเจน สงฺคหิตานิ.\csromanspacer{} เอกสจฺจํ ทฺวีหิ สจฺเจหิ สงฺคหิตํ.\csromanspacer{} ยํ ทุกฺขสจฺจํ อพฺยากตํ, นิโรธสจฺจํ อพฺยากตํ, เอวํ อพฺยากตฏฺเฐน เทฺว สจฺจานิ เอกสจฺเจน สงฺคหิตานิ.\csromanspacer{} เอกสจฺจํ ทฺวีหิ สจฺเจหิ สงฺคหิตํ.\csromanspacer{} เอวํ สิยา ตีณิ สจฺจานิ เอกสจฺเจน สงฺคหิตานิ, เอกสจฺจํ ตีหิ สจฺเจหิ สงฺคหิตํ วตฺถุวเสน ปริยาเยนาติ.}

\csromanmatikaanchor{mtk.90}
\csromantocmarkref{subsection}{2. ทุติยสุตฺตนฺตปาฬิ}{mtk.90}
\bookmark[dest={mtk.90},level=2]{2. ทุติยสุตฺตนฺตปาฬิ}
\csromanheader{2}{2. ทุติยสุตฺตนฺตปาฬิ}
\proseitem{13}{\csromanfolio{298}ปุพฺเพ เม ภิกฺขเว\footnote{สํ 2. 23 ปิฏฺเฐปิ.} สมฺโพธา อนภิสมฺพุทฺธสฺส โพธิสตฺตสฺเสว สโต เอตทโหสิ “โก นุ โข รูปสฺส อสฺสาโท, โก อาทีนโว, กิํ นิสฺสรณํ.\csromanspacer{} โก เวทนาย อสฺสาโท, โก อาทีนโว, กิํ นิสฺสรณํ.\csromanspacer{} โก สญฺญาย อสฺสาโท, โก อาทีนโว, กิํ นิสฺสรณํ.\csromanspacer{} โก สงฺขารานํ อสฺสาโท, โก อาทีนโว, กิํ นิสฺสรณํ.\csromanspacer{} โก วิญฺญาณสฺส อสฺสาโท, โก อาทีนโว, กิํ นิสฺสรณนฺ”ติ.\csromanspacer{} ตสฺส มยฺหํ ภิกฺขเว เอตทโหสิ “ยํ โข รูปํ ปฏิจฺจ อุปฺปชฺชติ สุขํ โสมนสฺสํ, อยํ รูปสฺส อสฺสาโท.\csromanspacer{} ยํ รูปํ อนิจฺจํ, ตํ ทุกฺขํ วิปริณามธมฺมํ, อยํ รูปสฺส อาทีนโว.\csromanspacer{} โย รูปสฺมิํ ฉนฺทราควินโย ฉนฺทราคปฺปหานํ, อิทํ รูปสฺส นิสฺสรณํ.\csromanspacer{}\csromanspacer{}ยํ เวทนํ ปฏิจฺจ ฯเปฯ.\csromanspacer{} ยํ สญฺญํ ปฏิจฺจ.\csromanspacer{} ยํ สงฺขาเร ปฏิจฺจ.\csromanspacer{} ยํ วิญฺญาณํ ปฏิจฺจ อุปฺปชฺชติ สุขํ โสมนสฺสํ อยํ วิญฺญาณสฺส อสฺสาโท.\csromanspacer{} ยํ วิญฺญาณํ อนิจฺจํ, ตํ ทุกฺขํ วิปริณามธมฺมํ, อยํ วิญฺญาณสฺส อาทีนโว.\csromanspacer{} โย วิญฺญาณสฺมิํ ฉนฺทราควินโย ฉนฺทราคปฺปหานํ, อิทํ วิญฺญาณสฺส นิสฺสรณํ.}

\prose{ยาว กีวญฺจาหํ ภิกฺขเว อิเมสํ ปญฺจนฺนํ อุปาทานกฺขนฺธานํ เอวํ อสฺสาทญฺจ อสฺสาทโต, อาทีนวญฺจ อาทีนวโต, นิสฺสรณญฺจ นิสฺสรณโต ยถาภูตํ นาพฺภญฺญาสิํ.\csromanspacer{} เนว ตาวาหํ ภิกฺขเว สเทวเก โลเก สมารเก สพฺรหฺมเก สสฺสมณพฺราหฺมณิยา ปชาย สเทวมนุสฺสาย “อนุตฺตรํ สมฺมาสมฺโพธิํ อภิสมฺพุทฺโธ”ติ ปจฺจญฺญาสิํ\footnote{อภิสมฺพุทฺโธ ปจฺจญฺญาสิํ (สฺยา) สํ 2. 24 ปิฏฺเฐ ปสฺสิตพฺพา.}.\csromanspacer{} ยโต จ ขฺวาหํ ภิกฺขเว อิเมสํ ปญฺจนฺนํ อุปาทานกฺขนฺธานํ เอวํ อสฺสาทญฺจ อสฺสาทโต, อาทีนวญฺจ อาทีนวโต, นิสฺสรณญฺจ นิสฺสรณโต ยถาภูตํ อพฺภญฺญาสิํ.\csromanspacer{} อถาหํ ภิกฺขเว สเทวเก โลเก สมารเก สพฺรหฺมเก สสฺสมณพฺราหฺมณิยา ปชาย สเทวมนุสฺสาย “อนุตฺตรํ สมฺมาสมฺโพธํ อภิสมฺพุทฺโธ”ติ ปจฺจญฺญาสิํ.\csromanspacer{} ญาณญฺจ ปน เม ทสฺสนํ อุทปาทิ, อกุปฺปา เม วิมุตฺติ\footnote{เจโตวิมุตฺติ (สฺยา, ก) สํ 2. 24 ปิฏฺเฐ ปสฺสิตพฺพา.}.\csromanspacer{} อยมนฺติมา ชาติ, นตฺถิ ทานิ ปุนพฺภโวติ.}

\csromanmatikaanchor{mtk.91}
\csromantocmarkref{subsection}{3. ทุติยสุตฺตนฺตนิทฺเทส}{mtk.91}
\bookmark[dest={mtk.91},level=2]{3. ทุติยสุตฺตนฺตนิทฺเทส}
\csromanheader{2}{2. สจฺจกถา \textbf{3. ทุติยสุตฺตนฺตนิทฺเทส}}
\proseitem{14}{\csromanfolio{299}ยํ รูปํ ปฏิจฺจ อุปฺปชฺชติ สุขํ โสมนสฺสํ, “อยํ รูปสฺส อสฺสาโท”ติ ปหานปฺปฏิเวโธ สมุทยสจฺจํ.\csromanspacer{} ยํ รูปํ อนิจฺจํ, ตํ ทุกฺขํ วิปริณามธมฺมํ, “อยํ รูปสฺส อาทีนโว”ติ ปริญฺญาปฏิเวโธ ทุกฺขสจฺจํ.\csromanspacer{} โย รูปสฺมิํ ฉนฺทราควินโย ฉนฺทราคปฺปหานํ, “อิทํ รูปสฺส นิสฺสรณนฺ”ติ สจฺฉิกิริยาปฏิเวโธ นิโรธสจฺจํ.\csromanspacer{} ยา อิเมสุ ตีสุ ฐาเนสุ ทิฏฺฐิ สงฺกปฺโป วาจา กมฺมนฺโต อาชีโว วายาโม สติ สมาธิ ภาวนาปฏิเวโธ มคฺคสจฺจํ.}

\prose{ยํ เวทนํ ปฏิจฺจ ฯเปฯ.\csromanspacer{} ยํ สญฺญํ ปฏิจฺจ.\csromanspacer{} ยํ สงฺขาเร ปฏิจฺจ.\csromanspacer{} ยํ วิญฺญาณํ ปฏิจฺจ อุปฺปชฺชติ สุขํ โสมนสฺสํ, “อยํ วิญฺญาณสฺส อสฺสาโท”ติ ปหานปฺปฏิเวโธ สมุทยสจฺจํ.\csromanspacer{} ยํ วิญฺญาณํ อนิจฺจํ, ตํ ทุกฺขํ วิปริณามธมฺมํ, “อยํ วิญฺญาณสฺส อาทีนโว”ติ ปริญฺญาปฏิเวโธ ทุกฺขสจฺจํ.\csromanspacer{} โย วิญฺญาณสฺมิํ ฉนฺทราควินโย ฉนฺทราคปฺปหานํ, “อิทํ วิญฺญาณสฺส นิสฺสรณนฺ”ติ สจฺฉิกิริยาปฏิเวโธ นิโรธสจฺจํ.\csromanspacer{} ยา อิเมสุ ตีสุ ฐาเนสุ ทิฏฺฐิ สงฺกปฺโป วาจา กมฺมนฺโต อาชีโว วายาโม สติ สมาธิ ภาวนาปฏิเวโธ มคฺคสจฺจํ.}

\proseitem{15}{\textbf{สจฺจนฺ}ติ กติหากาเรหิ สจฺจํ? เอสนฏฺเฐน ปริคฺคหฏฺเฐน ปฏิเวธฏฺเฐน.\csromanspacer{} กถํ เอสนฏฺเฐน สจฺจํ? ชรามรณํ กิํนิทานํ, กิํสมุทยํ, กิํชาติกํ, กิํปภวนฺติ เอวํ เอสนฏฺเฐน สจฺจํ.\csromanspacer{} ชรามรณํ ชาตินิทานํ, ชาติสมุทยํ, ชาติชาติกํ, ชาติปฺปภวนฺติ เอวํ ปริคฺคหฏฺเฐน สจฺจํ.\csromanspacer{} ชรามรณญฺจ ปชานาติ, ชรามรณสมุทยญฺจ ปชานาติ, ชรามรณนิโรธญฺจ ปชานาติ, ชรามรณนิโรธคามินิํ ปฏิปทญฺจ ปชานาติ, เอวํ ปฏิเวธฏฺเฐน สจฺจํ.}

\prose{ชาติ กิํนิทานา, กิํสมุทยา, กิํชาติกา, กิํปภวาติ เอวํ เอสนฏฺเฐน สจฺจํ.\csromanspacer{} ชาติ ภวนิทานา, ภวสมุทยา, ภวชาติกา, ภวปฺปภวาติ เอวํ ปริคฺคหฏฺเฐน สจฺจํ.\csromanspacer{} ชาติญฺจ ปชานาติ, ชาติสมุทยญฺจ ปชานาติ, ชาตินิโรธญฺจ ปชานาติ, ชาตินิโรธคามินิํ ปฏิปทญฺจ ปชานาติ, เอวํ ปฏิเวธฏฺเฐน สจฺจํ.}

\prose{ภโว กิํนิทาโน, กิํสมุทโย, กิํชาติโก, กิํปภโวติ เอวํ เอสนฏฺเฐน สจฺจํ, ภโว อุปาทานนิทาโน, อุปาทานสมุทโย, อุปาทานชาติโก, อุปาทานปฺปภโวติ เอวํ ปริคฺคหฏฺเฐน สจฺจํ.\csromanspacer{} ภวญฺจ ปชานาติ, \csromanfolio{300}ภวสมุทยญฺจ ปชานาติ, ภวนิโรธญฺจ ปชานาติ, ภวนิโรธคามินิํ ปฏิปทญฺจ ปชานาติ, เอวํ ปฏิเวธฏฺเฐน สจฺจํ.}

\prose{อุปาทานํ กิํนิทานํ, กิํสมุทยํ, กิํชาติกํ, กิํปภวนฺติ เอวํ เอสนฏฺเฐน สจฺจํ.\csromanspacer{} อุปาทานํ ตณฺหานิทานํ, ตณฺหาสมุทยํ, ตณฺหาชาติกํ, ตณฺหาปภวนฺติ เอวํ ปริคฺคหฏฺเฐน สจฺจํ.\csromanspacer{} อุปาทานญฺจ ปชานาติ, อุปาทานสมุทยญฺจ ปชานาติ, อุปาทานนิโรธญฺจ ปชานาติ, อุปาทานนิโรธคามินิํ ปฏิปทญฺจ ปชานาติ, เอวํ ปฏิเวธฏฺเฐน สจฺจํ.}

\prose{ตณฺหา กิํ นิทานา, กิํสมุทยา, กิํชาติกา, กิํปภวาติ เอวํ เอสนฏฺเฐน สจฺจํ.\csromanspacer{} ตณฺหา เวทนานิทานา, เวทนาสมุทยา, เวทนาชาติกา, เวทนาปภวาติ เอวํ ปริคฺคหฏฺเฐน สจฺจํ.\csromanspacer{} ตณฺหญฺจ ปชานาติ, ตณฺหาสมุทยญฺจ ปชานาติ, ตณฺหานิโรธญฺจ ปชานาติ, ตณฺหานิโรธคามินิํ ปฏิปทญฺจ ปชานาติ, เอวํ ปฏิเวธฏฺเฐน สจฺจํ.}

\prose{เวทนา กิํนิทานา, กิํสมุทยา, กิํชาติกา, กิํปภวาติ เอวํ เอสนฏฺเฐน สจฺจํ.\csromanspacer{} เวทนา ผสฺสนิทานา, ผสฺสสมุทยา, ผสฺสชาติกา, ผสฺสปฺปภวาติ เอวํ ปริคฺคหฏฺเฐน สจฺจํ.\csromanspacer{} เวทนญฺจ ปชานาติ, เวทนาสมุทยญฺจ ปชานาติ, เวทนานิโรธญฺจ ปชานาติ, เวทนานิโรธคามินิํ ปฏิปทญฺจ ปชานาติ, เอวํ ปฏิเวธฏฺเฐน สจฺจํ.}

\prose{ผสฺโส กิํนิทาโน, กิํสมุทโย, กิํชาติโก, กิํปภโวติ เอวํ เอสนฏฺเฐน สจฺจํ.\csromanspacer{} ผสฺโส สฬายตนนิทาโน, สฬายตนสมุทโย, สฬายตนชาติโก, สฬายตนปฺปภโวติ เอวํ ปริคฺคหฏฺเฐน สจฺจํ.\csromanspacer{} ผสฺสญฺจ ปชานาติ, ผสฺสสมุทยญฺจ ปชานาติ, ผสฺสนิโรธญฺจ ปชานาติ, ผสฺสนิโรธคามินิํ ปฏิปทญฺจ ปชานาติ, เอวํ ปฏิเวธฏฺเฐน สจฺจํ.}

\prose{สฬายตนํ กิํนิทานํ, กิํสมุทยํ, กิํชาติกํ, กิํปภวนฺติ เอวํ เอสนฏฺเฐน สจฺจํ.\csromanspacer{} สฬายตนํ นามรูปนิทานํ, นามรูปสมุทยํ, นามรูปชาติกํ, นามรูปปฺปภวนฺติ เอวํ ปริคฺคหฏฺเฐน สจฺจํ.\csromanspacer{} สฬายตนญฺจ ปชานาติ, สฬายตนสมุทยญฺจ ปชานาติ, สฬายตนนิโรธญฺจ ปชานาติ, สฬายตนนิโรธคามินิํ ปฏิปทญฺจ ปชานาติ, เอวํ ปฏิเวธฏฺเฐน สจฺจํ.}

\prose{นามรูปํ กิํนิทานํ, กิํสมุทยํ, กิํชาติกํ, กิํปภวนฺติ เอวํ เอสนฏฺเฐน สจฺจํ.\csromanspacer{} นามรูปํ วิญฺญาณนิทานํ, วิญฺญาณสมุทยํ, วิญฺญาณชาติกํ,}

\vspace{\tipitakaparskip}
\csromancenter{สจฺจกถา วิญฺญาณปฺปภวนฺติ เอวํ ปริคฺคหฏฺเฐน สจฺจํ.\csromanspacer{} นามรูปญฺจ ปชานาติ, นามรูปสมุทยญฺจ ปชานาติ, นามรูปนิโรธญฺจ ปชานาติ, นามรูปนิโรธคามินิํ ปฏิปทญฺจ ปชานาติ, เอวํ ปฏิเวธฏฺเฐน สจฺจํ.}
\prose{\csromanfolio{301}วิญฺญาณํ กิํนิทานํ, กิํสมุทยํ, กิํชาติกํ, กิํปภวนฺติ เอวํ เอสนฏฺเฐน สจฺจํ.\csromanspacer{} วิญฺญาณํ สงฺขารนิทานํ, สงฺขารสมุทยํ, สงฺขารชาติกํ, สงฺขารปฺปภวนฺติ เอวํ ปริคฺคหฏฺเฐน สจฺจํ.\csromanspacer{} วิญฺญาณญฺจ ปชานาติ, วิญฺญาณสมุทยญฺจ ปชานาติ, วิญฺญาณนิโรธญฺจ ปชานาติ, วิญฺญาณนิโรธคามินิํ ปฏิปทญฺจ ปชานาติ, เอวํ ปฏิเวธฏฺเฐน สจฺจํ.}

\prose{สงฺขารา กิํนิทานา, กิํสมุทยา, กิํชาติกา, กิํปภวาติ เอวํ เอสนฏฺเฐน สจฺจํ.\csromanspacer{} สงฺขารา อวิชฺชานิทานา, อวิชฺชาสมุทยา, อวิชฺชาชาติกา, อวิชฺชาปภาวาติ เอวํ ปริคฺคหฏฺเฐน สจฺจํ.\csromanspacer{} สงฺขาเร จ ปชานาติ, สงฺขารสมุทยญฺจ ปชานาติ, สงฺขารนิโรธญฺจ ปชานาติ, สงฺขารนิโรธคามินิํ ปฏิปทญฺจ ปชานาติ, เอวํ ปฏิเวธฏฺเฐน สจฺจํ.}

\proseitem{16}{ชรามรณํ ทุกฺขสจฺจํ, ชาติ สมุทยสจฺจํ, อุภินฺนมฺปิ นิสฺสรณํ นิโรธสจฺจํ, นิโรธปฺปชานนา มคฺคสจฺจํ.\csromanspacer{}\csromanspacer{}ชาติ ทุกฺขสจฺจํ, ภโว สมุทยสจฺจํ, อุภินฺนมฺปิ นิสฺสรณํ นิโรธสจฺจํ, นิโรธปฺปชานนา มคฺคสจฺจํ.\csromanspacer{}\csromanspacer{}ภโว ทุกฺขสจฺจํ, อุปาทานํ สมุทยสจฺจํ, อุภินฺนมฺปิ นิสฺสรณํ นิโรธสจฺจํ, นิโรธปฺปชานนา มคฺคสจฺจํ.\csromanspacer{}\csromanspacer{}อุปาทานํ ทุกฺขสจฺจํ, ตณฺหา สมุทยสจฺจํ, อุภินฺนมฺปิ นิสฺสรณํ นิโรธสจฺจํ, นิโรธปฺปชานนา มคฺคสจฺจํ.\csromanspacer{}\csromanspacer{}ตณฺหา ทุกฺขสจฺจํ, เวทนา สมุทยสจฺจํ, อุภินฺนมฺปิ นิสฺสรณํ นิโรธสจฺจํ, นิโรธปฺปชานนา มคฺคสจฺจํ.\csromanspacer{}\csromanspacer{}เวทนา ทุกฺขสจฺจํ, ผสฺโส สมุทยสจฺจํ, อุภินฺนมฺปิ นิสฺสรณํ นิโรธสจฺจํ, นิโรธปฺปชานนา มคฺคสจฺจํ.\csromanspacer{}\csromanspacer{}ผสฺโส ทุกฺขสจฺจํ, สฬายตนํ สมุทยสจฺจํ, อุภินฺนมฺปิ นิสฺสรณํ นิโรธสจฺจํ, นิโรธปฺปชานนา มคฺคสจฺจํ.\csromanspacer{}\csromanspacer{}สฬายตนํ ทุกฺขสจฺจํ, นามรูปํ สมุทยสจฺจํ, อุภินฺนมฺปิ นิสฺสรณํ นิโรธสจฺจํ, นิโรธปฺปชานนา มคฺคสจฺจํ.\csromanspacer{}\csromanspacer{}นามรูปํ ทุกฺขสจฺจํ, วิญฺญาณํ สมุทยสจฺจํ, อุภินฺนมฺปิ นิสฺสรณํ นิโรธสจฺจํ, นิโรธปฺปชานนา มคฺคสจฺจํ.\csromanspacer{}\csromanspacer{}วิญฺญาณํ ทุกฺขสจฺจํ, สงฺขารา สมุทยสจฺจํ, อุภินฺนมฺปิ นิสฺสรณํ นิโรธสจฺจํ, นิโรธปฺปชานนา มคฺคสจฺจํ.\csromanspacer{}\csromanspacer{}สงฺขารา ทุกฺขสจฺจํ, อวิชฺชา สมุทยสจฺจํ, อุภินฺนมฺปิ นิสฺสรณํ นิโรธสจฺจํ, นิโรธปฺปชานนา มคฺคสจฺจํ.}

\prose{\csromanfolio{302}ชรามรณํ สิยา ทุกฺขสจฺจํ, สิยา สมุทยสจฺจํ, อุภินฺนมฺปิ นิสฺสรณํ นิโรธสจฺจํ, นิโรธปฺปชานนา มคฺคสจฺจํ.\csromanspacer{}\csromanspacer{}ชาติ สิยา ทุกฺขสจฺจํ, สิยา สมุทยสจฺจํ ฯเปฯ.\csromanspacer{} ภโว สิยา ทุกฺขสจฺจํ, สิยา สมุทยสจฺจํ, อุภินฺนมฺปิ นิสฺสรณํ นิโรธสจฺจํ, นิโรธปฺปชานนา มคฺคสจฺจนฺติ.}

\nitthitam{สจฺจกถา นิฏฺฐิตา.}
\csromanheader{2}{ภาณวาโร.}
\csromansectionrule
\csromanmatikaanchor{mtk.92}
\csromantocmarknopage{section}{3. โพชฺฌงฺคกถา}{mtk.92}
\bookmark[dest={mtk.92},level=1]{3. โพชฺฌงฺคกถา}
\csromanheader{1}{3. โพชฺฌงฺคกถา}
\proseitem{17}{สาวตฺถินิทานํ.\csromanspacer{}\csromanspacer{}สตฺติเม ภิกฺขเว\footnote{สํ 3. 75 ปิฏฺเฐ.} โพชฺฌงฺคา, กตเม สตฺต? สติสมฺโพชฺฌงฺโค ธมฺมวิจยสมฺโพชฺฌงฺโค วีริยสมฺโพชฺฌงฺโค ปีติสมฺโพชฺฌงฺโค ปสฺสทฺธิสมฺโพชฺฌงฺโค สมาธิสมฺโพชฺฌงฺโค อุเปกฺขาสมฺโพชฺฌงฺโค.\csromanspacer{} อิเม โข ภิกฺขเว สตฺต โพชฺฌงฺคา.}

\prose{\textbf{โพชฺฌงฺคา}ติ เกนฏฺเฐน โพชฺฌงฺคา? โพธาย สํวตฺตนฺตีติ โพชฺฌงฺคา.\csromanspacer{} พุชฺฌนฺตีติ โพชฺฌงฺคา, อนุพุชฺฌนฺตีติ โพชฺฌงฺคา, ปฏิพุชฺฌนฺตีติ โพชฺฌงฺคา, สมฺพุชฺฌนฺตีติ โพชฺฌงฺคา. (1)}

\prose{พุชฺฌนฏฺเฐน โพชฺฌงฺคา, อนุพุชฺฌนฏฺเฐน โพชฺฌงฺคา, ปฏิพุชฺฌนฏฺเฐน โพชฺฌงฺคา, สมฺพุชฺฌนฏฺเฐน โพชฺฌงฺคา. (2)}

\prose{โพเธนฺตีติ โพชฺฌงฺคา, อนุโพเธนฺตีติ โพชฺฌงฺคา, ปฏิโพเธนฺตีติ โพชฺฌงฺคา, สมฺโพเธนฺตีติ โพชฺฌงฺคา. (3)}

\prose{โพธนฏฺเฐน โพชฺฌงฺคา, อนุโพธนฏฺเฐน โพชฺฌงฺคา, ปฏิโพธนฏฺเฐน โพชฺฌงฺคา, สมฺโพธนฏฺเฐน โพชฺฌงฺคา. (4)}

\prose{โพธิปกฺขิยฏฺเฐน โพชฺฌงฺคา, อนุโพธิปกฺขิยฏฺเฐน โพชฺฌงฺคา, ปฏิโพธิปกฺขิยฏฺเฐน โพชฺฌงฺคา, สมฺโพธิปกฺขิยฏฺเฐน โพชฺฌงฺคา. (5) (ปญฺจมจตุกฺกํ)}

\prose{พุทฺธิลภนฏฺเฐน โพชฺฌงฺคา, พุทฺธิปฏิลภนฏฺเฐน โพชฺฌงฺคา, พุทฺธิโรปนฏฺเฐน โพชฺฌงฺคา, พุทฺธิ-อภิโรปนฏฺเฐน โพชฺฌงฺคา, พุทฺธิปาปนฏฺเฐน โพชฺฌงฺคา, พุทฺธิสมฺปาปนฏฺเฐน โพชฺฌงฺคา. (ฉกฺกํ)}

\csromanmatikaanchor{mtk.93}
\csromantocmarkref{section}{มูลมูลกาทิทสก}{mtk.93}
\bookmark[dest={mtk.93},level=1]{มูลมูลกาทิทสก}
\csromanheader{1}{3. โพชฺฌงฺคกถา \textbf{มูลมูลกาทิทสก}}
\proseitem{18}{\csromanfolio{303}มูลฏฺเฐน โพชฺฌงฺคา, มูลจริยฏฺเฐน โพชฺฌงฺคา, มูลปริคฺคหฏฺเฐน โพชฺฌงฺคา, มูลปริวารฏฺเฐน โพชฺฌงฺคา, มูลปริปูรณฏฺเฐน\footnote{มูลปริปูรฏฺเฐน (สฺยา, ก) เอวมุปริปิ.} โพชฺฌงฺคา, มูลปริปากฏฺเฐน โพชฺฌงฺคา, มูลปฏิสมฺภิทฏฺเฐน โพชฺฌงฺคา, มูลปฏิสมฺภิทาปาปนฏฺเฐน โพชฺฌงฺคา, มูลปฏิสมฺภิทาย วสีภาวฏฺเฐน โพชฺฌงฺคา, มูลปฏิสมฺภิทาย วสีภาวปฺปตฺตานมฺปิ โพชฺฌงฺคา. (1)}

\prose{เหตุฏฺเฐน โพชฺฌงฺคา, เหตุจริยฏฺเฐน โพชฺฌงฺคา, เหตุปริคฺคหฏฺเฐน โพชฺฌงฺคา, เหตุปริวารฏฺเฐน โพชฺฌงฺคา, เหตุปริปูรณฏฺเฐน โพชฺฌงฺคา, เหตุปริปากฏฺเฐน โพชฺฌงฺคา, เหตุปฏิสมฺภิทฏฺเฐน โพชฺฌงฺคา, เหตุปฏิสมฺภิทาปาปนฏฺเฐน โพชฺฌงฺคา, เหตุปฏิสมฺภิทาย วสีภาวฏฺเฐน โพชฺฌงฺคา, เหตุปฏิสมฺภิทาย วสีภาวปฺปตฺตานมฺปิ โพชฺฌงฺคา. (2)}

\prose{ปจฺจยฏฺเฐน โพชฺฌงฺคา, ปจฺจยจริยฏฺเฐน โพชฺฌงฺคา, ปจฺจยปริคฺคหฏฺเฐน โพชฺฌงฺคา, ปจฺจยปริวารฏฺเฐน โพชฺฌงฺคา, ปจฺจยปริปูรณฏฺเฐน โพชฺฌงฺคา, ปจฺจยปริปากฏฺเฐน โพชฺฌงฺคา, ปจฺจยปฏิสมฺภิทฏฺเฐน โพชฺฌงฺคา, ปจฺจยปฏิสมฺภิทาปาปนฏฺเฐน โพชฺฌงฺคา, ปจฺจยปฏิสมฺภิทาย วสีภาวฏฺเฐน โพชฺฌงฺคา, ปจฺจยปฏิสมฺภิทาย วสีภาวปฺปตฺตานมฺปิ โพชฺฌงฺคา. (3)}

\prose{วิสุทฺธฏฺเฐน โพชฺฌงฺคา, วิสุทฺธิจริยฏฺเฐน โพชฺฌงฺคา, วิสุทฺธิปริคฺคหฏฺเฐน โพชฺฌงฺคา, วิสุทฺธิปริวารฏฺเฐน โพชฺฌงฺคา, วิสุทฺธิปริปูรณฏฺเฐน โพชฺฌงฺคา, วิสุทฺธิปริปากฏฺเฐน โพชฺฌงฺคา, วิสุทฺธิปฏิสมฺภิทฏฺเฐน โพชฺฌงฺคา, วิสุทฺธิปฏิสมฺภิทาปาปนฏฺเฐน โพชฺฌงฺคา, วิสุทฺธิปฏิสมฺภิทาย วสีภาวฏฺเฐน โพชฺฌงฺคา, วิสุทฺธิปฏิสมฺภิทาย วสีภาวปฺปตฺตานมฺปิ โพชฺฌงฺคา. (4)}

\prose{อนวชฺชฏฺเฐน โพชฺฌงฺคา, อนวชฺชจริยฏฺเฐน โพชฺฌงฺคา, อนวชฺชปริคฺคหฏฺเฐน โพชฺฌงฺคา, อนวชฺชปริวารฏฺเฐน โพชฺฌงฺคา, อนวชฺชปริปูรณฏฺเฐน โพชฺฌงฺคา, อนวชฺชปริปากฏฺเฐน โพชฺฌงฺคา, อนวชฺชปฏิสมฺภิทฏฺเฐน โพชฺฌงฺคา, อนวชฺชปฏิสมฺภิทาปาปนฏฺเฐน โพชฺฌงฺคา, อนวชฺชปฏิสมฺภิทาย วสีภาวฏฺเฐน โพชฺฌงฺคา, อนวชฺชปฏิสมฺภิทาย วสีภาวปฺปตฺตานมฺปิ โพชฺฌงฺคา. (5)}

\prose{เนกฺขมฺมฏฺเฐน โพชฺฌงฺคา, เนกฺขมฺมจริยฏฺเฐน โพชฺฌงฺคา, เนกฺขมฺมปริคฺคหฏฺเฐน โพชฺฌงฺคา, เนกฺขมฺมปริวารฏฺเฐน โพชฺฌงฺคา, เนกฺขมฺมปริปูรณฏฺเฐน โพชฺฌงฺคา, \csromanfolio{304}เนกฺขมฺมปริปากฏฺเฐน โพชฺฌงฺคา, เนกฺขมฺมปฏิสมฺภิทฏฺเฐน โพชฺฌงฺคา, เนกฺขมฺมปฏิสมฺภิทาปาปนฏฺเฐน โพชฺฌงฺคา, เนกฺขมฺมปฏิสมฺภิทาย วสีภาวฏฺเฐน โพชฺฌงฺคา, เนกฺขมฺมปฏิสมฺภิทาย วสีภาวปฺปตฺตานมฺปิ โพชฺฌงฺคา. (6)}

\prose{วิมุตฺตฏฺเฐน โพชฺฌงฺคา, วิมุตฺติจริยฏฺเฐน โพชฺฌงฺคา, วิมุตฺติปริคฺคหฏฺเฐน โพชฺฌงฺคา, วิมุตฺติปริวารฏฺเฐน โพชฺฌงฺคา, วิมุตฺติปริปูรณฏฺเฐน โพชฺฌงฺคา, วิมุตฺติปริปากฏฺเฐน โพชฺฌงฺคา, วิมุตฺติปฏิสมฺภิทฏฺเฐน โพชฺฌงฺคา, วิมุตฺติปฏิสมฺภิทาปาปนฏฺเฐน โพชฺฌงฺคา, วิมุตฺติปฏิสมฺภิทาย วสีภาวฏฺเฐน โพชฺฌงฺคา, วิมุตฺติปฏิสมฺภิทาย วสีภาวปฺปตฺตานมฺปิ โพชฺฌงฺคา. (7)}

\prose{อนาสวฏฺเฐน โพชฺฌงฺคา, อนาสวจริยฏฺเฐน โพชฺฌงฺคา, อนาสวปริคฺคหฏฺเฐน โพชฺฌงฺคา, อนาสวปริวารฏฺเฐน โพชฺฌงฺคา, อนาสวปริปูรณฏฺเฐน โพชฺฌงฺคา, อนาสวปริปากฏฺเฐน โพชฺฌงฺคา, อนาสวปฏิสมฺภิทฏฺเฐน โพชฺฌงฺคา, อนาสวปฏิสมฺภิทาปาปนฏฺเฐน โพชฺฌงฺคา, อนาสวปฏิสมฺภิทาย วสีภาวฏฺเฐน โพชฺฌงฺคา, อนาสวปฏิสมฺภิทาย วสีภาวปฺปตฺตานมฺปิ โพชฺฌงฺคา. (8)}

\prose{วิเวกฏฺเฐน โพชฺฌงฺคา, วิเวกจริยฏฺเฐน โพชฺฌงฺคา, วิเวกปริคฺคหฏฺเฐน โพชฺฌงฺคา, วิเวกปริวารฏฺเฐน โพชฺฌงฺคา, วิเวกปริปูรณฏฺเฐน โพชฺฌงฺคา, วิเวกปริปากฏฺเฐน โพชฺฌงฺคา, วิเวกปฏิสมฺภิทฏฺเฐน โพชฺฌงฺคา, วิเวกปฏิสมฺภิทาปาปนฏฺเฐน โพชฺฌงฺคา, วิเวกปฏิสมฺภิทาย วสีภาวฏฺเฐน โพชฺฌงฺคา, วิเวกปฏิสมฺภิทาย วสีภาวปฺปตฺตานมฺปิ โพชฺฌงฺคา. (9)}

\prose{โวสคฺคฏฺเฐน โพชฺฌงฺคา, โวสคฺคจริยฏฺเฐน โพชฺฌงฺคา, โวสคฺคปริคฺคหฏฺเฐน โพชฺฌงฺคา, โวสคฺคปริวารฏฺเฐน โพชฺฌงฺคา, โวสคฺคปริปูรณฏฺเฐน โพชฺฌงฺคา, โวสคฺคปริปากฏฺเฐน โพชฺฌงฺคา, โวสคฺคปฏิสมฺภิทฏฺเฐน โพชฺฌงฺคา, โวสคฺคปฏิสมฺภิทาปาปนฏฺเฐน โพชฺฌงฺคา, โวสคฺคปฏิสมฺภิทาย วสีภาวฏฺเฐน โพชฺฌงฺคา, โวสคฺคปฏิสมฺภิทาย วสีภาวปฺปตฺตานมฺปิ โพชฺฌงฺคา. (10)}

\proseitem{19}{มูลฏฺฐํ พุชฺฌนฺตีติ โพชฺฌงฺคา, เหตุฏฺฐํ พุชฺฌนฺตีติ โพชฺฌงฺคา, ปจฺจยฏฺฐํ พุชฺฌนฺตีติ โพชฺฌงฺคา, วิสุทฺธฏฺฐํ พุชฺฌนฺตีติ โพชฺฌงฺคา, อนวชฺชฏฺฐํ พุชฺฌนฺตีติ โพชฺฌงฺคา, เนกฺขมฺมฏฺฐํ พุชฺฌนฺตีติ โพชฺฌงฺคา, วิมุตฺตฏฺฐํ พุชฺฌนฺตีติ โพชฺฌงฺคา, อนาสวฏฺฐํ พุชฺฌนฺตีติ โพชฺฌงฺคา, วิเวกฏฺฐํ พุชฺฌนฺตีติ โพชฺฌงฺคา, โวสคฺคฏฺฐํ พุชฺฌนฺตีติ โพชฺฌงฺคา.}

\csromanheader{2}{3. โพชฺฌงฺคกถา}
\prose{\csromanfolio{305}มูลจริยฏฺฐํ พุชฺฌนฺตีติ โพชฺฌงฺคา, เหตุจริยฏฺฐํ พุชฺฌนฺตีติ โพชฺฌงฺคา, ปจฺจยจริยฏฺฐํ พุชฺฌนฺตีติ โพชฺฌงฺคา, วิสุทฺธิจริยฏฺฐํ พุชฺฌนฺตีติ โพชฺฌงฺคา, อนวชฺชจริยฏฺฐํ พุชฺฌนฺตีติ โพชฺฌงฺคา, เนกฺขมฺมจริยฏฺฐํ พุชฺฌนฺตีติ โพชฺฌงฺคา, วิมุตฺติจริยฏฺฐํ พุชฺฌนฺตีติ โพชฺฌงฺคา, อนาสวจริยฏฺฐํ พุชฺฌนฺตีติ โพชฺฌงฺคา, วิเวกจริยฏฺฐํ พุชฺฌนฺตีติ โพชฺฌงฺคา, โวสคฺคจริยฏฺฐํ พุชฺฌนฺตีติ โพชฺฌงฺคา.}

\prose{มูลปริคฺคหฏฺฐํ พุชฺฌนฺตีติ โพชฺฌงฺคา ฯเปฯ โวสคฺคปริคฺคหฏฺฐํ พุชฺฌนฺตีติ โพชฺฌงฺคา.\csromanspacer{} มูลปริวารฏฺฐํ พุชฺฌนฺตีติ โพชฺฌงฺคา ฯเปฯ โวสคฺคปริวารฏฺฐํ พุชฺฌนฺตีติ โพชฺฌงฺคา.\csromanspacer{} มูลปริปูรณฏฺฐํ พุชฺฌนฺตีติ โพชฺฌงฺคา ฯเปฯ โวสคฺคปริปูรฏฺฐํ พุชฺฌนฺตีติ โพชฺฌงฺคา.\csromanspacer{} มูลปริปากฏฺฐํ พุชฺฌนฺตีติ โพชฺฌงฺคา ฯเปฯ โวสคฺคปริปากฏฺฐํ พุชฺฌนฺตีติ โพชฺฌงฺคา.\csromanspacer{} มูลปฏิสมฺภิทฏฺฐํ พุชฺฌนฺตีติ โพชฺฌงฺคา ฯเปฯ โวสคฺคปฏิสมฺภิทฏฺฐํ พุชฺฌนฺตีติ โพชฺฌงฺคา.\csromanspacer{}\csromanspacer{}มูลปฏิสมฺภิทาปาปนฏฺฐํ พุชฺฌนฺตีติ โพชฺฌงฺคา ฯเปฯ โวสคฺคปฏิสมฺภิทาปาปนฏฺฐํ พุชฺฌนฺตีติ โพชฺฌงฺคา.\csromanspacer{}\csromanspacer{}มูลปฏิสมฺภิทาย วสีภาวฏฺฐํ พุชฺฌนฺตีติ โพชฺฌงฺคา ฯเปฯ โวสคฺคปฏิสมฺภิทาย วสีภาวฏฺฐํ พุชฺฌนฺตีติ โพชฺฌงฺคา ฯเปฯ.}

\prose{ปริคฺคหฏฺฐํ พุชฺฌนฺตีติ โพชฺฌงฺคา, ปริวารฏฺฐํ พุชฺฌนฺตีติ โพชฺฌงฺคา, ปริปูรณฏฺฐํ พุชฺฌนฺตีติ โพชฺฌงฺคา, เอกคฺคฏฺฐํ พุชฺฌนฺตีติ โพชฺฌงฺคา, อวิกฺเขปฏฺฐํ พุชฺฌนฺตีติ โพชฺฌงฺคา, ปคฺคหฏฺฐํ พุชฺฌนฺตีติ โพชฺฌงฺคา, อวิสารฏฺฐํ พุชฺฌนฺตีติ โพชฺฌงฺคา, อนาวิลฏฺฐํ พุชฺฌนฺตีติ โพชฺฌงฺคา, อนิญฺชนฏฺฐํ พุชฺฌนฺตีติ โพชฺฌงฺคา, เอกตฺตุปฏฺฐานวเสน จิตฺตสฺส ฐิตฏฺฐํ พุชฺฌนฺตีติ โพชฺฌงฺคา, อารมฺมณฏฺฐํ พุชฺฌนฺตีติ โพชฺฌงฺคา, โคจรฏฺฐํ พุชฺฌนฺตีติ โพชฺฌงฺคา, ปหานฏฺฐํ พุชฺฌนฺตีติ โพชฺฌงฺคา, ปริจฺจาคฏฺฐํ พุชฺฌนฺตีติ โพชฺฌงฺคา, วุฏฺฐานฏฺฐํ พุชฺฌนฺตีติ โพชฺฌงฺคา, วิวฏฺฏนฏฺฐํ พุชฺฌนฺตีติ โพชฺฌงฺคา, สนฺตฏฺฐํ พุชฺฌนฺตีติ โพชฺฌงฺคา, ปณีตฏฺฐํ พุชฺฌนฺตีติ โพชฺฌงฺคา, วิมุตฺตฏฺฐํ พุชฺฌนฺตีติ โพชฺฌงฺคา, อนาสวฏฺฐํ พุชฺฌนฺตีติ โพชฺฌงฺคา, ตรณฏฺฐํ พุชฺฌนฺตีติ โพชฺฌงฺคา, อนิมิตฺตฏฺฐํ พุชฺฌนฺตีติ โพชฺฌงฺคา, อปฺปณิหิตฏฺฐํ พุชฺฌนฺตีติ โพชฺฌงฺคา, สุญฺญตฏฺฐํ พุชฺฌนฺตีติ โพชฺฌงฺคา, เอกรสฏฺฐํ พุชฺฌนฺตีติ โพชฺฌงฺคา, อนติวตฺตนฏฺฐํ พุชฺฌนฺตีติ โพชฺฌงฺคา, ยุคนทฺธฏฺฐํ พุชฺฌนฺตีติ โพชฺฌงฺคา, นิยฺยานฏฺฐํ พุชฺฌนฺตีติ โพชฺฌงฺคา, เหตุฏฺฐํ พุชฺฌนฺตีติ โพชฺฌงฺคา, ทสฺสนฏฺฐํ พุชฺฌนฺตีติ โพชฺฌงฺคา, อาธิปเตยฺยฏฺฐํ พุชฺฌนฺตีติ โพชฺฌงฺคา.}

\prose{สมถสฺส อวิกฺเขปฏฺฐํ พุชฺฌนฺตีติ โพชฺฌงฺคา, วิปสฺสนาย อนุปสฺสนฏฺฐํ พุชฺฌนฺตีติ โพชฺฌงฺคา, สมถวิปสฺสนานํ เอกรสฏฺฐํ พุชฺฌนฺตีติ โพชฺฌงฺคา, ยุคนทฺธสฺส อนติวตฺตนฏฺฐํ พุชฺฌนฺตีติ โพชฺฌงฺคา.}

\prose{\csromanfolio{306}สิกฺขาย สมาทานฏฺฐํ พุชฺฌนฺตีติ โพชฺฌงฺคา, อารมฺมณสฺส โคจรฏฺฐํ พุชฺฌนฺตีติ โพชฺฌงฺคา, ลีนสฺส จิตฺตสฺส ปคฺคหฏฺฐํ พุชฺฌนฺตีติ โพชฺฌงฺคา, อุทฺธตสฺส จิตฺตสฺส นิคฺคหฏฺฐํ พุชฺฌนฺตีติ โพชฺฌงฺคา.\csromanspacer{} อุโภวิสุทฺธานํ อชฺฌุเปกฺขนฏฺฐํ พุชฺฌนฺตีติ โพชฺฌงฺคา.\csromanspacer{} วิเสสาธิคมฏฺฐํ พุชฺฌนฺตีติ โพชฺฌงฺคา.\csromanspacer{} อุตฺตริปฏิเวธฏฺฐํ พุชฺฌนฺตีติ โพชฺฌงฺคา.\csromanspacer{} สจฺจาภิสมยฏฺฐํ พุชฺฌนฺตีติ โพชฺฌงฺคา.\csromanspacer{} นิโรเธ ปติฏฺฐาปกฏฺฐํ พุชฺฌนฺตีติ โพชฺฌงฺคา.}

\prose{สทฺธินฺทฺริยสฺส อธิโมกฺขฏฺฐํ พุชฺฌนฺตีติ โพชฺฌงฺคา ฯเปฯ ปญฺญินฺทฺริยสฺส ทสฺสนฏฺฐํ พุชฺฌนฺตีติ โพชฺฌงฺคา.\csromanspacer{} สทฺธาพลสฺส อสฺสทฺธิเย อกมฺปิยฏฺฐํ พุชฺฌนฺตีติ โพชฺฌงฺคา ฯเปฯ ปญฺญาพลสฺส อวิชฺชาย อกมฺปิยฏฺฐํ พุชฺฌนฺตีติ โพชฺฌงฺคา.\csromanspacer{} สติสมฺโพชฺฌงฺคสฺส อุปฏฺฐานฏฺฐํ พุชฺฌนฺตีติ โพชฺฌงฺคา ฯเปฯ อุเปกฺขาสมฺโพชฺฌงฺคสฺส ปฏิสงฺขานฏฺฐํ พุชฺฌนฺตีติ โพชฺฌงฺคา.\csromanspacer{} สมฺมาทิฏฺฐิยา ทสฺสนฏฺฐํ พุชฺฌนฺตีติ โพชฺฌงฺคา ฯเปฯ สมฺมาสมาธิสฺส อวิกฺเขปฏฺฐํ พุชฺฌนฺตีติ โพชฺฌงฺคา.}

\prose{อินฺทฺริยานํ อาธิปเตยฺยฏฺฐํ พุชฺฌนฺตีติ โพชฺฌงฺคา.\csromanspacer{} พลานํ อกมฺปิยฏฺฐํ พุชฺฌนฺตีติ โพชฺฌงฺคา.\csromanspacer{} นิยฺยานฏฺฐํ พุชฺฌนฺตีติ โพชฺฌงฺคา.\csromanspacer{} มคฺคสฺส เหตุฏฺฐํ พุชฺฌนฺตีติ โพชฺฌงฺคา.\csromanspacer{} สติปฏฺฐานานํ อุปฏฺฐานฏฺฐํ พุชฺฌนฺตีติ โพชฺฌงฺคา.\csromanspacer{} สมฺมปฺปธานานํ ปทหนฏฺฐํ พุชฺฌนฺตีติ โพชฺฌงฺคา.\csromanspacer{} อิทฺธิปาทานํ อิชฺฌนฏฺฐํ พุชฺฌนฺตีติ โพชฺฌงฺคา.\csromanspacer{} สจฺจานํ ตถฏฺฐํ พุชฺฌนฺตีติ โพชฺฌงฺคา.\csromanspacer{} ปโยคานํ ปฏิปฺปสฺสทฺธฏฺฐํ พุชฺฌนฺตีติ โพชฺฌงฺคา.\csromanspacer{} ผลานํ สจฺฉิกิริยฏฺฐํ พุชฺฌนฺตีติ โพชฺฌงฺคา.}

\prose{วิตกฺกสฺส อภินิโรปนฏฺฐํ พุชฺฌนฺตีติ โพชฺฌงฺคา.\csromanspacer{} วิจารสฺส อุปวิจารฏฺฐํ พุชฺฌนฺตีติ โพชฺฌงฺคา.\csromanspacer{} ปีติยา ผรณฏฺฐํ พุชฺฌนฺตีติ โพชฺฌงฺคา.\csromanspacer{} สุขสฺส อภิสนฺทนฏฺฐํ พุชฺฌนฺตีติ โพชฺฌงฺคา.\csromanspacer{} จิตฺตสฺส เอกคฺคฏฺฐํ พุชฺฌนฺตีติ โพชฺฌงฺคา.}

\prose{อาวชฺชนฏฺฐํ พุชฺฌนฺตีติ โพชฺฌงฺคา.\csromanspacer{} วิชานนฏฺฐํ พุชฺฌนฺตีติ โพชฺฌงฺคา.\csromanspacer{} ปชานนฏฺฐํ พุชฺฌนฺตีติ โพชฺฌงฺคา.\csromanspacer{} สญฺชานนฏฺฐํ พุชฺฌนฺตีติ โพชฺฌงฺคา.\csromanspacer{} เอโกทฏฺฐํ พุชฺฌนฺตีติ โพชฺฌงฺคา.\csromanspacer{} อภิญฺญาย ญาตฏฺฐํ\footnote{อภิญฺเญยฺยตฺถํ (สฺยา) 17 ปิฏฺเฐ ปสฺสิตพฺพา.} พุชฺฌนฺตีติ โพชฺฌงฺคา.\csromanspacer{} ปริญฺญาย ตีรณฏฺฐํ พุชฺฌนฺตีติ โพชฺฌงฺคา.\csromanspacer{} ปหานสฺส ปริจฺจาคฏฺฐํ พุชฺฌนฺตีติ โพชฺฌงฺคา.\csromanspacer{} ภาวนาย เอกรสฏฺฐํ พุชฺฌนฺตีติ โพชฺฌงฺคา.\csromanspacer{} สจฺฉิกิริยาย ผสฺสนฏฺฐํ พุชฺฌนฺตีติ โพชฺฌงฺคา.\csromanspacer{} ขนฺธานํ ขนฺธฏฺฐํ พุชฺฌนฺตีติ โพชฺฌงฺคา.\csromanspacer{} ธาตูนํ ธาตุฏฺฐํ พุชฺฌนฺตีติ โพชฺฌงฺคา.\csromanspacer{} อายตนานํ อายตนฏฺฐํ พุชฺฌนฺตีติ โพชฺฌงฺคา.\csromanspacer{} สงฺขตานํ สงฺขตฏฺฐํ พุชฺฌตีติ โพชฺฌงฺคา.\csromanspacer{} อสงฺขตสฺส อสงฺขตฏฺฐํ พุชฺฌนฺตีติ โพชฺฌงฺคา.}

\csromanheader{2}{3. โพชฺฌงฺคกถา}
\prose{\csromanfolio{307}จิตฺตฏฺฐํ พุชฺฌนฺตีติ โพชฺฌงฺคา.\csromanspacer{} จิตฺตานนฺตริยฏฺฐํ พุชฺฌนฺตีติ โพชฺฌงฺคา.\csromanspacer{} จิตฺตสฺส วุฏฺฐานฏฺฐํ พุชฺฌนฺตีติ โพชฺฌงฺคา.\csromanspacer{} จิตฺตสฺส วิวฏฺฏนฏฺฐํ\footnote{วิวชฺชนฏฺฐํ (สฺยา)} พุชฺฌนฺตีติ โพชฺฌงฺคา.\csromanspacer{} จิตฺตสฺส เหตุฏฺฐํ พุชฺฌนฺตีติ โพชฺฌงฺคา.\csromanspacer{} จิตฺตสฺส ปจฺจยฏฺฐํ พุชฺฌนฺตีติ โพชฺฌงฺคา.\csromanspacer{} จิตฺตสฺส วตฺถุฏฺฐํ พุชฺฌนฺตีติ โพชฺฌงฺคา.\csromanspacer{} จิตฺตสฺส ภูมฏฺฐํ\footnote{ภุมฺมฏฺฐํ (สฺยา) 18 ปิฏฺเฐ ปสฺสิตพฺพา.} พุชฺฌนฺตีติ โพชฺฌงฺคา.\csromanspacer{} จิตฺตสฺส อารมฺมณฏฺฐํ พุชฺฌนฺตีติ โพชฺฌงฺคา.\csromanspacer{} จิตฺตสฺส โคจรฏฺฐํ พุชฺฌนฺตีติ โพชฺฌงฺคา.\csromanspacer{} จิตฺตสฺส จริยฏฺฐํ พุชฺฌนฺตีติ โพชฺฌงฺคา.\csromanspacer{} จิตฺตสฺส คตฏฺฐํ พุชฺฌนฺตีติ โพชฺฌงฺคา.\csromanspacer{} จิตฺตสฺส อภินีหารฏฺฐํ พุชฺฌนฺตีติ โพชฺฌงฺคา.\csromanspacer{} จิตฺตสฺส นิยฺยานฏฺฐํ พุชฺฌนฺตีติ โพชฺฌงฺคา.\csromanspacer{} จิตฺตสฺส นิสฺสรณฏฺฐํ พุชฺฌนฺตีติ โพชฺฌงฺคา.}

\prose{เอกตฺเต อาวชฺชนฏฺฐํ พุชฺฌนฺตีติ โพชฺฌงฺคา.\csromanspacer{} เอกตฺเต วิชานนฏฺฐํ พุชฺฌนฺตีติ โพชฺฌงฺคา.\csromanspacer{} เอกตฺเต ปชานนฏฺฐํ พุชฺฌนฺตีติ โพชฺฌงฺคา.\csromanspacer{} เอกตฺเต สญฺชานนฏฺฐํ พุชฺฌนฺตีติ โพชฺฌงฺคา.\csromanspacer{} เอกตฺเต เอโกทฏฺฐํ พุชฺฌนฺตีติ โพชฺฌงฺคา. ( )\footnote{(เอกตฺเต อุปนิพนฺธนฏฺฐํ พุชฺฌนฺตีติ โพชฺฌงฺคา) (สพฺพตฺถ)} เอกตฺเต ปกฺขนฺทนฏฺฐํ พุชฺฌนฺตีติ โพชฺฌงฺคา.\csromanspacer{} เอกตฺเต ปสีทนฏฺฐํ พุชฺฌนฺตีติ โพชฺฌงฺคา.\csromanspacer{} เอกตฺเต สนฺติฏฺฐนฏฺฐํ พุชฺฌนฺตีติ โพชฺฌงฺคา.\csromanspacer{} เอกตฺเต วิมุจฺจนฏฺฐํ พุชฺฌนฺตีติ โพชฺฌงฺคา.\csromanspacer{} เอกตฺเต “เอตํ สนฺตนฺ”ติ ปสฺสนฏฺฐํ พุชฺฌนฺตีติ โพชฺฌงฺคา.\csromanspacer{} เอกตฺเต ยานีกตฏฺฐํ พุชฺฌนฺตีติ โพชฺฌงฺคา.\csromanspacer{} เอกตฺเต วตฺถุกตฏฺฐํ พุชฺฌนฺตีติ โพชฺฌงฺคา.\csromanspacer{} เอกตฺเต อนุฏฺฐิตฏฺฐํ พุชฺฌนฺตีติ โพชฺฌงฺคา.\csromanspacer{} เอกตฺเต ปริจิตฏฺฐํ พุชฺฌนฺตีติ โพชฺฌงฺคา.\csromanspacer{} เอกตฺเต สุสมารทฺธฏฺฐํ พุชฺฌนฺตีติ โพชฺฌงฺคา.\csromanspacer{} เอกตฺเต ปริคฺคหฏฺฐํ พุชฺฌนฺตีติ โพชฺฌงฺคา.\csromanspacer{} เอกตฺเต ปริวารฏฺฐํ พุชฺฌนฺตีติ โพชฺฌงฺคา.\csromanspacer{} เอกตฺเต ปริปูรณฏฺฐํ พุชฺฌนฺตีติ โพชฺฌงฺคา.\csromanspacer{} เอกตฺเต สโมธานฏฺฐํ พุชฺฌนฺตีติ โพชฺฌงฺคา.\csromanspacer{} เอกตฺเต อธิฏฺฐานฏฺฐํ พุชฺฌนฺตีติ โพชฺฌงฺคา.\csromanspacer{} เอกตฺเต อาเสวนฏฺฐํ พุชฺฌนฺตีติ โพชฺฌงฺคา.\csromanspacer{} เอกตฺเต ภาวนฏฺฐํ พุชฺฌนฺตีติ โพชฺฌงฺคา.\csromanspacer{} เอกตฺเต พหุลีกมฺมฏฺฐํ พุชฺฌนฺตีติ โพชฺฌงฺคา.\csromanspacer{} เอกตฺเต สุสมุคฺคตฏฺฐํ พุชฺฌนฺตีติ โพชฺฌงฺคา.\csromanspacer{} เอกตฺเต สุวิมุตฺตฏฺฐํ พุชฺฌนฺตีติ โพชฺฌงฺคา.\csromanspacer{} เอกตฺเต พุชฺฌนฏฺฐํ พุชฺฌนฺตีติ โพชฺฌงฺคา.เอกตฺเต อนุพุชฺฌนฏฺฐํ พุชฺฌนฺตีติ โพชฺฌงฺคา.\csromanspacer{} เอกตฺเต ปฏิพุชฺฌนฏฺฐํ พุชฺฌนฺตีติ โพชฺฌงฺคา.\csromanspacer{} เอกตฺเต สมฺพุชฺฌนฏฺฐํ พุชฺฌนฺตีติ โพชฺฌงฺคา.\csromanspacer{} เอกตฺเต โพธนฏฺฐํ พุชฺฌนฺตีติ โพชฺฌงฺคา.\csromanspacer{} เอกตฺเต อนุโพธนฏฺฐํ พุชฺฌนฺตีติ โพชฺฌงฺคา.\csromanspacer{} เอกตฺเต ปฏิโพธนฏฺฐํ พุชฺฌนฺตีติ โพชฺฌงฺคา.\csromanspacer{} เอกตฺเต สมฺพุชฺฌนฏฺฐํ พุชฺฌนฺตีติ โพชฺฌงฺคา.\csromanspacer{} เอกตฺเต โพธิปกฺขิยฏฺฐํ พุชฺฌนฺตีติ โพชฺฌงฺคา.\csromanspacer{} เอกตฺเต อนุโพธิปกฺขิยฏฺฐํ พุชฺฌนฺตีติ โพชฺฌงฺคา.\csromanspacer{} เอกตฺเต ปฏิโพธิปกฺขิยฏฺฐํ พุชฺฌนฺตีติ โพชฺฌงฺคา.\csromanspacer{} เอกตฺเต สมฺโพธิปกฺขิยฏฺฐํ พุชฺฌนฺตีติ โพชฺฌงฺคา.\csromanspacer{} เอกตฺเต}

\prose{\csromanfolio{308}โชตนฏฺฐํ พุชฺฌนฺตีติ โพชฺฌงฺคา.\csromanspacer{} เอกตฺเต อุชฺโชตนฏฺฐํ พุชฺฌนฺตีติ โพชฺฌงฺคา.\csromanspacer{} เอกตฺเต อนุโชตนฏฺฐํ พุชฺฌนฺตีติ โพชฺฌงฺคา.\csromanspacer{} เอกตฺเต ปฏิโชตนฏฺฐํ พุชฺฌนฺตีติ โพชฺฌงฺคา.\csromanspacer{} เอกตฺเต สญฺโชตนฏฺฐํ พุชฺฌนฺตีติ โพชฺฌงฺคา.}

\prose{ปตาปนฏฺฐํ พุชฺฌนฺตีติ โพชฺฌงฺคา.\csromanspacer{} วิโรจนฏฺฐํ พุชฺฌนฺตีติ โพชฺฌงฺคา.\csromanspacer{} กิเลสานํ สนฺตาปนฏฺฐํ พุชฺฌนฺตีติ โพชฺฌงฺคา.\csromanspacer{} อมลฏฺฐํ พุชฺฌนฺตีติ โพชฺฌงฺคา.\csromanspacer{} วิมลฏฺฐํ พุชฺฌนฺตีติ โพชฺฌงฺคา.\csromanspacer{} นิมฺมลฏฺฐํ พุชฺฌนฺตีติ โพชฺฌงฺคา.\csromanspacer{} สมฏฺฐํ พุชฺฌนฺตีติ โพชฺฌงฺคา.\csromanspacer{} สมยฏฺฐํ พุชฺฌนฺตีติ โพชฺฌงฺคา, วิเวกฏฺฐํ พุชฺฌนฺตีติ โพชฺฌงฺคา, วิเวกจริยฏฺฐํ พุชฺฌนฺตีติ โพชฺฌงฺคา, วิราคฏฺฐํ พุชฺฌนฺตีติ โพชฺฌงฺคา, วิราคจริยฏฺฐํ พุชฺฌนฺตีติ โพชฺฌงฺคา, นิโรธฏฺฐํ พุชฺฌนฺตีติ โพชฺฌงฺคา, นิโรธจริยฏฺฐํ พุชฺฌนฺตีติ โพชฺฌงฺคา, โวสคฺคฏฺฐํ พุชฺฌนฺตีติ โพชฺฌงฺคา, โวสคฺคจริยฏฺฐํ พุชฺฌนฺตีติ โพชฺฌงฺคา, วิมุตฺตฏฺฐํ พุชฺฌนฺตีติ โพชฺฌงฺคา, วิมุตฺติจริยฏฺฐํ พุชฺฌนฺตีติ โพชฺฌงฺคา.}

\prose{ฉนฺทฏฺฐํ พุชฺฌนฺตีติ โพชฺฌงฺคา, ฉนฺทสฺส มูลฏฺฐํ พุชฺฌนฺตีติ โพชฺฌงฺคา, ฉนฺทสฺส ปาทฏฺฐํ พุชฺฌนฺตีติ โพชฺฌงฺคา, ฉนฺทสฺส ปธานฏฺฐํ พุชฺฌนฺตีติ โพชฺฌงฺคา, ฉนฺทสฺส อิชฺฌนฏฺฐํ พุชฺฌนฺตีติ โพชฺฌงฺคา, ฉนฺทสฺส อธิโมกฺขฏฺฐํ พุชฺฌนฺตีติ โพชฺฌงฺคา, ฉนฺทสฺส ปคฺคหฏฺฐํ พุชฺฌนฺตีติ โพชฺฌงฺคา, ฉนฺทสฺส อุปฏฺฐานฏฺฐํ พุชฺฌนฺตีติ โพชฺฌงฺคา, ฉนฺทสฺส อวิกฺเขปฏฺฐํ พุชฺฌนฺตีติ โพชฺฌงฺคา, ฉนฺทสฺส ทสฺสนฏฺฐํ พุชฺฌนฺตีติ โพชฺฌงฺคา.}

\prose{วีริยฏฺฐํ พุชฺฌนฺตีติ โพชฺฌงฺคา ฯเปฯ.\csromanspacer{} จิตฺตฏฺฐํ พุชฺฌนฺตีติ โพชฺฌงฺคา ฯเปฯ.\csromanspacer{} วีมํสฏฺฐํ พุชฺฌนฺตีติ โพชฺฌงฺคา, วีมํสาย มูลฏฺฐํ พุชฺฌนฺตีติ โพชฺฌงฺคา, วีมํสาย ปาทฏฺฐํ พุชฺฌนฺตีติ โพชฺฌงฺคา, วีมํสาย ปธานฏฺฐํ พุชฺฌนฺตีติ โพชฺฌงฺคา.\csromanspacer{} วีมํสาย อิชฺฌนฏฺฐํ พุชฺฌนฺตีติ โพชฺฌงฺคา, วีมํสาย อธิโมกฺขฏฺฐํ พุชฺฌนฺตีติ โพชฺฌงฺคา, วีมํสาย ปคฺคหฏฺฐํ พุชฺฌนฺตีติ โพชฺฌงฺคา, วีมํสาย อุปฏฺฐานฏฺฐํ พุชฺฌนฺตีติ โพชฺฌงฺคา, วีมํสาย อวิกฺเขปฏฺฐํ พุชฺฌนฺตีติ โพชฺฌงฺคา, วีมํสาย ทสฺสนฏฺฐํ พุชฺฌนฺตีติ โพชฺฌงฺคา.}

\prose{ทุกฺขฏฺฐํ พุชฺฌนฺตีติ โพชฺฌงฺคา.\csromanspacer{} ทุกฺขสฺส ปีฬนฏฺฐํ พุชฺฌนฺตีติ โพชฺฌงฺคา, ทุกฺขสฺส สงฺขตฏฺฐํ พุชฺฌนฺตีติ โพชฺฌงฺคา, ทุกฺขสฺส สนฺตาปฏฺฐํ พุชฺฌนฺตีติ โพชฺฌงฺคา, ทุกฺขสฺส วิปริณามฏฺฐํ พุชฺฌนฺตีติ โพชฺฌงฺคา.\csromanspacer{} สมุทยฏฺฐํ พุชฺฌนฺตีติ โพชฺฌงฺคา.\csromanspacer{} สมุทยสฺส อายูหนฏฺฐํ.\csromanspacer{} นิทานฏฺฐํ.\csromanspacer{} สญฺโญคฏฺฐํ.\csromanspacer{} ปลิโพธฏฺฐํ พุชฺฌนฺตีติ โพชฺฌงฺคา.\csromanspacer{} นิโรธฏฺฐํ พุชฺฌนฺตีติ โพชฺฌงฺคา, นิโรธสฺส นิสฺสรณฏฺฐํ.\csromanspacer{} วิเวกฏฺฐํ.\csromanspacer{} อสงฺขตฏฺฐํ.\csromanspacer{} อมตฏฺฐํ พุชฺฌนฺตีติ โพชฺฌงฺคา.\csromanspacer{} มคฺคฏฺฐํ พุชฺฌนฺตีติ โพชฺฌงฺคา, มคฺคสฺส นิยฺยานฏฺฐํ.\csromanspacer{} เหตุฏฺฐํ.\csromanspacer{} ทสฺสนฏฺฐํ.\csromanspacer{} อาธิปเตยฺยฏฺฐํ พุชฺฌนฺตีติ โพชฺฌงฺคา.}

\csromanheader{2}{3. โพชฺฌงฺคกถา}
\prose{\csromanfolio{309}ตถฏฺฐํ พุชฺฌนฺตีติ โพชฺฌงฺคา, อนตฺตฏฺฐํ พุชฺฌนฺตีติ โพชฺฌงฺคา, สจฺจฏฺฐํ พุชฺฌนฺตีติ โพชฺฌงฺคา, ปฏิเวธฏฺฐํ พุชฺฌนฺตีติ โพชฺฌงฺคา, อภิชานนฏฺฐํ พุชฺฌนฺตีติ โพชฺฌงฺคา, ปริชานนฏฺฐํ พุชฺฌนฺตีติ โพชฺฌงฺคา, ธมฺมฏฺฐํ พุชฺฌนฺตีติ โพชฺฌงฺคา, ธาตุฏฺฐํ พุชฺฌนฺตีติ โพชฺฌงฺคา, ญาตฏฺฐํ พุชฺฌนฺตีติ โพชฺฌงฺคา, สจฺฉิกิริยฏฺฐํ พุชฺฌนฺตีติ โพชฺฌงฺคา, ผสฺสนฏฺฐํ พุชฺฌนฺตีติ โพชฺฌงฺคา, อภิสมยฏฺฐํ พุชฺฌนฺตีติ โพชฺฌงฺคา.}

\prose{เนกฺขมฺมํ พุชฺฌนฺตีติ โพชฺฌงฺคา, อพฺยาปาทํ พุชฺฌนฺตีติ โพชฺฌงฺคา, อาโลกสญฺญํ พุชฺฌนฺตีติ โพชฺฌงฺคา, อวิกฺเขปํ พุชฺฌนฺตีติ โพชฺฌงฺคา, ธมฺมววตฺถานํ พุชฺฌนฺตีติ โพชฺฌงฺคา, ญาณํ พุชฺฌนฺตีติ โพชฺฌงฺคา, ปาโมชฺชํ พุชฺฌนฺตีติ โพชฺฌงฺคา, ปฐมํ ฌานํ พุชฺฌนฺตีติ โพชฺฌงฺคา ฯเปฯ อรหตฺตมคฺคํ พุชฺฌนฺตีติ โพชฺฌงฺคา.\csromanspacer{} อรหตฺตผลสมาปตฺติํ พุชฺฌนฺตีติ โพชฺฌงฺคา.}

\prose{อธิโมกฺขฏฺเฐน สทฺธินฺทฺริยํ พุชฺฌนฺตีติ โพชฺฌงฺคา ฯเปฯ.\csromanspacer{} ทสฺสนฏฺเฐน ปญฺญินฺทฺริยํ พุชฺฌนฺตีติ โพชฺฌงฺคา.\csromanspacer{}\csromanspacer{}อสฺสทฺธิเย อกมฺปิยฏฺเฐน สทฺธาพลํ พุชฺฌนฺตีติ โพชฺฌงฺคา ฯเปฯ.\csromanspacer{} อวิชฺชาย อกมฺปิยฏฺเฐน ปญฺญาพลํ พุชฺฌนฺตีติ โพชฺฌงฺคา.\csromanspacer{}\csromanspacer{}อุปฏฺฐานฏฺเฐน สติสมฺโพชฺฌงฺคํ พุชฺฌนฺตีติ โพชฺฌงฺคา ฯเปฯ.\csromanspacer{} ปฏิสงฺขานฏฺเฐน อุเปกฺขาสมฺโพชฺฌงฺคํ พุชฺฌนฺตีติ โพชฺฌงฺคา.}

\prose{ทสฺสนฏฺเฐน สมฺมาทิฏฺฐิํ พุชฺฌนฺตีติ โพชฺฌงฺคา ฯเปฯ.\csromanspacer{} อวิกฺเขปฏฺเฐน สมฺมาสมาธิํ พุชฺฌนฺตีติ โพชฺฌงฺคา.\csromanspacer{}\csromanspacer{}อาธิปเตยฺยฏฺเฐน อินฺทฺริยํ พุชฺฌนฺตีติ โพชฺฌงฺคา, อกมฺปิยฏฺเฐน พลํ พุชฺฌนฺตีติ โพชฺฌงฺคา, นิยฺยานฏฺฐํ พุชฺฌนฺตีติ โพชฺฌงฺคา, เหตุฏฺเฐน มคฺคํ พุชฺฌนฺตีติ โพชฺฌงฺคา, อุปฏฺฐานฏฺเฐน สติปฏฺฐานํ พุชฺฌนฺตีติ โพชฺฌงฺคา, ปทหนฏฺเฐน สมฺมปฺปธานํ พุชฺฌนฺตีติ โพชฺฌงฺคา, อิชฺฌนฏฺเฐน อิทฺธิปาทํ พุชฺฌนฺตีติ โพชฺฌงฺคา, ตถฏฺเฐน สจฺจํ พุชฺฌนฺตีติ โพชฺฌงฺคา, อวิกฺเขปฏฺเฐน สมถํ พุชฺฌนฺตีติ โพชฺฌงฺคา.\csromanspacer{} อนุปสฺสนฏฺเฐน วิปสฺสนํ ฯเปฯ เอกรสฏฺเฐน สมถวิปสฺสนํ พุชฺฌนฺตีติ โพชฺฌงฺคา, อนติวตฺตนฏฺเฐน ยุคนทฺธํ พุชฺฌนฺตีติ โพชฺฌงฺคา, สํวรฏฺเฐน สีลวิสุทฺธิํ พุชฺฌนฺตีติ โพชฺฌงฺคา, อวิกฺเขปฏฺเฐน จิตฺตวิสุทฺธิํ พุชฺฌนฺตีติ โพชฺฌงฺคา, ทสฺสนฏฺเฐน ทิฏฺฐิวิสุทฺธิํ พุชฺฌนฺตีติ โพชฺฌงฺคา, มุตฺตฏฺเฐน วิโมกฺขํ พุชฺฌนฺตีติ โพชฺฌงฺคา, ปฏิเวธฏฺเฐน วิชฺชํ พุชฺฌนฺตีติ โพชฺฌงฺคา, ปริจฺจาคฏฺเฐน วิมุตฺติํ พุชฺฌนฺตีติ โพชฺฌงฺคา, สมุจฺเฉทฏฺเฐน ขเย ญาณํ พุชฺฌนฺตีติ โพชฺฌงฺคา, ปฏิปฺปสฺสทฺธฏฺเฐน อนุปฺปาเท ญาณํ พุชฺฌนฺตีติ โพชฺฌงฺคา.}

\prose{ฉนฺทํ มูลฏฺเฐน พุชฺฌนฺตีติ โพชฺฌงฺคา, มนสิการํ สมุฏฺฐานฏฺเฐน พุชฺฌนฺตีติ โพชฺฌงฺคา, ผสฺสํ สโมธานฏฺเฐน พุชฺฌนฺตีติ โพชฺฌงฺคา, เวทนํ สโมสรณฏฺเฐน พุชฺฌนฺตีติ โพชฺฌงฺคา, สมาธิํ ปมุขฏฺเฐน พุชฺฌนฺตีติ โพชฺฌงฺคา, สติํ}

\prose{\csromanfolio{310}อาธิปเตยฺยฏฺเฐน พุชฺฌนฺตีติ โพชฺฌงฺคา, ปญฺญํ ตทุตฺตรฏฺเฐน พุชฺฌนฺตีติ โพชฺฌงฺคา, วิมุตฺติํ สารฏฺเฐน พุชฺฌนฺตีติ โพชฺฌงฺคา, อมโตคธํ นิพฺพานํ ปริโยสานฏฺเฐน พุชฺฌนฺตีติ โพชฺฌงฺคา.}

\proseitem{20}{สาวตฺถินิทานํ.\csromanspacer{}\csromanspacer{}ตตฺร โข อายสฺมา สาริปุตฺโต ภิกฺขู อามนฺเตสิ “อาวุโส ภิกฺขโว”ติ\footnote{อาวุโสติ (สฺยา) สํ 3. 64 ปิฏฺเฐ ปสฺสิตพฺพา.}. “อาวุโส”ติ โข เต ภิกฺขู อายสฺมโต สาริปุตฺตสฺส ปจฺจสฺโสสุํ.\csromanspacer{} อายสฺมา สาริปุตฺโต เอตทโวจ\csromandash{}}

\prose{สตฺติเม อาวุโส โพชฺฌงฺคา.\csromanspacer{} กตเม สตฺต? สติสมฺโพชฺฌงฺโค ธมฺมวิจยสมฺโพชฺฌงฺโค ฯเปฯ อุเปกฺขาสมฺโพชฺฌงฺโค, อิเม โข อาวุโส สตฺต โพชฺฌงฺคา.\csromanspacer{} อิเมสํ ขฺวาหํ อาวุโส สตฺตนฺนํ โพชฺฌงฺคานํ เยน เยน โพชฺฌงฺเคน อากงฺขามิ ปุพฺพณฺหสมยํ วิหริตุํ, เตน เตน โพชฺฌงฺเคน ปุพฺพณฺหสมยํ วิหรามิ.\csromanspacer{} เยน เยน โพชฺฌงฺเคน อากงฺขามิ มชฺฌนฺหิกสมยํ ฯเปฯ สายนฺหสมยํ วิหริตุํ, เตน เตน โพชฺฌงฺเคน สายนฺหสมยํ วิหรามิ.\csromanspacer{} สติสมฺโพชฺฌงฺโค อิติ เจ เม อาวุโส โหติ, “อปฺปมาโณ”ติ เม โหติ, “สุสมารทฺโธ”ติ เม โหติ, “ติฏฺฐนฺตํ จ นํ\footnote{ติฏฺฐนฺตํ จรํ (สฺยา) สํ 3. 64 ปิฏฺเฐ ปสฺสิตพฺพา.} ติฏฺฐตี”ติ ปชานามิ, สเจปิ เม จวติ, “อิทปฺปจฺจยา เม จวตี”ติ ปชานามิ.\csromanspacer{} ธมฺมวิจยสมฺโพชฺฌงฺโค ฯเปฯ อุเปกฺขาสมฺโพชฺฌงฺโค อิติ เจ เม อาวุโส โหติ, “อปฺปมาโณ”ติ เม โหติ, “สุสมารทฺโธ”ติ เม โหติ, “ติฏฺฐนฺตํ จ นํ ติฏฺฐตี”ติ ปชานามิ.\csromanspacer{} สเจปิ เม จวติ, “อิทปฺปจฺจยา เม จวตี”ติ ปชานามิ.}

\prose{เสยฺยถาปิ อาวุโส รญฺโญ วา ราชมหามตฺตสฺส วา นานารตฺตานํ ทุสฺสานํ ทุสฺสกรณฺฑโก ปูโร อสฺส, โส ยญฺญเทว ทุสฺสยุคํ อากงฺเขยฺย ปุพฺพณฺหสมยํ ปารุปิตุํ, ตํ ตเทว ทุสฺสยุคํ ปุพฺพณฺหสมยํ ปารุเปยฺย.\csromanspacer{} ยญฺญเทว ทุสฺสยุคํ อากงฺเขยฺย มชฺฌนฺหิกสมยํ ฯเปฯ สายนฺหสมยํ ปารุปิตุํ, ตํ ตเทว ทุสฺสยุคํ สายนฺหสมยํ ปารุเปยฺย, เอวเมว ขฺวาหํ อาวุโส อิเมสํ สตฺตนฺนํ โพชฺฌงฺคานํ เยน เยน โพชฺฌงฺเคน อากงฺขามิ ปุพฺพณฺหสมยํ วิหริตุํ, เตน เตน โพชฺฌงฺเคน ปุพฺพณฺหสมยํ วิหรามิ.\csromanspacer{} เยน เยน โพชฺฌงฺเคน อากงฺขามิ มชฺฌนฺหิกสมยํ ฯเปฯ}

\vspace{\tipitakaparskip}
\csromancenter{โพชฺฌงฺคกถา สายนฺหสมยํ วิหริตุํ, เตน เตน โพชฺฌงฺเคน สายนฺหสมยํ วิหรามิ.\csromanspacer{} สติสมฺโพชฺฌงฺโค อิติ เจ เม อาวุโส โหติ, “อปฺปมาโณ”ติ เม โหติ, “สุสมารทฺโธ”ติ เม โหติ, “ติฏฺฐนฺตํ จ นํ ติฏฺฐตี”ติ ปชานามิ, สเจปิ เม จวติ, “อิทปฺปจฺจยา เม จวตี”ติ ปชานามิ ฯเปฯ อุเปกฺขาสมฺโพชฺฌงฺโค อิติ เจ เม อาวุโส โหติ, “อปฺปมาโณ”ติ เม โหติ, “สุสมารทฺโธ”ติ เม โหติ, “ติฏฺฐนฺตํ จ นํ ติฏฺฐตี”ติ ปชานามิ, สเจปิ เม จวติ, “อิทปฺปจฺจยา เม จวตี”ติ ปชานามิ.}
\csromanmatikaanchor{mtk.94}
\csromantocmarkref{section}{สุตฺตนฺตนิทฺเทส}{mtk.94}
\bookmark[dest={mtk.94},level=1]{สุตฺตนฺตนิทฺเทส}
\csromanheader{1}{\textbf{สุตฺตนฺตนิทฺเทส}}
\proseitem{21}{\csromanfolio{311}กถํ \textbf{สติสมฺโพชฺฌงฺโค อิติ เจ เม โหตี}ติ โพชฺฌงฺโค? ยาวตา นิโรธู’ปฏฺฐาติ, ตาวตา สติสมฺโพชฺฌงฺโค, อิติ เจ เม โหตีติ โพชฺฌงฺโค.\csromanspacer{} เสยฺยถาปิ เตลปฺปทีปสฺส ฌายโต ยาวตา อจฺจิ, ตาวตา วณฺโณ.\csromanspacer{} ยาวตา วณฺโณ, ตาวตา อจฺจิ.\csromanspacer{} เอวเมว ยาวตา นิโรธู’ปฏฺฐาติ, ตาวตา สติสมฺโพชฺฌงฺโค.\csromanspacer{} อิติ เจ เม โหตีติ โพชฺฌงฺโค.}

\prose{กถํ \textbf{อปฺปมาโณ อิติ เจ เม โหตี}ติ โพชฺฌงฺโค? ปมาณพทฺธา\footnote{ปมาณวนฺตา (สฺยา), อฏฺฐกถา โอโลเกตพฺพา.} กิเลสา สพฺเพ จ ปริยุฏฺฐานา เย จ สงฺขารา โปโนภวิกา, อปฺปมาโณ นิโรโธ อจลฏฺเฐน อสงฺขตฏฺเฐน ยาวตา นิโรธู’ปฏฺฐาติ, ตาวตา \textbf{อปฺปมาโณ อิติ เจ เม โหตี}ติ โพชฺฌงฺโค.}

\prose{กถํ \textbf{สุสมารทฺโธ อิติ เจ เม โหตี}ติ โพชฺฌงฺโค? วิสมา กิเลสา สพฺเพ จ ปริยุฏฺฐานา เย จ สงฺขารา โปโนภวิกา, สมธมฺโม นิโรโธ สนฺตฏฺเฐน ปณีตฏฺเฐน.\csromanspacer{} ยาวตา นิโรธู’ปฏฺฐาติ, ตาวตา สุสมารทฺโธ, อิติ เจ เม โหตีติ โพชฺฌงฺโค.}

\prose{กถํ \textbf{“ติฏฺฐนฺตํ จ นํ ติฏฺฐตี”ติ ปชานามิ, สเจปิ จวติ, “อิทปฺปจฺจยา จวตี”ติ ปชานามิ?} กติหากาเรหิ สติสมฺโพชฺฌงฺโค ติฏฺฐติ, กติหากาเรหิ สติสมฺโพชฺฌงฺโค จวติ? อฏฺฐหากาเรหิ สติสมฺโพชฺฌงฺโค ติฏฺฐติ, อฏฺฐหากาเรหิ สติสมฺโพชฺฌงฺโค จวติ.}

\prose{\csromanfolio{312}กตเมหิ \textbf{อฏฺฐหากาเรหิ สติสมฺโพชฺฌงฺโค ติฏฺฐติ?} อนุปฺปาทํ อาวชฺชิตตฺตา สติสมฺโพชฺฌงฺโค ติฏฺฐติ, อุปฺปาทํ อนาวชฺชิตตฺตา สติสมฺโพชฺฌงฺโค ติฏฺฐติ, อปฺปวตฺตํ อาวชฺชิตตฺตา สติสมฺโพชฺฌงฺโค ติฏฺฐติ, ปวตฺตํ อนาวชฺชิตตฺตา สติสมฺโพชฺฌงฺโค ติฏฺฐติ, อนิมิตฺตํ อาวชฺชิตตฺตา สติสมฺโพชฺฌงฺโค ติฏฺฐติ, นิมิตฺตํ อนาวชฺชิตตฺตา สติสมฺโพชฺฌงฺโค ติฏฺฐติ, นิโรธํ อาวชฺชิตตฺตา สติสมฺโพชฺฌงฺโค ติฏฺฐติ, สงฺขาเร อนาวชฺชิตตฺตา สติสมฺโพชฺฌงฺโค ติฏฺฐติ, อิเมหิ อฏฺฐหากาเรหิ สติสมฺโพชฺฌงฺโค ติฏฺฐติ.}

\prose{กตเมหิ \textbf{อฏฺฐหากาเรหิ สติสมฺโพชฺฌงฺโค จวติ?} อุปฺปาทํ อาวชฺชิตตฺตา สติสมฺโพชฺฌงฺโค จวติ, อนุปฺปาทํ อนาวชฺชิตตฺตา สติสมฺโพชฺฌงฺโค จวติ, ปวตฺตํ อาวชฺชิตตฺตา สติสมฺโพชฺฌงฺโค จวติ, อปฺปวตฺตํ อนาวชฺชิตตฺตา สติสมฺโพชฺฌงฺโค จวติ, นิมิตฺตํ อาวชฺชิตตฺตา สติสมฺโพชฺฌงฺโค จวติ, อนิมิตฺตํ อนาวชฺชิตตฺตา สติสมฺโพชฺฌงฺโค จวติ, สงฺขาเร อาวชฺชิตตฺตา สติสมฺโพชฺฌงฺโค จวติ, นิโรธํ อนาวชฺชิตตฺตา สติสมฺโพชฺฌงฺโค จวติ.\csromanspacer{} อิเมหิ อฏฺฐหากาเรหิ สติสมฺโพชฺฌงฺโค จวติ.\csromanspacer{}\csromanspacer{}เอวํ “ติฏฺฐนฺตํ จ นํ ติฏฺฐตี”ติ ปชานามิ.\csromanspacer{} สเจปิ จวติ, “อิทปฺปจฺจยา เม จวตี”ติ ปชานามิ ฯเปฯ.}

\prose{กถํ \textbf{อุเปกฺขาสมฺโพชฺฌงฺโค อิติ เจ เม โหตี}ติ โพชฺฌงฺโค? ยาวตา นิโรธู’ปฏฺฐาติ, ตาวตา \textbf{อุเปกฺขาสมฺโพชฺฌงฺโค อิติ เจ เม โหตี}ติ โพชฺฌงฺโค.\csromanspacer{} เสยฺยถาปิ เตลปฺปทีปสฺส ฌายโต ยาวตา อจฺจิ, ตาวตา วณฺโณ.\csromanspacer{} ยาวตา วณฺโณ, ตาวตา อจฺจิ.\csromanspacer{} เอวเมว ยาวตา นิโรธู’ปฏฺฐาติ, ตาวตา อุเปกฺขาสมฺโพชฺฌงฺโค อิติ เจ โหตีติ โพชฺฌงฺโค.}

\prose{กถํ \textbf{อปฺปมาโณ อิติ เจ โหตี}ติ โพชฺฌงฺโค? ปมาณพทฺธา กิเลสา สพฺเพ จ ปริยุฏฺฐานา เย จ สงฺขารา โปโนภวิกา, อปฺปมาโณ นิโรโธ อจลฏฺเฐน อสงฺขตฏฺเฐน.\csromanspacer{} ยาวตา นิโรธู’ปฏฺฐาติ, ตาวตา \textbf{อปฺปมาโณ อิติ เจ โหตี}ติ โพชฺฌงฺโค.}

\prose{กถํ \textbf{สุสมารทฺโธ อิติ เจ โหตี}ติ โพชฺฌงฺโค? วิสมา กิเลสา สพฺเพ จ ปริยุฏฺฐานา เย จ สงฺขารา โปโนภวิกา, สมธมฺโม}

\vspace{\tipitakaparskip}
\csromancenter{เมตฺตากถา นิโรโธ สนฺตฏฺเฐน ปณีตฏฺเฐน.\csromanspacer{} ยาวตา นิโรธูปฏฺฐาติ, ตาวตา สุสมารทฺโธ, อิติ เจ โหตีติ โพชฺฌงฺโค.}
\prose{\csromanfolio{313}กถํ \textbf{“ติฏฺฐนฺตํ จ นํ ติฏฺฐตี”ติ ปชานามิ.}\csromanspacer{}\textbf{ สเจปิ จวติ, “อิทปฺปจฺจยา เม จวตี”ติ ปชานามิ?} กติหากาเรหิ อุเปกฺขาสมฺโพชฺฌงฺโค ติฏฺฐติ, กติหากาเรหิ อุเปกฺขาสมฺโพชฺฌงฺโค จวติ? อฏฺฐหากาเรหิ อุเปกฺขาสมฺโพชฺฌงฺโค ติฏฺฐติ, อฏฺฐหากาเรหิ อุเปกฺขาสมฺโพชฺฌงฺโค จวติ.}

\prose{กตเมหิ \textbf{อฏฺฐหากาเรหิ อุเปกฺขาสมฺโพชฺฌงฺโค ติฏฺฐติ?} อนุปฺปาทํ อาวชฺชิตตฺตา อุเปกฺขาสมฺโพชฺฌงฺโค ติฏฺฐติ, อุปฺปาทํ อนาวชฺชิตตฺตา อุเปกฺขาสมฺโพชฺฌงฺโค ติฏฺฐติ, อปฺปวตฺตํ อาวชฺชิตตฺตา อุเปกฺขาสมฺโพชฺฌงฺโค ติฏฺฐติ, ปวตฺตํ อนาวชฺชิตตฺตา อุเปกฺขาสมฺโพชฺฌงฺโค ติฏฺฐติ, อนิมิตฺตํ อาวชฺชิตตฺตา อุเปกฺขาสมฺโพชฺฌงฺโค ติฏฺฐติ, นิมิตฺตํ อนาวชฺชิตตฺตา อุเปกฺขาสมฺโพชฺฌงฺโค ติฏฺฐติ, นิโรธํ อาวชฺชิตตฺตา อุเปกฺขาสมฺโพชฺฌงฺโค ติฏฺฐติ, สงฺขาเร อนาวชฺชิตตฺตา อุเปกฺขาสมฺโพชฺฌงฺโค ติฏฺฐติ.\csromanspacer{} อิเมหิ อฏฺฐหากาเรหิ อุเปกฺขาสมฺโพชฺฌงฺโค ติฏฺฐติ.}

\prose{กตเมหิ \textbf{อฏฺฐหากาเรหิ อุเปกฺขาสมฺโพชฺฌงฺโค จวติ?} อุปฺปาทํ อาวชฺชิตตฺตา อุเปกฺขาสมฺโพชฺฌงฺโค จวติ, อนุปฺปาทํ อนาวชฺชิตตฺตา อุเปกฺขาสมฺโพชฺฌงฺโค จวติ, ปวตฺตํ อาวชฺชิตตฺตา อุเปกฺขาสมฺโพชฺฌงฺโค จวติ, อปฺปวตฺตํ อนาวชฺชิตตฺตา อุเปกฺขาสมฺโพชฺฌงฺโค จวติ, นิมิตฺตํ อาวชฺชิตตฺตา อุเปกฺขาสมฺโพชฺฌงฺโค จวติ, อนิมิตฺตํ อนาวชฺชิตตฺตา อุเปกฺขาสมฺโพชฺฌงฺโค จวติ, สงฺขาเร อาวชฺชิตตฺตา อุเปกฺขาสมฺโพชฺฌงฺโค จวติ, นิโรธํ อนาวชฺชิตตฺตา อุเปกฺขาสมฺโพชฺฌงฺโค จวติ.\csromanspacer{} อิเมหิ อฏฺฐหากาเรหิ อุเปกฺขาสมฺโพชฺฌงฺโค จวติ, เอวํ “ติฏฺฐนฺตํ จ นํ ติฏฺฐตี”ติ ปชานามิ.\csromanspacer{} สเจปิ จวติ, “อิทปฺปจฺจยา เม จวตี”ติ ปชานามิ.}

\nitthitamruled{โพชฺฌงฺคกถา นิฏฺฐิตา.}
\csromanmatikaanchor{mtk.95}
\csromantocmarknopage{section}{4. เมตฺตากถา}{mtk.95}
\bookmark[dest={mtk.95},level=1]{4. เมตฺตากถา}
\csromanheader{1}{4. เมตฺตากถา}
\proseitem{22}{สาวตฺถินิทานํ.\csromanspacer{}\csromanspacer{}เมตฺตาย ภิกฺขเว เจโตวิมุตฺติยา อาเสวิตาย ภาวิตาย พหุลีกตาย ยานีกตาย วตฺถุกตาย \csromanfolio{314}อนุฏฺฐิตาย ปริจิตาย สุสมารทฺธาย เอกาทสานิสํสา ปาฏิกงฺขา\csromandash{}กตเม เอกาทส? สุขํ สุปติ, สุขํ ปฏิพุชฺฌติ, น ปาปกํ สุปินํ ปสฺสติ, มนุสฺสานํ ปิโย โหติ, อมนุสฺสานํ ปิโย โหติ, เทวตา รกฺขนฺติ, นาสฺส อคฺคิ วา วิสํ วา สตฺถํ วา กมติ, ตุวฏํ จิตฺตํ สมาธิยติ, มุขวณฺโณ วิปฺปสีทติ, อสมฺมูโฬฺห กาลํ กโรติ, อุตฺตริ\footnote{อุตฺตริํ (สฺยา, อิ) อํ 3. 542 ปิฏฺเฐ ปสฺสิตพฺพา.} อปฺปฏิวิชฺฌนฺโต พฺรหฺมโลกูปโค โหติ.\csromanspacer{} เมตฺตาย ภิกฺขเว เจโตวิมุตฺติยา อาเสวิตาย ภาวิตาย พหุลีกตาย ยานีกตาย วตฺถุกตาย อนุฏฺฐิตาย ปริจิตาย สุสมารทฺธาย อิเม เอกาทสานิสํสา ปาฏิกงฺขา.}

\prose{อตฺถิ อโนธิโส ผรณา เมตฺตาเจโตวิมุตฺติ, อตฺถิ โอธิโส ผรณา เมตฺตาเจโตวิมุตฺติ, อตฺถิ ทิสาผรณา เมตฺตาเจโตวิมุตฺติ.\csromanspacer{} กติหากาเรหิ อโนธิโส ผรณา เมตฺตาเจโตวิมุตฺติ, กติหากาเรหิ โอธิโส ผรณา เมตฺตาเจโตวิมุตฺติ, กติหากาเรหิ ทิสาผรณา เมตฺตาเจโตวิมุตฺติ? ปญฺจหากาเรหิ อโนธิโส ผรณา เมตฺตาเจโตวิมุตฺติ, สตฺตหากาเรหิ โอธิโส ผรณา เมตฺตาเจโตวิมุตฺติ, ทสหากาเรหิ ทิสาผรณา เมตฺตาเจโตวิมุตฺติ.}

\prose{กตเมหิ ปญฺจหากาเรหิ อโนธิโส ผรณา เมตฺตาเจโตวิมุตฺติ? สพฺเพ สตฺตา อเวรา อพฺยาปชฺชา\footnote{อพฺยาปชฺฌา (สฺยา)} อนีฆา สุขี อตฺตานํ ปริหรนฺตุ.\csromanspacer{} สพฺเพ ปาณา ฯเปฯ.\csromanspacer{} สพฺเพ ภูตา ฯเปฯ.\csromanspacer{} สพฺเพ ปุคฺคลา ฯเปฯ.\csromanspacer{} สพฺเพ อตฺตภาวปริยาปนฺนา อเวรา อพฺยาปชฺชา อนีฆา สุขี อตฺตานํ ปริหรนฺตูติ.\csromanspacer{} อิเมหิ ปญฺจหากาเรหิ อโนธิโส ผรณา เมตฺตาเจโตวิมุตฺติ.}

\prose{กตเมหิ สตฺตหากาเรหิ โอธิโส ผรณา เมตฺตาเจโตวิมุตฺติ, สพฺพา อิตฺถิโย อเวรา อพฺยาปชฺชา อนีฆา สุขี อตฺตานํ ปริหรนฺตุ.\csromanspacer{} สพฺเพ ปุริสา ฯเปฯ.\csromanspacer{} สพฺเพ อริยา ฯเปฯ.\csromanspacer{} สพฺเพ อนริยา ฯเปฯ.\csromanspacer{} สพฺเพ เทวา ฯเปฯ.\csromanspacer{} สพฺเพ มนุสฺสา ฯเปฯ.\csromanspacer{} สพฺเพ วินิปาติกา อเวรา อพฺยาปชฺชา อนีฆา สุขี อตฺตานํ ปริหรนฺตูติ อิเมหิ สตฺตหากาเรหิ โอธิโส ผรณา เมตฺตาเจโตวิมุตฺติ.}

\prose{กตเมหิ ทสหากาเรหิ ทิสาผรณา เมตฺตาเจโตวิมุตฺติ, สพฺเพ ปุรตฺถิมาย ทิสาย สตฺตา อเวรา อพฺยาปชฺชา อนีฆา สุขี อตฺตานํ}

\vspace{\tipitakaparskip}
\csromancenter{เมตฺตากถา ปริหรนฺตุ.\csromanspacer{} สพฺเพ ปจฺฉิมาย ทิสาย สตฺตา ฯเปฯ.\csromanspacer{} สพฺเพ อุตฺตราย ทิสาย สตฺตา ฯเปฯ.\csromanspacer{} สพฺเพ ทกฺขิณาย ทิสาย สตฺตา ฯเปฯ.\csromanspacer{} สพฺเพ ปุรตฺถิมาย อนุทิสาย สตฺตา ฯเปฯ.\csromanspacer{} สพฺเพ ปจฺฉิมาย อนุทิสาย สตฺตา ฯเปฯ.\csromanspacer{} สพฺเพ อุตฺตราย อนุทิสาย สตฺตา ฯเปฯ.\csromanspacer{} สพฺเพ ทกฺขิณาย อนุทิสาย สตฺตา ฯเปฯ.\csromanspacer{} สพฺเพ เหฏฺฐิมาย ทิสาย สตฺตา ฯเปฯ.\csromanspacer{} สพฺเพ อุปริมาย ทิสาย สตฺตา อเวรา อพฺยาปชฺชา อนีฆา สุขี อตฺตานํ ปริหรนฺตุ.\csromanspacer{} สพฺเพ ปุรตฺถิมาย ทิสาย ปาณา ฯเปฯ.\csromanspacer{} ภูตา.\csromanspacer{} ปุคฺคลา.\csromanspacer{} อตฺตภาวปริยาปนฺนา.\csromanspacer{} สพฺพา อิตฺถิโย.\csromanspacer{} สพฺเพ ปุริสา.\csromanspacer{} สพฺเพ อริยา.\csromanspacer{} สพฺเพ อนริยา.\csromanspacer{} สพฺเพ เทวา.\csromanspacer{} สพฺเพ มนุสฺสา.\csromanspacer{} สพฺเพ วินิปาติกา อเวรา อพฺยาปชฺชา อนีฆา สุขี อตฺตานํ ปริหรนฺตุ.\csromanspacer{} สพฺเพ ปจฺฉิมาย ทิสาย วินิปาติกา ฯเปฯ.\csromanspacer{} สพฺเพ อุตฺตราย ทิสาย วินิปาติกา.\csromanspacer{} สพฺเพ ทกฺขิณาย ทิสาย วินิปาติกา.\csromanspacer{} สพฺเพ ปุรตฺถิมาย อนุทิสาย วินิปาติกา.\csromanspacer{} สพฺเพ ปจฺฉิมาย อนุทิสาย วินิปาติกา.\csromanspacer{} สพฺเพ อุตฺตราย อนุทิสาย วินิปาติกา.\csromanspacer{} สพฺเพ ทกฺขิณาย อนุทิสาย วินิปาติกา.\csromanspacer{} สพฺเพ เหฏฺฐิมาย ทิสาย วินิปาติกา.\csromanspacer{} สพฺเพ อุปริมาย ทิสาย วินิปาติกา อเวรา อพฺยาปชฺชา อนีฆา สุขี อตฺตานํ ปริหรนฺตูติ อิเมหิ ทสหากาเรหิ ทิสาผรณา เมตฺตาเจโตวิมุตฺติ.}
\csromanmatikaanchor{mtk.96}
\csromantocmarkref{subsection}{1. อินฺทฺริยวาร}{mtk.96}
\bookmark[dest={mtk.96},level=2]{1. อินฺทฺริยวาร}
\csromanheader{2}{1. อินฺทฺริยวาร}
\proseitem{23}{\csromanfolio{315}สพฺเพสํ สตฺตานํ ปีฬนํ วชฺเชตฺวา อปีฬนาย, อุปฆาตํ วชฺเชตฺวา อนุปฆาเตน, สนฺตาปํ วชฺเชตฺวา อสนฺตาเปน ปริยาทานํ วชฺเชตฺวา อปริยาทาเนน, วิเหสํ วชฺเชตฺวา อวิเหสาย สพฺเพ สตฺตา อเวริโน โหนฺตุ, มา เวริโน.\csromanspacer{} สุขิโน โหนฺตุ, มา ทุกฺขิโน.\csromanspacer{} สุขิตตฺตา โหนฺตุ, มา ทุกฺขิตตฺตาติ อิเมหิ อฏฺฐหากาเรหิ สพฺเพ สตฺเต เมตฺตายตีติ เมตฺตา.\csromanspacer{} ตํ ธมฺมํ เจตยตีติ เจโต.\csromanspacer{} สพฺพพฺยาปาทปริยุฏฺฐาเนหิ วิมุจฺจตีติ วิมุตฺติ.\csromanspacer{} เมตฺตา จ เจโต จ วิมุตฺติ จาติ เมตฺตาเจโตวิมุตฺติ.}

\prose{สพฺเพ สตฺตา อเวริโน โหนฺตุ, เขมิโน โหนฺตุ, สุขิโน โหนฺตูติ สทฺธาย อธิมุจฺจติ.\csromanspacer{} สทฺธินฺทฺริยปริภาวิตา โหติ เมตฺตาเจโตวิมุตฺติ.}

\prose{\csromanfolio{316}สพฺเพ สตฺตา อเวริโน โหนฺตุ, เขมิโน โหนฺตุ.\csromanspacer{} สุขิโน โหนฺตูติ วีริยํ ปคฺคณฺหาติ.\csromanspacer{} วีริยินฺทฺริยปริภาวิตา โหติ เมตฺตาเจโตวิมุตฺติ.}

\prose{สพฺเพ สตฺตา อเวริโน โหนฺตุ, เขมิโน โหนฺตุ, สุขิโน โหนฺตูติ สติํ อุปฏฺฐาเปติ.\csromanspacer{} สตินฺทฺริยปริภาวิตา โหติ เมตฺตาเจโตวิมุตฺติ.}

\prose{สพฺเพ สตฺตา อเวริโน โหนฺตุ, เขมิโน โหนฺตุ, สุขิโน โหนฺตูติ จิตฺตํ สมาทหติ.\csromanspacer{} สมาธินฺทฺริยปริภาวิตา โหติ เมตฺตาเจโตวิมุตฺติ.}

\prose{สพฺเพ สตฺตา อเวริโน โหนฺตุ, เขมิโน โหนฺตุ, สุขิโน โหนฺตูติ ปญฺญาย ปชานาติ.\csromanspacer{} ปญฺญินฺทฺริยปริภาวิตา โหติ เมตฺตาเจโตวิมุตฺติ.}

\prose{อิมานิ ปญฺจินฺทฺริยานิ เมตฺตาย เจโตวิมุตฺติยา อาเสวนา โหนฺติ.\csromanspacer{} อิเมหิ ปญฺจหิ อินฺทฺริเยหิ เมตฺตาเจโตวิมุตฺติ อาเสวียติ.\csromanspacer{} อิมานิ ปญฺจินฺทฺริยานิ เมตฺตาย เจโตวิมุตฺติยา ภาวนา โหนฺติ.\csromanspacer{} อิเมหิ ปญฺจหิ อินฺทฺริเยหิ เมตฺตาเจโตวิมุตฺติ ภาวียติ.\csromanspacer{} อิมานิ ปญฺจินฺทฺริยานิ เมตฺตาย เจโตวิมุตฺติยา พหุลีกตา โหนฺติ.\csromanspacer{} อิเมหิ ปญฺจหิ อินฺทฺริเยหิ เมตฺตาเจโตวิมุตฺติ พหุลีกรียติ.\csromanspacer{} อิมานิ ปญฺจินฺทฺริยานิ เมตฺตาย เจโตวิมุตฺติยา อลงฺการา โหนฺติ.\csromanspacer{} อิเมหิ ปญฺจหิ อินฺทฺริเยหิ เมตฺตาเจโตวิมุตฺติ สฺวาลงฺกตา โหติ.\csromanspacer{} อิมานิ ปญฺจินฺทฺริยานิ เมตฺตาย เจโตวิมุตฺติยา ปริกฺขารา โหนฺติ.\csromanspacer{} อิเมหิ ปญฺจหิ อินฺทฺริเยหิ เมตฺตาเจโตวิมุตฺติ สุปริกฺขตา โหติ.\csromanspacer{} อิมานิ ปญฺจินฺทฺริยานิ เมตฺตาย เจโตวิมุตฺติยา ปริวารา โหนฺติ.\csromanspacer{} อิเมหิ ปญฺจหิ อินฺทฺริเยหิ เมตฺตาเจโตวิมุตฺติ สุปริวุตา โหติ.\csromanspacer{} อิมานิ ปญฺจินฺทฺริยานิ เมตฺตาย เจโตวิมุตฺติยา อาเสวนา โหนฺติ, ภาวนา โหนฺติ, พหุลีกตา โหนฺติ, อลงฺการา โหนฺติ, ปริกฺขารา โหนฺติ, ปริวารา โหนฺติ, ปาริปูรี โหนฺติ, สหคตา โหนฺติ, สหชาตา โหนฺติ, สํสฏฺฐา โหนฺติ, สมฺปยุตฺตา โหนฺติ, ปกฺขนฺทนา โหนฺติ, ปสีทนา\footnote{สํสีทนา (ก, สี-ฏฺฐ)} โหนฺติ, สนฺติฏฺฐนา โหนฺติ, วิมุจฺจนา โหนฺติ, “เอตํ สนฺตนฺ”ติ ผสฺสนา โหนฺติ, ยานีกตา โหนฺติ, วตฺถุกตา โหนฺติ,}

\vspace{\tipitakaparskip}
\csromancenter{เมตฺตากถา อนุฏฺฐิตา โหนฺติ, ปริจิตา โหนฺติ, สุสมารทฺธา โหนฺติ, สุภาวิตา โหนฺติ, สฺวาธิฏฺฐิตา โหนฺติ, สุสมุคฺคตา โหนฺติ, สุวิมุตฺตา โหนฺติ นิพฺพตฺเตนฺติ โชเตนฺติ ปตาเปนฺติ.}
\csromanmatikaanchor{mtk.97}
\csromantocmarkref{subsection}{2. พลวาร}{mtk.97}
\bookmark[dest={mtk.97},level=2]{2. พลวาร}
\csromanheader{2}{2. พลวาร}
\proseitem{24}{\csromanfolio{317}สพฺเพ สตฺตา อเวริโน โหนฺตุ, เขมิโน โหนฺตุ, สุขิโน โหนฺตูติ อสฺสทฺธิเย น กมฺปติ.\csromanspacer{} สทฺธาพลปริภาวิตา โหติ เมตฺตาเจโตวิมุตฺติ.}

\prose{สพฺเพ สตฺตา อเวริโน โหนฺตุ, เขมิโน โหนฺตุ, สุขิโน โหนฺตูติ โกสชฺเชน กมฺปติ.\csromanspacer{} วีริยพลปริภาวิตา โหติ เมตฺตาเจโตวิมุตฺติ.}

\prose{สพฺเพ สตฺตา อเวริโน โหนฺตุ, เขมิโน โหนฺตุ, สุขิโน โหนฺตูติ ปมาเท น กมฺปติ.\csromanspacer{} สติพลปริภาวิตา โหติ เมตฺตาเจติวิมุตฺติ.}

\prose{สพฺเพ สตฺตา อเวริโน โหนฺตุ, เขมิโน โหนฺตุ, สุขิโน โหนฺตูติ อุทฺธจฺเจ น กมฺปติ.\csromanspacer{} สมาธิพลปริภาวิตา โหติ เมตฺตาเจติวิมุตฺติ.}

\prose{สพฺเพ สตฺตา อเวริโน โหนฺตุ, เขมิโน โหนฺตุ, สุขิโน โหนฺตูติ อวิชฺชาย น กมฺปติ.\csromanspacer{} ปญฺญาพลปริภาวิตา โหติ เมตฺตาเจโตวิมุตฺติ.}

\prose{อิมานิ ปญฺจ พลานิ เมตฺตาย เจโตวิมุตฺติยา อาเสวนา โหนฺติ.\csromanspacer{} อิเมหิ ปญฺจหิ พเลหิ เมตฺตาเจโตวิมุตฺติ อาเสวียติ.\csromanspacer{} อิมานิ ปญฺจ พลานิ เมตฺตาย เจโตวิมุตฺติยา ภาวนา โหนฺติ.\csromanspacer{} อิเมหิ ปญฺจหิ พเลหิ เมตฺตาเจโตวิมุตฺติ ภาวียติ.\csromanspacer{} อิมานิ ปญฺจ พลานิ เมตฺตาย เจโตวิมุตฺติยา พหุลีกตา โหนฺติ.\csromanspacer{} อิเมหิ ปญฺจหิ พเลหิ เมตฺตาเจโตวิมุตฺติ พหุลีกรียติ.\csromanspacer{} อิมานิ ปญฺจ พลานิ เมตฺตาย เจโตวิมุตฺติยา อลงฺการา โหนฺติ.\csromanspacer{} อิเมหิ ปญฺจหิ พเลหิ เมตฺตาเจโตวิมุตฺติ สฺวาลงฺกตา โหติ.\csromanspacer{} อิมานิ ปญฺจ พลานิ เมตฺตาย เจโตวิมุตฺติยา ปริกฺขารา โหนฺติ.\csromanspacer{} อิเมหิ ปญฺจหิ พเลหิ เมตฺตาเจโตวิมุตฺติ สุปริกฺขตา โหติ.\csromanspacer{} อิมานิ ปญฺจ พลานิ เมตฺตาย เจโตวิมุตฺติยา \csromanfolio{318}ปริวารา โหนฺติ.\csromanspacer{} อิเมหิ ปญฺจหิ พเลหิ เมตฺตาเจโตวิมุตฺติ สุปริวุตา โหติ.\csromanspacer{} อิมานิ ปญฺจ พลานิ เมตฺตาย เจโตวิมุตฺติยา อาเสวนา โหนฺติ, ภาวนา โหนฺติ, พหุลีกตา โหนฺติ, อลงฺการา โหนฺติ, ปริกฺขารา โหนฺติ, ปริวารา โหนฺติ, ปาริปูรี โหนฺติ, สหคตา โหนฺติ, สหชาตา โหนฺติ, สํสฏฺฐา โหนฺติ, สมฺปยุตฺตา โหนฺติ, ปกฺขนฺทนา โหนฺติ, ปสีทนา โหนฺติ, สนฺติฏฺฐนา โหนฺติ, วิมุจฺจนา โหนฺติ, “เอตํ สนฺตนฺ”ติ ผสฺสนา โหนฺติ, ยานีกตา โหนฺติ, วตฺถุกตา โหนฺติ, อนุฏฺฐิตา โหนฺติ, ปริจิตา โหนฺติ, สุสมารทฺธา โหนฺติ, สุภาวิตา โหนฺติ, สฺวาธิฏฺฐิตา โหนฺติ, สุสมุคฺคตา โหนฺติ, สุวิมุตฺตา โหนฺติ นิพฺพตฺเตนฺติ โชเตนฺติ ปตาเปนฺติ.}

\csromanmatikaanchor{mtk.98}
\csromantocmarkref{subsection}{3. โพชฺฌงฺควาร}{mtk.98}
\bookmark[dest={mtk.98},level=2]{3. โพชฺฌงฺควาร}
\csromanheader{2}{3. โพชฺฌงฺควาร}
\proseitem{25}{สพฺเพ สตฺตา อเวริโน โหนฺตุ, เขมิโน โหนฺตุ, สุขิโน โหนฺตูติ สติํ อุปฏฺฐาเปติ.\csromanspacer{} สติสมฺโพชฺฌงฺคปริภาวิตา โหติ เมตฺตาเจโตวิมุตฺติ.}

\prose{สพฺเพ สตฺตา อเวริโน โหนฺตุ, เขมิโน โหนฺตุ, สุขิโน โหนฺตูติ ปญฺญาย ปวิจินาติ.\csromanspacer{} ธมฺมวิจยสมฺโพชฺฌงฺคปริภาวิตา โหติ เมตฺตาเจโตวิมุตฺติ.}

\prose{สพฺเพ สตฺตา อเวริโน โหนฺตุ, เขมิโน โหนฺตุ, สุขิโน โหนฺตูติ วีริยํ ปคฺคณฺหาติ.\csromanspacer{} วีริยสมฺโพชฺฌงฺคปริภาวิตา โหติ เมตฺตาเจโตวิมุตฺติ.}

\prose{สพฺเพ สตฺตา อเวริโน โหนฺตุ, เขมิโน โหนฺตุ, สุขิโน โหนฺตูติ ปริฬาหํ ปฏิปฺปสฺสมฺเภติ.\csromanspacer{} ปีติสมฺโพชฺฌงฺคปริภาวิตา โหติ เมตฺตาเจโตวิมุตฺติ.}

\prose{สพฺเพ สตฺตา อเวริโน โหนฺตุ, เขมิโน โหนฺตุ, สุขิโน โหนฺตูติ ทุฏฺฐุลฺลํ ปฏิปฺปสฺสมฺเภติ.\csromanspacer{} ปสฺสทฺธิสมฺโพชฺฌงฺคปริภาวิตา โหติ เมตฺตาเจโตวิมุตฺติ.}

\prose{สพฺเพ สตฺตา อเวริโน โหนฺตุ, เขมิโน โหนฺตุ, สุขิโน โหนฺตูติ จิตฺตํ สมาทหติ.\csromanspacer{} สมาธิสมฺโพชฺฌงฺคปริภาวิตา โหติ เมตฺตาเจโตวิมุตฺติ.}

\csromanheader{2}{4. เมตฺตากถา}
\prose{\csromanfolio{319}สพฺเพ สตฺตา อเวริโน โหนฺตุ, เขมิโน โหนฺตุ, สุขิโน โหนฺตูติ ญาเณน กิเลเส ปฏิสงฺขาติ.\csromanspacer{} อุเปกฺขาสมฺโพชฺฌงฺคปริภาวิตา โหติ เมตฺตาเจโตวิมุตฺติ.}

\prose{อิเม สตฺต โพชฺฌงฺคา เมตฺตาย เจโตวิมุตฺติยา อาเสวนา โหนฺติ.\csromanspacer{} อิเมหิ สตฺตหิ โพชฺฌงฺเคหิ เมตฺตาเจโตวิมุตฺติ อาเสวียติ.\csromanspacer{} อิเม สตฺต โพชฺฌงฺคา เมตฺตาย เจโตวิมุตฺติยา ภาวนา โหนฺติ.\csromanspacer{} อิเมหิ สตฺตหิ โพชฺฌงฺเคหิ เมตฺตาเจโตวิมุตฺติ ภาวียติ.\csromanspacer{} อิเม สตฺต โพชฺฌงฺคา เมตฺตาย เจโตวิมุตฺติยา พหุลีกตา โหนฺติ.\csromanspacer{} อิเมหิ สตฺตหิ โพชฺฌงฺเคหิ เมตฺตาเจโตวิมุตฺติ พหุลีกรียติ.\csromanspacer{} อิเม สตฺต โพชฺฌงฺคา เมตฺตาย เจโตวิมุตฺติยา อลงฺการา โหนฺติ.\csromanspacer{} อิเมหิ สตฺตหิ โพชฺฌงฺเคหิ เมตฺตาเจโตวิมุตฺติ สฺวาลงฺกตา โหติ.\csromanspacer{} อิเม สตฺต โพชฺฌงฺคา เมตฺตาย เจโตวิมุตฺติยา ปริกฺขารา โหนฺติ.\csromanspacer{} อิเมหิ สตฺตหิ โพชฺฌงฺเคหิ เมตฺตาเจโตวิมุตฺติ สุปริกฺขตา โหติ.\csromanspacer{} อิเม สตฺต โพชฺฌงฺคา เมตฺตาย เจโตวิมุตฺติยา ปริวารา โหนฺติ.\csromanspacer{} อิเมหิ สตฺตหิ โพชฺฌงฺเคหิ เมตฺตาเจโตวิมุตฺติ สุปริวุตา โหติ.\csromanspacer{} อิเม สตฺต โพชฺฌงฺคา เมตฺตาย เจโตวิมุตฺติยา อาเสวนา โหนฺติ.\csromanspacer{} ภาวนา โหนฺติ, พหุลีกตา โหนฺติ, อลงฺการา โหนฺติ, ปริกฺขารา โหนฺติ, ปริวารา โหนฺติ, ปาริปูรี โหนฺตี, สหคตา โหนฺติ, สหชาตา โหนฺติ, สํสฏฺฐา โหนฺติ, สมฺปยุตฺตา โหนฺติ, ปกฺขนฺทนา โหนฺติ, ปสีทนา โหนฺติ, สนฺติฏฺฐนา โหนฺติ, วิมุจฺจนา โหนฺติ, “เอตํ สนฺตนฺ”ติ ผสฺสนา โหนฺติ, ยานีกตา โหนฺติ, วตฺถุกตา โหนฺติ, อนุฏฺฐิตา โหนฺติ, ปริจิตา โหนฺติ, สุสมารทฺธา โหนฺติ, สุภาวิตา โหนฺติ, สฺวาธิฏฺฐิตา โหนฺติ, สุสมุคฺคตา โหนฺติ, สุวิมุตฺตา โหนฺติ นิพฺพตฺเตนฺติ โชเตนฺติ ปตาเปนฺติ.}

\csromanmatikaanchor{mtk.99}
\csromantocmarkref{subsection}{4. มคฺคงฺควาร}{mtk.99}
\bookmark[dest={mtk.99},level=2]{4. มคฺคงฺควาร}
\csromanheader{2}{4. มคฺคงฺควาร}
\proseitem{26}{สพฺเพ สตฺตา อเวริโน โหนฺตุ, เขมิโน โหนฺตุ, สุขิโน โหนฺตูติ สมฺมา ปสฺสติ.\csromanspacer{} สมฺมาทิฏฺฐิปริภาวิตา โหติ เมตฺตาเจโตวิมุตฺติ.}

\prose{สพฺเพ สตฺตา อเวริโน โหนฺตุ, เขมิโน โหนฺตุ, สุขิโน โหนฺตูติ สมฺมา อภินิโรเปติ.\csromanspacer{} สมฺมาสงฺกปฺปปริภาวิตา โหติ เมตฺตาเจโตวิมุตฺติ.}

\prose{\csromanfolio{320}สพฺเพ สตฺตา อเวริโน โหนฺตุ, เขมิโน โหนฺตุ, สุขิโน โหนฺตูติ สมฺมา ปริคฺคณฺหาติ.\csromanspacer{} สมฺมาวาจาปริภาวิตา โหติ เมตฺตาเจโตวิมุตฺติ.}

\prose{สพฺเพ สตฺตา อเวริโน โหนฺตุ, เขมิโน โหนฺตุ, สุขิโน โหนฺตูติ สมฺมา สมุฏฺฐาเปติ.\csromanspacer{} สมฺมากมฺมนฺตปริภาวิตา โหติ เมตฺตาเจโตวิมุตฺติ.}

\prose{สพฺเพ สตฺตา อเวริโน โหนฺตุ, เขมิโน โหนฺตุ, สุขิโน โหนฺตูติ สมฺมา โวทาเปติ.\csromanspacer{} สมฺมา-อาชีวปริภาวิตา โหติ เมตฺตาเจโตวิมุตฺติ.}

\prose{สพฺเพ สตฺตา อเวริโน โหนฺตุ, เขมิโน โหนฺตุ, สุขิโน โหนฺตูติ สมฺมา ปคฺคณฺหาติ.\csromanspacer{} สมฺมาวายามปริภาวิตา โหติ เมตฺตาเจโตวิมุตฺติ.}

\prose{สพฺเพ สตฺตา อเวริโน โหนฺตุ, เขมิโน โหนฺตุ, สุขิโน โหนฺตูติ สมฺมา อุปฏฺฐาเปติ.\csromanspacer{} สมฺมาสติปริภาวิตา โหติ เมตฺตาเจโตวิมุตฺติ.}

\prose{สพฺเพ สตฺตา อเวริโน โหนฺตุ, เขมิโน โหนฺตุ, สุขิโน โหนฺตูติ สมฺมา สมาทหติ.\csromanspacer{} สมฺมาสมาธิปริภาวิตา โหติ เมตฺตาเจโตวิมุตฺติ.}

\prose{อิเม อฏฺฐ มคฺคงฺคา เมตฺตาย เจโตวิมุตฺติยา อาเสวนา โหนฺติ.\csromanspacer{} อิเมหิ อฏฺฐหิ มคฺคงฺเคหิ เมตฺตาเจโตวิมุตฺติ อาเสวียติ.\csromanspacer{} อิเม อฏฺฐ มคฺคงฺคา เมตฺตาย เจโตวิมุตฺติยา ภาวนา โหนฺติ.\csromanspacer{} อิเมหิ อฏฺฐหิ มคฺคงฺเคหิ เมตฺตาเจโตวิมุตฺติ ภาวียติ.\csromanspacer{} อิเม อฏฺฐ มคฺคงฺคา เมตฺตาย เจโตวิมุตฺติยา พหุลีกตา โหนฺติ.\csromanspacer{} อิเมหิ อฏฺฐหิ มคฺคงฺเคหิ เมตฺตาเจโตวิมุตฺติ พหุลีกรียติ.\csromanspacer{} อิเม อฏฺฐ มคฺคงฺคา เมตฺตาย เจโตวิมุตฺติยา อลงฺการา โหนฺติ.\csromanspacer{} อิเมหิ อฏฺฐหิ มคฺคงฺเคหิ เมตฺตาเจโตวิมุตฺติ สฺวาลงฺกตา โหนฺติ.\csromanspacer{} อิเม อฏฺฐ มคฺคงฺคา เมตฺตาย เจโตวิมุตฺติยา ปริกฺขารา โหนฺติ.\csromanspacer{} อิเมหิ อฏฺฐหิ มคฺคงฺเคหิ เมตฺตาเจโตวิมุตฺติ สุปริกฺขตา โหติ.\csromanspacer{} อิเม อฏฺฐ มคฺคงฺคา เมตฺตาย เจโตวิมุตฺติยา ปริวารา โหนฺติ.\csromanspacer{} อิเมหิ อฏฺฐหิ มคฺคงฺเคหิ เมตฺตาเจโตวิมุตฺติ สุปริวุตา โหติ.\csromanspacer{} อิเม อฏฺฐ มคฺคงฺคา เมตฺตาย เจโตวิมุตฺติยา}

\vspace{\tipitakaparskip}
\csromancenter{เมตฺตากถา อาเสวนา โหนฺติ, ภาวนา โหนฺติ, พหุลีกตา โหนฺติ, อลงฺการา โหนฺติ, ปริกฺขารา โหนฺติ, ปริวารา โหนฺติ, ปาริปูรี โหนฺติ, สหคตา โหนฺติ, สหชาตา โหนฺติ, สํสฏฺฐา โหนฺติ, สมฺปยุตฺตา โหนฺติ, ปกฺขนฺทนา โหนฺติ, ปสีทนา โหนฺติ, สนฺติฏฺฐนา โหนฺติ, วิมุจฺจนา โหนฺติ, “เอตํ สนฺตนฺ”ติ ผสฺสนา โหนฺติ, ยานีกตา โหนฺติ, วตฺถุกตา โหนฺติ, อนุฏฺฐิตา โหนฺติ, ปริจิตา โหนฺติ, สุสมารทฺธา โหนฺติ, สุภาวิตา โหนฺติ, สฺวาธิฏฺฐิตา โหนฺติ, สุสมุคฺคตา โหนฺติ, สุวิมุตฺตา โหนฺติ นิพฺพตฺเตนฺติ โชเตนฺติ ปตาเปนฺติ.}
\proseitem{27}{\csromanfolio{321}สพฺเพสํ ปาณานํ ฯเปฯ.\csromanspacer{} สพฺเพสํ ภูตานํ.\csromanspacer{} สพฺเพสํ ปุคฺคลานํ.\csromanspacer{} สพฺเพสํ อตฺตภาวปริยาปนฺนานํ.\csromanspacer{} สพฺพาสํ อิตฺถีนํ.\csromanspacer{} สพฺเพสํ ปุริสานํ.\csromanspacer{} สพฺเพสํ อริยานํ.\csromanspacer{} สพฺเพสํ อนริยานํ.\csromanspacer{} สพฺเพสํ เทวานํ.\csromanspacer{} สพฺเพสํ มนุสฺสานํ.\csromanspacer{} สพฺเพสํ วินิปาติกานํ ปีฬนํ วชฺเชตฺวา อปีฬนาย อุปฆาตํ วชฺเชตฺวา อนุปฆาเตน, สนฺตาปํ วชฺเชตฺวา อสนฺตาเปน, ปริยาทานํ วชฺเชตฺวา อปริยาทาเนน, วิเหสํ วชฺเชตฺวา อวิเหสาย สพฺเพ วินิปาติกา อเวริโน โหนฺตุ, มา เวริโน.\csromanspacer{} สุขิโน โหนฺตุ, มา ทุกฺขิโน.\csromanspacer{} สุขิตตฺตา โหนฺตุ, มา ทุกฺขิตตฺตาติ อิเมหิ อฏฺฐหากาเรหิ สพฺเพ วินิปาติเก เมตฺตายตีติ เมตฺตา, ตํ ธมฺมํ เจตยตีติ เจโต, สพฺพพฺยาปาทปริยุฏฺฐาเนหิ วิมุจฺจตีติ วิมุตฺติ, เมตฺตา จ เจโต จ วิมุตฺติ จาติ เมตฺตาเจโตวิมุตฺติ.}

\prose{สพฺเพ วินิปาติกา อเวริโน โหนฺตุ, เขมิโน โหนฺตุ, สุขิโน โหนฺตูติ สทฺธาย อธิมุจฺจติ.\csromanspacer{} สทฺธินฺทฺริยปริภาวิตา โหติ เมตฺตาเจโตวิมุตฺติ ฯเปฯ นิพฺพตฺเตนฺติ โชเตนฺติ ปตาเปนฺติ.}

\prose{สพฺเพสํ ปุรตฺถิมาย ทิสาย สตฺตานํ ฯเปฯ.\csromanspacer{} สพฺเพสํ ปจฺฉิมาย ทิสาย สตฺตานํ.\csromanspacer{} สพฺเพสํ อุตฺตราย ทิสาย สตฺตานํ.\csromanspacer{} สพฺเพสํ ทกฺขิณาย ทิสาย สตฺตานํ.\csromanspacer{} สพฺเพสํ ปุรตฺถิมาย อนุทิสาย สตฺตานํ.\csromanspacer{} สพฺเพสํ ปจฺฉิมาย อนุทิสาย สตฺตานํ.\csromanspacer{} สพฺเพสํ อุตฺตราย อนุทิสาย สตฺตานํ.\csromanspacer{} สพฺเพสํ ทกฺขิณาย อนุทิสาย สตฺตานํ.\csromanspacer{} สพฺเพสํ เหฏฺฐิมาย ทิสาย สตฺตานํ สพฺเพสํ อุปริมาย ทิสาย สตฺตานํ ปีฬนํ วชฺเชตฺวา อปีฬนาย, อุปฆาตํ วชฺเชตฺวา อนุปฆาเตน, สนฺตาปํ วชฺเชตฺวา อสนฺตาเปน, ปริยาทานํ วชฺเชตฺวา อปริยาทาเนน, วิเหสํ วชฺเชตฺวา อวิเหสาย สพฺเพ อุปริมาย ทิสาย สตฺตา อเวริโน โหนฺตุ, มา เวริโน.\csromanspacer{} สุขิโน \csromanfolio{322}โหนฺตุ, มา ทุกฺขิโน.\csromanspacer{} สุขิตตฺตา โหนฺตุ, มา ทุกฺขิตตฺตาติ อิเมหิ อฏฺฐหากาเรหิ สพฺเพ อุปริมาย ทิสาย สตฺเต เมตฺตายตีติ เมตฺตา, ตํ ธมฺมํ เจตยตีติ เจโต, สพฺพพฺยาปาทปริยุฏฺฐาเนหิ วิมุจฺจตีติ วิมุตฺติ, เมตฺตา จ เจโต จ วิมุตฺติ จาติ เมตฺตาเจโตวิมุตฺติ.}

\prose{สพฺเพ อุปริมาย ทิสาย สตฺตา อเวริโน โหนฺตุ, เขมิโน โหนฺตุ, สุขิโน โหนฺตูติ สทฺธาย อธิมุจฺจติ.\csromanspacer{} สทฺธินฺทฺริยปริภาวิตา โหติ, เมตฺตาเจโตวิมุตฺติ ฯเปฯ นิพฺพตฺเตนฺติ โชเตนฺติ ปตาเปนฺติ.}

\prose{สพฺเพสํ ปุรตฺถิมาย ทิสาย ปาณานํ.\csromanspacer{} ภูตานํ.\csromanspacer{} ปุคฺคลานํ.\csromanspacer{} อตฺตภาวปริยาปนฺนานํ.\csromanspacer{} สพฺพาสํ อิตฺถีนํ.\csromanspacer{} สพฺเพสํ ปุริสานํ.\csromanspacer{} สพฺเพสํ อริยานํ.\csromanspacer{} สพฺเพสํ อนริยานํ.\csromanspacer{} สพฺเพสํ เทวานํ.\csromanspacer{} สพฺเพสํ มนุสฺสานํ.\csromanspacer{} สพฺเพสํ วินิปาติกานํ.\csromanspacer{} สพฺเพสํ ปจฺฉิมาย ทิสาย วินิปาติกานํ.\csromanspacer{} สพฺเพสํ อุตฺตราย ทิสาย วินิปาติกานํ.\csromanspacer{} สพฺเพสํ ทกฺขิณาย ทิสาย วินิปาติกานํ.\csromanspacer{} สพฺเพสํ ปุรตฺถิมาย อนุทิสาย วินิปาติกานํ.\csromanspacer{} สพฺเพสํ ปจฺฉิมาย อนุทิสาย วินิปาติกานํ.\csromanspacer{} สพฺเพสํ อุตฺตราย อนุทิสาย วินิปาติกานํ.\csromanspacer{} สพฺเพสํ ทกฺขิณาย อนุทิสาย วินิปาติกานํ.\csromanspacer{} สพฺเพสํ เหฏฺฐิมาย ทิสาย วินิปาติกานํ.\csromanspacer{} สพฺเพสํ อุปริมาย ทิสาย วินิปาติกานํ ปีฬนํ วชฺเชตฺวา อปีฬนาย, อุปฆาตํ วชฺเชตฺวา อนุปฆาเตน, สนฺตาปํ วชฺเชตฺวา อสนฺตาเปน, ปริยาทานํ วชฺเชตฺวา อปริยาทาเนน, วิเหสํ วชฺเชตฺวา อวิเหสาย สพฺเพ อุปริมาย ทิสาย วินิปาติกา อเวริโน โหนฺตุ, มา เวริโน.\csromanspacer{} สุขิโน โหนฺตุ, มา ทุกฺขิโน.\csromanspacer{} สุขิตตฺตา โหนฺตุ, มา ทุกฺขิตตฺตาติ อิเมหิ อฏฺฐหากาเรหิ สพฺเพ อุปริมาย ทิสาย วินิปาติเก เมตฺตายตีติ เมตฺตา, ตํ ธมฺมํ เจตยตีติ เจโต, สพฺพพฺยาปาทปริยุฏฺฐาเนหิ วิมุจฺจตีติ วิมุตฺติ, เมตฺตา จ เจโต จ วิมุตฺติ จาติ เมตฺตาเจโตวิมุตฺติ.}

\prose{สพฺเพ อุปริมาย ทิสาย วินิปาติกา อเวริโน โหนฺตุ, เขมิโน โหนฺตุ, สุขิโน โหนฺตูติ สทฺธาย อธิมุจฺจติ.\csromanspacer{} สทฺธินฺทฺริยปริภาวิตา โหติ เมตฺตาเจโตวิมุตฺติ.}

\prose{สพฺเพ อุปริมาย ทิสาย วินิปาติกา อเวริโน โหนฺตุ, เขมิโน โหนฺตุ, สุขิโน โหนฺตูติ วีริยํ ปคฺคณฺหาติ.\csromanspacer{} วีริยินฺทฺริยปริภาวิตา โหติ เมตฺตาเจโตวิมุตฺติ.}

\csromanheader{2}{4. เมตฺตากถา}
\prose{\csromanfolio{323}สติํ อุปฏฺฐาเปติ, สตินฺทฺริยปริภาวิตา โหติ เมตฺตาเจโตวิมุตฺติ.\csromanspacer{} จิตฺตํ สมาทหติ, สมาธินฺทฺริยปริภาวิตา โหติ เมตฺตาเจโตวิมุตฺติ.\csromanspacer{} ปญฺญาย ปชานาติ, ปญฺญินฺทฺริยปริภาวิตา โหติ เมตฺตาเจโตวิมุตฺติ.}

\prose{อิมานิ ปญฺจินฺทฺริยานิ เมตฺตาย เจโตวิมุตฺติยา อาเสวนา โหนฺติ.\csromanspacer{} อิเมหิ ปญฺจหิ อินฺทฺริเยหิ เมตฺตาเจโตวิมุตฺติ อาเสวียติ ฯเปฯ นิพฺพตฺเตนฺติ โชเตนฺติ ปตาเปนฺติ.}

\prose{สพฺเพ อุปริมาย ทิสาย วินิปาติกา อเวริโน โหนฺตุ, เขมิโน โหนฺตุ, สุขิโน โหนฺตูติ อสฺสทฺธิเย น กมฺปติ, สทฺธาพลปริภาวิตา โหติ เมตฺตาเจโตวิมุตฺติ.\csromanspacer{} โกสชฺเช น กมฺปติ, วีริยพลปริภาวิตา โหติ เมตฺตาเจโตวิมุตฺติ.\csromanspacer{} ปมาเท น กมฺปติ, สติพลปริภาวิตา โหติ เมตฺตาเจโตวิมุตฺติ.\csromanspacer{} อุทฺธจฺเจ น กมฺปติ, สมาธิพลปริภาวิตา โหติ เมตฺตาเจโตวิมุตฺติ.\csromanspacer{} อวิชฺชาย น กมฺปติ, ปญฺญาพลปริภาวิตา โหติ เมตฺตาเจโตวิมุตฺติ.}

\prose{อิมานิ ปญฺจ พลานิ เมตฺตาย เจโตวิมุตฺติยา อาเสวนา โหนฺติ.\csromanspacer{} อิเมหิ ปญฺจหิ พเลหิ เมตฺตาเจโตวิมุตฺติ อาเสวียติ ฯเปฯ นิพฺพตฺเตนฺติ โชเตนฺติ ปตาเปนฺติ.}

\prose{สพฺเพ อุปริมาย ทิสาย วินิปาติกา อเวริโน โหนฺตุ, เขมิโน โหนฺตุ, สุขิโน โหนฺตูติ สติํ อุปฏฺฐาเปติ, สติสมฺโพชฺฌงฺคปริภาวิตา โหติ เมตฺตาเจโตวิมุตฺติ.}

\prose{ปญฺญาย ปวิจินาติ, ธมฺมวิจยสมฺโพชฺฌงฺคปริภาวิตา โหติ เมตฺตาเจโตวิมุตฺติ.\csromanspacer{} วีริยํ ปคฺคณฺหาติ, วีริยสมฺโพชฺฌงฺคปริภาวิตา โหติ เมตฺตาเจโตวิมุตฺติ.\csromanspacer{} ปริฬาหํ ปฏิปฺปสฺสมฺเภติ, ปีติสมฺโพชฺฌงฺคปริภาวิตา โหติ เมตฺตาเจโตวิมุตฺติ.\csromanspacer{} ทุฏฺฐุลฺลํ ปฏิปฺปสฺสมฺเภติ, ปสฺสทฺธิสมฺโพชฺฌงฺคปริภาวิตา โหติ เมตฺตาเจโตวิมุตฺติ.\csromanspacer{} จิตฺตํ สมาทหติ, สมาธิสมฺโพชฺฌงฺคปริภาวิตา โหติ เมตฺตาเจโตวิมุตฺติ.\csromanspacer{} ญาเณน กิเลเส ปฏิสงฺขาติ, อุเปกฺขาสมฺโพชฺฌงฺคปริภาวิตา โหติ เมตฺตาเจโตวิมุตฺติ.}

\prose{อิเม สตฺต โพชฺฌงฺคา เมตฺตาย เจโตวิมุตฺติยา อาเสวนา โหนฺติ.\csromanspacer{} อิเมหิ สตฺตหิ โพชฺฌงฺเคหิ เมตฺตาเจโตวิมุตฺติ อาเสวียติ ฯเปฯ นิพฺพตฺเตนฺติ โชเตนฺติ ปตาเปนฺติ.}

\csromanmatikaanchor{mtk.100}
\csromantocmarkref{section}{5. วิราคกถา}{mtk.100}
\bookmark[dest={mtk.100},level=1]{5. วิราคกถา}
\prose{\csromanfolio{324}สพฺเพ อุปริมาย ทิสาย วินิปาติกา อเวริโน โหนฺตุ, เขมิโน โหนฺตุ, สุขิโน โหนฺตูติ สมฺมา ปสฺสติ, สมฺมาทิฏฺฐิปริภาวิตา โหติ เมตฺตาเจโตวิมุตฺติ.\csromanspacer{} สมฺมา อภินิโรเปติ, สมฺมาสงฺกปฺปปริภาวิตา โหติ เมตฺตาเจโตวิมุตฺติ.\csromanspacer{} สมฺมา ปริคฺคณฺหาติ, สมฺมาวาจาปริภาวิตา โหติ เมตฺตาเจโตวิมุตฺติ.\csromanspacer{} สมฺมา สมุฏฺฐาเปติ, สมฺมากมฺมนฺตปริภาวิตา โหติ เมตฺตาเจโตวิมุตฺติ.\csromanspacer{} สมฺมา โวทาเปติ, สมฺมา-อาชีวปริภาวิตา โหติ เมตฺตาเจโตวิมุตฺติ.\csromanspacer{} สมฺมา ปคฺคณฺหาติ, สมฺมาวายามปริภาวิตา โหติ เมตฺตาเจโตวิมุตฺติ.\csromanspacer{} สมฺมา อุปฏฺฐาติ, สมฺมาสติปริภาวิตา โหติ เมตฺตาเจโตวิมุตฺติ.\csromanspacer{} สมฺมา สมาทหติ, สมฺมาสมาธิปริภาวิตา โหติ เมตฺตาเจโตวิมุตฺติ.}

\prose{อิเม อฏฺฐ มคฺคงฺคา เมตฺตาย เจโตวิมุตฺติยา อาเสวนา โหนฺติ.\csromanspacer{} อิเมหิ อฏฺฐหิ มคฺคงฺเคหิ เมตฺตาเจโตวิมุตฺติ อาเสวียติ ฯเปฯ.\csromanspacer{} อิเม อฏฺฐ มคฺคงฺคา เมตฺตาย เจโตวิมุตฺติยา ปริภาวิตา โหนฺติ.\csromanspacer{} อิเมหิ อฏฺฐหิ มคฺคงฺเคหิ เมตฺตาเจโตวิมุตฺติ สุปริวุตา โหติ.\csromanspacer{} อิเม อฏฺฐ มคฺคงฺคา เมตฺตาย เจโตวิมุตฺติยา อาเสวนา โหนฺติ, ภาวนา โหนฺติ, พหุลีกตา โหนฺติ, อลงฺการา โหนฺติ, ปริกฺขารา โหนฺติ, ปริวารา โหนฺติ, ปาริปูรี โหนฺติ, สหคตา โหนฺติ, สหชาตา โหนฺติ, สํสฏฺฐา โหนฺติ, สมฺปยุตฺตา โหนฺติ, ปกฺขนฺทนา โหนฺติ, ปสีทนา โหนฺติ, สนฺติฏฺฐนา โหนฺติ, วิมุจฺจนา โหนฺติ, “เอตํ สนฺตนฺ”ติ ผสฺสนา โหนฺติ, ยานีกตา โหนฺติ, วตฺถุกตา โหนฺติ, อนุฏฺฐิตา โหนฺติ, ปริจิตา โหนฺติ, สุสมารทฺธา โหนฺติ, สุภาวิตา โหนฺติ, สฺวาธิฏฺฐิตา โหนฺติ, สุสมุคฺคตา โหนฺติ, สุวิมุตฺตา โหนฺติ นิพฺพตฺเตนฺติ โชเตนฺติ ปตาเปนฺตีติ.}

\nitthitamruled{เมตฺตากถา นิฏฺฐิตา.}
\csromanheader{2}{5. วิราคตถา}
\proseitem{28}{วิราโค มคฺโค, วิมุตฺติ ผลํ.\csromanspacer{} กถํ \textbf{วิราโค มคฺโค?} โสตาปตฺติมคฺคกฺขเณ ทสฺสนฏฺเฐน สมฺมาทิฏฺฐิ มิจฺฉาทิฏฺฐิยา วิรชฺชติ, ตทนุวตฺตกกิเลเสหิ จ ขนฺเธหิ จ วิรชฺชติ.\csromanspacer{} พหิทฺธา จ สพฺพนิมิตฺเตหิ วิรชฺชติ.}

\vspace{\tipitakaparskip}
\csromancenter{วิราคกถา \textbf{วิราโค} วิราคารมฺมโณ วิราคโคจโร วิราเค สมุทาคโต\footnote{สมุปาคโต (สฺยา)} วิราเค ฐิโต วิราเค ปติฏฺฐิโต.}
\prose{\csromanfolio{325}\textbf{วิราโค}ติ เทฺว วิราคา นิพฺพานญฺจ วิราโค, เย จ นิพฺพานารมฺมณตาชาตา ธมฺมา สพฺเพ วิราคา โหนฺตีติ วิราคา.\csromanspacer{} สหชาตานิ สตฺตงฺคานิ วิราคํ คจฺฉนฺตีติ วิราโค มคฺโค.\csromanspacer{} เอเตน มคฺเคน พุทฺธา จ สาวกา จ อคตํ ทิสํ นิพฺพานํ คจฺฉนฺตีติ อฏฺฐงฺคิโก มคฺโค, ยาวตา ปุถุสมณพฺราหฺมณานํ ปรปฺปวาทานํ มคฺคา, อยเมว อริโย อฏฺฐงฺคิโก มคฺโค อคฺโค จ เสฏฺโฐ จ ปาโมกฺโข\footnote{วิโมกฺโข (สฺยา), โมกฺโข (ฏฺฐ)} จ อุตฺตโม จ ปวโร จาติ มคฺคานํ อฏฺฐงฺคิโก เสฏฺโฐ.}

\prose{อภินิโรปนฏฺเฐน สมฺมาสงฺกปฺโป มิจฺฉาสงฺกปฺปา วิรชฺชติ ปริคฺคหฏฺเฐน สมฺมาวาจา มิจฺฉาวาจาย วิรชฺชติ.\csromanspacer{} สมุฏฺฐานฏฺเฐน สมฺมากมฺมนฺโต มิจฺฉากมฺมนฺตา วิรชฺชติ.\csromanspacer{} โวทานฏฺเฐน สมฺมา-อาชีโว มิจฺฉา-อาชีวา วิรชฺชติ.\csromanspacer{} ปคฺคหฏฺเฐน สมฺมาวายาโม มิจฺฉาวายามา วิรชฺชติ.\csromanspacer{} อุปฏฺฐานฏฺเฐน สมฺมาสติ มิจฺฉาสติยา วิรชฺชติ.\csromanspacer{} อวิกฺเขปฏฺเฐน สมฺมาสมาธิ มิจฺฉาสมาธิโต วิรชฺชติ.\csromanspacer{} ตทนุวตฺตกกิเลเสหิ จ ขนฺเธหิ จ วิรชฺชติ.\csromanspacer{} พหิทฺธา จ สพฺพนิมิตฺเตหิ วิรชฺชติ, วิราโค วิราคารมฺมโณ วิราคโคจโร วิราเค สมุทาคโต วิราเค ฐิโต วิราเค ปติฏฺฐิโต.}

\prose{\textbf{วิราโค}ติ เทฺว วิราคา นิพฺพานญฺจ วิราโค, เย จ นิพฺพานารมฺมณตาชาตา ธมฺมา สพฺเพ วิราคา โหนฺตีติ วิราคา, สหชาตานิ สตฺตงฺคานิ วิราคํ คจฺฉนฺตีติ วิราโค มคฺโค.\csromanspacer{} เอเตน มคฺเคน พุทฺธา จ สาวกา จ อคตํ ทิสํ นิพฺพานํ คจฺฉนฺตีติ อฏฺฐงฺคิโก มคฺโค, ยาวตา ปุถุสมณพฺราหฺมณานํ ปรปฺปวาทานํ มคฺคา, อยเมว อริโย อฏฺฐงฺคิโก มคฺโค อคฺโค จ เสฏฺโฐ จ ปาโมกฺโข จ อุตฺตโม จ ปวโร จาติ มคฺคานํ อฏฺฐงฺคิโก เสฏฺโฐ.}

\prose{สกทาคามิมคฺคกฺขเณ ทสฺสนฏฺเฐน สมฺมาทิฏฺฐิ ฯเปฯ อวิกฺเขปฏฺเฐน สมฺมาสมาธิ โอฬาริกา กามราคสญฺโญชนา ปฏิฆสญฺโญชนา โอฬาริกา กามราคานุสยา ปฏิฆานุสยา วิรชฺชติ, ตทนุวตฺตกกิเลเสหิ จ \csromanfolio{326}ขนฺเธหิ จ วิรชฺชติ, พหิทฺธา จ สพฺพนิมิตฺเตหิ วิรชฺชติ, วิราโค วิราคารมฺมโณ วิราคโคจโร วิราเค สมุทาคโต วิราเค ฐิโต วิราเค ปติฏฺฐิโต.}

\prose{\textbf{วิราโค}ติ เทฺว วิราคา นิพฺพานญฺจ วิราโค, เย จ นิพฺพานารมฺมณตาชาตา ธมฺมา สพฺเพ วิราคา โหนฺตีติ วิราคา, สหชาตานิ สตฺตงฺคานิ วิราคํ คจฺฉนฺตีติ วิราโค มคฺโค, เอเตน มคฺเคน พุทฺธา จ สาวกา จ อคตํ ทิสํ นิพฺพานํ คจฺฉนฺตีติ อฏฺฐงฺคิโก มคฺโค, ยาวตา ปุถุสมณพฺราหฺมณานํ ปรปฺปวาทานํ มคฺคา.\csromanspacer{} อยเมว อริโย อฏฺฐงฺคิโก มคฺโค อคฺโค จ เสฏฺโฐ จ ปาโมกฺโข จ อุตฺตโม จ ปวโร จาติ มคฺคานํ อฏฺฐงฺคิโก เสฏฺโฐ.}

\prose{อนาคามิมคฺคกฺขเณ ทสฺสนฏฺเฐน สมฺมาทิฏฺฐิ ฯเปฯ อวิกฺเขปฏฺเฐน สมฺมาสมาธิ อณุสหคตา กามราคสญฺโญชนา ปฏิฆสญฺโญชนา อณุสหคตา กามราคานุสยา ปฏิฆานุสยา วิรชฺชติ, ตทนุวตฺตกกิเลเสหิ จ ขนฺเธหิ จ วิรชฺชติ, พหิทฺธา จ สพฺพนิมิตฺเตหิ วิรชฺชติ, วิราโค วิราคารมฺมโณ ฯเปฯ มคฺคานํ อฏฺฐงฺคิโก เสฏฺโฐ.}

\prose{อรหตฺตมคฺคกฺขเณ ทสฺสนฏฺเฐน สมฺมาทิฏฺฐิ ฯเปฯ อวิกฺเขปฏฺเฐน สมฺมาสมาธิ รูปราคา อรูปราคา มานา อุทฺธจฺจา อวิชฺชาย มานานุสยา ภวราคานุสยา อวิชฺชานุสยา วิรชฺชติ, ตทนุวตฺตกกิเลเสหิ จ ขนฺเธหิ จ วิรชฺชติ, พหิทฺธา จ สพฺพนิมิตฺเตหิ วิรชฺชติ, วิราโค วิราคารมฺมโณ วิราคโคจโร วิราเค สมุทาคโต วิราเค ฐิโต วิราเค ปติฏฺฐิโต.}

\prose{\textbf{วิราโค}ติ เทฺว วิราคา นิพฺพานญฺจ วิราโค, เย จ นิพฺพานารมฺมณตาชาตา ธมฺมา สพฺเพ วิราคา โหนฺตีติ วิราคา, สหชาตานิ สตฺตงฺคานิ วิราคํ คจฺฉนฺตีติ วิราโค มคฺโค, เอเตน มคฺเคน พุทฺธา จ สาวกา จ อคตํ ทิสํ นิพฺพานํ คจฺฉนฺตีติ อฏฺฐงฺคิโก มคฺโค, ยาวตา ปุถุสมณพฺราหฺมณานํ ปรปฺปวาทานํ มคฺคา, อยเมว อริโย อฏฺฐงฺคิโก มคฺโค อคฺโค จ เสฏฺโฐ จ ปาโมกฺโข จ อุตฺตโม จ ปวโร จาติ มคฺคานํ อฏฺฐงฺคิโก เสฏฺโฐ.}

\prose{ทสฺสนวิราโค สมฺมาทิฏฺฐิ, อภินิโรปนวิราโค สมฺมาสงฺกปฺโป, ปริคฺคหวิราโค สมฺมาวาจา, สมุฏฺฐานวิราโค สมฺมากมฺมนฺโต, โวทานวิราโค}

\vspace{\tipitakaparskip}
\csromancenter{วิราคกถา สมฺมา-อาชีโว, ปคฺคหวิราโค สมฺมาวายาโม, อุปฏฺฐานวิราโค สมฺมาสติ, อวิกฺเขปวิราโค สมฺมาสมาธิ.\csromanspacer{}\csromanspacer{}อุปฏฺฐานวิราโค สติสมฺโพชฺฌงฺโค, ปวิจยวิราโค ธมฺมวิจยสมฺโพชฺฌงฺโค, ปคฺคหวิราโค วีริยสมฺโพชฺฌงฺโค, ผรณวิราโค ปีติสมฺโพชฺฌงฺโค, อุปสมวิราโค ปสฺสทฺธิสมฺโพชฺฌงฺโค, อวิกฺเขปวิราโค สมาธิสโมชฺฌงฺโค, ปฏิสงฺขานวิราโค อุเปกฺขาสมฺโพชฺฌงฺโค.\csromanspacer{}\csromanspacer{}อสฺสทฺธิเย อกมฺปิยวิราโค สทฺธาพลํ, โกสชฺเช อกมฺปิยวิราโค วีริยพลํ, ปมาเท อกมฺปิยวิราโค สติพลํ, อุทฺธจฺเจ อกมฺปิยวิราโค สมาธิพลํ, อวิชฺชาย อกมฺปิยวิราโค ปญฺญาพลํ.\csromanspacer{}\csromanspacer{}อธิโมกฺขวิราโค สทฺธินฺทฺริยํ, ปคฺคหวิราโค วีริยินฺทฺริยํ, อุปฏฺฐานวิราโค สตินฺทฺริยํ, อวิกฺเขปวิราโค สมาธินฺทฺริยํ, ทสฺสนวิราโค ปญฺญินฺทฺริยํ.\csromanspacer{}\csromanspacer{}อาธิปเตยฺยฏฺเฐน อินฺทฺริยานิ วิราโค, อกมฺปิยฏฺเฐน พลํ วิราโค, นิยฺยานฏฺเฐน โพชฺฌงฺคา วิราโค, เหตุฏฺเฐน มคฺโค วิราโค, อุปฏฺฐานฏฺเฐน สติปฏฺฐานา วิราโค, ปทหนฏฺเฐน สมฺมปฺปธานา วิราโค, อิชฺฌนฏฺเฐน อิทฺธิปฺปาทา วิราโค, ตถฏฺเฐน สจฺจา วิราโค, อวิกฺเขปฏฺเฐน สมโถ วิราโค, อนุปสฺสนฏฺเฐน วิปสฺสนา วิราโค, เอกรสฏฺเฐน สมถวิปสฺสนา วิราโค, อนติวตฺตนฏฺเฐน ยุคนทฺธํ วิราโค, สํวรฏฺเฐน สีลวิสุทฺธิ วิราโค, อวิกฺเขปฏฺเฐน จิตฺตวิสุทฺธิ วิราโค, ทสฺสนฏฺเฐน ทิฏฺฐิวิสุทฺธิ วิราโค, วิมุตฺตฏฺเฐน วิโมกฺโข วิราโค, ปฏิเวธฏฺเฐน วิชฺชา วิราโค, ปริจฺจาคฏฺเฐน วิมุตฺติ วิราโค, สมุจฺเฉทฏฺเฐน ขเย ญาณํ วิราโค, ฉนฺโท มุลฏฺเฐน วิราโค, มนสิกาโร สมุฏฺฐานฏฺเฐน วิราโค, ผสฺโส สโมธานฏฺเฐน วิราโค, เวทนา สโมสรณฏฺเฐน วิราโค, สมาธิ ปมุขฏฺเฐน วิราโค, สติ อาธิปเตยฺยฏฺเฐน วิราโค, ปญฺญา ตทุตฺตรฏฺเฐน วิราโค, วิมุตฺติ สารฏฺเฐน วิราโค, อมโตคธํ นิพฺพานํ ปริโยสานฏฺเฐน มคฺโค.}
\prose{\csromanfolio{327}ทสฺสนมคฺโค สมฺมาทิฏฺฐิ, อภินิโรปนมคฺโค สมฺมาสงฺกปฺโป ฯเปฯ อมโตคธํ นิพฺพานํ ปริโยสานฏฺเฐน มคฺโค, เอวํ วิราโค มคฺโค.}

\proseitem{29}{กถํ \textbf{วิมุตฺติ ผลํ?} โสตาปตฺติผลกฺขเณ ทสฺสนฏฺเฐน สมฺมาทิฏฺฐิ มิจฺฉาทิฏฺฐิยา วิมุตฺตา โหติ, ตทนุวตฺตกกิเลเสหิ จ ขนฺเธหิ จ วิมุตฺตา โหติ, พหิทฺธา จ สพฺพนิมิตฺเตหิ วิมุตฺตา โหติ, วิมุตฺติ วิมุตฺตารมฺมณา วิมุตฺติโคจรา วิมุตฺติยา สมุทาคตา วิมุตฺติยา ฐิตา วิมุตฺติยา ปติฏฺฐิตา.\csromanspacer{} \textbf{วิมุตฺตี}ติ เทฺว วิมุตฺติโย นิพฺพานญฺจ วิมุตฺติ, เย จ นิพฺพานารมฺมณตาชาตา ธมฺมา สพฺเพ จ วิมุตฺตา โหนฺตีติ วิมุตฺติ ผลํ.}

\prose{\csromanfolio{328}อภินิโรปนฏฺเฐน สมฺมาสงฺกปฺโป มิจฺฉาสงฺกปฺปา วิมุตฺโต โหติ, ตทนุวตฺตกกิเลเสหิ จ ขนฺเธหิ จ วิมุตฺโต โหติ, พหิทฺธา จ สพฺพนิมิตฺเตหิ วิมุตฺโต โหติ, วิมุตฺติ วิมุตฺตารมฺมณา วิมุตฺติโคจรา วิมุตฺติยา สมุทาคตา วิมุตฺติยา ฐิตา วิมุตฺติยา ปติฏฺฐิตา.\csromanspacer{} \textbf{วิมุตฺตี}ติ เทฺว วิมุตฺติโย นิพฺพานญฺจ วิมุตฺติ, เย จ นิพฺพานารมฺมณตาชาตา ธมฺมา สพฺเพ วิมุตฺตา โหนฺตีติ วิมุตฺติ ผลํ.}

\prose{ปริคฺคหฏฺเฐน สมฺมาวาจา มิจฺฉาวาจาย วิมุตฺตา โหติ, สมุฏฺฐานฏฺเฐน สมฺมากมฺมนฺโต มิจฺฉากมฺมนฺตา วิมุตฺโต โหติ.\csromanspacer{} โวทานฏฺเฐน สมฺมา-อาชีโว มิจฺฉา-อาชีวา วิมุตฺโต โหติ, ปคฺคหฏฺเฐน สมฺมาวายาโม มิจฺฉาวายามา วิมุตฺโต โหติ, อุปฏฺฐานฏฺเฐน สมฺมาสติ มิจฺฉาสติยา วิมุตฺตา โหติ, อวิกฺเขปฏฺเฐน สมฺมาสมาธิ มิจฺฉาสมาธิโต วิมุตฺโต โหติ, ตทนุวตฺตกกิเลเสหิ จ ขนฺเธหิ จ วิมุตฺโต โหติ, พหิทฺธา จ สพฺพนิมิตฺเตหิ วิมุตฺโต โหติ, วิมุตฺติ วิมุตฺตารมฺมณา วิมุตฺติโคจรา วิมุตฺติยา สมุทาคตา วิมุตฺติยา ฐิตา วิมุตฺติยา ปติฏฺฐิตา.\csromanspacer{} \textbf{วิมุตฺตี}ติ เทฺว วิมุตฺติโย, นิพฺพานญฺจ วิมุตฺติ, เย จ นิพฺพานารมฺมณตาชาตา ธมฺมา สพฺเพ วิมุตฺตา โหนฺตีติ วิมุตฺติ ผลํ.}

\prose{สกทาคามิผลกฺขเณ ทสฺสนฏฺเฐน สมฺมาทิฏฺฐิ ฯเปฯ อวิกฺเขปฏฺเฐน สมฺมาสมาธิ โอฬาริกา กามราคสญฺโญชนา ปฏิฆสญฺโญชนา โอฬาริกา กามราคานุสยา ปฏิฆานุสยา วิมุตฺโต โหติ, ตทนุวตฺตกกิเลเสหิ จ ขนฺเธหิ จ วิมุตฺโต โหติ, พหิทฺธา จ สพฺพนิมิตฺเตหิ วิมุตฺโต โหติ, วิมุตฺติ วิมุตฺตารมฺมณา วิมุตฺติโคจรา วิมุตฺติยา สมุทาคตา วิมุตฺติยา ฐิตา วิมุตฺติยา ปติฏฺฐิตา.\csromanspacer{} \textbf{วิมุตฺตี}ติ เทฺว วิมุตฺติโย, นิพฺพานญฺจ วิมุตฺติ, เย จ นิพฺพานารมฺมณตาชาตา ธมฺมา สพฺเพ วิมุตฺตา โหนฺตีติ วิมุตฺติ ผลํ.}

\prose{อนาคามิผลกฺเขเณ ทสฺสนฏฺเฐน สมฺมาทิฏฺฐิ ฯเปฯ อวิกฺเขปฏฺเฐน สมฺมาสมาธิ อณุสหคตา กามราคสญฺโญชนา ปฏิฆสญฺโญชนา อณุสหคตา กามราคานุสยา ปฏิฆานุสยา วิมุตฺโต โหติ, ตทนุวตฺตกกิเลเสหิ จ ขนฺเธหิ จ วิมุตฺโต โหติ, พหิทฺธา จ สพฺพนิมิตฺเตหิ วิมุตฺโต โหติ, วิมุตฺติ วิมุตฺตารมฺมณา วิมุตฺติโคจรา วิมุตฺติยา สมุทาคตา วิมุตฺติยา ฐิตา วิมุตฺติยา ปติฏฺฐิตา ฯเปฯ.}

\csromanmatikaanchor{mtk.101}
\csromantocmarknopage{section}{6. ปฏิสมฺภิทากถา}{mtk.101}
\bookmark[dest={mtk.101},level=1]{6. ปฏิสมฺภิทากถา}
\csromanheader{1}{6. ปฏิสมฺภิทากถา}
\prose{\csromanfolio{329}อรหตฺตผลกฺขเณ ทสฺสนฏฺเฐน สมฺมาทิฏฺฐิ ฯเปฯ อวิกฺเขปฏฺเฐน สมฺมาสมาธิ รูปราคา อรูปราคา มานา อุทฺธจฺจา อวิชฺชาย มานานุสยา ภาวราคานุสยา อวิชฺชานุสยา วิมุตฺโต โหติ.\csromanspacer{} ตทนุวตฺตกกิเลเสหิ จ ขนฺเธหิ จ วิมุตฺโต โหติ, พหิทฺธา จ สพฺพนิมิตฺเตหิ วิมุตฺโต โหติ, วิมุตฺติ วิมุตฺตารมฺมณา วิมุตฺติโคจรา วิมุตฺติยา สมุทาคตา วิมุตฺติยา ฐิตา วิมุตฺติยา ปติฏฺฐิตา.\csromanspacer{} \textbf{วิมุตฺตี}ติ เทฺว วิมุตฺติโย, นิพฺพานญฺจ วิมุตฺติ, เย จ นิพฺพานารมฺมณตาชาตา ธมฺมา สพฺเพ วิมุตฺตา โหนฺตีติ วิมุตฺติ ผลํ.}

\prose{ทสฺสนวิมุตฺติ สมฺมาทิฏฺฐิ ฯเปฯ อวิกฺเขปวิมุตฺติ สมฺมาสมาธิ.\csromanspacer{}\csromanspacer{}อุปฏฺฐานวิมุตฺติ สติสมฺโพชฺฌงฺโค.\csromanspacer{} ปฏิสงฺขานวิมุตฺติ อุเปกฺขาสมฺโพชฺฌงฺโค.\csromanspacer{}\csromanspacer{}อสฺสทฺธิเย อกมฺปิยวิมุตฺติ สทฺธาพลํ ฯเปฯ อวิชฺชาย อกมฺปิยวิมุตฺติ ปญฺญาพลํ.\csromanspacer{}\csromanspacer{}อธิโมกฺขวิมุตฺติ สทฺธินฺทฺริยํ ฯเปฯ ทสฺสนวิมุตฺติ ปญฺญินฺทฺริยํ.}

\prose{อาธิปเตยฺยฏฺเฐน อินฺทฺริยา วิมุตฺติ.\csromanspacer{} อกมฺปิยฏฺเฐน พลา วิมุตฺติ, นิยฺยานฏฺเฐน โพชฺฌงฺคา วิมุตฺติ, เหตุฏฺเฐน มคฺโค วิมุตฺติ, อุปฏฺฐานฏฺเฐน สติปฏฺฐานา วิมุตฺติ, ปทหนฏฺเฐน สมฺมปฺปธานา วิมุตฺติ, อิชฺฌนฏฺเฐน อิทฺธิปาทา วิมุตฺติ, ตถฏฺเฐน สจฺจา วิมุตฺติ, อวิกฺเขปฏฺเฐน สมโถ วิมุตฺติ, อนุปสฺสนฏฺเฐน วิปสฺสนา วิมุตฺติ, เอกรสฏฺเฐน สมถวิปสฺสนา วิมุตฺติ, อนติวตฺตนฏฺเฐน ยุคนทฺธํ วิมุตฺติ, สํวรฏฺเฐน สีลวิสุทฺธิ วิมุตฺติ, อวิกฺเขปฏฺเฐน จิตฺตวิสุทฺธิ วิมุตฺติ, ทสฺสนฏฺเฐน ทิฏฺฐิวิสุทฺธิ วิมุตฺติ, วิมุตฺตฏฺเฐน วิโมกฺโข วิมุตฺติ, ปฏิเวธฏฺเฐน วิชฺชา วิมุตฺติ, ปริจฺจาคฏฺเฐน วิมุตฺติ วิมุตฺติ, ปฏิปฺปสฺสทฺธิยฏฺเฐน อนุปฺปาเท ญาณํ วิมุตฺติ, ฉนฺโท มูลฏฺเฐน วิมุตฺติ, มนสิกาโร สมุฏฺฐานฏฺเฐน วิมุตฺติ, ผสฺโส สโมธานฏฺเฐน วิมุตฺติ, เวทนา สโมสรณฏฺเฐน วิมุตฺติ, สมาธิ ปมุขฏฺเฐน วิมุตฺติ, สติ อาธิปเตยฺยฏฺเฐน วิมุตฺติ, ปญฺญา ตทุตฺตรฏฺเฐน วิมุตฺติ, วิมุตฺติ สารฏฺเฐน วิมุตฺติ, อมโตคธํ นิพฺพานํ ปริโยสานฏฺเฐน วิมุตฺติ, เอวํ วิมุตฺติ ผลํ.\csromanspacer{} เอวํ วิราโค มคฺโค วิมุตฺติ ผลนฺติ.}

\nitthitamruled{วิราคกถา นิฏฺฐิตา.}
\csromanheader{2}{6. ปฏิสมฺภิทากถา}
\csromanmatikaanchor{mtk.102}
\csromantocmarkref{subsection}{1. ธมฺมจกฺกปวตฺตนวาร}{mtk.102}
\bookmark[dest={mtk.102},level=2]{1. ธมฺมจกฺกปวตฺตนวาร}
\csromanheader{2}{1. ธมฺมจกฺกปวตฺตนวาร}
\proseitem{30}{เอวํ เม สุตํ\csromandash{}เอกํ สมยํ ภควา พาราณสิยํ วิหรติ อิสิปตเน มิคทาเย.\csromanspacer{} ตตฺร โข ภควา ปญฺจวคฺคิเย ภิกฺขู อามนฺเตสิ}

\prose{\csromanfolio{330}เทฺวเม ภิกฺขเว อนฺตา ปพฺพชิเตน น เสวิตพฺพา.\csromanspacer{} กตเม เทฺว? โย จายํ กาเมสุ กามสุขลฺลิกานุโยโค หีโน คมฺโม โปถุชฺชนิโก อนริโย อนตฺถสํหิโต, โย จายํ อตฺตกิลมถานุโยโค ทุกฺโข อนริโย อนตฺถสํหิโต.\csromanspacer{} เอเต โข\footnote{เอเต เต (สฺยา, ก, สี-ฏฺฐ) สํ 3. 369; วิ 3. 14 ปิฏฺเฐสุ ปสฺสิตพฺพา.} ภิกฺขเว อุโภ อนฺเต อนุปคมฺม มชฺฌิมา ปฏิปทา ตถาคเตน อภิสมฺพุทฺธา จกฺขุกรณี ญาณกรณี อุปสมาย อภิญฺญาย สมฺโพธาย นิพฺพานาย สํวตฺตติ.}

\prose{กตมา จ สา ภิกฺขเว มชฺฌิมา ปฏิปทา ตถาคเตน อภิสมฺพุทฺธา จกฺขุกรณี ญาณกรณี อุปสมาย อภิญฺญาย สมฺโพธาย นิพฺพานาย สํวตฺตติ? อยเมว อริโย อฏฺฐงฺคิโก มคฺโค.\csromanspacer{} เสยฺยถิทํ, สมฺมาทิฏฺฐิ ฯเปฯ สมฺมาสมาธิ.\csromanspacer{} อยํ โข สา ภิกฺขเว มชฺฌิมา ปฏิปทา ตถาคเตน อภิสมฺพุทฺธา จกฺขุกรณี ญาณกรณี อุปสมาย อภิญฺญาย สมฺโพธาย นิพฺพานาย สํวตฺตติ.}

\prose{อิทํ โข ปน ภิกฺขเว ทุกฺขํ อริยสจฺจํ\csromandash{}ชาติปิ ทุกฺขา, ชราปิ ทุกฺขา, พฺยาธิปิ ทุกฺโข, มรณมฺปิ ทุกฺขํ, อปฺปิเยหิ สมฺปโยโค ทุกฺโข, ปิเยหิ วิปฺปโยโค ทุกฺโข, ยมฺปิจฺฉํ น ลภติ ตมฺปิ ทุกฺขํ, สํขิตฺเตน ปญฺจุปาทานกฺขนฺธา ทุกฺขา.\csromanspacer{} อิทํ โข ปน ภิกฺขเว ทุกฺขสมุทยํ อริยสจฺจํ\csromandash{}ยายํ ตณฺหา โปโนภวิกา นนฺทิราคสหคตา ตตฺร ตตฺราภินนฺทินี.\csromanspacer{} เสยฺยถิทํ, กามตณฺหา ภวตณฺหา วิภวตณฺหา.\csromanspacer{} อิทํ โข ปน ภิกฺขเว ทุกฺขนิโรธํ อริยสจฺจํ\csromandash{}โย ตสฺสาเยว ตณฺหาย อเสสวิราคนิโรโธ จาโค ปฏินิสฺสคฺโค มุตฺติ อนาลโย.\csromanspacer{} อิทํ โข ปน ภิกฺขเว ทุกฺขนิโรธคามินี ปฏิปทา อริยสจฺจํ, อยเมว อริโย อฏฺฐงฺคิโก มคฺโค.\csromanspacer{} เสยฺยถิทํ, สมฺมาทิฏฺฐิ ฯเปฯ สมฺมาสมาธิ.}

\prose{“อิทํ ทุกฺขํ อริยสจฺจนฺ”ติ เม ภิกฺขเว ปุพฺเพ อนนุสฺสุเตสุ ธมฺเมสุ จกฺขุํ อุทปาทิ, ญาณํ อุทปาทิ, ปญฺญา อุทปาทิ, วิชฺชา อุทปาทิ, อาโลโก อุทปาทิ.\csromanspacer{}\csromanspacer{}“ตํ โข ปนิทํ ทุกฺขํ อริยสจฺจํ ปริญฺเญยฺยนฺ”ติ เม ภิกฺขเว ฯเปฯ “ปริญฺญาตนฺ”ติ เม ภิกฺขเว ปุพฺเพ อนนุสฺสุเตสุ ธมฺเมสุ จกฺขุํ อุทปาทิ, ญาณํ อุทปาทิ, ปญฺญา อุทปาทิ, วิชฺชา อุทปาทิ, อาโลโก อุทปาทิ.}

\csromanheader{2}{6. ปฏิสมฺภิทากถา}
\prose{\csromanfolio{331}“อิทํ ทุกฺขสมุทยํ อริยสจฺจนฺ”ติ เม ภิกฺขเว ปุพฺเพ อนนุสฺสุเตสุ ธมฺเมสุ จกฺขุํ อุทปาทิ, ญาณํ อุทปาทิ, ปญฺญา อุทปาทิ, วิชฺชา อุทปาทิ, อาโลโก อุทปาทิ.\csromanspacer{}\csromanspacer{}“ตํ โข ปนิทํ ทุกฺขสมุทยํ อริยสจฺจํ ปหาตพฺพนฺ”ติ เม ภิกฺขเว ฯเปฯ “ปหีนนฺ”ติ เม ภิกฺขเว ปุพฺเพ อนนุสฺสุเตสุ ธมฺเมสุ จกฺขุํ อุทปาทิ, ญาณํ อุทปาทิ, ปญฺญา อุทปาทิ, วิชฺชา อุทปาทิ, อาโลโก อุทปาทิ.}

\prose{“อิทํ ทุกฺขนิโรธํ อริยสจฺจนฺ”ติ เม ภิกฺขเว ปุพฺเพ อนนุสฺสุเตสุ ธมฺเมสุ จกฺขุํ อุทปาทิ, ญาณํ อุทปาทิ, ปญฺญา อุทปาทิ, วิชฺชา อุทปาทิ, อาโลโก อุทปาทิ.\csromanspacer{}\csromanspacer{}“ตํ โข ปนิทํ ทุกฺขนิโรธํ อริยสจฺจํ สจฺฉิกาตพฺพนฺ”ติ เม ภิกฺขเว ฯเปฯ “สจฺฉิกตนฺ”ติ เม ภิกฺขเว ปุพฺเพ อนนุสฺสุเตสุ ธมฺเมสุ จกฺขุํ อุทปาทิ, ญาณํ อุทปาทิ, ปญฺญา อุทปาทิ, วิชฺชา อุทปาทิ, อาโลโก อุทปาทิ.}

\prose{“อิทํ ทุกฺขนิโรธคามินี ปฏิปทา อริยสจฺจนฺ”ติ เม ภิกฺขเว ปุพฺเพ อนนุสฺสุเตสุ ธมฺเมสุ จกฺขุํ อุทปาทิ, ญาณํ อุทปาทิ, ปญฺญา อุทปาทิ, วิชฺชา อุทปาทิ, อาโลโก อุทปาทิ.\csromanspacer{}\csromanspacer{}“ตํ โข ปนิทํ ทุกฺขนิโรธคามินี ปฏิปทา อริยสจฺจํ ภาเวตพฺพนฺ”ติ เม ภิกฺขเว ฯเปฯ “ภาวิตนฺ”ติ เม ภิกฺขเว ปุพฺเพ อนนุสฺสุเตสุ ธมฺเมสุ จกฺขุํ อุทปาทิ, ญาณํ อุทปาทิ, ปญฺญา อุทปาทิ, วิชฺชา อุทปาทิ, อาโลโก อุทปาทิ.}

\prose{ยาว กีวญฺจ เม ภิกฺขเว อิเมสุ จตูสุ อริยสจฺเจสุ เอวํ ติปริวฏฺฏํ ทฺวาทสาการํ ยถาภูตํ ญาณทสฺสนํ น สุวิสุทฺธํ อโหสิ, เนว ตาวาหํ ภิกฺขเว สเทวเก โลเก สมารเก สพฺรหฺมเก สสฺสมณพฺราหฺมณิยา ปชาย สเทวมนุสฺสาย “อนุตฺตรํ สมฺมาสมฺโพธิํ อภิสมฺพุทฺโธ”ติ ปจฺจญฺญาสิํ.\csromanspacer{} ยโต จ โข เม ภิกฺขเว อิเมสุ จตูสุ อริยสจฺเจสุ เอวํ ติปริวฏฺฏํ ทฺวาทสาการํ ยถาภูตํ ญาณทสฺสนํ สุวิสุทฺธํ อโหสิ, อถาหํ ภิกฺขเว สเทวเก โลเก สมารเก สพฺรหฺมเก สสฺสมณพฺรหฺมณิยา ปชาย สเทวมนุสฺสาย “อนุตฺตรํ สมฺมาสมฺโพธํ อภิสมฺพุทฺโธ”ติ ปจฺจญฺญาสิํ.\csromanspacer{} ญาณญฺจ ปน เม ทสฺสนํ อุทปาทิ.\csromanspacer{} อกุปฺปา เม วิมุตฺติ, อยมนฺติมา ชาติ, นตฺถิ ทานิ ปุนพฺภโวติ.}

\prose{อิทมโวจ ภควา, อตฺตมนา ปญฺจวคฺคิยา ภิกฺขู ภควโต ภาสิตํ อภินนฺทุนฺติ.}

\prose{\csromanfolio{332}อิมสฺมิํ จ ปน เวยฺยากรณสฺมิํ ภญฺญมาเน อายสฺมโต โกณฺฑญฺญสฺส วิรชํ วีตมลํ ธมฺมจกฺขุํ อุทปาทิ “ยํ กิญฺจิ สมุทยธมฺมํ, สพฺพํ ตํ นิโรธธมฺมนฺ”ติ.}

\prose{ปวตฺติเต จ ปน ภควตา ธมฺมจกฺเก ภุมฺมา\footnote{ภูมา (ก)} เทวา สทฺทมนุสฺสาเวสุํ “เอตํ ภควตา พาราณสิยํ อิสิปตเน มิคทาเย อนุตฺตรํ ธมฺมจกฺกํ ปวตฺติตํ อปฺปฏิวตฺติยํ สมเณน วา พฺราหฺมเณน วา เทเวน วา มาเรน วา พฺรหฺมุนา วา เกนจิ วา โลกสฺมินฺ”ติ.\csromanspacer{} ภุมฺมานํ เทวานํ สทฺทํ สุตฺวา จาตุมหาราชิกา\footnote{จาตุมฺมหาราชิกา (สฺยา)} เทวา สทฺทมนุสฺสาเวสํ.\csromanspacer{} จาตุมหาราชิกานํ เทวานํ สทฺทํ สุตฺวา ตาวติํสา เทวา ฯเปฯ.\csromanspacer{} ยามา เทวา ฯเปฯ.\csromanspacer{} ตุสิตา เทวา ฯเปฯ.\csromanspacer{} นิมฺมานรตี เทวา ฯเปฯ, ปรนิมฺมิตวสวตฺตี เทวา ฯเปฯ.\csromanspacer{} พฺรหฺมกายิกา เทวา สทฺทมนุสฺสาเวสุํ “เอตํ ภควตา พาราณสิยํ อิสิปตเน มิคทาเย อนุตฺตรํ ธมฺมจกฺกํ ปวตฺติกํ อปฺปฏิวตฺติยํ สมเณน วา พฺราหฺมเณน วา เทเวน วา มาเรน วา พฺรหฺมุนา วา เกนจิ วา โลกสฺมินฺ”ติ.}

\prose{อิติห เตน ขเณน เตน ลเยน เตน มุหุตฺเตน ยาว พฺรหฺมโลกา สทฺโท อพฺภุคฺคจฺฉิ.\csromanspacer{} อยญฺจ ทสสหสฺสี โลกธาตุ สํกมฺปิ สมฺปกมฺปิ สมฺปเวธิ, อปฺปมาโณ จ อุฬาโร โอภาโส โลเก ปาตุรโหสิ อติกฺกมฺม\footnote{อติกฺกมฺเมว (สฺยา, ก) สํ 3. 371 ปิฏฺเฐ ปสฺสิตพฺพา.} เทวานํ เทวานุภาวนฺติ.}

\prose{อถ โข ภควา อิมํ อุทานํ อุทาเนสิ “อญฺญาสิ วต โภ โกณฺฑญฺโญ, อญฺญาสิ วต โภ โกณฺฑญฺโญ”ติ.\csromanspacer{} อิติ หิทํ อายสฺมโต โกณฺฑญฺญสฺส “อญฺญาสิโกณฺฑญฺโญ”เตฺวว\footnote{อญฺญาโกณฺฑญฺโญเตฺวว (สฺยา)} นามํ อโหสิ.}

\prose{(ก) “อิทํ ทุกฺขํ อริยสจฺจนฺ”ติ ปุพฺเพ อนนุสฺสุเตสุ ธมฺเมสุ จกฺขุํ อุทปาทิ, ญาณํ อุทปาทิ, ปญฺญา อุทปาทิ, วิชฺชา อุทปาทิ, อาโลโก อุทปาทิ.}

\prose{\textbf{จกฺขุํ อุทปาที}ติ เกนฏฺเฐน, \textbf{ญาณํ อุทปาที}ติ เกนฏฺเฐน, \textbf{ปญฺญา อุทปาที}ติ เกนฏฺเฐน, \textbf{วิชฺชา อุทปาที}ติ เกนฏฺเฐน, \textbf{อาโลโก อุทปาที}ติ เกนฏฺเฐน? \textbf{จกฺขุํ อุทปาที}ติ ทสฺสนฏฺเฐน, \textbf{ญาณํ อุทปาที}ติ ญาตฏฺเฐน, \textbf{ปญฺญา อุทปาที}ติ ปชานนฏฺเฐน, \textbf{วิชฺชา อุทปาที}ติ ปฏิเวธฏฺเฐน, \textbf{อาโลโก อุทปาที}ติ โอภาสฏฺเฐน.}

\csromanheader{2}{6. ปฏิสมฺภิทากถา}
\prose{\csromanfolio{333}จกฺขุํ ธมฺโม, ญาณํ ธมฺโม, ปญฺญา ธมฺโม, วิชฺชา ธมฺโม, อาโลโก ธมฺโม.\csromanspacer{} อิเม ปญฺจ ธมฺมา ธมฺมปฏิสมฺภิทาย อารมฺมณา เจว โหนฺติ โคจรา จ.\csromanspacer{} เย ตสฺสา อารมฺมณา, เต ตสฺสา โคจรา.\csromanspacer{} เย ตสฺสา โคจรา, เต ตสฺสา อารมฺมณา.\csromanspacer{} เตน วุจฺจติ ธมฺเมสุ ญาณํ ธมฺมปฏิสมฺภิทา.}

\prose{ทสฺสนฏฺโฐ อตฺโถ, ญาตฏฺโฐ อตฺโถ ปชานนฏฺโฐ อตฺโถ, ปฏิเวธฏฺโฐ อตฺโถ โอภาสฏฺโฐ อตฺโถ.\csromanspacer{} อิเม ปญฺจ อตฺถา อตฺถปฏิสมฺภิทาย อารมฺมณา เจว โหนฺติ โคจรา จ, เย ตสฺสา อารมฺมณา, เต ตสฺสา โคจรา.\csromanspacer{} เย ตสฺสา โคจรา, เต ตสฺสา อารมฺมณา.\csromanspacer{} เตน วุจฺจติ อตฺเถสุ ญาณํ อตฺถปฏิสมฺภิทา.}

\prose{ปญฺจ ธมฺเม สนฺทสฺเสตุํ พฺยญฺชนนิรุตฺตาภิลาปา, ปญฺจ อตฺเถ สนฺทสฺเสตุํ พฺยญฺชนนิรุตฺตาภิลาปา.\csromanspacer{} อิมา ทส นิรุตฺติโย นิรุตฺติปฏิสมฺภิทาย อารมฺมณา เจว โหนฺติ โคจรา จ, เย ตสฺสา อารมฺมณา, เต ตสฺสา โคจรา.\csromanspacer{} เย ตสฺสา โคจรา, เต ตสฺสา อารมฺมณา.\csromanspacer{} เตน วุจฺจติ นิรุตฺตีสุ ญาณํ นิรุตฺติปฏิสมฺภิทา.}

\prose{ปญฺจสุ ธมฺเมสุ ญาณานิ, ปญฺจสุ อตฺเถสุ ญาณานิ, ทสสุ นิรุตฺตีสุ ญาณานิ.\csromanspacer{} อิมานิ วีสติ ญาณานิ ปฏิภานปฏิสมฺภิทาย อารมฺมณา เจว โหนฺติ โคจรา จ, เย ตสฺสา อารมฺมณา, เต ตสฺสา โคจรา.\csromanspacer{} เย ตสฺสา โคจรา, เต ตสฺสา อารมฺมณา.\csromanspacer{} เตน วุจฺจติ ปฏิภาเนสุ ญาณํ ปฏิภานปฏิสมฺภิทา.}

\prose{“ตํ โข ปนิทํ ทุกฺขํ อริยสจฺจํ ปริญฺเญยฺยนฺ”ติ ฯเปฯ “ปริญฺญาตนฺ”ติ ปุพฺเพ อนนุสฺสุเตสุ ธมฺเมสุ จกฺขุํ อุทปาทิ, ญาณํ อุทปาทิ, ปญฺญา อุทปาทิ, วิชฺชา อุทปาทิ, อาโลโก อุทปาทิ.}

\prose{\textbf{จกฺขุํ อุทปาที}ติ เกนฏฺเฐน, \textbf{ญาณํ อุทปาที}ติ เกนฏฺเฐน, \textbf{ปญฺญา อุทปาที}ติ เกนฏฺเฐน, \textbf{วิชฺชา อุทปาที}ติ เกนฏฺเฐน, \textbf{อาโลโก อุทปาที}ติ เกนฏฺเฐน? \textbf{จกฺขุํ อุทปาที}ติ ทสฺสนฏฺเฐน, \textbf{ญาณํ อุทปาที}ติ ญาตฏฺเฐน, \textbf{ปญฺญา อุทปาที}ติ ปชานนฏฺเฐน, \textbf{วิชฺชา อุทปาที}ติ ปฏิเวธฏฺเฐน, \textbf{อาโลโก อุทปาที}ติ โอภาสฏฺเฐน.}

\prose{จกฺขุํ ธมฺโม, ญาณํ ธมฺโม, ปญฺญา ธมฺโม, วิชฺชา ธมฺโม, อาโลโก ธมฺโม.\csromanspacer{} อิเม ปญฺจ ธมฺมา ธมฺมปฏิสมฺภิทาย อารมฺมณา เจว โหนฺติ โคจรา}

\prose{\csromanfolio{334}จ, เย ตสฺสา อารมฺมณา, เต ตสฺสา โคจรา.\csromanspacer{} เย ตสฺสา โคจรา, เต ตสฺสา อารมฺมณา.\csromanspacer{} เตน วุจฺจติ ธมฺเมสุ ญาณํ ธมฺมปฏิสมฺภิทา.}

\prose{ทสฺสนฏฺโฐ อตฺโถ, ญาตฏฺโฐ อตฺโถ, ปชานนฏฺโฐ อตฺโถ, ปฏิเวธฏฺโฐ อตฺโถ, โอภาสฏฺโฐ อตฺโถ.\csromanspacer{} อิเม ปญฺจ อตฺถา อตฺถปฏิสมฺภิทาย อารมฺมณา เจว โหนฺติ โคจรา จ, เย ตสฺสา อารมฺมณา, เต ตสฺสา โคจรา.\csromanspacer{} เย ตสฺสา โคจรา, เต ตสฺสา อารมฺมณา.\csromanspacer{} เตน วุจฺจติ อตฺเถสุ ญาณํ อตฺถปฏิสมฺภิทา.}

\prose{ปญฺจ ธมฺเม สนฺทสฺเสตุํ พฺยญฺชนนิรุตฺตาภิลาปา, ปญฺจ อตฺเถ สนฺทสฺเสตุํ พฺยญฺชนนิรุตฺตาภิลาปา, อิมา ทส นิรุตฺติโย นิรุตฺติปฏิสมฺภิทาย อารมฺมณา เจว โหนฺติ โคจรา จ, เย ตสฺสา อารมฺมณา, เต ตสฺสา โคจรา.\csromanspacer{} เย ตสฺสา โคจรา, เต ตสฺสา อารมฺมณา.\csromanspacer{} เตน วุจฺจติ นิรุตฺตีสุ ญาณํ นิรุตฺติปฏิสมฺภิทา.}

\prose{ปญฺจสุ ธมฺเมสุ ญาณานิ, ปญฺจสุ อตฺเถสุ ญาณานิ, ทสสุ นิรุตฺตีสุ ญาณานิ.\csromanspacer{} อิมานิ วีสติ ญาณานิ ปฏิภานปฏิสมฺภิทาย อารมฺมณา เจว โหนฺติ โคจรา จ, เย ตสฺสา อารมฺมณา, เต ตสฺสา โคจรา.\csromanspacer{} เย ตสฺสา โคจรา, เต ตสฺสา อารมฺมณา.\csromanspacer{} เตน วุจฺจติ ปฏิภาเนสุ ญาณํ ปฏิภานปฏิสมฺภิทา.}

\prose{ทุกฺเข อริยสจฺเจ ปนฺนรส ธมฺมา, ปนฺนรส อตฺถา, ติํส นิรุตฺติโย, สฏฺฐิ ญาณานิ.}

\prose{(ข) “อิทํ ทุกฺขสมุทยํ อริยสจฺจนฺ”ติ ปุพฺเพ อนนุสฺสุเตสุ ธมฺเมสุ จกฺขุํ อุทปาทิ ฯเปฯ อาโลโก อุทปาทิ. “ตํ โข ปนิทํ ทุกฺขสมุทยํ อริยสจฺจํ ปหาตพฺพนฺ”ติ ฯเปฯ “ปหีนนฺ”ติ ปุพฺเพ อนนุสฺสุเตสุ ธมฺเมสุ จกฺขุํ อุทปาทิ ฯเปฯ อาโลโก อุทปาทิ ฯเปฯ.}

\prose{ทุกฺขสมุทเย อริยสจฺเจ ปนฺนรส ธมฺมา, ปนฺนรส อตฺถา, ติํส นิรุตฺติโย, สฏฺฐิ ญาณานิ.}

\prose{(ค) “อิทํ ทุกฺขนิโรธํ อริยสจฺจนฺ”ติ ปุพฺเพ อนนุสฺสุเตสุ ธมฺเมสุ จกฺขุํ อุทปาทิ ฯเปฯ อาโลโก อุทปาทิ. “ตํ โข ปนิทํ ทุกฺขนิโรธํ อริยสจฺจํ สจฺฉิกาตพฺพนฺ”ติ ฯเปฯ “สจฺฉิกตนฺ”ติ ปุพฺเพ อนนุสฺสุเตสุ ธมฺเมสุ จกฺขุํ อุทปาทิ ฯเปฯ อาโลโก อุทาปาทิ ฯเปฯ.}

\csromanheader{2}{6. ปฏิสมฺภิทากถา}
\prose{\csromanfolio{335}ทุกฺขนิโรเธ อริยสจฺเจ ปนฺนรส ธมฺมา, ปนฺนรส อตฺถา, ติํส นิรุตฺติโย, สฏฺฐิ ญาณานิ.}

\prose{(ฆ) “อิทํ ทุกฺขนิโรธคามินี ปฏิปทา อริยสจฺจนฺ”ติ ปุพฺเพ อนนุสฺสุเตสุ ธมฺเมสุ จกฺขุํ อุทปาทิ ฯเปฯ อาโลโก อุทปาทิ. “ตํ โข ปนิทํ ทุกฺขนิโรธคามินี ปฏิปทา อริยสจฺจํ ภาเวตพฺพนฺ”ติ ฯเปฯ “ภาวิตนฺ”ติ ปุพฺเพ อนนุสฺสุเตสุ ธมฺเมสุ จกฺขุํ อุทปาทิ ฯเปฯ อาโลโก อุทปาทิ ฯเปฯ.}

\prose{ทุกฺขนิโรธคามินิยา ปฏิปทาย อริยสจฺเจ ปนฺนรส ธมฺมา, ปนฺนรส อตฺถา, ติํส นิรุตฺติโย, สฏฺฐิ ญาณานิ.}

\prose{จตูสุ อริยสจฺเจสุ สฏฺฐิ ธมฺมา, สฏฺฐิ อตฺถา, วีสติสตนิรุตฺติโย, จตฺตาลีสญฺจ เทฺว จ ญาณสตานิ.}

\csromanmatikaanchor{mtk.103}
\csromantocmarkref{subsection}{2. สติปฏฺฐานวาร}{mtk.103}
\bookmark[dest={mtk.103},level=2]{2. สติปฏฺฐานวาร}
\csromanheader{2}{2. สติปฏฺฐานวาร}
\proseitem{31}{“อยํ กาเย กายานุปสฺสนา”ติ เม ภิกฺขเว ปุพฺเพ อนนุสฺสุเตสุ ธมฺเมสุ จกฺขุํ อุทปาทิ ฯเปฯ อาโลโก อุทปาทิ, “สา โข ปนายํ กาเย กายานุปสฺสนา ภาเวตพฺพา”ติ เม ภิกฺขเว ฯเปฯ “ภาวิตา”ติ เม ภิกฺขเว ปุพฺเพ อนนุสฺสุเตสุ ธมฺเมสุ จกฺขุํ อุทปาทิ ฯเปฯ อาโลโก อุทปาทิ.}

\prose{“อยํ เวทนาสุ ฯเปฯ.\csromanspacer{} อยํ จิตฺเต ฯเปฯ.\csromanspacer{} อยํ ธมฺเมสุ ธมฺมานุปสฺสนา”ติ เม ภิกฺขเว ปุพฺเพ อนนุสฺสุเตสุ ธมฺเมสุ จกฺขุํ อุทปาทิ ฯเปฯ อาโลโก อุทปาทิ, “สา โข ปนายํ ธมฺเมสุ ธมฺมานุปสฺสนา ภาเวตพฺพา”ติ ฯเปฯ “ภาวิตา”ติ เม ภิกฺขเว ปุพฺเพ อนนุสฺสุเตสุ ธมฺเมสุ จกฺขุํ อุทปาทิ ฯเปฯ อาโลโก อุทปาทิ.}

\prose{(ก) “อยํ กาเย กายานุปสฺสนา”ติ ปุพฺเพ อนนุสฺสุเตสุ ธมฺเมสุ จกฺขุํ อุทปาทิ ฯเปฯ อาโลโก อุทปาทิ ฯเปฯ “สา โข ปนายํ กาเย กายานุปสฺสนา ภาเวตพฺพา”ติ ฯเปฯ “ภาวิตา”ติ ปุพฺเพ อนนุสฺสุเตสุ ธมฺเมสุ จกฺขุํ อุทปาทิ, ญาณํ อุทปาทิ, ปญฺญา อุทปาทิ, วิชฺชา อุทปาทิ, อาโลโก อุทปาทิ.}

\prose{\textbf{จกฺขุํ อุทปาที}ติ เกนฏฺเฐน, \textbf{ญาณํ อุทปาที}ติ เกนฏฺเฐน, ปญฺญา อุทปาทีติ เกนฏฺเฐน, วิชฺชา อุทปาทีติ เกนฏฺเฐน, อาโลโก อุทปาทีติ เกนฏฺเฐน? \textbf{จกฺขุํ อุทปาที}ติ ทสฺสนฏฺฐน, \textbf{ญาณํ อุทปาที}ติ ญาตฏฺเฐน, \csromanfolio{336}\textbf{ปญฺญา อุทปาที}ติ ปชานนฏฺเฐน, \textbf{วิชฺชา อุทปาที}ติ ปฏิเวธฏฺเฐน, \textbf{อาโลโก อุทปาที}ติ โอภาสฏฺเฐน.}

\prose{จกฺขุํ ธมฺโม, ญาณํ ธมฺโม, ปญฺญา ธมฺโม, วิชฺชา ธมฺโม, อาโลโก ธมฺโม.\csromanspacer{} อิเม ปญฺจ ธมฺมา ธมฺมปฏิสมฺภิทาย อารมฺมณา เจว โหนฺติ โคจรา จ.\csromanspacer{} เย ตสฺสา อารมฺมณา, เต ตสฺสา โคจรา, เย ตสฺสา โคจรา, เต ตสฺสา อารมฺมณา, เตน วุจฺจติ ธมฺเมสุ ญาณํ ธมฺมปฏิสมฺภิทา.}

\prose{ทสฺสนฏฺโฐ อตฺโถ, ญาตฏฺโฐ อตฺโถ, ปชานนฏฺโฐ อตฺโถ, ปฏิเวธฏฺโฐ อตฺโถ, โอภาสฏฺโฐ อตฺโถ.\csromanspacer{} อิเม ปญฺจ อตฺถา อตฺถปฏิสมฺภิทาย อารมฺมณา เจว โหนฺติ โคจรา จ, เย ตสฺสา อารมฺมณา, เต ตสฺสา โคจรา, เย ตสฺสา โคจรา, เต ตสฺสา อารมฺมณา, เตน วุจฺจติ อตฺเถสุ ญาณํ อตฺถปฏิสมฺภิทา.}

\prose{ปญฺจ ธมฺเม สนฺทสฺเสตุํ พฺยญฺชนนิรุตฺตาภิลาปา, ปญฺจ อตฺเถ สนฺทสฺเสตุํ พฺยญฺชนนิรุตฺตาภิลาปา, อิมา ทส นิรุตฺติโย นิรุตฺติปฏิสมฺภิทาย อารมฺมณา เจว โหนฺติ โคจรา จ, เย ตสฺสา อารมฺมณา, เต ตสฺสา โคจรา.\csromanspacer{} เย ตสฺสา โคจรา, เต ตสฺสา อารมฺมณา.\csromanspacer{} เตน วุจฺจติ นิรุตฺตีสุ ญาณํ นิรุตฺติปฏิสมฺภิทา.}

\prose{ปญฺจสุ ธมฺเมสุ ญาณานิ, ปญฺจสุ อตฺเถสุ ญาณานิ, ทสสุ นิรุตฺตีสุ ญาณานิ, อิมานิ วีสติ ญาณานิ ปฏิภานปฏิสมฺภิทาย อารมฺมณา เจว โหนฺติ โคจรา จ, เย ตสฺสา อารมฺมณา, เต ตสฺสา โคจรา.\csromanspacer{} เย ตสฺสา โคจรา, เต ตสฺสา อารมฺมณา.\csromanspacer{} เตน วุจฺจติ ปฏิภาเนสุ ญาณํ ปฏิภานปฏิสมฺภิทา.}

\prose{กาเย กายานุปสฺสนาสติปฏฺฐาเน ปนฺนรส ธมฺมา, ปนฺนรส อตฺถา, ติํส นิรุตฺติโย, สฏฺฐิ ญาณานิ.}

\prose{(ข.\csromanspacer{} ฆ) “อยํ เวทนาสุ ฯเปฯ.\csromanspacer{} อยํ จิตฺเต ฯเปฯ.\csromanspacer{} อยํ ธมฺเมสุ ธมฺมานุปสฺสนา”ติ ปุพฺเพ อนนุสฺสุเตสุ ธมฺเมสุ จกฺขุํ อุทปาทิ ฯเปฯ อาโลโก อุทปาทิ ฯเปฯ “สา โข ปนายํ ธมฺเมสุ ธมฺมานุปสฺสนา ภาเวตพฺพา”ติ ฯเปฯ “ภาวิตา”ติ ปุพฺเพ อนนุสฺสุเตสุ ธมฺเมสุ จกฺขุํ อุทปาทิ ฯเปฯ อาโลโก อุทปาทิ ฯเปฯ.}

\csromanheader{2}{6. ปฏิสมฺภิทากถา}
\prose{\csromanfolio{337}ธมฺเมสุ ธมฺมานุปสฺสนา สติปฏฺฐาเน ปนฺนรส ธมฺมา, ปนฺนรส อตฺถา, ติํส นิรุตฺติโย, สฏฺฐิ ญาณานิ.}

\prose{จตูสุ สติปฏฺฐาเนสุ สฏฺฐิ ธมฺมา, สฏฺฐิ อตฺถา, วีสติสตนิรุตฺติโย, จตฺตาลีสญฺจ เทฺว จ ญาณสตานิ.}

\csromanmatikaanchor{mtk.104}
\csromantocmarkref{subsection}{3. อิทฺธิปาทวาร}{mtk.104}
\bookmark[dest={mtk.104},level=2]{3. อิทฺธิปาทวาร}
\csromanheader{2}{3. อิทฺธิปาทวาร}
\proseitem{32}{“อยํ ฉนฺทสมาธิปธานสงฺขารสมนฺนาคโต อิทฺธิปาโท”ติ เม ภิกฺขเว ปุพฺเพ อนนุสฺสุเตสุ ธมฺเมสุ จกฺขุํ อุทปาทิ ฯเปฯ อาโลโก อุทปาทิ. “โส โข ปนายํ ฉนฺทสมาธิปธานสงฺขารสมนฺนาคโต อิทฺธิปาโท ภาเวตพฺโพ”ติ เม ภิกฺขเว ฯเปฯ “ภาวิโต”ติ เม ภิกฺขเว ปุพฺเพ อนนุสฺสุเตสุ ธมฺเมสุ จกฺขุํ อุทปาทิ ฯเปฯ อาโลโก อุทปาทิ.}

\prose{“อยํ วีริยสมาธิ ฯเปฯ.\csromanspacer{} อยํ จิตฺตสมาธิ ฯเปฯ.\csromanspacer{} อยํ วีมํสาสมาธิปธานสงฺขารสมนฺนาคโต อิทฺธิปาโท”ติ เม ภิกฺขเว ปุพฺเพ อนนุสฺสุเตสุ ธมฺเมสุ จกฺขุํ อุทปาทิ ฯเปฯ อาโลโก อุทปาทิ. “โส โข ปนายํ วีมํสาสมาธิปธานสงฺขารสมนฺนาคโต อิทฺธิปาโท ภาเวตพฺโพ”ติ เม ภิกฺขเว ปุพฺเพ อนนุสฺสุเตสุ ธมฺเมสุ จกฺขุํ อุทปาทิ ฯเปฯ อาโลโก อุทปาทิ.}

\prose{(ก) “อยํ ฉนฺทสมาธิปธานสงฺขารสมนฺนาคโต อิทฺธิปาโท”ติ ปุพฺเพ อนนุสฺสุเตสุ ธมฺเมสุ จกฺขุํ อุทปาทิ ฯเปฯ อาโลโก อุทปาทิ ฯเปฯ “โส โข ปนายํ ฉนฺทสมาธิปธานสงฺขารสมนฺนาคโต อิทฺธิปาโท ภาเวตพฺโพ”ติ ฯเปฯ “ภาวิโต”ติ ปุพฺเพ อนนุสฺสุเตสุ ธมฺเมสุ จกฺขุํ อุทปาทิ, ญาณํ อุทปาปาทิ, ปญฺญา อุทปาทิ, วิชฺชา อุทปาทิ, อาโลโก อุทปาทิ.}

\prose{\textbf{จกฺขุํ อุทปาที}ติ เกนฏฺเฐน, ญาณํ อุทปาทีติ เกนฏฺเฐน, ปญฺญา อุทปาทีติ เกนฏฺเฐน, วิชฺชา อุทปาทีติ เกนฏฺเฐน, อาโลโก อุทปาทีติ เกนฏฺเฐน? \textbf{จกฺขุํ อุทปาที}ติ ทสฺสนฏฺเฐน.\csromanspacer{} \textbf{ญาณํ อุทปาที}ติ ญาตฏฺเฐน.\csromanspacer{} \textbf{ปญฺญา อุทปาที}ติ ปชานนฏฺเฐน.\csromanspacer{} \textbf{วิชฺชา อุทปาที}ติ ปฏิเวธฏฺเฐน.\csromanspacer{} \textbf{อาโลโก อุทปาที}ติ โอภาสฏฺเฐน.}

\prose{จกฺขุํ ธมฺโม, ญาณํ ธมฺโม, ปญฺญา ธมฺโม, วิชฺชา ธมฺมา, อาโลโก ธมฺโม.\csromanspacer{} อิเม ปญฺจ ธมฺมา ธมฺมปฏิสมฺภิทาย อารมฺมณา เจว โหนฺติ โคจรา จ. \csromanfolio{338}เย ตสฺสา อารมฺมณา, เต ตสฺสา โคจรา.\csromanspacer{} ย ตสฺสา โคจรา, เต ตสฺสา อารมฺมณา.\csromanspacer{} เตน วุจฺจติ ธมฺเมสุ ญาณํ ธมฺมปฏิสมฺภิทา.}

\prose{ทสฺสนฏฺโฐ อตฺโถ, ญาตฏฺโฐ อตฺโถ, ปชานนฏฺโฐ อตฺโถ, ปฏิเวธฏฺโฐ อตฺโถ, โอภาสฏฺโฐ อตฺโถ.\csromanspacer{} อิเม ปญฺจ อตฺถา อตฺถปฏิสมฺภิทาย อารมฺมณา เจว โหนฺติ โคจรา จ.\csromanspacer{} เย ตสฺสา อารมฺมณา, เต ตสฺสา โคจรา.\csromanspacer{} เย ตสฺสา โคจรา, เต ตสฺสา อารมฺมณา.\csromanspacer{} เตน วุจฺจติ อตฺเถสุ ญาณํ อตฺถปฏิสมฺภิทา.}

\prose{ปญฺจ ธมฺเม สนฺทสฺเสตุํ พฺยญฺชนนิรุตฺตาภิลาปา, ปญฺจ อตฺเถ สนฺทสฺเสตุํ พฺยญฺชนนิรุตฺตาภิลาปา.\csromanspacer{} อิมา ทส นิรุตฺติโย นิรุตฺติปฏิสมฺภิทาย อารมฺมณา เจว โหนฺติ โคจรา จ.\csromanspacer{} เย ตสฺสา อารมฺมณา, เต ตสฺสา โคจรา.\csromanspacer{} เย ตสฺสา โคจรา, เต ตสฺสา อารมฺมณา.\csromanspacer{} เตน วุจฺจติ นิรุตฺตีสุ ญาณํ นิรุตฺติปฏิสมฺภิทา.}

\prose{ปญฺจสุ ธมฺเมสุ ญาณานิ, ปญฺจสุ อตฺเถสุ ญาณานิ, ทสสุ นิรุตฺตีสุ ญาณานิ.\csromanspacer{} อิมานิ วีสติ ญาณานิ ปฏิภานปฏิสมฺภิทาย อารมฺมณา เจว โหนฺติ โคจรา จ.\csromanspacer{} เย ตสฺสา อารมฺมณา, เต ตสฺสา โคจรา.\csromanspacer{} เย ตสฺสา โคจรา, เต ตสฺสา อารมฺมณา.\csromanspacer{} เตน วุจฺจติ ปฏิภาเนสุ ญาณํ ปฏิภานปฏิสมฺภิทา.}

\prose{ฉนฺทสมาธิปธานสงฺขารสมนฺนาคเต อิทฺธิปาเท ปนฺนรส ธมฺมา, ปนฺนรส อตฺถา, ติํส นิรุตฺติโย, สฏฺฐิ ญาณานิ.}

\prose{(ข.\csromanspacer{} ฆ) “อยํ วีริยสมาธิ ฯเปฯ.\csromanspacer{} อยํ จิตฺตสมาธิ ฯเปฯ.\csromanspacer{} อยํ วีมํสาสมาธิปธานสงฺขารสมนฺนาคโต อิทฺธิปาโท”ติ ปุพฺเพ อนนุสฺสุเตสุ ธมฺเมสุ จกฺขุํ อุทปาทิ ฯเปฯ อาโลโก อุทปาทิ ฯเปฯ “โส โข ปนายํ วีมํสาสมาธิปธานสงฺขารสมนฺนาคโต อิทฺธิปาโท ภาเวตพฺโพ”ติ ฯเปฯ “ภาวิโต”ติ ปุพฺเพ อนนุสฺสุเตสุ ธมฺเมสุ จกฺขุํ อุทปาทิ ฯเปฯ อาโลโก อุทปาทิ ฯเปฯ.}

\prose{วีมํสาสมาธิปธานสงฺขารสมนฺนาคเต อิทฺธิปาเท ปนฺนรส ธมฺมา, ปนฺนรส อตฺถา, ติํส นิรุตฺติโย, สฏฺฐิ ญาณานิ.}

\prose{จตูสุ อิทฺธิปาเทสุ สฏฺฐิ ธมฺมา, สฏฺฐิ อตฺถา, วีสติสตนิรุตฺติโย, จตฺตาลีสญฺจ เทฺว จ ญาณสตานิ.}

\csromanheader{2}{6. ปฏิสมฺภิทากถา}
\csromanmatikaanchor{mtk.105}
\csromantocmarkref{subsection}{4. สตฺตโพธิสตฺตวาร}{mtk.105}
\bookmark[dest={mtk.105},level=2]{4. สตฺตโพธิสตฺตวาร}
\csromanheader{2}{4. สตฺตโพธิสตฺตวาร}
\proseitem{33}{\csromanfolio{339}“สมุทโย สมุทโย”ติ โข ภิกฺขเว วิปสฺสิสฺส โพธิสตฺตสฺส ปุพฺเพ อนนุสฺสุเตสุ ธมฺเมสุ จกฺขุํ อุทปาทิ ฯเปฯ อาโลโก อุทปาทิ. “นิโรโธ นิโรโธ”ติ โข ภิกฺขเว วิปสฺสิสฺส โพธิสตฺตสฺส ปุพฺเพ อนนุสฺสุเตสุ ธมฺเมสุ จกฺขุํ อุทปาทิ ฯเปฯ อาโลโก อุทปาทิ.\csromanspacer{} วิปสฺสิสฺส โพธิสตฺตสฺส เวยฺยากรเณ ทส ธมฺมา, ทส อตฺถา, วีสติ นิรุตฺติโย, จตฺตาลีสํ\footnote{จตฺตารีส (ก)} ญาณานิ.}

\prose{“สมุทโย สมุทโย”ติ โข ภิกฺขเว สิขิสฺส โพธิสตฺตสฺส ฯเปฯ เวสฺสภุสฺส โพธิสตฺตสฺส ฯเปฯ กกุสนฺธสฺส โพธิสตฺตสฺส ฯเปฯ โกณาคมนสฺส โพธิสตฺตสฺส ฯเปฯ กสฺสปสฺส โพธิสตฺตสฺส ปุพฺเพ อนนุสฺสุเตสุ ธมฺเมสุ จกฺขุํ อุทปาทิ ฯเปฯ อาโลโก อุทปาทิ. “นิโรโธ นิโรโธ”ติ โข ภิกฺขเว กสฺสปสฺส โพธิสตฺตสฺส ปุพฺเพ อนนุสฺสุเตสุ ธมฺเมสุ จกฺขุํ อุทปาทิ ฯเปฯ อาโลโก อุทปาทิ.\csromanspacer{} กสฺสปสฺส โพธิสตฺตสฺส เวยฺยากรเณ ทส ธมฺมา, ทส อตฺถา, วีสติ นิรุตฺติโย จตฺตาลีสํ ญาณานิ.}

\prose{“สมุทโย สมุทโย”ติ โข ภิกฺขเว โคตมสฺส โพธิสตฺตสฺส ปุพฺเพ อนนุสฺสุเตสุ ธมฺเมสุ จกฺขุํ อุทปาทิ ฯเปฯ อาโลโก อุทปาทิ. “นิโรโธ นิโรโธ”ติ โข ภิกฺขเว โคตมสฺส โพธิสตฺตสฺส ปุพฺเพ อนนุสฺสุเตสุ ธมฺเมสุ จกฺขุํ อุทปาทิ ฯเปฯ อาโลโก อุทปาทิ.\csromanspacer{} โคตมสฺส โพธิสตฺตสฺส เวยฺยากรเณ ทส ธมฺมา, ทส อตฺถา, วีสติ นิรุตฺติโย, จตฺตาลีสํ ญาณานิ.}

\prose{สตฺตนฺนํ โพธิสตฺตานํ สตฺตสุ เวยฺยากรเณสุ สตฺตติ ธมฺมา, สตฺตติ อตฺถา, จตฺตาลีสสตํ\footnote{จตฺตารีสสตา (สฺยา)} นิรุตฺติโย, อสีติ จ เทฺว จ ญาณสตานิ.}

\csromanmatikaanchor{mtk.106}
\csromantocmarkref{subsection}{5. อภิญฺญาทิวาร}{mtk.106}
\bookmark[dest={mtk.106},level=2]{5. อภิญฺญาทิวาร}
\csromanheader{2}{5. อภิญฺญาทิวาร}
\proseitem{34}{ยาวตา อภิญฺญาย อภิญฺญฏฺโฐ ญาโต ทิฏฺโฐ วิทิโต สจฺฉิกโต ผสฺสิโต ปญฺญาย.\csromanspacer{} อผสฺสิโต ปญฺญาย อภิญฺญฏฺโฐ นตฺถีติ จกฺขุํ อุทปาทิ, ญาณํ อุทปาทิ, ปญฺญา อุทปาทิ, วิชฺชา อุทปาทิ, อาโลโก อุทปาทิ.\csromanspacer{} อภิญฺญาย อภิญฺญฏฺเฐ ปญฺจวีสติ ธมฺมา, ปญฺจวีสติ อตฺถา, ปญฺญาสํ นิรุตฺติโย, สตํ ญาณานิ.}

\prose{\csromanfolio{340}ยาวตา ปริญฺญาย ปริญฺญฏฺโฐ ฯเปฯ.\csromanspacer{} ยาวตา ปหานสฺส ปหานฏฺโฐ ฯเปฯ.\csromanspacer{} ยาวตา ภาวนาย ภาวนฏฺโฐ ฯเปฯ.\csromanspacer{} ยาวตา สจฺฉิกิริยาย สจฺฉิกิริยฏฺโฐ ญาโต ทิฏฺโฐ วิทิโต สจฺฉิกโต ผสฺสิโต ปญฺญาย.\csromanspacer{} อผสฺสิโต ปญฺญาย สจฺฉิกิริยฏฺโฐ นตฺถีติ จกฺขุํ อุทปาทิ, ญาณํ อุทปาทิ, ปญฺญา อุทปาทิ, วิชฺชา อุทปาทิ, อาโลโก อุทปาทิ, สจฺฉิกิริยาย สจฺฉิกิริยฏฺเฐ ปญฺจวีสติ ธมฺมา, ปญฺจวีสติ อตฺถา, ปญฺญาสํ นิรุตฺติโย, สตํ ญาณานิ.}

\prose{อภิญฺญาย อภิญฺญฏฺเฐ ปริญฺญาย ปริญฺญฏฺเฐ ปหานาย ปหานฏฺเฐ.\csromanspacer{} ภาวนาย ภาวนฏฺเฐ.\csromanspacer{} สจฺฉิกิริยาย สจฺฉิกิริยฏฺเฐ ปญฺจวีสสตํ ธมฺมา ปญฺจวีสสตํ อตฺถา, อฑฺฒเตยฺยานิ นิรุตฺติสตานิ, ปญฺจ ญาณสตานิ.}

\csromanmatikaanchor{mtk.107}
\csromantocmarkref{subsection}{6. ขนฺธาทิวาร}{mtk.107}
\bookmark[dest={mtk.107},level=2]{6. ขนฺธาทิวาร}
\csromanheader{2}{6. ขนฺธาทิวาร}
\proseitem{35}{ยาวตา ขนฺธานํ ขนฺธฏฺโฐ ญาโต ทิฏฺโฐ วิทิโต สจฺฉิกโต ผสฺสิโต ปญฺญาย.\csromanspacer{} อผสฺสิโต ปญฺญาย ขนฺธฏฺโฐ นตฺถีติ จกฺขุํ อุทปาทิ.\csromanspacer{} ญาณํ อุทปาทิ, ปญฺญา อุทปาทิ, วิชฺชา อุทปาทิ, อาโลโก อุทปาทิ, ขนฺธานํ ขนฺธฏฺเฐ ปญฺจวีสติ ธมฺมา, ปญฺจวีสติ อตฺถา, ปญฺญาสํ นิรุตฺติโย, สตํ ญาณานิ.}

\prose{ยาวตา ธาตูนํ ธาตุฏฺโฐ ฯเปฯ.\csromanspacer{} ยาวตา อายตนานํ อายตนฏฺโฐ ฯเปฯ.\csromanspacer{} ยาวตา สงฺขตานํ สงฺขตฏฺโฐ ฯเปฯ.\csromanspacer{} ยาวตา อสงฺขตสฺส อสงฺขตฏฺโฐ ญาโต ทิฏฺโฐ วิทิโต สจฺฉิกโต ผสฺสิโต ปญฺญาย.\csromanspacer{} อผสฺสิโต ปญฺญาย อสงฺขตฏฺโฐ นตฺถีติ จกฺขุํ อุทปาทิ ฯเปฯ อาโลโก อุทปาทิ.\csromanspacer{} อสงฺขตสฺส อสงฺขตฏฺเฐ ปญฺจวีสติ ธมฺมา, ปญฺจวีสติ อตฺถา, ปญฺญาสํ นิรุตฺติโย, สตํ ญาณานิ.}

\prose{ขนฺธานํ ขนฺธฏฺเฐ ธาตูนํ ธาตุฏฺเฐ อายตนานํ อายตนฏฺเฐ สงฺขตานํ สงฺขตฏฺเฐ อสงฺขตสฺส อสงฺขตฏฺเฐ ปญฺจวีสติสตํ ธมฺมา, ปญฺจวีสติสตํ อตฺถา, อฑฺฒเตยฺยานิ นิรุตฺติสตานิ, ปญฺจ ญาณสตานิ.}

\csromanmatikaanchor{mtk.108}
\csromantocmarkref{subsection}{7. สจฺจวาร}{mtk.108}
\bookmark[dest={mtk.108},level=2]{7. สจฺจวาร}
\csromanheader{2}{7. สจฺจวาร}
\proseitem{36}{ยาวตา ทุกฺขสฺส ทุกฺขฏฺโฐ ญาโต ทิฏฺโฐ วิทิโต สจฺฉิกโต ผสฺสิโต ปญฺญาย.\csromanspacer{} อผสฺสิโต ปญฺญาย ทุกฺขฏฺโฐ นตฺถีติ จกฺขุํ อุทปาทิ, ญาณํ อุทปาทิ, ปญฺญา อุทปาทิ, วิชฺชา อุทปาทิ, อาโลโก อุทปาทิ.\csromanspacer{} ทุกฺขสฺส}

\csromanheader{2}{6. ปฏิสมฺภิทากถา}
\prose{\csromanfolio{341}ทุกฺขฏฺเฐ ปญฺจวีสติ ธมฺมา, ปญฺจวีสติ อตฺถา, ปญฺญาสํ นิรุตฺติโย, สตํ ญาณานิ.}

\prose{ยาวตา สมุทยสฺส สมุทยฏฺโฐ ฯเปฯ.\csromanspacer{} ยาวตา นิโรธสฺส นิโรธฏฺโฐ ฯเปฯ.\csromanspacer{} ยาวตา มคฺคสฺส มคฺคฏฺโฐ ญาโต ทิฏฺโฐ วิทิโต สจฺฉิกโต ผสฺสิโต ปญฺญาย.\csromanspacer{} อผสฺสิโต ปญฺญาย มคฺคฏฺโฐ นตฺถีติ จกฺขุํ อุทปาทิ ฯเปฯ อาโลโก อุทปาทิ.\csromanspacer{} มคฺคสฺส มคฺคฏฺเฐ ปญฺจวีสติ ธมฺมา, ปญฺจวีสติ อตฺถา.\csromanspacer{} ปญฺญาสํ นิรุตฺติโย, สตํ ญาณานิ.}

\prose{จตูสุ อริยสจฺเจสุ สตํ ธมฺมา, สตํ อตฺถา, เทฺว นิรุตฺติสตานิ, จตฺตาริ ญาณสตานิ.}

\csromanmatikaanchor{mtk.109}
\csromantocmarkref{subsection}{8. ปฏิสมฺภิทาวาร}{mtk.109}
\bookmark[dest={mtk.109},level=2]{8. ปฏิสมฺภิทาวาร}
\csromanheader{2}{8. ปฏิสมฺภิทาวาร}
\proseitem{37}{ยาวตา อตฺถปฏิสมฺภิทาย อตฺถปฏิสมฺภิทฏฺโฐ ญาโต ทิฏฺโฐ วิทิโต สจฺฉิกโต ผสฺสิโต ปญฺญาย.\csromanspacer{} อผสฺสิโต ปญฺญาย อตฺถปฏิสมฺภิทฏฺโฐ นตฺถีติ จกฺขุํ อุทปาทิ ฯเปฯ อาโลโก อุทปาทิ.\csromanspacer{} อตฺถปฏิสมฺภิทาย อตฺถปฏิสมฺภิทฏฺเฐ ปญฺจวีสติ ธมฺมา, ปญฺจวีสติ อตฺถา.\csromanspacer{} ปญฺญาสํ นิรุตฺติโย, สตํ ญาณานิ.}

\prose{ยาวตา ธมฺมปฏิสมฺภิทาย ธมฺมปฏิสมฺภิทฏฺโฐ ฯเปฯ.\csromanspacer{} ยาวตา นิรุตฺติปฏิสมฺภิทาย นิรุตฺติปฏิสมฺภิทฏฺโฐ ฯเปฯ.\csromanspacer{} ยาวตา ปฏิภานปฏิสมฺภิทาย ปฏิภานปฏิสมฺภิทฏฺโฐ ญาโต ทิฏฺโฐ วิทิโต สจฺฉิกโต ผสฺสิโต ปญฺญาย.\csromanspacer{} อผสฺสิโต ปญฺญาย ปฏิภานปฏิสมฺภิทฏฺโฐ นตฺถีติ จกฺขุํ อุทปาทิ ฯเปฯ อาโลโก อุทปาทิ.\csromanspacer{} ปฏิภานปฏิสมฺภิทฏฺเฐ ปญฺจวีสติ ธมฺมา ปญฺจวีสติ อตฺถา, ปญฺญาสํ นิรุตฺติโย, สตํ ญาณานิ.}

\prose{จตูสุ ปฏิสมฺภิทาสุ สตํ ธมฺมา, สตํ อตฺถา, เทฺว นิรุตฺติสตานิ, จตฺตาริ ญาณสตานิ.}

\csromanmatikaanchor{mtk.110}
\csromantocmarkref{subsection}{9. ฉพุทฺธธมฺมวาร}{mtk.110}
\bookmark[dest={mtk.110},level=2]{9. ฉพุทฺธธมฺมวาร}
\csromanheader{2}{9. ฉพุทฺธธมฺมวาร}
\proseitem{38}{ยาวตา อินฺทฺริยปโรปริยตฺตญาณํ ญาตํ ทิฏฺฐํ วิทิตํ สจฺฉิกตํ ผสฺสิตํ ปญฺญาย, อผสฺสิตํ ปญฺญาย อินฺทฺริยปโรปริยตฺตญาณํ นตฺถีติ จกฺขุํ อุทปาทิ ฯเปฯ อาโลโก อุทปาทิ, อินฺทฺริยปโรปริยตฺตญาเณ ปญฺจวีสติ ธมฺมา, ปญฺจวีสติ อตฺถา, ปญฺญาสํ นิรุตฺติโย, สตํ ญาณานิ.}

\prose{\csromanfolio{342}ยาวตา สตฺตานํ อาสยานุสเย ญาณํ ฯเปฯ.\csromanspacer{} ยาวตา ยมกปาฏิหีเร ญาณํ ฯเปฯ.\csromanspacer{} ยาวตา มหากรุณาสมาปตฺติยา ญาณํ ฯเปฯ.\csromanspacer{} ยาวตา สพฺพญฺญุตญฺญาณํ ฯเปฯ.\csromanspacer{} ยาวตา อนาวรณํ ญาณํ ญาตํ ทิฏฺฐํ วิทิตํ สจฺฉิกตํ ผสฺสิตํ ปญฺญาย, อผสฺสิตํ ปญฺญาย อนาวรณํ ญาณํ นตฺถีติ จกฺขุํ อุทปาทิ, ญาณํ อุทปาทิ, ปญฺญา อุทปาทิ, วิชฺชา อุทปาทิ, อาโลโก อุทปาทิ.\csromanspacer{} อนาวรเณ ญาเณ ปญฺจวีสติ ธมฺมา, ปญฺจวีสติ อตฺถา, ปญฺญาสํ นิรุตฺติโย, สตํ ญาณานิ.}

\prose{ฉสุ พุทฺธธมฺเมสุ ทิยฑฺฒสตํ ธมฺมา, ทิยฑฺฒสตํ อตฺถา, ตีณิ นิรุตฺติสตานิ, ฉ ญาณสตานิ.}

\prose{ปฏิสมฺภิทาธิกรเณ\footnote{ปฏิสมฺภิทาปกรเณ (สฺยา)} อฑฺฒนวธมฺมสตานิ\footnote{อฑฺฒนวมานิ ธมฺมสตานิ (สฺยา), อฑฺฒนวมธมฺมสตานิ (ก)}, อฑฺฒนว-อตฺถสตานิ, นิรุตฺติสหสฺสญฺจ สตฺต จ นิรุตฺติสตานิ, ตีณิ จ ญาณสหสฺสานิ, จตฺตาริ จ ญาณสตานีติ.}

\nitthitamruled{ปฏิสมฺภิทากถา นิฏฺฐิตา.}
\csromanmatikaanchor{mtk.111}
\csromanmatikaanchor{mtk.112}
\csromantocmarknopage{section}{7. ธมฺมจกฺกกถา}{mtk.111}
\bookmark[dest={mtk.111},level=1]{7. ธมฺมจกฺกกถา}
\csromantocmarkref{subsection}{1. สจฺจวาร}{mtk.112}
\bookmark[dest={mtk.112},level=2]{1. สจฺจวาร}
\csromanheader{2}{7. ธมฺมจกฺกกถา 1. สจฺจวาร}
\proseitem{39}{เอกํ สมยํ ภควา พาราณสิยํ วิหรติ ฯเปฯ.\csromanspacer{} อิติ หิทํ อายสฺมโต โกณฺฑญฺญสฺส “อญฺญาสิโกณฺฑญฺโญ”เตฺวว นามํ อโหสิ. (ก) “อิทํ ทุกฺขํ อริยสจฺจนฺ”ติ ปุพฺเพ อนนุสฺสุเตสุ ธมฺเมสุ จกฺขุํ อุทปาทิ, ญาณํ อุทปาทิ, ปญฺญา อุทปาทิ, วิชฺชา อุทปาทิ, อาโลโก อุทปาทิ.}

\prose{\textbf{จกฺขุํ อุทปาที}ติ เกนฏฺเฐน, ญาณํ อุทปาทีติ เกนฏฺเฐน, ปญฺญา อุทปาทีติ เกนฏฺเฐน, วิชฺชา อุทปาทีติ เกนฏฺเฐน, อาโลโก อุทาปาทีติ เกนฏฺเฐน? \textbf{จกฺขุํ อุทปาที}ติ ทสฺสนฏฺเฐน.\csromanspacer{} \textbf{ญาณํ อุทปาที}ติ ญาตฏฺเฐน.\csromanspacer{} \textbf{ปญฺญา อุทปาที}ติ ปชานนฏฺเฐน.\csromanspacer{} \textbf{วิชฺชา อุทปาที}ติ ปฏิเวธฏฺเฐน.\csromanspacer{} \textbf{อาโลโก อุทปาที}ติ โอภาสฏฺเฐน.}

\prose{จกฺขุํ ธมฺโม, ทสฺสนฏฺโฐ อตฺโถ.\csromanspacer{} ญาณํ ธมฺโม, ญาตฏฺโฐ อตฺโถ.\csromanspacer{} ปญฺญา ธมฺโม, ปชานนฏฺโฐ อตฺโถ.\csromanspacer{} วิชฺชา ธมฺโม, ปฏิเวธฏฺโฐ อตฺโถ, อาโลโก ธมฺโม, โอภาสฏฺโฐ อตฺโถ.\csromanspacer{} อิเม ปญฺจ ธมฺมา ปญฺจ อตฺถา}

\csromanheader{2}{7. ธมฺมจกฺกกถา}
\prose{\csromanfolio{343}ทุกฺขวตฺถุกา สจฺจวตฺถุกา สจฺจารมฺมณา สจฺจโคจรา สจฺจสงฺคหิตา สจฺจปริยาปนฺนา สจฺเจ สมุทาคตา สจฺเจ ฐิตา สจฺเจ ปติฏฺฐิตา.}

\proseitem{40}{\textbf{ธมฺมจกฺกนฺ}ติ เกนฏฺเฐน ธมฺมจกฺกํ? ธมฺมญฺจ ปวตฺเตติ จกฺกญฺจาติ ธมฺมจกฺกํ, จกฺกญฺจ ปวตฺเตติ ธมฺมญฺจาติ ธมฺมจกฺกํ, ธมฺเมน ปวตฺเตตีติ ธมฺมจกฺกํ, ธมฺมจริยาย ปวตฺเตตีติ ธมฺมจกฺกํ, ธมฺเม ฐิโต ปวตฺเตตีติ ธมฺมจกฺกํ, ธมฺเม ปติฏฺฐิโต ปวตฺเตตีติ ธมฺมจกฺกํ, ธมฺเม ปติฏฺฐาเปนฺโต ปวตฺเตตีติ ธมฺมจกฺกํ, ธมฺเม วสิปฺปตฺโต ปวตฺเตตีติ ธมฺมจกฺกํ, ธมฺเม วสิํ ปาเปนฺโต ปวตฺเตตีติ ธมฺมจกฺกํ, ธมฺเม ปารมิปฺปตฺโต ปวตฺเตตีติ ธมฺมจกฺกํ, ธมฺเม ปารมิํ ปาเปนฺโต ปวตฺเตตีติ ธมฺมจกฺกํ, ธมฺเม เวสารชฺชปฺปตฺโต ปวตฺเตตีติ ธมฺมจกฺกํ, ธมฺเม เวสารชฺชํ ปาเปนฺโต ปวตฺเตตีติ ธมฺมจกฺกํ, ธมฺมํ สกฺกโรนฺโต ปวตฺเตตีติ ธมฺมจกฺกํ, ธมฺมํ ครุํ กโรนฺโต\footnote{ครุกโรนฺโต (สฺยา)} ปวตฺเตตีติ ธมฺมจกฺกํ, ธมฺมํ มาเนนฺโต ปวตฺเตตีติ ธมฺมจกฺกํ, ธมฺมํ ปูเชนฺโต ปวตฺเตตีติ ธมฺมจกฺกํ, ธมฺมํ อปจายมาโน ปวตฺเตตีติ ธมฺมจกฺกํ, ธมฺมทฺธโช ปวตฺเตตีติ ธมฺมจกฺกํ, ธมฺมเกตุ ปวตฺเตตีติ ธมฺมจกฺกํ, ธมฺมาธิปเตยฺโย ปวตฺเตตีติ ธมฺมจกฺกํ.\csromanspacer{} ตํ โข ปน ธมฺมจกฺกํ อปฺปฏิวตฺติยํ สมเณน วา พฺราหฺมเณน วา เทเวน วา มาเรน วา พฺรหฺมุนา วา เกนจิ วา โลกสฺมินฺติ ธมฺมจกฺกํ.}

\prose{สทฺธินฺทฺริยํ ธมฺโม, ตํ ธมฺมํ ปวตฺเตตีติ ธมฺมจกฺกํ, วีริยินฺทฺริยํ ธมฺโม, ตํ ธมฺมํ ปวตฺเตตีติ ธมฺมจกฺกํ, สตินฺทฺริยํ ธมฺโม, ตํ ธมฺมํ ปวตฺเตตีติ ธมฺมจกฺกํ, สมาธินฺทฺริยํ ธมฺโม, ตํ ธมฺมํ ปวตฺเตตีติ ธมฺมจกฺกํ, ปญฺญินฺทฺริยํ ธมฺโม ตํ ธมฺมํ ปวตฺเตตีติ ธมฺมจกฺกํ.\csromanspacer{}\csromanspacer{}สทฺธาพลํ ธมฺโม, ตํ ธมฺมํ ปวตฺเตตีติ ธมฺมจกฺกํ, วีริยพลํ ธมฺโม, ตํ ธมฺมํ ปวตฺเตตีติ ธมฺมจกฺกํ, สติพลํ ธมฺโม, ตํ ธมฺมํ ปวตฺเตตีติ ธมฺมจกฺกํ, สมาธิพลํ ธมฺโม, ตํ ธมฺมํ ปวตฺเตตีติ ธมฺมจกฺกํ, ปญฺญาพลํ ธมฺโม, ตํ ธมฺมํ ปวตฺเตตีติ ธมฺมจกฺกํ.\csromanspacer{}\csromanspacer{}สติสมฺโพชฺฌงฺโค ธมฺโม, ตํ ธมฺมํ ปวตฺเตตีติ ธมฺมจกฺกํ, ธมฺมวิจยสมฺโพชฺฌงฺโค ธมฺโม, ตํ ธมฺมํ ปวตฺเตตีติ ธมฺมจกฺกํ, วีริยสมฺโพชฺฌงฺโค ธมฺโม, ตํ ธมฺมํ ปวตฺเตตีติ ธมฺมจกฺกํ, ปีติสมฺโพชฺฌงฺโค ธมฺโม, ตํ ธมฺมํ ปวตฺเตตีติ ธมฺมจกฺกํ, ปสฺสทฺธิสมฺโพชฺฌงฺโค ธมฺโม, ตํ ธมฺมํ ปวตฺเตตีติ ธมฺมจกฺกํ, สมาธิสมฺโพชฺฌงฺโค ธมฺโม, ตํ ธมฺมํ ปวตฺเตตีติ ธมฺมจกฺกํ, อุเปกฺขาสมฺโพชฺฌงฺโค ธมฺโม, ตํ ธมฺมํ \csromanfolio{344}ปวตฺเตตีติ ธมฺมจกฺกํ.\csromanspacer{}\csromanspacer{}สมฺมาทิฏฺฐิ ธมฺโม, ตํ ธมฺมํ ปวตฺเตตีติ ธมฺมจกฺกํ, สมฺมาสงฺกปฺโป ธมฺโม, ตํ ธมฺมํ ปวตฺเตตีติ ธมฺมจกฺกํ, สมฺมาวาจา ธมฺโม, ตํ ธมฺมํ ปวตฺเตตีติ ธมฺมจกฺกํ, สมฺมากมฺมนฺโต ธมฺโม, ตํ ธมฺมํ ปวตฺเตตีติ ธมฺมจกฺกํ, สมฺมา-อาชีโว ธมฺโม, ตํ ธมฺมํ ปวตฺเตตีติ ธมฺมจกฺกํ, สมฺมาวายาโม ธมฺโม, ตํ ธมฺมํ ปวตฺเตตีติ ธมฺมจกฺกํ, สมฺมาสติ ธมฺโม, ตํ ธมฺมํ ปวตฺเตตีติ ธมฺมจกฺกํ, สมฺมาสมาธิ ธมฺโม, ตํ ธมฺมํ ปวตฺเตตีติ ธมฺมจกฺกํ.}

\prose{อาธิปเตยฺยฏฺเฐน อินฺทฺริยํ ธมฺโม, ตํ ธมฺมํ ปวตฺเตตีติ ธมฺมจกฺกํ, อกมฺปิยฏฺเฐน พลํ ธมฺโม, ตํ ธมฺมํ ปวตฺเตตีติ ธมฺมจกฺกํ, นิยฺยานิกฏฺเฐน โพชฺฌงฺโค ธมฺโม, ตํ ธมฺมํ ปวตฺเตตีติ ธมฺมจกฺกํ, เหตุฏฺเฐน มคฺโค ธมฺโม, ตํ ธมฺมํ ปวตฺเตตีติ ธมฺมจกฺกํ, อุปฏฺฐานฏฺเฐน สติปฏฺฐานา ธมฺโม, ตํ ธมฺมํ ปวตฺเตตีติ ธมฺมจกฺกํ, ปทหนฏฺเฐน สมฺมปฺปธานา ธมฺโม, ตํ ธมฺมํ ปวตฺเตตีติ ธมฺมจกฺกํ, อิชฺฌนฏฺเฐน อิทฺธิปาทา ธมฺโม, ตํ ธมฺมํ ปวตฺเตตีติ ธมฺมจกฺกํ, ตถฏฺเฐน สจฺจา ธมฺโม, ตํ ธมฺมํ ปวตฺเตตีติ ธมฺมจกฺกํ, อวิกฺเขปฏฺเฐน สมโถ ธมฺโม, ตํ ธมฺมํ ปวตฺเตตีติ ธมฺมจกฺกํ, อนุปสฺสนฏฺเฐน วิปสฺสนา ธมฺโม, ตํ ธมฺมํ ปวตฺเตตีติ ธมฺมจกฺกํ, เอกรสฏฺเฐน สมถวิปสฺสนา ธมฺโม, ตํ ธมฺมํ ปวตฺเตตีติ ธมฺมจกฺกํ, อนติวตฺตนฏฺเฐน ยุคนทฺธํ ธมฺโม, ตํ ธมฺมํ ปวตฺเตตีติ ธมฺมจกฺกํ, สํวรฏฺเฐน สีลวิสุทฺธิ ธมฺโม, ตํ ธมฺมํ ปวตฺเตตีติ ธมฺมจกฺกํ อวิกฺเขปฏฺเฐน จิตฺตวิสุทฺธิ ธมฺโม, ตํ ธมฺมํ ปวตฺเตตีติ ธมฺมจกฺกํ, ทสฺสนฏฺเฐน ทิฏฺฐิวิสุทฺธิ ธมฺโม, ตํ ธมฺมํ ปวตฺเตตีติ ธมฺมจกฺกํ, มุตฺตฏฺเฐน วิโมกฺโข ธมฺโม, ตํ ธมฺมํ ปวตฺเตตีติ ธมฺมจกฺกํ, ปฏิเวธฏฺเฐน วิชฺชา ธมฺโม, ตํ ธมฺมํ ปวตฺเตตีติ ธมฺมจกฺกํ, ปริจฺจาคฏฺเฐน วิมุตฺติ ธมฺโม, ตํ ธมฺมํ ปวตฺเตตีติ ธมฺมจกฺกํ, สมุจฺเฉทฏฺเฐน ขเย ญาณํ ธมฺโม, ตํ ธมฺมํ ปวตฺเตตีติ ธมฺมจกฺกํ, ปฏิปฺปสฺสทฺธฏฺเฐน อนุปฺปาเท ญาณํ ธมฺโม, ตํ ธมฺมํ ปวตฺเตตีติ ธมฺมจกฺกํ, ฉนฺโท มูลฏฺเฐน ธมฺโม, ตํ ธมฺมํ ปวตฺเตตีติ ธมฺมจกฺกํ, มนสิกาโร สมุฏฺฐานฏฺเฐน ธมฺโม, ตํ ธมฺมํ ปวตฺเตตีติ ธมฺมจกฺกํ, ผสฺโส สโมธานฏฺเฐน ธมฺโม, ตํ ธมฺมํ ปวตฺเตตีติ ธมฺมจกฺกํ, เวทนา สโมสรณฏฺเฐน ธมฺโม, ตํ ธมฺมํ ปวตฺเตตีติ ธมฺมจกฺกํ, สมาธิ ปมุขฏฺเฐน ธมฺโม, ตํ ธมฺมํ ปวตฺเตตีติ ธมฺมจกฺกํ, สติ อาธิปเตยฺยฏฺเฐน ธมฺโม, ตํ ธมฺมํ ปวตฺเตตีติ ธมฺมจกฺกํ, ปญฺญา ตทุตฺตรฏฺเฐน ธมฺโม, ตํ ธมฺมํ ปวตฺเตตีติ ธมฺมจกฺกํ, วิมุตฺติ สารฏฺเฐน ธมฺโม, ตํ ธมฺมํ ปวตฺเตตีติ ธมฺมจกฺกํ, อมโตคธํ นิพฺพานํ ปริโยสานฏฺเฐน ธมฺโม, ตํ ธมฺมํ ปวตฺเตตีติ ธมฺมจกฺกํ.}

\csromanheader{2}{7. ธมฺมจกฺกกถา}
\prose{\csromanfolio{345}“ตํ โข ปนิทํ ทุกฺขํ อริยสจฺจํ ปริญฺเญยฺยนฺ”ติ ฯเปฯ “ปริญฺญาตนฺ”ติ ปุพฺเพ อนนุสฺสุเตสุ ธมฺเมสุ จกฺขุํ อุทปาทิ ฯเปฯ อาโลโก อุทปาทิ.}

\prose{\textbf{จกฺขุํ อุทปาที}ติ เกนฏฺเฐน ฯเปฯ อาโลโก อุทปาทีติ เกนฏฺเฐน? \textbf{จกฺขุํ อุทปาที}ติ ทสฺสนฏฺเฐน ฯเปฯ.\csromanspacer{} \textbf{อาโลโก อุทปาที}ติ โอภาสฏฺเฐน.\csromanspacer{} จกฺขุํ ธมฺโม, ทสฺสนฏฺโฐ อตฺโถ ฯเปฯ อาโลโก ธมฺโม โอภาสฏฺโฐ อตฺโถ.\csromanspacer{} อิเม ปญฺจ ธมฺมา ปญฺจ อตฺถา ทุกฺขวตฺถุกา สจฺจวตฺถุกา สจฺจารมฺมณา สจฺจโคจรา สจฺจสงฺคหิตา สจฺจปริยาปนฺนา สจฺเจ สมุทาคตา สจฺเจ ฐิตา สจฺเจ ปติฏฺฐิตา.}

\prose{\textbf{ธมฺมจกฺกนฺ}ติ เกนฏฺเฐน ธมฺมจกฺกํ? ธมฺมญฺจ ปวตฺเตติ จกฺกญฺจาติ ธมฺมจกฺกํ, จกฺกญฺจ ปวตฺเตติ ธมฺมญฺจาติ ธมฺมจกฺกํ, ธมฺเมน ปวตฺเตตีติ ธมฺมจกฺกํ, ธมฺมจริยาย ปวตฺเตตีติ ธมฺมจกฺกํ, ธมฺเม ฐิโต ปวตฺเตตีติ ธมฺมจกฺกํ, ธมฺเม ปติฏฺฐิโต ปวตฺเตตีติ ธมฺมจกฺกํ ฯเปฯ อมโตคธํ นิพฺพานํ ปริโยสานฏฺเฐน ธมฺโม, ตํ ธมฺมํ ปวตฺเตตีติ ธมฺมจกฺกํ.}

\prose{(ข.\csromanspacer{} ฆ) “อิทํ ทุกฺขสมุทยํ อริยสจฺจนฺ”ติ ปุพฺเพ อนนุสฺสุเตสุ ธมฺเมสุ จกฺขุํ อุทปาทิ ฯเปฯ อาโลโก อุทปาทิ ฯเปฯ “ตํ โข ปนิทํ ทุกฺขสมุทยํ อริยสจฺจํ ปหาตพฺพนฺ”ติ ฯเปฯ “ปหีนนฺ”ติ ปุพฺเพ อนนุสฺสุเตสุ ธมฺเมสุ จกฺขุํ อุทปาทิ ฯเปฯ อาโลโก อุทปาทิ.}

\prose{\textbf{จกฺขุํ อุทปาที}ติ เกนฏฺเฐน ฯเปฯ อาโลโก อุทปาทีติ เกนฏฺเฐน? \textbf{จกฺขุํ อุทปาที}ติ ทสฺสนฏฺเฐน ฯเปฯ.\csromanspacer{} \textbf{อาโลโก อุทปาที}ติ โอภาสฏฺเฐน.}

\prose{จกฺขุํ ธมฺโม, ทสฺสนฏฺโฐ อตฺโถ ฯเปฯ อาโลโก ธมฺโม, โอภาสฏฺโฐ อตฺโถ.\csromanspacer{} อิเม ปญฺจ ธมฺมา ปญฺจ อตฺถา สมุทยวตฺถุกา สจฺจวตฺถุกา ฯเปฯ นิโรธวตฺถุกา สจฺจวตฺถุกา ฯเปฯ มคฺควตฺถุกา สจฺจวตฺถุกา สจฺจารมฺมณา สจฺจโคจรา สจฺจสงฺคหิตา สจฺจปริยาปนฺนา สจฺเจ สมุทาคตา สจฺเจ ฐิตา สจฺเจ ปติฏฺฐิตา.}

\prose{\textbf{ธมฺมจกฺกนฺ}ติ เกนฏฺเฐน ธมฺมจกฺกํ? ธมฺมญฺจ ปวตฺเตติ จกฺกญฺจาติ ธมฺมจกฺกํ, จกฺกญฺจ ปวตฺเตติ ธมฺมญฺจาติ ธมฺมจกฺกํ, ธมฺเมน ปวตฺเตตีติ ธมฺมจกฺกํ, ธมฺมจริยาย ปวตฺเตตีติ ธมฺมจกฺกํ, ธมฺเม ฐิโต ปวตฺเตตีติ ธมฺมจกฺกํ, ธมฺเม ปติฏฺฐิโต ปวตฺเตตีติ ธมฺมจกฺกํ ฯเปฯ อมโตคธํ นิพฺพานํ ปริโยสานฏฺเฐน ธมฺโม, ตํ ธมฺมํ ปวตฺเตตีติ ธมฺมจกฺกํ.}

\csromanmatikaanchor{mtk.113}
\csromantocmarkref{subsection}{2. สติปฏฺฐานวาร}{mtk.113}
\bookmark[dest={mtk.113},level=2]{2. สติปฏฺฐานวาร}
\csromanheader{2}{2. สติปฏฺฐานวาร}
\proseitem{41}{\csromanfolio{346}“อยํ กาเย กายานุปสฺสนา”ติ เม ภิกฺขเว ปุพฺเพ อนนุสฺสุเตสุ ธมฺเมสุ จกฺขุํ อุทปาทิ ฯเปฯ อาโลโก อุทปาทิ. “สา โข ปนายํ กาเย กายานุปสฺสนา ภาเวตพฺพา”ติ เม ภิกฺขเว ฯเปฯ “ภาวิตา”ติ เม ภิกฺขเว ปุพฺเพ อนนุสฺสุเตสุ ธมฺเมสุ จกฺขุํ อุทปาทิ ฯเปฯ อาโลโก อุทปาทิ.}

\prose{“อยํ เวทนาสุ ฯเปฯ.\csromanspacer{} อยํ จิตฺเต.\csromanspacer{} อยํ ธมฺเมสุ ธมฺมานุปสฺสนา”ติ เม ภิกฺขเว ปุพฺเพ อนนุสฺสุเตสุ ธมฺเมสุ จกฺขุํ อุทปาทิ ฯเปฯ อาโลโก อุทปาทิ. “สา โข ปนายํ ธมฺเมสุ ธมฺมานุปสฺสนา ภาเวตพฺพา”ติ เม ภิกฺขเว ฯเปฯ “ภาวิตา”ติ เม ภิกฺขเว ปุพฺเพ อนนุสฺสุเตสุ ธมฺเมสุ จกฺขุํ อุทปาทิ ฯเปฯ อาโลโก อุทปาทิ.}

\prose{“อยํ กาเย กายานุปสฺสนา”ติ ปุพฺเพ อนนุสฺสุเตสุ ธมฺเมสุ จกฺขุํ อุทปาทิ ฯเปฯ อาโลโก อุทปาทิ ฯเปฯ “สา โข ปนายํ กาเย กายานุปสฺสนา ภาเวตพฺพา”ติ ฯเปฯ “ภาวิตา”ติ ปุพฺเพ อนนุสฺสุเตสุ ธมฺเมสุ จกฺขุํ อุทปาทิ ฯเปฯ อาโลโก อุทปาทิ.}

\prose{\textbf{จกฺขุํ อุทปาที}ติ เกนฏฺเฐน ฯเปฯ อาโลโก \textbf{อุทปาที}ติ เกนฏฺเฐน? \textbf{จกฺขุํ อุทปาที}ติ ทสฺสนฏฺเฐน ฯเปฯ.\csromanspacer{} \textbf{อาโลโก} \textbf{อุทปาที}ติ โอภาสฏฺเฐน.}

\prose{จกฺขุํ ธมฺโม, ทสฺสนฏฺโฐ อตฺโถ ฯเปฯ อาโลโก ธมฺโม, โอภาสฏฺโฐ อตฺโถ.\csromanspacer{} อิเม ปญฺจ ธมฺมา ปญฺจ อตฺถา กายวตฺถุกา สติปฏฺฐานวตฺถุกา ฯเปฯ เวทนาวตฺถุกา สติปฏฺฐานวตฺถุกา, จิตฺตวตฺถุกา สติปฏฺฐานวตฺถุกา.\csromanspacer{} ธมฺมวตฺถุกา สติปฏฺฐานวตฺถุกา สติปฏฺฐานารมฺมณา สติปฏฺฐานโคจรา สติปฏฺฐานสงฺคหิตา สติปฏฺฐานปริยาปนฺนา สติปฏฺฐาเน สมุทาคตา สติปฏฺฐาเน ฐิตา สติปฏฺฐาเน ปติฏฺฐิตา.}

\prose{\textbf{ธมฺมจกฺกนฺ}ติ เกนฏฺเฐน ธมฺมจกฺกํ? ธมฺมญฺจ ปวตฺเตติ จกฺกญฺจาติ ธมฺมจกฺกํ, จกฺกญฺจ ปวตฺเตติ ธมฺมญฺจาติ ธมฺมจกฺกํ, ธมฺเมน ปวตฺเตตีติ ธมฺมจกฺกํ, ธมฺมจริยาย ปวตฺเตตีติ ธมฺมจกฺกํ, ธมฺเม ฐิโต ปวตฺเตตีติ ธมฺมจกฺกํ, ธมฺเม ปติฏฺฐิโต ปวตฺเตตีติ ธมฺมจกฺกํ ฯเปฯ อมโตคธํ นิพฺพานํ ปริโยสานฏฺเฐน ธมฺโม, ตํ ธมฺมํ ปวตฺเตตีติ ธมฺมจกฺกํ.}

\csromanheader{2}{7. ธมฺมจกฺกกถา}
\csromanmatikaanchor{mtk.114}
\csromantocmarkref{subsection}{3. อิทฺธิปาทวาร}{mtk.114}
\bookmark[dest={mtk.114},level=2]{3. อิทฺธิปาทวาร}
\csromanheader{2}{3. อิทฺธิปาทวาร}
\proseitem{42}{\csromanfolio{347}“อยํ ฉนฺทสมาธิปธานสงฺขารสมนฺนาคโต อิทฺธิปาโท”ติ เม ภิกฺขเว ปุพฺเพ อนนุสฺสุเตสุ ธมฺเมสุ จกฺขุํ อุทปาทิ ฯเปฯ อาโลโก อุทปาทิ. “โส โข ปนายํ ฉนฺทสมาธิปธานสงฺขารสมนฺนาคโต อิทฺธิปาโท ภาเวตพฺโพ”ติ เม ภิกฺขเว ฯเปฯ “ภาวิโต”ติ เม ภิกฺขเว ปุพฺเพ อนนุสฺสุเตสุ ธมฺเมสุ จกฺขุํ อุทปาทิ ฯเปฯ อาโลโก อุทปาทิ.}

\prose{“อยํ วีริยสมาธิ ฯเปฯ.\csromanspacer{} อยํ จิตฺตสมาธิ ฯเปฯ.\csromanspacer{} อยํ วีมํสาสมาธิปธานสงฺขารสมนฺนาคโต อิทฺธิปาโท”ติ เม ภิกฺขเว ปุพฺเพ อนนุสฺสุเตสุ ธมฺเมสุ จกฺขุํ อุทปาทิ ฯเปฯ อาโลโก อุทปาทิ. “โส โข ปนายํ วีมํสาสมาธิปธานสงฺขารสมนฺนาคโต อิทฺธิปาโท ภาเวตพฺโพ”ติ เม ภิกฺขเว ฯเปฯ “ภาวิโต”ติ เม ภิกฺขเว ปุพฺเพ อนนุสฺสุเตสุ ธมฺเมสุ จกฺขุํ อุทปาทิ ฯเปฯ อาโลโก อุทปาทิ.}

\prose{“อยํ ฉนฺทสมาธิปธานสงฺขารสมนฺนาคโต อิทฺธิปาโท”ติ ปุพฺเพ อนนุสฺสุเตสุ ธมฺเมสุ จกฺขุํ อุทปาทิ ฯเปฯ อาโลโก อุทปาทิ ฯเปฯ “โส โข ปนายํ ฉนฺทสมาธิปธานสงฺขารสมนฺนาคโต อิทฺธิปาโท ภาเวตพฺโพ”ติ ฯเปฯ “ภาวิโต”ติ ปุพฺเพ อนนุสฺสุเตสุ ธมฺเมสุ จกฺขุํ อุทปาทิ, ญาณํ อุทปาทิ, ปญฺญา อุทปาทิ, วิชฺชา อุทปาทิ, อาโลโก อุทปาทิ.}

\prose{\textbf{จกฺขุํ อุทปาที}ติ เกนฏฺเฐน, \textbf{ญาณํ อุทปาที}ติ เกนฏฺเฐน, \textbf{ปญฺญา อุทปาที}ติ เกนฏฺเฐน \textbf{วิชฺชา อุทปาที}ติ เกนฏฺเฐน, \textbf{อาโลโก อุทปาที}ติ เกนฏฺเฐน? \textbf{จกฺขุํ อุทปาที}ติ ทสฺสนฏฺเฐน, \textbf{ญาณํ อุทปาที}ติ ญาตฏฺเฐน, \textbf{ปญฺญา อุทปาที}ติ ปชานนฏฺเฐน, \textbf{วิชฺชา อุทปาที}ติ ปฏิเวธฏฺเฐน, \textbf{อาโลโก อุทปาที}ติ โอภาสฏฺเฐน.}

\prose{จกฺขุํ ธมฺโม, ทสฺสนฏฺโฐ อตฺโถ.\csromanspacer{} ญาณํ ธมฺโม, ญาตตฺโถ อตฺโถ.\csromanspacer{} ปญฺญา ธมฺโม, ปชานนฏฺโฐ อตฺโถ.\csromanspacer{} วิชฺชา ธมฺโม, ปฏิเวธฏฺโฐ อตฺโถ.\csromanspacer{} อาโลโก ธมฺโม, โอภาสฏฺโฐ อตฺโถ.\csromanspacer{} อิเม ปญฺจ ธมฺมา ปญฺจ อตฺถา ฉนฺทวตฺถุกา อิทฺธิปาทวตฺถุกา อิทฺธิปาทารมฺมณา อิทฺธิปาทโคจรา อิทฺธิปาทสงฺคหิตา อิทฺธิปาทปริยาปนฺนา อิทฺธิปาเท สมุทาคตา อิทฺธิปาเท ฐิตา อิทฺธิปาเท ปติฏฺฐิตา.}

\prose{\textbf{ธมฺมจกฺกนฺ}ติ เกนฏฺเฐน ธมฺมจกฺกํ? ธมฺมญฺจ ปวตฺเตติ จกฺกญฺจาติ ธมฺมจกฺกํ, จกฺกญฺจ ปวตฺเตติ ธมฺมญฺจาติ ธมฺมจกฺกํ, ธมฺเมน ปวตฺเตตีติ ธมฺมจกฺกํ, ธมฺมจริยาย \csromanfolio{348}ปวตฺเตตีติ ธมฺมจกฺกํ, ธมฺเม ฐิโต ปวตฺเตตีติ ธมฺมจกฺกํ, ธมฺเม ปติฏฺฐิโต ปวตฺเตตีติ ธมฺมจกฺกํ, ธมฺเม ปติฏฺฐาเปนฺโต ปวตฺเตตีติ ธมฺมจกฺกํ, ธมฺเม วสิปฺปตฺโต ปวตฺเตตีติ ธมฺมจกฺกํ, ธมฺเม วสิํ ปาเปนฺโต ปวตฺเตตีติ ธมฺมจกฺกํ, ธมฺเม ปารมิปฺปตฺโต ปวตฺเตตีติ ธมฺมจกฺกํ, ธมฺเม ปารมิํ ปาเปนฺโต ปวตฺเตตีติ ธมฺมจกฺกํ ฯเปฯ.\csromanspacer{} ธมฺมํ อปจายมาโน ปวตฺเตตีติ ธมฺมจกฺกํ, ธมฺมทฺธโช ปวตฺเตตีติ ธมฺมจกฺกํ, ธมฺมเกตุ ปวตฺเตตีติ ธมฺมจกฺกํ, ธมฺมาธิปเตยฺโย ปวตฺเตตีติ ธมฺมจกฺกํ.\csromanspacer{} ตํ โข ปน ธมฺมจกฺกํ อปฺปฏิวตฺติยํ สมเณน วา พฺราหฺมเณน วา เทเวน วา มาเรน วา พฺรหฺมุนา วา เกนจิ วา โลกสฺมินฺติ ธมฺมจกฺกํ.}

\prose{สทฺธินฺทฺริยํ ธมฺโม, ตํ ธมฺมํ ปวตฺเตตีติ ธมฺมจกฺกํ ฯเปฯ อมโตคธํ นิพฺพานํ ปริโยสานฏฺเฐน ธมฺโม, ตํ ธมฺมํ ปวตฺเตตีติ ธมฺมจกฺกํ.}

\prose{“อยํ วีริยสมาธิปธานสงฺขารสมนฺนาคโต อิทฺธิปาโท”ติ ปุพฺเพ อนนุสฺสุเตสุ ธมฺเมสุ จกฺขุํ อุทปาทิ ฯเปฯ อาโลโก อุทปาทิ ฯเปฯ “โส โข ปนายํ วีริยสมาธิปธานสงฺขารสมนฺนาคโต อิทฺธิปาโท ภาเวตพฺโพ”ติ ฯเปฯ “ภาวิโต”ติ ปุพฺเพ อนนุสฺสุเตสุ ธมฺเมสุ จกฺขุํ อุทปาทิ ฯเปฯ อาโลโก อุทปาทิ.}

\prose{\textbf{จกฺขุํ อุทปาที}ติ เกนฏฺเฐน ฯเปฯ อาโลโก อุทปาทีติ เกนฏฺเฐน? \textbf{จกฺขุํ อุทปาที}ติ ทสฺสนฏฺเฐน.\csromanspacer{} \textbf{อาโลโก อุทปาที}ติ โอภาสฏฺเฐน.}

\prose{จกฺขุํ ธมฺโม, ทสฺสนฏฺโฐ อตฺโถ ฯเปฯ อาโลโก ธมฺโม, โอภาสฏฺโฐ อตฺโถ.\csromanspacer{} อิเม ปญฺจ ธมฺมา ปญฺจ อตฺถา วีริยวตฺถุกา อิทฺธิปาทวตฺถุกา ฯเปฯ จิตฺตวตฺถุกา อิทฺธิปาทวตฺถุกา.\csromanspacer{} วีมํสาวตฺถุกา อิทฺธิปาทวตฺถุกา อิทฺธิปาทารมฺมณา อิทฺธิปาทโคจรา อิทฺธิปาทสงฺคหิตา อิทฺธิปาทปริยาปนฺนา อิทฺธิปาเท สมุทาคตา อิทฺธิปาเท ฐิตา อิทฺธิปาเท ปติฏฺฐิตา.}

\prose{\textbf{ธมฺมจกฺกนฺ}ติ เกนฏฺเฐน ธมฺมจกฺกํ? ธมฺมญฺจ ปวตฺเตติ จกฺกญฺจาติ ธมฺมจกฺกํ, จกฺกญฺจ ปวตฺเตติ ธมฺมญฺจาติ ธมฺมจกฺกํ, ธมฺเมน ปวตฺเตตีติ ธมฺมจกฺกํ, ธมฺมจริยาย ปวตฺเตตีติ ธมฺมจกฺกํ, ธมฺเม ฐิโต ปวตฺเตตีติ ธมฺมจกฺกํ, ธมฺเม ปติฏฺฐิโต ปวตฺเตตีติ ธมฺมจกฺกํ ฯเปฯ อมโตคธํ นิพฺพานํ ปริโยสานฏฺเฐน ธมฺโม, ตํ ธมฺมํ ปวตฺเตตีติ ธมฺมจกฺกนฺติ.}

\nitthitam{ธมฺมจกฺกกถา นิฏฺฐิตา.}
\csromanmatikaanchor{mtk.115}
\csromantocmarkref{section}{8. โลกุตฺตรกถา}{mtk.115}
\bookmark[dest={mtk.115},level=1]{8. โลกุตฺตรกถา}
\csromanheader{1}{8. โลกุตฺตรกถา}
\csromanheader{2}{8. โลกุตฺตรกถา}
\proseitem{43}{\csromanfolio{349}กตเม ธมฺมา โลกุตฺตรา? จตฺตาโร สติปฏฺฐานา จตฺตาโร สมฺมปฺปธานา จตฺตาโร อิทฺธิปาทา ปญฺจินฺทฺริยานิ ปญฺจ พลานิ สตฺต โพชฺฌงฺคา อริโย อฏฺฐงฺคิโก มคฺโค จตฺตาโร อริยมคฺคา จตฺตาริ จ สามญฺญผลานิ นิพฺพานญฺจ.\csromanspacer{} อิเม ธมฺมา โลกุตฺตรา.}

\prose{\textbf{โลกุตฺตรา}ติ เกนฏฺเฐน โลกุตฺตรา? โลกํ ตรนฺตีติ โลกุตฺตรา, โลกา อุตฺตรนฺตีติ โลกุตฺตรา, โลกโต อุตฺตรนฺตีติ โลกุตฺตรา, โลกมฺหา อุตฺตรนฺตีติ โลกุตฺตรา, โลกํ อติกฺกมนฺตีติ โลกุตฺตรา, โลกํ สมติกฺกมนฺตีติ โลกุตฺตรา.\csromanspacer{} โลกํ สมติกฺกนฺตาติ โลกุตฺตรา, โลเกน อติเรกาหิ โลกุตฺตรา, โลกนฺตํ ตรนฺตีติ โลกุตฺตรา, โลกานิสฺสรนฺตีติ โลกุตฺตรา, โลกโต นิสฺสรนฺตีติ โลกุตฺตรา, โลกมฺหา นิสฺสรนฺตีติ โลกุตฺตรา, โลกา นิสฏาติ โลกุตฺตรา, โลเกน นิสฺสฏาติ โลกุตฺตรา, โลกมฺหา นิสฺสฏาติ โลกุตฺตรา.\csromanspacer{}\csromanspacer{}โลเก น ติฏฺฐนฺตีติ โลกุตฺตรา, โลกสฺมิํ น ติฏฺฐนฺตีติ โลกุตฺตรา, โลเก น ลิมฺปนฺตีติ โลกุตฺตรา, โลเกน น ลิมฺปนฺตีติ โลกุตฺตรา, โลเก อสํลิตฺตาติ โลกุตฺตรา, โลเกน อสํลิตฺตาติ โลกุตฺตรา, โลเก อนุปลิตฺตาติ โลกุตฺตรา, โลเกน อนุปลิตฺตาติ โลกุตฺตรา.\csromanspacer{} โลเก วิปฺปมุตฺตาติ โลกุตฺตรา โลเกน วิปฺปมุตฺตาติ โลกุตฺตรา, โลกา วิปฺปมุตฺตาติ โลกุตฺตรา, โลกโต วิปฺปมุตฺตาติ โลกุตฺตรา, โลกมฺหา วิปฺปมุตฺตาติ โลกุตฺตรา.\csromanspacer{} โลเก วิสญฺญุตฺตาติ โลกุตฺตรา, โลเกน วิสญฺญุตฺตาติ โลกุตฺตรา, โลกา วิสญฺญุตฺตาติ โลกุตฺตรา, โลกสฺมิํ วิสญฺญุตฺตาติ โลกุตฺตรา, โลกโต วิสญฺญุตฺตาติ โลกุตฺตรา, โลกมฺหา วิสญฺญุตฺตาติ โลกุตฺตรา.\csromanspacer{}\csromanspacer{}โลกา สุชฺฌนฺตีติ โลกุตฺตรา, โลกโต สุชฺฌนฺตีติ โลกุตฺตรา, โลกมฺหา สุชฺฌนฺตีติ โลกุตฺตรา, โลกา วิสุชฺฌนฺตีติ โลกุตฺตรา, โลกโต วิสุชฺฌนฺตีติ โลกุตฺตรา, โลกมฺหา วิสุชฺฌนฺตีติ โลกุตฺตรา.\csromanspacer{}\csromanspacer{}โลกา วุฏฺฐหนฺตีติ\footnote{อุทฺธรนฺตีติ (ก), อุฏฺฐหนฺตีติ (สี, ฏฺฐ)} โลกุตฺตรา, โลกโต วุฏฺฐหนฺตีติ โลกุตฺตรา, โลกมฺหา วุฏฺฐหนฺตีติ โลกุตฺตรา.\csromanspacer{}\csromanspacer{}โลกา วิวฏฺฏนฺตีติ \csromanfolio{350}โลกุตฺตรา, โลกโต วิวฏฺฏนฺตีติ โลกุตฺตรา, โลกมฺหา วิวฏฺฏนฺตีติ โลกุตฺตรา.\csromanspacer{}\csromanspacer{}โลเก น สชฺชนฺตีติ โลกุตฺตรา, โลเก น คยฺหนฺตีติ โลกุตฺตรา, โลเก น พชฺฌนฺตีติ โลกุตฺตรา.\csromanspacer{}\csromanspacer{}โลกํ สมุจฺฉินฺทนฺตีติ โลกุตฺตรา, โลกํ สมุจฺฉินฺนตฺตาติ โลกุตฺตรา.\csromanspacer{}\csromanspacer{}โลกํ ปฏิปฺปสฺสมฺเภนฺตีติ โลกุตฺตรา, โลกํ ปฏิปฺปสฺสมฺภิตตฺตาติ โลกุตฺตรา.\csromanspacer{}\csromanspacer{}โลกสฺส อปถาติ โลกุตฺตรา, โลกสฺส อคตีติ โลกุตฺตรา, โลกสฺส อวิสยาติ โลกุตฺตรา, โลกสฺส อสาธารณาติ โลกุตฺตรา.\csromanspacer{}\csromanspacer{}โลกํ วมนฺตีติ โลกุตฺตรา, โลกํ น ปจฺจาวมนฺตีติ โลกุตฺตรา.\csromanspacer{} โลกํ ปชหนฺตีติ โลกุตฺตรา.\csromanspacer{} โลกํ น อุปาทิยนฺตีติ โลกุตฺตรา.\csromanspacer{}\csromanspacer{}โลกํ วิสีเนนฺตีติ โลกุตฺตรา, โลกํ น อุสฺสีเนนฺตีติ โลกุตฺตรา.\csromanspacer{} โลกํ วิธูเปนฺตีติ โลกุตฺตรา, โลกํ น สํธูเปนฺตีติ โลกุตฺตรา.\csromanspacer{} โลกํ สมติกฺกมฺม อภิภุยฺย ติฏฺฐนฺตีติ โลกุตฺตรา.}

\nitthitamruled{โลกุตฺตรกถา นิฏฺฐิตา.}
\csromanmatikaanchor{mtk.116}
\csromantocmarkref{section}{9. พลกถา}{mtk.116}
\bookmark[dest={mtk.116},level=1]{9. พลกถา}
\csromanheader{1}{9. พลกถา}
\proseitem{44}{สาวตฺถินิทานํ\footnote{สํ 3. 218 ปิฏฺเฐ ปสฺสิตพฺพา.}.\csromanspacer{}\csromanspacer{}ปญฺจิมานิ ภิกฺขเว พลานิ.\csromanspacer{} กตมานิ ปญฺจ? สทฺธาพลํ วีริยพลํ สติพลํ สมาธิพลํ ปญฺญาพลํ.\csromanspacer{} อิมานิ โข ภิกฺขเว ปญฺจ พลานิ.}

\prose{อปิ จ อฏฺฐสฏฺฐิ พลานิ.\csromanspacer{} สทฺธาพลํ, วีริยพลํ, สติพลํ, สมาธิพลํ, ปญฺญาพลํ, หิริพลํ, โอตฺตปฺปพลํ, ปฏิสงฺขานพลํ, ภาวนาพลํ, อนวชฺชพลํ, สงฺคหพลํ, ขนฺติพลํ, ปญฺญตฺติพลํ, นิชฺฌตฺติพลํ, อิสฺสริยพลํ, อธิฏฺฐานพลํ, สมถพลํ, วิปสฺสนาพลํ, ทส เสขพลานิ, ทส อเสขพลานิ, ทส ขีณาสวพลานิ, ทส อิทฺธิพลานิ, ทส ตถาคตพลานิ.}

\prose{กตมํ \textbf{สทฺธาพลํ?} อสฺสทฺธิเย น กมฺปตีติ สทฺธาพลํ, สหชาตานํ ธมฺมานํ อุปตฺถมฺภนฏฺเฐน สทฺธาพลํ, กิเลสานํ ปริยาทานฏฺเฐน สทฺธาพลํ, ปฏิเวธาทิวิโสธนฏฺเฐน สทฺธาพลํ, จิตฺตสฺส อธิฏฺฐานฏฺเฐน สทฺธาพลํ, จิตฺตสฺส โวทานฏฺเฐน สทฺธาพลํ, วิเสสาธิคมฏฺเฐน สทฺธาพลํ, อุตฺตริ-}

\csromanheader{2}{9. พลกถา}
\prose{\csromanfolio{351}ปฏิเวธฏฺเฐน สทฺธาพลํ, สจฺจาภิสมยฏฺเฐน สทฺธาพลํ, นิโรเธ ปติฏฺฐาปกฏฺเฐน สทฺธาพลํ, อิทํ สทฺธาพลํ. (1)}

\prose{กตมํ \textbf{วีริยพลํ?} โกสชฺเช น กมฺปตีติ วีริยพลํ, สหชาตานํ ธมฺมานํ อุปตฺถมฺภนฏฺเฐน วีริยพลํ, กิเลสานํ ปริยาทานฏฺเฐน วีริยพลํ, ปฏิเวธาทิวิโสธนฏฺเฐน วีริยพลํ, จิตฺตสฺส อธิฏฺฐานฏฺเฐน วีริยพลํ, จิตฺตสฺส โวทานฏฺเฐน วีริยพลํ, วิเสสาธิคมฏฺเฐน วีริยพลํ, อุตฺตริปฏิเวธฏฺเฐน วีริยพลํ, สจฺจาภิสมยฏฺเฐน วีริยพลํ, นิโรเธ ปติฏฺฐาปกฏฺเฐน วีริยพลํ, อิทํ วีริยพลํ. (2)}

\prose{กตมํ \textbf{สติพลํ?} ปมาเท น กมฺปตีติ สติพลํ, สหชาตานํ ธมฺมานํ อุปตฺถมฺภนฏฺเฐน สติพลํ ฯเปฯ นิโรเธ ปติฏฺฐาปกฏฺเฐน สติพลํ, อิทํ สติพลํ. (3)}

\prose{กตมํ \textbf{สมาธิพลํ?} อุทฺธจฺเจ น กมฺปตีติ สมาธิพลํ, สหชาตานํ ธมฺมานํ อุปตฺถมฺภนฏฺเฐน สมาธิพลํ ฯเปฯ นิโรเธ ปติฏฺฐาปกฏฺเฐน สมาธิพลํ, อิทํ สมาธิพลํ. (4)}

\prose{กตมํ \textbf{ปญฺญาพลํ?} อวิชฺชาย น กมฺปตีติ ปญฺญาพลํ, สหชาตานํ ธมฺมานํ อุปตฺถมฺภนฏฺเฐน ปญฺญาพลํ ฯเปฯ นิโรเธ ปติฏฺฐาปกฏฺเฐน ปญฺญาพลํ, อิทํ ปญฺญาพลํ. (5)}

\prose{กตมํ \textbf{หิริพลํ?} เนกฺขมฺเมน กามจฺฉนฺทํ หิรียตีติ\footnote{หิริยตีติ (สฺยา, ก)} หิริพลํ, อพฺยาปาเทน พฺยาปาทํ หิรียตีติ หิริพลํ, อาโลกสญฺญาย ถินมิทฺธํ หิรียตีติ หิริพลํ, อวิกฺเขเปน อุทฺธจฺจํ หิรียตีติ หิริพลํ, ธมฺมววตฺถาเนน วิจิกิจฺฉํ หิรียตีติ หิริพลํ, ญาเณน อวิชฺชํ หิรียตีติ หิริพลํ, ปาโมชฺเชน อรติํ หิรียตีติ หิริพลํ, ปฐเมน ฌาเนน นีวรเณ หิรียตีติ หิริพลํ ฯเปฯ อรหตฺตมคฺเคน สพฺพกิเลเส หิรียตีติ หิริพลํ, อิทํ หิริพลํ. (6)}

\prose{กตมํ \textbf{โอตฺตปฺปพลํ?} เนกฺขมฺเมน กามจฺฉนฺทํ โอตฺตปฺปตีติ โอตฺตปฺปพลํ, อพฺยาปาเทน พฺยาปาทํ โอตฺตปฺปตีติ โอตฺตปฺปพลํ, อาโลกสญฺญาย ถินมิทฺธํ โอตฺตปฺปตีติ โอตฺตปฺปพลํ, อวิกฺเขเปน อุทฺธจฺจํ โอตฺตปฺปตีติ โอตฺตปฺปพลํ, ธมฺมววตฺถาเนน วิจิกิจฺฉํ โอตฺตปฺปตีติ โอตฺตปฺปพลํ, ญาเณน อวิชฺชํ โอตฺตปฺปตีติ}

\prose{\csromanfolio{352}โอตฺตปฺปพลํ, ปาโมชฺเชน อรติํ โอตฺตปฺปตีติ โอตฺตปฺปพลํ, ปฐเมน ฌาเนน นีวรเณ โอตฺตปฺปตีติ โอตฺตปฺปพลํ ฯเปฯ อรหตฺตมคฺเคน สพฺพกิเลเส โอตฺตปฺปตีติ โอตฺตปฺปพลํ, อิทํ โอตฺตปฺปพลํ. (7)}

\prose{กตมํ \textbf{ปฏิสงฺขานพลํ?} เนกฺขมฺเมน กามจฺฉนฺทํ ปฏิสงฺขาตีติ ปฏิสงฺขานพลํ.\csromanspacer{} อพฺยาปาเทน พฺยาปาทํ ปฏิสงฺขาตีติ ปฏิสงฺขานพลํ, อาโลกสญฺญาย ถินมิทฺธํ ปฏิสงฺขาตีติ ปฏิสงฺขานพลํ, อวิกฺเขเปน อุทฺธจฺจํ ปฏิสงฺขาตีติ ปฏิสงฺขานพลํ, ธมฺมววตฺถาเนน วิจิกิจฺฉํ ปฏิสงฺขาตีติ ปฏิสงฺขานพลํ, ญาเณน อวิชฺชํ ปฏิสงฺขาตีติ ปฏิสงฺขานพลํ, ปาโมชฺเชน อรติํ ปฏิสงฺขาตีติ ปฏิสงฺขานพลํ, ปฐเมน ฌาเนน นีวรเณ ปฏิสงฺขาตีติ ปฏิสงฺขานพลํ ฯเปฯ อรหตฺตมคฺเคน สพฺพกิเลเส ปฏิสงฺขาตีติ ปฏิสงฺขานพลํ, อิทํ ปฏิสงฺขานพลํ. (8)}

\prose{กตมํ \textbf{ภาวนาพลํ?} กามจฺฉนฺทํ ปชหนฺโต เนกฺขมฺมํ ภาเวตีติ ภาวนาพลํ, พฺยาปาทํ ปชหนฺโต อพฺยาปาทํ ภาเวตีติ ภาวนาพลํ, ถินมิทฺธํ ปชหนฺโต อาโลกสญฺญํ ภาเวตีติ ภาวนาพลํ, อุทฺธจฺจํ ปชหนฺโต อวิกฺเขปํ ภาเวตีติ ภาวนาพลํ, วิจิกิจฺฉํ ปชหนฺโต ธมฺมววตฺถานํ ภาเวตีติ ภาวนาพลํ, อวิชฺชํ ปชหนฺโต ญาณํ ภาเวตีติ ภาวนาพลํ, อรติํ ปชหนฺโต ปาโมชฺชํ ภาเวตีติ ภาวนาพลํ, นีวรเณ ปชหนฺโต ปฐมํ ฌานํ ภาเวตีติ ภาวนาพลํ ฯเปฯ สพฺพกิเลเส ปชหนฺโต อรหตฺตมคฺคํ ภาเวตีติ ภาวนาพลํ, อิทํ ภาวนาพลํ. (9)}

\prose{กตมํ \textbf{อนวชฺชพลํ?} กามจฺฉนฺทสฺส ปหีนตฺตา เนกฺขมฺเม นตฺถิ กิญฺจิ วชฺชนฺติ อนวชฺชพลํ, พฺยาปาทสฺส ปหีนตฺตา อพฺยาปาเท นตฺถิ กิญฺจิ วชฺชนฺติ อนวชฺชพลํ, ถินมิทฺธสฺส ปหีนตฺตา อาโลกสญฺญาย นตฺถิ กิญฺจิ วชฺชนฺติ อนวชฺชพลํ, อุทฺธจฺจสฺส ปหีนตฺตา อวิกฺเขเป นตฺถิ กิญฺจิ วชฺชนฺติ อนวชฺชพลํ, วิจิกิจฺฉาย ปหีนตฺตา ธมฺมววตฺถาเน นตฺถิ กิญฺจิ วชฺชนฺติ อนวชฺชพลํ, อวิชฺชาย ปหีนตฺตา ญาเณ นตฺถิ กิญฺจิ วชฺชนฺติ อนวชฺชพลํ, อรติยา ปหีนตฺตา ปาโมชฺเช นตฺถิ กิญฺจิ วชฺชนฺติ อนวชฺชพลํ, นีวรณานํ ปหีนตฺตา ปฐมชฺฌาเน นตฺถิ กิญฺจิ วชฺชนฺติ อนวชฺชพลํ ฯเปฯ สพฺพกิเลสานํ ปหีนตฺตา อรหตฺตมคฺเค นตฺถิ กิญฺจิ วชฺชนฺติ อนวชฺชพลํ, อิทํ อนวชฺชพลํ. (10)}

\csromanheader{2}{9. พลกถา}
\prose{\csromanfolio{353}กตมํ \textbf{สงฺคหพลํ?} กามจฺฉนฺทํ ปชหนฺโต เนกฺขมฺมวเสน จิตฺตํ สงฺคณฺหาตีติ สงฺคหพลํ, พฺยาปาทํ ปชหนฺโต อพฺยาปาทวเสน จิตฺตํ สงฺคณฺหาตีติ สงฺคหพลํ, ถินมิทฺธํ ปชหนฺโต อาโลกสญฺญาวเสน จิตฺตํ สงฺคณฺหาตีติ สงฺคหพลํ ฯเปฯ สพฺพกิเลเส ปชหนฺโต อรหตฺตมคฺควเสน จิตฺตํ สงฺคณฺหาตีติ สงฺคหพลํ, อิทํ สงฺคหพลํ. (11)}

\prose{กตมํ \textbf{ขนฺติพลํ?} กามจฺฉนฺทสฺส ปหีนตฺตา เนกฺขมฺมํ ขมตีติ ขนฺติพลํ, พฺยาปาทสฺส ปหีนตฺตา อพฺยาปาโท ขมตีติ ขนฺติพลํ, ถินมิทฺธสฺส ปหีนตฺตา อาโลกสญฺญา ขมตีติ ขนฺติพลํ, อุทฺธจฺจสฺส ปหีนตฺตา อวิกฺเขโป ขมตีติ ขนฺติพลํ, วิจิกิจฺฉาย ปหีนตฺตา ธมฺมววตฺถานํ ขมตีติ ขนฺติพลํ, อวิชฺชาย ปหีนตฺตา ญาณํ ขมตีติ ขนฺติพลํ, อรติยา ปหีนตฺตา ปาโมชฺชํ ขมตีติ ขนฺติพลํ, นีวรณานํ ปหีนตฺตา ปฐมํ ฌานํ ขมตีติ ขนฺติพลํ ฯเปฯ สพฺพกิเลสานํ ปหีนตฺตา อรหตฺตมคฺโค ขมตีติ ขนฺติพลํ, อิทํ ขนฺติพลํ. (12)}

\prose{กตมํ \textbf{ปญฺญตฺติพลํ?} กามจฺฉนฺทํ ปชหนฺโต เนกฺขมฺมวเสน จิตฺตํ ปญฺญาเปตีติ ปญฺญตฺติพลํ, พฺยาปาทํ ปชหนฺโต อพฺยาปาทวเสน จิตฺตํ ปญฺญาเปตีติ ปญฺญตฺติพลํ, ถินมิทฺธํ ปชหนฺโต อาโลกสญฺญาวเสน จิตฺตํ ปญฺญาเปตีติ ปญฺญตฺติพลํ ฯเปฯ สพฺพกิเลเส ปชหนฺโต อรหตฺตมคฺควเสน จิตฺตํ ปญฺญาเปตีติ ปญฺญตฺติพลํ, อิทํ ปญฺญตฺติพลํ. (13)}

\prose{กตมํ \textbf{นิชฺฌตฺติพลํ?} กามจฺฉนฺทํ ปชหนฺโต เนกฺขมฺมวเสน จิตฺตํ นิชฺฌาเปตีติ นิชฺฌตฺติพลํ, พฺยาปาทํ ปชหนฺโต อพฺยาปาทวเสน จิตฺตํ นิชฺฌาเปตีติ นิชฺฌตฺติพลํ, ถินมิทฺธํ ปชหนฺโต อาโลกสญฺญาวเสน จิตฺตํ นิชฺฌาเปตีติ นิชฺฌตฺติพลํ ฯเปฯ สพฺพกิเลเส ปชหนฺโต อรหตฺตมคฺควเสน จิตฺตํ นิชฺฌาเปตีติ นิชฺฌตฺติพลํ, อิทํ นิชฺฌตฺติพลํ. (14)}

\prose{กตมํ \textbf{อิสฺสริยพลํ?} กามจฺฉนฺทํ ปชหนฺโต เนกฺขมฺมวเสน จิตฺตํ วสํ วตฺเตตีติ อิสฺสริยพลํ, พฺยาปาทํ ปชหนฺโต อพฺยาปาทวเสน จิตฺตํ วสํ วตฺเตตีติ อิสฺสริยพลํ, ถินมิทฺธํ ปชหนฺโต อาโลกสญฺญาวเสน จิตฺตํ วสํ วตฺเตตีติ อิสฺสริยพลํ ฯเปฯ สพฺพกิเลเส ปชหนฺโต อรหตฺตมคฺควเสน จิตฺตํ วสํ วตฺเตตีติ อิสฺสริยพลํ, อิทํ อิสฺสริยพลํ. (15)}

\prose{กตมํ \textbf{อธิฏฺฐานพลํ?} กามจฺฉนฺทํ ปชหนฺโต เนกฺขมฺมวเสน จิตฺตํ อธิฏฺฐาตีติ อธิฏฺฐานพลํ, พฺยาปาทํ ปชหนฺโต อพฺยาปาทวเสน จิตฺตํ}

\prose{\csromanfolio{354}อธิฏฺฐาตีติ อธิฏฺฐานพลํ, ถินมิทฺธํ ปชหนฺโต อาโลกสญฺญาวเสน จิตฺตํ อธิฏฺฐาตีติ อธิฏฺฐานพลํ ฯเปฯ สพฺพกิเลเส ปชหนฺโต อรหตฺตมคฺควเสน จิตฺตํ อธิฏฺฐาตีติ อธิฏฺฐานพลํ, อิทํ อธิฏฺฐานพลํ. (16)}

\prose{กตมํ \textbf{สมถพลํ?} เนกฺขมฺมวเสน จิตฺตสฺส เอกคฺคตา อวิกฺเขโป สมถพลํ, อพฺยาปาทวเสน จิตฺตสฺส เอกคฺคตา อวิกฺเขโป สมถพลํ.\csromanspacer{} อาโลกสญฺญาวเสน จิตฺตสฺส เอกคฺคตา อวิกฺเขโป สมถพลํ ฯเปฯ ปฏินิสฺสคฺคานุปสฺสี อสฺสาสวเสน จิตฺตสฺส เอกคฺคตา อวิกฺเขโป สมถพลํ, ปฏินิสฺสคฺคานุปสฺสี ปสฺสาสวเสน จิตฺตสฺส เอกคฺคตา อวิกฺเขโป สมถพลํ.}

\prose{\textbf{สมถพลนฺ}ติ เกนฏฺเฐน \textbf{สมถพลํ?} ปฐเมน ฌาเนน นีวรเณ น กมฺปตีติ สมถพลํ, ทุติเยน ฌาเนน วิตกฺกวิจาเร น กมฺปตีติ สมถพลํ, ตติยน ฌาเนน ปีติยา น กมฺปตีติ สมถพลํ, จตุตฺเถน ฌาเนน สุขทุกฺเข น กมฺปตีติ สมถพลํ, อากาสานญฺจายตนสมาปตฺติยา รูปสญฺญาย ปฏิฆสญฺญาย นานตฺตสญฺญาย น กมฺปตีติ สมถพลํ, วิญฺญาณญฺจายตนสมาปตฺติยา อากาสานญฺจายตนสญฺญาย น กมฺปตีติ สมถพลํ, อากิญฺจญฺญายตนสมาปตฺติยา วิญฺญาณญฺจายตนสญฺญาย น กมฺปตีติ สมถพลํ, เนวสญฺญานาสญฺญายตนสมาปตฺติยา อากิญฺจญฺญายตนสญฺญาย น กมฺปตีติ สมถพลํ.\csromanspacer{} อุทฺธจฺเจ จ อุทฺธจฺจสหคตกิเลเส จ ขนฺเธ จ น กมฺปติ น จลติ น เวธตีติ สมถพลํ, อิทํ สมถพลํ. (17)}

\prose{กตมํ \textbf{วิปสฺสนาพลํ?} อนิจฺจานุปสฺสนา วิปสฺสนาพลํ, ทุกฺขานุปสฺสนา วิปสฺสนาพลํ ฯเปฯ ปฏินิสฺสคฺคานุปสฺสนา วิปสฺสนาพลํ.\csromanspacer{} รูเป อนิจฺจานุปสฺสนา วิปสฺสนาพลํ, รูเป ทุกฺขานุปสฺสนา วิปสฺสนาพลํ ฯเปฯ รูเป ปฏินิสฺสคฺคานุปสฺสนา วิปสฺสนาพลํ, เวทนาย ฯเปฯ สญฺญาย.\csromanspacer{} สงฺขาเรสุ.\csromanspacer{} วิญฺญาเณ จกฺขุสฺมิํ ฯเปฯ ชรามรเณ อนิจฺจานุปสฺสนา วิปสฺสนาพลํ, ชรามรเณ ทุกฺขานุปสฺสนา วิปสฺสนาพลํ ฯเปฯ ชรามรเณ ปฏินิสฺสคฺคานุปสฺสนา วิปสฺสนาพลํ.\csromanspacer{}\csromanspacer{}\textbf{วิปสฺสนาพลนฺ}ติ เกนฏฺเฐน \textbf{วิปสฺสนาพลํ?} อนิจฺจานุปสฺสนาย นิจฺจสญฺญาย น กมฺปตีติ วิปสฺสนาพลํ, ทุกฺขานุปสฺสนาย สุขสญฺญาย น กมฺปตีติ วิปสฺสนาพลํ, อนตฺตานุปสฺสนาย อตฺตสญฺญาย น กมฺปตีติ วิปสฺสนาพลํ, นิพฺพิทานุปสฺสนาย นนฺทิยา น กมฺปตีติ วิปสฺสนาพลํ, วิราคานุปสฺสนาย ราเค น กมฺปตีติ วิปสฺสนาพลํ,}

\csromanheader{2}{9. พลกถา}
\prose{\csromanfolio{355}นิโรธานุปสฺสนาย สมุทเย น กมฺปตีติ วิปสฺสนาพลํ, ปฏินิสฺสคฺคานุปสฺสนาย อาทาเน น กมฺปตีติ วิปสฺสนาพลํ, อวิชฺชาย จ อวิชฺชาสหคตกิเลเส จ ขนฺเธ จ น กมฺปติ น จลติ น เวธตีติ วิปสฺสนาพลํ, อิทํ วิปสฺสนาพลํ. (18)}

\prose{กตมานิ ทส เสขพลานิ, ทส อเสขพลานิ\textbf{?} สมฺมาทิฏฺฐิํ\footnote{สมฺมาทิฏฺฐิ (ก) เอวมีทิเสสุ นวสุ ปเทสุ ทุติยนฺตวจเนน. ที 3. 225 ปิฏฺเฐ ปสฺสิตพฺพา.} สิกฺขตีติ เสขพลํ, ตตฺถ สิกฺขิตตฺตา อเสขพลํ.\csromanspacer{} สมฺมาสงฺกปฺปํ สิกฺขตีติ เสขพลํ, ตตฺถ สิกฺขิตตฺตา อเสขพลํ.\csromanspacer{} สมฺมาวาจํ ฯเปฯ.\csromanspacer{} สมฺมากมฺมนฺตํ.\csromanspacer{} สมฺมา-อาชีวํ.\csromanspacer{} สมฺมาวายามํ.\csromanspacer{} สมฺมาสติํ.\csromanspacer{} สมฺมาสมาธิํ.\csromanspacer{} สมฺมาญาณํ ฯเปฯ.\csromanspacer{} สมฺมาวิมุตฺติํ สสิกฺขตีติ เสขพลํ, ตตฺถ สิกฺขิตตฺตา อเสขพลํ.\csromanspacer{} อิมานิ ทส เสขพลานิ, ทส อเสขพลานิ. (\textbf{2}0, 38)}

\prose{กตมานิ \textbf{ทส ขีณาสวพลานิ}\footnote{อํ 3. 397 ปิฏฺเฐ ปสฺสิตพฺพา.}\textbf{?} อิธ ขีณาสวสฺส ภิกฺขุโน อนิจฺจโต สพฺเพ สงฺขารา ยถาภูตํ สมฺมปฺปญฺญาย สุทิฏฺฐา โหนฺติ.\csromanspacer{} ยมฺปิ ขีณาสวสฺส ภิกฺขุโน อนิจฺจโต สพฺเพ สงฺขารา ยถาภูตํ สมฺมปฺปญฺญาย สุทิฏฺฐา โหนฺติ, อิทมฺปิ ขีณาสวสฺส ภิกฺขุโน พลํ โหติ.\csromanspacer{} ยํ พลํ อาคมฺม ขีณาสโว ภิกฺขุ อาสวานํ ขยํ ปฏิชานาติ “ขีณา เม อาสวา”ติ. (1)}

\prose{ปุน จปรํ ขีณาสวสฺส ภิกฺขุโน องฺคารกาสูปมา กามา ยถาภูตํ สมฺมปฺปญฺญาย สุทิฏฺฐา โหนฺติ.\csromanspacer{} ยมฺปิ ขีณาสวสฺส ภิกฺขุโน องฺคารกาสูปมา กามา ยถาภูตํ สมฺมปฺปญฺญาย สุทิฏฺฐา โหนฺติ, อิทมฺปิ ขีณาสวสฺส ภิกฺขุโน พลํ โหติ.\csromanspacer{} ยํ พลํ อาคมฺม ขีณาสโว ภิกฺขุ อาสวานํ ขยํ ปฏิชานาติ “ขีณา เม อาสวา”ติ. (\textbf{2})}

\prose{ปุน จปรํ ขีณาสวสฺส ภิกฺขุโน วิเวกนินฺนํ จิตฺตํ โหติ วิเวกโปณ วิเวกปพฺภารํ วิเวกฏฺฐํ เนกฺขมฺมาภิรตํ พฺยนฺตีภูตํ\footnote{พฺยนฺติภูตํ (ก)} สพฺพโส อาสวฏฺฐานิเยหิ ธมฺเมหิ.\csromanspacer{} ยมฺปิ ขีณาสวสฺส ภิกฺขุโน วิเวกนินฺนํ จิตฺตํ โหติ วิเวกโปณํ วิเวกปพฺภารํ วิเวกฏฺฐํ เนกฺขมฺมาภิรตํ พฺยนฺตีภูตํ สพฺพโส อาสวฏฺฐานิเยหิ ธมฺเมหิ, อิทมฺปิ ขีณาสวสฺส ภิกฺขุโน พลํ โหติ.\csromanspacer{} ยํ พลํ อาคมฺม ขีณาสโว ภิกฺขุ อาสวานํ ขยํ ปฏิชานาติ “ขีณา เม อาสวา”ติ. (3)}

\prose{\csromanfolio{356}ปุน จปรํ ขีณาสวสฺส ภิกฺขุโน จตฺตาโร สติปฏฺฐานา ภาวิตา โหนฺติ สุภาวิตา.\csromanspacer{} ยมฺปิ ขีณาสวสฺส ภิกฺขุโน จตฺตาโร สติปฏฺฐานา ภาวิตา โหนฺติ สุภาวิตา, อิทมฺปิ ขีณาสวสฺส ภิกฺขุโน พลํ โหติ.\csromanspacer{} ยํ พลํ อาคมฺม ขีณาสโว ภิกฺขุ อาสวานํ ขยํ ปฏิชานาติ “ขีณา เม อาสวา”ติ. (4)}

\prose{ปุน จปรํ ขีณาสวสฺส ภิกฺขุโน จตฺตาโร สมฺมปฺปธานา ภาวิตา โหนฺติ สุภาวิตา ฯเปฯ.\csromanspacer{} จตฺตาโร อิทฺธิปาทา ภาวิตา โหนฺติ สุภาวิตา.\csromanspacer{} ปญฺจินฺทฺริยานิ ภาวิตานิ โหนฺติ สุภาวิตานิ.\csromanspacer{} ปญฺจ พลานิ ภาวิตานิ โหนฺติ สุภาวิตานิ.\csromanspacer{} สตฺต โพชฺฌงฺคา ภาวิตา โหนฺติ สุภาวิตา ฯเปฯ.\csromanspacer{} อริโย อฏฺฐงฺคิโก มคฺโค ภาวิโต โหติ สุภาวิโต.\csromanspacer{} ยมฺปิ ขีณาสวสฺส ภิกฺขุโน อฏฺฐงฺคิโก มคฺโค ภาวิโต โหติ สุภาวิโต, อิทมฺปิ ขีณาสวสฺส ภิกฺขุโน พลํ โหติ.\csromanspacer{} ยํ พลํ อาคมฺม ขีณาสโว ภิกฺขุ อาสวานํ ขยํ ปฏิชานาติ “ขีณา เม อาสวา”ติ.\csromanspacer{} อิมานิ ทส ขีณาสวพลานิ. (5-10, 48)}

\prose{กตมานิ \textbf{ทส อิทฺธิพลานิ?} อธิฏฺฐานา อิทฺธิ, วิกุพฺพนา อิทฺธิ, มโนมยา อิทฺธิ, ญาณวิปฺผารา อิทฺธิ, สมาธิวิปฺผารา อิทฺธิ, อริยา อิทฺธิ, กมฺมวิปากชา อิทฺธิ, ปุญฺญวโต อิทฺธิ, วิชฺชามยา อิทฺธิ, ตตฺถ ตตฺถ สมฺมา ปโยคปฺปจฺจยา อิชฺฌนฏฺเฐน อิทฺธิ.\csromanspacer{} อิมานิ ทส อิทฺธิพลานิ. (10, 58)}

\prose{กตมานิ \textbf{ทส ตถาคตพลานิ?} อิธ ตถาคโต ฐานญฺจ ฐานโต อฏฺฐานญฺจ อฏฺฐานโต ยถาภูตํ ปชานาติ, ยมฺปิ ตถาคโต ฐานญฺจ ฐานโต อฏฺฐานญฺจ อฏฺฐานโต ยถาภูตํ ปชานาติ, อิทมฺปิ ตถาคตสฺส ตถาคตพลํ โหติ, ยํ พลํ อาคมฺม ตถาคโต อาสภํ ฐานํ ปฏิชานาติ, ปริสาสุ สีหนาทํ นทติ, พฺรหฺมจกฺกํ ปวตฺเตติ. (1)}

\prose{ปุน จปรํ ตถาคโต อตีตานาคตปจฺจุปฺปนฺนานํ กมฺมสมาทานานํ ฐานโส เหตุโส วิปากํ ยถาภูตํ ปชานาติ, ยมฺปิ ตถาคโต อตีตานาคตปจฺจุปฺปนฺนานํ กมฺมสมาทานานํ ฐานโส เหตุโส วิปากํ ยถาภูตํ ปชานาติ, อิทมฺปิ ตถาคตสฺส ตถาคตพลํ โหติ.\csromanspacer{} ยํ พลํ อาคมฺม ตถาคโต อาสภํ ฐานํ ปฏิชานาติ, ปริสาสุ สีหนาทํ นทติ, พฺรหฺมจกฺกํ ปวตฺเตติ. (2)}

\csromanheader{2}{9. พลกถา}
\prose{\csromanfolio{357}ปุน จปรํ ตถาคโต สพฺพตฺถคามินิํ ปฏิปทํ\footnote{สพฺพตฺถคามินีปฏิปทํ (สฺยา) ม 1. 100 ปิฏฺเฐ ปสฺสิตพฺพา.} ยถาภูตํ ปชานาติ, ยมฺปิ ตถาคโต สพฺพตฺถคามินิํ ปฏิปทํ ยถาภูตํ ปชานาติ, อิทมฺปิ ตถาคตสฺส ตถาคตพลํ โหติ, ยํ พลํ อาคมฺม ตถาคโต อาสภํ ฐานํ ปฏิชานาติ, ปริสาสุ สีหนาทํ นทติ, พฺรหฺมจกฺกํ ปวตฺเตติ. (3)}

\prose{ปุน จปรํ ตถาคโต อเนกธาตุํ นานาธาตุํ โลกํ ยถาภูตํ ปชานาติ.\csromanspacer{} ยมฺปิ ตถาคโต อเนกธาตุํ นานาธาตุํ โลกํ ยถาภูตํ ปชานาติ, อิทมฺปิ ตถาคตสฺส ฯเปฯ. (4)}

\prose{ปุน จปรํ ตถาคโต สตฺตานํ นานาธิมุตฺติกตํ ยถาภูตํ ปชานาติ, ยมฺปิ ตถาคโต สตฺตานํ นานาธิมุตฺติกตํ ยถาภูตํ ปชานาติ, อิทมฺปิ ตถาคตสฺส ฯเปฯ. (5)}

\prose{ปุน จปรํ ตถาคโต ปรสตฺตานํ ปรปุคฺคลานํ อินฺทฺริยปโรปริยตฺตํ ยถาภูตํ ปชานาติ, ยมฺปิ ตถาคโต ปรสตฺตานํ ปรปุคฺคลานํ อินฺทฺริยปโรปริยตฺตํ ยถาภูตํ ปชานาติ, อิทมฺปิ ตถาคตสฺส ฯเปฯ. (6)}

\prose{ปุน จปรํ ตถาคโต ฌานวิโมกฺขสมาธิสมาปตฺตีนํ สํกิเลสํ โวทานํ วุฏฺฐานํ ยถาภูตํ ปชานาติ, ยมฺปิ ตถาคโต ฌานวิโมกฺขสมาธิสมาปตฺตีนํ สํกิเลสํ โวทานํ วุฏฺฐานํ ยถาภูตํ ปชานาติ, อิทมฺปิ ตถาคตสฺส ฯเปฯ. (7)}

\prose{ปุน จปรํ ตถาคโต อเนกวิหิตํ ปุพฺเพนิวาสํ อนุสฺสรติ, เสยฺยถิทํ, เอกมฺปิ ชาติํ เทฺวปิ ชาติโย ฯเปฯ อิติ สาการํ ส-อุทฺเทสํ อเนกวิหิตํ ปุพฺเพนิวาสํ อนุสฺสรติ, ยมฺปิ ตถาคโต อเนกวิหิตํ ปุพฺเพนิวาสํ อนุสฺสรติ.\csromanspacer{} เสยฺยถิทํ, เอกมฺปิ ชาติํ เทฺวปิ ชาติโย ฯเปฯ อิทมฺปิ ตถาคตสฺส ฯเปฯ. (8)}

\prose{ปุน จปรํ ตถาคโต ทิพฺเพน จกฺขุนา วิสุทฺเธน อติกฺกนฺตมานุสเกน สตฺเต ปสฺสติ จวมาเน อุปปชฺชมาเน ฯเปฯ ยมฺปิ ตถาคโต ทิพฺเพน จกฺขุนา วิสุทฺเธน อติกฺกนฺตมานุสเกน สตฺเต ปสฺสติ จวมาเน อุปปชฺชมาเน ฯเปฯ อิทมฺปิ ตถาคตสฺส ฯเปฯ. (9)}

\prose{\csromanfolio{358}ปุน จปรํ ตถาคโต อาสวานํ ขยา อนาสวํ เจโตวิมุตฺติํ ปญฺญาวิมุตฺติํ ทิฏฺเฐว ธมฺเม สยํ อภิญฺญา สจฺฉิกตฺวา อุปสมฺปชฺช วิหรติ, ยมฺปิ ตถาคโต อาสวานํ ขยา อนาสวํ เจโตวิมุตฺติํ ปญฺญาวิมุตฺติํ ทิฏฺเฐว ธมฺเม สยํ อภิญฺญา สจฺฉิกตฺวา อุปสมฺปชฺช วิหรติ, อิทมฺปิ ตถาคตสฺส ตถาคตพลํ โหติ, ยํ พลํ อาคมฺม ตถาคโต อาสภํ ฐานํ ปฏิชานาติ, ปริสาสุ สีหนาทํ นทติ, พฺรหฺมจกฺกํ ปวตฺเตติ.\csromanspacer{} อิมานิ ทส ตถาคตพลานิ. (10, 68)}

\proseitem{45}{เกนฏฺเฐน สทฺธาพลํ.\csromanspacer{} เกนฏฺเฐน วีริยพลํ.\csromanspacer{} เกนฏฺเฐน สติพลํ.\csromanspacer{} เกนฏฺเฐน สมาธิพลํ.\csromanspacer{} เกนฏฺเฐน ปญฺญาพลํ.\csromanspacer{} เกนฏฺเฐน หิริพลํ.\csromanspacer{} เกนฏฺเฐน โอตฺตปฺปพลํ.\csromanspacer{} เกนฏฺเฐน ปฏิสงฺขานพลํ ฯเปฯ.\csromanspacer{} เกนฏฺเฐน ตถาคตพลํ?}

\prose{อสฺสทฺธิเย อกมฺปิยฏฺเฐน สทฺธาพลํ.\csromanspacer{} โกสชฺเช อกมฺปิยฏฺเฐน วีริยพลํ.\csromanspacer{} ปมาเท อกมฺปิยฏฺเฐน สติพลํ.\csromanspacer{} อุทฺธจฺเจ อกมฺปิยฏฺเฐน สมาธิพลํ.\csromanspacer{} อวิชฺชาย อกมฺปิยฏฺเฐน ปญฺญาพลํ.\csromanspacer{} หิรียติ ปาปเก อกุสเล ธมฺเมติ หิริพลํ.\csromanspacer{} โอตฺตปฺปติ ปาปเก อกุสเล ธมฺเมติ โอตฺตปฺปพลํ.\csromanspacer{} ญาเณน กิเลเส ปฏิสงฺขาตีติ ปฏิสงฺขานพลํ.\csromanspacer{} ตตฺถ ชาตา ธมฺมา เอกรสา โหนฺตีติ ภาวนาพลํ, ตตฺถ นตฺถิ กิญฺจิ วชฺชนฺติ อนวชฺชพลํ.\csromanspacer{} เตน จิตฺตํ สงฺคณฺหาตีติ สงฺคหพลํ.\csromanspacer{} ตํ ตสฺส ขมตีติ\footnote{ตํ ขมตีติ (ก)} ขนฺติพลํ.\csromanspacer{} เตน จิตฺตํ ปญฺญาเปตีติ ปญฺญตฺติพลํ, เตน จิตฺตํ นิชฺฌาเปตีติ นิชฺฌตฺติพลํ, เตน จิตฺตํ วสํ วตฺเตตีติ อิสฺสริยพลํ, เตน จิตฺตํ อธิฏฺฐาตีติ อธิฏฺฐานพลํ, เตน จิตฺตํ เอกคฺคนฺติ สมถพลํ, ตตฺถ ชาเต ธมฺเม อนุปสฺสตีติ วิปสฺสนาพลํ, ตตฺถ สิกฺขตีติ เสขพลํ, ตตฺถ สิกฺขิตตฺตา อเสขพลํ, เตน อาสวา ขีณาติ ขีณาสวพลํ.\csromanspacer{} ตสฺส อิชฺฌตีติ อิทฺธิพลํ.\csromanspacer{} อปฺปเมยฺยฏฺเฐน ตถาคตพลนฺติ.}

\nitthitamruled{พลกถา นิฏฺฐิตา.}
\csromanmatikaanchor{mtk.117}
\csromantocmarkref{section}{10. สุญฺญกถา}{mtk.117}
\bookmark[dest={mtk.117},level=1]{10. สุญฺญกถา}
\csromanheader{1}{10. สุญฺญกถา}
\proseitem{46}{เอวํ เม สุตํ\csromandash{}เอกํ สมยํ ภควา สาวตฺถิยํ วิหรติ เชตวเน อนาถปิณฺฑิกสฺส อาราเม.\csromanspacer{} อถ โข อายสฺมา อานนฺโท เยน}

\csromanheader{2}{10. สุญฺญกถา}
\prose{\csromanfolio{359}ภควา เตนุปสงฺกมิ, อุปสงฺกมิตฺวา ภควนฺตํ อภิวาเทตฺวา เอกมนฺตํ นิสีทิ, เอกมนฺตํ นิสินฺโน โข อายสฺมา อานนฺโท ภควนฺตํ เอตทโวจ\csromandash{}}

\prose{“สุญฺโญ โลโก สุญฺโญ โลโก”ติ ภนฺเต วุจฺจติ, กิตฺตาวตา นุ โข ภนฺเต “สุญฺโญ โลโก”ติ วุจฺจตีติ? ยสฺมา โข อานนฺท สุญฺญํ อตฺเตน วา อตฺตนิเยน วา, ตสฺมา “สุญฺโญ โลโก”ติ วุจฺจติ.\csromanspacer{} กิญฺจานนฺท\footnote{กิญฺจ อานนฺท (สํ 2. 280 ปิฏฺเฐ)} สุญฺญํ อตฺเตน วา อตฺตนิเยน วา? จกฺขุ โข อานนฺท สุญฺญํ อตฺเตน วา อตฺตนิเยน วา, รูปา สุญฺญา อตฺเตน วา อตฺตนิเยน วา, จกฺขุวิญฺญาณํ สุญฺญํ อตฺเตน วา อตฺตนิเยน วา, จกฺขุสมฺผสฺโส สุญฺโญ อตฺเตน วา อตฺตนิเยน วา, ยมฺปิทํ\footnote{ยทิทํ (ก) สํ 2. 280 ปิฏฺเฐ ปสฺสิตพฺพา. 3. สุญฺญํ สุญฺญํ (สฺยา)} จกฺขุสมฺผสฺสปจฺจยา อุปฺปชฺชติ เวทยิตํ สุขํ วา ทุกฺขํ วา อทุกฺขมสุขํ วา, ตมฺปิ สุญฺญํ อตฺเตน วา อตฺตนิเยน วา.}

\prose{โสตํ สุญฺญํ, สทฺทา สุญฺญา.\csromanspacer{} ฆานํ สุญฺญํ, คนฺธา สุญฺญา.\csromanspacer{} ชิวฺหา สุญฺญา, รสา สุญฺญา.\csromanspacer{} กาโย สุญฺโญ, โผฏฺฐพฺพา สุญฺญา.\csromanspacer{} มโน สุญฺโญ อตฺเตน วา อตฺตนิเยน วา.\csromanspacer{} ธมฺมา สุญฺญา อตฺเตน วา อตฺตนิเยน วา.\csromanspacer{} มโนวิญฺญาณํ สุญฺญํ อตฺเตน วา อตฺตนิเยน วา, มโนสมฺผสฺโส สุญฺโญ อตฺเตน วา อตฺตนิเยน วา, ยมฺปิทํ มโนสมฺผสฺสปจฺจยา อุปฺปชฺชติ เวทยิตํ สุขํ วา ทุกฺขํ วา อทุกฺขมสุขํ วา, ตมฺปิ สุญฺญํ อตฺเตน วา อตฺตนิเยน วา.\csromanspacer{} ยสฺมา โข อานนฺท สุญฺญํ อตฺเตน วา อตฺตนิเยน วา, ตสฺมา “สุญฺโญ โลโก”ติ วุจฺจตีติ.}

\csromanheader{2}{1. มาติกา}
\proseitem{47}{สุญฺญสุญฺญํ3, สงฺขารสุญฺญํ, วิปริณามสุญฺญํ, อคฺคสุญฺญํ, ลกฺขณสุญฺญํ, วิกฺขมฺภนสุญฺญํ, ตทงฺคสุญฺญํ, สมุจฺเฉทสุญฺญํ, ปฏิปฺปสฺสทฺธิสุญฺญํ, นิสฺสรณสุญฺญํ, อชฺฌตฺตสุญฺญํ, พหิทฺธาสุญฺญํ, ทุภโตสุญฺญํ, สภาคสุญฺญํ, วิสภาคสุญฺญํ, เอสนาสุญฺญํ, ปริคฺคหสุญฺญํ, ปฏิลาภสุญฺญํ, ปฏิเวธสุญฺญํ, เอกตฺตสุญฺญํ, นานตฺตสุญฺญํ, ขนฺติสุญฺญํ, อธิฏฺฐานสุญฺญํ, ปริโยคาหณสุญฺญํ\footnote{ปริโยคาหนสุญฺญํ (สฺยา)}, สมฺปชานสฺส ปวตฺตปริยาทานํ สพฺพสุญฺญตานํ ปรมตฺถสุญฺญํ. (25)}

\csromanheader{2}{2. นิทฺเทส}
\proseitem{48}{\csromanfolio{360}กตมํ \textbf{สุญฺญสุญฺญํ?} จกฺขุ สุญฺญํ อตฺเตน วา อตฺตนิเยน วา นิจฺเจน วา ธุเวน วา สสฺสเตน วา อวิปริณามธมฺเมน วา.\csromanspacer{}\csromanspacer{}โสตํ สุญฺญํ ฯเปฯ.\csromanspacer{} ฆานํ สุญฺญํ.\csromanspacer{} ชิวฺหา สุญฺญา.\csromanspacer{} กาโย สุญฺโญ, มโน สุญฺโญ อตฺเตน วา อตฺตนิเยน วา นิจฺเจน วา ธุเวน วา สสฺสเตน วา อวิปริณามธมฺเมน วา.\csromanspacer{} อิทํ สุญฺญสุญฺญํ. (1)}

\prose{กตมํ \textbf{สงฺขารสุญฺญํ?} ตโย สงฺขารา ปุญฺญาภิสงฺขาโร อปุญฺญาภิสงฺขาโร อาเนญฺชาภิสงฺขาโร.\csromanspacer{} ปุญฺญาภิสงฺขาโร อปุญฺญาภิสงฺขาเรน จ อาเนญฺชาภิสงฺขาเรน จ สุญฺโญ, อปุญฺญาภิสงฺขาโร ปุญฺญาภิสงฺขาเรน จ อาเนญฺชาภิสงฺขาเรน จ สุญฺโญ, อาเนญฺชาภิสงฺขาโร ปุญฺญาภิสงฺขาเรน จ อปุญฺญาภิสงฺขาเรน จ สุญฺโญ.\csromanspacer{} อิเม ตโย สงฺขารา.}

\prose{อปเรปิ ตโย สงฺขารา กายสงฺขาโร วจีสงฺขาโร จิตฺตสงฺขาโร.\csromanspacer{} กายสงฺขาโร วจีสงฺขาเรน จ จิตฺตสงฺขาเรน จ สุญฺโญ, วจีสงฺขาโร กายสงฺขาเรน จ จิตฺตสงฺขาเรน จ สุญฺโญ, จิตฺตสงฺขาโร กายสงฺขาเรน จ วจีสงฺขาเรน จ สุญฺโญ.\csromanspacer{} อิเม ตโย สงฺขารา.}

\prose{อปเรปิ ตโย สงฺขารา อตีตา สงฺขารา อนาคตา สงฺขารา ปจฺจุปฺปนฺนา สงฺขารา.\csromanspacer{} อตีตา สงฺขารา อนาคเตหิ จ ปจฺจุปฺปนฺเนหิ จ สงฺขาเรหิ สุญฺญา, อนาคตา สงฺขารา อตีเตหิ จ ปจฺจุปฺปนฺเนหิ จ สงฺขาเรหิ สุญฺญา, ปจฺจุปฺปนฺนา สงฺขารา อตีเตหิ จ อนาคเตหิ จ สงฺขาเรหิ สุญฺญา.\csromanspacer{} อิเม ตโย สงฺขารา.\csromanspacer{} อิทํ สงฺขารสุญฺญํ. (2)}

\prose{กตมํ \textbf{วิปริณามสุญฺญํ?} ชาตํ รูปํ สภาเวน สุญฺญํ, วิคตํ รูปํ วิปริณตญฺเจว สุญฺญญฺจ, ชาตา เวทนา สภาเวน สุญฺญา, วิคตา เวทนา วิปริณตา เจว สุญฺญา จ ฯเปฯ ชาตา สญฺญา.\csromanspacer{} ชาตา สงฺขารา.\csromanspacer{} ชาตํ วิญฺญาณํ.\csromanspacer{} ชาตํ จกฺขุ ฯเปฯ ชาโต ภโว สภาเวน สุญฺโญ, วิคโต ภโว วิปริณโต เจว สุญฺโญ จ.\csromanspacer{} อิทํ วิปริณามสุญฺญํ. (3)}

\prose{กตมํ \textbf{อคฺคสุญฺญํ?} อคฺคเมตํ ปทํ, เสฏฺฐเมตํ ปทํ, วิสิฏฺฐเมตํ\footnote{วิเสฏฺฐเมตํ (ก)} ปทํ, ยทิทํ สพฺพสงฺขารสมโถ สพฺพูปธิปฏินิสฺสคฺโค ตณฺหกฺขโย วิราโค นิโรโธ นิพฺพานํ.\csromanspacer{} อิทํ อคฺคสุญฺญํ. (4)}

\csromanheader{2}{10. สุญฺญกถา}
\prose{\csromanfolio{361}กตมํ \textbf{ลกฺขณสุญฺญํ?} เทฺว ลกฺขณานิ พาลลกฺขณญฺจ ปณฺฑิตลกฺขณญฺจ.\csromanspacer{} พาลลกฺขณํ ปณฺฑิตลกฺขเณน สุญฺญํ, ปณฺฑิตลกฺขณํ พาลลกฺขเณน สุญฺญํ.\csromanspacer{}\csromanspacer{}ตีณิ ลกฺขณานิ อุปฺปาทลกฺขณํ วยลกฺขณํ ฐิตญฺญถตฺตลกฺขณํ.\csromanspacer{} อุปฺปาทลกฺขณํ วยลกฺขเณน จ ฐิตญฺญถตฺตลกฺขเณน จ สุญฺญํ, วยลกฺขณํ อุปฺปาทลกฺขเณน จ ฐิตญฺญถตฺตลกฺขเณน จ สุญฺญํ, ฐิตญฺญถตฺตลกฺขณํ อุปฺปาทลกฺขเณน จ วยลกฺขเณน จ สุญฺญํ.}

\prose{รูปสฺส อุปฺปาทลกฺขณํ วยลกฺขเณน จ ฐิตญฺญถตฺตลกฺขเณน จ สุญฺญํ, รูปสฺส วยลกฺขณํ อุปฺปาทลกฺขเณน จ ฐิตญฺญถตฺตลกฺขเณน จ สุญฺญํ, รูปสฺส ฐิตญฺญถตฺตลกฺขณํ อุปฺปาทลกฺขเณน จ วยลกฺขเณน จ สุญฺญํ.\csromanspacer{}\csromanspacer{}เวทนาย ฯเปฯ.\csromanspacer{} สญฺญาย.\csromanspacer{} สงฺขารานํ.\csromanspacer{} วิญฺญาณสฺส.\csromanspacer{} จกฺขุสฺส ฯเปฯ.\csromanspacer{} ชรามรณสฺส อุปฺปาทลกฺขณํ วยลกฺขเณน จ ฐิตญฺญถตฺตลกฺขเณน จ สุญฺญํ, ชรามรณสฺส วยลกฺขณํ อุปฺปาทลกฺขเณน จ ฐิตญฺญถตฺตลกฺขเณน จ สุญฺญํ, ชรามรณสฺส ฐิตญฺญถตฺตลกฺขณํ อุปฺปาทลกฺขเณน จ วยลกฺขเณน จ สุญฺญํ.\csromanspacer{} อิทํ ลกฺขณสุญฺญํ. (5)}

\prose{กตมํ \textbf{วิกฺขมฺภนสุญฺญํ?} เนกฺขมฺเมน กามจฺฉนฺโท วิกฺขมฺภิโต เจว สุญฺโญ จ, อพฺยาปาเทน พฺยาปาโท วิกฺขมฺภิโต เจว สุญฺโญ จ, อาโลกสญฺญาย ถินมิทฺธํ วิกฺขมฺภิตญฺเจว สุญฺญญฺจ, อวิกฺเขเปน อุทฺธจฺจํ วิกฺขมฺภิตญฺเจว สุญฺญญฺจ, ธมฺมววตฺถาเนน วิจิกิจฺฉา วิกฺขมฺภิตา เจว สุญฺญา จ, ญาเณน อวิชฺชา วิกฺขมฺภิตา เจว สุญฺญา จ, ปาโมชฺเชน อรติ วิกฺขมฺภิตา เจว สุญฺญา จ, ปฐเมน ฌาเนน นีวรณา วิกฺขมฺภิตา เจว สุญฺญา จ ฯเปฯ อรหตฺตมคฺเคน สพฺพกิเลสา วิกฺขมฺภิตา เจว สุญฺญา จ.\csromanspacer{} อิทํ วิกฺขมฺภนสุญฺญํ. (6)}

\prose{กตมํ \textbf{ตทงฺคสุญฺญํ?} เนกฺขมฺเมน กามจฺฉนฺโท ตทงฺคสุญฺโญ, อพฺยาปาเทน พฺยาปาโท ตทงฺคสุญฺโญ, อาโลกสญฺญาย ถินมิทฺธํ ตทงฺคสุญฺญํ, อวิกฺเขเปน อุทฺธจฺจํ ตทงฺคสุญฺญํ, ธมฺมววตฺถาเนน วิจิกิจฺฉา ตทงฺคสุญฺญา, ญาเณน อวิชฺชา ตทงฺคสุญฺญา, ปาโมชฺเชน อรติ ตทงฺคสุญฺญา, ปฐเมน ฌาเนน นีวรณา ตทงฺคสุญฺญา ฯเปฯ วิวฏฺฏนานุปสฺสนาย สญฺโญคาภินิเวโส ตทงฺคสุญฺโญ.\csromanspacer{} อิทํ ตทงฺคสุญฺญํ. (7)}

\prose{กตมํ \textbf{สมุจฺเฉทสุญฺญํ?} เนกฺขมฺเมน กามจฺฉนฺโท สมุจฺฉินฺโน เจว สุญฺโญ จ.\csromanspacer{} อพฺยาปาเทน พฺยาปาโท สมุจฺฉินฺโน เจว สุญฺโญ จ,}

\prose{\csromanfolio{362}อาโลกสญฺญาย ถินมิทฺธํ สมุจฺฉินฺนญฺเจว สุญฺญญฺจ, อวิกฺเขเปน อุทฺธจฺจํ สมุจฺฉินฺนญฺเจว สุญฺญญฺจ, ธมฺมวตฺถาเนน วิจิกิจฺฉา สมุจฺฉินฺนา เจว สุญฺญา จ, ญาเณน อวิชฺชา สมุจฺฉินฺนา เจว สุญฺญา จ, ปาโมชฺเชน อรติ สมุจฺฉินฺนา เจว สุญฺญา จ, ปฐเมน ฌาเนน นีวรณา สมุจฺฉินฺนา เจว สุญฺญา จ ฯเปฯ อรหตฺตมคฺเคน สพฺพกิเลสา สมุจฺฉินฺนา เจว สุญฺญา จ.\csromanspacer{} อิทํ สมุจฺเฉทสุญฺญํ. (8)}

\prose{กตมํ \textbf{ปฏิปฺปสฺสทฺธิสุญฺญํ?} เนกฺขมฺเมน กามจฺฉนฺโท ปฏิปฺปสฺสทฺโธ เจว สุญฺโญ จ, อพฺยาปาเทน พฺยาปาโท ปฏิปฺปสฺสทฺโธ เจว สุญฺโญ จ.\csromanspacer{} อาโลกสญฺญาย ถินมิทฺธํ ปฏิปฺปสฺสทฺธญฺเจว สุญฺญญฺจ, อวิกฺเขเปน อุทฺธจฺจํ ปฏิปฺปสฺสทฺธญฺเจว สุญฺญญฺจ.\csromanspacer{} ธมฺมวตฺถาเนน วิจิกิจฺฉา ปฏิปฺปสฺสทฺธา เจว สุญฺญา จ, ญาเณน อวิชฺชา ปฏิปฺปสฺสทฺธา เจว สุญฺญา จ.\csromanspacer{} ปาโมชฺเชน อรติ ปฏิปฺปสฺสทฺธา เจว สุญฺญา จ, ปฐเมน ฌาเนน นีวรณา ปฏิปฺปสฺสทฺธา เจว สุญฺญา จ ฯเปฯ อรหตฺตมคฺเคน สพฺพกิเลสา ปฏิปฺปสฺสทฺธา เจว สุญฺญา จ.\csromanspacer{} อิทํ ปฏิปฺปสฺสทฺธิสุญฺญํ. (9)}

\prose{กตมํ \textbf{นิสฺสรณสุญฺญํ?} เนกฺขมฺเมน กามจฺฉนฺโท นิสฺสโฏ เจว สุญฺโญ จ, อพฺยาปาเทน พฺยาปาโท นิสฺสโฏ เจว สุญฺโญ จ, อาโลกสญฺญาย ถินมิทฺธํ นิสฺสฏญฺเจว สุญฺญญฺจ, อวิกฺเขเปน อุทฺธจฺจํ นิสฺสฏญฺเจว สุญฺญญฺจ, ธมฺมววตฺถาเนน วิจิกิจฺฉา นิสฺสฏา เจว สุญฺญา จ, ญาเณน อวิชฺชา นิสฺสฏา เจว สุญฺญา จ, ปาโมชฺเชน อรติ นิสฺสฏา เจว สุญฺญา จ, ปฐเมน ฌาเนน นีวรณา นิสฺสฏา เจว สุญฺญา จ ฯเปฯ อรหตฺตมคฺเคน สพฺพกิเลสา นิสฺสฏา เจว สุญฺญา จ.\csromanspacer{} อิทํ นิสฺสรณสุญฺญํ. (10)}

\prose{กตมํ \textbf{อชฺฌตฺตสุญฺญํ?} อชฺฌตฺตํ จกฺขุํ สุญฺญํ อตฺเตน วา อตฺตนิเยน วา นิจฺเจน วา ธุเวน วา สสฺสเตน วา อวิปริณามธมฺเมน วา.\csromanspacer{} อชฺฌตฺตํ โสตํ สุญฺญํ.\csromanspacer{} อชฺฌตฺตํ ฆานํ สุญฺญํ.\csromanspacer{} อชฺฌตฺตํ ชิวฺหา สุญฺญา.\csromanspacer{} อชฺฌตฺตํ กาโย สุญฺโญ.\csromanspacer{} อชฺฌตฺตํ มโน สุญฺโญ อตฺเตน วา อตฺตนิเยน วา นิจฺเจน วา ธุเวน วา สสฺสเตน วา อวิปริณามธมฺเมน วา.\csromanspacer{} อิทํ อชฺฌตฺตสุญฺญํ. (11)}

\prose{กตมํ \textbf{พหิทฺธาสุญฺญํ?} พหิทฺธา รูปา สุญฺญา ฯเปฯ.\csromanspacer{} พหิทฺธา ธมฺมา สุญฺญา อตฺเตน วา อตฺตนิเยน วา นิจฺเจน วา ธุเวน วา สสฺสเตน วา อวิปริณามธมฺเมน วา.\csromanspacer{} อิทํ พหิทฺธาสุญฺญํ. (12)}

\csromanheader{2}{10. สุญฺญกถา}
\prose{\csromanfolio{363}กตมํ \textbf{ทุภโตสุญฺญํ?} ยญฺจ อชฺฌตฺตํ จกฺขุ เย จ พหิทฺธา รูปา, อุภยเมตํ สุญฺญํ\footnote{อุภยโต ตํ สุญฺญา (สฺยา)} อตฺเตน วา อตฺตนิเยน วา นิจฺเจน วา ธุเวน วา สสฺสเตน วา อวิปริณามธมฺเมน วา.\csromanspacer{} ยญฺจ อชฺฌตฺตํ โสตํ, เย จ พหิทฺธา สทฺทา ฯเปฯ.\csromanspacer{} ยญฺจ อชฺฌตฺตํ ฆานํ, เย จ พหิทฺธา คนฺธา.\csromanspacer{} ยา จ อชฺฌตฺตํ ชิวฺหา, เย จ พหิทฺธา รสา.\csromanspacer{} โย จ อชฺฌตฺตํ กาโย, เย จ พหิทฺธา โผฏฺฐพฺพา.\csromanspacer{} โย จ อชฺฌตฺตํ มโน, เย จ พหิทฺธา ธมฺมา, อุภยเมตํ สุญฺญํ อตฺเตน วา อตฺตนิเยน วา นิจฺเจน วา ธุเวน วา สสฺสเตน วา อวิปริณามธมฺเมน วา.\csromanspacer{} อิทํ ทุภโตสุญฺญํ. (13)}

\prose{กตมํ \textbf{สภาคสุญฺญํ?} ฉ อชฺฌตฺติกานิ อายตนานิ สภาคานิ เจว สุญฺญานิ จ, ฉ พาหิรานิ อายตนานิ สภาคานิ เจว สุญฺญานิ จ, ฉ วิญฺญาณกายา สภาคา เจว สุญฺญา จ, ฉ ผสฺสกายา สภาคา เจว สุญฺญา จ, ฉ เวทนากายา สภาคา เจว สุญฺญา จ, ฉ สญฺญากายา สภาคา เจว สุญฺญา จ, ฉ เจตนากายา สภาคา เจว สุญฺญา จ.\csromanspacer{} อิทํ สภาคสุญฺญํ. (14)}

\prose{กตมํ \textbf{วิสภาคสุญฺญํ?} ฉ อชฺฌตฺติกานิ อายตนานิ ฉหิ พาหิเรหิ อายตเนหิ วิสภาคานิ เจว สุญฺญานิ จ, ฉ พาหิรานิ อายตนานิ ฉหิ วิญฺญาณกาเยหิ วิสภาคานิ เจว สุญฺญานิ จ, ฉ วิญฺญาณกายา ฉหิ ผสฺสกาเยหิ วิสภาคา เจว สุญฺญา จ, ฉ ผสฺสกายา ฉหิ เวทนากาเยหิ วิสภาคา เจว สุญฺญา จ, ฉ เวทนากายา ฉหิ สญฺญากาเยหิ วิสภาคา เจว สุญฺญา จ, ฉ สญฺญากายา ฉหิ เจตนากาเยหิ วิสภาคา เจว สุญฺญา จ.\csromanspacer{} อิทํ วิสภาคสุญฺญํ. (15)}

\prose{กตมํ \textbf{เอสนาสุญฺญํ?} เนกฺขมฺเมสนา กามจฺฉนฺเทน สุญฺญา, อพฺยาปาเทสนา พฺยาปาเทน สุญฺญา, อาโลกสญฺเญสนา ถินมิทฺเธน สุญฺญา, อวิกฺเขเปสนา อุทฺธจฺเจน สุญฺญา, ธมฺมววตฺถาเนสนา วิจิกิจฺฉาย สุญฺญา, ญาเณสนา อวิชฺชาย สุญฺญา, ปาโมชฺเชสนา อรติยา สุญฺญา, ปฐมชฺฌาเนสนา นีวรเณหิ สุญฺญา ฯเปฯ อรหตฺตมคฺเคสนา สพฺพกิเลเสหิ สุญฺญา.\csromanspacer{} อิทํ เอสนาสุญฺญํ. (16)}

\prose{กตมํ \textbf{ปริคฺคหสุญฺญํ?} เนกฺขมฺมปริคฺคโห กามจฺฉนฺเทน สุญฺโญ, อพฺยาปาทปริคฺคโห พฺยาปาเทน สุญฺโญ, อาโลกสญฺญาปริคฺคโห}

\prose{\csromanfolio{364}ถินมิทฺเธน สุญฺโญ, อวิกฺเขปปริคฺคโห อุทฺธจฺเจน สุญฺโญ, ธมฺมววตฺถานปริคฺคโห วิจิกิจฺฉาย สุญฺโญ ญาณปริคฺคโห อวิชฺชาย สุญฺโญ, ปาโมชฺชปริคฺคโห อรติยา สุญฺโญ, ปฐมชฺฌานปริคฺคโห นีวรเณหิ สุญฺโญ ฯเปฯ อรหตฺตมคฺคปริคฺคโต สพฺพกิเลเสหิ สุญฺโญ, อิทํ ปริคฺคหสุญฺญํ. (17)}

\prose{กตมํ \textbf{ปฏิลาภสุญฺญํ?} เนกฺขมฺมปฏิลาโภ กามจฺฉนฺเทน สุญฺโญ, อพฺยาปาทปฏิลาโภ พฺยาปาเทน สุญฺโญ, อาโลกสญฺญาปฏิลาโภ ถินมิทฺเธน สุญฺโญ, อวิกฺเขปปฏิลาโภ อุทฺธจฺเจน สุญฺโญ, ธมฺมววตฺถานปฏิลาโภ วิจิกิจฺฉาย สุญฺโญ, ญาณปฏิลาโภ อวิชฺชาย สุญฺโญ, ปาโมชฺชปฏิลาโภ อรติยา สุญฺโญ, ปฐมชฺฌานปฏิลาโภ นีวรเณหิ สุญฺโญ ฯเปฯ อรหตฺตมคฺคปฏิลาโภ สพฺพกิเลเสหิ สุญฺโญ.\csromanspacer{} อิทํ ปฏิลาภสุญฺญํ. (18)}

\prose{กตมํ \textbf{ปฏิเวธสุญฺญํ?} เนกฺขมฺมปฺปฏิเวโธ กามจฺฉนฺเทน สุญฺโญ, อพฺยาปาทปฺปฏิเวโธ พฺยาปาเทน สุญฺโญ, อาโลกสญฺญาปฺปฏิเวโธ ถินมิทฺเธน สุญฺโญ, อวิกฺเขปปฺปฏิเวโธ อุทฺธจฺเจน สุญฺโญ, ธมฺมววตฺถานปฺปฏิเวโธ วิจิกิจฺฉาย สุญฺโญ, ญาณปฺปฏิเวโธ อวิชฺชาย สุญฺโญ, ปาโมชฺชปฺปฏิเวโธ อรติยา สุญฺโญ, ปฐมชฺฌานปฺปฏิเวโธ นีวรเณหิ สุญฺโญ ฯเปฯ อรหตฺตมคฺคปฺปฏิเวโธ สพฺพกิเลเสหิ สุญฺโญ.\csromanspacer{} อิทํ ปฏิเวธสุญฺญํ. (19)}

\prose{กตมํ \textbf{เอกตฺตสุญฺญํ, นานตฺตสุญฺญํ?} กามจฺฉนฺโท นานตฺตํ, เนกฺขมฺมํ เอกตฺตํ, เนกฺขมฺเมกตฺตํ เจตยโต กามจฺฉนฺเทน สุญฺญํ.\csromanspacer{} พฺยาปาโท นานตฺตํ, อพฺยาปาโท เอกตฺตํ, อพฺยาปาเทกตฺตํ เจตยโต พฺยาปาเทน สุญฺญํ.\csromanspacer{} ถินมิทฺธํ นานตฺตํ อาโลกสญฺญา เอกตฺตํ, อาโลกสญฺเญกตฺตํ เจตยโต ถินมิทฺเธน สุญฺญํ.\csromanspacer{} อุทฺธจฺจํ นานตฺตํ, อวิกฺเขโป เอกตฺตํ, อวิกฺเขเปกตฺตํ เจตยโต อุทฺธจฺเจน สุญฺญํ, วิจิกิจฺฉา นานตฺตํ, ธมฺมววตฺถานํ เอกตฺตํ, ธมฺมววตฺถาเนกตฺตํ เจตยโต วิจิกิจฺฉาย สุญฺญํ.\csromanspacer{} อวิชฺชา นานตฺตํ, ญาณํ เอกตฺตํ, ญาเณกตฺตํ, เจตยโต อวิชฺชาย สุญฺญํ.\csromanspacer{} อรติ นานตฺตํ, ปโมชฺชํ เอกตฺตํ, ปาโมชฺเชกตฺตํ เจตยโต อรติยา สุญฺญํ.\csromanspacer{} นีวรณา นานตฺตํ, ปฐมชฺฌานํ เอกตฺตํ, ปฐมชฺฌาเนกตฺตํ เจตยโต นีวรเณหิ สุญฺญํ ฯเปฯ สพฺพกิเลสา นานตฺตํ, อรหตฺตมคฺโค เอกตฺตํ,}

\csromanheader{2}{10. สุญฺญกถา}
\prose{\csromanfolio{365}อรหตฺตมคฺเคกตฺตํ เจตยโต สพฺพกิเลเสหิ สุญฺญํ.\csromanspacer{} อิทํ เอกตฺตสุญฺญํ, นานตฺตสุญฺญํ. (20, 21)}

\prose{กตมํ \textbf{ขนฺติสุญฺญํ?} เนกฺขมฺมขนฺติ กามจฺฉนฺเทน สุญฺญา, อพฺยาปาทขนฺติ พฺยาปาเทน สุญฺญา, อาโลกสญฺญาขนฺติ ถินมิทฺเธน สุญฺญา, อวิกฺเขปขนฺติ อุทฺธจฺเจน สุญฺญา, ธมฺมววตฺถานขนฺติ วิจิกิจฺฉาย สุญฺญา, ญาณขนฺติ อวิชฺชาย สุญฺญา, ปาโมชฺชขนฺติ อรติยา สุญฺญา, ปฐมชฺฌานขนฺติ นีวรเณหิ สุญฺญา ฯเปฯ อรหตฺตมคฺคขนฺติ สพฺพกิเลเสหิ สุญฺญา.\csromanspacer{} อิทํ ขนฺติสุญฺญํ. (22)}

\prose{กตมํ \textbf{อธิฏฺฐานสุญฺญํ?} เนกฺขมฺมาธิฏฺฐานํ กามจฺฉนฺเทน สุญฺญํ, อพฺยาปาทาธิฏฺฐานํ พฺยาปาเทน สุญฺญํ, อาโลกสญฺญาธิฏฺฐานํ ถินมิทฺเธน สุญฺญํ, อวิกฺเขปาธิฏฺฐานํ อุทฺธจฺเจน สุญฺญํ, ธมฺมววตฺถานาธิฏฺฐานํ วิจิกิจฺฉาย สุญฺญํ, ญาณาธิฏฺฐานํ อวิชฺชาย สุญฺญํ, ปาโมชฺชาธิฏฺฐานํ อรติยา สุญฺญํ, ปฐมชฺฌานาธิฏฺฐานํ นีวรเณหิ สุญฺญํ ฯเปฯ อรหตฺตมคฺคาธิฏฺฐานํ สพฺพกิเลเสหิ สุญฺญํ.\csromanspacer{} อิทํ อธิฏฺฐานสุญฺญํ. (23)}

\prose{กตมํ \textbf{ปริโยคาหณสุญฺญํ?} เนกฺขมฺมปริโยคาหณํ กามจฺฉนฺเทน สุญฺญํ, อพฺยาปาทปริโยคาหณํ พฺยาปาเทน สุญฺญํ, อาโลกสญฺญาปริโยคาหณํ ถินมิทฺเธน สุญฺญํ, อวิกฺเขปปริโยคาหณํ อุทฺธจฺเจน สุญฺญํ, ธมฺมววตฺถานปริโยคาหณํ วิจิกิจฺฉาย สุญฺญํ, ญาณปริโยคาหณํ อวิชฺชาย สุญฺญํ, ปาโมชฺชปริโยคาหณํ อรติยา สุญฺญํ, ปฐมชฺฌานปริโยคาหณํ นีวรเณหิ สุญฺญํ ฯเปฯ อรหตฺตมคฺคปริโยคาหณํ สพฺพกิเลเสหิ สุญฺญํ.\csromanspacer{} อิทํ ปริโยคาหณสุญฺญํ. (24)}

\prose{กตมํ \textbf{สมฺปชานสฺส ปวตฺตปริยาทานํ สพฺพสุญฺญตานํ ปรมตฺถสุญฺญํ?} อิธ สมฺปชาโน เนกฺขมฺเมน กามจฺฉนฺทสฺส ปวตฺตํ ปริยาทิยติ, อพฺยาปาเทน พฺยาปาทสฺส ปวตฺตํ ปริยาทิยติ, อาโลกสญฺญาย ถินมิทฺธสฺส ปวตฺตํ ปริยาทิยติ, อวิกฺเขเปน อุทฺธจฺจสฺส ปวตฺตํ ปริยาทิยติ, ธมฺมววตฺถาเนน วิจิกิจฺฉาย ปวตฺตํ ปริยาทิยติ, ญาเณน อวิชฺชาย ปวตฺตํ ปริยาทิยติ, ปาโมชฺเชน อรติยา ปวตฺตํ ปริยาทิยติ, ปฐเมน ฌาเนน นีวรณานํ ปวตฺตํ ปริยาทิยติ ฯเปฯ อรหตฺตมคฺเคน สพฺพกิเลสานํ ปวตฺตํ ปริยาทิยติ.\csromanspacer{} อถ วา ปน สมฺปชานสฺส อนุปาทิเสสาย นิพฺพานธาตุยา ปรินิพฺพายนฺตสฺส อิทํ เจว จกฺขุปวตฺตํ ปริยาทิยติ, อญฺญญฺจ จกฺขุปวตฺตํ}

\prose{\csromanfolio{366}น อุปฺปชฺชติ.\csromanspacer{} อิทํ เจว โสตปวตฺตํ ฯเปฯ.\csromanspacer{} ฆานปวตฺตํ.\csromanspacer{} ชิวฺหาปวตฺตํ.\csromanspacer{} กายปวตฺตํ.\csromanspacer{} มโนปวตฺตํ ปริยาทิยติ, อญฺญญฺจ มโนปวตฺตํ น อุปฺปชฺชติ.\csromanspacer{} อิทํ สมฺปชานสฺส ปวตฺตปริยาทานํ สพฺพสุญฺญตานํ ปรมตฺถสุญฺญนฺติ. (25)}

\nitthitam{สุญฺญกถา นิฏฺฐิตา.}
\vspace{\tipitakaparskip}
\csromancenter{\textbf{ยุคนทฺธวคฺโค ทุติโย.}}
\titlehead{\textbf{ตสฺสุทฺทานํ}}
\csromangathameasure{ยุคนทฺธา สจฺจโพชฺฌงฺคา, เมตฺตา วิราคปญฺจมา. \\ ปฏิสมฺภิทา ธมฺมจกฺกํ, โลกุตฺตรพลสุญฺญาติ. \\ เอส นิกายธเรหิ ฐปิโต อสโม ทุติโย ปวโร วรวคฺโคติ.}
\csromangathagroup{\csromangathastanza{ยุคนทฺธา สจฺจโพชฺฌงฺคา, เมตฺตา วิราคปญฺจมา. \\ ปฏิสมฺภิทา ธมฺมจกฺกํ, โลกุตฺตรพลสุญฺญาติ. \\ เอส นิกายธเรหิ ฐปิโต อสโม ทุติโย ปวโร วรวคฺโคติ.}}

\csromanmatikaanchor{mtk.118}
\csromantocmarknopage{chapter}{3. ปญฺญาวคฺค}{mtk.118}
\bookmark[dest={mtk.118},level=0]{3. ปญฺญาวคฺค}
\chapterheadpageread{3. ปญฺญาวคฺค}
\csromanmatikaanchor{mtk.119}
\csromantocmarknopage{section}{1. มหาปญฺญากถา}{mtk.119}
\bookmark[dest={mtk.119},level=1]{1. มหาปญฺญากถา}
\csromanheader{1}{1. มหาปญฺญากถา}
\proseitem{1}{\csromanfolio{367}อนิจฺจานุปสฺสนา ภาวิตา พหุลีกตา กตมํ ปญฺญํ ปริปูเรติ, ทุกฺขานุปสฺสนา ภาวิตา พหุลีกตา กตมํ ปญฺญํ ปริปูเรติ, อนตฺตานุปสฺสนา ภาวิตา พหุลีกตา กตมํ ปญฺญํ ปริปูเรติ ฯเปฯ ปฏินิสฺสคฺคานุปสฺสนา ภาวิตา พหุลีกตา กตมํ ปญฺญํ ปริปูเรติ?}

\prose{อนิจฺจานุปสฺสนา ภาวิตา พหุลีกตา ชวนปญฺญํ ปริปูเรติ, ทุกฺขานุปสฺสนา ภาวิตา พหุลีกตา นิพฺเพธิกปญฺญํ ปริปูเรติ, อนตฺตานุปสฺสนา ภาวิตา พหุลีกตา มหาปญฺญํ ปริปูเรติ, นิพฺพิทานุปสฺสนา ภาวิตา พหุลีกตา ติกฺขปญฺญํ ปริปูเรติ, วิราคานุปสฺสนา ภาวิตา พหุลีกตา วิปุลปญฺญํ ปริปูเรติ, นิโรธานุปสฺสนา ภาวิตา พหุลีกตา คมฺภีรปญฺญํ ปริปูเรติ, ปฏินิสฺสคฺคานุปสฺสนา ภาวิตา พหุลีกตา อสามนฺตปญฺญํ\footnote{อสฺสามนฺตปญฺญํ (สฺยา)} ปริปูเรติ.\csromanspacer{} อิมา สตฺต ปญฺญา ภาวิตา พหุลีกตา ปณฺฑิจฺจํ ปริปูเรนฺติ, อิมา อฏฺฐ ปญฺญา ภาวิตา พหุลีกตา ปุถุปญฺญํ ปริปูเรนฺติ, อิมา นว ปญฺญา ภาวิตา พหุลีกตา หาสปญฺญํ ปริปูเรนฺติ.}

\prose{หาสปญฺญา ปฏิภานปฏิสมฺภิทา, ตสฺสา อตฺถววตฺถานโต อตฺถปฏิสมฺภิทา อธิคตา โหติ สจฺฉิกตา ผสฺสิตา ปญฺญาย, ธมฺมววตฺถานโต ธมฺมปฏิสมฺภีทา อธิคตา โหติ สจฺฉิกตา ผสฺสิตา ปญฺญาย, นิรุตฺติววตฺถานโต นิรุตฺติปฏิสมฺภิทา อธิคตา โหติ สจฺฉิกตา ผสฺสิตา ปญฺญาย, ปฏิภานววตฺถานโต ปฏิภานปฏิสมฺภิทา อธิคตา โหติ สจฺฉิกตา ผสฺสิตา ปญฺญาย.\csromanspacer{} ตสฺสิมา จตสฺโส ปฏิสมฺภิทาโย อธิคตา โหนฺติ สจฺฉิกตา ผสฺสิตา ปญฺญาย.}

\prose{รูเป อนิจฺจานุปสฺสนา ภาวิตา พหุลีกตา กตมํ ปญฺญํ ปริปูเรติ ฯเปฯ รูเป ปฏินิสฺสคฺคานุปสฺสนา ภาวิตา พหุลีกตา กตมํ ปญฺญํ ปริปูเรติ? รูเป อนิจฺจานุปสฺสนา ภาวิตา พหุลีกตา ชวนปญฺญํ ปริปูเรติ ฯเปฯ รูเป ปฏินิสฺสคฺคานุปสฺสนา ภาวิตา พหุลีกตา อสามนฺตปญฺญํ ปริปูเรติ.\csromanspacer{} อิมา สตฺต ปญฺญา ภาวิตา พหุลีกตา ปณฺฑิจฺจํ ปริปูเรนฺติ, อิมา อฏฺฐ \csromanfolio{368}ปญฺญา ภาวิตา พหุลีกตา ปุถุปญฺญํ ปริปูเรนฺติ, อิมา นว ปญฺญา ภาวิตา พหุลีกตา หาสปญฺญํ ปริปูเรนฺติ.}

\prose{หาสปญฺญา ปฏิภานปฏิสมฺภิทา, ตสฺสา อตฺถววตฺถานโต อตฺถปฏิสมฺภิทา อธิคตา โหติ สจฺฉิกตา ผสฺสิตา ปญฺญาย, ธมฺมววตฺถานโต ธมฺมปฏิสมฺภิทา อธิคตา โหติ สจฺฉิกตา ผสฺสิตา ปญฺญาย, นิรุตฺติววตฺถานโต นิรุตฺติปฏิสมฺภิทา อธิคตา โหติ สจฺฉิกตา ผสฺสิตา ปญฺญาย, ปฏิภานววตฺถานโต ปฏิภานปฏิสมฺภิทา อธิคตา โหติ สจฺฉิกตา ผสฺสิตา ปญฺญาย.\csromanspacer{} ตสฺสิมา จตสฺโส ปฏิสมฺภิทาโย อธิคตา โหนฺติ สจฺฉิกตา ผสฺสิตา ปญฺญาย.}

\prose{เวทนาย ฯเปฯ.\csromanspacer{} สญฺญาย.\csromanspacer{} สงฺขาเรสุ.\csromanspacer{} วิญฺญาเณ.\csromanspacer{} จกฺขุสฺมิํ ฯเปฯ.\csromanspacer{} ชรามรเณ อนิจฺจานุปสฺสนา ภาวิตา พหุลีกตา กตมํ ปญฺญํ ปริปูเรติ ฯเปฯ ชรามรเณ ปฏินิสฺสคฺคานุปสฺสนา ภาวิตา พหุลีกตา กตมํ ปญฺญํ ปริปูเรติ? ชรามรเณ อนิจฺจานุปสฺสนา ภาวิตา พหุลีกตา ชวนปญฺญํ ปริปูเรติ ฯเปฯ ชรามรเณ ปฏินิสฺสคฺคานุปสฺสนา ภาวิตา พหุลีกตา อสามนฺตปญฺญํ ปริปูเรติ.\csromanspacer{} อิมา สตฺต ปญฺญา ภาวิตา พหุลีกตา ปณฺฑิจฺจํ ปริปูเรนฺติ, อิมา อฏฺฐ ปญฺญา ภาวิตา พหุลีกตา ปุถุปญฺญํ ปริปูเรนฺติ.\csromanspacer{} อิมา นว ปญฺญา ภาวิตา พหุลีกตา หาสปญฺญํ ปริปูเรนฺติ.}

\prose{หาสปญฺญา ปฏิภานปฏิสมฺภิทา, ตสฺสา อตฺถววตฺถานโต อตฺถปฏิสมฺภิทา อธิคตา โหติ สจฺฉิกตา ผสฺสิตา ปญฺญาย.\csromanspacer{} ธมฺมววตฺถานโต ธมฺมปฏิสมฺภิทา อธิคตา โหติ สจฺฉิกตา ผสฺสิตา ปญฺญาย.\csromanspacer{} นิรุตฺติววตฺถานโต นิรุตฺติปฏิสมฺภิทา อธิคตา โหติ สจฺฉิกตา ผสฺสิตา ปญฺญาย, ปฏิภานววตฺถานโต ปฏิภานปฏิสมฺภิทา อธิคตา โหติ สจฺฉิกตา ผสฺสิตา ปญฺญาย, ตสฺสิมา จตสฺโส ปฏิสมฺภิทาโย อธิคตา โหนฺติ สจฺฉิกตา ผสฺสิตา ปญฺญาย.}

\proseitem{2}{รูเป อนิจฺจานุปสฺสนา ภาวิตา พหุลีกตา กตมํ ปญฺญํ ปริปูเรติ.\csromanspacer{} อตีตานาคตปจฺจุปฺปนฺเน รูเป อนิจฺจานุปสฺสนา ภาวิตา พหุลีกตา กตมํ ปญฺญํ ปริปูเรติ.\csromanspacer{} รูเป ทุกฺขานุปสฺสนา ภาวิตา พหุลีกตา กตมํ ปญฺญํ ปริปูเรติ.\csromanspacer{} อตีตานาคตปจฺจุปฺปนฺเน รูเป ทุกฺขานุปสฺสนา ภาวิตา พหุลีกตา กตมํ ปญฺญํ ปริปูเรติ.\csromanspacer{} รูเป}

\csromanheader{2}{1. มหาปญฺญากถา}
\prose{\csromanfolio{369}อนตฺตานุปสฺสนา ภาวิตา พหุลีกตา กตมํ ปญฺญํ ปริปูเรติ.\csromanspacer{} อตีตานาคตปจฺจุปฺปนฺเน รูเป อนตฺตานุปสฺสนา ภาวิตา พหุลีกตา กตมํ ปญฺญํ ปริปูเรติ.\csromanspacer{} รูเป นิพฺพิทานุปสฺสนา ภาวิตา พหุลีกตา กตมํ ปญฺญํ ปริปูเรติ.\csromanspacer{} อตีตานาคตปจฺจุปฺปนฺเน รูเป นิพฺพิทานุปสฺสนา ภาวิตา พหุลีกตา กตมํ ปญฺญํ ปริปูเรติ.\csromanspacer{} รูเป วิราคานุปสฺสนา ภาวิตา พหุลีกตา กตมํ ปญฺญํ ปริปูเรติ.\csromanspacer{} อตีตานาคตปจฺจุปฺปนฺเน รูเป วิราคานุปสฺสนา ภาวิตา พหุลีกตา กตมํ ปญฺญํ ปริปูเรติ.\csromanspacer{} รูเป นิโรธานุปสฺสนา ภาวิตา พหุลีกตา กตมํ ปญฺญํ ปริปูเรติ.\csromanspacer{} อตีตานาคตปจฺจุปฺปนฺเน รูเป นิโรธานุปสฺสนา ภาวิตา พหุลีกตา กตมํ ปญฺญํ ปริปูเรติ.\csromanspacer{} รูเป ปฏินิสฺสคฺคานุปสฺสนา ภาวิตา พหุลีกตา กตมํ ปญฺญํ ปริปูเรติ.\csromanspacer{} อตีตานาคตปจฺจุปฺปนฺเน รูเป ปฏินิสฺสคฺคานุปสฺสนา ภาวิตา พหุลีกตา กตมํ ปญฺญํ ปริปูเรติ?}

\prose{รูเป อนิจฺจานุปสฺสนา ภาวิตา พหุลีกตา ชวนปญฺญํ ปริปูเรติ.\csromanspacer{} อตีตานาคตปจฺจุปฺปนฺเน รูเป อนิจฺจานุปสฺสนา ภาวิตา พหุลีกตา ชวนปญฺญํ ปริปูเรติ.\csromanspacer{} รูเป ทุกฺขานุปสฺสนา ภาวิตา พหุลีกตา นิพฺเพธิกปญฺญํ ปริปูเรติ.\csromanspacer{} อตีตานาคตปจฺจุปฺปนฺเน รูเป ทุกฺขานุปสฺสนา ภาวิตา พหุลีกตา ชวนปญฺญํ ปริปูเรติ.\csromanspacer{} รูเป อนตฺตานุปสฺสนา ภาวิตา พหุลีกตา มหาปญฺญํ ปริปูเรติ.\csromanspacer{} อตีตานาคตปจฺจุปฺปนฺเน รูเป อนตฺตานุปสฺสนา ภาวิตา พหุลีกตา ชวนปญฺญํ ปริปูเรติ.\csromanspacer{} รูเป นิพฺพิทานุปสฺสนา ภาวิตา พหุลีกตา ติกฺขปญฺญํ ปริปูเรติ.\csromanspacer{} อตีตานาคตปจฺจุปฺปนฺเน รูเป นิพฺพิทานุปสฺสนา ภาวิตา พหุลีกตา ชวนปญฺญํ ปริปูเรติ.\csromanspacer{} รูเป วิราคานุปสฺสนา ภาวิตา พหุลีกตา วิปุลปญฺญํ ปริปูเรติ.\csromanspacer{} อตีตานาคตปจฺจุปฺปนฺเน รูเป วิราคานุปสฺสนา ภาวิตา พหุลีกตา ชวนปญฺญํ ปริปูเรติ.\csromanspacer{} รูเป นิโรธานุปสฺสนา ภาวิตา พหุลีกตา คมฺภีรปญฺญํ ปริปูเรติ.\csromanspacer{} อตีตานาคตปจฺจุปฺปนฺเน รูเป นิโรธานุปสฺสนา ภาวิตา พหุลีกตา ชาวนปญฺญํ ปริปูเรติ.\csromanspacer{} รูเป ปฏินิสฺสคฺคานุปสฺสนา ภาวิตา พหุลีกตา อสามนฺตปญฺญํ ปริปูเรติ.\csromanspacer{} อตีตานาคตปจฺจุปฺปนฺเน รูเป ปฏินิสฺสคฺคานุปสฺสนา ภาวิตา พหุลีกตา ชวนปญฺญํ ปริปูเรติ.\csromanspacer{} อิมา สตฺต ปญฺญา ภาวิตา พหุลีกตา ปณฺฑิจฺจํ ปริปูเรนฺติ.\csromanspacer{} อิมา อฏฺฐ ปญฺญา ภาวิตา พหุลีกตา ปุถุปญฺญํ ปริปูเรนฺติ.\csromanspacer{} อิมา นว ปญฺญา ภาวิตา พหุลีกตา หาสปญฺญํ ปริปูเรนฺติ.}

\prose{\csromanfolio{370}หาสปญฺญา ปฏิภานปฏิสมฺภิทา, ตสฺสา อตฺถววตฺถานโต อตฺถปฏิสมฺภิทา อธิคตา โหติ สจฺฉิกตา ผสฺสิตา ปญฺญาย, ธมฺมววตฺถานโต ธมฺมปฏิสมฺภิทา อธิคตา โหติ สจฺฉิกตา ผสฺสิตา ปญฺญาย, นิรุตฺติววตฺถานโต นิรุตฺติปฏิสมฺภิทา อธิคตา โหติ สจฺฉิกตา ผสฺสิตา ปญฺญาย, ปฏิภานววตฺถานโต ปฏิภานปฏิสมฺภิทา อธิคตา โหติ สจฺฉิกตา ผสฺสิตา ปญฺญาย.\csromanspacer{} ตสฺสิมา จตสฺโส ปฏิสมฺภิทาโย อธิคตา โหนฺติ สจฺฉิกตา ผสฺสิตา ปญฺญาย.}

\prose{เวทนาย ฯเปฯ.\csromanspacer{} สญฺญาย.\csromanspacer{} สงฺขาเรสุ.\csromanspacer{} วิญฺญาเณ.\csromanspacer{} จกฺขุสฺมิํ ฯเปฯ.\csromanspacer{} ชรามรเณ อนิจฺจานุปสฺสนา ภาวิตา พหุลีกตา กตมํ ปญฺญํ ปริปูเรติ.\csromanspacer{} อตีตานาคตปจฺจุปฺปนฺเน ชรามรเณ อนิจฺจานุปสฺสนา ภาวิตา พหุลีกตา กตมํ ปญฺญํ ปริปูเรติ ฯเปฯ ชรามรเณ ปฏินิสฺสคฺคานุปสฺสนา ภาวิตา พหุลีกตา กตมํ ปญฺญํ ปริปูเรติ.\csromanspacer{} อตีตานาคตปจฺจุปฺปนฺเน ชรามรเณ ปฏินิสฺสคฺคานุปสฺสนา ภาวิตา พหุลีกตา กตมํ ปญฺญํ ปริปูเรติ? ชรามรเณ อนิจฺจานุปสฺสนา ภาวิตา พหุลีกตา ชวนปญฺญํ ปริปูเรติ.\csromanspacer{} อตีตานาคตปจฺจุปฺปนฺเน ชรามรเณ อนิจฺจานุปสฺสนา ภาวิตา พหุลีกตา ชวนปญฺญํ ปริปูเรติ ฯเปฯ.\csromanspacer{} ตสฺสิมา จตสฺโส ปฏิสมฺภิทาโย อธิคตา โหนฺติ สจฺฉิกตา ผสฺสิตา ปญฺญาย.}

\proseitem{3}{จตฺตาโรเม ภิกฺขเว\footnote{สํ 3. 359 ปิฏฺเฐปิ.} ธมฺมา ภาวิตา พหุลีกตา โสตาปตฺติผลสจฺฉิกิริยาย สํวตฺตนฺติ.\csromanspacer{} กตเม จตฺตาโร? สปฺปุริสสํเสโว สทฺธมฺมสฺสวนํ โยนิโสมนสิกาโร ธมฺมานุธมฺมปฏิปตฺติ.\csromanspacer{} อิเม โข ภิกฺขเว จตฺตาโร ธมฺมา ภาวิตา พหุลีกตา โสตาปตฺติผลสจฺฉิกิริยาย สํวตฺตนฺติ.}

\prose{จตฺตาโรเม ภิกฺขเว ธมฺมา ภาวิตา พหุลีกตา สกทาคามิผลสจฺฉิกิริยาย สํวตฺตนฺติ ฯเปฯ อนาคามิผลสจฺฉิกิริยาย สํวตฺตนฺติ ฯเปฯ อรหตฺตผลสจฺฉิกิริยาย สํวตฺตนฺติ.\csromanspacer{} กตเม จตฺตาโร, สปฺปุริสสํเสโว สทฺธมฺมสฺสวนํ โยนิโสมนสิกาโร ธมฺมานุธมฺมปฏิปตฺติ.\csromanspacer{} อิเม โข ภิกฺขเว จตฺตาโร ธมฺมา ภาวิตา พหุลีกตา อรหตฺตมคฺคผลสจฺฉิกิริยาย สํวตฺตนฺติ.}

\csromanheader{2}{1. มหาปญฺญากถา}
\prose{\csromanfolio{371}จตฺตาโรเม ภิกฺขเว ธมฺมา ภาวิตา พหุลีกตา ปญฺญาปฏิลาภาย สํวตฺตนฺติ ฯเปฯ ปญฺญาพุทฺธิยา สํวตฺตนฺติ.\csromanspacer{} ปญฺญาเวปุลฺลาย สํวตฺตนฺติ.\csromanspacer{} มหาปญฺญตาย สํวตฺตนฺติ.\csromanspacer{} ปุถุปญฺญตาย สํวตฺตนฺติ.\csromanspacer{} วิปุลปญฺญตาย สํวตฺตนฺติ.\csromanspacer{} คมฺภีรปญฺญตาย สํวตฺตนฺติ.\csromanspacer{} อสามนฺตปญฺญตาย\footnote{อสฺสามนฺตปญฺญตาย (สฺยา) สํ 3. 361 ปิฏฺเฐ ปสฺสิตพฺพา.} สํวตฺตนฺติ.\csromanspacer{} ภูริปญฺญตาย สํวตฺตนฺติ.\csromanspacer{} ปญฺญาพาหุลฺลาย สํวตฺตนฺติ.\csromanspacer{} สีฆปญฺญตาย สํวตฺตนฺติ.\csromanspacer{} ลหุปญฺญตาย สํวตฺตนฺติ.\csromanspacer{} หาสปญฺญตาย สํวตฺตนฺติ.\csromanspacer{} ชวนปญฺญตาย สํวตฺตนฺติ.\csromanspacer{} ติกฺขปญฺญตาย สํวตฺตนฺติ.\csromanspacer{} นิพฺเพธิกปญฺญตาย สํวตฺตนฺติ.\csromanspacer{} กตเม จตฺตาโร? สปฺปุริสสํเสโว สทฺธมฺมสฺสวนํ โยนิโสมนสิกาโร ธมฺมานุธมฺมปฏิปตฺติ.\csromanspacer{} อิเม โข ภิกฺขเว จตฺตาโร ธมฺมา ภาวิตา พหุลีกตา ปญฺญาปฏิลาภาย สํวตฺตนฺติ.\csromanspacer{} ปญฺญาพุทฺธิยา สํวตฺตนฺติ ฯเปฯ นิพฺเพธิกปญฺญตาย สํวตฺตนฺติ.}

\csromanmatikaanchor{mtk.120}
\csromantocmarkref{subsection}{1. โสฬสปญฺญานิทฺเทส}{mtk.120}
\bookmark[dest={mtk.120},level=2]{1. โสฬสปญฺญานิทฺเทส}
\csromanheader{2}{1. โสฬสปญฺญานิทฺเทส}
\proseitem{4}{ปญฺญาปฏิลาภาย สํวตฺตนฺตีติ กตโม ปญฺญาปฏิลาโภ? จตุนฺนํ มคฺคญาณานํ จตุนฺนํ ผลญาณานํ จตุนฺนํ ปฏิสมฺภิทาญาณานํ ฉนฺนํ อภิญฺญาญาณานํ เตสตฺตตีนํ ญาณานํ สตฺตสตฺตตีนํ ญาณานํ ลาโภ ปฏิลาโภ ปตฺติ สมฺปตฺติ ผสฺสนา\footnote{ผุสนา (ก)} สจฺฉิกิริยา อุปสมฺปทา ปญฺญาปฏิลาภาย สํวตฺตนฺตีติ อยํ ปญฺญา ปฏิลาโภ. (1)}

\prose{\textbf{ปญฺญาพุทฺธิยา สํวตฺตนฺตี}ติ กตมา ปญฺญาพุทฺธิ? สตฺตนฺนํ จ เสกฺขานํ ปุถุชฺชนกลฺยาณกสฺส จ ปญฺญา วฑฺฒติ, อรหโต ปญฺญา วฑฺฒติ, วฑฺฒิภวฑฺฒนา ปญฺญาพุทฺธิยา สํวตฺตนฺติตี อยํ ปญฺญาพุทฺธิ. (2)}

\prose{\textbf{ปญฺญาเวปุลฺลาย สํวตฺตนฺตี}ติ กตมํ ปญฺญาเวปุลฺลํ? สตฺตนฺนํ เสกฺขานํ ปุถุชฺชนกลฺยาณกสฺส จ ปญฺญาเวปุลฺลํ คจฺฉติ.\csromanspacer{} อรหโต ปญฺญา เวปุลฺลคตา ปญฺญาเวปุลฺลาย สํวตฺตนฺตีติ อิทํ ปญฺญาเวปุลฺลํ. (3)}

\prose{\textbf{มหาปญฺญตาย สํวตฺตนฺตี}ติ กตมา มหาปญฺญา? มหนฺเต อตฺเถ ปริคฺคณฺหาตีติ มหาปญฺญา, มหนฺเต ธมฺเม ปริคฺคณฺหาตีติ มหาปญฺญา, มหนฺตา นิรุตฺติโย ปริคฺคณฺหาตีติ มหาปญฺญา, มหนฺตานิ ปฏิภานานิ ปริคฺคณฺหาตีติ มหาปญฺญา, มหนฺเต สีลกฺขนฺเธ ปริคฺคณฺหาตีติ มหาปญฺญา, มหนฺเต สมาธิกฺขนฺเธ ปริคฺคณฺหาตีติ มหาปญฺญา, มหนฺเต ปญฺญากฺขนฺเธ \csromanfolio{372}ปริคฺคณฺหาตีติ มหาปญฺญา, มหนฺเต วิมุตฺติกฺขนฺเธ ปริคฺคณฺหาตีติ มหาปญฺญา, มหนฺเต วิมุตฺติญาณทสฺสนกฺขนฺเธ ปริคฺคณฺหาตีติ มหาปญฺญา, มหนฺตานิ ฐานาฏฺฐานานิ ปริคฺคณฺหาตีติ มหาปญฺญา, มหนฺตา วิหารสมาปตฺติโย ปริคฺคณฺหาตีติ มหาปญฺญา, มหนฺตานิ อริยสจฺจานิ ปริคฺคณฺหาตีติ มหาปญฺญา, มหนฺเต สติปฏฺฐาเน ปริคฺคณฺหาตีติ มหาปญฺญา, มหนฺเต สมฺมปฺปธาเน ปริคฺคณฺหาตีติ มหาปญฺญา, มหนฺเต อิทฺธิปาเท ปริคฺคณฺหาตีติ มหาปญฺญา, มหนฺตานิ อินฺทฺริยานิ ปริคฺคณฺหาตีติ มหาปญฺญา, มหนฺตานิ พลานิ ปริคฺคณฺหาตีติ มหาปญฺญา, มหนฺเต โพชฺฌงฺเค ปริคฺคณฺหาตีติ มหาปญฺญา, มหนฺตํ อริยมคฺคํ\footnote{มหนฺเต อริยมคฺเค (ก)} ปริคฺคณฺหาตีติ มหาปญฺญา, มหนฺตานิ สามญฺญผลานิ ปริคฺคณฺหาตีติ มหาปญฺญา, มหนฺตา อภิญฺญาโย\footnote{มหาภิญฺญาโย (ก)} ปริคฺคณฺหาตีติ มหาปญฺญา, มหนฺตํ ปรมตฺถํ นิพฺพานํ ปริคฺคณฺหาตีติ มหาปญฺญา, มหาปญฺญตาย สํวตฺตนฺตีติ อยํ มหาปญฺญา. (4)}

\prose{\textbf{ปุถุปญฺญตาย สํวตฺตนฺตี}ติ กตมา ปุถุปญฺญา? ปุถุนานาขนฺเธสุ ญาณํ ปวตฺตตีติ ปุถุปญฺญา, ปุถุนานาธาตูสุ ญาณํ ปวตฺตตีติ ปุถุปญฺญา, ปุถุนานา-อายตเนสุ ญาณํ ปวตฺตตีติ ปุถุปญฺญา, ปุถุนานาปฏิจฺจสมุปฺปาเทสุ ญาณํ ปวตฺตตีติ ปุถุปญฺญา, ปุถุนานาสุญฺญตมนุปลพฺเภสุ ญาณํ ปวตฺตตีติ ปุถุปญฺญา, ปุถุนานา-อตฺเถสุ ญาณํ ปวตฺตตีติ ปุถุปญฺญา, ปุถุนานาธมฺเมสุ ญาณํ ปวตฺตตีติ ปุถุปญฺญา, ปุถุนานานิรุตฺตีสุ ญาณํ ปวตฺตตีติ ปุถุปญฺญา, ปุถุนานาปฏิภาเนสุ ญาณํ ปวตฺตตีติ ปุถุปญฺญา, ปุถุนานาสีลกฺขนฺเธสุ ญาณํ ปวตฺตตีติ ปุถุปญฺญา, ปุถุนานาสมาธิกฺขนฺเธสุ ญาณํ ปวตฺตตีติ ปุถุปญฺญา, ปุถุนานาปญฺญากฺขนฺเธสุ ญาณํ ปวตฺตตีติ ปุถุปญฺญา, ปุถุนานาวิมุตฺติกฺขนฺเธสุ ญาณํ ปวตฺตตีติ ปุถุปญฺญา.\csromanspacer{} ปุถุนานาวิมุตฺติญาณทสฺสนกฺขนฺเธสุ ญาณํ ปวตฺตตีติ ปุถุปญฺญา, ปุถุนานาฐานาฏฺฐาเนสุ ญาณํ ปวตฺตตีติ ปุถุปญฺญา, ปุถุนานาวิหารสมาปตฺตีสุ ญาณํ ปวตฺตตีติ ปุถุปญฺญา, ปุถุนานา-อริยสจฺเจสุ ญาณํ ปวตฺตตีติ ปุถุปญฺญา.\csromanspacer{} ปุถุนานาสติปฏฺฐาเนสุ ญาณํ ปวตฺตตีติ ปุถุปญฺญา, ปุถุนานาสมฺมปฺปธาเนสุ ญาณํ ปวตฺตตีติ ปุถุปญฺญา.\csromanspacer{} ปุถุนานา-อิทฺธิปาเทสุ ญาณํ ปวตฺตตีติ ปุถุปญฺญา.\csromanspacer{} ปุถุนานา-อินฺทฺริเยสุ ญาณํ ปวตฺตตีติ ปุถุปญฺญา.\csromanspacer{} ปุถุนานาพเลสุ ญาณํ ปวตฺตตีติ ปุถุปญฺญา.\csromanspacer{} ปุถุนานาโพชฺฌงฺเคสุ ญาณํ}

\csromanheader{2}{1. มหาปญฺญากถา}
\prose{\csromanfolio{373}ปวตฺตตีติ ปุถุปญฺญา.\csromanspacer{} ปุถุนานา-อริยมคฺเค ญาณํ ปวตฺตตีติ ปุถุปญฺญา, ปุถุนานาสามญฺญผเลสุ ญาณํ ปวตฺตตีติ ปุถุปญฺญา.\csromanspacer{} ปุถุนานา-อภิญฺญาสุ ญาณํ ปวตฺตตีติ ปุถุปญฺญา.\csromanspacer{} ปุถุชฺชนสาธารเณ ธมฺเม อติกฺกมฺม\footnote{สมติกฺกมฺม (สฺยา)} ปรมตฺเถ นิพฺพาเน ญาณํ ปวตฺตตีติ ปุถุปญฺญา.\csromanspacer{} ปุถุปญฺญตาย สํวตฺตนฺตีติ อยํ ปุถุปญฺญา. (5)}

\prose{\textbf{วิปุลปญฺญตาย สํวตฺตนฺตี}ติ กตมา วิปุลปญฺญา? วิปุเล อตฺเถ ปริคฺคณฺหาตีติ วิปุลปญฺญา, วิปุเล ธมฺเม ปริคฺคณฺหาตีติ วิปุลปญฺญา, วิปุลา นิรุตฺติโย ปริคฺคณฺหาตีติ วิปุลปญฺญา, วิปุลานิ ปฏิภานานิ ปริคฺคณฺหาตีติ วิปุลปญฺญา, วิปุเล สีลกฺขนฺเธ ปริคฺคณฺหาตีติ วิปุลปญฺญา, วิปุเล สมาธิกฺขนฺเธ ปริคฺคณฺหาตีติ วิปุลปญฺญา, วิปุเล ปญฺญากฺขนฺเธ ปริคฺคณฺหาตีติ วิปุลปญฺญา, วิปุเล วิมุตฺติกฺขนฺเธ ปริคฺคณฺหาตีติ วิปุลปญฺญา, วิปุเล วิมุตฺติญาณทสฺสนกฺขนฺเธ ปริคฺคณฺหาตีติ วิปุลปญฺญา, วิปุลานิ ฐานาฏฺฐานานิ ปริคฺคณฺหาตีติ วิปุลปญฺญา, วิปุลา วิหารสมาปตฺติโย ปริคฺคณฺหาตีติ วิปุลปญฺญา, วิปุลานิ อริยสจฺจานิ ปริคฺคณฺหาตีติ วิปุลปญฺญา, วิปุเล สติปฏฺฐาเน ปริคฺคณฺหาตีติ วิปุลปญฺญา, วิปุเล สมฺมปฺปธาเน ปริคฺคณฺหาตีติ วิปุลปญฺญา, วิปุเล อิทฺธิปาเท ปริคฺคณฺหาตีติ วิปุลปญฺญา, วิปุลานิ อินฺทฺริยานิ ปริคฺคณฺหาตีติ วิปุลปญฺญา, วิปุลานิ พลานิ ปริคฺคณฺหาตีติ วิปุลปญฺญา, วิปุเล โพชฺฌงฺเค ปริคฺคณฺหาตีติ วิปุลปญฺญา, วิปุลํ อริยมคฺคํ ปริคฺคณฺหาตีติ วิปุลปญฺญา, วิปุลานิ สามญฺญผลานิ ปริคฺคณฺหาตีติ วิปุลปญฺญา, วิปุลา อภิญฺญาโย ปริคฺคณฺหาตีติ วิปุลปญฺญา, วิปุลํ ปรมตฺถํ นิพฺพานํ ปริคฺคณฺหาตีติ วิปุลปญฺญา, วิปุลปญฺญตาย สํวตฺตนฺตีติ อยํ วิปุลปญฺญา. (6)}

\prose{\textbf{คมฺภีรปญฺญตาย สํวตฺตนฺตี}ติ กตมา คมฺภีรปญฺญา? คมฺภีเรสุ ขนฺเธสุ ญาณํ ปวตฺตตีติ คมฺภีรปญฺญา, คมฺภีราสุ ธาตูสุ ญาณํ ปวตฺตตีติ คมฺภีรปญฺญา, คมฺภีเรสุ อายตเนสุ ญาณํ ปวตฺตตีติ คมฺภีรปญฺญา, คมฺภีเรสุ ปฏิจฺจสมุปฺปาเทสุ ญาณํ ปวตฺตตีติ คมฺภีรปญฺญา.\csromanspacer{} คมฺภีเรสุ สุญฺญตมนุปลพฺเภสุ ญาณํ ปวตฺตตีติ คมฺภีรปญฺญา, คมฺภีเรสุ อตฺเถสุ ญาณํ ปวตฺตตีติ คมฺภีรปญฺญา, คมฺภีเรสุ ธมฺเมสุ ญาณํ ปวตฺตตีติ คมฺภีรปญฺญา, คมฺภีราสุ นิรุตฺตีสุ ญาณํ ปวตฺตตีติ คมฺภีรปญฺญา, คมฺภีเรสุ ปฏิภาเนสุ ญาณํ}

\prose{\csromanfolio{374}ปวตฺตตีติ คมฺภีรปญฺญา, คมฺภีเรสุ สีลกฺขนฺเธสุ ญาณํ ปวตฺตตีติ คมฺภีรปญฺญา, คมฺภีเรสุ สมาธิกฺขนฺเธสุ ญาณํ ปวตฺตตีติ คมฺภีรปญฺญา, คมฺภีเรสุ ปญฺญากฺขนฺเธสุ ญาณํ ปวตฺตตีติ คมฺภีรปญฺญา, คมฺภีเรสุ วิมุตฺติกฺขนฺเธสุ ญาณํ ปวตฺตตีติ คมฺภีรปญฺญา, คมฺภีเรสุ วิมุตฺติญาณทสฺสนกฺขนฺเธสุ ญาณํ ปวตฺตตีติ คมฺภีรปญฺญา, คมฺภีเรสุ ฐานาฏฺฐาเนสุ ญาณํ ปวตฺทตีติ คมฺภีรปญฺญา, คมฺภีราสุ วิหารสมาปตฺตีสุ ญาณํ ปวตฺตทีติ คมฺภีรปญฺญา, คมฺภีเรสุ อริยสจฺเจสุ ญาณํ ปวตฺตตีติ คมฺภีรปญฺญา, คมฺภีเรสุ สติปฏฺฐาเนสุ ญาณํ ปวตฺตตีติ คมฺภีรปญฺญา, คมฺภีเรสุ สมฺมปฺปธาเนสุ ญาณํ ปวตฺตตีติ คมฺภีรปญฺญา, คมฺภีเรสุ อิทฺธิปาเทสุ ญาณํ ปวตฺตตีติ คมฺภีรปญฺญา, คมฺภีเรสุ อินฺทฺริเยสุ ญาณํ ปวตฺตตีติ คมฺภีรปญฺญา, คมฺภีเรสุ พเลสุ ญาณํ ปวตฺตตีติ คมฺภีรปญฺญา, คมฺภีเรสุ โพชฺฌงฺเคสุ ญาณํ ปวตฺตตีติ คมฺภีรปญฺญา, คมฺภีเร อริยมคฺเค ญาณํ ปวตฺตตีติ คมฺภีรปญฺญา, คมฺภีเรสุ สามญฺญผเลสุ ญาณํ ปวตฺตตีติ คมฺภีรปญฺญา, คมฺภีราสุ อภิญฺญาสุ ญาณํ ปวตฺตตีติ คมฺภีรปญฺญา, คมฺภีเร ปรมตฺเถ นิพฺพาเน ญาณํ ปวตฺตตีติ คมฺภีรปญฺญา, คมฺภีรปญฺญตาย สํวตฺตนฺตีติ อยํ คมฺภีรปญฺญา. (7)}

\prose{\textbf{อสามนฺตปญฺญตาย สํวตฺตนฺตี}ติ กตมา อสามนฺตปญฺญา? ยสฺส ปุคฺคลสฺส อตฺถววตฺถานโต อตฺถปฏิสมฺภิทา อธิคตา โหติ สจฺฉิกตา ผสฺสิตา ปญฺญาย, ธมฺมววตฺถานโต ธมฺมปฏิสมฺภิทา อธิคตา โหติ สจฺฉิกตา ผสฺสิตา ปญฺญาย, นิรุตฺติววตฺถานโต นิรุตฺติปฏิสมฺภิทา อธิคตา โหติ สจฺฉิกตา ผสฺสิตา ปญฺญาย, ปฏิภานววตฺถานโต ปฏิภานปฏิสมฺภิทา อธิคตา โหติ สจฺฉิกตา ผสฺสิตา ปญฺญาย, ตสฺส อตฺเถ จ ธมฺเม จ นิรุตฺติยา จ ปฏิภาเน จ น อญฺโญ โกจิ สกฺโกติ อภิสมฺภวิตุํ, อนภิสมฺภวนีโย จ โส อญฺเญหีติ อสามนฺตปญฺโญ.}

\prose{ปุถุชฺชนกลฺยาณกสฺส ปญฺญา อฏฺฐมกสฺส ปญฺญาย ทูเร วิทูเร สุวิทูเร น สนฺติเก น สามนฺตา, ปุถุชฺชนกลฺยาณกํ อุปาทาย อฏฺฐมโก อสามนฺตปญฺโญ.\csromanspacer{} อฏฺฐมกสฺส ปญฺญา โสตาปนฺนสฺส ปญฺญาย ทูเร วิทูเร สุวิทูเร น สนฺติเก น สามนฺตา, อฏฺฐมกํ อุปาทาย โสตาปนฺโน อสามนฺตปญฺโญ.\csromanspacer{} โสตาปนฺนสฺส ปญฺญา สกทาคามิสฺส ปญฺญาย ทูเร วิทูเร สุวิทูเร น สนฺติเก น สามนฺตา, โสตาปนฺนํ อุปาทาย สกทาคามี อสามนฺตปญฺโญ.\csromanspacer{} สกทาคามิสฺส ปญฺญา อนาคามิสฺส ปญฺญาย}

\csromanheader{2}{1. มหาปญฺญากถา}
\prose{\csromanfolio{375}ทูเร วิทูเร สุวิทูเร น สนฺติเก น สามนฺตา, สกทาคามิํ อุปาทาย อนาคามี อสามนฺตปญฺโญ.\csromanspacer{} อนาคามิสฺส ปญฺญา อรหโต ปญฺญาย ทูเร วิทูเร สุวิทูเร น สนฺติเก น สามนฺตา, อนาคามิํ อุปาทาย อรหา อสามนฺตปญฺโญ.\csromanspacer{} อรหโต ปญฺญา ปจฺเจกสมฺพุทฺธสฺส\footnote{ปจฺเจกพุทฺธสฺส (สฺยา, ก) อฏฺฐกถา โอโลเกตพฺพา.} ปญฺญาย ทูเร วิทูเร สุวิทูเร น สนฺติเก น สามนฺตา, อรหนฺตํ อุปาทาย ปจฺเจกพุทฺโธ อสามนฺตปญฺโญ.\csromanspacer{} ปจฺเจกพุทฺธญฺจ สเทวกญฺจ โลกํ อุปาทาย ตถาคโต อรหํ สมฺมาสมฺพุทฺโธ อคฺโค อสามนฺตปญฺโญ.}

\proseitem{5}{ปญฺญาปเภทกุสโล ปภินฺนญาโณ อธิคตปฺปฏิสมฺภิโท จตุเวสารชฺชปฺปตฺโต ทสพลธารี ปุริสาสโภ ปุริสสีโห ปุริสนาโค ปุริสาชญฺโญ ปุริสโธรโยฺห อนนฺตญาโณ อนนฺตเตโช อนนฺตยโส อฑฺโฒ มหทฺธโน ธนวา เนตา วิเนตา อนุเนตา ปญฺญาเปตา นิชฺฌาเปตา เปกฺเขตา ปสาเทตา.\csromanspacer{} โส หิ ภควา อนุปฺปนฺนสฺส มคฺคสฺส อุปฺปาเทตา, อสญฺชาตสฺส มคฺคสฺส สญฺชเนตา\footnote{สญฺชาเนตา (สฺยา)}, อนกฺขาตสฺส มคฺคสฺส อกฺขาตา, มคฺคญฺญู มคฺควิทู มคฺคโกวิโท มคฺคานุคามี\footnote{มคฺคานุคา (สฺยา) ขุ 7. 138 ปิฏฺเฐ ปสฺสิตพฺพา.} จ ปนสฺส เอตรหิ สาวกา วิหรนฺติ ปจฺฉา สมนฺนาคตา.}

\prose{โส หิ ภควา ชานํ ชานาติ, ปสฺสํ ปสฺสติ, จกฺขุภูโต ญาณภูโต ธมฺมภูโต พฺรหฺมภูโต วตฺตา ปวตฺตา อตฺถสฺส นินฺเนตา อมตสฺส ทาตา ธมฺมสฺสามี ตถาคโต, นตฺถิ ตสฺส ภควโต อญฺญาตํ อทิฏฺฐํ อวิทิตํ อสจฺฉิกตํ อผสฺสิตํ ปญฺญาย.\csromanspacer{}\csromanspacer{}อตีตํ อนาคตํ ปจฺจุปฺปนฺนํ\footnote{อตีตานาคตปจฺจุปฺปนฺนํ (สฺยา) ขุ 7. 138 ปิฏฺเฐปิ.} อุปาทาย สพฺเพ ธมฺมา สพฺพากาเรน พุทฺธสฺส ภควโต ญาณมุเข อาปาถํ อาคจฺฉนฺติ, ยํ กิญฺจิ เนยฺยํ นาม อตฺถิ, ตํ สพฺพํ ชานิตพฺพํ\footnote{สพฺพํ ธมฺมํ ชานิตพฺพํ (ก)}, อตฺตตฺโถ วา ปรตฺโถ วา อุภยตฺโถ วา ทิฏฺฐธมฺมิโก วา อตฺโถ สมฺปรายิโก วา อตฺโถ อุตฺตาโน วา อตฺโถ คมฺภีโร วา อตฺโถ คูโฬฺห วา อตฺโถ ปฏิจฺฉนฺโน วา อตฺโถ เนยฺโย วา อตฺโถ นีโต วา อตฺโถ อนวชฺโช วา อตฺโถ นิกฺกิเลโส วา อตฺโถ โวทาโน วา อตฺโถ ปรมตฺโถ วา อตฺโถ, สพฺพํ ตํ อนฺโตพุทฺธญาเณ ปริวตฺตติ.}

\prose{\csromanfolio{376}สพฺพํ กายกมฺมํ พุทฺธสฺส ภควโต ญาณานุปริวตฺติ, สพฺพํ วจีกมฺมํ พุทฺธสฺส ภควโต ญาณานุปริวตฺติ, สพฺพํ มโนกมฺมํ พุทฺธสฺส ภควโต ญาณานุปริวตฺติ.\csromanspacer{} อตีเต พุทฺธสฺส ภควโต อปฺปฏิหตํ ญาณํ, อนาคเต พุทฺธสฺส ภควโต อปฺปฏิหตํ ญาณํ, ปจฺจุปฺปนฺเน พุทฺธสฺส ภควโต อปฺปฏิหตํ ญาณํ, ยาวตกํ เนยฺยํ, ตาวตกํ ญาณํ.\csromanspacer{} ยาวตกํ ญาณํ, ตาวตกํ เนยฺยํ, เนยฺยปริยนฺติกํ ญาณํ, ญาณปริยนฺติกํ เนยฺยํ, เนยฺยํ อติกฺกมิตฺวา ญาณํ นปฺปวตฺตติ, ญาณํ อติกฺกมิตฺวา เนยฺยปโถ นตฺถิ, อญฺญมญฺญปริยนฺตฏฺฐายิโน เต ธมฺมา, ยถา ทฺวินฺนํ สมุคฺคปฏลานํ สมฺมา ผุสิตานํ\footnote{สุผุสฺสิตานํ (สฺยา, ก)} เหฏฺฐิมํ สมุคฺคปฏลํ อุปริมํ นาติวตฺตติ, อุปริมํ สมุคฺคปฏลํ เหฏฺฐิมํ นาติวตฺตติ, อญฺญมญฺญปริยนฺตฏฺฐายิโน.\csromanspacer{} เอวเมว พุทฺธสฺส ภควโต เนยฺยญฺจ ญาณญฺจ อญฺญมญฺญปริยนฺตฏฺฐายิโน.\csromanspacer{} ยาวตกํ เนยฺยํ, ตาวตกํ ญาณํ, ยาวตกํ ญาณํ, ตาวตกํ เนยฺยํ, เนยฺยปริยนฺติกํ ญาณํ, ญาณปริยนฺติกํ เนยฺยํ, เนยฺยํ อติกฺกมิตฺวา ญาณํ นปฺปวตฺตติ, ญาณํ อติกฺกมิตฺวา เนยฺยปโถ นตฺถิ, อญฺญมญฺญปริยนฺตฏฺฐายิโน เต ธมฺมา.\csromanspacer{} สพฺพธมฺเมสุ พุทฺธสฺส ภควโต ญาณํ ปวตฺตติ.}

\prose{สพฺเพ ธมฺมา พุทฺธสฺส ภควโต อาวชฺชนปฺปฏิพทฺธา อากงฺขปฺปฏิพทฺทา มนสิการปฺปฏิพทฺธา จิตฺตุปฺปาทปฺปฏิพทฺธา, สพฺพสตฺเตสุ พุทฺธสฺส ภควโต ญาณํ ปวตฺตติ, สพฺเพสํ สตฺตานํ พุทฺโธ อาสยํ ชานาติ, อนุสยํ ชานาติ, จริตํ\footnote{จริยํ (สฺยา) ขุ 7. 139 ปิฏฺเฐ ปสฺสิตพฺพา.} ชานาติ, อธิมุตฺติํ ชานาติ, อปฺปรชกฺเข มหารชกฺเข ติกฺขินฺทฺริเย มุทินฺทฺริเย สฺวากาเร ทฺวากาเร สุวิญฺญาปเย ทุวิญฺญาปเย ภพฺพาภพฺเพ สตฺเต ปชานาติ, สเทวโก โลโก สมารโก สพฺรหฺมโก สสฺสมณพฺราหฺมณี ปชา สเทวมนุสฺสา อนฺโตพุทฺธญาเณ ปริวตฺตติ.}

\prose{ยถา เย เกจิ มจฺฉกจฺฉปา อนฺตมโส ติมิติมิงฺคลํ อุปาทาย อนฺโตมหาสมุทฺเท ปริวตฺตนฺติ, เอวเมวํ สเทวโก โลโก สมารโก สพฺรหฺมโก สสฺสมณพฺราหฺมณี ปชา สเทวมนุสฺสา อนฺโตพุทฺธญาเณ ปริวตฺตติ.\csromanspacer{}\csromanspacer{}ยถา เย เกจิ ปกฺขิโน อนฺตมโส ครุฬํ เวนเตยฺยํ อุปาทาย อากาสสฺส ปเทเส ปริวตฺตนฺติ.\csromanspacer{} เอวเมว เยปิ เต สาริปุตฺตสมา ปญฺญาย, เตปิ พุทฺธญาณสฺส ปเทเส ปริวตฺตนฺติ, พุทฺธญาณํ}

\csromanheader{2}{1. มหาปญฺญากถา}
\prose{\csromanfolio{377}เทวมนุสฺสานํ ปญฺญํ ผริตฺวา อติฆํสิตฺวา ติฏฺฐติ.\csromanspacer{} เยปิ เต ขตฺติยปณฺฑิตา พฺราหฺมณปณฺฑิตา คหปติปณฺฑิตา สมณปณฺฑิตา นิปุณา กตปรปฺปวาทา วาลเวธิรูปา, โวภินฺทนฺตา\footnote{เต ภินฺทนฺตา (สฺยา, ก)} มญฺเญ จรนฺติ ปญฺญาคเตน ทิฏฺฐิคตานิ, เต ปญฺหํ อภิสงฺขริตฺวา อภิสงฺขริตฺวา ตถาคตํ อุปสงฺกมิตฺวา ปุจฺฉนฺติ คูฬฺหานิ จ ปฏิจฺฉนฺนานิ จ, กถิตา วิสชฺชิตา จ เต ปญฺหา ภควตา โหนฺติ นิทฺทิฏฺฐการณา\footnote{นิทฺทิฏฺฐิการณา (สฺยา)}, อุปกฺขิตฺตกา จ เต ภควโต สมฺปชฺชนฺติ.\csromanspacer{} อถ โข ภควา ตตฺถ อติโรจติ ยทิทํ ปญฺญายาติ อคฺโค อสามนฺตปญฺโญ, อสามนฺตปญฺญตาย สํวตฺตนฺตีติ อยํ อสามนฺตปญฺญา. (8)}

\proseitem{6}{\textbf{ภูริปญฺญตาย สํวตฺตนฺตี}ติ กตมา ภูริปญฺญา? ราคํ อภิภุยฺยตีติ ภูริปญฺญา, อภิภวิตาติ ภูริปญฺญา.\csromanspacer{} โทสํ อภิภุยฺยตีติ ภูริปญฺญา, อภิภวิตาติ ภูริปญฺญา,.\csromanspacer{} โมหํ อภิภุยฺยตีติ ภูริปญฺญา, อภิภวิตาติ ภูริปญฺญา.\csromanspacer{} โกธํ ฯเปฯ.\csromanspacer{} อุปนาหํ.\csromanspacer{} มกฺขํ.\csromanspacer{} ปฬาสํ.\csromanspacer{} อิสฺสํ.\csromanspacer{} มจฺฉริยํ.\csromanspacer{} มายํ.\csromanspacer{} สาเฐยฺยํ.\csromanspacer{} ถมฺภํ.\csromanspacer{} สารมฺภํ.\csromanspacer{} มานํ.\csromanspacer{} อติมานํ.\csromanspacer{} มทํ.\csromanspacer{} ปมาทํ.\csromanspacer{} สพฺเพ กิเลเส.\csromanspacer{} สพฺเพ ทุจฺจริเต.\csromanspacer{} สพฺเพ อภิสงฺขาเร ฯเปฯ.\csromanspacer{} สพฺเพ ภวคามิกมฺเม อภิภุยฺยตีติ ภูริปญฺญา, อภิภวิตาติ ภูริปญฺญา.\csromanspacer{} ราโค อริ, ตํ อริํ มทฺทนิปญฺญาติ ภูริปญฺญา, โทโส อริ, ตํ อริํ มทฺทนิปญฺญาติ ภูริปญฺญา.\csromanspacer{} โมโห อริ, ตํ อริํ มทฺทนิปญฺญาติ ภูริปญฺญา.\csromanspacer{} โกโธ ฯเปฯ.\csromanspacer{} อุปนาโห.\csromanspacer{} มกฺโข.\csromanspacer{} ปฬาโส.\csromanspacer{} อิสฺสา.\csromanspacer{} มจฺฉริยํ.\csromanspacer{} มายา.\csromanspacer{} สาเฐยฺยํ.\csromanspacer{} ถมฺโภ.\csromanspacer{} สารมฺโภ.\csromanspacer{} มาโน.\csromanspacer{} อติมาโน.\csromanspacer{} มโท.\csromanspacer{} ปมาโท.\csromanspacer{} สพฺเพ กิเลสา.\csromanspacer{} สพฺเพ ทุจฺจริตา.\csromanspacer{} สพฺเพ อภิสงฺขารา ฯเปฯ.\csromanspacer{} สพฺเพ ภวคามิกมฺมา อริ, ตํ อริํ มทฺทนิปญฺญาติ ภูริปญฺญา.\csromanspacer{} ภูริ วุจฺจติ ปถวี\footnote{ปฐวี (สฺยา)}, ตาย ปถวิสมาย วิตฺถตาย วิปุลาย ปญฺญาย สมนฺนาคโตติ ภูริปญฺญา.\csromanspacer{} อปิ จ ปญฺญาย เมตํ อธิวจนํ “ภูริ เมธา ปริณายิกา”ติ ภูริปญฺญา, \textbf{ภูริปญฺญตาย สํวตฺตนฺตี}ติ อยํ ภูริปญฺญา. (9)}

\prose{\textbf{ปญฺญาพาหุลฺลาย สํวตฺตนฺตี}ติ กตมํ ปญฺญาพาหุลฺลํ? อิเธกจฺโจ ปญฺญาครุโก โหติ ปญฺญาจริโต ปญฺญาสโย ปญฺญาธิมุตฺโต ปญฺญาธโช ปญฺญาเกตุ ปญฺญาธิปเตยฺโย วิจยพหุโล ปวิจยพหุโล โอกฺขายนพหุโล สโมกฺขายนพหุโล\footnote{สมฺเปกฺขายนพหุโล (สฺยา, ก) ขุ 7. 393 ปิฏฺเฐ ปสฺสิตพฺพา.} สมฺเปกฺขายนธมฺโม วิภูตวิหารี ตจฺจริโต ตคฺครุโก ตพฺพหุโล ตนฺนินฺโน ตปฺโปโณ \csromanfolio{378}ตปฺปพฺภาโร ตทธิมุตฺโต ตทธิปเตยฺโย, ยถา คณครุโก วุจฺจติ คณพาหุลิโกติ, จีวรครุโก วุจฺจติ “จีวรพาหุลิโก”ติ, ปตฺตครุโก วุจฺจติ “ปตฺตพาหุลิโก”ติ, เสนาสนครุโก วุจฺจติ “เสนาสนพาหุลิโก”ติ.\csromanspacer{} เอวเมวํ อิเธกจฺโจ ปญฺญาครุโก โหติ ปญฺญาจริโต ปญฺญาสโย ปญฺญาธิมุตฺโต ปญฺญาธโช ปญฺญาเกตุ ปญฺญาธิปเตยฺโย วิจยพหุโล ปวิจยพหุโล โอกฺขายนพหุโล สโมกฺขายนพหุโล สมฺเปกฺขายนธมฺโม วิภูตวิหารี ตจฺจริโต ตคฺครุโก ตพฺพหุโล ตนฺนินฺโน ตปฺโปโณ ตปฺปพฺภาโร ตทธิมุตฺโต ตทธิปเตยฺโย, ปญฺญาพาหุลฺลาย สํวตฺตนฺตีติ อิทํ ปญฺญาพาหุลฺลํ. (10)}

\prose{\textbf{สีฆปญฺญตาย สํวตฺตนฺตี}ติ กตมา สีฆปญฺญา? สีฆํ สีฆํ สีลานิ ปริปูเรตีติ สีฆปญฺญา, สีฆํ สีฆํ อินฺทฺริยสํวรํ ปริปูเรตีติ สีฆปญฺญา, สีฆํ สีฆํ โภชเน มตฺตญฺญุตํ ปริปูเรตีติ สีฆปญฺญา, สีฆํ สีฆํ ชาคริยานุโยคํ ปริปูเรตีติ สีฆปญฺญา, สีฆํ สีฆํ สีลกฺขนฺธํ ปริปูเรตีติ สีฆปญฺญา, สีฆํ สีฆํ สมาธิกฺขนฺธํ ปริปูเรตีติ สีฆปญฺญา, สีฆํ สีฆํ ปญฺญากฺขนฺธํ ปริปูเรตีติ สีฆปญฺญา, สีฆํ สีฆํ วิมุตฺติกฺขนฺธํ ปริปูเรตีติ สีฆปญฺญา, สีฆํ สีฆํ วิมุตฺติญาณทสฺสนกฺขนฺธํ ปริปูเรตีติ สีฆปญฺญา, สีฆํ สีฆํ ฐานาฏฺฐานานิ ปฏิวิชฺฌตีติ สีฆปญฺญา, สีฆํ สีฆํ วิหารสมาปตฺติโย ปริปูเรตีติ สีฆปญฺญา, สีฆํ สีฆํ อริยสจฺจานิ ปฏิวิชฺฌตีติ สีฆปญฺญา, สีฆํ สีฆํ สติปฏฺฐาเน ภาเวตีติ สีฆปญฺญา, สีฆํ สีฆํ สมฺมปฺปธาเน ภาเวตีติ สีฆปญฺญา, สีฆํ สีฆํ อิทฺธิปาเท ภาเวตีติ สีฆปญฺญา, สีฆํ สีฆํ อินฺทฺริยานิ ภาเวตีติ สีฆปญฺญา, สีฆํ สีฆํ พลานิ ภาเวตีติ สีฆปญฺญา, สีฆํ สีฆํ โพชฺฌงฺเค ภาเวตีติ สีฆปญฺญา, สีฆํ สีฆํ อริยมคฺคํ ภาเวตีติ สีฆปญฺญา, สีฆํ สีฆํ สามญฺญผลานิ สจฺฉิกโรตีติ สีฆปญฺญา, สีฆํ สีฆํ อภิญฺญาโย ปฏิวิชฺฌตีติ สีฆปญฺญา, สีฆํ สีฆํ ปรมตฺถํ นิพฺพานํ สจฺฉิกโรตีติ สีฆปญฺญา, สีฆปญฺญตาย สํวตฺตนฺตีติ อยํ สีฆปญฺญา. (11)}

\prose{\textbf{ลหุปญฺญตาย สํวตฺตนฺตี}ติ กตมา ลหุปญฺญา? ลหุํ ลหุํ สีลานิ ปริปูเรตีติ ลหุปญฺญา, ลหุํ ลหุํ อินฺทฺริยสํวรํ ปริปูเรตีติ ลหุปญฺญา, ลหุํ ลหุํ โภชเน มตฺตญฺญุตํ ปริปูเรตีติ ลหุปญฺญา, ลหุํ ลหุํ ชาคริยานุโยคํ ปริปูเรตีติ ลหุปญฺญา, ลหุํ ลหุํ สีลกฺขนฺธํ ฯเปฯ}

\csromanheader{2}{1. มหาปญฺญากถา}
\prose{\csromanfolio{379}สมาธิกฺขนฺธํ.\csromanspacer{} ปญฺญากฺขนฺธํ.\csromanspacer{} วิมุตฺติกฺขนฺธํ.\csromanspacer{} วิมุตฺติญาณทสฺสนกฺขนฺธํ ปริปูเรตีติ ลหุปญฺญา, ลหุํ ลหุํ ฐานาฏฺฐานานิ ปฏิวิชฺฌตีติ ลหุปญฺญา, ลหุํ ลหุํ วิหารสมาปตฺติโย ปริปูเรตีติ ลหุปญฺญา, ลหุํ ลหุํ อริยสจฺจานิ ปฏิวิชฺฌตีติ ลหุปญฺญา, ลหุํ ลหุํ สติปฏฺฐาเน ภาเวตีติ ลหุปญฺญา, ลหุํ ลหุํ สมฺมปฺปธาเน ภาเวตีติ ลหุปญฺญา, ลหุํ ลหุํ อิทฺธิปาเท ภาเวตีติ ลหุปญฺญา, ลหุํ ลหุํ อินฺทฺริยานิ ภาเวตีติ ลหุปญฺญา, ลหุํ ลหุํ พลานิ ภาเวตีติ ลหุปญฺญา, ลหุํ ลหุํ โพชฺฌงฺเค ภาเวตีติ ลหุปญฺญา, ลหุํ ลหุํ อริยมคฺคํ ภาเวตีติ ลหุปญฺญา, ลหุํ ลหุํ สามญฺญผลานิ สจฺฉิกโรตีติ ลหุปญฺญา, ลหุํ ลหุํ อภิญฺญาโย ปฏิวิชฺฌตีติ ลหุปญฺญา, ลหุํ ลหุํ ปรมตฺถํ นิพฺพานํ สจฺฉิกโรตีติ ลหุปญฺญา, ลหุปญฺญตาย สํวตฺตนฺตีติ อยํ ลหุปญฺญา. (12)}

\prose{\textbf{หาสปญฺญตาย สํวตฺตนฺตี}ติ กตมา หาสปญฺญา? อิเธกจฺโจ หาสพหุโล เวทพหุโล ตุฏฺฐิพหุโล ปาโมชฺชพหุโล สีลานิ ปริปูเรตีติ หาสปญฺญา, หาสพหุโล เวทพหุโล ตุฏฺฐิพหุโล ปาโมชฺชพหุโล อินฺทฺริยสํวรํ ปริปูเรตีติ หาสปญฺญา, หาสพหุโล เวทพหุโล ตุฏฺฐิพหุโล ปาโมชฺชพหุโล โภชเน มตฺตญฺญุตํ ปริปูเรตีติ หาสปญฺญา, หาสพหุโล เวทพหุโล ตุฏฺฐิพหุโล ปาโมชฺชพหุโล ชาคริยานุโยคํ ปริปูเรตีติ หาสปญฺญา, หาสพหุโล เวทพหุโล ตุฏฺฐิพหุโล ปาโมชฺชพหุโล สีลกฺขนฺธํ ฯเปฯ สมาธิกฺขนฺธํ.\csromanspacer{} ปญฺญากฺขนฺธํ.\csromanspacer{} วิมุตฺติกฺขนฺธํ.\csromanspacer{} วิมุตฺติญาณทสฺสนกฺขนฺธํ ปริปูเรตีติ ฯเปฯ ฐานาฏฺฐานานิ ปฏิวิชฺฌตีติ.\csromanspacer{} วิหารสมาปตฺติโย ปริปูเรตีติ.\csromanspacer{} อริยสจฺจานิ ปฏิวิชฺฌตีติ.\csromanspacer{} สติปฏฺฐาเน ภาเวตีติ.\csromanspacer{} สมฺมปฺปธาเน ภาเวตีติ.\csromanspacer{} อิทฺธิปาเท ภาเวตีติ.\csromanspacer{} อินฺทฺริยานิ ภาเวตีติ.\csromanspacer{} พลานิ ภาเวตีติ.\csromanspacer{} โพชฺฌงฺเค ภาเวตีติ.\csromanspacer{} อริยมคฺคํ ภาเวตีติ ฯเปฯ สามญฺญผลานิ สจฺฉิกโรตีติ หาสปญฺญา, หาสพหุโล เวทพหุโล ตุฏฺฐิพหุโล ปาโมชฺชพหุโล อภิญฺญาโย ปฏิวิชฺฌตีติ หาสปญฺญา, หาสพหุโล เวทพหุโล ตุฏฺฐิพหุโล ปาโมชฺชพหุโล ปรมตฺถํ นิพฺพานํ สจฺฉิกโรตีติ หาสปญฺญา, หาสปญฺญตาย สํวตฺตนฺตีติ อยํ หาสปญฺญา. (13)}

\proseitem{7}{\textbf{ชวนปญฺญตาย สํวตฺตนฺตี}ติ กตมา ชวนปญฺญา? ยํ กิญฺจิ รูปํ อตีตานาคตปจฺจุปฺปนฺนํ อชฺฌตฺตํ วา พหิทฺธา วา โอฬาริกํ วา สุขุมํ วา หีนํ วา \csromanfolio{380}ปณีตํ วา ยํ ทูเร สนฺติเก วา, สพฺพํ รูปํ อนิจฺจโต ขิปฺปํ ชวตีติ ชวนปญฺญา, ทุกฺขโต ขปฺปํ ชวตีติ ชวนปญฺญา, อนตฺตโต ขิปฺปํ ชวตีติ ชวนปญฺญา.\csromanspacer{} ยา กาจิ เวทนา ฯเปฯ.\csromanspacer{} ยา กาจิ สญฺญา.\csromanspacer{} เย เกจิ สงฺขารา.\csromanspacer{} ยํ กิญฺจิ วิญฺญาณํ อตีตานาคตปจฺจุปฺปนฺนํ อชฺฌตฺตํ วา พหิทฺธา วา โอฬาริกํ วา สุขุมํ วา หีนํ วา ปณีตํ วา ยํ ทูเร สนฺติเก วา, สพฺพํ วิญฺญาณํ อนิจฺจโต ขิปฺปํ ชวตีติ ชวนปญฺญา, ทุกฺขโต ขิปฺปํ ชวตีติ ชวนปญฺญา, อนตฺตโต ขิปฺปํ ชวตีติ ชวนปญฺญา.\csromanspacer{} จกฺขุ ฯเปฯ.\csromanspacer{} ชรามรณํ อตีตานาคตปจฺจุปฺปนฺนํ อนิจฺจโต ขิปฺปํ ชวตีติ ชวนปญฺญา, ทุกฺขโต ขิปฺปํ ชวตีติ ชวนปญฺญา, อนตฺตโต ขิปฺปํ ชวตีติ ชวนปญฺญา.}

\prose{รูปํ อตีตานาคตปจฺจุปฺปนฺนํ อนิจฺจํ ขยฏฺเฐน ทุกฺขํ ภยฏฺเฐน อนตฺตา อสารกฏฺเฐนาติ ตุลยิตฺวา ตีรยิตฺวา วิภาวยิตฺวา วิภูตํ กตฺวา รูปนิโรเธ นิพฺพาเน ขิปฺปํ ชวตีติ ชวนปญฺญา.\csromanspacer{}\csromanspacer{}เวทนา ฯเปฯ.\csromanspacer{} สญฺญา.\csromanspacer{} สงฺขารา.\csromanspacer{} วิญฺญาณํ.\csromanspacer{} จกฺขุ ฯเปฯ.\csromanspacer{} ชรามรณํ อตีตานาคตปจฺจุปฺปนฺนํ อนิจฺจํ ขยฏฺเฐน ทุกฺขํ ภยฏฺเฐน อนตฺตา อสารกฏฺเฐนาติ ตุลยิตฺวา ตีรยิตฺวา วิภาวยิตฺวา วิภูตํ กตฺวา ชารามรณนิโรเธ\footnote{ชรามรณํ รูปนิโรเธ (สฺยา)} นิพฺพาเน ขิปฺปํ ชวตีติ ชวนปญฺญา.}

\prose{รูปํ อตีตานาคตปจฺจุปฺปนฺนํ อนิจฺจํ สงฺขตํ ปฏิจฺจสมุปฺปนฺนํ ขยธมฺมํ วยธมฺมํ วิราคธมฺมํ นิโรธธมฺมนฺติ ตุลยิตฺวา ตีรยิตฺวา วิภาวยิตฺวา วิภูตํ กตฺวา รูปนิโรเธ นิพฺพาเน ขิปฺปํ ชวตีติ ชวนปญฺญา.\csromanspacer{}\csromanspacer{}เวทนา ฯเปฯ.\csromanspacer{} สญฺญา.\csromanspacer{} สงฺขารา.\csromanspacer{} วิญฺญาณํ.\csromanspacer{} จกฺขุ ฯเปฯ.\csromanspacer{} ชรามรณํ อตีตานาคตปจฺจุปฺปนฺนํ อนิจฺจํ สงฺขตํ ปฏิจฺจสมุปฺปนฺนํ ขยธมฺมํ วยธมฺมํ วิราคธมฺมํ นิโรธธมฺมนฺติ ตุลยิตฺวา ตีรยิตฺวา วิภาวยิตฺวา วิภูตํ กตฺวา ชรามรณนิโรเธ นิพฺพาเน ขิปฺปํ ชวตีติ ชวนปญฺญา, ชวนปญฺญตาย สํวตฺตนฺตีติ อยํ ชวนปญฺญา. (14)}

\prose{\textbf{ติกฺขปญฺญตาย สํวตฺตนฺตี}ติ กตมา ติกฺขปญฺญา? ขิปฺปํ กิเลเส ฉินฺทตีติ ติกฺขปญฺญา, อุปฺปนฺนํ กามวิตกฺกํ นาธิวาเสติ ปชหติ วิโนเทติ พฺยนฺตีกโรติ\footnote{พฺยนฺติํกโรติ (ก)} อนภาวํ คเมตีติ ติกฺขปญฺญา, อุปฺปนฺนํ พฺยาปาทวิตกฺกํ นาธิวาเสติ ปชหติ วิโนเทติ พฺยนฺตีกโรติ อนภาวํ คเมตีติ ติกฺขปญฺญา, อุปฺปนฺนํ วิหิํสาวิตกฺกํ นาธิวาเสติ ฯเปฯ อุปฺปนฺนุปฺปนฺเน ปาปเก อกุสเล ธมฺเม นาธิวาเสติ ปชหติ วิโนเทติ พฺยนฺตีกโรติ}

\csromanheader{2}{1. มหาปญฺญากถา}
\prose{\csromanfolio{381}อนภาวํ คเมตีติ ติกฺขปญฺญา, อุปฺปนฺนํ ราคํ นาธิวาเสติ ปชหติ วิโนเทติ พฺยนฺตีกโรติ อนภาวํ คเมตีติ ติกฺขปญฺญา.\csromanspacer{} อุปฺปนฺนํ โทสํ ฯเปฯ.\csromanspacer{} อุปฺปนฺนํ โมหํ.\csromanspacer{} อุปฺปนฺนํ โกธํ.\csromanspacer{} อุปฺปนฺนํ อุปนาหํ.\csromanspacer{} มกฺขํ.\csromanspacer{} ปฬาสํ.\csromanspacer{} อิสฺสํ.\csromanspacer{} มจฺฉริยํ.\csromanspacer{} มายํ.\csromanspacer{} สาเฐยฺยํ.\csromanspacer{} ถมฺภํ.\csromanspacer{} สารมฺภํ.\csromanspacer{} มานํ.\csromanspacer{} อติมานํ.\csromanspacer{} มทํ.\csromanspacer{} ปมาทํ.\csromanspacer{} สพฺเพ กิเลเส.\csromanspacer{} สพฺเพ ทุจฺจริเต.\csromanspacer{} สพฺเพ อภิสงฺขาเร ฯเปฯ.\csromanspacer{} สพฺเพ ภวคามิกมฺเม นาธิวาเสติ ปชหติ วิโนเทติ พฺยนฺตีกโรติ อนภาวํ คเมตีติ ติกฺขปญฺญา.\csromanspacer{} เอกสฺมิํ อาสเน จตฺตาโร จ อริยมคฺคา จตฺตาริ จ สามญฺญผลานิ จตสฺโส ปฏิสมฺภิทาโย ฉ อภิญฺญาโย อธิคตา โหนฺติ สจฺฉิกตา ผสฺสิตา ปญฺญายาติ ติกฺขปญฺญา, ติกฺขปญฺญตาย สํวตฺตนฺตีติ อยํ ติกฺขปญฺญา. (15)}

\prose{\textbf{นิพฺเพธิกปญฺญตาย สํวตฺตนฺตี}ติ กตมา นิพฺเพธิกปญฺญา? อิเธกจฺโจ สพฺพสงฺขาเรสุ อุพฺเพคพหุโล โหติ อุตฺตาสพหุโล อุกฺกณฺฐนพหุโล อรติพหุโล อนภิรติพหุโล พหิมุโข น รมติ สพฺพสงฺขาเรสุ, อนิพฺพิทฺธปุพฺพํ อปฺปทาลิตปุพฺพํ โลภกฺขนฺธํ นิพฺพิชฺฌติ ปทาเลตีติ นิพฺเพธิกปญฺญา, อนิพฺพิทฺธปุพฺพํ อปฺปทาลิตปุพฺพํ โทสกฺขนฺธํ นิพฺพิชฺฌติ ปทาเลตีติ นิพฺเพธิกปญฺญา, อนิพฺพิทฺธปุพฺพํ อปฺปทาลิตปุพฺพํ โมหกฺขนฺธํ นิพฺพิชฺฌติ ปทาเลตีติ นิพฺเพธิกปญฺญา, อนิพฺพิทฺธปุพฺพํ อปฺปทาลิตปุพฺพํ โกธํ ฯเปฯ.\csromanspacer{} อุปนาหํ.\csromanspacer{} มกฺขํ.\csromanspacer{} ปฬาสํ.\csromanspacer{} อิสฺสํ.\csromanspacer{} มจฺฉริยํ.\csromanspacer{} มายํ.\csromanspacer{} สาเฐยฺยํ.\csromanspacer{} ถมฺภํ.\csromanspacer{} สารมฺภํ.\csromanspacer{} มานํ.\csromanspacer{} อติมานํ.\csromanspacer{} มทํ.\csromanspacer{} ปมาทํ.\csromanspacer{} สพฺเพ กิเลเส.\csromanspacer{} สพฺเพ ทุจฺจริเต.\csromanspacer{} สพฺเพ อภิสงฺขาเร ฯเปฯ.\csromanspacer{} สพฺเพ ภวคามิกมฺเม นิพฺพิชฺฌติ ปทาเลตีติ นิพฺเพธิกปญฺญา, นิพฺเพธิกปญฺญตาย สํวตฺตนฺตีติ อยํ นิพฺเพธิกปญฺญา. (16)}

\prose{อิมา โสฬส ปญฺญาโย.\csromanspacer{} อิมาหิ โสฬสหิ ปญฺญาหิ สมนฺนาคโต ปุคฺคโล ปฏิสมฺภิทปฺปตฺโต.}

\csromanmatikaanchor{mtk.121}
\csromantocmarkref{subsection}{2. ปุคฺคลวิเสสนิทฺเทส}{mtk.121}
\bookmark[dest={mtk.121},level=2]{2. ปุคฺคลวิเสสนิทฺเทส}
\csromanheader{2}{2. ปุคฺคลวิเสสนิทฺเทส}
\proseitem{8}{เทฺว ปุคฺคลา ปฏิสมฺภิทปฺปตฺตา.\csromanspacer{} เอโก ปุพฺพโยคสมฺปนฺโน, เอโก น ปุพฺพโยคสมฺปนฺโน.\csromanspacer{} โย ปุพฺพโยคสมฺปนฺโน, โส เตน อติเรโก โหติ, อธิโก โหติ, วิเสโส โหติ, ตสฺส ญาณํ ปภิชฺชติ.}

\prose{เทฺว ปุคฺคลา ปฏิสมฺภิทปฺปตฺตา.\csromanspacer{} เทฺวปิ ปุพฺพโยคสมฺปนฺนา.\csromanspacer{} เอโก พหุสฺสุโต, เอโก น พหุสฺสุโต.\csromanspacer{} โย พหุสฺสุโต, โส เตน \csromanfolio{382}อติเรโก โหติ, อธิโก โหติ, วิเสโส โหติ, ตสฺส ญาณํ ปภิชฺชติ.}

\prose{เทฺว ปุคฺคลา ปฏิสมฺภิทปฺปตฺตา.\csromanspacer{} เทฺวปิ ปุพฺพโยคสมฺปนฺนา.\csromanspacer{} เทฺวปิ พหุสฺสุตา.\csromanspacer{} เอโก เทสนาพหุโล, เอโก น เทสนาพหุโล.โย เทสนาพหุโล, โส เตน อติเรโก โหติ, อธิโก โหติ, วิเสโส โหติ, ตสฺส ญาณํ ปภิชฺชติ.}

\prose{เทฺว ปุคฺคลา ปฏิสมฺภิทปฺปตฺตา.\csromanspacer{} เทฺวปิ ปุพฺพโยคสมฺปนฺนา.\csromanspacer{} เทฺวปิ พหุสฺสุตา.\csromanspacer{} เทฺวปิ เทสนาพหุลา.\csromanspacer{} เอโก ครูปนิสฺสิโต, เอโก น ครูปนิสฺสิโต.\csromanspacer{} โย ครูปนิสฺสิโต, โส เตน อติเรโก โหติ, อธิโก โหติ, วิเสโส โหติ, ตสฺส ญาณํ ปภิชฺชติ.}

\prose{เทฺว ปุคฺคลา ปฏิสมฺภิทปฺปตฺตา.\csromanspacer{} เทฺวปิ ปุพฺพโยคสมฺปนฺนา.\csromanspacer{} เทฺวปิ พหุสฺสุตา.\csromanspacer{} เทฺวปิ เทสนาพหุลา.\csromanspacer{} เทฺวปิ ครูปนิสฺสิตา.\csromanspacer{} เอโก วิหารพหุโล, เอโก น วิหารพหุโล.\csromanspacer{} โย วิหารพหุโล, โส เตน อติเรโก โหติ, อธิโก โหติ, วิเสโส โหติ, ตสฺส ญาณํ ปภิชฺชติ.}

\prose{เทฺว ปุคฺคลา ปฏิสมฺภิทปฺปตฺตา.\csromanspacer{} เทฺวปิ ปุพฺพโยคสมฺปนฺนา.\csromanspacer{} เทฺวปิ พหุสฺสุตา.\csromanspacer{} เทฺวปิ เทสนาพหุลา.\csromanspacer{} เทฺวปิ ครูปนิสฺสิตา.\csromanspacer{} เทฺวปิ วิหารพหุลา.\csromanspacer{} เอโก ปจฺจเวกฺขณาพหุโล, เอโก น ปจฺจเวกฺขณาพหุโล.\csromanspacer{} โย ปจฺจเวกฺขณาพหุโล, โส เตน อติเรโก โหติ, อธิโก โหติ, วิเสโส โหติ, ตสฺส ญาณํ ปภิชฺชติ.}

\prose{เทฺว ปุคฺคลา ปฏิสมฺภิทปฺปตฺตา.\csromanspacer{} เทฺวปิ ปุพฺพโยคสมฺปนฺนา.\csromanspacer{} เทฺวปิ พหุสฺสุตา.\csromanspacer{} เทฺวปิ เทสนาพหุลา.\csromanspacer{} เทฺวปิ ครูปนิสฺสิตา.\csromanspacer{} เทฺวปิ วิหารพหุลา.\csromanspacer{} เทฺวปิ ปจฺจเวกฺขณาพหุลา.\csromanspacer{} เอโก เสขปฏิสมฺภิทปฺปตฺโต, เอโก อเสขปฏิสมฺภิทปฺปตฺโต.\csromanspacer{} โย อเสขปฏิสมฺภิทปฺปตฺโต, โส เตน อติเรโก โหติ, อธิโก โหติ, วิเสโส โหติ, ตสฺส ญาณํ ปภิชฺชติ.}

\prose{เทฺว ปุคฺคลา ปฏิสมฺภิทปฺปตฺตา.\csromanspacer{} เทฺวปิ ปุพฺพโยคสมฺปนฺนา.\csromanspacer{} เทฺวปิ พหุสฺสุตา.\csromanspacer{} เทฺวปิ เทสนาพหุลา.\csromanspacer{} เทฺวปิ ครูปนิสฺสิตา.\csromanspacer{} เทฺวปิ วิหารพหุลา.\csromanspacer{} เทฺวปิ ปจฺจเวกฺขณาพหุลา.\csromanspacer{} เทฺวปิ อเสขปฏิสมฺภิทปฺปตฺตา.\csromanspacer{} เอโก สาวกปารมิปฺปตฺโต, เอโก น สาวกปารมิปฺปตฺโต.\csromanspacer{} โย สาวกปารมิปฺปตฺโต,}

\csromanmatikaanchor{mtk.122}
\csromantocmarkref{section}{2. อิทฺธิกถา}{mtk.122}
\bookmark[dest={mtk.122},level=1]{2. อิทฺธิกถา}
\csromanheader{1}{2. อิทฺธิกถา}
\prose{\csromanfolio{383}โส เตน อติเรโก โหติ, อธิโก โหติ, วิเสโส โหติ, ตสฺส ญาณํ ปภิชฺชติ.}

\prose{เทฺว ปุคฺคลา ปฏิสมฺภิทปฺปตฺตา.\csromanspacer{} เทฺวปิ ปุพฺพโยคสมฺปนฺนา.\csromanspacer{} เทฺวปิ พหุสฺสุตา.\csromanspacer{} เทฺวปิ เทสนาพหุลา.\csromanspacer{} เทฺวปิ ครูปนิสฺสิตา.\csromanspacer{} เทฺวปิ วิหารพหุลา.\csromanspacer{} เทฺวปิ ปจฺจเวกฺขณาพหุลา.\csromanspacer{} เทฺวปิ อเสขปฏิสมฺภิทปฺปตฺตา.\csromanspacer{} เอโก สาวกปารมิปฺปตฺโต, เอโก ปจฺเจกสมฺพุทฺโธ.\csromanspacer{} โย ปจฺเจกสมฺพุทฺโธ, โส เตน อติเรโก โหติ, อธิโก โหติ, วิเสโส โหติ, ตสฺส ญาณํ ปภิชฺชติ.}

\prose{ปจฺเจกพุทฺธญฺจ สเทวกญฺจ โลกํ อุปาทาย ตถาคโต อรหํ สมฺมาสมฺพุทฺโธ อคฺโค ปฏิสมฺภิทปฺปตฺโต ปญฺญาปเภทกุสโล ปภินฺนญาโณ อธิคตปฏิสมฺภิโท จตุเวสารชฺชปฺปตฺโต ทสพลธารี ปุริสาสโภ ปุริสสีโห ฯเปฯ.\csromanspacer{} เยปิ เต ขตฺติยปณฺฑิตา พฺราหฺมณปณฺฑิตา คหปติปณฺฑิตา สมณปณฺฑิตา นิปุณา กตปรปฺปวาทา วาลเวธิรูปา โวภินฺทนฺตา มญฺเญ จรนฺติ ปญฺญาคเตน ทิฏฺฐิคตานิ.\csromanspacer{} เต ปญฺหํ อภิสงฺขริตฺวา อภิสงฺขาริตฺวา ตถาคตํ อุปสงฺกมิตฺวา ปุจฺฉนฺติ คูฬฺหานิ จ ปฏิจฺฉนฺนานิ จ, กถิตา วิสชฺชิตา จ เต ปญฺหา ภควตา โหนฺติ นิทฺทิฏฺฐการณา, อุปกฺขิตฺตกา จ เต ภควโต สมฺปชฺชนฺติ.\csromanspacer{} อถ โข ภควา ตตฺถ อติโรจติ ยทิทํ ปญฺญายาติ อคฺโค ปฏิสมฺภิทปฺปตฺโตติ.}

\nitthitamruled{มหาปญฺญากถา นิฏฺฐิตา.}
\csromanheader{2}{2. อิทฺธิกถา}
\proseitem{9}{กา อิทฺธิ.\csromanspacer{} \textbf{กติ อิทฺธิโย}.\csromanspacer{} \textbf{อิทฺธิยา กติ ภูมิโย}.\csromanspacer{} กติ ปาทา.\csromanspacer{} กติ ปทานิ.\csromanspacer{} กติ มูลานิ? \textbf{กา อิทฺธี}ติ อิชฺฌนฏฺเฐน อิทฺธิ.\csromanspacer{}\csromanspacer{}\textbf{กติ อิทฺธิโย}ติ ทส อิทฺธิโย.\csromanspacer{}\csromanspacer{}\textbf{อิทฺธิยา กติ ภูมิโย}ติ อิทฺธิยา จตสฺโส ภูมิโย.\csromanspacer{}\csromanspacer{}จตฺตาโร ปาทา.\csromanspacer{}\csromanspacer{}อฏฺฐ ปทานิ.\csromanspacer{}\csromanspacer{}โสฬส มูลานิ.}

\proseitem{10}{กตมา \textbf{ทส อิทฺธิโย?} อธิฏฺฐานา อิทฺธิ, วิกุพฺพนา อิทฺธิ, มโนมยา อิทฺธิ ญาณวิปฺผารา อิทฺธิ, สมาธิวิปฺผารา อิทฺธิ, อริยา อิทฺธิ, กมฺมวิปากชา อิทฺธิ, \csromanfolio{384}ปุญฺญวโต อิทฺธิ, วิชฺชามยา อิทฺธิ, ตตฺถ ตตฺถ สมฺมาปโยคปจฺจยา\footnote{สมฺมปฺปโยคปจฺจยา (สฺยา, ก)} อิชฺฌนฏฺเฐน อิทฺธิ.}

\prose{\textbf{อิทฺธิยา กตมา จตสฺโส ภูมิโย?} วิเวกชาภูมิ ปฐมํ ฌานํ, ปีติสุขภูมิ ทุติยํ ฌานํ, อุเปกฺขาสุขภูมิ ตติยํ ฌานํ, อทุกฺขมสุขาภูมิ จตุตฺถํ ฌานํ.\csromanspacer{} อิทฺธิยา อิมา จตสฺโส ภูมิโย อิทฺธิลาภาย อิทฺธิปฏิลาภาย อิทฺธิวิกุพฺพนตาย อิทฺธิวิสวิตาย อิทฺธิวสีภาวาย อิทฺธิเวสารชฺชาย สํวตฺตนฺตีติ.}

\prose{\textbf{อิทฺธิยา กตเม จตฺตาโร ปาทา?} อิธ ภิกฺขุ ฉนฺทสมาธิปธานสงฺขารสมนฺนาคตํ อิทฺธิปาทํ ภาเวติ, จิตฺตสมาธิปธานสงฺขารสมนฺนาคตํ อิทฺธิปาทํ ภาเวติ, วีริยสมาธิปธานสงฺขารสมนฺนาคตํ อิทฺธิปาทํ ภาเวติ, วีมํสาสมาธิปธานสงฺขารสมนฺนาคตํ อิทฺธิปาทํ ภาเวติ.\csromanspacer{} อิทฺธิยา อิเม จตฺตาโร ปาทา อิทฺธิลาภาย อิทฺธิปฏิลาภาย อิทฺธิวิกุพฺพนตาย อิทฺธิวิสวิตาย อิทฺธิวสีภาวาย อิทฺธิเวสารชฺชาย สํวตฺตนฺตีติ.}

\prose{\textbf{อิทฺธิยา กตมานิ อฏฺฐ ปทานิ?} ฉนฺทํ เจ ภิกฺขุ นิสฺสาย ลภติ สมาธิํ, ลภติ จิตฺตสฺส เอกคฺคตํ, ฉนฺโท น สมาธิ, สมาธิ น ฉนฺโท.\csromanspacer{} อญฺโญ ฉนฺโท, อญฺโญ สมาธิ.\csromanspacer{} วีริยํ เจ ภิกฺขุ นิสฺสาย ลภติ สมาธิํ, ลภติ จิตฺตสฺส เอกคฺคตํ.\csromanspacer{} วีริยํ น สมาธิ, สมาธิ น วีริยํ.\csromanspacer{} อญฺญํ วีริยํ, อญฺโญ สมาธิ.\csromanspacer{} จิตฺตํ เจ ภิกฺขุ นิสฺสาย ลภติ สมาธิํ, ลภติ จิตฺตสฺส เอกคฺคตํ.\csromanspacer{} จิตฺตํ น สมาธิ, สมาธิ น จิตฺตํ.\csromanspacer{} อญฺญํ จิตฺตํ, อญฺโญ สมาธิ.\csromanspacer{} วีมํสํ เจ ภิกฺขุ นิสฺสาย ลภติ สมาธิํ, ลภติ จิตฺตสฺส เอกคฺคตํ.\csromanspacer{} วีมํสา น สมาธิ, สมาธิ น วีมํสา.\csromanspacer{} อญฺญา วีมํสา, อญฺโญ สมาธิ.\csromanspacer{} อิทฺธิยา อิมานิ อฏฺฐ ปทานิ อิทฺธิลาภาย อิทฺธิปฏิลาภาย อิทฺธิวิกุพฺพนตาย อิทฺธิวิสวิตาย อิทฺธิวสีภาวาย อิทฺธิเวสารชฺชาย สํวตฺตนฺตีติ.}

\prose{\textbf{อิทฺธิยา กตมานิ โสฬส มูลานิ?} อโนนตํ\footnote{อโนณตํ (ก)} จิตฺตํ โกสชฺเช น อิญฺชตีติ อาเนญฺชํ, อนุนฺนตํ\footnote{อนุณฺณตํ (ก)} จิตฺตํ อุทฺธจฺเจ น อิญฺชตีติ อาเนญฺชํ, อนภินตํ จิตฺตํ ราเค น อิญฺชตีติ อาเนญฺชํ, อนปนตํ จิตฺตํ พฺยาปาเท น อิญฺชตีติ อาเนญฺชํ, อนิสฺสิตํ จิตฺตํ ทิฏฺฐิยา น อิญฺชตีติ อาเนญฺชํ, อปฺปฏิพทฺธํ จิตฺตํ ฉนฺทราเค น อิญฺชตีติ อาเนญฺชํ, วิปฺปมุตฺตํ จิตฺตํ กามราเค น อิญฺชตีติ อาเนญฺชํ, วิสญฺญุตฺตํ จิตฺตํ กิเลเส น อิญฺชตีติ อาเนญฺชํ, วิมริยาทิกตํ จิตฺตํ}

\csromanheader{2}{2. อิทฺธิกถา}
\prose{\csromanfolio{385}กิเลสมริยาเท น อิญฺชตีติ อาเนญฺชํ, เอกตฺตคตํ\footnote{เอกคฺคตํ (สฺยา)} จิตฺตํ นานตฺตกิเลเสหิ\footnote{นานตฺตกิเลเส (สฺยา)} น อิญฺชตีติ อาเนญฺชํ, สทฺธาย ปริคฺคหิตํ จิตฺตํ อสฺสทฺธิเย น อิญฺชตีติ อาเนญฺชํ, วีริเยน ปริคฺคหิตํ จิตฺตํ โกสชฺเช น อิญฺชตีติ อาเนญฺชํ, สติยา ปริคฺคหิตํ จิตฺตํ ปมาเท น อิญฺชตีติ อาเนญฺชํ, สมาธินา ปริคฺคหิตํ จิตฺตํ อุทฺธจฺเจ น อิญฺชตีติ อาเนญฺชํ, ปญฺญาย ปริคฺคหิตํ จิตฺตํ อวิชฺชาย น อิญฺชตีติ อาเนญฺชํ, โอภาสคตํ จิตฺตํ อวิชฺชนฺธกาเร น อิญฺชตีติ อาเนญฺชํ.\csromanspacer{} อิทฺธิยา อิมานิ โสฬส มูลานิ อิทฺธิลาภาย อิทฺธิปฏิลาภาย อิทฺธิวิกุพฺพนตาย อิทฺธิวิสวิตาย อิทฺธิวสีภาวาย อิทฺธิเวสารชฺชาย สํวตฺตนฺตีติ.}

\csromanmatikaanchor{mtk.123}
\csromantocmarkref{section}{ทส-อิทฺธินิทฺเทส}{mtk.123}
\bookmark[dest={mtk.123},level=1]{ทส-อิทฺธินิทฺเทส}
\csromanheader{1}{\textbf{ทส-อิทฺธินิทฺเทส}}
\proseitem{10}{กตมา \textbf{อธิฏฺฐานา อิทฺธิ?} อิธ ภิกฺขุ อเนกวิหิตํ อิทฺธิวิธํ ปจฺจนุโภติ, เอโกปิ หุตฺวา พหุธา โหติ, พหุธาปิ หุตฺวา เอโก โหติ, อาวิภาวํ ติโรภาวํ ติโรกุฏฺฏํ ติโรปาการํ ติโรปพฺพตํ อสชฺชมาโน คจฺฉติ เสยฺยถาปิ อากาเส.ปถวิยาปิ อุมฺมุชฺชนิมุชฺชํ กโรติ เสยฺยถาปิ อุทเก.\csromanspacer{} อุทเกปิ อภิชฺชมาเน\footnote{อภิชฺชมาโน (ก) ที 1. 73 ปิฏฺเฐ ปสฺสิตพฺพา.} คจฺฉติ เสยฺยถาปิ ปถวิยํ.\csromanspacer{} อากาเสปิ ปลฺลงฺเกน กมติ เสยฺยถาปิ ปกฺขี สกุโณ.\csromanspacer{} อิเมปิ จนฺทิมสูริเย เอวํมหิทฺธิเก เอวํมหานุภาเว ปาณินา ปรามสติ ปริมชฺชติ, ยาว พฺรหฺมโลกาปิ กาเยน วสํ วตฺเตตีติ.}

\prose{\textbf{อิธา}ติ อิมิสฺสา ทิฏฺฐิยา อิมิสฺสา ขนฺติยา อิมิสฺสา รุจิยา อิมสฺมิํ อาทาเย อิมสฺมิํ ธมฺเม อิมสฺมิํ วินเย อิมสฺมิํ ธมฺมวินเย อิมสฺมิํ ปาวจเน อิมสฺมิํ พฺรหฺมจริเย อิมสฺมิํ สตฺถุสาสเน, เตน วุจฺจติ อิธาติ.\csromanspacer{} \textbf{ภิกฺขู}ติ ปุถุชฺชนกลฺยาณโก วา โหติ ภิกฺขุ เสโข วา อรหา วา อกุปฺปธมฺโม.\csromanspacer{} \textbf{อเนกวิหิตํ อิทฺธิวิธํ ปจฺจนุโภตี}ติ นานปฺปการํ อิทฺธิวิธํ ปจฺจนุโภติ.\csromanspacer{} \textbf{เอโกปิ หุตฺวา พหุธา โหตี}ติ ปกติยา เอโก พหุกํ อาวชฺชติ, สตํ วา สหสฺสํ วา สตสหสฺสํ วา อาวชฺชติ, อาวชฺชิตฺวา ญาเณน อธิฏฺฐาติ “พหุโล โหมี”ติ\footnote{โหตูติ (ก)}, พหุโล โหติ.\csromanspacer{} ยถายสฺมา จูฬปนฺถโก\footnote{จุลฺลปนฺถโก (สฺยา)} เอโกปิ หุตฺวา พหุธา โหติ.\csromanspacer{} เอวเมวํ \csromanfolio{386}โส อิทฺธิมา เจโตวสิปฺปตฺโต เอโกปิ หุตฺวา พหุธา โหติ.\csromanspacer{} \textbf{พหุธาปิ หุตฺวา เอโก โหตี}ติ ปกติยา พหุโก เอกํ อาวชฺชติ, อาวชฺชิตฺวา ญาเณน อธิฏฺฐาติ “เอโก โหมี”ติ, เอโก โหติ.\csromanspacer{} ยถายสฺมา จูฬปนฺถโก พหุธาปิ หุตฺวา เอโก โหติ.\csromanspacer{} เอวเมวํ โส อิทฺธิมา เจโตวสิปฺปตฺโต พหุธาปิ หุตฺวา เอโก โหติ.}

\proseitem{11}{อา\textbf{วิภาว}นฺติ เกนจิ อนาวฏํ โหติ อปฺปฏิจฺฉนฺนํ วิวฏํ ปากฏํ.\csromanspacer{} \textbf{ติโรภาวนฺ}ติ เกนจิ อาวฏํ โหติ ปฏิจฺฉนฺนํ ปิหิตํ ปฏิกุชฺชิตํ.\csromanspacer{} \textbf{ติโรกุฏฺฏํ ติโรปาการํ ติโรปพฺพตํ อสชฺชมาโน คจฺฉติเสยฺยถาปิ อากาเส}ติ ปกติยา อากาสกสิณสมาปตฺติยา ลาภี โหติ, ติโรกุฏฺฏํ ติโรปาการํ ติโรปพฺพตํ อาวชฺชติ, อาวชฺชิตฺวา ญาเณน อธิฏฺฐาติ “อากาโส โหตู”ติ, อากาโส โหติ.\csromanspacer{} ติโรกุฏฺฏํ ติโรปาการํ ติโรปพฺพตํ อสชฺชมาโน คจฺฉติ.\csromanspacer{} ยถา มนุสฺสา ปกติยา อนิทฺธิมนฺโต เกนจิ อนาวเฏ อปริกฺขิตฺเต อสชฺชมานา คจฺฉนฺติ.\csromanspacer{} เอวเมวํ โส อิทฺธิมา เจโตวสิปฺปตฺโต ติโรกุฏฺฏํ ติโรปาการํ ติโรปพฺพตํ อสชฺชมาโน คจฺฉติ เสยฺยถาปิ อากาเส.}

\prose{\textbf{ปถวิยาปิ อุมฺมุชฺชนิมุชฺชํ กโรติ เสยฺยถาปิ อุทเก}ติ ปกติยา อาโปกสิณสมาปตฺติยา ลาภี โหติ.\csromanspacer{} ปถวิํ อาวชฺชติ, อาวชฺชิตฺวา ญาเณน อธิฏฺฐาติ “อุทกํ โหตู”ติ, อุทกํ โหติ.\csromanspacer{} โส ปถวิยา อุมฺมุชฺชนิมุชฺชํ กโรติ.\csromanspacer{} ยถา มนุสฺสา ปกติยา อนิทฺธิมนฺโต อุทเก อุมฺมุชฺชนิมุชฺชํ กโรนฺติ.\csromanspacer{} เอวเมวํ โส อิทฺธิมา เจโตวสิปฺปตฺโต ปถวิยา อุมฺมุชฺชนิมุชฺชํ กโรติ เสยฺยถาปิ อุทเก.}

\prose{\textbf{อุทเกปิ อภิชฺชมาเน คจฺฉติ เสยฺยถาปิ ปถวิยนฺ}ติ ปกติยา ปถวีกสิณสมาปตฺติยา ลาภี โหติ.\csromanspacer{} อุทกํ อาวชฺชติ, อาวชฺชิตฺวา ญาเณน อธิฏฺฐาติ “ปถวี โหตู”ติ, ปถวี โหติ.\csromanspacer{} โส อภิชฺชมาเน อุทเก คจฺฉติ.\csromanspacer{} ยถา มนุสฺสา ปกติยา อนิทฺธิมนฺโต อภิชฺชมานาย ปถวิยา คจฺฉนฺติ.\csromanspacer{} เอวเมวํ โส อิทฺธิมา เจโตวสิปฺปตฺโต อภิชฺชมาเน อุทเก คจฺฉติ เสยฺยถาปิ ปถวิยํ.}

\prose{\textbf{อากาเสปิ ปลฺลงฺเกน กมติ เสยฺยถาปิ ปกฺขี สกุโณ}ติ ปกติยา ปถวีกสิณสมาปตฺติยา ลาภี โหติ.\csromanspacer{} อากาสํ}

\csromanheader{2}{2. อิทฺธิกถา}
\prose{\csromanfolio{387}อาวชฺชติ, อาวชฺชิตฺวา ญาเณน อธิฏฺฐาติ “ปถวี โหตู”ติ, ปถวี โหติ.\csromanspacer{} โส อากาเส อนฺตลิกฺเข จงฺกมติปิ ติฏฺฐติปิ นิสีทติปิ เสยฺยมฺปิ กปฺเปติ.\csromanspacer{} ยถา มนุสฺสา ปกติยา อนิทฺธิมนฺโต ปถวิยา จงฺกมนฺติปิ ติฏฺฐนฺติปิ นิสีทนฺติปิ เสยฺยมฺปิ กปฺเปนฺติ.\csromanspacer{} เอวเมวํ โส อิทฺธิมา เจโตวสิปฺปตฺโต อากาเส อนฺตลิกฺเข จงฺกมติปิ ติฏฺฐติปิ นิสีทติปิ เสยฺยมฺปิ กปฺเปติ เสยฺยถาปิ ปกฺขี สกุโณ.}

\proseitem{12}{\textbf{อิเมปิ จนฺทิมสูริเย เอวํมหิทฺธิเก เอวํมหานุภาเว ปาณินา ปรามสติ ปริมชฺชตี}ติ อิธ โส อิทฺธิมา เจโตวสิปฺปตฺโต นิสินฺนโก วา นิปนฺนโก วา จนฺทิมสูริเย อาวชฺชติ, อาวชฺชิตฺวา ญาเณน อธิฏฺฐาติ “หตฺถปาเส โหตู”ติ, หตฺถปาเส โหติ.\csromanspacer{} โส นิสินฺนโก วา นิปนฺนโก วา จนฺทิมสูริเย ปาณินา อามสติ ปรามสติ ปริมชฺชติ.\csromanspacer{} ยถา มนุสฺสา ปกติยา อนิทฺธิมนฺโต กิญฺจิเทว รูปคตํ หตฺถปาเส อามสนฺติ ปรามสนฺติ ปริมชฺชนฺติ.\csromanspacer{} เอวเมวํ โส อิทฺธิมา เจโตวสิปฺปตฺโต นิสินฺนโก วา นิปนฺนโก วา จนฺทิมสูริเย ปาณินา อามสติ ปรามสติ ปริมชฺชติ.}

\prose{\textbf{ยาว พฺรหฺมโลกาปิ กาเยน วสํ วตฺเตตี}ติ สเจ โส อิทฺธิมา เจโตวสิปฺปตฺโต พฺรหฺมโลกํ คนฺตุกาโม โหติ, ทูเรปิ สนฺติเก อธิฏฺฐาติ “สนฺติเก โหตู”ติ, สนฺติเก โหติ.\csromanspacer{} สนฺติเกปิ ทูเร อธิฏฺฐาติ “ทูเร โหตู”ติ, ทูเร โหติ.\csromanspacer{} พหุกมฺปิ โถกํ อธิฏฺฐาติ “โถกํ โหตู”ติ, โถกํ โหติ.\csromanspacer{} โถกมฺปิ พหุกํ อธิฏฺฐาติ “พหุกํ โหตู”ติ, พหุกํ โหติ.\csromanspacer{} ทิพฺเพน จกฺขุนา ตสฺส พฺรหฺมุโน รูปํ ปสฺสติ, ทิพฺพาย โสตธาตุยา ตสฺส พฺรหฺมุโน สทฺทํ สุณาติ, เจโตปริยญาเณน ตสฺส พฺรหฺมุโน จิตฺตํ ปชานาติ, สเจ โส อิทฺธิมา เจโตวสิปฺปตฺโต ทิสฺสมาเนน กาเยน พฺรหฺมโลกํ คนฺตุกาโม โหติ, กายวเสน จิตฺตํ ปริณาเมติ, กายวเสน จิตฺตํ อธิฏฺฐาติ.\csromanspacer{} กายวเสน จิตฺตํ ปริณาเมตฺวา กายวเสน จิตฺตํ อธิฏฺฐหิตฺวา สุขสญฺญญฺจ ลหุสญฺญญฺจ โอกฺกมิตฺวา ทิสฺสมาเนน กาเยน พฺรหฺมโลกํ คจฺฉติ.\csromanspacer{} สเจ โส อิทฺธิมา เจโตวสิปฺปตฺโต อทิสฺสมาเนน กาเยน พฺรหฺมโลกํ คนฺตุกาโม โหติ, จิตฺตวเสน กายํ ปริณาเมติ, จิตฺตวเสน กายํ อธิฏฺฐาติ, จิตฺตวเสน กายํ ปริณาเมตฺวา จิตฺตวเสน กายํ อธิฏฺฐหิตฺวา สุขสญฺญญฺจ ลหุสญฺญญฺจ \csromanfolio{388}โอกฺกมิตฺวา อทิสฺสมาเนน กาเยน พฺรหฺมโลกํ คจฺฉติ.\csromanspacer{} โส ตสฺส พฺรหฺมุโน ปุรโต รูปิํ\footnote{รู (สฺยา, ก) ที 1. 73 ปิฏฺเฐ ปสฺสิตพฺพา.} อภินิมฺมินาติ มโนมยํ สพฺพงฺคปจฺจงฺคิํ\footnote{สพฺพงฺคปจฺจงฺคํ (สฺยา, ก)} อหีนินฺทฺริยํ, สเจ โส อิทฺธิมา จงฺกมติ, นิมฺมิโตปิ ตตฺถ จงฺกมติ.\csromanspacer{} สเจ โส อิทฺธิมา ติฏฺฐติ, นิมฺมิโตปิ ตตฺถ ติฏฺฐติ.\csromanspacer{} สเจ โส อิทฺธิมา นิสีทติ, นิมฺมิโตปิ ตตฺถ นิสีทติ.\csromanspacer{} สเจ โส อิทฺธิมา เสยฺยํ กปฺเปติ, นิมฺมิโตปิ ตตฺถ เสยฺยํ กปฺเปติ.\csromanspacer{} สเจ โส อิทฺธิมา ธูมายติ, นิมฺมิโตปิ ตตฺถ ธูมายติ.\csromanspacer{} สเจ โส อิทฺธิมา ปชฺชลติ, นิมฺมิโตปิ ตตฺถ ปชฺชลติ.\csromanspacer{} สเจ โส อิทฺธิมา ธมฺมํ ภาสติ, นิมฺมิโตปิ ตตฺถ ธมฺมํ ภาสติ.\csromanspacer{} สเจ โส อิทฺธิมา ปญฺหํ ปุจฺฉติ, นิมฺมิโตปิ ตตฺถ ปญฺหํ ปุจฺฉติ.\csromanspacer{} สเจ โส อิทฺธิมา ปญฺหํ ปุฏฺโฐ วิสชฺเชติ, นิมฺมิโตปิ ตตฺถ ปญฺหํ ปุฏฺโฐ วิสชฺเชติ.\csromanspacer{} สเจ โส อิทฺธิมา เตน พฺรหฺมุนา สทฺธิํ สนฺติฏฺฐติ สลฺลปติ สากจฺฉํ สมาปชฺชติ, นิมฺมิโตปิ ตตฺถ เตน พฺรหฺมุนา สทฺธิํ สนฺติฏฺฐติ สลฺลปติ สากจฺฉํ สมาปชฺชติ.\csromanspacer{} ยญฺญเทว โส อิทฺธิมา กโรติ, ตํตเทว หิ โส นิมฺมิโต กโรตีติ อยํ อธิฏฺฐานา อิทฺธิ. (1)}

\proseitem{13}{กตมา \textbf{วิกุพฺพนา อิทฺธิ?} สิขิสฺส ภควโต อรหโต สมฺมาสมฺพุทฺธสฺส อภิภู นาม สาวโก พฺรหฺมโลเก ฐิโต สหสฺสิโลกธาตุํ\footnote{สหสฺสีโลกธาตุํ (สฺยา) วตฺถุ สํ 1. 157 ปิฏฺเฐ อาคตํ.} สเรน วิญฺญาเปสิ.\csromanspacer{} โส ทิสฺสมาเนนปิ กาเยน ธมฺมํ เทเสสิ, อทิสฺสมาเนนปิ กาเยน ธมฺมํ เทเสสิ, ทิสฺสมาเนนปิ เหฏฺฐิเมน อุปฑฺฒกาเยน อทิสฺสมาเนนปิ อุปริเมน อุปฑฺฒกาเยน ธมฺมํ เทเสสิ.\csromanspacer{} ทิสฺสมาเนนปิ อุปริเมน อุปฑฺฒกาเยน, อทิสฺสมาเนนปิ เหฏฺฐิเมน อุปฑฺฒกาเยน ธมฺมํ เทเสสิ.\csromanspacer{} โส ปกติวณฺณํ วิชหิตฺวา กุมารกวณฺณํ วา ทสฺเสติ, นาควณฺณํ วา ทสฺเสติ, สุปณฺณวณฺณํ วา ทสฺเสติ, ยกฺขวณฺณํ วา ทสฺเสติ, อินฺทวณฺณํ วา\footnote{อสุรวณฺณํ วา ทสฺสติ, อินฺทวณฺณํ วา (สฺยา)} ทสฺเสติ, เทววณฺณํ วา ทสฺเสติ, พฺรหฺมวณฺณํ วา ทสฺเสติ, สมุทฺทวณฺณํ วา ทสฺเสติ, ปพฺพตวณฺณํ วา ทสฺเสติ, วนวณฺณํ วา ทสฺเสติ, สีหวณฺณํ วา ทสฺเสติ, พฺยคฺฆวณฺณํ วา ทสฺเสติ, ทีปิวณฺณํ วา ทสฺเสติ, หตฺถิมฺปิ ทสฺเสติ, อสฺสมฺปิ ทสฺเสติ, รถมฺปิ ทสฺเสติ, ปตฺติมฺปิ ทสฺเสติ, วิวิธมฺปิ เสนาพฺยูหํ ทสฺเสตีติ อยํ วิกุพฺพนา อิทฺธิ. (2)}

\csromanheader{2}{2. อิทฺธิกถา}
\proseitem{14}{\csromanfolio{389}กตมา \textbf{มโนมยา อิทฺธิ?} อิธ ภิกฺขุ อิมมฺหา กายา อญฺญํ กายํ อภินิมฺมินาติ รูปิํ มโนมยํ สพฺพงฺคปจฺจงฺคิํ อหีนินฺทฺริยํ, เสยฺยถาปิ ปุริโส มุญฺชมฺหา อีสิกํ ปวาเหยฺย.\csromanspacer{} ตสฺส เอวมสฺส “อยํ มุญฺโช, อยํ อีสิกา, อญฺโญ มุญฺโช, อญฺญา อีสิกา.\csromanspacer{} มุญฺชมฺหาเตฺวว อีสิกา ปวาฬฺหา”ติ.\csromanspacer{} เสยฺยถา วา ปน ปุริโส อสิํ โกสิยา ปวาเหยฺย.\csromanspacer{} ตสฺส เอวมสฺส “อยํ อสิ, อยํ โกสิ, อญฺโญ อสิ, อญฺญา โกสิ.\csromanspacer{} โกสิยาเตฺวว อสิ ปวาโฬฺห”ติ.\csromanspacer{} เสยฺยถา วา ปน ปุริโส อหิํ กรณฺฑา อุทฺธเรยฺย.\csromanspacer{} ตสฺส เอวมสฺส “อยํ อหิ, อยํ กรณฺโฑ.\csromanspacer{} อญฺโญ อหิ, อญฺโญ กรณฺโฑ.\csromanspacer{} กรณฺฑาเตฺวว อหิ อุพฺภโต”ติ.\csromanspacer{} เอวเมวํ ภิกฺขุ อิมมฺหา กายา อญฺญํ กายํ อภินิมฺมินาติ รูปิํ มโนมยํ สพฺพงฺคปจฺจงฺคิํ อหีนินฺทฺริยํ.\csromanspacer{} อยํ มโนมยา อิทฺธิ. (3)}

\proseitem{15}{กตมา \textbf{ญาณวิปฺผารา อิทฺธิ?} อนิจฺจานุปสฺสนาย นิจฺจสญฺญาย ปหานฏฺโฐ อิชฺฌตีติ ญาณวิปฺผารา อิทฺธิ.\csromanspacer{} ทุกฺขานุปสฺสนาย สุขสญฺญาย.\csromanspacer{} อนตฺตานุปสฺสนาย อตฺตสญฺญาย.\csromanspacer{} นิพฺพิทานุปสฺสนาย นนฺทิยา.\csromanspacer{} วิราคานุปสฺสนาย ราคสฺส.\csromanspacer{} นิโรธานุปสฺสนาย สมุทยสฺส.\csromanspacer{} ปฏินิสฺสคฺคานุปสฺสนาย อาทานสฺส ปหานฏฺโฐ อิชฺฌตีติ ญาณวิปฺผารา อิทฺธิ.\csromanspacer{} อายสฺมโต พากุลสฺส ญาณวิปฺผารา อิทฺธิ.\csromanspacer{} อายสฺมโต สํกิจฺจสฺส ญาณวิปฺผารา อิทฺธิ.\csromanspacer{} อายสฺมโต ภูตปาลสฺส ญาณวิปฺผารา อิทฺธิ.\csromanspacer{} อยํ ญาณวิปฺผารา อิทฺธิ. (4)}

\proseitem{16}{กตมา \textbf{สมาธิวิปฺผารา อิทฺธิ?} ปฐเมน ฌาเนน นีวรณานํ ปหานฏฺโฐ อิชฺฌตีติ สมาธิวิปฺผารา อิทฺธิ.\csromanspacer{} ทุติเยน ฌาเนน วิตกฺกวิจารานํ ปหานฏฺโฐ อิชฺฌตีติ สมาธิวิปฺผารา อิทฺธิ.\csromanspacer{} ตติเยน ฌาเนน ปีติยา ปหานฏฺโฐ อิชฺฌตีติ ฯเปฯ.\csromanspacer{} จตุตฺเถน ฌาเนน สุขทุกฺขานํ ปหานฏฺโฐ อิชฺฌตีติ ฯเปฯ.\csromanspacer{} อากาสานญฺจายตนสมาปตฺติยา รูปสญฺญาย ปฏิฆสญฺญาย นานตฺตสญฺญาย ปหานฏฺโฐ อิชฺฌตีติ ฯเปฯ.\csromanspacer{} วิญฺญาณญฺจายตนสมาปตฺติยา อากาสานญฺจายตนสญฺญาย ปหานฏฺโฐ อิชฺฌตีติ ฯเปฯ.\csromanspacer{} อากิญฺจญฺญายตนสมาปตฺติยา วิญฺญาณญฺจายตนสญฺญาย ปหานฏฺโฐ อิชฺฌตีติ ฯเปฯ.\csromanspacer{} เนวสญฺญานาสญฺญายตนสมาปตฺติยา อากิญฺจญฺญายตนสญฺญาย ปหานฏฺโฐ อิชฺฌตีติ สมาธิวิปฺผารา อิทฺธิ, อายสฺมโต สาริปุตฺตสฺส สมาธิวิปฺผารา อิทฺธิ, อายสฺมโต สญฺชีวสฺส สมาธิวิปฺผารา \csromanfolio{390}อิทฺธิ, อายสฺมโต ขาณุโกณฺฑญฺญสฺส สมาธิวิปฺผารา อิทฺธิ, อุตฺตราย อุปาสิกาย สมาธิวิปฺผารา อิทฺธิ, สามาวติยา อุปาสิกาย สมาธิวิปฺผารา อิทฺธิ, อยํ สมาธิวิปฺผารา อิทฺธิ. (5)}

\proseitem{117}{กตมา \textbf{อริยา อิทฺธิ?} อิธ ภิกฺขุ\footnote{ม 3. 351 ปิฏฺเฐ ปสฺสิตพฺพา.} สเจ อากงฺขติ “ปฏิกูเล อปฺปฏิกูลสญฺญี วิหเรยฺยนฺ”ติ, อปฺปฏิกูลสญฺญี ตตฺถ วิหรติ.\csromanspacer{} สเจ อากงฺขติ “อปฺปฏิกูเล ปฏิกูลสญฺญี วิหเรยฺยนฺ”ติ, ปฏิกูลสญฺญี ตตฺถ วิหรติ.\csromanspacer{} สเจ อากงฺขติ “ปฏิกูเล จ อปฺปฏิกูเล จ อปฺปฏิกูลสญฺญี วิหเรยฺยนฺ”ติ, อปฺปฏิกูลสญฺญี ตตฺถ วิหรติ.\csromanspacer{} สเจ อากงฺขติ “อปฺปฏิกูเล จ ปฏิกูเล จ ปฏิกูลสญฺญี วิหเรยฺยนฺ”ติ, ปฏิกูลสญฺญี ตตฺถ วิหรติ.\csromanspacer{} สเจ อากงฺขติ “ปฏิกูเล จ อปฺปฏิกูเล จ ตทุภยํ อภินิวชฺเชตฺวา อุเปกฺขโก วิหเรยฺยํ สโต สมฺปชาโน”ติ อุเปกฺขโก ตตฺถ วิหรติ สโต สมฺปชาโน.}

\prose{กถํ \textbf{ปฏิกูเล อปฺปฏิกูลสญฺญี วิหรติ?} อนิฏฺฐสฺมิํ วตฺถุสฺมิํ เมตฺตาย วา ผรติ, ธาตุโต วา อุปสํหรติ, เอวํ ปฏิกูเล อปฺปฏิกูลสญฺญี วิหรติ.}

\prose{กถํ \textbf{อปฺปฏิกูเล ปฏิกูลสญฺญี วิหรติ?} อิฏฺฐสฺมิํ วตฺถุสฺมิํ อสุภาย วา ผรติ, อนิจฺจโต วา อุปสํหรติ, เอวํ อปฺปฏิกูเล ปฏิกูลสญฺญี วิหรติ.}

\prose{กถํ \textbf{ปฏิกูเล จ อปฺปฏิกูเล จ อปฺปฏิกูลสญฺญี วิหรติ?} อนิฏฺฐสฺมิํ จ อิฏฺฐสฺมิํ จ วตฺถุสฺมิํ เมตฺตาย วา ผรติ, ธาตุโต วา อุปสํหรติ.\csromanspacer{} เอวํ ปฏิกูเล จ อปฺปฏิกูเล จ อปฺปฏิกูลญฺญี วิหรติ.}

\prose{กถํ \textbf{อปฺปฏิกูเล จ ปฏิกูเล จ ปฏิกูลสญฺญี วิหรติ?} อิฏฺฐสฺมิํ จ อนิฏฺฐสฺมิํ จ วตฺถุสฺมิํ อสุภาย วา ผรติ, อนิจฺจโต วา อุปสํหรติ, เอวํ อปฺปฏิกูเล จ ปฏิกูเล จ ปฏิกูลสญฺญี วิหรติ.}

\prose{กถํ \textbf{ปฏิกูเล จ อปฺปฏิกูเล จ ตทุภยํ อภินิวชฺเชตฺวา อุเปกฺขโก วิหรติ สโต สมฺปชาโน?} อิธ ภิกฺขุ จกฺขุนา รูปํ ทิสฺวา เนว สุมโน โหติ น ทุมฺมโน, อุเปกฺขโก วิหรติ สโต สมฺปชาโน.\csromanspacer{} โสเตน สทฺทํ สุตฺวา ฯเปฯ.\csromanspacer{} ฆาเนน คนฺธํ ฆายิตฺวา.\csromanspacer{} ชิวฺหาย รสํ สายิตฺวา.}

\vspace{\tipitakaparskip}
\csromancenter{อภิสมยกถา กาเยน โผฏฺฐพฺพํ ผุสิตฺวา.\csromanspacer{} มนสา ธมฺมํ วิญฺญาย เนว สุมโน โหติ น ทุมฺมโน, อุเปกฺขโก วิหรติ สโต สมฺปชาโน.\csromanspacer{} เอวํ ปฏิกูเล จ อปฺปฏิกูเล จ ตทุภยํ อภินิวชฺเชตฺวา อุเปกฺขโก วิหรติ สโต สมฺปชาโน, อยํ อริยา อิทฺธิ. (6)}
\proseitem{18}{\csromanfolio{391}กตมา \textbf{กมฺมวิปากชา อิทฺธิ?} สพฺเพสํ ปกฺขีนํ สพฺเพสํ เทวานํ เอกจฺจานํ มนุสฺสานํ เอกจฺจานํ วินิปาติกานํ, อยํ กมฺมวิปากชา อิทฺธิ. (7)}

\prose{กตมา \textbf{ปุญฺญวโต อิทฺธิ?} ราชา จกฺกวตฺตี\footnote{จกฺกวตฺติ (สฺยา)} เวหาสํ คจฺฉติ สทฺธิํ จตุรงฺคินิยา เสนาย อนฺตมโส อสฺสพนฺธโคปุริเส\footnote{อสฺสพนฺธโคปเก ปุริเส (สฺยา)} อุปาทาย.\csromanspacer{} โชติกสฺส คหปติสฺส ปุญฺญวโต อิทฺธิ.\csromanspacer{} ชฏิลสฺส คหปติสฺส ปุญฺญวโต อิทฺธิ.\csromanspacer{} เมณฺฑกสฺส คหปติสฺส ปุญฺญวโต อิทฺธิ.\csromanspacer{} โฆสิตสฺส คหปติสฺส ปุญฺญวโต อิทฺธิ.\csromanspacer{} ปญฺจนฺนํ มหาปุญฺญานํ ปุญฺญวโต อิทฺธิ, อยํ ปุญฺญวโต อิทฺธิ. (8)}

\prose{กตมา \textbf{วิชฺชามยา อิทฺธิ?} วิชฺชาธรา วิชฺชํ ปริชปฺเปตฺวา เวหาสํ คจฺฉนฺติ, อากาเส อนฺตลิกฺเข หตฺถิมฺปิ ทสฺเสนฺติ, อสฺสมฺปิ ทสฺเสนฺติ, รถมฺปิ ทสฺเสนฺติ, ปตฺติมฺปิ ทสฺเสนฺติ, วิวิธมฺปิ เสนาพฺยูหํ ทสฺเสนฺติ, อยํ วิชฺชามยา อิทฺธิ. (9)}

\prose{กถํ \textbf{ตตฺถ ตตฺถ สมฺมาปโยคปจฺจยา อิชฺฌนฏฺเฐน อิทฺธิ?} เนกฺขมฺเมน กามจฺฉนฺทสฺส ปหานฏฺโฐ อิชฺฌตีติ ตตฺถ ตตฺถ สมฺมาปโยคปจฺจยา อิชฺฌนฏฺเฐน อิทฺธิ, อพฺยาปาเทน พฺยาปาทสฺส ปหานฏฺโฐ อิชฺฌตีติ ตตฺถ ตตฺถ สมฺมาปโยคปจฺจยา อิชฺฌนฏฺเฐน อิทฺธิ ฯเปฯ อรหตฺตมคฺเคน สพฺพกิเลสานํ ปหานฏฺโฐ อิชฺฌตีติ ตตฺถ ตตฺถ สมฺมาปโยคปจฺจยา อิชฺฌนฏฺเฐน อิทฺธิ.\csromanspacer{} เอวํ ตตฺถ ตตฺถ สมฺมาปโยคปจฺจยา อิชฺฌนฏฺเฐน อิทฺธิ. (10) อิมา ทส อิทฺธิโย.}

\nitthitamruled{อิทฺธิกถา นิฏฺฐิตา.}
\csromanmatikaanchor{mtk.124}
\csromantocmarkref{section}{3. อภิสมยกถา}{mtk.124}
\bookmark[dest={mtk.124},level=1]{3. อภิสมยกถา}
\csromanheader{1}{3. อภิสมยกถา}
\proseitem{19}{อ\textbf{ภิสมโย}ติ เกน อภิสเมติ? จิตฺเตน อภิสเมติ.}

\prose{หญฺจิ จิตฺเตน อภิสเมติ, เตน หิ อญฺญาณี อภิสเมติ.\csromanspacer{} น อญฺญาณี อภิสเมติ, ญาเณน อภิสเมติ.}

\prose{\csromanfolio{392}หญฺจิ ญาเณน อภิสเมติ, เตน หิ อจิตฺเตน จ ญาเณน จ อจิตฺตโก\footnote{เตน หิ อจิตฺตโก (สฺยา)} อภิสเมติ.\csromanspacer{} น อจิตฺตโก อภิสเมติ, จิตฺเตน จ ญาเณน จ อภิสเมติ.}

\prose{หญฺจิ จิตฺเตน จ ญาเณน จ อภิสเมติ, เตน หิ กามาวจรจิตฺเตน จ ญาเณน จ อภิสเมติ, น กามาวจรจิตฺเตน จ ญาเณน จ อภิสเมติ.}

\prose{เตน หิ รูปาวจรจิตฺเตน จ ญาเณน จ อภิสเมติ, น รูปาวจรจิตฺเตน จ ญาเณน จ อภิสเมติ.}

\prose{เตน หิ อรูปาวจรจิตฺเตน จ ญาเณน จ อภิสเมติ.\csromanspacer{} น อรูปาวจรจิตฺเตน จ ญาเณน จ อภิสเมติ.}

\prose{เตน หิ กมฺมสฺสกตจิตฺเตน จ ญาเณน จ อภิสเมติ, น กมฺมสฺสกตจิตฺเตน จ ญาเณน จ อภิสเมติ.}

\prose{เตน หิ สจฺจานุโลมิกจิตฺเตน จ ญาเณน จ อภิสเมติ, น สจฺจานุโลมิกจิตฺเตน จ ญาเณน จ อภิสเมติ.}

\prose{เตน หิ อตีตจิตฺเตน จ ญาเณน จ อภิสเมติ, น อตีตจิตฺเตน จ ญาเณน จ อภิสเมติ.}

\prose{เตน หิ อนาคตจิตฺเตน จ ญาเณน จ อภิสเมติ, น อนาคตจิตฺเตน จ ญาเณน จ อภิสเมติ.}

\prose{เตน หิ ปจฺจุปฺปนฺนโลกิยจิตฺเตน จ ญาเณน จ อภิสเมติ, น ปจฺจุปฺปนฺนโลกิยจิตฺเตน จ ญาเณน จ อภิสเมติ.\csromanspacer{} โลกุตฺตรมคฺคกฺขเณ ปจฺจุปฺปนฺนจิตฺเตน จ ญาเณน จ อภิสเมติ.}

\prose{กถํ โลกุตฺตรมคฺคกฺขเณ ปจฺจุปฺปนฺนจิตฺเตน จ ญาเณน จ อภิสเมติ? โลกุตฺตรมคฺคกฺขเณ อุปฺปาทาธิปเตยฺยํ จิตฺตํ ญาณสฺส เหตุ ปจฺจโย จ ตํสมฺปยุตฺตํ นิโรธโคจรํ ทสฺสนาธิปเตยฺยํ ญาณํ จิตฺตสฺส เหตุ ปจฺจโย จ ตํสมฺปยุตฺตํ ญาณํ นิโรธโคจรํ, เอวํ โลกุตฺตรมคฺคกฺขเณ ปจฺจุปฺปนฺนจิตฺเตน จ ญาเณน จ อภิสเมติ.}

\csromanheader{2}{3. อภิสมยกถา}
\proseitem{20}{\csromanfolio{393}กิํ นุ เอตฺตโกเยว อภิสมโยติ? น หิ, โลกุตฺตรมคฺคกฺขเณ ทสฺสนาภิสมโย สมฺมาทิฏฺฐิ, อภินิโรปนาภิสมโย สมฺมาสงฺกปฺโป, ปริคฺคหาภิสมโย สมฺมาวาจา, สมุฏฺฐานาภิสมโย สมฺมากมฺมนฺโต, โวทานาภิสมโย สมฺมา-อาชีโว, ปคฺคหาภิสมโย สมฺมาวายาโม, อุปฏฺฐานาภิสมโย สมฺมาสติ, อวิกฺเขปาภิสมโย สมฺมาสมาธิ.\csromanspacer{}\csromanspacer{}อุปฏฺฐานาภิสมโย สติสมฺโพชฺฌงฺโค, ปวิจยาภิสมโย ธมฺมวิจยสมฺโพชฺฌงฺโค, ปคฺคหาภิสมโย วีริยสมฺโพชฺฌงฺโค, ผรณาภิสมโย ปีติสมฺโพชฺฌงฺโค, อุปสมาภิสมโย ปสฺสทฺธิสมฺโพชฺฌงฺโค, อวิกฺเขปาภิสมโย สมาธิสมฺโพชฺฌงฺโค, ปฏิสงฺขานาภิสมโย อุเปกฺขาสมฺโพชฺฌงฺโค.\csromanspacer{}\csromanspacer{}อสฺสทฺธิเย อกมฺปิยาภิสมโย สทฺธาพลํ, โกสชฺเช อกมฺปิยาภิสมโย วีริยพลํ, ปมาเท อกมฺปิยาภิสมโย สติพลํ, อุทฺธจฺเจ อกมฺปิยาภิสมโย สมาธิพลํ, อวิชฺชาย อกมฺปิยาภิสมโย ปญฺญาพลํ.\csromanspacer{}\csromanspacer{}อธิโมกฺขาภิสมโย สทฺธินฺทฺริยํ, ปคฺคหาภิสมโย วีริยินฺทฺริยํ, อุปฏฺฐานาภิสมโย สตินฺทฺริยํ, อวิกฺเขปาภิสมโย สมาธินฺทฺริยํ, ทสฺสนาภิสมโย ปญฺญินฺทฺริยํ.\csromanspacer{}\csromanspacer{}อาธิปเตยฺยฏฺเฐน อินฺทฺริยาภิสมโย, อกมฺปิยฏฺเฐน พลาภิสมโย, นิยฺยานฏฺเฐน โพชฺฌงฺคาภิสมโย, เหตุฏฺเฐน มคฺคาภิสมโย, อุปฏฺฐานฏฺเฐน สติปฏฺฐานาภิสมโย, ปทหฏฺเฐน สมฺมปฺปธานาภิสมโย, อิชฺฌนฏฺเฐน อิทฺธิปาทาภิสมโย, ตถฏฺเฐน สจฺจาภิสมโย, อวิกฺเขปฏฺเฐน สมถาภิสมโย, อนุปสฺสนฏฺเฐน วิปสฺสนาภิสมโย, เอกรสฏฺเฐน สมถวิปสฺสนาภิสมโย, อนติวตฺตนฏฺเฐน ยุคนทฺธาภิสมโย, สํวรฏฺเฐน สีลวิสุทฺธิ-อภิสมโย, อวิกฺเขปฏฺเฐน วิสุทฺธิ-อภิสมโย, ทสฺสนฏฺเฐน ทิฏฺฐิวิสุทฺธิ-อภิสมโย, มุตฺตฏฺเฐน วิโมกฺขาภิสมโย, ปฏิเวธฏฺเฐน วิชฺชาภิสมโย, ปริจฺจาคฏฺเฐน วิมุตฺติ-อภิสมโย, สมุจฺเฉทฏฺเฐน ขเย ญาณํ อภิสมโย.\csromanspacer{}\csromanspacer{}ฉนฺโท มูลฏฺเฐน อภิสมโย, มนสิกาโร สมุฏฺฐานฏฺเฐน อภิสมโย, ผสฺโส สโมธานฏฺเฐน อภิสมโย, เวทนา สโมสรณฏฺเฐน อภิสมโย, สมาธิ ปมุขฏฺเฐน อภิสมโย, สติ อาธิปเตยฺยฏฺเฐน อภิสมโย, ปญฺญา ตทุตฺตรฏฺเฐน อภิสมโย, วิมุตฺติ สารฏฺเฐน อภิสมโย, อมโตคธํ นิพฺพานํ ปริโยสานฏฺเฐน อภิสมโย.}

\proseitem{21}{\csromanfolio{394}กิํนุ เอตฺตโกเยว อภิสมโยติ? น หิ, โสตาปตฺติมคฺคกฺขเณ ทสฺสนาภิสมโย สมฺมาทิฏฺฐิ ฯเปฯ อมโตคธํ นิพฺพานํ ปริโยสานฏฺเฐน อภิสมโย.}

\prose{กิํนุ เอตฺตโกเยว อภิสมโยติ? น หิ, โสตาปตฺติผลกฺขเณ ทสฺสนาภิสมโย สมฺมาทิฏฺฐิ ฯเปฯ ปฏิปฺปสฺสทฺธฏฺเฐน อนุปฺปาเท ญาณํ อภิสมโย, ฉนฺโท มูลฏฺเฐน อภิสมโย ฯเปฯ อมโตคธํ นิพฺพานํ ปริโยสานฏฺเฐน อภิสมโย.}

\prose{กิํนุ เอตฺตโกเยว อภิสมโยติ? น หิ, สกทาคามิมคฺคกฺขเณ ฯเปฯ สกทาคามิผลกฺขเณ.\csromanspacer{} อนาคามิมคฺคกฺขเณ.\csromanspacer{} อนาคามิผลกฺขเณ.\csromanspacer{} อรหตฺตมคฺคกฺขเณ ฯเปฯ.\csromanspacer{} อรหตฺตผลกฺขเณ ทสฺสนาภิสมโย สมฺมาทิฏฺฐิ, อภินิโรปนาภิสมโย สมฺมาสงฺกปฺโป ฯเปฯ ปฏิปฺปสฺสทฺธฏฺเฐน อนุปฺปาเท ญาณํ อภิสมโย.\csromanspacer{} ฉนฺโท มูลฏฺเฐน อภิสมโย ฯเปฯ อมโตคธํ นิพฺพานํ ปริโยสานฏฺเฐน อภิสมโย.}

\prose{ยฺวายํ\footnote{สฺวายํ (สฺยา, อิ)} กิเลเส ปชหติ, อตีเต กิเลเส ปชหติ ฯเปฯ อนาคเต กิเลเส ปชหติ, ปจฺจุปฺปนฺเน กิเลเส ปชหติ? อตีเต กิเลเส ปชหตีติ.\csromanspacer{} หญฺจิ อตีเต กิเลเส ปชหติ, เตน หิ ขีณํ เขเปติ, นิรุทฺธํ นิโรเธติ, วิคตํ วิคเมติ, อตฺถงฺคตํ อตฺถงฺคเมติ, อตีตํ ยํ น อตฺถิ, ตํ ปชหตีติ น อตีเต กิเลเส ปชหตีติ.\csromanspacer{}\csromanspacer{}อนาคเต กิเลเส ปชหตีติ.\csromanspacer{} หญฺจิ อนาคเต กิเลเส ปชหติ, เตน หิ อชาตํ ปชหติ, อนิพฺพตฺตํ ปชหติ, อนุปฺปนฺนํ ปชหติ, อปาตุภูตํ ปชหติ, อนาคตํ ยํ น อตฺถิ, ตํ ปชหตีติ น อนาคเต กิเลเส ปชหตีติ.\csromanspacer{}\csromanspacer{}ปจฺจุปฺปนฺเน กิเลเส ปชหตีติ.\csromanspacer{} หญฺจิ ปจฺจุปฺปนฺเน กิเลเส ปชหติ, เตน หิ รตฺโต ราคํ ปชหติ, ทุฏฺโฐ โทสํ ปชหติ, มูโฬฺห โมหํ ปชหติ, วินิพทฺโธ มานํ ปชหติ, ปรามฏฺโฐ ทิฏฺฐิํ ปชหติ, วิกฺเขปคโต อุทฺธจฺจํ ปชหติ, อนิฏฺฐงฺคโต วิจิกิจฺฉํ ปชหติ, ถามคโต อนุสยํ ปชหติ, กณฺหสุกฺกธมฺมา ยุคนทฺธา สมเมว วตฺตนฺติ.\csromanspacer{} สํกิเลสิกา\footnote{ตํสํกิเลสิกา (สฺยา, ก)} มคฺคภาวนา โหติ.}

\prose{น หิ อตีเต กิเลเส ปชหติ, น อนาคเต กิเลเส ปชหติ, น ปจฺจุปฺปนฺเน กิเลเส ปชหตีติ.\csromanspacer{} หญฺจิ น อตีเต กิเลเส ปชหติ,}

\csromanmatikaanchor{mtk.125}
\csromantocmarknopage{section}{4. วิเวกกถา}{mtk.125}
\bookmark[dest={mtk.125},level=1]{4. วิเวกกถา}
\csromanheader{1}{4. วิเวกกถา}
\prose{\csromanfolio{395}น อนาคเต ฯเปฯ.\csromanspacer{} น ปจฺจุปฺปนฺเน กิเลเส ปชหติ, เตน หิ นตฺถิ มคฺคภาวนา, นตฺถิ ผลสจฺฉิกิริยา, นตฺถิ กิเลสปฺปหานํ, นตฺถิ ธมฺมาภิสมโยติ.\csromanspacer{} อตฺถิ มคฺคภาวนา, อตฺถิ ผลสจฺฉิกิริยา, อตฺถิ กิเลสปฺปหานํ, อตฺถิ ธมฺมาภิสมโย.\csromanspacer{} ยถา กถํ วิย, เสยฺยถาปิ ตรุโณ รุกฺโข อชาตผโล, ตเมนํ ปุริโส มูลํ ฉินฺเทยฺย, เย ตสฺส รุกฺขสฺส อชาตผลา, เต อชาตาเยว น ชายนฺติ.\csromanspacer{} อนิพฺพตฺตาเยว น นิพฺพตฺตนฺติ, อนุปฺปนฺนาเยว น อุปฺปชฺชนฺติ, อปาตุภูตาเยว น ปาตุภวนฺติ.\csromanspacer{} เอวเมวํ “อุปฺปาโท เหตุ อุปฺปาโท ปจฺจโย กิเลสานํ นิพฺพตฺติยา”ติ อุปฺปาเท อาทีนวํ ทิสฺวา อนุปฺปาเท จิตฺตํ ปกฺขนฺทติ, อนุปฺปาเท จิตฺตสฺส ปกฺขนฺทตฺตา เย อุปฺปาทปจฺจยา กิเลสา นิพฺพตฺเตยฺยุํ, เต อชาตาเยว น ชายนฺติ, อนิพฺพตฺตาเยว น นิพฺพตฺตนฺติ, อนุปฺปนฺนาเยว น อุปฺปชฺชนฺติ, อปาตุภูตาเยว น ปาตุภวนฺติ.\csromanspacer{} เอวํ “เหตุนิโรธา ทุกฺขนิโรโธ, ปวตฺตํ เหตุ, นิมิตฺตํ เหตุ, อายูหนา เหตุ, อายูหนา ปจฺจโย กิเลสานํ นิพฺพตฺติยา”ติ อายูหเน อาทีนวํ ทิสฺวา อนายูหเน จิตฺตํ ปกฺขนฺทติ, อนายูหเน จิตฺตสฺส ปกฺขนฺทตฺตา เย อายูหนปจฺจยา กิเลสา นิพฺพตฺเตยฺยุํ, เต อชาตาเยว น ชายนฺติ, อนิพฺพตฺตาเยว น นิพฺพตฺตนฺติ, อนุปฺปนฺนาเยว น อุปฺปชฺชนฺติ, อปาตุภูตาเยว น ปาตุภวนฺติ.\csromanspacer{} เอวํ เหตุนิโรธา ทุกฺขนิโรโธ.\csromanspacer{} เอวํ อตฺถิ มคฺคภาวนา, อตฺถิ ผลสจฺฉิกิริยา, อตฺถิ กิเลสปฺปหานํ, อตฺถิ ธมฺมาภิสมโยติ.}

\nitthitamruled{อภิสมยกถา นิฏฺฐิตา.}
\csromanheader{2}{4. วิเวกกถา}
\proseitem{22}{สาวตฺถินิทานํ\footnote{สํ 3. 42 ปิฏฺเฐ ปสฺสิตพฺพา.}.\csromanspacer{}\csromanspacer{}เสยฺยถาปิ ภิกฺขเว เย เกจิ พลกรณียา กมฺมนฺตา กรียนฺติ, สพฺเพ เต ปถวิํ นิสฺสาย ปถวิยํ ปติฏฺฐาย, เอวเมเต พลกรณียา กมฺมนฺตา กรียนฺติ.\csromanspacer{} เอวเมวํ ภิกฺขเว ภิกฺขุ สีลํ นิสฺสาย สีเล ปติฏฺฐาย อริยํ อฏฺฐงฺคิกํ มคฺคํ ภาเวติ, อริยํ อฏฺฐงฺคิกํ มคฺคํ พหุลีกโรติ.}

\prose{กถญฺจ ภิกฺขเว ภิกฺขุ สีลํ นิสฺสาย สีเล ปติฏฺฐาย อริยํ อฏฺฐงฺคิกํ มคฺคํ ภาเวติ, อริยํ อฏฺฐงฺคิกํ มคฺคํ พหุลีกโรติ? อิธ ภิกฺขเว \csromanfolio{396}ภิกฺขุ สมฺมาทิฏฺฐิํ ภาเวติ วิเวกนิสฺสิตํ วิราคนิสฺสิตํ นิโรธนิสฺสิตํ โวสคฺคปริณามิํ.\csromanspacer{} สมฺมาสงฺกปฺปํ ฯเปฯ.\csromanspacer{} สมฺมาวาจํ.\csromanspacer{} สมฺมากมฺมนฺตํ.\csromanspacer{} สมฺมา-อาชีวํ.\csromanspacer{} สมฺมาวายามํ.\csromanspacer{} สมฺมาสติํ.\csromanspacer{} สมฺมาสมาธิํ ภาเวติ วิเวกนิสฺสิตํ วิราคนิสฺสิตํ นิโรธนิสฺสิตํ โวสคฺคปริณามิํ.\csromanspacer{} เอวํ โข ภิกฺขเว ภิกฺขุ สีลํ นิสฺสาย สีเล ปติฏฺฐาย อริยํ อฏฺฐงฺคิกํ มคฺคํ ภาเวติ, อริยํ อฏฺฐงฺคิกํ มคฺคํ พหุลีกโรติ.}

\proseitem{23}{เสยฺยถาปิ ภิกฺขเว เย เกจิ\footnote{เยปิเม (สฺยา) สํ 3. 43 ปิฏฺเฐ ปสฺสิตพฺพา.} พีชคามภูตคามา วุทฺธิํ วิรูฬฺหิํ เวปุลฺลํ อาปชฺชนฺติ, สพฺเพ เต ปถวิํ นิสฺสาย ปถวิยํ ปติฏฺฐาย, เอวเมเต พีชคามภูตคามา วุทฺธิํ วิรุฬฺหิํ เวปุลฺลํ อาปชฺชนฺติ.\csromanspacer{} เอวเมวํ โข ภิกฺขเว ภิกฺขุ สีลํ นิสฺสาย สีเล ปติฏฺฐาย อริยํ อฏฺฐงฺคิกํ มคฺคํ ภาเวนฺโต อริยํ อฏฺฐงฺคิกํ มคฺคํ พหุลีกโรนฺโต วุทฺธิํ วิรุฬฺหิํ เวปุลฺลํ ปาปุณาติ ธมฺเมสุ.}

\prose{กถญฺจ ภิกฺขเว ภิกฺขุ สีลํ นิสฺสาย สีเล ปติฏฺฐาย อริยํ อฏฺฐงฺคิกํ มคฺคํ ภาเวนฺโต อริยํ อฏฺฐงฺคิกํ มคฺคํ พหุลีกโรนฺโต วุทฺธิํ วิรุฬฺหิํ เวปุลฺลํ ปาปุณาติ ธมฺเมสุ? อิธ ภิกฺขเว ภิกฺขุ สมฺมาทิฏฺฐิํ ภาเวติ วิเวกนิสฺสิตํ วิราคนิสฺสิตํ นิโรธนิสฺสิตํ โวสคฺคปริณามิํ, สมฺมาสงฺกปฺปํ ภาเวติ ฯเปฯ สมฺมาวาจํ ภาเวติ.\csromanspacer{} สมฺมากมฺมนฺตํ ภาเวติ.\csromanspacer{} สมฺมา-อาชีวํ ภาเวติ.\csromanspacer{} สมฺมาวายามํ ภาเวติ.\csromanspacer{} สมฺมาสติํ ภาเวติ.\csromanspacer{} สมฺมาสมาธิํ ภาเวติ วิเวกนิสฺสิตํ วิราคนิสฺสิตํ นิโรธนิสฺสิตํ โวสคฺคปริณามิํ.\csromanspacer{} เอวํ โข ภิกฺขเว ภิกฺขุ สีลํ นิสฺสาย สีเล ปติฏฺฐาย อริยํ อฏฺฐงฺคิกํ มคฺคํ ภาเวนฺโต อริยํ อฏฺฐงฺคิกํ มคฺคํ พหุลีกโรนฺโต วุฑฺฒิํ วิรุฬฺหิํ เวปุลฺลํ ปาปุณาติ ธมฺเมสูติ.}

\csromanmatikaanchor{mtk.126}
\csromantocmarkref{subsection}{1. มคฺคงฺคนิทฺเทส}{mtk.126}
\bookmark[dest={mtk.126},level=2]{1. มคฺคงฺคนิทฺเทส}
\csromanheader{2}{1. มคฺคงฺคนิทฺเทส}
\proseitem{24}{สมฺมาทิฏฺฐิยา ปญฺจ วิเวกา, ปญฺจ วิราคา, ปญฺจ นิโรธา, ปญฺจ โวสคฺคา, ทฺวาทส นิสฺสยา.\csromanspacer{} สมฺมาสงฺกปฺปสฺส ฯเปฯ.\csromanspacer{} สมฺมาวาจาย.\csromanspacer{} สมฺมากมฺมนฺตสฺส.\csromanspacer{} สมฺมา-อาชีวสฺส.\csromanspacer{} สมฺมาวายามสฺส.\csromanspacer{} สมฺมาสติยา.\csromanspacer{} สมฺมาสมาธิสฺส ปญฺจ วิเวกา, ปญฺจ วิราคา, ปญฺจ นิโรธา, ปญฺจ โวสคฺคา, ทฺวาทส นิสฺสยา.}

\csromanheader{2}{4. วิเวกกถา}
\prose{\csromanfolio{397}\textbf{สมฺมาทิฏฺฐิยา กตเม ปญฺจ วิเวกา?} วิกฺขมฺภนวิเวโก ตทงฺควิเวโก สมุจฺเฉทวิเวโก ปฏิปฺปสฺสทฺธิวิเวโก นิสฺสรณวิเวโก.\csromanspacer{} วิกฺขมฺภนวิเวโก จ นีวรณานํ ปฐมชฺฌานํ ภาวยโต, ตทงฺควิเวโก จ ทิฏฺฐิคตานํ นิพฺเพธภาคิยํ สมาธิํ ภาวยโต, สมุจฺเฉทวิเวโก จ โลกุตฺตรํ ขยคามิมคฺคํ ภาวยโต, ปฏิปฺปสฺสทฺธิวิเวโก จ ผลกฺขเณ, นิสฺสรณวิเวโก จ นิโรโธ นิพฺพานํ.\csromanspacer{} สมฺมาทิฏฺฐิยา อิเม ปญฺจ วิเวกา.\csromanspacer{} อิเมสุ ปญฺจสุ วิเวเกสุ ฉนฺทชาโต โหติ สทฺธาธิมุตฺโต, จิตฺตํ จสฺส สฺวาธิฏฺฐิตํ.}

\prose{\textbf{สมฺมาทิฏฺฐิยา กตเม ปญฺจ วิราคา?} วิกฺขมฺภนวิราโค ตทงฺควิราโค สมุจฺเฉทวิราโค ปฏิปฺปสฺสทฺธิวิราโค นิสฺสรณวิราโค.\csromanspacer{} วิกฺขมฺภนวิราโค จ นีวรณานํ ปฐมชฺฌานํ ภาวยโต, ตทงฺควิราโค จ ทิฏฺฐิคตานํ นิพฺเพธภาคิยํ สมาธิํ ภาวยโต, สมุจฺเฉทวิราโค จ โลกุตฺตรํ ขยคามิมคฺคํ ภาวยโต, ปฏิปฺปสฺสทฺธิวิราโค จ ผลกฺขเณ, นิสฺสรณวิราโค จ นิโรโธ นิพฺพานํ.\csromanspacer{} สมฺมาทิฏฺฐิยา อิเม ปญฺจ วิราคา.\csromanspacer{} อิเมสุ ปญฺจสุ วิราเคสุ ฉนฺทชาโต โหติ สทฺธาธิมุตฺโต, จิตฺตํ จสฺส สฺวาธิฏฺฐิตํ.}

\prose{\textbf{สมฺมาทิฏฺฐิยา กตเม ปญฺจ นิโรธา?} วิกฺขมฺภนนิโรโธ ตทงฺคนิโรโธ สมุจฺเฉทนิโรโธ ปฏิปฺปสฺสทฺธินิโรโธ นิสฺสรณนิโรโธ.\csromanspacer{} วิกฺขมฺภนนิโรโธ จ นีวรณานํ ปฐมชฺฌานํ ภาวยโต, ตทงฺคนิโรโธ จ ทิฏฺฐิคตานํ นิพฺเพธภาคิยํ สมาธิํ ภาวยโต, สมุจฺเฉทนิโรโธ จ โลกุตฺตรํ ขยคามิมคฺคํ ภาวยโต, ปฏิปฺปสฺสทฺธินิโรโธ จ ผลกฺขเณ, นิสฺสรณนิโรโธ จ อมตา ธาตุ.\csromanspacer{} สมฺมาทิฏฺฐิยา อิเม ปญฺจ นิโรธา.\csromanspacer{} อิเมสุ ปญฺจสุ นิโรเธสุ ฉนฺทชาโต โหติ สทฺธาธิมุตฺโต, จิตฺตํ จสฺส สฺวาธิฏฺฐิตํ.}

\prose{\textbf{สมฺมาทิฏฺฐิยา กตเม ปญฺจ โวสคฺคา?} วิกฺขมฺภนโวสคฺโค ตทงฺคโวสคฺโค สมุจฺเฉทโวสคฺโค ปฏิปฺปสฺสทฺธิโวสคฺโค นิสฺสรณโวสคฺโค.\csromanspacer{} วิกฺขมฺภนโวสคฺโค จ นีวรณานํ ปฐมชฺฌานํ ภาวยโต, ตทงฺคโวสคฺโค จ ทิฏฺฐิคตานํ นิพฺเพธภาคิยํ สมาธิํ ภาวยโต, สมุจฺเฉทโวสคฺโค จ โลกุตฺตรํ ขยคามิมคฺคํ ภาวยโต, ปฏิปฺปสฺสทฺธิโวสคฺโค จ ผลกฺขเณ, นิสฺสรณโวสคฺโค จ นิโรโธ นิพฺพานํ.\csromanspacer{} สมฺมาทิฏฺฐิยา อิเม ปญฺจ โวสคฺคา.\csromanspacer{} อิเมสุ ปญฺจสุ โวสคฺเคสุ ฉนฺทชาโต}

\prose{\csromanfolio{398}โหติ สทฺธาธิมุตฺโต, จิตฺตํ จสฺส สฺวาธิฏฺฐิตํ.\csromanspacer{} สมฺมาทิฏฺฐิยา อิเม ปญฺจ วิเวกา, ปญฺจ วิราคา, ปญฺจ นิโรธา, ปญฺจ โวสคฺคา, ทฺวาทส นิสฺสยา.}

\proseitem{25}{\textbf{สมฺมาสงฺกปฺปสฺส ฯเปฯ.}\csromanspacer{}\textbf{ สมฺมาวาจาย.}\csromanspacer{}\textbf{ สมฺมากมฺมนฺตสฺส.}\csromanspacer{}\textbf{ สมฺมา-อาชีวสฺส.}\csromanspacer{}\textbf{ สมฺมาวายามสฺส.}\csromanspacer{}\textbf{ สมฺมาสติยา.}\csromanspacer{}\textbf{ สมฺมาสมาธิสฺส กตเม ปญฺจ วิเวกา?} วิกฺขมฺภนวิเวโก ตทงฺควิเวโก สมุจฺเฉทวิเวโก ปฏิปฺปสฺสทฺธิวิเวโก นิสฺสรณวิเวโก.\csromanspacer{} วิกฺขมฺภนวิเวโก จ นีวรณานํ ปฐมชฺฌานํ ภาวยโต, ตทงฺควิเวโก จ ทิฏฺฐิคตานํ นิพฺเพธภาคิยํ สมาธิํ ภาวยโต, สมุจฺเฉทวิเวโก จ โลกุตฺตรํ ขยคามิมคฺคํ ภาวยโต, ปฏิปฺปสฺสทฺธิวิเวโก จ ผลกฺขเณ, นิสฺสรณวิเวโก จ นิโรโธ นิพฺพานํ.\csromanspacer{} สมฺมาสมาธิสฺส อิเม ปญฺจ วิเวกา.\csromanspacer{} อิเมสุ ปญฺจสุ วิเวเกสุ ฉนฺทชาโต โหติ สทฺธาธิมุตฺโต, จิตฺตํ จสฺส สฺวาธิฏฺฐิตํ.}

\prose{\textbf{สมฺมาสมาธิสฺส กตเม ปญฺจ วิราคา?} วิกฺขมฺภนวิราโค ตทงฺควิราโค สมุจฺเฉทวิราโค ปฏิปฺปสฺสทฺธิวิราโค นิสฺสรณวิราโค.\csromanspacer{} วิกฺขมฺภนวิราโค จ นีวรณานํ ปฐมชฺฌานํ ภาวยโต, ตทงฺควิราโค จ ทิฏฺฐิคตานํ นิพฺเพธภาคิยํ สมาธิํ ภาวยโต, สมุจฺเฉทวิราโค จ โลกุตฺตรํ ขยคามิมคฺคํ ภาวยโต, ปฏิปฺปสฺสทฺธิวิราโค จ ผลกฺขเณ, นิสฺสรณวิราโค จ นิโรโธ นิพฺพานํ.\csromanspacer{} สมฺมาสมาธิสฺส อิเม ปญฺจ วิราคา.\csromanspacer{} อิเมสุ ปญฺจสุ วิราเคสุ ฉนฺทชาโต โหติ สทฺธาธิมุตฺโต, จิตฺตํ จสฺส สฺวาธิฏฺฐิตํ.}

\prose{\textbf{สมฺมาสมาธิสฺส กตเม ปญฺจ นิโรธา?} วิกฺขมฺภนนิโรโธ ตทงฺคนิโรโธ สมุจฺเฉทนิโรโธ ปฏิปฺปสฺสทฺธินิโรโธ นิสฺสรณนิโรโธ.\csromanspacer{} วิกฺขมฺภนนิโรโธ จ นีวรณานํ ปฐมชฺฌานํ ภาวยโต, ตทงฺคนิโรโธ จ ทิฏฺฐิคตานํ นิพฺเพธภาคิยํ สมาธิํ ภาวยโต, สมุจฺเฉทนิโรโธ จ โลกุตฺตรํ ขยคามิมคฺคํ ภาวยโต, ปฏิปฺปสฺสทฺธินิโรโธ จ ผลกฺขเณ, นิสฺสรณนิโรโธ จ อมตา ธาตุ.\csromanspacer{} สมฺมาสมาธิสฺส อิเม ปญฺจ นิโรธา.\csromanspacer{} อิเมสุ ปญฺจสุ นิโรเธสุ ฉนฺทชาโต โหติ สทฺธาธิมุตฺโต, จิตฺตํ จสฺส สฺวาธิฏฺฐิตํ.}

\prose{\textbf{สมฺมาสมาธิสฺส กตเม ปญฺจ โวสคฺคา?} วิกฺขมฺภนโวสคฺโค ตทงฺคโวสคฺโค สมุจฺเฉทโวสคฺโค ปฏิปฺปสฺสทฺธิโวสคฺโค นิสฺสรณโวสคฺโค.\csromanspacer{} วิกฺขมฺภนโวสคฺโค จ นีวรณานํ ปฐมชฺฌานํ ภาวยโต, ตทงฺคโวสคฺโค จ ทิฏฺฐิคตานํ นิพฺเพธภาคิยํ สมาธิํ ภาวยโต, สมุจฺเฉทโวสคฺโค จ}

\csromanheader{2}{4. วิเวกกถา}
\prose{\csromanfolio{399}โลกุตฺตรํ ขยคามิมคฺคํ ภาวยโต, ปฏิปฺปสฺสทฺธิโวสคฺโค จ ผลกฺขเณ, นิสฺสรณโวสคฺโค จ นิโรโธ นิพฺพานํ.\csromanspacer{} สมฺมาสมาธิสฺส อิเม ปญฺจ โวสคฺคา.\csromanspacer{} อิเมสุ ปญฺจสุ โวสคฺเคสุ ฉนฺทชาโต โหติ สทฺธาธิมุตฺโต, จิตฺตํ จสฺส สฺวาธิฏฺฐิตํ.\csromanspacer{} สมฺมาสมาธิสฺส อิเม ปญฺจ วิเวกา ปญฺจ วิราคา ปญฺจ นิโรธา ปญฺจ โวสคฺคา ทฺวาทส นิสฺสยา.}

\proseitem{26}{เสยฺยถาปิ ภิกฺขเว เย เกจิ พลกรณียา กมฺมนฺตา กรียนฺติ, สพฺเพ เต ปถวิํ นิสฺสาย ปถวิยํ ปติฏฺฐาย, เอวเมเต พลกรณียา กมฺมนฺตา กรียนฺติ.\csromanspacer{} เอวเมวํ โข ภิกฺขเว ภิกฺขุ สีลํ นิสฺสาย สีเล ปติฏฺฐาย สตฺต โพชฺฌงฺเค ภาเวติ, สตฺต โพชฺฌงฺเค พหุลีกโรติ ฯเปฯ.\csromanspacer{} สตฺต โพชฺฌงฺเค ภาเวนฺโต สตฺต โพชฺฌงฺเค พหุลีกโรนฺโต วุทฺธิํ วิรุฬฺหิํ เวปุลฺลํ ปาปุณาติ ธมฺเมสุ ฯเปฯ.\csromanspacer{} ปญฺจ พลานิ ภาเวติ.\csromanspacer{} ปญฺจ พลานิ พหุลีกโรติ ฯเปฯ.\csromanspacer{} ปญฺจ พลานิ ภาเวนฺโต ปญฺจ พลานิ พหุลีกโรนฺโต วุทฺธิํ วิรุฬฺหิํ เวปุลฺลํ ปาปุณาติ ธมฺเมสุ ฯเปฯ.\csromanspacer{} ปญฺจินฺทฺริยานิ ภาเวติ, ปญฺจินฺทฺริยานิ พหุลีกโรติ ฯเปฯ.}

\prose{เสยฺยถาปิ ภิกฺขเว เย เกจิ พีชคามภูตคามา วุทฺธิํ วิรุฬฺหิํ เวปุลฺลํ อาปชฺชนฺติ, สพฺเพ เต ปถวิํ นิสฺสาย ปถวิยํ ปติฏฺฐาย, เอวเมเต พีชคามภูตคามา วุทฺธิํ วิรุฬฺหิํ เวปุลฺลํ อาปชฺชนฺติ.\csromanspacer{} เอวเมวํ โข ภิกฺขเว ภิกฺขุ สีลํ นิสฺสาย สีเล ปติฏฺฐาย ปญฺจินฺทฺริยานิ ภาเวนฺโต ปญฺจินฺทฺริยานิ พหุลีกโรนฺโต วุทฺธิํ วิรุฬฺหิํ เวปุลฺลํ ปาปุณาติ ธมฺเมสุ.}

\prose{กถญฺจ ภิกฺขเว ภิกฺขุ สีลํ นิสฺสาย สีเล ปติฏฺฐาย ปญฺจินฺทฺริยานิ ภาเวนฺโต ปญฺจินฺทฺริยานิ พหุลีกโรนฺโต วุทฺธิํ วิรุฬฺหิํ เวปุลฺลํ ปาปุณาติ ธมฺเมสุ? อิธ ภิกฺขเว ภิกฺขุ สทฺธินฺทฺริยํ ภาเวติ วิเวกนิสฺสิตํ วิราคนิสฺสิตํ นิโรธนิสฺสิตํ โวสคฺคปริณามิํ.\csromanspacer{} วีริยินฺทฺริยํ ภาเวติ ฯเปฯ.\csromanspacer{} สตินฺทฺริยํ ภาเวติ ฯเปฯ.\csromanspacer{} สมาธินฺทฺริยํ ภาเวติ ฯเปฯ.\csromanspacer{} ปญฺญินฺทฺริยํ ภาเวติ วิเวกนิสฺสิตํ วิราคนิสฺสิตํ นิโรธนิสฺสิตํ โวสคฺคปริณามิํ.\csromanspacer{} เอวํ โข ภิกฺขเว ภิกฺขุ สีลํ นิสฺสาย ฯเปฯ.\csromanspacer{} ปาปุณาติ ธมฺเมสูติ.}

\csromanmatikaanchor{mtk.127}
\csromantocmarkref{subsection}{2. อินฺทฺริยนิทฺเทส}{mtk.127}
\bookmark[dest={mtk.127},level=2]{2. อินฺทฺริยนิทฺเทส}
\csromanheader{2}{2. อินฺทฺริยนิทฺเทส}
\proseitem{27}{สทฺธินฺทฺริยสฺส ปญฺจ วิเวกา, ปญฺจ วิราคา, ปญฺจ นิโรธา, ปญฺจ โวสคฺคา, ทฺวาทส นิสฺสยา.\csromanspacer{} วีริยินฺทฺริยสฺส ฯเปฯ.\csromanspacer{} สตินฺทฺริยสฺส ฯเปฯ.\csromanspacer{} สมาธินฺทฺริยสฺส ฯเปฯ. \csromanfolio{400}ปญฺญินฺทฺริยสฺส ปญฺจ วิเวกา, ปญฺจ วิราคา, ปญฺจ นิโรธา, ปญฺจ โวสคฺคา, ทฺวาทส นิสฺสยา.}

\prose{\textbf{สทฺธินฺทฺริยสฺส กตเม ปญฺจ วิเวกา?} วิกฺขมฺภนวิเวโก ตทงฺควิเวโก สมุจฺเฉทวิเวโก ปฏิปฺปสฺสทฺธิวิเวโก นิสฺสรณวิเวโก.\csromanspacer{} วิกฺขมฺภนวิเวโก จ นีวรณานํ ปฐมชฺฌานํ ภาวยโต, ตทงฺควิเวโก จ ทิฏฺฐิคตานํ นิพฺเพธภาคิยํ สมาธิํ ภาวยโต, สมุจฺเฉทวิเวโก จ โลกุตฺตรํ ขยคามิมคฺคํ ภาวยโต, ปฏิปฺปสฺสทฺธิวิเวโก จ ผลกฺขเณ, นิสฺสรณวิเวโก จ นิโรโธ นิพฺพานํ.\csromanspacer{} สทฺธินฺทฺริยสฺส อิเม ปญฺจ วิเวกา.\csromanspacer{} อิเมสุ ปญฺจสุ วิเวเกสุ ฉนฺทชาโต โหติ สทฺธาธิมุตฺโต, จิตฺตํ จสฺส สฺวาธิฏฺฐิตํ ฯเปฯ.\csromanspacer{} สทฺธินฺทฺริยสฺส อิเม ปญฺจ วิเวกา, ปญฺจ วิราคา, ปญฺจ นิโรธา, ปญฺจ โวสคฺคา, ทฺวาทส นิสฺสยา.}

\prose{วีริยินฺทฺริยสฺส ฯเปฯ.\csromanspacer{} สตินฺทฺริยสฺส ฯเปฯ.\csromanspacer{} สมาธินฺทฺริยสฺส ฯเปฯ.\csromanspacer{} ปญฺญินฺทฺริยสฺส กตเม ปญฺจ วิเวกา.\csromanspacer{} วิกฺขมฺภนวิเวโก ตทงฺควิเวโก สมุจฺเฉทวิเวโก ปฏิปฺปสฺสทฺธิวิเวโก นิสฺสรณวิเวโก ฯเปฯ.\csromanspacer{} ปญฺญินฺทฺริยสฺส อิเม ปญฺจ วิเวกา, ปญฺจ วิราคา, ปญฺจ นิโรธา, ปญฺจ โวสคฺคา, ทฺวาทส นิสฺสยาติ.}

\nitthitamruled{วิเวกกถา นิฏฺฐิตา.}
\csromanmatikaanchor{mtk.128}
\csromantocmarkref{section}{5. จริยากถา}{mtk.128}
\bookmark[dest={mtk.128},level=1]{5. จริยากถา}
\csromanheader{1}{5. จฺ\textbf{อริยา}กถา}
\proseitem{28}{จฺ\textbf{อริยา}ติ อฏฺฐ จฺ\textbf{อริยา}โย อิริยาปถจฺ\textbf{อริยา} อายตนจฺ\textbf{อริยา} สติจฺ\textbf{อริยา} สมาธิจฺ\textbf{อริยา} ญาณจฺ\textbf{อริยา} มฺ\textbf{อคฺคจริยา} ปตฺติจฺ\textbf{อริยา} โลกตฺถจฺ\textbf{อริยา}ติ.}

\prose{\textbf{อิริยาปถจริยา}ติ จตูสุ อิริยาปเถสุ.\csromanspacer{} \textbf{อายตนจริยา}ติ ฉสุ อชฺฌตฺติกพาหิเรสุ อายตเนสุ.\csromanspacer{} \textbf{สติจริยา}ติ จตูสุ สติปฏฺฐาเนสุ.\csromanspacer{} \textbf{สมาธิจริยา}ติ จตูสุ ฌาเนสุ.\csromanspacer{} \textbf{ญาณจริยา}ติ จตูสุ อริยสจฺเจสุ.\csromanspacer{} มฺ\textbf{อคฺคจริยา}ติ จตูสุ อริยมคฺเคสุ.\csromanspacer{} \textbf{ปตฺติจริยา}ติ จตูสุ สามญฺญผเลสุ.\csromanspacer{} \textbf{โลกตฺถจริยา}ติ ตถาคเตสุ อรหนฺเตสุ สมฺมาสมฺพุทฺเธสุ ปเทเส\footnote{ปเทโส (สฺยา)} ปจฺเจกพุทฺเธสุ ปเทเส สาวเกสุ.}

\csromanmatikaanchor{mtk.129}
\csromantocmarkref{section}{6. ปาฏิหาริยกถา}{mtk.129}
\bookmark[dest={mtk.129},level=1]{6. ปาฏิหาริยกถา}
\csromanheader{1}{6. ปาฏิหาริยกถา}
\prose{\csromanfolio{401}อิริยาปถจริยา จ ปณิธิสมฺปนฺนานํ.\csromanspacer{} อายตนจริยา จ อินฺทฺริเยสุ คุตฺตทฺวารานํ.\csromanspacer{} สติจริยา จ อปฺปมาทวิหารีนํ.\csromanspacer{} สมาธิจริยา จ อธิจิตฺตมนุยุตฺตานํ.\csromanspacer{} ญาณจริยา จ พุทฺธิสมฺปนฺนานํ.\csromanspacer{} มคฺคจริยา จ สมฺมาปฏิปนฺนานํ.\csromanspacer{} ปตฺติจริยา จ อธิคตผลานํ.\csromanspacer{} โลกตฺถจริยา จ ตถาคตานํ อรหนฺตานํ สมฺมาสมฺพุทฺธานํ ปเทเส ปจฺเจกพุทฺธานํ ปเทเส สาวกานํ.\csromanspacer{} อิมา อฏฺฐ จริยาโย.}

\proseitem{29}{อปราปิ อฏฺฐ จริยาโย, อธิมุจฺจนฺโต สทฺธาย จรติ, ปคฺคณฺหนฺโต วีริเยน จรติ, อุปฏฺฐาเปนฺโต สติยา จรติ, อวิกฺเขปํ กโรนฺโต สมาธินา จรติ, ปชานนฺโต ปญฺญาย จรติ, วิชานนฺโต วิญฺญาณจริยาย จรติ, เอวํ ปฏิปนฺนสฺส กุสลา ธมฺมา อายาเปนฺตีติ อายตนจริยาย จรติ, เอวํ ปฏิปนฺโน วิเสสมธิคจฺฉตีติ วิเสสจริยาย จรติ.\csromanspacer{} อิมา อฏฺฐ จริยาโย.}

\prose{อปราปิ อฏฺฐ จริยาโย, ทสฺสนจริยา จ สมฺมาทิฏฺฐิยา, อภินิโรปนจริยา จ สมฺมาสงฺกปฺปสฺส, ปริคฺคหจริยา จ สมฺมาวาจาย, สมุฏฺฐานจริยา จ สมฺมากมฺมนฺตสฺส, โวทานจริยา จ สมฺมา-อาชีวสฺส, ปคฺคหจริยา จ สมฺมาวายามสฺส, อุปฏฺฐานจริยา จ สมฺมาสติยา, อวิกฺเขปจริยา จ สมฺมาสมาธิสฺส.\csromanspacer{} อิมา อฏฺฐ จริยาโยติ.}

\nitthitamruled{จริยากถา นิฏฺฐิตา.}
\csromanheader{2}{6. ปาฏิหาริยกถา}
\proseitem{30}{ตีณิมานิ ภิกฺขเว\footnote{อํ 1. 170 ปิฏฺเฐ ปสฺสิตพฺพา.} ปาฏิหาริยานิ.\csromanspacer{} กตมานิ ตีณิ? อิทฺธิปาฏิหาริยํ อาเทสนาปาฏิหาริยํ อนุสาสนีปาฏิหาริยํ.}

\prose{\textbf{กตมญฺจ ภิกฺขเว อิทฺธิปาฏิหาริยํ?} อิธ ภิกฺขเว เอกจฺโจ อเนกวิหิตํ อิทฺธิวิธํ ปจฺจนุโภติ, เอโกปิ หุตฺวา พหุธา โหติ, พหุธาปิ หุตฺวา เอโก โหติ, อาวิภาวํ ติโรภาวํ ฯเปฯ ยาว พฺรหฺมโลกาปิ กาเยน วสํ วตฺเตติ.\csromanspacer{} อิทํ วุจฺจติ ภิกฺขเว อิทฺธิปาฏิหาริยํ. (1)}

\prose{\csromanfolio{402}\textbf{กตมญฺจ ภิกฺขเว อาเทสนาปาฏิหาริยํ?} อิธ ภิกฺขเว เอกจฺโจ นิมิตฺเตน อาทิสติ “เอวมฺปิ เต มโน, อิตฺถมฺปิ เต มโน, อิติปิ เต จิตฺตนฺ”ติ.\csromanspacer{} โส พหุํ เจปิ อาทิสติ, ตเถว ตํ โหติ, โน อญฺญถา.\csromanspacer{}\csromanspacer{}อิธ ปน ภิกฺขเว เอกจฺโจ น เหว โข นิมิตฺเตน อาทิสติ, อปิ จ โข มนุสฺสานํ วา อมนุสฺสานํ วา เทวตานํ วา สทฺทํ สุตฺวา อาทิสติ “เอวมฺปิ เต มโน, อิตฺถมฺปิ เต มโน, อิติปิ เต จิตฺตนฺ”ติ.\csromanspacer{} โส พหุํ เจปิ อาทิสติ, ตเถว ตํ โหติ, โน อญฺญถา.\csromanspacer{}\csromanspacer{}อิธ ปน ภิกฺขเว เอกจฺโจ น เหว โข นิมิตฺเตน อาทิสติ, นปิ มนุสฺสานํ วา อมนุสฺสานํ วา เทวตานํ วา สทฺทํ สุตฺวา อาทิสติ, อปิ จ โข วิตกฺกยโต วิจารยโต วิตกฺกวิปฺผารสทฺทํ สุตฺวา อาทิสติ “เอวมฺปิ เต มโน, อิตฺถมฺปิ เต มโน, อิติปิ เต จิตฺตนฺ”ติ.\csromanspacer{} โส พหุํ เจปิ อาทิสติ, ตเถว ตํ โหติ, โน อญฺญถา.\csromanspacer{}\csromanspacer{}อิธ ปน ภิกฺขเว เอกจฺโจ นเหว โข นิมิตฺเตน อาทิสติ, นปิ มนุสฺสานํ วา อมนุสฺสานํ วา เทวตานํ วา สทฺทํ สุตฺวา อาทิสติ, นปิ วิตกฺกยโต วิจารยโต วิตกฺกวิปฺผารสทฺทํ สุตฺวา อาทิสติ, อปิ จ โข อวิตกฺกํ อวิจารํ สมาธิํ สมาปนฺนสฺส เจตสา เจโต ปริจฺจ ปชานาติ “ยถา อิมสฺส โภโต มโนสงฺขารา ปณิหิตา อิมสฺส\footnote{ตถา อิมสฺส (สฺยา), ยถา อิมสฺส (ก) อํ 1. 171 ปิฏฺเฐ ปสฺสิตพฺพา.} จิตฺตสฺส อนนฺตรา อมุกํ นาม วิตกฺกํ วิตกฺกยิสฺสตี”ติ\footnote{วิตกฺเกสฺสตีติ (สฺยา) อํ 1. 171 ปิฏฺเฐปิ.}.\csromanspacer{} โส พหุํ เจปิ อาทิสติ, ตเถว ตํ โหติ, โน อญฺญถา.\csromanspacer{} อิทํ วุจฺจติ ภิกฺขเว อาเทสนาปาฏิหาริยํ. (2)}

\prose{\textbf{กตมญฺจ ภิกฺขเว อนุสาสนีปาฏิหาริยํ?} อิธ ภิกฺขเว เอกจฺโจ เอวมนุสาสติ “เอวํ วิตกฺเกถ, มา เอวํ วิตกฺกยิตฺถ.\csromanspacer{} เอวํ มนสิ กโรถ, มา เอวํ มนสากตฺถ.\csromanspacer{} อิทํ ปชหถ, อิทํ อุปสมฺปชฺช วิหรถา”ติ, อิทํ วุจฺจติ ภิกฺขเว อนุสาสนีปาฏิหาริยํ.\csromanspacer{} อิมานิ โข ภิกฺขเว ตีณิ ปาฏิหาริยานิ. (3)}

\proseitem{31}{เนกฺขมฺมํ อิชฺฌตีติ อิทฺธิ, กามจฺฉนฺทํ ปฏิหรตีติ ปาฏิหาริยํ.\csromanspacer{} เย เตน เนกฺขมฺเมน สมนฺนาคตา, สพฺเพ เต วิสุทฺธจิตฺตา อนาวิลสงฺกปฺปาติ อาเทสนาปาฏิหาริยํ. “ตํ โข ปน เนกฺขมฺมํ เอวํ อาเสวิตพฺพํ เอวํ}

\csromanheader{2}{6. ปาฏิหาริยกถา}
\prose{\csromanfolio{403}ภาเวตพฺพํ, เอวํ พหุลีกาตพฺพํ, เอวํ ตทนุธมฺมตา สติ อุปฏฺฐาเปตพฺพา”ติ อนุสาสนีปาฏิหาริยํ. (1)}

\prose{อพฺยาปาโท อิชฺฌตีติ อิทฺธิ, พฺยาปาทํ ปฏิหรตีติ ปาฏิหาริยํ.\csromanspacer{} เย เตน อพฺยาปาเทน สมนฺนาคตา, สพฺเพ เต วิสุทฺธจิตฺตา อนาวิลสงฺกปฺปาติ อาเทสนาปาฏิหาริยํ. “โส โข ปน อพฺยาปาโท เอวํ อาเสวิตพฺโพ, เอวํ ภาเวตพฺโพ, เอวํ พหุลีกาตพฺโพ, เอวํ ตทนุธมฺมตา สติ อุปฏฺฐาเปตพฺพา”ติ อนุสาสนีปาฏิหาริยํ. (2)}

\prose{อาโลกสญฺญา อิชฺฌตีติ อิทฺธิ, ถินมิทฺธํ ปฏิหรตีติ ปาฏิหาริยํ.\csromanspacer{} เย ตาย อาโลกสญฺญาย สมนฺนาคตา, สพฺเพ เต วิสุทฺธจิตฺตา อนาวิลสงฺกปฺปาติ อาเทสนาปาฏิหาริยํ. “สา โข ปน อาโลกสญฺญา เอวํ อาเสวิตพฺพา, เอวํ ภาเวตพฺพา, เอวํ พหุลีกาตพฺพา, เอวํ ตทนุธมฺมตา สติ อุปฏฺฐาเปตพฺพา”ติ อนุสาสนีปาฏิหาริยํ. (3)}

\prose{อวิกฺเขโป อิชฺฌตีติ อิทฺธิ, อุทฺธจฺจํ ปฏิหรตีติ ปาฏิหาริยํ.\csromanspacer{} เย เตน อวิกฺเขเปน สมนฺนาคตา, สพฺเพ เต วิสุทฺธจิตฺตา อนาวิลสงฺกปฺปาติ อาเทสนาปาฏิหาริยํ. “โส โข ปน อวิกฺเขโป เอวํ อาเสวิตพฺโพ, เอวํ ภาเวตพฺโพ, เอวํ พหุลีกาตพฺโพ, เอวํ ตทนุธมฺมตา สติ อุปฏฺฐาเปตพฺพา”ติ อนุสาสนีปาฏิหาริยํ. (4)}

\prose{ธมฺมววตฺถานํ อิชฺฌตีติ อิทฺธิ, วิจิกิจฺฉํ ปฏิหรตีติ ปาฏิหาริยํ.\csromanspacer{} เย เตน ธมฺมววตฺถาเนน สมนฺนาคตา, สพฺเพ เต วิสุทฺธจิตฺตา อนาวิลสงฺกปฺปาติ อาเทสนาปาฏิหาริยํ. “ตํ โข ปน ธมฺมววตฺถานํ เอวํ อาเสวิตพฺพํ, เอวํ ภาเวตพฺพํ, เอวํ พหุลีกาตพฺพํ, เอวํ ตทนุธมฺมตา สติ อุปฏฺฐาเปตพฺพา”ติ อนุสาสนีปาฏิหาริยํ. (5)}

\prose{ญาณํ อิชฺฌตีติ อิทฺธิ, อวิชฺชํ ปฏิหรตีติ ปาฏิหาริยํ.\csromanspacer{} เย เตน ญาเณน สมนฺนาคตา, สพฺเพ เต วิสุทฺธจิตฺตา อนาวิลสงฺกปฺปาติ อาเทสนาปาฏิหาริยํ. “ตํ โข ปน ญาณํ เอวํ อาเสวิตพฺพํ, เอวํ ภาเวตพฺพํ, เอวํ พหุลีกาตพฺพํ, เอวํ ตทนุธมฺมตา สติ อุปฏฺฐาเปตพฺพา”ติ อนุสาสนีปาฏิหาริยํ. (6)}

\prose{ปาโมชฺชํ อิชฺฌตีติ อิทฺธิ, อรติํ ปฏิหรตีติ ปาฏิหาริยํ.\csromanspacer{} เย เตน ปาโมชฺเชน สมนฺนาคตา, สพฺเพ เต วิสุทฺธจิตฺตา อนาวิลสงฺกปฺปาติ อาเทสนาปาฏิหาริยํ. “ตํ โข ปน ปาโมชฺชํ เอวํ อาเสวิตพฺพํ, เอวํ}

\prose{\csromanfolio{404}ภาเวตพฺพํ, เอวํ พหุลีกาตพฺพํ, เอวํ ตทนุธมฺมตา สติ อุปฏฺฐาเปตพฺพา”ติ อนุสาสนีปาฏิหาริยํ ฯเปฯ. (7)}

\prose{ปฐมํ ฌานํ อิชฺฌตีติ อิทฺธิ, นีวรเณ ปฏิหรตีติ ปาฏิหาริยํ.\csromanspacer{} เย เตน ปฐเมน ฌาเนน สมนฺนาคตา, สพฺเพ เต วิสุทฺธจิตฺตา อนาวิลสงฺกปฺปาติ อาเทสนาปาฏิหาริยํ. “ตํ โข ปน ปฐมํ ฌานํ เอวํ อาเสวิตพฺพํ, เอวํ ภาเวตพฺพํ, เอวํ พหุลีกาตพฺพํ, เอวํ ตทนุธมฺมตา สติ อุปฏฺฐาเปตพฺพา”ติ อนุสาสนีปาฏิหาริยํ ฯเปฯ. (26, 33)}

\prose{อรหตฺตมคฺโค อิชฺฌตีติ อิทฺธิ, สพฺพกิเลเส ปฏิหรตีติ ปาฏิหาริยํ.\csromanspacer{} เย เตน อรหตฺตมคฺเคน สมนฺนาคตา, สพฺเพ เต วิสุทฺธจิตฺตา อนาวิลสงฺกปฺปาติ อาเทสนาปาฏิหาริยํ. “โส โข ปน อรหตฺตมคฺโค เอวํ อาเสวิตพฺโพ, เอวํ ภาเวตพฺโพ, เอวํ พหุลีกาตพฺโพ, เอวํ ตทนุธมฺมตา สติ อุปฏฺฐาเปตพฺพา”ติ อนุสาสนีปาฏิหาริยํ. (4, 37)}

\proseitem{32}{เนกฺขมฺมํ อิชฺฌตีติ อิทฺธิ, กามจฺฉนฺทํ ปฏิหรตีติ ปาฏิหาริยํ.\csromanspacer{} ยา จ อิทฺธิ ยญฺจ ปาฏิหาริยํ, อิทํ วุจฺจติ อิทฺธิปาฏิหาริยํ.\csromanspacer{}\csromanspacer{}อพฺยาปาโท อิชฺฌตีติ อิทฺธิ, พฺยาปาทํ ปฏิหรตีติ ปาฏิหาริยํ.\csromanspacer{} ยา จ อิทฺธิ ยญฺจ ปาฏิหาริยํ, อิทํ วุจฺจติ อิทฺธิปาฏิหาริยํ.\csromanspacer{}\csromanspacer{}อาโลกสญฺญา อิชฺฌตีติ อิทฺธิ, ถินมิทฺธํ ปฏิหรตีติ ปาฏิหาริยํ ฯเปฯ.\csromanspacer{} อรหตฺตมคฺโค อิชฺฌตีติ อิทฺธิ, สพฺพกิเลเส ปฏิหรตีติ ปาฏิหาริยํ.\csromanspacer{} ยา จ อิทฺธิ ยญฺจ ปาฏิหาริยํ, อิทํ วุจฺจติ อิทฺธิปาฏิหาริยนฺติ.}

\nitthitamruled{ปาฏิหาริยกถา นิฏฺฐิตา.}
\csromanmatikaanchor{mtk.130}
\csromantocmarkref{section}{7. สมสีสกถา}{mtk.130}
\bookmark[dest={mtk.130},level=1]{7. สมสีสกถา}
\csromanheader{1}{7. สมสีสกถา}
\proseitem{33}{สพฺพธมฺมานํ\footnote{สพฺเพสํ ธมฺมานํ (ก) 98 ปิฏฺเฐ ปสฺสิตพฺพา.} สมฺมา สมุจฺเฉเท นิโรเธ จ อนุปฏฺฐานตา ปญฺญา สมสีสฏฺเฐ ญาณํ.}

\prose{\textbf{สพฺพธมฺมานนฺ}ติ ปญฺจกฺขนฺธา ทฺวาทสายตนานิ อฏฺฐารส ธาตุโย กุสลา ธมฺมา อกุสลา ธมฺมา อพฺยากตา ธมฺมา กามาวจรา ธมฺมา รูปาวจรา ธมฺมา อรูปาวจรา ธมฺมา อปริยาปนฺนา ธมฺมา.\csromanspacer{}\csromanspacer{}\textbf{สมฺมา สมุจฺเฉเท}ติ}

\csromanheader{2}{7. สมสีสกถา}
\prose{\csromanfolio{405}เนกฺขมฺเมน กามจฺฉนฺทํ สมฺมา สมุจฺฉินฺทติ, อพฺยาปาเทน พฺยาปาทํ สมฺมา สมุจฺฉินฺทติ, อาโลกสญฺญาย ถินมิทฺธํ สมฺมา สมุจฺฉินฺทติ, อวิกฺเขเปน อุทฺธจฺจํ สมฺมา สมุจฺฉินฺทติ, ธมฺมววตฺถาเนน วิจิกิจฺฉํ สมฺมา สมุจฺฉินฺทติ, ญาเณน อวิชฺชํ สมฺมา สมุจฺฉินฺทติ, ปาโมชฺเชน อรติํ สมฺมา สมุจฺฉินฺทติ, ปฐเมน ฌาเนน นีวรเณ สมฺมา สมุจฺฉินฺทติ ฯเปฯ อรหตฺตมคฺเคน สพฺเพกิเลเส สมฺมา สมุจฺฉินฺทติ.}

\prose{\textbf{นิโรเธ}ติ เนกฺขมฺเมน กามจฺฉนฺทํ นิโรเธติ, อพฺยาปาเทน พฺยาปาทํ นิโรเธติ, อาโลกสญฺญาย ถินมิทฺธํ นิโรเธติ, อวิกฺเขเปน อุทฺธจฺจํ นิโรเธติ, ธมฺมววตฺถาเนน วิจิกิจฺฉํ นิโรเธติ, ญาเณน อวิชฺชํ นิโรเธติ ปาโมชฺเชน อรติํ นิโรเธติ, ปฐเมน ฌาเนน นีวรเณ นิโรเธติ ฯเปฯ อรหตฺตมคฺเคน สพฺพกิเลเส นิโรเธติ.}

\prose{\textbf{อนุปฏฺฐานตา}ติ เนกฺขมฺมํ ปฏิลทฺธสฺส กามจฺฉนฺโท น อุปฏฺฐาติ, อพฺยาปาทํ ปฏิลทฺธสฺส พฺยาปาโท น อุปฏฺฐาติ, อาโลกสญฺญํ ปฏิลทฺธสฺส ถินมิทฺธํ น อุปฏฺฐาติ, อวิกฺเขปํ ปฏิลทฺธสฺส อุทฺธจฺจํ น อุปฏฺฐาติ, ธมฺมววตฺถานํ ปฏิลทฺธสฺส วิจิกิจฺฉา น อุปฏฺฐาติ, ญาณํ ปฏิลทฺธสฺส อวิชฺชา น อุปฏฺฐาติ, ปาโมชฺชํ ปฏิลทฺธสฺส อรติ น อุปฏฺฐาติ, ปฐมํ ฌานํ ปฏิลทฺธสฺส นีวรณา น อุปฏฺฐนฺติ ฯเปฯ อรหตฺตมคฺคํ ปฏิลทฺธสฺส สพฺพกิเลสา น อุปฏฺฐนฺติ.}

\prose{\textbf{สมนฺ}ติ กามจฺฉนฺทสฺส ปหีนตฺตา เนกฺขมฺมํ สมํ, พฺยาปาทสฺส ปหีนตฺตา อพฺยาปาโท สมํ, ถินมิทฺธสฺส ปหีนตฺตา อาโลกสญฺญา สมํ, อุทฺธจฺจสฺส ปหีนตฺตา อวิกฺเขโป สมํ, วิจิกิจฺฉาย ปหีนตฺตา ธมฺมววตฺถานํ สมํ, อวิชฺชาย ปหีนตฺตา ญาณํ สมํ, อรติยา ปหีนตฺตา ปาโมชฺชํ สมํ, นีวรณานํ ปหีนตฺตา ปฐมํ ฌานํ สมํ ฯเปฯ สพฺพกิเลสานํ ปหีนตฺตา อรหตฺตมคฺโค สมํ.}

\prose{\textbf{สีสนฺ}ติ เตรส สีสานิ, ปลิโพธสีสญฺจ ตณฺหา, วินิพนฺธนสีสญฺจ มาโน, ปรามาสสีสญฺจ ทิฏฺฐิ, วิกฺเขปสีสญฺจ อุทฺธจฺจํ, สํกิเลสสีสญฺจ อวิชฺชา, อธิโมกฺขสีสญฺจ สทฺธา, ปคฺคหสีสญฺจ วีริยํ, อุปฏฺฐานสีสญฺจ สติ, อวิกฺเขปสีสญฺจ สมาธิ, ทสฺสนสีสญฺจ ปญฺญา, ปวตฺตสีสญฺจ ชีวิตินฺทฺริยํ, โคจรสีสญฺจ วิโมกฺโข, สงฺขารสีสญฺจ นิโรโธติ.}

\nitthitam{สมสีสกถา นิฏฺฐิตา.}
\csromanmatikaanchor{mtk.131}
\csromantocmarkref{section}{8. สติปฏฺฐานกถา}{mtk.131}
\bookmark[dest={mtk.131},level=1]{8. สติปฏฺฐานกถา}
\csromanheader{1}{8. สติปฏฺฐานกถา}
\proseitem{34}{\csromanfolio{406}สาวตฺถินิทานํ.\csromanspacer{}\csromanspacer{}จตฺตาโรเม ภิกฺขเว สติปฏฺฐานา.\csromanspacer{} กตเม จตฺตาโร? อิธ ภิกฺขเว ภิกฺขุ กาเย กายานุปสฺสี วิหรติ อาตาปี สมฺปชาโน สติมา วิเนยฺย โลเก อภิชฺฌาโทมนสฺสํ, เวทนาสุ.\csromanspacer{} จิตฺเต.\csromanspacer{} ธมฺเมสุ ธมฺมานุปสฺสี วิหรติ อาตาปี สมฺปชาโน สติมา วิเนยฺย โลเก อภิชฺฌาโทมนสฺสํ.\csromanspacer{} อิเม โข ภิกฺขเว จตฺตาโร สติปฏฺฐานาติ.}

\proseitem{35}{(ก) กถํ \textbf{กาเย กายานุปสฺสี วิหรติ?} อิเธกจฺโจ ปถวีกายํ อนิจฺจโต อนุปสฺสติ, โน นิจฺจโต, ทุกฺขโต อนุปสฺสติ, โน สุขโต, อนตฺตโต อนุปสฺสติ, โน อตฺตโต, นิพฺพินฺทติ, โน นนฺทติ, วิรชฺชติ, โน รชฺชติ, นิโรเธติ, โน สมุเทติ, ปฏินิสฺสชฺชติ, โน อาทิยติ.\csromanspacer{} อนิจฺจโต อนุปสฺสนฺโต นิจฺจสญฺญํ ปชหติ, ทุกฺขโต อนุปสฺสนฺโต สุขสญฺญํ ปชหติ, อนตฺตโต อนุปสฺสนฺโต อตฺตสญฺญํ ปชหติ, นิพฺพินฺทนฺโต นนฺทิํ ปชหติ, วิรชฺชนฺโต ราคํ ปชหติ, นิโรเธนฺโต สมุทยํ ปชหติ, ปฏินิสฺสชฺชนฺโต อาทานํ ปชหติ, อิเมหิ สตฺตหิ อากาเรหิ กายํ อนุปสฺสติ.\csromanspacer{} กาโย อุปฏฺฐานํ, โน สติ, สติ อุปฏฺฐานญฺเจว สติ จ, ตาย สติยา เตน ญาเณน ตํ กายํ อนุปสฺสติ, เตน วุจฺจติ กาเย กายานุปสฺสนาสติปฏฺฐานา.}

\prose{\textbf{ภาวนา}ติ จตสฺโส ภาวนา, ตตฺถ ชาตานํ ธมฺมานํ อนติวตฺตนฏฺเฐน ภาวนา, อินฺทฺริยานํ เอกรสฏฺเฐน ภาวนา, ตทุปควีริยวาหนฏฺเฐน ภาวนา, อาเสวนฏฺเฐน ภาวนา.}

\prose{อิเธกจฺโจ อาโปกายํ ฯเปฯ เตโชกายํ.\csromanspacer{} วาโยกายํ.\csromanspacer{} เกสกายํ.\csromanspacer{} โลมกายํ.\csromanspacer{} ฉวิกายํ.\csromanspacer{} จมฺมกายํ.\csromanspacer{} มํสกายํ.\csromanspacer{} รุธิรกายํ.\csromanspacer{} นฺหารุกายํ\footnote{นหารุกายํ (สฺยา)}.\csromanspacer{} อฏฺฐิกายํ.\csromanspacer{} อฏฺฐิมิญฺชกายํ อนิจฺจโต อนุปสฺสติ, โน นิจฺจโต, ทุกฺขโต อนุปสฺสติ, โน สุขโต, อนตฺตโต อนุปสฺสติ, โน อตฺตโต, นิพฺพินฺทติ, โน นนฺทติ, วิรชฺชติ, โน รชฺชติ, นิโรเธติ, โน สมุเทติ, ปฏินิสฺสชฺชติ, โน อาทิยติ.\csromanspacer{} อนิจฺจโต อนุปสฺสนฺโต นิจฺจสญฺญํ ปชหติ, ทุกฺขโต อนุปสฺสนฺโต สุขสญฺญํ ปชหติ, อนตฺตโต อนุปสฺสนฺโต อตฺตสญฺญํ ปชหติ,}

\csromanheader{2}{8. สติปฏฺฐานกถา}
\prose{\csromanfolio{407}นิพฺพินฺทนฺโต นนฺทิํ ปชหติ, วิรชฺชนฺโต ราคํ ปชหติ, นิโรเธนฺโต สมุทยํ ปชหติ, ปฏินิสฺสชฺชนฺโต อาทานํ ปชหติ.\csromanspacer{} อิเมหิ สตฺตหิ อากาเรหิ กายํ อนุปสฺสภิ.\csromanspacer{} กาโย อุปฏฺฐานํ, โน สติ, สติ อุปฏฺฐานญฺเจว สติ จ, ตาย สติยา เตเน ญาเณน ตํ กายํ อนุปสฺสติ, เตน วุจฺจติ กาเย กายานุปสฺสนาสติปฏฺฐานา.}

\prose{\textbf{ภาวนา}ติ จตสฺโส ภาวนา, ตตฺถ ชาตานํ ธมฺมานํ อนติวตฺตนฏฺเฐน ภาวนา, อินฺทฺริยานํ เอกรสฏฺเฐน ภาวนา, ตทุปควีริยวาหนฏฺเฐน ภาวนา, อาเสวนฏฺเฐน ภาวนา, เอวํ กาเย กายานุปสฺสี วิหรติ. (1)}

\prose{(ข) กถํ \textbf{เวทนาสุ เวทนานุปสฺสี วิหรติ?} อิเธกจฺโจ สุขํ เวทนํ อนิจฺจโต อนุปสฺสติ, โน นิจฺจโต ฯเปฯ ปฏินิสฺสชฺชติ, โน อาทิยติ.\csromanspacer{} อนิจฺจโต อนุปสฺสนฺโต นิจฺจสญฺญํ ปชหติ ฯเปฯ ปฏินิสฺสชฺชนฺโต อาทานํ ปชหติ.\csromanspacer{} อิเมหิ สตฺตหิ อากาเรหิ เวทนํ อนุปสฺสติ.\csromanspacer{} เวทนา อุปฏฺฐานํ, โน สติ, สติ อุปฏฺฐานญฺเจว สติ จ, ตาย สติยา เตน ญาเณน ตํ เวทนํ อนุปสฺสติ, เตน วุจฺจติ เวทนาสุ เวทนานุปสฺสนาสติปฏฺฐานา.}

\prose{\textbf{ภาวนา}ติ จตสฺโส ภาวนา ฯเปฯ อาเสวนฏฺเฐน ภาวนา ฯเปฯ.\csromanspacer{} อิเธกจฺโจ ทุกฺขํ เวทนํ ฯเปฯ อทุกฺขมสุขํ เวทนํ.\csromanspacer{} สามิสํ สุขํ เวทนํ.\csromanspacer{} นิรามิสํ สุขํ เวทนํ.\csromanspacer{} สามิสํ ทุกฺขํ เวทนํ.\csromanspacer{} นิรามิสํ ทุกฺขํ เวทนํ.\csromanspacer{} สามิสํ อทุกฺขมสุขํ เวทนํ.\csromanspacer{} นิรามิสํ อทุกฺขมสุขํ เวทนํ.\csromanspacer{} จกฺขุสมฺผสฺสชํ เวทนํ.\csromanspacer{} โสตสมฺผสฺสชํ เวทนํ.ฆานสมฺผสฺสชํ เวทนํ.\csromanspacer{} ชิวฺหาสมฺผสฺสชํ เวทนํ.\csromanspacer{} กายสมฺผสฺสชํ เวทนํ.\csromanspacer{} มโนสมฺผสฺสชํ เวทนํ อนิจฺจโต อนุปสฺสติ, โน นิจฺจโต ฯเปฯ ปฏินิสฺสชฺชติ, โน อาทิยติ.\csromanspacer{} อนิจฺจโต อนุปสฺสนฺโต นิจฺจสญฺญํ ปชหติ ฯเปฯ ปฏินิสฺสชฺชนฺโต อาทานํ ปชหติ.\csromanspacer{} อิเมหิ สตฺตหิ อากาเรหิ เวทนํ อนุปสฺสติ.\csromanspacer{} เวทนา อุปฏฺฐานํ, โน สติ, สติ อุปฏฺฐานญฺเจว สติ จ, ตาย สติยา เตน ญาเณน ตํ เวทนํ อนุปสฺสติ, เตน วุจฺจติ เวทนาสุ เวทนานุปสฺสนาสติปฏฺฐานา.}

\prose{\textbf{ภาวนา}ติ จตสฺโส ภาวนา ฯเปฯ เอวํ เวทนาสุ เวทนานุปสฺสี วิหรติ. (2)}

\prose{\csromanfolio{408}(ค) กถํ \textbf{จิตฺเต จิตฺตานุปสฺสี วิหรติ?} อิเธกจฺโจ สราคํ จิตฺตํ อนิจฺจโต อนุปสฺสติ, โน นิจฺจโต ฯเปฯ ปฏินิสฺสชฺชติ, โน อาทิยติ.\csromanspacer{} อนิจฺจโต อนุปสฺสนฺโต นิจฺจสญฺญํ ปชหติ ฯเปฯ ปฏินิสฺสชฺชนฺโต อาทานํ ปชหติ.\csromanspacer{} อิเมหิ สตฺตหิ อากาเรหิ จิตฺตํ อนุปสฺสติ, จิตฺตํ อุปฏฺฐานํ, โน สติ, สติ อุปฏฺฐานญฺเจว สติ จ, ตาย สติยา เตน ญาเณน ตํ จิตฺตํ อนุปสฺสติ, เตน วุจฺจติ จิตฺเต จิตฺตานุปสฺสนาสติปฏฺฐานา.}

\prose{\textbf{ภาวนา}ติ จตสฺโส ภาวนา ฯเปฯ อาเสวนฏฺเฐน ภาวนา.}

\prose{อิเธกจฺโจ วีตราคํ จิตฺตํ ฯเปฯ สโทสํ จิตฺตํ.\csromanspacer{} วีตโทสํ จิตฺตํ.\csromanspacer{} สโมหํ จิตฺตํ.\csromanspacer{} วีตโมหํ จิตฺตํ.\csromanspacer{} สํขิตฺตํ จิตฺตํ.\csromanspacer{} วิกฺขิตฺตํ จิตฺตํ.\csromanspacer{} มหคฺคตํ จิตฺตํ.\csromanspacer{} อมหคฺคตํ จิตฺตํ.\csromanspacer{} ส-อุตฺตรํ จิตฺตํ.\csromanspacer{} อนุตฺตรํ จิตฺตํ.\csromanspacer{} สมาหิตํ จิตฺตํ.\csromanspacer{} อสมาหิตํ จิตฺตํ.\csromanspacer{} วิมุตฺตํ จิตฺตํ.\csromanspacer{} อวิมุตฺตํ จิตฺตํ.\csromanspacer{} จกฺขุวิญฺญาณํ.\csromanspacer{} โสตวิญฺญาณํ.\csromanspacer{} ฆานวิญฺญาณํ.\csromanspacer{} ชิวฺหาวิญฺญาณํ.\csromanspacer{} กายวิญฺญาณํ.\csromanspacer{} มโนวิญฺญาณํ อนิจฺจโต อนุปสฺสติ, โน นิจฺจโต ฯเปฯ ปฏินิสฺสชฺชติ, โน อาทิยติ.\csromanspacer{} อนิจฺจโต อนุปสฺสนฺโต นิจฺจสญฺญํ ปชหติ ฯเปฯ ปฏินิสฺสชฺชนฺโต อาทานํ ปชหติ.\csromanspacer{} อิเมหิ สตฺตหิ อากาเรหิ จิตฺตํ อนุปสฺสติ.\csromanspacer{} จิตฺตํ อุปฏฺฐานํ โน สติ, สติ อุปฏฺฐานญฺเจว สติ จ, ตาย สติยา เตน ญาเณน ตํ จิตฺตํ อนุปสฺสติ, เตน วุจฺจติ จิตฺเต จิตฺตานุปสฺสนาสติปฏฺฐานา.}

\prose{\textbf{ภาวนา}ติ จตสฺโส ภาวนา ฯเปฯ อาเสวนฏฺเฐน ภาวนา.\csromanspacer{} เอวํ จิตฺเต จิตฺตานุปสฺสี วิหรติ. (3)}

\prose{(ฆ) กถํ ธมฺมฺ\textbf{เอสุ ธมฺมานุปสฺสี วิหรติ?} อิเธกจฺโจ ฐเปตฺวา กายํ ฐเปตฺวา เวทนํ ฐปตฺวา จิตฺตํ ตทวเสเส ธมฺเม อนิจฺจโต อนุปสฺสติ, โน นิจฺจโต, ทุกฺขโต อนุปสฺสติ, โน สุขโต, อนตฺตโต อนุปสฺสติ, โน อตฺตโต, นิพฺพินฺทติ, โน นนฺทติ, วิรชฺชติ, โน รชฺชติ, นิโรเธติ, โน สมุเทติ, ปฏินิสฺสชฺชติ, โน อาทิยติ.\csromanspacer{} อนิจฺจโต อนุปสฺสนฺโต นิจฺจสญฺญํ ปชหติ, ทุกฺขโต อนุปสฺสนฺโต สุขสญฺญํ ปชหติ, อนตฺตโต อนุปสฺสนฺโต อตฺตสญฺญํ ปชหติ, นิพฺพินฺทนฺโต นนฺทิํ ปชหติ, วิรชฺชนฺโต ราคํ ปชหติ, นิโรเธนฺโต สมุทยํ ปชหติ, ปฏินิสฺสชฺชนฺโต อาทานํ ปชหติ.\csromanspacer{} อิเมหิ สตฺตหิ อากาเรหิ เต ธมฺเม อนุปสฺสติ.\csromanspacer{} ธมฺมา อุปฏฺฐานํ, โน สติ, สติ อุปฏฺฐานญฺเจว สติ จ, ตาย สติยา เตน ญาเณน}

\csromanmatikaanchor{mtk.132}
\csromantocmarkref{section}{9. วิปสฺสนากถา}{mtk.132}
\bookmark[dest={mtk.132},level=1]{9. วิปสฺสนากถา}
\csromanheader{1}{9. วิปสฺสนากถา}
\prose{\csromanfolio{409}เต ธมฺเม อนุปสฺสติ, เตน วุจฺจติ ธมฺเมสุ ธมฺมานุปสฺสนาสติปฏฺฐานา.}

\prose{\textbf{ภาวนา}ติ จตสฺโส ภาวนา, ตตฺถ ชาตานํ ธมฺมานํ อนติวตฺตนฏฺเฐน ภาวนา, อินฺทฺริยานํ เอกรสฏฺเฐน ภาวนา, ตทุปควีริยวาหนฏฺเฐน ภาวนา อาเสวนฏฺเฐน ภาวนา.\csromanspacer{} เอวํ ธมฺเมสุ ธมฺมานุปสฺสี วิหรตีติ. (4)}

\nitthitamruled{สติปฏฺฐานกถา นิฏฺฐิตา.}
\csromanheader{2}{9. วิปสฺสนากถา}
\proseitem{36}{เอวํ เม สุตํ\csromandash{}เอกํ สมยํ ภควา สาวตฺถิยํ วิหรติ เชตวเน อนาถปิณฺฑิกสฺส อาราเม.\csromanspacer{} ตตฺร โข ภควา ภิกฺขู อามนฺเตสิ “ภิกฺขโว”ติ. “ภทนฺเต”ติ เต ภิกฺขู ภควโต ปจฺจสฺโสสุํ.\csromanspacer{} ภควา เอตทโวจ\csromandash{}}

\prose{โส วต ภิกฺขเว ภิกฺขุ กญฺจิ\footnote{กิญฺจิ (ก) อํ 2. 385 ปิฏฺเฐ ปสฺสิตพฺพา.} สงฺขารํ นิจฺจโต สมนุปสฺสนฺโต อนุโลมิกาย ขนฺติยา สมนฺนาคโต ภวิสฺสตีติ เนตํ ฐานํ วิชฺชติ, อนุโลมิกาย ขนฺติยา อสมนฺนาคโต สมฺมตฺตนิยามํ โอกฺกมิสฺสตีติ เนตํ ฐานํ วิชฺชติ, สมฺมตฺตนิยามํ อโนกฺกมมาโน โสตาปตฺติผลํ วา สกทาคามิผลํ วา อนาคามิผลํ วา อรหตฺตํ วา\footnote{อรหตฺตผลํ วา (สฺยา, ก) อํ 2. 384 ปิฏฺเฐ ปสฺสิตพฺพา.} สจฺฉิกริสฺสตีติ เนตํ ฐานํ วิชฺชติ. (1)}

\prose{โส วต ภิกฺขเว ภิกฺขุ สพฺพสงฺขาเร อนิจฺจโต สมนุปสฺสนฺโต อนุโลมิกาย ขนฺติยา สมนฺนาคโต ภวิสฺสตีติ ฐานเมตํ วิชฺชติ, อนุโลมิกาย ขนฺติยา สมนฺนาคโต สมฺมตฺตนิยามํ โอกฺกมิสฺสตีติ ฐานเมตํ วิชฺชติ, สมฺมตฺตนิยามํ โอกฺกมมาโน โสตาปตฺติผลํ วา สกทาคามิผลํ วา อนาคามิผลํ วา อรหตฺตํ วา สจฺฉิกริสฺสตีติ ฐานเมตํ วิชฺชติ. (2)}

\prose{โส วต ภิกฺขเว ภิกฺขุ กญฺจิ สงฺขารํ สุขโต สมนุปสฺสนฺโต อนุโลมิกาย ขนฺติยา สมนฺนาคโต ภวิสฺสตีติ เนตํ ฐานํ วิชฺชติ, อนุโลมิกาย ขนฺติยา อสมนฺนาคโต สมฺมตฺตนิยามํ โอกฺกมิสฺสตีติ \csromanfolio{410}เนตํ ฐานํ วิชฺชติ, สมฺมตฺตนิยามํ อโนกฺกมมาโน โสตาปตฺติผลํ วา สกทาคามิผลํ วา อนาคามิผลํ วา อรหตฺตํ วา สจฺฉิกริสฺสตีติ เนตํ ฐานํ วิชฺชติ. (3)}

\prose{โส วต ภิกฺขเว ภิกฺขุ สพฺพสงฺขาเร ทุกฺขโต สมนุปสฺสนฺโต อนุโลมิกาย ขนฺติยา สมนฺนาคโต ภวิสฺสตีติ ฐานเมตํ วิชฺชติ, อนุโลมิกาย ขนฺติยา สมนฺนาคโต สมฺมตฺตนิยามํ โอกฺกมิสฺสตีติ ฐานเมตํ วิชฺชติ, สมฺมตฺตนิยามํ โอกฺกมมาโน โสตาปตฺติผลํ วา สกทาคามิผลํ วา อนาคามิผลํ วา อรหตฺตํ วา สจฺฉิกริสฺสตีติ ฐานเมตํ วิชฺชติ. (4)}

\prose{โส วต ภิกฺขเว ภิกฺขุ กญฺจิ ธมฺมํ อตฺตโต สมนุปสฺสนฺโต อนุโลมิกาย ขนฺติยา สมนฺนาคโต ภวิสฺสตีติ เนตํ ฐานํ วิชฺชติ, อนุโลมิกาย ขนฺติยา อสมนฺนาคโต สมฺมตฺตนิยามํ โอกฺกมิสฺสตีติ เนตํ ฐานํ วิชฺชติ, สมฺมตฺตนิยามํ อโนกฺกมมาโน โสตาปตฺติผลํ วา สกทาคามิผลํ วา อนาคามิผลํ วา อรหตฺตํ วา สจฺฉิกริสฺสตีติ เนตํ ฐานํ วิชฺชติ. (5)}

\prose{โส วต ภิกฺขเว ภิกฺขุ สพฺพธมฺเม\footnote{ตญฺจิ ธมฺมํ (สฺยา), กิญฺจิ ธมฺมํ (ก) อํ 2. 385 ปิฏฺเฐ ปสฺสิตพฺพา.} อนตฺตโต สมนุปสฺสนฺโต อนุโลมิกาย ขนฺติยา สมนฺนาคโต ภวิสฺสตีติ ฐานเมตํ วิชฺชติ, อนุโลมิกาย ขนฺติยา สมนฺนาคโต สมฺมตฺตนิยามํ โอกฺกมิสฺสตีติ ฐานเมตํ วิชฺชติ, สมฺมตฺตนิยามํ โอกฺกมมาโน โสตาปตฺติผลํ วา สกทาคามิผลํ วา อนาคามิผลํ วา อรหตฺตํ วา สจฺฉิกริสฺสตีติ ฐานเมตํ วิชฺชติ. (6)}

\prose{โส วต ภิกฺขเว ภิกฺขุ นิพฺพานํ ทุกฺขโต สมนุปสฺสนฺโต อนุโลมิกาย ขนฺติยา สมนฺนาคโต ภวิสฺสตีติ เนตํ ฐานํ วิชฺชติ, อนุโลมิกาย ขนฺติยา อสมนฺนาคโต สมฺมตฺตนิยามํ โอกฺกมิสฺสตีติ เนตํ ฐานํ วิชฺชติ, สมฺมตฺตนิยามํ อโนกฺกมมาโน โสตาปตฺติผลํ วา สกทาคามิผลํ วา อนาคามิผลํ วา อรหตฺตํ วา สจฺฉิกริสฺสตีติ เนตํ ฐานํ วิชฺชติ. (7)}

\prose{โส วต ภิกฺขเว ภิกฺขุ นิพฺพานํ สุขโต สมนุปสฺสนฺโต อนุโลมิกาย ขนฺติยา สมนฺนาคโต ภวิสฺสตีติ ฐานเมตํ วิชฺชติ,}

\csromanheader{2}{9. วิปสฺสนากถา}
\prose{\csromanfolio{411}อนุโลมิกาย ขนฺติยา สมนฺนาคโต สมฺมตฺตนิยามํ โอกฺกมิสฺสตีติ ฐานเมตํ วิชฺชติ, สมฺมตฺตนิยามํ โอกฺกมมาโน โสตาปตฺติผลํ วา สกทาคามิผลํ วา อนาคามิผลํ วา อรหตฺตํ วา สจฺฉิกริสฺสตีติ ฐานเมตํ วิชฺชติ. (8)}

\proseitem{37}{กติหากาเรหิ อนุโลมิกํ ขนฺติํ ปฏิลภติ, กติหากาเรหิ สมฺมตฺตนิยามํ โอกฺกมติ? จตฺตารีสาย อากาเรหิ อนุโลมิกํ ขนฺติํ ปฏิลภติ, จตฺตารีสาย อากาเรหิ สมฺมตฺตนิยามํ โอกฺกมติ.}

\prose{กตเมหิ จตฺตารีสาย อากาเรหิ อนุโลมิกํ ขนฺติํ ปฏิลภติ, กตเมหิ จตฺตารีสาย อากาเรหิ สมฺมตฺตนิยามํ โอกฺกมติ? ปญฺจกฺขนฺเธ อนิจฺจโต ทุกฺขโต โรคโต คณฺฑโต สลฺลโต อฆโต อาพาธโต ปรโต ปโลกโต อีติโต อุปทฺทวโต ภยโต อุปสคฺคโต จลโต ปภงฺคุโต\footnote{ปภงฺคโต (สฺยา)} อทฺธุวโต อตาณโต\footnote{อตฺตาณโต (สฺยา)} อเลณโต อสรณโต ริตฺตโต ตุจฺฉโต สุญฺญโต อนตฺตโต อาทีนวโต วิปริณามธมฺมโต อสารกโต อฆมูลโต วธกโต วิภวโต สาสวโต สงฺขตโต มารามิสโต ชาติธมฺมโต ชราธมฺมโต พฺยาธิธมฺมโต มรณธมฺมโต โสกธมฺมโต ปริเทวธมฺมโต อุปายาสธมฺมโต สํกิเลสิกธมฺมโต.}

\proseitem{38}{ปญฺจกฺขนฺเธ อนิจฺจโต ปสฺสนฺโต อนุโลมิกํ ขนฺติํ ปฏิลภติ, “ปญฺจนฺนํ ขนฺธานํ นิโรโธ นิจฺจํ นิพฺพานนฺ”ติ ปสฺสนฺโต สมฺมตฺตนิยามํ โอกฺกมติ.\csromanspacer{} ปญฺจกฺขนฺเธ ทุกฺขโต ปสฺสนฺโต อนุโลมิกํ ขนฺติํ ปฏิลภติ, “ปญฺจนฺนํ ขนฺธานํ นิโรโธ สุขํ นิพฺพานนฺ”ติ ปสฺสนฺโต สมฺมตฺตนิยามํ โอกฺกมติ.\csromanspacer{} ปญฺจกฺขนฺเธ โรคโต ปสฺสนฺโต อนุโลมิกํ ขนฺติํ ปฏิลภติ, “ปญฺจนฺนํ ขนฺธานํ นิโรโธ อาโรคฺยํ นิพฺพานนฺ”ติ ปสฺสนฺโต สมฺมตฺตนิยามํ โอกฺกมติ.\csromanspacer{} ปญฺจกฺขนฺเธ คณฺฑโต ปสฺสนฺโต อนุโลมิกํ ขนฺติํ ปฏิลภติ, “ปญฺจนฺนํ ขนฺธานํ นิโรโธ อคณฺฑํ\footnote{นิคณฺโฑ (สฺยา)} นิพฺพานนฺ”ติ ปสฺสนฺโต สมฺมตฺตนิยามํ โอกฺกมติ.\csromanspacer{} ปญฺจกฺขนฺเธ สลฺลโต ปสฺสนฺโต อนุโลมิกํ ขนฺติํ ปฏิลภติ, “ปญฺจนฺนํ ขนฺธานํ นิโรโธ วิสลฺลํ\footnote{นิสฺสลฺลํ (สฺยา)} นิพฺพานนฺ”ติ ปสฺสนฺโต สมฺมตฺตนิยามํ โอกฺกมติ. (5)}

\prose{\csromanfolio{412}ปญฺจกฺขนฺเธ อฆโต ปสฺสนฺโต อนุโลมิกํ ขนฺติํ ปฏิลภติ, “ปญฺจนฺนํ ขนฺธานํ นิโรโธ อนโฆ นิพฺพานนฺ”ติ ปสฺสนฺโต สมฺมตฺตนิยามํ โอกฺกมติ.\csromanspacer{} ปญฺจกฺขนฺเธ อาพาธโต ปสฺสนฺโต อนุโลมิกํ ขนฺติํ ปฏิลภติ, “ปญฺจนฺนํ ขนฺธานํ นิโรโธ อนาพาธํ นิพฺพานนฺ”ติ ปสฺสนฺโต สมฺมตฺตนิยามํ โอกฺกมติ.\csromanspacer{} ปญฺจกฺขนฺเธ ปรโต ปสฺสนฺโต อนุโลมิกํ ขนฺติํ ปฏิลภติ, “ปญฺจนฺนํ ขนฺธานํ นิโรโธ อปรปฺปจฺจยํ นิพฺพานนฺ”ติ ปสฺสนฺโต สมฺมตฺตนิยามํ โอกฺกมติ.\csromanspacer{} ปญฺจกฺขนฺเธ ปโลกโต ปสฺสนฺโต อนุโลมิกํ ขนฺติํ ปฏิลภติ, “ปญฺจนฺนํ ขนฺธานํ นิโรโธ อปโลกธมฺโม นิพฺพานนฺ”ติ ปสฺสนฺโต สมฺมตฺตนิยามํ โอกฺกมติ.\csromanspacer{} ปญฺจกฺขนฺเธ อีติโต ปสฺสนฺโต อนุโลมิกํ ขนฺติํ ปฏิลภติ, “ปญฺจนฺนํ ขนฺธานํ นิโรโธ อนีติกํ นิพฺพานนฺ”ติ ปสฺสนฺโต สมฺมตฺตนิยามํ โอกฺกมติ. (10)}

\prose{ปญฺจกฺขนฺเธ อุปทฺทวโต ปสฺสนฺโต อนุโลมิกํ ขนฺติํ ปฏิลภติ, “ปญฺจนฺนํ ขนฺธานํ นิโรโธ อนุปทฺทวํ นิพฺพานนฺ”ติ ปสฺสนฺโต สมฺมตฺตนิยามํ โอกฺกมติ.\csromanspacer{} ปญฺจกฺขนฺเธ ภยโต ปสฺสนฺโต อนุโลมิกํ ขนฺติํ ปฏิลภติ, “ปญฺจนฺนํ ขนฺธานํ นิโรโธ อภยํ นิพฺพานนฺ”ติ ปสฺสนฺโต สมฺมตฺตนิยามํ โอกฺกมติ.\csromanspacer{} ปญฺจกฺขนฺเธ อุปสคฺคโต ปสฺสนฺโต อนุโลมิกํ ขนฺติํ ปฏิลภติ, “ปญฺจนฺนํ ขนฺธานํ นิโรโธ อนุปสคฺคํ นิพฺพานนฺ”ติ ปสฺสนฺโต สมฺมตฺตนิยามํ โอกฺกมติ.\csromanspacer{} ปญฺจกฺขนฺเธ จลโต ปสฺสนฺโต อนุโลมิกํ ขนฺติํ ปฏิลภติ, “ปญฺจนฺนํ ขนฺธานํ นิโรโธ อจลํ นิพฺพานนฺ”ติ ปสฺสนฺโต สมฺมตฺตนิยามํ โอกฺกมติ.\csromanspacer{} ปญฺจกฺขนฺเธ ปภงฺคุโต ปสฺสนฺโต อนุโลมิกํ ขนฺติํ ปฏิลภติ, “ปญฺจนฺนํ ขนฺธานํ นิโรโธ อปภงฺคุ\footnote{อปฺปภงฺคํ (สฺยา)} นิพฺพานนฺ”ติ ปสฺสนฺโต สมฺมตฺตนิยามํ โอกฺกมติ. (15)}

\prose{ปญฺจกฺขนฺเธ อทฺธุวโต ปสฺสนฺโต อนุโลมิกํ ขนฺติํ ปฏิลภติ, “ปญฺจนฺนํ ขนฺธานํ นิโรโธ ธุวํ นิพฺพานนฺ”ติ ปสฺสนฺโต สมฺมตฺตนิยามํ โอกฺกมติ.\csromanspacer{} ปญฺจกฺขนฺเธ อตาณโต ปสฺสนฺโต อนุโลมิกํ ขนฺติํ ปฏิลภติ, “ปญฺจนฺนํ ขนฺธานํ นิโรโธ ตาณํ นิพฺพานนฺ”ติ ปสฺสนฺโต สมฺมตฺตนิยามํ โอกฺกมติ.\csromanspacer{} ปญฺจกฺขนฺเธ อเลณโต ปสฺสนฺโต อนุโลมิกํ ขนฺติํ ปฏิลภติ, “ปญฺจนฺนํ ขนฺธานํ นิโรโธ เลณํ นิพฺพานนฺ”ติ ปสฺสนฺโต สมฺมตฺตนิยามํ โอกฺกมติ.\csromanspacer{} ปญฺจกฺขนฺเธ อสรณโต ปสฺสนฺโต อนุโลมิกํ ขนฺติํ ปฏิลภติ, “ปญฺจนฺนํ ขนฺธานํ นิโรโธ สรณํ นิพฺพานนฺ”ติ ปสฺสนฺโต สมฺมตฺตนิยามํ โอกฺกมติ.}

\csromanheader{2}{9. วิปสฺสนากถา}
\prose{\csromanfolio{413}ปญฺจกฺขนฺเธ ริตฺตโต ปสฺสนฺโต อนุโลมิกํ ขนฺติํ ปฏิลภติ, “ปญฺจนฺนํ ขนฺธานํ นิโรโธ อริตฺตํ นิพฺพานนฺ”ติ ปสฺสนฺโต สมฺมตฺตนิยามํ โอกฺกมติ. (20)}

\prose{ปญฺจกฺขนฺเธ ตุจฺฉโต ปสฺสนฺโต อนุโลมิกํ ขนฺติํ ปฏิลภติ, “ปญฺจนฺนํ ขนฺธานํ นิโรโธ อตุจฺฉํ นิพฺพานนฺ”ติ ปสฺสนฺโต สมฺมตฺตนิยามํ โอกฺกมติ.\csromanspacer{} ปญฺจกฺขนฺเธ สุญฺญโต ปสฺสนฺโต อนุโลมิกํ ขนฺติํ ปฏิลภติ, “ปญฺจนฺนํ ขนฺธานํ นิโรโธ ปรมสุญฺญํ นิพฺพานนฺ”ติ ปสฺสนฺโต สมฺมตฺตนิยามํ โอกฺกมติ.\csromanspacer{} ปญฺจกฺขนฺเธ อนตฺตโต ปสฺสนฺโต อนุโลมิกํ ขนฺติํ ปฏิลภติ, “ปญฺจนฺนํ ขนฺธานํ นิโรโธ ปรมตฺถํ นิพฺพานนฺ”ติ ปสฺสนฺโต สมฺมตฺตนิยามํ โอกฺกมติ.\csromanspacer{} ปญฺจกฺขนฺเธ อาทีนวโต ปสฺสนฺโต อนุโลมิกํ ขนฺติํ ปฏิลภติ, “ปญฺจนฺนํ ขนฺธานํ นิโรโธ อนาทีนวํ นิพฺพานนฺ”ติ ปสฺสนฺโต สมฺมตฺตนิยามํ โอกฺกมติ.\csromanspacer{} ปญฺจกฺขนฺเธ วิปริณามธมฺมโต ปสฺสนฺโต อนุโลมิกํ ขนฺติํ ปฏิลภติ, “ปญฺจนฺนํ ขนฺธานํ นิโรโธ อวิปริณามธมฺมํ นิพฺพานนฺ”ติ ปสฺสนฺโต สมฺมตฺตนิยามํ โอกฺกมติ. (25)}

\prose{ปญฺจกฺขนฺเธ อสารกโต ปสฺสนฺโต อนุโลมิกํ ขนฺติํ ปฏิลภติ, “ปญฺจนฺนํ ขนฺธานํ นิโรโธ สารํ นิพฺพานนฺ”ติ ปสฺสนฺโต สมฺมตฺตนิยามํ โอกฺกมติ.\csromanspacer{} ปญฺจกฺขนฺเธ อฆมูลโต ปสฺสนฺโต อนุโลมิกํ ขนฺติํ ปฏิลภติ, “ปญฺจนฺนํ ขนฺธานํ นิโรโธ อนฆมูลํ นิพฺพานนฺ”ติ ปสฺสนฺโต สมฺมตฺตนิยามํ โอกฺกมติ.\csromanspacer{} ปญฺจกฺขนฺเธ วธกโต ปสฺสนฺโต อนุโลมิกํ ขนฺติํ ปฏิลภติ, “ปญฺจนฺนํ ขนฺธานํ นิโรโธ อวธกํ นิพฺพานนฺ”ติ ปสฺสนฺโต สมฺมตฺตนิยามํ โอกฺกมติ.\csromanspacer{} ปญฺจกฺขนฺเธ วิภวโต ปสฺสนฺโต อนุโลมิกํ ขนฺติํ ปฏิลภติ, “ปญฺจนฺนํ ขนฺธานํ นิโรโธ อวิภวํ นิพฺพานนฺ”ติ ปสฺสนฺโต สมฺมตฺตนิยามํ โอกฺกมติ.\csromanspacer{} ปญฺจกฺขนฺเธ สาสวโต ปสฺสนฺโต อนุโลมิกํ ขนฺติํ ปฏิลภติ, “ปญฺจนฺนํ ขนฺธานํ นิโรโธ อนาสวํ นิพฺพานนฺ”ติ ปสฺสนฺโต สมฺมตฺตนิยามํ โอกฺกมติ. (30)}

\prose{ปญฺจกฺขนฺเธ สงฺขตโต ปสฺสนฺโต อนุโลมิกํ ขนฺติํ ปฏิลภติ, “ปญฺจนฺนํ ขนฺธานํ นิโรโธ อสงฺขตํ นิพฺพานนฺ”ติ ปสฺสนฺโต สมฺมตฺตนิยามํ โอกฺกมติ.\csromanspacer{} ปญฺจกฺขนฺเธ มารามิสโต ปสฺสนฺโต อนุโลมิกํ ขนฺติํ ปฏิลภติ, “ปญฺจนฺนํ ขนฺธานํ นิโรโธ นิรามิสํ นิพฺพานนฺ”ติ ปสฺสนฺโต สมฺมตฺตนิยามํ โอกฺกมติ.\csromanspacer{} ปญฺจกฺขนฺเธ ชาติธมฺมโต ปสฺสนฺโต อนุโลมิกํ ขนฺติํ ปฏิลภติ, “ปญฺจนฺนํ ขนฺธานํ นิโรโธ อชาตํ นิพฺพานนฺ”ติ ปสฺสนฺโต สมฺมตฺตนิยามํ โอกฺกมติ.\csromanspacer{} ปญฺจกฺขนฺเธ ชราธมฺมโต ปสฺสนฺโต อนุโลมิกํ ขนฺติํ ปฏิลภติ, “ปญฺจนฺนํ}

\prose{\csromanfolio{414}ขนฺธานํ นิโรโธ อชรํ นิพฺพานนฺ”ติ ปสฺสนฺโต สมฺมตฺตนิยามํ โอกฺกมติ.\csromanspacer{} ปญฺจกฺขนฺเธ พฺยาธิธมฺมโต ปสฺสนฺโต อนุโลมิกํ ขนฺติํ ปฏิลภติ, “ปญฺจนฺนํ ขนฺธานํ นิโรโธ อพฺยาธิ\footnote{อพฺยาธิธมฺมํ (สฺยา)} นิพฺพานนฺ”ติ ปสฺสนฺโต สมฺมตฺตนิยามํ โอกฺกมติ. (35)}

\prose{ปญฺจกฺขนฺเธ มรณธมฺมโต ปสฺสนฺโต อนุโลมิกํ ขนฺติํ ปฏิลภติ, “ปญฺจนฺนํ ขนฺธานํ นิโรโธ อมตํ นิพฺพานนฺ”ติ ปสฺสนฺโต สมฺมตฺตนิยามํ โอกฺกมติ.\csromanspacer{} ปญฺจกฺขนฺเธ โสกธมฺมโต ปสฺสนฺโต อนุโลมิกํ ขนฺติํ ปฏิลภติ, “ปญฺจนฺนํ ขนฺธานํ นิโรโธ อโสกํ นิพฺพานนฺ”ติ ปสฺสนฺโต สมฺมตฺตนิยามํ โอกฺกมติ.\csromanspacer{} ปญฺจกฺขนฺเธ ปริเทวธมฺมโต ปสฺสนฺโต อนุโลมิกํ ขนฺติํ ปฏิลภติ, “ปญฺจนฺนํ ขนฺธานํ นิโรโธ อปริเทวํ นิพฺพานนฺ”ติ ปสฺสนฺโต สมฺมตฺตนิยามํ โอกฺกมติ.\csromanspacer{} ปญฺจกฺขนฺเธ อุปายาสธมฺมโต ปสฺสนฺโต อนุโลมิกํ ขนฺติํ ปฏิลภติ, “ปญฺจนฺนํ ขนฺธานํ นิโรโธ อนุปายาสํ นิพฺพานนฺ”ติ ปสฺสนฺโต สมฺมตฺตนิยามํ โอกฺกมติ.\csromanspacer{} ปญฺจกฺขนฺเธ สํกิเลสิกธมฺมโต ปสฺสนฺโต อนุโลมิกํ ขนฺติํ ปฏิลภติ, “ปญฺจนฺนํ ขนฺธานํ นิโรโธ อสํกิลิฏฺฐํ นิพฺพานนฺ”ติ ปสฺสนฺโต สมฺมตฺตนิยามํ โอกฺกมติ. (40)}

\proseitem{39}{อนิจฺจโตติ อนิจฺจานุปสฺสนา.\csromanspacer{} ทุกฺขโตติ ทุกฺขานุปสฺสนา.\csromanspacer{} โรคโตติ ทุกฺขานุปสฺสนา.\csromanspacer{} คณฺฑโตติ ทุกฺขานุปสฺสนา.\csromanspacer{} สลฺลโตติ ทุกฺขานุปสฺสนา.\csromanspacer{} อฆโตติ ทุกฺขานุปสฺสนา.\csromanspacer{} อาพาธโตติ ทุกฺขานุปสฺสนา.\csromanspacer{} ปรโตติ อนตฺตานุปสฺสนา.\csromanspacer{} ปโลกโตติ อนิจฺจานุปสฺสนา. ¢ติโตติ ทุกฺขานุปสฺสนา. (10)}

\prose{อุปทฺทวโตติ ทุกฺขานุปสฺสนา.\csromanspacer{} ภยโตติ ทุกฺขานุปสฺสนา.\csromanspacer{} อุปสคฺคโตติ ทุกฺขานุปสฺสนา.\csromanspacer{} จลโตติ อนิจฺจานุปสฺสนา.\csromanspacer{} ปภงฺคุโตติ อนิจฺจานุปสฺสนา.\csromanspacer{} อทฺธุวโตติ อนิจฺจานุปสฺสนา.\csromanspacer{} อตาณโตติ ทุกฺขานุปสฺสนา.\csromanspacer{} อเลณโตติ ทุกฺขานุปสฺสนา.\csromanspacer{} อสรณโตติ ทุกฺขานุปสฺสนา.\csromanspacer{} ริตฺตโตติ อนตฺตานุปสฺสนา. (10, 20)}

\prose{ตุจฺฉโตติ อนตฺตานุปสฺสนา.\csromanspacer{} สุญฺญโตติ อนตฺตานุปสฺสนา.\csromanspacer{} อนตฺตโตติ อนตฺตานุปสฺสนา.\csromanspacer{} อาทีนวโตติ ทุกฺขานุปสฺสนา.\csromanspacer{} วิปริณามธมฺมโตติ อนิจฺจานุปสฺสนา.\csromanspacer{} อสารกโตติ อนตฺตานุปสฺสนา.}

\csromanmatikaanchor{mtk.133}
\csromantocmarkref{section}{10. มาติกากถา}{mtk.133}
\bookmark[dest={mtk.133},level=1]{10. มาติกากถา}
\csromanheader{1}{10. มาติกากถา}
\prose{\csromanfolio{415}อฆมูลโตติ ทุกฺขานุปสฺสนา.\csromanspacer{} วธกโตติ ทุกฺขานุปสฺสนา.\csromanspacer{} วิภวโตติ อนิจฺจานุปสฺสนา.\csromanspacer{} สาสวโตติ ทุกฺขานุปสฺสนา. (10, 30)}

\prose{สงฺขตโตติ อนิจฺจานุปสฺสนา.\csromanspacer{} มารามิสโตติ ทุกฺขานุปสฺสนา.\csromanspacer{} ชาติธมฺมโตติ ทุกฺขานุปสฺสนา.\csromanspacer{} ชราธมฺมโตติ ทุกฺขานุปสฺสนา.\csromanspacer{} พฺยาธิธมฺมโตติ ทุกฺขานุปสฺสนา.\csromanspacer{} มรณธมฺมโตติ อนิจฺจานุปสฺสนา.\csromanspacer{} โสกธมฺมโตติ ทุกฺขานุปสฺสนา.\csromanspacer{} ปริเทวธมฺมโตติ ทุกฺขานุปสฺสนา.\csromanspacer{} อุปายาสธมฺมโตติ ทุกฺขานุปสฺสนา.\csromanspacer{} สํกิเลสิกธมฺมโตติ ทุกฺขานุปสฺสนา. (10, 40)}

\prose{อิเมหิ จตฺตาลีสาย อากาเรหิ อนุโลมิกํ ขนฺติํ ปฏิลภติ.\csromanspacer{} อิเมหิ จตฺตาลีสาย อากาเรหิ สมฺมตฺตนิยามํ โอกฺกมติ.}

\prose{อิเมหิ จตฺตาลีสาย อากาเรหิ อนุโลมิกํ ขนฺติํ ปฏิลภนฺตสฺส, อิเมหิ จตฺตาลีสาย อากาเรหิ สมฺมตฺตนิยามํ โอกฺกมนฺตสฺส กติ อนิจฺจานุปสฺสนา, กติ ทุกฺขานุปสฺสนา, กติ อนตฺตานุปสฺสนา.}

\csromangathameasure{ปญฺจวีสติ อนตฺตานุปสฺสนา, \\ ปญฺญาส อนิจฺจานุปสฺสนา. \\ สตํ ปญฺจวีสติ เจว, \\ ยานิ ทุกฺเข ปวุจฺจเรติ.}
\csromangathagroup{\csromangathastanza{ปญฺจวีสติ อนตฺตานุปสฺสนา, \\ ปญฺญาส อนิจฺจานุปสฺสนา. \\ สตํ ปญฺจวีสติ เจว, \\ ยานิ ทุกฺเข ปวุจฺจเรติ.}}

\nitthitamruled{วิปสฺสนากถา นิฏฺฐิตา.}
\csromanheader{2}{10. มาติกากถา}
\proseitem{40}{นิจฺฉาโต\footnote{นิจฺฉาโต มุจฺจตีติ (สฺยา)}, โมกฺโข วิโมกฺโข, วิชฺชาวิมุตฺติ, อธิสีลํ, อธิจิตฺตํ, อธิปญฺญา, ปสฺสทฺธิ, ญาณํ, ทสฺสนํ, วิสุทฺธิ\footnote{สุทฺธิ (สฺยา)}, เนกฺขมฺมํ, นิสฺสรณํ, ปวิเวโก, โวสคฺโค, จริยา, ฌานวิโมกฺโข, ภาวนา, อธิฏฺฐานํ ชีวิตํ. (19)}

\proseitem{41}{\csromanfolio{416}\textbf{นิจฺฉาโต}ติ เนกฺขมฺเมน กามจฺฉนฺทโต นิจฺฉาโต, อพฺยาปาเทน พฺยาปาทโต นิจฺฉาโต ฯเปฯ ปฐเมน ฌาเนน นีวรเณหิ นิจฺฉาโต ฯเปฯ อรหตฺตมคฺเคน สพฺพกิเลเสหิ นิจฺฉาโต. (1)}

\prose{\textbf{โมกฺโข วิโมกฺโข}ติ เนกฺขมฺเมน กามจฺฉนฺทโต มุจฺจตีติ โมกฺโข วิโมกฺโข, อพฺยาปาเทน พฺยาปาทโต มุจฺจตีติ โมกฺโข วิโมกฺโข ฯเปฯ ปฐเมน ฌาเนน นีวรเณหิ มุจฺจตีติ โมกฺโข วิโมกฺโข ฯเปฯ อรหตฺตมคฺเคน สพฺพกิเลเสหิ มุจฺจตีติ โมกฺโข วิโมกฺโข. (2)}

\prose{\textbf{วิชฺชาวิมุตฺตี}ติ เนกฺขมฺมํ วิชฺชตีติ วิชฺชา, กามจฺฉนฺทโต มุจฺจตีติ วิมุตฺติ.\csromanspacer{} วิชฺชนฺโต มุจฺจติ, มุจฺจนฺโต วิชฺชตีติ วิชฺชาวิมุตฺติ.\csromanspacer{} อพฺยาปาโท\footnote{อพฺยาปาทํ (สฺยา)} วิชฺชตีติ วิชฺชา, พฺยาปาทโต วิมุจฺจตีติ วิมุตฺติ.\csromanspacer{} วิชฺชนฺโต มุจฺจติ, มุจฺจนฺโต วิชฺชตีติ วิชฺชาวิมุตฺติ ฯเปฯ อรหตฺตมคฺโค วิชฺชตีติ วิชฺชา, สพฺพกิเลเสหิ มุจฺจตีติ วิมุตฺติ.\csromanspacer{} วิชฺชนฺโต มุจฺจติ, มุจฺจนฺโต วิชฺชตีติ วิชฺชาวิมุตฺติ. (3)}

\prose{\textbf{อธิสีลํ อธิจิตฺตํ อธิปญฺญา}ติ เนกฺขมฺเมน กามจฺฉนฺทํ สํวรฏฺเฐน สีลวิสุทฺธิ, อวิกฺเขปฏฺเฐน จิตฺตวิสุทฺธิ, ทสฺสนฏฺเฐน ทิฏฺฐิวิสุทฺธิ.\csromanspacer{} โย ตตฺถ สํวรฏฺโฐ, อยํ อธิสีลสิกฺขา.\csromanspacer{} โย ตตฺถ อวิกฺเขปฏฺโฐ, อยํ อธิจิตฺตสิกฺขา.\csromanspacer{} โย ตตฺถ ทสฺสนฏฺโฐ, อยํ อธิปญฺญาสิกฺขา.\csromanspacer{} อพฺยาปาเทน พฺยาปาทํ สํวรฏฺเฐน สีลวิสุทฺธิ ฯเปฯ อรหตฺตมคฺเคน สพฺพกิเลเส สํวรฏฺเฐน สีลวิสุทฺธิ, อวิกฺเขปฏฺเฐน จิตฺตวิสุทฺธิ.\csromanspacer{} ทสฺสนฏฺเฐน ทิฏฺฐิวิสุทฺธิ.\csromanspacer{} โย ตตฺถ สํวรฏฺโฐ, อยํ อธิสีลสิกฺขา.\csromanspacer{} โย ตตฺถ อวิกฺเขปฏฺโฐ, อยํ อธิจิตฺตสิกฺขา.\csromanspacer{} โย ตตฺถ ทสฺสนฏฺโฐ, อยํ อธิปญฺญาสิกฺขา. (3, 6)}

\prose{\textbf{ปสฺสทฺธี}ติ เนกฺขมฺเมน กามจฺฉนฺทํ ปฏิปฺปสฺสมฺเภติ, อพฺยาปาเทน พฺยาปาทํ ปฏิปฺปสฺสมฺเภติ ฯเปฯ อรหตฺตมคฺเคน สพฺพกิเลเส ปฏิปฺปสฺสมฺเภติ. (7)}

\prose{\textbf{ญาณนฺ}ติ กามจฺฉนฺทสฺส ปหีนตฺตา เนกฺขมฺมํ ญาตฏฺเฐน ญาณํ, พฺยาปาทสฺส ปหีนตฺตา อพฺยาปาโท ญาตฏฺเฐน ญาณํ ฯเปฯ สพฺพกิเลสานํ ปหีนตฺตา อรหตฺตมคฺโค ญาตฏฺเฐน ญาณํ. (8)}

\prose{\textbf{ทสฺสนนฺ}ติ กามจฺฉนฺทสฺส ปหีนตฺตา เนกฺขมฺมํ ทิฏฺฐตฺตา ทสฺสนํ, พฺยาปาทสฺส ปหีนตฺตา อพฺยาปาโท ทิฏฺฐิตฺตา ทสฺสนํ ฯเปฯ สพฺพกิเลสานํ ปหีนตฺตา อรหตฺตมคฺโค ทิฏฺฐตฺตา ทสฺสนํ. (9)}

\csromanheader{2}{10. มาติกากถา}
\prose{\csromanfolio{417}\textbf{วิสุทฺธี}ติ กามจฺฉนฺทํ ปชหนฺโต เนกฺขมฺเมน วิสุชฺฌติ, พฺยาปาทํ ปชหนฺโต อพฺยาปาเทน วิสุชฺฌติ ฯเปฯ สพฺพกิเลเส ปชหนฺโต อรหตฺตมคฺเคน วิสุชฺฌติ. (10)}

\prose{\textbf{เนกฺขมฺมนฺ}ติ กามานเมตํ นิสฺสรณํ, ยทิทํ เนกฺขมฺมํ.\csromanspacer{} รูปานเมตํ นิสฺสรณํ, ยทิทํ อารุปฺปํ.\csromanspacer{} ยํ โข ปน กิญฺจิ ภูตํ สงฺขตํ ปฏิจฺจสมุปฺปนฺนํ, นิโรโธ ตสฺส เนกฺขมฺมํ.\csromanspacer{} พฺยาปาทสฺส อพฺยาปาโท เนกฺขมฺมํ.\csromanspacer{} ถินมิทฺธสฺส อาโลกสญฺญา เนกฺขมฺมํ ฯเปฯ สพฺพกิเลสานํ อรหตฺตมคฺโค เนกฺขมฺมํ. (11)}

\prose{\textbf{นิสฺสรณนฺ}ติ กามานเมตํ นิสฺสรณํ, ยทิทํ เนกฺขมฺมํ.\csromanspacer{} รูปานเมตํ นิสฺสรณํ, ยทิทํ อารุปฺปํ.\csromanspacer{} ยํ โข ปน กิญฺจิ ภูตํ สงฺขตํ ปฏิจฺจสมุปฺปนฺนํ, นิโรโธ ตสฺส นิสฺสรณํ.\csromanspacer{} กามจฺฉนฺทสฺส เนกฺขมฺมํ นิสฺสรณํ, พฺยาปาทสฺส อพฺยาปาโท นิสฺสรณํ ฯเปฯ สพฺพกิเลสานํ อรหตฺตมคฺโค นิสฺสรณํ. (12)}

\prose{\textbf{ปวิเวโก}ติ กามจฺฉนฺทสฺส เนกฺขมฺมํ ปวิเวโก ฯเปฯ สพฺพกิเลสานํ อรหตฺตมคฺโค ปวิเวโก. (13)}

\prose{\textbf{โวสคฺโค}ติ เนกฺขมฺเมน กามจฺฉนฺทํ โวสชฺชตีติ โวสคฺโค, อพฺยาปาเทน พฺยาปาทํ โวสชฺชตีติ โวสคฺโค ฯเปฯ อรหตฺตมคฺเคน สพฺพกิเลเส โวสชฺชตีติ โวสคฺโค. (14)}

\prose{\textbf{จริยา}ติ กามจฺฉนฺทํ ปชหนฺโต เนกฺขมฺเมน จรติ, พฺยาปาทํ ปชหนฺโต อพฺยาปาเทน จรติ ฯเปฯ สพฺพกิเลเส ปชหนฺโต อรหตฺตมคฺเคน จรติ. (15)}

\prose{\textbf{ฌานวิโมกฺโข}ติ เนกฺขมฺมํ ฌายตีติ ฌานํ, กามจฺฉนฺทํ ฌาเปตีติ ฌานํ, ฌายนฺโต มุจฺจตีติ ฌานวิโมกฺโข, ฌาเปนฺโต มุจฺจตีติ ฌานวิโมกฺโข, ฌายนฺตีติ ธมฺมา, ฌาเปนฺตีติ กิเลเส, ฌาเต จ ฌาเป จ ชานาตีติ ฌานฌายี\footnote{ฌานวิโมกฺโข (สฺยา) อฏฺฐกถา โอโลเกตพฺพา.}.\csromanspacer{} อพฺยาปาโท ฌายตีติ ฌานํ, พฺยาปาทํ ฌาเปตีติ ฌานํ ฯเปฯ อาโลกสญฺญา ฌายตีติ ฌานํ, ถินมิทฺธํ ฌาเปตีติ ฌานํ ฯเปฯ}

\prose{\csromanfolio{418}อรหตฺตมคฺโค ฌายตีติ ฌานํ, สพฺพกิเลเส ฌาเปตีติ ฌานํ, ฌายนฺโต มุจฺจตีติ ฌานวิโมกฺโข, ฌาเปนฺโต มุจฺจตีติ ฌานวิโมกฺโข, ฌายนฺตีติ ธมฺมา, ฌาเปนฺตีติ กิเลเส, ฌาเต จ ฌาเป จ ชานาตีติ ฌานฌายี. (16)}

\proseitem{42}{\textbf{ภาวนา อธิฏฺฐานํ ชีวิตนฺ}ติ กามจฺฉนฺทํ ปชหนฺโต เนกฺขมฺมํ ภาเวตีติ ภาวนาสมฺปนฺโน, เนกฺขมฺมวเสน จิตฺตํ อธิฏฺฐาตีติ อธิฏฺฐานสมฺปนฺโน, สฺวายํ เอวํ ภาวนาสมฺปนฺโน อธิฏฺฐานสมฺปนฺโน สมํ ชีวติ, โน วิสมํ.\csromanspacer{} สมฺมา ชีวติ, โน มิจฺฉา.\csromanspacer{} วิสุทฺธํ ชีวติ, โน กิลิฏฺฐนฺติ อาชีวสมฺปนฺโน.\csromanspacer{} สฺวายํ เอวํ ภาวนาสมฺปนฺโน อธิฏฺฐานสมฺปนฺโน อาชีวสมฺปนฺโน ยญฺญเทว ปริสํ อุปสงฺกมติ ยทิ ขตฺติยปริสํ ยทิ พฺราหฺมณปริยํ ยทิ คหปติปริสํ ยทิ สมณปริสํ วิสารโท อุปสงฺกมติ อมงฺกุภูโต, ตํ กิสฺส เหตุ, ตถา หิ โส ภาวนาสมฺปนฺโน อธิฏฺฐานสมฺปนฺโน อาชีวสมฺปนฺโน.}

\prose{พฺยาปาทํ ปชหนฺโต อพฺยาปาทํ ภาเวตีติ ภาวนาสมฺปนฺโน ฯเปฯ.\csromanspacer{} ถินมิทฺธํ ปชหนฺโต อาโลกสญฺญํ ภาเวตีติ ภาวนาสมฺปนฺโน ฯเปฯ.\csromanspacer{} อุทฺธจฺจํ ปชหนฺโต อวิกฺเขปํ ภาเวตีติ ภาวนาสมฺปนฺโน ฯเปฯ.\csromanspacer{} วิจิกิจฺฉํ ปชหนฺโต ธมฺมววตฺถานํ ภาเวตีติ ภาวนาสมฺปนฺโน ฯเปฯ.\csromanspacer{} อวิชฺชํ ปชหนฺโต วิชฺชํ ภาเวตีติ ภาวนาสมฺปนฺโน ฯเปฯ.\csromanspacer{} อรติํ ปชหนฺโต ปาโมชฺชํ ภาเวตีติ ภาวนาสมฺปนฺโน ฯเปฯ.\csromanspacer{} นีวรเณ ปชหนฺโต ปฐมํ ฌานํ ภาเวตีติ ภาวนาสมฺปนฺโน ฯเปฯ.\csromanspacer{} สพฺพกิเลเส ปชหนฺโต อรหตฺตมคฺคํ ภาเวตีติ ภาวนาสมฺปนฺโน.\csromanspacer{} อรหตฺตมคฺควเสน จิตฺตํ อธิฏฺฐาตีติ อธิฏฺฐานสมฺปนฺโน.\csromanspacer{} สฺวายํ เอวํ ภาวนาสมฺปนฺโน อธิฏฺฐานสมฺปนฺโน สมํ ชีวติ, โน วิสมํ.\csromanspacer{} สมฺมา ชีวติ, โน มิจฺฉา.\csromanspacer{} วิสุทฺธํ ชีวติ, โน กิลิฏฺฐนฺติ อาชีวสมฺปนฺโน.\csromanspacer{} สฺวายํ เอวํ ภาวนาสมฺปนฺโน อธิฏฺฐานสมฺปนฺโน อาชีวสมฺปนฺโน ยญฺญเทว ปริสํ อุปสงฺกมติ ยทิ ขตฺติยปริสํ ยทิ พฺราหฺมณปริสํ ยทิ คหปติปริสํ ยทิ สมณปริสํ วิสารโท อุปสงฺกมติ อมงฺกุภูโต, ตํ กิสฺส เหตุ, ตถา หิ โส ภาวนาสมฺปนฺโน อธิฏฺฐานสมฺปนฺโน อาชีวสมฺปนฺโนติ. (3, 19)}

\nitthitam{มาติกากถา นิฏฺฐิตา.}
\vspace{\tipitakaparskip}
\csromancenter{ปญฺญาวคฺโค ตติโย.}
\csromanheader{2}{อุทฺทานคาถา}
\titlehead{\textbf{ตสฺสุทฺทานํ}}
\csromangathameasure{ปญฺญา อิทฺธิ อภิสมโย, วิเวโก จริยปญฺจโม. \\ ปาฏิหาริ สมสีสิ, สติปฏฺฐานา วิปสฺสนา. \\ ตติเย ปญฺญาวคฺคมฺหิ, มาติกาย จ เต ทสาติ. \\ มหาวคฺโค ยุคนทฺโธ, ปญฺญาวคฺโค จ นามโต. \\ ตโยว วคฺคา อิมมฺหิ, ปฏิสมฺภิทาปกรเณ. \\ อนนฺตนยมคฺเคสุ, คมฺภีโร สาครูปโม. \\ นภํ จ ตารกากิณฺณํ, ถูโล ชาตสฺสโร ยถา. \\ กถิกานํ วิสาลาย, โยคีนํ ญาณโชตนนฺติ.}
\csromangathagroup{\csromangathastanza{\csromanfolio{419}ปญฺญา อิทฺธิ อภิสมโย, วิเวโก จริยปญฺจโม. \\ ปาฏิหาริ สมสีสิ, สติปฏฺฐานา วิปสฺสนา.}\csromangathastanza{ตติเย ปญฺญาวคฺคมฺหิ, มาติกาย จ เต ทสาติ. \\ มหาวคฺโค ยุคนทฺโธ, ปญฺญาวคฺโค จ นามโต.}\csromangathastanza{ตโยว วคฺคา อิมมฺหิ\footnote{ติวคฺโค ยสฺส วิกฺเขโป (สฺยา, ก)}, ปฏิสมฺภิทาปกรเณ. \\ อนนฺตนยมคฺเคสุ, คมฺภีโร สาครูปโม.}\csromangathastanza{นภํ จ ตารกากิณฺณํ, ถูโล ชาตสฺสโร ยถา. \\ กถิกานํ วิสาลาย, โยคีนํ ญาณโชตนนฺติ.}}

\nitthitam{\textbf{ปฏิสมฺภิทามคฺคปาฬิ นิฏฺฐิตา.}}
\orphannote{อาวชฺชนปฏิพนฺธา (ก) เอวมีทิเสสุ ปเทสุ.}

